\documentclass[11pt]{article}
%
%
%
%
%
%
%
\usepackage{fix-cm}
\usepackage[a4paper,margin=1.80cm]{geometry}
\usepackage{fontspec}
\usepackage{xeCJK}
%
\setCJKmainfont[
  BoldFont={Microsoft JhengHei Bold},
  ItalicFont={Microsoft JhengHei},
  ItalicFeatures={FakeSlant=0.2},
  BoldItalicFont={Microsoft JhengHei Bold},
  BoldItalicFeatures={FakeSlant=0.2}
]{Microsoft JhengHei}
\setCJKsansfont[BoldFont={Microsoft JhengHei Bold}]{Microsoft JhengHei}
\setCJKmonofont[BoldFont={Microsoft JhengHei Bold}]{Microsoft JhengHei}
\setmainfont{Noto Serif}
\setsansfont{Noto Sans}
\setmonofont{Noto Sans Mono}
\usepackage{amsmath,amssymb,mathtools,array,booktabs,longtable,pdflscape,adjustbox,diagbox,amsthm,xparse,hyperref,enumitem}
\hypersetup{unicode=true,hidelinks,pdfauthor={Emmy Noether},pdftitle={埃米·諾特數學著作集：繁體字版——基本正文}}
\setlength{\parindent}{2em}
\setlength{\parskip}{0.35em}
\setlist{nosep,leftmargin=2.5em}
\setlist[itemize]{leftmargin=*,itemsep=0pt,topsep=2pt}
\setlist[enumerate]{leftmargin=*,itemsep=2pt,topsep=2pt}
\allowdisplaybreaks
\sloppy
\hbadness=10000
\vbadness=10000
\hfuzz=5pt
\vfuzz=12pt
\emergencystretch=240em
\raggedbottom
\makeatletter
\renewcommand\footnotesize{\@setfontsize\footnotesize{8.2}{9.2}\selectfont}
\@ifundefined{selectlanguage}{\newcommand{\selectlanguage}[1]{}}{\renewcommand{\selectlanguage}[1]{}}
\makeatother
\theoremstyle{plain}
\newtheorem{satz}{定理}[section]
\newtheorem{hilfssatz}{引理}[section]
\newtheorem{behauptung}{斷言}[section]
\theoremstyle{definition}
\newtheorem{definition}{定義}[section]
\theoremstyle{remark}
\newtheorem{bemerkung}{注}[section]
\renewcommand{\proofname}{證明}

\providecommand{\A}{\Delta}
\providecommand{\D}{\Delta}
\providecommand{\R}{\varrho}
\providecommand{\eps}{\varepsilon}
\providecommand{\pair}[2]{(#1\,|\,#2)}
\providecommand{\modu}[1]{\pmod{#1}}
\providecommand{\hata}{\widehat{a}}
\providecommand{\hatnu}{\widehat{\nu}}
\providecommand{\mcell}[1]{\ensuremath{\begin{gathered}#1\end{gathered}}}
\providecommand{\Kfield}{\mathfrak{K}}
\providecommand{\Ssys}{\mathfrak{S}}
\providecommand{\Mfield}{\mathfrak{M}}
\providecommand{\Lsys}{\mathfrak{L}}
\providecommand{\Jdom}{\mathfrak{J}}
\providecommand{\Gdom}{\mathfrak{G}}
\providecommand{\Hgrp}{\mathfrak{H}}
\providecommand{\GaloisRes}{\Phi}
\providecommand{\pol}[2]{\left(#1\frac{\partial}{\partial #2}\right)}
\providecommand{\Deltaop}{\Delta\!\left(\frac{\partial}{\partial x},\frac{\partial}{\partial y},\frac{\partial}{\partial z}\right)}
\providecommand{\Omegaop}{\Omega}
\providecommand{\Afield}{\mathfrak{A}}
\providecommand{\Bfield}{\mathfrak{B}}
\providecommand{\Cfield}{\mathfrak{C}}
\providecommand{\Rfield}{\mathfrak{R}}
\providecommand{\Sdomain}{\mathfrak{S}}
\providecommand{\Yreal}{\mathfrak{Y}}
\providecommand{\bdelta}{\bar{\delta}}
\providecommand{\dd}{\mathrm d}
\providecommand{\Div}{\operatorname{Div}}
\providecommand{\G}{\mathfrak G}
\providecommand{\bd}{\bar{\delta}}
\providecommand{\p}{\partial}
\providecommand{\frS}{\mathfrak S}
\providecommand{\frp}{\mathfrak p}
\providecommand{\fra}{\mathfrak a}
\providecommand{\frb}{\mathfrak b}
\providecommand{\frc}{\mathfrak c}
\providecommand{\frf}{\mathfrak f}
\providecommand{\frD}{\mathfrak D}
\providecommand{\frr}{\mathfrak r}
\providecommand{\frm}{\mathfrak m}
\providecommand{\frn}{\mathfrak n}
\providecommand{\calA}{\mathcal A}
\providecommand{\calB}{\mathcal B}
\providecommand{\calN}{\mathcal N}
\providecommand{\frH}{\mathfrak H}
\providecommand{\frK}{\mathfrak K}
\providecommand{\frG}{\mathfrak G}
\providecommand{\tmod}[1]{\;(\operatorname{mod} #1)}
\providecommand{\mM}{\mathfrak{M}}
\providecommand{\mN}{\mathfrak{N}}
\providecommand{\mL}{\mathfrak{L}}
\providecommand{\mG}{\mathfrak{G}}
\providecommand{\mA}{\mathfrak{A}}
\providecommand{\mB}{\mathfrak{B}}
\providecommand{\mx}{\mathfrak{x}}
\providecommand{\my}{\mathfrak{y}}
\providecommand{\mz}{\mathfrak{z}}
\providecommand{\me}{\mathfrak{e}}
\providecommand{\mbarA}{\overline{\mathfrak{A}}}
\providecommand{\mbarB}{\overline{\mathfrak{B}}}
\providecommand{\bara}{\overline{\mathfrak{a}}}
\providecommand{\widebar}[1]{\overline{#1}}
\providecommand{\mU}{\mathfrak{U}}
\providecommand{\mP}{\mathfrak{P}}
\providecommand{\mQ}{\mathfrak{Q}}
\providecommand{\mS}{\mathfrak{S}}
\providecommand{\mD}{\mathfrak{D}}
\providecommand{\mR}{\mathfrak{R}}
\providecommand{\mV}{\mathfrak{V}}
\providecommand{\modu}[1]{\;\operatorname{mod} #1}
\providecommand{\ideal}[1]{\mathfrak{#1}}
\providecommand{\Amod}{\mathfrak A}
\providecommand{\Bmod}{\mathfrak B}
\providecommand{\Gmod}{\mathfrak G}
\providecommand{\Mmod}{\mathfrak M}
\providecommand{\Nmod}{\mathfrak N}
\providecommand{\Rmod}{\mathfrak R}
\providecommand{\Smod}{\mathfrak S}
\providecommand{\mideal}{\mathfrak m}
\providecommand{\nideal}{\mathfrak n}
\providecommand{\gideal}{\mathfrak g}
\providecommand{\hideal}{\mathfrak h}
\providecommand{\jideal}{\mathfrak j}
\providecommand{\tideal}{\mathfrak t}
\providecommand{\aideal}{\mathfrak a}
\providecommand{\bideal}{\mathfrak b}
\providecommand{\bb}{\mathfrak b}
\providecommand{\cc}{\mathfrak c}
\providecommand{\zero}[1]{\equiv 0\pmod{#1}}
\providecommand{\Norm}{N}
\providecommand{\mJ}{\mathfrak J}
\providecommand{\mS}{\mathfrak S}
\providecommand{\mo}{\mathfrak o}
\providecommand{\ma}{\mathfrak a}
\providecommand{\mb}{\mathfrak b}
\providecommand{\Om}{\Omega}
\providecommand{\dx}{\dd x}
\providecommand{\dy}{\dd y}
\providecommand{\fraka}{\mathfrak a}
\providecommand{\frakb}{\mathfrak b}
\providecommand{\frakc}{\mathfrak c}
\providecommand{\frakg}{\mathfrak g}
\providecommand{\frakh}{\mathfrak h}
\providecommand{\frakm}{\mathfrak m}
\providecommand{\frakn}{\mathfrak n}
\providecommand{\frako}{\mathfrak o}
\providecommand{\frakp}{\mathfrak p}
\providecommand{\frakq}{\mathfrak q}
\providecommand{\frakr}{\mathfrak r}
\providecommand{\frakt}{\mathfrak t}
\providecommand{\frakS}{\mathfrak S}
\providecommand{\barfrakm}{\bar{\mathfrak m}}
\providecommand{\mC}{\mathfrak{C}}
\providecommand{\mH}{\mathfrak{H}}
\providecommand{\mf}{\mathfrak{f}}
\providecommand{\mo}{\mathfrak{o}}
\providecommand{\mpideal}{\mathfrak{p}}
\providecommand{\mK}{\mathfrak{K}}
\providecommand{\mT}{\mathfrak{T}}
\providecommand{\bP}{\overline{P}}
\providecommand{\bR}{\overline{\mathfrak{R}}}
\providecommand{\ba}{\overline{\mathfrak{a}}}
\providecommand{\bp}{\overline{\mathfrak{p}}}
\providecommand{\bt}{\overline{\mathfrak{t}}}
\providecommand{\mat}[1]{\begin{pmatrix}#1\end{pmatrix}}
\providecommand{\OO}{\Omega}
\providecommand{\Th}{\Theta}
\providecommand{\iso}{\simeq}
\providecommand{\normlhd}{\triangleleft}
\providecommand{\Trdeg}{\operatorname{Trdeg}}
\providecommand{\const}{\operatorname{const.}}
\providecommand{\frakO}{\mathfrak{O}}
\providecommand{\frakp}{\mathfrak{p}}
\providecommand{\Gg}{\mathfrak{G}}
\providecommand{\frakC}{\mathfrak{C}}
\providecommand{\frakq}{\mathfrak{q}}
\providecommand{\frakr}{\mathfrak{r}}
\providecommand{\frakA}{\mathfrak{A}}
\providecommand{\frakK}{\mathfrak{K}}
\providecommand{\frakG}{\mathfrak{G}}
\providecommand{\N}{\operatorname{N}}
\providecommand{\frG}{\mathfrak{G}}
\providecommand{\Gcal}{\mathcal{G}}
\providecommand{\frA}{\mathfrak{A}}
\providecommand{\frB}{\mathfrak{B}}
\providecommand{\frC}{\mathfrak{C}}
\providecommand{\frL}{\mathfrak{L}}
\providecommand{\frM}{\mathfrak{M}}
\providecommand{\frN}{\mathfrak{N}}
\providecommand{\frR}{\mathfrak{R}}
\providecommand{\frS}{\mathfrak{S}}
\providecommand{\frT}{\mathfrak{T}}
\providecommand{\Sp}{\operatorname{Sp}}
\providecommand{\Ii}{\mathfrak{I}}
\providecommand{\PP}{\mathfrak{P}}
\providecommand{\Norm}{\operatorname{N}}
\providecommand{\Hh}{\mathfrak{H}}
\providecommand{\oo}{\mathfrak{o}}
\providecommand{\glqq}{``}
\providecommand{\grqq}{''}
\providecommand{\ZZ}{\mathbb{Z}}
\providecommand{\QQ}{\mathbb{Q}}
\providecommand{\KK}{\mathfrak K}
\providecommand{\ida}{\mathfrak a}
\providecommand{\cO}{\mathfrak O}
\providecommand{\cD}{\mathfrak D}
\providecommand{\Tr}{\operatorname{Tr}}
\providecommand{\Rang}{\operatorname{Rang}}
\providecommand{\tensor}{\otimes}
\providecommand{\wzbw}{證畢}
\providecommand{\foreign}[1]{#1}

%
\providecommand{\frakA}{\mathfrak{A}}
\providecommand{\frakB}{\mathfrak{B}}
\providecommand{\frakC}{\mathfrak{C}}
\providecommand{\frakD}{\mathfrak{D}}
\providecommand{\frakE}{\mathfrak{E}}
\providecommand{\frakF}{\mathfrak{F}}
\providecommand{\frakG}{\mathfrak{G}}
\providecommand{\frakH}{\mathfrak{H}}
\providecommand{\frakI}{\mathfrak{I}}
\providecommand{\frakJ}{\mathfrak{J}}
\providecommand{\frakK}{\mathfrak{K}}
\providecommand{\frakL}{\mathfrak{L}}
\providecommand{\frakM}{\mathfrak{M}}
\providecommand{\frakN}{\mathfrak{N}}
\providecommand{\frakO}{\mathfrak{O}}
\providecommand{\frakP}{\mathfrak{P}}
\providecommand{\frakQ}{\mathfrak{Q}}
\providecommand{\frakR}{\mathfrak{R}}
\providecommand{\frakS}{\mathfrak{S}}
\providecommand{\frakT}{\mathfrak{T}}
\providecommand{\frakU}{\mathfrak{U}}
\providecommand{\frakV}{\mathfrak{V}}
\providecommand{\frakW}{\mathfrak{W}}
\providecommand{\frakX}{\mathfrak{X}}
\providecommand{\frakY}{\mathfrak{Y}}
\providecommand{\frakZ}{\mathfrak{Z}}
\providecommand{\fraka}{\mathfrak{a}}
\providecommand{\frakb}{\mathfrak{b}}
\providecommand{\frakc}{\mathfrak{c}}
\providecommand{\frakd}{\mathfrak{d}}
\providecommand{\frake}{\mathfrak{e}}
\providecommand{\frakf}{\mathfrak{f}}
\providecommand{\frakg}{\mathfrak{g}}
\providecommand{\frakh}{\mathfrak{h}}
\providecommand{\fraki}{\mathfrak{i}}
\providecommand{\frakj}{\mathfrak{j}}
\providecommand{\frakk}{\mathfrak{k}}
\providecommand{\frakl}{\mathfrak{l}}
\providecommand{\frakm}{\mathfrak{m}}
\providecommand{\frakn}{\mathfrak{n}}
\providecommand{\frako}{\mathfrak{o}}
\providecommand{\frakp}{\mathfrak{p}}
\providecommand{\frakq}{\mathfrak{q}}
\providecommand{\frakr}{\mathfrak{r}}
\providecommand{\fraks}{\mathfrak{s}}
\providecommand{\frakt}{\mathfrak{t}}
\providecommand{\fraku}{\mathfrak{u}}
\providecommand{\frakv}{\mathfrak{v}}
\providecommand{\frakw}{\mathfrak{w}}
\providecommand{\frakx}{\mathfrak{x}}
\providecommand{\fraky}{\mathfrak{y}}
\providecommand{\frakz}{\mathfrak{z}}
%
\providecommand{\mC}{\mathfrak{C}}
\providecommand{\mF}{\mathfrak{F}}
\providecommand{\mI}{\mathfrak{I}}
\providecommand{\mJ}{\mathfrak{J}}
\providecommand{\mK}{\mathfrak{K}}
\providecommand{\mO}{\mathfrak{O}}
\providecommand{\mQ}{\mathfrak{Q}}
\providecommand{\mW}{\mathfrak{W}}
\providecommand{\mX}{\mathfrak{X}}
\providecommand{\mY}{\mathfrak{Y}}
\providecommand{\mZ}{\mathfrak{Z}}
\providecommand{\ma}{\mathfrak{a}}
\providecommand{\mb}{\mathfrak{b}}
\providecommand{\mc}{\mathfrak{c}}
\providecommand{\md}{\mathfrak{d}}
\providecommand{\me}{\mathfrak{e}}
\providecommand{\mf}{\mathfrak{f}}
\providecommand{\mg}{\mathfrak{g}}
\providecommand{\mh}{\mathfrak{h}}
\providecommand{\mi}{\mathfrak{i}}
\providecommand{\mj}{\mathfrak{j}}
\providecommand{\mk}{\mathfrak{k}}
\providecommand{\ml}{\mathfrak{l}}
\providecommand{\mm}{\mathfrak{m}}
\providecommand{\mn}{\mathfrak{n}}
\providecommand{\mo}{\mathfrak{o}}
\providecommand{\mr}{\mathfrak{r}}
\providecommand{\ms}{\mathfrak{s}}
\providecommand{\mt}{\mathfrak{t}}
\providecommand{\muideal}{\mathfrak{u}}
\providecommand{\mv}{\mathfrak{v}}
\providecommand{\mw}{\mathfrak{w}}
\providecommand{\mx}{\mathfrak{x}}
\providecommand{\my}{\mathfrak{y}}
\providecommand{\mz}{\mathfrak{z}}
\providecommand{\mq}{\mathfrak{q}}
%
\providecommand{\mA}{\mathfrak{A}}
\providecommand{\mB}{\mathfrak{B}}
\providecommand{\mC}{\mathfrak{C}}
\providecommand{\mD}{\mathfrak{D}}
\providecommand{\mE}{\mathfrak{E}}
\providecommand{\mF}{\mathfrak{F}}
\providecommand{\mG}{\mathfrak{G}}
\providecommand{\mH}{\mathfrak{H}}
\providecommand{\mI}{\mathfrak{I}}
\providecommand{\mJ}{\mathfrak{J}}
\providecommand{\mK}{\mathfrak{K}}
\providecommand{\mL}{\mathfrak{L}}
\providecommand{\mM}{\mathfrak{M}}
\providecommand{\mN}{\mathfrak{N}}
\providecommand{\mO}{\mathfrak{O}}
\providecommand{\mP}{\mathfrak{P}}
\providecommand{\mQ}{\mathfrak{Q}}
\providecommand{\mR}{\mathfrak{R}}
\providecommand{\mS}{\mathfrak{S}}
\providecommand{\mT}{\mathfrak{T}}
\providecommand{\mU}{\mathfrak{U}}
\providecommand{\mV}{\mathfrak{V}}
\providecommand{\mW}{\mathfrak{W}}
\providecommand{\mX}{\mathfrak{X}}
\providecommand{\mY}{\mathfrak{Y}}
\providecommand{\mZ}{\mathfrak{Z}}
\NewDocumentEnvironment{zusatz}{o}{\par\smallskip\noindent\IfNoValueTF{#1}{\textbf{補充.}}{\textbf{補充#1.}}\ }{\par\smallskip}
\RequirePackage{tikz}
\newcommand{\NoetherRosetteGlyph}{%
  \begin{tikzpicture}[x=0.1363584905660377pt,y=-0.1363584905660377pt,baseline=(current bounding box.south)]
    \path[use as bounding box] (0,0) rectangle (108,108);
    \fill[even odd rule]
    (33.5,84) -- (32.5,83) -- (30.5,83) -- (29,81.5) -- (29,76.5) -- (30.5,75) -- (31.5,75) -- (32.5,74) -- (34.5,74) -- (36,72.5) -- (36,71.5) -- (35.5,71) -- (33.5,71) -- (33,70.5) -- (33,69.5) -- (30.5,67) -- (28.5,69) -- (26.5,69) -- (24.5,71) -- (23.5,71) -- (22.5,72) -- (21.5,72) -- (20.5,73) -- (18.5,73) -- (17.5,72) -- (16.5,73) -- (16,72.5) -- (16,68.5) -- (18.5,66) -- (19.5,66) -- (20,65.5) -- (20,64.5) -- (26,58.5) -- (26,56.5) -- (24,54.5) -- (24,53.5) -- (17,46.5) -- (17,45.5) -- (16.5,45) -- (15.5,45) -- (15,44.5) -- (15,42.5) -- (15.5,42) -- (17.5,42) -- (19.5,44) -- (20.5,44) -- (21,44.5) -- (20,45.5) -- (21,46.5) -- (21,47.5) -- (22,48.5) -- (22,49.5) -- (29,56.5) -- (29,57.5) -- (27,59.5) -- (27,60.5) -- (26,61.5) -- (26,62.5) -- (24,64.5) -- (24,65.5) -- (22,67.5) -- (22.5,68) -- (24.5,68) -- (25.5,67) -- (28.5,67) -- (29.5,66) -- (31.5,66) -- (32.5,65) -- (33.5,65) -- (34.5,64) -- (36.5,64) -- (37.5,63) -- (38.5,63) -- (39.5,62) -- (41.5,62) -- (42.5,61) -- (44.5,61) -- (45.5,60) -- (46.5,60) -- (47.5,59) -- (48.5,59) -- (49,58.5) -- (49,57.5) -- (49.5,57) -- (51.5,57) -- (52,56.5) -- (50.5,55) -- (49.5,55) -- (48.5,54) -- (47.5,54) -- (46.5,53) -- (45.5,53) -- (43,50.5) -- (43,49.5) -- (42,48.5) -- (42,47.5) -- (41.5,47) -- (40.5,47) -- (39.5,48) -- (33.5,48) -- (32,46.5) -- (32,44.5) -- (34,42.5) -- (34,40.5) -- (33.5,40) -- (32.5,40) -- (31.5,39) -- (26.5,39) -- (26,38.5) -- (26,35.5) -- (27.5,34) -- (31.5,34) -- (33,35.5) -- (33,36.5) -- (34.5,38) -- (37.5,38) -- (38,37.5) -- (38,36.5) -- (41,33.5) -- (41,31.5) -- (40,30.5) -- (40,27.5) -- (39,26.5) -- (39,25.5) -- (38,24.5) -- (38,23.5) -- (37,22.5) -- (37,19.5) -- (37.5,19) -- (38.5,20) -- (39.5,20) -- (41,21.5) -- (41,24.5) -- (42,25.5) -- (42,27.5) -- (43,28.5) -- (43,30.5) -- (44,31.5) -- (44,35.5) -- (45,36.5) -- (45,37.5) -- (45.5,38) -- (46.5,38) -- (47,38.5) -- (47,40.5) -- (48,41.5) -- (48,42.5) -- (49,43.5) -- (49,44.5) -- (50,45.5) -- (50,48.5) -- (53,51.5) -- (52,52.5) -- (52,55.5) -- (53.5,57) -- (56.5,57) -- (58.5,59) -- (59.5,59) -- (60.5,60) -- (61.5,60) -- (62,60.5) -- (62,61.5) -- (63,62.5) -- (63,64.5) -- (63.5,65) -- (64.5,65) -- (65.5,64) -- (68.5,64) -- (69.5,65) -- (73.5,65) -- (74,65.5) -- (74,66.5) -- (75,67.5) -- (75,68.5) -- (75.5,69) -- (76,68.5) -- (76,66.5) -- (76.5,66) -- (78.5,68) -- (80.5,68) -- (81.5,69) -- (84.5,69) -- (85,68.5) -- (84,67.5) -- (84,65.5) -- (83.5,65) -- (82.5,65) -- (82,64.5) -- (82.5,64) -- (84.5,64) -- (86,65.5) -- (86,66.5) -- (88.5,69) -- (89.5,69) -- (90,69.5) -- (90,70.5) -- (91,71.5) -- (91,72.5) -- (93,74.5) -- (91.5,76) -- (90.5,76) -- (88.5,74) -- (87.5,74) -- (86.5,73) -- (85.5,73) -- (83.5,71) -- (82.5,71) -- (81.5,70) -- (80.5,70) -- (79.5,69) -- (77.5,69) -- (74.5,72) -- (73.5,72) -- (73,72.5) -- (73,74.5) -- (73.5,75) -- (78.5,75) -- (79.5,74) -- (80.5,74) -- (81,74.5) -- (81,78.5) -- (80,79.5) -- (80,80.5) -- (78.5,82) -- (73.5,82) -- (73,81.5) -- (73,79.5) -- (72,78.5) -- (72,77.5) -- (71.5,77) -- (70.5,77) -- (69.5,76) -- (68.5,76) -- (68,76.5) -- (68,77.5) -- (67,78.5) -- (67,79.5) -- (66.5,80) -- (65.5,80) -- (65,80.5) -- (66,81.5) -- (64.5,83) -- (63,81.5) -- (63,79.5) -- (62,78.5) -- (62,77.5) -- (61,76.5) -- (61,74.5) -- (60,73.5) -- (60,72.5) -- (59,71.5) -- (59,68.5) -- (58,67.5) -- (58,64.5) -- (56,62.5) -- (56,61.5) -- (55.5,61) -- (54.5,62) -- (53.5,62) -- (52,63.5) -- (52,64.5) -- (51,65.5) -- (51,66.5) -- (50,67.5) -- (50,68.5) -- (49,69.5) -- (49,70.5) -- (48,71.5) -- (48,75.5) -- (46,77.5) -- (46,79.5) -- (44.5,81) -- (42.5,81) -- (42,80.5) -- (44,78.5) -- (44,77.5) -- (43.5,77) -- (41.5,77) -- (39.5,75) -- (38.5,76) -- (37.5,76) -- (37,76.5) -- (37,78.5) -- (36,79.5) -- (38,81.5) -- (38,82.5) -- (37.5,83) -- (34.5,83) -- cycle
    (65.5,22) -- (63.5,22) -- (63,21.5) -- (63.5,21) -- (65.5,21) -- (66,21.5) -- cycle
    (55.5,28) -- (53.5,28) -- (53,27.5) -- (53.5,27) -- (54.5,27) -- (57.5,24) -- (59.5,24) -- (60,24.5) -- (57.5,27) -- (56.5,27) -- cycle
    (64.5,26) -- (64,25.5) -- (64.5,25) -- (65,25.5) -- cycle
    (63.5,29) -- (63,28.5) -- (63,27.5) -- (63.5,27) -- (64.5,27) -- (65,27.5) -- cycle
    (58.5,55) -- (57,53.5) -- (57,52.5) -- (56,51.5) -- (56,49.5) -- (57,48.5) -- (57,46.5) -- (60,43.5) -- (60,42.5) -- (59.5,42) -- (58.5,42) -- (58,41.5) -- (59,40.5) -- (59,39.5) -- (61.5,37) -- (62.5,37) -- (63.5,36) -- (65.5,36) -- (66.5,37) -- (68.5,37) -- (70,35.5) -- (69,34.5) -- (69,32.5) -- (68,31.5) -- (68.5,31) -- (74.5,31) -- (75,31.5) -- (75,32.5) -- (76,33.5) -- (76,35.5) -- (75,36.5) -- (75,37.5) -- (72,40.5) -- (72,41.5) -- (71,42.5) -- (72,43.5) -- (72,44.5) -- (72.5,45) -- (73.5,45) -- (74.5,46) -- (75.5,46) -- (76.5,45) -- (78.5,45) -- (79.5,44) -- (80,44.5) -- (80,47.5) -- (79.5,48) -- (77.5,48) -- (76.5,49) -- (74.5,49) -- (73.5,50) -- (71.5,50) -- (70.5,51) -- (68.5,51) -- (67.5,52) -- (65.5,52) -- (64.5,53) -- (63.5,53) -- (62.5,54) -- (59.5,54) -- cycle
    (86.5,50) -- (85.5,50) -- (85,49.5) -- (86,48.5) -- (86,46.5) -- (85.5,46) -- (82.5,46) -- (81.5,47) -- (81,46.5) -- (82,45.5) -- (82,44.5) -- (82.5,44) -- (84.5,44) -- (86.5,42) -- (88.5,42) -- (89.5,41) -- (92.5,41) -- (93,41.5) -- (93,42.5) -- (94,43.5) -- (94,44.5) -- (93.5,45) -- (92.5,45) -- (88.5,49) -- (87.5,49) -- cycle
    (44,42.5) -- (43.5,42) -- (43,42.5) -- (43.5,43) -- cycle
    (45,43.5) -- (44.5,43) -- (44,43.5) -- (44.5,44) -- cycle
    (80.5,60) -- (79.5,60) -- (79,59.5) -- (79,56.5) -- (81,54.5) -- (81,53.5) -- (82.5,52) -- (83,52.5) -- (83,53.5) -- (82,54.5) -- (82,55.5) -- (81.5,56) -- (80.5,56) -- (80,56.5) -- (80,57.5) -- (81,58.5) -- (81,59.5) -- cycle
    (55.5,54) -- (55,53.5) -- (55.5,53) -- (56,53.5) -- cycle
    (43,65.5) -- (42.5,65) -- (42,65.5) -- (42.5,66) -- cycle
    (61,67.5) -- (60.5,67) -- (60,67.5) -- (60.5,68) -- cycle
    (47,69.5) -- (47,68.5) -- (46.5,68) -- (46,68.5) -- (46,69.5) -- (46.5,70) -- cycle
    (62,68.5) -- (61.5,68) -- (61,68.5) -- (61.5,69) -- cycle
    (67.5,97) -- (65.5,97) -- (64,95.5) -- (64,94.5) -- (61.5,92) -- (60.5,92) -- (54.5,86) -- (53.5,86) -- (49.5,90) -- (48.5,90) -- (47.5,91) -- (46.5,91) -- (45.5,92) -- (45,91.5) -- (45,90.5) -- (45.5,90) -- (46.5,90) -- (48.5,88) -- (49.5,88) -- (54.5,83) -- (55.5,83) -- (60.5,88) -- (61.5,88) -- (62.5,89) -- (63.5,89) -- (64,89.5) -- (64,90.5) -- (64.5,91) -- (66.5,91) -- (67,90.5) -- (66,89.5) -- (67,88.5) -- (65,86.5) -- (65,85.5) -- (65.5,85) -- (66.5,85) -- (67,85.5) -- (67,86.5) -- (68,87.5) -- (68,90.5) -- (69,91.5) -- (69,94.5) -- (70,95.5) -- (69.5,96) -- (68.5,96) -- cycle
    (40.5,96) -- (36.5,96) -- (36,95.5) -- (37,94.5) -- (37,91.5) -- (40.5,88) -- (41.5,88) -- (42,88.5) -- (41,89.5) -- (41,92.5) -- (42,93.5) -- (42,94.5) -- cycle;
  \end{tikzpicture}%
}
\providecommand{\noethpIIrosette}{\makebox[9.44444pt][c]{\raisebox{-3.85pt}[0pt][0pt]{\NoetherRosetteGlyph}}}
\providecommand{\NoetherRosetteInline}{\makebox[17.72671698pt][c]{\raisebox{-3.85pt}[10.87671698pt][3.85pt]{\NoetherRosetteGlyph}}}
\begin{document}
\thispagestyle{empty}
\pdfbookmark[0]{書名與版本說明}{edition-notice}
\begin{center}
{\LARGE\bfseries 埃米·諾特數學著作集}\par
\bigskip
{\Large 繁體字版——基本正文}\par
\bigskip
第1—43篇論文、1929/30年冬季學期講義、\par
Kapferer論文及與Noether合作的補充、書目與相關著作目錄
\end{center}
\vfill
\begingroup\small
\noindent\textbf{作者與中文文本}\par
原作者為Emmy Noether及各篇署名的共同作者；1929/30年講義由M. Deuring記錄。
中文底本由AI翻譯，本繁體字版由簡體中文文本作繁簡字形轉換而來，沿用其原有詞彙；未作臺灣、香港或澳門的區域語言本地化。
早期中文文本的製作方為Manuscript Typesetting Project / OpenAI Codex；現存出處資料不足以確定其準確模型。
R9中文彙編及繁體字轉換程式：OpenAI Codex——GPT-5.6 Sol，Ultra推理強度，2026-08-12。
字形轉換採用opencc-python-reimplemented 0.1.7的s2t設定及固定正字規則；繁體字形式與簡體文本屬於同一個中文譯本。
本版AI訂正與排版：OpenAI Codex——GPT-6 Astra，Ultra推理強度，2026-10-23至2026-10-24；Claude Opus 5.5（Anthropic），2026-10-26。
最初逐段翻譯的準確模型、推理設定及工作日期尚未確定；2026-06-12是早期彙編的版本日期。
\par\medskip
\noindent\textbf{版本來源}\par
本繁體字版繼承R9 / ED0011中文文本；其德文底本為ED0011/v047（2026-08-08）。
早期中文文本見2026-06-12彙編：\url{https://doi.org/10.5281/zenodo.20836874}。
歷史多語種彙編：\url{https://doi.org/10.5281/zenodo.20412587}。
各篇原刊資訊、作者署名與文獻引用保留在正文中；上述DOI並非當前GitHub版的新DOI。
\par\medskip
\noindent\textbf{權利與使用}\par
原著及受保護的衍生貢獻，其權利均歸各自權利人所有。
本版免費提供，僅供非商業、純教育用途；出版者推定權利人不會反對。這是推定，並非已獲證實的許可；免費可讀本身並不授權衍生出版。
權利人可透過\url{https://github.com/KokunoYumeto/modern-latex-manuscripts/issues/new}提出異議；出版者承諾依權利人的請求撤下相關材料。
本版不作新的整體許可或公有領域奉獻；既有的真實權利聲明仍在其實際授權範圍內適用。
\par\medskip
\noindent PDF書籤可導覽至各篇論文、講義及書目。完整可編輯正文與排版說明隨原始檔提供。
\par\medskip
\noindent\href{https://github.com/KokunoYumeto/modern-latex-manuscripts/blob/main/editions/noether/methods/chinese-script-2026-10-24-r2/README.zh-Hant.md}{字形轉換方法、程式與分項署名}。
\endgroup
\clearpage
\setcounter{page}{1}

%
%
%

%

%<unit id="noether.p01.title.001" type="title" scan="scan:p01:041">
\begin{center}
\pdfbookmark[0]{1. 關於三元雙二次型的型系統的構造}{paper-01}
{\Large\bfseries 1. 關於三元雙二次型的型系統的構造}\par
\vspace{1.5em}
埃爾朗根物理—醫學學會會報 39 (1907)，第 176--179 頁
\end{center}
%</unit>

\vspace{3em}
%<unit id="noether.p01.note.001" type="front-note" scan="scan:p01:041">
\begin{center}
（摘自作者的博士論文。）
\end{center}
%</unit>

%<unit id="noether.p01.par.001" type="paragraph" scan="scan:p01:041">
關於三元雙二次型的型系統，Gordan、Maisano 和 Pascal 的若干工作\footnote{P. Gordan, “Über das volle Formensystem der ternären biquadratischen Form” $f=x_1^3x_2+x_2^3x_3+x_3^3x_1$ (Math. Annalen, Bd. XVII, S. 217--233, 1880); G. Maisano, 1. “Sistemi completi dei primi cinque gradi della forma ternaria biquadratica e degl' invarianti, covarianti e contravarianti di sesto grado” (Giorn. di Battaglini XIX). 2. “Sui covarianti indipendenti di $6^\circ$ grado nei coefficienti della forma biquadratica ternaria” (Rend. Circ. Mat. di Palermo I, 1887); E. Pascal, “Contributo alla teoria della forma ternaria biquadratica e delle sue varie decomposizioni in fattori” (Memoria premiata dalla R. Accademia delle scienze fisiche e matematiche di Napoli, 1905).} 已有所論述。Gordan 在與他為二元領域中的型系統所給出的原則相類似的原則基礎上，對特殊自同構型
\[
f=x_1^3x_2+x_2^3x_3+x_3^3x_1
\]
建立了由 54 個構成式組成的完全型系統。
%</unit>

%<unit id="noether.p01.par.002" type="paragraph" scan="scan:p01:041-042">
在 Maisano 的工作中，對於一般雙二次型，直到第五階\footnote{這裡“階”應理解為相對於係數的維數，而“次數”應理解為相對於變量的維數。}（含第五階）的型都已列出，並且還列出了若干較高階的不變式、協變式和逆變式；所用方法是 Gordan 在 \emph{Mathematische Annalen} 第一卷中處理三元三次型時採用的方法。Pascal 利用 Maisano 的結果，主要研究雙二次型分解為因子的問題\footnote{在第二篇短札中，Maisano 試圖證明三個 6 階且 6 次的協變式之間存在線性依賴。與這個尚不完全的證明相反，Pascal 認為他已經證明了這三個 6 階協變式的線性無關性，也證明了三個 9 階不變式的線性無關性。然而，在三個協變式一方與三個不變式另一方之間事實上存在一條線性關係，這一點是在我的計算過程中顯現出來的；用 Pascal 的記號，顯式公式為
\[
0=20\Omega_1+6\Omega_2-3\Omega_3+4fC_1-12fC_2-4A\cdot\Delta,
\]
\[
0=4A^3-15AB+30C-90D+9E.
\]}。
%</unit>

%<unit id="noether.p01.par.003" type="paragraph" scan="scan:p01:042">
\textbf{我的研究目標，是為一般三元雙二次型建立型系統；首先只給出主要基礎，並建立所謂的“相對完備系統”\footnote{關於這一名稱，參見 Gordan--Kerschensteiner, \emph{Vorlesungen über Invariantentheorie}, Bd. II, S. 227.}。}
%</unit>

%<unit id="noether.p01.par.004" type="paragraph" scan="scan:p01:042">
三元型的系統構造的基本思想與二元領域中的思想相同。從第一個相對完備系統——即從二元情形繼承而來的型所構成的系統——出發，按照某一確定規律，過渡到模越來越高的系統，直到某一個模的系統——以該模作為基本型——成為有限並且已知，或者直到某個模可以化歸為以不變式為因子的型。通過把相對完備系統同該模的系統作轉積（德語 \emph{Überschiebung}），在第一種情形中得到絕對完備系統，而在第二種情形中，相對完備系統與絕對完備系統相同。由於 Hilbert 一般性地證明了型系統的有限性，這一過程必然會終止。
%</unit>

%<unit id="noether.p01.par.005" type="paragraph" scan="scan:p01:042-043">
在我們的情形中，第一個相對完備系統的模 $(abc)$ 可以歸結為下列兩個模
\[
\Delta=(abc)^2a_x^2b_x^2c_x^2\quad\text{以及}\quad \nu=(abu)^4;
\]
由此很容易得到模 $\nu$ 的相對完備系統。現在我們選取以下型作為模的序列：
\[
\nu;\qquad \nu^{(\nu)}=(\nu\nu_1x)^4=s_x^4,
\]
\[
\nu(s)=(ss'u)^4\ldots,
\]
並且把在模 $s$ 的相對完備系統的構造中出現的兩個二次型 $u_\varrho^2$ 和 $t_x^2$ 作為進一步的模加入。隨後可以證明，跟在 $s$ 後面的模 $(ss'u)^4$ 可化歸為模 $(\varrho,t)$。但兩個二次型的同時系統是有限且已知的，所以這樣就達到了上面所說的終點。模 $(\varrho,t)$ 的相對完備系統給出 331 個構成式。---
%</unit>

%<unit id="noether.p01.par.006" type="paragraph" scan="scan:p01:043">
我還要補充說明所用方法的若干細節。
%</unit>

%<unit id="noether.p01.par.007" type="paragraph" scan="scan:p01:043">
生成型的基本過程，即縮並過程（Faltungsprozeß），在三元領域中可以如下定義：
%</unit>

%<unit id="noether.p01.eq.001" type="formula" scan="scan:p01:043">
若在一個符號乘積
\[
s_x^m t_x^n u_\sigma^\mu u_\tau^\nu
\]
%</unit>
%<unit id="noether.p01.tbl.001" type="table" scan="scan:p01:043">
中把因子對
\[
\begin{array}{c|c|c|c|c}
 & s_xt_x & u_\sigma u_\tau & s_xu_\sigma\ \text{或}\ s_xu_\tau & t_xu_\tau\ \text{或}\ t_xu_\sigma\\
\text{分別替換為} & (stu) & (\sigma\tau x) & s_\sigma\ \text{或}\ s_\tau & t_\tau\ \text{或}\ t_\sigma\\
\text{縮並} & \mathrm{I} & \mathrm{II} & \mathrm{III} & \mathrm{IV}
\end{array}
\]
%</unit>
%<unit id="noether.p01.par.008" type="paragraph" scan="scan:p01:043">
則由此產生的型就是從原來的型經縮並而得到的\footnote{Gordan, Math. Annalen, Bd. XVII, S. 219.}。
%</unit>

%<unit id="noether.p01.par.009" type="paragraph" scan="scan:p01:043">
這裡還總可以令 $s_\sigma=0;\ t_\tau=0$。關於各個縮並之間的關係，有如下命題：
%</unit>

%<unit id="noether.p01.thm.001" type="theorem" scan="scan:p01:043-044">
\textbf{縮並 I 和 II 是基本縮並；縮並 III 和 IV 可以由它們組合而成，且與組合的次序無關。換言之：要形成由一個給定表達式經縮並所產生的全部型，只需施行縮並 I 和 II}\footnote{當 $n$ 與 $\mu$，或相應地 $m$ 與 $\nu$ 同時為零時，該命題有一個例外；此時“型列”化為對角項。}。
%</unit>

%<unit id="noether.p01.par.010" type="paragraph" scan="scan:p01:044">
所謂“型列”，是指一個初始型以及所有由它通過對自身縮並而產生的型。由上面的命題，一個型列 $s_x^m t_x^n u_\sigma^\mu u_\tau^\nu$ 可以排列成一個矩形圖式。在前進時，可以得到：
%</unit>

%<unit id="noether.p01.list.001" type="list" scan="scan:p01:044">
\begin{enumerate}
\item 向右移一列，得到所有由相鄰的型經縮並 I 產生的型；
\item 向下移一行，得到所有由上方的型經縮並 II 產生的型；因而得到所有由縮併產生的型。
\end{enumerate}
%</unit>

%<unit id="noether.p01.par.011" type="paragraph" scan="scan:p01:044">
在這些前提下，關於化簡子的定理可以在最一般的形式中表述出來；定義為
\begin{center}
“化簡子就是一個可化簡的型列”
\end{center}
時，有：
%</unit>

%<unit id="noether.p01.thm.002" type="theorem" scan="scan:p01:044">
\begin{center}
\textbf{若一個型列的初始型之所以可化簡，是因為它的某一項由同一個化簡子縮並而來，並且該型列的終型也是由這一同一項經縮並而產生的，則整個型列可化簡。}
\end{center}
%</unit>

%<unit id="noether.p01.par.012" type="paragraph" scan="scan:p01:044">
還應提到以下進一步的化簡方法：
%</unit>

%<unit id="noether.p01.list.002" type="list" scan="scan:p01:044">
\begin{enumerate}
\item 所謂“雙重化簡”，即以兩種方式化簡同一個型，以獲得高階型之間的關係；
\item “同分解型作縮並”，它一方面具有藉助化簡子進行化簡的特徵，另一方面又具有“雙重化簡”的特徵。
\end{enumerate}
%</unit>
%
\clearpage
%
\begingroup
\renewcommand{\modu}[1]{\;\operatorname{mod} #1}
\renewcommand{\glqq}{“}
\renewcommand{\grqq}{”}
\newcommand{\srcfn}[2]{%
  \begingroup\renewcommand{\thefootnote}{#1}\footnote{#2}\addtocounter{footnote}{-1}\endgroup}
\newcommand{\srcfnmark}[1]{%
  \begingroup\renewcommand{\thefootnote}{#1}\footnotemark\addtocounter{footnote}{-1}\endgroup}
\newcommand{\srcfntext}[2]{%
  \begingroup\renewcommand{\thefootnote}{#1}\footnotetext{#2}\endgroup}
\newcommand{\srcnumdisplay}[3][3.4em]{%
  \par\smallskip\noindent
  \begin{minipage}[t]{#1}\vspace*{0pt}#2\end{minipage}%
  \begin{minipage}[t]{\dimexpr\linewidth-#1\relax}\vspace*{0pt}\centering
  \(\displaystyle #3\)
  \end{minipage}\par\smallskip}
\setcounter{footnote}{0}
\begin{center}
\pdfbookmark[0]{2. 關於三元雙二次型之型系統的構造}{paper-02}
{\Large\bfseries 2. 關於三元雙二次型之型系統的構造\srcfn{*)}{引言和第一章的一小段摘錄曾刊於《埃朗根物理醫學學會會議報告》1907 年卷，第 176 至 179 頁。}.}\par
\vspace{0.8em}
埃朗根的 \emph{埃米·諾特} 小姐撰。\par
\vspace{1.2em}
《純粹與應用數學雜誌》134（1908），第 23--90 頁及兩張表
\end{center}

\vspace{3em}

\begin{center}
\textbf{引言。}
\end{center}

\emph{戈爾丹}、\emph{邁薩諾} 和 \emph{帕斯卡爾} 的著作均研究了三元雙二次型的型系統\srcfn{**)}{\emph{P. 戈爾丹}，《論三元雙二次型的完全型系統》
\[
f=x_1^3x_2+x_2^3x_3+x_3^3x_1.
\]
（《數學年刊》第 XVII 卷（1880），第 217--233 頁。）
\emph{G. 邁薩諾}，1) \emph{Sistemi completi dei primi cinque gradi della forma ternaria biquadratica e degl' invarianti, covarianti e contravarianti di sesto grado}.（Giorn. di Battaglini XIX.）
2) \emph{Sui covarianti indipendenti di $6^\circ$ grado nei coefficienti della forma biquadratica ternaria}.（Rend. Circ. Mat. di Palermo I, 1887.）
\emph{E. 帕斯卡爾}，\emph{Contributo alla teoria della forma ternaria biquadratica e delle sue varie decomposizioni in fattori}.（Memoria premiata dalla R. Accademia delle scienze fisiche e matematiche di Napoli. 1905.）}。戈爾丹先生以他為二元領域的型系統所提出的類似原則為基礎，為特殊自同型 \(f\equiv x_1^3x_2+x_2^3x_3+x_3^3x_1\) 建立了由 54 個構成式組成的完全型系統。

邁薩諾先生依照戈爾丹先生在《數學年刊》第一卷中用於三元三次型的方法，列出了一般雙二次型直至並包括第 5 階\srcfn{***)}{下文所謂“階”是指關於係數的次數，所謂“次數”則指關於變量的次數。}的各型，以及若干較高階的不變式、協變式和逆變式。帕斯卡爾先生利用邁薩諾的結果，主要研究雙二次型分解為因子的問題\srcfn{*)}{在第二篇短札中，邁薩諾先生試圖證明三個 6 階、6 次協變式之間存在線性相關。與這一尚不完備的證明相反，帕斯卡爾先生認為自己已經證明這三個 6 階協變式線性無關，也證明三個 9 階不變式線性無關。然而，這三個協變式之間以及三個不變式之間事實上分別存在一條線性關係；本文公式 (13)、§ 11 以及公式 (22)、§ 17 注中明確給出了這些關係。據帕斯卡爾先生來函說明，這一矛盾源於他論文第 46 頁逆變式 $p$（本文記作 $\varrho$）的表達式中有一個數值錯誤。}。

\emph{以下研究的目標是建立一般三元雙二次型的型系統；本文只給出主要基礎，並建立一個所謂的“相對完備系統”}\srcfn{**)}{術語參見 § 3a。}\emph{。}

本文與\emph{戈爾丹}的工作緊密相承；不過，他在那裡僅以十分概括的方式給出的原則，尚須逐項具體展開。藉助一個關於各縮並之間聯繫的定理，並引入“型列”（§ 1），我們對化簡定理作出嚴格表述並加以補全（§ 3）；而特殊級數展開的遞歸構造（§ 2）則提供了實際實施化簡的計算手段。

三元型系統構造的基本思想與二元領域中的相同。從第一個相對完備系統——即從二元情形承襲而來的各型所成的系統——出發，按照確定的規律得到模越來越高的系統，直至某個模的系統——把該模當作基型——成為有限且已知的系統，或者直至某個模可以化簡為以不變式為因子的型。在第一種情形中，將這個相對完備系統與該模的系統作轉積，便得到絕對完備系統；在第二種情形中，相對完備系統與絕對完備系統相同。由於\emph{希爾伯特}已經一般地證明了型系統的有限性，這一過程必然終止。

在我們的情形中，第一個相對完備系統（§ 4）的模 $(abc)$ 被化歸為模
\[
\Delta=(abc)^2a_x^2b_x^2c_x^2\quad\text{和}\quad \nu=(abu)^4
\]
（§ 5），繼而求得模 $\nu$ 下的相對完備系統（§ 6）。對於一般雙二次型，這些結果已經見於\emph{戈爾丹}的工作；不過，我們的推導方式更容易推廣到次數更高的型。

現在我們選取下列各型作為模的序列（§ 7）：
\[
\nu,\qquad \nu(\nu)=(\nu\nu_1x)^4=s_x^4,\qquad \nu(s)=(ss'u)^4,\ldots .
\]

在構造模 $s$ 下的相對完備系統時，把系統中出現的兩個二次型 $u_\varrho^2$ 和 $t_x^2$ 添作模（§ 9 和 10），因而構造模 $(s,\varrho,t)$ 下的相對完備系統；隨後證明（§ 17），模 $(ss'u)^4$ 可以化簡為模 $\varrho$ 和 $t$。由此，再高一級的系統轉化為模 $(\varrho,t)$ 下的相對完備系統。可是，兩個二次型的聯合系統是有限且已知的，在一般情形下由 20 個構成式組成\srcfn{*)}{參見 \emph{克萊佈施--林德曼}，《幾何學講義》（第 I 卷，第 288 頁及以下）：兩個同變型的系統。\emph{戈爾丹}，《論圓錐曲線束》（《數學年刊》第 XIX 卷，第 530 頁）：兩個逆變型的系統。}，所以隨着模 $(\varrho,t)$ 下相對完備系統的建立，便達到了上述終點。為構造絕對完備系統而將此係統與 $(\varrho,t)$ 的系統作轉積，則留待以後完成。

模 $(\varrho,t)$ 下的相對完備系統共得到 331 個構成式；附表按它們關於變量 $x$ 和 $u$ 的次數排列。

最後再列出所引入的符號：
\begin{align*}
f&=a_x^4,\\
\theta&=\theta_x^4u_\vartheta^2=(abu)^2a_x^2b_x^2,\\
K&=K_x^6u_k^3=(a\theta u)u_\vartheta^2a_x^3\theta_x^3,\\
j&=u_j^6=(a\theta u)^4u_\vartheta^2,\\
\Delta&=\Delta_x^6=a_\vartheta^2a_x^2\theta_x^4=(abc)^2a_x^2b_x^2c_x^2,\\
N&=N_x^8u_n=(a\Delta u)a_x^3\Delta_x^5,\\
\nu&=u_\nu^4=(abu)^4,\\
H&=H_x^2u_\eta^4=(\nu\nu_1x)^2u_\nu^2u_{\nu_1}^2,\\
L&=L_x^3u_l^6=(\nu\eta x)H_x^2u_\nu^3u_\eta^3,\\
g&=g_x^6=(\nu\eta x)^4H_x^2,\\
\sigma&=u_\sigma^6=H_\nu^2u_\nu^2u_\eta^4,\\
s&=s_x^4=(\nu\nu_1x)^4,\\
Z&=Z_x^4u_\zeta^2=(ss'u)^2s_x^2s_x'^2,\\
\varrho&=u_\varrho^2=\theta_\nu^4u_\vartheta^2,\\
t&=t_x^2=a_\eta^4H_x^2,\\
i&=a_\nu^4,\\
J&=s_\nu^4.
\end{align*}

\begin{center}
{\Large\bfseries 第一章。\quad 關於三元型的一般定理。}
\end{center}

\begin{center}
\textbf{§ 1.}\quad \emph{縮並過程。型列。}
\end{center}

生成型的基本過程是\emph{縮並過程}，在三元領域中可定義如下。給定一個符號乘積：
\[
s_x^m t_x^n u_\sigma^\mu u_\tau^\nu .
\]
若將各因子對
\[
s_xt_x\qquad u_\sigma u_\tau\qquad s_xu_\sigma\ \text{或}\ s_xu_\tau
\qquad t_xu_\tau\ \text{或}\ t_xu_\sigma
\]
分別替換為
\[
(stu)\qquad (\sigma\tau x)\qquad s_\sigma\ \text{或}\ s_\tau
\qquad t_\tau\ \text{或}\ t_\sigma,
\]
即縮並
\[
\text{I}\qquad\text{II}\qquad\text{III}\qquad\text{IV},
\]
則所得各型便由原式經縮併產生\srcfn{*)}{\emph{戈爾丹}，《數學年刊》第 XVII 卷，第 219 頁。}。

這裡，表達式 $s_x^m t_x^n u_\sigma^\mu u_\tau^\nu$ 既可以理解為兩個型 $S\cdot T$ 的實際乘積，也可以理解為一個具有若干符號列的單一型：
\[
s_x^m t_x^n u_\sigma^\mu u_\tau^\nu=A_x^{m+n}u_\lambda^{\mu+\nu}.
\]
在第一種情形中，我們稱之為型 $S$ 與 $T$ 的縮並；在第二種情形中，則稱為“型 $A$ 的自縮並”。為簡便起見，下文總假設（後文考察的所有型事實上均如此）
\srcnumdisplay{(a)}{\begin{gathered}
s_\sigma^\lambda=0,\\
t_\tau^\varkappa=0
\end{gathered}}
對於任意非零的 $\lambda$ 和 $\varkappa$ 都成立，並且與任意其他縮並相結合時也成立。這表示型 $s_x^m t_y^n u_\sigma^\mu v_\tau^\nu$ 是所謂的“正規型”，即無論按變量 $x$ 和 $u$，還是按變量 $y$ 和 $v$ 施行 $\Omega$ 過程，它都消失。

因此，條件 (a) 並不構成限制，而總能達到\srcfn{*)}{參見 \emph{戈爾丹}，《數學年刊》第 V 卷，第 104 頁。}。

關於各個縮並之間的聯繫，有如下定理：

\emph{定理 I。縮並 I 和 II 是基本縮並；縮並 III 和 IV 可以由它們複合而成，且與複合的先後次序無關。換言之：要構造由一個給定表達式經縮並所產生的全部型，只須施行縮並 I 和 II。}

\textbf{證明：}由恆等式定理以及矩陣的乘積定理，並考慮 (a)，可得：
\[
(\widehat{st}\sigma x)=-t_\sigma,\qquad
(\widehat{\sigma\tau}su)=-s_\tau,
\]
\[
(\widehat{st}\tau x)=s_\tau,\qquad
(\widehat{\sigma\tau}tu)=t_\sigma,
\]
\[
(stu)(\sigma\tau x)=s_\tau+t_\sigma-u_x\cdot s_\tau t_\sigma .
\]
因此，複合縮並 I 和 II 即產生縮並 III 和 IV；無論將縮並 II 施於縮並 I 所得的因子 $(stu)u_\sigma u_\tau$，還是將縮並 I 施於縮並 II 所得的因子 $(\sigma\tau x)s_xt_x$，結論都相同。

我們把一個初始型連同由它作自縮並所得的全部型定義為一個\emph{型列}\srcfn{**)}{參見 \emph{克萊佈施}，《論不變量理論的一項基本問題》（Abh. der Gött. Ges. d. Wiss.，第 XVII 卷）：§ 17 以及 § 18 末尾。這裡的“型列”與該文所引入的“真正化簡的等價系統”不同，其區別在於按較高型排序的原則。}\setcounter{footnote}{9}，並用初始型來表示該型列。這裡，任何含有更多自縮並的型都稱為“較高型”；而在地位相同的縮並 $s_\tau$ 與 $t_\sigma$ 之間，必須規定其中一個為較高者。由型列的構造律可得：

1) 型列中的每一個型都關於初始型的符號呈線性。

2) 對於各有兩列逆變符號的初始型，型列唯一確定；對於符號列更多的初始型，則須在由恆等式定理聯繫起來的各縮並中指定某些特殊縮並，才能使其唯一規範化。

由定理 I，型列 $s_x^m t_x^n u_\sigma^\mu u_\tau^\nu$（$m>n;\ \mu>\nu$）可按如下矩形圖式排列：
\begingroup\small
\[
\begin{array}{c|c|c|c|c}
s_x^m t_x^n u_\sigma^\mu u_\tau^\nu
& (stu)
& (stu)^2
& \cdots
& \noethpIIrosette\\[0.3em]
(\sigma\tau x)
& t_\sigma,\ s_\tau
& t_\sigma(stu),\ s_\tau(stu)
& \cdots
& \\[0.3em]
(\sigma\tau x)^2
& t_\sigma(\sigma\tau x),\ s_\tau(\sigma\tau x)
& t_\sigma^2,\ t_\sigma s_\tau,\ s_\tau^2
& \cdots
& \vdots\\[0.3em]
\vdots&\vdots&\vdots&\ddots&\vdots\\[0.3em]
(\sigma\tau x)^\nu
& t_\sigma(\sigma\tau x)^{\nu-1},\ s_\tau(\sigma\tau x)^{\nu-1}
& t_\sigma^2(\sigma\tau x)^{\nu-2},\ t_\sigma s_\tau(\sigma\tau x)^{\nu-2},\ s_\tau^2(\sigma\tau x)^{\nu-2}
& \cdots
& \vdots\\[0.3em]
\vdots
& t_\sigma(\sigma\tau x)^\nu
& t_\sigma^2(\sigma\tau x)^{\nu-1},\ t_\sigma s_\tau(\sigma\tau x)^{\nu-1}
& \cdots
& \vdots\\[0.3em]
\vdots
& \vdots
& t_\sigma^2(\sigma\tau x)^\nu
& \cdots
& \vdots
\end{array}
\]
\endgroup
\begingroup\small
\[
\begin{array}{r@{\quad}c}
\noethpIIrosette\srcfnmark{*)}
&
\begin{array}{c|c|c}
(stu)^n & \cdots & \cdots\\[0.3em]
t_\sigma(stu)^{n-1},\ s_\tau(stu)^{n-1}
& s_\tau(stu)^n & \cdots\\[0.3em]
t_\sigma^2(stu)^{n-2},\ t_\sigma s_\tau(stu)^{n-2},\ s_\tau^2(stu)^{n-2}
& t_\sigma s_\tau(stu)^{n-1},\ s_\tau^2(stu)^{n-1}
& s_\tau^2(stu)^n .
\end{array}
\end{array}
\]
\endgroup
\srcfntext{*)}{此處以及下文類似之處，標有\NoetherRosetteInline{}的下部應接到圖式右上方同樣標記的位置。}

在這裡，若

1) 向右前進一列，便得到由相鄰各型經縮並 I 所產生的全部型；

2) 向下前進一行，便得到由上方各型經縮並 II 所產生的全部型；

從而得到經縮併產生的全部型。

總型列中的任意一個型 $t_\sigma^\kappa s_\tau^\lambda(stu)^\mu$ 或 $t_\sigma^\kappa s_\tau^\lambda(\sigma\tau x)^\nu$，都構成一個新型列的開端；這個新型列由矩形圖式中以 $t_\sigma^\kappa s_\tau^\lambda(stu)^\mu$ 或 $t_\sigma^\kappa s_\tau^\lambda(\sigma\tau x)^\nu$ 為左上角、並含有因子 $t_\sigma^\kappa s_\tau^\lambda$ 的所有型組成。

補充說明。定理 I 在數 $n$ 和 $\mu$（相應地，$m$ 和 $\nu$）等於零時有一個例外：這時縮並 I、II 和 IV 不再存在，而縮並 III（相應地，IV）仍然存在。在這種情形中，我們把圖式的對角項稱為型列：
\[
\begin{array}{c@{\qquad}c@{\qquad}c}
\begin{array}{ccccc}
s_x^m u_\tau^\nu &&&&\\
&s_\tau&&&\\
&&s_\tau^2&&\\
&&&\ddots&\\
&&&&s_\tau^\nu
\end{array}
&\text{相應地}&
\begin{array}{ccccc}
t_x^n u_\sigma^\mu &&&&\\
&t_\sigma&&&\\
&&t_\sigma^2&&\\
&&&\ddots&\\
&&&&t_\sigma^\mu
\end{array}
\end{array}
\]
與縮並 I 和 II 的缺失相一致，在這種情形中，按型列極式所作的級數展開總是退化為初始項；但關於化簡子的各定理仍然成立。

\begin{center}
\textbf{§ 2.}\quad \emph{按型列極式所作的級數展開}\srcfn{*)}{參見 Study，《三元型理論的方法》（Teubner 1889）II，§ 7。與該書的推導不同，我們的推導給出了快速計算數值係數 $C_{pr\varrho}$ 的方法。克萊佈施--Study 展開經特殊化 a)、a') 後將寫成：
\[
(stu)^\mu u_\tau^{\nu-\varkappa}v_\tau^\varkappa s_x^m t_x^{\,n-\lambda}(tuv)^\lambda
=\sum_{p,r,\varrho} C'_{pr\varrho}
\bigl[s_\tau^\varrho(stu)^{\mu+r-\varrho}\bigr]_{\widehat{uv}^{\lambda-r+\varrho}v^{\varkappa-\varrho},\,w^p,\,(w=x\widehat{uv})},
\]
所以它不能作出我們的展開在計算過程中已經出現的簡化，因為我們知道所有含因子 $u_x$ 的項都應略去。}。
\end{center}

對於含 $n$ 個變量的型，已有兩類級數展開被明確寫出\srcfn{**)}{參見 戈爾丹，《論組合量》，§ 2 和 5（《數學年刊》第 V 卷）。}：

1) 對於各有一列逆變變量的型 $s_x^m u_\tau^\nu$，有一個按 $u_x$ 的冪遞進的級數，其係數為“正規型”；

2) 對於有兩列同變變量的型 $s_x^m t_x^n$，有一個按基本協變式的極式（即型列 $s_x^m t_x^n$ 的極式）所作的展開；在三元情形下，它寫作：
\[
\tag{I}
s_x^m t_x^{n-\lambda}t_y^\lambda
=\sum_{\varrho}
\frac{\binom{m}{\varrho}\binom{\lambda}{\varrho}}
 {\binom{m+n-\varrho+1}{\varrho}}\cdot
\bigl[(st\widehat{xy})^\varrho s_x^{m-\varrho}t_x^{n-\varrho}\bigr]_{y^{\lambda-\varrho}}.
\]

我們把這兩類展開結合起來：對於各有兩列逆變變量的型
\[
s_x^m t_x^{n-\lambda}t_y^\lambda u_\sigma^\mu u_\tau^{\nu-\varkappa}v_\tau^\varkappa
\]
建立按 $u_x$ 的冪所作的展開，其係數是型列 $s_x^m t_x^n u_\sigma^\mu u_\tau^\nu$ 關於變量 $x$ 和 $u$ 的複合極式。

只須對各有兩列逆變符號（相應地，變量）的情形建立級數，因為每次把兩列符號（變量）合併成一個新的單列，就可將符號列（變量列）更多的型化歸到這種情形。

為了使型列的所有型都只含各兩列符號或變量，我們作如下特殊化：
\[
\begin{array}{ll}
\text{a) } \sigma=\widehat{st}, & \text{a') } y=\widehat{uv},\\[0.3em]
\text{b) } s=\widehat{\sigma\tau}, & \text{b') } v=\widehat{xy}
\end{array}
\]
並對特殊化 a)、a') 進行展開。

直接轉用展開 I，可得到型 $(stu)^\mu s_x^m t_x^{n-\lambda}t_y^\lambda$ 的複合極式；而對於型 $(stu)^\mu u_\tau^{\nu-\varkappa}v_\tau^\varkappa s_x^m t_x^{n-\lambda}t_y^\lambda$，所得各項不再是複合極式，對它們求值需要不斷引入新的輔助變量，並在最後才令這些變量相等，這會使計算變得不必要地複雜。

因此，我們採用一條間接途徑：把型列所有型的複合極式，即表達式
\[
\bigl[s_\tau^\varrho(stu)^{\mu-\varrho+r}\bigr]_{y^\lambda}^{v^\varkappa}
\]
用這些極式的初始項表示，再將所得遞歸方程組反演，從而得到所求的按極式展開。

根據轉移原理，對於型 $s_x^m t_x^n$ 的第 $\lambda$ 個極式 $s_x^m t_x^n:[s_x^m t_x^n]_{y^\lambda}$（$\lambda\le n$），有如下展開\srcfn{*)}{戈爾丹，《二元型的結式》（第 I 章，§ 5）。Rend. Circ. Mat. di Palermo XXII. 1906.}（當 $\lambda>n$ 時，可通過展開 $[s_y^m t_y^n]_{x^{m+n-\lambda}}$ 化歸到前一種情形）：
\[
\bigl[s_x^m t_x^n\bigr]_{y^\lambda}
=\sum_r (-1)^r\,
\frac{\binom{m}{r}\binom{\lambda}{r}}{\binom{m+n}{r}}\cdot
(st\widehat{xy})^r s_x^{m-r}t_x^{n-\lambda}t_y^{\lambda-r},
\]
相應地，對於 $u_\sigma^\mu u_\tau^\nu$ 的第 $\varkappa$ 個極式 $u_\sigma^\mu u_\tau^\nu:[u_\sigma^\mu u_\tau^\nu]^{v^\varkappa}$，有
\[
\bigl[u_\sigma^\mu u_\tau^\nu\bigr]^{v^\varkappa}
=\sum_\varrho (-1)^\varrho\,
\frac{\binom{\mu}{\varrho}\binom{\varkappa}{\varrho}}{\binom{\mu+\nu}{\varrho}}\,
(\sigma\tau\widehat{uv})^\varrho u_\sigma^{\mu-\varrho}u_\tau^{\nu-\varkappa}v_\tau^{\varkappa-\varrho}.
\]

相乘並引入特殊化 a) 和 a')，得到
\[
\begin{aligned}
\bigl[s_x^m t_x^n(stu)^\mu u_\tau^\nu\bigr]_{\widehat{uv}^{\lambda}}^{v^\varkappa}
={}&
\sum_{r,\varrho}c_{r\varrho}\cdot
(st\widehat{x\widehat{uv}})^r(\widehat{st}\tau\widehat{uv})^\varrho
s_x^{m-r}t_x^{n-\lambda}(tuv)^{\lambda-r}(stu)^{\mu-\varrho}
u_\tau^{\nu-\varkappa}v_\tau^{\varkappa-\varrho},
\end{aligned}
\]
由此按 $u_x$ 的冪展開，單列出首項（$\sum'_{r,\varrho}$ 表示略去 $r=0$、$\varrho=0$ 的項），並考慮第 26 頁的 (a)，得到
\[
\tag{II}
\begin{aligned}
&s_x^m t_x^{n-\lambda}(tuv)^\lambda(stu)^\mu u_\tau^{\nu-\varkappa}v_\tau^\varkappa\\
&\quad =
\bigl[s_x^m t_x^n(stu)^\mu u_\tau^\nu\bigr]_{\widehat{uv}^{\lambda}}^{v^\varkappa}\\
&\qquad
-\sum_{r,\varrho}' c_{r\varrho}\cdot s_\tau^\varrho(stu)^{\mu-\varrho+r}
u_\tau^{\nu-\varkappa}v_\tau^{\varkappa-\varrho}
s_x^{m-r}t_x^{n-\lambda}(tuv)^{\lambda-r+\varrho}\cdot v_x^r\\
&\qquad
+u_x\cdot\sum_{\varrho}\cdot\sum_{r=1}^{\lambda} c_{r\varrho}\cdot
s_\tau^\varrho(stu)^{\mu-\varrho+r-1}(stv)
u_\tau^{\nu-\varkappa}v_\tau^{\varkappa-\varrho}
s_x^{m-r}t_x^{n-\lambda}(tuv)^{\lambda-r+\varrho}\cdot v_x^{r-1}\\
&\qquad
\vdots\\
&\qquad
-(-1)^\lambda u_x^\lambda\cdot\sum_\varrho c_{\lambda\varrho}\cdot
s_\tau^\varrho(stu)^{\mu-\varrho}(stv)^\lambda
u_\tau^{\nu-\varkappa}v_\tau^{\varkappa-\varrho}
s_x^{m-\lambda}t_x^{n-\lambda}(tuv)^\varrho,
\end{aligned}
\]
其中
\[
c_{r\varrho}=(-1)^{r+\varrho}\cdot
\frac{\binom{m}{r}\binom{\lambda}{r}}{\binom{m+n}{r}}\cdot
\frac{\binom{\mu}{\varrho}\binom{\varkappa}{\varrho}}{\binom{\mu+\nu}{\varrho}}
\qquad
\begin{array}{l}
\text{當 }\lambda\le n,\\
\varkappa\le\nu.
\end{array}
\]
對於數
\[\lambda,n;\varkappa,\nu.\]
的其餘組合，情形相應。因此，表達式
\[s_x^m t_x^{n-\lambda}(tuv)^\lambda(stu)^\mu u_\tau^{\nu-\varkappa}v_\tau^\varkappa\]
被化歸為一個複合極式；再應用恆等式
\[
(stv)u_\tau=(stu)v_\tau-s_\tau(tuv),
\]
還可化為一組構造方式相同、但對應型列中較高型的表達式。

反覆迭代便得到展開：
\[
\tag{III}
s_x^m t_x^{n-\lambda}(tuv)^\lambda(stu)^\mu u_\tau^{\nu-\varkappa}v_\tau^\varkappa
=
\sum_{p,r,\varrho}u_x^p\cdot C_{pr\varrho}\cdot
\bigl[s_\tau^\varrho(stu)^{\mu-\varrho+r}\bigr]_{\widehat{uv}^{\lambda-r+\varrho}v^{\varkappa-\varrho+p}\cdot v_x^{r-p}}.
\]
數值係數 $C_{pr\varrho}$ 可由迭代 (II.) 所得方程組的係數 $c_{r\varrho}$、$c'_{r\varrho}$ 唯一地遞歸算出（參見第 42 頁的例子）。其餘特殊化給出類似的展開。

\begin{center}
\textbf{§ 3.}\quad \emph{化簡定理。}
\end{center}

現在把關於型和型系統化簡的已知定理同第 1、2 節聯繫起來；為此需要若干取自二元型理論的定義。

a) 對於預先給定的一列型，如果一個型系統具有如下性質，我們就稱它為\emph{模該列型的相對完備系統}：由這些型的任意乘積經縮並所得的一切構成式，都能表示為系統中各型的整有理函數，至多相差一個加性項，而該加性項由系統型與模的系統作縮並而產生\srcfn{*)}{戈爾丹--Kerschensteiner，《不變量理論講義》，第 II 卷，第 227 頁。}。

b) 如果一個型能夠用以不變式為因子的型或“較高型”表示，我們就稱它\emph{可化簡}；這裡所謂較高型，是指它所屬總型列中的較高型，或是在模的符號中含有較高階的型。

現在，關於化簡子\srcfn{**)}{戈爾丹，《數學年刊》第 XVII 卷，第 222 頁。}的定理可以用最一般的形式表述如下：

\textbf{定義：}\emph{化簡子是一個可化簡的型列。}

\textbf{定理 II。}\emph{設一個型列的初始型之所以可化簡，是由於它的某一項由同一個化簡子作縮並而產生（即含有一個化簡子作為因子），並且該型列的終型也正由這一項經縮併產生，則整個型列可化簡}\srcfn{*)}{定理 II 所帶來的簡化可由下例說明：對於特殊型 $f=x_1^3x_2+x_2^3x_3+x_3^3x_1$，$a_\nu^2a_x^2u_\nu^2$ 是化簡子。因此，由定理 II 可得型列
\[
\theta_\nu(\vartheta\nu x)\theta_x^3u_\vartheta u_\nu^2
=
a_\nu^2(abu)-\underbrace{a_\nu b_\nu(abu)}_{0},
\]
的化簡，因為 $\theta_\nu^4=a_\nu^2b_\nu^2(abu)^2$ 是由項 $a_\nu^2(abu)$ 經縮併產生的。反之，在戈爾丹前引文第 231 頁中，各型
\[
\theta_\nu(\vartheta\nu x),\quad \theta_\nu(\vartheta\nu x)^2,\quad
\theta_\nu^2,\quad \theta_\nu^2(\vartheta\nu x),\quad
\theta_\nu^2(\vartheta\nu x)^2
\]
則是逐一化簡的。}。

\textbf{證明：}該型列由初始項作自縮並而生，也就是說，一方面由化簡子的較高縮並而生，另一方面由化簡子的自縮並而生。在兩種情形中，由 § 2，型列的所有型都可表示為化簡子型列上的極式（轉積）。

由初始型推及整個型列的結論，對其餘化簡方法不再成立。這些方法是：

1) 所謂\emph{“雙重化簡”}。

所謂雙重化簡，是把一個含有兩個彼此獨立的化簡子作為因子的表達式用兩種方式化簡，由此得到被化簡至的那些較高型之間的關係。

系統地應用“雙重化簡”，一方面必定得到全部關係\srcfn{**)}{例如，通過“雙重化簡”得到型 $a_\eta^3$（§ 10）的化簡；它是邁薩諾系統中唯一多餘列入的型。此外，還得到引言注中所述三個協變式之間以及三個不變式之間的關係。}；另一方面，如果將“雙重化簡”用於不同表達式時不再產生新關係，就既得到判別較高型獨立性的準則，也得到對計算的檢驗。

2) \emph{與可分解型作縮並。}

這裡所謂“可分解型”——較通常概念略作推廣——是指能夠用次數較低的型的乘積以及“較高”型表示的型。型的乘積經一次縮並仍化為乘積；同樣，在特殊化 a') 和 b')（§ 2）中，\emph{非線性}型的乘積經二次縮並後，依據乘積定理也化為乘積：
\[
a_yb_y a_xb_x=\frac12a_y^2\cdot b_x^2+\frac12b_y^2\cdot a_x^2-\frac12(ab\widehat{yx})^2.
\]
因此，可分解型的一階極式（轉積）以及某些二階極式仍是可分解型。於是有：

一重（二重）轉積中，由與一個可分解型作一次（兩次）縮並所產生的一個\emph{項}，本身也是可分解型：

a) 如果可分解型的型列中緊接着的較高型分解或變為可化簡，或者在一次縮並時出現特殊化 $y=\widehat{uv}$ 或 $v=\widehat{xy}$；

b) 如果其餘的一次縮並導向較高型（參見 § 12，B. $H_k$ 和 $H_k(k\eta x)$）。

這樣的轉積項也可能分解：

c) 如果其餘的一次縮並導向較低型，而這些較低型又獨立於同可分解型的縮並而分解，即可以用乘積和待化簡的型表示；不過，每次都必須先通過計算確認有關項的數值係數並不恆等於零。這不過是較低型的“雙重化簡”（參見 § 23，C. $H_\vartheta(\theta H u)(\theta s u)$）。

與函數行列式所作的一切一次縮並都會產生可分解型\srcfn{*)}{參見戈爾丹前引文 § 3。}，某些較高縮並也會產生可分解型。有
\[
(st\widehat{yx})s_y=\frac12\,t\cdot s_y^2-\frac12\,s\cdot t_y^2+\frac12(st\widehat{yx})^2,
\]
\[
(st\widehat{yx})s_y^2=\frac13\,t\cdot s_y^3-\frac13\,s\cdot t_y^3+s_y(st\widehat{yx})^2-\frac13(st\widehat{yx})^3.
\]
對於下列對偶型，也有類似公式：
\[
(\sigma\tau\widehat{vu})v_\sigma
\quad\text{和}\quad
(\sigma\tau\widehat{vu})v_\sigma^2.
\]

根據本節各定理，為列出由 $S$ 與 $T$ 縮併產生的全部不可化簡型，可得如下規則：

從最低的縮並開始，沿任意一條事先確定的路徑（例如逐行或關於對角項對稱）構造型列 $ST$，直至遇到一個以化簡子為因子的型。只要終型滿足定理 II 所列的條件，就應略去由這個型確定的型列。剔除所有可化簡型列以後，再應用雙重化簡，除去其餘所有可化簡型以及所有可分解型。總型列中被彼此獨立的可化簡型列雙重覆蓋的位置，會引起較高型的化簡（參見 § 11，D）。

\begin{center}
\textbf{第二章。\quad 相對完備的基本系統。}

\textbf{§ 4.}\quad\emph{第一個相對完備系統。}
\end{center}

下文除第一個定理以外，只討論雙二次型。不過，只要適當選擇各模，§§ 5 和 6 的結果就很容易推廣到次數更高的型，正如它們也適用於三次型一樣\srcfn{*)}{對於三次型，相應的、以類似方式導出的公式為：
\[
(abc)a_x^2b_y^2c_z^2=\frac13(sxy)(sxz)(syz)+\frac23\Delta_x\Delta_y\Delta_z\cdot(xyz),\qquad (\mathrm{IV}.)
\]
\[
\begin{gathered}
s=(abc)(abu)(acu)(bcu),\qquad \Delta=(abc)^2a_xb_xc_x,\\
t=\Delta_\vartheta(\Delta\theta u)^2u_\vartheta=\theta_s^2u_su_\vartheta^2,
\qquad T=a_t^3,
\end{gathered}
\]
\[
\left.
\begin{aligned}
(a\theta u)^2&=\frac13 u_xs,\\
a_\vartheta&=\frac13u_x\cdot\Delta,
& a_\vartheta(a\theta u)&=0,
& a_\vartheta(a\theta u)^2&=s,\\
a_\vartheta^2&=\Delta,
& a_\vartheta^2(a\theta u)&=0.
\end{aligned}
\right\}\tag{1.}
\]
}。

根據下列已知定理，可以得到任意次數型的第一個相對完備系統：

\clearpage
\noindent\textbf{定理 III。}
\emph{從二元領域承襲而來的型，即根據克萊佈施轉移原理、由相應二元型的系統型經添邊所得的那些型，構成模 $(abc)$ 下的一個相對完備系統。}

\noindent\textbf{證明：}
設一個二元關係表明各型 $Q'_1,\ldots,Q'_n$ 構成某個二元基型的完全系統：
\[
Q'-F(Q')=0=\sum_\lambda\bigl\{(ab)c_x+(bc)a_x+(ca)b_x\bigr\}\varphi_\lambda\srcfnmark{*)}
\]
\srcfntext{*)}{戈爾丹--Kerschensteiner，前引書，第 134 頁。}
經添邊後，它對應於三元關係：
\[
Q-F(Q_i)=\sum_\lambda u_x\cdot(abc)\varphi_\lambda .
\]
而這正是模 $(abc)$ 下相對完備系統的定義方程。對於雙二次型，承襲而來的型的系統由下列各型組成：
\[
\begin{gathered}
f=a_x^4,\qquad
\theta=\theta_x^4u_\vartheta^2=(abu)^2a_x^2b_x^2,
\qquad
K=K_x^6u_k^3=(a\theta u)u_\vartheta^2a_x^3\theta_x^3,\\
j=u_j^6=(a\theta u)^4u_\vartheta^2,\qquad
\nu=u_\nu^4=(abu)^4,
\end{gathered}
\]
其中前四個型構成模 $((abc),\nu)$ 下的相對完備系統。

\begin{center}
\textbf{§ 5.}\quad\emph{把模 $(abc)$ 化歸為模 $\Delta$ 和 $\nu$。}
\end{center}

僅由一個括號因子給出的模 $(abc)$，按照定義（§ 3，a）應化歸為系統中的型。換言之：型列 $(abc)$ 必須唯一規範化（參見 § 1）。

\[
a_s^2a_xu_s=\frac13a_s^3\cdot u_x=\frac13S\cdot u_x. \tag{2.}
\]
\[
\left.
\begin{aligned}
(a\Delta u)^2&=a_s-\frac{S}{6}u_x^2,\\
(a\Delta u)^3&=0.
\end{aligned}
\right\} \tag{3.}
\]
\[
\left.
\begin{aligned}
\theta(s)=(ss_1x)^2&=2a_t+\frac13S\cdot\theta-\frac13Tu_x^2,\\
\theta(t)=(tt_1x)^2&=-\frac13S\cdot a_t+\frac23T\cdot a_s+\frac1{18}S^2\cdot\theta-\frac1{18}u_x^2\cdot S\cdot T.
\end{aligned}
\right\} \tag{4.}
\]

模的序列：$(abc)$；$\Delta$；$s$；$t$（模 $t$ 的系統只含型 $t$）。

我們將證明，下列各型：
\[
\Delta=(abc)^2a_x^2b_x^2c_x^2,\qquad \nu=(abu)^4
\]
作為模已經足夠，也就是說，各型
\[
\begin{gathered}
\Delta=(abc)^2a_x^2b_x^2c_x^2,
\qquad a_\nu=(abc)(bcu)^3a_x^3,\\
a_\nu^2=(abc)^2(bcu)^2a_x^2,
\qquad i=a_\nu^4=(abc)^4
\end{gathered}
\]
確定了型列 $(abc)$。在型列 $a_\nu$ 中，根據恆等式
\[
a_\nu^3=(abc)^3a_x(bcu)=\frac13u_x\cdot(abc)^4=\frac13i\cdot u_x.
\]
不存在型 $a_\nu^3$。

為了把一個模化歸為較高的模，按定義，我們必須把與前一模所作的最一般縮並，或者把一切需要考慮的特殊縮並，化歸為與較高模所作的縮並。

在我們的情形中，最一般的縮併產生於表達式：
\[
(abc)a_x^3b_y^3c_z^3,
\qquad
(abc)a_x^2b_y^2c_z^2a_xb_xc_x
\]
（從三次型承襲而來），
\[
(abc)a_xb_yc_za_x^2b_x^2c_x^2
\]
（從二次型承襲而來），以及這些表達式的極化。

為了用 $\Delta$ 的極式和型列 $a_\nu$ 的極式（或者說它們的初始項）計算這些表達式，我們同 § 2 一樣採用間接途徑。將各極式項
\[
(abc)^2a_x^2b_y^2c_z^2,
\qquad
 a_\nu(\nu xy)(\nu yz)(\nu zx)a_xa_ya_z=(abc)(bc\widehat{xy})(bc\widehat{yz})(bc\widehat{zx})a_xa_ya_z\quad\text{等等}
\]
展開為含符號因子 $(abc)$ 的表達式的線性組合；將所得方程組反演，便得到所求關係，以及按變量的不同組合所取各極式之間的一個關係。求值時使用行列式的恆等式定理和乘積定理。

對於 $(abc)a_x^3b_y^3c^3$，方程組如下：
\begingroup\small
\[
\begin{aligned}
1)\quad
\sum_3a_\nu(\nu yz)^3a_x^3
&=6(abc)a_x^3b_y^3c_z^3-6\sum_3(abc)a_x^3b_y^2b_zc_z^2c_y\srcfnmark{*)},\\[0.4em]
2)\quad
3(abc)^2a_x^2b_y^2c_z^2\cdot(xyz)-\frac12\sum_3a_\nu^2(\nu yz)^2a_x^2\cdot(xyz)
&=2\sum_3(abc)a_x^3b_y^2b_zc_z^2c_y\\
&\quad-4\sum_3(abc)a_xa_ya_zb_y^2b_xc_z^2c_x,\\[0.4em]
3)\quad
a_\nu(\nu xy)(\nu yz)(\nu zx)a_xa_ya_z
&=2\sum_3(abc)a_xa_ya_zb_y^2b_xc_z^2c_x,\\[0.4em]
4)\quad
\sum_3a_\nu(\nu yz)^3a_x^3-\frac32\sum_3a_\nu^2(\nu yz)^2a_x^2\cdot(xyz)+\frac16i\cdot(xyz)^3
&=6\sum_3(abc)a_xa_ya_zb_y^2b_xc_z^2c_x.
\end{aligned}
\]
\endgroup
\srcfntext{*)}{$\sum_3$ 表示把變量作循環置換後由初始項產生的三項之和。各方程對和式的每一項也分別成立。}

由此，對於 $(abc)a_x^3b_y^3c_z^3$，並經類似計算對於 $(abc)a_x^2b_y^2c_z^2a_xb_xc_x$、$(abc)a_xb_yc_za_x^2b_x^2c_x^2$，可得：
\[
\begin{aligned}
(abc)a_x^3b_y^3c_z^3
&=\frac32(abc)^2a_x^2b_y^2c_z^2\cdot(xyz)
+\frac32a_\nu(\nu xy)(\nu yz)(\nu zx)a_xa_ya_z
-\frac1{12}i\cdot(xyz)^3,\\[0.4em]
(\mathrm{IV}.)\quad
(abc)a_x^2b_y^2c_z^2a_xb_xc_x
&=\frac23(abc)^2a_x^2b_yb_xc_zc_x\cdot(xyz)
+\frac13a_\nu(\nu xy)(\nu yz)(\nu zx)a_x^3\\
&\quad-\frac16a_\nu^2(\nu xy)(\nu zx)a_x^2\cdot(xyz),\\[0.4em]
(abc)a_xb_yc_za_x^2b_x^2c_x^2
&=\frac16(abc)^2a_x^2b_x^2c_x^2\cdot(xyz).
\end{aligned}
\]

這樣，我們得到一個模 $(\Delta,\nu)$ 下的相對完備系統，它由四個型 $f,\theta,K,j$ 組成。

由此可知，$f$ 與 $\theta$ 的所有並非承襲自二元領域的轉積，以及在二元情形中可化簡為 $(ab)^4$ 的轉積 $(a\theta u)^2$，都可以用符號 $\Delta$ 和 $\nu$ 表示。

對於 $a_\vartheta(a\theta u)^2$、$a_\vartheta^2(a\theta u)^2$、$(a\theta u)^2$，通過將展開 IV 特殊化，或者更簡便地直接計算（交換符號），我們由如下設式得到下文所用的基本化簡公式：
\[
\begin{aligned}
a_\vartheta(a\theta u)^2
&=(abc)(bcu)(abu)^2-\frac13(abc)(bcu)\{(bcu)a_x-u_x\cdot(abc)\}^2,\\
a_\vartheta^2(a\theta u)^2
&=(abc)^2(abu)^2-\frac13(abc)^2\{(bcu)a_x-u_x\cdot(abc)\}^2
\end{aligned}
\]
對於 $(a\theta u)^2$\srcfn{*)}{戈爾丹--Kerschensteiner，前引書，第 92 頁。}：
\[
\begin{aligned}
(abu)a_y^3b_x^3&=\frac32u_\vartheta(\vartheta yx)\theta_y^2+\frac14(\nu yx)^3,\\
(abu)(acu)^3b_x^3&=\frac32u_\vartheta(\vartheta\widehat{au}x)(\theta au)^2+\frac14(\nu\widehat{au}x)^3:
\end{aligned}
\]
\begingroup\small
\[
\begin{array}{r@{\quad}|l@{\qquad}c|}
& \underline{(a\theta u)^2}=\frac16\nu\cdot f-\frac23u_x\cdot a_\nu+\frac23u_x^2\cdot a_\nu^2-\frac{i}{18}u_x^4
& \bigstar\\[0.55em]
\underline{a_\vartheta}=\frac13u_x\cdot\Delta
& \underline{a_\vartheta(a\theta u)}=0
& \\[0.55em]
& \underline{a_\vartheta^2}=\Delta
&
\end{array}
\]
\vspace{-0.5em}
\[
\begin{array}{r@{\quad}c@{\quad}|l@{\qquad}|l}
(1.) & \bigstar
& \underline{(a\theta u)^3}=0
& \underline{(a\theta u)^4}=j\\[0.55em]
&
& \underline{a_\vartheta(a\theta u)^2}=-\frac56a_\nu+\frac76u_x\cdot a_\nu^2-\frac{i}{9}u_x^3
& \underline{a_\vartheta(a\theta u)^3}=0\\[0.55em]
&
& \underline{a_\vartheta^2(a\theta u)}=0
& \underline{a_\vartheta^2(a\theta u)^2}=\frac23a_\nu^2-\frac{i}{9}u_x^2
\end{array}
\]
\endgroup

再添上同樣由模 $(abc)$ 的化簡得到的公式：
\[
\begin{array}{r@{\qquad}l}
(2.) & a_\nu^3a_xu_\nu=\frac13i\cdot u_x.
\end{array}
\]

\emph{除關於可分解型的關係外，後來所有化簡公式均由公式 (1.)、(2.) 以及對基型 $\nu$ 類似構造的公式經極化導出。}由公式 (1.) 可見：

下列型列分別導向：
\[
\begin{array}{rcl}
a_\vartheta(a\theta u)^2 &\text{導向}& \text{僅含符號 }\nu,\\[0.2em]
a_\vartheta &\text{導向}& \text{符號 }\Delta\text{ 和 }\nu,\\[0.2em]
(a\theta u)^2 &\text{導向}& \text{符號 }j,\Delta\text{ 和 }\nu,\\[0.2em]
(a\theta u)^2\text{ 經縮並 I} &\text{導向}& \text{符號 }j\text{ 和 }\nu,\\[0.2em]
(a\theta u)^2\text{ 經縮並 II} &\text{導向}& \text{符號 }\Delta\text{ 和 }\nu.
\end{array}
\]

因此可以作如下替換：
\begingroup
\renewcommand{\labelenumi}{\arabic{enumi})}
\begin{enumerate}
\item 模 $(\Delta,\nu)$：將與 $j$ 所作的縮並替換為與 $(a\theta u)^2$ 或 $(a\theta u)^3$ 所作的相應縮並；
\item 模 $(\nu)$：將與 $\Delta$ 所作的縮並替換為與 $a_\vartheta$ 或 $a_\vartheta(a\theta u)$ 所作的相應縮並，或者替換為與 $(a\theta u)^2$ 所作的特殊縮並（它只產生一次縮並 I）。
\end{enumerate}
\endgroup
特別有：
\[
\begin{aligned}
(a\theta u)^2a_y^2\theta_x^2&=\frac13(jyx)^2\modu{\nu},&
(a\theta u)^2a_y\theta_y&=-\frac16(jyx)^2\modu{\nu},\\[0.4em]
(a\theta u)^2v_\vartheta^2&=\frac13(\Delta uv)^2\modu{\nu},&
(a\theta u)(a\theta v)u_\vartheta v_\vartheta&=-\frac16(\Delta uv)^2\modu{\nu},\\[0.4em]
(a\theta u)^2a_yv_\vartheta^2&=\frac13(\Delta uv)^2\Delta_y\modu{\nu}.
\end{aligned}
\]

\begin{center}
\textbf{§ 6.}\quad\emph{把模 $(\Delta,\nu)$ 化歸為模 $(\nu)$。}
\end{center}

上一節的化簡公式使我們能夠直接建立模 $\nu$ 下的相對完備系統；我們將看到，只須在承襲而來的型的系統中加入型 $\Delta$ 和 $(a\Delta u)=N$（與三次型的情形類似，而且一般而言也成立）。

為了化簡模 $\Delta$，我們考察一切需要考慮的特殊縮並，即模 $(\Delta,\nu)$ 下相對完備系統與 $\Delta$ 的系統之間的縮並，並從關於係數階數最低的型——即 $f$ 與 $\Delta$ 的縮並——開始。

為了化簡型 $(a\Delta u)^2$、$(a\Delta u)^3$、$(a\Delta u)^4$，我們依 § 5 用型列中較低的型替換符號 $\Delta$，也就是展開表達式
\[
a_\vartheta b_\vartheta(b\theta u)^2,
\qquad a_\vartheta(a\theta u)b_\vartheta(b\theta u)^2,
\qquad a_\vartheta(a\theta u)b_\vartheta(b\theta u)^3
\]
1) 按型列 $a_\vartheta$ 的極式展開（符號 $\Delta$ 和 $\nu$），

2) 按型列 $b_\vartheta(b\theta u)^2$ 的極式展開（符號 $\nu$），

從而得到化簡公式（計算見下）：
\[
\begin{aligned}
\text{a)}\quad \underline{(a\Delta u)^2}
&=\frac12(\vartheta \nu x)^2-\frac25f\cdot a_\nu^2
-\frac45u_x\cdot\theta_\nu(\vartheta \nu x)^2
+u_x^2\left\{\frac16i f-\frac1{40}s\right\},\\
(3.)\quad\text{b)}\quad \underline{(a\Delta u)^3}
&=-\frac3{10}\theta_\nu(\vartheta \nu x),\\
\text{c)}\quad \underline{(a\Delta u)^4}
&=\frac{21}{10}\theta_\nu^2-\frac3{10}H
-\frac65u_x\cdot\theta_\nu^3+\frac1{10}u_x^2\cdot\varrho.
\end{aligned}
\]

（將表達式 $a_\vartheta^2(b\theta u)^2$、$a_\vartheta^2(b\theta u)^3$、$a_\vartheta^2(a\theta u)^2(b\theta u)^2$ 和 $a_\vartheta^2(a\theta u)(b\theta u)^3$ 作雙重化簡，再消去 $a_j^2$，也會得到相同結果。）
由公式 (3.) 可知，$(a\Delta u)^2$ 是關於模 $\nu$ 的化簡子。由此得到：

1) $\Delta$ 的系統的化簡：型列 $(\Delta\Delta' u)^2\equiv a_\vartheta^2(a\Delta u)^2$ 含有化簡子 $(a\Delta u)^2$ 作為因子。

2) 由與模 $\Delta$ 作縮並所產生的系統的化簡：

型列 $(\theta\Delta u)^2$ 含有化簡子 $(a\Delta u)^2$ 作為因子。

對於型 $(\theta\Delta u)$ 和 $\Delta_\vartheta$，有：
\[
(a\Delta u)^2(abu)=(\theta\Delta u)-u_x\cdot\Delta_\vartheta(\theta\Delta u),
\]
\[
(a\Delta u)(a\Delta b)=-\frac12\Delta_\vartheta+\frac12u_x\cdot\Delta_\vartheta^2.
\]

而化簡子 $(\theta\Delta u)$ 又給出由與模 $\Delta$ 作縮並所產生的系統的化簡。

\bigskip

因此，模 $\nu$ 下的相對完備系統由下列六個型組成：
\[
f,\ \theta,\ K,\ j,\ \A,\ N.
\]

作為 § 2 中級數展開的例子，我們詳細給出公式 (3.)a 的推導；在後續計算中將直接寫出最終公式 III。（消失型的極式在計算過程中已經略去；每組的第一行給出按展開 II 代入的值，第二行給出從方程組最後一個方程開始遞歸求得的最終值。）

由展開 II 可得：
\[
\begin{aligned}
\text{1)}\quad&m=3.\quad n=4.\quad \lambda=2.\quad \mu=0.\quad \nu=\varkappa=1,\\[0.25em]
\underline{a_\vartheta b_\vartheta(b\theta u)^2}
&=[a_\vartheta]_{\widehat{ub^2}}b+\frac67\{a_\vartheta b_\vartheta(a\theta u)(b\theta u)-u_x\cdot a_\vartheta b_\vartheta(a\theta b)(b\theta u)\}\\
&\quad-\frac17\{a_\vartheta(a\theta u)^2b_\vartheta-2u_x\cdot a_\vartheta(a\theta u)(a\theta b)b_\vartheta+u_x^2\cdot a_\vartheta(a\theta b)^2b_\vartheta\}\\
&=\frac23(a\A u)^2+\frac15[a_\vartheta(a\theta u)^2]b+\frac1{30}f\cdot[a_\vartheta^2(a\theta u)^2]-\frac25u_x\cdot[a_\vartheta(a\theta u)^2]b^2\\
&\quad-\frac1{15}u_x\cdot[a_\vartheta^2(a\theta u)^2]b+\frac15u_x^2\cdot[a_\vartheta(a\theta u)^2]b^3+\frac1{30}u_x^2\cdot[a_\vartheta^2(a\theta u)^2]b^2;\\[0.7em]
\text{2)}\quad&m=2.\quad n=3.\quad \lambda=1.\quad \mu=1.\quad \nu=\varkappa=1,\\[0.25em]
\underline{a_\vartheta(a\theta u)b_\vartheta(b\theta u)}
&=\frac25\{a_\vartheta(a\theta u)^2b_\vartheta-u_x\cdot a_\vartheta(a\theta u)(a\theta b)b_\vartheta\}+\frac12a_\vartheta^2(b\theta u)^2\\
&\quad-\frac15\{a_\vartheta^2(b\theta u)(a\theta u)-u_x\cdot a_\vartheta^2(a\theta b)(b\theta u)\}\\
&=\frac12(a\A u)^2+\frac25[a_\vartheta(a\theta u)^2]b+\frac1{15}[a_\vartheta^2(a\theta u)^2]\cdot f\\
&\quad-\frac25u_x\cdot[a_\vartheta(a\theta u)^2]b^2-\frac1{10}u_x\cdot[a_\vartheta^2(a\theta u)^2]b+\frac1{30}u_x^2\cdot[a_\vartheta^2(a\theta u)^2]b^2;\\[0.7em]
\text{3)}\quad&m=2.\quad n=3.\quad \lambda=1.\quad \mu=1.\quad \nu=1.\quad \varkappa=2,
\end{aligned}
\]
\[
\begin{aligned}
\underline{a_\vartheta(a\theta b)b_\vartheta(b\theta u)}
&=\frac25\{a_\vartheta(a\theta b)(a\theta u)b_\vartheta-u_x\cdot a_\vartheta(a\theta b)^2b_\vartheta\}\\
&=\frac25[a_\vartheta(a\theta u)^2]b^2+\frac1{30}[a_\vartheta^2(a\theta u)^2]b-\frac25u_x\cdot[a_\vartheta(a\theta u)^2]b^3\\
&\quad-\frac1{30}u_x\cdot[a_\vartheta^2(a\theta u)^2]b^2;\\[0.7em]
\text{4)}\quad&m=1.\quad n=2.\quad \lambda=0.\quad \mu=2.\quad \nu=\varkappa=1,\\[0.25em]
\underline{a_\vartheta(a\theta u)^2b_\vartheta}
&=[a_\vartheta(a\theta u)^2]b+\frac23a_\vartheta^2(a\theta u)(b\theta u)\\
&=[a_\vartheta(a\theta u)^2]b+\frac16[a_\vartheta^2(a\theta u)^2]\cdot f-\frac16u_x\cdot[a_\vartheta^2(a\theta u)^2]b;\\[0.7em]
\text{5)}\quad&m=1.\quad n=2.\quad \lambda=0.\quad \mu=2.\quad \nu=1.\quad \varkappa=2.\quad(\varkappa>\nu),\\[0.25em]
\underline{a_\vartheta(a\theta u)(a\theta b)b_\vartheta}
&=[a_\vartheta(a\theta u)^2]b^2+\frac13a_\vartheta^2(a\theta b)(b\theta u)\\
&=[a_\vartheta(a\theta u)^2]b^2+\frac1{12}[a_\vartheta^2(a\theta u)^2]b-\frac1{12}u_x\cdot[a_\vartheta^2(a\theta u)^2]b^2;\\[0.7em]
\text{6)}\quad&m=1.\quad n=2.\quad \lambda=0.\quad \mu=2.\quad \nu=1.\quad \varkappa=3,\\[0.25em]
\underline{a_\vartheta(a\theta b)^2b_\vartheta}&=[a_\vartheta(a\theta u)^2]b^3;\\[0.7em]
\text{7)}\quad&m=2.\quad n=4.\quad \lambda=2.\quad \mu=\nu=\varkappa=0.\quad(\text{按展開 I}),\\[0.25em]
\underline{a_\vartheta^2(b\theta u)^2}&=[a_\vartheta^2]_{\widehat{ub^2}}+\frac1{10}\{[a_\vartheta^2(a\theta u)^2]\cdot f-2u_x\cdot[a_\vartheta^2(a\theta u)^2]b+u_x^2\cdot[a_\vartheta^2(a\theta u)^2]b^2\};\\[0.7em]
\text{8)}\quad&m=1.\quad n=3.\quad \lambda=1.\quad \mu=1.\quad \nu=\varkappa=0.\quad(\text{展開 I}),\\[0.25em]
\underline{a_\vartheta^2(a\theta u)(b\theta u)}&=\frac14[a_\vartheta^2(a\theta u)^2]\cdot f-\frac14u_x\cdot[a_\vartheta^2(a\theta u)^2]b;\\[0.7em]
\text{9)}\quad&m=1.\quad n=3.\quad \lambda=1.\quad \mu=\varkappa=1.\quad \nu=0.\quad(\text{展開 I}),\\[0.25em]
\underline{a_\vartheta^2(a\theta b)(b\theta u)}&=\frac14[a_\vartheta^2(a\theta u)^2]b-\frac14u_x\cdot[a_\vartheta^2(a\theta u)^2]b^2.
\end{aligned}
\]

由 1) 和 4) 得：
\[
\begin{aligned}
(a\A u)^2
&=\frac65[a_\vartheta(a\theta u)^2]b+\frac15f\cdot[a_\vartheta^2(a\theta u)^2]+\frac35u_x\cdot[a_\vartheta(a\theta u)^2]b^2-\frac3{20}u_x\cdot[a_\vartheta^2(a\theta u)^2]b\\
&\quad-\frac3{10}u_x^2\cdot[a_\vartheta(a\theta u)^2]b^3-\frac1{20}u_x^2\cdot[a_\vartheta^2(a\theta u)^2]b^2,\\[0.4em]
(a\A u)^2&=-a_\nu b_\nu+\frac35f\cdot a_\nu^2+\frac45u_x\cdot a_\nu^2b_\nu-\frac3{20}u_x^2\cdot a_\nu^2b_\nu^2-\frac1{12}u_x^2\cdot i\cdot f.
\end{aligned}
\]

（關於改寫為符號 $\theta$，參見 § 8 和 § 9。）
公式 (3.)a 還可通過直接轉用展開 I，並引入輔助變量 $w=x\widehat{ub}$，作更簡短的計算；然而，從公式 (3.)b 和 (3.)c 開始，我們所採用的途徑就更為簡便。對於 (3.)b，由 § 9 可預先知道，一切含因子 $u_x$ 的項都應略去。為計算 (3.)b 和 (3.)c，可得：
\[
\begin{aligned}
(3.)\mathrm{b)}\quad 1)&\quad a_\vartheta(a\theta u)b_\vartheta(b\theta u)^2=\frac12(a\A u)^3+\frac45[a_\vartheta(a\theta u)^2]_{\widehat{ub}}b+\frac{77}{360}[a_\vartheta^2(a\theta u)^2]_{\widehat{ub}},\\
2)&\quad a_\vartheta(a\theta u)b_\vartheta(b\theta u)^2=[a_\vartheta(a\theta u)^2]_{\widehat{ub}}b+\frac{13}{36}[a_\vartheta^2(a\theta u)^2]_{\widehat{ub}};\\[0.5em]
(3.)\mathrm{c)}\quad 1)&\quad a_\vartheta(a\theta u)b_\vartheta(b\theta u)^3=\frac12(a\A u)^4+\frac65[a_\vartheta(a\theta u)^2]_{\widehat{ub^2}}b+\frac{53}{120}[a_\vartheta^2(a\theta u)^2]_{\widehat{ub^2}}\\
&\qquad-\frac65u_x\cdot[a_\vartheta(a\theta u)^2]_{\widehat{ub^2}}b^2-\frac35u_x\cdot[a_\vartheta^2(a\theta u)^2]_{\widehat{ub^2}}b+\frac{19}{120}u_x^2\cdot[a_\vartheta^2(a\theta u)^2]_{\widehat{ub^3}}b^2,\\
2)&\quad a_\vartheta(a\theta u)b_\vartheta(b\theta u)^3=\frac34[a_\vartheta^2(a\theta u)^2]_{\widehat{ub^2}}.
\end{aligned}
\]

補充說明：按 § 5 可化簡的 $f$ 與 $j$ 的轉積也可類似計算。可得
\[
\tag{4.}
\begin{aligned}
a_j&=-\theta_\nu+u_x\cdot\left\{\frac92\theta_\nu^2-\frac12H\right\}-3u_x^2\cdot\theta_\nu^3+\frac12u_x^3\cdot\varrho,\\
a_j^2&=\frac{21}{10}\theta_\nu^2-\frac3{10}H-2u_x\cdot\theta_\nu^3+\frac12u_x^2\cdot\varrho,\\
a_j^3&=-\frac7{10}\theta_\nu^3+\frac12u_x\cdot\varrho,\\
a_j^4&=\frac35\varrho.
\end{aligned}
\]

由 (3.) 和 (4.)，通過從二元領域轉移，可得
\[
(\theta\theta'u)^2=\frac13 j\cdot f\modu{\nu}\srcfn{*)}{戈尔丹--Kerschensteiner，前引书，第 181 页。}\text{。}
\]
型列 $(\theta\theta'u)^2$ 的其餘各型以及型 $\theta'_\vartheta$ 均可化簡為符號 $\nu$。因此可以作如下替換：
\[
\modu{\nu}:\quad \text{將与 } f\cdot j \text{ 所作的縮並替換為与 }(\theta\theta'u)^2\text{ 所作的相應縮並。}
\]
此外有：
\[
(\vartheta\vartheta'x)^2=\frac43 f\cdot\A,\qquad
\theta_{\vartheta'}(\vartheta\vartheta'x)=-N.
\]
\begin{center}
{\Large\bfseries 第三章。\quad 模 $(s,\varrho,t)$ 的相對完備系統。}
\end{center}

\begin{center}
\textbf{\S\ 7.}\quad\emph{系統構造概覽。}
\end{center}

為了從模 $\nu$ 的相對完備系統過渡到模更高一級的相對完備系統，按照 \S\ 3 的定義，須相對於這一較高模構造 $\nu$ 的型系統，並把它同模 $\nu$ 的相對完備系統中的各型作縮並。不過，$\nu=u_\nu^4$ 作為型 $f$ 的四次逆變式，在對偶意義下與 $f$ 相對應。因此，若把 $\nu$ 視為基型，便得到一個由六個型組成的模 $\nu(\nu)=(\nu\nu_1x)^4=s_x^4$ 的相對完備系統。

因此，為構造模 $s$ 的相對完備系統（系統 III），須作如下轉積：
\[
\begin{array}{ll}
\text{系統 I：} & f,\ \theta,\ K,\ j,\ \A,\ N\quad \text{同}\vphantom{\dfrac11}\\[0.3em]
\text{系統 II：} & \nu,\ H=(\nu\nu_1x)^2u_\nu^2u_{\nu_1}^2,\ L=(\nu\eta x)H_x^2u_\nu^3u_\eta^3,\ g=(\nu\eta x)^4H_x^2,\\
& \sigma=H_\nu^2u_\nu^2u_\eta^4,\ (\nu\sigma x).
\end{array}
\]
下文將表明（公式 (13.)），型 $g$ 可化簡，因而系統 II 實際只含五個型。

在計算過程中，將二次型
\[
\varrho=u_\varrho^2=\theta_\nu^4u_\vartheta^2,\qquad t=t_x^2=a_\eta^4H_x^2
\]
作為模加入，從而得到模 $(s,\varrho,t)$ 的相對完備系統；並規定模 $(\varrho,t)$ 高於模 $s$。

各型按其關於係數的階排列。

我們約定縮並
\[
\theta_\eta(K_\eta,\theta_\iota,\ldots)
\]
高於縮並
\[
H_\vartheta(H_\kappa,L_\vartheta,\ldots).
\]
於是依照 \S\ 1 與 \S\ 3b，同階各型的次序便唯一確定。

同階各型的構造按表格進行：

I. 指明系統 I 與系統 II 中哪些型須彼此縮並，才能得到指定的關於係數的階。

II. 按型列方案列出不可化簡型。

III. 化簡可化簡或可分解的型列或型，也就是處理總型列中斷之處，從而證明所列構造已窮盡全部不可化簡構造。

IV. 推論，即：1) 指明新產生的化簡子；2) 指明哪些型在同一系統各型的乘積中不參與同另一系統各型的縮並。

下文將表明，只會出現：

1) 系統 I 的一個型同系統 II 的一個型的縮並。

2) $\theta$（相應地 $H$）的冪，或一個系統中的兩個型，同另一系統中一個型的縮並。

我們將得到如下縮並方案：
\[
\begin{array}{r@{\ }c@{\ }l@{\qquad}c@{\qquad}r@{\ }c@{\ }l}
\multicolumn{3}{c}{\text{系統 I}} & \text{同之縮並} & \multicolumn{3}{c}{\text{系統 II}}\\[0.25em]
1. & \text{階} & f & & 2. & \text{階} & \nu\\
2. & ,, & \theta & & 4. & ,, & H\\
3. & ,, & K,j,\A & & 6. & ,, & L,\sigma\\
4. & ,, & N,\theta^2,f\cdot j & & 8. & ,, & (\nu\sigma x),H^2\\
5. & ,, & \theta\cdot K & & 10. & ,, & H\cdot L\\
6. & ,, & \theta^3,\A\cdot j & & 12. & ,, & H^3.
\end{array}
\]

除\glqq{}雙重化簡\grqq{}外，還用下列兩個特殊型覈對計算：
\[
\begin{array}{ll}
\text{I.}&\begin{array}{rlrlrl}
f&=\displaystyle\sum_3x_1^3x_2,& a_\nu^2&=\dfrac{i}{6}u_x^2,& s&=\dfrac{i}{3}f,\\
a_\eta^2&=-\dfrac19i\theta,& \A_\nu^2&=\dfrac{i}{6}\theta,& \sigma&=-\dfrac{i}{6}j,\\
g&=-\dfrac23i\A,& \varrho&=0,& t&=0,
\end{array}\\[1.2em]
\text{II.}&\begin{array}{rlrlrl}
f&=\displaystyle\sum_3x_1^4,& s&=\dfrac43if,& a_\eta^2&=\dfrac29i\theta,\\
\A_\nu^2&=\dfrac1{15}i\theta,& \sigma&=\dfrac43ij,& g&=\dfrac43i\A,\\
\varrho&=0,& t&=0.
\end{array}
\end{array}
\]

附註：對於基型 $\nu$，與公式 (1.)、(2.)、(3.)、(4.) 相對應的是下列公式：
\[
\tag{5.}
\begin{array}{@{}c|c|c@{}}
\cdot & \underline{H_\nu}=\dfrac13u_x\cdot\sigma & \cdot\\[0.45em]
\underline{(\nu\eta x)^2}=A & \underline{H_\nu(\nu\eta x)}=0 & \underline{H_\nu^2}=\sigma\\[0.45em]
\underline{(\nu\eta x)^3}=0 & \underline{H_\nu(\nu\eta x)^2}=B & \underline{H_\nu^2(\nu\eta x)}=0\\[0.45em]
\underline{(\nu\eta x)^4}=g & \underline{H_\nu(\nu\eta x)^3}=0 & \underline{H_\nu^2(\nu\eta x)^2}=C,
\end{array}
\]
\[
\begin{aligned}
A&=\frac16s\cdot\nu-\frac23u_x\cdot s_\nu+\frac23u_x^2\cdot s_\nu^2-\frac{J}{18}u_x^4,\\
B&=-\frac56s_\nu+\frac76u_x\cdot s_\nu^2-\frac{J}{9}u_x^3,\\
C&=\frac23s_\nu^2-\frac{J}{9}u_x^2.
\end{aligned}
\]
\[
\tag{6.}
s_\nu^3u_\nu s_x=\frac13u_x\cdot s_\nu^4=\frac13J\cdot u_x.
\]

\[
\tag{7.}
\begin{aligned}
\text{a)}\quad (\nu\sigma x)^2
&=\frac12(Hsu)^2-\frac25\nu\cdot s_\nu^2-\frac45u_x\cdot s_\eta(Hsu)^2+u_x^2\left\{\frac16J\cdot\nu-\frac1{40}(ss'u)^4\right\},\\
\text{b)}\quad (\nu\sigma x)^3&=-\frac3{10}s_\eta(Hsu),\\
\text{c)}\quad (\nu\sigma x)^4&=\frac{21}{10}s_\eta^2-\frac3{10}Z-\frac65u_x\cdot s_\eta^3+\frac1{10}u_x^2\cdot s_\eta^4.
\end{aligned}
\]
\[
\tag{8.}
\begin{aligned}
\text{a)}\quad g_\nu&=-s_\eta+u_x\left\{\frac92s_\eta^2-\frac12Z\right\}-3u_x^2\cdot s_\eta^3+\frac12u_x^3\cdot s_\eta^4,\\
\text{b)}\quad g_\nu^2&=\frac{21}{10}s_\eta^2-\frac3{10}Z-2u_x\cdot s_\eta^3+\frac12u_x^2\cdot s_\eta^4,\\
\text{c)}\quad g_\nu^3&=-\frac7{10}s_\eta^3+\frac12u_x\cdot s_\eta^4,\\
\text{d)}\quad g_\nu^4&=\frac35s_\eta^4.
\end{aligned}
\]

\begin{center}
\textbf{\S\ 8.}\quad\emph{三階型（系統 III）。}
\end{center}

\noindent\textbf{$f$ 與 $\nu$ 的縮並。}

\noindent\emph{不可化簡型：}
\[
\begin{array}{cccc}
a_\nu&\cdot&\cdot&\cdot\\
\cdot&a_\nu^2&\cdot&\cdot\\
\cdot&\cdot&\cdot&\cdot\\
\cdot&\cdot&\cdot&a_\nu^4=i.
\end{array}
\]

\noindent\emph{化簡：}
\[
a_\nu^3=\frac13iu_x\qquad\text{（公式 (2.)）。}
\]
\noindent\emph{推論：}1) $a_\nu^3$ 是化簡子。2) $f$ 只有在與逆變型 $(j)$ 成積時才參與同 $\nu$ 的縮並。$\nu$ 相對於 $f$ 亦然。

\begin{center}
\textbf{\S\ 9.}\quad\emph{四階型（系統 III）。}
\end{center}

\noindent\textbf{$\theta$ 與 $\nu$ 的縮並。}

\noindent\emph{不可化簡型：}
\[
\begin{array}{cccc}
(\vartheta\nu x)&\theta_\nu&\cdot&\cdot\\
(\vartheta\nu x)^2&\theta_\nu(\vartheta\nu x)&\theta_\nu^2&\cdot\\
\cdot&\theta_\nu(\vartheta\nu x)^2&\cdot&\theta_\nu^3.
\end{array}
\]

\noindent\emph{化簡：}由化簡子 $a_\nu^3$，按照下列設式對表達式
$a_\nu^3(abu)$、$a_\nu^3b_\nu$、$a_\nu^3b_\nu(abu)$ 作雙重化簡：
\[
\begin{aligned}
\theta_\nu^2(\vartheta\nu x)
&=a_\nu^2(\widehat{ab}\nu x)(abu)-\frac13(\nu\nu_1x)^3
=a_\nu b_\nu(\widehat{ab}\nu x)(abu)+\frac16(\nu\nu_1x)^3\\
&=a_\nu^3(abu)-a_\nu^2b_\nu(abu)=2a_\nu^2b_\nu(abu)=\frac23a_\nu^3(abu),\\[.4em]
\theta_\nu^2(\vartheta\nu x)^2
&=a_\nu^2(\widehat{ab}\nu x)^2-\frac13(\nu\nu_1x)^4
=a_\nu b_\nu(\widehat{ab}\nu x)^2+\frac16(\nu\nu_1x)^4\\
&=i\cdot f-2a_\nu^3b_\nu+a_\nu^2b_\nu^2-\frac13s
=2a_\nu^3b_\nu-2a_\nu^2b_\nu^2+\frac16s
=\frac23i\cdot f-\frac23a_\nu^3b_\nu-\frac16s,\\[.4em]
\theta_\nu^3(\vartheta\nu x)&=a_\nu^2b_\nu(\widehat{ab}\nu x)(abu)=a_\nu^3b_\nu(abu).
\end{aligned}
\]

由此得到：
\[
\tag{9}
\begin{array}{c|c|c}
\underline{\theta_\nu^2(\vartheta\nu x)}=0
&\underline{\theta_\nu^3(\vartheta\nu x)}=0
&\underline{\theta_\nu^4u_\vartheta^2}=u_\varrho^2=\varrho\quad(\text{所加入的模，\S\ 7}),\\[0.6em]
\underline{\theta_\nu^2(\vartheta\nu x)^2}=\dfrac49 i f-\dfrac16s
&&
\end{array}
\]

\noindent\emph{推論：}1) $\theta_\nu^2(\vartheta\nu x)^2$ 是化簡子。
2) $\nu$ 只有在與逆變型 $(g)$ 成積時才參與同 $\theta$ 的縮並。

\begin{center}
\textbf{\S\ 10.}\quad\emph{五階型（系統 III）。}
\end{center}

\[
\text{下列縮並：}\left\{\begin{array}{l}
K,j,\A \text{ 与 }\nu,\\
f\text{ 与 }H.
\end{array}\right.
\]

\noindent\emph{不可化簡型：}
\[
\begin{array}{c@{\qquad}c@{\qquad}c@{\qquad}c}
\text{A)} & \begin{array}{cccc}
\cdot&K_\nu&\cdot&\cdot\\
(k\nu x)^2&\cdot&K_\nu^2&\cdot\\
(k\nu x)^3&K_\nu(k\nu x)^2&\cdot&\cdot\\
\cdot&K_\nu(k\nu x)^3&\cdot&\cdot,
\end{array}
&
\text{B)}\ \begin{array}{c}(j\nu x)\\(j\nu x)^2\\(j\nu x)^3\\ \cdot\end{array}
&
\text{C)}\ \begin{array}{c}\A_\nu\\\cdot\\\cdot\\\cdot\end{array}
\end{array}
\]
\[
\text{D)}\qquad
\begin{array}{cccc}
(aHu)&(aHu)^2&\cdot&\\
a_\eta&a_\eta(aHu)&a_\eta(aHu)^2&\\
\cdot&a_\eta^2&\cdot&\cdot
\end{array}
\]

\noindent\emph{化簡：}

\noindent A) 型列 $K_\nu^3$：化簡子 $a_\nu^3$。

型 $K_\nu^2(k\nu x)^2$：化簡子 $\theta_\nu^2(\vartheta\nu x)^2$。
\[
(k\nu x),\quad K_\nu(k\nu x),\quad K_\nu^2(k\nu x)
\]
依 \S\ 3 均為可分解型。

\noindent B) 型
\[
\begin{aligned}
(j\nu x)^4&\equiv(\widehat{a\theta}\nu x)^4
=a_\nu^4\cdot\theta+\theta_\nu^4\cdot f-4a_\nu^3\theta_\nu
-4a_\nu\{a_\nu\cdot\theta_x-(\widehat{a\theta}\nu x)\}^3\\
&\quad+6a_\nu^2\{a_\nu\theta_x-(\widehat{a\theta}\nu x)\}^2
=\R\cdot f+\frac53 i\theta-6a_\nu^2(\widehat{a\theta}\nu x)^2
+4a_\nu(\widehat{a\theta}\nu x)^3,
\end{aligned}
\]
因而
\[
(j\nu x)^4=\R\cdot f+\frac53 i\theta\cdot\modu{(\A,\nu)},\qquad (\text{參見 \S\ 5 末尾}).
\]

\noindent C) 型 $\A_\nu^4$：化簡子 $\theta_\nu^4$。

型 $\A_\nu^2$ 與 $\A_\nu^3$：將型 $\A$ 替換為型列中的較低型（\S\ 5 末尾），亦即對下列表達式作雙重化簡，便可由化簡子 $a_\nu^3$ 或 $\theta_\nu^4$ 得到：
\[
a_\vartheta a_\nu^3,\qquad
a_\vartheta a_\nu^3\theta_\nu,\qquad
a_\vartheta\theta_\nu^4.
\]
對 $\A_\nu^3$ 的雙重計算引出較高型 $a_\eta^3$ 的化簡。

\noindent\emph{化簡公式：}
\[
\tag{10}
\begin{aligned}
\text{a)}\quad \A_\nu^2
&=\frac35a_\eta^2-\frac3{10}(asu)^2+\frac13 i\theta
+\frac15u_x^2\cdot(2a_\R^2-t),\\
\text{b)}\quad \A_\nu^3
&=-\frac3{10}a_\R+\frac35u_x\left\{\frac{13}{10}a_\R^2-\frac25t\right\},\\
\text{c)}\quad \A_\nu^4&=a_\R^2-\frac25t.
\end{aligned}
\]

\noindent\emph{公式的推導：}
\[
\begin{aligned}
\text{a)}\quad
 a_\vartheta a_\nu^3
&=\frac13 i\theta
=[a_\vartheta]_{\nu^3}+\frac67[a_\vartheta^2]_{\nu^2}
+\frac65[a_\vartheta(a\theta u)^2]_{\nu}\,\widehat{\nu x}^2
+\frac{11}{30}[a_\vartheta^2(a\theta u)^2]\,\widehat{\nu x}^2\\
&\quad-\frac67u_x\cdot[a_\vartheta^2]_\nu^3
-\frac16u_x\cdot[a_\vartheta^2(a\theta u)^2]_{\nu}\,\widehat{\nu x}^2,\\
\frac13 i\theta
&=\A_\nu^2-\frac23u_x\cdot\A_\nu^3-a_\nu a_{\nu_1}(\nu\nu_1x)^2
+\frac25a_\nu^2(\nu\nu_1x)^2
+\frac15u_x\cdot a_\nu^2a_{\nu_1}(\nu\nu_1x)^2;
\end{aligned}
\]
\[
\begin{aligned}
\text{b)}\quad
 a_\vartheta a_\nu^3\theta_\nu
&=0=[a_\vartheta]_{\nu^4}+\frac9{14}[a_\vartheta^2]_{\nu^3}
+\frac35[a_\vartheta(a\theta u)^2]_{\nu^2}\,\widehat{\nu x}^2
+\frac{17}{120}[a_\vartheta^2(a\theta u)^2]_{\nu}\,\widehat{\nu x}^2\\
&\quad-\frac9{14}u_x\cdot[a_\vartheta^2]_{\nu^4}
-\frac1{24}u_x\cdot[a_\vartheta^2(a\theta u)^2]_{\nu^2}\,\widehat{\nu x}^2,\\
0&=\frac56\A_\nu^3-\frac12u_x\cdot\A_\nu^4-\frac14a_\eta^3+\frac1{20}u_x\cdot a_\eta^4;
\end{aligned}
\]
\[
\begin{aligned}
\text{b')}\quad
 a_\vartheta\theta_\nu^4
&=a_\R=a_\vartheta\theta_\nu\{a_\nu-(a\theta\widehat{\nu x})\}^3\\
&=-\frac32[a_\vartheta^2]_{\nu^3}
+\frac35[a_\vartheta(a\theta u)^2]_{\nu^2}\,\widehat{\nu x}^2
-\frac{37}{120}[a_\vartheta^2(a\theta u)^2]_{\nu}\,\widehat{\nu x}^2\\
&\quad+\frac32u_x\cdot[a_\vartheta^2]_{\nu^4}
+\frac{49}{120}u_x\cdot[a_\vartheta^2(a\theta u)^2]_{\nu^2}\,\widehat{\nu x}^2,\\
a_\R&=-\frac32\A_\nu^3+\frac32u_x\cdot\A_\nu^4-\frac{11}{20}a_\eta^3+\frac7{20}u_x\cdot a_\eta^4;
\end{aligned}
\]
\[
\begin{aligned}
\text{c)}\quad
 a_\vartheta^2\theta_\nu^4
&=a_\R^2=[a_\vartheta^2]_{\nu^4}+\frac35[a_\vartheta^2(a\theta u)^2]_{\nu^2}\,\widehat{\nu x}^2,\\
a_\R^2&=\A_\nu^4+\frac25a_\eta^4.
\end{aligned}
\]

\noindent D) 對 $\A_\nu^3$（C）作雙重化簡，便得到 $a_\eta^3$ 的化簡。

其餘化簡由與 \S\ 9 中的計算對偶相反的計算得到。

\noindent\emph{化簡公式：}
\[
\tag{11}
\begin{array}{rl|rl}
\underline{a_\eta^2(aHu)}&=-\dfrac16(asu)^3,&
\underline{a_\eta^2(aHu)^2}&=\dfrac49 i \nu-\dfrac16(asu)^4,\\[0.7em]
\underline{a_\eta^3}&=-a_\R+\dfrac35u_x\cdot a_\R^2+\dfrac15u_x\cdot t,
&\underline{a_\eta^3(aHu)}&=0,\\[0.7em]
&&\underline{a_\eta^4}&=t\quad(\text{作為模加入}).
\end{array}
\]

\noindent\emph{推論：}

1) $(j\nu x)^4$、$\A_\nu^2$、$a_\eta^2(aHu)$ 是化簡子。

2) $\nu$ 只有在與逆變型成積時才參與同 $K,j,\A$ 的縮並；$f$ 相對於 $H$ 亦然，$j$ 相對於 $\nu$ 亦然，因為依 \S\ 11 B，$\theta_\nu(\vartheta\cdot j,\nu,x)^3$ 可化簡。

\begin{center}
\textbf{\S\ 11.\quad 六階型（系統 III）。}
\end{center}
\[
\text{下列縮並：}\left\{\begin{array}{l}
\theta^2,\ f\cdot j,\ N \text{ 与 }\nu,\\
\theta\text{ 与 }H.
\end{array}\right.
\]


\noindent\emph{不可化簡型：}
\[
\renewcommand{\arraystretch}{1.25}
\begin{array}{@{}r@{\quad}c@{\qquad\qquad}r@{\quad}c@{\qquad\qquad}r@{\quad}c@{}}
\text{A)}&(\vartheta^2\nu x)^3&\text{B)}&a_\nu(j\nu x)&\text{C)}&N_\nu,\\[0.45em]
&\cdot&&\cdot&&\\[-0.1em]
&\theta_\nu(\vartheta^2\nu x)^3,&&a_\nu^2(j\nu x),&&
\end{array}
\]
\[
\renewcommand{\arraystretch}{1.25}
\begin{array}{@{}r@{\quad}c|c|c|c@{}}
&\cdot&(\theta Hu)&(\theta Hu)^2&\cdot\\
\text{D)}&(\vartheta\eta x)&\theta_\eta&H_\vartheta(\theta Hu),\ \theta_\eta(\theta Hu)&\theta_\eta(\theta Hu)^2\\
&(\vartheta\eta x)^2&H_\vartheta(\vartheta\eta x),\ \theta_\eta(\vartheta\eta x)&\theta_\eta^2&\theta_\eta^2(\theta Hu)\\
&\cdot&\theta_\eta(\vartheta\eta x)^2&\cdot&\cdot
\end{array}
\]
\noindent\emph{化簡：}

\noindent A) $(\vartheta^2\nu x)^4$ 依公式 (3.)a 分解：
\[
(a\A\widehat{\vartheta}x)^2=\frac12(\vartheta^2\nu x)^4-\frac35 f\cdot a_\nu^2(\nu\vartheta x)^2
=\frac13\A^2+f\cdot\A_\vartheta^2.
\]

\noindent B) $a_\nu(j\nu x)^2$ 依公式 (9.)a 分解：按照 \S\ 6 末尾，以 $(\theta\theta'u)^2$ 替換 $f\cdot j$，並注意 $\theta_\vartheta'$ 的可化簡性：
\[
\frac13a_\nu(j\nu x)^2=(\widehat{\theta\theta'}\nu x)^2\theta_\nu
=\theta_\nu^3\cdot\theta-\theta_\nu^2\theta_\nu'
=\theta_\nu^3\cdot\theta+\theta_\nu^2(\widehat{\vartheta\nu}\theta'u)
=\theta_\nu^3\cdot\theta-\frac23\theta_\nu^3\cdot\theta.
\]
型 $a_\nu(j\nu x)^3$ 與 $a_\nu^2(j\nu x)^2$：化簡子為 $a_j$ 與 $(j\nu x)^4$：
\[
a_\nu(j\nu x)^3=a_j(j\nu x)^3-(j\nu x)^3(j\nu\widehat{au}),
\]
\[
a_\nu^2(j\nu x)^2=a_j^2(j\nu x)^2-2a_j(j\nu x)^2(j\nu\widehat{au})+(j\nu x)^2(j\nu\widehat{au})^2.
\]

\noindent C) 型列 $N_\nu^2$；化簡子為 $\A_\nu^2$。
依 \S\ 3，$(n\nu x)$ 與 $N_\nu(n\nu x)$ 作為經由函數行列式的轉積而分解。

\noindent D) \emph{化簡概覽：}

a) 型列 $\theta_\eta^2H_\vartheta$：化簡子 $a_\eta^2(aHu)$。（導向含較高符號的型。）

b) 型列 $\theta_\eta H_\vartheta$：通過雙重化簡，以 $a_\eta^2(aHu)$ 為化簡子。
（導向含較高符號的型以及總型列中的較高型，即型列 $\theta_\eta^2$。）我們用下列型列代替 $\theta_\eta H_\vartheta$ 加以考察：
\[
\theta_\eta(\vartheta\eta x)(\theta Hu)=\theta_\eta H_\vartheta+\theta_\eta^2,
\]
將其中的符號 $\theta$ 替換為較低型 $(abu)$，便得到型列
\[
a_\eta^2(aHu)(abu)=\theta_\eta H_\vartheta+\frac32\theta_\eta^2\modu{s}
\]
的雙重化簡，一直進行到末型
\[
a_\eta^3b_\eta(bHu)(abu)=\theta_\eta^3H_\vartheta+\frac12\theta_\eta^4.
\]

c) 型列 $\theta_\eta^3$：同時屬於型列 $\theta_\eta H_\vartheta$ 與 $\theta_\eta^2H_\vartheta$ 的那些型——也就是型列 $\theta_\eta^2H_\vartheta$ 中的型——一方面化為含較高符號的型，另一方面化為含較高符號的型以及較高型列 $\theta_\eta^3$；對這些型作雙重化簡。

d) 型 $\theta_\eta^2(\vartheta\eta x)$、$\theta_\eta^2(\vartheta\eta x)^2$、$\theta_\eta^3(\vartheta\eta x)$：以 $a$ 為化簡子，仿照 \S\ 9 的計算作雙重化簡。

e) 型 $\theta_\eta^2(\theta Hu)^2$ 與 $\theta_\eta^2(\vartheta\eta x)^2$：分別以 $\A_\nu^2$ 與 $(a\A u)^4$ 為化簡子，對表達式 $\A_\nu^2(a\A u)^4$ 與 $(a\A\widehat{\nu}x)^4$ 作雙重化簡。

f) 型 $H_\vartheta$ 與 $H_\vartheta^2$：可分解型。對 $H_\vartheta^2$，這一點可由表達式 $\A_\nu^2(a\A u)^2$ 的雙重化簡得到，也可用系統中的各型直接計算乘積 $(a_\nu^2)^2$、$a_\nu^2\cdot b_{\nu_1}$ 得到。（類似地計算 $(a_\nu)^2$，會給出可分解型之間的一個關係。）

g) 較高型之所以能夠化簡，是因為型列 $\theta\cdot H$ 中的型容許兩種化簡（即屬於兩個可化簡型列），具體為：
\[
\begin{array}{rl}
g&\text{由 }\theta_\eta^2(\vartheta\eta x)^2\text{ 的雙重化簡：可採用化簡 d) 与 e)；}\\[0.3em]
s_\vartheta^2(\theta s u)&\text{由 }\theta_\eta^3(\vartheta\eta x)\text{ 的雙重化簡：可採用化簡 c) 与 d)；}\\[0.3em]
s_\vartheta^2(\theta s u)^2&\text{由 }\theta_\eta^2H_\vartheta^2\text{ 的雙重化簡：可採用化簡 a)，並且}\\
&\text{以化簡子 }\theta_\nu^2(\vartheta\nu x)^2\text{ 為因子；}\\[0.3em]
s_\nu^2&\text{由 }\theta_\eta^3H_\vartheta\text{ 的雙重化簡：可採用化簡 a) 与 c)。}
\end{array}
\]
其餘各型的獨立性可由型列 $\A_\nu^2(a\A u)^2\equiv\theta_\eta^2$ 的雙重化簡證明。

\noindent\emph{化簡的實施：}

化簡 a)、b)、c)、d) 的存在性先驗即明，因為按照上述構造法則，雙重化簡所產生的關係彼此獨立。對於化簡 e)、f)、g)，則須通過計算證明不會產生恆等式。

我們以型列 $\theta\cdot H$ 的對角項為對稱軸，推進到型 $\theta_\eta$ 為止；同時也列出由化簡 a)、b)、c)、d) 得到的公式，並對部分計算略作提示，因為這些公式一部分用於化簡 e)、f)、g)，一部分將在後文使用。

\noindent 1) 按 f) 處理 $H_\vartheta$ 與 $H_\vartheta^2$。按照設式\srcfn{*)}{以符號 $\sim$ 代替等號，表示省略 $u_x$ 與比縮並 III、IV 所產生的型更高的型的乘積。}
\[
\begin{aligned}
a_\nu\cdot b_{\nu_1}^2
&=\nu\cdot a_\nu b_\nu^2-(aHu)(bHu)b_\eta
\sim \nu\cdot a_\nu b_\nu^2-f\cdot a_\eta(aHu)^2-(abu)(aHu)b_\eta+(abu)^2b_\eta,\\
a_\nu^2\cdot b_{\nu_1}^2
&=b_{\nu_1}^2\{a_{\nu_1}u_\nu-(a\widehat{\nu\nu_1}u)\}^2
\sim \nu\cdot a_\nu^2b_\nu^2-2a_\eta b_\eta(aHu)(bHu)-\frac16(asu)^2(bsu)^2\\
&\sim \nu\cdot a_\nu^2b_\nu^2-2a_\eta^2(aHu)(abu)-\frac16(asu)^2(bsu)^2+(\vartheta\eta x)^2(\theta Hu)^2
\end{aligned}
\]
代入 $\theta_\eta H_\vartheta$ 的值，便得到：
\[
\tag{12}
\begin{aligned}
\text{a)}\quad H_\vartheta&\sim 2a_\nu\cdot a_\nu^2+2\nu\cdot\theta_\nu(\vartheta\nu x)^2+2f\cdot a_\eta(aHu)^2-2\theta_\eta,\\
\text{b)}\quad H_\vartheta^2&\sim (a_\nu^2)^2+2\theta_\eta^2+\frac23(\theta su)^2+\frac1{12}s\cdot\nu-\frac19 i f\cdot\nu-\frac16f\cdot(asu)^4.
\end{aligned}
\]

\noindent 2) 按 b) 處理 $\theta_\eta H_\vartheta$：
\[
\begin{aligned}
a_\eta^2(aHu)(abu)
&\sim [a_\eta^2(aHu)]_{\widehat{bu}}+\frac12 f\cdot[a_\eta^2(aHu)^2]
=a_\eta b_\eta(aHu)(abu)\\
&\quad+a_\eta(\widehat{ab}\eta x)(abu)(aHu)\sim\theta_\eta(\vartheta\eta x)(\theta Hu)+\frac12\theta_\eta^2+\frac14(\nu\eta x)^2.
\end{aligned}
\]
由公式 (5.) 與 (11.) 得：
\[
\theta_\eta H_\vartheta\sim-\frac32\theta_\eta^2-\frac14(\theta su)^2-\frac1{12}s\cdot\nu+\frac29 i f\cdot\nu.
\]

\noindent 3) 按 b)，$\theta_\eta H_\vartheta(\theta Hu)$ 的計算與 $\theta_\eta H_\vartheta$ 類似：
\[
\theta_\eta H_\vartheta(\theta Hu)\sim-\theta_\eta^2(\theta Hu)-\frac16(\theta su)^3.
\]

\noindent 4) 按 b) 處理 $\theta_\eta H_\vartheta(\vartheta\eta x)$；按 d) 處理 $\theta_\eta^2(\vartheta\eta x)$（參見 \S\ 9）：
\[
\theta_\eta H_\vartheta(\vartheta\eta x)\sim-\theta_\eta^2(\vartheta\eta x)-\frac16s_\vartheta(\theta su)
\sim\frac13(\vartheta\R x)-\frac16s_\vartheta(\theta su),
\]
\[
\theta_\eta^2(\vartheta\eta x)\sim\frac23a_\eta^3(abu)\sim-\frac13(\vartheta\R x).
\]


\noindent 5)
\[
\begin{array}{ccc}
\theta_\eta H_\vartheta^2\text{ 按 b)},&
\theta_\eta^2H_\vartheta\text{ 按 a)},&
\theta_\eta^2H_\vartheta\text{ 按 b)，並由此按 c) 處理 }\theta_\eta^3,\\
\text{由 }a_\eta^2(aHb)(bHu)\text{；}&
\text{由 }a_\eta^2(Hab)(abu)\text{；}&
\text{由 }a_\eta^3(bHu)(abu)\text{：}
\end{array}
\]
\[
\theta_\eta H_\vartheta^2\sim-\frac32\theta_\eta^2H_\vartheta-\frac14s_\vartheta(\theta su)^2+\frac16s_\nu-\frac49ia_\nu
\sim-\theta_\varrho-\frac16s_\nu-\frac19ia_\nu,
\]
\[
\theta_\eta^2H_\vartheta\sim\frac23\theta_\varrho-\frac16s_\vartheta(\theta su)^2+\frac29s_\nu-\frac29ia_\nu,
\]
\[
\theta_\eta^2H_\vartheta\sim\frac23\theta_\varrho-\frac23\theta_\eta^3+\frac5{36}s_\nu\quad\text{並由此得到}
\]
\[
\theta_\eta^3\sim\frac14s_\vartheta(\theta su)^2-\frac18s_\nu+\frac13i\cdot a_\nu.
\]

\noindent 6) 按 e) 處理 $\theta_\eta^2(\theta Hu)^2$：
\[
\A_\nu^2(a\A u)^4=[(a\A u)^4]_{\nu^2}=-\frac{12}{5}\theta_\eta^2(\theta Hu)^2-\frac3{10}\sigma-\frac1{20}(\theta su)+\nu\cdot\varrho
\quad(\text{依公式 (3.) c}),
\]
\[
\A_\nu^2(a\A u)^4=[\A_\nu^2]^4_{\widehat{u a}}=\frac35\theta_\eta^2(\theta Hu)^2+\frac15\sigma-\frac3{10}(\theta su)^4+\frac13ij
\quad(\text{依公式 (10.) a}),
\]
\[
\theta_\eta^2(\theta Hu)^2=-\frac16\sigma+\frac1{12}(\theta su)^4-\frac19ij+\frac13\nu\cdot\varrho.
\]

\noindent 7) 按 d) 與 e) 處理 $\theta_\eta^2(\vartheta\eta x)^2$。由此化簡較高型 $g$。

仿照 \S\ 9 的計算可得：
\[
\begin{aligned}
\theta_\eta^2(\vartheta\eta x)^2
&=\frac12t\cdot f-\frac12a_\eta^2b_\eta-\frac1{12}g=\frac23t\cdot f-\frac23a_\eta^3b_\eta-\frac16g\\
&=-\frac16g-\frac13(\vartheta\varrho x)^2+\frac4{15}f\cdot a_\varrho^2+\frac8{15}f\cdot t\quad(\text{依公式 (11.)}).
\end{aligned}
\]
作與上述相應的計算可得：
\[
(a\A\widehat{\nu x})^4=[(a\A u)^4]\widehat{\nu x}^4=\frac{21}{10}\theta_\eta^2(\vartheta\eta x)^2+\frac7{10}s_\vartheta^2-\frac3{10}g
\quad(\text{依公式 (3.) c}),
\]
\[
\begin{aligned}
(a\A\widehat{\nu x})^4&=-\frac13i\A+6[\A_\nu^2]a^2-4[\A_\nu^3]a+\A_\nu^4\cdot f\\
&=-\frac{36}{5}\theta_\eta^2(\vartheta\eta x)^2-\frac95s_\vartheta^2-\frac35g-\frac35(\vartheta\varrho x)^2\\
&\quad+\frac53i\A+\frac{37}{25}f\cdot a_\varrho^2+\frac{74}{25}f\cdot t,\quad(\text{依公式 (10.)}).
\end{aligned}
\]
因此
\[
\begin{aligned}
\frac{93}{10}\theta_\eta^2(\vartheta\eta x)^2
&=-\frac3{10}g-\frac52s_\vartheta^2-\frac35(\vartheta\varrho x)^2\\
&\quad+\frac53i\A+\frac{37}{25}f\cdot a_\varrho^2+\frac{74}{25}f\cdot t,
\end{aligned}
\]
再依 g) 比較 $\theta_\eta^2(\vartheta\eta x)^2$ 的兩個值，得到：
\[
\tag{13.)\srcfnmark{*)}}
0=g-2s_\vartheta^2+2(\vartheta\varrho x)^2+\frac43i\A-\frac45f\cdot a_\varrho^2-\frac85f\cdot t.
\]
\srcfntext{*)}{這就是引言註釋中提到的三個六階六次協變式之間的關係。}

\noindent 8) 按 a) 與 b) 處理 $\theta_\eta^2H_\vartheta(\theta Hu)$，由此再按 c) 得到 $\theta_\eta^3(\theta Hu)$（含因子 $u_x$ 的型消失）：
\[
\theta_\eta^2H_\vartheta(\theta Hu)=-\frac16s_\vartheta(\theta su)^3-\frac13(\nu\varrho x),
\]
\[
\theta_\eta^2H_\vartheta(\theta Hu)=-\frac13\theta_\eta^3(\theta Hu)-\frac13(\nu\varrho x),
\]
\[
\theta_\eta^3(\theta Hu)=\frac12s_\vartheta(\theta su)^3.
\]

\noindent 9) 按 a) 與 b) 處理 $\theta_\eta^2H_\vartheta(\vartheta\eta x)$，由此按 c) 得到 $\theta_\eta^3(\vartheta\eta x)$。再按 d) 處理 $\theta_\eta^3(\vartheta\eta x)$，由此按 g) 化簡較高型 $s_\vartheta^2(\theta su)$（含因子 $u_x$ 的型消失）：
\[
\theta_\eta^2H_\vartheta(\vartheta\eta x)=\frac13(atu),
\]
\[
\theta_\eta^2H_\vartheta(\vartheta\eta x)=\frac13(atu)+\frac13\theta_\eta^3(\vartheta\eta x)-\frac16s_\vartheta^2(\theta su),
\]
\[
\theta_\eta^3(\vartheta\eta x)=\frac12s_\vartheta^2(\theta su),
\]
\[
\theta_\eta^3(\vartheta\eta x)=\frac25\theta_\varrho(\vartheta\varrho x)+\frac15(atu):
\]
\srcnumdisplay{(14.)}{%
s_\vartheta^2(\theta su)=\frac45\theta_\varrho(\vartheta\varrho x)+\frac25(atu).
}

\noindent 10) 按 a) 並藉助化簡子 $\theta_\nu^2(\vartheta\nu x)^2$ 處理 $\theta_\eta^2H_\vartheta^2$；按 a) 與 c) 處理 $\theta_\eta^3H_\vartheta$；按 c) 處理 $\theta_\eta^4$。由此按 g) 化簡較高型 $s_\vartheta^2(\theta su)^2$ 與 $s_\nu^2$：
\[
\begin{aligned}
\text{按 a)}\quad
\theta_n^2H_\vartheta^2&=[a_\eta^2(aHu)^2]b^2+\frac13[a_\eta^4]_{\widehat{bu^2}}-\frac13H_\nu^2(\nu\eta x)^2\\
&=\frac49i\cdot a_\nu^2-\frac16s_\vartheta^2(\theta su)^2-\frac5{18}s_\nu^2+\frac13(atu)^2+\frac1{54}J\cdot u_x^2.
\end{aligned}
\]
\[
\begin{aligned}
\text{由 }\theta_\nu^2(\vartheta\nu x)^2\quad
\theta_\eta^2H_\vartheta^2
&=[\theta_\nu^2(\vartheta\nu x)^2]_{\nu_1^2}+\frac13[\theta_\nu^4]\widehat{\nu_1}x^2-\frac13s_\vartheta^2(\theta su)^2\\
&=\frac49i\cdot a_\nu^2-\frac16s_\nu^2-\frac13s_\vartheta^2(\theta su)^2+\frac13(\nu\varrho x)^2.
\end{aligned}
\]
\[
\begin{aligned}
\text{按 a)}\quad
\theta_\eta^3H_\vartheta&=-[a_\eta^2(aHu)]_{\widehat{bu}}b^2-\frac12[a_\eta^2(aHu)^2]b^2-\frac13[a_\eta^3]_{\widehat{bu^2}}b+\frac1{18}[a_\eta^4]_{\widehat{bu^2}}\\
&\quad+\frac12u_x\cdot[a_\eta^2(aHu)^2]b^3\\
&=-\frac29ia_\nu^2+\frac1{12}s_\nu^2+\frac4{15}\theta_\varrho^2-\frac7{90}(\nu\varrho x)^2+\frac1{30}(atu)^2+\frac2{27}i^2\cdot u_x^2-\frac1{36}J\cdot u_x^2.
\end{aligned}
\]
\[
\begin{aligned}
\text{按 c)}\quad
\theta_\eta^3H_\vartheta:\quad a_\eta^3(Hab)(bHu)&=\frac12(atu)^2
=\theta_\eta^2H_\vartheta(\theta Hu)(\vartheta\eta x)+\frac14H_\nu^2(\nu\eta x)^2\\
&\quad+\frac12\theta_\eta^2H_\vartheta^2
=\frac32\theta_\eta^2H_\vartheta^2+\theta_\eta^3H_\vartheta-u_x\cdot[a_\eta^2(aHu)^2]b^3\\
&\quad+\frac14H_\nu^2(\nu\eta x)^2,
\end{aligned}
\]
因而
\[
\begin{aligned}
\theta_\eta^3H_\vartheta&=-\frac23i\cdot a_\nu^2+\frac5{12}s_\nu^2+\frac12(\nu\varrho x)^2-\frac12(atu)^2\\
&\quad+\frac4{27}i^2\cdot u_x^2-\frac1{12}J\cdot u_x^2.
\end{aligned}
\]
按 c)
\[
\theta_\eta^4=\frac49ia_\nu^2-\frac16s_\nu^2+\frac{16}{15}\theta_\varrho^2-\frac{14}{45}(\nu\varrho x)^2+\frac2{15}(atu)^2.
\]
按 g)，由 $\theta_\eta^2H_\vartheta^2$ 的兩個值得到：
\par\smallskip\noindent
\begin{minipage}[c]{3.6em}
\centering (15.)
\end{minipage}%
\begin{minipage}[c]{\dimexpr\linewidth-3.6em\relax}
\centering
\(\displaystyle
\begin{aligned}
s_\vartheta^2(\theta su)^2
&=\frac23s_\nu^2-2(atu)^2+2(\nu\varrho x)^2-\frac19J\cdot u_x^2\\
&=\frac89i\cdot a_\nu^2+\frac8{15}\theta_\varrho^2-\frac{14}{15}(atu)^2+\frac{38}{45}(\nu\varrho x)^2-\frac4{27}i^2\cdot u_x^2.
\end{aligned}
\)
\end{minipage}
\par\smallskip
按 g)，由 $\theta_\eta^3H_\vartheta$ 的兩個值得到：
\par\smallskip\noindent
\begin{minipage}[c]{3.6em}
\centering (16.)
\end{minipage}%
\begin{minipage}[c]{\dimexpr\linewidth-3.6em\relax}
\centering
\(\displaystyle
s_\nu^2=\frac43i\cdot a_\nu^2+\frac45\theta_\varrho^2+\frac85(atu)^2-\frac{26}{15}(\nu\varrho x)^2+\frac16J\cdot u_x^2-\frac29i^2\cdot u_x^2.
\)
\end{minipage}
\par\smallskip
公式 (16.) 是化簡模 $(ss'u)^4$ 的基礎，公式 (13.) 是化簡 $s$ 的型系統的基礎。（\S\ 17。）

\noindent\emph{推論：}

1) $a_\nu(j\nu x)^3$、$\theta_\eta H_\vartheta$、$\theta_\eta^2(\vartheta\eta x)$、$\theta_\eta^2(\theta Hu)^2$ 是化簡子；$a_\nu(j\nu x)^2$、$H_\vartheta$ 與 $H_\vartheta^2$ 是可分解型。

2) 系統 II 中的型 $j(\nu)=g=(\nu\eta x)^4H_x^2$ 可化簡。

3) 由於唯一的逆變型 $g$ 已成為可化簡型，$\nu$ 在同系統 I 縮並時，完全不以冪或與本系統各型的乘積形式出現。

$\theta$ 只有在作為逆變式看待時才以乘積或冪的形式參與縮並；相應地，$H$ 則作為協變式看待（參見 \S\ 13 B）。

\begin{center}
\textbf{\S\ 12.}\quad\emph{七階型。}
\end{center}
\[
\text{下列縮並：}\left\{\begin{array}{l}
\theta\cdot K\text{ 与 }\nu,\\
K,j,\A\text{ 与 }H,\\
f\text{ 与 }L,\sigma.
\end{array}\right.
\]


\noindent\emph{不可化簡型：}
\[
\renewcommand{\arraystretch}{1.35}
\begin{array}{@{}r@{\quad}c|c|c|c@{}}
&\cdot&\cdot&(KHu)^2&\cdot\\
&\cdot&K_\eta&\cdot&K_\eta(KHu)^2\\
\text{B)}&(k\eta x)^2&\cdot&K_\eta^2&\cdot\\
&(k\eta x)^3&H_k(k\eta x)^2,\ K_\eta(k\eta x)^2&\cdot&\cdot\\
&\cdot&K_\eta(k\eta x)^3&\cdot&\cdot
\end{array}\, ,
\]
\[
\renewcommand{\arraystretch}{1.35}
\begin{array}{@{}r@{\quad}c@{\qquad}c@{\qquad}c@{\qquad\qquad}r@{\quad}c@{}}
\text{C)}&(j\eta x)&H_j&\cdot&\text{D)}&(\A Hu)\\
&\cdot&H_j(j\eta x)&H_j^2,&&\A_\eta,
\end{array}
\]
\[
\renewcommand{\arraystretch}{1.35}
\begin{array}{@{}r@{\quad}c@{\qquad}c@{\qquad}c@{\qquad}c@{}}
&\cdot&(aLu)^2&(aLu)^3&\cdot\\
\text{E)}&a_l&\cdot&a_l(aLu)^2&a_l(aLu)^3\\
&\cdot&a_l^2&\cdot&\cdot
\end{array}\, .
\]
\noindent\emph{化簡：}

\noindent A) $(\vartheta\cdot k,\nu,x)^4$ 作為對可分解型 $(\vartheta^2\nu x)^4$ 的一次轉積而分解。

\noindent B) 型列 $H_kK_\eta$：化簡子 $H_\vartheta\theta_\eta$。

型列 $K_\eta^2(KHu)$：化簡子 $a_\eta^2(aHu)$。

型列 $K_\eta^2(k\eta x)$：化簡子 $\theta_\eta^2(\vartheta\eta x)$。

由此對型列 $K_\eta^3$ 作雙重化簡，又引出較高型列 $(j\eta x)^3$ 的化簡。

依 \S\ 3，$(KHu)$、$(k\eta x)$、$K_\eta(k\eta x)$ 是可分解型。

$H_k(KHu)$、$K_\eta(KHu)$ 是通過同可分解型 $K_\nu(k\nu x)$ 及函數行列式作一次縮並而產生的，因而可分解：
\[
K_\nu^2\equiv\theta_\nu^2(a\theta u)\equiv a_\nu^2(a\theta u)\modu{\nu^2}.
\]
$K_\nu^2$ 的雙重表示對應於可分解型之間的一個關係。可得：
\[
H_k(KHu)=2K_{\nu_1}(\nu\nu_1k)=2\theta_{\nu_1}(\nu\nu_1\widehat{a\theta})
=2\theta_\nu^2\cdot a_\nu-a_\nu^2\cdot\theta_\nu+f\cdot\theta_\eta(\theta Hu)^2-f\cdot \nu\cdot\theta_\nu^3\quad(\text{乘積定理})\modu{(\nu^3)},
\]
\[
\begin{aligned}
K_\eta(KHu)&=-K_\nu^2(\nu\nu_1x)=-a_\nu^2(a\theta u)(\nu\nu_1x)=-\theta_\nu^2(a\theta u)(\nu\nu_1x)\\
&=a_\eta(aHu)^2\cdot\theta+a_\nu^2\cdot\theta_\nu
=-\theta_\eta(\theta Hu)^2\cdot f-\theta_\nu^2\cdot a_\nu:
\end{aligned}
\]
\par\smallskip\noindent
\begin{minipage}[c]{4.8em}
\centering (17.)\srcfnmark{*)}
\end{minipage}%
\begin{minipage}[c]{\dimexpr\linewidth-4.8em\relax}
\centering
\(\displaystyle
0=a_\nu^2\cdot\theta_{\nu_1}+\theta_\nu^2\cdot a_{\nu_1}+f\cdot\theta_\eta(\theta Hu)^2+\theta\cdot a_\eta(aHu)^2\modu{(\nu^3)}.
\)
\end{minipage}
\par\smallskip
\srcfntext{*)}{在可分解型之間的所有關係中，只要型的乘積中含有一個逆變式——例如這裡的 $f\cdot\theta_\nu^3\cdot \nu$——就將其省略，因為在把這些公式應用於型系統 $s_{IV}$ 時，它們不起作用。}
型 $H_k$、$H_k(k\eta x)$、$H_k^2$、$H_k^2(k\eta x)$ 是通過同可分解型 $H_\vartheta$ 與 $H_\vartheta^2$ 作一次縮並而產生的，因而可分解（參見 \S\ 3. 2) a) 與 b)）。

有
\[
H_k=H_\vartheta(a\theta u)\sim[H_\vartheta]_{\widehat{ua}}-f\cdot H_\vartheta(\theta Hu).\qquad(\text{將 2) a 中的 }y=\widehat{uv}\text{ 特化}).
\]
\[
H_k(k\eta x)=H_\vartheta a_\eta-H_\vartheta\theta_\eta\cdot f,
\qquad
H_\vartheta a_\eta=\frac54[H_\vartheta]a-\frac14H_\vartheta a_\vartheta
\quad(\S\ 3. 2)b),
\]
\[
H_\vartheta^2(a\theta u)=[H_\vartheta^2]_{\widehat{ua}},
\qquad
H_\vartheta^2a_\eta=[H_\vartheta^2]a=H_k^2(k\eta x)+H_\vartheta^2\theta_\eta\cdot f.
\]

\noindent C) 型列 $(j\eta x)^3$：由 B) 中型列 $K_\eta^3$ 的雙重化簡可知其可化簡；所用設式為：
\[
\begin{aligned}
0&=a_\eta^3(a\theta u)=\theta_\eta^3(a\theta u)
=\{\theta_\eta+(\widehat{a\theta}\eta x)\}^3(a\theta u)-\theta_\eta^3(a\theta u)\\
&=\frac12(j\eta x)^3\modu{(\nu^3,s,\R,t,i,J)}.\qquad(\text{參見 \S\ 5 末尾。})
\end{aligned}
\]
同樣的設式適用到該型列的末型：
\[
0=H_\vartheta^2(\widehat{a\theta}\eta x)a_\eta^3
=H_\vartheta^2(\widehat{a\theta}\eta x)\theta_\eta^3=\frac12H_j^2(j\eta x)^4\modu{(\nu^3,s,\R,t,i,J)}.
\]
型 $(j\eta x)^2$ 的分解有兩種途徑：1) 對可分解型 $H_\vartheta^2$ 的關係作兩次極化（(12.) b）；

2) 直接計算乘積 $a_\nu\cdot\theta_\nu^3$；由此較高型 $(\A Hu)^2$ 也分解。可得：
\par\smallskip\noindent
\begin{minipage}[c]{4.8em}
\centering (18.)
\end{minipage}%
\begin{minipage}[c]{\dimexpr\linewidth-4.8em\relax}
\centering
\(\displaystyle
\left.\begin{array}{ll}
\text{a)}&\dfrac13(\A Hu)^2+\dfrac13(j\eta x)^2=\theta_\nu^2\cdot a_{\nu_1}^2,\\[0.6em]
\text{b)}&(j\eta x)^2=\dfrac43\theta_\nu^3\cdot a_{\nu_1},
\end{array}\right\}\modu{(\nu^3,s,\R,t,i,J)}
\)
\end{minipage}
\par\smallskip
所用設式為：
\[
\begin{aligned}
H_\vartheta^2(a\theta u)^2=\frac13(\A Hu)^2
&=2\theta_\eta^2(a\theta u)(aHu)+\theta_\nu^2\cdot a_{\nu_1}^2\quad(\text{乘積定理})\\
&=(a\theta u)(aHu)\{a_\eta-(\widehat{a\theta}\eta x)\}^2+\theta_\nu^2\cdot a_{\nu_1}^2,
\end{aligned}
\]
\[
\begin{aligned}
a_{\nu_1}\cdot\theta_\nu^3
&=-(a\widehat{\nu\nu_1}u)\theta_\nu^2(\theta\widehat{\nu\nu_1}u)-\frac12(a\widehat{\nu\nu_1}u)\theta_\nu\theta_{\nu_1}(\theta\widehat{\nu\nu_1}u)\\
&=-\frac32\theta_\eta^2(\theta Hu)(aHu)=-\frac32\theta_\eta^2(\theta Hu)(a\theta u)\\
&=-\frac32(a_\eta-(\widehat{a\theta}\eta x))^2\{(aHu)-(a\theta u)\}(a\theta u).
\end{aligned}
\]

\noindent D) 型列 $\A_\eta(\A Hu)$：化簡子 $\A_\nu^2$。

依 C)，$(\A Hu)^2$ 是可分解型。

\noindent E) 型列 $a_l^2(aLu)$：化簡子 $a_\eta^2(aHu)$。
型 $(aLu)$ 與 $a_l(aLu)$：依 \S\ 3，作為經由函數行列式的轉積而分解。

\noindent F) 型列 $a_\sigma^3$：化簡子 $a_\eta^3$。

型 $a_\sigma^2$ 與 $a_\sigma^3$：通過對相應表達式作雙重化簡，以 $a_\nu^3$ 與 $a_\eta^3$ 為化簡子：
\[
H_\nu a_\nu^3\quad\text{与}\quad H_\nu a_\eta^3,
\qquad
H_\nu a_\nu^3a_\eta\quad\text{与}\quad H_\nu a_\eta^4
\]
（參見 \S\ 10. C) 中 $\A_\nu^2$ 與 $\A_\nu^3$ 的化簡）。

由此，在顧及 (16.) 後，可化簡較高型 $a_\nu s_\nu(asu)^2$、$a_\nu^2s_\nu(asu)^2$ 以及含符號 $\R,t(a_\nu,a_\varrho^2)$ 的型：
\par\smallskip\noindent
\begin{minipage}[c]{4.8em}
\centering (19.)
\end{minipage}%
\begin{minipage}[c]{\dimexpr\linewidth-4.8em\relax}
\centering
\(\displaystyle
\left.\begin{array}{rcl}
a_\nu s_\nu(asu)^2&=&\dfrac13 i\theta_\nu^2-\dfrac1{18}iH,\\[0.5em]
a_\nu^2s_\nu(asu)^2&=&0,
\end{array}\right\}\modu{(\R,t)}
\)
\end{minipage}
\par\smallskip
型 $a_\sigma$：可由 (12.)b 的極化而分解；更簡潔地，也可依下列設式直接展開乘積 $a_\nu^2\cdot\theta_\nu^3$：
\[
a_\nu^2\cdot\theta_{\nu_1}^3=a_\nu^2\theta_{\nu_1}\{a_{\nu_1}-(\widehat{a\theta}\nu_1x)\}^2
=-2a_\eta(aHu)(a\theta H)(\widehat{a\theta}\eta x)+(\widehat{a\theta}\nu_1x)^2a_\nu^2\theta_{\nu_1}
\]
並顧及 $Hj(j\eta x)^2$ 的化簡（C），得到：
\par\smallskip\noindent
\begin{minipage}[c]{4.8em}
\centering (20.)
\end{minipage}%
\begin{minipage}[c]{\dimexpr\linewidth-4.8em\relax}
\centering
\(\displaystyle a_\sigma=2a_\nu^2\cdot\theta_{\nu_1}^3\modu{(s,\R,t,i)}.\)
\end{minipage}
\par\smallskip

\noindent\emph{推論：}

1) $(j\eta x)^3$、$\A_\eta(\A Hu)$、$a_\sigma^2$ 是化簡子；$(j\eta x)^2$、$(\A Hu)^2$、$a_\sigma$ 是可分解型。

2) $f$ 只有在與逆變型成積時才參與同 $L$ 的縮並；$f$，因而 $K$ 與 $N$，完全不參與同 $\sigma$、從而也不參與同 $(\nu\sigma x)$ 的縮並。$j$ 只有在與逆變型成積時才出現；$\A$ 在同 $H$ 與 $L$ 縮並時完全不以乘積形式出現（由 $\A_\eta(\A Hu)$ 與 $(\A Hu)^2$ 可知）。同樣，$\sigma$，從而 $(\nu\sigma x)$，在同系統 I 的任何型縮並時都不以乘積形式出現。顧及 \S\ 11 的推論，系統 II 的乘積中只剩 $H$ 的冪以及 $H$ 與 $L$ 的乘積。

\begin{center}
\textbf{\S\ 13.}\quad\emph{八階型（系統 III）。}
\end{center}
\[
\text{下列縮並：}\left\{\begin{array}{l}
\A\cdot j\text{ 与 }\nu,\\
\theta^2,
 f\cdot j,
 N\text{ 与 }H,\\
\theta\text{ 与 }L,\sigma.
\end{array}\right.
\]


\noindent\emph{不可化簡型：}
\[
\renewcommand{\arraystretch}{1.28}
\begin{array}{@{}r@{\quad}c@{\qquad\qquad}r@{\quad}c@{\qquad}c@{}}
\text{A)}&\A_\nu(j\nu x),&\text{B)}&(\vartheta^2\eta x)^3&\cdot\\
&&&(\vartheta^2\eta x)^4&H_\vartheta(\vartheta^2\eta x)^3,\ \theta_\eta(\vartheta^2\eta x)^3,
\end{array}
\]
\[
\renewcommand{\arraystretch}{1.28}
\begin{array}{@{}r@{\quad}c@{\qquad}c@{\qquad\qquad}r@{\quad}c@{}}
\text{C)}&(aHu)H_j&\cdot&\text{D)}&H_n;\ N_\eta,\\
&\cdot&a_\eta(aHu)H_j,&&
\end{array}
\]
\[
\renewcommand{\arraystretch}{1.30}
\begin{array}{@{}r@{\quad}c@{\qquad}c@{\qquad}c@{\qquad}c@{\qquad}c@{}}
&\cdot&\cdot&(\theta Lu)^2&(\theta Lu)^3&\cdot\\
&\cdot&\theta_l&\cdot&L_\vartheta(\theta Lu)^2,\ \theta_l(\theta Lu)^2&\theta_l(\theta Lu)^3\\
\text{E)}&(\vartheta lx)^2&\cdot&\theta_l^2&\cdot&\cdot\\
&\cdot&\theta_l(\vartheta lx)^2&\cdot&\cdot&\cdot
\end{array}
\]
\noindent\emph{化簡：}

\noindent A) $\A_\nu(j\nu x)^3$：化簡子 $a_\nu(j\nu x)^3$。

顧及化簡子 $\theta_{\vartheta'}$，依 \S\ 3. 2)b，可借可分解型 $a_\nu(j\nu x)^2$ 化簡 $\A_\nu(j\nu x)^2$。事實上，由 \S\ 11 B) 得：
\[
\frac3{10}\A_\nu(j\nu x)^2+\frac35a_\nu(j\nu\vartheta)a_\vartheta(j\nu x)
+\frac1{10}a_\nu(j\nu\vartheta)^2
=\frac35\theta_\nu^3\cdot\theta_{\vartheta''}^{\prime2}+\frac25\theta_\nu^3\theta_{\vartheta''}\theta_{\vartheta''}'.
\]
（化簡子為 $a_j$ 與 $\theta_{\vartheta'}$）。

\noindent B) 依 \S\ 3. 2)b) 及下列關係，型 $H_\vartheta(\vartheta'\eta x)^2$、$H_\vartheta^2(\vartheta'\eta x)$、$H_\vartheta^2(\vartheta'\eta x)^2$ 分別是通過依次同可分解型 $H_\vartheta$ 與 $H_\vartheta^2$，相應地同 $H_\vartheta a_\eta$ 與 $H_\vartheta^2a_\eta$ 作一次縮並而產生的，因而可分解：
\[
H_\vartheta(\vartheta'\eta x)^2=2H_\vartheta a_\eta^2\cdot f-2H_\vartheta a_\eta b_\eta.
\]

\noindent C) 型列 $a_\eta(j\eta x)^2$：化簡子為 $a_j$ 與 $(j\eta x)^3$。

型 $a_\eta(j\eta x)$ 可分解：以 $a_j$ 為化簡子，並依 \S\ 3. 2)a 同可分解型 $(j\eta x)^2$ 作一次縮並（特化 $y=uv$）。

型 $a_\eta H_j$ 可分解：1) 同可分解型 $a_\nu(j\nu x)^2$ 作一次縮並；

2) 同可分解型 $(j\eta x)^2$ 作特殊的二重縮並；由此得到可分解型之間的一個關係。

由 $a_\nu(j\nu x)^2$ 得：
\[
0=\theta_\nu^3\cdot\theta_\nu+2a_\eta H_j\modu{(s,\R,t,i,J)}.
\]
由 $(j\eta x)^2$ 得：
\[
0=\theta_\nu^3\cdot\theta_\nu-\frac32a_\eta H_j\modu{(s,\R,t,i,J)}
\]
（省略含逆變式的乘積），所用設式為：
\[
0=\frac12a_\nu(abu)^2\theta_\nu^3
+\frac12a_\nu(abu)\theta_{\nu_1}^3(\theta bu)
-\frac34(\widehat{j\eta}bu)^2+\frac34H_j(\widehat{j\eta}bu)
-\frac18f\cdot H_j^2
\]
並在 $a_\nu(abu)(\theta bu)\theta_{\nu_1}^3$ 中交換符號。

\noindent D) 型列 $N_\eta(NHu)$：化簡子 $\A_\nu^2$。

$(NHu)$、$(n\eta x)$、$N_\eta(n\eta x)$、$H_n(NHu)$ 是可分解型（參見 \S\ 12 B）；$(NHu)^2$ 則通過同型 $(\A Hu)^2$ 作一次縮並而分解。

\noindent E) 型列 $\theta_lL_\vartheta$、$\theta_l^2(\theta Lu)^2$、$\theta_l^2(\vartheta lx)$：

化簡子為 $\theta_\eta H_\vartheta$、$\theta_\eta^2(\theta Hu)^2$、$\theta_\eta^2(\vartheta\eta x)$。

型 $(\theta Lu)$、$(\vartheta lx)$、$L_\vartheta$、$L_\vartheta(\theta Lu)$、$\theta_l(\theta Lu)$、$L_\vartheta(\vartheta lx)$、$\theta_l(\vartheta lx)$、$L_\vartheta^2$、$L_\vartheta^2(\theta Lu)$、$\theta_l^2(\theta Lu)$ 均可分解。（它們與 \S\ 12 B) 的可分解型在對偶意義下相反。）

\noindent F) 型列 $\theta_\sigma(\vartheta\sigma x)$：化簡子 $a_\sigma^2$。

型 $(\vartheta\sigma x)$ 與 $(\vartheta\sigma x)^2$ 是通過同可分解型 $a_\sigma$ 作一次縮並而產生的，因而可分解；$\theta_\sigma$ 由二重縮併產生，也可分解，因為型列
\[
a_\nu^2(a\theta u)\theta_{\nu_1}^3\equiv H_j^2(j\eta x)\equiv0\;\text{模 }(s,\R,t,i,J):
\]
\[
\theta_\sigma=a_\sigma(abu)^2=a_\nu^2(abu)^2\cdot\theta_\nu^3.
\]

\noindent\emph{推論：}

$\theta^3$ 或 $\theta^2K$ 不參與同 $H$ 的縮並；$\theta$ 只有在作為逆變式看待時才以冪或乘積形式參與同 $L$ 的縮並，並且完全不參與同 $\sigma$ 與 $(\nu\sigma x)$ 的縮並。

$N$ 在同系統 II 的任何型縮並時都不以乘積形式出現，也不與系統 II 各型的冪或乘積一同出現。顧及前幾節的推論，系統 I 的乘積中只剩 $f\cdot j$、$\A\cdot j$、$\theta$ 的冪以及 $\theta\cdot K$ 的乘積。

\begin{center}
\textbf{\S\ 14.}\quad\emph{九階型（系統 III）。}
\end{center}
\[
\text{下列縮並：}\left\{\begin{array}{l}
\theta\cdot K\text{ 与 }H,\\
K,j,\A\text{ 与 }L,\\
j,\A\text{ 与 }\sigma,\\
f\text{ 与 }H^2.
\end{array}\right.
\]


\noindent\emph{不可化簡型：}
\[
\renewcommand{\arraystretch}{1.35}
\begin{array}{@{}r@{\quad}c@{\qquad}c@{\qquad\qquad}r@{\quad}c@{\qquad}c@{\qquad}c@{\qquad}c@{\qquad}c@{}}
&&&&\cdot&\cdot&\cdot&(KLu)^3&\cdot\\
\text{A)}&(\vartheta\cdot k,\eta,x)^4&\cdot&\text{B)}&\cdot&\cdot&\cdot&\cdot&K_l(KLu)^3\\
&\cdot&H_k(\vartheta\cdot k,\eta,x)^4,&&\cdot&\cdot&K_l^2&\cdot&\cdot\\
&&&& (klx)^3&\cdot&\cdot&\cdot&\cdot\\
&&&&\cdot&K_l(Klx)^3&\cdot&\cdot&\cdot,
\end{array}
\]
\[
\renewcommand{\arraystretch}{1.35}
\begin{array}{@{}r@{\quad}c@{\qquad}c@{\qquad}r@{\quad}c@{\qquad}r@{\quad}c@{\qquad}r@{\quad}c@{\qquad}c@{}}
\text{C)}&L_j&\cdot&\text{D)}&\A_l,&\text{E)}&(j\sigma x),&\text{G)}&(aH^2u)^3&(aH^2u)^4\\
&\cdot&L_j^2,&&&&&&\cdot&a_\eta(aH^2u)^3.
\end{array}
\]
\noindent\emph{化簡：}

A) 型 $H_\vartheta(k\eta x)^3$、$H_\vartheta^2(k\eta x)^2$、$H_\vartheta^2(k\eta x)^3$ 是通過依次同可分解型作一次縮並而產生的，因而可分解。

B) 型列 $K_lL_k$、$K_l^2(KLu)$、$K_l^2(klx)$：

化簡子為 $\theta_\eta H_\vartheta$、$a_\eta^2(aHu)$、$\theta_\eta^2(\vartheta\eta x)$。

型 $K_l$、$K_l(KLu)$、$K_l(klx)$、$(KLu)^2$、$(klx)^2$ 作為經由函數行列式的轉積而分解（\S\ 3）。

型 $L_k$、$L_k(KLu)$、$L_k(KLu)^2$、$L_k(klx)$、$L_k(klx)^2$、$L_k^2$、$L_k^2(KLu)$、$L_k^2(klx)$、$L_k^3$、$K_l(KLu)^2$、$K_l(klx)^2$ 是通過同下列可分解型作一次縮並而產生的，因而可分解：
\[
L_\vartheta,
\ H_k(KHu),
\ L_\vartheta(\vartheta lx),
\ H_k^2,
\ L_\vartheta^2,
\ L_\vartheta^2(\theta Lu),
\ K_\eta(KHu),
\ \theta_l(\vartheta lx)
\]
依 \S\ 3. 2) a) 與 b)。

C) 型列 $(jlx)^3$：化簡子 $(j\eta x)^3$。

型 $(jlx)^2$ 與 $L_j(jlx)$：由同可分解型 $(j\eta x)^2$ 作一次縮並而產生。

D) 型列 $\A_l(\A Lu)$：化簡子 $\A_\nu^2$。

型 $(\A Lu)^2$、$(\A Lu)^3$：同可分解型 $(\A Hu)^2$ 作一次縮並。

E) 型列 $(j\sigma x)^2$：將 $j$ 替換為型列中的較低型（\S\ 5），便以 $a_\sigma^2$ 為化簡子：
\[
a_\sigma^2(a\theta u)^2=\frac13(j\sigma x)^2\ldots,
\qquad
 a_\sigma^2(a\theta\sigma x)^2=\frac13(j\sigma x)^4.
\]
F) 型列 $\A_\sigma^2$：化簡子 $a_\sigma^2$。

型 $\A_\sigma$：將 $\A$ 替換為型列中的較低型，便以 $a_\sigma^2$ 為化簡子：
\[
a_\vartheta a_\sigma^2=\frac23\A_\sigma\modu{\nu}.
\]

\noindent\emph{推論：}只有型 $j$ 參與同 $\sigma$ 與 $(\nu\sigma x)$ 的縮並。

$N$ 不參與同 $L$ 的縮並，因為與 $\A_l$ 對應的型 $N_l$ 可分解。

\begin{center}
\textbf{\S\ 15.}\quad\emph{十階型（系統 III）。}
\end{center}
\[
\text{下列縮並：}\left\{\begin{array}{l}
\theta^2,
\ f\cdot j\text{ 与 }L,\\
\theta\text{ 与 }H^2.
\end{array}\right.
\]
\noindent\emph{不可化簡型：}
\[
\begin{array}{c@{\quad}l@{\qquad}c@{\quad}l}
\text{A)}&\begin{array}{cc}(\vartheta^2lx)^4&\cdot\\ \cdot&
\theta_l(\vartheta^2lx)^4,\end{array}&
\text{B)}&\begin{array}{cc}(aLu)^2L_j&\cdot\\ \cdot&
a_l(aLu)^2L_j,\end{array}\\[0.8em]
\text{C)}&\begin{array}{ccc}(\theta H^2u)^3
&(\theta H^2u)^4&\cdot\\ \cdot
&H_\vartheta(\theta H^2u)^3,
&\theta_\eta(\theta H^2u)^3.\end{array}
\end{array}
\]

\noindent\emph{化簡：}

A) 與 B)：其餘各型均通過同可分解型作一次縮並而分解。

C) 依 \S\ 13 B)，由對偶相反的考察可知其餘各型均可分解。

\noindent\emph{推論：}

$\theta^2\cdot K$ 不參與同 $L$ 的縮並；相應地，$H^2L$ 不參與同 $K$ 的縮並。$H^3$ 與 $H^2L$ 不參與同 $\theta$ 的縮並。由此，\S\ 7 的縮並方案得證。

\begin{center}
\textbf{\S\ 16.}\quad\emph{十一、十二、十三、十五階型（系統 III）。}
\end{center}

\noindent\emph{十一階。}
\[
\text{下列縮並：}\left\{\begin{array}{l}
\theta\cdot K\text{ 与 }L,\\
K,j\text{ 与 }H^2,\\
j\text{ 与 }(\nu\sigma x),\\
f\text{ 与 }H\cdot L.
\end{array}\right.
\]
\noindent\emph{不可化簡型：}
\[
\begin{array}{c@{\quad}l@{\qquad}c@{\quad}l}
\text{A)}&\begin{array}{cc}(\vartheta\cdot k,l,x)^5&\cdot\\ \cdot&
\theta_l(\vartheta\cdot k,l,x)^5,\end{array}&
\text{B)}&\begin{array}{cc}(KH^2u)^4&\cdot\\ \cdot&
K_\eta(KH^2u)^4,\end{array}\\[0.8em]
\text{C)}&(\nu\sigma j),&
\text{D)}&(a,H\cdot L,u)^4.
\end{array}
\]

\noindent\emph{十二階。}
\[
\text{下列縮並：}\left\{\begin{array}{l}
\theta^3\text{ 与 }L,\\
\theta\text{ 与 }H\cdot L.
\end{array}\right.
\]
\noindent\emph{不可化簡型：}
\[
\text{E)}\ (\vartheta^3lx)^6,
\qquad
\text{F)}\ (\theta,H\cdot L,u)^4,
\qquad L_\vartheta(\theta,H\cdot L,u^4).
\]

\noindent\emph{十三階。}
\[
\text{$K$ 与 $H\cdot L$ 的縮並。}
\]
\noindent\emph{不可化簡型：}
\[
\text{G)}\ (K,H\cdot L,u)^5,
\qquad K_\eta(K,H\cdot L,u)^5.
\]

\noindent\emph{十五階。}
\[
\text{$K$ 与 $H_3$ 的縮並。}
\]
\noindent\emph{不可化簡型：}
\[
\text{H)}\ (KH^3u)^6.
\]

\noindent\emph{化簡：}

型 $H_{j^3}^2$、$H_{j^4}^2$。由下列關係，以 $L_j^3$ 為化簡子：
\[
(\nu lx)L_x^3=-\frac12H^2\modu{(\sigma,s)}.
\]
由前幾節的考察可知其餘各型均可分解或可化簡。

\noindent\emph{推論：}

依照 \S\ 7 的縮並方案以及 \S\S\ 8 至 15 的推論，系統 III 的不可化簡型（共 117 個）至此已經窮盡。



\clearpage
\begin{center}
{\Large\bfseries 第四章。\quad 模 $(\varrho,t)$ 的相對完備系統。}
\end{center}

\begin{center}
\textbf{\S\ 17.}\quad\emph{模 $(ss'u)^4$ 與 $s$ 的型系統的化簡。}
\end{center}

為構造再高一級的相對完備系統，仿照 \S\ 7，須關於下一個模
\[
\nu(s)=(ss'u)^4
\]
構造 $s$ 的型系統，並將它同模 $s$ 的相對完備系統中的諸型作縮並。現將證明，若引入模 $(\varrho,t)$ 作為較高的模，則
\[
\begin{array}{ll}
\text{A)}&\text{模 }(ss'u)^4\text{ 可化簡，}\vphantom{\dfrac11}\\
\text{B)}&s\text{ 的型系統僅由型 }s\text{ 組成。}
\end{array}
\]

A) 模 $(ss'u)^4$ 的化簡直接由公式 (16.)、\S\ 11 中所得的化簡子 $s_\nu^2$ 得出。由
\[
 s_\nu^2=\frac43 i a_\nu^2+\frac45\theta_\varrho^2+\frac85(atu)^2-
 \frac{26}{15}(\nu\varrho x)^2+\frac16J\cdot u_x^2-
 \frac29i^2\cdot u_x^2
\]
將 $x$ 向 $\nu_1$ 極化，得到：
\srcnumdisplay{(21.)\srcfnmark{*)}}{%
\begin{aligned}
(ss'u)^4&=-\frac23J\cdot \nu+6s_\nu^2s_{\nu_1}^2\\
&=\frac{24}{5}\theta_\nu^2\theta_\varrho^2+
\frac{48}{5}a_\nu^2(atu)^2-
\frac{52}{5}H_\varrho^2+
\frac43i\cdot(asu)^4+\frac13J\cdot \nu-
\frac49i^2\cdot \nu,
\end{aligned}
}
從而得到模 $(ss'u)^4$ 的化簡。

\srcfntext{*)}{由公式 (21.)，應用公式 (10.)c，即得引言註釋中提到的三個九階不變量之間的關係：
\srcnumdisplay{(22.)}{%
\begin{aligned}
(ass')^4&=\frac{24}{5}\theta_\nu^2\theta_\varrho^2a_\nu^2a_\vartheta^2+
\frac{48}{5}a_\nu^2b_\nu^2(tab)^2-
\frac{52}{5}H_\varrho^2a_\eta^4+
\frac53J\cdot i-\frac49i^3,\\
(ass')^4&=\frac{24}{5}a_\varrho^2a_{\varrho'}^2-
\frac{12}{5}t_\varrho^2+\frac53J\cdot i-\frac49i^3.
\end{aligned}
}
}

B) 對
\[
Z=(ss'u)^2=\theta(s)
\]
的化簡，以及由此對 $s$ 的型系統的化簡，是依公式 (8.)b、\S\ 7，以較低型 $g_\nu^2$ 代替型 $Z$ 而得到的；此處還要利用關係
\[
 s_\eta^2=s_\nu^2(\nu\nu_1x)^2-\frac13Z.
\]
由公式 (13.)、\S\ 11，型 $g_\nu^2$ 以化簡子 $s_\nu^2$ 為因子。依公式 (8.) 與 (16.)，有：
\[
 g_\nu^2=-Z+\frac{14}{5}ia_\eta^2+\frac{14}{15}i(asu)^2
 \modu{(\varrho,t)},
\]
依公式 (13.)、(14.)、(15.)、(16.)，有：
\[
\begin{aligned}
g_\nu^2&=2[s_\vartheta^2]_{\nu^2}-\frac43 i\A_\nu^2
      =2s_\vartheta^2s_\nu^2-\frac65[s_\vartheta^2(\theta su)^2]\widehat{\nu x}^{2}
        -\frac43i\A_\nu^2\\
&=\frac45i\cdot a_\eta^2-\frac25i\cdot(asu)^2+\frac13J\cdot\theta
\modu{(\varrho,t)}:
\end{aligned}
\]
於是
\srcnumdisplay{(23.)}{%
Z=2i\cdot a_\eta^2+\frac43i\cdot(asu)^2-\frac13J\cdot\theta
 \modu{(\varrho,t)}.
}

模 $((ss'u)^4,\varrho,t)$ 的相對完備系統因而轉化為模 $(\varrho,t)$ 的相對完備系統；再將它同兩個二次型的已知系統作轉積，便得到絕對完備系統。為構造模 $(\varrho,t)$ 的相對完備系統，由於依公式 (23.)，$s$ 的型系統化簡為型 $s$，我們必須把模 $(s,\varrho,t)$ 的相對完備系統同 $s$ 及 $s$ 的冪作縮並。下文將表明，$s$ 的冪只同系統 I 發生縮並。

在同階、同次數的諸型中，我們將這樣一些型定義為\glqq{}較高型\grqq{}：它們由模 $(s,\varrho,t)$ 的相對完備系統中的較高型同 $s$ 縮並而來，而同 $s$ 的縮並僅視為原型的極化。這樣，各型的排列次序便唯一確定（參看 \S\ 1，型列 2）。例如：
\[
(\theta_\eta(\theta Hu),s,u)^2\text{ 高於 }s_\eta((\theta Hu)^2,s,u)\srcfnmark{*)},
\]
\[
(\theta_\eta(\theta Hu)^2,s,u)\text{ 高於 }(\theta_\eta(\theta Hu),s,u)^2.
\]

我們把所得系統分成四個子系統：
\[
\begin{array}{ll}
\text{系統 }s_I:&\text{由 }s\text{ 及 }s\text{ 的冪同系統 I 縮並而成，}\\
\text{系統 }s_{II}:&\text{由 }s\text{ 同系統 II 縮並而成，}\\
\text{系統 }s_{III}:&\text{由 }s\text{ 同系統 III 的各個型}\\
&\text{（不含積）以及同系統 I 和 II 各取}\\
&\text{一型所成的積縮並而成，}\\
\text{系統 }s_{IV}:&\text{由 }s\text{ 同系統 I 和 III、II 和 III、III 和 III}\\
&\text{各取一型所成的積縮並而成。}
\end{array}
\]
還須指出，從現在起，所有公式均應理解為模 $(\varrho,t)$；從公式 (27.) 起，則均應理解為模 $(\varrho,t,i,J,\text{ 較高型})$，不再逐次明言。

為進行化簡，我們彙集先前所得的公式：
\srcnumdisplay{(24.)}{%
\begin{aligned}
s_\nu^2&=\frac43i\,a_\nu^2+\frac16J\cdot u_x^2-\frac29i^2\cdot u_x^2,\\
s_\nu^3&=\frac13J\cdot u_x,
& s_\nu^4&=J,\\
g&=2s_\vartheta^2-\frac43i\Delta,\\
s_\vartheta^2(\theta su)&=0,\\
s_\vartheta^2(\theta su)^2&=\frac89i\cdot a_\nu^2-\frac4{27}i^2\cdot u_x^2,\\
Z&=2i\cdot a_\eta^2+\frac43i(asu)^2-\frac13J\cdot\theta,\\
g_\nu&=-s_\eta+u_x\left\{2ia_\eta^2-\frac23i(asu)^2+\frac23J\cdot\theta\right\},\\
g_\nu^2&=\frac45ia_\eta^2-\frac25i(asu)^2+\frac13J\cdot\theta,
&g_\nu^3&=0,\quad g_\nu^4=0.
\end{aligned}
}
\srcfntext{*)}{為使寫法簡短，我們把轉積式的首項作為型列中的型，以代替完整轉積式 $[(\theta Hu)^2]_{\widehat{su}}s$。列寫公式時，所有轉積式都以其首項代替；例如，以 $\theta_\eta(\theta Hu)(\theta su)^2$ 代替 $(\theta_\eta(\theta Hu),s,u)^2$。}

\noindent\emph{推論：}
\[
 s_\nu^2,\quad s_\vartheta^2(\theta su),\quad s_\vartheta^2s_\nu,\quad s_\vartheta^2\theta_\nu
\]
都是化簡子。

\begin{center}
\textbf{\S\ 18.}\quad\emph{系統 $s_I$。}
\end{center}

\[
\text{作如下縮並：}\left\{\begin{array}{l}
s\text{ 同 } f;\ \theta;\ K,j,\Delta;\ f\cdot j,N;\ \Delta\cdot j,\\
s^2\text{ 同 }\theta,K.
\end{array}\right.
\]

\noindent\emph{不可化簡型。}

\noindent 五階：
\[
(asu);\qquad (asu)^2;\qquad (asu)^3;\qquad (asu)^4.
\]

\noindent 六階：
\[
\begin{array}{cccc}
(\theta su)&(\theta su)^2&(\theta su)^3&(\theta su)^4\\
s_\vartheta&s_\vartheta(\theta su)&s_\vartheta(\theta su)^2&s_\vartheta(\theta su)^3\\
\cdot&s_\vartheta^2&\cdot&\cdot
\end{array}
\]

\noindent 七階：
\[
\begin{array}{c@{\quad}c@{\quad}c@{\quad}c@{\qquad}l}
\cdot&(Ksu)^2&(Ksu)^3&(Ksu)^4&\\
\text{A)}\ s_k&s_k(Ksu)&s_k(Ksu)^2&s_k(Ksu)^3&
\text{B)}\ s_j,\\
\cdot&s_k^2&\cdot&\cdot,&
\text{C)}\ (\Delta su),\ (\Delta su)^2.
\end{array}
\]

\noindent 八階：
\[
\begin{array}{ll}
\text{A)}\ s_j(asu),\ s_j(asu)^2,\ s_j(asu)^3,
&
\text{B)}\ \begin{array}[t]{cc}
\cdot&(Nsu)^2\\
s_n&s_n(Nsu).
\end{array}
\end{array}
\]

\noindent 十階：
\[
\text{A)}\ s_j(\Delta su),\qquad \text{B)}\ s_\vartheta(\theta s'u)^4.
\]

\noindent 十一階：
\[
\begin{array}{cc}
(Ks^2u)^6&\cdot\\
s_k(Ks^2u)^5&s_k(Ks^2u)^6.
\end{array}
\]

\noindent\emph{化簡：}

型列
\[
 s_\vartheta^2(\theta su),\quad s_k^2(Ksu),\quad s_j^2,\quad (\Delta su)^3,\quad (Nsu)^3
\]
均以化簡子 $s_\vartheta^2(\theta su)$ 為因子；對於 $s_j^2$ 與 $(\Delta su)^3$，這是通過把 $j$ 和 $\Delta$ 還原到型列中的較低型而得到的：
\[
\begin{aligned}
s_\vartheta^2(\theta su)(a\theta u)^3&=-\frac12s_j^2,\\
s_\vartheta^2(\theta su)(a\theta u)^2&=\frac13(\Delta su)^3,\\
s_\vartheta^2(\theta su)^2(a\theta u)^2&=\frac13(\Delta su)^4+\frac13s_j^2.
\end{aligned}
\]
型 $s_\vartheta^2s_{\vartheta'}$ 與 $s_\vartheta^2s_{\vartheta'}^2$：化簡子為 $s_\vartheta^2(\theta su)$ 與 $\theta_{\vartheta'}$：
\[
\begin{aligned}
s_\vartheta^2s_{\vartheta'}&=s_\vartheta^2\theta_{\vartheta'}-s_\vartheta^2(\theta s\widehat{\vartheta' x}),\\
s_\vartheta^2s_{\vartheta'}^2&=s_\vartheta^2s_{\vartheta'}\theta_{\vartheta'}-s_\vartheta^2s_{\vartheta'}(\theta s\widehat{\vartheta' x}).
\end{aligned}
\]
型列 $s_j(\Delta su)^2$：化簡子為 $\Delta_j$ 與 $(\Delta su)^3$。

\noindent\emph{推論：}

$s$ 在同 $f,j,\Delta,N$ 的縮並中不以冪出現。$f,j,\Delta,N$ 只以同逆變型的積進入縮並；不過，同 $f$ 所成的積仍可關於 $x$ 為二次，同 $\Delta$ 所成的積仍可關於 $x$ 為一次。也就是說，應考察如下型的積：
\[
 f\cdot(0,\lambda),\quad f\cdot(1,\lambda),\quad f\cdot(2,\lambda),\quad
 j\cdot(\mu,0),\quad N,\quad
 \Delta\cdot(0,\lambda)\quad \Delta\cdot(1,\lambda)\quad(\lambda>0),
\]
其中 $(\mu,\lambda)$ 表示關於 $x$ 為 $\mu$ 次、關於 $u$ 為 $\lambda$ 次的型。

$\theta$ 或 $K$ 在同 $s$ 或 $s^2$ 的縮並中，既不以冪出現，也不同任意型的積出現。事實上，對於含函數行列式 $K$ 的積 $K\cdot\varphi_x$（$\varphi\ne f$ 或 $\theta$），分解型之間有關係：
\[
K\cdot\varphi_x=(a\theta u)\varphi_x=f\cdot(\varphi\theta u)-\theta\cdot(\varphi a u).
\]
至此，上述系統 $s_I$ 的縮並方案即得證明。

\begin{center}
\textbf{\S\ 19.}\quad\emph{系統 $s_{II}$。}
\end{center}

$s$ 同 $\nu;\ H;\ L;\ H^2$ 的縮並。

\noindent\emph{不可化簡型：}
\[
\begin{array}{c@{\qquad}c@{\qquad}c@{\qquad}c}
\text{六階} & \text{八階} & \text{十階} & \text{十二階}\\[0.3em]
s_\nu & (Hsu),\ (Hsu)^2 & (Lsu)^2,\ (Lsu)^3 & (H^2su)^3\\
\cdot & s_\eta & s_l & \cdot\\
\cdot & \cdot & \cdot & \cdot\\
s_\nu^4=J & \cdot & \cdot & \cdot
\end{array}
\]

\noindent\emph{化簡：}

型列 $s_\eta(Hsu)$ 與 $s_l(Lsu)$：化簡子為 $s_\nu^2$。

型列 $s_\sigma$：由 $H$ 得到化簡子 $s_\nu^2$，且 $s_\nu^2=\dfrac23s_\sigma$。

$(H^2su)^4$ 由公式 (7.)a 分解（參看 \S\ 11 A）。由此，$(H\cdot L,s,u)^4$ 分解。

\noindent\emph{推論：}

$s^2$ 不進入同系統 II 的縮並。

$\nu$ 只以同逆變型的積進入縮並；把 $H$ 看作協變式時，它以同型 $(1,\lambda)$、$(2,\lambda)$、$(3,\lambda)$（$\lambda$ 任意）的積進入縮並。

$\sigma$ 與 $(\nu\sigma x)$ 完全不進入縮並；$L$ 作為函數行列式，也不以積進入縮並（對系統 III 的協變式所作的完全轉積式均可化簡）。

系統 I 與 II 的積中只餘下：
\[
f\cdot \nu,\qquad \Delta\cdot \nu.
\]

\begin{center}
\textbf{\S\ 20.}\quad\emph{七階型（系統 $s_{III}$）。}
\end{center}

$s$ 同 $f\cdot \nu$、$a_\nu$、$a_\nu^2$ 的縮並。

\noindent\emph{不可化簡型：}
\[
\begin{gathered}
s_\nu(asu),\quad a_\nu(asu),\quad s_\nu(asu)^2,\quad a_\nu(asu)^2,\quad
s_\nu(asu)^3,\quad a_\nu(asu)^3,\\
a_\nu s_\nu,\qquad a_\nu s_\nu(asu),\qquad a_\nu^2(asu),\qquad a_\nu^2(asu)^2.
\end{gathered}
\]

\noindent\emph{化簡：}

由直接同化簡子 $s_\nu^2$ 與 $s_\eta(Hsu)$ 縮並而產生的顯然化簡不再提及；同樣也不再提及含函數行列式的轉積式的分解。

型 $a_\nu s_\nu(asu)^2$ 與 $a_\nu^2s_\nu(asu)^2$：參看公式 (19.)、\S\ 12；
\[
a_\nu s_\nu(asu)^3=a_\nu(a\widehat{\nu_1\nu_2}u)^3(\nu_1\nu_2\nu)=0.
\]
型列 $a_\nu^2s_\nu$：化簡子 $s_\vartheta^2(\theta su)$ 由把 $a_\nu^2$ 還原到較低型而得，也就是說，按如下設式對表達式
\[
s_\vartheta^2(a\theta s)(a\theta u),\qquad
s_\vartheta^2(a\theta s)(a\theta u)^2
\]
作雙重化簡：
\[
\begin{aligned}
s_\vartheta^2(a\theta s)(a\theta u)
&=[(a\theta u)^2]s^3+\frac15[a_\vartheta(a\theta u)^2]s^2+\frac1{60}[a_\vartheta^2(a\theta u)^2]s\\
&=\frac12a_\nu^2s_\nu-\frac23a_\nu s_\nu^2,\\
s_\vartheta^2(a\theta s)(a\theta u)^2
&=\frac45[a_\vartheta(a\theta u)^2]_{\widehat{us}}s^2+\frac{23}{180}[a_\vartheta^2(a\theta u)^2]_{\widehat{us}}s\\
&=-\frac12a_\nu^2s_\nu(asu).
\end{aligned}
\]
此處 $(a_\nu^2b_\nu(abu)=0)$。

由此得到公式：
\srcnumdisplay{(25.)}{%
\begin{array}{rcl@{\qquad}rcl}
\underline{a_\nu^2s_\nu}&=&\dfrac43i\cdot a_\nu^2b_\nu,
& \underline{a_\nu s_\nu(asu)^2}&=&\dfrac13i\cdot\theta_\nu^2-\dfrac1{18}iH,\\[0.7em]
\underline{a_\nu s_\nu(asu)^3}&=&0,
& \underline{a_\nu^2s_\nu(asu)}&=&0,\\[0.7em]
\underline{a_\nu^2s_\nu(asu)^2}&=&0.
\end{array}
}

\noindent\emph{推論：}

1) $a_\nu s_\nu(asu)^2$、$a_\nu^2s_\nu$ 是化簡子。

2) 由此得到在 \S\ 19 中已判定為可化簡的型 $(\Delta su)^3$、$(\Delta su)^4$ 的值；相應的型列 $(agu)^2$ 亦然，它在同 $\nu$ 縮並時，由於化簡子 $g_\nu$ 而引起雙重化簡。
\srcnumdisplay{(26.)}{%
\begin{aligned}
\text{a)}\quad (\Delta su)^3&=-\frac65a_\nu^2(asu)+3a_\nu s_\nu(asu)+2i\cdot\theta_\nu(\vartheta\nu x),\\
\text{b)}\quad (\Delta su)^4&=-\frac25a_\nu^2(asu)^2+\frac43i\theta_\nu^2+\frac19iH.
\end{aligned}
}

由公式 (24.) 與 (26.)，模 $(i,J)$（參看第 71 頁）：
\srcnumdisplay{(27.)}{%
\begin{aligned}
\text{a)}\quad (agu)^2&=\frac23(\Delta su)^2-\frac43a_\nu s_\nu+\frac35s\cdot a_\nu^2,\\
\text{b)}\quad (agu)^3&=2a_\nu s_\nu(asu)-a_\nu^2(asu)^2,\\
\text{c)}\quad (agu)^4&=-a_\nu^2(asu)^2.
\end{aligned}
}


\begin{center}
\textbf{\S\ 21.}\quad\emph{八階型（系統 $s_{III}$）。}
\end{center}

$s$ 同四階型（系統 III）的縮並。（\S\ 9。）

\noindent\emph{不可化簡型：}
\[
\begin{array}{c|c|c}
\cdot
%
& ((\vartheta \nu x),s,u)^2,\ (\theta_\nu su)
& \bigstar\\[0.45em]
(\vartheta \nu x)s_\nu,\ ((\vartheta \nu x^2),s,u)
& ((\vartheta \nu x)^2,s,u)^2,\ \theta s_\nu
& \\[0.45em]
(\vartheta \nu x)^2s_\nu & ((\vartheta \nu x)^2,s,u)s_\nu,\ (\theta_\nu(\vartheta \nu x)^2,s,u)
&
\end{array}
\]
\[
\begin{array}{c|c|c|c}
\bigstar
& ((\vartheta \nu x),s,u)^3,\ (\theta_\nu su)^2
& ((\vartheta \nu x),s,u)^4,\ (\theta_\nu su)^3
& \cdot\\[0.45em]
& ((\vartheta \nu x)^2,s,u)^3,\ (\theta_\nu(\vartheta \nu x),s,u)^2,\ (\theta_\nu^2su)
& (\theta_\nu(\vartheta \nu x),s,u)^3,\ (\theta_\nu^2su)^2
& (\theta_\nu(\vartheta\nu x),s,u)^4.\\[0.45em]
& \cdot
& (\theta_\nu^3su)
& \cdot
\end{array}
\]

\noindent\emph{化簡：}

$((\vartheta \nu x),s,u)s_\nu$ 依 \S\ 3 分解為函數行列式上的轉積式。

型列 $((\vartheta \nu x),s,u)^2s_\nu$：化簡子為 $s_\nu^2$ 與 $s_\vartheta^2\theta_\nu$：
\[
(\widehat{\vartheta \nu}su)s_\nu(\theta su)=s_\nu s_\vartheta(\theta su)=-\frac12(\widehat{\vartheta \nu}su)^2(\theta su),
\]
\[
\theta_\nu(\widehat{\vartheta \nu}su)s_\nu=\theta_\nu s_\nu s_\vartheta=-\frac12\theta_\nu(\widehat{\vartheta \nu}su)^2.
\]
型列 $\theta_\nu s_\nu(\theta su)$：化簡子 $a_\nu s_\nu(asu)^2$ 由下列表達式的雙重化簡得到：
\[
a_\nu s_\nu(asu)^2(abu)\quad\text{和}\quad a_\nu s_\nu(asu)^2(abu)(sbu);
\]
$\theta_\nu s_\nu(\theta su)^3$ 則由 $(agu)^3(abu)(bgu)^3$ 的雙重化簡得到。

型列 $\theta_\nu(\vartheta \nu x)s_\nu$：化簡子 $a_\nu^2s_\nu$ 來自 $\theta_\nu(\vartheta \nu x)s_\nu=a_\nu^2(abu)s_\nu$。

型列 $(\theta_\nu(\vartheta \nu x)^2,s,u)$：對型列 $\theta_\nu(\widehat{\vartheta \nu}su)s_\nu$ 中的諸型作雙重化簡；這些型同時屬於型列 $((\vartheta \nu x)su)^2s_\nu$ 與 $\theta_\nu(\vartheta \nu x)s_\nu$。

型 $(\theta_\nu(\vartheta \nu x),s,u)$ 分解為函數行列式 $\theta_\nu(\vartheta \nu x)=a_\nu^2(abu)$ 上的一重轉積式。

型 $(\vartheta \nu x)^2,s,u)^4$：對表達式 $(a\A u)^2(\A su)^4$ 作雙重化簡；也可依據公式 (27.) 對 $(\theta gu)^4$ 作雙重化簡，仿照 $(j\eta x)^4$ 的化簡（\S\ 12 C）。

雙重化簡給出公式：
\srcnumdisplay{(28.)}{%
\begin{aligned}
0&=\theta_\nu s_\nu(\theta su)-\frac16(\widehat{\vartheta \nu}su)^2(\theta su)-\frac1{12}(Hsu),\\
0&=\theta_\nu s(\theta su)^2-\frac14(\widehat{\vartheta \nu}su)^2(\theta su)^2-\frac1{24}(Hsu)^2,\\
0&=(\widehat{\vartheta \nu}su)^2(\theta su)^2+2\theta_\nu(\widehat{\vartheta \nu}su)(\theta su)^2
  +3\theta_\nu^2(\theta su)^2-\frac13H(su)^2,\\
0&=\theta_\nu s_\nu(\theta su)^3+\frac12\theta_\nu(\widehat{\vartheta \nu}su)(\theta su)^3,\\
0&=\theta_\nu s_\nu s_\vartheta-\frac14s_\eta,
\quad 0=\theta_\nu s_\nu s_\vartheta(\theta su),
\quad 0=\theta_\nu s_\nu s_\vartheta(\theta su)^2.
\end{aligned}
}
最後一組對應於 $(\theta_\nu(\vartheta \nu x)^2,s,u)^2\ldots$。

\noindent\emph{推論：}

1) $\theta(\vartheta \nu x)s_\nu$、$(\widehat{\vartheta \nu}su)(\theta su)s_\nu$、$\theta_\nu(\widehat{\vartheta \nu}su)^2$、$(\widehat{\vartheta \nu}su)^2(\theta su)^2$ 是化簡子。

2) 四階型中的 $(5,\lambda)$、$(6,\lambda)$ 等不進入同 $s^2$ 的縮並。

3) 對型列 $(\theta gu)^2$，由 (27.) 並計及本節的化簡，得到如下公式：
\srcnumdisplay{(29.)}{%
\begin{aligned}
\text{a)}\quad (\theta gu)^2&=\frac43\theta_\nu s_\nu+\frac23s_\nu s_\vartheta-\frac13s\cdot\theta_\nu^2-\frac19s\cdot H-\frac13f\cdot a_\nu^2(asu)^2+\frac13a_\nu^2\cdot(asu)^2,\\
\text{b)}\quad (\theta gu)^3&=-\frac13(\widehat{\vartheta \nu}su)^2(\theta su)-\theta_\nu(\widehat{\vartheta \nu}su)(\theta su)-\theta_\nu^2(\theta su),\\
\text{c)}\quad (\theta gu)^4&=2\theta_\nu^2(\theta su)^2-\frac13(Hsu)^2,\\
g\vartheta(\theta gu)&=(\widehat{\vartheta \nu}su)(\vartheta \nu x)s_\nu-\theta_\nu(\vartheta \nu x)(\widehat{\vartheta \nu}su),\\
g\vartheta(\theta gu)^2&=\frac23s\cdot\theta_\nu^3,
\quad g\vartheta(\theta gu)^3=\frac13\theta_\nu^3(\theta su),
\quad g\vartheta(\theta gu)^4=0.
\end{aligned}
}

\begin{center}
\textbf{\S\ 22.}\quad\emph{九階型（系統 $s_{III}$）。}
\end{center}

$s$ 同五階型（系統 III）及積 $\A\cdot \nu$ 的縮並（\S\ 10）。

\noindent\emph{不可化簡型：}
\[
\begin{array}{c@{\quad}c|c|c|c}
\text{A)}& \cdot & \cdot & (K_\nu su)^2 & \bigstar\\[0.45em]
& \cdot & ((k\nu x)^2,s,u)^2,\ K_\nu s_\nu & ((k\nu x)^2,s,u)^3 & \\[0.45em]
& (k\nu x)^2s_\nu,\ ((k\nu x)^3,s,u) & ((k\nu x)^3,s,u)^2 & \cdot & \cdot\\[0.45em]
& (k\nu x)^3s_\nu & ((k\nu x)^3,s,u)s_\nu,\ (K_\nu(k\nu x)^3,s,u) & \cdot & \cdot\\[0.45em]
\bigstar & (K_\nu su)^3 & (K_\nu su)^4 & \cdot & \\[0.45em] & \cdot & \cdot & (K_\nu^2su)^4, &
\end{array}
\]
\[
\begin{array}{c@{\quad}c|c@{\qquad}c@{\quad}c}
\text{B)}& (j\nu x)s_\nu,\ ((j\nu x)^2,s,u) & ((j\nu x)^2,s,u)^2 && \\
& (j\nu x)^3,s,u) & \cdot, && \text{C)}\quad s_\nu(\A su),\ (\A_\nu su),
\end{array}
\]
\[
\begin{array}{c@{\quad}c|c|c}
\text{D)}& \cdot & ((aHu),s,u)^2,\ ((aHu)^2,s,u) & \bigstar\\[0.45em]
& (aHu)s_\eta,\ (a_\eta su) & (aHu)^2s_\eta,\ (a_\eta su)^2 & \\[0.45em]
& \cdot & (a_\eta^2su) & \\
\bigstar& ((aHu),s,u)^3,\ (aHu)^2,s,u)^2 & ((aHu),s,u)^4 & \\
& (a_\eta su)^3,\ (a_\eta(aHu),s,u)^2,\ (a_\eta(aHu)^2,s,u) & (a_\eta(aHu),s,u)^3. &
\end{array}
\]

\noindent\emph{化簡：}

A) 這些化簡由 \S\ 21 的化簡經一重縮並直接得到。型列 $K_\nu s_\nu(Ksu)^2$ 以化簡子 $a_\nu s_\nu(asu)^2$ 與 $\theta_\nu s_\nu(\theta su)^2$ 為因子。由此，較高型 $K_\nu^2(Ksu)^2$、$K_\nu^2(Ksu)^3$ 按如下設式分解：
\[
0=a_\nu s_\nu(asu)^2(a\theta u)=K_\nu s_\nu(Ksu)^2=\theta_\nu s_\nu(\theta su)^2(a\theta u).
\]

B) 型列 $(j\nu x)^2s_\nu$：化簡子 $a_\nu^2s_\nu$ 來自
\[
(j\nu x)^2s_\nu=3a_\nu^2s_\nu(a\theta u)^2.
\]
$((j\nu x)^3,su)^2$：化簡子 $a_\nu^2s_\nu$ 來自
\[
((j\nu x)^3,s,u)^2=-6a_\nu^2s_\nu(a\theta u)^2(\theta su).
\]

C) 型 $\A_\nu(\A su)^2$ 與 $s_\nu(\A su)^2$：由公式 (27.)a 得化簡子 $g_\nu$。

型 $\A_\nu s_\nu$：1) 由 $a_\vartheta a_\nu^2s_\nu=\frac23\A_\nu s_\nu$ 得化簡子 $a_\nu^2s_\nu$；

2) 依公式 (27.)a，通過表達式 $(\A s\widehat{\nu x})^2$ 的雙重化簡而分解。

由此，較高型 $a_\eta s_\eta$ 分解。

D) 型列 $(aHu)^2s_\eta(asu)$：通過對型列 $[a_\nu s_\nu(asu)^2]\widehat{\nu_1x}$ 作雙重化簡，由化簡子 $a_\nu s_\nu(asu)^2$ 得到；$(aHu)^2s_\eta(asu)^2$ 的化簡由 $a_\nu s_\nu(asu)^2$ 與 $a_\nu s_\nu(asu)^3$ 得到。由此得到 $(a_\eta,s,u)^4$ 的化簡。

型列 $a_\eta(aHu)s_\eta$：化簡子為 $a_\nu^2s_\nu$；$a_\eta^2s_\eta$ 通過表達式 $(\A s\widehat{\nu x})^3$ 的雙重化簡得到，因為 $a_\eta^2s_\eta$ 並非由型 $a_\eta(aHu)s_\eta$ 的可化簡項 $a_\nu^2s_\nu(\nu\nu_1x)$ 經縮併產生。

型列 $(a_\eta^2su)^2$：化簡子 $a_\nu^2s_\nu$ 來自
\[
a_\nu^2s_\nu s_{\nu_1}=-\frac12a_\eta^2(Hsu)^2.
\]
型 $(a_\eta(aHu),s,u)$ 分解為函數行列式上的轉積式。

由此得到化簡公式：
\srcnumdisplay{(30.)}{%
\begin{aligned}
\text{a)}\quad 0&=(aHu)^2(asu)s_\eta-a_\eta(asu)(Hsu)^2-a_\eta(aHu)(asu)(Hsu),\\
\text{b)}\quad 0&=(aHu)^2(asu)^2s_\eta-a_\eta(aHu)(asu)^2(Hsu),\\
\text{c)}\quad 0&=a_\eta(asu)^2(Hsu)^2,\\
0&=a_\eta s_\eta+\frac14a_\nu^2\cdot g-\frac14s\cdot a_\eta^2,\\
0&=a_\eta(aHu)s_\eta-\frac12a_\eta^2(Hsu),
\quad 0=a_\eta^2(Hsu)^2,\ldots,
\quad 0=a_\eta^2s_\eta.
\end{aligned}
}

\noindent\emph{推論：}

1) $(j\nu x)^2s_\nu$、$\A_\nu s_\nu$、$\A_\nu(\A su)^2$，以及因而 $N_\nu(Nsu)^2$、$(aHu)^2(asu)s_\eta$、$a_\eta(aHu)s_\eta$、$a_\eta^2(Hsu)^2$，都是化簡子；$a_\eta s_\eta$ 是分解型，並在所有一重和二重縮並下化為分解型。

2) 五階型中的 $(5,\lambda)$、$(6,\lambda)$ 等不進入同 $s^2$ 的縮並。


\begin{center}\textbf{\S\ 23.}\quad\emph{十階型（系統 $s_{III}$）。}\end{center}
$s$ 同六階型（系統 III）的縮並（\S\ 11）。

\noindent\emph{不可化簡型：}
\[
\begin{array}{r@{\quad}c}
\text{A)}&
\begin{array}{c|c}
(\vartheta^2\nu x)^3s_\nu&\cdot\\[0.45em]
\cdot&(\vartheta^2\nu x)^3s_\nu s_\vartheta,\ \theta_\nu(\vartheta^2\nu x)^3s_\vartheta,
\end{array}
\end{array}
\]
\[
\begin{array}{r@{\quad}c}
\text{B)}&
\begin{array}{c|ccc}
\cdot&(a_\nu(j\nu x),s,u)^2&(a_\nu(j\nu x),s,u)^3&(a_\nu(j\nu x),s,u)^4\\[0.45em]
a_\nu(j\nu x)s_\nu&\cdot&(a_\nu^2(j\nu x),s,u)^2&(a_\nu^2(j\nu x),s,u)^3,
\end{array}
\end{array}
\]
\[
\begin{array}{r@{\quad}c}
\text{C)}&
\begin{array}{c|c|c|c}
\cdot&\cdot&((\theta Hu),s,u)^2,\ ((\theta Hu)^2,s,u)&\bigstar\\[0.35em]
\cdot&((\vartheta\eta x),s,u)^2,\ (\theta Hu)s_\eta,\ (\theta_\eta su)&((\vartheta\eta x),s,u)^3,\ (\theta_\eta su)^2&\\[0.35em]
((\vartheta\eta x)^2,s,u)&\cdot&(\theta_\eta^2su)&\\[0.35em]
\cdot&(\theta_\eta(\vartheta\eta x)^2,s,u)&\cdot&\cdot
\end{array}
\end{array}
\]
\[
\begin{array}{c@{\quad}|c|c}
\bigstar&((\theta Hu),s,u)^3,\ ((\theta Hu)^2,s,u)^2&((\theta Hu),s,u)^4\\[0.45em]
&(H_\vartheta(\theta Hu),s,u)^2,\ (\theta_\eta(\theta Hu),s,u)^2,\ (\theta_\eta(\theta Hu)^2,s,u)&(\theta_\eta su)^4,\ (\theta_\eta(\theta Hu),s,u)^3\\[0.45em]
&\cdot&(\theta_\eta^2(\theta Hu),s,u)^2.
\end{array}
\]

\noindent\emph{化簡：}

A) 型 $((\vartheta^2\nu x)^3,s,u)^2$、$((\vartheta^2\nu x)^3,s,u)^3$ 分解：
\[
(\vartheta \nu x)(\widehat{\vartheta\nu}su)(\widehat{\vartheta'}\nu su)
=(\vartheta \nu x)s_\vartheta s_{\vartheta'}-(\vartheta\nu x)s_\nu s_{\vartheta'}
=-2\theta\cdot(\vartheta\nu x)s_\nu s_\vartheta+(\vartheta\nu x)\cdot s_\vartheta^2.
\]
其餘諸型或者以化簡子為因子，或者是同分解型的一重縮並。

B) 化簡子為 $a_\nu^2s_\nu$ 與 $(\widehat{j\nu}su)s_\nu$。

C) 1. 型列 $(\theta Hu)^2s_\eta$：a) 化簡子為 $(aHu)^2(asu)s_\eta$；b) 對型列 $(\theta gu)^3g_\nu$ 與 $g_\nu(agu)^3(bgu)^2(abu)$ 作雙重化簡。

由此得到較高型 $(\theta_\eta su)^3$、$(H_\vartheta(\theta Hu),s,u)^3$ 與 $(a_\nu\cdot b_{\nu_1}^2,s,u)^4$ 的化簡［後一型用來代替 $(H_\vartheta su)^4$］。

2. $((\vartheta\eta x),s,u)^4$：化簡子為 $(a_\eta su)^4$。

3. $((\vartheta\eta x)^2,s,u)^3$：化簡子 $s_\vartheta^2s_\nu$ 來自 $(\widehat{\vartheta\eta}su)^2(Hsu)$。
型 $(\vartheta\eta x)s_\eta$、$(\vartheta\eta x)^2s_\eta$、$((\vartheta\eta x)^2,s,u)^2$、$\theta_\eta s_\eta$ 通過同分解型 $a_\eta s_\eta$ 縮並而分解。

4. 型列 $H_\vartheta(\theta Hu)s_\eta$、$\theta_\eta(\theta Hu)s_\eta$、$H_\vartheta(\vartheta\eta x)s_\eta$、$\theta_\eta(\vartheta\eta x)s_\eta$：
化簡子為 $a_\eta(aHu)s_\eta$ 與 $a_\eta^2s_\eta$。

型列 $(\theta_\eta(\vartheta\eta x),s,u)^2$：化簡子為 $a_\eta^2(Hsu)^2$［但 $(\theta_\eta^2(\theta Hu),s,u)^2$ 除外，因為它由項 $a_\eta^2(bsu)(Hsu)$ 經縮併產生］。

5. $(H_\vartheta(\vartheta\eta x),s,u)$、$(H_\vartheta(\vartheta\eta x),s,u)^2$：對 $a_\eta(aHu)s_\eta(abu)$ 與 $g_\vartheta(\theta gu)g_\nu$ 作雙重化簡。

型列 $(H_\vartheta(\vartheta\eta x),s,u)^3$：化簡子為 $((\vartheta \nu x)^2,s,u)^2s_\nu$。

6. $(H_\vartheta(\theta Hu),s,u)$ 依 \S\ 3. 2), c)，通過同分解型 $(\theta_\nu(\vartheta \nu x),s,u)$ 作一重縮並而分解；也就是說，通過對較低的分解型 $(H_\vartheta su)^2$ 作雙重化簡：\srcfn{*)}{符號 $\equiv$ 表示略去了較低次數型的積。}
\[
\begin{aligned}
0&\equiv \theta_\nu(\vartheta\nu\nu_1)(\theta su)+\frac23\theta_\nu^2(\vartheta \nu_1 x)(\theta su)+\theta_\nu(\vartheta \nu x)(\theta su)s_{\nu_1}\\
&\equiv -\frac12H_\vartheta(\theta Hu)(\theta su)+\theta_\eta(\theta su)(Hsu)+\frac12H_\vartheta(\theta su)(Hsu)\\
&\equiv -\frac12H_\vartheta(\theta Hu)(\theta su)+\theta_\eta(\theta su)(Hsu)+\frac12[H_\vartheta]_{\widehat{su}^2}-\frac12\cdot\frac35H_\vartheta(\theta Hu)(\theta su)\\
&\equiv -\frac35H_\vartheta(\theta Hu)(\theta su).
\end{aligned}
\]
依 1. 與 5. 作雙重化簡，得到公式：
\[
\begin{array}{@{}r@{\quad}l@{}}
&\displaystyle
\begin{aligned}
0={}&\theta_\eta(\theta su)(Hsu)^2+\frac14H_\vartheta(\theta Hu)(\theta su)^2-\frac34\theta_\eta(\theta Hu)(\theta su)(Hsu)\\
&-\frac54\theta_\eta(\theta Hu)^2(\theta su)-(asu)^3\cdot b_\eta(bHu)^2-a_\nu(asu)^3\cdot b_\nu^2,
\end{aligned}\\[0.55em]
(31.)&\displaystyle 0=H_\vartheta(\theta Hu)(\theta su)^3-7\theta_\eta(\theta Hu)(\theta su)^2(Hsu),\\[0.55em]
&\displaystyle 0=a_\nu(asu)^2b_{\nu_1}^{2}(bsu)^2-\theta_\eta(\theta su)^2(Hsu)^2+6\theta_\eta(\theta Hu)(\theta su)^2(Hsu),\\[0.55em]
&\displaystyle 0\equiv (H_\vartheta(\vartheta\eta x),s,u),\qquad
0\equiv H_\vartheta(\widehat{\vartheta\eta}su)(Hsu)-4\theta_\eta^2(Hsu).
\end{array}
\]

\noindent\emph{推論：}

1) $(\theta Hu)^2s_\eta$、$H_\vartheta(\theta Hu)s_\eta$、$\theta_\eta(\theta Hu)s_\eta$、$H_\vartheta(\vartheta\eta x)s_\eta$、$\theta_\eta(\vartheta\eta x)s_\eta$、$(H_\vartheta(\vartheta\eta x),s,u)$、$(\theta_\eta(\vartheta\eta x),s,u)^2$ 是化簡子。

2) $s^2$ 不進入同六階型中的 $(5,\lambda)$ 等的縮並。

\begin{center}\textbf{\S\ 24.}\quad\emph{十一階型（系統 $s_{III}$）。}\end{center}
$s$ 同七階型（系統 III）的縮並（\S\ 12）。

\noindent\emph{不可化簡型：}
\[
\text{A)}\quad
\begin{array}{c|c|c|c|c}
\cdot&\cdot&\cdot&((KHu)^2su)^2&\bigstar\\[0.35em]
\cdot&\cdot&(K_\eta su)^2&\cdot&\\[0.35em]
\cdot&\cdot&\cdot&\cdot&\cdot\\[0.35em]
((k\eta x)^3,s,u)&\cdot&\cdot&\cdot&\cdot\\[0.35em]
\cdot&(K_\eta(k\eta x)^3,s,u)&\cdot&\cdot&\cdot
\end{array}
\]
\[
\begin{array}{c@{\quad}|c|c}
\bigstar&((KHu)^2,s,u)^3&((KHu)^2,s,u)^4\\[0.45em]
&\cdot&(K_\eta(KHu)^2,s,u)^3,
\end{array}
\]
\[
\text{B)}\quad ((j\eta x),s,u)^2,\ (H_jsu)\ \big|\ ((j\eta x),s,u)^3,
\qquad \text{C)}\quad (\D_\eta su),
\]
\[
\text{D)}\quad
\begin{array}{c|c|c}
\cdot&((aLu)^2,s,u)^2,\ ((aLu)^3,s,u)&((aLu)^2,s,u)^3\\[0.5em]
(aLu)^2s_l,\ (a_lsu)^2&(a_lsu)^3&(a_l(aLu)^2,s,u)^2.
\end{array}
\]

\noindent\emph{化簡：}

A) 通過同以符號 $\theta\cdot H$ 表示的相應諸型作一重縮並而化簡或分解。

$(K_\eta(KHu)^2,s,u)$ 與 $(K_\eta(KHu)^2,s,u)^2$：依 \S\ 3. 2), a)，通過同 $K_\nu^2(Ksu)$、$K_\nu^2(Ksu)^2$ 縮並而分解。

$(K_\eta(KHu),s,u)^4\equiv(a_\nu^2\cdot\theta_{\nu_1},s,u)^4$：依 \S\ 3. 2), b)，由分解型 $K_\nu^2(Ksu)^3$ 分解。

B) 型列 $(j\eta x)s_\eta$：化簡子為 $(j\nu x)^2s_\nu$。

型 $H_j(\widehat{j\eta}su)(Hsu)$：化簡子為 $(j\nu x)(\widehat{j\nu}su)^2$。

型 $H_j(j\eta x)(Hsu)$：通過用化簡子 $(j\nu x)^2s_\nu$ 化簡分解型 $((j\eta x)^2,s,u)^2$（公式 (18.)a）而分解：
\[
0=-2(j\nu x)^2s_\nu s_{\nu_1}=[(j\eta x)^2]_{\widehat{su}^2}+H_j(j\eta x)(Hsu).
\]
C) 化簡子為 $\D_\nu s_\nu$ 與 $\D_\nu(\D su)^2$。

D) 通過同以符號 $f\cdot H$ 表示的相應諸型作一重縮並而化簡和分解。

E) 由 (30.)c 得到分解型列 $(a_\sigma su)^2$ 的化簡：
\[
\begin{aligned}
0&=a_\eta(Hsu)^2(as\widehat{\nu x})^2
=a_\eta(Hsu)^2(a\widehat{\nu\eta}u)^2+2a_\eta(a\widehat{\nu\eta}u)(\widehat{\nu\eta}su)(Hsu)^2
=\frac13a_\sigma(asu)^2,\\
0&=a_\eta a_\nu(Hsu)^2(as\widehat{\nu x})(asu)
=a_\eta(a\widehat{\nu\eta}u)^2(Hsu)^2(asu)+a_\eta(a\widehat{\nu\eta}u)(\widehat{\nu\eta}su)(Hsu)^2(asu)
=\frac13a_\sigma(asu)^3.
\end{aligned}
\]
\noindent\emph{推論：} $s^2$ 不進入同七階型的縮並。


\begin{center}\textbf{\S\ 25.}\quad\emph{十二、十三、十四、十五階型（系統 $s_{III}$）。}\end{center}
$s$ 同八、九、十、十一階型（系統 III）的縮並（\S\S\ 13--16）。

\noindent\emph{不可化簡型：}

\noindent 十二階。
\[
\begin{adjustbox}{max width=0.98\linewidth}
$\displaystyle
\begin{array}{c@{\qquad}c@{\qquad}c}
\text{A)}&
\begin{array}{c|c}
((\vartheta^2\eta x)^4,s,u)&\cdot\\[0.35em]
\cdot&\theta_\eta(\vartheta^2\eta x)^3s_\vartheta,
\end{array}
&
\text{B)}\quad ((aHu)H_j,s,u)^2\;\big|\;((aHu)H_j,s,u)^3,
\\[1.0em]
\text{C)}&
\begin{array}{c|c|c}
\cdot&((\theta Lu)^2,s,u)^2,\ ((\theta Lu)^3,s,u)&((\theta Lu)^2,s,u)^3\\[0.35em]
(\theta_lsu)^2&\cdot&(\theta_l(\theta Lu)^2,s,u)^2.
\end{array}
&
\end{array}$
\end{adjustbox}
\]
\noindent\emph{化簡：} 由直接同化簡子或分解型縮並而產生的化簡不再提及。

B) $((aHu)H_j,s,u)$ 與 $(a_\eta(aHu)H_j,s,u)$ 分解為函數行列式上的轉積式。

$(a_\eta(aHu)H_j,s,u)^2$ 的分解：對分解型 $[\theta_\nu\cdot\theta_{\nu_1}^{\prime 3}]_{\widehat{su}^3}$（\S\ 13 C）作雙重化簡：
\begin{align*}
0&\equiv[\theta_{\nu_1}^{\prime 3}\cdot\theta_\nu]_{\widehat{su}^3}
\equiv-\frac34\theta_\nu(\theta su)^2(\theta\theta'u)\theta_{\nu_1}^{\prime 3}\\
&\equiv-\frac98(\widehat{\theta\theta'}\eta\widehat{su})(\theta\theta'u)(\theta Hu)(\theta'Hu)\theta_\eta'(\theta su)
\equiv-\frac38a_\eta(aHu)H_j(asu)^2.
\end{align*}

C) 由公式 (31.) 得到分解型
\[
 [L_\vartheta(\theta Lu)]_{\widehat{su}^3}\quad\hbox{与}\quad [\theta_l(\theta Lu)]_{\widehat{su}^3},
\]
即積 $[\theta_\nu^2\cdot H]_{\widehat{su}^3}$ 與 $[a_\nu^2\cdot(bHu)^2]_{\widehat{su}^3}$ 的化簡：
\[
\begin{array}{@{}r@{\quad}l@{}}
&\displaystyle 0\equiv\theta_\nu^2(\theta su)^2(Hsu),\qquad
0\equiv [L_\vartheta(\theta Lu)]_{\widehat{su}^4}-7[\theta_l(\theta Lu)]_{\widehat{su}^4},\\[0.55em]
&\displaystyle 0\equiv a_\nu^2(asu)^2(bHu)^2(bsu)+6\theta_l(\theta Lu)^2(\theta su)(Lsu),\\[0.55em]
(32.)&\displaystyle 0\equiv\theta_\nu^3(\theta su)(Hsu)^2
\quad\hbox{由 }0\equiv\theta_l^2(\theta Lu)(\theta su)(Lsu)^2\\[-0.05em]
&\displaystyle\hfill(\hbox{化简子 }a_\eta^2(Hsu)^2).
\end{array}
\]

\noindent\emph{不可化簡型。}

\noindent 十三階。
\[
\begin{array}{rl}
\text{A)}&((KLu)^3,s,u)^2;\ ((KLu)^3,s,u)^3,\\[0.25em]
\text{B)}&(L_jsu)^2,\\[0.25em]
\text{C)}&((aH^2u)^3,s,u)^2.
\end{array}
\]
\noindent 十四階。
\[
\begin{array}{rl}
\text{A)}&((aLu)^2L_j,s,u)^2,\\[0.25em]
\text{B)}&((\theta H^2u)^3,s,u)^2.
\end{array}
\]
\noindent 十五階。
\[
 ((KH^2u)^4,s,u)^2.
\]
\noindent\emph{化簡：} $(K_l(KLu)^3,s,u)^2$、$(H_\vartheta(\theta H^2u)^3,s,u)$、$((K,H\cdot L,u)^5,s,u)$ 依 \S\ 3. 2), c) 分解；也就是說，同時考察分解型 $(K_l(KLu)^2,s,u)^3$、$(H_\vartheta(\theta H'u)^2,s,u)^2$、$((K,H\cdot L,u)^4,s,u)^2$。

\noindent\emph{推論：} $s^2$ 不進入同系統 III 的各個型的縮並。

\begin{center}\textbf{\S\ 26.}\quad\emph{系統 $s_{IV}$。}\end{center}
\noindent 1) \emph{$s$ 同系統 I$\cdot$III 的縮並。}

\noindent\emph{化簡：} 依前幾節的化簡（$(a_\nu^2s_\nu,\ (aHu)^2(asu)s_\eta$ 等），型 $(0,\lambda)$、$(1,\lambda)$、$(2,\lambda)$ 不容許同時作 I 與 III 的縮並；因此依 \S\ 18，為構造系統 $s_{IV}$，$f,\D,N$ 不以積進入縮並。

$j$ 不進入縮並，因為依同樣的化簡，與關於 $j$ 逆變的縮並
\[
(K_\nu(k\nu x)^3,s,u),\qquad ((\vartheta^2\eta x)^4,s,u)\quad\hbox{等等}
\]
相對應的是關於 $j$ 同變的縮並：
\[
K_\nu(k\nu x)^2s_k,
\qquad (\vartheta^2\eta x)^3s_\vartheta\quad\hbox{等等。}
\]

\noindent 2) \emph{$s$ 同系統 II$\cdot$III 的縮並}
也就是說，$s$ 同 $H\cdot(1,\lambda)$、$H\cdot(2,\lambda)$、$H\cdot(3,\lambda)$ 的縮並。

\noindent\emph{不可化簡型：}
\[
\begin{array}{rl}
\text{A)}&(a_\nu\cdot H,s,u)^4,
((aHu)^2\cdot H,s,u)^3,
((aHu)^2\cdot H,s,u)^4,\\[0.25em]
&((aLu)^2\cdot H,s,u)^4,
((aLu)^3\cdot H,s,u)^3,
((aH^2u)^3\cdot H,s,u)^4.\\[0.45em]
\text{B)}&(a_\nu^2\cdot H,s,u)^3,
(a_\eta(aHu)^2\cdot H,s,u)^3,
(a_l(aLu)^2\cdot H,s,u)^4.\\[0.45em]
\text{C)}&(\theta_\nu\cdot H,s,u)^4,
((\theta Hu)^2\cdot H,s,u)^3,
((\theta Hu)^2\cdot H,s,u)^4,\\[0.25em]
&((\theta Lu)^2\cdot H,s,u)^4,
((\theta Lu)^3\cdot H,s,u)^3,
((\theta H^2u)^3\cdot H,s,u)^4,\\[0.45em]
\text{D)}&(\theta_\eta(\theta Hu)^2\cdot H,s,u)^3,
(\theta_l(\theta Lu)^2\cdot H,s,u)^4,\\[0.45em]
\text{E)}&((KLu)^3\cdot H,s,u)^4,
((KH^2u)^4\cdot H,s,u)^4,\\[0.45em]
\text{F)}&((j\nu x)^2\cdot H,s,u)^3,
((j\nu x)^2\cdot H,s,u)^4,
((j\eta x)\cdot H,s,u)^4,\\[0.25em]
&(H_j\cdot H,s,u)^3,
(L_j\cdot H,s,u)^4.
\end{array}
\]
\noindent\emph{化簡：} $\nu$ 與 $j$ 類似，不進入縮並。

B) 與 D)。$(a_\nu^2\cdot H,s,u)^4$ 和 $(\theta_\nu^2\cdot H,s,u)^4$ 由公式 (27.)c) 與 (29.)c) 分解，並須計及 $(gHu)^2=-\frac13s\cdot\sigma$：
\[
(a_\nu^2\cdot H,s,u)^4\equiv\frac13(asu)^4\cdot\sigma,
\qquad
(\theta_\nu^2\cdot H,s,u)^4\equiv-\frac16(\theta su)^4\cdot\sigma;
\]
由此，$(a_\eta(aHu)\cdot H,s,u)^4$、$(\theta_\eta(\theta Hu)\cdot H,s,u)^4$ 分解。

$(\theta_\nu^2\cdot H,s,u)^3$、$(\theta_\nu^3\cdot H,s,u)^3$，以及由此 $(\theta_\eta^2(\theta Hu)\cdot H,s,u)^4$，依 \S\ 25 十二階型 C) 分解。

\noindent 3) \emph{$s$ 同系統 III$\cdot$III 的縮並}，也就是說，$s$ 同系統 III 中如下型的縮並：
\[
(3,\lambda)\cdot(3,\lambda),\quad (3,\lambda)\cdot(2,\lambda),\quad (3,\lambda)\cdot(1,\lambda),\quad
(2,\lambda)\cdot(2,\lambda),\quad (2,\lambda)\cdot(1,\lambda),\quad
(1,\lambda)\cdot(1,\lambda).
\]
\noindent\emph{不可化簡型：}
\[
\begin{array}{rl}
\text{A)}&(a_\nu^2\cdot a_\nu^2,s,u)^3,
(a_\nu^2\cdot a_\nu^2,s,u)^4,
(a_\nu^2\cdot a_\eta(aHu)^2,s,u)^3,\\[0.45em]
\text{B)}&(a_\nu\cdot H_j,s,u)^4,
((aHu)^2\cdot H_j,s,u)^3,
((aLu)^2\cdot H_j,s,u)^4,\\[0.25em]
&((aLu)^3\cdot H_j,s,u)^2,
((aH^2u)^3\cdot H_j,s,u)^3.
\end{array}
\]
\noindent\emph{化簡：}

在符號 $\nu$ 與 $f$ 中階數相同時，積 II$\cdot$III 應視為高於積 III$\cdot$III。

這些化簡一部分由積之間的關係（合式）產生，一部分則通過以相應的分解型代替這些積，並化簡這些型同 $s$ 的縮並而產生。（含逆變式的積一律略去，因為它們會轉化為積。）

A) $f$ 為二次，符號 $\nu$ 為任意階。依 \S\ 11 公式 (12.)，並經類似計算，有：
\begin{align*}
1)\quad 0&=a_\nu\cdot b_{\nu_1}+\frac12f\cdot(aHu)^2-\frac14\theta\cdot H+\frac12u_xH_\vartheta-\frac14u_x^2\cdot H_\vartheta^2,\\
2)\quad 0&=a_\nu\cdot b_{\nu_1}^2+f\cdot a_\eta(aHu)^2-\theta_\eta-\frac12H_\vartheta,\\
3)\quad 0&=a_\nu^2\cdot b_{\nu_1}^2-H_\vartheta^2+2\theta_\eta^2.
\end{align*}
由此得：
\begin{align*}
4)\quad 0&=a_\nu^2\cdot b_{\nu_1}^2(j\nu_1x)\quad(\text{由 }3)),\\
5)\quad L_\vartheta(\theta Lu)&=-a_\nu\cdot b_\eta(bHu)^2+\frac34a_\nu^2\cdot(bHu)^2+\frac34\theta_\nu^2\cdot H\quad(\text{由 }2)),\\
6)\quad L_\vartheta(\theta Lu)&=-6a_\nu\cdot b_\eta(bHu)^2+2a_\nu^2\cdot(bHu)^2-\frac12\theta_\nu^2\cdot H\quad(\text{由 }1)),\\
7)\quad 0&=a_\nu\cdot(bHu)^2+f\cdot(aLu)^3+\theta_\nu\cdot H\quad(\text{由 }1)),\\
\end{align*}
\[
\left.
\begin{aligned}
8)\quad 0&=a_\nu\cdot(bLu)^3+\frac34(\theta Hu)^2\cdot H\\
9)\quad 0&=(aHu)^2\cdot(bHu)^2+2(\theta Hu)^2\cdot H
\end{aligned}
\right\}\quad\text{（由 7）），}
\]
\begin{align*}
10)\quad 0&=a_\eta(aHu)^2\cdot b_\eta(bHu)^2\quad(\text{由 }5)\text{ 與 }6)),\\
11)\quad 0&= [a_\nu^2\cdot b_\eta(bHu)]_{\widehat{su}^4}\quad(\text{由公式 (30.)b}).
\end{align*}
對
\[
(H_\vartheta su)^4,
\qquad (L_\vartheta(\theta Lu),s,u)^3,
\qquad (L_\vartheta(\theta Lu),s,u)^4
\]
的化簡（公式 (31.) 與 (32.)），連同公式 1) 至 11)，給出 $s$ 同如下各積縮並時的全部化簡：這些積關於符號 $f$ 為二次，而關於 $\nu$ 為任意階。

B) 符號 $f,\theta,j$ 中每次取兩個所成的積；$\nu$ 為任意階。依公式 (17.)、(18.)、(20.) 並經直接計算，有如下關係：
\[
\left.
\begin{aligned}
1)\quad 0&=a_\nu\cdot\theta_{\nu_1}+\frac14\theta\cdot(aHu)^2+\frac14f\cdot(\theta Hu)^2,\\
2)\quad 0&=\theta_\nu\cdot\theta_{\nu_1}+\frac12\theta\cdot(\theta Hu)^2
\end{aligned}
\right\}\quad\text{（直接計算，類似於 A 1））。}
\]
由此得：
\[
\left.
\begin{aligned}
3)\quad 0&=3a_\nu\cdot(\theta Hu)^2+\theta\cdot(aLu)^3+2f\cdot(\theta Lu)^3\\
4)\quad 0&=3\theta_\nu\cdot(aHu)^2+f\cdot(\theta Lu)^3+2\theta\cdot(aLu)^3
\end{aligned}
\right\}\quad\text{（由 1) 与 A 6)）），}
\]
\begin{align*}
5)\quad 0&=a_\nu\cdot(\theta Lu)^3,\quad 0=\theta_\nu\cdot(aLu)^3,\quad 0=(aHu)^2\cdot(\theta Hu)^2
\quad\text{（由 3)、4) 與 A 8)）），}\\
6)\quad 0&=a_\nu\cdot\theta_\nu^2+\theta_\nu\cdot a_\nu^2+f\cdot\theta_\eta(\theta Hu)^2+\theta\cdot a_\eta(aHu)^2\quad(\text{公式 (17.)}),\\
7)\quad 0&=\theta_\nu\cdot\theta_\nu^2+\theta\cdot\theta_\eta(\theta Hu)^2+\frac16f\cdot H_j\\
&\qquad [\text{仿照 }6)\text{，由下式推得：}\theta_\nu^2(\theta\theta'\widehat{\nu}_1x)],\\
8)\quad 0&=a_\nu\cdot(j\nu x)^2+f\cdot H_j\quad(\text{由 }6)),\\
9)\quad 0&=(aHu)^2\cdot(j\nu x)^2-2a_\nu\cdot H_j\quad(\text{由 }8)),\\
10)\quad 0&=\theta_\nu\cdot(j\nu x)^2,\quad 0=\theta\cdot H_j\quad(\text{由 }7)\text{ 與 }8)),
\end{align*}
\[
\left.
\begin{aligned}
11)\quad 0&=a_\nu\cdot\theta_\nu^3-\frac34(j\eta x)^2,\\
12)\quad 0&=a_\nu^2\cdot\theta_\nu^2-\frac13(j\eta x)^2-\frac13(\D Hu)^2
\end{aligned}
\right\}\quad\text{（公式 (18.)），}
\]
\begin{align*}
13)\quad 0&=a_\nu^2\cdot\theta_\nu^3-\frac12a_\sigma\quad(\text{公式 (20.)}),\\
14)\quad 0&=\theta_\nu\cdot\theta_\nu^3,
&0&=a_\eta H_j\quad(\S\ 13. C)).
\end{align*}
這些公式一直推導到諸積可以極化為止；也就是說，縮並只會產生\emph{一個}新積，這是因為
\begin{center}
\begin{tabular}{rl}
a)&積同自身的縮並可化簡，\\
b)&縮並所產生的其餘積可化簡
\end{tabular}
\end{center}
或導向逆變式（參看 \S\ 3. 2), a) 與 b)）。

公式 1) 至 5) 給出 $s$ 同由 $a_\nu\cdot\theta_\nu$ 經縮併產生的積作縮並時的全部化簡，

公式 6) 至 10) 給出由 $a_\nu\cdot\theta_\nu^2$ 經縮併產生的諸積的化簡。

依 \S\ 24 A)，並計及公式 6) 至 10)，$s$ 同由 $a_\nu^2\cdot\theta_\nu$ 產生的諸積的縮並均分解。

有：
\begin{align*}
0&\equiv a_\nu^2(asu)^2\theta_{\nu_1}(\theta su)^2,
&0&\equiv a_\nu^2(asu)\theta_{\nu_1}(\theta su)^3-K_\eta(KHu)^2(Ksu)^3,\\
0&\equiv a_\nu^2(asu)^2\theta_{\nu_1}(\theta su),
&0&\equiv a_\nu^2(asu)\theta_{\nu_1}(\theta su)^2,\\
0&\equiv a_\nu^2(asu)\theta_{\nu_1}(\theta su).
\end{align*}
化簡子
\[
 ((j\eta x)^2,s,u)^3
 \ [\hbox{由 }(\vartheta\eta\widehat{su})^2(Hsu).\ \S\ 23.\ C\ 3)],
\]
\[
(H_j(j\eta x),s,u)^2\ [\S\ 24,\ B],
\quad ((\D Hu)^2,s,u)^2\ [\S\ 24,\ C],
\]
\[
(a_\sigma su)^2\ [\S\ 24,\ E]
\]
連同公式 11) 至 14)，給出由
$a_\nu\cdot\theta_{\nu_1}^3$、$a_\nu^2\cdot\theta_{\nu_1}^2$、$a_\nu^2\cdot\theta_{\nu_1}^3$ 產生的全部積的化簡。

以上計算可以概括為如下
\emph{結論：}

相對完備系統 $\modu{(\R,t)}$ 由以下 331 個型與子系統組成；我們隨後還將它們按表格排列列出：
\[
\left.
\begin{array}{lr}
u_x&1\ \text{個型}\\
\text{系統 I}&6\ \text{個型}\\
\text{系統 II}&5\ \text{個型}\\
\text{系統 III}&117\ \text{個型}
\end{array}
\right\}
\quad
\begin{gathered}
\text{相對完備系統，模 }(s,\varrho,t)\\
[\text{見表 I}]\text{。}
\end{gathered}
\]
\[
\left.
\begin{array}{lr}
s&1\ \text{個型}\\
\text{系統 }s_I&35\ \text{個型}\\
\text{系統 }s_{II}&9\ \text{個型}\\
\text{系統 }s_{III}&125\ \text{個型}\\
\text{系統 }s_{IV}&32\ \text{個型}
\end{array}
\right\}
\quad [\text{見表 II}]\text{。}
\]

\clearpage
\begin{landscape}
\thispagestyle{plain}
\begin{center}\bfseries 表 I（見第 90 頁）。\end{center}
\vspace{-0.6em}
\begingroup
\setlength{\tabcolsep}{2pt}
\renewcommand{\arraystretch}{1.35}
\tiny
\begin{center}
\begin{adjustbox}{max width=0.98\linewidth,max totalheight=0.72\textheight,center}
\begin{tabular}{|c|c|c|c|c|c|c|c|c|c|c|c|c|c|c|c|c|c|c|c|c|c|c|}
\hline
{} & 0 & 1 & 2 & 3 & 4 & 5 & 6 & 7 & 8 & 9 & 10 & 11 & 12 & 13 & 14 & 15 & 16 & 17 & 18 & 19 & 20 & 21 \\ \hline
0 & \mcell{i} &  &  &  & \mcell{f} &  & \mcell{\A} &  & \mcell{K_{\nu}(k\nu x)^3} &  & \mcell{K_{\eta}(k\eta x)^3} &  &  &  & \mcell{(\vartheta^2\eta x)^4} & \mcell{H_k(k\cdot\vartheta,\eta x)^4} &  & \mcell{\theta_l(\vartheta\cdot k,lx)^5} &  &  &  & \mcell{(\vartheta^3lx)^6} \\ \hline
1 &  & \mcell{u_x} &  &  &  & \mcell{\theta_\nu(\vartheta\nu x)^2} &  & \mcell{\theta_\eta(\vartheta\eta x)^2} & \mcell{N} & \mcell{(k\nu x)^3} & \mcell{\theta_\nu(\vartheta^2\nu x)^3} & \mcell{(k\eta x)^3} & \mcell{H_{\vartheta}(\vartheta^2\eta x)^3\\ \theta_\eta(\vartheta^2\eta x)^3} &  & \mcell{\theta_l(\vartheta^2lx)^4} &  & \mcell{(\vartheta k,\eta,x)^4} &  & \mcell{(\vartheta k,l,x)^5} &  &  &  \\ \hline
2 &  &  & \mcell{a_\nu^2} &  & \mcell{\theta\\ a_\eta^2} &  & \mcell{(\vartheta\nu x)^2} & \mcell{K_\nu(k\nu x)^2} & \mcell{(\vartheta\eta x)^2} & \mcell{H_k(k\eta x)^2\\K_\eta(k\eta x)^2} &  & \mcell{(\vartheta^2\nu x)^3\\ K_l(klx)^3} &  & \mcell{(\vartheta^2\eta x)^3} &  & \mcell{(\vartheta^2lx)^4} &  &  &  &  &  &  \\ \hline
3 &  & \mcell{\theta_\nu^3} &  & \mcell{a_\nu} & \mcell{\theta_\nu(\vartheta\nu x)} & \mcell{\A_\nu\\a_\eta} & \mcell{K\\H_{\vartheta}(\vartheta\eta x)\\\theta_\eta(\vartheta\eta x)} & \mcell{\A_\eta} & \mcell{(k\nu x)^2\\\theta_l(\vartheta lx)^2} &  & \mcell{(k\eta x)^2} &  & \mcell{(klx)^3} &  &  &  &  &  &  &  &  &  \\ \hline
4 & \mcell{\nu} &  & \mcell{H\\\theta_\nu^2} & \mcell{(j\nu x)^3\\a_\eta(aHu)} & \mcell{\theta_\eta^2} & \mcell{(\vartheta\nu x)\\a_l^2} &  & \mcell{N_\nu\\(\vartheta\eta x)} &  & \mcell{H_n\\N_\eta\\(\vartheta lx)^2} &  &  &  &  &  &  &  &  &  &  &  &  \\ \hline
5 &  & \mcell{a_\eta(aHu)^2} & \mcell{\theta_\eta^2(\theta Hu)} & \mcell{\theta_\nu} & \mcell{K_\nu^2\\(aHu)} & \mcell{\theta_\eta} & \mcell{K_\eta^2\\(\A Hu)\\a_l} &  & \mcell{\A_l} &  &  &  &  &  &  &  &  &  &  &  &  &  \\ \hline
6 & \mcell{j\\\sigma} &  & \mcell{(j\nu x)^2\\(aHu)^2} & \mcell{L\\a_\nu^2(j\nu x)\\H_{\vartheta}(\theta Hu)\\\theta_\eta(\theta Hu)} &  & \mcell{K_\nu\\\theta_l^2} &  & \mcell{K_\eta} &  &  &  &  &  &  &  &  &  &  &  &  &  &  \\ \hline
7 &  & \mcell{\theta_\eta(\theta Hu)^2} & \mcell{H_j(j\eta x)\\a_l(aLu)^2} &  & \mcell{a_\nu(j\nu x)\\(\theta Hu)} &  & \mcell{\A_\nu(j\nu x)\\\theta_l} & \mcell{K_l^2} &  &  &  &  &  &  &  &  &  &  &  &  &  &  \\ \hline
8 & \mcell{H_j^2\\a_l(aLu)^3} & \mcell{(\nu\sigma x)\\(j\nu x)} & \mcell{(\theta Hu)^2} & \mcell{K_\eta(KHu)^2\\(j\eta x)\\(aLu)^2} &  &  &  &  &  &  &  &  &  &  &  &  &  &  &  &  &  &  \\ \hline
9 &  & \mcell{H_j\\(aLu)^3} & \mcell{a_\eta(aHu)H_j\\L_{\vartheta}(\theta Lu)^2\\\theta_l(\theta Lu)^2} &  & \mcell{(KHu)} &  &  &  &  &  &  &  &  &  &  &  &  &  &  &  &  &  \\ \hline
10 & \mcell{\theta_l(\theta Lu)^3} & \mcell{L_j^2\\(j\sigma x)\\a_\eta(aH^2u)^3} &  & \mcell{H_j(aHu)\\(\theta Lu)^2} &  &  &  &  &  &  &  &  &  &  &  &  &  &  &  &  &  &  \\ \hline
11 &  & \mcell{(\theta Lu)^3} & \mcell{K_l(KLu)^3\\L_j\\(aH^2u)^3} &  &  &  &  &  &  &  &  &  &  &  &  &  &  &  &  &  &  &  \\ \hline
12 & \mcell{(aH^2u)^4} & \mcell{a_l(aLu)^2L_j\\H_{\vartheta}(\theta H^2u)^3\\\theta_\eta(\theta H^2u)^3} &  & \mcell{(KLu)^3} &  &  &  &  &  &  &  &  &  &  &  &  &  &  &  &  &  &  \\ \hline
13 & \mcell{(\nu\sigma j)} &  & \mcell{(aLu)^2L_j\\(\theta H^2u)^3} &  &  &  &  &  &  &  &  &  &  &  &  &  &  &  &  &  &  &  \\ \hline
14 & \mcell{(\theta H^2u)^4} & \mcell{K_\eta(KH^2u)^4\\(a,H\!\cdot L,u)^4} &  &  &  &  &  &  &  &  &  &  &  &  &  &  &  &  &  &  &  &  \\ \hline
15 & \mcell{L_\vartheta(\theta,H\!\cdot L,u)^4} &  & \mcell{(KH^2u)^4} &  &  &  &  &  &  &  &  &  &  &  &  &  &  &  &  &  &  &  \\ \hline
16 &  & \mcell{(\theta,H\!\cdot L,u)^4} &  &  &  &  &  &  &  &  &  &  &  &  &  &  &  &  &  &  &  &  \\ \hline
17 & \mcell{K_\eta(K,H\!\cdot L,u)^5} &  &  &  &  &  &  &  &  &  &  &  &  &  &  &  &  &  &  &  &  &  \\ \hline
18 &  & \mcell{(K,H\!\cdot L,u)^5} &  &  &  &  &  &  &  &  &  &  &  &  &  &  &  &  &  &  &  &  \\ \hline
19 &  &  &  &  &  &  &  &  &  &  &  &  &  &  &  &  &  &  &  &  &  &  \\ \hline
20 &  &  &  &  &  &  &  &  &  &  &  &  &  &  &  &  &  &  &  &  &  &  \\ \hline
21 & \mcell{(KH^3u)^6} &  &  &  &  &  &  &  &  &  &  &  &  &  &  &  &  &  &  &  &  &  \\ \hline
\end{tabular}
\end{adjustbox}
\end{center}
\vspace{-0.2em}
\begin{center}\footnotesize（最上方橫行中的數字表示關於 $x$ 的次數，最左側豎行中的數字表示關於 $u$ 的次數。）\end{center}
\endgroup
\end{landscape}
\clearpage
\begin{landscape}
\thispagestyle{plain}
\begin{center}\bfseries 表 II（見第 90 頁）。\end{center}
\vspace{-0.6em}
\begingroup
\setlength{\tabcolsep}{2pt}
\renewcommand{\arraystretch}{1.35}
\tiny
\begin{center}
\begin{adjustbox}{max width=0.98\linewidth,max totalheight=0.72\textheight,center}
\begin{tabular}{|c|c|c|c|c|c|c|c|c|c|c|c|c|c|c|c|c|c|}
\hline
{} & 0 & 1 & 2 & 3 & 4 & 5 & 6 & 7 & 8 & 9 & 10 & 11 & 12 & 13 & 14 & 15 & 16 \\ \hline
0 & \mcell{J} &  &  &  & \mcell{s} &  & \mcell{s_\vartheta^2} &  &  &  &  & \mcell{s_n} & \mcell{(k\nu x)^3s_\nu} & \mcell{(\vartheta^2\nu x)^3s_\nu s_\vartheta\\\theta_\nu(\vartheta^2\nu x)^3s_\vartheta} &  & \mcell{\theta_\eta(\vartheta^2\eta x)^3s_\vartheta} &  \\ \hline
1 &  &  &  &  &  &  & \mcell{(asu)} & \mcell{s_\vartheta} & \mcell{s_k^2\\(\Delta su)} & \mcell{s_n(Nsu)\\(\vartheta\nu x)^2s_\nu} & \mcell{((k\nu x)^3,s,u)s_\nu\\(K_\nu(k\nu x)^3,s,u)} &  & \mcell{(K_\eta(k\eta x)^3,s,u)} &  & \mcell{(\vartheta^2\nu x)^3s_\nu} &  & \mcell{((\vartheta^2\eta x)^4,s,u)} \\ \hline
2 &  &  &  &  & \mcell{(asu)^2} & \mcell{s_\vartheta(\theta su)} & \mcell{(\Delta su)^2\\a_\nu s_\nu} & \mcell{((\vartheta\nu x)^2,s,u)s_\nu\\(\theta_\nu(\vartheta\nu x)^2,s,u)} &  & \mcell{s_k\\(\theta_\eta(\vartheta\eta x)^2,s,u)} &  & \mcell{(k\nu x)^2s_\nu\\((k\nu x)^3,s,u)} &  & \mcell{((k\eta x)^3,s,u)} &  &  &  \\ \hline
3 &  &  & \mcell{(asu)^3} & \mcell{s_\vartheta(\theta su)^2\\s_\nu} & \mcell{a_\nu s_\nu(asu)\\a_\nu^2(asu)} & \mcell{s_\eta} & \mcell{(\theta su)\\(a_\eta^2su)} & \mcell{s_k(Ksu)} & \mcell{(Nsu)^2\\(\vartheta\nu x)s_\nu\\((\vartheta\nu x)^2,s,u)} & \mcell{((k\nu x)^3,s,u)^2} & \mcell{((\vartheta\eta x)^2,s,u)} &  &  &  &  &  &  \\ \hline
4 & \mcell{(asu)^4} & \mcell{s_\vartheta(\theta su)^3} & \mcell{a_\nu^2(asu)^2} & \mcell{(\theta_\nu^3su)} & \mcell{(\theta su)^2} & \mcell{s_k(Ksu)^2\\a_\nu(asu)\\s_\nu(asu)} & \mcell{((\vartheta\nu x)^2,s,u)^2\\\theta_\nu s_\nu} & \mcell{s_\nu(\Delta su)\\(\Delta_\nu su)\\(aHu)s_\eta\\(a_\eta su)} &  & \mcell{(\Delta_\eta su)} &  &  &  &  &  &  &  \\ \hline
5 &  &  & \mcell{(\theta su)^3} & \mcell{s_k(Ksu)^3,s_j,\\s_\vartheta(\theta s^\prime u)^4\\s_\nu(asu)^2\\a_\nu(asu)^2} & \mcell{(Hsu)\\((\vartheta\nu x)^2,s,u)^3\\(\theta_\nu(\vartheta\nu x),s,u)^2\\(\theta_\nu^2su)} & \mcell{((j\nu x)^3,s,u)\\(aHu)^2s_\eta\\(a_\eta su)^2} & \mcell{(Ksu)^2\\s_l\\(\theta_\eta^2su)} &  & \mcell{((k\nu x)^2,s,u)^2\\K_\nu s_\nu} &  &  &  &  &  &  &  &  \\ \hline
6 & \mcell{(\theta su)^4} & \mcell{s_\nu(asu)^3\\a_\nu(asu)^3} & \mcell{(Hsu)^2\\(\theta_\nu(\vartheta\nu x),s,u)^3\\(\theta_\nu^2su)^2} & \mcell{(a_\eta su)^3\\(a_\eta(aHu),s,u)^2\\(a_\eta(aHu)^2,s,u)} & \mcell{(Ksu)^3} & \mcell{s_j(asu)\\((\vartheta\nu x),s,u)^2\\(\theta_\nu su)} & \mcell{((k\nu x)^2,s,u)^3} & \mcell{s_j(\Delta su)\\a_\nu(j\nu x)s_\nu\\((\vartheta\eta x),s,u)^2\\(\theta Hu)s_\eta\\(\theta_\eta su)} &  &  &  &  &  &  &  &  &  \\ \hline
7 & \mcell{(\theta_\nu(\vartheta\nu x),s,u)^4} & \mcell{(a_\eta(aHu),s,u)^3} & \mcell{(Ksu)^4\\(\theta_\eta^2(\theta Hu),s,u)^2\\(a_\nu^2\cdot a_\nu^2,s,u)^3} & \mcell{s_j(asu)^2\\s_k(Ks^2u)^5\\((\vartheta\nu x),s,u)^3\\(\theta_\nu su)^2} & \mcell{(j\nu x)s_\nu\\((j\nu x)^2,s,u)\\((aHu),s,u)^2\\((aHu)^2,s,u)} & \mcell{((\vartheta\eta x),s,u)^3\\(\theta_\eta su)^2} & \mcell{(aLu)^2s_l\\(a_lsu)^2} &  &  &  &  &  &  &  &  &  &  \\ \hline

8 & \mcell{(a_\nu^2\cdot a_\nu^2,s,u)^4} & \mcell{s_j(asu)^3\\s_k(Ks^2u)^6\\((\vartheta\nu x),s,u)^4\\(\theta_\nu su)^3} & \mcell{((j\nu x)^2,s,u)^2\\((aHu),s,u)^3\\((aHu)^2,s,u)^2} & \mcell{(Lsu)^2\\(a_\nu^2(j\nu x),s,u)^2\\(H_\vartheta(\theta Hu),s,u)^2\\(\theta_\eta(\theta Hu),s,u)^2\\(\theta_\eta(\theta Hu)^2,s,u)} & \mcell{(a_lsu)^3} & \mcell{(K_\nu su)^2} &  & \mcell{(K_\eta su)^2} &  &  &  &  &  &  &  &  &  \\ \hline
9 & \mcell{(K_\nu^2su)^4\\((aHu),s,u)^4} & \mcell{(Lsu)^3\\(a_\nu^2(j\nu x),s,u)^3\\(\theta_\eta su)^4\\(\theta_\eta(\theta Hu),s,u)^3} & \mcell{(Ks^2u)^6\\(a_l(aLu)^2,s,u)^2\\(a_\nu^2\cdot H,s,u)^3} & \mcell{(K_\nu su)^3} & \mcell{(a_\nu(j\nu x),s,u)^2\\((\theta Hu),s,u)^2\\((\theta Hu)^2,s,u)} &  & \mcell{(\theta_lsu)^2} &  &  &  &  &  &  &  &  &  &  \\ \hline
10 &  & \mcell{(K_\nu su)^4\\(a_\nu^2\cdot a_\eta(aHu)^2,s,u)^3} & \mcell{(a_\nu(j\nu x),s,u)^3\\((\theta Hu),s,u)^3\\((\theta Hu),s,u)^2} & \mcell{((j\eta x),s,u)^2\\(H_jsu)\\((aLu)^2,s,u)^2\\((aLu)^3,s,u)} &  &  &  &  &  &  &  &  &  &  &  &  &  \\ \hline
11 & \mcell{(a_\nu(j\nu x),s,u)^4\\((\theta Hu),s,u)^4} & \mcell{(K_\eta(KHu)^2,s,u)^3\\((j\eta x),s,u)^3\\((aLu)^2,s,u)^3\\(a_\nu\cdot H,s,u)^4} & \mcell{(H^2su)^3\\(\theta_l(\theta Lu)^2,s,u)^2} &  & \mcell{((KHu)^2,s,u)^2} &  &  &  &  &  &  &  &  &  &  &  &  \\ \hline
12 &  & \mcell{(a_\eta(aHu)^2\cdot H,s,u)^3} & \mcell{((KHu)^2,s,u)^3} & \mcell{((aHu)H_j,s,u)^2\\((\theta Lu)^2,s,u)^2\\((\theta Lu)^3,s,u)} &  &  &  &  &  &  &  &  &  &  &  &  &  \\ \hline
13 & \mcell{((KHu)^2,s,u)^4} & \mcell{((aHu)H_j,s,u)^3\\((\theta Lu)^2,s,u)^3\\(\theta_\nu\cdot H,s,u)^4} & \mcell{(L_jsu)^2\\((aH^2u)^3,s,u)^2\\((j\nu x)^2\cdot H,s,u)^3\\((aHu)^2\cdot H,s,u)^3} &  &  &  &  &  &  &  &  &  &  &  &  &  &  \\ \hline
14 & \mcell{((j\nu x)^2\cdot H,s,u)^4\\((aHu)^2\cdot H,s,u)^4} & \mcell{(\theta_\eta(\theta Hu)^2\cdot H,s,u)^3} &  & \mcell{((KLu)^3,s,u)^2} &  &  &  &  &  &  &  &  &  &  &  &  &  \\ \hline
15 & \mcell{(a_l(aLu)^2\cdot H,s,u)^4} & \mcell{((KLu)^3,s,u)^3} & \mcell{((aLu)^2L_j,s,u)^2\\((\theta H^2u)^3,s,u)^2\\((\theta Hu)^2\cdot H,s,u)^3} &  &  &  &  &  &  &  &  &  &  &  &  &  &  \\ \hline
16 & \mcell{((\theta Hu)^2\cdot H,s,u)^4\\(a_\nu\cdot H_j,s,u)^4} & \mcell{((j\eta x)\cdot H,s,u)^4\\(H_j\cdot H,s,u)^3\\((aLu)^2\cdot H,s,u)^4\\((aLu)^3\cdot H,s,u)^3} &  &  &  &  &  &  &  &  &  &  &  &  &  &  &  \\ \hline
17 & \mcell{(\theta_l(\theta Lu)^2\cdot H,s,u)^4} &  & \mcell{((KH^2u)^4,s,u)^2} &  &  &  &  &  &  &  &  &  &  &  &  &  &  \\ \hline
18 &  & \mcell{((\theta Lu)^2\cdot H,s,u)^4\\((\theta Lu)^3\cdot H,s,u)^3\\((aHu)^2\cdot H_j,s,u)^3} &  &  &  &  &  &  &  &  &  &  &  &  &  &  &  \\ \hline
19 & \mcell{((aH^2u)^3\cdot H,s,u)^4\\(L_j\cdot H,s,u)^4} &  &  &  &  &  &  &  &  &  &  &  &  &  &  &  &  \\ \hline
20 &  & \mcell{((KLu)^3\cdot H,s,u)^4} & \mcell{((aLu)^3\cdot H_j,s,u)^2} &  &  &  &  &  &  &  &  &  &  &  &  &  &  \\ \hline
21 & \mcell{((\theta H^2u)^3\cdot H,s,u)^4\\((aLu)^2\cdot H_j,s,u)^4} &  &  &  &  &  &  &  &  &  &  &  &  &  &  &  &  \\ \hline
22 &  &  &  &  &  &  &  &  &  &  &  &  &  &  &  &  &  \\ \hline
23 & \mcell{((KH^2u)^4\cdot H,s,u)^4} & \mcell{((aH^2u)^3\cdot H,s,u)^3} &  &  &  &  &  &  &  &  &  &  &  &  &  &  &  \\ \hline
\end{tabular}
\end{adjustbox}
\end{center}
\vspace{-0.2em}
\begin{center}\footnotesize（最上方橫行中的數字表示關於 $x$ 的次數，最左側豎行中的數字表示關於 $u$ 的次數。）\end{center}
\endgroup
\end{landscape}
\endgroup
%
\clearpage
%

%<unit id="noether.p03.title.001" type="title" scan="scan:p03:114">
\begin{center}
\pdfbookmark[0]{3. 關於 n 個變量的型的不變式論}{paper-03}
{\Large\bfseries 3. 關於 $n$ 個變量的型的不變式論\footnote{1909 年在薩爾茨堡自然科學家集會上所作的報告。}}\par
\vspace{1.2em}
Jahresbericht der DMV 19 (1910), pp. 101--104
\end{center}
%</unit>

\vspace{1em}
%<unit id="noether.p03.par.001" type="paragraph" scan="scan:p03:114">
在 $n$ 個變量的型的射影不變式論中，主問題已經解決：型系統的有限性已經得到證明。然而，還有一系列進一步的問題並未被處理，即那些涉及研究型之間聯繫的方法的問題。我願意簡短地報告我在這些問題上得到的結果；但為清楚起見，我先在已知的三元領域中概述這些問題。
%</unit>

%<unit id="noether.p03.par.002" type="paragraph" scan="scan:p03:114">
我首先想到的是符號方法的兩個基本定理：其一，所有不變構成式都可以用符號表示；其二，所有存在於不變式之間的關係，都由少數符號恆等式的逐次應用得到。因此，符號方法給出全部關係，而不必離開不變式的領域。\footnote{Study: \emph{Methoden zur Theorie der ternären Formen}. II. \S\ 6. (Leipzig, Teubner, 1889.)} 建立在這些定理之上的，是生成型的各種過程、符號性的摺疊過程以及非符號性的微分過程、級數展開理論，等等。在這裡我還應提到我在博士論文中證明的一個定理，\footnote{\emph{Über die Bildung des Formensystems der ternären biquadratischen Form}. \S\ 1. (Crelle's Journal Bd. 134. 1908.)} 即所有摺疊都可以由兩種可能的“同變”摺疊逐次應用而生成；對於符號乘積 $s_x t_x u_\sigma u_\tau$，這裡所謂的兩種摺疊，是指 $(stu)$ 與 $(\sigma\tau x)$。在這個定理之上，可以像在二元領域中那樣建立化簡子的理論以及型系統化簡的理論；我在博士論文中已經說明了這一點。
%</unit>

%<unit id="noether.p03.par.003" type="paragraph" scan="scan:p03:114-115">
現在轉到 $n$ 元領域時，甚至符號方法的第一個定理也只在有條件的意義下成立。眾所周知，已經在四元型中，除點座標與平面座標 $x$ 和 $u$ 之外，還出現線座標 $p_{ik}$，即由兩行點座標形成的子行列式。類似地，在 $n$ 元領域中有變量行 $p_\varrho$，其各個元素是由 $\varrho$ 行 $x$ 形成的子行列式；也有符號行 $S^\varrho$，其元素像由 $\varrho$ 行與 $u$ 同變的行 $s$ 形成的子行列式那樣運作。由行列式給出的符號表示，只在把符號行 $S^\varrho$ 分解為行 $s$ 並把這些行分配到不同的行列式中時才已知成立；只有那些依賴於 $S^\varrho$ 的行列式組合才有真正意義。但哪些行列式組合具有這種性質並不能一覽而知；由此，三元領域中提到的所有其他定理也就失效了。
%</unit>

%<unit id="noether.p03.par.004" type="paragraph" scan="scan:p03:115">
我現在要陳述一個簡單原則；藉助它可以得到顯含行 $S^\varrho$ 的符號表示，並由此也重新獲得其餘定理。這是熟知的對應矩陣定理的一個應用：“每一個變量行 $p_\varrho$ 都與另一個行 $q_{n-\varrho}$ 一一對應；後者的各個元素是由 $n-\varrho$ 行 $u$ 形成的子行列式，而這些行 $u$ 通過方程 $(u\mid x)=0$ 與 $\varrho$ 行 $x$ 相聯繫。”相應地，每一個符號行 $S^\varrho$ 都配以另一個 $S_{n-\varrho}$，它的行為彷彿由 $n-\varrho$ 行與 $x$ 同變的行組成。
%</unit>

%<unit id="noether.p03.par.005" type="paragraph" scan="scan:p03:115">
這個定理還容許一個適合應用的第二表述：“對於每一個矩陣積 $(v^{(1)}\ldots v^{(\varrho)}\mid y^{(1)}\ldots y^{(\varrho)})$，\footnote{也就是說，對於由 $\varrho$ 行 $v$ 組成的矩陣和由 $\varrho$ 行 $y$ 組成的矩陣，取相應行列式乘積之和。} 若第二個矩陣的行與第一個矩陣的行反變，則有一個行列式與之對應；反過來，每一個行列式也可以變換為具有預先指定行數 $\varrho$ 的矩陣積。”於是，行列式與矩陣積表現為完全等價，符號方法的第一個定理獲得如下形式：“所有不變構成式都可以由矩陣積作符號表示。”由兩個符號行 $S^\sigma$, $T^\tau$，藉助變量行，現在很容易形成顯含這些符號行的矩陣積，從而表示形如 $(S^\sigma\mid p_\sigma)(T^\tau\mid p_\tau)$ 的不變構成式，即：
\begin{equation}
 \pair{q_{n-\sigma-\lambda}S^\sigma}{T_{n-\tau}p_{\tau-\lambda}},
 \qquad (\lambda=1,2\ldots\tau\ \text{或}\ n-\sigma).
\tag{1}
\end{equation}
%</unit>

%<unit id="noether.p03.par.006" type="paragraph" scan="scan:p03:115-116">
更重要的是逆命題：所有不變構成式都能顯式地由矩陣積表示，因此上述過程就是二元和三元摺疊過程的推廣。為了證明這個逆命題，必須把符號方法的第二個定理轉移到 $n$ 元領域；也就是說，必須建立所有獨立不可分解的恆等式的全體，在這些恆等式中所有行 $S^\varrho$, $p_\varrho$ 都被看作不可分解的。這些恆等式來源於只含行 $x$ 與 $u$ 的已知恆等式，\footnote{E. Pascal in \emph{Memorie d. R. Acc. d. Lincei}. (4) 5 (1888), p. 375.} 其方式是用矩陣的若干乘積定理取代行列式的乘積定理，並且恰恰藉助這些乘積定理把不同的矩陣積縮併為一個矩陣積。這樣得到兩個對偶公式；這裡只給出其中一個：
\begin{equation}
\begin{aligned}
 \pair{S^{\varrho_1}S^{\varrho_2}\ldots S^{\varrho_k}}{x p_{\varrho-1}}
 &=\sum_{i=1}^{k}\eps\,\pair{S_1^{\rho_1}\cdots S_{i-1}^{\rho_{i-1}}S_{i+1}^{\rho_{i+1}}\cdots S_k^{\rho_k}q_{n-\rho_i+1}}{S_{n-\rho_i}x},\\
 \varrho&=\sum_{i=1}^{k}\varrho_i\leqq n;\qquad \eps=\pm1.
\end{aligned}
\tag{2}
\end{equation}
%</unit>

%<unit id="noether.p03.par.007" type="paragraph" scan="scan:p03:116">
這些公式中唯一包含的特殊化，即出現一個行 $x$，是不可避免的；因為對於完全一般的符號行與變量行，我們會得到一個“分解恆等式”，即一個把行 $p_\varrho$ 分解為各個行 $x$ 的恆等式。分解恆等式最簡單的特殊情形，Clebsch 已在他的“不變式論的基本問題”中用於從級數展開理論導出初等協變式。藉助這些恆等式——它們調解從非顯式行列式表達式到矩陣積的過渡——事實上可以證明，所要求的顯式符號表示正是由矩陣表示得到的。
%</unit>

%<unit id="noether.p03.par.008" type="paragraph" scan="scan:p03:116-117">
對於摺疊過程，卻有一個與三元領域完全類似的定理；化簡定理也同樣可以以類似方式建立起來。如果在 (1) 中把數 $\lambda$ 稱為摺疊的虧格，並稱 $\sigma=\tau$ 的摺疊為同變摺疊，那麼所有摺疊都可以由虧格為 1 的同變摺疊逐次應用而生成。這個定理的證明本質上基於這樣一個事實：最一般的型可以由“正規型”代替；所謂正規型，就是在應用所有不變微分運算時恆等為零的型。在四元領域中，這後一事實已由 Mertens 證明；\footnote{Über invariante Gebilde quaternärer Formen. Wiener Berichte. Bd. 98. 1889.} 我們對於最一般微分運算的知識——眾所周知，它們是把摺疊過程中的符號行替換為微分符號而產生的——使我們能夠藉助分解恆等式，幾乎逐字地把 Mertens 的證明轉移到 $n$ 元領域。
%</unit>

\clearpage

%<unit id="noether.p04.title.001" type="title" scan="scan:p04:118">
\begin{center}
\pdfbookmark[0]{4. 關於 n 個變量的型的不變式論}{paper-04}
{\Large\bfseries 4. 關於 $n$ 個變量的型的不變式論}\par
\vspace{1.2em}
Journal für die reine und angewandte Mathematik 139 (1911), pp. 118--154
\end{center}
%</unit>

%<unit id="noether.p04.sec.intro" type="section" scan="scan:p04:118-121">
\section*{引言}
在 $n$ 個變量的型的射影不變式論中，一旦越出二元和三元領域，與主問題——型系統有限性的證明——相鄰接的問題幾乎還沒有被處理。這些問題涉及型之間的相互依賴關係，以及控制這種型間聯繫的方法。正如在二元和三元領域中已經完全證明的那樣，這些方法本質上依賴於 Study 稱為“符號方法的基本定理”的兩個定理，即型的符號可表示性定理與符號恆等式定理。\footnote{E. Study, \emph{Methoden zur Theorie der ternären Formen} (Leipzig, Teubner, 1889). II. \S\ 6. 關於這些基本定理也出現在特殊變換群的不變式論中，參見 Study, \emph{Das Apollonische Problem} (Math. Ann. Bd. 49 [1897]) 以及 \emph{Zur Differentialgeometrie der analytischen Kurven} (Transactions of the American Math. Society, Bd. 1 [1909]).} 後一個定理說，在整不變構成式之間存在的一切關係，都是由少數符號恆等式的逐次應用得到的；因此，符號方法並且只有符號方法，能在不離開不變式領域的情況下認識型之間的聯繫。

Gordan 在他的埃爾朗根綱領中\footnote{P. Gordan, \emph{Über das Formensystem binärer Formen}. S. 46. (Leipzig, Teubner, 1875).} 已經指出，向 $n$ 個變量過渡的困難在於：符號方法的第一個定理不再以其一般形式成立。在 $n$ 元領域中，眾所周知，會出現較高階變量行 $p_\rho$\footnote{記號見 \S\ 1。}，並且相應地出現較高階符號行 $S^\rho$；無論是像 Clebsch 最初那樣\footnote{A. Clebsch, über eine Fundamentalaufgabe der Invariantentheorie. (Abh. der Ges. d. Wiss. zu Göttingen, Bd. XVII, 1872).} 從含有任意多個行 $p_\rho$ 的型出發，還是按照 F. Deruyts 與 A. Capelli\footnote{參見 A. Capelli, \emph{Lezioni sulla teoria delle forme algebriche} (Napoli 1902), \S\ 22--24 以及那裡所引文獻。} 從只含任意多個同變行 $x$ 的型出發，都會到達這些行。顯含行 $p_\rho$, $S^\rho$ 的符號表示\footnote{“顯式表示”的較精確定義見 \S\ 2。} 並不存在；必須把 $p_\rho$, $S^\rho$ 分解為各個行 $x$ 以及與之反變的行 $s$，再把這些行分配到不同的行列式中。確實，可以藉助微分條件（參見公式 (19.)）來檢驗一個給定的行列式組合是否只依賴於行 $p_\rho$, $S^\rho$；但這樣做並不能概覽不變構成式的全體以及生成它們的過程。\footnote{Study 所作的一個在計算上很有價值的符號學發展（由 F. Engel 通報，Ber. d. K. Sächs. Ges. d. Wiss., math.-phys. Kl., 1900, p. 126；四元領域見 E. Study, \emph{Geometrie der Dynamen}, p. 135）在我們的意義下也不是顯式表示。例如，它很難導向生成這些型的非符號微分過程。}

下面我們將在 \S\S\ 2 和 6 中說明，藉助“對應矩陣定理”\footnote{除 Grassmann 1862 年的 \emph{Ausdehnungslehre} 第 112 號外，最早見 A. Brill, über zwei Berührungsprobleme, \S\ 2 (Math. Ann. Bd. 4. 1871).}，通過用“矩陣積”表示來取代行列式表示，第一個基本定理如何恢復其一般性。在 \S\ 2 中證明，每一個顯含行 $S^\rho,p_\rho$ 的矩陣積都是基型的一個不變構成式；而 \S\ 6 中給出的逆命題——每一個不變構成式都可以顯式地由矩陣積表示——只有通過第二個基本定理的轉移（\S\S\ 3--5）才得以實現。

對於只含變量行 $x$ 的型，這個定理是已知的。\footnote{E. Pascal, Giorn. di mat. 26 (1888), pp. 33, 102; Rendic. d. Accad. d. Linc. (4) 4 (1888), p. 119; ib. Mem. (4) 5 (1888), p. 375.} 一般問題，即建立把所有行 $S^\rho,p_\rho$ 視為不可分解時的不可分解恆等式全體，是由 Study 提出的；\footnote{``Methoden \dots'', p. 204, note 17).} 同時，他把出現的恆等式數目看作 $n$ 元領域中方法困難程度的一種度量。既然現在可以把所有恆等式合併在一個公式 (36.) 中，這一困難至少被壓縮到最小。在應用中，經常出現的卻是 (36.) 的若干突出的特殊情形，例如 (20.)、(21.)，其特徵在於它們不必代入新行便容許顯式表示；或者公式 (31.)，稱為“分解恆等式”，因為在其中一行 $p_\rho$ 似乎被分解為若干行 $y$。

在這兩個基本定理之上，進一步的理論便可以容易地建立起來。第一個定理直接給出生成型的過程、符號摺疊過程和非符號微分過程，\footnote{在論文 “Invariante Gebilde von Nullsystemen” (Sitzungsber. d. Ak. d. Wiss. Wien, Bd. 97, Abt. IIa, 1888) 中，Mertens 通過另一條不能直接推廣的路徑得到了四元摺疊過程。那裡出現的 6 個行 $S^2$ 的交錯不變式，與我們這裡得到的不變式相差一個由代換 (11.) 的一次應用所產生的偶不變式組合。四元摺疊過程的一個特殊情形也見於 P. Gordan, \emph{Das simultane System von zwei quadratischen quaternären Formen} (Math. Ann. Bd. 56. 1903).} 而分解恆等式則在這些過程中消去所有等價的過程（\S\ 6）。對微分過程的認識進一步給出“正規型”概念向 $n$ 元領域的轉移；並且藉助 Mertens\footnote{F. Mertens, Über invariante Gebilde quaternärer Formen. (Sitzber. d. Ak. d. Wiss. Wien. Bd. 98. Abt. IIa. 1889.)} 在四元領域給出的證明方法，得到一個定理：一個不變構成式系統可以由一個等價的正規型系統代替（\S\ 7）。在後一假設下，我在博士論文中\footnote{Über die Bildung des Formensystems der ternären biquadratischen Form. (Dieses Journal, Bd. 134. 1908.)} 證明了三元摺疊過程可以由兩個基本摺疊逐次應用而生成。在這個定理以及級數展開理論的基礎上，我又引入“型行”概念，發展了型系統的主要化簡定理。完全類似地，在 $n$ 元領域中，摺疊過程可以由 $n-1$ 個基本摺疊生成（\S\ 8），並且可以建立化簡定理（\S\ 9）。
%</unit>

%<unit id="noether.p04.afterword.001" type="afterword" scan="scan:p04:121">
\paragraph{附記。} 在薩爾茨堡通報之際，維也納的 E. Müller 教授提醒我注意，本論文所依據的矩陣積計算，在原則上與 Grassmann 擴張論的演算相同。\footnote{Hermann Graßmanns gesammelte mathematische und physikalische Werke. Ersten Bandes zweiter Teil: Die Ausdehnungslehre von 1862. In Gemeinschaft mit H. Graßmann d. Jüngeren herausgegeben v. Fr. Engel (Leipzig, Teubner 1896).} 的確，Grassmann 的“組合積” $[S^{\rho_1}S^{\rho_2}\cdots S^{\rho_i}]$ 可以看作由這些行給出的矩陣（參見 \S\ 2）；更具體地說，“遞進積”就是由這些行本身組成的矩陣，“逆進積”是由所配行 $S_{n-\rho_1},S_{n-\rho_2},\ldots S_{n-\rho_i}$ 組成的矩陣；而對應於“混合積”的矩陣，則在若干行 $S$ 通過代換 (9.) 合併為新的行 $P,Q$ 時產生。我們的“所配行”對應於 Grassmann 的“補”，我們的“矩陣積”對應於“階數為 0 的混合積”。

在這種對應下，1862 年 \emph{Ausdehnungslehre} 第 113 號給出的展開，便與我們公式 (32.) 的對偶公式相同。在特殊情形 $\rho=1$ 下，(32.) 化為 (17.)，對偶公式化為 (16.)。Müller\footnote{E. Müller, Beiträge zur Graßmannschen Ausdehnungslehre. 1. Mitteilung: Einige allgemeine Sätze. (Sitzungsber. d. Ak. d. Wiss. Wien; Bd. 118. Abt. IIa. Juli 1909), Satz 16 und 18.} 最近正是從 Grassmann 展開的這些特殊情形導出了上面提到的恆等式 (20.)、(21.)。

與 Grassmann--Müller 的推導相比，我們的推導同時還給出了恆等式完備性的證明；對於這裡所考察的問題，這一點似乎是本質所在。而且，我們的推導從 (20.)、(21.) 進一步到達最一般的不可分解恆等式 (36.)；這些恆等式此前似乎尚未被注意到。
%</unit>

%<unit id="noether.p04.sec.001" type="section" scan="scan:p04:121-124">
\section*{\S\ 1. 記號和定義。已知結果概述。}
照通常做法，我們用變量行 $x$ 表示元素行 $x_1,x_2,\ldots,x_n$；類似地，用變量行 $p_\rho$（$\rho=1,2,\ldots,n-1$）表示 $\binom n\rho$ 個元素 $p_{i_1i_2\ldots i_\rho}$ 組成的行，其中 $p_{i_1i_2\ldots i_\rho}$ 表示由 $\rho$ 行 $x^{(1)},x^{(2)},\ldots,x^{(\rho)}$ 組成的矩陣所形成的 $\binom n\rho$ 個行列式，
\[
\left|\begin{array}{cccc}
 x_{i_1}^{(1)}&x_{i_2}^{(1)}&\cdots&x_{i_\rho}^{(1)}\\
 x_{i_1}^{(2)}&x_{i_2}^{(2)}&\cdots&x_{i_\rho}^{(2)}\\
 \vdots&\vdots&&\vdots\\
 x_{i_1}^{(\rho)}&x_{i_2}^{(\rho)}&\cdots&x_{i_\rho}^{(\rho)}
\end{array}\right|,
\qquad (i_1<i_2<\cdots<i_\rho=1,2,\ldots,n).
\]
按照對應矩陣定理，\footnote{例如參見 Gordan--Kerschensteiner, \emph{Vorlesungen über Invariantentheorie}, Bd. I, p. 99.} 每一個行 $p_\rho$ 都通過下列對該行每個元素有效的關係，一一對應地配以第二個行 $q_{n-\rho}$：
\begin{equation}
 p_{i_1i_2\ldots i_\rho}=(-1)^{\frac{\rho(\rho+1)}{2}+i_1+i_2+\cdots+i_\rho}q_{j_1j_2\ldots j_{n-\rho}},
\tag{1}
\end{equation}
其中 $i$ 和 $j$ 各自按順序排列，並且合在一起構成數字 $1$ 到 $n$ 的一個完全置換。行 $q_{n-\rho}$ 由 $\binom n\rho$ 個 $n-\rho$ 階行列式 $q_{j_1j_2\ldots j_{n-\rho}}$ 組成，這些行列式由 $n-\rho$ 行 $u^{(1)},u^{(2)},\ldots,u^{(n-\rho)}$ 組成的矩陣形成，而這些行與 $x$ 反變。它們滿足 $\rho\cdot(n-\rho)$ 個條件方程
\begin{equation}
 (u^{(k)}\mid x^{(l)})=\sum_{\alpha=1}^{n}u_\alpha^{(k)}x_\alpha^{(l)}=0;
 \qquad (k=1,2,\ldots,n-\rho;\ l=1,2,\ldots,\rho).
\tag{2}
\end{equation}
我們將假定它們按 Clebsch\footnote{同上，\S\ 5。} 的方式歸一化，使得在 (1.) 中本來會出現的比例因子等於一。

按照 Clebsch\footnote{同上，\S\ 12。}（也同樣按照 F. Deruyts 和 A. Capelli 後來的研究，參見引言），最一般的型可以化歸為這樣的型：每一個 $n-1$ 個行 $p_\rho$ 至多隻出現一次；因而它們可以用符號表示為
\begin{equation}
 (S^1\mid p_1)^{m_1}(S^2\mid p_2)^{m_2}\cdots(S^{n-1}\mid p_{n-1})^{m_{n-1}},
\tag{3}
\end{equation}
其中，類似於 (2.)，我們置
\[
 (S^\rho\mid p_\rho)=\sum S^{i_1i_2\ldots i_\rho}p_{i_1i_2\ldots i_\rho}.
\]
由於行 $p_\rho$ 的元素之間存在非線性的恆等關係，符號行 $S^\rho$ 可以被歸一化，使其滿足完全相同的恆等式；因此它們的行為就像由一個 $\rho$ 行矩陣的行列式組成，而該矩陣的行 $s^{(1)},s^{(2)},\ldots s^{(\rho)}$ 與 $x$ 反變。\footnote{Clebsch，同上，\S\ 4。} 因而，每一個行 $S^\rho$ 都可以通過關係
\begin{equation}
 S_{i_1i_2\ldots i_{n-\rho}}=(-1)^{\frac{(n-\rho)(n-\rho+1)}2+i_1+i_2+\cdots+i_{n-\rho}}S^{j_1j_2\ldots j_\rho},
\tag{4}
\end{equation}
一一對應地配以另一個 $S_{n-\rho}$，其中 $i$ 與 $j$ 仍各自按順序排列，並且恰好窮盡數字 $1$ 至 $n$。行 $S_{n-\rho}$ 的元素 $S_{i_1i_2\ldots i_{n-\rho}}$ 的行為，像是由 $(n-\rho)$ 行 $\sigma^{(1)},\sigma^{(2)},\ldots\sigma^{(n-\rho)}$ 形成的行列式；這些行與 $x$ 同變，而行 $s$ 與 $\sigma$ 之間由 $\rho\cdot(n-\rho)$ 個條件相連：
\begin{equation}
 (s^{(k)}\mid \sigma^{(l)})=0.\qquad(k=1,2,\ldots\rho;\ l=1,2,\ldots n-\rho).
\tag{5}
\end{equation}
由關係 (1.) 與 (4.)，(3.) 的每個因子都有以下等價的四重符號表示：
\begin{equation}
 (S^\rho\mid p_\rho)=(S^\rho q_{n-\rho})=(S_{n-\rho}p_\rho)=(-1)^{\rho(n-\rho)}(q_{n-\rho}\mid S_{n-\rho}),
\tag{6}
\end{equation}
其中，在以下一切地方，$(S^\rho q_{n-\rho})$ 與 $(S_{n-\rho}p_\rho)$ 總是行列式 $(s^{(1)}\cdots s^{(\rho)}u^{(1)}\cdots u^{(n-\rho)})$ 與 $(\sigma^{(1)}\cdots\sigma^{(n-\rho)}x^{(1)}\cdots x^{(\rho)})$ 的縮寫；由 Laplace 展開可知，它們只依賴於行 $S^\rho,q_{n-\rho}$，或分別依賴於 $S_{n-\rho},p_\rho$，正如上述記號所意圖表示的那樣。根據前面的討論，表達式 $(S^\rho\mid p_\rho)$ 與 $(q_{n-\rho}\mid S_{n-\rho})$ 是行 $s$ 與 $x$ 的矩陣的乘積，或行 $u$ 與 $\sigma$ 的矩陣的乘積；這裡所謂矩陣積，是指相應行列式乘積之和，也就是乘積矩陣的行列式值。

一個符號行（或變量行）所屬的特徵數 $\rho\cdot(n-\rho)$，按照 Clebsch 的說法，稱為該行的權；\footnote{同上，\S\ 11。} 數 $\rho$ 給出行 $s$ 的數目，稱為上權指標；數 $n-\rho$ 給出行 $\sigma$ 的數目，稱為下權指標。由關係 (4.) 可知，一個符號行可以在同一個關係中一次帶上權指標出現，另一次帶下權指標出現；相應行 $p_\rho$ 與 $q_{n-\rho}$ 的出現也是如此。最後還要說明：一般地，我們把其配屬矩陣由行 $x$ 組成的較高階變量行記為 $p$，把其矩陣由行 $u$ 組成的記為 $q$；與 $x$ 同變的符號行用小希臘字母 $\sigma,\tau,\alpha,\beta$ 表示；與 $u$ 同變的符號行用小拉丁字母 $s,t,a,b$ 表示；其餘所有符號行則用帶權指標的大拉丁字母表示，如 $S^\rho$, $T^\tau$, $A^\rho$ 或 $S_{n-\sigma}$, $T_{n-\tau}$, $A_{n-\rho}$。與上權指標或下權指標相對應，我們也把一個行的各個元素寫成帶上標或下標的形式，如 (4.) 中所示。
%</unit>
%
%
%<unit id="noether.p04.sec.002" type="section" lang="zh-Hans" scan="scan:p04:124-128">
\section*{\S\ 2. 用矩陣積表示。}
以下所謂“矩陣積”，始終是指兩個行數相同的矩陣的相應行列式乘積之和；第二個矩陣的行與第一個矩陣的行互為反變。每個矩陣的行可以是單獨給出的，也可以如 (6.) 中那樣合併成一個行 $S^\rho$，或合併成若干不同行 $S^\rho,S^\sigma,\ldots$。矩陣積仍用 (6.) 中的符號 $(\ \mid\ )$ 表示。表示式 (6.) 是對應矩陣定理的另一種表述的特殊應用：每一個 $\rho$ 行矩陣積 $(v^{(1)}v^{(2)}\cdots v^{(\rho)}\mid y^{(1)}\cdots y^{(\rho)})$ 都配以一個行列式；反過來，每個行列式也可配以一個具有指定行數 $\rho$ 的矩陣積，$\rho=1,2,\ldots,n-1$。\footnote{當 $\rho>n$ 時矩陣積已知為零；當 $\rho=n$ 時它化為兩個行列式的乘積。} 實際上，只須把各行 $y$ 合併為一行 $p_\rho$（或把各行 $v$ 合併為一行 $q_\rho$），再執行 (1.) 給出的代換，便可在給定的矩陣積和行列式之間來回過渡。行列式的值並非由單個元素立即給出，而是經 Laplace 展開才得到。因此行列式表示與矩陣積表示等價，於是關於型的符號可表示性的定理（第一基本定理）也可寫成如下形式：

\emph{定理 I.} “所有不變構成式都可以符號地表示為矩陣積；反過來，所有矩陣積都是不變構成式。”\footnote{當然，對給定基型的不變構成式，只是那些依賴於符號行 $S^\rho\ldots$ 的矩陣積聚合。}

矩陣積表示克服了行列式表示所造成的非顯式表示的根本困難。\footnote{參見引言。} “顯式表示”是指最終表達式只含行 $S^\sigma,T^\tau,\ldots$ 本身，或只含由這些行唯一導出的行 $R^\rho$，而不再含有各個組成行 $s,t,\ldots$。在 \S\ 6 中，對於只依賴於符號行 $S^\sigma,T^\tau,\ldots$ 的所有不變構成式，我們將得到關於這些符號行的顯式表示，並從而把生成過程、摺疊過程以及微分過程轉移到 $n$ 元領域。顯式表示之所以可能，是因為只有最簡單的矩陣積才逐一對應於行列式表示中的行列式；其餘矩陣積則把若干行列式組合成一個單一表達式。若把符號行和變量行分解為單個行 $s,u,x$，則按通常形式的第一基本定理，必然出現行列式聚合。反過來，表示給定基型的不變構成式的一切行列式聚合都可組合為顯式矩陣積；這一點要用第二基本定理來證明（\S\ 6）。

為了由兩個符號行 $S^\sigma,T^\tau$，藉助變量行，得到基型 $(S^\sigma\mid p_\sigma)(T^\tau\mid p_\tau)$ 的一個顯含所有行的不變構成式（摺疊過程），只須按照本段開頭的定義構造第三種矩陣積：
\begin{equation}
 \pair{q_{n-\sigma-\lambda}S^\sigma}{T_{n-\tau}p_{\tau-\lambda}},
 \qquad (\lambda=1,2,\ldots,\tau,n-\sigma).
\tag{7}
\end{equation}
展開 (7.) 得到
\begin{equation}
\begin{aligned}
&\pair{v^{(1)}\cdots v^{(n-\sigma-\lambda)}s^{(1)}\cdots s^{(\sigma)}}
       {\alpha^{(1)}\cdots\alpha^{(n-\tau)}y^{(1)}\cdots y^{(\tau-\lambda)}} \\
&=\sum_{i,k,l}\eps\,q_{k_1k_2\ldots k_{n-\sigma-\lambda}}
 S^{k_{n-\sigma-\lambda+1}\ldots k_{n-\lambda}}
 T_{l_1\ldots l_{n-\tau}}
 p_{l_{n-\tau+1}\ldots l_{n-\lambda}},
\end{aligned}
\tag{8}
\end{equation}
也就是說，它只依賴於行 $S,T,q,p$。這裡 $\eps=\pm1$，\footnote{此記號以下始終保留。} 求和的意思是：每一行 $S\ldots$ 的指標按正序排列；對於每個組合 $i_1i_2\ldots i_{n-\lambda}$，$k$ 與 $l$ 分別遍歷這些數 $i$ 的全部可能置換。進入 (7.) 的數 $\lambda$，在 $\sigma>\tau$ 的情形下稱為摺疊的缺陷；這一限制的理由將在 \S\ 6 中出現。\footnote{數 $\lambda$ 總是取到最後所示兩個數中較小者。}

出現在 (7.) 中的兩個矩陣的行列式，可以看作新行 $Q^{n-\lambda},P_{n-\lambda}$ 的元素；其配屬行分別為 $Q_\lambda,P^\lambda$。由 (8.)，這些行 $Q,P$ 可分別由原來的 $S$ 與 $q$、$T$ 與 $p$ 表出。我們用下列等式表示：
\begin{equation}
\begin{aligned}
 q_{n-\sigma-\lambda}S^\sigma&=Q^{n-\lambda}\sim Q_\lambda, &\quad\text{亦即}\quad Q_\lambda&=[q_{n-\sigma-\lambda}S^\sigma]_\lambda,\\
 T_{n-\tau}p_{\tau-\lambda}&=P_{n-\lambda}\sim P^\lambda, &\quad\text{亦即}\quad P^\lambda&=[T_{n-\tau}p_{\tau-\lambda}]^\lambda.
\end{aligned}
\tag{9}
\end{equation}
考慮到 (4.)，與 (6.) 類似，矩陣積 (7.) 有等價的四重符號表示：
\begin{equation}
\pair{q_{n-\sigma-\lambda}S^\sigma}{T_{n-\tau}p_{\tau-\lambda}}
=(q_{n-\sigma-\lambda}S^\sigma P^\lambda)=(Q_\lambda T_{n-\tau}p_{\tau-\lambda})=(-1)^{\lambda(n-\lambda)}\pair{P^\lambda}{Q_\lambda}.
\tag{10}
\end{equation}
還可對 (7.) 進行進一步的新行代換，令
\begin{equation}
 \pair{q_{n-\sigma-\lambda}S^\sigma}{T_{n-\tau}p_{\tau-\lambda}}
 =\pair{R^{\sigma+\lambda}}{p_{\sigma+\lambda}}(R^{\tau-\lambda}p_{\tau-\lambda}).
\tag{11}
\end{equation}
這裡同樣由 (8.) 知道，行 $R$ 的元素的乘積可唯一地通過行 $S$ 與 $T$ 表出。代換 (9.) 和 (11.) 將特別在對 (7.) 再同其他行摺疊時使用。

矩陣表示使我們能夠以不變形式寫出行 $p_\rho$ 的元素之間的二次關係（眾所周知，其餘關係都由此推出），從而也能按 \S\ 1 寫出符號行 $S^\rho$ 滿足的、關於基型係數線性的關係。這些恆等式是
\begin{equation}
 \pair{q_{n-\rho-\lambda}S^\rho}{S_{n-\rho}p_{\rho-\lambda}}=0,
 \qquad (\lambda=1,2,\ldots,\rho,n-\rho),
\tag{12}
\end{equation}
其中 $p$ 與 $q$ 看作任意量。我們將在 \S\ 5 中看到，缺陷數 $\lambda$ 為奇數的恆等式是較高缺陷恆等式的後果。關係 (12.) 直接由關係 (5.) 推出，因為有
\[
 \pair{q_{n-\rho-\lambda}S^\rho}{S_{n-\rho}p_{\rho-\lambda}}
 =(v^{(1)}\cdots v^{(n-\rho-\lambda)}s^{(1)}\cdots s^{(\rho)}\mid \sigma^{(1)}\cdots\sigma^{(n-\rho)}y^{(1)}\cdots y^{(\rho-\lambda)}),
\]
應用乘積定理後，由 (5.) 得到一個消失的 $(n-\lambda)$ 行行列式。恆等式 (12.) 說明：為求出一切摺疊型，只要摺疊在每個變量行中線性的型的乘積即可；或者，按 \S\ 6，歸一化型 $(S^\rho\mid p_\rho)^m$ 不具有係數線性的不變構成式。\footnote{參見 Study 在 Grassmann 著作第一卷第二部 p. 510 所作附錄（量 $S^\rho$ 為簡單量的條件）。}

條件 (12.) 也足以刻畫行 $p_\rho$。因為 $p_\rho$ 的各元素是某矩陣行列式這一性質是不變性質，所以必須能不變地表達；而由定理 VI，條件 (12.) 給出對 $p_\rho$ 可能有的最大不變限制，即以行 $p_\rho$ 為係數構成的線性型完全沒有不變構成式。
%</unit>

%<unit id="noether.p04.sec.003" type="section" lang="zh-Hans" scan="scan:p04:128-134">
\section*{\S\ 3. 符號恆等式。}
現在必須轉移符號方法的第二基本定理，也就是列出所有不可約、不可分解的恆等式；在這些恆等式中，所有行 $S^\rho,p_\rho$ 都被看作不可分解的量。還要約定：按照符號行和變量行的含義，凡由把變量行 $p_\rho$ 特化為 $p_{\rho-1}y$ 並把某個系統恆等式應用到所得表達式而得到的恆等式，均視為複合的；另一方面，含有 $k$ 個符號行的恆等式，即使由上述過程從含較少符號行的恆等式導出，也仍算作不可分解。行 $S^\rho,p_\rho$ 不必已經在恆等式中顯式出現；只要證明出現的矩陣積僅依賴於這些行，就足以把它們解釋為不可分解量。不過我們將證明，在這裡也總能通過（部分地）使用代換 (11.) 得到顯式表示。一般恆等式由矩陣表示原則，從 Pascal 給出的特殊恆等式得到；這些恆等式只含行 $x$ 與 $u$，並且直接轉移了 Study 在三元領域所給的恆等式：\footnote{問題的表述及文獻參見引言。}
\begin{align}
\sum_{i=1}^n(-1)^{i+1}\pair{a^{(i)}}{x}(a^{(1)}\cdots a^{(i-1)}a^{(i+1)}\cdots a^{(n)}u)&=(-1)^{n+1}\pair{u}{x}(a^{(1)}\cdots a^{(n)}),\tag{13}\\
\sum\pm\pair{a^{(1)}}{x^{(1)}}\pair{a^{(2)}}{x^{(2)}}\cdots\pair{a^{(n)}}{x^{(n)}}&=(a^{(1)}\cdots a^{(n)})(x^{(1)}\cdots x^{(n)}),\tag{14}\\
\sum_{i=1}^n(-1)^{i+1}\pair{u}{\alpha^{(i)}}(\alpha^{(1)}\cdots\alpha^{(i-1)}\alpha^{(i+1)}\cdots \alpha^{n}x)&=(-1)^{n+1}\pair{u}{x}(\alpha^{(1)}\cdots\alpha^{(n)}).\tag{15}
\end{align}
由 (13.) 與 (15.) 還可通過代換 $(a^{(i)}\mid x)=(a^{(i)}q_{n-1})$ 以及 $(u\mid \alpha^{(i)})=(p_{n-1}\alpha^{(i)})$ 得到兩個恆等式；按我們的解釋，它們不應看成實質上不同。(14.) 表示行列式的乘積定理；現在說明如何用適當構成的矩陣乘積定理系統代替 (13.) 與 (14.)，以及對偶地代替 (15.) 與 (14.)。一次對 $\rho$ 行矩陣、一次對 $(\rho-1)$ 行矩陣寫出乘積定理，有
\[
\pair{a^{(1)}a^{(2)}\cdots a^{(\rho)}}{xp_{\rho-1}}
=(a^{(1)}a^{(2)}\cdots a^{(\rho)}\mid xy^{(2)}\cdots y^{(\rho)})
=\sum\pm\pair{a^{(1)}}{x}\pair{a^{(2)}}{y^{(2)}}\cdots\pair{a^{(\rho)}}{y^{(\rho)}},
\]
\[
\sum\pm\pair{a^{(2)}}{y^{(2)}}\cdots\pair{a^{(\rho)}}{y^{(\rho)}}
=(a^{(2)}\cdots a^{(\rho)}\mid y^{(2)}\cdots y^{(\rho)})
=(a^{(2)}\cdots a^{(\rho)}\mid p_{\rho-1})=(a^{(2)}\cdots a^{(\rho)}q_{n-\rho-1});
\]
於是得到下列乘積定理系統，其中 $\rho=2,3,\ldots n$：
\begin{equation}
\begin{aligned}
\pair{a^{(1)}a^{(2)}\cdots a^{(\rho)}}{xp_{\rho-1}}
 &=\sum_{i=1}^{\rho}(-1)^{i+1}\pair{a^{(i)}}{x}(a^{(1)}\cdots a^{(i-1)}a^{(i+1)}\cdots a^{(\rho)}\mid p_{\rho-1})\\
 &=\sum_{i=1}^{\rho}(-1)^{i+1}\pair{a^{(i)}}{x}(a^{(1)}\cdots a^{(i-1)}a^{(i+1)}\cdots a^{(\rho)}q_{n-\rho+1}).
\end{aligned}
\tag{16}
\end{equation}
當 $\rho=n$ 時，(16.) 化為 (13.)；若繼續特化 $p_{n-1}=y^{(2)}p_{n-2}$、$p_{n-2}=y^{(3)}p_{n-3}$ 等，並把恆等式 (16.) 應用於表達式 $(a^{(1)}\cdots a^{(i-1)}a^{(i+1)}\cdots a^{(n)}\mid y^{(2)}p_{n-2})$、$(a^{(1)}\cdots a^{(i-1)}a^{(i+1)}\cdots a^{(k-1)}a^{(k+1)}\cdots a^{(n)}\mid y^{(3)}p_{n-3})$ 等，就從 (13.) 得到 (14.)。因此，按照我們的定義，恆等式 (14.) 由恆等式 (16.) 複合而成；也就是說，系統 (16.) 與恆等式 (13.) 和 (14.) 等價。系統 (16.) 滿足一般恆等式所要求的條件之一：它把行 $p_{\rho-1}$ 作為不可分解的量包含在內，同時又將其顯式表示出來。它是滿足這一要求且僅滿足這一要求的唯一系統。因為我們不能由行 $a,x,p$ 構成其他行列式或矩陣的乘積；而這些式子之間也不可能存在獨立於 (16.) 的關係，否則消去 $(a^{(1)}\cdots a^{(\rho)}\mid xp_{\rho-1})$ 後，就會得到 $\rho$ 個（$\rho\leqq n$）表達式 $(a^{(i)}\mid x)(a^{(1)}\cdots a^{(i-1)}a^{(i+1)}a^{(\rho)}q_{n-\rho+1})$ 之間的一個關係，但根據 (13.)，至少要有 $n+1$ 個這樣的表達式才能構成一個關係。由類似的考察，可得到與 (15.) 和 (14.) 等價的系統：
\begin{equation}
\begin{aligned}
\pair{u q_{\rho-1}}{\alpha^{(1)}\cdots \alpha^{(\rho)}}
 &=\sum_{i=1}^{\rho}(-1)^{i+1}\pair{u}{\alpha^{(i)}}(q_{\rho-1}\mid \alpha^{(1)}\cdots \alpha^{(i-1)}\alpha^{(i+1)}\cdots \alpha^{(\rho)})\\
 &=\sum_{i=1}^{\rho}(-1)^{i+1}\pair{u}{\alpha^{(i)}}(p_{n-\rho+1}\alpha^{(1)}\cdots \alpha^{(i-1)}\alpha^{(i+1)}\cdots \alpha^{(\rho)}).
\end{aligned}
\tag{17}
\end{equation}
為了從 (16.) 過渡到任意行 $A^{\rho_1},B^{\rho_2},\ldots,x,p$ 之間的恆等式，注意：一旦把 $A^{\rho_1}=a^{(1)}a^{(2)}\cdots a^{(\rho_1)}$、$B^{\rho_2}=b^{(1)}b^{(2)}\cdots b^{(\rho_2)}$ 等分解，這些恆等式必須與 (16.) 相同；並且只要各個表達式屬於 (16.) 中的類型，最終表達式也只能藉助 (16.) 所規定的運算組合而成。因此，若設
\begin{equation}
 \rho=\rho_1+\rho_2+\cdots\rho_k\leqq n,
\tag{18}
\end{equation}
則有
\begin{align*}
&\pair{a^{(1)}\cdots a^{(\rho_1)}b^{(1)}\cdots b^{(\rho_2)}\cdots k^{(1)}\cdots k^{(\rho_k)}}{xp_{\rho-1}}
  =\pair{A^{\rho_1}B^{\rho_2}\cdots K^{\rho_k}}{xp_{\rho-1}} \\
&\quad=\sum_{i=1}^{\rho_1}(-1)^{i+1}\pair{a^{(i)}}{x}(a^{(1)}\cdots a^{(i-1)}a^{(i+1)}\cdots a^{(\rho_1)}B^{\rho_2}\cdots K^{\rho_k}q_{n-\rho+1})\\
&\qquad+(-1)^{\rho_1\rho_2}\sum_{i=1}^{\rho_2}(-1)^{i+1}\pair{b^{(i)}}{x}(b^{(1)}\cdots b^{(i-1)}b^{(i+1)}b^{(\rho_2)}A^{\rho_1}\cdots K^{\rho_k}q_{n-\rho+1})\\
&\qquad\vdots\\
&\qquad+(-1)^{(\rho_1+\cdots\rho_{k-1})\rho_k}\sum_{i=1}^{\rho_k}(-1)^{i+1}\pair{k^{(i)}}{x}(k^{(1)}\cdots k^{(i-1)}k^{(i+1)}\cdots k^{(\rho_k)}A^{\rho_1}\cdots (K-1)^{\rho_{k-1}}q_{n-\rho+1}).
\end{align*}
每一個單獨的和（但不是其中一部分）確實把行 $A^{\rho_1},B^{\rho_2}\ldots$ 作為不可分解量，因為它滿足必要且充分的條件\footnote{Capelli，同上，\S\ 23 給出充分條件：$\bigl(a^{(k)}\mid \frac{\partial}{\partial a^{(i)}}\bigr)=0$, $(k>i,\ i=1,2,\ldots,\rho_1-1)$。通過 Wellstein 的 Poisson 括號過程（Von den Differentialgleichungen der projektiven Invarianten, Math. Ann. 67 (1909) \S\ 3），這些條件最簡單地推出必要且充分條件。}
\begin{equation}
 \left(a^{(i+1)}\middle|\frac{\partial}{\partial a^{(i)}}\right)=0;\ \ldots\ ;
 \left(k^{(j+1)}\middle|\frac{\partial}{\partial k^{(j)}}\right)=0.
 \quad (i=1,2,\ldots,\rho_1-1;\ j=1,2,\ldots,\rho_k-1)
\tag{19}
\end{equation}
還須說明它也顯式包含所有行。若置 $B^{\rho_2}\cdots K^{\rho_k}q_{n-\rho+1}=q_{n-\rho_1+1}$，則由 (16.) 與 (10.) 得
\begin{align*}
&\sum(-1)^{i+1}\pair{a^{(i)}}{x}\,
 (a^{(1)}\cdots a^{(i-1)}a^{(i+1)}\cdots a^{(\rho_1)}q_{n-\rho_1+1})\\
&\quad=\pair{A^{\rho_1}}{xp_{\rho_1-1}}\\
&\quad=(A_{n-\rho_1}xp_{\rho_1-1})=(-1)^{(\rho_1+1)(n+1)}\pair{q_{n-\rho_1+1}}{A_{n-\rho_1}x}\\
&\quad=(-1)^{(\rho_1+1)(n+1)}\pair{B^{\rho_2}\cdots K^{\rho_k}q_{n-\rho+1}}{A_{n-\rho_1}x}.
\end{align*}
對其餘各和作同樣計算，繼而因各行已由權指標區別，把 $B^{\rho_2}=A^{\rho_2}$、$C^{\rho_3}=A^{\rho_3},\ldots$ 等代入，即得所需恆等式：
\begin{equation}
\begin{aligned}
\pair{A^{\rho_1}A^{\rho_2}\cdots A^{\rho_k}}{xp_{\rho-1}}
&=\sum_{i=1}^{k}(-1)^{\delta_i}\pair{A^{\rho_1}\cdots A^{\rho_{i-1}}A^{\rho_{i+1}}\cdots A^{\rho_k}q_{n-\rho+1}}{A_{n-\rho_i}x},\\
\delta_i&=(\rho_1+\rho_2+\cdots+\rho_{i-1})\rho_i+(\rho_i+1)(n+1),
\end{aligned}
\tag{20}
\end{equation}
並由 (17.) 得到對偶恆等式：
\begin{equation}
\begin{aligned}
\pair{u q_{\sigma-1}}{A_{\sigma_1}A_{\sigma_2}\cdots A_{\sigma_k}}
&=\sum_{i=1}^{k}(-1)^{\eps_i}\pair{u A^{n-\sigma_i}}{p_{n-\sigma+1}A_{\sigma_1}\cdots A_{\sigma_{i-1}}A_{\sigma_{i+1}}\cdots A_{\sigma_k}},\\
\eps_i&=(\sigma_1+\sigma_2+\cdots+\sigma_{i-1})\sigma_i+(\sigma_i+1)(n+\sigma).
\end{aligned}
\tag{21}
\end{equation}
\emph{定理 II.} 由 (20.) 與 (21.) 已窮盡行 $A^{\rho_i},x,p$ 之間，分別行 $A_{\sigma_i},u,q$ 之間的所有不可分解恆等式；同時也窮盡那些不用執行代換 (11.) 便容許顯式表示的不可分解恆等式。

第一部分由導出 (20.) 與 (21.) 的討論而來；第二部分則來自兩個符號行之間不存在如下恆等式：
\[
\pair{q_\rho A^\sigma}{B_{\rho+\sigma-\tau}p_\tau}
=\eps_1\pair{q_\rho B^{n-\rho-\sigma+\tau}}{A_{n-\sigma}p_\tau}
+\eps_2\pair{q_\rho q_{n-\tau}}{B_{\rho+\sigma-\tau}A_{n-\sigma}},
\]
其中 $\rho\leqq\tau\leqq\sigma$ 且沒有一個行的權為 $1\cdot(n-1)$。其餘關於 $\rho,\tau,\sigma$ 的假設可通過交換行名歸約到這一情形。該不可能性可通過把 (8.) 在特化 $q_\rho=lM^{\rho-1}$、$q_{n-\tau}=lN^{n-\tau-1}$ 下展開來驗證，例如取 $l_1=1$、$M^{2,3,\ldots\rho}=1$、$A^{\rho+1,\ldots\rho+\sigma}=1$，而行 $l,M,A$ 的其餘元素為零，$N$ 與 $B$ 任意。把任意行間恆等式中的 $p_\tau$ 分解為 $y^{(1)}y^{(2)}\cdots y^{(\tau)}$ 後，它們化為由 (20.) 複合而成的恆等式；因此只要說明 (20.) 可由僅含兩個符號行的恆等式生成，第二部分即證。事實上，由
\begin{equation}
\pair{A^{\rho_1}A^{\rho_2}}{xp_{\rho-1}}
=\eps\pair{A^{\rho_1}q_{n-\rho+1}}{A_{n-\rho_2}x}
+\pair{A^{\rho_1}}{xp_{\rho_1-1}},
\quad \text{其中 }p_{\rho_1-1}\sim A^{\rho_2}q_{n-\rho+1},
\tag{20a}
\end{equation}
再特化 $A^{\rho_1}=B^{\sigma_1}B^{\sigma_2}$，即用
\[
\pair{B^{\sigma_1}B^{\sigma_2}}{xp_{\rho_1-1}}
=\eps\pair{B^{\sigma_1}q_{n-\rho_1+1}}{B_{n-\sigma_2}x}
+\pair{B^{\sigma_1}}{xp_{\sigma_1-1}},
\quad \text{其中 } p_{\sigma_1-1}\sim B^{\sigma_2}q_{n-\rho_1+1},
\]
便得到含三個符號行的 (20.) 型恆等式；繼續特化 $B^{\sigma_1}=C^{\tau_1}C^{\tau_2}$ 等，即得一般的 (20.)。兩個符號行與兩個任意變量行之間不存在顯式表示的恆等式，因此在任意多個符號行的情形也不存在這樣的恆等式，只要變量行中不含行 $x$ 或 $u$。
%</unit>

%<unit id="noether.p04.sec.004" type="section" lang="zh-Hans" scan="scan:p04:134-140">
\section*{\S\ 4. 分解恆等式。}
現在考慮任意符號行和變量行之間最一般的恆等式：在其最終表達式中，不執行代換 (11.)，某個行 $p_\tau$ 分解為個別行 $y^{(1)}\cdots y^{(\tau)}$ 出現；我們稱這些恆等式為“分解恆等式”。依照上一段末尾的說明，先在只有兩個符號行的情形中確定它們，也就是展開表達式 $\pair{q_\rho A^\sigma}{B_{\rho+\sigma-\tau}p_\tau}$。令
\begin{equation}
\begin{aligned}
q_{\rho-\alpha}^{(i_1\ldots i_\alpha)}&\sim p_{n-\rho}\,y^{(i_1)}\cdots y^{(i_\alpha)},\qquad
p_{\tau-\alpha}^{(i_1\ldots i_\alpha)}=y^{(i_{\alpha+1})}\cdots y^{(i_\tau)},\qquad
p_\tau=y^{(1)}\cdots y^{(\tau)}.
\end{aligned}
\tag{22}
\end{equation}
關於數 $\rho,\sigma,\tau$ 必須區分四種情形：
\[
\begin{array}{llll}
1)&\rho+\sigma\leqq n,& \tau\leqq\rho,& \tau\leqq\sigma;\\[2pt]
2)&\rho+\sigma\leqq n,& \tau\geqq\rho,& \tau\leqq\sigma;\\[2pt]
3)&\rho+\sigma\leqq n,& \tau\leqq\rho,& \tau\geqq\sigma;\\[2pt]
4)&\rho+\sigma\leqq n,& \tau\geqq\rho,& \tau\geqq\sigma,
\end{array}
\]
先取第一種情形；其餘情形隨後由此簡述。由 (20.)、(10.) 與 (9.)，當把行 $y^{(1)}y^{(2)}\cdots y^{(\tau)}B_{\rho+\sigma-\tau}$ 分解為 $y^{(1)}$ 和 $y^{(2)}\cdots y^{(\tau)}B_{\rho+\sigma-\tau}$ 時，得到
\begin{equation}
\begin{aligned}
\pair{q_\rho A^\sigma}{y^{(1)}y^{(2)}\cdots y^{(\tau)}B_{\rho+\sigma-\tau}}
&=\pair{q_{\rho-1}^{(1)}A^\sigma}{y^{(2)}\cdots y^{(\tau)}B_{\rho+\sigma-\tau}}\\
&\quad+(-1)^\rho\pair{q_\rho[A_{n-\sigma}y^{(1)}]^{\sigma-1}}{y^{(2)}\cdots y^{(\tau)}B_{\rho+\sigma-\tau}}.
\end{aligned}
\tag{23}
\end{equation}
(23.) 的最後一項仍恰有原表達式的形式，只是 $\sigma$ 被 $\sigma-1$、$\tau$ 被 $\tau-1$ 代替。因此在 $\tau\leqq\sigma$ 時可連續應用運算 (23.) 共 $\tau$ 次，並結合 (10.)，得到
\begin{equation}
\begin{aligned}
\pair{q_\rho A^\sigma}{p_\tau B_{\rho+\sigma-\tau}}
&=(-1)^{\rho\sigma+(\sigma+\tau)(n+1)}
  \pair{q_\rho B^{n-\rho-\sigma+\tau}}{A_{n-\sigma}p_\tau}\\
&\quad+\sum_{i_1=1}^{\tau}(-1)^{\rho(i_1-1)}
 \pair{q_{\rho-1}^{(i_1)}[A_{n-\sigma}y^{(1)}y^{(2)}\cdots y^{(i_1-1)}]^{\sigma-i_1+1}}
 {y^{(i_1+1)}\cdots y^{(\tau)}B_{\rho+\sigma-\tau}}.
\end{aligned}
\tag{24}
\end{equation}
求和號下每項又是原來類型的表達式；其中 $\rho$ 換為 $\rho-1$，$\sigma$ 換為 $\sigma-i_1+1$，$\tau$ 換為 $\tau-i_1$，因而條件 1) 更強地滿足。既然 $\tau\leqq\rho$，可把運算 (24.) 再應用 $\tau$ 次，得到所需的分解恆等式：
\begin{equation}
\begin{aligned}
\pair{q_\rho A^\sigma}{p_\tau B_{\rho+\sigma-\tau}}
&=\eps_{i_0}\pair{q_\rho B^{n-\rho-\sigma+\tau}}{A_{n-\sigma}p_\tau}\\
&\quad+\sum_{i_1=1}^{\tau}\eps_{i_1}
 \pair{q_{\rho-1}^{(i_1)}B^{n-\rho-\sigma+\tau}}{A_{n-\sigma}p_{\tau-1}^{(i_1)}}\\
&\quad+\sum_{i_2=i_1+1}^{\tau}\eps_{i_2}
 \pair{q_{\rho-2}^{(i_1i_2)}B^{n-\rho-\sigma+\tau}}{A_{n-\sigma}p_{\tau-2}^{(i_1i_2)}}\\ &\quad\vdots\\
&\quad+\eps_{i_\tau}\pair{q_{\rho-\tau}^{(i_1i_2\ldots i_\tau)}B^{n-\rho-\sigma+\tau}}{A_{n-\sigma}}\\
&=\sum_{\alpha=0}^{\tau}\sum_{i_\alpha=i_{\alpha-1}+1}^{\tau}\eps_{i_\alpha}
 \pair{q_{\rho-\alpha}^{(i_1\ldots i_\alpha)}B^{n-\rho-\sigma+\tau}}{A_{n-\sigma}p_{\tau-\alpha}^{(i_1\ldots i_\alpha)}}.
\end{aligned}
\tag{25}
\end{equation}
符號因子為
\begin{equation}
\begin{aligned}
\eps_{i_\alpha}&=(-1)^{\delta_{i_\alpha}},\qquad
\delta_{i_\alpha}=\sum_{\beta=1}^{\alpha}(\rho-\beta+1)(i_{\beta-1}+i_\beta+1)\\
&\quad +(\rho-\alpha)(\sigma+i_\alpha+\alpha)+(\sigma+\tau+\alpha)(n+1),
\end{aligned}
\tag{26}
\end{equation}
其中置 $i_0=0$。(25.) 中求和號的含義是：對每個組合 $i_1i_2\ldots i_{\alpha-1}$，其中 $i_1<i_2\cdots<i_{\alpha-1}$，附加所有滿足 $i_{\alpha-1}<i_\alpha\le\tau$ 的 $i_\alpha$。由 (26.) 容易驗證，(25.) 中每個單獨的和滿足與 (19.) 類似的微分方程
\begin{equation}
 \left(\frac{\partial}{\partial y^{(i)}}\middle|y^{(i+1)}\right)=0,
 \qquad (i=1,2,\ldots\tau-1),
\tag{27}
\end{equation}
因而只依賴於行 $p_\tau$。於是分解恆等式是行 $A,B,q,p$ 之間的不可分解恆等式；其顯式表示將在下一段再討論。

在情形 2) 中，運算 (23.) 同樣可應用 $\tau$ 次，但運算 (24.) 在第 $\rho$ 次應用後停止。因此 (25.) 終表達式中的第一個求和號必須換為 $\sum_{\alpha=0}^{\rho}$。在情形 3) 與 4) 中，運算 (23.) 只能進行 $\sigma$ 次；此時代替 (24.) 的關係是
\begin{equation}
\begin{aligned}
\pair{q_\rho A^\sigma}{p_\tau B_{\rho+\sigma-\tau}}
&=(-1)^{\rho\sigma}
 (q_{\rho-(\tau-\sigma)}^{(\sigma+1\ldots\tau)}B^{n-\rho+\tau-\sigma})(A_{n-\sigma}p_{\tau-(\tau-\sigma)}^{(\sigma+1\ldots\tau)})\\
&\quad+\sum_{i_1=1}^{\sigma}(-1)^{\rho(i_1-1)}
 \pair{q_{\rho-1}^{(i_1)}[A_{n-\sigma}y^{(1)}\cdots y^{(i_1-1)}]^{\sigma-i_1+1}}
 {y^{(i_1+1)}\cdots y^{(\tau)}B_{\rho+\sigma-\tau}};
\end{aligned}
\tag{28}
\end{equation}
重複 (28.) 共 $(\tau-\sigma)$ 次，且求和號相應變為 $\sum_{i_\beta=i_{\beta-1}+1}^{\sigma+\beta-1}$（這由對 (28.) 中求和項作 $(\sigma-i_1+1)$ 次 (23.) 可見），便給出 (25.) 中屬於指標 $\alpha=\tau-\sigma$ 的單獨和的全部項以及相應符號。於是情形 3) 與 4) 過渡為情形 1) 與 2)，因為對應於 $\sigma,\tau$ 的數變為 $\sigma'=\sigma-i_{\tau-\sigma}+\tau-\sigma$、$\tau'=\tau-i_{\tau-\sigma}=\sigma'$；程序繼續後有 $\tau'<\sigma'$。因此在 (25.) 的最終表達式中，第一個求和號須換成 $\sum_{\alpha=\tau-\sigma}^{\tau,\rho}$。

類似地，由 (21.) 得到對偶的分解恆等式，其中行 $q_\rho$ 分解為各行 $v^{(1)},\ldots,v^{(\rho)}$。若置
\begin{equation}
 \dot p_{\tau-\beta}^{(i_1\ldots i_\beta)}\sim v^{(i_1)}\cdots v^{(i_\beta)}q_{n-\tau};\qquad
 \dot q_{\rho-\beta}^{(i_1\ldots i_\beta)}=v^{(i_{\beta+1})}\cdots v^{(i_\rho)};\qquad
 q_\rho=v^{(1)}\cdots v^{(\rho)},
\tag{29}
\end{equation}
則
\begin{equation}
 \pair{q_\rho A^\sigma}{p_\tau B_{\rho+\sigma-\tau}}
 =\sum_{\beta=\max(0,\tau-\sigma)}^{\min(\tau,\rho)}
   \sum_{i_\beta=i_{\beta-1}+1}^{\rho}\eps
\pair{\dot q_{\rho-\beta}^{(i_1\ldots i_\beta)}B^{n-\rho-\sigma+\tau}}{A_{n-\sigma}\dot p_{\tau-\beta}^{(i_1\ldots i_\beta)}} ,
\tag{30}
\end{equation}
其中求和指標 $\beta$ 從兩個下限中較大者取到兩個上限中較小者。(25.) 與 (30.) 中每個單獨和的各項，都是同一個型經由極化過程
\[
\left(\frac{\partial}{\partial p_{\tau-\alpha}}\middle|p_{\tau-\alpha}^{(i_1\ldots i_\alpha)}\right),\qquad
\left(q_{\rho-\alpha}^{(i_1\ldots i_\alpha)}\middle|\frac{\partial}{\partial q_{\rho-\alpha}}\right)
\]
生成的極化式。為表達這一點，用 $\Pi$ 表示乘以常數的這類極化運算的聚合，於是代替 (25.) 與 (30.) 可得關係
\begin{equation}
 \pair{q_\rho A^\sigma}{p_\tau B_{\rho+\sigma-\tau}}
 =\sum_{\alpha=\max(0,\tau-\sigma)}^{\min(\tau,\rho)}
 \Pi\pair{q_{\rho-\alpha}B^{n-\rho-\sigma+\tau}}{A_{n-\sigma}p_{\tau-\alpha}}.
\tag{31}
\end{equation}
%</unit>

%<unit id="noether.p04.sec.005" type="section" lang="zh-Hans" scan="scan:p04:140-144">
\section*{\S\ 5. 顯式形式的分解恆等式。}
在 (25) 的情形 4) 中把 $\tau$ 特化為 $\sigma+\rho$，便得
\begin{equation}
 \pair{q_\rho A^\sigma}{y^{(1)}\cdots y^{(\rho+\sigma)}}
 =\sum_{1\le i_1<\cdots<i_\rho\le \rho+\sigma}
 \eps_{i_\rho}\pair{q_\rho}{y^{(i_1)}\cdots y^{(i_\rho)}}
 \pair{A^\sigma}{y^{(i_{\rho+1})}\cdots y^{(i_{\rho+\sigma})}},
\tag{32}
\end{equation}
其中 $\eps_{i_\rho}=(-1)^{\delta_{i_\rho}}$，$\delta_{i_\rho}=\rho(\rho+1)/2+i_1+\cdots+i_\rho$。這是比 (16.)、(17.) 更一般的矩陣乘積定理形式，並可由乘積矩陣的行列式的 Laplace 展開得到：
\[
 \sum\pm(v^{(1)}y^{(1)})\cdots(v^{(\rho)}y^{(\rho)})(a^{(1)}y^{(\rho+1)})\cdots(a^{(\sigma)}y^{(\rho+\sigma)}).
\]
因此更一般的乘積定理由 (20.) 與 (21.) 組成，特殊形式 (16.)、(17.) 的選擇也由此得到解釋。

關係 (32.) 給出了顯式表示分解恆等式的手段。為此，把 (31.) 中出現的型看作基型；也就是說，執行代換 (11.)。這並不改變關於行 $A,B$ 的顯式表示，因為由 (8.)，新引入的行 $R$ 可由行 $A,B$ 本身表示，而不依賴於其組成行 $a^{(1)}a^{(2)}\ldots$、$b^{(1)}b^{(2)}\ldots$。於是令
\begin{equation}
 \pair{q_{\rho-\alpha}B^{n-\rho-\sigma+\tau}}{A_{n-\sigma}p_{\tau-\alpha}}
 =\pair{R_{\rho-\alpha}p_{n-\rho+\alpha}}{R^{\tau-\alpha}p_{\tau-\alpha}}.
\tag{33}
\end{equation}
那麼，對應於指標 $\alpha$ 的 (25.) 中的單獨和給出
\begin{equation}
\begin{aligned}
&\sum \eps_{i_\alpha}\pair{R_{\rho-\alpha}p_{n-\rho}y^{(i_1)}\cdots y^{(i_\alpha)}}{R^{\tau-\alpha}y^{(i_{\alpha+1})}\cdots y^{(i_\tau)}}\\
&\quad=\sum \eps_{i_\alpha}\pair{[R_{\rho-\alpha}p_{n-\rho}]^\alpha y^{(i_1)}\cdots y^{(i_\alpha)}}{R^{\tau-\alpha}y^{(i_{\alpha+1})}\cdots y^{(i_\tau)}}\\
&\quad=\eps\pair{[R_{\rho-\alpha}p_{n-\rho}]^\alpha R^{\tau-\alpha}}{p_\tau}
 =\eps\pair{R^{\tau-\alpha}q_{n-\tau}}{R_{\rho-\alpha}p_{n-\rho}}.
\end{aligned}
\tag{34}
\end{equation}
這樣確實得到顯式終表達式；$q_\rho$ 與 $p_\tau$ 在其中對稱出現，說明 (30.) 導向同一終表達式，因而與原來的分解無關。只有當行 $y$ 與 $v$ 仍有獨立意義時，(25.) 與 (30.) 才有差異；此時分解恆等式按定義過渡為可分解恆等式。要到達含更多符號行的恆等式，只需按 \S\ 3 末尾，將 $q_\rho$ 分解為 $q_{\rho_1}C^{\sigma_1}$，或將 $p_\tau$ 分解為 $p_{\tau_1}C_{\tau-\tau_1}$。所得到的表達式
\begin{equation}
 \pair{q_{\rho_1}C^{\sigma_1}}{[R^{\tau-\alpha}q_{n-\tau}]_\alpha R_{\rho_1+\sigma_1-\alpha}},
 \qquad
 \pair{[R_{\rho-\alpha}p_{n-\rho}]^\alpha R^{\tau-\alpha}}{p_{\tau_1}C_{\tau-\tau_1}}
\tag{35}
\end{equation}
恰屬於兩個符號行時已出現的類型，因此在作相應於 (33.) 的代換後，又有顯式表示。重複此過程便得到任意多個行的顯式表示。還應指出，(25.) 與 (30.) 中的第一個和、最後一個和分別不用 (33.)，僅靠公式 (32.)，便給出顯式表示，或者完全化為一個已顯含所有行的單項。把 (32.) 中的 $q_\rho$ 分解為 $q_{\rho_1}C^{\sigma_1}$ 等所得到的關係，可以看作乘積定理的最一般形式。

綜上得到：

\emph{定理 III.} 恆等式
\begin{equation}
 \pair{q_\rho A^\sigma}{p_\tau B_{\rho+\sigma-\tau}}
 =\sum_{\alpha=\max(0,\tau-\sigma)}^{\min(\tau,\rho)}
 \eps\pair{R^{\tau-\alpha}q_{n-\tau}}{R_{\rho-\alpha}p_{n-\rho}},
\tag{36}
\end{equation}
其中行 $R$ 由 (33) 決定；再加上由把 $q_\rho,p_\tau$ 分解為 $q_{\rho_1}C^{\sigma_1}$ 等所得的恆等式，已經窮盡所有不可分解恆等式。

確實，由此得到的恆等式是同類中最一般的，因為各個行不再受 (20.) 與 (21.) 中仍有的權或數目的限制。也不存在更多行間恆等式；因為把 $p_\tau$ 分解為 $y^{(1)}\cdots y^{(\tau)}$ 時，最一般的恆等式由 (20.) 組成，因而由 \S\ 3 末尾又由 (20a.) 組成；而 (16.) 之後所討論的複合正是 \S\ 4 所進行的複合。向任意多個行過渡，由 (20a.) 的討論可知，是對由分解 $q_\rho$ 為 $q_{\rho_1}C^{\sigma_1}$ 等產生的表達式不斷應用同一運算。因此，直接由 (20.) 而非 (20a.) 過渡並不給出新東西；通過計算也容易驗證：一次在 (25.) 中特化一個行 $q_\rho$，另一次按 \S\ 4 的方法從 (20.) 過渡，二者得到同樣的和；這些和經一般乘積定理化為從 (36.) 得到的恆等式。(20.)、(21.) 對應於特殊情形 $\tau=1$、$\rho=1$；此時只出現兩個單獨和，因而按照前面說明，無須代換 (33.) 的顯式表示與先前所得相一致。

最後略述兩個分解恆等式的特殊情形。在 (31) 中置 $\tau=\sigma$、$\rho=n-\sigma$，並把行 $p_\tau$ 記為 $p'_\sigma\sim q'_{n-\sigma}$。則
\begin{equation}
\begin{aligned}
\pair{A^\sigma}{p_\sigma}\pair{B^\sigma}{p'_\sigma}
-\pair{A^\sigma}{p'_\sigma}\pair{B^\sigma}{p_\sigma}
=\sum_{\alpha=1}^{\sigma,n-\sigma}
\Pi\pair{q_{n-\sigma-\alpha}B^\sigma}{A_{n-\sigma}p_{\sigma-\alpha}}.
\end{aligned}
\tag{37}
\end{equation}
這是 Clebsch 用以把含兩個互伴行 $p_\sigma,p'_\sigma$ 的型展開為初等同變式的公式。\footnote{Clebsch, loc. cit., pp. 25--32. Clebsch 證明右邊各項只依賴於行 $A,B$；顯式表示可通過把 (32.) 作用於其表達式 $\Phi_h$ 得到。} 屬於型 $\pair{A^\sigma}{p_\sigma}^m\pair{B^\sigma}{p'_\sigma}^n$ 的初等同變式為
\begin{equation}
 \left\{\sum_{\alpha=1}^{\sigma,n-\sigma}
 \pair{q_{n-\sigma-\alpha}B^\sigma}{A_{n-\sigma}p_{\sigma-\alpha}}
 \right\}^{\rho}
 \pair{A^\sigma}{p_\sigma}^{m-\rho}\pair{B^\sigma}{p_\sigma}^{n-\rho},
 \qquad (\rho=0,1,\ldots m,n).
\tag{38}
\end{equation}
我們將簡短地說，它們由基型經“缺陷 $\alpha$ 的互伴摺疊”產生（參見 \S\ 2）。

其次，置 $\tau=\sigma-\lambda$、$\rho=n-\sigma-\lambda$、$B^\sigma=A^\sigma$。則
\begin{equation}
\begin{aligned}
\pair{q_{n-\sigma-\lambda}A^\sigma}{A_{n-\sigma}p_{\sigma-\lambda}}
&=(-1)^\lambda\pair{q_{n-\sigma-\lambda}A^\sigma}{A_{n-\sigma}p_{\sigma-\lambda}}\\
&\quad+(-1)^{(n-\sigma)(\sigma-\lambda)}
 \sum_{\alpha=1}^{\sigma-\lambda,n-\sigma-\lambda}
 \Pi\pair{q_{n-\sigma-\lambda-\alpha}A^\sigma}{A_{n-\sigma}p_{\sigma-\lambda}}.
\end{aligned}
\tag{39}
\end{equation}
因而，如 \S\ 2 所述，當缺陷 $\lambda$ 為奇數時，刻畫符號行的條件 (12.) 是更高缺陷條件的後果。對於線性型 $\pair{A^\sigma}{p_\sigma}\equiv\pair{B^\sigma}{p_\sigma}$，(39.) 蘊含其在係數中二次的不可約不變構成式化歸為偶缺陷構成式；這與二元、三元領域中的已知事實相符：線性型沒有不變構成式。

還須注意，一般恆等式原是在各個行滿足條件 (12.) 的假設下導出的。但這些行在恆等式中只線性出現，所以這些恆等式同樣適用於不滿足條件 (12.) 的更一般行。
%</unit>

%<unit id="noether.p04.sec.006" type="section" lang="zh-Hans" scan="scan:p04:144-150">
\section*{\S\ 6. 不變構成式的顯式表示。\\摺疊過程和微分過程。}
前幾段的結果使我們能夠完成 \S\ 2 所述的符號可表示性定理（第一基本定理），即補上顯式可表示性：

\emph{定理 IV.} “任意基型的所有不變構成式都可顯式地表示為矩陣積；也就是說，在最終表達式中，除變量行 $p_\rho$ 外，只出現原型的行 $S^\sigma,\ldots T^\tau$，或由代換 (9.)、(11.) 從它們導出的行 $P,Q,R$。”

證明如下。設一個不變構成式除變量行外還依賴於行 $S^\sigma,\ldots T^\tau$；並把這些行充分分解為個別行 $S^\sigma=S^{\sigma_1}\cdots S^{\sigma_i},\ldots T^\tau=T^{\tau_1}\cdots T^{\tau_k}$，使行列式表示成為可能。將這些行按某個任意但固定的次序 $S^\sigma,\ldots T^\tau$ 排列。先取出一個依賴於第一個總行 $S^\sigma=S^{\sigma_1}\cdots S^{\sigma_i}$、而不僅依賴於其個別組成行的行列式聚合。按 \S\ 3--\S\ 5，這種依賴是一般恆等式的後果。若分解變量行 $p_\rho$ 為行 $y,v$，或分解後面的某個行 $T$ 為個別行 $t$ 或 $\tau$，則最一般恆等式由 (20.) 與 (21.) 組成；而這些恆等式總導向一個顯含行 $S^\sigma=S^{\sigma_1}\cdots S^{\sigma_i}$ 的表達式，即 (20.)、(21.) 的左邊。因為由 (19.) 可檢驗，只有右邊完整的和，而非對應於 (20a.) 某項的一部分，才具有僅依賴於 $S^\sigma$ 的性質。程序繼續時，凡已顯式進入的行都保持顯式；恆等式的應用至多要求把若干行合併為一個行，即使用代換 (9.)。若最後只剩單個行 $T^\tau$ 需要顯式化，則有兩種情況：若仍有自由變量行，即未通過代換 (9.) 與其他行相連的行，則把它分解為組成行 $y$ 或 $v$，如 (31.) 那樣得到一個或多個已顯含所有行的型的極化式；若沒有自由變量行，則由其生成方式可知，在含個別行 $t$ 或 $\tau$、卻依賴於行 $T^\tau$ 的表達式全體中，正出現一個分解恆等式的各項；按 \S\ 5，這些在應用代換 (11.) 後總有關於所有行的顯式表示。這便證明了定理的完全一般性。還應注意：只有兩個符號行時，代換 (11.) 從不出現。因為 $\sigma,n-\sigma,\tau,n-\tau<n$，無變量行出現只可能在兩個符號行合併成單個行列式並因而顯式出現時發生；而一旦變量行出現，(34.) 與 (33.) 就說明不需要代換 (11.)。

由此可知，為生成所有不變構成式，只需在附加變量行後，把符號行合併成所有可能的矩陣積。最一般矩陣積有形式
\begin{equation}
 \pair{P^{\mu_1}\cdots P^{\mu_i}}{Q_{\nu_1}\cdots Q_{\nu_k}},\qquad \sum\mu=\sum\nu,
\tag{40}
\end{equation}
其中 $P$ 與 $Q$ 可以是原型的行，也可以是由原來的行通過一串代換 (9.) 或 (11.) 產生的行，最後也可以是變量行。但 (40.) 可由每次只把兩個符號行藉助變量行合併為一個矩陣積而逐次生成；因為由
\[
 \pair{q_{n-\mu-\lambda}P^\mu}{Q_{\nu_1}p_{n-\mu-\nu_1-\lambda}}=\pair{R_\lambda Q_{\nu_1}}{p_{n-\mu-\nu_1-\lambda}}
\]
再同 $Q_{\nu_2}$ 合併，得到
\[
 \pair{[R_\lambda Q_{\nu_1}]^{n-\lambda-\nu_1}}{Q_{\nu_2}p_{n-\mu-\nu_1-\nu_2-\lambda}}
 =\pair{q_{n-\mu-\lambda}P^\mu}{Q_{\nu_1}Q_{\nu_2}p_{n-\mu-\nu_1-\nu_2-\lambda}},
\]
繼續下去便得 (40.)。若 $P,Q$ 不是原型中的行，則由 (9.)、(11.) 與剛才所證，它們本身也是由每次合併兩個符號行的一串過程產生。因此：

\emph{定理 V.} 生成所有不變構成式的過程，即摺疊過程，就是每次藉助變量行把兩個符號行合併為一個矩陣積並逐次進行。

由兩個符號行 $S^\sigma,T^\tau$（其中 $\sigma\geqq\tau$）可形成下列矩陣積：
\begin{align}
&\pair{q_{n-\sigma-\lambda}S^\sigma}{T_{n-\tau}p_{\tau-\lambda}}; &&(\lambda=1,2,\ldots\tau,n-\sigma)\tag{41}\\
&\pair{q_{n-\tau-\mu}T^\tau}{S_{n-\sigma}p_{\sigma-\mu}}; &&(\mu=\sigma-\tau+\lambda)\tag{42}\\
&\pair{q_{n-\tau-\nu}T^\tau}{S_{n-\sigma}p_{\sigma-\nu}}. &&(\nu=1,2,\ldots\sigma-\tau)
\tag{43}
\end{align}
然而由 (31.) 可得 (42.) 與 (43.) 的值：
\begin{align}
&\pair{q_{n-\sigma-\lambda}S^\sigma}{T_{n-\tau}p_{\tau-\lambda}}
 +\sum_{\alpha}\Pi\pair{q_{n-\sigma-\lambda-\alpha}S^\sigma}{T_{n-\tau}p_{\tau-\lambda-\alpha}};\tag{42a}\\
&\Pi(q_{n-\sigma}S^\sigma)(T_{n-\tau}p_\tau)
 +\sum_{\beta}\Pi\pair{q_{n-\sigma-\beta}S^\sigma}{T_{n-\tau}p_{\tau-\beta}}.
\tag{43a}
\end{align}
因此，型 (42.) 可線性地由 (41.) 及其極化式表示，型 (43.) 可由基型及 (41.) 的極化式表示。同樣，(31.) 給出 (41.) 對 (42.) 及其極化式的線性依賴。所以按 Clebsch 的定義，\footnote{Clebsch, loc. cit., \S\ 7.} (42.) 與 (41.) 等價，而 (43.) 回到基型。於是生成過程同樣可以用 (41.) 或 (42.) 給出；我們把關於數 $\lambda$ 較簡單的過程 (41.) 定義為摺疊過程。數 $\lambda$，即摺疊缺陷（參見 \S\ 2），給出行列式乘積中缺少的行數，指的是對給定摺疊含最多行的那個矩陣積。由於 $\lambda$ 只取到 $\tau,n-\sigma$ 中較小者，且 $\sigma\geqq\tau$，有
\begin{equation}
 \lambda\leqq\frac n2,\ n\text{ 為偶數};\qquad
 \lambda\leqq\frac{n-1}{2},\ n\text{ 為奇數}.
\tag{44}
\end{equation}
即使通過進一步摺疊，行 $p,q$ 已經變為符號行，(42.) 與 (43.) 到 (41.) 的化歸仍然有效；因為由 (36.) 有
\[
 \pair{A^{n-\tau-\varkappa}T^\tau}{S_{n-\sigma}B_{\sigma-\varkappa}}
 =\sum_{\alpha=0,\varkappa-\sigma+\tau}^{\tau,n-\sigma}
 \eps\pair{R^{\tau-\alpha}B^{n-\sigma+\varkappa}}{A_{\tau+\varkappa}R_{n-\sigma-\alpha}},
\]
其中行 $R$ 由 $S$ 與 $T$ 按 (33.) 得來。因而這些只含符號行的型，可通過把 $\pair{q_\tau A^{n-\tau-\varkappa}}{B_{\sigma-\varkappa}p_{n-\sigma}}$ 或 $\pair{q_{\tau-\alpha}B^{n-\sigma+\varkappa}}{A_{\tau+\varkappa}p_{n-\sigma-\alpha}}$（按 $\sigma-\tau\gtreqless 2\varkappa$ 分別取之）同基型及型 (41.) 摺疊而得。

由符號摺疊過程按通常方式得到不變微分過程：只須把符號行 $S^\sigma,T_{n-\tau}$ 換成微分符號行 $\frac{\partial}{\partial p_\sigma}$、$\frac{\partial}{\partial q_{n-\tau}}$；眾所周知，乘積 $\frac{\partial}{\partial p_{i_1\ldots i_\sigma}}\cdot\frac{\partial}{\partial q_{j_1\ldots j_{n-\tau}}}$ 的實際意義是 $\frac{\partial^2}{\partial p_{i_1\ldots i_\sigma}\partial q_{j_1\ldots j_{n-\tau}}}$。行 $\frac{\partial}{\partial p_\sigma}$ 與 $\frac{\partial}{\partial q_{n-\tau}}$ 同行 $S^\sigma,T_{n-\tau}$ 互為同變，所以滿足相同恆等式，特別滿足 (42.) 與 (42a.) 之間、(43.) 與 (43a.) 之間的恆等式。綜上得到：

\emph{定理 VI.} 任意基型的所有不變構成式都由反覆應用符號摺疊過程 (41.) 而來，或者也可由反覆應用不變微分過程
\begin{equation}
 \left(q_{n-\sigma-\lambda}\frac{\partial}{\partial p_\sigma}\middle|\frac{\partial}{\partial q_{n-\tau}}p_{\tau-\lambda}\right),
 \qquad(\sigma\geqq\tau,\ \lambda=1,2,\ldots\sigma,n-\tau).
\tag{45}
\end{equation}
得到。還須指出，在每一次摺疊 (41.) 中，型的總權，即各變量行權的和（參見 \S\ 1），減少正數 $2(\lambda^2+\lambda(\sigma-\tau))$。這不依賴於一般有限性證明，已經推出型行的有限性，也就是給定基型的係數線性不變構成式全體的有限性。\footnote{參見 Clebsch, loc. cit., \S\ 11 和 12：初等同變式系統的有限性。}

還應指出，定理 VI 對不滿足條件 (12.) 的型仍成立。因為 (3) 中每個因子 $\pair{S^\rho}{p_\rho}^{m_\rho}$ 都可替換為 $m_\rho$ 個線性因子的乘積 $\pair{A^\rho}{p_\rho}\pair{B^\rho}{p_\rho}\cdots\pair{M^\rho}{p_\rho}$；於是其不變構成式變為一個線性型系統的不變構成式，最後再把這些線性型彼此相等即可。對於每一個單獨線性型，條件 (12.)——引入微分符號時只含二階微分商——總是滿足的。
%</unit>

%<unit id="noether.p04.sec.007" type="section" lang="zh-Hans" scan="scan:p04:150-156">
\section*{\S\ 7. 一般型按正規型的展開。}
以下（\S\ 8）將導出摺疊過程 (41.) 的一個主要性質，即它可由基本摺疊複合而成。為此需要把 Mertens\footnote{F. Mertens, \emph{\"Uber invariante Gebilde quatern\"arer Formen}（見引言），下稱 “Mertens”。} 在四元領域給出的一個重要定理轉移到 $n$ 元領域。該定理說：任意有限型系統都可由一個等價的正規型系統代替。所謂“正規型”，是對二元、三元定義\footnote{Study, loc. cit., p. 54.}的推廣，指除基型本身外，其所有係數線性的不變構成式都恆等為零的型。因此，由定理 VI（\S\ 6），正規型 $\varphi$ 滿足微分方程組
\begin{equation}
 \left(q_{n-\sigma-\lambda}\frac{\partial}{\partial p_\sigma}\middle|\frac{\partial}{q_{n-\tau}}p_{\tau-\lambda}\right)=0,
 \qquad(\sigma\geqq\tau;\ \lambda=1,2,\ldots\tau,n-\sigma),
\tag{46}
\end{equation}
其中行 $q_{n-\sigma-\lambda},p_{\tau-\lambda}$ 看作任意量。在四元領域中，考慮到 $\sigma=\tau=2$ 時的 (39.)，這些微分方程與 Mertens 在不引入任意行 $p,q$ 的情況下給出的十個微分方程完全一致（loc. cit., formula 28）。這些方程同定理 VI 一起，使其證明幾乎逐字轉移到 $n$ 個變量。唯一要補充的是：用分解恆等式 (30.) 驗證所構造的型確為正規型；在四元領域，這一點尚無需進一步變形。因此只須給出證明的簡短說明，並且為簡單起見，只在下面需要的情形中討論，即單個基型 $f$ 及其型行（參見 \S\ 6 末尾）；對於任意多個基型的係數任意次數不變構成式，證明相同。

證明的基本思想是：引入變換後的係數，把含 $n-1$ 個行 $p_\rho$ 的型替換為含 $n$ 個與 $x$ 同變的行 $\xi$ 的型；這些行就是變換量的行。所求正規型由這些型通過不變微分過程直接得到。一般型按正規型展開，則由含 $\rho$ 個（$\rho\leqq n$）互伴行 $\xi$ 的型的級數展開來中介；不變微分過程再把它們帶回原來的行 $p_\rho$ 中。

令 $\xi^{(1)},\ldots\xi^{(n)}$ 為變換量的行，於是
\begin{equation}
\begin{aligned}
x_i&=\xi_i^{(1)}x'_1+\xi_i^{(2)}x'_2+\cdots+\xi_i^{(n)}x'_n,\\
p_{i_1\ldots i_\rho}&=\sum_j(\xi_{i_1}^{(j_1)}\cdots \xi_{i_\rho}^{(j_\rho)})p'_{j_1\ldots j_\rho}.
\end{aligned}
\tag{47}
\end{equation}
型 $f,\varphi,\ldots$ 的變換後係數記為 $(f),(\varphi),\ldots$，並令
\[
 f\equiv\pair{S^1}{p_1}^{m_1}\pair{S^2}{p_2}^{m_2}\cdots\pair{S^{n-1}}{p_{n-1}}^{m_{n-1}}.
\]
於是
\begin{equation}
\begin{aligned}
(f)&=\pair{S^1}{\xi^{(1)}}^{m_{11}}\cdots\pair{S^1}{\xi^{(n)}}^{m_{1n}}
\pair{S^2}{\xi^{(1)}\xi^{(2)}}^{m_{21}}\cdots
\pair{S^{n-1}}{\xi^{(2)}\cdots\xi^{(n)}}^{m_{n-1,n}}\\
&=\pair{s^{(1)}}{\xi^{(1)}}^{k_1}\pair{s^{(2)}}{\xi^{(2)}}^{k_2}\cdots\pair{s^{(n)}}{\xi^{(n)}}^{k_n},
\end{aligned}
\tag{48}
\end{equation}
其中 $m_{11}+\cdots+m_{1n}=m_1$，等等。對於每一個滿足
\begin{equation}
 k_1\ge k_2\ge k_3\ge\cdots\ge k_n,
\tag{49}
\end{equation}
的係數 $(f)$，正好耗盡行 $\xi$ 的微分過程
\begin{equation}
\left(\frac{\partial}{\partial\xi^{(1)}}\middle|p_1\right)^{k_1-k_2}
\left(\frac{\partial}{\partial\xi^{(1)}}\frac{\partial}{\partial\xi^{(2)}}\middle|p_2\right)^{k_2-k_3}
\cdots
\left(\frac{\partial}{\partial\xi^{(1)}}\cdots\frac{\partial}{\partial\xi^{(n-1)}}\middle|p_{n-1}\right)^{k_{n-1}-k_n}
\left(\frac{\partial}{\partial\xi^{(1)}}\cdots\frac{\partial}{\partial\xi^{(n)}}\right)^{k_n}
\tag{50}
\end{equation}
給出型
\begin{equation}
 \varphi=\bigl(s^{(1)}\mid p_1\bigr)^{k_1-k_2}
 \bigl(s^{(1)}s^{(2)}\mid p_2\bigr)^{k_2-k_3}\cdots
 \bigl(s^{(1)}\cdots s^{(n-1)}\mid p_{n-1}\bigr)^{k_{n-1}-k_n}
 \bigl(s^{(1)}s^{(2)}\cdots s^{(n)}\bigr)^{k_n}.
\tag{51}
\end{equation}
這些 $\varphi$ 就是所求正規型。它們滿足定義正規型的性質。

第一，它們是基型 $f$ 的、關於係數線性的不變構成式。這是因為它們由 $f$ 通過不變微分過程 $\left(\frac{\partial}{\partial p_\rho}\middle|\xi^{(i_1)}\xi^{(i_2)}\cdots\xi^{(i_\rho)}\right)$ 和 (50.) 得到。因此按 \S\ 6 末尾，它們可由 $f$ 的有限型行的極化式線性表示。這些極化式的取得方式是：把進入 $\varphi$ 的變量 $p_\rho$ 看作一個分解為變量行 $\mathfrak{p}_\sigma$ 的線性型的型的係數：
\begin{equation}
 H=(q_{n-1}\mid \mathfrak{p}_{n-1})^{k_1-k_2}(q_{n-2}\mid \mathfrak{p}_{n-2})^{k_2-k_3}\cdots(q_1\mid \mathfrak{p}_1)^{k_{n-1}-k_n}.
\tag{52}
\end{equation}
型行 $f$ 與 $H$ 的同時不變量給出所有相關極化式，因為後者關於各行 $p_\rho$ 中變數的次數必須與 $\varphi$ 本身相同。

第二，它們滿足微分方程 (46.)。若把微分過程 (45.) 作用於 $\varphi$，則
\[
 \varphi'=\bigl(q_{n-\sigma-\lambda}s^{(1)}\cdots s^{(\sigma)}\mid [s^{(1)}\cdots s^{(\tau)}]_{n-\tau}p_{\tau-\lambda}\bigr).\qquad (\sigma\geqq\tau).
\]
由分解恆等式 (30.)，若把由行 $s$ 構成的 $q_\sigma=s^{(1)}\cdots s^{(\sigma)}$、$p_{n-\tau}=[s^{(1)}\cdots s^{(\tau)}]_{n-\tau}$ 同那裡出現的 $q_\rho,p_\tau$ 識別，則
\[
 \varphi'=\sum_{\beta=\sigma-\tau+\lambda}^{n-\tau,\sigma}\sum_{i_\beta}\eps\,
 \pair{\dot{q}_{\sigma-\beta}^{(i_1\ldots i_\beta)}q_{n-\tau+\lambda}}{p_{\sigma+\lambda}\dot{p}_{n-\tau-\beta}^{(i_1\ldots i_\beta)}}.
\]
這裡
\[
 \dot{p}_{n-\tau-\beta}^{(i_1\ldots i_\beta)}\sim s^{(1)}\cdots s^{(\tau)}s^{(i_1)}\cdots s^{(i_\beta)},
\]
而 $i_1,\ldots i_\beta$ 表示從 $1$ 到 $\sigma$ 中任取的 $\beta$ 個不同數。由於 $\beta>\sigma-\tau$，這些數中必有某個介於 $1$ 與 $\tau$ 之間；故 $\dot{p}_{n-\tau-\beta}^{(i_1\ldots i_\beta)}$ 恆等為零，每個矩陣積以及 $\varphi'$ 均為零。

為證明正規型系統與型行 $f$ 等價，還需說明型行中的所有型都可按第一項所規定的意義展開為正規型的極化式。令 $\Pi$ 表示極化運算
\begin{equation}
 \left(\frac{\partial}{\partial\xi^{(i)}}\middle|\xi^{(k)}\right),
 \qquad (i\ne k;\ i=1,2,\ldots,\rho;\ k=1,2,\ldots,\rho),
\tag{53}
\end{equation}
的一個聚合，令 $\Pi^\sigma$ 表示滿足 $i=\sigma$ 且 $k<\sigma$ 的特殊運算 (53.) 的一個聚合。若 $F$ 是 $\rho$ 個行 $\xi^{(1)},\ldots,\xi^{(\rho)}$ 的型，在最後一行 $\xi^{(\rho)}$ 中次數為 $k_\rho$，則有展開\footnote{F. Mertens, \emph{\"Uber eine Formel der Determinantentheorie}. Sitzungsberichte der Akademie der Wissenschaften, Vienna, Math.-Naturw. Kl., vol. 91, 2nd section (1885), formula 20).}
\begin{equation}
 F=\sum_{\alpha=0}^{k_\rho}\Pi F_\alpha,
 \qquad(\rho\leqq n),
\tag{54}
\end{equation}
其中 $F_\alpha$ 在最後一行 $\xi^{(\rho)}$ 中次數為 $\alpha$，並由 $F$ 經不變微分運算
\begin{equation}
 \left(\frac{\partial}{\partial\xi^{(1)}}\cdots\frac{\partial}{\partial\xi^{(\rho)}}\middle|\xi^{(1)}\cdots\xi^{(\rho)}\right)^\alpha
\tag{55}
\end{equation}
以及 $\Pi^\rho$ 得到。若在構造 $F_\alpha$ 時，把過程 (55.) 換成
\begin{equation}
 \left(\frac{\partial}{\partial\xi^{(1)}}\cdots\frac{\partial}{\partial\xi^{(\rho)}}\middle|p_\rho\right)^\alpha,
\tag{56}
\end{equation}
則得到只含 $\rho-1$ 個行 $\xi^{(1)},\ldots,\xi^{(\rho-1)}$ 的型 $F_\alpha(p_\rho)$，可以再次應用 (54.)。繼續此程序得到型 $G(p_\rho,p_{\rho-1},\ldots,p_1)$。為了由這些型得到 $F$ 的展開，再把各行 $p_\sigma$ 替換為 $\xi^{(1)},\ldots,\xi^{(\sigma)}$，並對所得型施行極化過程 $\Pi$ (53.)；這樣得到一組變換後係數 $(G)$ 之和（參見 (48.)）。在 $\rho=n$ 時，過程 (55.) 在第一步保留下來。若 $c$ 表示數值常數，並簡記
\begin{equation}
 (\xi^{(1)}\xi^{(2)}\cdots\xi^{(n)})=\Delta;
 \qquad
 \left(\frac{\partial}{\partial\xi^{(1)}}\frac{\partial}{\partial\xi^{(2)}}\cdots\frac{\partial}{\partial\xi^{(n)}}\right)=\nabla,
\tag{57}
\end{equation}
則在 $\rho=n$ 時有
\begin{equation}
 F=\sum c\Delta^\alpha(G),
\tag{58}
\end{equation}
而在 $\rho<n$ 時右邊只出現 $\alpha=0$，即 $\sum c(G)$。

把 (58.) 作用於變換後係數 $(f)$ (48.)，並利用極化過程 $\Pi^\sigma$ 與所有 $\rho\ge\sigma$ 的運算 (56.) 可交換，以及變換後係數的極化式仍給出變換後係數，得到
\[
 G=\left(\frac{\partial}{\partial\xi^{(1)}}\middle|p_1\right)^{\beta_1}
 \left(\frac{\partial}{\partial\xi^{(1)}}\frac{\partial}{\partial\xi^{(2)}}\middle|p_2\right)^{\beta_2}
 \cdots
 \left(\frac{\partial}{\partial\xi^{(1)}}\cdots\frac{\partial}{\partial\xi^{(n-1)}}\middle|p_{n-1}\right)^{\beta_{n-1}}
 \nabla^{\beta_n}\Pi(f)=\sum c\varphi,
\]
因而由 (58.) 得
\begin{equation}
 (f)=\sum c\Delta^\alpha(\varphi)\qquad\text{（Mertens, formula 23）}.
\tag{59}
\end{equation}
為了從 (59.) 過渡到 $f$ 本身的展開，再把變量 $p_\rho$ 看作型 $H$ (52.) 的係數，其次數 $m_1,\ldots m_{n-1}$ 對應於 $f$，並置 $m=m_1+\cdots+m_{n-1}$。由不變量的定義，
\[
 f\cdot\Delta^m=\sum(f)(H)=\sum c\Delta^\alpha(\varphi)(H),
\]
施行正好耗盡行 $\xi$ 的過程 $\nabla^m$，便得
\begin{equation}
 f=\sum c\Phi,
\tag{60}
\end{equation}
其中型 $\Phi$ 由 $\varphi$ 與 $H$ 經關於變量的不變微分過程導出，因而是型行 $\varphi$ 與 $H$ 的同時不變量。\footnote{Mertens 在 \S\ 7 中不用這個直接結論，而引入一串進一步的微分過程，其本質作用是構造型行 $H$；這在這裡由定理 VI 完成。} 但 $\varphi$ 是正規型，故型行 $\varphi$ 化為 $\varphi$ 本身；型行 $H$ 含所有可能的不變變量組合。因此 $\varphi$ 與 $H$ 的同時不變量就是第一項所說意義下的 $\varphi$ 的極化式。換言之，(60.) 給出 $f$ 按正規型極化式的展開。對型行 $f$ 的任意型 $\theta$，還需驗證同樣展開。令 $\Gamma$ 為由 $(\theta)$ 按 (50.) 導出的正規型，則
\begin{equation}
 (\theta)=\sum c\Delta^\beta(\Gamma).
\tag{61}
\end{equation}
型 $\theta$ 由對每個變換後係數成立的關係刻畫：
\[
 (\theta)\Delta^\sigma=\sum c(f)\qquad\text{（Mertens, formula 9）}.
\]
但每個含因子 $\Delta^\sigma$ 的 $n$ 個行 $\xi$ 的型 $F$ 也含第二個 $\nabla^\sigma\Pi F$（Mertens, formula 20）；故
\[
 (\theta)=\sum c\nabla^\sigma(f),\qquad\text{因此}\qquad \Gamma=\sum c\varphi,
\]
由 (61.) 得到與 (59.) 相應的關係
\[
 (\theta)=\sum c\Delta^\beta(\varphi),
\]
於是對 $\theta$ 也推出按正規型 $\varphi$ 的極化式展開 (60)。因此：

\emph{定理 VII.} 每個型行 $f$ 都可由一個等價的正規型系統代替。
%</unit>

%<unit id="noether.p04.sec.008" type="section" lang="zh-Hans" scan="scan:p04:156-169">
\section*{\S\ 8. 摺疊過程化歸為基本摺疊。}
按照定理 VII，現在來實施 \S\ 7 開頭所說的、把摺疊過程化歸為基本摺疊。我們把兩個同變行 $S^\rho,T^\rho$ 的任意摺疊稱為“同變摺疊”，並把兩個行 $S^\rho,T^\rho$ 的缺陷為 1 的同變摺疊稱為“類型 $\rho$ 的基本摺疊”。於是證明如下命題，它補足定理 VI：

\emph{定理 VIII.} 一般摺疊過程 (41.) 可由類型 $\rho$ 的 $(n-1)$ 個基本摺疊複合而成，其中 $\rho=1,2,\ldots n-1$。換言之，為生成所有不變構成式，只需逐次應用這 $n-1$ 個基本摺疊。

在二元領域中，摺疊過程 (41.) 直接就是唯一可能的基本摺疊；在三元領域中，我已在博士論文 \S\ 1 證明定理 VIII（參見引言）。在那裡，由於 (44.)，一般同變摺疊仍等同於基本摺疊。值得注意的是，$n$ 個變量的定理 VIII 給出的不是較弱的“化歸為同變摺疊”，而是真正化歸為基本摺疊。證明在原則上仿照三元情形：藉助符號恆等式，由基本摺疊構造最一般的摺疊。我們將生成：
\begin{enumerate}
\item 缺陷 1 的任意摺疊，由缺陷 1 的同變摺疊，即由基本摺疊生成（三元情形的類似物）；
\item 缺陷 $\lambda$ 的任意摺疊，由缺陷 $\lambda$ 的同變摺疊和缺陷 1 的任意摺疊生成；
\item 缺陷 $\lambda$ 的同變摺疊，由缺陷 $\lambda-1$ 的同變摺疊和缺陷 1 的任意摺疊生成。
\end{enumerate}
關鍵的是，根據定理 VII，假定要摺疊的兩個型 $S$ 與 $T$ 都是正規型。因此由 (46.) 有關係
\begin{equation}
\begin{aligned}
\pair{q_{n-\sigma_1-\lambda}S^{\sigma_1}}{S_{n-\sigma_2}p_{\sigma_2-\lambda}}&=0,
&& (\sigma_1\ge\sigma_2,\ \lambda=1,2,\ldots,\sigma_2,n-\sigma_1),\\
\pair{q_{n-\tau_1-\lambda}T^{\tau_1}}{T_{n-\tau_2}p_{\tau_2-\lambda}}&=0,
&& (\tau_1\ge\tau_2,\ \lambda=1,2,\ldots,\tau_2,n-\tau_1).
\end{aligned}
\tag{62}
\end{equation}
此外，按 \S\ 2 末尾，只需假定型 $S$ 與 $T$ 在所有變量行中都是線性的；也就是說，所構造的型屬於型行
\[
 f=\pair{S^1}{p_1}\pair{S^2}{p_2}\cdots\pair{S^{n-1}}{p_{n-1}}
   \pair{T^1}{p_1}\pair{T^2}{p_2}\cdots\pair{T^{n-1}}{p_{n-1}}.
\]
1) 由類型 $\tau+\alpha$ 的基本摺疊（$\alpha=0,1,\ldots\sigma-\tau$）生成 $\pair{q_{n-\sigma-1}S^\sigma}{T_{n-\tau}p_{\tau-1}}$，其中 $\sigma>\tau$。從類型 $\tau$ 的基本摺疊
\[
 \pair{q_{n-\tau-1}S^\tau}{T_{n-\tau}p_{\tau-1}}
\]
經代換 $w\sim T_{n-\tau}p_{\tau-1}$，得到型行 $f$ 中一個具有三個上權指標（參見 \S\ 1）為 $\tau+1$ 的行的型：
\begin{equation}
 \pair{wS^\tau}{p_{\tau+1}}\pair{S^{\tau+1}}{p_{\tau+1}}\pair{T^{\tau+1}}{p_{\tau+1}}.
\tag{63}
\end{equation}
對 (63.) 施以類型 $\tau+1$ 的基本摺疊，得到
\[
 \pair{q_{n-\tau-2}wS^\tau}{S_{n-\tau-1}p_\tau};
\]
由 (31.)，作代換 $q_{n-\tau-2}w=q'_{n-\tau-1}$，其值為
\[
 \eps\bigl(q_{n-\tau-2}wS^{\tau+1}\bigr)\pair{S^\tau}{p_\tau}
 +\sum_{\alpha=1}^{\tau,n-\tau-1}\Pi\pair{q'_{n-\tau-1-\alpha}S^{\tau+1}}{S_{n-\tau}p_{\tau-\alpha}}.
\]
求和號下的表達式由 (62.) 恆等為零；於是到達
\[
 \pair{wS^{\tau+1}}{p_{\tau+2}}\pair{S^{\tau+2}}{p_{\tau+2}}\pair{T^{\tau+2}}{p_{\tau+2}},
\]
即由 (63.) 把 $\tau$ 改為 $\tau+1$ 所得的型。重複此過程 $(\sigma-\tau)$ 次，遂到達所需摺疊：
\[
 \pair{wS^\sigma}{p_{\sigma+1}}=\eps\pair{q_{n-\sigma-1}S^\sigma}{T_{n-\tau}p_{\tau-1}}.
\]
剛才的過程記為
\begin{equation}
 T^\tau S^\tau,
\quad S^{\tau+1},\quad S^{\tau+2},\ldots,S^\sigma;
\tag{64}
\end{equation}
可見只用了 $S$ 的正規型性質，而沒有用到 $T$ 的。對偶於 (64.) 的過程只用 $T$ 的正規型性質，以相反次序進行：
\begin{equation}
 S^\sigma T^\sigma,
\quad T^{\sigma-1},\quad T^{\sigma-2},\ldots,T^\tau,
\tag{65}
\end{equation}
相應關係為
\begin{align*}
\pair{q_{n-\sigma-1}S^\sigma}{T_{n-\sigma}p_{\sigma-1}}&=\eps\pair{T_{n-\sigma}y}{p_{\sigma-1}},\quad y\sim q_{n-\sigma-1}S^\sigma,\\
\pair{q_{n-\sigma}T^{\sigma-1}}{T_{n-\sigma}yp_{\sigma-2}}&=\eps\pair{T_{n-\sigma+1}y}{p_{\sigma-2}},\\
&\vdots\\
\pair{q_{n-\tau-1}T^\tau}{T_{n-\tau-1}yp_{\tau-1}}&=\eps\pair{q_{n-\sigma-1}S^\sigma}{T_{n-\tau}p_{\tau-1}}.
\end{align*}
還應指出這同 \S\ 6 末尾的總權下降相聯繫：缺陷 1 的摺疊的下降量為 $2(1+(\sigma-\tau))$，基本摺疊的下降量為 $2\cdot1$。因此，若確實能由基本摺疊複合而成，則必需且正好需要 $\sigma-\tau+1$ 個這樣的摺疊，這正是上面完成的過程。

2) 由同變摺疊和基本摺疊生成一般摺疊。一般摺疊 (41) 的權下降為 $2(\lambda^2+\lambda(\sigma-\tau))=2[\lambda^2+(\sigma-\tau)(1+(\lambda-1))]$。這提示可把它分解為一個缺陷 $\lambda$ 的同變摺疊（權下降 $2\lambda^2$）以及 $\sigma-\tau$ 個缺陷 1、權指標差為 $\lambda-1$ 的摺疊（每個下降 $2\{1+(\lambda-1)\}$）。當 $\lambda=1$ 時此分解與 1) 中的分解相同；現在說明它確實可實現。

從缺陷 $\lambda$ 的同變摺疊
\[
 \pair{q_{n-\tau-\lambda}S^\tau}{T_{n-\tau}p_{\tau-\lambda}}
\]
經代換 $R^\lambda\sim T_{n-\tau}p_{\tau-\lambda}$，得到具有上權指標 $\tau+\lambda$ 和 $\tau+1$ 的兩個行的型
\begin{equation}
 \pair{R^\lambda S^\tau}{p_{\tau+\lambda}}\pair{S^{\tau+1}}{p_{\tau+1}}.
\tag{66}
\end{equation}
下權指標為 $\tau+1$ 的行屬於正規型；因此 (66.) 容許一個缺陷 1 的摺疊，它由 $\lambda$ 個基本摺疊按 (65) 的次序組成：
\[
 (R^\lambda S^\tau)S^{\tau+\lambda},\quad S^{\tau+\lambda-1},\quad S^{\tau+\lambda-2},\ldots S^{\tau+1}.
\]
考慮 (31.) 與 (62.)，得到
\[
 \pair{q_{n-\tau-\lambda-1}R^\lambda S^\tau}{S_{n-\tau-1}p_\tau}=\eps\pair{R^\lambda S^{\tau+1}}{p_{\tau+\lambda+1}},
\]
即由 (66.) 把 $\tau$ 替換為 $\tau+1$ 得到的型，正如 1) 中的 (63.)。重複 $(\sigma-\tau)$ 次得到所需摺疊：
\[
 \pair{R^\lambda S^\sigma}{p_{\sigma+\lambda}}=\eps\pair{q_{n-\sigma-\lambda}S^\sigma}{T_{n-\tau}p_{\tau-\lambda}}.
\]
整個過程中只使用 $S$ 的正規型性質。把 $T$ 看作正規型，則得到對偶過程，次序完全相反：
\begin{align*}
\pair{q_{n-\sigma-\lambda}S^\sigma}{T_{n-\sigma}p_{\sigma-\lambda}}&=\eps\pair{T_{n-\sigma}R_\lambda}{p_{\sigma-\lambda}},\quad R_\lambda\sim q_{n-\sigma-\lambda}S^\sigma,\\
\pair{q_{n-\sigma}T^{\sigma-1}}{T_{n-\sigma}R_\lambda p_{\sigma-\lambda-1}}&=\eps\pair{T_{n-\sigma+1}R_\lambda}{p_{\sigma-\lambda-1}},\\
&\vdots\\
\pair{q_{n-\tau-1}T^\tau}{T_{n-\tau-1}R_\lambda p_{\tau-\lambda}}&=\eps\pair{q_{n-\sigma-\lambda}S^\sigma}{T_{n-\tau}p_{\tau-\lambda}}.
\end{align*}
3) 由缺陷 $\lambda-1$ 的同變摺疊和基本摺疊生成缺陷 $\lambda$ 的同變摺疊。同變摺疊缺陷 $\lambda$ 的權下降為 $2\lambda^2=2((\lambda-1)^2+2\lambda-1)$，因而可分解為一個缺陷 $\lambda-1$ 的同變摺疊和 $2\lambda-1$ 個基本摺疊。最容易先取 $n=2\lambda$。缺陷 $\lambda$ 的唯一可能同變摺疊是 $\pair{S^\lambda}{T_\lambda}=\pair{S^\lambda}{T_{n/2}}$；它比同一兩行的缺陷 $\lambda-1$ 摺疊 $\pair{uS^\lambda}{T_\lambda x}$ 少含一個行 $u$ 與一個行 $x$。反過來，用導致一個行 $u$ 與一個行 $x$ 消失的基本摺疊（權下降 $2(n-1)=2(2\lambda-1)$）來複合缺陷 $\lambda$ 的摺疊，同時生成新的符號行 $R^\lambda$，再加上一個缺陷 $\lambda-1$ 的同變摺疊。最簡單的這種缺陷 1 摺疊為
\[
 \pair{sT^\lambda}{\sigma p_\lambda}=\eps\pair{S^{2\lambda-1}}{R_{\lambda-1}p_\lambda}=\eps\pair{R^\lambda}{p_\lambda},
\]
其中
\[
 R^{\lambda+1}=sT^\lambda,
 \quad s=S^1,
 \quad \sigma=S_1,
 \quad R^\lambda\sim\sigma R_{\lambda-1}.
\]
由 (65.) 與 (64.)，它由 $2\lambda-1$ 個基本摺疊組成：
\begin{equation}
 T^\lambda S^\lambda,S^{\lambda-1},\ldots S^1;
 \qquad R^{\lambda+1}S^{\lambda+1},S^{\lambda+2},\ldots S^{2\lambda-1}.
\tag{67}
\end{equation}
考慮 (21.) 與 (62.)，缺陷 $\lambda-1$ 的摺疊給出
\[
 \pair{uS^\lambda}{\sigma R_{\lambda-1}x}=\eps\pair{R^{\lambda+1}}{S_\lambda x}
 =\eps\pair{sT^\lambda}{S^\lambda x}=\eps\pair{S^\lambda}{T_\lambda}.
\]
任意 $n$ 與兩個行 $S^\rho,T^\rho$ 的一般情形可直接由該特殊情形得到。代替 (67.) 有下列 $2\lambda-1$ 個基本摺疊：
\begin{equation}
 T^\rho S^\rho,S^{\rho-1},\ldots S^{\rho-\lambda+1};
 \qquad R^{\rho+1}S^{\rho+1},S^{\rho+2},\ldots S^{\rho+\lambda-1}.
\tag{68}
\end{equation}
前半摺疊給出型
\[
 \pair{q_{n-\rho-1}T^\rho}{S_{n-\rho+\lambda-1}p_{\rho-\lambda}},
\]
由代換
\[
 S_{n-\rho+\lambda-1}p_{\rho-\lambda}\sim w,
 \qquad T^\rho w=R^{\rho+1},
\]
並應用後半摺疊，得到
\[
 \pair{q_{n-\rho-\lambda}S^{\rho+\lambda-1}}{R_{n-\rho-1}p_\rho}.
\]
再代入
\[
 q_{n-\rho-\lambda}S^{\rho+\lambda-1}\sim y,
\]
用一個缺陷 $\lambda-1$ 的同變摺疊得到
\begin{equation}
 \pair{q_{n-\rho-\lambda+1}S^\rho}{yR_{n-\rho-1}p_{\rho-\lambda+1}}.
\tag{69}
\end{equation}
現在說明這就是所需的缺陷 $\lambda$ 摺疊。代入
\[
 R_{n-\rho-1}p_{\rho-\lambda+1}=R_{n-\lambda}
\]
並注意分解 $y$ 後有
\[
 \pair{q_{n-\rho-\lambda+1}R^\lambda}{S_{n-\rho}y}
 =\eps\pair{q_{n-\rho-\lambda}S^{\rho+\lambda-1}}{S_{n-\rho}p'_{\rho-1}}=0.
\]
於是由 (20.)，(69.) 的值為
\[
 \eps\pair{S^\rho R^\lambda}{p_{\rho+\lambda-1}y}
 =\eps\pair{q_{n-\rho-\lambda}S^{\rho+\lambda-1}}{[S^\rho R^\lambda]_{n-\rho-\lambda}p_{\rho+\lambda-1}}.
\]
這一折疊屬於 (43.) 型，因而等價於
\[
 \pair{q_{n-\rho-\lambda}S^\rho}{R_{n-\rho-1}p_{\rho-\lambda+1}}
 =\eps\pair{wT^\rho}{R_\lambda p_{\rho-\lambda+1}},
 \qquad R_\lambda\sim q_{n-\rho-\lambda}S^\rho.
\]
這裡 $R_\lambda$ 已是所需終形式；該表達式對應於特殊情形中的 $\pair{sT^\lambda}{S_\lambda x}$。再注意分解 $w$ 與 $R_\lambda$ 後
\[
 \pair{wR^{n-\lambda}}{T_{n-\rho}p_{\rho-\lambda+1}}
 =\eps\pair{q'_{\lambda-1}R^{n-\lambda}}{S_{n-\rho+\lambda-1}p_{\rho-\lambda}}
 =\eps\pair{q''_{n-\rho-1}S^\rho}{S_{n-\rho+\lambda-1}p_{\rho-\lambda}}=0,
\]
於是由 (21.) 得其值
\[
 \pair{wq_{n-\rho+\lambda-1}}{T_{n-\rho}R_\lambda}
 =\eps\pair{q_{n-\rho+\lambda-1}[T_{n-\rho}R_\lambda]^{\rho-\lambda}}{S_{n-\rho+\lambda-1}p_{\rho-\lambda}}
 \quad\text{等價於}\quad
 \pair{q_{n-\rho-\lambda}S^\rho}{T_{n-\rho}p_{\rho-\lambda}}.
\]
整個過程中只使用 $S$ 而非 $T$ 的正規型性質。因此，生成 (69.) 時使用的缺陷 $\lambda-1$ 的同變摺疊，可由 $2\lambda-3$ 個基本摺疊和一個缺陷 $\lambda-2$ 的同變摺疊代替，後者又可繼續分解為基本摺疊，如此下去。類似地，若使用 $T$ 的正規型性質並應用對偶過程，則得到基本摺疊
\[
 S^\rho T^\rho,\quad T^{\rho+1},\ldots T^{\rho+\lambda-1};
 \qquad R^{\rho-1}T^{\rho-1},\quad T^{\rho-2},\ldots T^{\rho-\lambda+1},
\]
以及同 (69.) 對偶的摺疊。於是得到同前一折疊等價的摺疊
\[
\pair{q_{n-\rho-\lambda}T^\rho}{S_{n-\rho}p_{\rho-\lambda}}.
\]

把 2) 中所得分解加上去，確實得到最一般摺疊由基本摺疊生成。還需說明這種生成除了排列次序外是唯一的，即每個摺疊都對應於唯一的一組類型 $\rho$ 的基本摺疊個數 $\alpha_\rho$，其中 $\rho=1,2,\ldots n-1$。

設初始型 $f$ 在變量行 $p_\rho$ 中的次數為 $m_\rho$，摺疊所得型的次數為 $l_\rho$。類型 $\rho$ 的一個基本摺疊分別把次數 $m_{\rho-1},m_\rho,m_{\rho+1}$ 變為 $m_{\rho-1}+1,m_\rho-2,m_{\rho+1}+1$。因此 $\alpha$ 滿足方程組
\begin{equation}
 m_\rho-l_\rho=-\alpha_{\rho-1}+2\alpha_\rho-\alpha_{\rho+1},
 \qquad (\rho=1,2,\ldots n-1),
\tag{70}
\end{equation}
其中置 $\alpha_0$ 與 $\alpha_n=0$，其行列式的值為 $n$。\footnote{若 $\Delta_n$ 表示 $n$ 行行列式，則有遞推公式 $\Delta_n=2\Delta_{n-1}-\Delta_{n-2}$；也就是說，$\Delta_n-\Delta_{n-1}=\Delta_{n-1}-\Delta_{n-2}=\cdots=\Delta_2-\Delta_1=1$，因為 $\Delta_2=3$、$\Delta_1=2$。} 因而 $\alpha$ 被唯一確定，而且由定理 VIII，它們是正整數；這給出差 $m_\rho-l_\rho$ 的條件。

本段結果依賴於關係 (62)；然而若關係 (62) 不成立，終表達式並不等價於初始表達式。\S\ 6 中的等價性定義也可表述為：\footnote{參見下一段末尾的考察。} “兩個型等價，當且僅當它們相差一個較高摺疊的型，即顯示出更大權下降的型。” 這種較大的權下降總在進入摺疊的兩個行 $q_{n-\sigma-\lambda},p_{\tau-\lambda}$ 真正為變量行時發生，如 (41.)、(42.)、本段出現為等價的摺疊以及定理 VII 的等價型中那樣。但當這些行由符號行經代換導出時，不一定發生；本段中能憑 (62.) 相互表示的型就是這種情況。容易看到，消失的型甚至部分具有更高的總權。
%</unit>

%<unit id="noether.p04.sec.009" type="section" lang="zh-Hans" scan="scan:p04:169-174">
\section*{\S\ 9. 型行。化簡定理。}
現在把 \S\ 6 末尾引入的型行更精確地定義為：“給定初始型的、關於係數線性的不變構成式全體，也就是由初始型通過自身摺疊而得到的所有型，並按更高型排列。”

為說明更高型，按定理 VII 假設初始型是兩個正規型 $S\cdot T$ 的乘積。所謂更高型，是指具有更高自身摺疊的型；在自身摺疊相同的型之間，則指經過特別歸一化的型。更具體地說，設型 $A$ 由 $\alpha_\rho$ 個類型 $\rho$ 的基本摺疊產生，型 $B$ 由 $\beta_\rho$ 個這樣的摺疊產生。則 $B$ 高於 $A$：
\begin{enumerate}
\item 若在不等式組 $\beta_\rho\geqq\alpha_\rho$（$\rho=1,2,\ldots n-1$）中，至少對一個 $\rho$ 有嚴格不等號；\footnote{若在不等式組中一部分為 $>$、一部分為 $<$，則這些型不可比較，因為一個由另一個通過摺疊得到的型，總具有一個正的基本摺疊值系統，因而在這裡有正或為零的值 $\beta_\rho-\alpha_\rho$。}
\item 若所有 $\rho$ 都為等號，但合併在單個矩陣積中的符號行更少。這種歸一化因 \S\ 8 的化歸而必要。例如，$\pair{S^{n-1}}{T_{n-1}}$ 高於 $\pair{S^{n-1}}{[S^1T^\rho]_{n-\rho-1}p_\rho}$；
\item 若所有 $\rho$ 都為等號，且合併到一個矩陣積中的符號行數也相同，則通過一個本身任意但一次確定的歸一化把 $B$ 定為高於 $A$。這種歸一化總可行，因為屬於一個值系統 $\alpha$ 的型只有有限多個；例如可優先考慮一個型 $S$（符號行 $S$ 的數目較大、行 $S$ 的上權指標之和較大，等等）。
\end{enumerate}
把型行想成一個 $(n-1)$ 維格點流形很有用，其中 $n-1$ 個方向對應於 $n-1$ 個基本摺疊。於是每個型對應唯一的值系統 $\alpha$，即一個正整數格點；屬於同一值系統 $\alpha$ 的不同型，則由上述歸一化唯一確定其次序。型行的任意一個型又給出一個新型行的初始型，該新型行作為總型行的一部分包含在其中；事實上它正是由新的初始型沿 $n-1$ 個正方向繼續取得的那一部分。

由定理 VIII 和級數展開的一般理論，對型行成立的正是二元、三元領域中的化簡定理（參見博士論文 \S\ 3）。我們簡要推出主要定理。先給出定義：

“在一個給定型系統中，設有有限多個型被合併為一個模。若一個型能由含有不變量為因子的型表示，或由‘更高型’表示，即由該型所屬總型行的更高型表示，或由以更高階含有該模符號的型表示，則稱該型可化簡。一個可化簡的型行稱為化簡子。” 於是有：

\emph{定理 IX.} 若一個型行的初始型由於同一個化簡子摺疊而可化簡，則由該初始型產生的整個型行都可化簡。

證明時注意，由型 $f$ 與第二個型 $g$ 摺疊而產生的型 $h$，可因 (45.) 看作含有若干互伴變量行 $p_\rho,p'_\rho$ 的型 $f$。這類型可由級數展開的一般理論表示為 $f$ 的初等同變式的極化式，方法是對每一對互伴變量行 $p_\rho,p'_\rho$ 依次應用級數展開。由 (38.)，出現的初等同變式正是由互伴摺疊生成的型。不過，級數展開的應用次序仍是任意的；因而可特別選取一個與一般摺疊由基本摺疊生成相對應的次序。所得初等同變式系統就與型行 $f$ 相同；由較高缺陷的互伴摺疊產生的型，被排在由基本摺疊產生的型之中。但是對級數展開的任意次序也得到型行 $f$ 的極化式，並且可能對應另一個初等同變式系統。由於型行窮盡所有不變構成式，第二系統的每一個型都可由型行的極化式線性表示，並且事實上由顯示出同一或更大權下降的型的極化式表示。後一點是因為（參見 \S\ 7）這些極化式可看作一個型行 $H$ (52.)（其次數字對應於型 $h$）和型行 $f$ 的同時不變量；$H$ 的每個自身摺疊都對應一個權下降，在執行代換 (11.) 後，該下降轉移到符號行並因此轉移到 $f$ 的各型上。\footnote{這一考察也支撐等價概念的第二種表述（\S\ 8 末尾）。關於三元領域中不同初等同變式系統的進一步說明，可見 Study, loc. cit., II, \S\ 7, Theorems 7) and 8)。由我們的定理 VII，同樣結果也適用於 $n$ 個變量。}

型行 $h$ 的任意一個型由 $h$ 的自身摺疊而來；也就是說，一方面由 $g$ 與 $f$ 的更高摺疊而來，另一方面由 $f$ 或 $g$ 自身的摺疊而來。在兩種情形中，級數展開都導向型行 $f$ 的極化式，正如上面對初始型 $h$ 所證明的。若特別地 $f$ 是一個化簡子，則得到按該化簡子的極化式展開；因此型行 $h$ 可化簡。

關於完全型系統化簡的應用，與二元和三元領域完全類似。
%</unit>
%
\clearpage
%
% unit: noether.p05.title.001
\begin{center}
\pdfbookmark[0]{5. 有理函數域}{paper-05}
{\Large\bfseries 5. 有理函數域\footnote{1913 年在維也納 \emph{Naturforscherversammlung} 上作的報告。}}\par
\vspace{1.2em}
\emph{Jahresbericht der DMV} 22 (1913), pp. 316--319
\end{center}

\vspace{1.5em}

% unit: noether.p05.par.001
以下這些問題最初來自與 E. Fischer 的談話。此外，其中若干問題，對於一個不定元的特殊情形，並且在對係數域作更一般假設的條件下，已經由 E. Steinitz 提出並解決。\footnote{E. Steinitz: \emph{Algebraische Theorie der Körper}, 尤其是 \S\ 24. Crelle's Journal 137, 1910.}

% unit: noether.p05.par.002
我所謂“有理函數域”，是指其元素為 $n$ 個不定元的有理函數、且係數取自一個給定數域的域；這個數域特別也可以包含所有複數。這類域的例子包括 $n$ 個量的對稱函數域，或更一般地說，允許某個群的置換的 $n$ 個不定元的有理函數全體（Lagrange 的 \emph{Gattungsbereiche}）；此外還有不變量域。

% unit: noether.p05.par.003
前面最先提到的兩個域與最後一個域之間存在一個特徵性的差別：前者含有 $n$ 個代數無關函數，即初等對稱函數；而代數無關不變量的個數總是小於不定元的個數。因此，重要的是，一般地可以通過把若干不定元代以數，把這種第二種的域與第一種的域一一對應起來。所以下文我將限於第一種的域，也就是具有 $n$ 個代數無關函數的域；藉助這種對應，所得結果隨後便在一般情形下成立。

% unit: noether.p05.par.004
這些問題圍繞三個基概念而聚集：有理基、最小基和整基底。

% unit: noether.p05.par.005
1. 我所謂“有理基”，是指該域中的有限個函數，使得域中的每一個函數都能用給定係數域中的係數表示為這有限個函數的有理組合。通過簡單的考察可以證明，每一個有理函數域都有有理基；在一個不定元的情形下，這些考察也可見於 E. Steinitz 前引處。例如，按 Lagrange 定理，初等對稱函數以及一個屬於該群的函數，構成前面提到的 Lagrange \emph{Gattungsbereiche} 的一個有理基。

% unit: noether.p05.par.006
2. 每一個有理基（對於第一種的域而言）至少必須含有 $n$ 個函數；如果它恰好含有 $n$ 個函數，那麼這些函數必然代數無關，我把這樣的基稱作“最小基”。最小基存在性的問題，可以由 Lüroth、Castelnuovo 和 Enriques 的定理部分地回答。\footnote{J. Lüroth: \emph{Beweis eines Satzes über rationale Kurven}. Math. Ann. 9, 1875. --- G. Castelnuovo: \emph{Sulla razionalità delle involuzioni piane}. Math. Ann. 44, 1893. --- F. Enriques: \emph{Sopra una involuzione non razionale dello spazio}. Rend. Acc. Linc. vol. 21, 21 Jan. 1912.} 由這些定理可知，對於一個不定元的域，總存在一個最小基：即該域的 Lüroth 函數，它在分式線性變換之下唯一確定（參見 Steinitz \S\ 24）。對於兩個不定元的域，當且一般僅當允許係數域作代數擴張時，最小基總是存在；而對於三個及更多不定元，最小基一般根本不存在。對那些具有最小基的特殊域作特徵刻畫，目前尚未有人嘗試。

% unit: noether.p05.par.007
我進一步考察了 Lagrange \emph{Gattungsbereiche} 中最小基的問題。在這裡它具有如下意義：“如果屬於群 $G$ 的 Lagrange \emph{Gattungsbereich} 具有一個最小基，那麼就可以通過參數表示有理地構造出具有給定群 $G$ 的全部方程；並且這是相對於每一個包含該最小基的係數域的數域而言的。因此，特別地，如果這個係數域就是有理數域，那麼對於任意數域都可以這樣做。”

% unit: noether.p05.par.008
最簡單的例子由對稱群給出；這裡初等對稱函數構成一個具有有理數係數的最小基。由此得到的參數表示無非就是帶不定係數的方程。如何從這種方程過渡到每一個數域上的任意多的互不受影響的（\emph{affektlos}）方程，則由 Hilbert 的不可約性定理說明。

% unit: noether.p05.par.009
現在，從 Lüroth 和 Castelnuovo 的定理直接得到：對於三次和四次方程中出現的所有群，最小基都存在；同時還進一步表明，這個最小基具有有理數係數。因此，對於任意數域，都可以有理地構造出具有給定群的三次和四次方程全體。對於循環群和二面體群，存在一個以 $n$ 次單位根所在域為係數域的最小基；對於若干其他亞循環群也是如此。最小基是否根本就是某些群的特徵性對象，這個問題必須暫且懸而未決。

% unit: noether.p05.par.010
3. 為了達到最後一個基概念，我考察域中所含多項式的全體。我所謂“整基底”，是指這些多項式中的有限個，使得域中的每一個多項式都能用給定係數域中的係數表示為這有限個多項式的整有理組合。整基底存在性的問題，例如把不變量系的有限性作為特殊情形包含在內；Hilbert 在他的“數學問題”中\footnote{D. Hilbert: \emph{Mathematische Probleme}. 1900 年巴黎報告，第 14 問題。\emph{Göttinger Nachr.} 1900.}以稍有不同的形式，即作為相對整函數的問題提出了它。

% unit: noether.p05.par.011
目前我只能指出一類整基底存在的域。確切地說，如果該域含有一個由 $n$ 個多項式組成的系統，使得這些多項式最高維項的齊次結式不為零，那麼整基底存在；它由線性形式
% unit: noether.p05.eq.001
\[
  u_0=u_1x_1+u_2x_2+\cdots+u_nx_n.\footnotemark
\]
\footnotetext{在一個不定元的情形下，這個不可約方程見於 E. Steinitz 對 Lüroth 定理的證明。}
在該域中不可約的方程裡所有 $u$ 的係數給出。

% unit: noether.p05.par.012
特別地，如果結式的值等於係數域的一個單位，那麼表示的整性也得到保證。Lagrange 的 \emph{Gattungsbereiche} 給出了這一類域的一個例子，因為初等對稱函數的結式具有值 $\pm 1$。另一個例子（但沒有整性）由所有隻含有一個代數無關多項式的域給出；這裡整基底與包含所有多項式的最小子域的 Lüroth 函數相同。因此，眾所周知，$n$ 個變量的二次型的所有整有理射影不變量，都是判別式的整函數。

% unit: noether.p05.par.013
上述條件只是整基底存在的充分條件，絕不是必要條件。人們可以容易地構造這樣的域：其中任取 $n$ 個多項式的結式均為零，而這些域仍然具有由上述不可約方程的係數給出的整基底。於是問題產生了：是否即使在最一般的情形下，那個方程的係數也給出整基底？
%
\clearpage
%
% unit: noether.p06.title.001
\begin{center}
\pdfbookmark[0]{6. 有理函數的域與系統}{paper-06}
{\Large\bfseries 6. 有理函數的域與系統}\par
\vspace{0.6em}
Emmy Noether (Erlangen).\par
\vspace{1.2em}
\emph{Math. Ann.} 76 (1915), pp. 161--196
\end{center}

\vspace{1.5em}

% unit: noether.p06.par.001
本文研究任意的有理函數系統和整有理函數系統中的基問題；並且，由於所採用的域論方法，使得對於由有理函數構成的域——有理函數域\footnote{參看在維也納自然科學家會議時所作的預告通信：\emph{Rationale Funktionenkörper}，\emph{Jahresber. d. D. Math.-Ver.} 22 (1913)，其中給出了問題以及有關函數域的結果的概覽。}——解決這些問題顯得是本質性的，而把結果推廣到任意系統則作為其後果而得到。

% unit: noether.p06.par.002
關於一般系統的基問題，到目前為止只知道由 Hilbert 定理（\emph{Math. Ann.} 36）所保證的模基存在性。下文中，有理可表示性問題由每一個任意系統的有理基的存在性完全解決（\S\ 7）；域的有理基已在 \S\ 4 中得到。有理基的這種存在性使我們處處可以從抽象定義的域或系統出發，從而避免那些只由有理基的特殊選擇、而不是由系統本身造成的困難，例如特殊分母或基函數的基本點的出現。

% unit: noether.p06.par.003
對於域，還處理最小基的問題，即由代數無關函數構成的有理基的問題（\S\ 6）。域論方法還導向一個特殊的有理基，即對合基\footnote{這個名稱的選擇，是因為在幾何解釋中，有理函數域代表線性空間中最一般的對合。}（\S\ 5）；它對於由多項式構成的整性域（\S\ 8）變得特別重要。在那裡它給出一個具有固定分母的表示（這個分母是由整性域所確定的某個函數的冪），例如在不變量的典型表示這個特殊情形中已知的那樣。

% unit: noether.p06.par.004
現在，從有理可表示性出發，要儘可能廣泛地推出關於整有理可表示性的後果，也就是關於狹義有限性問題的後果。

% unit: noether.p06.par.005
迄今為止已知的有限性定理都依賴於這樣一個事實：Hilbert 的模基定理對所有這樣的系統保證有限性，即在其中一個表示
% unit: noether.p06.eq.001
\[
  F=A_1f_1+A_2f_2+\cdots+A_kf_k
\]
會導致第二個表示
% unit: noether.p06.eq.002
\[
  F=B_1f_1+B_2f_2+\cdots+B_kf_k,
\]
其中 $B_i$ 屬於該系統，並且次數低於 $F$。

% unit: noether.p06.par.006
與此相反，有理基定理和域論方法允許在性質完全不同的假設下推出有限性。

% unit: noether.p06.par.007
因此，作為對 Hilbert 一個問題（“數學問題”第 14 問題）的部分回答，可以得到一類相對整函數（第一類相對整域），它們是有限整性域，並且其整基底由上述對合基給出（\S\ 10）。與 Hilbert 對 Kronecker 關於代數整數基本系統的定理的證明（“完備不變量系”，\S\ 2；\emph{Math. Ann.} 42）相類似，還可以證明所有正則多項式系統——即所有能夠映到沒有基本點的系統上的系統——都具有整基底（\S\ 12）；而沒有整基底的非正則系統則很容易給出。作為這個事實的一個簡單例子，可以提到如下定理：“在一個不定元的任意無限多項式系統中，存在其中有限多個多項式，使得該系統中的每個多項式都能表示為這有限多個多項式的整有理組合。”進一步的例子見於 \S\ 13。

% unit: noether.p06.par.008
最後，末尾幾節（\S\S\ 14 和 15）給出關於整性的基定理的加強。

% unit: noether.p06.par.009
研究中的一個本質工具是 \S\ 3 的傳遞原理。按照該原理，對於這裡提出的所有基問題，只需考慮那些系統，其中不定元個數與該系統中代數無關函數的個數相一致。

% unit: noether.p06.par.010
本文工作的推動來自與 E. Fischer 先生的談話，特別是來自他關於 Lagrange 的 \emph{Gattungsbereiche} 的最小基的問題。與此相關的研究——它們導向具有預定群的方程的構造（參見所引關於“有理函數域”的通信第二部分）——將留待另一篇發表。

% unit: noether.p06.par.011
本文相對於原先通信的本質擴展在於：那裡只對域或特殊整性域陳述的基定理，在這裡被轉移到任意系統上。這個轉移藉助於任意系統的最小包含域這一概念以簡單方式實現；這個概念是整性域的商域的推廣（\S\ 7）。促使我形成這一概念的是這樣一個事實：K. Hentzelt 先生在得知域的有理基以後，能夠繞道 Hilbert 的模基來證明任意系統有理基的存在性。我關於正則系統整基底的定理，也來自 K. Hentzelt 的觀察：我用於整性域的論證，在同樣假設下對任意系統也有效；此外，我還應感謝 Hentzelt 先生的一些零星的小評論。

% unit: noether.p06.par.012
最後還應提到，在一個不定元的情形下，域的有理基和最小基見於 E. Steinitz 的《\emph{Algebraische Theorie der Körper}》\S\ 24（Crelle's Journal 137）；在那裡係數域取為最一般意義下的域。這裡給出的定理在這些更一般的假設下能在多大程度上保持成立，這個問題下文不予討論；無論如何，證明需要修改，因為例如函數行列式理論的工具對於特徵 $p$ 的域失效。

% unit: noether.p06.sec.001
\section*{\S\ 1. 域 $\Kfield_{n\rho}$ 與系統 $\Ssys_{n\rho}$。}

% unit: noether.p06.s01.par.001
下文中我們研究 $n$ 個不定元的有理函數系統，特別是有理函數域。

% unit: noether.p06.s01.par.002
因此，$f(x_1\cdots x_n)$、$g(x_1\cdots x_n)$、……，或簡寫為 $f(x)$、$g(x)$、……，總是指 $x_1\cdots x_n$ 的有理函數；並且假定它們處於既約表示中，即分子與分母互素。整有理函數（多項式）用大寫字母表示，$F(x)$、$G(x)$、……。係數域 $\Omega$ 假定為某個數域；它尤其也可以包含所有複數。域 $\Omega$ 也可以包含有限多個參數。

% unit: noether.p06.s01.par.003
來自 $\Omega$ 的量，即所有零次函數，總是算作屬於該系統。

% unit: noether.p06.s01.par.004
在這些約定下，由有理函數組成的域——有理函數域——可以抽象地定義。

% unit: noether.p06.def.001
\emph{定義 I：一個有理函數系統若滿足下列條件，則稱為一個域：}
\begin{enumerate}
\item \emph{連同 $f(x)$，並且對 $\Omega$ 中任意量 $c$，它也含有 $c\cdot f(x)$。}
\item \emph{連同 $f(x)$ 和 $g(x)$，它總也含有 $f(x)+g(x)$、$f(x)\cdot g(x)$，並且在 $g(x)\ne0$ 時含有商 $f(x):g(x)$。}\footnote{這裡 $f(x)$ 與 $g(x)$ 也可以是零次的，即 $\Omega$ 中的量。}
\end{enumerate}

% unit: noether.p06.s01.par.005
特殊的域是那些通過向 $\Omega$ 添加有限多個有理函數
% unit: noether.p06.eq.003
\[
  f_1(x),\ f_2(x),\ \cdots,\ f_k(x)
\]
而產生的域；它們記作
% unit: noether.p06.eq.004
\[
  \Omega\bigl(f_1(x),\ldots,f_k(x)\bigr)
  \quad\text{或}\quad
  \Omega(f_1,\ldots,f_k).
\]
\footnote{這些“特殊域”已經窮盡所有域的總體，這一點在 \S\ 4 中得出。}
在這些特殊域中，我們還提到由向 $\Omega$ 添加 $x_1\cdots x_n$ 所產生的域
% unit: noether.p06.eq.005
\[
  \Omega(x_1,\ldots,x_n),
\]
它等同於所有以 $\Omega$ 中元素為係數的 $x_1\cdots x_n$ 的有理函數全體。

% unit: noether.p06.s01.par.006
我們還（沿用 E. Steinitz）引入中間域的概念：
\begin{quote}
“如果一個域 $\Mfield$ 含有 $\mathfrak{L}_1$ 且包含於 $\mathfrak{L}_2$，則稱它位於 $\mathfrak{L}_1$ 與 $\mathfrak{L}_2$ 之間；”
\end{quote}
於是定義 I 也容許如下表述：

% unit: noether.p06.s01.par.007
$n$ 個不定元的有理函數域，是 $\Omega$ 與 $\Omega(x_1\cdots x_n)$ 之間的一個中間域，並且它在至少一個函數中實際含有這 $n$ 個不定元。

% unit: noether.p06.s01.par.008
除了不定元的個數外，對於有理函數系統還出現另一個特徵數，即代數無關函數的個數，或代數秩，由如下定義給出。

% unit: noether.p06.def.002
\emph{定義 II：如果一個有理函數系統中可以指定 $\rho$ 個代數無關函數，而該系統中任意 $(\rho+1)$ 個函數都是代數相關的，則稱該系統具有代數秩 $\rho$。}\footnote{稱 $\rho$ 個函數 $f_1(x),\ldots,f_\rho(x)$ 代數無關，是指從任何關係
\[
F\bigl(f_1(x),\ldots,f_\rho(x)\bigr)=0\quad\text{在 }x_1,\ldots,x_n\text{ 中恆等}
\]
都推出
\[
F(\lambda_1,\ldots,\lambda_\rho)=0\quad\text{在 }\lambda_1,\ldots,\lambda_\rho\text{ 中恆等}.
\]
反之，則稱它們代數相關。}

% unit: noether.p06.s01.par.009
$n$ 個不定元、代數秩為 $\rho$ 的有理函數系統（有限個或無限多個）用雙重下標表示為
% unit: noether.p06.eq.006
\[
  \Ssys_{n\rho}.
\]
特別地，記號
% unit: noether.p06.eq.007
\[
  \Kfield_{n\rho}
\]
表示 $n$ 個不定元、代數秩為 $\rho$ 的域。有不等式
% unit: noether.p06.eq.008
\[
  1\leq \rho\leq n.
\]

% unit: noether.p06.s01.par.010
我們再給出一些域 $\Kfield_{nn}$ 和 $\Kfield_{n\rho}$ 的例子：

% unit: noether.p06.s01.par.011
1. 域 $\Kfield_{nn}$ 包括 Lagrange 的 \emph{Gattungsbereiche}，它們可以定義為
\begin{quote}
“$x_1,\ldots,x_n$ 的所有有理（以及整有理）函數的總體，這些函數允許 $x_1,\ldots,x_n$ 的某個置換群 $G$ 的置換。”
\end{quote}
它們總是包含對稱函數域，因而也包含 $n$ 個代數無關的初等對稱函數。

% unit: noether.p06.s01.par.012
2. 域 $\Kfield_{n\rho}$ 包括射影不變量域，它們由一個或多個基本型的所有有理（以及整有理）不變量構成。\footnote{只考慮行列式為 $1$ 的變換是適宜的；這樣，任意多個不變量的任意有理組合仍然允許該變換，並且確實得到一個域。單獨取出的齊次、等重不變量不構成域，但它們構成一個系統 $\Ssys_{n\rho}$。} 這裡 $\rho<n$，因為代數無關不變量的個數總是小於係數的個數。

% unit: noether.p06.s01.par.013
對於射影群的子群所對應的不變量域，也有相應的說法。

% unit: noether.p06.sec.002
\section*{\S\ 2. 關於代數相關函數的引理。}

% unit: noether.p06.s02.par.001
本節的引理將在 \S\ 3 中說明：系統的代數秩就是它真正的特徵數。確切地說，該引理說明，通過依次把若干不定元專門化為數，使得 $\rho$ 個代數無關函數仍保持代數無關，則任意一個代數依賴於這 $\rho$ 個函數的函數，既不能恆等為零，也不能變為不定。

% unit: noether.p06.s02.par.002
因此，設系統
% unit: noether.p06.eq.009
\begin{equation}
  g_1(x),\ g_2(x),\ldots,g_\rho(x)
\tag{1}
\end{equation}
具有代數秩 $\rho$，其中 $\rho<n$，並設 $h(x)$ 是一個代數依賴於函數 (1) 的有理函數。由已知定理，(1) 的函數矩陣的秩等於 $\rho$。因此，例如可以給不定元重新命名，使得
% unit: noether.p06.eq.010
\begin{equation}
 \frac{\partial(g_1,g_2,\ldots,g_\rho)}{\partial(x_1,x_2,\ldots,x_\rho)}
 =\frac{\Gamma(x)}{G(x)^{\rho+1}}\ne0,
\tag{2}
\end{equation}
其中 $G(x)$ 照例表示函數 (1) 的公分母。

% unit: noether.p06.s02.par.003
\emph{現在選擇數 $a$，使得乘積 $G(x)\cdot\Gamma(x)$ 不被 $(x_i-a)$ 整除}，這裡 $i$ 被理解為 $\rho+1,\rho+2,\ldots,n$ 中的一個數。這樣，當 $x_i=a$ 時，函數 (1) 的公分母不可能消失，並且 $\rho$ 個函數
% unit: noether.p06.eq.011
\begin{equation}
  [g_1(x)]_{x_i=a},\ldots,[g_\rho(x)]_{x_i=a}
\tag{3}
\end{equation}
同樣代數無關。按我們的說法，系統 (1) 的代數秩在代入 $x_i=a$ 後保持不變。

% unit: noether.p06.s02.par.004
設 $h(x)$ 關於函數 (1) 所滿足的不可約方程為
% unit: noether.p06.eq.012
\begin{equation}
\begin{aligned}
&A_0(g_1,\ldots,g_\rho)\cdot\,h(x)^\tau
 +A_1(g_1,\ldots,g_\rho)\cdot\,h(x)^{\tau-1}
 +\cdots \\
&\hspace{6em}+A_\tau(g_1,\ldots,g_\rho)=0,
\end{aligned}
\tag{4}
\end{equation}
其中 $A_0(g_1,\ldots,g_\rho)$ 與 $A_\tau(g_1,\ldots,g_\rho)$ 不恆等為零。為了由此過渡到多項式之間的恆等式，我們置
% unit: noether.p06.eq.013
\begin{equation}
  g_i(x)=G_i(x):G(x);\qquad h(x)=H_1(x):H_2(x),
\tag{5}
\end{equation}
並注意到——由於假定 $h(x)$ 為既約表示——$H_1(x)$ 與 $H_2(x)$ 互素。由 (5)，(4) 變為
% unit: noether.p06.eq.014
\begin{equation}
\begin{aligned}
&B_0(G,G_1,\ldots,G_\rho)\cdot\,G(x)^{\sigma_0\cdot}H_1(x)^\tau \\
&\quad+B_1(G,G_1,\ldots,G_\rho)\cdot\,G(x)^{\sigma_1\cdot}H_1(x)^{\tau-1}H_2(x)+\cdots \\
&\quad+B_\tau(G,G_1,\ldots,G_\rho)\cdot\,G(x)^{\sigma_\tau\cdot}H_2(x)^\tau=0,
\end{aligned}
\tag{6}
\end{equation}
其中 $B_k$ 是 $G,G_i$ 中的齊次形式，並且非負指數 $\sigma_k$ 中至少有一個必須為零。

% unit: noether.p06.s02.par.005
現在假設 $H_1(x)$ 可被 $(x_i-a)$ 整除。那麼由 (6)，
% unit: noether.p06.eq.015
\[
  B_\tau(G\,G_1\cdots G_\rho)\cdot G(x)^{\sigma_\tau}\cdot H_2(x)^\tau
\]
也會被 $(x_i-a)$ 整除。按假設，這對於 $G(x)$ 是排除的；於是由 $B_\tau$ 的可整除性還會推出
% unit: noether.p06.eq.016
\[
  A_\tau=B_\tau:G^\lambda
\]
的可整除性，並且函數 (3) 之間就會存在關係 $A_\tau=0$，這與所假定的代數無關性矛盾。最後，由於 $H_2(x)$ 與 $H_1(x)$ 互素，它也不能被 $(x_i-a)$ 整除；\footnote{這個最後結論依賴於既約表示；若作代入 $x_i=a$, $x_k=b$ 而其餘假設保持，則不再可能。那時 $H_1(x)$ 與 $H_2(x)$ 可能同時消失，商在代入下會變成不定的 $=\frac00$；例如
\[
 h(x)=\frac{x_3(\alpha x_1+\beta x_2)+x_4(\gamma x_1+\delta x_2)}
 {x_3(\alpha' x_1+\beta' x_2)+x_4(\gamma' x_1+\delta' x_2)}
\]
在代入 $x_3=0$, $x_4=0$ 時。} 因而 $H_1(x)$ 可被 $(x_i-a)$ 整除的假設導致矛盾。以同樣方式可證明 $H_2(x)$ 也不能被 $(x_i-a)$ 整除。因此函數 $[h(x)]_{x_i=a}=k(x)$ 既不能恆等為零，也不能變為不定。

% unit: noether.p06.s02.par.006
上述結論現在可以再次應用於系統 (3) 和函數 $k(x)$，後者也再次假定為既約表示；繼續這個過程，最終得到如下引理。

% unit: noether.p06.lemma.001
\emph{引理：兩個系統}
% unit: noether.p06.eq.017
\[
  g_1(x),\ g_2(x),\ldots,g_\rho(x),
\]
% unit: noether.p06.eq.018
\[
  g_1(x),\ g_2(x),\ldots,g_\rho(x),\ h(x)
\]
\emph{各自具有代數秩 $\rho$，其中 $\rho<n$，並且例如設}
% unit: noether.p06.eq.019
\[
 \frac{\partial(g_1,\ldots,g_\rho)}{\partial(x_1,\ldots,x_\rho)}
 =\frac{\Gamma(x)}{G(x)^{\rho+1}}\ne0.
\]
\emph{數}
% unit: noether.p06.eq.020
\[
  a_n,\ a_{n-1},\ldots,a_{\rho+1}
\]
\emph{應當選得使得在代入}
% unit: noether.p06.eq.021
\[
  x_n=a_n,\quad x_{n-1}=a_{n-1},\ldots,\quad x_{\rho+1}=a_{\rho+1},
\]
\emph{時，乘積 $G(x)\cdot\Gamma(x)$ 不消失，也就是系統 (1) 的代數秩被保持。在連續代入}
% unit: noether.p06.eq.022
\[
\begin{aligned}
 [h(x)]_{x_n=a_n}&=h^{(n-1)}(x),\qquad
 [h^{(n-1)}(x)]_{x_{n-1}=a_{n-1}}=h^{(n-2)}(x),\ldots,\\
 [h^{(\rho+1)}(x)]_{x_{\rho+1}=a_{\rho+1}}&=h^*(x_1,x_2,\ldots,x_\rho),
\end{aligned}
\]
\emph{中，函數 $h(x)$、$h^{(n-1)}(x)$、……、$h^{(\rho+1)}(x)$ 每次都置為既約表示。在這些假設下，在函數 $h^*(x_1,\ldots,x_\rho)$ 中，分子和分母都不能恆等為零，並且該函數也不能變為 $\frac{0}{0}$。}\footnote{例如，在前注中的函數 $h(x)$ 中，連續代入給出
\[
 [h(x)]_{x_4=0}=h^{(3)}(x)=\frac{(\alpha x_1+\beta x_2)}{(\alpha' x_1+\beta' x_2)}.\qquad
 h^*(x_1 x_2)=\frac{\alpha x_1+\beta x_2}{\alpha' x_1+\beta' x_2}.
\]}

% unit: noether.p06.sec.003
\section*{\S\ 3. 系統 $\Ssys_{n\rho}$ 到系統 $\Ssys^*_{\rho\rho}$ 的一一映射。}

% unit: noether.p06.s03.par.001
由引理，我們將立即得到把系統 $\Ssys_{n\rho}$ 映到系統 $\Ssys^*_{\rho\rho}$ 的映射，即通過把若干不定元替換為數；由此，代數秩 $\rho$ 顯示為真正的特徵數。

% unit: noether.p06.s03.par.002
由於 $\Ssys_{n\rho}$ 具有代數秩 $\rho$，所以存在 $\Ssys_{n\rho}$ 中的 $\rho$ 個函數
% unit: noether.p06.eq.023
\begin{equation}
  g_1(x),\ldots,g_\rho(x),
\tag{1}
\end{equation}
它們代數無關。此外，如果 $f(x)$ 表示 $\Ssys_{n\rho}$ 中任意函數，則系統
% unit: noether.p06.eq.024
\begin{equation}
  g_1(x),\ldots,g_\rho(x),\ f(x)
\tag{2}
\end{equation}
也具有代數秩 $\rho$，並且系統 (1) 與 (2) 滿足引理的條件。\footnote{如有必要，不定元應重新編號。}

% unit: noether.p06.s03.par.003
因此，若按照引理確定數
% unit: noether.p06.eq.025
\[
  a_n,\ a_{n-1},\ldots,a_{\rho+1}
\]
使得系統 (1) 的代數秩 $\rho$ 在代入下保持不變——這可以用任意多種方式做到——則由
% unit: noether.p06.eq.026
\[
\begin{aligned}
 [f(x)]_{(x_n=a_n)}&=f^{(n-1)}(x),\qquad
 [f^{(n-1)}(x)]_{(x_{n-1}=a_{n-1})}=f^{(n-2)}(x),\\
 [f^{(\rho+1)}(x)]_{(x_{\rho+1}=a_{\rho+1})}&=f^*(x_1,\ldots,x_\rho)
\end{aligned}
\]
所定義的函數，即
% unit: noether.p06.eq.027
\begin{equation}
  f^*(x_1,\ldots,x_\rho),
\tag{3}
\end{equation}
其分子與分母都不能恆等為零。

% unit: noether.p06.s03.par.004
現在讓 $f(x)$ 遍歷 $\Ssys_{n\rho}$ 中所有（有限或無限多個）函數。我們考慮由所有函數 (3) 的總體組成的系統；它具有以下性質：

% unit: noether.p06.s03.par.005
1. 它具有代數秩 $\rho$；因為它包含由 (1) 得到的函數
% unit: noether.p06.eq.028
\[
  g_1^*(x_1\cdots x_\rho)\cdots g_\rho^*(x_1\cdots x_\rho),
\]
這些函數由於 $a_n,a_{n-1},\ldots,a_{\rho+1}$ 的選擇而代數無關。因此該系統記為
% unit: noether.p06.eq.029
\[
  \Ssys^*_{\rho\rho}.
\]
\footnote{相應地，$f^*(x)$、$g^*(x)$、……應始終理解為映射函數，即 $x_1\cdots x_\rho$ 的函數。}

% unit: noether.p06.s03.par.006
2. 函數 $f(x)$ 與 $f^*(x)$ 的對應毫無例外地是一一對應。

% unit: noether.p06.s03.par.007
事實上，如果 $f_1(x)$ 和 $f_2(x)$ 屬於 $\Ssys_{n\rho}$，則它們的差
% unit: noether.p06.eq.030
\[
  d(x)=f_1(x)-f_2(x)
\]
也代數依賴於系統 (1)；並且
% unit: noether.p06.eq.031
\[
  d^*(x)=f_1^*(x)-f_2^*(x).
\]
由於 $d(x)$ 連同系統 (1) 滿足引理的條件，由
% unit: noether.p06.eq.032
\[
  d(x)\ne0
\]
必然推出
% unit: noether.p06.eq.033
\[
  d^*(x)\ne0;
\]
換言之，$\Ssys_{n\rho}$ 中不同的函數對應於 $\Ssys^*_{\rho\rho}$ 中不同的函數。

% unit: noether.p06.s03.par.008
3. 從 $\Ssys^*_{\rho\rho}$ 中函數之間的每一個關係
% unit: noether.p06.eq.034
\[
  F\bigl(f_1^*(x),f_2^*(x),\ldots,f_k^*(x)\bigr)=0
  \quad\text{在 }x_1,\ldots,x_\rho\text{ 中恆等},
\]
對於 $\Ssys_{n\rho}$ 中唯一對應的函數，推出
% unit: noether.p06.eq.035
\[
  F\bigl(f_1(x),f_2(x),\ldots,f_k(x)\bigr)=0
  \quad\text{在 }x_1,\ldots,x_n\text{ 中恆等}.
\]
因為連同 $f_1(x),f_2(x),\ldots,f_k(x)$，這些函數的每一個有理組合也代數依賴於系統 (1)。

% unit: noether.p06.s03.par.009
這還由
% unit: noether.p06.eq.036
\begin{equation}
\begin{aligned}
 h(x)&=F\bigl(f_1(x),\ldots,f_k(x)\bigr):\\
 h^*(x)&=F\bigl(f_1^*(x),\ldots,f_k^*(x)\bigr)
\end{aligned}
\tag{4}
\end{equation}
推出。因此，由引理，
% unit: noether.p06.eq.037
\[
  h(x)\ne0\quad\text{蘊涵}\quad h^*(x)\ne0;
\]
證畢。

% unit: noether.p06.thm.001
綜上，我們得到：

\emph{定理 I：對每一個由（有限或無限多個）有理函數組成的系統 $\Ssys_{n\rho}$，可以用任意多種方式配以一個一一對應的系統 $\Ssys^*_{\rho\rho}$，使得 $\Ssys^*_{\rho\rho}$ 的函數之間存在的所有關係也在 $\Ssys_{n\rho}$ 中保持，反過來也如此。特別地，由 (4)，一個域 $\Kfield_{n\rho}$ 又對應於一個域 $\Kfield^*_{\rho\rho}$。}

% unit: noether.p06.s03.par.010
由定理 I 我們得出一個簡化所有後續證明的結論：\emph{如果系統 $\Ssys_{n\rho}$ 的某些性質由系統函數之間的有限或無限多個關係來刻畫，那麼只要這些性質對一個映射系統 $\Ssys^*_{\rho\rho}$ 成立，它們也對系統 $\Ssys_{n\rho}$ 成立。}

% unit: noether.p06.s03.par.011
我們簡稱這個結論為“傳遞原理”。

% unit: noether.p06.s03.par.012
這個映射最簡單的例子由代數不變量理論提供。事實上，通過線性變換，總可以把基本型的某些係數數值地固定，使得對特殊型成立的所有關係在一般情形下也仍然成立。

% unit: noether.p06.sec.004
\section*{\S\ 4. 域 $\Kfield_{n\rho}$ 的有理基。}

% unit: noether.p06.s04.par.001
在下面圍繞基概念展開的研究中，我們將始終使用 \S\ 3 的“傳遞原理”。我們首先證明域的基的存在性，但把定義作一般表述：

% unit: noether.p06.def.003
\emph{定義 III：$\Ssys_{n\rho}$ 中有限多個函數稱為有理基，如果 $\Ssys_{n\rho}$ 中每個函數都能用 $\Omega$ 中係數表示為這有限多個函數的有理組合。}

% unit: noether.p06.s04.par.002
有理基必須含有 $\rho$ 個代數無關函數，因為從有理基導出的系統的代數秩等於有理基的代數秩。

% unit: noether.p06.s04.par.003
我們首先說明，有理基的存在性證明容許傳遞原理。確切地說，定義 III 可以表述如下：

% unit: noether.p06.s04.par.004
$\Ssys_{n\rho}$ 中的函數
% unit: noether.p06.eq.038
\[
  g_1(x),\ g_2(x),\ldots, g_k(x)
\]
構成一個有理基，當且僅當滿足無限多個關係
% unit: noether.p06.eq.039
\begin{equation}
  f(x)=\gamma\bigl(g_1(x),\ldots,g_k(x)\bigr),
\tag{1}
\end{equation}
其中 $f(x)$ 遍歷 $\Ssys_{n\rho}$ 中所有函數，$\gamma$ 表示 $k$ 個變元的有理函數，並隨 $f$ 而改變。

% unit: noether.p06.s04.par.005
於是，如果在映射系統 $\Ssys^*_{\rho\rho}$ 中，一個有理基由
% unit: noether.p06.eq.040
\[
  g_1^*(x),\ g_2^*(x),\ldots,g_k^*(x)
\]
給出，並因而有關係
% unit: noether.p06.eq.041
\[
  f^*(x)=\gamma\bigl(g_1^*(x)\cdots g_k^*(x)\bigr),
\]
那麼由定理 I，關係 (1) 對 $\Ssys_{n\rho}$ 成立。因此，\emph{關於有理基的傳遞原理}說：

% unit: noether.p06.s04.par.006
\emph{由映射系統 $\Ssys^*_{\rho\rho}$ 的有理基的存在性，推出原系統 $\Ssys_{n\rho}$ 的有理基。}\footnote{這裡處理的每一個基問題，都相應地表述傳遞原理。}

% unit: noether.p06.s04.par.007
為了證明所有域 $\Kfield_{n\rho}$ 的有理基存在，我們因此可以限於域 $\Kfield^*_{\rho\rho}$；這裡，對 E. Steinitz 所用論證的一個直接推廣即可達到目的。\footnote{E. Steinitz, \emph{Algebraische Theorie der Körper}, \S\ 24, Crelle's Journal 137 (1909).}

% unit: noether.p06.s04.par.008
設
% unit: noether.p06.eq.042
\begin{equation}
  g_1^*(x_1,\ldots,x_\rho),\ldots,
  g_\rho^*(x_1,\ldots,x_\rho)
\tag{2}
\end{equation}
是 $\Kfield^*_{\rho\rho}$ 中 $\rho$ 個代數無關函數組成的系統。通過消去 $x_1,\ldots,x_\rho$，得到 $\rho$ 個關係：
% unit: noether.p06.eq.043
\begin{equation}
\begin{aligned}
  F_1\bigl(g_1^*(x),\ldots,g_\rho^*(x),x_1\bigr)&=0,\ldots,\\
  F_\rho\bigl(g_1^*(x),\ldots,g_\rho^*(x),x_\rho\bigr)&=0,
\end{aligned}
\tag{3}
\end{equation}
其中在每個關係 $F_i=0$ 中，不定元 $x_i$ 都實際出現；否則，與假設相反，函數 (2) 之間會存在代數依賴。關係 (3) 說明
% unit: noether.p06.eq.044
\[
  x_1,x_2,\ldots,x_\rho
\]
相對於
% unit: noether.p06.eq.045
\[
  \Omega\bigl(g_1^*(x),\ldots,g_\rho^*(x)\bigr)
\]
是代數元素。因此，通過添加 $x_1,\ldots,x_\rho$ 而產生的域
% unit: noether.p06.eq.046
\[
  \Omega\bigl(g_1^*,\ldots,g_\rho^*,x_1,
  \ldots,x_\rho\bigr)=\Omega(x_1,
  \ldots,x_\rho),
\]
是 $\Omega(g_1^*,\ldots,g_\rho^*)$ 的有限代數擴張。由一個已知定理，\footnote{例如比較 Weber, \emph{Lehrbuch d. Algebra} 1, \S\ 144（第一版）或 \S\ 55（小版）。} $\Kfield^*_{\rho\rho}$ 作為 $\Omega(g_1^*\cdots g_\rho^*)$ 與 $\Omega(x_1\cdots x_\rho)$ 之間的中間域，本身是 $\Omega(g_1^*\cdots g_\rho^*)$ 的代數擴張，而 $\Omega(x_1\cdots x_\rho)$ 是 $\Kfield^*_{\rho\rho}$ 的代數擴張。因此，$\Kfield^*_{\rho\rho}$ 由 $\Omega(g_1^*\cdots g_\rho^*)$ 通過添加 $\Kfield^*_{\rho\rho}$ 中有限多個元素得到，或者也可以通過添加一個適當選擇的函數得到；也就是說，有
% unit: noether.p06.eq.047
\[
  \Kfield^*_{\rho\rho}
  =\Omega(g_1^*\cdots g_\rho^*g_{\rho+1}^*\cdots g_k^*)
  =\Omega(g_1^*\cdots g_\rho^*g_0^*).
\]

% unit: noether.p06.thm.002
由傳遞原理，我們因此有：

\emph{定理 II：每一個域 $\Kfield_{n\rho}$ 都具有一個代數秩為 $\rho$ 的有理基；特別地，有理基可以選得使它至多含有 $\rho+1$ 個函數。}

% unit: noether.p06.s04.par.009
設 $\Kfield_{n\rho}$ 中任意 $\rho$ 個代數無關函數組成的系統為 $h_1(x),\ldots,h_\rho(x)$，並假設它能由上述方法擴充為有理基
% unit: noether.p06.eq.048
\[
  h_1(x),\ldots,h_\rho(x);
  h_{\rho+1}(x),\ldots,h_\sigma(x).
\]
由定義性質，存在關係
% unit: noether.p06.eq.049
\[
\begin{aligned}
  h_i(x)&=\chi_i\bigl(g_1(x),\ldots,g_k(x)\bigr)
    &&(i=1,2,\ldots,\sigma),\\
  g_j(x)&=\psi_j\bigl(h_1(x),\ldots,h_\sigma(x)\bigr)
    &&(j=1,2,\ldots,k),
\end{aligned}
\]
其中 $\chi_i$ 和 $\psi_j$ 表示其變元的有理函數。因此：

% unit: noether.p06.s04.par.010
從任意一個有理基，通過雙有理變換可得到任何另一個有理基；並且：

% unit: noether.p06.s04.par.011
如果
% unit: noether.p06.eq.050
\[
  h_1(x),\ldots,h_\rho(x)
\]
是 $\Kfield_{n\rho}$ 中任意 $\rho$ 個代數無關函數組成的系統，那麼 $\Kfield_{n\rho}$ 是
% unit: noether.p06.eq.051
\[
  \Omega\bigl(h_1(x),\ldots,h_\rho(x)\bigr)
\]
上的有限代數域。

% unit: noether.p06.sec.005
\section*{\S\ 5. 域 $\Kfield_{n\rho}$ 的對合形式與對合基。}

% unit: noether.p06.s05.par.001
在 \S\ 4 中構造的每個有理基都有一個缺點，即依賴任意選擇的起始函數；現在我們證明存在一個\emph{由域唯一確定的有理基}。

% unit: noether.p06.s05.par.002
由 \S\ 4，$\Omega(x_1,\ldots,x_\rho)$ 是 $\Kfield^*_{\rho\rho}$ 的代數擴張，設其次數為 $i$；並且由於
% unit: noether.p06.eq.052
\[
  \Kfield^*_{\rho\rho}(x_1,\ldots,x_\rho)=\Omega(x_1,\ldots,x_\rho),
\]
它由添加元素 $x_1,\ldots,x_\rho$ 而產生，這些元素相對於 $\Kfield^*_{\rho\rho}$ 是代數的。因此我們考慮相對於 $\Kfield^*_{\rho\rho}$ 不可約的整函數 $\Phi^*(z,u)$，其零點為
% unit: noether.p06.eq.053
\[
  z=u_1x_1+u_2x_2+\cdots+u_\rho x_\rho=u(x),
\]
並且它關於 $z$ 是 $i$ 次的。此外，作為與 $[z-u(x)]$ 共軛的量的乘積，$\Phi^*(z,u)$ 是 $z,u_1,\ldots,u_\rho$ 中維數為 $i$ 的齊次形式；因此，作為 $z^i$ 的係數，它含有一個不含 $u$ 的 $\Kfield^*_{\rho\rho}$ 中的函數。如果我們通過除法使這個最高係數等於 $1$，則不可約函數 $\Phi^*(z,u)$ 被唯一確定。這樣定義的齊次形式
% unit: noether.p06.eq.054
\begin{equation}
\begin{aligned}
  \Phi^*(z,u)
  &=z^i+
    \sum{}' g^*_{\alpha\alpha_1\ldots\alpha_\rho}(x_1,
    \ldots,x_\rho)\cdot
    z^\alpha u_1^{\alpha_1}\cdots u_\rho^{\alpha_\rho},\\
  &\hspace{8em}(\alpha+\alpha_1+\cdots+\alpha_\rho=i),
\end{aligned}
\tag{1}
\end{equation}
稱為 $\Kfield^*_{\rho\rho}$ 的\emph{對合形式}；\footnote{這個名稱的理由將在本節末尾變得清楚；$\sum'$ 表示求和遍及除 $z^i$ 以外的所有項。} 由 (1) 得到
% unit: noether.p06.eq.055
\begin{equation}
\begin{aligned}
  \Phi(z,u)
  &=z^i+
    \sum{}' g_{\alpha\alpha_1\ldots\alpha_\rho}(x)\cdot
    z^\alpha u_1^{\alpha_1}\cdots u_\rho^{\alpha_\rho},\\
  &\hspace{8em}(\alpha+\alpha_1+\cdots+\alpha_\rho=i),
\end{aligned}
\tag{2}
\end{equation}
作為 $\Kfield_{n\rho}$ 的\emph{對合形式}。

% unit: noether.p06.s05.par.003
對合形式 $\Phi^*(z,u)$ 的係數將會表明是 $\Kfield^*_{\rho\rho}$ 的一個有理基；也就是說，有關係
% unit: noether.p06.eq.056
\begin{equation}
  \Omega\bigl(g^*_{\alpha\alpha_1\ldots\alpha_\rho}(x_1,
  \ldots,x_\rho)\bigr)=\Kfield^*_{\rho\rho}
  \quad(\alpha+\alpha_1+\cdots+\alpha_\rho=i).
\tag{3}
\end{equation}
證明如下。注意到 $\Phi^*(z,u)$ 的係數屬於 $\Omega(g^*_{\alpha\alpha_1\ldots\alpha_\rho}(x))$，並且由於它相對於 $\Kfield^*_{\rho\rho}$ 不可約，它對於這個子域更加不可約。因此 $\Phi^*(z,u)$ 給出 $u(x)$ 相對於 $\Omega(g^*_{\alpha\alpha_1\ldots\alpha_\rho}(x))$ 的不可約方程。於是 $\Omega(x_1,\ldots,x_\rho)$ 是 $\Omega(g^*_{\alpha\alpha_1\ldots\alpha_\rho}(x))$ 的一個代數擴張，次數為 $i$。由於 $\Kfield^*_{\rho\rho}$ 位於 $\Omega(g^*_{\alpha\alpha_1\ldots\alpha_\rho}(x))$ 與 $\Omega(x_1\cdots x_\rho)$ 之間，次數相等按已知事實蘊涵域相等，從而得到關係 (3)。由傳遞原理，$\Phi(z,u)$ 的係數相應地給出 $\Kfield_{n\rho}$ 的一個有理基；即我們有：

% unit: noether.p06.thm.003
\emph{定理 III：對合形式的係數構成一個由域確定的有理基，即對合基；對於域 $\Kfield_{\rho\rho}$，這個基是唯一確定的。}\footnote{對於域 $\Kfield_{n\rho}$，對合形式可能依賴於映射；唯一性可以通過帶不定係數的映射獲得。}

% unit: noether.p06.s05.par.004
我們補充這個事實的一個非有理解釋，它建立了與對合的幾何理論的聯繫。

% unit: noether.p06.s05.par.005
在一個擴張域中，$\Phi^*(z,u)$ 分解為線性因子如下：
% unit: noether.p06.eq.057
\begin{equation}
\begin{aligned}
 \Phi^*(z,u)&=(z-u_1x_1-\cdots-u_\rho x_\rho)
 (z-u_1x_1^{(1)}-\cdots-u_\rho x_\rho^{(1)})\cdots\\
 &\qquad\cdot
 (z-u_1x_1^{(i-1)}-\cdots-u_\rho x_\rho^{(i-1)}).
\end{aligned}
\tag{4}
\end{equation}
與 (1) 比較可知，函數
% unit: noether.p06.eq.058
\[
  g^*_{\alpha\alpha_1\ldots\alpha_\rho}(x_1,
  \ldots,x_\rho)\quad(\alpha+\alpha_1+\cdots+\alpha_\rho=i)
\]
等同於下列各行量的初等對稱函數：
% unit: noether.p06.eq.059
\begin{equation}
\begin{array}{cccc}
  x_1 & x_2 & \cdots & x_\rho\\
  x_1^{(1)} & x_2^{(1)} & \cdots & x_\rho^{(1)}\\
  \vdots & \vdots & & \vdots\\
  x_1^{(i-1)} & x_2^{(i-1)} & \cdots & x_\rho^{(i-1)}.
\end{array}
\tag{5}
\end{equation}
因此，$\Kfield^*_{\rho\rho}$ 中每一個函數都是這些在行 (5) 中對稱的初等函數的有理組合；反過來，行 (5) 的每個對稱函數都屬於 $\Kfield^*_{\rho\rho}$。

% unit: noether.p06.s05.par.006
於是，當
% unit: noether.p06.eq.060
\[
  f^*(x)=\frac{F^*(x)}{H^*(x)}
\]
遍歷 $\Kfield^*_{\rho\rho}$ 中所有函數時，存在一個無限多關係組成的系統：
% unit: noether.p06.eq.061
\begin{equation}
  \frac{F^*(x)}{H^*(x)}=
  \frac{F^*(x^{(1)})}{H^*(x^{(1)})}=\cdots=
  \frac{F^*(x^{(i-1)})}{H^*(x^{(i-1)})},
\tag{6}
\end{equation}
或概括為
% unit: noether.p06.eq.062
\[
  F^*(x^{(k)})H^*(x^{(l)})-
  F^*(x^{(l)})H^*(x^{(k)})=0
  \quad(k,l=0,1,\cdots,i-1),
\]
其中 $F^*(x^{(k)})$ 與 $H^*(x^{(k)})$ 都不恆等為零，因為它們與 $F^*(x)$ 和 $H^*(x)$ 共軛，並且形式上是對稱的。

% unit: noether.p06.s05.par.007
反過來，對於滿足無限多關係
% unit: noether.p06.eq.063
\begin{equation}
  F^*(x)H^*(\xi)-H^*(x)F^*(\xi)=0,
  \qquad H^*(\xi)\ne0,
\tag{7}
\end{equation}
的每一行量 $\xi$，都可得到其屬於量行 (5)。

% unit: noether.p06.s05.par.008
事實上，如果 $G^*(x)$ 表示 $\Phi^*(z,u)$ 的係數的公分母，那麼由假設 (7)，函數
% unit: noether.p06.eq.064
\[
  G^*(\xi)\cdot\Phi^*(z,u;\xi),
\]
它關於 $\xi$ 也為整函數，並不恆等為零——也就是說，$\Phi^*(z,u;\xi)$ 的係數不變成不定——且具有零點
% unit: noether.p06.eq.065
\[
  z=u_1\xi_1+u_2\xi_2+\cdots+u_\rho\xi_\rho.
\]
由於 (7)，$\Phi^*(z,u;x)$ 因而也具有零點 $z=u(\xi)$；也就是說，$\xi$ 屬於量行 (5)。若 $\xi$ 滿足關於某個共軛量行的方程系統，則由 (6) 有相應斷言：
% unit: noether.p06.eq.066
\[
  F^*(x^{(k)})H^*(\xi)-H^*(x^{(k)})F^*(\xi)=0;
  \qquad H(\xi)\ne0.
\]
因此，量行 (5) 被定義為方程系統
% unit: noether.p06.eq.067
\[
  F^*(x)H^*(\xi)-H^*(x)F^*(\xi)=0;
  \qquad H^*(\xi)\ne0,
\]
的解 $\xi$，其中 $F^*(x):H^*(x)$ 遍歷 $\Kfield^*_{\rho\rho}$ 中所有函數；這裡的行 $x$ 也可以由任何共軛行來替換。

% unit: noether.p06.s05.par.009
從幾何上說，若把每一行 $x$ 解釋為一個點，則 (7) 的解在一個 $\rho$ 維線性空間中表示 $\infty^\rho$ 個點群，每個點群有 $i$ 個點，並且這個群中的任意一個點唯一決定該個別群。這樣的點群系統稱為“\emph{線性空間中的對合}”。\footnote{(7) 中被排除的那些解，即 $F^*(\xi)=0$, $H^*(\xi)=0$ 的解，以及與對合基相應的方程系統中被排除的解，稱為對合的基本點；其中還包括由齊次化產生的點，即“位於無窮遠處”的點。}

% unit: noether.p06.s05.par.010
相應地，一個域 $\Kfield_{n\rho}$ 刻畫了一個由曲線、曲面等等組成的線性空間中的對合。

% unit: noether.p06.s05.par.011
根據上面的討論，$\Omega(x_1,\ldots,x_\rho)$ 關於 $\Kfield^*_{\rho\rho}$ 的次數 $i$ 等於滿足附加條件 $H^*(\xi)\ne0$ 的 (7) 的解系統個數——我們可以把它們稱為 $\Kfield^*_{\rho\rho}$ 的解系統。因此，用於證明對合基存在性的論證也允許如下非有理的表述：

% unit: noether.p06.s05.par.012
“如果 $\mathfrak{L}^*$ 是 $\Kfield^*_{\rho\rho}$ 的一個因子，並且域 $\mathfrak{L}^*$ 的每一個解系統也是 $\Kfield^*_{\rho\rho}$ 的一個解系統，那麼 $\mathfrak{L}^*$ 與 $\Kfield^*_{\rho\rho}$ 相同。”

% unit: noether.p06.s05.par.013
由傳遞原理，對於 $\Kfield_{n\rho}$ 的因子 $\mathfrak{L}$ 也有相應說法。

% unit: noether.p06.sec.006
\section*{\S\ 6. 域 $\Kfield_{n\rho}$ 的最小基。}

% unit: noether.p06.s06.par.001
本節給出已知定理的概覽，但採用一種由傳遞原理和有理基存在性所可能的擴展表述。

% unit: noether.p06.def.004
\emph{定義 IV：$\Kfield_{n\rho}$ 的有理基若只含有最小數目 $\rho$ 個函數，則稱為最小基；這些函數於是必須是代數無關的。}

% unit: noether.p06.s06.par.002
由 Lüroth、Castelnuovo 和 Enriques 的定理\footnote{J. Lüroth: \emph{Beweis eines Satzes über rationale Kurven}, \emph{Math. Ann.} 9 (1875). --- G. Castelnuovo: \emph{Sulla razionalità delle involuzioni piane}, \emph{Math. Ann.} 44 (1893). --- F. Enriques: \emph{Sopra una involuzione non razionale dello spazio}, \emph{Rend. Acc. Linc.}, Vol. 21, 21 Jan. 1912.}得到如下事實：

% unit: noether.p06.s06.par.003
I. 每一個域 $\Kfield_{n1}$ 都具有一個最小基，即該域的 Lüroth 函數，它在分式線性變換之下唯一確定。

% unit: noether.p06.s06.par.004
II. 對於域 $\Kfield_{n2}$，如果允許係數域 $\Omega$ 作代數擴張，則最小基總是存在，並且一般也只有在這種情況下存在。

% unit: noether.p06.s06.par.005
III. 對於域 $\Kfield_{n\rho}$（$\rho\ge3$），一般不存在最小基。對於具有最小基的特殊域 $\Kfield_{n\rho}$，尚未嘗試作出特徵刻畫。

% unit: noether.p06.s06.par.006
我們簡要說明 E. Steinitz 給出的域 $\Kfield^*_{11}$ 情形下 Lüroth 定理的證明：\footnote{同前引，\S\ 24。}

% unit: noether.p06.s06.par.007
設 $\Phi^*(z,u)$ 為 $\Kfield^*_{11}$ 的對合形式，$\psi^*(x)$ 為 $\Phi^*(z,u)$ 的任意一個實際含有不定元 $x$ 的係數。可以看出，$\Omega(x)$ 關於 $\Omega(\psi^*(x))$ 的次數等於 $\Omega(x)$ 關於 $\Kfield^*_{11}$ 的次數，由此推出域相等。此外，如果 $\Omega(\chi^*(x))$ 等於 $\Omega(\psi^*(x))$，則 $\chi^*(x)$ 與 $\psi^*(x)$ 的相互有理依賴蘊涵這兩個函數之間的分式線性依賴。

% unit: noether.p06.s06.par.008
因此，由傳遞原理，Lüroth 定理也容許如下表述：

% unit: noether.p06.s06.par.009
\emph{域 $\Kfield_{n1}$ 的最小基由對合基中的任意一個函數構成；並且 $\Kfield_{n1}$ 的對合基中任意兩個函數彼此分式線性地依賴。}

% unit: noether.p06.s06.par.010
G. Castelnuovo 藉助對合的幾何理論，對由三個函數的有理基導出的域 $\Kfield^*_{22}$ 證明事實 II。由於傳遞原理在係數域 $\Omega$ 的有限代數擴張下也保持有效，有理基的存在性便使這個證明對所有域 $\Kfield_{n2}$ 完成。

% unit: noether.p06.s06.par.011
關於 III，應當指出，Enriques 構造了三維空間中的一個非有理對合；也就是說，一個沒有最小基的域 $\Kfield_{33}$。
%
\clearpage
%
% unit: noether.p06.s07.h.001
\section*{\S\ 7. 任意系統 $\Ssys_{n\rho}$、線性族 $\Lsys_{n\rho}$ 以及有理函數的整性域 $\Jdom_{n\rho}$。有理基的存在。}

% unit: noether.p06.s07.par.001
由域 $\Kfield_{n\rho}$ 的有理基的存在性，現在可以推出任意有理函數系統和整有理函數系統的有理基。我們按照這些系統由加法、乘法和除法這些基本運算相互逐次產生的順序來排列它們。

% unit: noether.p06.s07.par.002
I. 如往常一樣，$\Ssys_{n\rho}$ 表示 $n$ 個不定元的有理函數（或整有理函數）的任意系統，其代數秩為 $\rho$。來自 $\Omega$ 的量也算作屬於這個系統。

% unit: noether.p06.s07.par.003
II. 如果系統 $\Lsys_{n\rho}$ 滿足下列條件，則稱它為\emph{線性族}：

1. 與 $f(x)$ 同時，$\Lsys_{n\rho}$ 總還包含 $c\cdot f(x)$，其中 $c$ 是 $\Omega$ 中的任意量；

2. 與 $f(x)$ 和 $g(x)$ 同時，$\Lsys_{n\rho}$ 總還包含 $f(x)+g(x)$。

% unit: noether.p06.s07.par.004
III. 線性族 $\Jdom_{n\rho}$ 在進一步滿足下列條件時稱為\emph{整性域}：

3. 與 $f(x)$ 和 $g(x)$ 同時，$\Jdom_{n\rho}$ 總還包含 $f(x)\cdot g(x)$。

% unit: noether.p06.s07.par.005
IV. 整性域 $\Kfield_{n\rho}$ 在進一步滿足下列條件時稱為\emph{域}：

4. 與 $f(x)$ 和 $g(x)$ 同時——其中 $g(x)\ne0$——$\Kfield_{n\rho}$ 總還包含商 $f(x):g(x)$。

% unit: noether.p06.s07.par.006
條件 1 到 4 就是 \S\ 1 中對域給出的條件。\footnote{$f(x)$ 和 $g(x)$ 也可以是零次的，即為 $\Omega$ 中的量。}

% unit: noether.p06.s07.par.007
現在要把這些域依次由彼此推出。

% unit: noether.p06.s07.par.008
a) 任意系統 $\Ssys_{n\rho}$ 由下列要求擴充為一個線性族 $\Lsys_{n\rho}$：

% unit: noether.p06.s07.par.009
除 $\Ssys_{n\rho}$ 之外，$\Lsys_{n\rho}$ 包含所有且僅包含那些函數 $h(x)$，這些函數可用 $\Ssys_{n\rho}$ 中有限個函數並以 $\Omega$ 中的係數線性表示：
% unit: noether.p06.s07.eq.001
\[
  h(x)=c_1f_1(x)+c_2f_2(x)+\cdots+c_kf_k(x),
\]
其中 $f_1(x),\ldots,f_k(x)$ 遍歷 $\Ssys_{n\rho}$ 中的全部函數，而 $c_1,
\ldots,c_k$ 遍歷 $\Omega$ 中的全部量。

% unit: noether.p06.s07.par.010
因為把系統中函數的有理組合添加進去，會保持一個系統的代數秩不變。$\Lsys_{n\rho}$ 滿足條件 1 和 2，因而是線性族。並且，任何包含 $\Ssys_{n\rho}$ 的線性族，按 1 和 2 也包含 $\Lsys_{n\rho}$；所以 $\Lsys_{n\rho}$ 是包含 $\Ssys_{n\rho}$ 的最小線性族。

% unit: noether.p06.s07.par.011
b) 相應地，從線性族 $\Lsys_{n\rho}$ 由下列要求產生包含它的最小整性域 $\Jdom_{n\rho}$：

% unit: noether.p06.s07.par.012
除 $\Lsys_{n\rho}$ 之外，$\Jdom_{n\rho}$ 包含所有且僅包含那些函數 $h(x)$，這些函數能夠表示為 $\Lsys_{n\rho}$ 中有限個函數 $f_1(x),\ldots,f_k(x)$ 的、以 $\Omega$ 中的係數作成的\emph{整有理組合}；這裡 $f_1(x),\ldots,f_k(x)$ 遍歷 $\Lsys_{n\rho}$ 中的全部函數。

% unit: noether.p06.s07.par.013
c) 最後，從整性域 $\Jdom_{n\rho}$ 由下列要求產生包含它的最小域 $\Kfield_{n\rho}$：除 $\Jdom_{n\rho}$ 之外，$\Kfield_{n\rho}$ 包含所有且僅包含那些函數 $h(x)$，這些函數可表示為 $\Jdom_{n\rho}$ 中兩個函數的商。

% unit: noether.p06.s07.par.014
把這三個擴充合併為一步，就得到：

% unit: noether.p06.s07.thm.001
\emph{每一個系統 $\Ssys_{n\rho}$ 都可以由下列要求擴充為包含它的最小域 $\Kfield_{n\rho}$：}

% unit: noether.p06.s07.thm.002
\emph{除 $\Ssys_{n\rho}$ 之外，$\Kfield_{n\rho}$ 包含所有且僅包含那些函數 $h(x)$，這些函數能夠表示為 $\Ssys_{n\rho}$ 中有限個函數 $f_1(x),\ldots,f_k(x)$ 的、以 $\Omega$ 中的係數作成的有理組合；這裡 $f_1(x),\ldots,f_k(x)$ 遍歷 $\Ssys_{n\rho}$ 中的全部函數。}

% unit: noether.p06.s07.par.015
由後一事實，任意域的有理基存在性立即得出：

% unit: noether.p06.s07.par.016
令 $\Kfield_{n\rho}$ 為包含 $\Ssys_{n\rho}$ 的最小域，並設 $\Kfield_{n\rho}$ 的一個有理基由
% unit: noether.p06.s07.eq.002
\begin{equation}
  h_1(x),\ h_2(x),\ldots,h_\sigma(x)
\tag{1}
\end{equation}
給出。按照 $\Kfield_{n\rho}$ 的定義，存在關係
% unit: noether.p06.s07.eq.003
\begin{equation}
\begin{aligned}
 h_1(x)&=\gamma_1\bigl(g_1(x),\ldots,g_\lambda(x)\bigr);\ldots;\\
 h_\sigma(x)&=\gamma_\sigma\bigl(g_\mu(x),\ldots,g_\nu(x)\bigr),
\end{aligned}
\tag{2}
\end{equation}
其中 $r_1,
\ldots,r_\sigma$ 是其自變量的有理函數，而
% unit: noether.p06.s07.eq.004
\begin{equation}
  g_1(x),\ldots,g_\lambda(x),\ldots,
  g_\mu(x),\ldots,g_\nu(x)
\tag{3}
\end{equation}
全都屬於 $\Ssys_{n\rho}$。由於關係 (2)，(3) 除了 (1) 之外也表示 $\Kfield_{n\rho}$ 的一個有理基；又因為函數 (3) 屬於 $\Ssys_{n\rho}$，它同時也是 $\Ssys_{n\rho}$ 的一個有理基。因此有

% unit: noether.p06.s07.thm.003
\emph{定理 IV：有理函數（或整有理函數）的任意系統 $\Ssys_{n\rho}$ 都有一個有理基，它由該系統中有限個函數組成，代數秩為 $\rho$。}

% unit: noether.p06.s07.par.017
現在特別令 $\Lsys_{n\rho}$ 為一個線性族，並令
% unit: noether.p06.s07.eq.005
\begin{equation}
  g_1(x),\ldots,g_\rho(x),\ g_{\rho+1}(x),\ldots,g_\nu(x)
\tag{4}
\end{equation}
為 $\Lsys_{n\rho}$ 的一個有理基。假設例如
% unit: noether.p06.s07.eq.006
\[
  g_1(x)\cdots g_\rho(x)
\]
是代數無關的。如果於是置
% unit: noether.p06.s07.eq.007
\[
  g(x)=c_1g_{\rho+1}(x)+c_2g_{\rho+2}(x)+\cdots+c_{\nu-\rho}g_\nu(x),
\]
那麼可以把 $c_i$ 選為 $\Omega$ 中的數，使得
% unit: noether.p06.s07.eq.008
\begin{equation}
  g_1(x),\ldots,g_\rho(x),\ g(x)
\tag{5}
\end{equation}
成為包含 $\Lsys_{n\rho}$ 的最小域 $\Kfield_{n\rho}$ 的一個有理基。但由於 $\Lsys_{n\rho}$ 是線性族，$g(x)$ 屬於 $\Lsys_{n\rho}$，而 (5) 表示 $\Lsys_{n\rho}$ 的一個有理基；也就是說：

% unit: noether.p06.s07.thm.004
\emph{定理 V：有理函數（或整有理函數）的每一個線性族 $\Lsys_{n\rho}$ 的有理基——因而特別地，每一個整性域 $\Jdom_{n\rho}$ 和每一個域 $\Kfield_{n\rho}$ 的有理基——都可以特別選成至多含有 $\rho+1$ 個函數（且至少含有 $\rho$ 個函數）。}

% unit: noether.p06.s07.par.018
因此，\S\ 4 中對域的有理基函數個數所得的界，依賴於域作為線性族這一性質；而在最小基存在時把函數個數限制為 $\rho$，則使用了域的全部假設。

% unit: noether.p06.s07.par.019
還需注意，按照 \S\ 3 (4)，系統以及包含它們的最小系統的全部特徵性質，在映射之下都保持不變。

% unit: noether.p06.s08.h.001
\section*{\S\ 8. 由多項式組成的整性域 $\Jdom_{n\rho}$ 的對合基。}

% unit: noether.p06.s08.par.001
在以下各節中，我們處理元素為整有理函數（多項式）的系統 $\Ssys_{n\rho}$；首先處理由多項式組成的整性域 $\Jdom_{n\rho}$。對於這些域 $\Jdom_{n\rho}$，須先固定整除概念。

% unit: noether.p06.s08.par.002
設 $F(x),G(x)$ 為 $\Jdom_{n\rho}$ 中的兩個多項式，並設 $L(x)$ 為這些多項式的一個公因子，使得
% unit: noether.p06.s08.eq.001
\[
  F(x)=L(x)\cdot F_1(x),\qquad G(x)=L(x)\cdot G_1(x).
\]
當且僅當三個多項式 $L(x),F_1(x),G_1(x)$ 都屬於 $\Jdom_{n\rho}$ 時，稱 $L(x)$ 為\emph{$\Jdom_{n\rho}$ 中的公因子}。

% unit: noether.p06.s08.par.003
如果 $\Jdom_{n\rho}$ 中的兩個多項式在 $\Jdom_{n\rho}$ 中沒有公因子，則稱它們在 $\Jdom_{n\rho}$ 中\emph{互素}。

% unit: noether.p06.s08.par.004
設 $\Kfield_{n\rho}$ 為包含 $\Jdom_{n\rho}$ 的最小域。由定義可知，$\Kfield_{n\rho}$ 中的每個函數都有表示
% unit: noether.p06.s08.eq.002
\[
  f(x)=F_1(x):F_2(x),
\]
其中 $F_1(x),F_2(x)$ 屬於 $\Jdom_{n\rho}$，並且在 $\Jdom_{n\rho}$ 中互素。相應地，$\Kfield_{n\rho}$ 中的若干函數可以化為 $\Jdom_{n\rho}$ 中的公分母：
% unit: noether.p06.s08.eq.003
\[
  f_1(x)=\frac{F_1(x)}{F(x)},\ldots,
  f_k(x)=\frac{F_k(x)}{F(x)}.
\]
當且僅當 $\Jdom_{n\rho}$ 中的多項式
% unit: noether.p06.s08.eq.004
\[
  F(x),F_1(x),\ldots,F_k(x)
\]
在 $\Jdom_{n\rho}$ 中互素時，稱 $F(x)$ 為\emph{$\Jdom_{n\rho}$ 中的最小公分母}。

% unit: noether.p06.s08.par.005
這樣的表示可能以不同方式實現，因為在 $\Jdom_{n\rho}$ 中一般不成立分解為不可約函數的唯一分解律。

% unit: noether.p06.s08.par.006
現在，我們把對合形式和對合基的概念轉移到整性域，來研究 $\Kfield_{n\rho}$ 的對合基對於 $\Jdom_{n\rho}$ 的意義。設
% unit: noether.p06.s08.eq.005
\[
  \Phi(z,u)=z^i+\sum{}' g_{\alpha\alpha_1\ldots\alpha_\rho}(x)
  z^\alpha u_1^{\alpha_1}\cdots u_\rho^{\alpha_\rho},
  \qquad
  (\alpha+\alpha_1+\cdots+\alpha_\rho=i)
\]
為 $\Kfield_{n\rho}$ 的一個對合形式，並設 $G(x)$ 為 $\Phi(z,u)$ 各系數在 $\Jdom_{n\rho}$ 中的一個最小公分母。乘以 $G(x)$ 後，$\Phi(z,u)$ 變為一個形式 $\Psi(z,u)$，其係數為 $\Jdom_{n\rho}$ 中的多項式，並且在 $\Jdom_{n\rho}$ 中互素：
% unit: noether.p06.s08.eq.006
\[
  G(x)\cdot\Phi(z,u)=\Psi(z,u)
  =G(x)z^i+
  \sum{}' G_{\alpha\alpha_1\ldots\alpha_\rho}(x)
  z^\alpha u_1^{\alpha_1}\cdots u_\rho^{\alpha_\rho},
  \qquad
  (\alpha+\alpha_1+\cdots+\alpha_o=i).
\]

% unit: noether.p06.s08.par.007
凡是由 $\Kfield_{n\rho}$ 的對合形式 $\Phi(z,u)$ 乘以 $\Jdom_{n\rho}$ 中的多項式而得到的形式 $\Psi(z,u)$，若其係數為 $\Jdom_{n\rho}$ 中的多項式並且在 $\Jdom_{n\rho}$ 中互素，就稱為 $\Jdom_{n\rho}$ 的一個\emph{對合形式}；$\Psi(z,u)$ 的係數稱為相應的\emph{對合基}。

% unit: noether.p06.s08.par.008
由 $\Kfield_{n\rho}$ 的對合基的意義首先清楚可見，$\Jdom_{n\rho}$ 的每一個對合基同時也是有理基。為進一步研究，我們考察一個映射域 $\Jdom^*_{\rho\rho}$，其最小包含域 $\Kfield^*_{\rho\rho}$ 對零
% unit: noether.p06.s08.eq.007
\[
  z=u_1x_1+\cdots+u_\rho x_\rho
\]
具有不可約方程 $\Phi^*(z,u)=0$。

% unit: noether.p06.s08.par.009
按照 \S\ 5，
% unit: noether.p06.s08.eq.008
\[
  \Phi^*(z,u)=z^i+\sum{}' g^*_{\alpha\alpha_1\ldots\alpha_\rho}(x)
  z^\alpha u_1^{\alpha_1}\cdots u_\rho^{\alpha_\rho},
  \qquad
  (\alpha+\alpha_1+\cdots+\alpha_\rho=i)
\]
的係數表示相對於 $\Kfield^*_{\rho\rho}$ 共軛的若干量行
% unit: noether.p06.s08.eq.009
\[
  x,\ x^{(1)}\cdots x^{(i-1)}
\]
的初等對稱函數。令 $F^*(x)$ 為 $\Kfield^*_{\rho\rho}$ 中任意一個多項式。由於
% unit: noether.p06.s08.eq.010
\[
  F^*(x)=\frac1i\Bigl(F^*(x)+F^*(x^{(1)})+
  \cdots+F^*(x^{(i-1)})\Bigr),
\]
可知 $F^*(x)$ 是這些量行的整對稱函數；因此它能夠表示為初等對稱函數的整有理組合，也就是表示為函數
% unit: noether.p06.s08.eq.011
\[
  g^*_{\alpha\alpha_1\ldots\alpha_\rho}(x)
\]
的整有理組合。如果現在
% unit: noether.p06.s08.eq.012
\[
  G^*(x)\cdot\Phi^*(z,u)=\Psi^*(z,u)
  =G^*(x)z^i+
  \sum{}'G^*_{\alpha\alpha_1\ldots\alpha_\rho}(x)
  z^\alpha u_1^{\alpha_1}\cdots u_\rho^{\alpha_\rho},
  \qquad
  (\alpha+\alpha_1+\cdots+\alpha_\rho=i)
\]
是 $\Jdom^*_{\rho\rho}$ 的一個對合形式，那麼 $\Kfield^*_{\rho\rho}$ 中的每個多項式，特別是 $\Jdom^*_{\rho\rho}$ 中的每個多項式，都是商
% unit: noether.p06.s08.eq.013
\[
  G^*_{\alpha\alpha_1\ldots\alpha_\rho}(x):G^*(x)
\]
的整有理組合。

% unit: noether.p06.s08.par.010
考慮到轉移原則，就得到

% unit: noether.p06.s08.thm.001
\emph{定理 VI：對於每一個多項式的整性域 $\Jdom_{n\rho}$，至少存在一個特殊的有理基，即對合基；在用對合基的函數對 $\Jdom_{n\rho}$ 中所有多項式作有理表示時，分母中只出現某個固定函數（相應對合形式的最高係數）的冪。}

% unit: noether.p06.s09.h.001
\section*{\S\ 9. 由多項式組成的相對整域 $\Gdom_{n\rho}$。}

% unit: noether.p06.s09.par.001
下面考察由多項式組成的特殊整性域，它們構成整性域與域之間的一個中間階段。

% unit: noether.p06.s09.par.002
由多項式組成的整性域 $\Gdom_{n\rho}$ 在下列條件下稱為“相對整域”：

% unit: noether.p06.s09.def.001
\emph{4a) 與 $F(x)$ 和 $G(x)$ 同時——其中 $G(x)\ne0$——$\Gdom_{n\rho}$ 總是且僅在商 $F(x):G(x)$ 關於不定元為整、也就是為多項式時，包含這個商。}

% unit: noether.p06.s09.par.003
首先要證明，相對整域與一個有理函數域 $\Kfield_{n\sigma}$（$\sigma\ge \rho$）中全部多項式的集合相同。

% unit: noether.p06.s09.par.004
設 $\Kfield_{n\rho}$ 是包含 $\Gdom_{n\rho}$ 的最小域。對於 $\Kfield_{n\rho}$ 中的每個函數 $f(x)$，存在一個作為 $\Gdom_{n\rho}$ 中兩個多項式的商的表示：
% unit: noether.p06.s09.eq.001
\[
  f(x)=F_1(x):F_2(x),
\]
這裡也允許零次多項式。只要 $f(x)$ 成為多項式，它就按條件 4a) 屬於 $\Gdom_{n\rho}$；而由於 $\Gdom_{n\rho}$ 不能包含其他多項式，\emph{每一個相對整域 $\Gdom_{n\rho}$ 都與其最小包含域中的全部多項式相同。}

% unit: noether.p06.s09.par.005
反過來，令 $\Kfield_{n\sigma}$ 為任意一個有理函數域（因此不一定是某個給定系統的最小包含域），並令 $\Jdom$ 為 $\Kfield_{n\rho}$ 中所含多項式的全體。$\Jdom$ 滿足整性域的條件 1、2、3 以及條件 4a)，所以是一個相對整域；也就是說，\emph{一個域 $\Kfield_{n\sigma}$ 中所含全部多項式構成一個相對整域 $\Gdom_{n\rho}$（$\rho\le\sigma$）。}\footnote{也可以有 $\rho=0$；也就是說，$\Kfield_{n\sigma}$ 中的全部多項式由係數域 $\Omega$ 的元素組成。}

% unit: noether.p06.s09.par.006
我們還要說明，相對整域與 Hilbert 的“相對整函數的整性域”相同。\footnote{D. Hilbert: \emph{Mathematische Probleme}. 1900 年巴黎報告，第 14 問題（\emph{Göttinger Nachrichten} 1900）。}

% unit: noether.p06.s09.par.007
設
% unit: noether.p06.s09.eq.002
\begin{equation}
  G_1(x),\ldots,G_k(x)
\tag{1}
\end{equation}
為 $\Gdom_{n\rho}$ 的一個有理基。按條件 1、2、3、4a，多項式 (1) 的任何有理組合，只要關於不定元為整、即成為多項式，就屬於 $\Gdom_{n\rho}$。反過來，$\Gdom_{n\rho}$ 中的每一個多項式，按有理基的概念，都可以由多項式 (1) 有理表示；因此 $\Gdom_{n\rho}$ \emph{正是那些多項式 (1) 的有理組合的全體，其中這些組合關於不定元為整，也就是說為多項式。} 於是我們得到從特殊有理基出發的 Hilbert 定義；而我們的定義只使用域的抽象性質。

% unit: noether.p06.s09.par.008
對於相對整域，\S\ 8 中為一般整性域給出的公因子定義可以簡化：

% unit: noether.p06.s09.par.009
設 $F(x),G(x)$ 為 $\Gdom_{n\rho}$ 中的兩個多項式，並設 $L(x)$ 為這些多項式的一個公因子。那麼，當且僅當 $L(x)$ 屬於 $\Gdom_{n\rho}$ 時，稱 $L(x)$ 為\emph{$\Gdom_{n\rho}$ 中的公因子}。

% unit: noether.p06.s09.par.010
因為商 $F(x):L(x)$ 與 $G(x):L(x)$ 的歸屬隨後由條件 4a) 得出。

% unit: noether.p06.s09.par.011
相對整域的另一個本質性質是，“關於 $\Gdom_{n\rho}$ 代數整”的概念可以藉助任意方程來定義，正如代數數域或域 $\Omega(x_1,
\ldots,x_n)$ 中所熟知的那樣。

% unit: noether.p06.s09.def.002
\emph{定義 V：如果一個量 $\xi$ 滿足某個方程，其最高係數為單位，而其餘多項式來自 $\Gdom_{n\rho}$，則稱 $\xi$ 關於 $\Gdom_{n\rho}$ 代數整，或也稱相對代數整。}

% unit: noether.p06.s09.par.012
事實上，假設相對代數整的量 $\xi$ 例如滿足方程
% unit: noether.p06.s09.eq.003
\[
  L(z)=z^\lambda+G_1(x)z^{\lambda-1}+\cdots+G_\lambda(x)=0.
\]
$L(z)$ 可被關於最小包含域 $\Kfield_{n\rho}$ 不可約的整函數
% unit: noether.p06.s09.eq.004
\[
  M(z)=z^\mu+H_1(x)z^{\mu-1}+\cdots+H_\mu(x)
\]
整除。但由這種整除性，根據 Gauss 定理的已知推廣\footnote{例如比較 J. König: \emph{Einleitung in die allgemeine Theorie der algebraischen Größen}（Leipzig, B.~G. Teubner, 1903），第 IX 章，\S\ 2，或 Weber（小版），\S\ 20。}，可知有理函數 $H_1(x),\ldots,H_\mu(x)$ 是 $\Kfield_{n\rho}$ 中的多項式，因此屬於 $\Kfield_{n\rho}$。於是相對代數整量的定義可以從不可約方程得到。特別地，$\Gdom_{n\rho}$ 中的每個多項式 $F(x)$ 也關於 $\Gdom_{n\rho}$ 代數整，因為它滿足方程 $z-F(x)=0$。

% unit: noether.p06.s09.par.013
還應指出，對於多項式的整性域 $\Jdom_{n\rho}$，類似於最小包含域，包含它的最小相對整域 $\Gdom_{n\rho}$ 也由下列要求定義：

% unit: noether.p06.s09.par.014
與 $\Jdom_{n\rho}$ 同時，$\Gdom_{n\rho}$ 包含所有且僅包含那些多項式 $H(x)$，它們可表示為 $\Jdom_{n\rho}$ 中兩個多項式的商。

% unit: noether.p06.s09.par.015
由這個定義還得到：從 $\Jdom_{n\rho}$ 的每個對合形式和對合基，通過至多分離掉基中所有多項式在 $\Gdom_{n\rho}$ 中的一個公因子，就產生 $\Gdom_{n\rho}$ 的相應對合形式和對合基。

% unit: noether.p06.s09.par.016
因為 $\Jdom_{n\rho}$ 和 $\Gdom_{n\rho}$ 具有同一個最小包含域 $\Kfield_{n\rho}$；而 $\Jdom_{n\rho}$ 的每個對合形式，都由 $\Kfield_{n\rho}$ 的一個對合形式乘以 $\Jdom_{n\rho}$ 中的多項式得到，因此具有 $\Gdom_{n\rho}$ 中的係數；不過這些係數可能在 $\Gdom_{n\rho}$ 中具有公因子。\footnote{例如，對於由 $x_1^2$ 和 $x_1x_2$ 的所有整有理組合所組成的整性域 $\Jdom_{22}$，得到如下對合形式：
\[
\Psi(z,u)=x_1^2\cdot z^2-x_1^4u_1^2-2x_1^3x_2u_1u_2-x_1^2x_2^2u_2^2.
\]
由於 $x_1^2$ 屬於 $\Jdom_{22}$，因而更屬於最小包含相對整域 $\Gdom_{22}$，於是它是 $\Gdom_{22}$ 中的公因子，所以得到 $\Gdom_{22}$ 的如下對合形式：
\[
\Psi(z,u)=z^2-x_1^2u_1^2-2x_1x_2u_1u_2-x_2^2u_2^2.
\]}

% unit: noether.p06.s09.par.017
最後還要指出，相對整域 $\Gdom_{n\rho}$ 的\emph{映射域}，按 \S\ 7 又是一個多項式的整性域 $\Jdom^*_{\rho\rho}$，但它本身不一定是相對整域。因為如果 $F(x)$ 和 $G(x)$ 是 $\Gdom_{n\rho}$ 中的兩個多項式，而它們的商不是多項式、因而不屬於 $\Gdom_{n\rho}$，那麼商 $F^*(x):G^*(x)$ 也不屬於 $\Jdom^*_{\rho\rho}$。可是，這個由若干不定元的特殊化而來的商，很可能是一個多項式，於是對於 $\Jdom^*_{\rho\rho}$，條件 4a) 就不滿足。\footnote{例如令 $\Gdom_{32}$ 由所有可表示為 $F=x_1^2x_2$ 和 $G=x_1^2(1+x_3)$ 的有理組合的多項式組成。可允許的特殊化 $x_3=0$ 給出 $F^*:G^*=x_2$，而 $F:G$ 不是多項式。因此 $\Jdom^*_{22}$ 不是相對整域。}
%
\clearpage
%
% unit: noether.p06.s10.title.001
\section*{\S\ 10. 第一種相對整域及其整性基底。}

% unit: noether.p06.s10.par.001
\S\ 9 中引入的概念（定義 V）“相對代數整”，使我們能夠給出一類相對整域，它們是有限整性域；從而，至少對於這一類，可以肯定地回答 D. Hilbert 提出的一個問題（《數學問題》第 14 問題）。我們給出如下定義。

% unit: noether.p06.s10.def.006
\emph{定義 VI：$\Gdom_{n\rho}$ 中有限個多項式稱為一個整性基底，如果 $\Gdom_{n\rho}$ 中的每個多項式都能用來自 $\Omega$ 的係數表示為這有限個多項式的整有理組合。一個具有整性基底的多項式整性域稱為有限整性域。}

% unit: noether.p06.s10.par.002
對於將要考察的相對整域類，其對合基將證明是整性基底；按以下定義，我們把它們稱為第一種相對整域。

% unit: noether.p06.s10.def.007
\emph{定義 VII：若相對整域 $\Gdom_{n\rho}$ 作為映射域具有一個相對整域 $\Gdom^*_{\rho\rho}$，使得 $x_1,\ldots,x_\rho$ 的每個任意多項式（係數在 $\Omega$ 中）都關於 $\Gdom^*_{\rho\rho}$ 代數整，則稱 $\Gdom_{n\rho}$ 為第一種。}

% unit: noether.p06.s10.par.003
按照這個定義，不定元 $x_1,\ldots,x_\rho$，因而線性形式 $u(x)=u_1x_1+\cdots+u_\rho x_\rho$，都關於 $\Gdom^*_{\rho\rho}$ 代數整。以 $z=u(x)$ 為零點的不可約整函數，即最小包含域 $\Kfield^*_{\rho\rho}$ 的對合形式，因此具有單位作為最高係數，並以 $\Gdom^*_{\rho\rho}$ 中的多項式作為其餘係數；所以它同時也是 $\Gdom^*_{\rho\rho}$ 的對合形式，而且在 $\Gdom^*_{\rho\rho}$ 中不可能存在其他對合形式。由轉移原理，$\Gdom_{n\rho}$ 也具有一個最高係數為單位的對合形式；又因按定理 VI，用相應對合基對 $\Gdom_{n\rho}$ 中所有多項式作有理表示時，只有這個最高係數進入分母，所以對合基成為整性基底。於是：

% unit: noether.p06.s10.thm.007
\emph{定理 VII。每個第一種相對整域都是有限整性域；整性基底由從其特徵映射域得到的對合基給出。特別地，第一種相對整域 $\Gdom_{\rho\rho}$ 的整性基底由 $\Gdom_{\rho\rho}$ 的唯一確定的對合基給出。}

% unit: noether.p06.s10.par.004
我們還給出第一種相對整域的一個判別準則。設 $\Gdom_{n\rho}$ 為第一種，且 $\Gdom_{n\rho}$ 的一個整性基底為
% unit: noether.p06.s10.eq.001
\begin{equation}
  F_1(x),F_2(x),\ldots,F_k(x).
\tag{1}
\end{equation}
那麼，由 Hilbert 的一個輔助引理，\footnote{D. Hilbert: \emph{Über die vollen Invariantensysteme}, Math. Ann. 42 (1893), \S\ 1。那裡只對齊次多項式陳述的輔助引理，同樣適用於非齊次多項式。} 存在 $\Gdom_{n\rho}$ 中 $\rho$ 個代數無關多項式
% unit: noether.p06.s10.eq.002
\begin{equation}
  G_1(x),G_2(x),\ldots,G_\rho(x),
\tag{2}
\end{equation}
使得多項式 (1)，因而 $\Gdom_{n\rho}$ 中的每個多項式，都關於多項式 (2) 代數整。對於映射域 $\Gdom^*_{\rho\rho}$ 中的每個多項式，相對於相應的多項式
% unit: noether.p06.s10.eq.003
\begin{equation}
  G_1^*(x),G_2^*(x),\ldots,G_\rho^*(x)
\tag{3}
\end{equation}
也有同樣的結論。

% unit: noether.p06.s10.par.005
因此，按定義 VII，$x_1,\ldots,x_\rho$ 中的每個任意多項式也關於多項式 (3) 代數整。

% unit: noether.p06.s10.par.006
反過來，假設某個任意相對整域 $\Gdom_{n\rho}$ 的一個映射域 $\Jdom^*_{\rho\rho}$ 含有 $\rho$ 個多項式 (3)，而 $x_1,\ldots,x_\rho$ 的每個多項式都關於它們代數整。於是我們說明，由此 $\Jdom^*_{\rho\rho}$ 本身成為一個相對整域 $\Gdom^*_{\rho\rho}$，並進一步說明 $\Gdom^*_{\rho\rho}$，因而 $\Gdom_{n\rho}$，都是第一種。

% unit: noether.p06.s10.par.007
事實上，設 $F^*(x)$ 是最小包含 $\Jdom^*_{\rho\rho}$ 的域 $\Kfield^*_{\rho\rho}$ 中的一個多項式。按假設，它關於多項式 (3) 代數整。由於 $\Kfield_{n\rho}$ 中相應的函數於是關於 $\Gdom_{n\rho}$ 中 $\rho$ 個多項式代數整，它就是一個多項式 $F(x)$，因而屬於 $\Gdom_{n\rho}$；所以 $F^*(x)$ 屬於 $\Jdom^*_{\rho\rho}$。於是映射域 $\Jdom^*_{\rho\rho}$ 含有 $\Kfield^*_{\rho\rho}$ 中全部多項式，是相對整域 $\Gdom^*_{\rho\rho}$。對於 $\Gdom^*_{\rho\rho}$，因而對於 $\Gdom_{n\rho}$，由定義 V、定義 VII 和假設立即推出它們為第一種；也就是說：

% unit: noether.p06.s10.par.008
\emph{第一種相對整域 $\Gdom_{n\rho}$ 的特徵是：在一個映射域 $\Jdom^*_{\rho\rho}$ 中，存在 $\rho$ 個多項式，$x_1,\ldots,x_\rho$ 的每個任意多項式都關於它們代數整。從這個假設推出映射域本身就是相對整域。}\footnote{例如，在 \S\ 9 末尾註釋中提到的域 $\Gdom_{32}$ 中，可允許的特殊化 $x_1=1$ 給出兩個多項式 $F^*=x_2$, $G^*=1+x_3$，而 $x_2,x_3$ 中的每個任意多項式都關於它們代數整。因此 $\Gdom_{32}$ 為第一種，並且以 $F(x)$ 和 $G(x)$ 為整性基底。}

% unit: noether.p06.s11.title.001
\section*{\S\ 11. 關於代數整依賴的輔助引理。}

% unit: noether.p06.s11.par.001
現在我們給出一個判別準則，用來判斷 $x_1,\ldots,x_\rho$ 中的每個任意多項式是否關於 $x_1,\ldots,x_\rho$ 中的 $\rho$ 個多項式
% unit: noether.p06.s11.eq.001
\begin{equation}
  G_1^*(x),\ldots,G_\rho^*(x)
\tag{1}
\end{equation}
代數整。

% unit: noether.p06.s11.par.002
如果多項式 (1) 特別是齊次形式，那麼對於使 (1) 消失的每一個特殊值系，所有代數整依賴的形式也同時消失；因而特別是 $x_1,\ldots,x_\rho$ 也消失。

% unit: noether.p06.s11.par.003
所以形式 (1) 除了
% unit: noether.p06.s11.eq.002
\[
  x_1=0,
  \ldots,
  x_\rho=0
\]
這個值系以外，不可能對任何值系同時消失；它們的結式必須不為零。

% unit: noether.p06.s11.par.004
反過來，這個條件也是充分的，並且按照下列輔助引理可推廣到非齊次多項式。

% unit: noether.p06.s11.lem.001
\emph{輔助引理：使 $x_1,\ldots,x_\rho$ 中每個任意多項式關於 $x_1,\ldots,x_\rho$ 中 $\rho$ 個多項式}
% unit: noether.p06.s11.eq.003
\[
  \tag{1} G_1^*(x),\ldots,G_\rho^*(x)
\]
\emph{代數整的一個充分條件，是這些多項式 (1) 的最高維項}
% unit: noether.p06.s11.eq.004
\[
  H_1^*(x),\ldots,H_\rho^*(x)
\]
\emph{關於 $x_1,\ldots,x_\rho$ 的結式不消失：}
% unit: noether.p06.s11.eq.005
\begin{equation}
  R_0=\left[
  \begin{array}{ccc}
    H_1^*&\cdots&H_\rho^*\\
    x_1&\cdots&x_\rho
  \end{array}
  \right]\ne0.
\tag{2}
\end{equation}
\emph{對於齊次形式 (1)，條件 (2) 是必要的。}\footnote{關於結式理論，參見 F. Mertens, \emph{Zur Theorie der Elimination}（I. 和 II. 部分），Sitzungsber. d. k. Ak. d. Wiss. Wien, Math.-naturw. Kl. Abt. IIa, Bd. 108 (1899)。J. König, \emph{Theorie d. algebr. Größen}, 第 VI 章中作了部分轉載。}

% unit: noether.p06.s11.par.005
為證明起見，我們構造多項式
% unit: noether.p06.s11.eq.006
\begin{equation}
\begin{gathered}
  \bar G_1^*(x)=G_1^*(x)-\lambda_1,
  \ldots,
  \bar G_\rho^*(x)=G_\rho^*(x)-\lambda_\rho,\\
  \bar\Phi^*(x)=\Phi^*(x)-\mu,
\end{gathered}
\tag{3}
\end{equation}
其中 $\lambda_1,\ldots,\lambda_\rho,\mu$ 表示不定元，$\Phi^*(x)$ 表示 $x_1,\ldots,x_\rho$ 中的任意多項式。令多項式 (3) 關於 $x_1,\ldots,x_\rho,1$\footnote{也就是說，如果藉助另一個不定元 $x_0$ 把多項式 (3) 齊次化，使得關於 $x$ 的次數得以保持，則這是關於 $x_1,\ldots,x_\rho,x_0$ 的齊次結式。} 的結式為
% unit: noether.p06.s11.eq.007
\begin{equation}
  R(\lambda,\mu)=\left[
  \begin{array}{cccc}
   \bar G_1^*&\cdots&\bar G_\rho^*&\bar\Phi^*\\
   x_1&\cdots&x_\rho&1
  \end{array}
  \right];
\tag{4}
\end{equation}
它是 $\lambda_1,\ldots,\lambda_\rho,\mu$ 的整有理函數。根據 F. Mertens，在按 $\mu$ 的冪展開 (4) 時，可以明確給出最高係數；\footnote{A. a. O. \S\ 3（定理 I 的證明），並且明確地在 \S\ 6（確定結式 $R$ 的所有項中含有給定冪積 $a_{\lambda\lambda}^{p_\lambda}a_{\mu\mu}^{p_\mu}\cdots a_{\sigma\sigma}^{p_\sigma}$ 作為因子的項的全體）。在我們的情形中，取 $p_\lambda=0$, $p_\mu=0,\ldots,p_\sigma=\tau$, $a_{\sigma\sigma}=(-\mu)$。König, 第 VI 章, \S\ 9 中有轉載。} 得到
% unit: noether.p06.s11.eq.008
\begin{equation}
  (-1)^\tau\cdot R(\lambda,\mu)
  =R_0^\sigma\cdot\mu^\tau+A_1(\lambda)\mu^{\tau-1}+\cdots+A_\tau(\lambda),
\tag{5}
\end{equation}
其中 $A_1(\lambda),\ldots,A_\tau(\lambda)$ 表示 $\lambda_1,\ldots,\lambda_\rho$ 以及多項式 (3) 的係數的整值整函數。這個展開首先表明，在假設 (2) $R_0\ne0$ 下，$R(\lambda,\mu)$ 也不可能恆等地消失。關於 $x_1,\ldots,x_\rho,\lambda_1,\ldots,\lambda_\rho,\mu$ 的恆等式
% unit: noether.p06.s11.eq.009
\[
  R(\lambda,
\mu)\equiv 0\pmod{\bar G_1^*(x),\ldots,
\bar G_\rho^*(x),\bar\Phi^*(x)},
\]
再通過代換
% unit: noether.p06.s11.eq.010
\[
  \lambda_i=G_i^*(x),\qquad \mu=\Phi^*(x)
\]
化為
% unit: noether.p06.s11.eq.011
\begin{equation}
\begin{aligned}
  R(G_i^*(x),\Phi^*(x))
  &=R_0^\sigma\cdot\Phi^{*\tau}(x)
    +A_1(G_i^*(x))\Phi^{*\tau-1}(x)+\cdots \\
  &\qquad +A_\tau(G_i(x))=0,
\end{aligned}
\tag{6}
\end{equation}
這就證明了輔助引理，因為 $R_0$ 按假設是 $\Omega$ 中的一個數。

% unit: noether.p06.s12.title.001
\section*{\S\ 12. 由多項式組成的正則系統 $\Ssys_{n\rho}$ 的整性基底。}

% unit: noether.p06.s12.par.001
\S\ 11 的輔助引理結合 \S\ 10 末尾的結論，立即給出下列事實：

% unit: noether.p06.s12.par.002
\emph{如果在 $\Gdom_{n\rho}$ 的某個映射域 $\Jdom^*_{\rho\rho}$ 中存在 $\rho$ 個多項式，使其最高維項的齊次結式不恆等消失，即 $R_0\ne0$，那麼 $\Gdom_{n\rho}$ 是第一種，因而是有限整性域，並且其整性基底為對合基。}

% unit: noether.p06.s12.par.003
條件 $R_0\ne0$ 現在藉助輔助引理給出整性基底存在性的第二個證明，這個證明本質上由 D. Hilbert 給出；\footnote{\emph{Über die vollen Invariantensysteme}, \S\ 2。在 Hilbert 的情形中，在不變量域的整性基底已由其他方式獲得存在性的假設下，要點只是通過一個域的代數整量的 Kronecker 基本系統來解釋這個基底。} 這個證明誠然不顯示對合基本身，但它統一地適用於所有具有所述性質的系統 $\Ssys_{n\rho}$。

% unit: noether.p06.s12.par.004
事實上，設 $\Jdom^*_{\rho\rho}$ 中滿足 $R_0\ne0$、因而代數無關的 $\rho$ 個多項式為
% unit: noether.p06.s12.eq.001
\begin{equation}
  G_1^*(x),\ldots,G_\rho^*(x).
\tag{1}
\end{equation}
按輔助引理，$x_1,\ldots,x_\rho$ 中的每個任意多項式，因而特別是最小包含 $\Ssys^*_{\rho\rho}$ 的域 $\Kfield^*_{\rho\rho}$ 中的每個多項式，都關於多項式 (1) 代數整。如果在 $\Ssys_{n\rho}$ 中相應的多項式為
% unit: noether.p06.s12.eq.002
\begin{equation}
  G_1(x),\ldots,G_\rho(x),
\tag{2}
\end{equation}
那麼由轉移原理，最小包含 $\Ssys_{n\rho}$ 的域 $\Kfield_{n\rho}$ 中的每個多項式，因而特別是 $\Ssys_{n\rho}$ 中的每個多項式，都關於多項式 (2) 代數整。反過來，$\Kfield_{n\rho}$ 中關於多項式 (2) 代數整的每個函數也是多項式。因此，如果把 $\Kfield_{n\rho}$（按 \S\ 4）看作 $\Omega(G_1,\ldots,G_\rho)$ 上的代數域，那麼 $\Kfield_{n\rho}$ 中多項式的全體，即相對整域 $\Gdom_{n\rho}$，與這個域中的代數整量相同。設 $G(x)$ 是 $\Kfield_{n\rho}$ 中把 (2) 補成有理基的一個多項式，$D(x)$ 是 $G(x)$ 關於 $\Omega(G_1,\ldots,G_\rho)$ 滿足的 $k$ 次不可約方程的判別式。那麼 $\Kfield_{n\rho}$ 中的每個多項式，特別是 $\Ssys_{n\rho}$ 中的每個多項式 $F(x)$，作為代數整量，都有表示
% unit: noether.p06.s12.eq.003
\begin{equation}
  F(x)=\frac{\Gamma_1(G_i)G(x)^{k-1}+\Gamma_2(G_i)G(x)^{k-2}+\cdots+\Gamma_k(G_i)}{D(x)}.
\tag{3}
\end{equation}

% unit: noether.p06.s12.par.005
令 $F(x)$ 遍歷 $\Ssys_{n\rho}$ 中所有多項式。把 Hilbert 關於模基的定理 I（Math. Ann. 36）應用於由相應函數
% unit: noether.p06.s12.eq.004
\[
  v_1\Gamma_1(G_i)+v_2\Gamma_2(G_i)+\cdots+v_k\Gamma_k(G_i)
\]
組成的系統，可知存在 $\Ssys_{n\rho}$ 中有限多個多項式
% unit: noether.p06.s12.eq.005
\begin{equation}
  F_1(x),F_2(x),\ldots,F_\lambda(x),
\tag{4}
\end{equation}
使得 $\Ssys_{n\rho}$ 中每個多項式都允許表示為
% unit: noether.p06.s12.eq.006
\begin{equation}
  F(x)=A_1(G_i)F_1(x)+A_2(G_i)F_2(x)+\cdots+A_\lambda(G_i)F_\lambda(x),
\tag{5}
\end{equation}
其中 $A_1,\ldots,A_\lambda$ 表示 $G_i(x)$ 的多項式。因此多項式 (2) 與 (4) 構成 $\Ssys_{n\rho}$ 的一個整性基底；也就是說：

% unit: noether.p06.s12.thm.008
\emph{定理 VIII：如果在多項式系統 $\Ssys_{n\rho}$ 的一個映射系統 $\Jdom^*_{\rho\rho}$ 中存在 $\rho$ 個多項式，使其最高維項的結式不恆等消失，即 $R_0\ne0$，那麼 $\Ssys_{n\rho}$ 具有一個整性基底。}

% unit: noether.p06.s12.par.006
這個定理還有一個輕微推廣。只要滿足 $R_0\ne0$ 的 $\rho$ 個多項式 (1) 屬於最小包含 $\Ssys^*_{\rho\rho}$ 的整性域 $\Jdom^*_{\rho\rho}$ 即可。事實上，導出定理 VIII 的推理表明，此時表示 (5) 仍然成立。而按定義，最小包含整性域 $\Jdom_{n\rho}$ 的每個多項式 $L(x)$ 都能由 $\Ssys_{n\rho}$ 中有限個（依賴於 $L(x)$）多項式作整有理表示；因此存在 $\Ssys_{n\rho}$ 中的 $\sigma$ 個多項式
% unit: noether.p06.s12.eq.007
\begin{equation}
  K_1(x),K_2(x),\ldots,K_\sigma(x),
\tag{6}
\end{equation}
使得多項式 $G_i(x)$ 成為 (6) 的整有理組合。於是 (6) 與 (4) 給出整性基底。引入下列術語：

% unit: noether.p06.s12.def.008
\emph{定義 VIII：如果在最小包含整性域的一個映射域 $\Jdom^*_{\rho\rho}$ 中，存在 $\rho$ 個多項式，使其最高維項的結式不恆等消失，即 $R_0\ne0$，則稱多項式系統 $\Ssys_{n\rho}$ 為正則。}

% unit: noether.p06.s12.par.007
於是有：

% unit: noether.p06.s12.thm.009
\emph{定理 IX：每個由多項式組成的正則系統 $\Ssys_{n\rho}$ 都具有一個整性基底。}\footnote{Hilbert 用一個例子說明了存在沒有整性基底的非正則系統（Gött. Nachr. 1891, p. 233）。下面這個例子還要簡單：
\[
  \Ssys_{22}=x_1,
  x_1x_2,
  x_1x_2^2,
  \ldots,
  x_1x_2^\nu,
  \ldots\quad\text{無窮地繼續。}
\]
這裡整性基底的存在等價於存在一個正整數 $\nu$，使恆等式
\[
  x_1x_2^{\nu+\sigma}=\sum c\cdot x_1^{i_0}(x_1x_2)^{i_1}(x_1x_2)^{i_2}\cdots(x_1x_2^\nu)^{i_\nu}
  \quad(\sigma=1,2,\ldots)
\]
成立。只需對 $\sigma=1$ 比較 $x_1,x_2$ 中的維數，就得到指數的兩個條件方程
\[
  1=i_0+i_1+\cdots+i_\nu,
  \qquad
  \nu+1=i_1+2i_2+\cdots+\nu\cdot i_\nu,
\]
它們顯然在非負整數中沒有解。因此 $\Ssys_{22}$ 沒有整性基底。}

% unit: noether.p06.s12.par.008
我們還像 \S\ 5 對對合所作的那樣，給出正則系統的一個幾何解釋。為此，如那裡 (7) 中一樣，考察方程組
% unit: noether.p06.s12.eq.008
\begin{equation}
  L^*(x)-L^*(\xi)=0,
\tag{7}
\end{equation}
其中 $L^*(x)$ 遍歷 $\Jdom^*_{\rho\rho}$ 中所有多項式。用 $x_0,\xi_0$ 齊次化後，令 (7) 變為齊次形式之間的方程組
% unit: noether.p06.s12.eq.009
\begin{equation}
  \xi_0^{m\cdot}M^*(x)-x_0^{m\cdot}M^*(\xi)=0,
\tag{8}
\end{equation}
並且若 $H^*(x)$ 表示 $L^*(x)$ 的最高維項，則有關係
% unit: noether.p06.s12.eq.010
\begin{equation}
  H^*(x)\equiv M^*(x)\pmod{x_0}.
\tag{9}
\end{equation}

% unit: noether.p06.s12.par.009
我們稱 $\Jdom^*_{\rho\rho}$ 的\emph{基本點}（比較 \S\ 5 的註釋）為 (8) 中與 $x$ 無關的解，也就是不同於與 $x$ 共軛的那些行的解。因此基本點就是不同於 $\xi_1=0,\ldots,\xi_\rho=0$ 的值系，並且對這些值系同時有
% unit: noether.p06.s12.eq.011
\[
  \xi_0^m=0,
  \qquad M^*(\xi)=0,
\]
或者由 (9) 得
% unit: noether.p06.s12.eq.012
\[
  \xi_0=0,
  \qquad H^*(\xi)=0.
\]

% unit: noether.p06.s12.par.010
如果現在存在 $\rho$ 個形式 $H^*(\xi)$，其 $R_0\ne0$，也就是沒有共同解系，那麼 $\Jdom^*_{\rho\rho}$ 不可能有基本點。特別地，如果 $\Ssys^*_{\rho\rho}$ 由齊次形式組成，那麼當 $R_0\ne0$ 時，$\Ssys^*_{\rho\rho}$ 本身也不可能有基本點，因為這樣的點同時也是 $\Jdom^*_{\rho\rho}$ 的基本點。但反過來也成立。為此考察由所有形式 $H^*(x)$ 組成的系統 $\Jdom(H^*)$；它又是一個整性域。

% unit: noether.p06.s12.par.011
由 Hilbert 的模基定理 I 和 Hilbert 的輔助引理（比較 \S\ 10 的引用）可知，存在 $\Jdom(H^*)$ 中的 $\rho$ 個形式，其消失導致 $\Jdom(H^*)$ 中所有形式的消失。如果 $\Jdom^*_{\rho\rho}$ 沒有基本點，那麼這 $\rho$ 個形式不可能同時消失；其結式 $R_0\ne0$。

% unit: noether.p06.s12.par.012
於是：

% unit: noether.p06.s12.par.013
\emph{多項式的正則系統 $\Ssys_{n\rho}$ 的特徵在於：其最小包含整性域的一個映射域 $\Jdom^*_{\rho\rho}$ 沒有基本點。對於由齊次形式組成的正則系統 $\Ssys^*_{\rho\rho}$，系統本身已經沒有基本點。}\footnote{前一注中給出的非正則系統 $\Ssys_{22}$ 具有基本點
\[
  \xi_1:\xi_2:\xi_0=0:1:0.
\]}

% unit: noether.p06.s13.title.001
\section*{\S\ 13. 第一種相對整域和正則系統的例子。}

% unit: noether.p06.s13.par.001
1. 作為第一種相對整域 $\Gdom_{nn}$ 的第一個例子，我們考察\emph{所有允許給定置換群 $G$ 的 $x_1,\ldots,x_n$ 中多項式的全體}，也就是一個 Lagrange 屬域的多項式全體。由於 $x_1,\ldots,x_n$ 藉助恆等式
% unit: noether.p06.s13.eq.001
\begin{equation}
  x_k^n-\sigma_1(x)x_k^{n-1}+\sigma_2(x)x_k^{n-2}+\cdots\pm\sigma_n(x)=0
  \qquad(k=1,2,
\cdots,n),
\tag{1}
\end{equation}
關於屬於 $\Gdom_{nn}$ 的初等對稱函數代數整，因而按定義 V 關於 $\Gdom_{nn}$ 代數整，所以 $\Gdom_{nn}$ 是第一種。此外，$\Gdom_{nn}$ 作為第一種域，由於其元素為齊次形式且 $n=\rho$，按 \S\ 10 的結論構成正則系統。事實上，由 (1) 可知，初等對稱函數的結式 $R_0$ 不可能恆等消失；按結式理論的計算規則，$R_0$ 算得為 $\pm1$。$\Gdom_{nn}$ 的\emph{對合形式}唯一確定——因為 $G_{nn}$ 為第一種且 $n=\rho$——它就是相應 Lagrange 屬域的對合形式。由於關於該域與 $x_1,\ldots,x_n$ 共軛的量由群 $G$ 的置換產生，這個對合形式因此由
% unit: noether.p06.s13.eq.002
\begin{equation}
\begin{aligned}
  \Psi(z,u)&=\prod_G
  (z-u_1x_{i_1}-u_2x_{i_2}-\cdots-u_nx_{i_n})\\
  &=z^i+\sum{}'G_{\alpha\alpha_1\ldots\alpha_n}(x)
  z^\alpha u_1^{\alpha_1}\cdots u_n^{\alpha_n},
\end{aligned}
\tag{2}
\end{equation}
給出，其中乘積擴及 $G$ 的所有置換。公式 (2) 表示“群 $G$ 的 Galois 形式”，因而有以下事實：

% unit: noether.p06.s13.par.002
\emph{所有允許置換群 $G$ 的 $x_1,\ldots,x_n$ 中多項式，都是該群的 Galois 形式的係數 $G_{\alpha\alpha_1\ldots\alpha_n}(x)$ 的整有理組合。}

% unit: noether.p06.s13.par.003
2. 作為\emph{正則系統}的例子，考察多項式系統 $\Ssys_{n1}$，即只含有一個代數無關多項式的系統。

% unit: noether.p06.s13.par.004
事實上，設 $\Ssys^*_{11}$ 是 $\Ssys_{n1}$ 的映射系統，並設
% unit: noether.p06.s13.eq.003
\[
  F^*(x)=a_0x^k+a_1x^{k-1}+\cdots+a_k\qquad(a_0\ne0)
\]
是 $\Ssys^*_{11}$ 中的一個多項式。那麼
% unit: noether.p06.s13.eq.004
\begin{equation}
  R_0=a_0\ne0\footnotemark.
\tag{3}
\end{equation}
\footnotetext{Mertens，前引處，\S\ 5。}
因此 $\Ssys_{n1}$ 按定義 VIII 是正則的，故：

% unit: noether.p06.s13.par.005
\emph{每個多項式系統 $\Ssys_{n1}$，特別是一個不定元的無限多個多項式的每個系統，都有一個整性基底。}\footnote{最簡單的無整性基底的非正則系統因此是系統 $\Ssys_{22}$。Hilbert 的例子以及 \S\ 12 註釋中的例子表明這樣的系統確實存在。}

% unit: noether.p06.s13.par.006
3. 現在把系統 $\Ssys_{n1}$ 特殊化為相對整域 $\Gdom_{n1}$；按 \S\ 12，作為正則系統，它們是第一種相對整域。因此其整性基底由一個對合基給出，而這個對合基同時也是最小包含域 $\Kfield_{n1}$ 的對合基。但按 \S\ 6，$\Kfield_{n1}$ 的對合基中的任意函數同時也是一個最小基——我們稱其為該域的 Lüroth 函數——並且對合基中的每個函數都通過分式線性變換依賴於這個 Lüroth 函數。在當前情形中，Lüroth 函數因為屬於 $\Gdom_{n1}$，所以是一個多項式；對合基中的每個其他多項式因此分式線性地，並且由於 (3) 還代數整地，亦即整且線性地，依賴於這個 Lüroth 函數 $L(x)$，後者從而成為整性基底。$\Gdom^*_{11}$ 的對合形式，在同時取 $(u_1=1)$ 時，為
% unit: noether.p06.s13.eq.005
\[
  \Psi^*(z)=L^*(z)-L^*(x),
\]
由此 $L(x)$ 又顯示為整性基底。因此：

% unit: noether.p06.s13.par.007
\emph{相對整域 $\Gdom_{n1}$ 是第一種；其整性基底由一個多項式組成，即最小包含域的 Lüroth 函數。}

% unit: noether.p06.s13.par.008
4. 作為相對整域 $\Gdom_{n1}$ 的例子，還特別提到 $s$ 個變量的二次型的整有理射影不變量。與前述相一致，眾所周知，它們可以表示為該形式的行列式 $|a_{ik}|$ 的整有理組合；而這個行列式同時也是相應不變量域的 Lüroth 函數。

% unit: noether.p06.s14.title.001
\section*{\S\ 14. 由整數多項式組成的系統。}

% unit: noether.p06.s14.par.001
現在要在整數性方面加強已經得到的基定理。因此以下把係數域 $\Omega$ 假定為一個（有限）代數數域；並以 $[\Omega]$ 表示 $\Omega$ 中代數整數的全體，我們將簡稱其為整數。所謂整數多項式，是指係數屬於 $[\Omega]$ 的多項式；從現在起，$F(x),G(x),\ldots$ 只表示整數多項式。

% unit: noether.p06.s14.par.002
我們按照 \S\ 7 的類比，考察由整數多項式組成的各種系統。

% unit: noether.p06.s14.par.003
一個由整數多項式組成的任意系統稱為整數系統 $\Ssys_{n\rho}$。

% unit: noether.p06.s14.par.004
整數線性族 $\Lsys_{n\rho}$ 是一個整數系統，它在含有 $F(x)$ 和 $F_1(x)$ 的同時，總也含有 $c_1F(x)+c_2F_1(x)$，其中 $c_1,c_2$ 是 $[\Omega]$ 中任意整數。

% unit: noether.p06.s14.par.005
整數整性域 $\Jdom_{n\rho}$ 是一個整數線性族，它在含有 $F(x)$ 和 $G(x)$ 的同時，總也含有 $F(x)\cdot G(x)$。

% unit: noether.p06.s14.par.006
整數相對整域 $\Gdom_{n\rho}$ 是一個整性域，它在含有 $F(x)$ 和 $G(x)$（其中 $G(x)\ne0$）的同時，當且僅當商 $F(x):G(x)$ 是整數多項式時，含有這個商。

% unit: noether.p06.s14.par.007
包含系統 $\Ssys_{n\rho}$ 的最小整數整性域 $\Jdom_{n\rho}$ 由所有這樣的多項式 $H(x)$ 組成：它們是 $\Ssys_{n\rho}$ 中有限個多項式的整有理、整數係數組合；

% unit: noether.p06.s14.par.008
包含它的最小整數相對整域 $\Gdom_{n\rho}$ 由所有這樣的整數多項式 $H(x)$ 組成：它們是 $\Ssys_{n\rho}$ 中有限個多項式的有理組合。

% unit: noether.p06.s14.par.009
現在轉向基定理的轉移。關於有理基，需要注意的是，僅依賴“有理性”的域、最小包含域以及有理基這些概念沒有發生變化。此外，整數系統 $\Ssys_{n\rho}$ 的映射系統 $\Ssys^*_{\rho\rho}$ 可以以任意多種方式再選成整數系統。由於對於整數線性族，導出生成元個數限制的推理也仍然有效，得到：

% unit: noether.p06.s14.thm.005p
\emph{定理 V$'$：每個整數系統 $\Ssys_{n\rho}$ 都有有理基；特別地，整數線性族 $\Lsys_{n\rho}$ 的有理基可以選成至多含有 $\rho+1$ 個多項式。}

% unit: noether.p06.s14.par.010
現在考察多項式整性域的對合基（\S\ 8）的類似物。為此，需要把關於\emph{量行的對稱函數}的定理在\emph{整數上的整性}方面擴展為下面的輔助引理。

% unit: noether.p06.s14.lem.001
\emph{$n$ 行量 $(yz\cdots t)$ 的整整數對稱函數}
% unit: noether.p06.s14.eq.001
\begin{equation}
\begin{array}{cccc}
  y_1&z_1&\cdots&t_1\\
  y_2&z_2&\cdots&t_2\\
  \vdots&\vdots&&\vdots\\
  y_n&z_n&\cdots&t_n
\end{array}
\tag{1}
\end{equation}
\emph{是其中有限個函數的整整數函數；確切地說，是初等對稱函數}
% unit: noether.p06.s14.eq.002
\begin{equation}
\begin{gathered}
  \sigma_1(y),\sigma_2(y),\ldots,\sigma_n(y);\quad
  \sigma_1(z),\sigma_2(z),\ldots,\sigma_n(z);\quad\ldots;\\
  \sigma_1(t),\ldots,\sigma_n(t);
\end{gathered}
\tag{2}
\end{equation}
\emph{以及函數}
% unit: noether.p06.s14.eq.003
\begin{equation}
  \chi_1(y,z,\ldots,t),\ldots,\chi_k(y,z,\ldots,t)
\tag{3}
\end{equation}
\emph{的整整數函數，而這些函數關於函數 (2) 代數整且整數。}

% unit: noether.p06.s14.par.011
事實上，$n$ 個量 $y_1,\ldots,y_n$ 關於 $\sigma_1(y),\ldots,\sigma_n(y)$ 在整數上代數整且整數；對於 $z_1,\ldots,z_n$ 相對於 $\sigma_1(z),\ldots,\sigma_n(z)$，以及諸如此類，也有相應斷言。因此，量行 (1) 的對稱函數域在由函數 (2) 確定的域上形成一個代數域，並且這個域中代數整且整數的量，正好就是量行 (1) 的整整數對稱函數。如果取 (3) 為 Kronecker 的基本系統，輔助引理即得證。

% unit: noether.p06.s14.par.012
現在令 $\Jdom_{n\rho}$ 為整數整性域，$\Jdom^*_{\rho\rho}$ 為整數映射域，並令
% unit: noether.p06.s14.eq.004
\[
  \Psi^*(z,u)=G^*(x)\cdot z^i+
  \sum{}'G^*_{\alpha\alpha_1\ldots\alpha_\rho}(x)\cdot
  z^\alpha u_1^{\alpha_1}\cdots u_\rho^{\alpha_\rho}
\]
為 $\Jdom^*_{\rho\rho}$ 的一個對合形式；這裡 $G^*(x)$ 也可以為零維，也就是 $[\Omega]$ 中的一個數。與函數 (2) 相應的初等對稱函數於是由若干商
% unit: noether.p06.s14.eq.005
\begin{equation}
  G^*_{\alpha\alpha_1\ldots\alpha_\rho}(x):G^*(x)
\tag{4}
\end{equation}
給出。由於按輔助引理，函數 (3) 關於 (2) 在整數上代數整且整數，所以與 (3) 對應的 $\Kfield^*_{\rho\rho}$ 中的函數，由 Gauss 定理對帶代數整數係數多項式的推廣給出為
% unit: noether.p06.s14.eq.006
\begin{equation}
  K_\lambda^*(x):(G^*(x))^\sigma,
\tag{5}
\end{equation}
其中 $K_\lambda^*(x)$ 是 $\Kfield^*_{\rho\rho}$ 中的整數多項式。因此，在相對整域情形中，$K_\lambda^*(x)$ 屬於該域；一般地，它們是 $\Jdom^*_{\rho\rho}$ 中多項式的商。

% unit: noether.p06.s14.par.013
如果 $F^*(x)$ 是 $\Jdom^*_{\rho\rho}$ 中的任意整數多項式，則有
% unit: noether.p06.s14.eq.007
\[
  F^*(x)=\frac1i\Bigl(F^*(x)+F^*(x^{(1)})+\cdots+F^*(x^{(i-1)})\Bigr),
\]
所以 $i\cdot F^*(x)$ 成為量行 $x,x^{(1)},\ldots,x^{(i-1)}$ 的整整數對稱函數。

% unit: noether.p06.s14.par.014
於是，由輔助引理和轉移原理得到：

% unit: noether.p06.s14.thm.006p
\emph{定理 VI$'$：對於整數整性域 $\Jdom_{n\rho}$，每一個對合形式確定一個特殊有理基，使得 $\Jdom_{n\rho}$ 中每個多項式 $F(x)$ 都允許表示為}
% unit: noether.p06.s14.eq.008
\[
  F(x)=\frac{\Gamma(H(x),H_1(x),\ldots,H_\sigma(x))}{i\cdot H(x)^k},
\]
\emph{其中 $\Gamma$ 表示一個整整數函數。對於整數相對整域 $\Gdom_{n\rho}$，若其映射域為整數相對整域 $\Gdom^*_{\rho\rho}$，則分母 $H(x)$ 由相應對合形式的最高係數給出。}

% unit: noether.p06.s15.title.001
\section*{\S\ 15. 第一種整數相對整域和正則系統。}

% unit: noether.p06.s15.par.001
對於整數系統，整數整性基底取代整性基底；也就是說，它是系統中有限個整數多項式，使得系統中的每個整數多項式都能表示為這有限個多項式的整整數係數組合。

% unit: noether.p06.s15.par.002
關於整數整性基底的定理本質上與前述定理平行，只有證明需要輕微修改。

% unit: noether.p06.s15.par.003
首先，定義 V（\S\ 9）直接轉移。

% unit: noether.p06.s15.def.005p
\emph{定義 V$'$：若一個量 $\xi$ 滿足某個方程，其最高係數是 $[\Omega]$ 中的單位，其餘多項式來自整數相對整域 $\Gdom_{n\rho}$，則稱 $\xi$ 關於 $\Gdom_{n\rho}$ 代數整且整數。}

% unit: noether.p06.s15.par.004
Gauss 定理對帶代數整數係數的多項式的推廣表明，如 \S\ 9 中一樣，這一定義可以從不可約方程得到。

% unit: noether.p06.s15.par.005
由定義 V$'$ 進一步得到定義 VII 的轉移。

% unit: noether.p06.s15.def.007p
\emph{定義 VII$'$：若整數相對整域 $\Gdom_{n\rho}$ 作為映射域具有一個相對整域 $\Gdom^*_{\rho\rho}$，使得 $x_1,\ldots,x_\rho$ 中每個任意整數多項式都關於 $\Gdom^*_{\rho\rho}$ 代數整且整數，則稱它為第一種。}

% unit: noether.p06.s15.par.006
為證明這些第一種域的有限性，我們首先注意到，如 \S\ 10 中一樣，由定義推出，對合形式的最高係數是 $[\Omega]$ 中的一個單位。

% unit: noether.p06.s15.par.007
於是按定理 VI$'$（\S\ 14），$\Gdom_{n\rho}$ 中每個多項式 $F(x)$ 都允許表示為
% unit: noether.p06.s15.eq.001
\[
  F(x)=\frac{\Gamma(H_1(x),\ldots,H_\sigma(x))}{i}.
\]
在這裡，把 Hilbert 關於整數模基的定理 II（Math. Ann. 36）——原本為整有理數係數所陳述——在其推廣到代數整數多項式的形式下，應用於由整數多項式 $\Gamma(H_1,\ldots,H_\sigma)$ 形成的系統，\footnote{若 $w_1,\ldots,w_k$ 表示 $\Omega$ 中代數整數的一個基本系統，並設
\[
  \Gamma(H_1,\ldots,H_\sigma)=\Gamma_1(H)w_1+\cdots+\Gamma_k(H)w_k,
\]
則 $\Gamma_1(H),\ldots,\Gamma_k(H)$ 具有整有理數係數。把定理 II 應用於由形式 $v_1\Gamma_1(H)+\cdots+v_k\Gamma_k(H)$ 構成的系統，就直接說明了這個推廣的有效性。} 即得整數整性基底的存在；也就是說：

% unit: noether.p06.s15.thm.007p
\emph{定理 VII$'$：每個第一種整數相對整域都具有一個整數整性基底。}

% unit: noether.p06.s15.par.008
再轉移正則系統的定義：

% unit: noether.p06.s15.def.008p
\emph{定義 VIII$'$：若在最小包含整數整性域的一個映射域 $\Jdom^*_{\rho\rho}$ 中存在 $\rho$ 個多項式，使得最高維項的結式是 $[\Omega]$ 中的單位，即 $R_0=\eps$，則稱整數系統 $\Ssys_{n\rho}$ 為整數正則。}

% unit: noether.p06.s15.par.009
為轉移整性基底的存在性證明，我們注意到 \S\ 11 的輔助引理具有如下更強形式；這由 \S\ 11 中的方程 (5)、(6) 以及 $A_1(\lambda),\ldots,A_\tau(\lambda)$ 是 $\lambda_1,\ldots,\lambda_\rho$ 的整整數函數這一事實得出：

% unit: noether.p06.s15.par.010
\emph{使 $x_1,\ldots,x_\rho$ 中每個任意整數多項式關於 $\rho$ 個整數多項式代數整且整數的一個充分條件，是這些多項式最高維項的結式為 $[\Omega]$ 中的單位，即 $R_0=\eps$。}

% unit: noether.p06.s15.par.011
由這個輔助引理，完全像 \S\ 12 中一樣，對於整數正則系統中的每個多項式 $F(x)$，得到 \S\ 12 的表示 (3)，其中現在 $\Gamma_1(G_i),\ldots,\Gamma_k(G_i)$ 是 $G_i$ 中的整數多項式。進一步，如 \S\ 12 中一樣，把 Hilbert 的定理 II（Math. Ann. 36）以其推廣到代數整數多項式的形式應用，並結合最小包含整數整性域的定義，得到整數整性基底的存在；也就是說：

% unit: noether.p06.s15.thm.009p
\emph{定理 IX$'$：每個整數正則系統都具有一個整數整性基底。}

% unit: noether.p06.s15.par.012
作為第一種整數相對整域的例子，按 \S\ 13 中例 1 的類比，提及一個 Lagrange 屬域的整數多項式全體 $\Gdom_{nn}$。$\Gdom_{nn}$ 為整數第一種，來自恆等式
% unit: noether.p06.s15.eq.002
\[
  x_k^n-\sigma_1(x)x_k^{n-1}+\cdots\pm\sigma_n(x)=0
  \qquad(k=1,2,
\cdots,n);
\]
由於初等對稱函數的結式具有值 $\pm1$，$\Gdom_{nn}$ 也是一個整數正則系統。

% unit: noether.p06.s15.par.013
另一個例子由 \S\ 14 的輔助引理給出，即 $n$ 行量的整整數對稱函數的全體。

% unit: noether.p06.closing.001
Erlangen，1914 年 5 月。
%
\clearpage
%
% unit: noether.p07.title.001
\pdfbookmark[0]{7. 有限群不變量的有限性定理}{paper-07}
\section*{7. 有限群不變量的有限性定理}
\begin{center}
Math. Ann. 77 (1916), pp. 89--92
\end{center}

% unit: noether.p07.par.001
以下將給出一個關於有限群不變量有限性的完全初等的證明；它只依賴於對稱函數理論，同時還給出完整系統的實際指定。通常的證明則依靠 Hilbert 的模基定理（Ann. 36），因而只是存在性證明。\footnote{參見例如 Weber, \emph{Lehrbuch der Algebra} (2nd ed.), vol. 2, \S\ 57.}

% unit: noether.p07.par.002
設有限群 $\Hgrp$ 由 $h$ 個線性變換（行列式不為零）$A_1\cdots A_h$ 構成，其中 $A_k$ 表示線性變換
% unit: noether.p07.eq.001
\[
  x_1^{(k)}=\sum_{\nu=1}^{n}a_{1\nu}^{(k)}x_\nu,\cdots,
  x_n^{(k)}=\sum_{\nu=1}^{n}a_{n\nu}^{(k)}x_\nu,
\]
或簡寫為 $(x^{(k)})=A_k(x)$。因此，群 $\Hgrp$ 把具有元素 $x_1\cdots x_n$ 的行 $(x)$ 變為具有元素 $x_1^{(k)}\cdots x_n^{(k)}$ 的各行 $(x^{(k)})$。由於恆等變換必須包含在 $A_1\cdots A_h$ 之中，行 $(x)$ 本身也包含在這些行 $(x^{(k)})$ 之中。所謂該群的整有理（絕對）不變量，是指 $x_1\cdots x_n$ 的這樣的整有理函數：在施用 $A_1\cdots A_h$ 時恆等地保持不變。因此，對於這樣的不變量 $f(x)$，有
% unit: noether.p07.eq.002
\begin{equation}
  f(x)=f(x^{(1)})=\cdots=f(x^{(k)})=
  \frac1h\cdot\sum_{k=1}^{h}f(x^{(k)}).
\tag{1}
\end{equation}

% unit: noether.p07.sec.001
\subsection*{1.}
% unit: noether.p07.par.003
公式 (1) 表明，$f(x)$ 是諸數量行 $(x^{(1)})\cdots(x^{(h)})$ 的整有理對稱函數；而且，由於每一項 $f(x^{(k)})$ 只含有一個行 $(x^{(k)})$，這裡涉及的是最簡單的情形，通常稱為單式情形。由關於數量行的對稱函數的已知定理，\footnote{參見 2. 的注。}$f(x)$ 因而能由這些行的初等對稱函數整有理地表示，也就是說，由“Galois 預解式”的係數 $G_{\alpha\alpha_1\cdots\alpha_n}(x)$ 表示：
% unit: noether.p07.eq.003
\[
\begin{aligned}
  \Phi(z,u)
  &=\prod_{k=1}^{h}\bigl(z+u_1x_1^{(k)}+\cdots+u_nx_n^{(k)}\bigr) \\
  &=z^h+
    \sum G_{\alpha\alpha_1\cdots\alpha_n}(x)
    z^{\alpha}u_1^{\alpha_1}\cdots u_n^{\alpha_n},
\end{aligned}
\qquad
\left(\begin{array}{c}
(\alpha+\alpha_1+
\cdots+
\alpha_n=h)\\
\alpha\ne h
\end{array}\right),
\]
其中 $G_{\alpha\alpha_1\cdots\alpha_n}(x)$ 是關於 $x$ 的次數為 $\alpha_1+\cdots+\alpha_n$ 的不變量。於是得到證明：

% unit: noether.p07.thm.001
\emph{Galois 預解式的係數 $G_{\alpha\alpha_1\cdots\alpha_n}(x)$ 構成該群的不變量完整系統，使每一個不變量都能由這有限多個不變量整有理地表示。}

% unit: noether.p07.sec.002
\subsection*{2.}
% unit: noether.p07.par.004
從 (1) 出發，還可以作如下更為初等的考察，\footnote{這等同於在上述單式情形中證明關於數量行對稱函數的定理。}它同時導向第二個完整系統。設
% unit: noether.p07.eq.004
\[
  f(x)=a+b x_1^{\mu_1}\cdots x_n^{\mu_n}
       +\cdots+c x_1^{\nu_1}\cdots x_n^{\nu_n},
\]
其中 $a,b,\cdots,c$ 表示常數。由 (1) 得
% unit: noether.p07.eq.005
\[
\begin{aligned}
 h\cdot f(x)
 &=h\cdot a+b\cdot\sum_{k=1}^{h}x_1^{(k)^{\mu_1}}\cdots x_n^{(k)^{\mu_n}}
     +\cdots
     +c\cdot\sum_{k=1}^{h}x_1^{(k)^{\nu_1}}\cdots x_n^{(k)^{\nu_n}} \\
 &=h\cdot a+b\cdot J_{\mu_1\cdots\mu_n}+\cdots+c\cdot J_{\nu_1\cdots\nu_n}.
\end{aligned}
\]
因此，每個不變量都可以由特殊不變量
% unit: noether.p07.eq.006
\[
  J_{\mu_1\cdots\mu_n}
  =\sum_{k=1}^{h}x_1^{(k)^{\mu_1}}\cdots x_n^{(k)^{\mu_n}}
\]
整線性地合成，所以只需對這些特殊不變量證明斷言。但除去一個數值因子外，$J_{\mu_1\cdots\mu_n}$ 是表達式
% unit: noether.p07.eq.007
\[
  S_\mu=\sum_{k=1}^{h}
  \bigl(u_1x_1^{(k)}+\cdots+u_nx_n^{(k)}\bigr)^{\mu}
\]
中 $u_1^{\mu_1}\cdots u_n^{\mu_n}$ 的係數，其中 $\mu=\mu_1+\cdots+\mu_n$；而該表達式是 $h$ 個線性形式
% unit: noether.p07.eq.008
\[
  \xi_1=u_1x_1^{(1)}+\cdots+u_nx_n^{(1)},\cdots,
  \xi_h=u_1x_1^{(h)}+\cdots+u_nx_n^{(h)}
\]
的第 $\mu$ 個冪和。眾所周知，無窮多個冪和 $S_\mu$ 都是
% unit: noether.p07.eq.009
\[
  S_1,S_2,\cdots,S_h
\]
的整有理組合，其係數由不變量
% unit: noether.p07.eq.010
\[
  J_{\mu_1\cdots\mu_n}\qquad(\mu_1+\cdots+\mu_n\leqq h)
\]
給出。於是得到第二個完整系統：

% unit: noether.p07.thm.002
\emph{該群的不變量完整系統由所有滿足 $\mu_1+\cdots+\mu_n\leqq h$ 的不變量 $J_{\mu_1\cdots\mu_n}$ 給出，其中 $h$ 表示群的階。}

% unit: noether.p07.par.005
與 1. 的聯繫由如下觀察建立：冪和 $S_1\cdots S_h$ 可以用 $\xi_1\cdots\xi_h$ 的初等對稱函數替代，這就導向那裡給出的 $G_{\alpha\alpha_1\cdots\alpha_n}(x)$ 的完整系統。兩個結果都表明，所有不變量都可以由那些關於 $x$ 的次數不超過群階 $h$ 的不變量整有理地表示。

% unit: noether.p07.sec.003
\subsection*{3.}
% unit: noether.p07.par.006
由剛剛證明的結果可以推出關於有理表示的結論。每一個有理絕對不變量都可以表示為兩個整有理絕對不變量的商，這兩個不變量不必互素；這是通過把分子和分母乘以分母關於該群的共軛即可直接看出的。因此，每一個有理絕對不變量都可以由 Galois 預解式 $\Phi(z,u)$ 的係數 $G_{\alpha\alpha_1\cdots\alpha_n}(x)$ 有理地表示。這個定理也可以不經由有限性定理而以多種方式容易地證明；它在 Weber II, \S\ 58 中已經表述。\footnote{然而，那裡給出的證明並不可靠。它只表明公式 (7) 中的函數 $\Psi(t)$ 以不變量為係數，卻沒有表明這些不變量能由 $\Phi(t)$ 的係數有理表示。若在用 $\xi_k$ 表示 $x_i^{(k)}$ 時，不使用 Lagrange 插值公式，而使用一個已知的微分方法，則可避免這個缺口。這就是對恆等式 $\Phi(-\xi_k,u)=0$ 關於所有 $u$ 微分得到的關係：
\[
 \left[\frac{\partial\Phi}{\partial u_i}
      -x_i^{(k)}\cdot\frac{\partial\Phi}{\partial z}\right]_{z=-\xi_k}=0.
\]
相應地，對於每一個有理函數 $\omega(x)$，Weber 中的公式 (7) 和 (8) 應由下式代替：
\[
 \omega(x_1^{(k)}\cdots x_n^{(k)})=
 \left.
 \omega\!\left(
   \frac{\frac{\partial\Phi}{\partial u_1}}{\frac{\partial\Phi}{\partial z}}\cdots,
   \frac{\frac{\partial\Phi}{\partial u_n}}{\frac{\partial\Phi}{\partial z}}
 \right)\right|_{z=-\xi_k},
\]
它只含有 $\Phi(z,u)$ 的係數；當 $\omega(x)$ 為不變量時，對所有 $k$ 求和即給出所要求的表示。}

% unit: noether.p07.sec.004
\subsection*{4.}
% unit: noether.p07.par.007
最後要指出，這裡推出的結果已隱含於我的論文“有理函數的域與系統”之中；\footnote{Math. Ann. 76, p. 161 (1915).}更準確地說，1. 中給出的完整系統包含在定理 VI 和 VII 中，而獨立於這個有限性定理的有理表示證明包含在定理 III 中。\footnote{對於相對不變量，定理 VIII 和 IX 包含了一個有限性證明，它同樣建立在不同於通常證明的基礎上，但也只是存在性證明；原因在於相對不變量並不構成一個域。}不過，在目前這個特殊情形中，證明過程和結果相對於一般理論都大為簡化。我之所以想到把一般研究應用於有限群不變量，是由於與 E. Fischer 的談話；2. 中給出的證明在內容上也來自他，3. 的註釋亦然；在一般情形中，那裡出現的完整系統並不能有理定義。

% unit: noether.p07.close.001
Erlangen, 1915 年 5 月。

\clearpage
% unit: noether.p08.title.001
\pdfbookmark[0]{8. 關於任意多個基本形式的系統的不變量的整有理表示}{paper-08}
\section*{8. 關於任意多個基本形式的系統的不變量的整有理表示}
\begin{center}
Math. Ann. 77 (1916), pp. 93--102
\end{center}

% unit: noether.p08.par.001
在他為 Schwarz 紀念文集所寫的《任意多個基本形式的系統的不變量》一文的結尾，D. Hilbert 提出猜想：這些不變量可以由系統 $(J,PJ)$ 的有限多個不變量整有理地表示，其中 $J$ 表示那些線性無關的基本形式的完整不變量系統，而不變量 $PJ$ 則由 $J$ 通過極化過程產生。

% unit: noether.p08.par.002
下面我在 I 中說明，這個猜想確實成立，而且它只是如下歸約定理的簡單推論：任意多個、每個含 $n$ 個變量的行的任一形式，都可表示為若干極化形式的和；這些極化形式只含有 $n$ 個固定行，並且它們本身由原始形式通過極化過程導出。這個歸約定理是 F. Mertens 首先表述的一個特殊情形，屬於 A. Capelli、F. Mertens 和 J. Deruyts 給出的 Clebsch--Gordan 級數展開推廣；該推廣適用於任意多個、每個含 $n$ 個變量的行的形式，並通過對微分過程的形式變換而證明。\footnote{F. Mertens: ``Über eine Formel der Determinantentheorie.'' (公式 5), Sitzb. d. Ak. d. Wiss. Wien, vol. 91, II. Abt. (1885), 以及更詳盡的（含應用）``Über ganze Funktionen von $m$ Systemen von je $n$ Unbestimmten''. Monatsh. f. Math. u. Phys. 第 4 年 (1893)。關於 Capelli 自 1882 年以來的若干證明，參見例如 Math. Ann. 37，或其影印講義 ``Lezioni sulla teoria delle forme algebriche'' (1902)。J. Deruyts: \emph{Essai d'une théorie générale des formes algébriques}, Mém. d. l. Soc. Roy. des Sciences de Liège (2) 17 (1892).}

% unit: noether.p08.par.003
在 II 中，我還給出歸約定理的一個更概念化的證明，把這個問題看作線性形式族的等價問題。這種觀點，以及所用論證中的一個實質部分，取自 E. Fischer 的論文“論代數模系統”\footnote{E. Fischer: Über algebraische Modulsysteme und lineare homogene partielle Differentialgleichungen mit konstanten Koeffizienten, \S\ 1--4, J. f. M. 140 (1911).}以及與之相連的 E. Fischer 未發表定理；他友好地允許我使用這些結果。

% unit: noether.p08.par.004
最後，在 III 中我還說明，對應的論證如何導向最一般的級數展開。

% unit: noether.p08.sec.001
\subsection*{I.}
% unit: noether.p08.par.005
設 $\theta$ 是在每一個 $N>n$ 行
% unit: noether.p08.eq.001
\[
  A_1,A_2,\ldots,A_N,
\]
中都齊次的形式，其中一般地，行 $A_k$ 由 $n$ 個元素
% unit: noether.p08.eq.002
\[
  A_k^{(1)},A_k^{(2)},\ldots,A_k^{(n)}
\]
組成。$\theta$ 的係數可以是不定元，也可以是某個給定有理性域中的量。又令 $P$ 表示極化過程
% unit: noether.p08.eq.003
\[
  P_{hk}=\left(A_h\frac{\partial}{\partial A_k}\right)
  =A_h^{(1)}\frac{\partial}{\partial A_k^{(1)}}+
   A_h^{(2)}\frac{\partial}{\partial A_k^{(2)}}+
   \cdots+
   A_h^{(n)}\frac{\partial}{\partial A_k^{(n)}}
\]
的一個整有理函數，其係數為有理整數；其中乘積 $P_{hk}P_{ij}\theta$ 應理解為把 $P_{hk}$ 施加於 $P_{ij}\theta$。

% unit: noether.p08.par.006
剛才提到的歸約定理由恆等式
% unit: noether.p08.eq.004
\begin{equation}
  \theta=\sum PZ,
\tag{1}
\end{equation}
給出，其中形式 $Z$ 只含有行
% unit: noether.p08.eq.005
\[
  A_1,A_2,\ldots,A_n
\]
並且由 $\theta$ 通過極化過程 $P$ 導出。

% unit: noether.p08.par.007
現在考慮基本形式系統 $(F)$，它由 $N$ 個同階且變量數相同的形式
% unit: noether.p08.eq.006
\[
  F_1,F_2,\ldots,F_N
\]
組成，其係數分別取作 $N$ 個行
% unit: noether.p08.eq.007
\[
  A_1,A_2,\ldots,A_N.
\]
令 $(J)$ 表示 $F_1,\ldots,F_n$ 的完整系統，它由有限多個不變量組成，並且使每個 $F_1,\ldots,F_n$ 的不變量都可由 $(J)$ 中的不變量整有理地表示。要證明的 Hilbert 猜想便是：對於 $F_1,\ldots,F_N$ 的同時不變量 $S$，一個完整系統由有限多個不變量 $(J,PJ)$ 給出，其中 $PJ$ 表示所有由 $J$ 通過極化過程 $P$ 導出的不變量。

% unit: noether.p08.par.008
事實上，$(F)$ 的每一個帶有有理整數係數的同時不變量 $S$，都是在 $A_1,\ldots,A_N$ 的每一個行中齊次的形式，且係數為有理整數，因此可對它應用恆等式 (1)。由於對 $S$ 施加極化過程 $P$ 已知仍產生不變量，所以形式 $Z$ 也都是不變量；並且由於它們只含有行 $A_1,\ldots,A_n$，它們是 $F_1,\ldots,F_n$ 的同時不變量；於是按假設，它們是 $(J)$ 中不變量的整有理函數。因此恆等式 (1) 變成
% unit: noether.p08.eq.008
\[
  S=\sum PY,
\]
其中 $Y$ 表示 $(J)$ 中不變量的冪積。可是由 $P$ 的定義可知，實際在 $Y$ 上施行過程 $P$，就是有限次相繼施行簡單的極化操作 $P_{hk}$；按乘積的微分規則，這導致 $(J)$ 與 $(PJ)$ 中不變量的乘積之和。對來自 $(J)$ 與 $(PJ)$ 的乘積施加 $P_{hk}$ 亦然；因此 $PY$ 也只給出 $(J,PJ)$ 的整有理組合。\emph{由此證明了所有不變量 $S$ 都可由 $(J,PJ)$ 整有理地表示。}

% unit: noether.p08.sec.002
\subsection*{II.}
% unit: noether.p08.par.009
在證明歸約定理以前，先回憶關於線性形式族的一些事實。所謂 \emph{$K$ 上的線性形式族}，是指同一維數的形式系統，它是完備的，意即只要含有 $\theta_1$ 和 $\theta_2$，就總含有 $c_1\theta_1+c_2\theta_2$，其中 $c_1,c_2$ 是給定有理性域 $K$ 中的量，並且 $K$ 包含各 $\theta$ 的係數域。在這些形式中，總有有限個（設為 $\rho$ 個）在 $K$ 上線性無關，而任意 $\rho+1$ 個形式都滿足一個係數取自 $K$ 的線性關係；這個形式族關於 $K$ 具有\emph{秩 $\rho$}。\footnote{若去掉同維數的限制，或者去掉 $K$ 包含 $\theta$ 的係數域這一限制，就得到更一般的線性形式族；它們的秩不必是有限的。}我們將使用一個顯然的事實：若 $\mathfrak L$ 與 $\mathfrak T$ 是 $K$ 上兩個同秩 $\rho$ 的線性形式族，並且 $\mathfrak T$ 是 $\mathfrak L$ 的子族，則 $\mathfrak L$ 與 $\mathfrak T$ 相同（等價）。

% unit: noether.p08.par.010
有了這些，我們轉向歸約定理。I 中的恆等式
% unit: noether.p08.eq.009
\[
  \theta=\sum PZ
\]
在對兩邊施加同一個極化過程 $P$ 時，變為 $P\theta$ 的類似恆等式；因此歸約定理同時給出了 $\theta$ 的所有極化形式用特殊極化形式 $PZ$ 的和來表示。反過來，這些特殊極化形式當然包含在所有極化形式中，所以歸約定理斷言這兩個線性形式族等價；特別也斷言其中每個形式在各個行中的維數均與 $\theta$ 相同的子族等價。為證明這種等價，我們插入由形式 $Z$ 決定的中間形式族。為簡單起見，先給出三行二元的情形 $(N=3,n=2)$ 的證明，然後簡要指出完全平行的一般證明。

% unit: noether.p08.par.011
於是令 $\theta(x,y,z)$ 在行
% unit: noether.p08.eq.010
\[
  x_1x_2;\quad y_1y_2;\quad z_1z_2;
\]
中分別具有維數 $\alpha,\beta,p$；並且先把 $\theta$ 的係數看作不定元 $a$（或有理數）。我們考慮以下三個線性形式族。

% unit: noether.p08.par.012
1. 形式族 $\mathfrak L$，由所有從 $\theta$ 通過極化過程 $P$ 得到、且在各個行 $x,y,z$ 中具有與 $\theta$ 相同維數的形式 $f(x,y,z)$ 構成；特別地，$\mathfrak L$ 包含 $\theta$。若把 $f(x,y,z)$ 看作 $x,y,z$ 以及不定元 $a$ 中的形式，則係數域是有理數域 $R$；$K$ 也必須取為 $R$，因為只有當 $c_1,c_2$ 為有理數時，$c_1P_1+c_2P_2$ 才仍是一個過程 $P$；設 $\mathfrak L$ 關於 $R$ 的秩為 $\rho$。

% unit: noether.p08.par.013
2. 形式族 $\mathfrak S$，由所有按
% unit: noether.p08.eq.011
\[
  g(\lambda;x,y)=f(x,y,\lambda_1x+\lambda_2y)
\]
定義的形式 $g(\lambda;x,y)$ 構成，或用極化過程表示為
% unit: noether.p08.eq.012
\[
\begin{aligned}
  g(\lambda;xy)
  &=\frac1{p!}\left[(\lambda_1x+\lambda_2y)\frac{\partial}{\partial z}\right]^p f(x,y,z) \\
  &=g_0(xy)\lambda_1^p+g_1(x,y)\lambda_1^{p-1}\lambda_2+\cdots+g_p(xy)\lambda_2^p,
\end{aligned}
\]
其中
% unit: noether.p08.eq.013
\[
  i!(p-i)!\,g_i(xy)=
  \left(y\frac{\partial}{\partial z}\right)^i
  \left(x\frac{\partial}{\partial z}\right)^{p-i}f(xyz).
\]
\footnote{如 I 中一樣，
\[
 \left(t\frac{\partial}{\partial z}\right)=t_1\frac{\partial}{\partial z_1}+t_2\frac{\partial}{\partial z_2},
\]
因而特別有
\[
 \left[(\lambda_1x+\lambda_2y)\frac{\partial}{\partial z}\right]
 = (\lambda_1x_1+\lambda_2y_1)\frac{\partial}{\partial z_1}
 + (\lambda_1x_2+\lambda_2y_2)\frac{\partial}{\partial z_2},
\]
\[
 \left[z\left(\frac{\partial}{\partial\lambda_1}\frac{\partial}{\partial x}
 +\frac{\partial}{\partial\lambda_2}\frac{\partial}{\partial y}\right)\right]
 =z_1\left(\frac{\partial}{\partial\lambda_1}\frac{\partial}{\partial x_1}
 +\frac{\partial}{\partial\lambda_2}\frac{\partial}{\partial y_1}\right)
 +z_2\left(\frac{\partial}{\partial\lambda_1}\frac{\partial}{\partial x_2}
 +\frac{\partial}{\partial\lambda_2}\frac{\partial}{\partial y_2}\right).
\]}
這樣，$g_i(xy)$ 給出恆等式 (1) 中的形式 $Z$。各 $g$ 是 $x,y,\lambda$ 以及不定元 $a$ 中的有理整係數形式；設 $\mathfrak S$ 關於 $R$ 的秩為 $\sigma$。

% unit: noether.p08.par.014
3. 形式族 $\mathfrak T$，它由 $\mathfrak S$ 經逆極化過程得到，正如 $\mathfrak S$ 由 $\mathfrak L$ 得到一樣；$\mathfrak T$ 中的形式 $h$ 定義為
% unit: noether.p08.eq.014
\[
\begin{aligned}
 h(x,y,z)
 &=\frac1{p!}\left[z\left(\frac{\partial}{\partial\lambda_1}\frac{\partial}{\partial x}
 +\frac{\partial}{\partial\lambda_2}\frac{\partial}{\partial y}\right)\right]^p g(\lambda;xy)\\
 &=h_0(xyz)+h_1(xyz)+\cdots+h_p(xyz),
\end{aligned}
\]
其中由 2. 得
% unit: noether.p08.eq.015
\[
  h_i(xyz)=
  \left(z\frac{\partial}{\partial x}\right)^{p-i}
  \left(z\frac{\partial}{\partial y}\right)^i g_i(xy).
\]
各個 $h_i$，從而 $h$ 本身，都是形式 $PZ$ 以及它們的和；設 $\mathfrak T$ 關於 $R$ 的秩為 $\tau$。

% unit: noether.p08.par.015
由 $\mathfrak T$ 的形式 $h(x,y,z)$ 的構造可見，這些形式由 $\theta$ 通過極化過程 $P$ 得到，並且在行 $x,y,z$ 中分別具有與 $\theta$ 相同的維數；因此 $\mathfrak T$ 是 $\mathfrak L$ 的子族。由上面提到的事實，為證明等價只剩下證明秩相等。在此證明中，不再使用 $f(xyz)$ 的特殊結構，即它們是同一形式 $\theta$ 的極化形式。

% unit: noether.p08.par.016
\emph{a) 為比較 $\mathfrak L$ 與 $\mathfrak S$ 的秩，使用引理 a)：由 $f(x,y,\lambda_1x+\lambda_2y)=0$ 必然推出 $f(x,y,z)=0$。} 這直接來自三個二元行之間的恆等關係
% unit: noether.p08.eq.016
\[
  (xy)z+(yz)x+(zx)y=0,
  \qquad (xy)=(x_1y_2-x_2y_1)\cdots,
\]
它推出
% unit: noether.p08.eq.017
\[
  f[x,y,(zy)x+(xz)y]=(xy)^p f(x,y,z).
\]
因此，由 $f(x,y,\lambda_1x+\lambda_2y)=0$（關於 $\lambda$ 恆等），通過代入 $\lambda_1=(zy)$、$\lambda_2=(xz)$，並且由於 $(xy)\ne0$，必然得 $f(x,y,z)=0$。

% unit: noether.p08.par.017
由這個引理，$\mathfrak S$ 中的形式與 $\mathfrak L$ 中的形式之間存在一一對應。事實上，由
% unit: noether.p08.eq.018
\[
  f_1(x,y,\lambda_1x+\lambda_2y)=g(\lambda;xy),\qquad
  f_2(x,y,\lambda_1x+\lambda_2y)=g(\lambda;xy)
\]
必然推出
% unit: noether.p08.eq.019
\[
  f_1(xyz)=f_2(xyz).
\]
此外，$\mathfrak S$ 中的每一個關係都保留在 $\mathfrak L$ 中唯一對應的形式之間（反之亦然）。因為由
% unit: noether.p08.eq.020
\[
  c_1g^{(1)}(\lambda;xy)+c_2g^{(2)}(\lambda;xy)+\cdots+c_\varkappa g^{(\varkappa)}(\lambda;xy)=0
\]
通過引理必然得到
% unit: noether.p08.eq.021
\[
  c_1f^{(1)}(x,y,z)+c_2f^{(2)}(x,y,z)+\cdots+c_\varkappa f^{(\varkappa)}(x,y,z)=0.
\]
\emph{由此證明了 $\mathfrak L$ 與 $\mathfrak S$ 的秩相等；$\rho=\sigma$。}

% unit: noether.p08.par.018
\emph{b) 為比較 $\mathfrak S$ 與 $\mathfrak T$ 的秩，使用對應的引理 b)：由 $h(x,y,z)=0$ 必然推出 $g(\lambda;xy)=0$。} 證明時注意，由 3.，$h(x,y,z)=0$ 等價於各個冪積
% unit: noether.p08.eq.022
\[
 \left(\frac{\partial}{\partial\lambda_1}\frac{\partial}{\partial x_1}
 +\frac{\partial}{\partial\lambda_2}\frac{\partial}{\partial y_1}\right)^{p_1}
 \left(\frac{\partial}{\partial\lambda_1}\frac{\partial}{\partial x_2}
 +\frac{\partial}{\partial\lambda_2}\frac{\partial}{\partial y_2}\right)^{p_2}g(\lambda;x,y),
 \qquad (p_1+p_2=p).
\]
的消失。進一步微分還得到冪積
% unit: noether.p08.eq.023
\[
\begin{aligned}
&\left(\frac{\partial}{\partial x_1}\right)^{\alpha_1}
 \left(\frac{\partial}{\partial x_2}\right)^{\alpha_2}
 \left(\frac{\partial}{\partial y_1}\right)^{\beta_1}
 \left(\frac{\partial}{\partial y_2}\right)^{\beta_2} \\
&\quad\cdot
 \left(\frac{\partial}{\partial\lambda_1}\frac{\partial}{\partial x_1}
 +\cdots\right)^{p_1}
 \left(\frac{\partial}{\partial\lambda_1}\frac{\partial}{\partial x_2}
 +\cdots\right)^{p_2}
 g(\lambda,xy)
\end{aligned}
\]
的消失。因此，這些冪積的任意線性組合都消失，從而每一個形式
% unit: noether.p08.eq.024
\[
  f\left(\frac{\partial}{\partial x},\frac{\partial}{\partial y},
  \frac{\partial}{\partial\lambda_1}\frac{\partial}{\partial x}
  +\frac{\partial}{\partial\lambda_2}\frac{\partial}{\partial y}\right)g(\lambda,xy)
\]
也消失。特別選擇 $f$ 使得
% unit: noether.p08.eq.025
\[
  f(x,y,\lambda_1x+\lambda_2y)=g(\lambda;xy),
\]
便還得到
% unit: noether.p08.eq.026
\[
  g\left(\frac{\partial}{\partial\lambda};
    \frac{\partial}{\partial x},
    \frac{\partial}{\partial y}\right)g(\lambda,x,y)=0;
\]
或者，如果
% unit: noether.p08.eq.027
\[
  g(\lambda,xy)=\sum G_i\lambda_1^{\varkappa_1}\lambda_2^{\varkappa_2}
  x_1^{\varkappa_3}x_2^{\varkappa_4}y_1^{\varkappa_5}y_2^{\varkappa_6},
\]
其中 $G_i$ 是不定元 $a$ 中的有理整係數線性形式，則
% unit: noether.p08.eq.028
\[
  \sum \nu_1!\nu_2!\alpha_1!\alpha_2!\beta_1!\beta_2!\,G_i^2=0,
\]
於是 $g(\lambda;xy)=0$。

% unit: noether.p08.par.019
由引理 b)，完全如 a) 中那樣推出：\emph{$\mathfrak S$ 與 $\mathfrak T$ 的秩相等；$\sigma=\tau$。}

% unit: noether.p08.par.020
因此由 a) 和 b) 得 $\rho=\tau$；於是，由於 $\mathfrak T$ 是 $\mathfrak L$ 的子族，這兩個線性形式族的等價性得到證明。於是 $\mathfrak L$ 的每一個形式，特別是 $\theta$，都可表示為 $\rho$ 個線性無關形式 $h(x,y,z)$ 的有理整係數線性組合：
% unit: noether.p08.eq.029
\[
  \theta=\sum c_i h^{(i)}(x,y,z)=\sum PZ.
\]
這個對 $\theta$ 的不定係數 $a$ 推出的恆等式，在把不定元 $a$ 替換為任意有理性域中的量時仍然成立；因此，\emph{三行二元情形下的歸約定理已一般地得到證明。}

% unit: noether.p08.par.021
對於 $N$ 個、每個含 $n$ 個變量的行，一般證明完全平行。設 $\theta(A)$ 在行 $A_i$ 中具有維數 $\alpha_i$。我們再次考慮三個線性形式族：

% unit: noether.p08.par.022
1. 形式族 $\mathfrak L$，由所有從 $\theta(A)$ 通過極化過程 $P$ 得到、並在每個單獨行 $A_i$ 中具有與 $\theta(A)$ 相同維數的形式 $f(A)$ 構成；

% unit: noether.p08.par.023
2. 形式族 $\mathfrak S$，由所有通過代入
% unit: noether.p08.eq.030
\[
\begin{aligned}
 A_{n+1}&=\lambda_{n+1}^{(1)}A_1+\lambda_{n+1}^{(2)}A_2+
        \cdots+\lambda_{n+1}^{(n)}A_n\\
 &\vdotswithin{=}\\
 A_N&=\lambda_N^{(1)}A_1+\lambda_N^{(2)}A_2+
        \cdots+\lambda_N^{(n)}A_n
\end{aligned}
\]
從 $f(A)$ 得到的形式 $g(\lambda;A_1\cdots A_n)$ 構成；在此過程中，$g(\lambda;A_1\cdots A_n)$ 中各個 $\lambda$ 冪積的係數給出恆等式 (1) 的形式 $Z$；

% unit: noether.p08.par.024
3. 形式族 $\mathfrak T$，由所有按
% unit: noether.p08.eq.031
\[
\begin{aligned}
 h(A)=&\left[A_{n+1}\left(
 \frac{\partial}{\partial\lambda_{n+1}^{(1)}}\frac{\partial}{\partial A_1}
 +\cdots+
 \frac{\partial}{\partial\lambda_{n+1}^{(n)}}\frac{\partial}{\partial A_n}\right)\right]^{\alpha_{n+1}} \\
&\cdots
\left[A_N\left(
 \frac{\partial}{\partial\lambda_N^{(1)}}\frac{\partial}{\partial A_1}
 +\cdots+
 \frac{\partial}{\partial\lambda_N^{(n)}}\frac{\partial}{\partial A_n}\right)\right]^{\alpha_N}
 g(\lambda;A_1\cdots A_n)
\end{aligned}
\]
定義的形式 $h(A)$ 構成。這裡 $h(A)$ 是形式 $PZ$ 的和。

% unit: noether.p08.par.025
由於任意 $n+1$ 個行之間存在恆等關係，引理 a) 仍然成立；引理 b) 同樣仍然成立，因為變量數並未進入它的證明。因此 $\mathfrak L$ 與 $\mathfrak T$ 具有相同的秩，並且由於 $\mathfrak T$ 是 $\mathfrak L$ 的子族，它們相同。所以恆等式
% unit: noether.p08.eq.032
\[
  \theta=\sum PZ
\]
對不定係數成立，從而也對任意有理性域中的量作為係數成立；\emph{歸約定理由此在一般情形下得到證明。}

% unit: noether.p08.sec.003
\subsection*{III.}
% unit: noether.p08.par.026
我們還簡要指出，類似的考察如何也給出\emph{一般級數展開}（對於 $N\le n$）。設 $Z$ 是 $n$ 個行 $A_1,\ldots,A_n$ 中分別齊次的形式，每個行含 $n$ 個量；設 $\Delta$ 為這些行的行列式。先證明與 (1) 對應的恆等式：
% unit: noether.p08.eq.033
\begin{equation}
  Z=\sum PH \mod{\Delta},
\tag{2}
\end{equation}
其中形式 $H$ 只含有行 $A_1\cdots A_{n-1}$，並且由 $Z$ 通過過程 $P$ 導出。

% unit: noether.p08.par.027
證明在於把 II 中推出的關係轉化為模 $\Delta$ 的同餘。為簡單起見，取三個三元行 $x,y,z$；係數假定為不定元。

% unit: noether.p08.par.028
三個線性形式族 $\mathfrak L,\mathfrak S,\mathfrak T$ 如 II 中二元行情形那樣定義；設其秩分別為 $\rho,\sigma,\tau$，而 $\rho^*$ 和 $\tau^*$ 表示 $\mathfrak L$ 和 $\mathfrak T$ 模 $\Delta$ 的秩（也就是說，有 $\rho^*$ 個，但沒有 $\rho^*+1$ 個，$\mathfrak L$ 中的形式模 $\Delta$ 線性無關）。那麼如 II 中一樣，有：

% unit: noether.p08.par.029
\emph{引理 a)：由 $f(x,y,\lambda_1x+\lambda_2y)=0$ 必然推出 $f(xyz)\equiv0\pmod{\Delta}$（反之亦然）。} 因為由四個三元行之間的關係
% unit: noether.p08.eq.034
\[
  \Delta_1x-\Delta_2y+\Delta_3z=\Delta t,
  \qquad \Delta_1=(yzt)\cdots,
\]
總可得到三個三元行之間的一個模 $\Delta$ 的線性關係：
% unit: noether.p08.eq.035
\[
  \Delta_1x-\Delta_2y+\Delta_3z\equiv0\mod{\Delta},
\]
因此如 II 中一樣，
% unit: noether.p08.eq.036
\[
  f(x,y,-\Delta_1x+\Delta_2y)\equiv\Delta_3^p f(x,y,z)\pmod{\Delta}.
\]
由 $f(x,y,\lambda_1x+\lambda_2y)=0$（關於 $\lambda$ 恆等），由於 $\Delta$ 不可約且 $\Delta_3$ 不為 $\Delta$ 所整除，必然推出 $f(x,y,z)\equiv0\pmod{\Delta}$。反向由 $\Delta(x,y,\lambda_1x+\lambda_2y)=0$ 顯然。於是如 II 中那樣得 $\rho^*=\sigma$。

% unit: noether.p08.par.030
引理 b) 及其證明仍然有效，故有 $\sigma=\tau$。

% unit: noether.p08.par.031
為證明 $\tau=\tau^*$，使用

% unit: noether.p08.par.032
\emph{引理 c)：由 $h(x,y,z)\equiv0\pmod{\Delta}$ 必然推出 $h(x,y,z)=0$。} 事實上，作為極化形式，$h(x,y,z)$ 在施用已知的 $\Omega$ 過程時為零；這由微分方程
% unit: noether.p08.eq.037
\[
  \Delta\!\left(\frac{\partial}{\partial x},\frac{\partial}{\partial y},\frac{\partial}{\partial z}\right)h(xyz)
  -\left(\frac{\partial}{\partial x}\frac{\partial}{\partial y}\frac{\partial}{\partial z}\right)h(x,y,z)=0.\begingroup\renewcommand{\thefootnote}{*)}\footnotemark\addtocounter{footnote}{-1}\endgroup
\]
\begingroup\renewcommand{\thefootnote}{*)}\footnotetext{按 3.，有：
\[
 c\cdot\Omega(h)=\left(
\frac{\partial}{\partial \xi}\frac{\partial}{\partial \eta}
\left(\frac{\partial}{\partial\lambda_1}\frac{\partial}{\partial \xi}
+\frac{\partial}{\partial\lambda_2}\frac{\partial}{\partial \eta}\right)
\right)
\left(x\frac{\partial}{\partial \xi}\right)^{\alpha-1}
\left(y\frac{\partial}{\partial \eta}\right)^{\beta-1}
\left[
z\left(\frac{\partial}{\partial\lambda_1}\frac{\partial}{\partial \xi}
+\frac{\partial}{\partial\lambda_2}\frac{\partial}{\partial \eta}\right)
\right]^{p-1}\cdot g(\lambda;\xi,\eta)
\]
故 $\Omega(h)$ 因行列式因子的消失而消失。}\endgroup
表示。若現在 $h(x,y,z)=\Delta(x,y,z)\cdot k(x,y,z)$，則進一步微分也給出
% unit: noether.p08.eq.038
\[
  k\!\left(\frac{\partial}{\partial x},\frac{\partial}{\partial y},\frac{\partial}{\partial z}\right)
  \Delta\!\left(\frac{\partial}{\partial x},\frac{\partial}{\partial y},\frac{\partial}{\partial z}\right)
  =
  h\!\left(\frac{\partial}{\partial x},\frac{\partial}{\partial y},\frac{\partial}{\partial z}\right)h(x,y,z),
\]
的消失，因而 $h(x,y,z)$ 也消失。因此 $\rho^*=\sigma=\tau=\tau^*$；由於 $\mathfrak T$ 是 $\mathfrak L$ 的子族，恆等式 (2) 得證。同樣的論證給出 $n$ 個變量情形的證明。

% unit: noether.p08.par.033
把恆等式 (2) 應用於 $\Delta$ 的乘子，得到展開
% unit: noether.p08.eq.039
\[
  Z=\varphi_0+\Delta\varphi_1+\Delta^2\varphi_2+\cdots+\Delta^\mu\varphi_\mu,
\]
其中每個 $\varphi_i$ 在施用 $\Omega$ 過程時都消失。該過程重複 $i$ 次給出，因為一般有 $\Omega(\Delta\cdot\psi)=c_1\psi+c_2\Delta\cdot\Omega(\psi)$，
% unit: noether.p08.eq.040
\[
  c\cdot \Omega^i(Z)\equiv\varphi_i\pmod{\Delta};
\]
另一方面，由 (2) 有
% unit: noether.p08.eq.041
\[
  c\cdot\Omega^i(Z)\equiv\sum PH_i\pmod{\Delta},
\]
其中 $H_i$ 由形式 $\Omega^i(Z)$ 以 $H$ 由 $Z$ 產生的同樣方式導出。由引理 c)（推廣到 $n$ 個變量），便得 $\varphi_i=\sum PH_i$，從而
% unit: noether.p08.eq.042
\begin{equation}
  Z=\sum PH+\Delta\cdot\sum PH_1+\Delta^2\cdot\sum PH_2+
  \cdots\Delta^\mu\cdot\sum PH_\mu
\tag{3}
\end{equation}
\emph{這是一個關於 $n$ 個行、每行 $n$ 個變量的形式的最一般級數展開，即按只含 $n-1$ 個行的形式的極化形式來展開。}

% unit: noether.p08.par.034
對於 $N$ 個行、每行 $n$ 個變量的形式 $(N<n)$，其級數展開可由一個簡單代入得到，這是已知的。我們設（$N=3$, $n>3$）
% unit: noether.p08.eq.043
\[
  u=\xi_1x+\xi_2y+\xi_3z,
  \quad v=\eta_1x+\eta_2y+\eta_3z,
  \quad w=\zeta_1x+\zeta_2y+\zeta_3z,
\]
並按 (3) 展開 $H(u,v,w)=Z(\xi,\eta,\zeta)$。於是，由於 $\xi,\eta,\zeta$ 只以組合 $u,v,w$ 出現，有
% unit: noether.p08.eq.044
\[
  \left(\eta\frac{\partial}{\partial \xi}\right)=
  \left(v\frac{\partial}{\partial u}\right)\cdots,
  \qquad
  \Omega=\sum\left(\frac{\partial}{\partial u}
  \frac{\partial}{\partial v}
  \frac{\partial}{\partial w}\right)_{ikl}(xyz)_{ikl}
  =\nabla_{uvw}(xyz).
\]
在特化 $\xi_1=\eta_2=\zeta_3=1$；$\xi_2=\xi_3=\cdots=\zeta_2=0$ 下，$H(u,v,w)$ 變為 $H(x,y,z)$，而 $\nabla_{uvw}(xyz)$ 在相差一個數值因子下變為 $\nabla_{xyz}(xyz)=\nabla_{xyz}$，因為把 $(y\partial/\partial x)$ 等作用於行列式 $(xyz)_{ikl}$ 會使它們消失。由於相應說法對 $N$ 個行也成立，於是由 (3) 得到展開
% unit: noether.p08.eq.045
\begin{equation}
  H=\sum P\Phi+\sum P\Phi_1+\cdots+\sum P\Phi_\mu,
\tag{4}
\end{equation}
其中 $\Phi_i$ 由 $\nabla^i_{A_1\ldots A_N}H$ 通過過程 $P$ 產生，並且除 $\nabla^i$ 中出現的行列式組合外，只含有行 $A_1,\ldots,A_{N-1}$；如前所述，這些行列式組合在極化時不需考慮。

% unit: noether.p08.par.035
如果接着把 $\nabla^i$ 中這些行列式替換為不定元 $p$，得到的形式又可以按 (4) 展開；這樣最後得到一個展開
% unit: noether.p08.eq.046
\begin{equation}
  H=\left[\sum P\psi(p)\right]_{p=\text{$A_1\cdots A_N$ 的行列式}A_1\cdots A_N},
\tag{5}
\end{equation}
其中 $\psi$ 由初始形式通過過程 $P$ 以及 $\nabla_{A_1\ldots A_N}(p)$、$\nabla_{A_1\ldots A_{N-1}}(p)$ 等產生。藉助這些展開及其組合，例如把 (3) 然後 (4) 與 (5) 應用於 (1) 中出現的形式 $Z$，對於任意多個、每個含 $n$ 個同變變量的行的形式，所有已知的展開都已窮盡。

% unit: noether.p08.close.001
Erlangen, 1915 年 1 月 5 日。
%
\clearpage
%
% unit: noether.p09.title.001
\begin{center}
\pdfbookmark[0]{9. 由整超越數構成的最一般領域}{paper-09}
{\Large\bfseries 9. 由整超越數構成的最一般領域}\par
\vspace{1.2em}
\emph{Math. Ann.} 77 (1916), pp. 103--128
\end{center}

\vspace{1.5em}

% unit: noether.p09.par.001
藉助良序定理，E. Zermelo 構造了一個“整超越數”的領域，也就是說，一個數的領域 $\mathfrak G$，它具有下列抽象性質：
\begin{enumerate}
\item[I.] 來自 $\mathfrak G$ 的兩個數的和、差與積仍然是 $\mathfrak G$ 中的數。
\item[II.] 每一個實數或複數都是 $\mathfrak G$ 中兩個數的商。
\item[III.] 每一個整有理數（相應地，整代數數）都是 $\mathfrak G$ 的元素。
\item[IV.] 沒有一個非整有理數（相應地，非整代數數）屬於 $\mathfrak G$。\footnote{E. Zermelo, \emph{Über ganze transzendente Zahlen}, Math. Ann. 75, p. 434 (1914).}
\end{enumerate}

% unit: noether.p09.par.002
這個構造基於如下事實：良序定理顯示出所有數的一個\emph{超越基}的存在；也就是說，存在一個由數 $\eta$ 組成的\emph{系統} $\mathsf{H}$，這些數之間沒有代數關係，而所有其餘數都能借助它們代數地表示出來。由於這些基元素 $\eta$ 代數無關，它們可以看作不定元；於是所要求的領域成為不定元 $\eta$ 的有理函數與代數函數所成的整域，\footnote{因此，本文的考慮部分地與作者以前的論文 \emph{Körper und Systeme rationaler Funktionen}, Math. Ann. 76, p. 161 (1915) 平行。}其係數屬於所有代數數的域 $\mathsf{K}$。於是 Zermelo 所構造的特殊領域 $\mathfrak G_\eta$ 由下列基性質刻畫：

% unit: noether.p09.par.003
$\mathfrak G_\eta$ 包含：
\begin{enumerate}
\item[1)] 所有基元素 $\eta$；
\item[2)] 所有以 $\eta$ 為變量且具有整係數的多項式；\footnote{全文中，“整”總是指係數屬於所有代數整數的整域 $[\mathsf{K}]$。對於 $\mathfrak G_\eta$ 的定義，在 1)--5) 中僅考慮有理整數係數也已足夠；不過對於更一般的領域，基性質就會變得不那麼簡單。}
\item[3)] 所有以 $\eta$ 為自變量且具有整係數的整代數函數。
\end{enumerate}
$\mathfrak G_\eta$ 排除：
\begin{enumerate}
\item[4)] 所有以 $\eta$ 為自變量、具有整係數的有理分式函數；
\item[5)] 所有以 $\eta$ 為自變量、具有整係數的代數分式函數。\footnote{Zermelo, loc. cit., § 3.}
\end{enumerate}

% unit: noether.p09.par.004
現在很容易看出（參見 §§ 2 和 3 中的例子），存在一些領域 $\mathfrak G$ 與 Zermelo 領域 $\mathfrak G_\eta$ 本質上不同；也就是說，存在這樣的領域 $\mathfrak G$，在任何良序下，基性質 1)--5) 都不能同時成立，因而它們在任何良序下都不能由 Zermelo 的構造生成。換言之，存在一些領域 $\mathfrak G$ 不能一一且同構地映到 $\mathfrak G_\eta$ 上。因此，由整超越數構成的最一般領域的問題就產生了；下文將對此作完全回答。

% unit: noether.p09.par.005
為此，必須先確定哪些基性質是抽象定義性質的結果，哪些則由特殊構造造成。將證明（§ 1）：對每一個給定的領域 $\mathfrak G$，都存在一個良序，使得基性質 1) 和 2) 成立；因此每一個領域 $\mathfrak G$ 都可以藉助所有數的一個超越基構造出來。進一步的基性質沒有一個必須成立（§§ 2 和 3）。特別地，由此得出，Zermelo 領域的另一個抽象性質並不屬於所有領域 $\mathfrak G$，即代數整閉性。它可以表述如下：
\begin{enumerate}
\item[V.] 每一個以整係數代數整地依賴於 $\mathfrak G$ 中有限多個數的數，都屬於 $\mathfrak G$。
\end{enumerate}

% unit: noether.p09.par.006
除 I--IV 之外還滿足 V 的特殊領域，將稱為由代數整超越數構成的領域 $\mathfrak H$。對於最一般的領域 $\mathfrak H$，得到如下簡單結果（§ 4）：
\begin{enumerate}
\item[a)] 藉助同一個基 $\mathsf{H}$ 可構造的所有領域 $\mathfrak H$ 的交，由該基的所有整代數函數給出，也就是由 Zermelo 領域 $\mathfrak G_\eta$ 給出；而 $\mathfrak G_\eta$ 本身也是一個領域 $\mathfrak H$。這同時也給出了 Zermelo 領域的一個抽象刻畫。
\item[b)] 每一個領域 $\mathfrak H$ 都由把一個任意系統 $\mathfrak S(\eta)$ 代數整地附加到 $\mathfrak G_\eta$ 而得到；\footnote{此表達式的解釋見 § 4。}這裡 $\mathfrak S(\eta)$ 只須滿足唯一的條件：沒有代數分式數通過這個附加進入該領域。
\end{enumerate}

% unit: noether.p09.par.007
當以超越基為基礎時，關於最一般領域 $\mathfrak G$ 的結果不那麼簡單。這裡所有領域 $\mathfrak G$ 的交本身並不是一個領域 $\mathfrak G$，因此最一般領域的構造也較難一目瞭然（§ 5）。

% unit: noether.p09.par.008
為了得到與領域 $\mathfrak H$ 類似的結果，超越基被一個\emph{有理基} $\Theta$ 取代；也就是說，由數 $\vartheta$ 構成的一個\emph{系統} $\Theta$，它允許所有數作有理表示，而在至少一個良序下，沒有一個基元素能夠由有限多個在它之前的基元素有理地表示。這個有理基 $\Theta$ 還必須服從一個附加條件：由所有 $\vartheta$ 的整多項式構成的整域 $[\Theta]$ 不含任何代數分式數。具有這個附加條件的這種基的構造見 § 7。於是，對於最一般的領域 $\mathfrak G$——也可以把它們稱為由有理整超越數構成的領域——又得到如下簡單結果（§ 6）：
\begin{enumerate}
\item[a)] 藉助同一個有理基 $\Theta$ 可構造的所有領域 $\mathfrak G$ 的交，由該基的所有整有理函數給出，也就是由領域 $[\Theta]$ 給出；而 $[\Theta]$ 本身就是一個領域 $\mathfrak G$。
\item[b)] 每一個領域 $\mathfrak G$ 都由把一個任意系統 $\mathfrak S(\vartheta)$ 有理整地附加到 $[\Theta]$ 而得到；這裡 $\mathfrak S(\vartheta)$ 只須滿足唯一的條件：沒有代數分式數通過這個附加進入該領域。
\end{enumerate}

% unit: noether.p09.par.009
為了與領域 $\mathfrak H$ 完全類比，必須在領域 $\mathfrak G$ 的定義性質 III 和 IV 中刪去代數數。作了這種 III 和 IV 的修改以後，用有理基構造整元素的方法可以推廣到任意域（§ 10）；而 Zermelo 的構造則預設該域代數閉。

% unit: noether.p09.par.010
當 $\mathsf{H}$ 遍歷每一個超越基，或者當 $\Theta$ 遍歷每一個滿足附加條件的有理基時，結果 a) 和 b) 給出了全部領域 $\mathfrak H$ 和 $\mathfrak G$。如果把所有同構領域歸入一個類，那麼可以從這個全體中選出一個子集，它代表所有類的全體——按 § 9，至少有可數無窮多個類（§ 8）。

% unit: noether.p09.sec.001
\subsection*{§ 1. 對所有領域 $\mathfrak G$ 證明基性質 1) 和 2).}

% unit: noether.p09.par.011
設 $\mathfrak G$ 是任意給定的由整超越數構成的領域，因而它具有抽象性質 I--IV。按照 E. Zermelo 對所有複數的域所採用的程序，我們從 $\mathfrak G$ 中選取一個\emph{所有 $\mathfrak G$ 中數的超越基} $\mathsf{H}$。由於按 II，每一個實數或複數都能表示為 $\mathfrak G$ 中兩個數的商，$\mathsf{H}$ 同時也是\emph{所有}數的一個超越基。因此，對於每一個給定的領域 $\mathfrak G$，都能把所有數的超越基選成屬於 $\mathfrak G$；所以基性質 1) 成立。此外，由 III，$\mathfrak G$ 包含所有代數整數的整域 $[\mathsf{K}]$，因而由 I，也包含所有以基元素 $\eta$ 為變量、係數在 $[\mathsf{K}]$ 中的多項式，也就是所有 $\eta$ 的整多項式；這些多項式構成整域 $[\mathsf{H}]$。所以基性質 2) 也總是成立；換言之，每一個領域 $\mathfrak G$ 都可以藉助所有數的一個超越基 $\mathsf{H}$ 來構造，使得 $\mathfrak G$ 包含整域 $[\mathsf{H}]$ 作為一個子環。

% unit: noether.p09.par.012
\textbf{註記。} 然而，從一個基 $\mathsf{H}$ 出發，當然可以構造一個領域 $\mathfrak G$，使得基性質 1) 和 2) 對 $\mathsf{H}$ 本身不成立，卻對另一個基 $\mathsf{H}'$ 成立。例如，令 $\mathfrak G'$ 由 Zermelo 領域 $\mathfrak G_\eta$ 得到，方式是在整個構造中把某個基元素 $\eta$ 替換為一個多項式 $f(\eta)$。那麼 $\mathfrak G'$ 既不包含 $\eta$，也不包含 $[\mathsf{H}]$；但是如果在基 $\mathsf{H}$ 中把基元素 $\eta$ 替換為 $\zeta=f(\eta)$——由此 $\mathsf{H}$ 又變為一個超越基 $\mathsf{H}'$——則 $\mathfrak G'$ 對 $\mathsf{H}'$ 滿足基性質 1) 和 2)。

% unit: noether.p09.sec.002
\subsection*{§ 2. 排除基性質 4) 和 5).}

% unit: noether.p09.par.013
現在我們表明，領域 $\mathfrak G$ 並不必滿足任何進一步的基性質。首先，通過在 Zermelo 的構造中把整域 $[\mathsf{H}]$ 替換為整有理“函數泛函”的領域，可以排除 4) 和 5)。

% unit: noether.p09.par.014
按照 Weber，\footnote{H. Weber, \emph{Lehrbuch der Algebra}. Kleine Ausgabe, §§ 96 和 97.}所謂\emph{有理函數泛函}，是指以有理數為係數的每一個有理函數，其表示唯一確定為
% unit: noether.p09.eq.001
\[
 A(\eta_1\cdots \eta_t)
   = a\cdot\frac{E_1(\eta_1\cdots \eta_t)}{E_2(\eta_1\cdots \eta_t)},
\]
其中 $E_1,E_2$ 是具有有理整數係數的互素本原多項式，而 $a$ 是一個正有理數，即 $A$ 的絕對值。特別地，如果 $a$ 是一個\emph{有理整數}，那麼 $A$ 稱為一個\emph{整有理函數泛函}。作為特殊情形，整有理函數泛函包括有理整數以及具有有理整數係數的多項式；而有理分式數以及具有有理分式係數的多項式則應歸入分式有理函數泛函之中。

% unit: noether.p09.par.015
令 $\mathfrak J$ 為每次只涉及有限多個基元素 $\eta$ 的所有整有理函數泛函的全體，並令領域 $\mathfrak G$ 由 $\mathfrak J$ 以及所有代數整地依賴於 $\mathfrak J$ 的函數組成，也就是說，後者滿足某個首項係數為一、其餘係數為整有理函數泛函的方程。

% unit: noether.p09.par.016
對 $\mathfrak G$ 證明性質 I--IV 的過程與 Zermelo 的相應證明完全平行，這裡只簡要指出。首先 II 和 III 是清楚的，因為 $\mathfrak G$ 包含 Zermelo 領域 $\mathfrak G_\eta$ 作為子環。由關於有理整數係數多項式的 Gauss 定理，整有理函數泛函的和、差與積仍是整有理函數泛函；因此 $\mathfrak J$ 形成一個整域，並由熟知的論證\footnote{例如比較 Weber, § 97, 8.}可知 $\mathfrak G$ 也成為一個整域；所以 I 成立。進一步，由擴展的 Gauss 定理\footnote{Weber, § 20, 5.} 得出兩個輔助命題：“如果 $z$ 通過\emph{某個}方程代數整地依賴於 $\mathfrak J$，那麼它也通過它相對於所有有理函數泛函的域所滿足的\emph{不可約}方程這樣依賴”，\footnote{比較 Weber, § 97, 4.}以及“一個代數數在所有有理數的域上不可約並滿足的方程，在所有有理函數泛函的域上也不可約”。因此 $\mathfrak G$ 不能包含任何分式代數數；於是 IV，因而所有抽象性質 I--IV，都已得到驗證。

% unit: noether.p09.par.017
另一方面，在任何良序下，都不能從 $\mathfrak G$ 中選出一個基，使得 $\mathfrak G$ 排除所有基元素的有理分式函數和代數分式函數。事實上，設 $z(\eta)$ 是領域中任一將被選作基元素的函數，因此它由一個方程
% unit: noether.p09.eq.002
\[
 z^k+A_1(\eta)z^{k-1}+\cdots+A_k(\eta)=0
\]
定義，其中
% unit: noether.p09.eq.003
\[
 A_i(\eta)=a_i\cdot\frac{E_1(\eta)}{E_2(\eta)}
\]
是整有理函數泛函。那麼函數 $\frac{a_k}{z}$ 屬於 $\mathfrak G$，同樣 $\frac{a_k}{\sqrt z}$、$\frac{a_k}{\sqrt[3]z}$ 等也屬於 $\mathfrak G$。\footnote{完全類似地，每一個 $\eta$ 的代數函數都可通過乘以一個有理整數而變成函數泛函領域 $\mathfrak G$ 的元素。因此，這個領域具有如下性質：每一個複數都能表示為 $\mathfrak G$ 中一個元素與一個有理整數的商，這與代數數的情形相類似。} 因而，對於領域 $\mathfrak G$ 的全體，基性質 4) 和 5) 被排除了。

% unit: noether.p09.par.018
\textbf{註記。} § 3 中的一個例子將表明，$\mathfrak G$ 在每一個良序下都可以包含有理分式函數泛函，而不包含有理分式數或代數分式數。

% unit: noether.p09.sec.003
\subsection*{§ 3. 排除基性質 3).}

% unit: noether.p09.par.019
為了排除基性質 3)，我們使用一個廣義的\emph{模領域}。在這個領域中，領域內多項式的自變量不再是不定元，而是不定元的所有整代數函數；\footnote{所謂“整代數函數”總是指具有\emph{整}係數的函數；也就是說，滿足某個方程的函數，該方程的首項係數為一，其餘係數是以 $\eta$ 為變量、具有代數整數係數的多項式，即來自 $[\mathsf{H}]$ 的多項式。特別地，$[\mathsf{H}]$ 中所有多項式都屬於整代數函數。}而不定元本身承擔模基的角色。

% unit: noether.p09.par.020
令 $\mathfrak G$ 由\emph{所有元素}
% unit: noether.p09.eq.004
\begin{equation}
 z=\alpha+g_1(\eta)\cdot\eta_1+g_2(\eta)\cdot\eta_2+\cdots+g_t(\eta)\cdot\eta_t,
\tag{1}
\end{equation}
\emph{組成，其中 $\alpha$ 遍歷所有代數整數，$\eta_1\cdots\eta_t$ 遍歷所有基元素，而對於每一個固定組合 $\eta_1\cdots\eta_t$，函數 $g_1(\eta)\cdots g_t(\eta)$ 分別遍歷所有 $\eta$ 的整代數函數。}\footnote{模性質屬於元素 $z-\alpha$。}

% unit: noether.p09.par.021
容易驗證，性質 I--IV 對這個模領域成立。首先 I 成立，因為兩個元素 $z$ 的和、差與積仍具有形式 (1)。領域 $\mathfrak G$ 也包含所有基元素 $\eta$，並且還包含所有元素 $\eta\cdot g(\eta)$，其中 $g(\eta)$ 遍歷所有 $\eta$ 的整代數函數。因此所有 $\eta$ 的整代數函數，進而所有 $\eta$ 的代數函數，都能表示為 $\mathfrak G$ 中兩個元素的商，這證明了 II。在表示式 (1) 中，所有代數整數也特別包含在內——取 $g_1=0\cdots g_t=0$ 即可；但 $z$ 不可能等於一個分式代數數。因為如果 $z$ 表示一個代數數 $\beta$，那麼由關於 $\eta$ 恆等成立的關係
% unit: noether.p09.eq.005
\[
 \beta=\alpha+g_1(\eta)\eta_1+\cdots+g_t(\eta)\eta_t
\]
必然得出 $\beta=\alpha$。這可由特殊化
% unit: noether.p09.eq.006
\[
 \eta_1=0\cdots \eta_t=0
\]
看出；在這個特殊化下，乘子 $g_1(\eta)\cdots g_t(\eta)$ 作為整代數函數仍保持有限，因為它們變為整代數函數或整代數數。\footnote{這裡給出 IV 的較詳細證明，是為了本段末尾的擴展例子。在此只需指出，由構造可知 $z$ 是整代數函數即可。} 因而條件 I--IV 對 $\mathfrak G$ 成立。

% unit: noether.p09.par.022
另一方面，在任何良序下，都不能從 $\mathfrak G$ 中選出一個基，使得所有基元素的整代數函數都屬於 $\mathfrak G$。設 $z(\eta)$ 是領域中任一將被選作基元素的函數——因此它實際上必須含有不定元 $\eta$。我們將證明，從某個依賴於 $z$ 的有限值 $\nu$ 起，函數
% unit: noether.p09.eq.007
\[
 \xi=\sqrt[\nu]{z-\alpha}
\]
不再能屬於 $\mathfrak G$。

% unit: noether.p09.par.023
事實上，假設 $\xi$ 屬於 $\mathfrak G$。那麼由 (1)，它具有形式
% unit: noether.p09.eq.008
\[
 \zeta=\gamma+h_1(\eta)\eta_{i_1}+h_2(\eta)\eta_{i_2}
      +\cdots+h_\tau(\eta)\eta_{i_\tau},
\]
其中 $h_1(\eta)\cdots h_\tau(\eta)$ 仍是 $\eta$ 的整代數函數，而 $\eta_{i_1}\cdots\eta_{i_\tau}$ 是任意基元素；它們特別也可以與 $\eta_1\cdots\eta_t$ 重合。首先必須證明代數整數 $\gamma$ 為零。由於 $h_1(\eta)\cdots h_\tau(\eta)$ 是整代數函數，當 $\eta_{i_1}=0\cdots\eta_{i_\tau}=0$ 且其餘 $\eta$ 任意時，$\xi$ 等於 $\gamma$；特別地，在
% unit: noether.p09.eq.009
\[
 \eta_{i_1}=0\cdots \eta_{i_\tau}=0; \qquad
 \eta_1=0\cdots \eta_t=0
\]
時，$\xi$ 等於 $\gamma$。但在這個特殊化下，函數 $(z-\alpha)$ 由 (1) 等於零，因此 $\xi$ 也等於零；故 $\gamma=0$，於是有
% unit: noether.p09.eq.010
\[
 \zeta=h_1(\eta)\eta_{i_1}+h_2(\eta)\eta_{i_2}
      +\cdots+h_\tau(\eta)\eta_{i_\tau}.
\]
將其升到第 $\nu$ 次冪，得到
% unit: noether.p09.eq.011
\[
 \zeta^\nu=z-\alpha=F_\nu(\eta),
\]
其中 $F_\nu(\eta)$ 表示在 $\eta_{i_1}\cdots\eta_{i_\tau}$ 中一個 $\nu$ 次非恆零齊次形式，其係數是 $\eta$ 的整代數函數。現在，$z-\alpha$ 由 (1) 是 $\eta$ 的一個整代數函數，它滿足相對於 $[\mathsf{H}]$ 不可約的方程：
% unit: noether.p09.eq.012
\[
 (z-\alpha)^\varkappa+B(\eta)(z-\alpha)^{\varkappa-1}+\cdots+C(\eta)=0,
\]
其中最後一個係數 $C(\eta)$——即 $(z-\alpha)$ 與其共軛的乘積——是一個多項式；設它的次數為 $\lambda$。另一方面，由 $z-\alpha=F_\nu(\eta)$，把 $F_\nu(\eta)$ 與相應共軛相乘，得到
% unit: noether.p09.eq.013
\[
 C(\eta)=G_{\nu\varkappa}(\eta),
\]
其中 $G_{\nu\varkappa}(\eta)$ 表示在 $\eta_{i_1}\cdots\eta_{i_\tau}$ 中一個 $\nu\varkappa$ 次齊次形式，其係數是 $\eta$ 的多項式。因此 $G_{\nu\varkappa}$ 關於 $\eta$ 至少有次數 $\nu\varkappa$，即 $\lambda\geqq\nu\varkappa$；也就是說，當 $\nu>\frac{\lambda}{\varkappa}$ 時，函數 $\sqrt[\nu]{z-\alpha}$ 不屬於 $\mathfrak G$。所以，對於領域 $\mathfrak G$ 的全體，基性質 3) 被排除了。

% unit: noether.p09.par.024
\textbf{註記。} 對剛構造的領域作一個輕微推廣，可以得到一個領域 $\mathfrak G$，它在每一個良序下也包含分式函數泛函。令 $\mathfrak G$ 由所有元素
% unit: noether.p09.eq.014
\[
 z=\alpha+\gamma_1(\eta)\eta_1+\gamma_2(\eta)\eta_2
      +\cdots+\gamma_t(\eta)\eta_t,
\]
組成，其中現在 $\gamma(\eta)$ 遍歷 $\eta$ 的所有代數整但不必整的函數。性質 I--IV 的證明與上面相同。但是由於 $\mathfrak G$ 除每個函數 $z$ 以外，還包含 $a\cdot(z-\alpha)$，其中 $a$ 表示任意有理分式數或代數分式數，於是在每一個良序下也會出現分式函數泛函。

% unit: noether.p09.sec.004
\subsection*{§ 4. 由代數整超越數構成的最一般領域.}

% unit: noether.p09.par.025
為了更方便地表述下文，引入說法：“元素 $z$ 代數整地依賴於 $\mathfrak T$”；這表示 $z$ 滿足某個方程，其首項係數為一，其餘係數是 $\mathfrak T$ 的元素的整多項式；這裡 $\mathfrak T$ 表示任意給定領域。\footnote{對於整域 $\mathfrak T$，首項係數為一，其餘係數是 $\mathfrak T$ 的元素。}

% unit: noether.p09.par.026
如果一個領域 $\mathfrak T$ 滿足條件（比較引言）
\begin{enumerate}
\item[V.] 每一個代數整地依賴於 $\mathfrak T$ 的元素都屬於 $\mathfrak T$，
\end{enumerate}
則稱 $\mathfrak T$ 為\emph{代數整閉}。

% unit: noether.p09.par.027
我們先證明抽象性質 V 屬於 Zermelo 領域 $\mathfrak G_\eta$。設元素 $z$ 通過方程 $f(z)=0$ 代數整地依賴於 $\mathfrak G_\eta$。那麼，把 $f(z)$ 與它相對於 $[\mathsf{H}]$ 的共軛相乘，便表明 $z$ 也代數整地依賴於 $[\mathsf{H}]$，從而屬於 $\mathfrak G_\eta$。§ 2 中考慮的函數泛函領域的代數整閉性也以同樣方式證明。

% unit: noether.p09.par.028
性質 V 並不屬於每一個領域 $\mathfrak G$，這由 § 3 中考察的模領域顯示出來；它確實不具有基性質 3)。更準確地說，§ 3 表明，模領域對於該領域中的任何超越元素都不是代數整閉的；而由 III 和 IV 可知，對於該領域的每一個代數元素，情況當然成立。

% unit: noether.p09.par.029
這使我們首先從領域 $\mathfrak G$ 的全體中選出那些也具有抽象性質 V 的領域；這些領域將稱為“由代數整超越數構成的領域 $\mathfrak H$”。

% unit: noether.p09.par.030
設 $\mathfrak H$ 是這樣的領域，且 $\mathsf{H}$ 是 $\mathfrak H$ 中所有數的一個超越基；由 II，$\mathsf{H}$ 同時也是所有數的一個超越基。由 III、I 和 V，$\mathfrak H$ 除了 $\mathsf{H}$ 以外，還包含所有 $\eta$ 的整代數函數，也就是 Zermelo 領域 $\mathfrak G_\eta$。所以，與 § 1 類似：每一個由代數整超越數構成的領域 $\mathfrak H$ 都可以藉助所有數的一個超越基來構造，使得 $\mathfrak H$ 包含 Zermelo 領域 $\mathfrak G_\eta$ 作為子環。

% unit: noether.p09.par.031
現在令 $\mathsf{H}$ 是所有數的任意超越基，並令 $\mathfrak M(H)$ 表示所有包含 $H$ 的領域 $\mathfrak H$ 的全體。剛剛證明過，$\mathfrak M(H)$ 中的每一個領域 $\mathfrak H$ 都包含 Zermelo 領域 $\mathfrak G_\eta$ 作為子環；而由於 $\mathfrak G_\eta$ 本身就是一個領域 $\mathfrak H$，$\mathfrak G_\eta$ 是來自 $\mathfrak M(H)$ 的所有領域 $\mathfrak H$ 的最大公子環，或交。因此 $\mathfrak G_\eta$ 已被抽象地定義。還可直接推出 $\mathfrak M(H)$ 對取交封閉：也就是說，$\mathfrak M(H)$ 中任意子系統的所有領域 $\mathfrak H$ 的交仍屬於 $\mathfrak M(H)$。因為作為交，I 和 V 成立；II 和 III 由 $\mathfrak G_\eta$ 是子環而得；IV 由該交是 $\mathfrak H$ 的子環而得。

% unit: noether.p09.par.032
如 § 2 中函數泛函領域的例子所示，領域 $\mathfrak H$ 一般除了 $\mathfrak G_\eta$ 之外還含有其他函數。設 $\psi(\eta)$ 是 $\eta$ 的這樣一個有理函數或代數函數。於是由 I，$\mathfrak H$ 還包含 $\psi$ 與每次有限多個來自 $\mathfrak G_\eta$ 的函數的一切整有理組合。由此產生的整域——由所有以 $\psi$ 為變量、係數來自 $\mathfrak G_\eta$ 的多項式構成——記為 $[\mathfrak G_\eta,\psi]$，導致它的過程稱為“把 $\psi$ 有理整地附加到 $\mathfrak G_\eta$”。由 V，$\mathfrak H$ 進一步包含所有代數整地依賴於 $[\mathfrak G_\eta,\psi]$ 的函數；由此產生的領域記為 $\{\mathfrak G_\eta,\psi\}$，導致它的過程稱為“把 $\psi$ 代數整地附加到 $\mathfrak G_\eta$”。領域 $\{\mathfrak G_\eta,\psi\}$ 由所有以 $\psi$ 為自變量、係數來自 $\mathfrak G_\eta$ 的整代數函數組成，因此由熟知的論證可知它是一個整域。\footnote{例如比較 Weber, § 97, 6 和 7.}

% unit: noether.p09.par.033
我們證明 $\{\mathfrak G_\eta,\psi\}$ 也是代數整閉的，也就是說滿足條件 V。事實上，設 $z$ 通過方程
% unit: noether.p09.eq.015
\[
 f(z;a\cdots c)=z^\varkappa+a z^{\varkappa-1}+\cdots+c=0,
\]
代數整地依賴於 $\{\mathfrak G_\eta,\psi\}$，其中 $a\cdots c$ 作為 $\{\mathfrak G_\eta,\psi\}$ 的元素，滿足係數來自 $[\mathfrak G_\eta,\psi]$ 的方程：
% unit: noether.p09.eq.016
\[
\begin{gathered}
 a^\lambda+A_1a^{\lambda-1}+\cdots+A_\lambda=0,\\
 \vdots\\
 c^\mu+C_1c^{\mu-1}+\cdots+C_\mu=0.
\end{gathered}
\]
把它們的根記為 $a,a'\cdots a^{(\lambda-1)}\cdots c,c'\cdots c^{(\mu-1)}$。如果現在對根的所有組合取乘積
% unit: noether.p09.eq.017
\[
 g(z)=f(z;a\cdots c)f(z;a'\cdots c')\cdots
      f(z;a^{(\lambda-1)}\cdots c^{(\mu-1)}),
\]
那麼 $g(z)$ 的首項係數為一，其餘係數是 $[\mathfrak G_\eta,\psi]$ 的元素。因此 $z$ 屬於 $\{\mathfrak G_\eta,\psi\}$。\footnote{因此，只要用\emph{某個}方程定義代數整性，在證明 I 和 V 時就不再需要相對於基域的不可約性概念。另參 § 2，在那裡只有為了證明 IV，才需要輔助命題，把由任意方程給出的代數整性定義過渡到不可約方程。由於這個輔助命題在當前情形下一般不再成立，IV 必須作為特殊條件加在函數 $\psi$ 上。}

% unit: noether.p09.par.034
因此領域 $\{\mathfrak G_\eta,\psi\}$ 具有性質 I 和 V；並且由於 $\mathfrak G_\eta$ 是該領域的子環，它也具有 II 和 III。IV 是否成立取決於 $\psi$ 的選擇；例如，對於 $\psi=\frac{1}{n\cdot\eta}$，IV 不成立，因為這時 $\frac1n$ 屬於該領域。然而由 § 2 可知，對於 $\psi=\frac1\eta$，或者對於 $\psi=\frac{\eta}{n}$，IV 成立。因為在後一種情形，如果 $\{\mathfrak G_\eta,\psi\}$ 中的一個函數表示一個代數數，那麼在 $\eta=0$ 時它的值保持不變；由此它變為 $\mathfrak G_\eta$ 中的一個函數，因而變為一個代數整數。

% unit: noether.p09.par.035
所以，只要對 $\psi$ 加上條件：沒有代數分式數通過代數整附加進入該領域，$\{\mathfrak G_\eta,\psi\}$ 就成為一個領域 $\mathfrak H$。

% unit: noether.p09.par.036
完全類似地，可以定義把 $\eta$ 的有理函數或代數函數的任意系統 $\mathfrak S$ 有理整地、相應地代數整地附加到 $\mathfrak G_\eta$。有理整附加導向整域 $[\mathfrak G_\eta,\mathfrak S]$，它由所有每次取有限多個來自 $\mathfrak G_\eta$ 和 $\mathfrak S$ 的元素而形成的整多項式組成。代數整附加給出領域 $\{\mathfrak G_\eta,\mathfrak S\}$，它由所有代數整地依賴於 $[\mathfrak G_\eta,\mathfrak S]$ 的元素組成。正如上面一樣，可以證明整域 $\{\mathfrak G_\eta,\mathfrak S\}$ 是代數整閉的，並且也滿足 II 和 III。因此，如上所述：

% unit: noether.p09.par.037
$\{\mathfrak G_\eta,\mathfrak S\}$ 一旦對 $\mathfrak S$ 施加條件：沒有代數分式數通過代數整附加進入該領域，也就是說一旦 $\mathfrak S$ 成為一個可容許系統，就成為一個領域 $\mathfrak H$。\footnote{更一般地，同樣可證明：如果 $\mathfrak T$ 是任意領域，且 $\{\mathfrak T\}$ 表示所有代數整地依賴於 $\mathfrak T$ 的元素全體，那麼 $\{\mathfrak T\}$ 是代數整閉的。}

% unit: noether.p09.par.038
現在重要的是，反過來也成立；因此當 $\mathfrak S$ 遍歷所有可容許系統時，領域 $\{\mathfrak G_\eta,\mathfrak S\}$ 確實窮盡 $\mathfrak M(H)$ 中所有領域 $\mathfrak H$。事實上，設 $\mathfrak H$ 是 $\mathfrak M(\mathfrak H)$ 中任一領域，並記 $\mathfrak L=\mathfrak H-\mathfrak G_\eta$ 為由所有屬於 $\mathfrak H$ 但不屬於 $\mathfrak G_\eta$ 的函數組成的系統。由 V，$\mathfrak H$ 包含領域 $\{\mathfrak G_\eta,\mathfrak L\}$ 作為子環；反過來，$\{\mathfrak G_\eta,\mathfrak L\}$ 也包含 $\mathfrak H$，因為它包含 $\mathfrak G_\eta$ 與 $\mathfrak L$，而二者沒有公共元素。於是 $\mathfrak H=\{\mathfrak G_\eta,\mathfrak L\}$；並且由於 IV 對 $\mathfrak H$ 成立，$\mathfrak L$ 當然是可容許系統。\footnote{$\mathfrak L$ 的一個子系統的附加可能已經足夠；例如 § 2 的函數泛函領域，就是把所有整有理函數泛函——除多項式以外——代數整地附加到 $\mathfrak G_\eta$ 而產生的。}

% unit: noether.p09.par.039
當 $\mathsf{H}$ 遍歷每一個可能的超越基時，$\mathfrak M(\mathfrak H)$ 如開頭所證明，遍歷全部領域 $\mathfrak H$。因此，由代數整超越數構成的最一般領域 $\mathfrak H$ 的問題，由本節結果完全回答：
\begin{enumerate}
\item[a)] $\mathfrak M(H)$ 中所有領域 $\mathfrak H$ 的交由 Zermelo 領域 $\mathfrak G_\eta$ 給出，後者本身就是一個領域 $\mathfrak H$；$\mathfrak M(H)$ 對形成交封閉。
\item[b)] $\mathfrak M(H)$ 中每一個領域 $\mathfrak H$ 都由把一個可容許系統 $\mathfrak S$ 代數整地附加到 $\mathfrak G_\eta$ 而得到。當 $\mathsf{H}$ 遍歷所有可能的超越基時，$\mathfrak M(\mathfrak H)$ 遍歷全部領域 $\mathfrak H$。\footnote{可容許系統 $\mathfrak S$ 可如下找到。把所有 $\eta$ 的代數分式函數置於某個良序 $\Omega$ 之下，並把所有函數納入 $\mathfrak S$，但排除那些在代數整地附加到 $\mathfrak G_\eta$ 以及此前已納入的函數之後會生成一個代數分式數的函數。這個過程可以在任意地方中斷。當 $\Omega$ 遍歷所有良序，並且對每個固定的 $\Omega$，過程在每個任意地方中斷時，所有可容許系統 $\mathfrak S$ 由此窮盡。}
\end{enumerate}
%
\clearpage
%
% unit: noether.p09.sec.005
\section*{§ 5. 以超越基為基礎、由整超越數形成的最一般領域。}

% unit: noether.p09.par.041
令 $\mathsf{H}$ 是所有數的任意一個超越基，並令 $\mathfrak M(\mathfrak G)$ 表示所有包含 $H$、因而也包含 $[\mathsf{H}]$ 的領域 $\mathfrak G$ 的全體。當 $H$ 遍歷每一個可能的基時，$\mathfrak M(\mathfrak G)$ 遍歷全部領域 $\mathfrak G$；所以，和 § 4 一樣，我們可以限於考察屬於 $\mathfrak M(\mathfrak G)$ 的領域 $\mathfrak G$。

% unit: noether.p09.par.042
屬於 $\mathfrak M(\mathfrak G)$ 的領域 $\mathfrak G$ 全都含有 $[\mathsf{H}]$ 作為子環；但是，由於 $[\mathsf{H}]$ 本身並不是一個領域 $\mathfrak G$，不能像 § 4 那樣得出 $[\mathsf{H}]$ 是最大公子環的結論。儘管如此，下面由一列可數的、屬於 $\mathfrak M(\mathfrak G)$ 的領域 $\mathfrak G$ 表明這確是事實；這些領域由 § 3 的模領域作輕微推廣得到，並且除了 $[\mathsf{H}]$ 之外沒有共同元素。

% unit: noether.p09.par.043
\emph{對每個固定的 $\nu$，令領域 $\mathfrak G_\nu$ $(\nu=1,2,\ldots\text{ 無窮})$ 由所有量}
% unit: noether.p09.eq.018
\begin{equation}
z=f(\eta)+g_1(\eta)\eta_1^\nu+g_2(\eta)\eta_2^\nu+\cdots+g_t(\eta)\eta_t^\nu,
\tag{1}
\end{equation}
\emph{組成；其中 $f(\eta)$ 遍歷所有 $\eta$ 的整多項式，$\eta_1\cdots\eta_t$ 遍歷基量，而對每個固定組合 $\eta_1^\nu\cdots\eta_t^\nu$，函數 $g_1(\eta)\cdots g_t(\eta)$ 分別獨立地遍歷所有 $\eta$ 的整代數函數。}

% unit: noether.p09.par.044
與 § 3 一樣可見，領域 $\mathfrak G_\nu$ 具有抽象性質 I--IV；$\mathfrak G_\nu$ 只含有整代數函數，並且當 $\nu=1$ 時它與 § 3 的領域相同。現在我們證明，這些領域 $\mathfrak G_\nu$ 除了來自 $[\mathsf{H}]$ 的多項式 $f(\eta)$ 以外沒有其他共同元素。事實上，設 $z$ 是 $\eta$ 的一個整代數但非有理的函數，滿足相對於 $[\mathsf{H}]$ 不可約的方程
% unit: noether.p09.eq.019
\begin{equation}
z^\varkappa+a_1(\eta)z^{\varkappa-1}+a_2(\eta)z^{\varkappa-2}+\cdots+a_\varkappa(\eta)=0,
\tag{2}
\end{equation}
其中 $\varkappa>1$；並令 $\lambda$ 為多項式 $a_1(\eta)\cdots a_\varkappa(\eta)$ 的次數中的最大值。如果 $z$ 屬於 $\mathfrak G_\nu$，那麼量 $\zeta=z-f(\eta)$ 滿足同樣相對於 $[\mathsf{H}]$ 不可約的方程
% unit: noether.p09.eq.020
\begin{equation}
\zeta^\varkappa+b_1(\eta)\zeta^{\varkappa-1}+b_2(\eta)\zeta^{\varkappa-2}+\cdots+b_\varkappa(\eta)=0,
\tag{3}
\end{equation}
其中，由 (1) 可知，$b_1(\eta),b_2(\eta)\cdots b_\varkappa(\eta)$ 作為 $\zeta$ 的對稱函數，最低次項的次數至少分別為 $\nu,2\nu,\cdots,\varkappa\nu$。比較 (2) 與 (3) 得
% unit: noether.p09.eq.021
\[
 b_1(\eta)=a_1(\eta)+\varkappa\cdot f(\eta).
\]
因此，如果量 $\xi=z+\frac{a_1(\eta)}{\varkappa}$ 滿足
% unit: noether.p09.eq.022
\begin{equation}
\xi^\varkappa+c_2(\eta)\xi^{\varkappa-2}+\cdots+c_\varkappa(\eta)=0,
\tag{4}
\end{equation}
則
% unit: noether.p09.eq.023
\begin{equation}
 \xi-\zeta=\dfrac{b_1(\eta)}{\varkappa};\qquad;
 \qquad
 c_i(\eta)\equiv b_i(\eta)\pmod{\dfrac{b_1(\eta)}{\varkappa}},
\tag{5}
\end{equation}
其中 $b_1(\eta)$ 以次數至少為 $\nu$ 的項開始。另一方面，如由 (2) 與 (4) 的比較所示，$c_i(\eta)$ 是 (2) 的係數 $a_i(\eta)$ 的 $\chi$ 次式，因而在 $\eta$ 中次數至多為 $\varkappa\lambda$；而所有 $b_i(\eta)$ 都以次數至少為 $\nu$ 的項開始。於是若取 $\nu>\varkappa\lambda$，(5) 給出
% unit: noether.p09.eq.024
\[
 c_i(\eta)=0;\qquad
 \xi^\varkappa=0\quad\text{或}\quad
 \left(z+\frac{a_1(\eta)}{\varkappa}\right)^\varkappa=0.
\]
所以 $z$ 與假設相反，成為 $\eta$ 的有理函數。因此：

% unit: noether.p09.par.045
\emph{如果 $z$ 是 $\eta$ 的整代數而非有理函數，那麼當 $\nu>\chi\lambda$ 時，$z$ 不能屬於任何領域 $\mathfrak G_\nu$；或者說：}

% unit: noether.p09.par.046
\emph{屬於 $\mathfrak M(\mathfrak G)$ 的所有領域 $\mathfrak G$ 的交，由整域 $[\mathsf{H}]$ 給出，而 $[\mathsf{H}]$ 本身並不是一個領域 $\mathfrak G$；因此 $\mathfrak M(\mathfrak G)$ 對取交併不封閉。}

% unit: noether.p09.par.047
於是，用超越基構造最一般領域 $\mathfrak G$ 也比構造領域 $\mathfrak H$ 更難把握。領域 $[\mathsf{H}]$ 具有性質 I、III 和 IV；若把任意系統 $\mathfrak S$ 有理整地附加上去，則對所得整域 $[\mathsf{H},\mathfrak S]$，性質 I 和 III 仍然成立，而 IV 成為對 $\mathfrak S$ 的一個特殊條件。為了得到一個領域 $\mathfrak G$，還必須把 II 作為條件加在 $\mathfrak S$ 上；換言之：對每個在 $(\mathsf{H})$ 上有限的域 $\mathfrak K$，\footnote{$(\mathsf{H})$ 指由所有以 $\eta$ 為變量、以代數數為係數的有理函數組成的域。} 必須存在一個在 $\mathfrak K$ 之上的域 $\mathfrak L$，它在 $\mathfrak K$ 上有限，並使 $[\mathsf{H},\mathfrak S]$ 含有 $\mathfrak L$ 的一個本原元。\footnote{參見 § 7。} 反過來，屬於 $\mathfrak M(\mathfrak G)$ 的每個領域 $\mathfrak G$ 也能表示為 $[\mathsf{H},\mathfrak S]$ 的形式；所以當 $\mathfrak S$ 遍歷所有以這種方式受限的系統時，就得到 $\mathfrak M(\mathfrak G)$ 中的所有領域 $\mathfrak G$。系統 $\mathfrak S$ 的結構將在下一節藉助“有理基”更簡單地顯現；同時也由此得到與 § 4 的定理的完全類比。

% unit: noether.p09.sec.006
\section*{§ 6. 以有理基為基礎、由整超越數形成的最一般領域。}

% unit: noether.p09.par.048
與超越基類比，有理基定義如下：“由數 $\vartheta$ 組成的系統 $\Theta$ 稱為所有數的有理基，如果 $\Theta$ 能使每一個數都以 $\mathsf{K}$ 中係數有理地表示出來\footnote{所謂有理函數，總應理解為具有代數數係數的這類函數；由於 III 和 IV，這些代數數係數被納入定義。更適當的做法是在 III 和 IV 中以及有理基的定義中都要求有理數係數。若僅把所有代數數的域 $\mathsf{K}$ 替換為所有有理數的域，§§ 6 和 7 的論證逐字保持不變。}；同時，在至少一個良序下，沒有一個基元素能由有限多個在它之前的基元素以 $\mathsf{K}$ 中係數有理地表示。”\footnote{與線性基和超越基相反，這裡元素的次序起作用。例如，在次序 $\eta,\sqrt{\eta}$ 中，兩個元素都必須收入；而在相反次序中，只須收入 $\sqrt{\eta}$。}

% unit: noether.p09.par.049
把由整超越數形成的一個領域 $\mathfrak G$ 置於良序 $\Omega$ 之下。於是可能發生這樣的情形：$\mathfrak G$ 中的某個量 $z$ 能由 $\mathfrak G$ 中有限多個在它之前的量有理表示；也就是說，它滿足方程 $z=\varphi(\zeta)$，其中 $\varphi(\zeta)$ 是 $\xi_1,\ldots,\xi_t$ 的、以 $K$ 為係數的有理函數。特別地，該領域中的每個代數數 $\alpha$ 都滿足方程 $z=\alpha$。$\mathfrak G$ 中剩餘的量 $\vartheta$ 組成一個系統 $\Theta$，需要證明它是 $\mathfrak G$ 中所有量的有理基。因為按假設，每個量 $\vartheta$ 都與 $\Theta$ 中在它之前的量有理無關。歸納論證\footnote{參見 Zermelo, loc. cit., § 1。}還表明，每個關係 $z=\varphi(\zeta)$ 都導致一個關係 $z=\psi(\vartheta)$；否則就必須存在第一個元素 $z_0=\varphi_0(\zeta_1\cdots\zeta_t)$，它不能由 $\vartheta$ 有理表示，而在它之前的 $\xi_1,\ldots,\xi_t$ 又都是 $\vartheta$ 的有理函數；這就矛盾。由 II，$\Theta$ 同時是所有數的有理基；由 III 和 I，除 $\Theta$ 外，$\mathfrak G$ 還含有由所有 $\vartheta$ 的整多項式構成的整域 $[\Theta]$ 作為子環；而由 IV，$[\Theta]$ 不含任何代數分式數。因此，所有數的有理基 $\Theta$ 滿足一個附加條件：由 $\Theta$ 導出的整域 $[\Theta]$ 不含任何代數分式數；\footnote{這確實是一個條件，例如由下述事實可見：兩個量 $\eta,\frac1{2\sqrt{\eta}}$ 不能作為基元素，而 $\eta,\frac1{\sqrt{\eta}}$ 可以。}$\Theta$ 是所有數的一個\emph{帶附加條件的有理基}。於是，與 § 2 類比：\emph{每個領域 $\mathfrak G$ 都可藉助一個帶附加條件的有理基來構造，使得 $\mathfrak G$ 含有整域 $[\Theta]$ 作為子環。}

% unit: noether.p09.par.050
令 $\Theta$ 為所有數的任意一個帶附加條件的有理基——其存在性因而由領域 $\mathfrak G$ 的存在得到保證。令 $\mathfrak N(\mathfrak G)$ 表示所有包含 $\Theta$、因而也包含 $[\Theta]$ 的領域 $\mathfrak G$ 的全體。當 $\Theta$ 遍歷每個可能的帶附加條件的基時，由上文可知，$\mathfrak N(\mathfrak G)$ 遍歷全部領域 $\mathfrak G$；所以我們又可限於考察 $\mathfrak N(\mathfrak G)$ 中的領域 $\mathfrak G$。

% unit: noether.p09.par.051
現在 $[\Theta]$ 本身就是一個領域 $\mathfrak G$；因為每個數都可由 $\Theta$ 有理表示，因而可表示為 $[\Theta]$ 中兩個多項式的商；並且按其構造，$[\Theta]$ 也滿足其餘條件 I、III、IV。\emph{因此 $[\Theta]$ 是 $\mathfrak N(\mathfrak G)$ 中所有領域 $\mathfrak G$ 的最大公子環，或者說交；而且 $\mathfrak N(\mathfrak G)$ 對取交封閉，即 $\mathfrak N(\mathfrak G)$ 的任一子系統中所有領域 $\mathfrak G$ 的交仍屬於 $\mathfrak N(\mathfrak G)$。}因為 I 對交成立；II 和 III 由 $[\Theta]$ 是子環得出；IV 則由交領域作為子環包含在一個領域 $\mathfrak G$ 中得出。

% unit: noether.p09.par.052
把 $\vartheta$ 的有理函數的任意系統 $\mathfrak S$ 有理整地附加到 $[\Theta]$——由於有理基的性質，每個 $\vartheta$ 的代數函數都能通過某些其他 $\vartheta$ 有理表示——得到一個具有性質 I、II、III 的整域 $[\Theta,\mathfrak S]$；只要 $\mathfrak S$ 是“可容許系統”，也就是說，只要它滿足沒有代數分式數通過這種有理整附加進入該領域這個條件，它就成為一個領域 $\mathfrak G$。反過來，$\mathfrak N(\mathfrak G)$ 中的每個領域 $\mathfrak G$ 都允許表示為 $[\Theta,\mathfrak S]$，只須令 $\mathfrak S$ 表示剩餘系統 $\mathfrak G-[\Theta]$。於是，對最一般領域 $\mathfrak G$，完全類似於 § 4，有：
\begin{enumerate}
\item[a)] \emph{$\mathfrak N(\mathfrak G)$ 中所有領域 $\mathfrak G$ 的交由整域 $[\Theta]$ 給出，而 $[\Theta]$ 本身是一個領域 $\mathfrak G$；$\mathfrak N(\mathfrak G)$ 對取交封閉。}
\item[b)] \emph{$\mathfrak N(\mathfrak G)$ 中的每個領域 $\mathfrak G$ 都由把可容許系統 $\mathfrak S$ 有理整地附加到 $[\Theta]$ 而產生。當 $\Theta$ 遍歷每個可能的帶附加條件的有理基時，$\mathfrak N(\mathfrak G)$ 遍歷全部領域 $\mathfrak G$。}\footnote{對於最一般的可容許系統 $\mathfrak S$，§ 4 所說的話仍適用，只須把“代數”換成“有理”。}
\end{enumerate}

% unit: noether.p09.par.053
\emph{註記。} 若從 $\Theta$ 中抽取出 $\Theta$ 中所有量的一個超越基 $\mathsf{H}$，那麼 $\mathsf{H}$ 同時也是所有數的一個超越基；反過來，每個超越基都可通過把 $\mathsf{H}$ 置於良序之首而擴張為一個有理基。相應地，§ 5 中藉助超越基構造領域 $\mathfrak G$ 可以解釋如下。把要附加到 $[\mathsf{H}]$ 的系統 $\mathfrak S$ 分解為兩個子系統 $\mathfrak S_1$ 和 $\mathfrak S_2$；附加 $\mathfrak S_1$ 相當於把 $\mathsf{H}$ 擴張為一個帶附加條件的有理基，隨後再把一個可容許系統 $\mathfrak S_2$ 有理整地附加上去。

% unit: noether.p09.sec.007
\section*{§ 7. 帶附加條件的最一般有理基的構造。}

% unit: noether.p09.par.054
前一節的結果還需要補充，因為那裡證明了所有數的帶附加條件有理基的\emph{存在}，但沒有給出由一個給定良序構造它的方法。由於有附加條件，當給出的不是領域 $\mathfrak G$ 而是全體複數的域時，這個構造具有更復雜的形式。\footnote{在 § 6 b) 中，只讓 $\Theta$ 遍歷那些由超越基和基量的整代數函數組成的特殊有理基就足夠了；這時附加條件自動滿足。為看出這一點，只需按 § 6 把從 $\mathfrak G$ 中抽出的超越基擴張為有理基，然後在所得基中，把每個 $\eta$ 的代數分式函數乘以一個適當選擇的 $\eta$ 的多項式，由此保持基性質。這個構造可不加改變地轉移到全體複數的域。從 § 6 產生的關於帶附加條件的最一般有理基的問題，不過也有獨立的興趣。}

% unit: noether.p09.par.055
把全體複數的域置於一個良序 $\Omega$ 之下。那麼可能出現三類量：
\begin{enumerate}
\item[1.] 量 $y$，其對在它之前的基量——在 3. 中遞歸定義\footnote{如歸納論證所示，這樣的遞歸定義是可能的。}——的有理整附加，迫使一個代數分式數進入由此得到的整域；如果在它之前沒有基量，則這個整域由 $y$ 的所有整多項式組成。特別地，所有代數分式數都是量 $y$，因為所得整域含有量 $y=\beta$。
\item[2.] 量 $z$，它們不具有性質 1，但可由有限多個在它們之前且不同於 $y$ 的數有理表示；因此它們滿足關係 $z=\varphi(\zeta_1\cdots\zeta_t)$，其中 $\varphi$ 是以 $\mathsf{K}$ 為係數的有理函數，並且沒有一個 $\xi$ 是 $y$。特別地，所有代數整數 $\alpha$ 都屬於此類，因為性質 1 對它們不成立，而它們滿足關係 $z=\alpha$。
\item[3.] 量 $\vartheta$，1 和 2 對它們都不成立，並且它們應稱作基量。特別地，良序中的第一個超越數 $\vartheta_0$ 就是這樣的基量；因為 $\vartheta_0$ 的整多項式不能表示代數分式數，而 $\vartheta_0$ 也不能有理依賴於在它之前的代數數。
\end{enumerate}

% unit: noether.p09.par.056
\emph{現在要證明，所有基量 $\vartheta$ 的系統 $\Theta$ 是所有數的一個帶附加條件的有理基。}

% unit: noether.p09.par.057
由於沒有量 $y$ 出現在 $\vartheta$ 之中，整域 $[\Theta]$ 不可能含有任何代數分式數；所以附加條件成立。又由於沒有量 $z$ 出現在 $\vartheta$ 之中，每個 $\vartheta$ 都與在它之前的基量有理無關。歸納論證完全如 § 6 所示：每個關係 $z=\varphi(\zeta)$——其中沒有 $y$ 出現在 $\xi$ 之中——都導致一個關係 $z=\psi(\vartheta)$，所以所有量 $z$ 都能由 $\vartheta$ 有理表示。

% unit: noether.p09.par.058
剩下只需證明較困難的斷言：$y$ 也能由 $\vartheta$ 有理表示。首先證明 $y$ 代數依賴於 $\vartheta$。按假設，對每個 $y$ 至少存在一個代數分式數 $\beta$，它具有表示
% unit: noether.p09.eq.025
\begin{equation}
\beta=f_0(\vartheta)+f_1(\vartheta)y+\cdots+f_\varkappa(\vartheta)y^\varkappa,
\tag{1}
\end{equation}
其中 $f_0(\vartheta)\cdots f_\varkappa(\vartheta)$ 是來自 $[\Theta]$ 的多項式，因而不能表示代數分式數。如果 $y$ 與 $\vartheta$ 代數無關，那麼 (1) 將作為 $y$ 的恆等式成立，並且會有 $\beta=f_0(\vartheta)$，這與 $[\Theta]$ 的性質矛盾。

% unit: noether.p09.par.059
為了由代數依賴推出有理依賴，我們從 $\Theta$ 中抽取出 $\Theta$ 中所有量的一個超越基 $\mathsf{H}$。於是每個 $y$ 作為 $\vartheta$ 的代數函數，也代數依賴於 $\eta$；由 $y=y(\eta)$ 因而確定了一個在 $(\mathsf{H})$ 上的代數域 $(\mathsf{H},y(\eta))$。由於所有 $\eta$ 同時也是量 $\vartheta$，證明有理依賴就等於證明對每個這樣的域，都存在一些本原元，它們是 $\vartheta$ 的有理函數；更準確地說，將證明 $y(\eta)$ 在乘以一個 $\eta$ 的多項式以後失去性質 1，因而變成這樣的一個元素。

% unit: noether.p09.par.060
已知在 $(\mathsf{H})$ 與 $(\mathsf{H},y)$ 之間只有有限多箇中間域。令 $\Omega_0=(H),\Omega_1,\ldots,\Omega_\sigma$ 為其中那些具有可由 $\vartheta$ 有理表示的本原元的中間域，從而該域的每個元素都成為 $\vartheta$ 的有理函數；這個條件至少總是對 $\mathfrak L_0=(\mathsf{H})$ 滿足。設屬於 $\mathfrak L_i$ 的本原函數由 $\frac{g_i(\vartheta)}{h_i(\vartheta)}$ 給出，其中 $g_i,h_i$ 是 $[\Theta]$ 中的多項式。那麼——$\alpha_i$ 表示一個適當選擇的整數——$[\Theta]$ 中的多項式
% unit: noether.p09.eq.026
\[
 \zeta_i=g_i(\vartheta)+\alpha_i h_i(\vartheta)
\]
是 $\mathfrak L_i$ 的本原函數，或是 $\mathfrak L_i$ 上某個代數域的本原函數，\footnote{參見 § 5 末尾。}並且 $\mathfrak L_i$ 中的每個函數 $\psi(\eta)$ 都有表示
% unit: noether.p09.eq.027
\begin{equation}
\psi(\eta)=
\frac{a_1(\eta)\zeta_i^{\varkappa-1}+a_2(\eta)\zeta_i^{\varkappa-2}+\cdots+a_\varkappa(\eta)}
     {a(\eta)}
=\frac{A(\eta,\zeta_i)}{a(\eta)},
\tag{2}
\end{equation}
其中 $a(\eta),a_1(\eta),\ldots,a_\varkappa(\eta)$ 是 $\eta$ 的整多項式，因此 $A(\eta,\zeta_i)$ 和 $a(\eta)$ 是 $[\Theta]$ 中的量。

% unit: noether.p09.par.061
$y$ 相對於 $\mathfrak L_i$ 所滿足的不可約方程，按 (2) 具有形式
% unit: noether.p09.eq.028
\[
 f^{(i)}(\eta)y^{\lambda_i}+f_1^{(i)}(\eta,\zeta_i)y^{\lambda_i-1}
+\cdots+f_{\lambda_i}^{(i)}(\eta,\zeta_i)=0,
\]
並且該方程的所有係數都屬於 $[\Theta]$。於是函數 $y_i=f^{(i)}(\eta)\cdot y$ 滿足相對於 $\mathfrak L_i$ 同樣不可約的方程
% unit: noether.p09.eq.029
\[
 y_i^{\lambda_i}+f_1^{(i)}(\eta,\zeta_i)y_i^{\lambda_i-1}
+\cdots+\bigl(f^{(i)}(\eta)\bigr)^{\lambda_i-1}
 f_{\lambda_i}(\eta,\zeta_i)=0,
\]
由此它代數整地依賴於 $[\mathsf{H},\zeta_i]$；$y_i$ 與任意 $\eta$ 的多項式的乘積也是如此。特別地，函數
% unit: noether.p09.eq.030
\begin{equation}
Y=f^{(0)}(\eta)f^{(1)}(\eta)\cdots f^{(\sigma)}(\eta)\cdot y
\tag{3}
\end{equation}
通過在每種情形下相對於 $\mathfrak L_i$ 不可約的方程
% unit: noether.p09.eq.031
\begin{equation}
Y^{\lambda_i}+F_1^{(i)}(\eta,\zeta_i)Y^{\lambda_i-1}
+\cdots+F_{\lambda_i}^{(i)}(\eta,\zeta_i)=0
\qquad (i=0,1,\ldots,\sigma)
\tag{4}
\end{equation}
代數整地依賴於所有領域 $[\mathsf{H},\zeta_i]$。現在我們證明，$Y$ 就是所要求的、可由 $\vartheta$ 有理表示的本原元。比如，假設代數數 $\gamma$ 由 $Y$ 和 $[\Theta]$ 表示為
% unit: noether.p09.eq.032
\begin{equation}
\gamma=k_0(\vartheta)+k_1(\vartheta)Y+\cdots+k_\nu(\vartheta)Y^\nu,
\tag{5}
\end{equation}
其中 $k_0(\vartheta)\cdots k_\nu(\vartheta)$ 是 $[\Theta]$ 中的多項式。現在確定 $Y$ 相對於由 (5) 的係數域產生的域 $\mathfrak K=(\mathsf{H};k_0(\vartheta)\cdots k_\nu(\vartheta))$ 所滿足的不可約方程。已知它由 $Y$ 相對於 $\mathfrak K$ 與 $(\mathsf{H},Y)$ 的交域的不可約方程給出，該交域位於 $(\mathsf{H})$ 和 $(\mathsf{H},Y)$ 之間。由 (3)，$(\mathsf{H},Y)$ 與 $(\mathsf{H},y)$ 相同；又因為 $\mathfrak K$ 的所有元素，因而交域的所有元素，都能由有限多個 $\vartheta$ 有理表示，所以這個交域必然是 $\mathfrak L_0,\mathfrak L_1,\cdots,\mathfrak L_\sigma$ 中的一個，設為 $\mathfrak L_\tau$。因此所求不可約方程按 (4) 具有次數 $\lambda_\tau$，併為
% unit: noether.p09.eq.033
\begin{equation}
Y^{\lambda_\tau}+F_1^{(\tau)}(\eta,\zeta_\tau)Y^{\lambda_\tau-1}
+\cdots+F_{\lambda_\tau}^{(\tau)}(\eta,\zeta_\tau)=0
\tag{6}
\end{equation}
其係數是 $[\Theta]$ 中的多項式。利用 (6)，(5) 可改寫為
% unit: noether.p09.eq.034
\begin{equation}
\gamma=K_0(\vartheta)+K_1(\vartheta)Y+\cdots+
K_{\lambda_\tau-1}(\vartheta)Y^{\lambda_\tau-1},
\tag{7}
\end{equation}
其中 $K_0(\vartheta),K_1(\vartheta),\cdots$ 屬於 $\mathfrak K$，另一方面又是 $[\Theta]$ 中的多項式。但由於 $\gamma$ 是代數數，關係 (7) 在 $Y$ 中達不到次數 $\lambda_\tau$，因而作為 $Y$ 的\emph{恆等式}成立。所以
% unit: noether.p09.eq.035
\[
 \gamma=K_0(\vartheta),
\]
即 $\gamma$ 屬於 $[\Theta]$，因而是代數整數。

% unit: noether.p09.par.062
當 $\gamma$ 遍歷所有可由 $Y$ 和 $[\Theta]$ 表示的代數數時，$(\mathsf{H},Y)$ 與相應域 $\mathfrak K$ 的交域始終只遍歷 $\mathfrak L_0,\mathfrak L_1,\cdots,\mathfrak L_\sigma$，前面的所有論證仍然有效。也就是說，\emph{通過把 $Y$ 有理整地附加到 $[\Theta]$，不能表示出任何代數分式數}。因此 $Y$ 不再具有性質 1；按照 2) 或 3)，它是一個量 $z$ 或 $\vartheta$，所以能由有限多個 $\vartheta$ 有理表示；由 (3)，$y$ 也如此：\emph{$\Theta$ 已被證明是所有數的有理基。}反過來，由於每個帶附加條件的有理基都可以通過把 $\Theta$ 置於良序之首而按 1), 2), 3) 構造出來，我們證明了：\emph{由 1), 2), 3) 給出的構造導向所有數的最一般帶附加條件的有理基。}

% unit: noether.p09.sec.008
\section*{§ 8. 領域 $\mathfrak G$ 和 $\mathfrak H$ 的類劃分。\footnote{參見 E. Steinitz, Algebraische Theorie der Körper. Crelles Journal 137 (1910)，尤其是 §§ 23、24。（§ 23 已證明超越基的存在性。）}}

% unit: noether.p09.par.063
在最一般領域 $\mathfrak G$ 和 $\mathfrak H$ 的問題之外，還出現了\emph{本質上不同}的領域的問題；這些領域在下文將更精確說明的意義下，不能由同一個構造原則生成。這導致把這些領域劃分為類。

% unit: noether.p09.par.064
\emph{如果兩個領域的元素能夠一一對應，使得一個領域中任意兩個元素的和與積總是對應於另一個領域中相應元素的和與積，那麼這兩個領域稱為屬於同一類，或稱為同構。}由此定義可知，在每個同構關係下，零和單位元都對應自身，有限多個元素之間的一切代數關係都對相應元素保持，反之亦然。特別地，除單位元外，所有有理數都對應自身；而代數數全體——同樣代數整數全體——變到自身，雖然某個個別代數數完全可能對應於它的一個共軛。

% unit: noether.p09.par.065
首先注意，包含固定超越基 $H$ 的全部領域 $\mathfrak G$ 的全體 $\mathfrak M(\mathfrak G)$ 中的每個領域，都同構於屬於另一個超越基 $\mathsf{H}^*$ 的全體 $\mathfrak M^*(\mathfrak G)$ 中的一個領域。因為全體複數的域由每個超越基經代數擴張得到，因而具有與基相同的勢；\footnote{Steinitz, §§ 23 和 24。}所以 $\mathsf{H}$ 和 $\mathsf{H}^*$ 具有同樣的勢，而所需同構映射通過把基量 $\eta_i$ 對應到 $\eta_i^*$ 得到。\footnote{這個結論當然不適用於屬於兩個有理基 $\Theta$ 和 $\Theta^*$ 的集合 $\mathfrak N(\mathfrak G)$ 和 $\mathfrak N^*(\mathfrak G)$；也不能由 $\Theta$ 和 $\Theta^*$ 的相互有理可表示性推出 $\mathfrak N(\mathfrak G)$ 與 $\mathfrak N^*(\mathfrak G)$ 的同構。不過得到的是域 $(\Theta)$ 與 $(\Theta^*)$ 之間的同構關係。}另一方面，$\mathfrak M(\mathfrak G)$ 含有非同構的、或者說本質上不同的領域 $\mathfrak G$，這由 §§ 2、3 和 5 的例子得出。因為在同構映射下所有代數關係都保持，超越基必須再對應於超越基；但是 §§ 2 和 3 表明，函數泛函領域和模領域沒有一個超越基能使 Zermelo 領域的所有基性質都得到滿足。由於 § 2 的函數泛函領域是代數整閉的，子集 $\mathfrak M(\mathfrak H)$ 也含有本質上不同的領域 $\mathfrak H$。

% unit: noether.p09.par.066
現在我們仿照基的構造說明，怎樣從 $\mathfrak M(\mathfrak G)$ 中抽取一個子集 $\mathfrak L(\mathfrak G)$，使得 $\mathfrak L(\mathfrak G)$ 中所有領域 $\mathfrak G$ 都本質上不同，而 $\mathfrak M(\mathfrak G)$ 中的每個領域——從而任意一個領域 $\mathfrak G$——都同構於 $\mathfrak L(\mathfrak G)$ 中的一個領域。把 $\mathfrak M(\mathfrak G)$ 置於良序 $\Omega$ 下；那麼 $\mathfrak M(\mathfrak G)$ 中的領域 $\mathfrak G$ 可能同構於某個在它之前的領域，也可能不然。那些不同構於任何在先領域的領域 $\mathfrak G$ 組成集合 $\mathfrak L(\mathfrak G)$，按構造它只含有本質上不同的領域 $\mathfrak G$。歸納論證還表明，$\mathfrak M(\mathfrak G)$ 中的每個領域都同構於 $\mathfrak L(\mathfrak G)$ 中的一個領域。因為如果 $\mathfrak G_1$ 是第一個不滿足這一點的領域，那麼 $\mathfrak G_1$ 要麼屬於 $\mathfrak L(\mathfrak G)$，要麼同構於一個在先領域 $\mathfrak G_0$，而按假設 $\mathfrak G_0$ 又同構於 $\mathfrak L(\mathfrak G)$ 中的某個領域；所以 $\mathfrak G_1$ 也同樣如此。由於每個領域 $\mathfrak G$ 屬於某個集合 $\mathfrak M^*(\mathfrak G)$，因而同構於 $\mathfrak M(\mathfrak G)$ 中的一個領域，\emph{領域 $\mathfrak G$ 的全部類由 $\mathfrak L(\mathfrak G)$ 表示。}

% unit: noether.p09.par.067
若在 $\mathfrak L(\mathfrak G)$ 中的一個領域 $\mathfrak G$ 內，把基量 $\eta$ 換成另一個基的那些 $\eta^*$，則如上所示，得到一個與 $\mathfrak G$ 同構的領域。反過來，令 $\mathfrak G^*$ 為任意領域，並假設它同構於 $\mathfrak L(\mathfrak G)$ 中的一個領域 $\mathfrak G$。如果基量 $\eta$ 對應於 $\mathfrak G^*$ 中的量 $\eta^*$，那麼這些量又組成一個超越基，而 $\mathfrak G^*$ 含有與 $\mathfrak G$ 所含的 $\eta$ 的代數函數相同的 $\eta^*$ 的代數函數——或其共軛。\emph{因此，與 $\mathfrak G$ 同構的每個領域，都通過使超越基 $H$ 遍歷 $\mathfrak G$ 中及其相對於 $(\mathsf{H})$ 的共軛中的每個可能的超越基而得到}\footnote{這些與 $\mathfrak G$ 共軛的領域的元素，固然可由 $\mathfrak G$ 的元素按 II 有理表示，但不一定能整有理地表示；因此這些共軛可以不同於 $\mathfrak G$。}\emph{——如果這些共軛與 $\mathfrak G$ 不同。}由 $\mathfrak L(\mathfrak G)$ 中每個固定領域 $\mathfrak G_i$，以這種方式得到類 $\mathfrak K_i$，它包含全部且僅包含同構於 $\mathfrak G_i$ 的領域 $\mathfrak G$，雖然每個領域可能出現多次，甚至無限多次。由 $\mathfrak K_i$ 又可按上述程序抽取一個子集 $\mathfrak K_i^*$，其中每個同構於 $\mathfrak G_i$ 的領域恰出現一次。當 $\mathfrak K_i^*$ 遍歷所有類時，得到全部領域 $\mathfrak G$，每個恰出現一次。為了刻畫類 $\mathfrak K_i$，當然也可以不用 $\mathfrak G_i$，而用該類中任一其他領域，只要 $\mathfrak K_i$ 的所有領域都可由它按上法導出。綜上：

% unit: noether.p09.par.068
\emph{可以從 $\mathfrak M(\mathfrak G)$ 中抽取一個子集 $\mathfrak L(\mathfrak G)$，使得 $\mathfrak L(\mathfrak G)$ 中的領域 $\mathfrak G$ 與全部類一一對應（無重複）。由 $\mathfrak G_i$ 出發，相應類 $\mathfrak K_i$ 通過使 $H$ 遍歷 $\mathfrak G_i$ 中及其相對於 $(\mathsf{H})$ 的共軛中的所有可能超越基而產生。如果 $\mathfrak K_i^*$ 表示 $\mathfrak K_i$ 中兩兩不同的全部領域，並且 $\mathfrak K_i^*$ 遍歷所有類，那麼全部領域 $\mathfrak G$ 出現，每個恰出現一次。}\footnote{這裡 $\mathfrak M(\mathfrak G)$ 應按 §§ 6 和 7 構造：按 § 7，把超越基 $\mathsf{H}$ 擴張為每個由它產生的、帶附加條件的有理基 $\Theta$；並對每個 $\Theta$，按 § 6 b) 形成包含 $\Theta$ 的全部領域 $\mathfrak G$ 的全體 $\mathfrak N(\mathfrak G)$。}

% unit: noether.p09.par.069
對於由代數整超越數形成的領域 $\mathfrak H$，有相應的陳述。

% unit: noether.p09.sec.009
\section*{§ 9. 領域 $\mathfrak G$ 的可數無窮多個類的一個例子。}

% unit: noether.p09.par.070
§ 4 b) 和 § 6 b) 給出的構造的勢，由相應註記可看出等於所有良序的勢；若 $c$ 表示連續統的勢，則它為 $2^c$。但由此不能對只計一次的全部領域 $\mathfrak G$ 的勢作出結論，更不能對全部類的勢作出結論。現在將通過一個可數無窮多個領域 $\mathfrak G$ 的例子——類似於 § 5 的那些領域——說明類的數目當然不是有限的；這些領域將全部證明為本質上不同，從而給出可數無窮多個類。\footnote{§ 6 的領域 $\mathfrak G$ 同樣給出可數無窮多個類，但證明稍複雜。}

% unit: noether.p09.par.071
與 § 5, (1) 類比，令領域 $\mathfrak G_\sigma$ 由全部量
% unit: noether.p09.eq.036
\begin{equation}
 z\equiv a_0+a_1(\eta)+\cdots+a_{\sigma-1}(\eta)
 \;\mathrm{mod.}\;\mathfrak M_\sigma(\eta),
\tag{1}
\end{equation}
組成；\emph{其中 $a_i(\eta)$ 表示 $\eta$ 的 $i$ 次齊次型；而模 $\mathfrak M_\sigma$ 以 $\eta$ 中所有 $\sigma$ 次冪積為基，}\footnote{自然，在每個單獨的 $z$ 中，模裡只出現有限多個 $\eta$ 的冪積。}\emph{並以所有 $\eta$ 的整代數函數為乘子。}我們證明，在兩個領域 $\mathfrak G_\sigma$ 與 $\mathfrak G_\tau$ $(\sigma\ne\tau)$ 之間的同構關係下，模 $\mathfrak M_\sigma$ 與 $\mathfrak M_\tau$ 必然同構地對應；於是矛盾很容易推出。

% unit: noether.p09.par.072
記 $\mathfrak G_\tau$ 中的不定元為 $\xi$；在同構關係下，令 $\mathfrak G_\sigma$ 中的量 $f(\eta)$ 對應於這些 $\xi$。由於同構關係在取商時保持，由所有 $\xi$ 的整代數函數組成的領域 $\mathfrak G_\zeta$ 對應於一個代數整閉領域 $\mathfrak T$，它由所有在 $f(\eta)$ 上整代數的函數組成，因而也在 $\eta$ 中整代數。

% unit: noether.p09.par.073
假設 (1) 對應於 $\mathfrak M_\tau$ 中的一個量。那麼，由 $\mathfrak M_\tau$ 的模性質，(1) 與 $\mathfrak T$ 中任意量 $\psi(\eta)$ 的乘積也必須屬於 $\mathfrak G_\sigma$；用公式表示為：
% unit: noether.p09.eq.037
\begin{equation}
(a_0+a_1(\eta)+\cdots+a_{\sigma-1}(\eta))\cdot\psi(\eta)
\equiv b_0+b_1(\eta)+\cdots+b_{\sigma-1}(\eta)
\;\mathrm{mod.}\;\mathfrak M_\sigma.
\tag{2}
\end{equation}
令所有 $\eta$ 等於零，得 $a_0\cdot\psi(0)=b_0$，於是 (2) 變為
% unit: noether.p09.eq.038
\begin{equation}
(a_0+a_1(\eta)+\cdots+a_{\sigma-1}(\eta))(\psi(\eta)-\psi(0))
\equiv c_1(\eta)+\cdots+c_{\sigma-1}(\eta)
\;\mathrm{mod.}\;\mathfrak M_\sigma.
\tag{3}
\end{equation}
現在，除 $\psi(\eta)$ 外，對每個 $\nu$，函數
% unit: noether.p09.eq.039
\[
 \chi_\nu(\eta)=\sqrt[\nu]{\psi(\eta)-\psi(0)}
\]
也屬於 $\mathfrak T$；且由於 $\chi_\nu(0)=0$，與 (3) 類似的公式為
% unit: noether.p09.eq.040
\begin{equation}
(a_0+a_1(\eta)+\cdots+a_{\sigma-1}(\eta))\chi_\nu(\eta)
\equiv c_1^{(\nu)}(\eta)+\cdots+c_{\sigma-1}^{(\nu)}(\eta)
\pmod{\mathfrak M_\sigma}
\quad(\nu=1,2,\cdots).
\tag{4}
\end{equation}
令 $A(\eta)$，次數為 $\lambda$，表示 $\chi_1=\psi(\eta)-\psi(0)$ 與其 $\varkappa-1$ 個共軛的乘積。由於 $\chi_1(0)=0$，$A(\eta)$ 必須以次數至少為一的項開始。由 (4) 得
% unit: noether.p09.eq.041
\begin{equation}
a_0\cdot\chi_\nu(\eta)\equiv0\pmod{\mathfrak M_1};\qquad
a_0^\nu\cdot\chi_1(\eta)\equiv0\pmod{\mathfrak M_\nu};\qquad
a_0^{\nu\varkappa\cdot}A(\eta)\equiv0\pmod{\mathfrak M_{\nu\varkappa}},
\tag{5}
\end{equation}
因而當 $a_0\ne0$ 時，如 § 3 那樣有 $\lambda>\nu\varkappa$。但 $\nu<\frac{\lambda}{\varkappa}$ 與 $\mathfrak T$ 的代數整閉性矛盾，所以必有 $a_0=0$。由 (4) 現在得到
% unit: noether.p09.eq.042
\begin{equation}
a_1(\eta)\cdot\chi_\nu(\eta)\equiv c_1^{(\nu)}(\eta)\pmod{\mathfrak M_2};\qquad
[a_1(\eta)]^{\nu\varkappa}\cdot A(\eta)
\equiv[c_1^{(\nu)}(\eta)]^{\nu\varkappa}\pmod{\mathfrak M_{\nu\varkappa+1}},
\tag{6}
\end{equation}
由此，由於 $A(\eta)$ 以次數至少為一的項開始，得 $c_1^{(\nu)}(\eta)=0$。於是，完全類似於 (5)，對每個 $\nu$ 有
% unit: noether.p09.eq.043
\[
 a_1(\eta)\cdot\chi_\nu(\eta)\equiv0\pmod{\mathfrak M_2};\qquad
 [a_1(\eta)]^{\nu\varkappa}A(\eta)\equiv0\pmod{\mathfrak M_{2\nu\varkappa}},
\]
從而 $a_1(\eta)=0$。繼續這樣做，交替使用 (5) 與 (6)，得到
% unit: noether.p09.eq.044
\[
 a_0=0,\qquad a_1(\eta)=0,\quad \ldots,\quad a_{\sigma-1}(\eta)=0.
\]
由於關於 $\mathfrak G_\tau$ 的類似考察也成立：\emph{在 $\mathfrak G_\sigma$ 與 $\mathfrak G_\tau$ $(\sigma\ne\tau)$ 之間的每個同構關係下，模 $\mathfrak M_\sigma$ 與 $\mathfrak M_\tau$ 都同構地對應。}

% unit: noether.p09.par.074
例如設 $\sigma>\tau$。與不定元 $\xi$ 對應的是 $\mathfrak G_\sigma$ 中的量 $f(\eta)$；又因 $\mathfrak G_\tau$ 中所有量都能表示為 $\xi$ 的多項式模 $\mathfrak M_\tau$，$\mathfrak M_\tau$ 與 $\mathfrak M_\sigma$ 的對應意味着 $\mathfrak G_\sigma$ 中所有量都可由 $f(\eta)$ 的多項式模 $\mathfrak M_\sigma$ 得到。因為 $\sigma>\tau\geqq1$，不定元 $\eta$ 當然不屬於 $\mathfrak M_\sigma$，所以既然它們能由 $f(\eta)$ 模 $\mathfrak M_\sigma$ 整有理地表示出來，就必有某個 $f(\eta)$ 具有非零線性項。因此令
% unit: noether.p09.eq.045
\[
 \zeta:f(\eta)\equiv a_0+a_1(\eta)\pmod{\mathfrak M_2}\quad[a_1(\eta)\ne0],
\]
% unit: noether.p09.eq.046
\[
 \zeta^\tau:[f(\eta)]^\tau\equiv a_0^\tau\;\mathrm{mod.}\;\mathfrak M_1\quad\text{當 }a_0\ne0,
\]
% unit: noether.p09.eq.047
\[
 \hphantom{\zeta^\tau:[f(\eta)]^\tau}\equiv [a_1(\eta)]^\tau
 \;\mathrm{mod.}\;\mathfrak M_{\tau+1}\quad\text{當 }a_0=0.
\]
現在 $\zeta^\tau$ 屬於 $\mathfrak M_\tau$，而 $[f(\eta)]^\tau$ 以非零的零次項或 $\tau$ 次項開始，並且由於 $\tau<\sigma$，它不可能屬於 $\mathfrak M_\sigma$，這與上面得到的 $\mathfrak M_\sigma$ 與 $\mathfrak M_\tau$ 的對應矛盾。情形 $\tau>\sigma$ 同時通過交換 $\eta$ 與 $\xi$ 處理。因此證明了領域 $\mathfrak G_\sigma$ 與 $\mathfrak G_\tau$ \emph{不同構，或者說本質上不同}。於是，\emph{領域 $\mathfrak G_i$ $(i=1,2,\cdots\text{ 無窮})$ 給出了領域 $\mathfrak G$ 的可數無窮多個類的一個例子。}

% unit: noether.p09.par.075
\emph{註記。}雖然任意兩個領域 $\mathfrak G_i$ 都不同構，但對兩個領域來說，完全可能每一個都同構於另一個的一部分。\footnote{這種現象也可以發生在域中；參見 Steinitz, loc. cit., 結尾。}例如令 $\tau=1$，$\sigma$ 任意：
% unit: noether.p09.eq.048
\[
\mathfrak G_1:a_0+\mathfrak M_1(\zeta);\qquad
G_\sigma:a_0+a_1(\eta)+\cdots+a_{\sigma-1}(\eta)+\mathfrak M_\sigma(\eta).
\]
賦值 $\zeta=\eta$ 使 $\mathfrak G_\sigma$ 成為 $\mathfrak G_1$ 的真子整環；反過來，在 $\mathfrak G_1$ 與 $\mathfrak G_\sigma$ 中由 $\zeta=\eta^\sigma$、$\eta=\sqrt[\sigma]{\zeta}$ 給出的一一對應，使 $\mathfrak G_1$ 成為 $\mathfrak G_\sigma$ 的真子整環。因此每個 $\mathfrak G_\sigma$ 也同構於自身的一個真子整環。與交的概念相反，交類的概念——即這樣一種類：其中每個領域都同構於每個其他類的每個領域的真子整環——並不存在，至少不是一個唯一確定的概念。

% unit: noether.p09.sec.010
\section*{§ 10. 任意域的整量。}

% unit: noether.p09.par.076
前面各節的展開，對領域 $\mathfrak G$ 來說——只要不把代數數包括在領域 $\mathfrak G$ 的定義性質 III 和 IV 中——只依賴於全體複數形成一個域這一事實；\footnote{只有在 § 7 中還使用了有理數“素域”包含無窮多個元素這一附加假設；然而在只含有限多個元素的素域情形下，如我們將看到的，這一節簡單地消失。}對領域 $\mathfrak H$ 來說，還依賴於這個域是代數閉的。所得結果因而直接允許\emph{推廣到最一般的域}；這種域依 Weber 和 E. Steinitz\footnote{loc. cit., 引言。}抽象地定義為“具有加法和乘法兩種運算的元素系統，這兩種運算服從結合律和交換律，由分配律相聯繫，並且允許不受限制且唯一的逆運算。”\footnote{只須排除除以零。}

% unit: noether.p09.par.077
每個這樣的域都含有一個單位元，並且因而含有由加法、乘法及其逆運算從單位元導出的“素域”。這個素域或者是有理數型（特徵 $0$），或者是某個素數 $p$ 的剩餘類系統型（特徵 $p$）。\footnote{Steinitz, § 4。}\emph{素域的整量}於是明確定義為由單位元通過加法和減法得到的量；對於特徵 $p$ 的素域，整量已經窮盡素域，並且不存在素域的分式量。每個代數閉域都含有一個“絕對代數域”，它由素域附加所有在素域上代數的元素而得到；對於特徵 $0$ 的域，它因而是所有代數數的域這一類型。\emph{絕對代數域的整代數量}於是明確定義為所有在素域的整量上（或者等價地，在單位元上）整代數的量；對於特徵 $p$ 的域，因此也不存在絕對代數域的代數分式量。在這些預備說明之後，“整”量和“代數整”量的定義可分別轉移到任意域和代數閉域：\emph{域 $\mathfrak K$ 的整量定義為 $\mathfrak K$ 中量的一個整域 $\mathfrak G$，其商域\footnote{也就是由 $\mathfrak G$ 中成對量的所有商組成的域。}窮盡 $\mathfrak K$，並且它含有單位元但不含素域的任何分式量。}

% unit: noether.p09.par.078
\emph{代數閉域 $\mathfrak L$ 的代數整量定義為 $\mathfrak L$ 中量的一個代數整閉整域 $\mathfrak H$，其商域窮盡 $\mathfrak L$，並且它含有單位元但不含絕對代數域的任何代數分式量。}

% unit: noether.p09.par.079
對於特徵零的域，只要把“代數數”分別替換為素域的量和絕對代數域的量，前幾節所得結果逐字適用。對於特徵 $p$ 的域，結果還要進一步簡化，因為素域不含任何分式量，所以有理基以及待附加系統的附加條件簡單地消失。\footnote{參見本節第一條注。}於是結果為：\emph{抽象定義允許作為特徵 $p$ 的域的整量的是：由任意有理基導出的每個整域、域本身，以及處於這些領域和該域之間的每個整域。相應地，作為素特徵代數閉域的代數整量，有：處於由超越基導出的領域和該域本身之間的每個代數整閉整域——包括兩個邊界領域。}因此，對於最一般的素特徵域，和相應素域一樣，整與分式的區別消失。

% unit: noether.p09.close.001
\medskip
\noindent Erlangen，1915 年 3 月 30 日。
%
\clearpage
%
% unit: noether.p10.title.001
\begin{center}
\pdfbookmark[0]{10. 同構映射的函數方程}{paper-10}
{\Large\bfseries 10. 同構映射的函數方程}\par
\vspace{1.25em}
\emph{Math. Ann.} 77 (1916), pp. 536--545
\end{center}

\vspace{2em}

% unit: noether.p10.par.001
按照 Dedekind，\footnote{參見例如 Dirichlet-Dedekind，\emph{Vorlesungen über Zahlentheorie}（第 4 版），補編 XI，\S\ 161。仿照群論而來的“同構映射”或“同構”這一名稱，只見於較晚的文獻；Dedekind 說的是“一個域的置換”。}所謂域 $\Afield$ \emph{到一個系統} $\Bfield$ \emph{上的同構映射} --- 該系統事後也證明為一個域 --- 是指這樣一個單值對應：$\Afield$ 的每一個元素都對應 $\Bfield$ 的一個且僅一個元素；並且 $\Afield$ 中任意兩個元素的和、差、積、商，總是又被對應為 $\Bfield$ 中那兩個唯一對應元素的和、差、積、商。

% unit: noether.p10.par.002
設這樣的映射由一個 --- 依上文為單值的 --- 函數 $f(z)$ 來實現。若 $z$ 遍歷域 $\Afield$ 的所有元素 $x,y,\cdots$，那麼 $f(z)$ 由下列函數方程來刻畫：
% unit: noether.p10.eq.001
\[
\begin{array}{ll}
(1)\quad f(x+y)=f(x)+f(y); & (2)\quad f(x-y)=f(x)-f(y);\\[0.45em]
(3)\quad f(x\cdot y)=f(x)\cdot f(y); & (4)\quad f\!\left(\dfrac{x}{y}\right)=\dfrac{f(x)}{f(y)}.
\end{array}
\]

% unit: noether.p10.par.003
\emph{下文將給出這些函數方程的最一般單值解} $f(z)$，\emph{其中} $x,y,\cdots$ \emph{遍歷所有實數和複數；也就是給出實現全體複數之域的同構映射的最一般函數} $f(z)$。\footnote{這個問題是 Landau 先生偶然向我提出的。正如我後來注意到的，解的一部分已經見於 H. Lebesgue, \emph{Sur les transformations ponctuelles, transformant les plans en plans}, Atti Acad. Torino, 1906/07；並且也隱含地見於 A. Ostrowski, \emph{Über einige Fragen der allgemeinen Körpertheorie}, Crelle's Journal 143 (1913), \S\ 2。}對於任意域，解法完全類似。

% unit: noether.p10.par.004
G. Hamel\footnote{G. Hamel, \emph{Eine Basis aller Zahlen und die unstetigen Lösungen der Funktionalgleichung:} $f(x+y)=f(x)+f(y)$, Math. Ann. 60, p. 459 (1905).}藉助所有數的一個\emph{線性}基，構造了函數方程 (1) 的最一般解。這裡給出的 (1) 至 (4) 的解的構造 --- 因而也是使有理運算和代數運算保持不變的映射函數的構造 --- 使用所有數的\emph{有理}基和\emph{超越}基。在第 1 節中較詳細地討論同構映射之後，第 2 節給出 $f(z)$ 的必要條件；第 3 節中這些條件被看出也是充分的，並導向實際構造。\footnote{參見我論文中的平行考慮：\emph{Die allgemeinsten Bereiche aus ganzen transzendenten Zahlen}. Math. Ann. 77, p. 103 (1915)，特別是 \S\S\ 8 和 10。}除已知的例外 $f(z)=z$ 或 $\overline z$ 外，(1) 至 (4) 的解是不連續的；第 4 節還將進一步說明它們是“極端不連續的” --- 這一事實由 G. Hamel 對 (1) 的實解證明過，但在複數範圍內並不一定成立。\footnote{Hamel 的術語“完全不連續”在其餘文獻中以另一種意義使用。}

% unit: noether.p10.sec.001
\medskip
\noindent\textbf{1.}\quad 我們先從函數方程 (1) 至 (4) 推出若干後果，並且直接針對一般域 $\Afield$。由 (1)，單值性給出 $f(0)=0$。若 $f(y)$ 不為零，則原來的 $y$ 也不可能為零，於是函數方程 (1) 至 (4) 表明，所有值 $f(z)$ 構成的\emph{映射系統} $\Bfield$ \emph{又形成一個域}。反過來，如果預先規定初始條件 $f(0)=0$，那麼 (2) 推出函數 $f(z)$ 的單值性：\emph{單值性要求與初始條件} $f(0)=0$ \emph{因而是等價的}。由 (4) 可知 $f(z)$ 不可能恆等為零；由 (3) 還可知 $f(z)$ 只在 $z=0$ 時為零。因為若 $f(y)=0$ 且 $y\ne0$，由於 $\frac{z}{y}$ 也屬於 $\Afield$，就會推出
% unit: noether.p10.eq.002
\[
f(z)=f\!\left(\frac{z}{y}\cdot y\right)=f\!\left(\frac{z}{y}\right)\cdot f(y)=0,
\]
這與 (4) 相矛盾，因為它使 $f(z)$ 恆等為零。因此由 $f(x)=f(y)$ 得出 $x=y$；也就是說，對 $f(z)$ 的每一個值也對應着一個且僅一個 $z$ 的值：\emph{$f(z)$ 是} $z$ \emph{的可逆單值函數}。由函數方程可見，由反函數實現的映射又是同構的。由於每一個有理運算都由有限次加法、乘法及其逆運算組成，也可以這樣說：\emph{同構映射在} $\Afield$ \emph{與} $\Bfield$ \emph{的元素之間建立了一個一一關係，使得} $\Afield$ \emph{中任意有限多個元素之間成立的所有有理關係}
% unit: noether.p10.eq.003
\[
F(x,y,\cdots,t)=0,
\]
\emph{在} $\Bfield$ \emph{中保持，反之亦然}。

% unit: noether.p10.par.005
還應指出，對於一個域，$f(z)$ 已經由函數方程 (1) 和 (3)、$f(z)$ 不恆等為零這一要求以及初始條件 $f(0)=0$ 所刻畫。事實上，(2) 由 (1) 得出，只需把 $x$ 替換為屬於 $\Afield$ 的元素 $(x-y)$；然後 (2) 和 $f(0)=0$ 給出 $f(z)$ 的單值性。如上所示，再由 (3)，把 $x$ 替換為屬於 $\Afield$ 的元素 $\frac{x}{y}$，可知當 $y\ne0$ 時 $f(y)$ 不可能為零，從而得到 (4) 的有效性，同時也得到一一性。如果以整域 $\Sdomain$ 而不是域為基礎，則 $\frac{x}{y}$ 未必屬於 $\Sdomain$，而且不能由 (3) 推出一一性。然而在這裡，下述最熟悉的表述已經足夠：同構映射是 $\Sdomain$ 與 $\Sdomain'$ 的元素之間的\emph{一一}對應，使得和與積總是分別對應和與積。\footnote{關於第 1 節，參見 Dedekind，前引處，\S\ 161。}

% unit: noether.p10.sec.002
\medskip
\noindent\textbf{2.}\quad 從現在起，為簡單起見，取所有實數和複數的域 $\Cfield$ 代替 $\Afield$ 為基礎。首先要討論的是 $f(z)$ 所取的值系統。由 (3) 或 (4) 得出 $f(1)=1$，於是由函數方程，對於每個\emph{有理數} $\alpha$，$f(\alpha)=\alpha$；\emph{因此每個有理數在映射下對應它自身}。由於\emph{代數數} $\beta$ 滿足一個有有理數係數的不可約方程 $F(\beta,\alpha)=0$，第 1 節和剛才的觀察給出關於 $f(\beta)$ 的方程 $F(f(\beta),\alpha)=0$；\emph{因此每個代數數變為它自身或它的一個共軛}，並且由於一一性，代數數全體又對應於所有代數數的域。為了認識\emph{超越}數的行為，同時分離共軛值，我們取所有數的一個\emph{有理基}\footnote{參見我所引關於“整個超越數”的論文，\S\ 6。}為基礎，也就是一個由數 $\vartheta$ 組成的系統 $\Theta$，它允許所有數以有理方式 --- 帶有理數係數\footnote{我所謂“有理函數”總是指其係數也為有理數的函數。} --- 表示出來；同時在至少一種良序下，沒有任何基元素能由有限多個在它之前的基元素有理地表示出來。這個基的存在性完全像線性基和超越基一樣由良序定理得出。

% unit: noether.p10.par.006
令 $\Cfield$ 服從一個良序 $\Omega$。於是每一個數或者能由有限多個前面的數有理表示，或者與前面的數有理無關。後者形成基 $\Theta$，而歸納推理表明確實每個數都能由有限多個 $\vartheta$ 有理表示。

% unit: noether.p10.par.007
對每一個基元素 $\vartheta$，定義“截段域 $\Kfield(\vartheta)$”：它由在良序中先於 $\vartheta$ 的那些基元素的所有有理組合組成。第一個基元素的截段域應視為全體有理數的域 $\Rfield$。數 $\vartheta$ 相對於 $\Kfield(\vartheta)$ 可表現出兩種行為：它可以代數依賴於 $\Kfield(\vartheta)$，也可以代數獨立於它。由於按第 1 節，$\Cfield$ 中所有成立的關係在映射域中也成立，反之亦然，所以對應於各個 $\vartheta$ 的值 $f(\vartheta)$ 全體形成映射域的一個基，並藉助 $\Theta$ 自身的良序而良序化。特別地，截段域 $\Kfield(\vartheta)$ 對應於截段域 $\Kfield(f(\vartheta))$；而 $\vartheta$ 相對於 $\Kfield(\vartheta)$ 是代數依賴還是代數獨立，$f(\vartheta)$ 相對於 $\Kfield(f(\vartheta))$ 也同樣如此。在依賴的情形中，$\vartheta$ 與 $f(\vartheta)$ 的不可約方程也相互對應。每次相對於 $\Kfield(\vartheta)$ 代數獨立的那些值 $\vartheta$ 的全體 --- 如歸納推理再次表明 --- 形成\emph{所有數的一個超越基}，\footnote{參見 E. Zermelo, \emph{Über ganze transzendente Zahlen}, \S\ 1, Math. Ann. 75 (1914)。}即一個由數 $\eta$ 組成的系統 $\mathsf{H}$，它允許所有數代數地表示出來，而 $\eta$ 本身之間不存在代數關係。在映射域中，與這個超越基 $\mathsf{H}$ 對應的又是一個超越基 $\mathsf{Z}$；由於一一性，它與 $\mathsf{H}$ 有相同的勢，即連續統的勢，因為眾所周知代數擴張不會增大勢。\footnote{參見 E. Steinitz, \emph{Algebraische Theorie der Körper}, Crelle's Journal 137 (1910), \S\ 24。}然而，一旦知道各個值 $f(\vartheta)$，便立即可得所有值系統上的 $f(z)$，因為由 $z=\psi(\vartheta)$ 總是推出 $f(z)=\psi(f(\vartheta))$。

% unit: noether.p10.sec.003
\medskip
\noindent\textbf{3.}\quad 現在我們證明，上面得到的必要條件實際上也給出了 $f(z)$ 的構造。與超越基 $\mathsf{H}$ 一一對應的基 $\mathsf{Z}$ 按如下方式構造並指定給 $\mathsf{H}$。設置第二個良序 $\Omega'$，它給出所有數的一個超越基 $\mathsf{Z}'$。令以 $\xi$ 為元素的 $Z$ 是 $\mathsf{Z}'$ 的一個子集，其勢與 $\mathsf{Z}'$ 相同，因而也與 $\mathsf{H}$ 相同。對 $H$ 的每一個元素 $\eta_i$，按如下規則指定一個且僅一個具有相同編號的元素 $\xi_i$：$\xi_i$ 是良序 $\Omega'$ 下 $\mathsf{Z}$ 中所有尚未指定給 $H$ 中先於 $\eta_i$ 的任何元素的那些元素中的\emph{第一個}。

% unit: noether.p10.par.008
現在如下定義 $f(z)$：
\begin{enumerate}
\item[(a)] \emph{對所有有理數} $\alpha$，\emph{由} $f(0)=0$, $f(\alpha)=\alpha$ \emph{定義}；
\item[(b)] \emph{對超越基} $\mathsf{H}$ \emph{的所有量，由} $f(\eta_i)=\xi_i$ \emph{定義}；
\item[(c)] \emph{對有理基} $\Theta$ \emph{的所有其餘量} $\vartheta$ --- \emph{即那些代數依賴於截段域} $\Kfield(\vartheta)$ \emph{、因而滿足在} $\Kfield(\vartheta)$ \emph{中不可約的方程} $F(\vartheta;\vartheta_{i_1}\cdots\vartheta_{i_\tau})=0$ \emph{的量} --- \emph{作為方程} $F(f(\vartheta);f(\vartheta_{i_1})\cdots f(\vartheta_{i_\tau}))=0$ \emph{的一個根來定義}；
\item[(d)] \emph{對所有其餘的值} $z$，\emph{這些值按} $\Theta$ \emph{的定義可表示為} $z=\psi(\vartheta_{k_1}\cdots\vartheta_{k_\sigma})$，\emph{由} $f(z)=\psi(f(\vartheta_{k_1})\cdots f(\vartheta_{k_\sigma}))$ \emph{定義}。\footnote{在 (d) 中，(a) 作為特殊情形包含在內。}
\end{enumerate}

% unit: noether.p10.par.009
為了證明這樣定義的 $f(z)$ 確實滿足函數方程 (1) 至 (4)，根據第 1 節，只需證明 (1) 和 (3)，因為 $\Cfield$ 是一個域，並且 $f(z)$ 由於 (a) 和 (b) 不能恆等為零，同時還滿足初始條件 $f(0)=0$。藉助有理基 $\Theta$，把 $x,y,x+y$ 表示為
% unit: noether.p10.eq.004
\[
x=\psi_1(\vartheta_1\cdots\vartheta_t),\quad
 y=\psi_2(\vartheta_1\cdots\vartheta_t),\quad
 (x+y)=\psi_3(\vartheta_1\cdots\vartheta_t).
\]
於是證明 $f(z)$ 滿足函數方程 (1)，即
% unit: noether.p10.eq.005
\[
f(x)+f(y)-f(x+y)=0,
\]
按 (d) 等同於證明：由
% unit: noether.p10.eq.006
\[
\psi_1(\vartheta)+\psi_2(\vartheta)-\psi_3(\vartheta)
 =\frac{H_1(\vartheta)}{H_2(\vartheta)}=0
\]
總是推出
% unit: noether.p10.eq.007
\[
\psi_1(f(\vartheta))+\psi_2(f(\vartheta))-\psi_3(f(\vartheta))
 =\frac{H_1(f(\vartheta))}{H_2(f(\vartheta))}=0,
\]
其中 $H_1,H_2$ 表示以其變元為整有理函數。函數方程 (3) 也導向完全相同形式的條件；因此所有函數方程被滿足，等同於下述條件的成立 --- 該條件按第 1 和第 2 節也是必要的：對於\emph{每一個}整有理函數 $H(\vartheta)$，由 $H(\vartheta)=0$ 總是推出 $H(f(\vartheta))=0$，並且由 $H(\vartheta)\ne0$ 推出 $H(f(\vartheta))\ne0$，後一項是由於分母的出現。

% unit: noether.p10.par.010
這確實成立，可由歸納證明。設在 $H(\vartheta_1\cdots\vartheta_t)$ 中，基量 $\vartheta_t$ 是在良序中位於 $\vartheta_1\cdots\vartheta_t$ 中最後的一個。\footnote{那些在各 $\vartheta$ 中次數為零的表達式 $H(\vartheta)$，在映射下變為其自身，因此無需進一步研究。}於是 $H(\vartheta_1\cdots\vartheta_t)=0$ 可以看作 $\vartheta_t$ 相對於截段域 $\Kfield(\vartheta_t)$ 所滿足的一個關係；同樣，$H(f(\vartheta_1)\cdots f(\vartheta_t))=0$ 可以看作 $f(\vartheta_t)$ 相對於 $\Kfield(f(\vartheta_t))$ 所滿足的一個關係。如果並非所有關係 $H(\vartheta)=0$ 對 $f(\vartheta)$ 也成立，且反之也成立，那麼必定存在第一個基量，比如 $\vartheta_0$，對於它，至少有一個相對於 $\Kfield(\vartheta_0)$ 成立的關係沒有在映射域中保持，或者反過來沒有保持，而在它之前的截段域 $\Kfield(\vartheta_0)$ 和 $\Kfield(f(\vartheta_0))$ 已經是同構的。特別地，由 (a)，$\Theta$ 的第一個基量的截段域，即全體有理數的域 $\Rfield$，與它的映射域同構。現在設 $H(\vartheta_0)$ 具有形式
% unit: noether.p10.eq.008
\[
H(\vartheta_0)=a_0(\vartheta)\cdot\vartheta_0^{\varkappa}
+a_1(\vartheta)\cdot\vartheta_0^{\varkappa-1}+\cdots+a_{\varkappa}(\vartheta),
\]
其中 $a_i(\vartheta)$ 是 $\Kfield(\vartheta_0)$ 中的量。於是有兩種可能。

% unit: noether.p10.par.011
1) $\vartheta_0$ 代數獨立於 $\Kfield(\vartheta_0)$。於是 $H(\vartheta_0)=0$ 關於 $\vartheta_0$ 恆等地成立，\footnote{因為 $x,y,\cdots$ 可以用各 $\vartheta$ 表示，所以這種形式的關係完全可能出現。}從而所有係數都有 $a_i(\vartheta)=0$，並且由於 $\Kfield(\vartheta_0)$ 與 $\Kfield(f(\vartheta_0))$ 的同構，也有 $a_i(f(\vartheta))=0$。因此由 $H(\vartheta_0)=0$ 總是推出 $H(f(\vartheta_0))=0$。反過來，設 $H(f(\vartheta_0))=0$。由於 $\vartheta_0$ 代數獨立於 $\Kfield(\vartheta_0)$，它屬於超越基 $\mathsf{H}$；由 (b)，$f(\vartheta_0)$ 屬於超越基 $\mathsf{Z}$，並且代數獨立於 $\Kfield(f(\vartheta_0))$。因此所有 $a_i(f(\vartheta))$ 都為零，並且由同構，所有 $a_i(\vartheta)$ 都為零。由 $H(f(\vartheta_0))=0$ 推出 $H(\vartheta_0)=0$，或者由 $H(\vartheta_0)\ne0$ 也推出 $H(f(\vartheta_0))\ne0$。

% unit: noether.p10.par.012
2) $\vartheta_0$ 代數依賴於 $\Kfield(\vartheta_0)$。於是 $H(\vartheta_0)$ 可被相對於 $\Kfield(\vartheta_0)$ 不可約的方程 $F(\vartheta_0)$ 整除；也就是說，關於 $t$ 恆等地有
% unit: noether.p10.eq.009
\[
H(t;a_i(\vartheta))=F(t,\vartheta)\cdot G(t,\vartheta_i).
\]
由於出現的所有係數都屬於 $\Kfield(\vartheta_0)$，在同構域 $\Kfield(f(\vartheta_0))$ 中相應的關係也成立：
% unit: noether.p10.eq.010
\[
H(t;a_i(f(\vartheta)))=F(t;f(\vartheta))\cdot G(t;f(\vartheta)).
\]
但按 (c)，$f(\vartheta_0)$ 被選為 $F(t;f(\vartheta))$ 的一個根。因此由 $H(\vartheta_0)=0$ 推出 $H(f(\vartheta_0))=0$。

% unit: noether.p10.par.013
反過來，若 $H(f(\vartheta_0))=0$，那麼 $H(t;a_i(f(\vartheta)))$ 被函數 $F(t;f(\vartheta))$ 整除；該函數相對於 $\Kfield(f(\vartheta_0))$ 不可約。這又對同構域 $\Kfield(\vartheta_0)$ 推出 $H(t;a_i(\vartheta))$ 被 $F(t;\vartheta)$ 整除。並且由於 $F(t;f(\vartheta))$ 是由 $\vartheta_0$ 的不可約方程通過同構關係得到的，而此關係是一一的，$F(t;\vartheta)$ 必然表示以 $\vartheta_0$ 為根的不可約函數。因此由 $H(f(\vartheta_0))=0$ 推出 $H(\vartheta_0)=0$，或者由 $H(\vartheta_0)\ne0$ 也推出 $H(f(\vartheta_0))\ne0$。

% unit: noether.p10.par.014
因此，從同構的截段域 $\Kfield(\vartheta_0)$ 和 $\Kfield(f(\vartheta_0))$ 出發，分別添入 $\vartheta_0$ 和 $f(\vartheta_0)$ 後，也得到同構的域；假定這對某個 $\vartheta_0$ 失敗，就會導致矛盾。於是證明了：\emph{由 (a) 至 (d) 定義的函數} $f(z)$ \emph{確實表示函數方程 (1) 至 (4) 的一個解，並且按第 2 節，是最一般的解}。

% unit: noether.p10.par.015
用於構造 $f(z)$ 的所有考慮，只用到了全體 $\Cfield$ 數形成一個域這一事實；因此它們適用於每一個抽象定義的域 $\Afield$。這裡應注意，有理數域 $\Rfield$ 要由 $\Afield$ 的“素域”來代替，即由 $\Afield$ 的單位元通過加法、乘法及其逆運算所導出的域。因此，如果在 (a) 至 (d) 中把有理數替換為素域的元素，那麼\emph{由此得到的函數} $f(z)$ \emph{就實現任意抽象定義的域的最一般同構映射}。

% unit: noether.p10.sec.004
\medskip
\noindent\textbf{4.}\quad 還需證明，為全體實數和複數的域 $\Cfield$ 所定義的函數 $f(z)$ --- 除 $f(z)=z$ 或 $\overline z$ 外，已知它們是不連續的 --- 是\emph{極端不連續的}；也就是說，在任意實數或複數 $z_0$ 的每一個鄰域中，都有一些值 $z$ 使 $f(z)$ 任意接近任意預先給定的值 $Z_0$。如 G. Hamel 前引處所證明的，這種極端不連續性屬於定義在所有\emph{實}數上的函數方程 (1) 的\emph{實}解 $\varphi(x)$。但是，如下列例子所示，它不一定屬於定義在複數上的解。Hamel 的論證如下。如果 $\varphi(x)$ 不同於 $C\cdot x$，那麼線性基中至少有兩個值，比如 $x_1$ 和 $x_2$，\footnote{所有（實）數的線性基由一個（實）數系統給出，它允許所有（實）數以有理係數線性表示，而在任意有限多個基元素之間不存在有理係數的線性關係。}使得 $\varphi(x_1):x_1\ne\varphi(x_2):x_2$。於是兩個方程
% unit: noether.p10.eq.011
\[
\alpha x_1+\beta x_2=a;\qquad \alpha\varphi(x_1)+\beta\varphi(x_2)=b
\]
在實數 $\alpha,\beta$ 中有解；因而 $a$ 和 $b$ 可以由有理數 $\alpha,\beta$ 任意逼近。由於 $\alpha,\beta$ 是有理數，$\alpha\varphi(x_1)+\beta\varphi(x_2)$ 於是實際上就是 $\varphi(\alpha x_1+\beta x_2)$。這個推理對復值系統失效。事實上，在複數域中也可以給出 (1) 的一些解，它們不連續但並非極端不連續；例如
% unit: noether.p10.eq.012
\[
\varphi(z)=\varphi(x+iy)=\varphi(x)+i(y+\varphi(x)),
\]
其中 $\varphi(x)$ 表示一個在實數域中極端不連續的實解；或者更簡單地，$\varphi(z)=\varphi(x)$。在兩種情形中，$\varphi(z)$ 的實部都是極端不連續的，而虛部由實部共同決定。

% unit: noether.p10.par.016
如果引入秩的概念，這種行為，同時也是函數 $f(z)$ 的行為，就會完全清楚。令 $z=x+iy$，$\varphi(z)=X+iY$，並按照不存在、存在一個、兩個或三個線性無關關係
% unit: noether.p10.eq.013
\[
c_1x+c_2y+c_3X+c_4Y=0,
\]
分別稱 $\varphi(z)$ 具有四秩、三秩、二秩或一秩；這裡 $c$ 當然是實數。

% unit: noether.p10.par.017
1) 設 $\varphi(z)$ 具有\emph{四秩}。那麼線性基中至少存在四個值，例如 $z_1,z_2,z_3,z_4$，使得行列式
% unit: noether.p10.eq.014
\[
\begin{vmatrix}
x_1&x_2&x_3&x_4\\
y_1&y_2&y_3&y_4\\
X_1&X_2&X_3&X_4\\
Y_1&Y_2&Y_3&Y_4
\end{vmatrix}
\]
不為零。因此方程
% unit: noether.p10.eq.015
\[
\begin{aligned}
\alpha x_1+\beta x_2+\gamma x_3+\delta x_4&=a,\\
\alpha y_1+{}\cdot{}\ \cdot{}\ \cdot{}+\delta y_4&=b,\\
\alpha X_1+{}\cdot{}\ \cdot{}\ \cdot{}+\delta X_4&=A,\\
\alpha Y_1+{}\cdot{}\ \cdot{}\ \cdot{}+\delta Y_4&=B
\end{aligned}
\]
在實數 $\alpha,\cdots,\delta$ 中有解，並且 $a,b,A,B$ 可以由有理數 $\alpha,\beta,\gamma,\delta$ 任意逼近。此外有
% unit: noether.p10.eq.016
\[
\alpha(X_1+iY_1)+\beta(X_2+iY_2)+\gamma(X_3+iY_3)+\delta(X_4+iY_4)
 =\varphi(\alpha z_1+\beta z_2+\gamma z_3+\delta z_4).
\]
因此，“一般地說”，解 $\varphi(z)$ 在複數域中也是極端不連續的；$z_0=a+ib$ 鄰域中的 $\varphi(z)$ 的\emph{不連續值}形成一個\emph{二維線性流形}。

% unit: noether.p10.par.018
2) 設 $\varphi(z)$ 具有\emph{三秩}。那麼線性基中至少存在三個值系統，例如 $z_1,z_2,z_3$，使得上述行列式的秩為三。滿足關係
% unit: noether.p10.eq.017
\[
c_1a+c_2b+c_3A+c_4B=0,
\]
的那些值 $a,b,A,B$，且只有這些值，能夠被任意逼近；不連續值形成一個由 $z_0=a+ib$ 確定的\emph{一維線性流形}。如所給例子所示，這種情形確實可能發生。

% unit: noether.p10.par.019
3) 設 $\varphi(z)$ 具有\emph{二秩}。由於總可以選擇兩個複數 $z_1,z_2$，使得
% unit: noether.p10.eq.018
\[
\begin{vmatrix}x_1&x_2\\ y_1&y_2\end{vmatrix}\ne0,
\]
各關係具有形式
% unit: noether.p10.eq.019
\[
X=\lambda_1x+\lambda_2y;\qquad Y=\mu_1x+\mu_2y.
\]
由此得到的表達式
% unit: noether.p10.eq.020
\[
\varphi(z)=(\lambda_1x+\lambda_2y)+i\cdot(\mu_1x+\mu_2y)
\]
恰好給出複數範圍內函數方程 (1) 的最一般\emph{連續}解。

% unit: noether.p10.par.020
4) 設 $\varphi(z)$ 具有\emph{一秩}。這種情形在複數範圍內不可能發生；在實數範圍內，它意味着\emph{連續}解 $\varphi(x)=C\cdot x$。

% unit: noether.p10.par.021
現在我們證明，定義在所有實數和複數上的函數方程 (1) 至 (4) 的\emph{解} $f(z)$ \emph{只能具有二秩或四秩}。這裡二秩只給出兩個\emph{連續}函數 $f(z)=z$ 或 $\overline z$，這是由特化 $f(1)=1$ 和 $f(i)=\pm i$ 結合 3) 看出的。由於一秩已經被排除，還需排除三秩。因此我們假設至少存在一個關係
% unit: noether.p10.eq.021
\[
c_1x+c_2y+c_3X+c_4Y=0,
\]
使得 $f(z)$ 具有三秩，或者可能具有更低的秩，並證明在這個假設下總是後一種情形發生。如特化 $f(1)=1$、$f(i)=\pm i$ 再次表明，關係變為
% unit: noether.p10.eq.022
\[
c_3(X-x)+c_4(Y\mp y)=0,
\]
其中 $c_3$ 和 $c_4$ 不同時為零。先設，例如，$c_4=0$。關係 $(X-x)=0$ 意味着純虛的 $z$ 也對應純虛的 $f(z)$。因此對每個實數 $\mathfrak{y}$，有 $f(i\mathfrak{y})=i\cdot\Yreal$，其中 $\Yreal$ 仍為實數。由此並且由於
% unit: noether.p10.eq.023
\[
f(iy)=\pm i f(y),
\]
還得到 $f(y)=\pm\Yreal$；於是每個實的 $z$ 也對應一個實的 $f(z)$，並且由於 $(X-x)$，所有實值都映為自身。因此只有 $f(z)=z$ 或 $\overline z$，也就是實際上為二秩。當 $c_3=0$, $c_4\ne0$ 時，結論完全類似。

% unit: noether.p10.par.022
現在設 $c_3$ 和 $c_4$ 都不為零，於是關係可以寫成形式
% unit: noether.p10.eq.024
\[
(Y\mp y)=c\cdot(X-x)
\]
其中 $c$ 是實的、非零且有限的。這意味着當 $y=\pm cx$ 時，也有 $Y=c\cdot X$。因此，考慮到函數方程，對每個實數 $\mathfrak{x}$ 有
% unit: noether.p10.eq.025
\[
f\{\mathfrak{x}\cdot(1\pm ic)\}=X(1+ic)=f(\mathfrak{x})\cdot(1+i f(c)),
\]
其中 $X$ 仍為實數。但這裡按照第 1 節，$f(\mathfrak{x})$ 的因子不可能恆等為零，作除法得到
% unit: noether.p10.eq.026
\[
f(\mathfrak{x})=X(\sigma+i\tau)
\]
其中 $\sigma$ 和 $\tau$ 為實數，並且與 $\mathfrak{x}$ 和 $X$ 無關。特別地，對 $\mathfrak{x}=1$ 也有一個實的 $\mathfrak{X}_0$ 與之相屬，而條件 $1=\mathfrak{X}_0(\sigma+i\tau)$ 表明必然 $\tau=0$，於是 $f(\mathfrak{x})=\sigma\cdot X$。因此每個實的 $z$ 也對應一個實的 $f(z)$；等價地，由 $y=0$ 必然推出 $Y=0$。於是對於實的 $z$，線性關係變為 $X-x=0$；\footnote{這裡使用了 $c\ne0$ 這一事實，因而 $c_3$ 和 $c_4$ 都不為零，而前面只用了 $c_4\ne0$；所以必須預先處理這兩個數中有一個為零的情形。}每個實值對應它自身。再次只有 $f(z)=z$ 或 $\overline z$，所以實際上是二秩。由此在所有情形中排除了三秩。但是，正如前幾節的研究所表明的，不連續解確實存在，而這些解因而必然具有四秩；故由 1) 得到定理：\emph{“函數方程 (1) 至 (4) 的所有解} $f(z)$，\emph{除連續解} $f(z)=z$ \emph{或} $\overline z$ \emph{之外，在實數和複數範圍內都是極端不連續的。”}\footnote{Lebesgue 在前引處表明，$f(z)$ 的實部和虛部都是 $x$ 與 $y$ 的“不可測”函數；如第 4 節開頭的例子 $\varphi(z)=\varphi(x)+i(y+\varphi(x))$ 所示，這還不能推出在複數範圍內的極端不連續性。}

% unit: noether.p10.close.001
\medskip
\noindent Erlangen，1915 年 10 月 30 日。
%
\clearpage
%
% unit: noether.p11.title.001
\begin{center}
\pdfbookmark[0]{11. 具有給定群的方程}{paper-11}
{\Large\bfseries 11. 具有給定群的方程}\par
\vspace{1.2em}
\emph{Math. Ann.} 78 (1918), pp. 221--229
\end{center}

\vspace{1.5em}

% unit: noether.p11.par.001
構造具有給定群的方程這一問題，可以從兩個方向來處理；可簡要稱為“無理的”、即刻畫根的方向，和“有理的”、即刻畫係數的方向。

% unit: noether.p11.par.002
在“無理的”方向中，即以函數論和算術方式工作的方向中，有 Kronecker 定理：有理數域上的所有 Abel 域都是圓分域；以及關於二次數域的相對 Abel 域的相應定理。這些定理一方面給出所有 Abel 群相對於這些特殊有理性域的可實現性；另一方面，它們給出方程的全體，並同時對由這些方程定義的數域的結構給出更深的洞察。但是每一個單獨的定理都限於一個特殊的有理性域；每一個新的有理性域都要求完全重新處理。

% unit: noether.p11.par.003
“有理的”、代數地工作的方向，則以 Hilbert 的不可約性定理為基礎；按該定理，在係數表示中可以限於參數表示。例如，關於任意有理性域，存在任意多的具有交錯群的方程，就屬於這裡。當然，在這裡缺少上面所說的關於由方程定義的域的結構的洞察；不過這個“有理”方向的優點在於：群的可實現性是同時對所有有理性域證明的，雖然迄今已知的參數表示一般並不給出方程的全體。\footnote{對於可解方程，係數的參數表示可以由根的參數表示取代。最近 F. Mertens 對某些 8 次群給出了最一般的根表示：“Gleichungen mit Quaternionengruppen”, Sitzb. d. Ak. d. Wiss., Wien, Abt. IIa, vol. 125, p. 735。據我從這篇札記中得知，G. Bucht 已經更早給出了三次、四次方程以及具有上述群的八次方程的總能構造的標準形的最一般根表達式：“Über einige algebraische Körper achten Grades”, Arkiv for Math., Astron. och Physik, vol. 6, no. 30. [校樣時補記。]}

% unit: noether.p11.par.004
以下是對“有理”方向的一個貢獻：我們給出充分條件，使得對於任意有理性域，存在參數表示來給出具有給定群的方程全體。這些條件是：屬於該群的不變量域，即 Lagrange 的 \emph{Gattungsbereich}，可以由該域的獨立函數有理地表示，也就是說它具有一個“最小基”；換言之，它同構於 $n$ 個不定元的所有有理函數的域。

% unit: noether.p11.par.005
如在 \S 2 中還要說明的，這一點例如對三次和四次方程中出現的所有群都成立。關於這些三次和四次方程的實際構造和更詳細討論，可見 F. Seidelmann 的論文\footnote{\emph{Die Gesamtheit der kubischen und biquadratischen Gleichungen mit Affekt bei beliebigem Rationalitätsbereich}. Erlangen, 1916.}；該論文的摘錄緊接在本通信之後。

% unit: noether.p11.par.006
Lagrange 屬域的最小基問題，正如在論文“Körper und Systeme rationaler Funktionen”（Math. Ann. 76）中已經提到的，是 E. Fischer 在同我的談話中獨立於這裡闡述的聯繫而提出的，並給了本研究以推動。\footnote{參見預告通信“Rationale Funktionenkörper”的第 2 部分，Jahresb. d. d. Math. Vereinig., vol. 22, 1913.}

% unit: noether.p11.sec.001
\subsection*{\S 1. 參數表示的充分條件。}

% unit: noether.p11.par.007
\textbf{1.} 把具有給定群的方程的構造化歸為參數表示的可能性，來自 Hilbert 的不可約性定理；它說如下。設方程
% unit: noether.p11.eq.001
\begin{equation}\tag{1}
f(x)=x^n+F_1(\lambda_1,\ldots,\lambda_r)x^{n-1}+\cdots+F_n(\lambda_1,\ldots,\lambda_r)=0
\end{equation}
其中 $F_1(\lambda_1,
\ldots,\lambda_r),\ldots,F_n(\lambda_1,
\ldots,\lambda_r)$ 是獨立參數 $\lambda_1,
\ldots,\lambda_r$ 的有理函數，係數在一個數域\footnote{也可以用一般的有理性域代替這個數域。} $\Omega$ 中，並且該方程相對於由 $\Omega$ 中係數的 $\lambda_1,
\ldots,\lambda_r$ 的所有有理函數組成的域 $\Omega(\lambda_1,
\ldots,\lambda_r)$ 具有群 $\Gamma$。那麼，可以把參數 $\lambda$ 以任意多種方式代為有理數，使得所得數值方程
% unit: noether.p11.eq.002
\begin{equation}\tag{2}
g(x)=x^n+a_1x^{n-1}+\cdots+a_n=0
\end{equation}
相對於 $\Omega$ 正好具有群 $\Gamma$。這樣的 $g(x)$ 將更精確地記作 $g_\Gamma$。

% unit: noether.p11.par.008
現在的問題是：是否能給出這樣的參數表示 (1)，使得當 $\lambda_1,
\ldots,\lambda_r$ 遍歷 $\Omega$ 中所有量、但排除那些導致群降低的 $\Omega$ 中的量時，所得數值方程 (2) 遍歷全部 $g_\Gamma$。\footnote{這裡還應把 $g(x)=0$ 出現重根的情形也計算在內。} 換句話說，非線性方程組
% unit: noether.p11.eq.003
\begin{equation}\tag{3}
F_1(\lambda_1,
\ldots,\lambda_r)=a_1;\quad \ldots;\quad F_n(\lambda_1,
\ldots,\lambda_r)=a_n
\end{equation}
應當對於每一個對應於一個 $g_\Gamma$ 的值系 $a_1,
\ldots,a_n$，在 $\Omega$ 中的量 $\lambda_1,
\ldots,\lambda_r$ 中有解。由於系統 (3) 是非線性的，要求通過一個參數表示得到全部 $g_\Gamma$，從一開始就必須加以限制。事實上，對於這樣的非線性系統，總可能出現奇異的值系 $a_1,
\ldots,a_n$，它們滿足某個代數關係 $H(a_1,
\ldots,a_n)=0$，但沒有解；\footnote{這一點不久將在一篇關於“非線性方程組的替代”的論文中一般地加以闡明。} 於是參數表示 (1) 失效，而需要一個補充表示。不過在當前情形下，這些奇異值 $a_1,
\ldots,a_n$ 可以簡單地刻畫。

% unit: noether.p11.par.009
\textbf{2.} 為了得到這樣的參數表示，我們進一步要求：$\lambda_i$ 在把根視為不定元的函數時，恆等地是屬於該群的自然無理量。這意思如下。若在 (1) 中，把參數 $\lambda_1,
\ldots,\lambda_r$ 替換為適當選擇的、關於不定元 $x_1,
\ldots,x_n$ 的有理函數 $\varphi_1(x),\ldots,\varphi_r(x)$，且其係數在 $\Omega$ 中，那麼 (1) 應轉變為
% unit: noether.p11.eq.004
\begin{equation}\tag{4}
h(x)=x^n-\sigma_1(x)x^{n-1}+\sigma_2(x)x^{n-2}\cdots \pm \sigma_n(x),
\end{equation}
其中 $\sigma_1(x),\ldots,\sigma_n(x)$ 是 $x_1,
\ldots,x_n$ 的初等對稱函數；並且 $h(x)$ 應相對於由 $\Omega(\lambda_1,
\ldots,\lambda_r)$ 這樣產生的有理性域 $\Omega(\varphi_1(x),\ldots,\varphi_r(x))$ 具有群 $\Gamma$。由於參數獨立，函數 $\varphi_1(x),\ldots,\varphi_r(x)$ 也應代數無關。

% unit: noether.p11.par.010
為了使這一要求精確化，必須闡明 $\Omega(\varphi_1(x),\ldots,\varphi_r(x))$ 與群 $\Gamma$ 的關係。我們先確定最一般的有理性域 $K$，使得 $h(x)$ 相對於它成為 $h_\Gamma$。令 $\mathfrak L_\Gamma$ 表示 $x_1,
\ldots,x_n$ 中所有具有有理整數係數、且允許 $\Gamma$ 的有理函數的域；這也稱為屬於 $\Gamma$ 的 Lagrange \emph{Gattungsbereich}，或 $\Gamma$ 的不變量域。令 $\Omega(\Omega_\Gamma)$ 表示把 $\Omega_\Gamma$ 代入 $\Omega$ 後所得的相應域。

% unit: noether.p11.par.011
因此，$\mathfrak L_\Gamma$ 是 $\mathfrak R(x_1,
\ldots,x_n)$ 的一個除子，後者是 $x_1,
\ldots,x_n$ 中有理整數係數的所有有理函數的域。按群的定義，最一般的域 $K$ 由任一滿足如下條件的域給出：它包含 $\Omega_\Gamma$ 作為除子，並且同 $\mathfrak R(x_1,
\ldots,x_n)$ 的交正好是 $\Omega_\Gamma$。相應地，對於每個含有 $\Omega$ 的域 $K$，同 $\Omega(x_1,
\ldots,x_n)$ 的交由 $\Omega(\Omega_\Gamma)$ 給出。特別地，由於按我們的要求 $\Omega(\varphi_1(x),\ldots,\varphi_r(x))$ 成為一個域 $K$，它同 $\Omega(x_1,
\ldots,x_n)$ 的交由 $\Omega(\Omega_\Gamma)$ 給出；另一方面，由於 $\Omega(\varphi_1(x),\ldots,\varphi_r(x))$ 是 $\Omega(x_1,
\ldots,x_n)$ 的一個除子，它就正好等於這個交。這個要求因此意味着，Lagrange \emph{Gattungsbereich} $\Omega(\Omega_\Gamma)$ 應由 $\Omega$ 通過添入代數無關函數 $\varphi_1(x),\ldots,\varphi_r(x)$ 得到。這裡必須有 $r=n$，因為 $\Omega_\Gamma$ 含有 $n$ 個代數無關的初等函數，而不可能有超過 $n$ 個函數代數無關。我們說 $\varphi_1(x),\ldots,\varphi_n(x)$ 構成 $\Omega(\Omega_\Gamma)$ 的一個\emph{最小基}；也說 $\Omega_\Gamma$ 具有一個係數在 $\Omega$ 中的最小基。

% unit: noether.p11.par.012
\textbf{3.} 現在關鍵的是，整個考慮可以反過來，並且確實給出所要求的參數表示的充分條件。\footnote{第 2 中的要求相對於第 1 已經擴大，因此這些條件不必是必要條件。} 設 $\varphi_1(x),\ldots,\varphi_n(x)$ 是 $\Omega(\mathfrak L_\Gamma)$ 的一個最小基，並令 $\Omega$ 為包含 $\varphi_1(x)\ldots\varphi_n(x)$ 的係數的最小域。那麼 $\Omega$ 必然是一個代數數域；特別地，它可以是所有有理數的域 $\mathfrak R$。因為 $\sigma_1(x),\ldots,\sigma_n(x)$ 屬於 $\Omega(\mathfrak L_\Gamma)$，所以存在關於 $x_1,
\ldots,x_n$ 的恆等式，其中 $F_i$ 表示其自變量的有理函數：
% unit: noether.p11.eq.005
\begin{equation}\tag{5}
\sigma_1(x)=F_1(\varphi_1(x),\ldots,\varphi_n(x));\quad\ldots;\quad
\sigma_n(x)=F_n(\varphi_1(x),\ldots,\varphi_n(x)).
\end{equation}
反過來，$\varphi_1(x),\ldots,\varphi_n(x)$，因而還有 $\varphi(x)=\mu_1\varphi_1(x)+\cdots+\mu_n\varphi_n(x)$，通過相對於 $\Omega(\sigma_1(x),\ldots,\sigma_n(x))$ 不可約的方程
% unit: noether.p11.eq.006
\begin{equation}\tag{6}
G(\varphi(x);\sigma_1(x),\ldots,\sigma_n(x))=0,
\end{equation}
代數地依賴於 $\sigma_1(x),\ldots,\sigma_n(x)$；這仍是關於 $x$ 的恆等式。

% unit: noether.p11.par.013
由 (5)，這個恆等式 (6) 變為
% unit: noether.p11.eq.007
\begin{equation}\tag{7}
G\bigl(\varphi(x);F_1(\varphi_1,
\ldots,\varphi_n),\ldots,F_n(\varphi_1,
\ldots,\varphi_n)\bigr)
=H(\varphi_1(x),\ldots,\varphi_n(x))=0.
\end{equation}
由於 $\varphi_1(x),\ldots,\varphi_n(x)$ 的代數無關性，(7) 不僅對 $x$ 恆等成立，也對 $\varphi_i$ 恆等成立；也就是說，對獨立參數 $\lambda$ 恆等成立：
% unit: noether.p11.eq.008
\begin{equation}\tag{8}
H(\lambda_1,
\ldots,\lambda_n)=G\bigl(\mu_1\lambda_1+\cdots+\mu_n\lambda_n;
F_1(\lambda_1,
\ldots,\lambda_n),\ldots,F_n(\lambda_1,
\ldots,\lambda_n)\bigr)=0.
\end{equation}
由恆等式 (8) 可推出，方程
% unit: noether.p11.eq.009
\begin{equation}\tag{9}
f(x)=x^n-F_1(\lambda_1,
\ldots,\lambda_n)x^{n-1}+F_2(\lambda_1,
\ldots,\lambda_n)x^{n-2}\cdots \pm F_n(\lambda_1,
\ldots,\lambda_n)=0
\end{equation}
相對於 $\Omega(\lambda_1,
\ldots,\lambda_n)$ 具有群 $\Gamma$。因為這個恆等式表明，已知的有理量 $\lambda=\mu_1\lambda_1+\cdots+\mu_n\lambda_n$ 是屬於 $\Gamma$ 的根的函數；但是不可能還有屬於 $\Gamma$ 的某個子群的其他函數被有理地知道，這可由代入 $\lambda_i=\varphi_i(x)$ 看出；而在同一代入之後，$\lambda$ 也不同於其所有共軛。此外，如果 $\Omega^*$ 是任一包含 $\Omega$ 的數域，則 $f(x)$ 相對於 $\Omega^*(\lambda_1,
\ldots,\lambda_n)$ 也具有群 $\Gamma$；因為按第 2，在代入 $\lambda_i=\varphi_i(x)$ 時，即使添入 $\Omega^*$，群也不降低。

% unit: noether.p11.par.014
還要證明的是，在第 1 中修正過的意義下，$f(x)$ 確實遍歷全部 $g_\Gamma$。為此注意到，由 (5) 可知，$F_i(\lambda)$ 是其自變量 $\lambda$ 的代數無關函數。因此方程組
% unit: noether.p11.eq.010
\begin{equation}\tag{10}
F_1(\lambda_1,
\ldots,\lambda_n)=-a_1;\quad F_2(\lambda_1,
\ldots,\lambda_n)=a_2;\quad\ldots;\quad F_n(\lambda_1,
\ldots,\lambda_n)=\pm a_n
\end{equation}
對於任意值系 $a_1,
\ldots,a_n$ 都有解，除了還要指出的那些奇異值之外。對於相對於 $\Omega^*$ 的某個 $g_\Gamma$ 所對應的每個值系 $a_1,
\ldots,a_n$，相應的 $\lambda_i$，作為屬於群的自然無理量，\footnote{亦參見下文。}成為 $\Omega^*$ 中的量。因此如果 $\lambda_1,
\ldots,\lambda_n$ 遍歷所有值系，那麼 $g(x)$ 在修正意義下遍歷方程的全體；其中包含了那些對應於 $\Omega^*$ 中值的全部 $g_\Gamma$。

% unit: noether.p11.par.015
反過來，由於按 Hilbert 的不可約性定理，$\Omega^*$ 中的 $\lambda$ 值也“一般地”對應於方程 $g_\Gamma$，所以在第 1 中修正過的意義下，參數表示 (9) 的確給出相對於 $\Omega^*$ 的全部 $g_\Gamma$。

% unit: noether.p11.par.016
現在還要確定使參數表示失效的奇異值系。設最小基中的函數分為分子和分母，寫成
% unit: noether.p11.eq.011a
\[
\varphi_i(x)=\frac{\psi_i(x)}{\chi_i(x)}.
\]
把它代入恆等式 (5)，設這些恆等式在不約去因子的情況下變為
% unit: noether.p11.eq.011
\begin{equation}\tag{11}
\sigma_k(x)=\frac{G'_k(x)}{H_k(x)}\quad\text{或}\quad H_k(x)\cdot\sigma_k(x)=G_k(x),
\end{equation}
於是所有 $H_k(x)$ 的乘積可被所有分母 $\chi_i(x)$ 的乘積整除。對於所有根系 $\alpha_1,
\ldots,\alpha_n$，若
% unit: noether.p11.eq.011b
\[
H_1(\alpha)\cdot H_2(\alpha)\cdots H_n(\alpha)\ne0
\]
並因而沒有分母 $\chi_i(\alpha)$ 消失，則可以在 (11) 中除以 $H_k(\alpha)$，得到
% unit: noether.p11.eq.012
\begin{equation}\tag{12}
\sigma_k(\alpha)=\frac{G_k(\alpha)}{H_k(\alpha)}=F_k(\varphi_1(\alpha),\ldots,\varphi_n(\alpha)),
\end{equation}
這裡因為分母不為零，$\varphi_i(\alpha)$ 取有限值；對於方程 $g_\Gamma$，它們成為 $\Omega^*$ 中的量。一旦 $\alpha_1,
\ldots,\alpha_n$ 被確定為使 $a_k=(-1)^k\sigma_k(\alpha)$，(12) 就表示 (10) 的解系。\footnote{由此，方程的求解被化歸為由獨立函數組成的系統的求解：
$\varphi_1(\alpha_1,
\ldots,\alpha_n)=\lambda_1;\ldots;\varphi_n(\alpha_1,
\ldots,\alpha_n)=\lambda_n$.}

% unit: noether.p11.par.017
因此，只有那些滿足
% unit: noether.p11.eq.012a
\[
H_1(\alpha)H_2(\alpha)\cdots H_n(\alpha)=0;
\]
的係數系 $a_k=(-1)^k\sigma_k(\alpha)$ 才可能成為奇異的；也就是說，正是那些使 (11) 變為表示 $0=0$ 的係數系。於是需要特別研究，才能決定在這樣定義的奇異 $a_1,
\ldots,a_n$ 中是否有對應於方程 $g_\Gamma$ 的情形，以及對於哪些數域有這種情形。\footnote{例如，如 F. Seidelmann 在前引處所示，這樣的方程 --- 對應於奇異值系 --- 只在三次和四次方程的交錯群中出現，並且只對包含三次單位根域的數域 $\Omega^*$ 出現。每一種情況下都可以給出一個簡單的補充表示。} 總結起來得到如下

% unit: noether.p11.thm.001
\smallskip
\noindent\textbf{定理。} \emph{如果屬於群 $\Gamma$ 的 Lagrange \emph{Gattungsbereich} 具有一個係數在 $\Omega$ 中的最小基，\footnote{由一個最小基的存在立刻推出任意多最小基的存在；它們可由原來的最小基通過可逆的有理變換得到。}那麼對於每個包含 $\Omega$ 的數域 $\Omega^*$，相對於 $\Omega^*$ 的具有給定群 $\Gamma$ 的方程全體，可以通過參數表示 (2) 有理地構造出來；至多例外的是那些相對於該參數表示為奇異的方程，它們可以另行給出。參數必須遍歷 $\Omega^*$ 中所有不導致群降低的值。如果最小基的係數域 $\Omega$ 等於所有有理數的域 $\mathfrak R$，則這種構造對每個數域都是可能的。}

% unit: noether.p11.par.018
此外，由 Hilbert 的不可約性定理，參數流形在排除 $\Omega^*$ 中那些導致群降低的值之後並不降低。因此還得到：\emph{由最小基的存在可知，相對於 $\Omega^*$ 的具有群 $\Gamma$ 的方程的流形，等於無 affect 方程的流形；流形並不因 affect 而降低。}

% unit: noether.p11.sec.002
\subsection*{\S 2. 特殊方程的應用。}

% unit: noether.p11.par.019
現在我們證明，對三次和四次方程中出現的所有群，最小基都存在；首先完全一般地說明，對 $n$ 次方程，問題可以化歸為含 $n-2$ 個不定元的問題。第一個化歸對應於去掉次高項的線性 Tschirnhaus 變換。我們取第一基函數 $\varphi_1(x)=\sigma_1(x)$，並進一步注意到，$\mathfrak L_\Gamma$ 可以由 $x_1,
\ldots,x_n$ 中允許 $\Gamma$ 的所有齊次形式有理地導出。設 $p(x_1,
\ldots,x_n)$ 是這樣的一個形式；那麼
% unit: noether.p11.eq.013
\[
q(x)=p\left(x_1-\frac{\sigma_1(x)}{n},\ldots,x_n-\frac{\sigma_1(x)}{n}\right)=p\left(x_i-\frac{\sigma_1(x)}{n}\right)
\]
也屬於 $\mathfrak L_\Gamma$，並且
% unit: noether.p11.eq.014
\[
p(x)\equiv p\left(x_i-\frac{\sigma_1(x)}{n}\right)\pmod{\sigma_1(x)},
\quad\text{或}\quad
p(x)=q(x)+\sigma_1(x)\cdot p_1(x),
\]
其中 $p_1(x)$ 屬於 $\mathfrak L_\Gamma$；它的維數也低於 $p(x)$。繼續這一過程可知，$\mathfrak L_\Gamma$ 中的所有齊次形式 --- 從而 $\mathfrak L_\Gamma$ 中所有有理函數本身 --- 都能由 $\varphi_1(x)=\sigma_1(x)$ 以及齊次形式
% unit: noether.p11.eq.015
\[
q(x)=p\left(x_i-\frac{\sigma_1(x)}{n}\right)=p(y_i).
\]
有理地表示。在這些組合 $y_i$ 之間，卻存在一個線性關係 $\sum y_i=0$；所以 $q(x)$ 只依賴於 $(n-1)$ 個自變量。

% unit: noether.p11.par.020
第二個化歸基於這樣一個事實：$q(x)$ 是其自變量的齊次形式。我們取第二基函數
% unit: noether.p11.eq.016
\[
\varphi_2(x)=\frac{\tau_3(x)}{\tau_2(x)},
\]
其中
% unit: noether.p11.eq.017
\[
\tau_3(x)=\sigma_3\left(x_i-\frac{\sigma_1(x)}{n}\right)=\sigma_3(y_i),
\qquad
\tau_2(x)=\sigma_2\left(x_i-\frac{\sigma_1(x)}{n}\right)=\sigma_2(y_i).
\]
於是 $\varphi_2(x)$ 是關於 $(n-1)$ 個自變量，即滿足 $\sum y_i=0$ 的 $y_i$，的一次齊次分式函數。藉助 $\varphi_2(x)$，所有 $q(x)$ 都能由同樣屬於 $\mathfrak L_\Gamma$ 的函數
% unit: noether.p11.eq.018
\[
r(x)=\frac{q(x)}{\varphi_2(x)^\lambda}
=\frac{q(x)\cdot\tau_2(x)^\lambda}{\tau_3(x)^\lambda}
=\frac{p(y_i)\cdot\sigma_2(y_i)^\lambda}{\sigma_3(y_i)^\lambda},
\]
有理地表示，其中 $\lambda$ 表示 $q(x)$ 的次數。於是 $r(x)$ 成為零次齊次，因而只依賴於 $(n-2)$ 個自變量，即在 $\sum y_i=0$ 條件下的比值 $y_1:y_2:\cdots:y_n$。因此 $\mathfrak L_\Gamma$ 可以由
% unit: noether.p11.eq.019
\[
\varphi_1(x)=\sigma_1(x),
\qquad
\varphi_2(x)=\frac{\tau_3(x)}{\tau_2(x)},
\]
以及由 $\mathfrak L_\Gamma$ 中所有依賴於這 $(n-2)$ 個自變量的函數 $r(x)$ 所成的域有理地表示，\footnote{對於沒有第二項的 $n$ 次方程
\[
f(x)=x^n+a_2x^{n-2}+a_3x^{n-3}+\cdots+a_n=0,
\]
這個第二化歸對應於在 $a_2\ne0$, $a_3\ne0$ 時總能進行的代換
\[
y=\frac{x\cdot a_2}{a_3},
\]
通過這個代換，$f(x)$ 在差一個因子的意義下變為
\[
g(y)=y^n+\frac{a_2^3}{a_3^2}y^{n-2}+\frac{a_2^3}{a_3^2}y^{n-3}
+\frac{a_4\cdot a_2^4}{a_3^4}y^{n-4}+\cdots+a_n\cdot\frac{a_2^n}{a_3^n}=0,
\]
也就是說，變為只含有 $n-2$ 個參數的方程：
\[
g(y)=y^n+c\,y^{n-2}+c\,y^{n-3}+c_4y^{n-4}+\cdots+c_n=0.
\]} 這就完成了所要求的化歸。

% unit: noether.p11.par.021
對於三次和四次方程，這特別給出了分別化歸為一個或兩個自變量的函數域。對於這些域，由 Lüroth 和 Castelnuovo 的定理，總存在最小基。\footnote{J. Lüroth: “Beweis eines Satzes über rationale Kurven”, Math. Ann. 9 (1875); G. Castelnuovo: “Sulla razionalità delle involuzioni piane”, Math. Ann. 44 (1893). 對於超過兩個不定元的域，最小基不一定存在，這由 F. Enriques 通過一個反例說明，Rend. Acc. Linc., vol. 21, 21 Jan. 1912.} 在一個不定元的情形中，最小基可以由有理運算導出，因此特別對於 $\Omega_\Gamma$ 只含有有理數係數。對於兩個不定元，按 Castelnuovo，基仍可能具有代數數域 $\Omega$ 中的係數；實際研究（參見 F. Seidelmann 前引處）表明，在這裡每個 $\Omega_\Gamma$ 也都得到一個有理係數基。幾乎無需計算就能從下列事實看出這一點：對於 Klein 四元群，可以選擇這樣一個有理係數基，即
% unit: noether.p11.eq.020
\[
\varphi_1(x)=\sigma_1(x);\quad
\varphi_2(x)=\frac{vw}{u};\quad
\varphi_3(x)=\frac{uw}{v};\quad
\varphi_4(x)=\frac{u\cdot v}{w},
\]
其中
% unit: noether.p11.eq.021
\[
u=x_1+x_2-x_3-x_4;\quad v=x_1-x_2+x_3-x_4;\quad w=x_1-x_2-x_3+x_4,
\]
並且交錯群轉化為三個基函數 $\varphi_2,\varphi_3,\varphi_4$ 的循環群，而每個八階群轉化為這幾個函數中某兩個的對稱群，由此產生向三次和二次方程的群的化歸。至於循環群，一個有理係數基已由 Abel 所知，雖然沒有以這種形式表述。\footnote{見 Weber, \emph{Algebra}（小版），\S 93，公式 (12)。} 因而已經證明：\emph{具有給定群的三次和四次方程的全體 --- 至多例外的是那些相對於參數表示為奇異的方程 --- 可以對任意有理性域通過參數表示有理地構造；它們的流形等於無 affect 方程的流形。}

% unit: noether.p11.par.022
實際研究表明（參見 F. Seidelmann 前引處），只有在有理性域包含三次單位根域時，對於交錯群才需要補充表示；而在所有其他情形下，參數表示直接給出全部方程。

% unit: noether.p11.par.023
還應指出，所有 Abel 群的最小基都是已知的。\footnote{E. Fischer: “Die Isomorphie der Invariantenkörper der endlichen Abelschen Gruppen linearer Transformationen”, Gött. Nachr. 1915, 以及 “Zur Theorie der endlichen Abelschen Gruppen”, Math. Ann. 77 (1915).} 不過在這裡，它是以由群的特徵標導出的圓分域為係數的基。因此，在這種情形下，對於構造具有給定群的方程不能得到新的結果。

% unit: noether.p11.closing.001
\medskip
\noindent Göttingen，1916 年 7 月。
%
\clearpage
%
% unit: noether.p12.title.001
\providecommand{\bdelta}{\bar{\delta}}
\begin{center}
\pdfbookmark[0]{12. 任意微分表達式的不變量}{paper-12}
{\Large\bfseries 12. 任意微分表達式的不變量}\par
\vspace{1.5em}
\emph{Nachrichten von der Gesellschaft der Wissenschaften zu Göttingen} 1918, pp. 37--44
\end{center}

\vspace{3em}
\begin{center}
1918 年 1 月 25 日會議中由 F. Klein 提交。
\end{center}

% unit: noether.p12.par.001
我把一個“微分表達式”理解為函數
% unit: noether.p12.eq.000
\[
f(x,dx)=f(x_1,\ldots,x_n;\,dx_1,\ldots,dx_n)
\]
它依賴於 \(n\) 個變量 \(x_1\ldots x_n\) 以及它們的微分
\(dx_1\ldots dx_n\)，並假定它在這兩個各含 \(n\) 個自變量的系統中解析，但不必是實的。現在我取所有變量的解析變換所成的群為基礎；與它相應的是所有微分的線性變換所成的群：
% unit: noether.p12.eq.001
\begin{equation}
x_i=x_i(y_1\ldots y_n);\qquad
dx_i=\sum\frac{\partial x_i}{\partial y_k}\,dy_k;\qquad
\delta x_i=\sum\frac{\partial x_i}{\partial y_k}\,\delta y_k,
\tag{1}
\end{equation}
藉此 \(f(x,dx)\) 應過渡為 \(g(y,dy)\)。所謂 \(f(x,dx)\) 的一個不變量，我理解為相對於這個群、並且相對於由此在高階微分 \(d^2x,d\delta x\ldots\) 以及 \(f(x,dx)\) 的導數上誘導出的群而言的不變量；也就是這樣一個解析函數 \(J\)：對於任意 \(f\)，並藉助 (1)，在變量和所有出現的微分中恆等地有
% unit: noether.p12.eq.002
\[
\begin{aligned}
&J\!\left(f,\frac{\partial f}{\partial dx}\cdots
\frac{\partial^{\varrho+\sigma}f}{\partial x^\varrho\partial dx^\sigma}
\cdots dx,\delta x,d^2x,\ldots\right)\\
&\qquad =
J\!\left(g,\frac{\partial g}{\partial dy}\cdots
\frac{\partial^{\varrho+\sigma}g}{\partial y^\varrho\partial dy^\sigma}
\cdots dy,\delta y,d^2y\ldots\right)\footnotemark
\end{aligned}
\]
\footnotetext{\(\displaystyle
\frac{\partial^{\rho+\sigma}f}{\partial x^\rho\partial dx^\sigma}\)
在全文中作為下式的縮寫使用：
\[
\frac{\partial^{\varrho+\sigma}f}
{\partial x_{i_1}\ldots\partial x_{t_\varrho}
 \partial dx_{k_1}\ldots\partial dx_{k_\sigma}},
\]
其中指標表示從 \(1\) 到 \(n\) 的任意數。}

% unit: noether.p12.par.002
完全同樣地，可以定義一個微分表達式同時系統的不變量。特別地，如果這樣一個不變量只含一階微分 \(dx,\delta x\)，且只含關於這些微分的導數，而不含關於變量的導數，那麼函數中變量的出現以及微分的線性變換都不再需要考慮；這些特殊不變量就是相對於所有以未定係數給出的 \(dx,\delta x\) 的線性變換群的不變量，並稱為該同時系統的射影不變量。

% unit: noether.p12.par.003
對於二次齊次微分形式，Christoffel（Crelle 70）和 Ricci（\emph{Math. Annalen} 54）已經建立了這樣一個不變構成的系統，使所有不變量都成為這個同時系統的射影不變量；從而關於不變量全體以及關於等價性的問題，都化歸為線性不變量論的問題。這就是我想稱作“約化定理”的東西。完整系統由可數無窮多個形式組成；但是對於任意有限階 \(\varrho\) 的不變量，即那些不含關於變量的高於 \(\varrho\) 階導數的不變量，它只由有限多個形式組成。

% unit: noether.p12.par.004
以下我證明“約化定理”對於任意微分表達式成立；這裡同樣涉及一個由可數無窮多個不變構成組成的完整系統，但對於每一個有限階只涉及有限多個構成。不過，Christoffel 和 Ricci 是把不變量的生成以及約化定理的證明建立在難以總覽的消去之上，而我則直接通過不變微分和變分過程來生成不變量；後一種過程可以簡單地看成關於其他參數的微分過程。這些不變微分過程的算法，繼 Lagrange（\emph{Mécanique analytique}，第二部，第四節）之後，由 Riemann（《全集》XXII）和 Lipschitz（Crelle 70, 72）為二次微分形式發展出來，不過他們並未推進到約化定理。證明約化定理的輔助工具由 Riemann 為二次形式引入的“正規座標”（《全集》XIII）提供；這些座標把從一點發出的、屬於相應變分問題的極值線變為直線，但它們只存在於齊次微分表達式
\(f(x,\varkappa\cdot dx)=\varPhi(\varkappa)\cdot f(x,dx)\footnote{容易看出 \(\varPhi(\varkappa)=\varkappa^c\)，其中 \(c\) 表示任意實量或復量。}\) 的情形。然而，非齊次情形可以化歸到這一情形。

% unit: noether.p12.heading.001
\begin{center}
\textbf{I. 不變量的生成。}
\end{center}

% unit: noether.p12.par.005
\(f(x,dx)\) 的一個通常不終止的射影不變量序列由極化式給出。由於微分變換的線性性，由 \(f(x,dx)=g(y,dy)\) 還可推出
\(f(x,dx+\lambda\delta x)=g(y,dy+\lambda\delta y)\)；於是按 \(\lambda\) 展開，便得到表徵不變性的關係：
% unit: noether.p12.eq.003
\begin{equation}
\begin{gathered}
f(dx)=g(dy),\\
f_\delta=\sum\frac{\partial f(dx)}{\partial dx}\,\delta x
       =\sum\frac{\partial g(dy)}{\partial dy}\,\delta y,\\
\vdots\\
f_{\delta^\sigma}=
\sum\frac{\partial^\sigma f(dx)}{\partial dx^\sigma}\,\delta x^\sigma
=
\sum\frac{\partial^\sigma g(dy)}{\partial dy^\sigma}\,\delta y^\sigma,\\
\vdots
\end{gathered}
\tag{2}
\end{equation}

% unit: noether.p12.par.006
每一個極化式，通過應用一個變分過程 \(\bdelta\)，產生一個非終止的一般不變量序列：
% unit: noether.p12.eq.004
\begin{equation}
\begin{gathered}
\bdelta^\varrho f(x,dx)=\bdelta^\varrho g(y,dy),\\
\vdots\\
\bdelta^\varrho f_{\delta^\sigma}
=\bdelta^\varrho\sum\frac{\partial^\sigma f(dx)}{\partial dx^\sigma}\,
\delta x^\sigma
=\bdelta^\varrho\sum\frac{\partial^\sigma g(dy)}{\partial dy^\sigma}\,
\delta y^\sigma,\\
\vdots\\[0.2em]
(\varrho=0,1,2\ldots\ldots).
\end{gathered}
\tag{3}
\end{equation}
藉助 (1)，關係 (2) 和 (3) 可以讀出 \(f(x,dx)\) 的導數的變換律；不過下文並不需要這一點。由不變量 (3)，現在可以通過線性組合得到“第 \(\varrho\) 次變分的正規形式”；在 \(\varrho=1\) 時它見於 Lagrange，在 \(\varrho=2\) 時見於 Riemann。它的特徵在於階數為 \(\varrho+1\) 的混合微分，即 \(d^\varrho\delta,d^{\varrho-1}\delta^2\ldots\)，被消去。得到：
% unit: noether.p12.eq.005
\begin{equation}
\begin{gathered}
\varOmega_1=\delta f(dx)-d f_\delta(dx),\\
\varOmega_2=\delta^2 f-\delta d f_\delta+\frac12 d^2 f_{\delta^2},\\
\vdots\\
\varOmega_\varrho=\delta^\varrho f-\delta^{\varrho-1}d f_\delta
+\frac12\delta^{\varrho-2}d^2 f_{\delta^2}
+\cdots+\frac{(-1)^\varrho}{\varrho!}\,d^\varrho f_{\delta^\varrho},\\
\vdots
\end{gathered}
\tag{4}
\end{equation}

% unit: noether.p12.par.007
由不變量 (4)，通過推廣 Riemann 的一個設法，現在可以得到只含一階微分的不變量；這樣的不變量稱為“基本函數”。確切地說，如果在 \(\varOmega_1\) 中把微分替換為微商，那麼 \(\varOmega_1=0\)（關於 \(\delta x\) 恆等）就正好給出屬於
% unit: noether.p12.eq.006
\[
f\!\left(\frac{dx}{dt}\right).
\]
的變分問題的 Lagrange 方程。現在我加上這樣一個不變條件：方向
\((d+\lambda\delta)\) 滿足這些 Lagrange 方程：
% unit: noether.p12.eq.007
\begin{equation}
\varOmega_1(d+\lambda\delta)=
\bdelta f(dx+\lambda\delta x)
-(d+\lambda\delta)f_{\bdelta}(dx+\lambda\delta x)=0
\quad(\text{關於 }\delta x\text{ 恆等}),
\tag{5}
\end{equation}
由此，通過代入 \(\bdelta=d+\lambda\delta\)，對於齊次 \(f(dx)\) 還得到
% unit: noether.p12.eq.008
\begin{equation}
(d+\lambda\delta)f(dx+\lambda\delta x)=0,
\tag{6}
\end{equation}
但一階齊次的情形除外。在後一種情形中，(6) 必須作為新的條件添加，因為此時方程組 (5) 只表示 \((n-1)\) 個獨立條件。如果在 (5)，相應地在 (5) 和 (6) 中，把 \(\lambda\) 的零次、一次和二次冪的係數置零，就得到三個不變方程組
\(\varphi=0\)、\(\psi=0\)、\(\chi=0\)（關於 \(\delta x\) 恆等），藉此 \(d^2x,ddx,\delta^2x\) 可用一階微分表示。進一步的不變方程組
% unit: noether.p12.eq.009
\begin{equation}
\begin{gathered}
d^\sigma\varphi=0,\qquad d^\sigma\psi=0,\qquad
d^\sigma\chi=0,\qquad d^{\sigma-1}\delta\chi=0,\\
d^{\sigma-2}\delta^2\chi=0\ldots\delta^\sigma\chi=0
\qquad(\sigma=0,1,2\ldots)
\end{gathered}
\tag{7}
\end{equation}
於是恰好足以把所有高階微分表示為一階微分。\footnote{只有線性形式
\(\sum A_i\,dx_i\) 需要特殊處理，因為這裡方程組 (5) 不含二階微分。} 因此，函數 (4) 與不變方程組 (7) 結合，導出只含一階微分的基本函數 \([\varOmega_\varrho]\)；而且 \([\varOmega_1]\) 恆等為零。

% unit: noether.p12.par.008
同樣，從一個基本函數 \(h\) 的微分 \(\delta h(dx,\delta x)\) 出發，可藉助 (7) 再得到一個基本函數，稱為 \(h\) 的“協變導數” \([h^{(1)}]\)；一般地，令 \([h^{(\sigma)}]\) 表示 \([h^{(\sigma-1)}]\) 的協變導數。這兩個過程於是導出一個雙重無窮的基本函數序列：
% unit: noether.p12.eq.010
\begin{equation}
\begin{array}{cccccc}
[\varOmega_2], & [\varOmega_2^{(1)}], & [\varOmega_2^{(2)}] & \ldots,
& [\varOmega_2^{(\nu)}] & \ldots\\
\vdots &&&&&\\
{}[\varOmega_\varrho], & [\varOmega_\varrho^{(1)}], & \ldots & \ldots,
& [\varOmega_\varrho^{(\nu)}] & \ldots\\
\vdots & \vdots &&& \vdots&
\end{array}
\tag{8}
\end{equation}

% unit: noether.p12.par.009
還可看出，把高階微分用一階微分表示的那些方程的左端，與一階微分同變，即經歷同樣的線性變換。因此，如果不把高階微分本身，而把這些左端 \(p,q,r\ldots\) 作為不變量的自變量引入，那麼撇開導數不論，它只依賴於彼此同變的量。上面所指出的函數 (8) 的形成，正是這些新量的引入；每一個不變量由此過渡為一項不變量之和，而從這和式中選取恰好含有 \(dx,\delta x\)、但不含 \(p,q,r\ldots\) 的那一項。

% unit: noether.p12.par.010
下面將證明，在齊次性假設
% unit: noether.p12.eq.011
\[
f(\varkappa\cdot dx)=\varPhi(\varkappa)\cdot f(dx)
\]
下，基本函數 (8) 窮盡所尋求的完整系統。

% unit: noether.p12.heading.002
\begin{center}
\textbf{II. 正規座標、等價性與約化定理。}
\end{center}

% unit: noether.p12.par.011
我通過積分屬於 \(f\!\left(\frac{dx}{dt}\right)\) 的變分問題而達到正規座標：
% unit: noether.p12.eq.012
\begin{equation}
\begin{gathered}
\varOmega_1=\delta f\!\left(\frac{dx}{dt}\right)
-\frac{d}{dt}f_\delta\!\left(\frac{dx}{dt}\right)=0
\quad(\text{關於 }\delta x\text{ 恆等}),\\
\frac{d}{dt}f\!\left(\frac{dx}{dt}\right)=0
\quad(\text{作為結果或附加條件}).
\end{gathered}
\tag{9}
\end{equation}
由於假設 \(f(dx)\) 齊次，這個方程組在參數代換 \(t=\varkappa\cdot\tau\) 下變為自身。因此，如果藉助 (9) 把
\(\frac{d^\varrho x_i}{dt^\varrho}\) 表示為第一導數的函數
\(\varphi_\varrho^{(i)}\!\left(\frac{dx}{dt}\right)\)，那麼 \(\varphi_\varrho^{(i)}\) 是 \(\varrho\) 階齊次的：
% unit: noether.p12.eq.013
\begin{equation}
\varphi_\varrho^{(i)}\!\left(\varkappa\cdot\frac{dx}{dt}\right)
=\varkappa\varphi_\varrho^{(i)}\!\left(\frac{dx}{dt}\right).
\tag{10}
\end{equation}
如果在積分 (9) 時，規定 \(t=t_0\) 時的初值
% unit: noether.p12.eq.014
\[
x=(x)_0,\qquad \frac{dx}{dt}=\left(\frac{dx}{dt}\right)_0
\]
並置
% unit: noether.p12.eq.015
\begin{equation}
u_i=(t-t_0)\left(\frac{dx_i}{dt}\right)_0,
\tag{11}
\end{equation}
這裡由於 \(f\!\left(\frac{dx}{dt}\right)=\mathrm{const.}\)，參數
\((t-t_0)\) 與“極值線的長度”成比例，那麼按 \((t-t_0)\) 的冪展開，依 (10) 得
% unit: noether.p12.eq.016
\begin{equation}
x_i-(x_i)_0=u_i+\frac12\varphi_2^{(i)}(u)+\cdots+
\frac{1}{\varrho!}\varphi_\varrho^{(i)}(u)+\cdots.
\tag{12}
\end{equation}
這個展開可看成變量變換 \(x_i=x_i(u)\)，其中 \(u\) 正是 Riemann 正規座標。於是由 (11) 和 (12)，這些座標把經過 \((x_0)\) 的極值線變成直線。此外，對任意變量變換 \(x=x(y)\)，由 (11) 對應有相關正規座標的一個線性變換 \(u=u(v)\)，其係數只依賴於任意給定的值系 \((x)_0\)；因此微分 \(du,\delta u\) 服從與 \(u\) 本身相同的變換。

% unit: noether.p12.par.012
藉助 (12)，令 \(f(x,dx)\) 過渡為 \(F(u,du)\)，或者按 \(u\) 的冪展開為
% unit: noether.p12.eq.017
\begin{equation}
F(u,du)=F_0(u,du)+F_1(u,du)+\cdots F_\varrho(u,du)+\cdots,
\tag{13}
\end{equation}
其中 \(F_\varrho(u,du)\) 表示關於 \(u\) 的維數為 \(\varrho\) 的齊次形式。由於變換 \(u=u(v)\) 的線性性，恆等式 \(F(u,du)=G(v,dv)\) 也把各個齊次分量帶到相應的分量：
% unit: noether.p12.eq.018
\begin{equation}
F_\varrho(u,du)=G_\varrho(v,dv).
\tag{14}
\end{equation}
因此，如果對 \((x)_0\) 取任意固定的常值系統，則 (13) 和 (14) 表明，\(f(dx)\) 與 \(g(dy)\) 的等價性等價於相應函數系統 \(F_\varrho(u,du)\) 在一個線性變換下的等價性。此外，\(F_\varrho(u,du)\)，或相應的 \(F_\varrho(\delta u,du)\)，表示在點 \(x=(x)_0\) 處取的 \(f(dx)\) 的不變量；事實上它們構成該點 \(x=(x)_0\) 的一個完整基本函數系統。因為
% unit: noether.p12.eq.019
\[
F_\varrho(\delta u,du)=\frac1{\varrho!}\sum
\frac{\partial F_\varrho}{\partial \delta u^\varrho}\,\delta u^\varrho;
\qquad
\frac{\partial F_\varrho}{\partial \delta u^\varrho}
=\left(\frac{\partial F(du)}{\partial u^\varrho}\right)_0,
\]
而這個後一導數可以藉助 (12)，以完全同樣的方式，通過在 \(x=(x)_0\) 處取的 \(f(dx)\) 的導數表示出來；正如 \(G(dv)\) 的相應導數可通過 \(g(dy)\) 的導數表示出來一樣。由於
% unit: noether.p12.eq.020
\[
\left(\frac{\partial^{\varrho+\sigma}F(du)}
{\partial u^\varrho\partial du^\sigma}\right)_0
=
\frac{\partial^{\varrho+\sigma}F_\varrho(\delta u,du)}
{\partial\delta u^\varrho\partial du^\sigma},
\]
\(f(dx)=F(du)\) 的每一個在點 \(x=(x)_0\) 處取的不變量，一旦只把高階微分替換為與 \(dx,\delta x\) 同變的 \(p,q,r\ldots\)，也就成為 \(F_\varrho(\delta u,du)\) 的一個射影不變量。

% unit: noether.p12.par.013
但現在 \((x)_0\) 可以看成一個參數；每一個在點 \(x=(x)_0\) 處的不變量，只要把 \((x)_0\) 重新替換為 \(x\)，便對應於 \(f(dx)\) 的一個不變量。於是，如果 \(F_\varrho(\delta u,du)\) 變成在 \(x=(x)_0\) 處取的不變量 \(\varPsi_\varrho(\delta x,dx)\)——因為當 \(x=(x)_0\) 時，\(du,\delta u\) 過渡為 \(dx,\delta x\)——那麼 \(\varPsi_\varrho\) 本身構成一個完整系統的基本函數。剩下的只是把這個基本函數系統用 \(f(dx)\) 的顯式不變量表達出來，即用函數 (8) 表達出來。其可能性由它們通過微分過程的生成來說明。對於最高階項，這些過程變成簡單的極化過程；再用一個類似於 Clebsch--Gordan 級數展開的變換，並由於被忽略的項而配合歸納，便證明了這個斷言。如果處理的是一個同時系統的不變量，那麼還必須把其餘微分表達式的協變導數加入基本函數 (8) 中；這一點由把屬於 \(f(dx)\) 的正規座標引入這些表達式可見。

% unit: noether.p12.par.014
因此，已經被化歸為 (14) 的 \(f(dx)\) 與 \(g(dy)\) 的等價性，也等價於 \(\varPsi_\varrho\) 的等價性，或同樣地等價於函數 (8) 在微分線性變換下的等價性，而不需要任何特殊的可積性條件；這由可化歸為 (14) 得出。由此，約化定理在其全部部分中得證。特別地，函數序列
\([\varOmega_2],[\varOmega_3]\ldots[\varOmega_\varrho]\ldots\) 恆等消失，是 \(f(dx)\) 可化為一個常係數表達式的充要條件。

% unit: noether.p12.par.015
最後，如果不以所有解析變換的群為基礎，而取其一個子群為基礎，那麼相應的 \(u\) 的線性變換群也過渡為射影群的一個子群；\(f(dx)\) 的不變量再次成為函數 (8) 關於線性變換的不變量，不過現在是關於這個子群的不變量。因此，通過齊次化，非齊次函數 \(f(dx)\) 的情形可以化歸為仿射群；這裡的完整系統可由對多一個變量形成的函數 (8) 導出。

% unit: noether.p12.par.016
由約化定理可以得到二次形式的特殊情形，更一般地得到 \(p\) 維形式的特殊情形，其原因是這裡函數系統在 \([\varOmega_2]\) 處終止，更一般地在 \([\varOmega_p]\) 及其協變導數處終止，因為所有後續基本函數都成為這些函數的射影不變量。這解釋了 Riemann 曲率形式 \([\varOmega_2]\) 對二次形式所具有的突出地位。二次形式，分別 \(p\) 維形式，的等價性問題，正好化為 \([\varOmega_2]\)，分別 \([\varOmega_2]\ldots[\varOmega_p]\)，以及它們的協變導數在一個線性變換下的等價性；而 \(]\varOmega_2]\)，分別 \([\varOmega_2]\ldots[\varOmega_p]\)，的恆等消失，是變換為常係數形式的充要條件；對於二次形式，這正說明了最著名的結果之一。

% unit: noether.p12.par.017
更詳細的敘述將很快發表於 \emph{Mathematische Annalen}。
%
\clearpage
%
% unit: noether.p13.title.001
\begin{center}
\pdfbookmark[0]{13. 不變變分問題}{paper-13}
{\Large\bfseries 13. 不變變分問題}\par
\vspace{1.5em}
\emph{Nachrichten von der Gesellschaft der Wissenschaften zu Göttingen} 1918, pp. 235--257
\end{center}

\vspace{3em}
\begin{center}
由 F. Klein 在 1918 年 7 月 26 日的會議上提交。\footnote{手稿的最終稿直到九月底才提交。}
\end{center}

% unit: noether.p13.par.001
這裡討論的是容許一個連續群（按 Lie 的意義）的變分問題；由此對相應微分方程所產生的結果，在 \S{}1 所表述並在以下各節證明的定理中獲得其最一般的形式。對於這些由變分問題產生的微分方程，可以作出比對容許一個群的任意微分方程更精確得多的斷言；後者是 Lie 研究的對象。因此，以下內容建立在形式變分法的方法同 Lie 群論的方法相結合之上。對於特殊的群和變分問題，這種方法組合並不是新的；對於特殊有限群，我提到 Hamel 和 Herglotz；對於特殊無限群，我提到 Lorentz 及其學生（例如 Fokker）、Weyl 和 Klein。\footnote{Hamel: \emph{Math. Ann.} vol. 59 和 \emph{Zeitschrift f. Math. u. Phys.} vol. 50。Herglotz: \emph{Ann. d. Phys.} (4) vol. 36，特別是 \S{}9, p. 511。Fokker, Verslag d. Amsterdamer Akad., 1917 年 1 月 27 日。關於進一步文獻，可比較 Klein 的第二篇札記: \emph{Göttinger Nachrichten}, 1918 年 7 月 19 日。在 Kneser 剛發表的一篇論文（\emph{Math. Zeitschrift} vol. 2）中，問題是用類似方法構造不變量。} 特別地，Klein 的第二篇札記和本文的展開彼此影響；對此我可以參照 Klein 札記的結尾說明。

% unit: noether.p13.sec.001
\section*{\S{}1. 預備知識和定理的表述}

% unit: noether.p13.par.002
以下出現的所有函數都假定為解析的，或者至少在所考察的區域中是連續的、可連續微分有限次的，並且是單值的。

% unit: noether.p13.par.003
所謂“變換群”，通常是指這樣一個變換系統：對每一個變換，都有一個包含在該系統中的逆變換；並且該系統中任意兩個變換的複合仍屬於該系統。若一個群的各變換包含在一個最一般的變換中，而該變換解析地依賴於 \(\rho\) 個本質參數 \(\eps\)（也就是說，這 \(\rho\) 個參數不能表示為較少參數的 \(\rho\) 個函數），則稱該群為有限連續 \(\G_\varrho\)。相應地，無窮連續 \(\G_{\infty\varrho}\) 是指這樣一個群：其最一般的變換依賴於 \(\varrho\) 個本質任意函數 \(p(x)\) 及其導數，這種依賴是解析的，或者至少是連續的並具有有限個連續導數。介於兩者之間的情形，是依賴於無窮多個參數但不依賴任意函數的群。最後，混合群是指同時依賴任意函數和參數的群。\footnote{Lie 在 \emph{Grundlagen für die Theorie der unendlichen kontinuierlichen Transformationsgruppen}（皇家薩克森科學院報告，1891）中，把無窮連續群定義為這樣一個變換群：其變換由某個偏微分方程組的最一般解給出，只要這些解並不只依賴有限個參數。這給出了上面指出的、不同於有限群的類型之一；反過來說，具有無窮多個參數的極限情形卻不一定必須滿足一個微分方程組。}

% unit: noether.p13.par.004
設 \(x_1,\ldots,x_n\) 為獨立變量，設 \(u_1(x),\ldots,u_\mu(x)\) 為依賴於它們的函數。如果把 \(x\) 和 \(u\) 置於一個群的變換之下，那麼由於變換假定為可逆，在變換後的量中必須重新正好有 \(n\) 個獨立量，記為 \(y_1,\ldots,y_n\)；依賴於這些量的其餘量記為 \(v_1(y),\ldots,v_\mu(y)\)。這些變換還可以包含 \(u\) 關於 \(x\) 的導數，即 \(\frac{\p u}{\p x},\ \frac{\p^2u}{\p x^2},\ \cdots\)。\footnote{我儘可能省去指標，在求和中也如此；例如 \(\frac{\p^2u}{\p x^2}\) 表示 \(\p^2u_\nu/\p x_i\p x_k\)，等等。} 如果存在關係
% unit: noether.p13.eq.001
\[
P\!\left(x,u,\frac{\p u}{\p x},\frac{\p^2u}{\p x^2},\cdots\right)
 =
P\!\left(y,v,\frac{\p v}{\p y},\frac{\p^2v}{\p y^2},\cdots\right),
\]
則稱函數 \(P\) 為該群的不變量。特別地，如果存在關係
% unit: noether.p13.eq.002
\begin{equation}
\begin{aligned}
I&=\int\!\cdots\!\int f\!\left(x,u,\frac{\p u}{\p x},\frac{\p^2u}{\p x^2},\cdots\right)d x\\
 &=\int\!\cdots\!\int f\!\left(y,v,\frac{\p v}{\p y},\frac{\p^2v}{\p y^2},\cdots\right)d y,
\end{aligned}
\tag{1}
\end{equation}
則積分 \(I\) 因而是該群的不變量，其中積分是在任意實 \(x\)-區域及其對應的 \(y\)-區域上進行。\footnote{我簡寫 \(d x,d y\) 以表示 \(d x_1\ldots d x_n,d y_1\ldots d y_n\)。變換中出現的一切自變量 \(x,u,\eps,p(x)\) 都取為實的，而係數可以是復的。但是，由於最終結果是關於 \(x,u\)、參數和任意函數的恆等式，只要所有出現的函數都假定為解析，它們也對復值成立。此外，相當一部分結果可以不借助積分而建立，所以即便在證明上，對實數的限制也不是必要的。相反，\S{}2 末尾和 \S{}5 開頭的考察似乎不能不用積分來進行。}

% unit: noether.p13.par.005
另一方面，對於一個任意積分 \(I\)，不必假定它是不變的，我形成第一變分 \(\delta I\)，並按變分法規則用分部積分對它進行變換。眾所周知，只要 \(\delta u\) 連同它在邊界上出現的所有導數都在邊界上取零、而在其他地方任意，就得到
% unit: noether.p13.eq.003
\begin{equation}
\delta I=\int\!\cdots\!\int \delta f\,d x
=\int\!\cdots\!\int\left(\sum\psi_i\!\left(x,u,\frac{\p u}{\p x},\cdots\right)\delta u_i\right)d x,
\tag{2}
\end{equation}
其中 \(\psi\) 表示 Lagrange 表達式，即相應變分問題 \(\delta I=0\) 的 Lagrange 方程左端。這個積分關係對應於一個不含積分的、關於 \(\delta u\) 及其導數的恆等式；當把邊界項也寫出來時就得到它。分部積分表明，這些邊界項是關於散度的積分，也就是關於表達式
% unit: noether.p13.eq.004
\[
\Div A=\frac{\p A_1}{\p x_1}+\cdots+\frac{\p A_n}{\p x_n},
\]
的積分，其中 \(A\) 關於 \(\delta u\) 及其導數是線性的。於是得到
% unit: noether.p13.eq.005
\begin{equation}
\sum_i\psi_i\,\delta u_i=\delta f+\Div A.
\tag{3}
\end{equation}
特別地，如果 \(f\) 只含有 \(u\) 的一階導數，那麼在單積分的情形下，恆等式 (3) 與 Herr 所謂“Lagrange 中心方程”相同：
% unit: noether.p13.eq.006
\begin{equation}
\sum_i \psi_i\,\delta u_i
=\delta f-\frac{d}{d x}\left(\sum_i\frac{\p f}{\p u_i'}\,\delta u_i\right),
\qquad \left(u_i'=\frac{d u_i}{d x}\right),
\tag{4}
\end{equation}
而對於 \(n\) 重積分，(3) 變為
% unit: noether.p13.eq.007
\begin{equation}
\sum_i\psi_i\,\delta u_i
=
\delta f
-\frac{\p}{\p x_1}\left(\sum_i\frac{\p f}{\p \left(\frac{\p u_i}{\p x_1}\right)}\,\delta u_i\right)
-\cdots
-\frac{\p}{\p x_n}\left(\sum_i\frac{\p f}{\p \left(\frac{\p u_i}{\p x_n}\right)}\,\delta u_i\right).
\tag{5}
\end{equation}
對於單積分以及 \(u\) 的 \(\varkappa\) 階導數，(3) 給出為
% unit: noether.p13.eq.008
\begin{equation}
\begin{aligned}
\sum_i\psi_i\,\delta u_i
={}&\delta f
-\frac{d}{d x}\Bigg\{\sum_i\left(
\binom{1}{1}\frac{\p f}{\p u_i^{(1)}}\delta u_i
+\binom{2}{1}\frac{\p f}{\p u_i^{(2)}}\delta u_i^{(1)}
+\cdots+
\binom{\varkappa}{1}\frac{\p f}{\p u_i^{(\varkappa)}}\delta u_i^{(\varkappa-1)}
\right]\Bigg\}\\
&+\frac{d^2}{d x^2}\Bigg\{\sum_i\left(
\binom{2}{2}\frac{\p f}{\p u_i^{(2)}}\delta u_i
+\binom{3}{2}\frac{\p f}{\p u_i^{(3)}}\delta u_i^{(1)}
+\cdots+
\binom{\varkappa}{2}\frac{\p f}{\p u_i^{(\varkappa)}}\delta u_i^{(\varkappa-2)}
\right]\Bigg\}\\
&+\cdots
+(-1)^\varkappa\frac{d^\varkappa}{d x^\varkappa}
\left\{\sum\binom{\varkappa}{\varkappa}
\frac{\p f}{\p u_i^{(\varkappa)}}\delta u_i\right\}.
\end{aligned}
\tag{6}
\end{equation}
對於 \(n\) 重積分也有相應的恆等式；特別地，\(A\) 含有 \(\delta u\) 直到 \((\varkappa-1)\) 階導數。(4)、(5)、(6) 確實定義 Lagrange 表達式 \(\psi_i\)，這是因為在右端的組合中，\(\delta u\) 的所有較高導數都被消去；同時關係 (2) 被滿足，而分部積分唯一地導向這個關係。

% unit: noether.p13.par.006
以下涉及兩個定理：

% unit: noether.p13.thm.001
I. 如果積分 \(I\) 關於一個 \(\G_\varrho\) 不變，那麼 Lagrange 表達式的 \(\varrho\) 個線性無關組合成為散度；反過來，由此推出 \(I\) 關於一個 \(\G_\varrho\) 的不變性。這個定理在無窮多個參數的極限情形下也成立。

% unit: noether.p13.thm.002
II. 如果積分 \(I\) 關於一個 \(\G_{\infty\varrho}\) 不變，而其中任意函數出現到 \(\sigma\) 階導數，那麼在 Lagrange 表達式及其直到 \(\sigma\) 階的導數之間存在 \(\rho\) 個恆等關係；這裡反過來也是成立的。\footnote{關於某些平凡例外情形，比較 \S{}2 的第二個注。}

% unit: noether.p13.par.007
對於混合群，這兩個定理的斷言都適用；因此既出現依賴關係，也出現與這些依賴關係無關的散度關係。

% unit: noether.p13.par.008
從這些恆等式過渡到相應的變分問題，即令 \(\psi=0\)，\footnote{稍微更一般地，也可以令 \(\psi_i=T_i\)；比較 \S{}3 的第一個注。} 定理 I 在一維情形中說——此時散度成為全微分——存在 \(\varrho\) 個一階積分，儘管它們之間可以出現非線性依賴；\footnote{比較 \S{}3 的結尾。} 在多維情形中得到最近常被稱為“守恆律”的散度方程。定理 II 說，\(\varrho\) 個 Lagrange 方程是其餘方程的結果。

% unit: noether.p13.par.009
定理 II 的最簡單例子——沒有反命題——是 Weierstrass 參數表示；這裡在一階齊次性的條件下，眾所周知，如果把獨立變量 \(x\) 代以 \(x\) 的任意函數，而保持 \(u\) 不變 \((y=p(x);\ v_i(y)=u_i(x))\)，則積分是不變的。因此出現一個任意函數，但沒有導數；與之相應的是 Lagrange 表達式本身之間的著名線性關係
% unit: noether.p13.eq.009
\[
\sum_i \psi_i\,\frac{\dd u_i}{\dd x}=0.
\]
物理學家的“廣義相對論”提供了另一個例子。這裡涉及所有 \(x\) 的變換所成的群，\(y_i=p_i(x)\)，而 \(u\)（記作 \(g_{\mu\nu}\) 和 \(q\)）則服從由此對二次微分形式和一次微分形式的係數所誘導的變換；這些變換含有任意函數 \(p(x)\) 的一階導數。與此相應的是 Lagrange 表達式及其一階導數之間已知的 \(n\) 個依賴關係。\footnote{例如比較 Klein 的表述。}

% unit: noether.p13.par.010
如果特別地把群專門化：不允許 \(u(x)\) 的導數出現在變換中，並且還使變換後的獨立量只依賴於 \(x\) 而不依賴於 \(u\)，那麼，如 \S{}5 將要證明的，\(I\) 的不變性蘊含 \(\sum\psi_i\delta u_i\) 的相對不變性，\footnote{即 \(\sum \psi_i\delta u_i\) 在變換下獲得一個因子。} 也蘊含定理 I 中出現的那些散度的相對不變性，只要對參數施以適當的變換。還進一步得到，上述一階積分也容許該群。對於定理 II，同樣得到通過任意函數組合起來的依賴關係左端的相對不變性；並由此還得到一個散度恆等為零且容許該群的函數——在物理學家的相對論中，正是這個函數在依賴關係和能量定理之間起中介作用。\footnote{比較 Klein 的第二篇札記。} 最後，以群論形式來看，定理 II 給出了 Hilbert 關於“廣義相對論”中真正能量定理失效的一個相關斷言的證明。連同這些補充說明，定理 I 包含了力學等等中已知的所有關於一階積分的定理，而定理 II 可以說是“廣義相對論”在群論上的最大可能推廣。

% unit: noether.p13.sec.002
\section*{\S{}2. 散度關係和依賴關係}

% unit: noether.p13.par.011
設 \(\G\) 為一個連續群，可以是有限的也可以是無限的。總可以安排使恆等變換對應於參數 \(\eps\) 的零值，分別對應於任意函數 \(p(x)\) 的零值。\footnote{例如比較 Lie, \emph{Grundlagen}, p. 331。若有任意函數參與，則參數的特殊值 \(a^\sigma\) 應由固定函數 \(p^\sigma,\ \frac{\p p^\sigma}{\p x},\ \cdots\) 取代；相應地，值 \(a^\sigma+\eps\) 應由 \(p^\sigma+p(x),\ \frac{\p p^\sigma}{\p x}+\frac{\p p}{\p x}\) 等取代。} 因而最一般的變換具有形式
% unit: noether.p13.eq.010
\[
y_i=A_i\!\left(x,u,\frac{\p u}{\p x},\cdots\right)=x_i+\varDelta x_i+\cdots,\qquad
v_i(y)=B_i\!\left(x,u,\frac{\p u}{\p x},\cdots\right)=u_i+\varDelta u_i+\cdots,
\]
其中 \(\varDelta x_i,\varDelta u_i\) 表示 \(\eps\) 中，分別 \(p(x)\) 及其導數中的最低維項；它們將被假定為關於這些量線性的。後面將說明，這並不限制一般性。

% unit: noether.p13.par.012
現在設積分 \(I\) 關於 \(\G\) 不變，於是關係 (1) 成立。特別地，\(I\) 也關於包含在 \(\G\) 中的無窮小變換不變：
% unit: noether.p13.eq.011
\[
y_i=x_i+\varDelta x_i;\qquad v_i(y)=u_i+\varDelta u_i.
\]
對於它，關係 (1) 變為
% unit: noether.p13.eq.012
\begin{equation}
\begin{aligned}
0=\varDelta I
&=\int\!\cdots\!\int f\!\left(y,v(y),\frac{\p v}{\p y},\cdots\right)d y\\
&\qquad-\int\!\cdots\!\int f\!\left(x,u(x),\frac{\p u}{\p x},\cdots\right)d x.
\end{aligned}
\tag{7}
\end{equation}
第一個積分應擴展到同 \(x\)-區域對應的 \(x+\varDelta x\)-區域上。但是這種積分也可以藉助於對無窮小 \(\varDelta x\) 有效的變換，化為 \(x\)-區域上的積分：
% unit: noether.p13.eq.013
\begin{equation}
\begin{aligned}
\int\!\cdots\!\int f\!\left(y,v(y),\frac{\p v}{\p y},\cdots\right)d y
&=\int\!\cdots\!\int f\!\left(x,v(x),\frac{\p v}{\p x},\cdots\right)d x\\
&\quad+\int\!\cdots\!\int \Div(f\cdot \varDelta x)\cdot d x.
\end{aligned}
\tag{8}
\end{equation}
因此，如果以變分
% unit: noether.p13.eq.014
\begin{equation}
\bd u_i=v_i(x)-u_i(x)=\varDelta u_i-\sum\frac{\p u_i}{\p x_\lambda}\varDelta x_\lambda,
\tag{9}
\end{equation}
代替無窮小變換 \(\varDelta u\)，那麼 (7) 與 (8) 變為
% unit: noether.p13.eq.015
\begin{equation}
0=\int\!\cdots\!\int\{\bd f+\Div(f\cdot\varDelta x)\}d x.
\tag{10}
\end{equation}
由於關係 (10) 在任意區域上積分時都成立，被積函數必須恆等為零。於是 \(I\) 的不變性的 Lie 微分方程轉化為關係
% unit: noether.p13.eq.016
\begin{equation}
\bd f+\Div(f\cdot\varDelta x)=0.
\tag{11}
\end{equation}
如果在這裡按照 (3) 用 Lagrange 表達式表示 \(\bd f\)，就得到
% unit: noether.p13.eq.017
\begin{equation}
\sum\psi_i\bd u_i=\Div B\qquad (B=A-f\cdot\varDelta x),
\tag{12}
\end{equation}
因此，對於每一個不變積分 \(I\)，這個關係都表示為所有出現的自變量中的恆等式；這就是 \(I\) 的 Lie 微分方程所尋求的形式。\footnote{在平凡情形中，(12) 變成 \(0=0\)；這種情形只有在 \(\varDelta x,\varDelta u\) 也依賴於 \(u\) 的導數時才可能發生，即 \(\Div(f\cdot\varDelta x)=0,\ \bd u=0\)。因此，這些無窮小變換必須總是從群中分離出來，而在定理的表述中只計算剩下的參數或任意函數的個數。剩下的無窮小變換是否仍然總是構成一個群，這一點必須暫且不論。}

% unit: noether.p13.par.013
先假定 \(\G\) 是一個有限連續群 \(\G_\varrho\)。由於按假設 \(\varDelta u\) 和 \(\varDelta x\) 關於參數 \(\eps_1,\ldots,\eps_\varrho\) 是線性的，根據 (9)，\(\bd u\) 及其導數也是如此；因此 \(A\) 和 \(B\) 關於 \(\eps\) 是線性的。所以我令
% unit: noether.p13.eq.018
\[
B=B^{(1)}\eps_1+\cdots+B^{(\varrho)}\eps_\varrho;\qquad
\bd u=\bd u^{(1)}\eps_1+\cdots+\bd u^{(\varrho)}\eps_\varrho,
\]
其中 \(\bd u^{(1)},\ldots\) 是 \(x,u,\frac{\p u}{\p x},\cdots\) 的函數。由 (12) 得到所求的散度關係：
% unit: noether.p13.eq.019
\begin{equation}
\sum\psi_i\bd u_i^{(1)}=\Div B^{(1)};\ldots\qquad
\sum\psi_i\bd u_i^{(\varrho)}=\Div B^{(\varrho)}.
\tag{13}
\end{equation}
於是 Lagrange 表達式的 \(\varrho\) 個線性無關組合成為散度。線性無關性由如下事實推出：按照 (9)，\(\bd u=0,\ \varDelta x=0\) 將推出 \(\varDelta u=0,\ \varDelta x=0\)，因而推出無窮小變換之間的依賴關係。但按假設，對任何參數值都沒有這樣的依賴；否則，由無窮小變換重新積分得到的 \(\G_\varrho\) 將依賴於少於 \(\varrho\) 個本質參數。另一種可能性 \(\bd u=0,\ \Div(f\varDelta x)=0\) 已被排除。這些結論在無窮多個參數的極限情形中仍然有效。

% unit: noether.p13.par.014
現在設 \(\G\) 為一個無窮連續群 \(\G_{\infty\varrho}\)。於是 \(\bd u\) 及其導數，因而 \(B\)，又關於任意函數 \(p(x)\) 及其導數線性。\footnote{假定 \(p\) 獨立於 \(u,\frac{\p u}{\p x},\cdots\) 並不是限制，這一點由反命題說明。} 即使暫時不利用 (12)，把 \(\bd u\) 的值代入也給出
% unit: noether.p13.eq.020
\[
\sum\psi_i\bd u_i
=\sum_{\lambda,i}\psi_i\left\{
a_i^{(\lambda)}(x,u,\ldots)p^{(\lambda)}(x)
+b_i^{(\lambda)}(x,u,\ldots)\frac{\p p^{(\lambda)}}{\p x}
+\cdots+c_i^{(\lambda)}(x,u,\ldots)
\frac{\p^\sigma p^{(\lambda)}}{\p x^\sigma}\right\}.
\]
現在，類似於分部積分公式，藉助恆等式
% unit: noether.p13.eq.021
\[
\varphi(x,u,\ldots)\frac{\p^\tau p(x)}{\p x^\tau}
=(-1)^\tau\cdot\frac{\p^\tau\varphi}{\p x^\tau}\cdot p(x)
\quad \bmod \text{ divergences},
\]
可以把 \(p\) 的導數換成 \(p\) 本身以及關於 \(p\) 及其導數線性的散度。因此：
% unit: noether.p13.eq.022
\begin{equation}
\begin{aligned}
\sum\psi_i\delta u_i
={}&\sum_\lambda\left\{
(a_i^{(\lambda)}\psi_i)-\frac{\p}{\p x}\bigl(b_i^{(\lambda)}\psi_i\bigr)
+\cdots+(-1^\sigma)
\frac{\p^\sigma}{\p x^\sigma}\bigl(c_i^{(\lambda)}\psi_i\bigr)
\right\}p^{(\lambda)}+\Div\varGamma,
\end{aligned}
\tag{14}
\end{equation}
其中大括號內對 \(i\) 求和。與 (12) 結合，這給出
% unit: noether.p13.eq.023
\begin{equation}
\sum\left\{
(a_i^{(\lambda)}\psi_i)-\frac{\p}{\p x}\bigl(b_i^{(\lambda)}\psi_i\bigr)
+\cdots+(-1)^\sigma
\frac{\p^\sigma}{\p x^\sigma}\bigl(c_i^{(\lambda)}\psi_i\bigr)
\right\}p^{(\lambda)}
=\Div(B-\varGamma).
\tag{15}
\end{equation}
我現在對 (15) 形成 \(n\) 重積分，積分擴展到任意區域，並選擇 \(p(x)\)，使它們連同 \((B-\varGamma)\) 中出現的所有導數都在邊界上消失。由於散度的積分化為邊界積分，(15) 左端的積分對於在邊界上連同足夠多導數一起消失的任意 \(p(x)\) 也為零。由通常的論證推出，被積函數對於每個 \(p(x)\) 都為零，於是得到 \(\varrho\) 個關係
% unit: noether.p13.eq.024
\begin{equation}
\sum\left\{
(a_i^{(\lambda)}\psi_i)-\frac{\p}{\p x}\bigl(b_i^{(\lambda)}\psi_i\bigr)
+\cdots+(-1)^\sigma
\frac{\p^\sigma}{\p x^\sigma}\bigl(c_i^{(\lambda)}\psi_i\bigr)
\right\}=0\qquad(\lambda=1,2\ldots\varrho).
\tag{16}
\end{equation}
當 \(I\) 關於 \(\G_{\infty\varrho}\) 不變時，這些就是所求的 Lagrange 表達式及其導數之間的依賴關係。線性無關性如上證明，因為反命題又引回 (12)；並且人們可以再次由無窮小變換推出有限變換，如 \S{}4 將更詳細地進行。因此，對於一個 \(\G_{\infty\varrho}\)，在無窮小變換中已經出現 \(\varrho\) 個任意變換。由 (15) 和 (16) 還推出 \(\Div(B-\varGamma)=0\)。

% unit: noether.p13.par.015
如果按照一個“混合群”的情形取 \(\varDelta x\) 和 \(\varDelta u\) 關於 \(\eps\) 和 \(p(x)\) 線性，那麼通過一次令 \(p(x)\) 為零、一次令 \(\eps\) 為零，就看出散度關係 (13) 和依賴關係 (16) 都存在。

% unit: noether.p13.sec.003
\section*{\S{}3. 有限群情形下的反命題}

% unit: noether.p13.par.016
為了證明反命題，首先必須把前面的考察本質上反向進行。由 (13) 的有效性，乘以各個 \(\eps\) 並相加，得到 (12) 的有效性；並且藉助恆等式 (3) 得到一個關係
% unit: noether.p13.eq.025
\[
\bd f+\Div(A-B)=0.
\]
因此，如果令 \(\varDelta x=\frac1f(A-B)\)，就由此到達 (11)；經積分最終得到 (7)，即 \(\varDelta I=0\)，所以 \(I\) 關於由 \(\varDelta x,\varDelta u\) 確定的無窮小變換不變，其中 \(\varDelta u\) 由 (9) 從 \(\varDelta x\) 與 \(\bd u\) 確定，而且 \(\varDelta x,\varDelta u\) 關於參數成為線性的。但 \(\varDelta I=0\) 以已知方式蘊含 \(I\) 關於由同時系統的積分產生的有限變換不變：
% unit: noether.p13.eq.026
\begin{equation}
\frac{d x}{d t}=\varDelta x_i;\qquad
\frac{d u_i}{d t}=\varDelta u_i;\qquad
\left(\begin{array}{l}
x_i=y,\\
u_i=v_i
\end{array}\right\}
\quad\text{當 }t=0.
\tag{17}
\end{equation}
這些有限變換含有 \(\varrho\) 個參數 \(a_1\ldots a_\partial\)，即組合 \(t\eps_1,\ldots,t\eps_\varrho\)。由“應有且僅有 \(\rho\) 個線性無關的散度關係 (13)”這一假設進一步推出，只要這些有限變換不含導數 \(\frac{\p u}{\p x}\)，它們總是構成一個群。在相反情形中，至少一個通過 Lie 括號過程產生的無窮小變換不是其餘 \(\varrho\) 個的線性組合；而由於 \(I\) 也容許這個變換，就會有多於 \(\varrho\) 個線性無關的散度關係。或者這個無窮小變換具有特殊形式 \(\bd u=0,\ \Div(f\cdot\varDelta x)=0\)。但這時 \(\varDelta x\) 或 \(\varDelta u\) 就會依賴於導數，違反假設。當導數出現在 \(\varDelta x\) 或 \(\varDelta u\) 中時，這種情形是否可能發生，必須暫且不論；那麼還必須把所有使 \(\Div(f\cdot\varDelta x)=0\) 的函數 \(\varDelta x\) 加到上面已確定的 \(\varDelta x\) 上，以便重新得到群性質，但按照約定，由此增加的參數不予計算。反命題由此證明。

% unit: noether.p13.par.017
從這個反命題還推出，\(\varDelta x\) 和 \(\varDelta u\) 確實可以假定為關於參數線性的。如果 \(\varDelta u\) 和 \(\varDelta x\) 是關於 \(\eps\) 的較高次數形式，那麼由於 \(\eps\) 的冪乘積線性無關，會得到完全相應的關係 (13)，只是數量更多；由反命題，這些關係意味着 \(I\) 關於一個群不變，而這個群的無窮小變換線性地含有參數。如果這個群正好含有 \(\varrho\) 個參數，那麼原來由 \(\eps\) 的高次項得到的散度關係之間必須存在線性依賴。

% unit: noether.p13.par.018
還應注意，在 \(\varDelta x\) 和 \(\varDelta u\) 也含有 \(u\) 的導數的情形中，有限變換可以依賴於 \(u\) 的無窮多個導數；因為在這種情形中，對 (17) 積分時，在確定
% unit: noether.p13.eq.027
\[
\frac{\dd^2x_i}{\dd t^2},\quad \frac{\dd^2u_i}{\dd t^2},
\]
時會導致
% unit: noether.p13.eq.028
\[
\D\!\left(\frac{\p u}{\p x_\lambda}\right)
=\frac{\p\,\D u}{\p x_\lambda}
-\sum_\nu\frac{\p u}{\p x_\nu}
\frac{\p\,\D x_\nu}{\p x_\lambda},
\]
因此一般說來，\(u\) 的導數個數在每一步都增加。一個例子是：
% unit: noether.p13.eq.029
\[
f=\frac12 u'^2,\qquad
\psi=-u'',\qquad
\psi\cdot x=\frac{d}{d x}(u-u'x);\qquad
\bd u=x\cdot\eps,
\]
% unit: noether.p13.eq.030
\[
\varDelta x=\frac{-2u}{u'^2}\eps;\qquad
\varDelta u=\left(x-\frac{2u}{u'}\right)\cdot\eps.
\]
由於一個散度的 Lagrange 表達式恆等為零，反命題最後說明如下：如果 \(I\) 容許一個 \(\G_\varrho\)，那麼任何只與 \(I\) 相差一個邊界積分，也就是相差一個散度積分的積分，也容許一個具有相同 \(\bd u\) 的 \(\G_\varrho\)，而它的無窮小變換一般會含有 \(u\) 的導數。因此，對應於前面的例子，
% unit: noether.p13.eq.031
\[
f^*=\frac12\left\{u'^2-\frac{\dd}{\dd x}\left(\frac{u^2}{x}\right)\right\}
\]
容許無窮小變換 \(\varDelta u=x\eps,\ \varDelta x=0\)；而與 \(f\) 對應的無窮小變換中則出現 \(u\) 的導數。

% unit: noether.p13.par.019
如果過渡到變分問題，也就是說令 \(\psi_i=0\)，\footnote{\(\psi_i=0\)，或者稍更一般地，\(\psi_i=T_i\)，其中 \(T_i\) 是新添入的函數，在物理學中稱為“場方程”。在 \(\psi_i=T_i\) 的情形中，恆等式 (13) 變為方程 \(\Div B^{(\lambda)}=\sum T_i\delta u_i^{(\lambda)}\)，它們在物理學中仍稱為守恆律。} 那麼 (13) 變為方程
% unit: noether.p13.eq.032
\[
\Div B^{(1)}=0,\ldots,\qquad \Div B^{(\rho)}=0,
\]
它們常被稱為“守恆律”。在一維情形中，這給出
% unit: noether.p13.eq.033
\[
B^{(1)}=\mathrm{const.},\ldots,\qquad B^{(\rho)}=\mathrm{const.}.
\]
這裡，只要 \(\varDelta u\) 和 \(\varDelta x\) 不含高於 \(f\) 中出現的 \(\varkappa\) 階導數的導數，\(B\) 至多含有 \(u\) 的 \((2\kappa-1)\) 階導數（由 (6)）。由於 \(\psi\) 一般含有 \(2\varkappa\) 階導數，\footnote{只要 \(f\) 在 \(\varkappa\) 階導數中不是線性的。} 因而就有 \(\varrho\) 個一階積分存在。它們之間可以有非線性依賴，這也由上面的 \(f\) 說明。與線性無關的 \(\varDelta u=\eps_1,\ \varDelta x=\eps_2\) 相應的是線性無關關係
% unit: noether.p13.eq.034
\[
u''=\frac{\dd}{\dd x}u',\qquad
u''u'=\frac12\frac{\dd}{\dd x}(u')^2,
\]
而在一階積分 \(u'=\mathrm{const.},\ u'^2=\mathrm{const.}\) 之間存在非線性依賴。這是 \(\varDelta u,\varDelta x\) 不含 \(u\) 的導數的初等情形。\footnote{否則，對於每個 \(\lambda\)，還另外有 \(u''u'^{\lambda-1}=\mathrm{const.}\)，對應於
\[
u''(u')^{\lambda-1}=\frac1\lambda\frac{\dd}{\dd x}(u')^\lambda.
\]}
%
%
% unit: noether.p13.sec.004
\section*{\S{}4. 無限群情形中的逆命題}

% unit: noether.p13.par.061
首先要說明，假定 \(\D x\) 與 \(\D u\) 的線性並不構成限制。在這裡，即使不借助逆命題，也可由下述事實推出這一點：\(\G_{\infty\rho}\) 在形式上依賴於 \(\rho\) 個、且僅依賴於 \(\rho\) 個任意函數。確切地說，在非線性情形下，變換相複合時，最低階項相加，而任意函數的個數就會增加。比如設
% unit: noether.p13.eq.018a
\[
\begin{aligned}
y&=A\!\left(x,u,\frac{\p u}{\p x},\cdots;p\right)\\
&=x+\sum a(x,u,\ldots)p^\nu
+b(x,u,\ldots)p^{\nu-1}\frac{\p p}{\p x}
+c\,p^{\nu-2}\left(\frac{\p p}{\p x}\right)^2+\cdots
+d\left(\frac{\p p}{\p x}\right)^\nu+\cdots,
\end{aligned}
\]
其中
% unit: noether.p13.eq.018b
\[
p^\nu=(p^{(1)})^{\nu_1}\cdots(p^{(\rho)})^{\nu_\rho},
\]
並且相應地有
% unit: noether.p13.eq.018c
\[
v=B\!\left(x,u,\frac{\p u}{\p x},\ldots;p\right).
\]
於是同
% unit: noether.p13.eq.018d
\[
z=A\!\left(y,v,\frac{\p v}{\p y},\ldots;q\right)
\]
複合後，最低階項給出
% unit: noether.p13.eq.018e
\[
z=x+\sum a(p^\nu+q^\nu)+
b\left\{p^{\nu-1}\frac{\p p}{\p x}+q^{\nu-1}\frac{\p q}{\p x}\right\}
+c\left\{p^{\nu-2}\left(\frac{\p p}{\p x}\right)^2
+q^{\nu-2}\left(\frac{\p q}{\p x}\right)^2\right\}+\cdots.
\]
如果這裡有某個不同於 \(a\) 和 \(b\) 的係數不為零，從而對於 \(\sigma>1\) 實際出現一項
% unit: noether.p13.eq.018f
\[
p^{\nu-\sigma}\left(\frac{\p p}{\p x}\right)^\sigma
+q^{\nu-\sigma}\left(\frac{\p q}{\p x}\right)^\sigma,
\]
那麼它不能寫成單個函數的微分商，也不能寫成這種微分商的冪乘積。因此任意函數的個數違反假設而增加了。如果所有不同於 \(a\) 與 \(b\) 的係數均為零，那麼依指數 \(\nu_1\ldots\nu_\varrho\) 的取值，第二項或者就是第一項的微分商（例如對 \(\G_{\infty1}\) 總是如此），於是線性事實上出現；或者任意函數的個數在這裡也會增加。因此，由於關於 \(p(x)\) 的線性， 這些無窮小變換即滿足一個線性偏微分方程組；又因為群性質成立，它們就按 Lie 的定義形成一個“無限的無窮小變換群”（\emph{Grundlagen}, \S{}10）。

% unit: noether.p13.par.062
逆命題現在按與有限群情形相似的方式得到。依賴關係 (16) 成立時，把它們乘以 \(p^{(\lambda)}(x)\) 並相加，藉助恆等變形 (14) 得到
% unit: noether.p13.eq.019
\[
\sum_i\psi_i\bd u_i=\Div\varGamma.
\]
由此如 \S{}3 中一樣得到 \(\varDelta x\) 和 \(\varDelta u\) 的確定，以及 \(I\) 對這些無窮小變換的不變性；這些變換事實上線性地依賴於 \(\varrho\) 個任意函數及其直到 \(\sigma\) 階的導數。當這些無窮小變換不含 \(\frac{\p u}{\p x},\cdots\) 的導數時，它們確實形成一個群，這也如 \S{}3 中一樣：否則通過複合會出現更多任意函數，而按假設只應有 \(\varrho\) 個依賴關係 (16)。所以它們形成一個“無窮小變換的無限群”。但這樣的群（\emph{Grundlagen}, 定理 VII, p. 391）由某個由此定義的、按 Lie 意義的“有限變換的無限群 \(\G\)”的最一般無窮小變換組成。每個有限變換於是由無窮小變換生成（\emph{Grundlagen}, \S{}7），\footnote{由此特別可知，由 \(\G_{\infty\varrho}\) 的無窮小變換 \(\varDelta x,\varDelta u\) 生成的群 \(\G\) 又回到 \(\G_{\infty\varrho}\)。因為 \(\G_{\infty\varrho}\) 不含依賴於任意函數且不同於 \(\varDelta x,\varDelta u\) 的無窮小變換，也不能含有與之無關而依賴於參數的變換；否則它就是混合群。而按上文，有限變換由無窮小變換確定。} 因而通過積分下面的聯立系統而產生:
% unit: noether.p13.eq.020
\[
\frac{d x_i}{d t}=\varDelta x_i;\qquad
\frac{d u_i}{d t}=\varDelta u_i\quad
\left(\begin{array}{l}
x_i=y_i,\\
u_i=v
\end{array}\right\}
\quad\text{當 }t=0.
\]
不過可能有必要仍把任意的 \(p(x)\) 看作依賴於 \(t\)。因此 \(\G\) 實際上依賴於 \(\varrho\) 個任意函數。特別地，如果取 \(p(x)\) 與 \(t\) 無關已經足夠，那麼這種依賴性在任意函數 \(q(x)=t\cdot p(x)\) 中是解析的。\footnote{後一種情形是否也許總會發生，這個問題 Lie 曾以另一種表述提出（\emph{Grundlagen}, \S{}7 和 \S{}13 末尾）。} 如果出現 \(\frac{\p u}{\p x},\cdots\) 的導數，那麼在作同樣結論之前，可能還需要附加無窮小變換 \(\bd u=0,\ \Div(f\cdot\varDelta x)=0\)。

% unit: noether.p13.par.063
接着 Lie 的一個例子（\emph{Grundlagen}, \S{}7），現在指出一個相當一般的情形，在其中可以達到顯式公式；這些公式同時顯示，任意函數的導數出現到 \(\sigma\) 階為止，因此在這種情形下逆命題是完全的。這裡指的是這樣一些無窮小變換群：與它們相對應的是所有 \(x\) 的變換群，以及由此“誘導”的 \(u\) 的變換；也就是說，在這些 \(u\) 的變換中，\(\varDelta u\)，因而 \(u\)，只依賴於 \(\varDelta x\) 中出現的任意函數。還假設 \(\varDelta u\) 中不出現 \(\frac{\p u}{\p x},\cdots\) 的導數。於是有:
% unit: noether.p13.eq.021
\[
\varDelta x_i=p^{(i)}(x);\qquad
\varDelta u_i=
\sum_{\lambda=1}^{n}\left\{
a^{(\lambda)}(x,u)p^{(\lambda)}
+b^{(\lambda)}\frac{\p p^{(\lambda)}}{\p x}
+\cdots+
c^{(\lambda)}\frac{\p^\sigma p^{(\lambda)}}{\p x^\sigma}
\right\}.
\]
因為無窮小變換 \(\varDelta x=p(x)\) 生成每個帶任意 \(g(y)\) 的變換 \(x=y+g(y)\)，所以特別可以把 \(p(x)\) 作為 \(t\) 的函數來確定，使它生成一參數群:
% unit: noether.p13.eq.022
\begin{equation}
x_i=y_i+t\cdot g_i(y),
\tag{18}
\end{equation}
該群在 \(t=0\) 時化為恆等，在 \(t=1\) 時化為所求的 \(x=y+g(y)\)。對 (18) 微分，便得到:
% unit: noether.p13.eq.023
\begin{equation}
\frac{d x_i}{d t}=g_i(y)=p^{(i)}(x,t),
\tag{19}
\end{equation}
其中 \(p(x,t)\) 由 \(g(y)\) 通過反解 (18) 來確定；反過來，(18) 由 (19) 以及附加條件 \(x_i=y_i\) 當 \(t=0\) 得到，該條件唯一確定積分。藉助 (18)，\(\varDelta u\) 中的 \(x\) 可以由“積分常數” \(y\) 和 \(t\) 代替；於是 \(g(y)\) 恰好出現到其 \(\sigma\) 階導數為止，因為在
% unit: noether.p13.eq.024
\[
\frac{\p p}{\p x}=\sum_k\frac{\p g}{\p y_k}\frac{\p y_k}{\p x}
\]
中，\(\frac{\p y}{\p x}\) 由 \(\frac{\p x}{\p y}\) 表示，而一般地，\(\frac{\p^\sigma p}{\p x^\sigma}\) 被它在
% unit: noether.p13.eq.025
\[
\frac{\p g}{\p y},\ldots,\frac{\p^\sigma g}{\p y^\sigma},\ldots,\frac{\p^\sigma x}{\p y^\sigma}
\]
中的值所代替。為了確定 \(u\)，於是得到方程組
% unit: noether.p13.eq.026
\[
\frac{d u_i}{d t}
=
F_i\!\left(g(y),\frac{\p g}{\p y},\cdots\frac{\p^\sigma g}{\p y^\sigma},u,t\right),
\qquad (u_i=v_i\ \text{當 }t=0),
\]
其中只有 \(t\) 和 \(u\) 是變量，而 \(g(y),\ldots\) 屬於係數域；所以積分給出
% unit: noether.p13.eq.027
\[
u_i=v_i+
B_i\!\left(v,g(y),\frac{\p g}{\p y},\cdots\frac{\p^\sigma g}{\p y^\sigma},t\right)_{t=1}.
\]
由此得到恰好依賴於任意函數的 \(\sigma\) 階導數的變換。按 (18)，當 \(g(y)=0\) 時恆等變換包含在其中；群性質則來自下述事實：所述方法給出每個變換 \(x=y+g(y)\)，而由此 \(u\) 的誘導變換被唯一確定，所以群 \(\G\) 被窮盡。

% unit: noether.p13.par.064
從逆命題還順帶得出：假定任意函數只依賴於 \(x\)，而不依賴於 \(u,\frac{\p u}{\p x},\cdots\)，並不構成限制。在後一種情形下，在恆等變形 (14) 中，因而也在 (15) 中，除了 \(p^{(\lambda)}\) 以外，還會出現 \(\frac{\p p^{(\lambda)}}{\p u},\ \frac{\p p^{(\lambda)}}{\p\frac{\p u}{\p x}},\cdots\)。如果現在依次把 \(p^{(\lambda)}\) 取為關於 \(u,\frac{\p u}{\p x},\cdots\) 的零次、一次、……式，而以 \(x\) 的任意函數為係數，那麼又會產生依賴關係 (16)，只是數目更多；但按上述逆命題，通過同只依賴於 \(x\) 的任意函數合併，它們又可歸結到先前情形。同樣可以證明，同時出現依賴關係和與之無關的散度關係時，對應於混合群。\footnote{如 \S{}3 中一樣，這裡也由逆命題得出：除 \(I\) 以外，凡同它相差一個關於散度的積分的積分 \(I^*\)，也都容許一個具有相同 \(\bd u\) 的無限群；不過這裡 \(\varDelta x\) 和 \(\varDelta u\) 一般將含有 \(u\) 的導數。Einstein 在廣義相對論中引入了這樣一個積分 \(I^*\)，以便得到能量定理的較簡單表述；我按照 Klein 第二篇札記的記號，給出這個 \(I^*\) 所容許的無窮小變換。積分 \(I=\int\!\ldots\!\int K\,d\omega=\int\!\ldots\!\int \mathfrak{R}\,d S\) 容許所有 \(w\) 的變換群以及由此為 \(g_{\mu\nu}\) 誘導出的群；與之相應的是依賴關係（Klein 的 (30)）:
\[
\sum \mathfrak{R}_{\mu\nu}g_\tau^{\mu\nu}
+2\sum\frac{\p g^{\mu\nu}\mathfrak{R}_{\mu\tau}}{\p w^\sigma}=0.
\]
現在 \(I^*=\int\!\ldots\!\int \mathfrak{R}^*\,d S\)，其中 \(\mathfrak{R}^*=\mathfrak{R}+\Div\)，因而 \(\mathfrak{R}_{\mu\nu}^*=\mathfrak{R}_{\mu\nu}\)，這裡 \(\mathfrak{R}_{\mu\nu}^*,\mathfrak{R}_{\mu\nu}\) 分別表示 Lagrange 表達式。所給依賴關係因此也是關於 \(\mathfrak{R}_{\mu\nu}^*\) 的依賴關係；乘以 \(p^\tau\) 並相加之後，反向使用乘積微分的變形，得到
\[
\sum \mathfrak{R}_{\mu\nu}p^{\mu\nu}
+2\Div\!\left(\sum g^{\mu\sigma}\mathfrak{R}_{\mu\tau}p^\tau\right)=0,
\]
\[
\delta\mathfrak{R}^*+\Div\!\left(\sum 2g^{\mu\sigma}\mathfrak{R}_{\mu\tau}p^\tau
-\frac{\p\mathfrak{R}^*}{\p g^{\mu\nu}_{\sigma}}p^{\mu\nu}\right)=0.
\]
同 Lie 微分方程 \(\delta\mathfrak{R}^*+\Div(\mathfrak{R}^*\varDelta w)=0\) 比較，給出
\[
\varDelta w^\sigma=\frac1{\mathfrak{R}^*}\cdot\left(\sum (2g^{\mu\sigma}\mathfrak{R}_{\mu\tau}p^\tau
-\frac{\p\mathfrak{R}^*}{\p g^{\mu\nu}_{\sigma}}p^{\mu\nu}\right),
\qquad
\varDelta g^{\mu\nu}=p^{\mu\nu}+\sum g^{\mu\nu}_{\sigma}\varDelta w^\sigma,
\]
作為 \(I^*\) 所容許的無窮小變換。因此這些無窮小變換依賴於 \(g^{\mu\nu}\) 的一階和二階導數，並且含有直到一階導數為止的任意 \(p\)。}

% unit: noether.p13.sec.005
\section*{\S{}5. 關係各個組成部分的不變性}

% unit: noether.p13.par.065
如果把群 \(\G\) 特化為最簡單、通常考慮的情形，即在變換中不允許 \(u\) 的導數，並且使變換後的自變量只依賴於 \(x\) 而不依賴於 \(u\)，那麼就可以推出公式中各個組成部分的不變性。首先，按已知的推理，可得
% unit: noether.p13.eq.028
\[
\int\!\cdots\!\int\left(\sum_i\psi_i\delta u_i\right)\dd x
\]
的不變性；因而得到 \(\sum_i\psi_i\delta u_i\) 的相對不變性\footnote{即 \(\sum_i\psi_i\delta u_i\) 在變換下獲得一個因子；在代數不變量論中，這一情形一向稱為相對不變性。}，其中 \(\delta\) 表示任意變分。確實，一方面
% unit: noether.p13.eq.029
\[
\delta I
=\int\!\cdots\!\int \delta f\!\left(x,u,\frac{\p u}{\p x},\cdots\right)d x
=\int\!\cdots\!\int \delta f\!\left(y,v,\frac{\p v}{\p y},\cdots\right)d y.
\]
另一方面，對於在邊界上消失的 \(\delta u,\delta\frac{\p u}{\p x},\cdots\)，由於 \(\delta u,\delta\frac{\p u}{\p x},\cdots\) 的線性齊次變換，與之相應的 \(\delta v,\delta\frac{\p v}{\p y},\cdots\) 也在邊界上消失，於是有
% unit: noether.p13.eq.030
\[
\int\!\cdots\!\int \delta f\!\left(x,u,\frac{\p u}{\p x},\cdots\right)d x
=
\int\!\cdots\!\int\left(\sum\psi_i(u,\ldots)\delta u_i\right)d x,
\]
% unit: noether.p13.eq.031
\[
\int\!\cdots\!\int \delta f\!\left(y,v,\frac{\p v}{\p x},\cdots\right)d y
=
\int\!\cdots\!\int\left(\sum\psi_i(v,\ldots)\delta v_i\right)d y.
\]
於是，對於在邊界上消失的 \(\delta u,\delta\frac{\p u}{\p x},\cdots\)，
% unit: noether.p13.eq.032
\[
\int\!\cdots\!\int\left(\sum\psi_i(u,\ldots)\delta u_i\right)d x
=
\int\!\cdots\!\int\left(\sum\psi_i(v,\ldots)\delta v_i\right)d y
\]
% unit: noether.p13.eq.033
\[
=
\int\!\cdots\!\int\left(\sum\psi_i(v,\ldots)\delta v_i\right)
\left|\frac{\p y_i}{\p x_\varkappa}\right|d x.
\]
把第三個積分中的 \(y,v,\delta v\) 用 \(x,u,\delta u\) 表示，並使它等於第一個積分，就得到一個關係
% unit: noether.p13.eq.034
\[
\int\!\cdots\!\int\left(\sum\chi_i(u,\ldots)\delta u_i\right)d x=0
\]
其中 \(\delta u\) 在邊界上消失而除此以外任意。由此照常推出，對任意 \(\delta u\)，被積函數為零；所以在 \(\delta u\) 中恆有關係
% unit: noether.p13.eq.035
\[
\sum_i\psi_i(u,\ldots)\delta u_i
=
\left|\frac{\p y_i}{\p x_k}\right|
\left(\sum_i\psi_i(v,\ldots)\delta v_i\right).
\]
這表明 \(\sum_i\psi_i\delta u_i\) 的相對不變性，並因而表明
% unit: noether.p13.eq.036
\[
\int\!\cdots\!\int\left(\sum_i\psi_i\delta u_i\right)\dd x
\]
的不變性。\footnote{如果 \(y\) 也依賴於 \(u\)，這些推理便失效，因為此時 \(\delta f(y,v,\p v/\p y,\ldots)\) 還含有項 \(\sum(\p f/\p y)\delta y\)，所以散度變形不再導向 Lagrange 表達式。如果允許 \(u\) 的導數，它們同樣失效；這時 \(\delta v\) 成為 \(\delta u,\delta(\p u/\p x),\ldots\) 的線性組合，因而只有經過進一步的散度變形，才導向一個恆等式 \(\int\!\cdots\!\int(\sum\chi_i(u,\ldots)\delta u_i)\dd x=0\)，於是右邊又不出現 Lagrange 表達式。根據逆命題，是否可以從 \(\int\!\cdots\!\int(\sum\psi_i\delta u_i)\dd x\) 的不變性就推出散度關係的存在，等價於是否可以由此推出 \(I\) 對某個群的不變性；這個群不必導向相同的 \(\D u,\D x\)，但必須導向相同的 \(\bd u\)。在簡單積分且 \(f\) 中只有一階導數的特殊情形下，對有限群，可以從 Lagrange 表達式的不變性推出第一積分的存在（例如參見 Engel, Gött. Nachr. 1916, p. 270）。}

% unit: noether.p13.par.066
為了把這一點應用到導出的散度關係和依賴關係上，首先要證明，由 \(\varDelta u,\varDelta x\) 導出的 \(\bd u\) 確實滿足變分 \(\delta u\) 的變換律；只要在 \(\bd v\) 中的參數或任意函數，按它們在 \(y,v\) 中相似的無窮小變換群所對應的方式來確定即可。記 \(T_q\) 為把 \(x,u\) 變為 \(y,v\) 的變換；記 \(T_p\) 為 \(x,u\) 中的一個無窮小變換。則 \(y,v\) 中相似的變換由 \(T_r=T_qT_pT_q^{-1}\) 給出，其中參數或任意函數 \(r\) 由 \(p\) 與 \(q\) 確定。用公式表示就是:
% unit: noether.p13.eq.037
\[
T_p:\ \xi=x+\D x(x,p),\qquad u^*=u+\D u(x,u,p);
\]
% unit: noether.p13.eq.038
\[
T_q:\ y=A(x,q),\qquad v=B(x,u,q);
\]
% unit: noether.p13.eq.039
\[
T_qT_p:\ \eta=A(x+\D x(x,p),q),\quad
v^*=B(x+\varDelta x(p),\,u+\varDelta u(p),q).
\]
由此產生 \(T_r=T_qT_pT_q^{-1}\)，於是
% unit: noether.p13.eq.040
\[
\eta=y+\varDelta y(r);\qquad v^*=v+\varDelta v(r),
\]
其中藉助 \(T_q\) 的逆，把 \(x\) 看作 \(y\) 的函數，並只保留無窮小項。因此有恆等式
% unit: noether.p13.eq.041
\begin{equation}
\eta=y+\varDelta y(r)=y+\sum\frac{\p A(x,q)}{\p x}\varDelta x(p),
\tag{20a}
\end{equation}
% unit: noether.p13.eq.042
\begin{equation}
v^*=v+\varDelta v(r)=v+\sum\frac{\p B(x,u,q)}{\p x}\varDelta x(p)
+\sum\frac{\p B(x,u,q)}{\p u}\varDelta u(p).
\tag{20b}
\end{equation}
在其中把 \(\xi=x+\varDelta x\) 代以 \(\xi-\varDelta \xi\)，使 \(\xi\) 又回到 \(x\)，因而 \(\varDelta x\) 消失；那麼按照第一式 (20)，\(\eta\) 也又回到 \(y=\eta-\varDelta\eta\)。如果在這種代換下 \(\varDelta u(p)\) 變為 \(\bd u(p)\)，那麼 \(\varDelta v(r)\) 也變為 \(\bd v(r)\)，第二式 (20) 給出
% unit: noether.p13.eq.043
\[
v+\bd v(y,v,\ldots r)=v+\sum\frac{\p B(x,u,q)}{\p u}\,\bd u(p),
\]
% unit: noether.p13.eq.044
\[
\bd v(y,v,\ldots r)=\sum\frac{\p B}{\p u_\varkappa}\,\bd u_\varkappa(x,u,p).
\]
於是，只要把 \(\bd v\) 看作依賴於參數或任意函數 \(r\)，變分的變換公式就確實得到滿足。\footnote{再次可見，為了使這些結論成立，必須假定 \(y\) 與 \(u\) 無關，等等。作為例子，可以提到 Klein 給出的 \(\delta g^{\mu\nu}\) 和 \(\delta q_\varrho\)：只要 \(p\) 經受一個向量變換，它們就滿足關於變分的變換。}

% unit: noether.p13.par.067
因此特別得到 \(\sum\psi_i\bd u_i\) 的相對不變性；又因 (12)，由於散度關係在 \(y,v\) 中也成立，也得到 \(\Div B\) 的相對不變性；再由 (14) 和 (13)，得到 \(\Div\varGamma\) 以及同 \(p^{(\lambda)}\) 合併起來的依賴關係左邊的相對不變性，其中在變換後的公式中，任意 \(p(x)\)（或參數）處處要由 \(r\) 代替。由此還得到 \(\Div(B-\varGamma)\) 的相對不變性，也就是說，得到一個並非恆等消失、但其散度恆等消失的函數系 \(B-\varGamma\) 的散度的相對不變性。

% unit: noether.p13.par.068
由 \(\Div B\) 的相對不變性，在一維情形和有限群情形下，還可以推出第一積分的不變性。對應於無窮小變換的參數變換，由 (20) 變為線性齊次；並且由於所有變換可逆，\(\eps\) 也線性齊次地依賴於變換後的參數 \(\eps^*\)。當置 \(\psi=0\) 時，這種可逆性當然仍保留，因為 (20) 中不出現 \(u\) 的導數。把
% unit: noether.p13.eq.045
\[
\Div B(x,u,\ldots\eps)=\frac{d y}{d x}\cdot\Div B(y,v,\ldots\eps^*)
\]
中 \(\eps^*\) 的係數相等起來，則 \(\frac{d}{d y}B^{\lambda)}(y,v,\ldots)\) 也成為 \(\frac{d}{d x}B^{(\lambda)}(x,u,\ldots)\) 的線性齊次函數。因此由
% unit: noether.p13.eq.046
\[
\frac{\dd}{\dd x}B^{(\lambda)}(x,u,\ldots)=0
\quad\text{或}\quad B^{(\lambda)}(x,u)=\mathrm{const.}
\]
也推出
% unit: noether.p13.eq.047
\[
\frac{\dd}{\dd y}B^{(\lambda)}(y,v,\ldots)=0
\quad\text{或}\quad B^{(\lambda)}(y,v)=\mathrm{const.}
\]
於是對應於一個 \(\G_\varrho\) 的 \(\rho\) 個第一積分各自容許這個群，所以進一步積分也得到簡化。最簡單的例子是 \(f\) 與 \(x\) 無關，或與某個 \(u\) 無關；這分別對應於無窮小變換 \(\varDelta x=\eps,\varDelta u=0\)，以及 \(\varDelta x=0,\varDelta u=\eps\)。這時 \(\bd u=-\eps\frac{d u}{d x}\)，相應地為 \(\eps\)；而 \(B\) 由 \(f\) 和 \(\bd u\) 通過微分與有理組合導出，所以它也與 \(x\) 或相應的 \(u\) 無關，並容許相應的群。\footnote{在那些已能由 \(\int(\sum\psi_i\delta u_i)d x\) 的不變性推出第一積分存在的情形中，這些第一積分並不容許完整的群 \(\G_\varrho\)。例如，\(\int(u''\delta u)d x\) 容許無窮小變換 \(\varDelta x=\eps_2;\ \varDelta u=\eps_1+x\eps_3\)；而第一積分 \(u-u'x=\mathrm{const.}\)，對應於 \(\varDelta x=0,\varDelta u=x\eps_3\)，並不容許另外兩個無窮小變換，因為它顯式含有 \(u\) 和 \(x\)。與這個第一積分相對應的是 \(f\) 的含有導數的無窮小變換。因此可見，\(\int\!\ldots\!\int(\sum\psi_i\,d u_i)d x\) 的不變性無論如何比 \(I\) 的不變性所能完成的要少；這一點是對前一條注中提出的問題所應指出的。}

% unit: noether.p13.sec.006
\section*{\S{}6. Hilbert 的一個斷言}

% unit: noether.p13.par.069
由前述內容最後還得到 Hilbert 關於真正能量定理失效同“廣義相對論”之間聯繫的一個斷言的證明（Klein 第一篇札記, Göttinger Nachr. 1917, Reply, 第一段），並且這是在推廣的群論表述下的證明。

% unit: noether.p13.par.070
設積分 \(I\) 容許一個 \(\G_{\infty\varrho}\)，並令 \(\G_\sigma\) 為通過任意函數的特化而產生的任意有限群，也就是 \(\G_{\infty\varrho}\) 的子群。則無限群 \(\G_{\infty\varrho}\) 對應於依賴關係 (16)，有限群 \(\G_\sigma\) 對應於散度關係 (13)；反過來，由任意散度關係的成立推出 \(I\) 對某個有限群的不變性，而這個有限群當且僅當 \(\bd u\) 是由 \(\G_\sigma\) 得到的那些 \(\bd u\) 的線性組合時才同 \(\G_\sigma\) 相同。因此，對 \(\G_\sigma\) 的不變性不能導出不同於 (13) 的散度關係。但是由於從 (16) 的成立可以推出 \(I\) 對 \(\G_{\infty\varrho}\) 的無窮小變換 \(\varDelta u,\varDelta x\) 在任意 \(p(x)\) 下的不變性，所以特別已經推出了對由這些變換經特化產生的 \(\G_\sigma\) 的無窮小變換的不變性，因而推出了對 \(\G_\sigma\) 的不變性。散度關係
% unit: noether.p13.eq.048
\[
\sum_i\psi_i\bd u_i^{(\lambda)}=\Div B^{(\lambda)}
\]
因此必須是依賴關係 (16) 的後果；後者也可寫作
% unit: noether.p13.eq.049
\[
\sum_i\psi_i a_i^{(\lambda)}=\Div\chi^{(\lambda)},
\]
其中 \(\chi^{(\lambda)}\) 是 Lagrange 表達式及其導數的線性組合。由於 \(\psi\) 在 (13) 和 (16) 中都是線性進入的，散度關係尤其必須是依賴關係 (16) 的線性組合。於是
% unit: noether.p13.eq.050
\[
\Div B^{(\lambda)}=\Div\left(\sum \alpha_\mu\chi^{(\mu)}\right),
\]
而 \(B^{(\lambda)}\) 本身就由 \(\chi\)，也就是由 Lagrange 表達式及其導數，再加上其散度恆等消失的函數線性組成；例如 \S{}2 末尾出現的 \(B-\varGamma\)，其中 \(\Div(B-\varGamma)=0\)，並且該散度同時具有不變量性質。凡是 \(B^{(\;)}\) 可以按上述方式由 Lagrange 表達式及其導數組成的散度關係，我稱為“非真正的”；所有其餘的稱為“真正的”。

% unit: noether.p13.par.071
反過來，如果散度關係是依賴關係 (16) 的線性組合，因而是“非真正的”，那麼對 \(\G_\sigma\) 的不變性由對 \(\G_{\infty\varrho}\) 的不變性推出；\(\G\) 成為 \(\G_{\infty\varrho}\) 的子群。對應於有限群 \(\G_\sigma\) 的散度關係，當且僅當 \(\G_\sigma\) 是某個無限群的子群且 \(I\) 對該無限群不變時，才變成非真正的。

% unit: noether.p13.par.072
通過群的特化，Hilbert 的原始斷言由此得出。所謂“平移群”，是指有限群
% unit: noether.p13.eq.051
\[
y_i=x_i+\eps_i;\qquad v_i(y)=u_i(x);
\]
因此
% unit: noether.p13.eq.052
\[
\varDelta x_i=\eps_i,\qquad \varDelta u_i=0,\qquad
\bd u_i=-\sum_\lambda\frac{\p u_i}{\p x_\lambda}\eps_\lambda.
\]
如所周知，對平移群的不變性就是說，在
% unit: noether.p13.eq.053
\[
I=\int\!\cdots\!\int f\!\left(x,u,\frac{\p u}{\p x},\ldots\right)\dd x
\]
中，\(x\) 不在 \(f\) 中顯式出現。相應的 \(n\) 個散度關係
% unit: noether.p13.eq.054
\[
\sum\psi_i\frac{\p u_i}{\p x_\lambda}=\Div B^{(\lambda)}
\qquad(\lambda=1,2,\ldots n)
\]
稱為“能量關係”，因為同變分問題相應的“守恆律” \(\Div B^{(\lambda)}=0\) 對應於“能量定理”，而 \(B^{(\lambda)}\) 對應於“能量分量”。於是得到：如果 \(I\) 容許平移群，那麼能量關係當且僅當 \(I\) 對一個包含平移群為子群的無限群不變時才成為非真正的。\footnote{經典力學中的能量定理，以及舊“相對論”（其中 \(\sum d x^2\) 變到自身）中的能量定理，都是“真正的”，因為這裡不出現無限群。}

% unit: noether.p13.par.073
這種無限群的一個例子由所有 \(x\) 的變換以及 \(u(x)\) 的相應誘導變換的群給出，在這些誘導變換中只出現任意函數 \(p(x)\) 的導數。平移群通過特化 \(p^{(i)}(x)=\eps_i\) 產生；不過必須懸而未決的是，是否由此以及由把 \(I\) 改變一個邊界積分所產生的群，就已經給出了這類群中最一般者。所述類型的誘導變換例如可由下述方式產生：使 \(u\) 服從一個“全微分形式”的係數變換，也就是說，服從形如
% unit: noether.p13.eq.055
\[
\sum a\,\dd^\lambda x_i+\sum b\,\dd^{\lambda-1}x_i\,\dd x_\varkappa+\cdots
\]
的形式的係數變換；該形式除了 \(d x\) 之外還含有更高微分。更特殊的誘導變換，其中 \(p(x)\) 只以一階導數出現，由通常微分形式 \(\sum c\,d x_{i_1}\ldots d x_{i_\lambda}\) 的係數變換給出；通常所考慮的正是這些。

% unit: noether.p13.par.074
所述類型的另一個群——由於出現對數項，它不能是一個係數變換——例如如下:
% unit: noether.p13.eq.056
\[
y=x+p(x),\qquad
v_i=u_i+\lg(1+p'(x))=u_i+\lg\frac{d y}{d x};
\]
% unit: noether.p13.eq.057
\[
\varDelta x=p(x);\qquad
\varDelta u_i=p'(x)\footnotemark;\qquad
\bd u_i=p'(x)-u_i'p(x).
\]
\footnotetext{由這些無窮小變換，按照 \S{} 4 末尾所述的方法反向求出有限變換。}
這裡依賴關係 (16) 變為
% unit: noether.p13.eq.058
\[
\sum_i\left(\psi_i u_i'+\frac{d\psi_i}{d x}\right)=0,
\]
非真正的能量關係變為
% unit: noether.p13.eq.059
\[
\sum\left(\psi_i u_i'+
\frac{d(\psi_i+\mathrm{const.})}{d x}\right)=0.
\]
該群的一個最簡單的不變積分是
% unit: noether.p13.eq.060
\[
I=\int \frac{e^{-2u_1}}{u_1'-u_2'}\,d x.
\]
最一般的 \(I\) 由 Lie 微分方程 (11)
% unit: noether.p13.eq.061
\[
\bd f+\frac{d}{d x}(f\cdot \varDelta x)=0
\]
的積分確定；當把 \(\varDelta x\) 與 \(\bd u\) 的值代入，並且假定 \(f\) 只依賴於 \(u\) 的一階導數時，它變為
% unit: noether.p13.eq.062
\[
\frac{\p f}{\p x}p(x)
+\left\{\sum\frac{\p f}{\p u_i}
-\frac{\p f}{\p u_i'}u_i'+f\right\}p'(x)
+\left\{\sum\frac{\p f}{\p u_i''}\right\}p''(x)=0
\]
（在 \(p(x),p'(x),p''(x)\) 中恆等）。這個方程組已經在兩個函數 \(u(x)\) 的情形下有實際含有導數的解，即
% unit: noether.p13.eq.063
\[
f=(n_1'-u_2')\,
\varPhi\!\left(u_1-u_2,\frac{e^{-u_1}}{u_1'-u_2'}\right),
\]
其中 \(\varPhi\) 表示所示自變量的任意函數。

% unit: noether.p13.par.075
Hilbert 表述他的斷言時說，真正能量定理的失效是“廣義相對論”的一個特徵性標誌。為了使這個斷言按字面成立，就必須把“廣義相對論”這個稱謂理解得比通常更廣，並且也擴展到上述這些依賴於 \(n\) 個任意函數的群。\footnote{這再次確認了 Klein 一個評語的正確性：物理學中通常使用的“相對論”這一稱謂，應當由“相對於一個群的不變性”來取代。（\emph{Über die geometrischen Grundlagen der Lorentzgruppe}. \emph{Jhrber. d. d. Math. Vereinig.} vol. 19, p. 287, 1910；重印於 \emph{Physikalische Zeitschrift}。）}
%
\clearpage
%
% unit: noether.p14.title.001
\pdfbookmark[0]{14. 一變量代數函數的算術理論及其與其他理論和數域理論的關係}{paper-14}
\section*{14. 一變量代數函數的算術理論及其與其他理論和數域理論的關係}
\begin{center}
\emph{Jahresbericht der Deutschen Mathematiker-Vereinigung} 38 (1919), pp. 182--203
\end{center}

% unit: noether.p14.par.001
下面這篇報告，是出於在代數函數的各個理論之間建立一個聯繫環節的願望而產生的；同時它也可以看作對 Brill--Noether 報告的一個補充，\footnote{代數函數理論在較早和較近時期的發展。\emph{Jahresbericht der Deutschen Mathematiker-Vereinigung} vol. 3 (1894) [引用為 Brill--Noether]。}因為在那篇報告中，算術理論基本上沒有被討論。\footnote{除了對 Kronecker 關於判別式的研究的一段討論之外。（\emph{J. f. M.} vol. 91。）}

% unit: noether.p14.par.002
關於各個理論的論述，除 Brill--Noether 外，還應考慮下列百科全書條目及其中所列的文獻：

% unit: noether.p14.par.003
關於 Riemann 理論：Wirtinger (II B 2), nos. 1--19；

% unit: noether.p14.par.004
關於 Weierstrass 理論：Wirtinger (II B 2), nos. 23, 32, 33, 34；

% unit: noether.p14.par.005
關於 Brill--Noether 的代數幾何理論：Wirtinger (II B 2), nos. 21--31；Berzolari (III C 4), nos. 23--38；

% unit: noether.p14.par.006
關於 Dedekind--Weber 和 Hensel--Landsberg 的算術理論：Hensel (II C 5), [引用為 Hensel]；\footnote{涉及 Dedekind--Weber 的部分歸 A. Ostrowski。}

% unit: noether.p14.par.007
更一般地，關於代數量的算術理論：Landsberg (I C 5)，以及關於若干變量代數函數的算術理論：Jung (II C 6)。

% unit: noether.p14.par.008
關於代數數論，可比較 Hilbert 的《Zahlbericht》\footnote{代數數域理論。\emph{Jahresbericht der Deutschen Mathematiker-Vereinigung} vol. 4 (1894/95)。}以及 Dirichlet--Dedekind 的數論講義 [引用為 Dedekind--數論]。\footnote{Braunschweig, Vieweg，第 4 版，1894。}

% unit: noether.p14.toc.001
\subsection*{目錄}
\begin{enumerate}
\item 各種理論的基礎。
\item 代數數和代數函數之間的平行性（共軛元素、基定理、模概念）。
\item 代數數和代數函數的理想理論中的平行性和差異。
\item 級數展開的算術基礎與理想理論的關係。
\item 代數函數的算術理論與幾何理論的比較。不同的不變性概念。
\item 比較的繼續。奇異點。
\item 剩餘定理與理想理論諸定理之間的比較。
\item 代數函數和代數數中的超越問題。
\end{enumerate}

% unit: noether.p14.sec.001
\subsection*{1. 各種理論的基礎}

% unit: noether.p14.par.009
Riemann 的理論是唯一一個純粹超越的理論；它建立在存在定理之上（Dirichlet 原理，Schwarz--Neumann 方法）。在這裡，代數函數的積分是第一位的，而代數函數本身是第二位的。Riemann 曲面上的函數由其零點和極點來刻畫，對應於算術理論中的多邊形（相應地，除子）；而 Weierstrass 也談論函數本身以及它們的零點和極點。除存在定理外，主要問題是建立所有具有給定極點的函數，這由第一類和第二類積分來完成；進一步還有一個問題：當無窮遠點被規定時，仍然進入的任意常數有多少。後一問題由 Riemann--Roch 定理回答；在算術理論中與之相應的是多邊形類（相應地，除子類）的維數問題，在代數幾何理論中則是完全線性系的維數問題，而此處答案同樣由 Riemann--Roch 定理給出。函數的實際構造，在算術理論中通過基定理完成；在代數幾何理論中通過剩餘定理完成。特別地，第一類被積函數的零點，在算術理論中對應於微分類，在代數幾何理論中對應於特殊群理論（也就是由 \((n-3)\) 次伴隨曲線，即 \(\varphi\)-曲線，截出的點群鉛筆）。

% unit: noether.p14.par.010
正如 Riemann 理論在歷史上先於其他理論，儘管存在定理的嚴格證明只是後來才得到；同樣，後來藉助它的超越工具，人們解決了一些代數問題，而這些問題直到今天還不能由純代數或算術處理來達到。這裡首先要提到 Hurwitz 的奇異對應理論。\footnote{A. Hurwitz, Über algebraische Korrespondenzen und das erweiterte Korrespondenzprinzip. \emph{Math. Ann.} vol. 28.（比較 Brill--Noether, section 10。）}

% unit: noether.p14.par.011
與代數數域理論中超越問題的類比，將在最後一節討論。

% unit: noether.p14.par.012
對於其餘理論，代數函數是第一位的，積分是第二位的。在這方面，Weierstrass 理論建立在冪級數展開和剩餘定理的超越基礎之上；但隨後它以代數方式進行，即明確指定所有被考察的函數。

% unit: noether.p14.par.013
Hensel--Landsberg 形式的算術理論也使用級數展開的超越基礎。但在這個基礎上給各個點賦予除子之後，它便純算術地進行，因此可以看作 Weierstrass 與 Dedekind--Weber 的融合。與後者相比，它同時給出實際構造所有有關基函數的方法。細節上，它比 Dedekind--Weber 有若干簡化，主要是通過引入“分數”除子。最近 Hensel\footnote{\emph{Mathematische Zeitschrift} vol. 4；這個基礎也被用作他的報告的基礎。}給出了級數展開理論的算術基礎（除去收斂性證明中的一個優函數方法），由此該理論幾乎成為純算術的。

% unit: noether.p14.par.014
Brill--Noether 的代數幾何理論只預設一個分支的“後繼點”概念；這個概念被歸結為普通點處的級數展開，\footnote{關於“奇異”的後繼點，比較 Brill--Noether, section 6，特別是 nos. 10 和 11。普通後繼點例如是切線與曲線的交點，或更高階接觸點（分支點）。}並對應於 Weierstrass 的元素概念；此後它以代數有理方法繼續，實質上是三元消去理論的方法；它也達到顯式表示。

% unit: noether.p14.par.015
Dedekind--Weber 的原始純算術理論\footnote{\emph{J. f. M.} 92（引用為 D.--W.）。}與 Riemann 理論相近，因為它的定理又具有存在證明的性質；它保持了方法上的完全純粹性。\footnote{因此，與其他理論相比，它也比 Hensel--Landsberg 理論更適於比較；後者在下面的比較中只會偶爾出現。--- Abel 積分及其反演理論，在它預設連續性的範圍內，並沒有在 D.--W. 中處理，也沒有在代數理論的進一步發展中處理（M. Noether, \emph{Ann.} 37）。}該理論的算術基礎，其中變量只起不定元的作用，與代數數理論平行，如將在第 2 節中更詳細說明。\footnote{Hensel 中有 Hensel--Landsberg 理論與數論，尤其與 \(p\)-進數理論的比較。}這個基礎給出以算術方式引入點概念的方法（比較 Hensel no. 10a），並因而在不考慮連續性的情況下建立絕對 Riemann 曲面的概念。該理論由此獲得形式性質；例如，微分的概念也以純形式方式引入。在由點概念獲得的基礎上，可以證明它與代數幾何理論的方法和概念有密切親緣關係，因為後者除了元素概念外，也不含超越輔助工具；這將在以下各節中進行。

% unit: noether.p14.sec.002
\subsection*{2. 代數數和代數函數之間的平行性（共軛元素、基定理、模概念）}

% unit: noether.p14.par.016
一個代數數域被定義為有理數域 \(R\) 上次數為 \(n\) 的域；一變量代數函數域則被定義為 \(K(z)\) 上次數為 \(n\) 的域，其中 \(K(z)\) 是係數為任意複數的一個不定元 \(z\) 的所有有理函數所成的域。因此，對於這兩種域，一切完全抽去基域特殊性質的定理共同成立；這些是關於基、範數、跡和判別式的定理。\footnote{D.--W., § 1 和 2；Hensel no. 4（帶注 5）。}所有這些概念由 Dedekind--Weber 在不使用共軛元素的情況下引入。然而應當注意，共軛域的概念可以完全形式地定義，正如 Steinitz 繼 Kronecker 之後所說明的那樣。\footnote{Steinitz, Algebraische Theorie der Körper. (\emph{J. f. M.} 137, § 6 和 8。)}也就是說，若 \(K\) 是任意域，\(\varphi(x)\) 是在 \(K\) 上不可約的 \(x\) 中多項式，則模 \(\varphi(x)\) 的剩餘類系統形成一個域。如果把 \(x\) 的剩餘類指定為一個元素 \(\xi\)，那麼在同構指定下，元素 \(f(\xi)\) 對應於由多項式 \(f(x)\) 表示的剩餘類。由於所有能被 \(\varphi(x)\) 整除的多項式都表示零類，\(\varphi(\xi)\) 被指定給零類；元素 \(\xi\) 因而可以符號地看作 \(\varphi(x)=0\) 的一個零點，並生成代數域 \(K(\xi)\)。重複這個過程得到 \(n\) 個共軛域 \(K(j), K(j'),\ldots,K(j^{(n-1)})\)。因此，人們也可以特別地在代數函數中談論共軛元素，而不離開純算術理論。

% unit: noether.p14.par.017
兩個基域 \(R\) 和 \(K(z)\) 共有這樣一個性質：在其中，所有整量的系統分別被定義為有理整數和一個不定元的整有理函數。這些整量具有如下特徵性質：任意兩個 \(a\) 和 \(b\) 都有最大公約數 \(d\)，它可表示為 \(ma+nb\) 的形式，其中 \(d,m,n\) 也都是來自 \(R\) 或 \(K(z)\) 的整量；並且 \(d\) 正是可由這種形式表示的絕對最小的數，相應地是最低次數的多項式。正是僅在這一性質之上，藉助它對有理整數和一個不定元函數推出唯一分解為素因子，才建立了關於代數域或上域中整量的諸定理。證明最大公約數存在的論證，也給出代數域中整量的基本系統或基的存在性。若只考慮由整量組成的域的基，則其判別式成為一個有理整數，相應地成為一個不定元的有理函數。每一個對應於絕對最小數、相應地對應於最低次數函數的基，都形成一個基本系統。

% unit: noether.p14.par.018
對基問題的更精確認識，在更一般的域中也是如此，是由模概念的引入給出的。這個概念應當這樣表述，使它包括文獻中依元素性質不同而出現的各種定義。相對於一個基礎整域 \(\mathfrak J\) 而言，一個模是一個元素系統，其中任意兩個元素的差仍屬於該系統，並且一個元素與來自 \(\mathfrak J\) 的任意量的乘積也仍屬於該系統。（若 \(\mathfrak J\) 由有理整數構成，則第二個要求是第一個要求的結果。）\footnote{模概念歸功於 Dedekind，他首先在數論第二版中引入數模（\(\mathfrak J=\) 有理整數）；整線性形式的模也隱含地包含在那裡。在 Dedekind--Weber 中出現“函數模”，其元素由代數函數構成，而 \(\mathfrak J\) 由一個不定元的所有整有理函數構成。--- \(n\) 個不定元中的齊次型的模，更一般地說多項式的模，首先出現在 Hilbert（\emph{Annalen} 36）中；\(\mathfrak J\) 由所有多項式構成，並且由於元素本身也是多項式，這個模落入我們的理想定義之下。Kronecker 的多項式“模系統”從基出發，而不是從抽象定義出發。}

% unit: noether.p14.par.019
凡是具有有限個元素的基、並且每個元素都能用 \(\mathfrak J\) 中的係數通過該基線性表示的模，稱為有限的（有限生成的）；特別地，每一個一項模於是具有形式 \(c\cdot\alpha\)，其中 \(\alpha\) 表示基元素，而 \(c\) 遍歷 \(\mathfrak J\) 中的所有量。\emph{若一個模的元素本身屬於 \(\mathfrak J\)，則稱它為 \(\mathfrak J\) 中的一個理想；}由此得到有限理想的概念，特別是一項理想（主理想）的概念。

% unit: noether.p14.par.020
現在，關於整量和理想的所有基定理都建立在下列定理之上：

% unit: noether.p14.thm.001
\emph{若 \(M\) 是相對於 \(\mathfrak J\) 的一個模，其元素由 \(n\) 個不定元中、係數在 \(\mathfrak J\) 中的線性形式組成，並且 \(\mathfrak J\) 中的每個理想都是有限的，則 \(M\) 也是有限模。若特別地，\(\mathfrak J\) 中每個理想都是主理想，則 \(M\) 具有由線性無關形式組成的一個基。}\footnote{該定理在文獻中針對不同的係數域 \(\mathfrak J\) 分別出現。關於有理整數，可比較例如 Hilbert, Zahlbericht § 3 和 4；關於若干不定元中的多項式：Hilbert, Über die vollen Invariantensysteme: § 2 末尾（\emph{Math. Ann.} 42）。}

% unit: noether.p14.par.021
實際上，如果 \(M\) 的元素由 \(l(x)=a_1x_1+\cdots+a_nx_n\) 給出，那麼係數 \(a_1\) 的全體遍歷 \(\mathfrak J\) 中的一個理想；按假設，它具有形式 \(b_1a^{(1)}+\cdots+b_\varrho a^{(\varrho)}\)。屬於 \(M\) 的線性形式
% unit: noether.p14.eq.001
\[
m(x)=l(x)-b_1l^{(1)}(x)-\cdots-b_\varrho l^{(\varrho)}(x),
\]
% unit: noether.p14.par.022
因此只依賴於 \((n-1)\) 個不定元；經過有限次重複，斷言隨之得出。若每一次都有 \(\varrho=1\)，那麼該基由分別含 \(n,\ n-1\ldots1\) 個不定元的 \(n\) 個形式組成，這給出非零形式的線性無關性。

% unit: noether.p14.par.023
關於基本系統的定理，在各種情況下由此如下得出：

% unit: noether.p14.par.024
如果代數域由把整代數元素 \(\delta\) 添加到基域而得到，並且 \(\mathfrak J\) 表示基域中所有整量的全體，那麼上域中的每個整量都容許表示為（比較例如 Zahlbericht § 3）
% unit: noether.p14.eq.002
\[
\omega=\frac{a_1\delta^{\,n-1}+a_2\delta^{\,n-2}\cdots+a_n}{\Delta(\delta)},
\]
% unit: noether.p14.par.025
其中 \(a_i\) 是來自 \(\mathfrak J\) 的量，而 \(\Delta(\delta)\) 表示 \(\delta\) 的判別式。通過代換
% unit: noether.p14.eq.003
\[
x_i=\frac{\delta^{\,n-i}}{\Delta(\delta)}
\]
% unit: noether.p14.par.026
所有整量的模被變換為線性形式的模；因而在 \(\mathfrak J\) 中每個理想都是有限的所有情形中，基本系統的存在性隨之成立。同一論證更一般地說明，上域整量的域中的每個理想都是有限的。現在特別地，對於有理整數以及對於一個不定元的函數，每個理想都是主理想；於是得到：在代數數域和一變量代數函數域中，當 \(n\) 表示次數時，基本系統恰由 \(n\) 個量組成。並且由於在這些域中，每個理想按上面的論證都是有限的，於是同時得到相對域的基本系統；一般來說，它將由比相對次數所指示的更多的量組成。\footnote{按照 Steinitz 對線性形式模的研究，其中 \(\mathfrak J\) 由一個（有限）代數數域的整數構成，當 \(\nu\) 表示相對次數時，\(\nu+1\) 個基元素已經足夠（\emph{Math. Ann.} 71, § 4）。}此外，如果 \(\mathfrak J\) 表示若干不定元中所有多項式的域，那麼 \(\mathfrak J\) 中每個理想也是有限的；\footnote{由 Hilbert 的模基定理（\emph{Math. Ann.} 36）。這裡為了一致性而稱為理想的多項式模，通常被稱作“型模”或簡單稱作“模”。Hilbert 關於這種模 \(N\) 是有限的定理，也可歸約為關於線性形式的定理。事實上，設 \(f\) 是包含在 \(N\) 中、關於 \(r+1\) 個不定元 \(x_1,\ldots,x_r,x_{r+1}\) 的次數為 \(n\) 的多項式，並且含有項 \(x_{r+1}^{\,n}\)（這總可通過線性變換達到）；那麼 \(N\) 中的所有多項式都模 \(f\) 同餘於可表示為 \(a_1(x)x_{r+1}^{\,n-1}+\ldots a_n(x)\) 形式的多項式；於是通過令 \(x_{r+1}^{\,n-i}=X_i\)，它們形成線性形式的模，而對於這個模，\(r\) 個不定元的域中的每個理想可以假定為有限，因為一個不定元時情形如此。}因而，基本系統的存在性也為若干不定元的代數函數域得到，並且它又將由比次數所指示的更多的量組成。

% unit: noether.p14.sec.003
\subsection*{3. 代數數和代數函數的理想理論中的平行性和差異}

% unit: noether.p14.par.027
理想理論基本定理的證明，即每個理想都可以唯一地表示為素理想冪的乘積，可以對代數數和代數函數共同進行；它只基於這樣一個性質：對於有理整數和一個不定元的函數，每個理想都是主理想。\footnote{比較 Hurwitz: Über die Theorie der Ideale (\emph{Gött. Nachr.} 1894) 引言中的一個說明。在 Weber 的代數學教科書中，證明按照 Kronecker 的方法，通過引入泛函，對兩個域平行地進行。[Kronecker 理論和理想理論之間的聯繫由“擴展的 Gauss 定理”來中介；Hurwitz, loc. cit.]}證明的要點在於表明：由一個理想 \(\frc\) 被一個理想 \(\fra\) 整除（也就是說，\(\frc\) 的每個元素都包含在 \(\fra\) 中），可推出乘積表示 \(\frc=\fra\cdot\frb\)。\footnote{Dedekind（數論）特別強調了這一點；補篇 XI，定理 VII，§ 177；比較 p. 554 的注。該定理利用的是所涉對象為代數域這一事實；對於多項式的理想（模）它不再成立，那裡最小公倍式取代乘積。}在 Hurwitz 那裡，這個證明依賴於“擴展的 Gauss 定理”，該定理說：如果一個整量 \(\omega\) 整除兩個多項式 \(\varphi\) 與 \(\psi\) 的乘積的每個係數，那麼它也整除 \(\varphi\) 和 \(\psi\) 的所有係數的乘積；在 Dedekind 那裡，它依賴於一個與該定理等價的模定理。\footnote{數論，補篇 XI；定理 VI，§ 173。Dedekind 指出了這兩個定理的等價性（Über die Begründung der Idealtheorie, \emph{Göttinger Nachr.} 1895）。}

% unit: noether.p14.par.028
進一步，補模的概念（Hensel no. 12），以及由此得到的分歧理想（基本理想，不同）\(\mathfrak z\)，可以共同定義（它的範數成為域判別式 \(D\)）；並且基本定理成立：分歧理想 \(\mathfrak z\) 是\emph{所有主理想 \(F'(\delta)\) 的最大公約數}；這裡 \(\delta\) 遍歷該域的所有整量，而 \(F(\delta)=0\) 每次表示相應的不可約方程。由此得到：\emph{域判別式包含所有且僅包含那些能被一個素理想的平方整除的有理素數，相應地 \(z\) 中的線性因子}。對於 \(F'(\delta)\) 本身得到表示
% unit: noether.p14.eq.004
\[
F'(\delta)=\frf\cdot\frD,
\]
% unit: noether.p14.par.029
並取範數得
% unit: noether.p14.eq.005
\[
\Delta(\delta)=R^2\cdot D;
\]
% unit: noether.p14.par.030
這裡 \(\frf\) 表示由 \(\delta\) 導出的環（序，整域）的導子，也就是說，\(\frf\) 由所有這樣的（整）量組成：它們與任意整量相乘之後，都能被模 \([1,\delta,\ldots\delta^{n-1}]\) 整除。\(R^2=\operatorname{Norm}\frf\) 表示把基 \(1,\delta\ldots\delta^{n-1}\) 變到一個整量基的代換行列式的平方。\footnote{Dedekind, Über die Diskriminante endlicher Körper (\emph{Abhandl. der Gött. Ges. d. Wissensch.} vol. 29, 1882)。那裡為代數數給出的定義可以直接轉移到代數函數；證明也幾乎可以完全轉移。D.--W. 中的證明順序不同：見下文。關於所引定理，也比較 Zahlbericht, § 31 和 32。（\(F'(\delta)\) 在那裡被稱為整數 \(\delta\) 的“不同”。）}

% unit: noether.p14.par.031
進一步發展中的\emph{差異}由基域的特殊性質所決定。因此，在代數數的情形下，\emph{基域的元素同時也是指數}這一事實導致一種可能性：分歧理想可以被比 \((e-1)\) 次冪更高的冪整除，其中 \((e-1)\) 次冪來自在 \(p\) 中以 \(e\) 次冪出現的素理想 \(\frp\)（也就是當 \(e\) 被 \(p\) 整除時）；這又導致慣性域和分歧域的進一步概念，\footnote{比較 Zahlbericht, § 39--45。}而這些在代數函數理論中沒有類似物。最重要的是，由於基域中\emph{存在無窮多個常數}，對於代數函數而言，理想類數的有限性消失了，與確定類數相關的一切問題也隨之消失。不過，在某種意義上，所有（超越的）素形式的域\footnote{Klein, Zur Theorie der Abelschen Funktionen (\emph{Ann.} 36, 1890)。}可以看作類域的一個超越類似物；因為在這裡每個點由一個單獨的形式，即一個主理想，生成。

% unit: noether.p14.par.032
另一方面，基域 \(K(z)\) 中存在無窮多個常數這一事實，對於代數函數導致了代數數中沒有類似物的定理和問題。首先，由此在證明安排上得到某種簡化；例如在 D.--W. 所給出的理想唯一分解為素理想的證明中，可以避免“擴展的 Gauss 定理”或與之等價的定理，而使用這樣一個事實：每個線性形式 \((z-c)\) 有無窮多個剩餘類，即所有常數（而模一個素數的剩餘類數是有限的）。\footnote{D.--W. § 9。比較 pp. 212 和 213 的注。}

% unit: noether.p14.par.033
進一步的事實是，所有以 \(z\) 的整有理函數為係數的多項式，模 \((z-\alpha)\) 都分解成線性因子（這對一個整多項式模 \(p\) 並非如此），它導致真正新的結果；即使不取所有常數而只取所有代數數，這個事實也仍然成立。\footnote{D.--W. 在引言中指出，在這種情況下整個理論仍保持完整。}由此可得，每個素理想都是一度理想；\footnote{D.--W., § 9, no. 7。}進一步，域判別式 \(D\) 成為所有判別式 \(\Delta(\delta)\) 的交，其中 \(\delta\) 遍歷所有整量；\footnote{D.--W., p. 224。}這兩點都與代數數相反。隨後又由此得出，分歧理想是所有主理想 \(F'(\delta)\) 的交。

% unit: noether.p14.par.034
進一步的差異還由如下事實決定：對於代數數，有理數域 \(R\) 這個基域是絕對特別的東西，即包含在該域中的素域（也就是說，由單位導出的域）；而基域 \(K(z)\) 則是相對的東西。任意非常數函數 \(\xi\) 都可以選作獨立變量；於是該域成為 \(K(\xi)\) 上的代數域。每個素理想 \(\frp\) 都是一度理想、因而每個函數模 \(\frp\) 都同餘於一個常數這一定理，與可把任意函數作為獨立變量的可能性結合起來，導致超出代數數的最重要概念，即純算術形式的點概念。（比較 Hensel 10a。）“每個點在這裡由一個與函數域同構的常數域來表示。”\footnote{眾所周知，若兩個域可以一一對應，並且和、差、積、商也相互對應，則稱它們“同構”。}這樣定義的點的全體形成“絕對 Riemann 曲面”。這個點定義只依賴於域，而不依賴於一個特殊變量。然而存在性證明是通過把某個函數區別為獨立變量 \(z\) 來進行的；一個點 \(P\) 由一個素理想 \(\frp\) 生成，方法是把每個函數指定給它模 \(\frp\) 的常數剩餘；因此特別地，能被 \(\frp\) 整除的函數在 \(P\) 處消失，反之亦然。當 \(\frp\) 遍歷 \(z\) 中的所有素理想時，由此得到所有 \(z\) 有限的點；其餘點通過引入獨立變量 \(\zeta=\frac{1}{z}\) 而出現，並對應於整除主理想 \(\xi\) 的素理想。自然地，附屬於點概念的概念也不能在代數數中有類似物；例如多邊形、多邊形類、多邊形類的維數等概念。

% unit: noether.p14.par.035
還應注意，獨立變量的自由選擇導致對分解 \(F'(\vartheta)=\frf\cdot\mathfrak z\) 的另一種解釋。也就是說，如果把 \(\vartheta\) 看作獨立變量，那麼相應地得到
% unit: noether.p14.eq.006
\[
\left(\frac{\p F}{\p z}\right)=\frf\cdot\fra,
\]
% unit: noether.p14.par.036
因為由 \(\vartheta\) 導出的環是相對於 \(z\) 和 \(\vartheta\) 對稱定義的；於是 \(\frf\) 成為兩個主理想 \(\frac{\p F}{\p z}\) 和 \(\frac{\p F}{\p\vartheta}\) 的最大公約數；導子 \(\frf\) 同時成為相對於 \(\vartheta,z\) 的“雙重點理想”。\footnote{D.--W. p. 263。}

% unit: noether.p14.sec.004
\subsection*{4. 級數展開的算術基礎與理想理論的關係}

% unit: noether.p14.par.037
在 Hensel--Landsberg 中，級數展開的超越輔助工具以及從函數論借來的點概念取代理想理論的位置。Hensel 最近為代數數和代數函數的級數展開給出的算術基礎\footnote{Eine neue Theorie der algebraischen Zahlen (\emph{Math. Zeitschr.} vol. 2, 1918) 和 Neue Begründung der arithmetischen Theorie der algebraischen Funktionen einer Variablen (\emph{Math. Zeitschr.} vol. 4, 1919)。更詳細的說明可比較 Hensel no. 44 及以下。}因此可以看作理想理論的等價物。事實上，這個算術基礎等於對每個素理想轉到這樣一個（超越）擴域中，在其中這個素理想成為主理想，從而通常的分解定理成立；並且只需對出現在一個素數 \(p\)，相應地一個線性函數 \(z-a\)，中的所有素理想實施這個考察。

% unit: noether.p14.par.038
首先引入一個超越擴域，它由按 \(z-a\) 的整數冪、相應地按 \(p\)（\(p\)-進數）的整數冪的所有級數展開組成，並且每次至多有有限多個負冪。在這個域中，定義代數數域或函數域的不可約方程分解為 \(\varrho\) 個不可約因子的乘積，對應於 \(p\)，相應地 \(z-a\)，分解為 \(\varrho\) 個“準素”理想 \(\mathfrak q_1,\ldots\mathfrak q_\varrho\) 的乘積。每個準素理想進一步表示為一個素理想的冪，
% unit: noether.p14.eq.007
\[
\mathfrak q_i=\frp_i^{\,l_i},
\]
% unit: noether.p14.par.039
這在 Hensel 那裡對應於引入另一個相對於超越擴域為代數的域。並且在代數函數的情形中，這個域可以由一個純方程定義；在代數數的情形中，更一般地，由一個代數可解方程定義。代數函數中的簡化來自這樣一個事實：這裡沒有慣性域和分歧域（比較 no. 3）。

% unit: noether.p14.par.040
在這些研究中出現的、關於一個素元素 \(\pi\) 的“終止級數”\footnote{Neue Begründung \(\ldots\), 公式 (11)。}（而且它還不需要超越的收斂證明）
% unit: noether.p14.eq.008
\[
v=c_0+c_1\pi+\ldots+c_{\sigma-1}\pi^{\sigma-1}-v_\sigma\pi^\sigma,
\]
% unit: noether.p14.par.041
其中 \(c_0\ldots c_{\sigma-1}\) 是常數，\(v_\sigma\) 表示該域的一個整元素，在 D.--W. 中以完全對應的形式出現。\footnote{D.--W., p. 231 和 § 15。}在那裡，它從理想理論導出，並構成定義一個函數在一點處零化階數（相應地變為無窮的階數）的基礎。
%
\clearpage
%
% unit: noether.p14.s05.title
\subsection*{5. 代數函數的算術理論與幾何理論的比較。不同的不變性概念}

% unit: noether.p14.s05.par.001
“不變量”是指任何僅由該域所刻畫的概念。在算術理論中，這種不變性一般已經由定義本身給出；定義是相對於域的所有元素而說的，並不是相對於某個基或某個被突出出來的變量而說的。上面給出的點概念就是一個例子。相反，其餘理論則從某個特殊變量或基出發，然後再證明該概念不依賴於這種特殊選擇；這裡的不變性變成了對於雙有理變換的不變性。

% unit: noether.p14.s05.par.002
也就是說，設該域一方面由把代數元素 \(s\) 添加到 \(K(z)\) 而得到，另一方面由把 \(\sigma\) 添加到 \(K(\zeta)\) 而得到，其中 \(f(z,s)=0\)，相應地 \(F(\zeta,\sigma)=0\)，分別表示 \(s\) 和 \(\sigma\) 的不可約方程；或者更一般地，由把有限多個元素添加到 \(K(\eta)\) 而得到，並具有相應的不可約方程
% unit: noether.p14.s05.eq.001
\[
g_1(\eta,\tau_1)=0,\qquad g_2(\eta,\tau_1,\tau_2)=0,\ldots .
\]
那麼，按定義，所有元素都既可以由 \(z\) 和 \(s\) 有理地表示，也可以由 \(\xi\) 和 \(\sigma\) 有理地表示，還可以由 \(\eta,\tau_1,\tau_2\ldots\) 有理地表示。特別地，\(z\) 和 \(s\) 因而可由 \(\xi\) 和 \(\sigma\) 有理地表示，反過來也一樣；同樣，\(z\) 和 \(s\) 可由 \(\eta,\tau_1,\tau_2\ldots\) 有理地表示，反過來也一樣。平面曲線 \(f(z,s)=0\) 因此通過“雙有理變換”變為平面曲線 \(F(\zeta,\sigma)=0\)，或者也變為變量 \(\eta,\tau_1,\tau_2\ldots\) 的空間中由 \(g_1(\eta,\tau_1)=0; g_2(\eta,\tau_1,\tau_2)=0,\ldots\) 所定義的曲線。不變性於是表現為：為一條特殊曲線所作出的定義，對於由它經雙有理變換得到的所有曲線（在任意多維空間中）都保持有效；也就是對於 Riemann 意義下同一個“類”的所有曲線都保持有效。這裡，“點”被定義為一個相互關聯的值系 \(z=a,s=b\)，因而 \(f(a,b)=0\)；在雙有理變換下，這個點一般又變為 \(F(\zeta,\sigma)=0\) 的一個相互關聯的值系 \(\alpha,\beta\)，但在特殊情形中，當 \(a,b\) 是一個“奇異”點時，它可以對應於 \(F(\zeta,\sigma)=0\) 上的若干“點”；並且總可以指出一些曲線，使得 \(f(z,s)=0\) 的一個“奇異”點對應於有限多個簡單點，這些簡單點可以彼此分離，也可以相鄰（奇點消解）。\footnote{比較例如 Brill--Noether，第 6 節。關於“奇異”點的更精確信息見下一節。}同樣，一個點群又變為一個點群；而且，一旦把一個奇異點按適當變換下與之對應的簡單點個數來計數，它的點數就是不變的。由於每條曲線都可以雙有理地變換成只具有“普通重”點的曲線，\footnotemark[\value{footnote}]幾何理論便只對這種特殊曲線發展其概念，然後再證明這些概念對於任何雙有理變換都是不變的，即使這種變換會導致最一般的奇異點。

% unit: noether.p14.s05.par.003
特別地，作為該“類”的一條特出曲線，除超橢圓情形外，可以引入 \((p-1)\) 維空間中的所謂“\(\varphi\) 的正規曲線”，它沒有奇異點；這裡 \(\varphi\) 是使 \(f(z,s)=0\) 齊次化後得到的那 \(p\) 條“\((n-3)\) 次伴隨曲線”。把 \(f(z,s)=0\) 雙有理變換為 \(F(\zeta,\sigma)=0\)，對應於 \(\varphi\) 的一個線性變換，所以這裡涉及的是射影意義下的不變性。把域的一切函數表示為 \(\varphi\) 的有理函數，在幾何理論中稱為“不變表示”。\footnote{比較例如 Brill--Noether，第 8 節。對於若干變量的代數函數，尚未處理不變性是否可歸結為對於線性變換的不變性這一問題。}

% unit: noether.p14.s05.par.004
因此，\(\varphi\) 的線性束，或者由它們所截出的“特殊群”的線性束，在每一個雙有理變換下都變為自身，因而僅由該域所刻畫。它對應於算術理論中的微分類；在那裡，不變性又直接顯現在定義中，而定義是由相對於域的每一個函數作為獨立變量的性質給出的。

% unit: noether.p14.s06.title
\subsection*{6. 比較的繼續。奇異點}

% unit: noether.p14.s06.par.001
由於點概念的理解不同，曲線的“奇異點”在代數幾何理論中造成特殊困難，卻在算術理論的奠基中根本不出現；因而“伴隨函數”也不像在幾何理論中那樣是構造算術理論所必需的，而是屬於理論中特殊的較高部分。另一方面，在幾何理論的奠基中，分歧點並不出現；事實上，在 Aronhold 正規形中表示積分時（Hensel no. 27，公式 80），它們也消失了。

% unit: noether.p14.s06.par.002
第 5 節給出的奇異點定義可以表述如下：\(f(z,s)=0\)（相應地，高維空間中的一條曲線）在 \(z=a,s=b\)（相應地 \(z=a; s_i=b_i\)）處有一個奇異點，如果這個值系由絕對 Riemann 曲面的多於一個點（算術意義下的點）取到；或者如果它在一個或若干點處以較高階取到。在後一種情形中，這是 \(z\) 的分歧多邊形以及 \(s\) 的分歧多邊形上的一個點；這對應於幾何理論中的“相繼點”；這個點成為 \(f(z,s)=0\) 的一個（高階）尖點（相應地，高維空間中曲線的尖點）。

% unit: noether.p14.s06.par.003
由於按假設，域的任意函數 \(\eta\) 都可以由 \(z\) 和 \(s\) 有理地表示，因此，由同構性可知，在 \(z=a,s=b\) 的絕對 Riemann 曲面上的每個點，都是這樣得到的：把每個由 \(z\) 和 \(s\) 表示的函數 \(\eta\)，指定為它在 \(z=a,s=b\) 時的值。因此，為了使這個點是奇異的，必須至少存在一個函數 \(\eta\)，它在 \(z=a,s=b\) 時取多個值，或者多次取同一個值系。而由於可由 \(z\) 和 \(s\) 有理表示，這只能在下列情形發生：在該表示中，分子和分母在 \(z=a,s=b\) 時同時消失。由此可知，上述定義與通常的定義相同，即
% unit: noether.p14.s06.eq.001
\[
f,\quad \frac{\p f}{\p z},\quad \frac{\p f}{\p s}
\]
在 \(z=a,s=b\) 時同時消失。因為在這種情形中
% unit: noether.p14.s06.eq.002
\[
\eta=\frac{\frac{\p f}{\p z}}{\frac{\p f}{\p s}}
\]
就是所說類型的函數，它在 \(z=a,s=b\) 時為多值；而在相反情形中，每個函數，即使成為 \(0/0\)，也只取一個值。

% unit: noether.p14.s06.par.004
在算術點概念下，奇異點不能出現，因為只要有一個函數被指定了不同的值系，兩個點就被看作不同。但是，即使在作為算術點概念基礎的理想理論中，奇異點概念也不會出現，這是因為人們總是同時考察該域所有整量的全體。設 \(\delta\) 關於 \(z\) 是代數整的，則 \(F(z,\delta)=0\) 在有限部分中的所有奇異點，都（按第二個定義以及第 3 節末尾）由“雙重點理想” \(\frf\) 生成。現在，由基本定理可知，分歧理想 \(\mathfrak z\) 是所有主理想 \(F'(\delta)\) 的最大公約數；於是對於每個素理想 \(\frp\)，都可以給出一個 \(\delta\)，使得相應的 \(\frf\) 不被 \(\frp\) 整除；從而，由於 \(\frf\) 是 \(F'(\delta)\) 和 \(F'(z)\) 的最大公約數，兩個主理想 \(F'(\delta)\) 與 \(F'(z)\) 中至少有一個不被 \(\frp\) 整除。因此，\emph{不存在一個值系，使所有 \(F'(\delta)\) 和 \(F'(z)\) 同時消失}；因而值系 \(z=a,\delta_i=b_i\) 只在絕對 Riemann 曲面的一個點處被取到；\emph{該域所有整代數函數的全體}（可把它解釋為空間中一條無限維曲線）\emph{在有限部分沒有奇異點}；而在理想理論中，考慮的又只有有限值。

% unit: noether.p14.s06.par.005
所有整量的基函數 \(\omega_1\ldots\omega_n\) 也同樣在有限部分沒有奇異點。按照上文，一個與所有整函數同構地對應的常數系，已經唯一地確定了算術意義下的每個點。因此，如果由任意整函數 \(s_i\) 定義的一條“空間曲線”在 \(s_i=\alpha_i\) 時在有限部分有一個奇異點，那麼必定至少有一個整函數 \(\eta\)，在 \(s_i=\alpha_i\) 時取多個值系。現在若 \(\omega_1,
\ldots,\omega_n\) 表示所有整量的一個基，則由於 \(\delta=a_1\omega_1\ldots+a_n\omega_n\)，其中 \(a_i\) 是有理整數，\(\omega\) 的每個值系只對應於 \(\delta\) 的一個值系；因此，由 \(\omega_1,
\ldots,\omega_n\) 定義的空間曲線在有限部分不能有奇異點；並且\emph{絕對 Riemann 曲面上使 \(z\) 為有限的每個點，已經由一個同構地指定給各 \(\omega\) 的常數系唯一地定義。基函數正是那些在 \(s_i=\alpha_i\) 時成為多值的函數 \(\eta\)。}因此，通過引入整量的基，例如代替一個基 \(1,\delta\ldots\delta^{n-1}\)，奇異點的困難就被繞開了。

% unit: noether.p14.s06.par.006
在幾何理論中，唯一地生成絕對 Riemann 曲面的所有點這一任務，更一般地說，構造在給定點處變為無窮大的所有函數這一任務，是藉助伴隨形式及其商來完成的。若 \(\psi=0\) 在基曲線 \(f(z,s)=0\) 的每個 \(i\) 重點處至少有一個 \((i-1)\) 重點，則形式 \(\psi\) 稱為伴隨的；較高的奇異點首先要被消解為普通多重點。上面已經說明，對於最一般的函數 \(\eta=\frac{\psi}{\chi}\)，在 \(f(z,s)=0\) 的一個奇異點處，分子和分母總要消失；至於它們必須是伴隨的，則由處處有限的微分來考察可知。反過來，伴隨條件對於構造最一般函數也是充分的，這由以後還要更詳細討論的“剩餘定理”說明。因此，伴隨函數代替了算術理論中的基函數。

% unit: noether.p14.s06.par.007
在算術上，一種伴隨形式對應於該域中一個函數的分子，這個函數可被“雙重點理想” \(\frf\) 整除；這裡假定奇異點全都位於有限部分，這總可以通過變換達到。由此可以形式地看出，基量 \(\omega_1,
\ldots,\omega_n\) 可以表示為伴隨形式的商，並且由此得到不同表示方式之間的關係，如下：

% unit: noether.p14.s06.par.008
令 \(R(z)\) 表示把一個基 \(1,\delta\ldots\delta^{n-1}\) 代換到一個基 \(\omega_1,
\ldots,\omega_n\) 中的代換行列式。反過來，每個 \(\omega\) 都可表示為 \(z\) 與 \(\delta\) 的一個整函數除以 \(R(z)\) 所得的商：
% unit: noether.p14.s06.eq.003
\[
\omega_i=\frac{G_i(z,\delta)}{R(z)}.
\]
因此，由於 \(R(z)\omega_i\) 屬於由 \(\delta\) 導出的環，按導子的定義，\(R(z)\) 可被 \(\frf\) 整除 [從 \((R(z))^2=\operatorname{Norm}\frf\) 只能推出 \((R(z))^2\) 的這種可整除性]。現在，由於 \(\omega_i\) 作為一個整函數，對於所有有限的 \(z\) 都保持有限，所以分子函數 \(G_i(z,\delta)\) 也必須可被 \(\frf\) 整除；因此所有 \(\omega_i\) 的分子和分母都是伴隨的。而由 \(\omega\) 的這個表示，又得到該域一切函數都可由伴隨形式的商來表示。

% unit: noether.p14.s06.par.009
按照前文，幾何理論可以刻畫為\emph{由一個量 \(\delta\) 導出的環（序）的理論}，這與同時考察\emph{所有整量}的算術理論相對。因此很清楚，環的導子——也就是奇異點——起着決定性的作用。還可以指出環（序）的算術理論與幾何理論之間的更多類比。例如，在與伴隨曲線相交時，幾何理論從來不把落入奇異點的交點計入，而只考察與這些點不同的點群。與此對應的是，Dedekind\footnote{Über die Anzahl der Idealklassen in den verschiedenen Ordnungen eines endlichen Körpers. (Festschrift Braunschweig 1877); § 3--5。（這些段落中為代數數陳述的定理，逐字適用於一變量代數函數。）}只考察那些與導子理想 \(\frf\) 互素的環理想，並且對於這個理想範圍建立所有整除律，包括唯一分解為素理想。

% unit: noether.p14.s06.par.010
作為幾何理論基礎的“代數函數基本定理”\footnote{M. Noether, Über einen Satz aus der Theorie der algebraischen Funktionen. \emph{Math. Ann.} vol. 6 (1873).}，或者至少它的一個直接推論，在某種意義上也在環的算術理論中有一個類似物。這個定理給出可表示性
% unit: noether.p14.s06.eq.004
\[
\psi(z,s)=A(z,s)f(z,s)+B(z,s)\varphi(z,s),
\]
的條件，其中所有出現的量都是 \(s,z\) 中的整有理函數（更一般地，是 \(x_1x_2x_3\) 中整且齊次的）。由這些條件可知，\footnote{Brill--Noether, no. 55.}由於 \(f(z,s)=0\)，當商 \(\frac{\psi}{\varphi}\) 相對於 \(f\) 是伴隨的時，它總等於一個整有理函數 \(B(z,s)\)。與此對應的是第 3 節的定理，它是所有這些比較的基礎：由 \(\delta\) 導出的環的導子 \(\frf\) 等於 \(f(z,\delta)=0\) 的雙重點理想。因為按導子的定義，任何可被 \(\frf\) 整除的代數整量都屬於該環，因而可以整地並且有理地由 \(z\) 與 \(\delta\) 表示；而所引定理則說，這樣的量是一個伴隨函數。\footnote{這裡依照 Dedekind 始終假定 \(\delta\) 是代數整量（這總可以通過變換達到）。與幾何更精確對應的情形 \(f(z,s)=0\)，其中 \(s\) 也可以代數分式地依賴於 \(z\)，由 Hensel--Landsberg 處理；比較 Hensel nos. 26 和 29。}
%
\clearpage
%
% unit: noether.p14.sec.007
\subsection*{7. 剩餘定理與理想理論諸定理之間的比較}

% unit: noether.p14.par.196
剛才提到的“代數函數基本定理”的推論，在幾何理論中直接導向“剩餘定理”，從而導向整個幾何理論賴以建立的基礎。然而，在某種意義上，剩餘定理又比其餘基礎更初等，因為它可以在算術上表述為理想唯一分解為素理想的直接推論，或者更確切地說，表述為這個分解定理所依據的事實的直接推論；因而它並不依賴於較高的環理論。人們如下達到這種算術表述。

% unit: noether.p14.par.197
假設基曲線 \(f(z,\delta)=0\)——這總能通過變量的分式線性變換達到——具有如下性質：奇異點和所考察的有限多個點群都位於有限部分；並且還假設 \(\delta\) 關於 \(z\) 代數整，這也總能通過隨後在變量中作整的線性變換來達到。於是，對應於生成絕對 Riemann 曲面上一點 \(P\) 的素理想 \(\frp\)，在此點處 \(z=a,\delta=b\)，有 \(f(z,\delta)=0\) 上的點 \(z=a,\delta=b\)；更一般地，一個理想 \(\fra=\frp_1^{\,e_1}\cdots\frp_\varrho^{\,e_\varrho}\) 生成由點 \(z=a_i,\delta=b_i\) 以及相應重數組成的點群；\(f(z,\delta)=0\) 上位於有限部分的所有點都以這種方式生成。由於在一點 \(P\) 處，所有可被 \(\frp\) 整除的函數 \(g_1,g_2\ldots\) 都消失，而反過來，這些函數 \(g\) 的全體生成理想 \(\frp\)，所以 \(f\) 上相應的點就是
% unit: noether.p14.eq.046
\[
g_1(z,\delta)=0,\quad g_2(z,\delta)=0,\ldots
\]
與 \(f=0\) 的交；對應於一個理想 \(\fra\) 的點群也是如此。事實上，按每個理想都是兩個主理想的最大公約數這一定理，總有兩條這樣的曲線足以截出一個點群。特別地，對應於一個主理想的點群由一條曲線 \(g(z,\delta)=0\) 截出，反過來也一樣；主理想對應於完全交。

% unit: noether.p14.par.198
在幾何理論中，兩個點群 \(A\) 和 \(B\) 稱為\emph{餘殘的}，如果它們可由同一個點群 \(R\) 補成一個完全交。如果這些點群分別由理想 \(\fra,\frb,\frr\) 生成，那麼這些理想之間有關係
% unit: noether.p14.eq.047
\[
\fra\cdot\frr=(\alpha),\qquad \frb\cdot\frr=(\beta),
\]
這表明理想 \(\fra\) 和 \(\frb\) 等價。反過來，從兩個理想的等價 \(\fra=\eta\frb\)，根據每個理想都可以通過乘以另一個理想而變為主理想這一定理，總能得到這個關係；\footnote{比較 Dedekind，數論 § 181；也就是說，若選擇 \(\frr\) 使 \(\fra\frr\) 成為主理想 \((\alpha)\)，則 \(\frb\frr=\eta\cdot(\alpha)\)；並且由於 \(\frb\frr\) 只含整量，\(\eta\cdot(\alpha)\) 也如此，所以它也成為一個主理想 \((\beta)\)。}\emph{等價理想和餘殘點群因此是同一概念。}

% unit: noether.p14.par.199
剩餘定理現在說：\emph{如果 \(A\) 和 \(B\) 是餘殘的，那麼它們可以由具有相同剩餘（在奇異點之外）的伴隨曲線截出。}由於一個函數的零點和極點總是餘殘的，這同時證明了該域最一般的函數可表示為伴隨形式的商。用理想理論的語言來說，剩餘定理簡單地歸結為：除了 \(\fra\) 和 \(\frb\) 之外，\(\fra\frf\) 和 \(\frb\frf\) 也等價，其中 \(\frf\) 表示雙重點理想。因為如此一來，按上文，總存在一個理想 \(\frm\)，使得
% unit: noether.p14.eq.048
\[
\fra\frf\frm=(\alpha),\qquad \frb\frf\frm=(\beta)
\]
成為主理想；\(\alpha=0\) 和 \(\beta=0\) 就是分別截出點群 \(A\) 和 \(B\)，或者說截出它們不同於奇異點的子群 \(A'\) 和 \(B'\) 的伴隨曲線。\footnote{若 \(\fra=\frf'\fra_1,\ \frb=\frf'\frb_1,\ \frf=\frf'\frf_1\)，其中 \(\fra_1,\frb_1,\frf_1\) 兩兩互素，則已經有 \(\fra_1\frf\mathfrak n,\ \frb_1\frf\mathfrak n\) 成為主理想，其中 \(\mathfrak n\) 與 \(\frf\) 互素。因為每個理想都可通過乘以一個與給定理想互素的理想而變為主理想（D.--W. p. 214；或數論，§ 178，定理 XI）。}因此，剩餘定理的證明在算術上包含在把點指派給素理想的可能性中，也包含在如下證明中：等價理想可通過乘以同一個理想而變為主理想。

% unit: noether.p14.par.200
餘殘點群不必由同樣數目的點組成，正如等價理想不必具有同樣的次數。例如，所有主理想都是等價的，而由完全交生成的所有點群都是餘殘的。與等價多邊形（除子）\footnote{一個多邊形 \(A\) 由絕對 Riemann 曲面的有限多個點組成；兩個多邊形 \(A\) 和 \(B\) 稱為等價的，如果它們是某個函數 \(\eta\) 的零點和極點；此時 \(\eta\) 允許用多邊形商表示為 \(\eta=\frac{A}{B}\)。（D.--W. § 17 和 18。）}相比的這種差異，可由如下事實解釋：那些使 \(z\) 變為無窮的點，並不由 \(z\) 中的素理想生成。在把函數表示為理想商 \(\eta=\frac{\fra}{\frb}\) 時，這種表示陳述了 \(\fra\) 與 \(\frb\) 的等價，因此 \(\eta\) 的那些 \(z\) 變為無窮處的零點或極點並不出現。比如主理想 \((\alpha)\) 成為對應於整函數 \(\alpha\) 的零點的那些素理想的乘積；\(\alpha\) 的極點不出現，因為 \(\alpha\) 只在 \(z\) 變為無窮的點處變為無窮。在這個意義上，理想理論正是幾何理論中形式理論的精確類似物，在後者中也只有零點而沒有極點。

% unit: noether.p14.par.201
從形式到函數的過渡，在幾何上通過取同階形式的商來實現；在算術上，這對應於從理想類過渡到不變地定義的多邊形（除子）類。因為這裡不再有任何變量被區別出來，函數由多邊形商表示時便給出所有零點和極點。理想類（餘殘點群的類）與多邊形類的關係如下：若 \(\calA\) 和 \(\calB\) 是由理想 \(\fra\) 和 \(\frb\) 生成的多邊形，\(\calN\) 是 \(z\) 變為無窮的點的多邊形，即 \(z\) 和 \(z\) 的整函數的分母多邊形，那麼由 \(\fra\) 與 \(\frb\) 的等價推出 \(\calA\) 與 \(\calB\) 的等價，當且僅當 \(\fra\) 與 \(\frb\) 具有同樣的次數；一般地，只推出 \(\calA\calN^\nu\) 與 \(\calB\calN^{\nu'}\) 的等價，\((\nu\ne\nu')\)。因此，把理想（相應地餘殘點群）劃分為類，就是把無窮多個多邊形類合併為一個類；也可以把它理解為按 \(\calN\) 的一個冪取模來劃分多邊形。\footnote{Hensel--Landsberg，第 25 講，p. 438。--- 若把兩個變量 \(z\) 和 \(s\) 表示為多邊形商，就可以把這理解為分別對這些變量作二元齊次化；如果特別選擇 \(z\) 和 \(s\)，使它們具有同一個分母多邊形，那麼就得到幾何中通常的三元齊次化。}

% unit: noether.p14.par.202
對於由不同數目的點組成的餘殘點群，剩餘定理只出現在純幾何問題中。在應用於代數函數時，事實上只出現對應於多邊形類的情形，也就是同數目點的餘殘點群情形。由一個點群生成的“完全線性系”，即與這個點群餘殘且具有同樣點數的點群全體，正好對應於由相應多邊形生成的多邊形類；“兩個完全系之和”和“兩個多邊形類之積”這些概念也相互一致。這樣的完全系只含有限多個線性無關的點群，是剩餘定理的直接後果；該定理把這個系的維數，即線性無關點群的個數，化歸為所有通過剩餘群、任意但固定階數的伴隨曲線的維數。為了達到對這個維數的實際確定，即 Riemann--Roch 定理，幾何理論首先從剩餘定理推出“約化定理”（關於固定點的定理）；\footnote{Brill--Noether, no. 58.}再由此推出 Riemann--Roch 定理。

% unit: noether.p14.par.203
完全相應地，算術理論表明每一個多邊形類都是有限維的，並通過整量的一種特殊基，即由補模定義的正規基，來確定其維數；因此它把 Riemann--Roch 定理作為補模與分歧多邊形之間聯繫的結果而得到。在 D.--W. 的表述中，也出現了約化定理，\footnote{D.--W., § 30, no. 2.}不過並不像在幾何理論中那樣作為基礎。由於第 6 節所提到的基與伴隨函數之間的對應，在這裡也可以確立兩種理論中輔助工具的某種對應。概括地說，兩種理論中的概念形成在本質上是相合的，儘管除表述外，它們在安排上有所不同。這兩種理論彼此獨立地產生，幾何理論早了十年。還應指出，Riemann--Roch 定理最初在 Riemann 和 Roch 那裡，也在 Weierstrass 那裡，是作為秩定理出現的；所涉及的是矩陣 \((\varphi_i(x_\varkappa))\) 的秩，其中 \(\varphi\) 遍歷 \(p\) 條 \(n-3\) 次伴隨曲線，而 \(x\) 遍歷剩餘群的點。

% unit: noether.p14.sec.008
\subsection*{8. 代數函數和代數數中的超越問題}

% unit: noether.p14.par.204
Hilbert 特別指出了代數函數和代數數中的超越問題之間的類似；\footnote{Mathematische Probleme, problem 12.（\emph{Göttinger Nachrichten} 1900）。}這些本質上是存在性問題。因此，Riemann 關於給定 Riemann 曲面（因而也給定分支點）時代數函數存在性的問題，對應於具有給定群和判別式的代數數域存在性問題。\footnote{由此提示的、從群出發處理代數函數的方法，由 R. König 作為更一般問題的特殊情形完成；特別比較：Riemannsche Funktionen- und Differentialsysteme in der Ebene（\emph{J. f. M.} vol. 148）以及 \emph{Math. Ann.} 78, 79。在 König 那裡，取代代數域的是一個相對於 \(K(z)\) 的特殊有限模：這裡只有“類性質”。}在代數函數的情況下，存在性證明可以一般地進行；而在代數數域範圍內，它只對特殊基域（有理數、二次數域）並且只對 Abel 群成功。\footnote{只有三次和四次方程中出現的群的簡單情形，以及八次方程中的若干極特殊群，能對任意基域處理：G. Bucht, Über einige algebraische Körper 8. Grades（\emph{Arkiv f. Math., Astron. och Physik}, vol. 6, no. 30）。也比較 E. Noether, Gleichungen mit vorgeschriebener Gruppe 以及其中引用的 Seidelmann 論文（\emph{Math. Ann.} 78）。}這裡涉及 Kronecker 定理，即每個有理數域上的 Abel 數域都可由單位根域組合而成，以及 Fueter\footnote{\emph{Math. Ann.} 75。}和 Hecke\footnote{\emph{Math. Ann.} 76。}關於相對 Abel 數域的相應定理。因此，所尋求的數域實際上由一種超越手段，即指數函數、模函數，提供出來；這與代數函數情形下的純存在性證明相反。

% unit: noether.p14.par.205
代數函數的存在性證明之前，是處處有限積分的存在性。與此類似的是，代數數域中相對於一個素數 \(\ell\) 定義的“奇異主數”（\emph{singuläre Primärzahlen}）的存在性；它們的 \(\ell\) 次根相對不分歧，並且提供類域的第一批構件。\footnote{Furtwängler, Reziprozitätsgesetze für Potenzreste mit Primzahlexponenten in algebraischen Zahlkörpern. \emph{Math. Ann.} 67 和 72。}這個存在性證明藉助一個定理來進行，而且這種幫助是本質的：在數域中總存在具有給定剩餘特徵的素理想；因此，這個定理可以看作邊值問題的類似物，而處處有限積分的存在性正是以邊值問題為基礎。

% unit: noether.p14.par.206
處處有限積分給出一種超越判準，用以判定兩個多邊形是否屬於同一類，也就是說，兩個同數目的點群是否餘殘；這是通過 Abel 定理給出的，該定理說，第一類積分沿“餘殘路徑”延拓時為零；或者也可以通過相應的微分定理及其逆定理給出。\footnote{例如比較 M. Noether（\emph{Ann.} 37）中的表述。由於剩餘定理在歷史上是由 Abel 定理發展出來的，幾何理論沿用積分和的說法，談論點群和線性系的“和”；這與算術理論中類的“積”相對。}

% unit: noether.p14.par.207
作為代數數域中 Abel 定理的類似物，必須取“\(\ell\) 次冪剩餘的互反律”；它對於該域本身，或至少對於一個適當的上域，給出兩個理想屬於同一類的判準。其內容也可以表述為：在這個上域中，奇異主數相對於同一類的所有理想具有相同的剩餘特徵。後一事實在該域本身中也成立，但在那裡它並不與互反律重合。\footnote{Furtwängler, loc. cit.}

% unit: noether.p14.par.208
最後，Hensel 的 \(p\)-進數理論提供了代數函數冪級數展開的類似物；關於這一點，請參看 Hensel 的報告。
%
\clearpage
%
% unit: noether.p15.title.001
\begin{center}
\pdfbookmark[0]{15. 二元型的整不變量系統的有限性}{paper-15}
{\Large\bfseries 15. 二元型的整不變量系統的有限性}\par
\vspace{1em}
\emph{Nachrichten von der Gesellschaft der Wissenschaften zu Göttingen} 1919, pp. 138--156
\end{center}

\vspace{2em}
\begin{center}
由 F. Klein 於 1919 年 3 月 27 日會議上提交。
\end{center}

% unit: noether.p15.par.001
以下作為已知有限性定理的推廣，證明如下定理：\emph{一個二元基本型系統的所有多項式整不變量，都是其中有限多個不變量的、具有整數係數的多項式有理函數。} 這裡照通常用法，“不變量”始終指一個基本型或一個基本型系統的多項式有理不變量，也就是相對於由所有線性變換構成的群所誘導的係數變換而言的不變量。

% unit: noether.p15.par.002
證明依賴於對 Mertens--Hilbert 有限性證明作出的整性強化。\footnote{F. Mertens, \emph{Wiener Sitzungsberichte}, 1889 年 1 月；D. Hilbert, \emph{Math. Ann.} vol. 33 (1889) [署期 1888 年 3 月]。這兩個證明是在 Mertens, Crelle vol. 100 之後彼此獨立產生的，在內容上相一致。} 這個有限性證明可作如下刻畫。把基本型分解成它們的線性因子；更準確地說，給每一個基本型配上一個線性型的乘積。這樣，每個不變量 $I$ 就變成這些線性型的不變量，從而變成齊次化根的函數。這個函數必須滿足三個條件：1) 不變性條件，2) 權條件，3) 對稱性條件。條件 2) 和 3) 對於 $I$ 在基本型的係數中成為多項式且有理是必要的。條件 1) 通過把 $I$ 表示為線性型的行列式的多項式有理函數來滿足；條件 2) 通過不取單個行列式，而取它們的某些冪的乘積作為新的自變量來滿足。條件 3) 隨後只是說 $I$ 成為若干列量的多項式有理對稱函數，並且反過來，每一個這樣的列量對稱函數都成為基本型的不變量。因此，這些列量的初等對稱函數給出完整系統。

% unit: noether.p15.par.003
如果只考慮整不變量，那麼條件 1) 可以像上面一樣滿足；只要對於線性型的不變量證明了更精確的定理，即\emph{二元線性型的所有整不變量，都是這些線性型的行列式的多項式整函數}。我在 \S\ 3 中證明事實確實如此，其形式是：這些行列式之間模任意素數的所有關係的全體，與所有恆等關係的全體相一致；也就是說，與這些關係相應的模，在模每一個素數之後仍為素模。這個驗證通過相對於由行列式之間的二次關係所對應的二次型模來規範化剩餘類而完成；這些二次型由此證明在模每一個素數時都構成該模的一個整基底，而此前在忽略整性時它們只被知道是一個模基。

% unit: noether.p15.par.004
為滿足條件 2)，不需要任何改變，因為引入冪的乘積並不影響整數係數。條件 3) 現在說 $n_1!\ldots n_\mu! I$（其中 $n_1,\ldots n_\mu$ 是各個基本型的次數）成為列量的整對稱函數；反過來，每一個這些列量的整對稱函數都成為整不變量。但是，容易看出，這些列量的整對稱函數具有一個整基底，而僅用初等函數已不再足夠。於是 Hilbert 關於整模基的定理的應用，也給出這些 $I$ 本身的整基底。

% unit: noether.p15.par.005
\S\ 1 彙集與條件 1)、2)、3) 相對應的特殊有限性定理，有限性證明將在 \S\ 2 中由這些定理推出。與條件 3) 相對應的特殊有限性定理還給出由整數或代數整數替換構成的有限群的整不變量的有限性，這在 \S\ 4 中作簡要說明。

% unit: noether.p15.sec.001
\section*{\S\ 1. 特殊有限性定理}

% unit: noether.p15.par.006
我首先陳述以下特殊有限性定理；它們將用於滿足引言中提到的不變性、權和對稱性條件。

% unit: noether.p15.thm.001
\textbf{I.} 二元線性型系統的每一個整不變量，都是這些線性型的行列式的、具有整數係數的多項式有理函數。證明將在 \S\ 3 中給出。

% unit: noether.p15.thm.002
\textbf{II.} 設 $\frS$ 是 $n$ 個變量的非負指數冪乘積系統；若 $\frS$ 與任意兩個冪乘積 $f$ 和 $g$ 同時也含有它們的商（只要這個商在變量中是多項式），則它有一個有限乘法基 $f_1,\ldots,f_k$，由 $\frS$ 中的函數組成，使得對於 $\frS$ 中每個 $f$ 都有表示
% unit: noether.p15.eq.001
\[
        f=f_1^{\lambda_1}\cdots f_k^{\lambda_k}
\]
其中指數 $\lambda_1,\ldots,\lambda_k$ 非負。

% unit: noether.p15.par.007
事實上，由 Hilbert 的模基定理，存在 $\frS$ 中有限多個函數，記為 $f_1,\ldots,f_k$，其中每一個確實含有 $x$，使得 $\frS$ 中每個非常數 $f$ 都屬於由它們生成的模：
% unit: noether.p15.eq.002
\[
        f\equiv0\pmod{(f_1,\ldots,f_k)}.
\]
由此，對於每個 $f=x_1^{i_1}\cdots x_n^{i_n}$，至少有一個 $f_\nu=x_1^{\nu_1}\cdots x_n^{\nu_n}$ 使得 $\nu_1\leqq i_1;\ldots\nu_n\leqq i_n$。在表示
% unit: noether.p15.eq.003
\[
        f=x_1^{i_1-\nu_1}\cdots x_n^{i_n-\nu_n} f_\nu = A f_\nu
\]
中，因假設，因子 $A$ 屬於 $\frS$，但它在變量中的次數低於 $f$。所以有限次重複這個論證即得結論。為了使表示也適用於 $f=1$，只須把所有指數 $\lambda$ 都置為零。\footnote{這個表示也可以直接證明；反過來 Hilbert 定理也可由此推出。參見 P. Gordan, \emph{Neuer Beweis des Hilbertschen Satzes über homogene Funktionen}, Gött. Nachr. 1899。另一個直接證明見 A. Ostrowski, \emph{Über die Existenz einer endlichen Basis bei gewissen Funktionensystemen}, \emph{Math. Ann.} 78 (1916), \S\ 2,1 和 \S\ 3,2；其中也推出了定理 II。\S\ 2 所用定理 II 的特殊情形，即指數遍歷一個丟番圖方程組的所有非負解的情形，已經見於 Gordan--Kerschensteiner, \emph{Vorlesungen über Invariantentheorie}, vol. I, p. 199。}

% unit: noether.p15.thm.003
\textbf{III.} $n$ 列量的所有多項式整對稱函數，都是其中有限多個函數的多項式整函數。

% unit: noether.p15.par.008
令 $\tau_1(x),\ldots,\tau_n(x)$ 表示不定元 $x_1,\ldots,x_n$ 的初等對稱函數。則對每一個指數 $k$ 都有恆等式
% unit: noether.p15.eq.004
\[
        x_i^k = G_0^{(k)}(\tau)+x_iG_1^{(k)}(\tau)+\cdots+x_i^{n-1}G_{n-1}^{(k)}(\tau),
        \qquad (i=1,2,\ldots,n),
\]
其中 $G_0^{(k)},\ldots,G_{n-1}^{(k)}$ 是 $\tau(x)$ 的多項式整函數。若 $x_i,y_i,\ldots,z_i$ 是第 $i$ 列的元素，則藉助這些恆等式以及對 $y,\ldots,z$ 的相應恆等式，所有特殊的一列對稱函數
% unit: noether.p15.eq.005
\[
        \sum_{i=1}^n x_i^{a}y_i^{b}\cdots z_i^{c}
\]
都成為以下有限多個函數
% unit: noether.p15.eq.006
\[
        \sum_{i=1}^n x_i^{a}y_i^{b}\cdots z_i^{c}
        \qquad
        (0\le a<n,\;0\le b<n,\;\ldots,\;0\le c<n)
\]
以及初等對稱函數 $\tau(x),\tau(y),\ldots,\tau(z)$ 的多項式整函數。\footnote{Hilbert 前引處把冪和而不是 $\tau(x),\ldots,\tau(z)$ 納入基中；這樣基只由一列函數組成，但表示的整性便喪失了。} 若有 $n$ 列最一般的對稱函數--這是下面不會出現的情形--同樣可證明，函數
% unit: noether.p15.eq.007
\[
        \sum x_{i_1}^{a_1}y_{i_1}^{b_1}\cdots z_{i_1}^{c_1}
              x_{i_2}^{a_2}y_{i_2}^{b_2}\cdots z_{i_2}^{c_2}
              \cdots
              x_{i_n}^{a_n}y_{i_n}^{b_n}\cdots z_{i_n}^{c_n}
\]
% unit: noether.p15.eq.008
\[
        (0\le a_\lambda<n,\;0\le b_\lambda<n,\;\ldots,\;0\le c_\lambda<n),
\]
其中 $i_1,i_2,\ldots,i_n$ 遍歷所有置換；這些函數與 $\tau(x),\tau(y),\ldots,\tau(z)$ 一起構成一個整基底。

% unit: noether.p15.thm.004
\textbf{IV.} 設 $F_1(x),\ldots,F_k(x)$ 是 $x_1,\ldots,x_n$ 的多項式整函數，並設 $a$ 是固定整數。那麼所有函數
% unit: noether.p15.eq.009
\[
        \frac{G(F)}{a},
\]
其中 $G$ 表示多項式整函數，並且這些函數在 $x$ 中成為整函數，它們都是其中有限多個函數的多項式整函數。

% unit: noether.p15.par.009
事實上，由 Hilbert 的整模基定理，對每個這樣的 $G(F)/a$ 都有表示
% unit: noether.p15.eq.010
\[
 \frac{G(F)}{a}
 =
 A_1(F_1\ldots F_k)\frac{G_1(F)}{a}
 +\cdots+
 A_\varrho(F_1\ldots F_k)\frac{G_\varrho(F)}{a},
\]
其中 $A_1,\ldots,A_\varrho$ 是 $F$ 的多項式整函數，並且
% unit: noether.p15.eq.011
\[
        \frac{G_1(F)}{a},\ldots,\frac{G_\varrho(F)}{a}
\]
屬於該系統。因此函數
% unit: noether.p15.eq.012
\[
        F_1=\frac{aF_1}{a},\ldots,F_k=\frac{aF_k}{a};
        \quad
        \frac{G_1(F)}{a},\ldots,\frac{G_\varrho(F)}{a}
\]
構成所求的基。

% unit: noether.p15.sec.002
\setcounter{equation}{0}
\section*{\S\ 2. 二元整不變量的有限性}

% unit: noether.p15.par.010
令 $x,y$ 為服從帶不定係數的變換的二元變量
% unit: noether.p15.eq.013
\begin{equation}
        x=s_{11}x'+s_{12}y';\qquad y=s_{21}x'+s_{22}y',
\end{equation}
並令 $f$ 為一個二元基本型，其係數是不定元。於是由 (1)，
% unit: noether.p15.eq.014
\begin{equation}
        f=a_0x^n+a_1x^{n-1}y+\cdots+a_ny^n
        =a'_0x'^n_n+a'_1x'^{\,n-1}y'+\cdots+a'_ny'^n.
\end{equation}
如所周知，$f$ 的一個整不變量是指一個多項式整函數 $I(a)$，滿足
% unit: noether.p15.eq.015
\begin{equation}
        I(a')=|s_{ik}|^\varrho I(a),\qquad \text{在 $a,s_{ik}$ 中恆等}.
\end{equation}
基本型系統的整不變量相應地定義；可以把相對於每一個基本型的係數的齊次性納入定義，因為其餘的不變量都由這些分別齊次的不變量按整線性方式組合而成。

% unit: noether.p15.par.011
除 (2) 外，再考慮一個特殊的二元型 $g$，即 $n$ 個線性型的乘積，其係數也仍是不定元：\footnote{如果像不變量論中常做的那樣用二項式係數書寫 $f$，即 $f=b_0x^n+\binom n1b_1x^{n-1}y+\cdots+b_ny^n$，那麼所有整不變量 $I(b)$ 的區域相對於 $I(a)$ 的區域擴大，而這裡使用的有限性論證失效；$I(b)$ 不再是根的整函數。}
% unit: noether.p15.eq.016
\begin{align}
        g&=(\alpha_1x+\beta_1y)(\alpha_2x+\beta_2y)\cdots(\alpha_nx+\beta_ny)\notag\\
         &=\sigma_0(\alpha,\beta)x^n+\sigma_1(\alpha,\beta)x^{n-1}y+\cdots\sigma_n(\alpha,\beta)y^n\notag\\
         &=(\alpha_1'x'+\beta_1'y')(\alpha_2'x'+\beta_2'y')\cdots(\alpha_n'x'+\beta_n'y')\notag\\
         &=\sigma_0'x'^n+\sigma_1'x'^{\,n-1}y'+\cdots\sigma_n'y'^n.
\end{align}
因此 $\sigma_0(\alpha,\beta),\sigma_1(\alpha,\beta),\ldots,\sigma_n(\alpha,\beta)$ 表示齊次化的初等對稱函數，而變換後的係數 $\sigma'_i$ 等於 $\sigma_i(\alpha',\beta')$。

% unit: noether.p15.par.012
由於這些初等函數代數無關，每一個關係
% unit: noether.p15.eq.017
\begin{equation}
        G(\sigma_0,\sigma_1,\ldots,\sigma_n)=0
        \quad \text{在 $\alpha,\beta$ 中恆等}
\end{equation}
都對應於另一個關係
% unit: noether.p15.eq.018
\begin{equation}
        G(\sigma_0,\sigma_1,\ldots,\sigma_n)=0
        \quad \text{在 $\sigma_0,\ldots,\sigma_n$ 中恆等},
\end{equation}
從而對於不定元 $a_0,\ldots,a_n$，有
% unit: noether.p15.eq.019
\begin{equation}
        G(a_0,a_1,\ldots,a_n)=0
        \quad \text{在 $a_0,\ldots,a_n$ 中恆等}.
\end{equation}

% unit: noether.p15.par.013
現在，$f$ 的每個整不變量 $I(a)$，通過代換 $a_i=\sigma_i(\alpha,\beta)$，給出 $g$ 的一個整不變量 $I(\sigma)=H(\alpha,\beta)$，其中 $H$ 成為 $\alpha,\beta$ 的整函數。對這個不變量，關係 (3) 成為
% unit: noether.p15.eq.020
\begin{equation}
        I(\sigma')=|s_{ik}|^\varrho I(\sigma)
        \quad \text{在 $\sigma,s_{ik}$ 中恆等},
\end{equation}
而由於 $\sigma'=\sigma(\alpha',\beta')$，(8) 給出
% unit: noether.p15.eq.021
\begin{equation}
        H(\alpha',\beta')=|s_{ik}|^\varrho H(\alpha,\beta)
        \quad \text{在 $\alpha,\beta;s_{ik}$ 中恆等}.
\end{equation}
因此 $I(\sigma)=H(\alpha,\beta)$ 成為 $n$ 個線性型 $(\alpha_i x+\beta_i y)$ 的整不變量。反過來，設 $H(\alpha,\beta)$ 是這 $n$ 個線性型的整不變量，因而滿足 (9)，同時又等於 $\sigma$ 的一個整函數 $I(\sigma)$。由於這時
% unit: noether.p15.eq.022
\[
        H(\alpha',\beta)=I(\sigma(\alpha',\beta'))=I(\sigma'),
\]
方程 (9) 也可寫成
% unit: noether.p15.eq.023
\begin{equation*}
        I(\sigma')=|s_{ik}|^\varrho I(\sigma)
        \quad \text{在 $\alpha,\beta;s_{ik}$ 中恆等}.\tag{9a}
\end{equation*}
由 (5) 和 (6) 得出 (8) 成立，由 (7) 得出 (3) 成立。因此，每個整不變量 $I(a)$ 都對應於 $n$ 個線性型的整不變量 $H(\alpha,\beta)=I(\sigma)$，反之亦然。此外，若 $I_1(a),\ldots,I_k(a)$ 是所有 $I(\alpha)$ 的整基底，則 $H_1(\alpha,\beta),\ldots,H_k(\alpha,\beta)$ 是所有同時為整函數 $I(\sigma)$ 的 $H(\alpha,\beta)$ 的整基底；反過來這裡也由 (5)、(6)、(7) 成立。因此，要證明所有 $I(a)$ 的整基底存在，只須對所有同時為整函數 $I(\sigma)$ 的整不變量 $H(\alpha,\beta)$ 證明；而基
% unit: noether.p15.eq.024
\[
        H_1(\alpha,\beta)=I_1(\sigma),\ldots,H_k(\alpha,\beta)=I_k(\sigma)
\]
隨後給出 $I(a)$ 的基 $I_1(a),\ldots I_k(a)$。

% unit: noether.p15.par.014
令 $H(\alpha,\beta)=I(\sigma)$ 遍歷 (4) 的所有整不變量。由於由 (3) 每個不變量在 $\sigma$ 中是齊次的，而由 (4) 每個 $\sigma$ 在每個線性型的係數中是齊次且線性的，$H(\alpha,\beta)$ 在每一列 $\alpha_i\beta_i$ 中齊次，並且在每列中維數相同；此外它當然是 $n$ 列的對稱函數。反過來，由對稱函數基本定理，$n$ 列的每個整對稱函數 $H(\alpha,\beta)$，若分別齊次且在每列中維數相同，都是 $\sigma$ 的多項式整函數。因此，依據前面的論證，只要 $H(\alpha,\beta)$ 是不變量，它就是一個不變量 $I(\sigma)$。於是 (4) 的整不變量全體被刻畫為具有以下性質的整函數 $H(\alpha,\beta)$ 全體：
% unit: noether.p15.eq.025
\[
\begin{array}{ll}
1) & \text{$n$ 個線性型 $\alpha_i x+\beta_i y$ 的整不變量；}\\
2) & \text{在 $n$ 列 $\alpha_i,\beta_i$ 的每一列中分別齊次且維數相同；}\\
3) & \text{關於 $n$ 列 $\alpha_i,\beta_i$ 對稱。}
\end{array}
\]
由特殊有限性定理 I，$H(\alpha,\beta)$ 因 1) 成為行列式
% unit: noether.p15.eq.026
\[
        (ik)=\alpha_i\beta_k-\beta_i\alpha_k;
\]
的多項式整函數；於是
% unit: noether.p15.eq.027
\[
        H(\alpha,\beta)=G((ik)).
\]
在這個表示中，也可通過把 $n$ 列的所有置換作用於 $G((ik))$ 來形式地表示對稱性：
% unit: noether.p15.eq.028
\begin{equation}
        n!\,H(\alpha,\beta)
        =
        G^{(1)}((ik))+G^{(2)}((ik))+\cdots+G^{(n!)}((ik))
        =G^*((ik)).
\end{equation}
因此，伴隨 $(ik)$ 的每一個冪乘積 $P$，$G^*$ 也包含 $P$ 的所有置換之和 $Q=P^{(1)}+\cdots+P^{(n!)}$。事實上，$G^*$ 成為有限多個 $Q$ 的整線性組合，$G^*=\sum cQ$。由於這些 $Q$ 滿足條件 1)、2)、3)，因而本身成為不變量，故只須考慮 $Q$ 而非 $G^*((ik))$。

% unit: noether.p15.par.015
現在令 $P$ 遍歷所有出現在 $Q$ 中的 $(ik)$ 的冪乘積所成系統 $\frS$，也就是滿足條件 2) 的那些冪乘積。如果兩個 $P$ 的商在 $(ik)$ 中是多項式，那麼它同樣滿足條件 2)，因而屬於 $\frS$。所以由定理 II，存在有限多個這樣的 $P$，記為 $P_1,\ldots,P_\varrho$，使得每個 $P$ 都可表示為
% unit: noether.p15.eq.029
\[
        P=P_1^{\lambda_1}\cdots P_\varrho^{\lambda_\varrho}
\]
其中指數 $\lambda$ 非負。作用所有置換得到
% unit: noether.p15.eq.030
\[
        Q=
        P_1^{(1)\lambda_1}\cdots P_\varrho^{(1)\lambda_\varrho}
        +\cdots+
        P_1^{(n!)\lambda_1}\cdots P_\varrho^{(n!)\lambda_\varrho}.
\]
於是 $Q$ 成為 $n!$ 列量的整對稱函數，每列含有 $\varrho$ 個元素：
% unit: noether.p15.eq.031
\[
        P_1^{(1)}\ldots P_\varrho^{(1)};
        \quad
        P_1^{(2)}\ldots P_\varrho^{(2)};
        \quad \ldots;\quad
        P_1^{(n!)}\ldots P_\varrho^{(n!)};
\]
因此，由定理 III，它是這些列的有限多個對稱函數的多項式整函數。這些基函數本身就是 $Q$，或者是 $n!$ 個元素 $P_i^{(1)},P_i^{(2)},\ldots,P_i^{(n!)}$ 的初等對稱函數，因而滿足條件 1)、2)、3)，併成為整不變量
% unit: noether.p15.eq.032
\begin{equation}
        H_1(\alpha,\beta)=I_1(\sigma);\ldots;\quad
        H_l(\alpha,\beta)=I_l(\sigma).
\end{equation}
由 (10)，由於 $G^*=\sum cQ$，每個 $H(\alpha,\beta)=I(\sigma)$ 因此具有形式
% unit: noether.p15.eq.033
\[
        H(\alpha,\beta)=I(\sigma)=\frac1{n!}\Phi(I_1,\ldots,I_l),
\]
其中 $\Phi$ 表示多項式整函數；反過來，只要 $\frac1{n!}\Phi$ 在 $\alpha,\beta$ 中是多項式，它就是整不變量。因此由定理 IV，可以把有限多個進一步的不變量 $I_{l+1}(\sigma),\ldots,I_k(\sigma)$ 添入 (11)，使得
% unit: noether.p15.eq.034
\[
        H_1(\alpha,\beta)=I_1(\sigma);\ldots;H_l(\alpha,\beta)=I_l(\sigma);
\]
% unit: noether.p15.eq.035
\[
        H_{l+1}(\alpha,\beta)=I_{l+1}(\sigma);\ldots;H_k(\alpha,\beta)=I_k(\sigma)
\]
構成所有不變量 $H(\alpha,\beta)=I(\sigma)$ 的整基底。這就證明了一個基本型情形下的定理。

% unit: noether.p15.par.016
如果處理的是一個基本型系統，那麼如上所述，可以假設不變量在每一個基本型的係數中齊次。仍然把一般基本型分別替換為線性型的乘積。它們的不變量於是成為：1) 這些線性型的不變量；2) 在由線性型係數組成的列中分別齊次，並且在屬於每一個單獨基本型的列中維數相等；3) 關於屬於每一個單獨基本型的列對稱。反過來，所有滿足這些條件的線性型的每個整不變量，又對應於該基本型系統的一個整不變量。證明與剛才給出的證明完全平行。對於滿足權條件（定理 I 和 II）的各個行列式的每個冪乘積 $P=P_1^{\lambda_1}\cdots P_\varrho^{\lambda_\varrho}$，必須作用所有
% unit: noether.p15.eq.036
\[
        N=n_1!n_2!\cdots n_\mu!
\]
個置換，其中 $n_1,\ldots n_\mu$ 表示基本型的次數。這樣，$N\cdot I$ 被化為 $N$ 列量的整對稱函數，每列有 $\varrho$ 個元素。由此得到所有 $N\cdot I$ 的整基底，因而由定理 IV 得到所有 $I$ 的整基底。
%
%
% unit: noether.p15.sec.003
\setcounter{equation}{0}
\section*{\S\ 3. 二元線性型的整不變量的基}

% unit: noether.p15.par.017
現在還需要補上特殊有限性定理 I 的證明。根據一個已知定理，它最早由 Clebsch 證明，$n$ 個線性型 $\alpha_i x+\beta_i y$ 的所有多項式有理不變量，都是行列式 $(ik)$ 的多項式有理函數。因此，特別地，這些線性型的所有整不變量都成為這些行列式的多項式有理函數，雖然不一定是整函數。

% unit: noether.p15.par.018
下面將證明，由於 $(ik)$ 之間存在恆等關係，這種表示總是可以整地完成。

% unit: noether.p15.par.019
用公式並用模論的語言來說，這表述如下。令 $\xi_{ik}$, $i<k$, 為不定元。所有具有性質 $F((ik))$ 在 $\alpha,\beta$ 中恆等消失的整多項式 $F(\xi_{ik})$ 的全體形成一個模 $\frM$；兩個這樣的 $F$ 的和，以及一個這樣的 $F$ 與任意一個關於 $\xi_{ik}$ 的整多項式的乘積，仍屬於這個全體。公式地說：
% unit: noether.p15.eq.037
\[
        F((ik))=0\quad \text{在 $\alpha,\beta$ 中恆等}
\]
推出
% unit: noether.p15.eq.038
\[
        F(\xi_{ik})\equiv0 \bmod \frM\quad [\text{在 $\xi_{ik}$ 中恆等}],
\]
反之亦然。

% unit: noether.p15.par.020
我現在證明下列完全類似的命題。從
% unit: noether.p15.eq.039
\[
        F((ik))\equiv0\pmod p\quad \text{在 $\alpha,\beta$ 中恆等}
\]
% unit: noether.p15.eq.040
\begin{equation}
        \text{推出}\qquad
        F(\xi_{ik})\equiv0\pmod{(\frM,p)}
        \quad \text{在 $\xi_{ik}$ 中恆等},
\end{equation}
並且反之亦然，對於每個素數 $p$ 都如此。因此 $\frM$ 也與對應於 $(ik)$ 之間所有模 $p$ 關係的全體的模 $\frM_p$ 相同；這個模由所有整多項式 $F(\xi_{ik})$ 組成，使得 $F((ik))$ 在 $\alpha,\beta$ 中模 $p$ 恆等消失。

% unit: noether.p15.par.021
(1) 表示整可表示性，可以這樣看出。它也可寫成如下形式：若 $F((ik))=p\cdot H(\alpha,\beta)$ 在 $\alpha,\beta$ 中恆等成立，則
% unit: noether.p15.eq.041
\[
        F(\xi_{ik})-p\cdot G(\xi_{ik})\equiv0\pmod{\frM}
        \quad \text{在 $\xi_{ik}$ 中恆等},
\]
因而，由 $\frM$ 的定義，
% unit: noether.p15.eq.042
\[
        F((ik))=p\cdot G((ik))\quad [\text{在 $\alpha,\beta$ 中恆等}]\footnote{这个最后的结论同时证明了 (1) 的逆。}
\]
於是 $p\cdot H(\alpha,\beta)=p\cdot G((ik))$，或者
% unit: noether.p15.eq.043
\[
        H(\alpha,\beta)=G((ik)).
\]
因此，如果給出表示
% unit: noether.p15.eq.044
\[
        H(\alpha,\beta)
        =
        \frac1{p_1^{\lambda_1}\cdots p_\varrho^{\lambda_\varrho}}
        F((ik)),
\]
其中 $p_1,\ldots p_\varrho$ 是素數，那麼也得到
% unit: noether.p15.eq.045
\[
        H(\alpha,\beta)
        =
        \frac1{p_1^{\lambda_1-1}\cdots p_\varrho^{\lambda_\varrho}}
        G((ik)),
\]
有限次重複這個過程，就得到表示的整性。

% unit: noether.p15.par.022
為證明 (1)，記
% unit: noether.p15.eq.046
\[
        f_{ik,rs}=\xi_{ik}\xi_{rs}-\xi_{ir}\xi_{ks}+\xi_{is}\xi_{kr},
        \qquad i<k<r<s,
\]
為不定元 $\xi_{ik}$ 中對應於 $(ik)$ 之間二次關係
% unit: noether.p15.eq.047
\[
        (ik)(rs)-(ir)(ks)+(is)(kr)=0
\]
的型，並記 $\frN=(f_{ik,rs})$ 為由它們生成的整多項式模。由
% unit: noether.p15.eq.048
\[
        F(\xi_{ik})\equiv0\pmod{\frN}
\]
推出
% unit: noether.p15.eq.049
\[
        F((ik))=0\quad [\text{在 $\alpha,\beta$ 中恆等}],
\]
並且
% unit: noether.p15.eq.050
\begin{equation}
\begin{array}{rcl}
        \text{由 }F(\xi_{ik})&\equiv&0\pmod{(\frN,p)}\\
        \text{推出 }F((ik))&\equiv&0\pmod p\quad[\text{在 $\alpha,\beta$ 中恆等}].
\end{array}
\end{equation}
現在我將通過相對於 $\frN$ 規範化剩餘類來證明 (2) 的逆也成立。於是 $\frN=\frM=\frM_p$，從而 (1) 得證。等式 $\frN=\frM$ 作為 $p=0$ 時的特殊情形 $\frN=\frM_p$ 得到；在以下考察中總是包含 $p=0$ 的情形。換言之，$f_{ik,rs}$ 構成 $\frM_p$ 的一個整模基，特別也構成 $\frM_0=\frM$ 的整模基。

% unit: noether.p15.subsec.001
\subsection*{a) 兩行或三行 $\alpha_i\beta_i$ 的情形}

% unit: noether.p15.par.023
由於尚未出現四個不同的指標，$\frN$ 是零模。因此必須證明：由
% unit: noether.p15.eq.051
\begin{equation}
        F((ik))\equiv0\pmod p\quad \text{在 $\alpha,\beta$ 中恆等}
\end{equation}
推出
% unit: noether.p15.eq.052
\[
        F(\xi_{ik})\equiv0\pmod p\quad \text{在 $\xi_{ik}$ 中恆等}.
\]
這就以 $\frM=\frM_p=0$ 證明了 (1)。

% unit: noether.p15.par.024
由於條件 $F(ik)=p\cdot H(\alpha,\beta)$ 必須由各個行中齊次的組成部分分別滿足，所以始終只需假定 $F((ik))$ 在每一行中齊次，因而 $F(\xi_{ik})$ 在每個指標中齊次。在兩行的情形下，只有一個行列式 $(12)$ 和相應的不定元 $\xi_{12}$。由所假定的齊次性，
% unit: noether.p15.eq.053
\[
        F=c\cdot\,\xi_{12}^\nu.
\]
由
% unit: noether.p15.eq.054
\[
        c\cdot\,(12)^\varkappa\equiv0\pmod p\quad \text{在 $\alpha,\beta$ 中恆等}
\]
並且由於 $(12)\not\equiv0\pmod p$ 不在 $\alpha,\beta$ 中恆等為零，推出 $c\equiv0\pmod p$；故
% unit: noether.p15.eq.055
\[
        F(\xi_{ik})\equiv0\pmod p\quad [\text{在 $\xi_{ik}$ 中恆等}].
\]
這證明了 (3)，從而證明了 (1)。

% unit: noether.p15.par.025
在三行的情形下，有行列式 $(12),(13),(23)$，對應的不定元為 $\xi_{12},\xi_{13},\xi_{23}$。因此
% unit: noether.p15.eq.056
\[
        F=\sum c\,\xi_{12}^{\tau_{12}}\xi_{13}^{\tau_{13}}\xi_{23}^{\tau_{23}}.
\]
設 $F$ 在指標 $1,2,3$ 中的次數分別為 $\sigma_1,\sigma_2,\sigma_3$。於是
% unit: noether.p15.eq.057
\[
        \sigma_1=\tau_{12}+\tau_{13},\qquad
        \sigma_2=\tau_{12}+\tau_{23},\qquad
        \sigma_3=\tau_{13}+\tau_{23},
\]
其唯一逆變換為
% unit: noether.p15.eq.058
\[
        \tau_{12}=\frac12(\sigma_1+\sigma_2-\sigma_3),\quad
        \tau_{13}=\frac12(\sigma_1-\sigma_2+\sigma_3),\quad
        \tau_{23}=\frac12(-\sigma_1+\sigma_2+\sigma_3).
\]
因此每個分別齊次的 $F$ 至多含有一個項，
% unit: noether.p15.eq.059
\[
        F=c\,\xi_{12}^{\tau_{12}}\xi_{13}^{\tau_{13}}\xi_{23}^{\tau_{23}},
\]
由此如上得到 (3)，從而得到 (1)。

% unit: noether.p15.subsec.002
\subsection*{b) 四行 $\alpha_i,\beta_i$ 的情形}

% unit: noether.p15.par.026
出現六個不定元 $\xi_{12},\xi_{13},\xi_{14},\xi_{23},\xi_{24},\xi_{34}$ 及其相應行列式；並且 $\frN=(f_{12,34})$。由於
% unit: noether.p15.eq.060
\begin{equation}
        \xi_{14}\xi_{23}\equiv \xi_{13}\xi_{24}-\xi_{12}\xi_{34}\pmod{\frN},
\end{equation}
每一個 $\xi_{ik}$ 的冪乘積，模 $\frN$ 都同餘於若干冪乘積的整數線性組合，在這些冪乘積中不出現乘積 $\xi_{14}\xi_{23}$。於是對每個多項式 $F(\xi_{ik})$，有
% unit: noether.p15.eq.061
\begin{equation}
        F(\xi_{ik})\equiv F_0(\xi_{ik})\pmod{\frN},
\end{equation}
其中 $F_0$ 不含乘積 $\xi_{14}\xi_{23}$。多項式 $F_0$ 稱為 $F$ 的剩餘類的規範化多項式。由於 (4) 是分別齊次的，分別齊次的 $F$ 又對應於一個分別齊次的 $F_0$。

% unit: noether.p15.par.027
現在置
% unit: noether.p15.eq.062
\begin{equation}
        F_0(\xi_{ik})=G(\xi_{ik})+\xi_{34}\cdot F_1(\xi_{ik}),
\end{equation}
其中 $G$ 不含可被 $\xi_{34}$ 整除的冪乘積；並且再置
% unit: noether.p15.eq.063
\begin{equation}
        [G]_{\xi_{i4}=\xi_{i3}}=K(\xi_{12},\xi_{13},\xi_{23}).
\end{equation}
由
% unit: noether.p15.eq.064
\[
        F((ik))\equiv0\pmod p\quad \text{在 $\alpha,\beta$ 中恆等}
\]
按照 (5)，並顧及 (2)，推出
% unit: noether.p15.eq.065
\begin{equation}
        F_0((ik))\equiv0\pmod p\quad \text{在 $\alpha,\beta$ 中恆等}.
\end{equation}
通過特殊化 $\alpha_4=\alpha_3,\;\beta_4=\beta_3$，由 (6) 和 (7) 得到
% unit: noether.p15.eq.066
\[
        K(ik)\equiv0\pmod p\quad [\text{在 $\alpha,\beta$ 中恆等}].
\]
因此由 a) 中已證明的結果，
% unit: noether.p15.eq.067
\[
        K(\xi_{ik})\equiv0\pmod p\quad \text{在 $\xi_{ik}$ 中恆等}.
\]
由此推出下面還需證明的結論：
% unit: noether.p15.eq.068
\[
        G(\xi_{ik})\equiv0\pmod p\quad \text{在 $\xi_{ik}$ 中恆等}.
\]
於是 (6) 給出
% unit: noether.p15.eq.069
\[
        F_0(\xi_{ik})\equiv \xi_{34\cdot}F_1(\xi_{ik})\pmod p
        \quad \text{在 $\xi_{ik}$ 中恆等}.
\]
由於 $(34)\not\equiv0\pmod p$ 在 $\alpha,\beta$ 中不恆等為零，由 (8) 又推出
% unit: noether.p15.eq.070
\[
        F_1((ik))\equiv0\pmod p\quad [\text{在 $\alpha,\beta$ 中恆等}],
\]
其中 $F_1(\xi_{ik})$ 又是規範化的，並且在 $\xi_{34}$ 中次數較低。有限次重複這個過程，得到
% unit: noether.p15.eq.071
\[
        F_0(\xi_{ik})\equiv\xi_{34}^{\lambda}F_\lambda(\xi_{ik})\pmod p,
\]
其中 $F_\lambda$ 在 $\xi_{34}$ 中次數為零，因而由剛才證明的結果知其模 $p$ 消失。於是得到
% unit: noether.p15.eq.072
\begin{equation}
\begin{array}{rcl}
        F_0(\xi_{ik})&\equiv&0\pmod p\quad [\text{在 $\xi_{ik}$ 中恆等}],\quad\text{或者}\\
        F(\xi_{ik})&\equiv&0\pmod{(\frN,p)}\quad [\text{在 $\xi_{ik}$ 中恆等}].
\end{array}
\end{equation}
這同時以關於 $F_0(\xi_{ik})$ 的更強表述證明了 (2) 的逆，從而證明了 (1)。\footnote{這個更強表述意謂每個剩餘類只含有一個規範化多項式。}

% unit: noether.p15.par.028
還只需補充說明，$G(\xi_{ik})\equiv0\pmod p$ 由 $K(\xi_{ik})\equiv0\pmod p$ 推出；等價地，由 $G(\xi_{ik})\not\equiv0\pmod p$ 推出 $K(\xi_{ik})\not\equiv0\pmod p$。規範化只在歸納結論的這一部分中需要。在替換 $\xi_{i4}=\xi_{i3}$ 下，$G$ 的各個冪乘積並不消失。因此 $K$ 的消失只能由於 $G$ 中不同的冪乘積對應於 $K$ 中相同的冪乘積；而這又只有在這些冪乘積僅僅相差指標 3 和 4 的互換時才可能發生。設例如這樣的規範化冪乘積為
% unit: noether.p15.eq.073
\[
        \xi_{12}^{\tau_{12}}\xi_{13}^{\tau_{13}}\xi_{14}^{\tau_{14}}\xi_{24}^{\tau_{24}}.
\]
由於假定在每個指標中分別齊次，在某個 $\xi_{ik}$ 中把 3 替換為 4 必須伴隨在另一個 $\xi_{ik}$ 中把 4 替換為 3；一個相應項就會含有乘積 $\xi_{14}\xi_{23}$ 而不是 $\xi_{13}\xi_{24}$，這被規範化排除了。完全同樣地，對於
% unit: noether.p15.eq.074
\[
        \xi_{12}^{\tau_{12}}\xi_{13}^{\tau_{13}}\xi_{14}^{\tau_{14}}\xi_{24}^{\tau_{24}}
\]
能夠相抵的也只能是含有因子 $\xi_{14}\xi_{23}$ 的項，而這同樣被規範化排除了。因此 $K(\xi_{ik})\equiv0\pmod p$ 蘊含 $G(\xi_{ik})\equiv0\pmod p$。

% unit: noether.p15.subsec.003
\subsection*{c) $n$ 行 $\alpha_i,\beta_i$ 的情形}

% unit: noether.p15.par.029
$n$ 行的情形現在通過完全歸納來處理，正如剛才處理的四行情形的嚴格推廣。\footnote{證明也仍適用於 $n=3$。} 首先我仍證明：對於每個多項式 $F$，有一個同餘式
% unit: noether.p15.eq.075
\begin{equation}
        F(\xi_{ik})\equiv F_0(\xi_{ik})\pmod{\frN}
        \quad \text{在 $\xi_{ik}$ 中恆等},
\end{equation}
其中 $\frN=(f_{ik,rs})$，而 $F_0(\xi_{ik})$ 是規範化的。規範化規定如下：若 $i,k,r,s$ 是四個不同指標且 $i<k<r<s$，則 $F_0$ 不應含有乘積 $\xi_{is}\xi_{kr}$；也就是說，一個 $\xi$ 的兩個指標不應同時位於另一個 $\xi$ 的兩個指標之間。對於兩行和三行的情形，這種規範化還沒有任何限制，因而每個多項式已經是規範化的；在四行的情形下，如已證明，每個剩餘類也存在一個規範化多項式。因此可以假定，對於 $(n-1)$ 行，每個剩餘類中規範化多項式的存在已經證明。

% unit: noether.p15.par.030
把 $F$ 中任意一個冪乘積寫成
% unit: noether.p15.eq.076
\[
        \xi_{i_1n}\cdots\xi_{i_\sigma n}\,P,
\]
其中 $P$ 不再含指標 $n$。由歸納假設，$P$ 模 $\frN$ 同餘於一個整的規範化多項式 $P_0$：
% unit: noether.p15.eq.077
\[
        \xi_{i_1n}\cdots\xi_{i_\sigma n}P
        \equiv
        \xi_{i_1n}\cdots\xi_{i_\sigma n}P_0
        \pmod{\frN}.
\]
如果 $P_0$ 不含任何 $\xi_{rs}$，使得對至少一個指標 $i$，兩個指標 $r,s$ 位於 $i_\lambda$ 與 $n$ 之間，即 $i_\lambda<r<s<n$，那麼規範化已經達到，因為對於來自 $P_0$ 的任意兩個 $\xi$ 的乘積，條件都已滿足。否則，取這樣的一個指標 $i_\lambda$，使 $i_\lambda<r<s<n$，並令 $P_0^*$ 為 $P_0$ 中含有 $\xi_{rs}$ 的冪乘積。由於
% unit: noether.p15.eq.078
\[
        \xi_{i_\lambda n}\xi_{rs}\equiv
        \xi_{i_\lambda s}\xi_{rn}-\xi_{i_\lambda r}\xi_{sn}
        \pmod{\frN},
\]
於是得到
% unit: noether.p15.eq.079
\begin{align*}
 &\xi_{i_1n}\cdots\xi_{i_\lambda n}\cdots\xi_{i_\sigma n}P_0^*\\
 &\equiv
 \xi_{i_1n}\cdots\xi_{rn}\cdots\xi_{i_\sigma n}Q
 +
 \xi_{i_1n}\cdots\xi_{sn}\cdots\xi_{i_\sigma n}R
 \pmod{\frN},\\
 &\equiv
 \xi_{i_1n}\cdots\xi_{rn}\cdots\xi_{i_\sigma n}Q_0
 +
 \xi_{i_1n}\cdots\xi_{sn}\cdots\xi_{i_\sigma n}R_0
 \pmod{\frN},
\end{align*}
其中 $Q_0$ 和 $R_0$ 又表示 $Q$ 和 $R$ 的剩餘類的規範化多項式。它們由歸納假設存在，因為 $Q$ 和 $R$ 不含指標 $n$；並且它們當然是整多項式，因為 $Q$ 和 $R$ 只是冪乘積。由於 $i_\lambda<r<s<n$，還得到
% unit: noether.p15.eq.080
\[
 ((n-i_1)+\cdots(n-r)+\cdots(n-i_\sigma))
 <
 ((n-i_1)+\cdots(n-i_\lambda)+\cdots(n-i_\sigma)),
\]
% unit: noether.p15.eq.081
\[
 ((n-i_1)+\cdots(n-s)+\cdots(n-i_\sigma))
 <
 ((n-i_1)+\cdots(n-i_\lambda)+\cdots(n-i_\sigma)).
\]
所以從 $P_0^*$ 過渡到 $Q_0,R_0$ 時，這個差的和減小了；由於該和必然 $\ge \sigma$，對於每個冪乘積 $P_0^*$，這個過程經過有限步必須終止。於是得到
% unit: noether.p15.eq.082
\[
        F(\xi_{ik})
        =
        \sum \xi_{i_1n}\cdots\xi_{i_\sigma n}\cdot P^{(\nu)}
        \equiv
        \sum \xi_{j_1n}\cdots\xi_{j_\sigma n}\cdot S_0^{(\nu)}
        =
        F_0(\xi_{ik})\pmod{\frN},
\]
其中整多項式 $F_0(\xi_{ik})$ 必然是規範化的，否則差的和還可以繼續降低。這證明了 (10)。又因為按 $(n-1)$ 個指標的歸納假設，每個分別齊次的 $F$ 都有一個分別齊次的規範化 $F_0$，剛給出的過程說明，對 $n$ 個指標也是如此。故以下只需考慮分別齊次的多項式。

% unit: noether.p15.par.031
如 b) 中一樣，置
% unit: noether.p15.eq.083
\begin{equation}
        F_0(\xi_{ik})=G(\xi_{ik})+\xi_{n-1,n}F_1(\xi_{ik}),\footnote{對應於 (10) 和 (11) 的恆等式
        \[
          F((ik))=F_0((ik))=G((ik))+(n-1,n)\cdot F_1((ik))
        \]
        一旦把兩行 $(n-1)$ 和 $(n)$ 替換為變量行 $x,y$，從而把 $(n-1,n)$ 替換為 $(xy)$，就表示 Clebsch--Gordan 級數展開的一個類似物。但是，Clebsch--Gordan 是按極化式展開，而 $F_0$ 的規範化對應於極化式的“首項”，由此迫使 Clebsch--Gordan 中並不成立的整性出現。不考慮整性時，$f_{ik,rs}$ 形成 $\frM$ 的一個基這一已知證明，正是藉助 Clebsch--Gordan 級數展開進行的（Gordan--Kerschensteiner, \emph{Vorlesungen über Invariantentheorie}, p. 132）。}
\end{equation}
其中 $G$ 不含任何可被 $\xi_{n-1,n}$ 整除的冪乘積；並且再置
% unit: noether.p15.eq.084
\begin{equation}
        [G]_{(\xi_{in}=\xi_{i,n-1})}=K(\xi_{ik}).
\end{equation}
由規範化的定義可見，$K$ 和 $G$ 一樣，也是規範化的。

% unit: noether.p15.par.032
在 $G$ 和 $K$ 之間又有一一對應。$G$ 中不同的冪乘積只有在僅僅相差指標 $n$ 與 $(n-1)$ 的互換時，才可能對應於 $K$ 中同一個乘積。設例如 $G$ 在指標 $(n-1)$ 中次數為 $\rho$，在指標 $n$ 中次數為 $\sigma$，並設
% unit: noether.p15.eq.085
\[
        \chi_1\le\chi_2,\ldots,\le\chi_\rho\le\chi_{\rho+1},\ldots,\le\chi_{\rho+\sigma}
\]
是一個 $G$ 的冪乘積中與 $(n-1)$ 和 $n$ 相連的指標。它們的個數必須為 $\varrho+\sigma$，因為由 (11) 知不定元 $\xi_{n-1,n}$ 不出現在 $G$ 中。由規範化，冪乘積因此具有形式
% unit: noether.p15.eq.086
\[
        \xi_{\varkappa_1,n-1}\cdots\xi_{\varkappa_\varrho,n-1}
        \xi_{\varkappa_{\varrho+1},n}\cdots\xi_{\varkappa_{\varrho+\sigma},n}\cdot P(\xi_{ik}),
\]
其中 $P(\xi_{ik})$ 不再含指標 $n$ 和 $(n-1)$。每次把 $n$ 與 $(n-1)$ 互換，若 $i\ne j$，就會給出因子 $\xi_{in}\xi_{j,n-1}$ 且 $i<j$；於是 $j,n-1$ 位於 $i$ 與 $n$ 之間，該項就不是規範化的。因此，在固定 $\varkappa_1,\varkappa_\varrho,\ldots,\varkappa_{r+\sigma}$ 和固定 $P$ 的情況下，$G$ 只含一個冪乘積。這證明了 $G$ 中不同的冪乘積對應於 $K$ 中不同的乘積。因此
% unit: noether.p15.eq.087
\begin{equation}
        K(\xi_{ik})\equiv0\pmod p\quad\Longrightarrow\quad
        G(\xi_{ik})\equiv0\pmod p.
\end{equation}

% unit: noether.p15.par.033
由
% unit: noether.p15.eq.088
\[
        F((ik))\equiv0\pmod p\quad \text{在 $\alpha,\beta$ 中恆等}
\]
按照 (10)，並顧及 (2)，現在得到
% unit: noether.p15.eq.089
\begin{equation}
        F_0((ik))\equiv0\pmod p\quad \text{在 $\alpha,\beta$ 中恆等}.
\end{equation}
特殊化 $\alpha_n=\alpha_{n-1},\;\beta_n=\beta_{n-1}$，由 (11) 和 (12) 得到
% unit: noether.p15.eq.090
\begin{equation}
        K((ik))\equiv0\pmod p\quad [\text{在 $\alpha,\beta$ 中恆等}].
\end{equation}
由於 $K(\xi_{ik})$ 是規範化的並且只含有 $(n-1)$ 個指標，由 a) 和 b) 中公式 (9) 所證明的結果可推出
% unit: noether.p15.eq.091
\[
        K(\xi_{ik})\equiv0\pmod p\quad \text{在 $\xi_{ik}$ 中恆等}.
\]
於是由 (13)
% unit: noether.p15.eq.092
\[
        G(\xi_{ik})\equiv0\pmod p\quad \text{在 $\xi_{ik}$ 中恆等}.
\]
從而 (11) 給出
% unit: noether.p15.eq.093
\[
        F_0(\xi_{ik})\equiv \xi_{n-1,n}F_1(\xi_{ik})\pmod p
        \quad \text{在 $\xi_{ik}$ 中恆等}.
\]
由於 $(n,n-1)\not\equiv0\pmod p$ 在 $\alpha,\beta$ 中不恆等為零，由 (14) 推出
% unit: noether.p15.eq.094
\[
        F_1((ik))\equiv0\pmod p\quad [\text{在 $\alpha,\beta$ 中恆等}],
\]
其中 $F_1(\xi_{ik})$ 又是規範化的，並且在 $\xi_{n-1,n}$ 中次數低於 $F_0$。有限次重複這個過程，得到
% unit: noether.p15.eq.095
\[
        F_0(\xi_{ik})\equiv\xi_{n-1,n}^{\lambda}F_\lambda(\xi_{ik})\pmod p,
\]
其中 $F_\lambda$ 在 $\xi_{n-1,n}$ 中次數為零，因而由剛才證明的結果知其模 $p$ 消失。由此得到
% unit: noether.p15.eq.096
\begin{equation}
\begin{array}{rcl}
        F_0(\xi_{ik})&\equiv&0\pmod p\quad \text{在 $\xi_{ik}$ 中恆等},\quad\text{或者}\\
        F(\xi_{ik})&\equiv&0\pmod{(\frN,p)}\quad \text{在 $\xi_{ik}$ 中恆等}.
\end{array}
\end{equation}
這就證明了 (2) 的逆，從而證明了 (1)。同時也證明了在 $n$ 個指標情形下、因而在一般情形下，歸納步驟中對 $F_0(\xi_{ik})$ 所使用的更強斷言。

% unit: noether.p15.sec.004
\setcounter{equation}{0}
\section*{\S\ 4. 有限群的整不變量的有限性}

% unit: noether.p15.par.034
我在短文“有限群不變量的有限性定理”\footnote{\emph{Math. Ann.} 77, p. 89 (1915).}中給出的有限性的初等證明，藉助 \S\ 1 的特殊有限性定理 III 和 IV，可以強化為整性；而通常依靠 Hilbert 模基定理的證明並不給出這種整性。即便使用整模基，也只給出一種表示，其中在分母中出現整數 $h$ 的任意高次冪；這裡 $h$ 是群的階。

% unit: noether.p15.par.035
我全篇沿用上述短文的記號。設有限群 $\frH$ 由 $h$ 個行列式不為零的線性變換 $A_1,\ldots,A_h$ 組成，\footnote{只需假設 $\frH$ 與一個抽象有限群單同構；也就是說，如果某個 $A$ 的變換行列式為零，則 $\frH$ 應與一個群 $\begin{pmatrix}\frS&0\\0&0\end{pmatrix}$ 相似，其中 $\frS$ 的替換行列式都不為零。在此情況下，$\frH$ 的不變量與 $\frS$ 的不變量相一致，因此行列式不為零的條件實際上並不限制一般性。} 其中 $A_k$ 表示線性變換
% unit: noether.p15.eq.097
\[
        x_1^{(k)}=\sum_{\nu=1}^{n}a_{1\nu}^{(k)}x_\nu,
        \quad\ldots,\quad
        x_n^{(k)}=\sum_{\nu=1}^{n}a_{n\nu}^{(k)}x_\nu,
\]
或簡寫為 $(x^{(k)})=A_k(x)$。於是群 $\frH$ 把元素為 $x_1,\ldots,x_n$ 的行 $(x)$ 變為元素為 $x_1^{(k)},\ldots,x_n^{(k)}$ 的各行 $(x^{(k)})$。係數 $a_{i\nu}^{(k)}$ 假定為代數整數；這並不限制一般性，因為根據 J. Schur 的一個定理，\footnote{\emph{Über Gruppen linearer Substitutionen mit Koeffizienten aus einem algebraischen Zahlkörper}. Satz VII. \emph{Math. Ann.} 71, p. 555 (1912).} 每一個線性變換有限群都相似於一個群，使其中的 $a_{i\nu}^{(k)}$ 成為代數整數。

% unit: noether.p15.par.036
所謂群的整（絕對）不變量，是指 $x_1,\ldots,x_n$ 的一個多項式有理函數，其係數為由 $a_{i\nu}^{(k)}$ 生成的數域 $\frK$ 中的代數整數，並且在施加 $A_1,\ldots,A_h$ 時恆等不變。對於這樣的不變量 $f(x)$ 有
% unit: noether.p15.eq.098
\begin{equation}
        f(x)=f(x^{(1)})=\cdots=f(x^{(h)})
        =\frac1h\cdot\sum_{k=1}^{h}f(x^{(k)}).
\end{equation}
反過來，$h$ 行量 $(x^{(k)})$ 的每一個整對稱函數也成為該群的一個整不變量。

% unit: noether.p15.par.037
現在令 $\omega_1,\omega_2,\ldots,\omega_\varrho$ 表示 $\frK$ 的所有整數的一個基。那麼每個不變量都可表示為
% unit: noether.p15.eq.099
\[
        f(x)=\omega_1 f_1(x)+\cdots+\omega_\varrho f_\varrho(x),
\]
其中 $f_1(x),\ldots,f_\varrho(x)$ 有有理整數係數。於是由 (1) 得到
% unit: noether.p15.eq.100
\[
        h.f(x)
        =
        \omega_1\sum_{k=1}^{h}f_1(x^{(k)})
        +\cdots+
        \omega_\varrho\sum_{k=1}^{h}f_\varrho(x^{(k)}).
\]
這裡各個和是 $h$ 行量 $(x^{(k)})$ 的對稱函數，並且具有有理整數係數。於是根據定理 III，它們可用這些行的有限多個整有理對稱函數，以有理整數係數，多項式且有理地表示。由前面的論證，這些基函數 $I_1(x),\ldots,I_\sigma(x)$ 成為該群的整不變量；一般地，它們具有代數整係數，因為 $\frH$ 不必與 $A$ 一起包含它的所有代數共軛變換。

% unit: noether.p15.par.038
因此存在表示
% unit: noether.p15.eq.101
\begin{equation}
        h\cdot f(x)=\omega_1G_1(I)+\cdots+\omega_\varrho G_\varrho(I),
\end{equation}
其中 $G_1,\ldots,G_\varrho$ 有有理整數係數。

% unit: noether.p15.par.039
現在令 $w_1,\ldots,w_\varrho$ 為不定元，並使
% unit: noether.p15.eq.102
\[
        L=w_1G_1(I)+\cdots+w_\varrho G_\varrho(I)
\]
遍歷所有關於 $w$ 的線性型，這些線性型通過代換 $w_i=\omega_i/h$，在表示 (2) 中對應於不變量 $f(x)$。Hilbert 關於整模基的定理給出
% unit: noether.p15.eq.103
\begin{equation}
        L=A_1(I)L_1+\cdots+A_\tau(I)L_\tau,
\end{equation}
其中 $A_i$ 有有理整數係數；並且由於所有 $L$ 對 $w$ 都是線性的，$A_i$ 不再含 $w$。如果代換 $\omega_i=\frac{w_i}{h}$ 把 $L_1,\ldots,L_\tau$ 變為不變量 $\mathfrak I_1(x),\ldots,\mathfrak I_\tau(x)$，則 (3) 變為
% unit: noether.p15.eq.104
\begin{equation}
        f(x)=A_1(I)\cdot\mathfrak I_1(x)+\cdots+A_\tau(I)\mathfrak I_\tau(x).
\end{equation}
這個表示說明
% unit: noether.p15.eq.105
\[
        I_1(x),\ldots,I_\sigma(x);\mathfrak I_1(x),\ldots,\mathfrak I_\tau(x)
\]
構成群 $\frH$ 的所有整不變量的一個基，使得每個不變量都可表示為這有限多個不變量的、具有有理整數係數的多項式有理函數。

% unit: noether.p15.par.040
如果特別地，群 $\frH$ 具有如下性質：它與每個變換 $A$ 一起還包含所有具有代數共軛數係數的 $A',A'',\ldots$，那麼 $h$ 行 $(x^{(k)})$ 的整有理對稱函數成為 $x$ 的具有有理整數係數的函數。特別地，對稱函數的基函數也如此。因此，如果把自己限制於具有有理整數係數的不變量 $f(x)$，則基函數
% unit: noether.p15.eq.106
\[
        I_1(x),\ldots,I_\sigma(x);\mathfrak I_1(x),\ldots,\mathfrak I_\tau(x)
\]
也有有理整數係數。在這個特殊假設下——一般對群 $\frH$ 並不滿足——整有限性因而也適用於具有有理整數係數的不變量的子域。
%
\clearpage
%
% unit: noether.p16.title.001
\pdfbookmark[0]{16. 形式論中的級數展開}{paper-16}
\section*{16. 形式論中的級數展開}
\begin{center}
\emph{Math. Ann.} 81 (1920), pp. 25--30
\end{center}

% unit: noether.p16.par.001
所謂“形式論的級數展開”，眾所周知，一方面是指 Clebsch--Gordan 級數展開的推廣；也就是說，把一個含有兩個二元系列的形式，按這兩個系列的行列式的冪展開，其中係數是隻含一個變量系列的形式的極化式，或者等價地說，在應用 $\Omega$ 過程時為零。這樣的展開被推廣到任意多個系列、每個系列含 $n$ 個變量的形式。另一方面，作為對“復形座標” $t_{ik}$ 中出現的標準型的推廣，它也是對同時含有不同階層的量 $t_i, t_{ik}, t_{ikl}\ldots$ 的形式按標準型作展開。不過，這第二種展開可以由第一種通過不變微分過程得到，Mertens 和 Deruyts 尤其完成了這一點。\footnote{F. Mertens, \emph{Über invariante Gebilde quaternärer Formen} (Sitzber. d. Ak. d. Wiss. Wien 98, Abt. IIa (1889))，處理四元情形。一般標準型見 J. Deruyts: \emph{Sur la réduction des fonctions invariantes dans le système des variables géométriques}, Bulletins de l'Ac. R. de Belgique 24 (1892)。在不知道這篇工作的情況下，我緊密依循 Mertens，在論文 \emph{Zur Invariantentheorie der Formen von $n$ Variabeln}, J. f. M. 139 (1910) 的 \S~7 中給出了一般的按標準型的展開。}

% unit: noether.p16.par.002
第一種的各種展開，正如例如在 \emph{Math. Ann.} 77 前引處 III 中簡略勾畫的那樣，都作為如下展開的推論而出現：把一個由 $n$ 個系列、每個系列含 $n$ 個變量組成的形式，按這些系列的行列式的冪展開，其係數是隻含 $(n-1)$ 個系列的形式的極化式，而這些係數本身又由原始形式通過極化形成和 $\Omega$ 過程導出。\footnote{參見前引處的文獻說明。此後又有一篇同樣基於形式方法的論文發表: Schouten, \emph{Over reeksontwikkelingen van algebraische vormen met verschillende rijen van variabelen van verschillenden graad}, Verslag Ak. Amsterdam 27 (1919), p. 1481.} 原始的 Clebsch--Gordan 展開 $(n=2)$ 還更進一步：其中出現的特殊極化過程被明確給出，甚至包括數值係數。

% unit: noether.p16.par.003
我現在想給所指出的那條札記補充的是一種精確到數值係數的顯式表示；在那裡，級數展開較為概念性地被理解為線性形式束的問題。這種表示由 E. Fischer 關於一般模系的研究作專門化而得到。\footnote{E. Fischer, \emph{Über die Differentiationsprozesse der Algebra}, J. f. M. 148 (1917).}

% unit: noether.p16.par.004
通向這種歸入的，是對 $t_{ik}$ 中標準型的解釋。事實上，如果把“復形座標” $t_{ik}$ 替換為“線座標”
% unit: noether.p16.eq.001
\[
        p_{ik}=(x_i y_k-x_k y_i),
\]
那麼由於 $p_{ik}$ 之間存在恆等關係，一個形式 $F(p_{ik})$ 的表示就不再唯一。如果要求 $F$ 的係數之間滿足與 $F$ 中出現的 $p_{ik}$ 的冪積之間相同的線性關係，就可以取得唯一性。\footnote{參見 Clebsch, \emph{Über eine Fundamentalaufgabe der Invariantentheorie}, \S~4, Abh. d. K. G. d. Wissensch. zu Göttingen 17 (1872).} 因而對於這個標準型 $\Phi$ 有:
% unit: noether.p16.eq.002
\begin{equation}
\begin{aligned}
        F(p_{ik})&=\Phi(p_{ik}) &&[\text{在 }x,y\text{ 中恆等}],\quad\text{或者}\notag\\
        F(t_{ik})&\equiv\Phi(t_{ik}) &&\modu{M},
\end{aligned}
\tag{1}
\end{equation}
其中 $M$ 表示 $p_{ik}$ 之間所有恆等關係所成的模，也就是說，它由所有這樣的形式 $G(t_{ik})$ 組成：$G(p_{ik})$ 在 $x,y$ 中恆等地消失。$\Phi(t_{ik})$ 是復形座標中的“標準型”，而 $M$ 的基函數則成為 $t_{ik}$ 中的二次形式。對這些基函數的因子反覆施行這一過程，就得到完整的級數展開。當 $t_i,t_{ik},t_{ikl},\ldots$ 同時出現時，也有相應的考察。

% unit: noether.p16.par.005
按照類似的思路，也可以把按極化式的展開加以歸約。例如，在一個形式 $F(x,y,z)$ 中，把三個三元系列 $x=(x_1,x_2,x_3),y,z$ 替換為只由其中兩個線性組成的系列 $\xi,\eta,\zeta$，例如 $\xi=x,\eta=y;\zeta=\lambda_1x+\lambda_2y$，那麼 $F(\xi,\eta,\zeta)$ 的表示又不再唯一。仍然可以規定，係數之間滿足與 $F$ 中出現的 $\xi,\eta,\zeta$ 的冪積之間相同的線性關係。由於在這裡所有恆等關係所成的模 $M$ 以三個系列的行列式 $\Delta=(xyz)$ 為基函數（前引處 III，輔助引理 a)），[1] 被替換為:
% unit: noether.p16.eq.003
\begin{equation}
\begin{aligned}
        F(\xi,\eta,\zeta)&=\Phi(\xi,\eta,\zeta)&&[\text{在 }x,y\text{ 中恆等}],\quad\text{或者}\notag\\
        F(x,y,z)&\equiv\Phi(x,y,z)&&\modu{\Delta}.
\end{aligned}
\tag{2}
\end{equation}
可以看出，$\Phi$ 由極化形成而產生（參見前引處公式 (2)）；完整的級數展開又由迭代得到。

% unit: noether.p16.par.006
因此，在這兩種情形中，標準型 $\Phi$ 的形成都從屬於如下問題：若在 $F(z)$ 中把不定元 $z_i$ 替換為參數 $\xi$ 的多項式 $f_i(\xi)$，使得表示 $F(f_i(\xi))$ 因而不再唯一，則應通過要求係數之間滿足與 $F$ 中出現的 $f_i(\xi)$ 的冪積之間相同的線性關係來確定唯一性。於是有
% unit: noether.p16.eq.004
\begin{equation}
\begin{aligned}
        F(f_i(\xi))&=\Phi(f_i(\xi))&&[\text{在 }\xi\text{ 中恆等}],\quad\text{或者}\notag\\
        F(z)&\equiv\Phi(z)&&\modu{M},
\end{aligned}
\tag{3}
\end{equation}
其中 $M$ 表示 $f_i(\xi)$ 之間所有恆等關係所成的模。

% unit: noether.p16.par.007
對於這個標準型 $\Phi(z)$，E. Fischer 在所引論文的 \S~11, I 中給出了精確到數值係數的顯式展開，即藉助線性運算給出。同時，他在導言 III 和 \S~6 II 中，也對屬於 $M$ 的差 $F-\Phi$ 給出了同一意義下的顯式表達；這裡假定已經知道一個模基。又通過在 \S~11 II 中合併這些展開，他於是用一個顯式表示替換了公式 [3]。這裡仍未解決的數值係數問題，按照同一論文的 \S~12，等同於確定線性運算的特徵值和特徵函數（公式 82）。

% unit: noether.p16.par.008
現在我在 1. 中把 Fischer 的公式專門化到按極化式的級數展開的情形，並由此通過迭代得到級數展開。確定 $\Phi$ 的運算由此變成確定的極化過程，這比上面給出的已知表述要精確得多。在確定 $F-\Phi$ 的運算中，乘以行列式 $\Delta$ 與 $\Omega$ 過程組合地出現。為了得到通常的非顯式形式，還必須注意到 $\Omega\Delta$ 可以用極化過程來表示。

% unit: noether.p16.par.009
這裡還應補上一條修正。在上面引用的工作中（p. 101，自上第 6 行），我曾無意中把恆等變換
% unit: noether.p16.eq.005
\[
        \Omega(\Delta\psi)=c_1\psi+c_2\Delta\cdot\Omega(\psi),
\]
當作普遍有效；但它只在二元領域中成立。我因而從給出 [2] 中標準型 $\Phi$ 表達式的基本公式 (2) 過渡到了級數展開 (3)。因此，那裡的考察只給出公式 (2)，並作為其特殊情形給出歸約定理 (1)；但由於缺少差 $F-\Phi$ 的顯式表達，它們不足以奠定級數展開。

% unit: noether.p16.par.010
在 2. 中，我同樣通過專門化得到 [1] 的標準型 $\Phi(t_{ik})$；更一般地，得到標準型 $\Phi(t_i,t_{ik},t_{ikl},\ldots)$，這導向已知的顯式公式。不過，由於建立完整的級數展開需要知道 $M$ 的一個模基，這裡不變論的道路（參見所引文獻）更為可取；它不需要這種知識，而反過來通過級數展開給出一個模基。

% unit: noether.p16.par.011
1. 按 Fischer, \S~11（公式 (19) 及以下各行），[3] 的標準型 $\Phi(z)$ --- 在那裡記為 $N(\zeta)$ --- 由下式給出:
% unit: noether.p16.eq.006
\begin{equation}
        \Phi(z)=(a_1S+a_2S^2+\cdots+a_rS^r)F(z),
\tag{4}
\end{equation}
其中常數 $a_1\ldots a_r$ 屬於 $f(\xi)$ 的係數域，並且只依賴於 $F(z)$ 的次數；算子 $S$ 定義為
% unit: noether.p16.eq.007
\begin{equation}
\begin{gathered}
        S=AB;\qquad BF(z)=F(f(\xi));\\
        AF(f(\xi))=\left\{\sum z_i f_i\!\left(\frac{\p}{\p\xi}\right)\right\}^{p}F(f(\xi))\footnotemark.
\end{gathered}
\tag{5}
\end{equation}
\footnotetext{參見公式 (75)。在 [3] 的情形中，求和 $\alpha$ 遍歷所有 $p$ 維的冪積，因而 (75) 專門化為 [5]。}
如果在 [2] 中把 $F$ 看作僅關於 $z$ 的 $p$ 維形式，那麼 $f_i(\xi)$ 專門化為 $\lambda_1x_i+\lambda_2y_i$；因此得到:
% unit: noether.p16.eq.008
\[
\begin{gathered}
        BF=F(x,y,\lambda_1x+\lambda_2y),\\
        SF=\left\{\sum z_i\left(\frac{\p}{\p\lambda_1}\frac{\p}{\p x_i}+
        \frac{\p}{\p\lambda_2}\frac{\p}{\p y_i}\right)\right\}^{p}
        F(x,y,\lambda_1x+\lambda_2y),
\end{gathered}
\]
或寫成極化過程的形式（參見前引處 II, 2 和 3；縮寫記號亦同）:
% unit: noether.p16.eq.009
\begin{equation}
        SF(x,y,z)=\frac1{p!}\left[z\left(\frac{\p}{\p\lambda_1}\frac{\p}{\p x}+
        \frac{\p}{\p\lambda_2}\frac{\p}{\p y}\right)\right]^p
        \left[(\lambda_1x+\lambda_3y)\frac{\p}{\p z}\right]^p F(x,y,z).
\tag{6}
\end{equation}
\emph{方程} [4] \emph{和} [6] \emph{給出所要求的} $\Phi(x,y,z)$ \emph{的顯式表達}；同時 $\Phi(x,y,z)$ 歸入開頭提到的一般形式（前引處 (2) 中的 $\sum PH$）--- 即只依賴於兩個系列、並由 $F$ 通過極化過程得到的表達式的極化式。尚未確定的係數 $a_i$ 是有理數，因為這裡 $f_i(\xi)$ 具有整有理的數值係數。若把 [2] 寫成
% unit: noether.p16.eq.010
\begin{equation}
        F=\Phi+\Delta F_1,
\tag{2a}
\end{equation}
則由 Fischer 在導言 III 和 \S~6 II（p. 41）給出的公式專門化得到
% unit: noether.p16.eq.011
\begin{equation}
        F_1=(\alpha_1\Omega+\alpha_2\Omega\Delta\Omega+\cdots+\alpha_\lambda\Omega\Delta\Omega\cdots\Delta\Omega)F\footnotemark,
        \qquad \left[\Omega=\left(\frac{\p}{\p x}\frac{\p}{\p y}\frac{\p}{\p z}\right)\right],
\tag{7}
\end{equation}
\footnotetext{Kerim 在其即將出版的 Erlangen 博士論文中，將這個專門化公式用於“慣性形式”。}
其中 $c_i$ 又是隻依賴於 $F$ 的次數的有理數。

% unit: noether.p16.par.012
迭代導致:
% unit: noether.p16.eq.012
\[
        F_1=\Phi_1+\Delta F_2;\quad F_2=\Phi_2+\Delta F_3;\quad\ldots;\quad
        F_i=\Phi_i+\Delta F_{i+1};\quad\ldots,
\]
其中每一個 $\Phi_i$ 都是對應於 $F_i$ 的標準型。由於 $F_i$ 在 $z$ 中的次數為 $(p-i)$，由 [4] 和 [6] 得:
% unit: noether.p16.eq.013
\begin{equation}
\begin{gathered}
        \Phi_i=(a_{i1}S_i+a_{i2}S_i^2+\cdots)F_i,\\
        S_iF_i=\frac1{(p-i)!}
        \left[z\left(\frac{\p}{\p\lambda_1}\frac{\p}{\p x}+
        \frac{\p}{\p\lambda_2}\frac{\p}{\p y}\right)\right]^{p-i}
        \left[(\lambda_1x+\lambda_2y)\frac{\p}{\p z}\right]^{p-i}F_i.
\end{gathered}
\tag{8}
\end{equation}
對應於公式 [7] 有
% unit: noether.p16.eq.014
\[
        F_i=(\alpha_{i1}\Omega+\alpha_{i2}\Omega\Delta\Omega+\cdots)F_{i-1},
        \quad\text{並且通過迭代}
\]
% unit: noether.p16.eq.015
\begin{equation}
        F_i=(\beta_1\Omega^i+\beta_2\Omega^i\Delta\Omega+\cdots
        \beta_\rho\Omega^{k_1}\Delta\Omega^{k_2}\cdots\Delta\Omega^{k_\lambda}+\cdots)F
        \quad(k_1+k_2+\cdots+k_\lambda-\lambda+1=i).
\tag{9}
\end{equation}
\emph{藉助} [8] \emph{和} [9]，\emph{展開}
% unit: noether.p16.eq.016
\[
        F=\Phi+\Delta\Phi_1+\Delta^2\Phi_2+\cdots+\Delta^i\Phi_i+\cdots\Delta^p\Phi_p
\]
\emph{已經達到上述意義下的顯式表示。}

% unit: noether.p16.par.013
若對 [7] 應用恆等變換 $\Omega\Delta F=PF$，其中 $P$ 表示極化過程，\footnote{例如參見 F. Mertens, \emph{Über ganze Funktionen von $m$ Systemen von je $n$ Unbestimmten}, Monatsh. f. Math. u. Phys. 4, 公式 (5)。} 那麼，因為這些過程與 $\Omega$ 過程可交換，就得到已知公式
% unit: noether.p16.eq.017
\[
        F_1=P\Omega F;\quad\ldots;\quad F_i=P\Omega^iF;
\]
由此結合 [8] 或前引處的考察，便得到級數展開的已知表述，例如前引處 (3) 中的形式。如果處理的是 $n$ 個系列、每個系列含 $n$ 個變量，$x^{(1)},x^{(2)},\ldots,x^{(n-1)},z$，那麼相應地有
% unit: noether.p16.eq.018
\[
        f_i(\xi)=\lambda_1x_i^{(1)}+\lambda_2x_i^{(2)}+\cdots+\lambda_{n-1}x_i^{(n-1)};
\]
並且 $\Delta$ 與 $\Omega$ 也應按 $n$ 個系列來定義。

% unit: noether.p16.par.014
2. 對於 $F(t_i,t_{ik},t_{ikl},\ldots)$ 的標準型，函數 $f_i(\xi),f_{ik}(\xi),\ldots$ 由一系列、二系列、一直到 $(n-1)$ 系列的行列式給出:
% unit: noether.p16.eq.019
\[
        \xi_i^{(1)},\quad (\xi^{(1)}\xi^{(2)})_{ik},\quad\ldots,\quad
        (\xi^{(1)}\xi^{(2)}\cdots\xi^{(n-1)})_{i_1i_2\cdots i_{n-1}}.
\]
於是，作為 [5] 的專門化，得到:
% unit: noether.p16.eq.020
\[
\begin{gathered}
        S=AB;\qquad
        BF=F\bigl(\xi_i^{(1)},(\xi^{(1)}\xi^{(2)})_{ik},\ldots\bigr),\\
        ABF=\left(\sum t_i\frac{\p}{\p\xi_i^{(1)}}\right)^{\lambda_1}
        \left(\sum t_{ik}\left(\frac{\p}{\p\xi^{(1)}}\frac{\p}{\p\xi^{(2)}}\right)_{ik}\right)^{\lambda_2}
        \cdots F\bigl(\xi_i^{(1)},(\xi^{(1)}\xi^{(2)})_{ik},\ldots\bigr),
\end{gathered}
\]
其中 $\lambda_1,\lambda_2\ldots$ 表示在 $t_i,t_{ik},\ldots$ 中的次數。

% unit: noether.p16.par.015
由不變論的考察（也就是說，$F$ 的所有關於係數為線性的那些不變量，都可以由 $\Phi$ 和來自 $M$ 的形式線性組合而成，因而特別也可以由不屬於 $M$ 的 $SF$ 組成）\footnote{這可例如由我所引論文 \emph{Zur Invariantentheorie der Formen von $n$ Variabeln} 的 \S~6 和 \S~7, 2. 得出。} 得到這裡的同餘 $\Phi\equiv a_1SF\modu{M}$；並且由於標準型性質，由此取代 [4] 得到簡單關係
% unit: noether.p16.eq.021
\begin{equation}
        \Phi=a_1SF=a_1S\Phi.
\tag{10}
\end{equation}
第二個等式成立，是因為應用 $S$ 會使來自 $M$ 的每一個形式消失。公式 [10] 表明標準型 $\Phi$ 本身同時也是特徵函數；而且按照上面的不變論考察，並不存在其他同時又是 $F$ 的不變量的特徵函數（關於 [10]，參見 Deruyts，公式 (10) 和 (10')）。

% unit: noether.p16.par.016
\begin{center}
\textbf{更正。}
\end{center}

% unit: noether.p16.par.017
1. 在論文 “Die allgemeinsten Bereiche aus ganzen transzendenten Zahlen” (\emph{Math. Ann.} 77) 中，第 121 頁自上第 8 行、自上第 9 行和自下第 11 行，域 $(H,Y)$ 應替換為“域 $(H,Y,Y',\ldots,Y^{(\lambda_0-1)})$，其中 $Y',\ldots,Y^{(\lambda_0-1)}$ 表示 $Y$ 關於 $H$ 的共軛。”相應地，第 120 頁自上第 5 行以及第 121 頁自上第 9 行，$(H,y)$ 應替換為 $(H,y,y',\ldots,y^{(\lambda_0-1)})$。

% unit: noether.p16.par.018
2. 在論文 “Gleichungen mit vorgeschriebener Gruppe” (\emph{Math. Ann.} 78) 中，在 Hilbert 不可約性定理的表述裡 --- 第 222 頁自下第 3 行 --- 代替“數域 $\Omega$”，更準確地必須說：“有限代數數域 $\Omega$”；Ostrowski 先生向我指出了這一點。事實上，例如相對於所有代數數的域，任何一個具有代數數係數的方程的群都是恆等群。--- 相應地，在導言中，“任意有理性域”應理解為一個有限代數數域，還可以向其附加有限多個參數，或者理解為這樣產生的中間域之一。

% unit: noether.p16.closing.001
\begin{center}
（1919 年 9 月接受。）
\end{center}
%
\clearpage
%
% unit: noether.p17.title.001
\pdfbookmark[0]{17. 與 W. Schmeidler 合作：非交換域中的模，特別是來自差分表達式的模}{paper-17}
\section*{17. 與 W. Schmeidler 合作：非交換域中的模，特別是來自差分表達式的模}
\begin{center}
\emph{Mathematische Zeitschrift} 8 (1920), pp. 1--35
\end{center}

\begin{center}
\textbf{目錄。}
\end{center}
\begin{center}
\begin{minipage}{0.84\linewidth}
引言。\\
\S~1. 關於微分表達式和差分表達式的形式事項。\\
\S~2. 一般多項式域。\\
\S~3. 單側模及其剩餘群。\\
\S~4. 剩餘群分解為兩個子群的可約性及其與模的關係。\\
\S~5. 用單位進行計算以及分解為 $k$ 個子群的可約性。\\
\S~6. 有限性定理。\\
\S~7. 唯一分解的一個情形。\\
\S~8. 完全可約群到素群中的加法分解。同構定理。\\
\S~9. “同構”與“同類”這兩個概念之間的聯繫。\\
\S~10. 一個群存在無限多個分解。\\
\S~11. 與微分方程組及其積分的關係。\\
\S~12. 例子。
\end{minipage}
\end{center}

% unit: noether.p17.intro.001
\begin{center}\textbf{引言。}\end{center}

% unit: noether.p17.par.001
一元微分表達式的形式理論導向分解為不可約因子，而且這種分解在某種意義上是唯一的。\footnote{參見 E. Landau, \emph{Ein Satz über die Zerlegung homogener linearer Differentialausdrücke in irreduzible Faktoren}, \emph{Journal für die reine und angewandte Mathematik} 124 (1902), pp. 115--120；以及隨後 A. Loewy, \emph{Über reduzible lineare homogene Differentialgleichungen}, \emph{Mathematische Annalen} 56 (1903), pp. 549--584.} 但是，雖然代數中相應的定理可以轉移到若干變量的多項式上，對偏微分表達式卻不再可能。\footnote{參見 H. Blumberg, \emph{Über algebraische Eigenschaften von linearen homogenen Differentialausdrücken}, Göttingen 1912 年博士論文，52 頁，第 51--52 頁，那裡發表了 Landau 的一個反例。} 然而，對於一元微分表達式，除乘積分解之外，還出現了作為最小公倍式的表示，並且類似的規律成立；因為在兩個不同分解中，各個組成部分可以兩兩相配：它們是“同類”的。\footnote{除前引論文外，參見 A. Loewy, \emph{Über vollständig reduzible lineare homogene Differentialgleichungen}, \emph{Mathematische Annalen} 62 (1906), pp. 89--117；下文引作 Loewy。}

% unit: noether.p17.par.002
現在，這個最小公倍式的概念可以解釋為一個模的概念，並因而轉移到 $n$ 個變量；在本文中這就給出了上述定理的自然推廣。除此之外，我們還注意到，對微分表達式實際只用到下列事實：它們可以看作具有非交換乘法的多項式，並且一個乘積的次數等於因子次數之和。這些條件例如也由差分表達式滿足，因此定理也對差分表達式成立；此前它們在那裡同樣只對一個變量為人所知。\footnote{參見 G. Wallenberg--A. Guldberg, \emph{Theorie der linearen Differenzengleichungen}, Leipzig and Berlin 1911.}

% unit: noether.p17.par.003
因此，我們一般地發展一種模理論，其元素是具有非交換乘法的多項式；在這樣做時，我們把模的研究化歸為剩餘類的研究。\footnote{對於交換域的情形，參見 W. Schmeidler, \emph{Über Moduln und Gruppen hyperkomplexer Größen}, \emph{Mathematische Zeitschrift} 3 (1919), pp. 29--42.} 與交換的特殊情形不同，在這裡雖然也可以談剩餘類之和，卻不能談剩餘類之積；只能談一個剩餘類與一個多項式的乘積。因此，“剩餘群”的概念相對於交換情形必須加以修改。在一元多項式的特殊情形中，剩餘群還是“有限”的，即它只具有有限個線性無關剩餘類。一般地說，這樣的有限“階”並不存在；群是“無限”的。不過，對我們的方法來說，這個差別並不重要。

% unit: noether.p17.par.004
如同在交換情形中一樣，把一個模表示為互素模的最小公倍式，對應於剩餘群的加法分解這一較簡單過程。在這裡，剩餘群的同構性顯得本質重要；當專門化到一元微分表達式時，這種同構性變成相關模的同類概念，而當專門化到交換模時，它變成恆等。順便說， 同類的概念也可以轉移到 $n$ 個變量，並可證明它與同構概念一致。剩餘群的加法分解中總只出現有限多個加項，這是通過轉移 Hilbert 的模基定理而得出的；其證明基於 Gordan 的證明方法，而 Hilbert 原來的證明方法則會要求對乘法規律作限制（在公式 (6) 中，右邊為 $a$ 而不是 $a_i$，這對微分表達式成立，但對差分表達式不成立）。

% unit: noether.p17.par.005
為了得到一個唯一性定理，我們必須像 Loewy 所作的那樣，把模專門化為“完全可約”的模，即那些可以表示為素模的最小公倍式的模。給定兩個不同的分解，則各組成部分要麼兩兩相同，要麼至少對一個分解的每個組成部分，都有另一個分解的一個同構組成部分與之對應，反之亦然；同時，每個分解於是至少包含兩個同構的組成部分。此外，從同一個分解中出現兩個同構組成部分，可以推出存在無限多個分解；這甚至在不限制為素模的情況下也成立，而這一點此前即使在一元微分表達式的特殊情形中也未為人所知。

% unit: noether.p17.par.006
最後，對於來自微分表達式的模，我們討論它們與積分的聯繫，並證明：對於有限剩餘群，階等於線性無關積分的最大個數。因此，任意函數恰好在剩餘群不是有限時才出現在通積分中。兩個模同類這一概念，在這裡如在一元情形中一樣，意味着積分之間的一種“雙有理”關係。若干例子結束本文。

% unit: noether.p17.par.007
本文最初的推動來自數年前 E. Landau 向 E. Noether 提出的關於推廣乘積定理的問題。當時 E. Noether 只是注意到這是非交換多項式域中的模理論問題，但尚不能處理。那時已知的 Lasker 模定理不能直接轉移到這個域，因為 Lasker 的證明方法依賴於多項式可唯一分解為不可約因子的定理。\footnote{最近 E. Noether 注意到，Lasker 定理也可以不用這個定理，而用與本文方法相關的方法來建立。} 只有 Schmeidler 的方法提供了轉移的可能性，隨後我們共同實行了這種轉移。具體地說，需要指出：剩餘群的推廣、Hilbert 有限性定理的轉移和應用以及無限多個分解的存在，歸於 E. Noether；Loewy 的完全可約性概念的推廣、同構概念及其與同類概念的聯繫，歸於 W. Schmeidler。

% unit: noether.p17.sec.001
\subsection*{\S~1. 關於微分表達式和差分表達式的形式事項。}

% unit: noether.p17.s01.par.001
1. 對於一元微分表達式，人們已經引入一種符號乘積的構造。設
\[
 a_0(x)\frac{\dd^n y}{\dd x^n}+a_1(x)\frac{\dd^{n-1}y}{\dd x^{n-1}}+\cdots+a_n(x)y=A(y),
\]
\[
 b_0(x)\frac{\dd^m y}{\dd x^m}+b_1(x)\frac{\dd^{m-1}y}{\dd x^{m-1}}+\cdots+b_m(x)y=B(y).
\]
於是把符號乘積 $AB(y)$ 定義為微分表達式
\[
 a_0(x)\frac{\dd^n B(y)}{\dd x^n}+a_1(x)\frac{\dd^{n-1}B(y)}{\dd x^{n-1}}+\cdots+a_n(x)B(y).
\]
這種乘積構造實際上是兩個算子的乘法；最簡單的寫法，是像通常那樣把 $\frac{\dd^n}{\dd x^n}$ 替換為 $\left(\frac{\dd}{\dd x}\right)^n$。於是 $AB$ 變為
\[
 \left(a_0\left(\frac{\dd}{\dd x}\right)^n+\cdots+a_n\right)
 \left(b_0\left(\frac{\dd}{\dd x}\right)^m+\cdots+b_m\right)y,
\]
其中在進行乘法時採用對和與積求導的通常規則。現在，每個這樣的算子都可以看作單個變量 $\frac{\dd}{\dd x}=\xi$ 中的一個多項式；辦法是不顯式寫出算子所作用的函數，並給變量 $\xi$ 固定與乘積求導相應的乘法規則：
\[
 \frac{\dd}{\dd x}\bigl(a(x)\cdot y\bigr)=a(x)\frac{\dd}{\dd x}(y)+a'(x)y,
\]
或者，用我們的記號 $(n=1,m=0)$，
\begin{equation}
 \xi\cdot a=a\cdot\xi+a'.
\tag{1}
\end{equation}

% unit: noether.p17.s01.par.002
對於線性偏微分表達式也完全同樣處理：
\[
 \left(\sum a_{\alpha_1\ldots\alpha_n}(x_1,\ldots,x_n)
 \frac{\p^{\alpha_1+\cdots+\alpha_n}}
 {\p x_1^{\alpha_1}\cdots\p x_n^{\alpha_n}}\right)y
\]
其中 $y$ 是一個函數。對於 $n$ 個算子 $\frac{\p}{\p x_1},\ldots,\frac{\p}{\p x_n}$，我們引入變量 $\xi_1,\ldots,\xi_n$，由此得到多項式
\[
 \sum a_{\alpha_1\ldots\alpha_n}(x_1,\ldots,x_n)
 \xi_1^{\alpha_1}\cdots\xi_n^{\alpha_n}.
\]
兩個多項式的乘積於是按下列規律計算：
\begin{equation}
 \xi_i a=a\xi_i+a_i \qquad (i=1,\ldots,n),
\tag{2}
\end{equation}
其中 $a_i$ 表示 $a$ 關於 $x_i$ 的偏導數；各個 $\xi$ 相互之間的乘積是交換的。除了乘法交換律之外，這樣定義的多項式滿足加法和乘法的一切規律，正如通常代數中一樣。特別地，對於兩個這樣的多項式的乘積，總次數以及在每一個單獨變量中的次數，都等於兩個因子的相應次數之和。

% unit: noether.p17.s01.par.003
2. 對於形如
\[
 a_x^{(0)}y_{x+n}+a_x^{(1)}y_{x+n-1}+\cdots+a_x^{(n)}y_x=A_x,
\]
\[
 b_x^{(0)}y_{x+m}+b_x^{(1)}y_{x+m-1}+\cdots+b_x^{(m)}y_x=B_x
\]
的差分表達式（參見 Wallenberg 前引處），有相應的乘積定義
\[
 (AB)_x=a_x^{(0)}B_{x+n}+a_x^{(1)}B_{x+n-1}+\cdots+a_x^{(n)}B_x,
\]
它可以藉助乘法規則
\[
 (ay)_{x+1}=a_{x+1}y_{x+1}
\]
來實行。如果我們引入算子 $\xi$ 表示從 $x$ 到 $x+1$ 的過渡，那麼得到 $\xi\{ay\}=a^*\xi\{y\}$，其中 $a_x^*=a_{x+1}$。同樣可以略去算子所作用的函數 $y$，並得到規律
\begin{equation}
 \xi a=a^*\xi,
\tag{3}
\end{equation}
其中，隨 $a$ 一起，$a^*$ 也不恆等於零。對於 $n$ 個變量，完全類似地有
\begin{equation}
 \xi_i a=a_i\xi_i\qquad (i=1,\ldots,n),
\tag{4}
\end{equation}
其中 $\xi_i$ 表示從 $x_i$ 到 $x_i+1$ 的過渡，$a_i$ 表示由 $a$ 把 $x_i$ 替換為 $x_i+1$ 所得到的函數；因此，如果 $a$ 不恆等於零，它也不恆等於零。在這些約定下，除交換律外，加法和乘法的一切規則，包括關於乘積次數的規則，又都成立。

% unit: noether.p17.s01.par.004
在下一節中，這兩個情形都將被納入一個以任意係數定義的十分一般的多項式域。

% unit: noether.p17.sec.002
\subsection*{\S~2. 一般多項式域。}

% unit: noether.p17.s02.par.001
我們從一個任意抽象定義的域 $\mathsf{P}$ 出發（其元素用小拉丁字母表示），在其中代數的一切規則都成立，特別是乘法交換律以及唯一取逆也成立。我們在這個域上添加 $n$ 個不定元 $\xi_1,\ldots,\xi_n$；也就是說，我們考慮由乘積 $a\xi_{\nu_1}b\xi_{\nu_2}\ldots c\xi_{\nu_\rho}$ 及其有限和構成的整域 $J$，其中加法和乘法的一切規律都應成立，唯一例外是乘法交換律。特別地，我們把形如
\[
 F(\xi_1,\ldots,\xi_n)=\sum a_{\nu_1\ldots\nu_\rho}\xi_{\nu_1}\cdots\xi_{\nu_\rho}
\]
的有限和稱為“多項式”（本文始終用大拉丁字母表示）。現在，我們用這樣的乘積規律來代替交換律，即使得\emph{兩個多項式的乘積仍然是一個多項式}；因此，$J$ 中所有多項式的域已經窮盡了 $J$。為此，顯然必要且充分的是：對於元素 $\xi_1,\ldots,\xi_n$ 和 $\mathsf{P}$ 的任意元素 $a$，有形如
\begin{equation}
 \xi_i a=F_i(\xi_1,\ldots,\xi_n)\qquad (i=1,\ldots,n)
\tag{5}
\end{equation}
的乘積規則，其中 $F_i(\xi_1,\ldots,\xi_n)$ 表示某個自然依賴於 $\xi_i$ 和 $a$ 的多項式。（特別地，微分表達式和差分表達式的規則 (2) 與 (4) 分別屬於這種形式。）因為 $\xi_i a$ 本身不是多項式，而按假設屬於該域，所以這個等式是必要的。反過來，通過有限次反覆應用公式 (5)，域 $J$ 的每個元素都可以化為一個多項式。

% unit: noether.p17.s02.par.002
對於域的單位，我們記作 1，還要求 $\xi_i\cdot1=1\cdot\xi_i=\xi_i$。我們還規定 $\xi_1,\ldots,\xi_n$ 相互之間的乘法是交換的。\footnote{這隻有從 \S~6 起，為有限性定理及其推論才需要。} 於是 $J$ 中的每個多項式都像通常的 $\xi_1,\ldots,\xi_n$ 上、係數在 $\mathsf{P}$ 中的多項式那樣寫出；兩個多項式只有在所有係數都相同時才相同。

% unit: noether.p17.s02.par.003
更具體地，在下文中\footnote{這同樣只從 \S~6 起才使用。}我們要求用下列規則代替 (5)：
\begin{equation}
 \xi_i a=a_i\xi_i+b_i\qquad (a_i\ne0)\quad(i=1,\ldots,n).
\tag{6}
\end{equation}
這裡也特別包括規律 (2) 和 (4)。於是，一個乘積的次數等於因子次數之和；這對於總次數以及對於每個單獨變量中的次數都成立，因此兩個不恆等於零的多項式的乘積也不同於 0。

% unit: noether.p17.sec.003
\subsection*{\S~3. 單側模及其剩餘群。}

% unit: noether.p17.s03.par.001
由於在我們的域 $J$ 中乘法是非交換的，可以區分下列種類的模\footnote{我們用大德文字母表示模。}:\footnote{單側理想已經由 A. Hurwitz 考慮過；參見 \emph{Vorlesungen über die Zahlentheorie der Quaternionen} (Berlin, Springer 1919)，第六講。不過那裡本質上涉及的是主理想；該書注 1 中所引 Du Pasquier 的著作也是如此。}

% unit: noether.p17.s03.par.002
1. \emph{右側模}，它隨 $M$ 一起也包含任意多項式 $F$ 所給的 $FM$，並且隨 $M_1$ 與 $M_2$ 一起也包含 $M_1+M_2$。

% unit: noether.p17.s03.par.003
2. \emph{左側模}，它隨 $M$ 一起也包含 $MF$，並且隨 $M_1$ 與 $M_2$ 一起也包含 $M_1+M_2$。

% unit: noether.p17.s03.par.004
我們把這兩種統稱為\emph{單側模}；既是右側又是左側的模稱為\emph{雙側模}。下文中我們限於右側模，並指出左側模的理論可以完全以同樣方式建立。

% unit: noether.p17.s03.par.005
若兩個多項式 $F$ 和 $G$ 的差被 $\mM$ 整除，也就是屬於 $\mM$，則稱它們模 $\mM$ 同餘，記作 $F\equiv G(\mM)$。所有與給定多項式同餘的多項式全體構成模 $\mM$ 的一個\emph{剩餘類}。（剩餘類用小德文字母表示。）由於從
\[
 F_1\equiv F_2(\mM)\qquad\text{和}\qquad G_1\equiv G_2(\mM)
\]
推出 $F_1+G_1\equiv F_2+G_2(\mM)$，兩個剩餘類 $\mx$ 與 $\my$ 的和是定義好的，並且仍是剩餘類 $\mx+\my$。相反，由於模是單側的，不能談兩個剩餘類的乘積。事實上，從
\[
 F_1=F_2+M_1,\qquad G_1=G_2+M_2
\]
只推出
\[
 F_1G_1=F_2G_2+F_2M_2+M_1G_2+M_1M_2;
\]
但當 $M_1$ 屬於 $\mM$ 時，$M_1G_2$ 不一定屬於 $\mM$，所以 $F_1G_1$ 與 $F_2G_2$ 一般不必同餘。另一方面，如果 $M_1=0$，因而 $F_1=F_2=F$，則可見 $G_1\equiv G_2$ 蘊含 $FG_1\equiv FG_2$。因此，任意剩餘類 $\my$ 可以從前面乘以一個多項式 $F$，這給出一個良定義的剩餘類 $F\my=\mz$。

% unit: noether.p17.s03.par.006
在剩餘類的計算中，加法的交換律和結合律成立，加法的唯一可逆性規律也成立；此外，對於剩餘類乘以多項式以及加法，分配律 $F(\my_1+\my_2)=F\my_1+F\my_2$ 成立，這直接來自模 $\mM$ 下多項式的計算。

% unit: noether.p17.s03.par.007
單位所屬的剩餘類稱為單位類或單位，記作 $\me$。於是 $F\me=\mx$ 是多項式 $F$ 所屬的剩餘類。與一般剩餘類不同，這裡的乘積 $\mx\me$ 也良定義並且等於 $\mx$；因為從
\[
 F_1=F_2+M_1,\qquad G_1=1+M_2
\]
以及 $M_1\cdot1=M_1$ 可得
\[
 F_1G_1=F_2+M.
\]
模 $\mM$ 的所有剩餘類的全體稱為 $\mM$ 的\emph{剩餘群}。因此，一個剩餘群有下列性質：1. 它隨每個剩餘類 $\mx$ 和任意多項式 $F$ 一起包含剩餘類 $F\mx$；隨剩餘類 $\mx_1$ 和 $\mx_2$ 一起包含剩餘類 $\mx_1+\mx_2$。\footnote{這個性質在更一般的意義上可以稱為剩餘群的“模性質”。} 2. 它有一個單位 $\me$，使得對於每個多項式 $F$，$F\me$ 表示 $F$ 的剩餘類 $\mx$。\emph{因此，當} $F$ \emph{遍歷所有多項式時，} $F\me$ \emph{遍歷整個剩餘群。}

% unit: noether.p17.sec.004
\subsection*{\S~4. 剩餘群分解為兩個子群的可約性及其與模的關係。}

% unit: noether.p17.s04.par.001
對於單側模，\emph{整除性}、\emph{最大公約模}和\emph{最小公倍模}這些概念仍如雙側情形一樣，這從下列定義可以看出。

% unit: noether.p17.s04.par.002
若由 $M\equiv0(\mM)$ 也推出 $M\equiv0(\mN)$，則稱模 $\mM$ 被另一個模 $\mN$ 整除。

% unit: noether.p17.s04.par.003
對於兩個模 $\mM$ 和 $\mN$，\emph{最大公約模} $(\mM,\mN)$ 定義為所有可表示成 $M+N$ 的多項式全體，其中 $M\equiv0(\mM)$ 且 $N\equiv0(\mN)$；它仍然是一個模。多個模的情形也一樣。

% unit: noether.p17.s04.par.004
所有同時被 $\mM$ 和 $\mN$ 整除的多項式全體形成一個模 $[\mM,\mN]$，即 $\mM$ 與 $\mN$ 的\emph{最小公倍模}。這個定義也相應地適用於多個模，甚至適用於任意基數的模集合。

% unit: noether.p17.s04.par.005
若兩個模的最大公約模為單位模，也就是全體多項式，則稱這兩個模互素。

% unit: noether.p17.s04.par.006
現在我們從一個模 $\mM$ 出發，它可以表示為兩個互素模 $\mN$ 和 $\mL$ 的最小公倍模，並且這兩個模都假定為真因子：
\begin{equation}
 \mM=[\mN,\mL].
\tag{7}
\end{equation}
我們考察 $\mM$ 的剩餘群 $\mG$。按假設，存在兩個多項式 $N\equiv0(\mN)$ 和 $L\equiv0(\mL)$，使得
\begin{equation}
 1\equiv N+L(\mM).
\tag{8}
\end{equation}
由此我們證明
\begin{equation}
 \mL=(\mM,L),\qquad \mN=(\mM,N),
\tag{9}
\end{equation}
其中例如 $(\mM,L)$ 表示由所有形如 $M+FL$ 的多項式組成的模。事實上，首先有 $(\mM,L)\equiv0(\mL)$。反過來，由 (8) 對任意多項式 $F$ 推出
\begin{equation}
 F\equiv FN(\mL).
\tag{10}
\end{equation}
現在若 $F\equiv0(\mL)$，則 $FN\equiv0(\mL)$，因而 $\equiv0(\mM)$，於是 $F\equiv FL(\mM)$。同樣可證明 $\mN=(\mM,N)$。還順便得到，由於 $\mN$ 與 $\mL$ 是 $\mM$ 的真因子，既不能有 $L\equiv0$，也不能有 $N\equiv0(\mM)$。

% unit: noether.p17.s04.par.007
設 $\ma$ 為 $N$ 模 $\mM$ 的剩餘類，$\mb$ 為 $L$ 模 $\mM$ 的剩餘類。於是由 (8) 得
\[
 \me=\ma+\mb\qquad(\ma\ne0,\ \mb\ne0).
\]
由於每個剩餘類都是 $F\me$ 的形式，分配律對每個剩餘類給出
\begin{equation}
 F\me=F\ma+F\mb\qquad(\ma\ne0,\ \mb\ne0).
\tag{11}
\end{equation}
然而，把任意剩餘類表示為兩個形如 $G\ma$ 與 $H\mb$ 的剩餘類之和，是\emph{唯一}的。因為從關係 $0=G\ma+H\mb$ 出發，若 $K$ 是剩餘類 $G\ma=-H\mb$ 中的一個多項式，則顯然 $K\equiv GN(\mM)$，$K\equiv-HL(\mM)$，故 $K\equiv0(\mN)$，$K\equiv0(\mL)$，因而 $K\equiv0(\mM)$；也就是說，$G\ma=0$，$H\mb=0$。

% unit: noether.p17.s04.par.008
反過來，現在假設關係 (11) 對每個多項式 $F$ 成立，並且在上述意義下唯一。令 $N$ 為 $\ma$ 中的一個多項式，$L$ 為 $\mb$ 中的一個多項式。我們構造兩個模 (9)，並斷言 (7) 和 (8)。後者由 (11) 在 $F=1$ 時得出。此外，按 $\mN$ 與 $\mL$ 的定義，$\mM$ 是這兩個模的公倍模，而且是最小者。因為由 $F\equiv0(\mL)$ 與 $F\equiv0(\mN)$ 得 $F\equiv HL\equiv GN(\mM)$；故 $G\ma=H\mb$，$G\ma-H\mb=0$，於是由表示的唯一性，$G\ma=0$，$H\mb=0$，$F\equiv0(\mM)$。因此已經證明，\emph{把一個模} $\mM$ \emph{表示為兩個互素模的最小公倍模，等價於} $\mM$ \emph{的剩餘群} $\mG$ \emph{按公式} (11) \emph{給出的唯一加法分解。}

% unit: noether.p17.s04.par.009
現在我們考察由形如 $F\ma$ 的剩餘類組成的系統 $\mA$，並希望把它同 $\mL$ 的剩餘群 $\mbarA$ 聯繫起來。按 (11)，若 $\bara$ 是 $N$ 模 $\mL$ 的剩餘類，並且因而是 $\mbarA$ 的單位類，則 $\mL$ 中每個多項式 $F$ 屬於剩餘類 $F\ma$。若現在把 $\mA$ 中的剩餘類 $F\ma$ 與 $\mbarA$ 中的剩餘類 $F\bara$ 對應起來，那麼 $\mA$ 的每個剩餘類都被賦予 $\mbarA$ 的一個剩餘類，使得 $\mA$ 中兩個剩餘類之和對應於 $\mbarA$ 中被賦予的剩餘類之和，而 $\mA$ 中一個剩餘類與一個多項式之積對應於 $\mbarA$ 中被賦予的剩餘類與同一個多項式之積。這個對應還是一一的：由 $F\ma=0$ 按 (11) 推出 $F\equiv0(\mL)$，故 $F\bara=0$；反之，$F\bara=0$，故 $F\equiv0(\mathfrak L)$，也按 (10) 說明 $FN\equiv0(\mL)$，故 $FN\equiv0(\mM)$，也就是 $F\ma=0$。

% unit: noether.p17.s04.par.010
這把我們引向一個基本概念，定義如下：

% unit: noether.p17.def.001
\emph{兩個剩餘類系統 $\mA$ 和 $\mbarA$ 之間的一一對應稱為同構的，如果 $\mA$ 中兩個剩餘類之和被賦予 $\mbarA$ 中相應剩餘類之和，並且 $\mA$ 中任意剩餘類與一個多項式之積被賦予 $\mbarA$ 中相應剩餘類與同一個多項式之積。}

% unit: noether.p17.s04.par.011
$\mG$ 中一個剩餘類子系統 $\mA$，如果同構於某個模的剩餘群，就稱為 $\mG$ 的一個\emph{子群}。如上面的證明所示，這裡的賦值更具體地是這樣的：$\bar{\mathfrak A}$ 的每個剩餘類都由 $\mbarA$ 的相應剩餘類通過添入某些多項式而產生，即添入所有屬於 $\mL$ 但不屬於 $\mM$ 的多項式。

% unit: noether.p17.s04.par.012
類似地可證明，所有形如 $F\mb$ 的剩餘類全體構成 $\mG$ 的一個子群 $\mB$，它同構於 $\mN$ 的剩餘群 $\mbarB$。

% unit: noether.p17.s04.par.013
若一個剩餘群 $\mG$ 的剩餘類允許唯一的加法分解 (11)，則稱它可約為兩個子群 $\mA$ 和 $\mB$，記作 $\mG=\mA+\mB$；不可約的群稱為\emph{不可約}。我們有時也把相應的模稱為可約或不可約。於是有下列結果。

% unit: noether.p17.thm.001
\textbf{定理 1.}\footnote{參見 Schmeidler 前引處，第 39 頁。} \emph{模 $\mM$ 可以表示為兩個互素真因子 $\mL$ 與 $\mN$ 的最小公倍模，當且僅當它的剩餘群 $\mG$ 可約為兩個子群 $\mA$ 和 $\mB$。在這種情形下，$\mA$ 和 $\mB$ 分別同構於 $\mL$ 和 $\mN$ 的剩餘群。}
%
%
% unit: noether.p17.s05.title
\subsection*{\S~5. 用單位計算以及可約為 $k$ 個子群.}

% unit: noether.p17.s05.par.001
$\mG$ 的剩餘類 $\ma$ 和 $\mb$ 與單位 $\me$ 具有共同的性質：$\ma$ 和 $\mb$ 同任意剩餘類 $\mx$ 的乘積都有定義；因此我們也稱 $\ma$ 和 $\mb$ 為 $\mG$ 的單位。事實上，由 (11)，令 $F=M$，得到
% unit: noether.p17.s05.eq.001
\[
 0=M\me=M\ma+M\mb,
\]
從而由唯一性得 $M\ma=0$、$M\mb=0$，於是
% unit: noether.p17.s05.eq.002
\[
 (F+M)\ma=F\ma=\mx\ma,
\]
其中 $\mx$ 表示 $F$ 的剩餘類。因此，我們可以用等價的類關係
% unit: noether.p17.s05.eq.003
\[
 \mx\me=\mx\ma+\mx\mb
\]
來代替關係 (11)。特別地，令 $\mx=\ma$，得到
% unit: noether.p17.s05.eq.004
\[
 \ma=\ma^2+\ma\mb,
\]
因此由唯一性得 $\ma^2=\ma$、$\ma\mb=0$。同樣得到 $\mb^2=\mb$、$\mb\ma=0$。

% unit: noether.p17.s05.par.002
對應於把一個剩餘群表示為兩個組成部分之和，我們現在也可以定義把它表示為 $k$ 個組成部分之和。也就是說，設給定一個唯一表示
% unit: noether.p17.s05.eq.005
\[
 F\me=F\ma_1+F\ma_2+\cdots+F\ma_k\qquad(\ma_1\ne0,\ldots,\ma_k\ne0),
\]
其中從 $0=G_1\ma_1+\cdots+G_k\ma_k$ 也推出 $G_i\ma_i=0$。於是特別地，對 $F=M$ 得到關係 $M\ma_i=0$，因此該表示也可以寫成形式\footnote{對於剩餘群的加法分解，交換律和結合律成立；但是，與可交換情形相反，$\mU_1+\mU_2=\mU_1+\mU_3$ 並不推出 $\mU_2=\mU_3$；參見 \S~12 中的例 1。}
% unit: noether.p17.s05.eq.006
\begin{equation}
 \mx\me=\mx\ma_1+\cdots+\mx\ma_k\qquad(\ma_1\ne0,\ldots,\ma_k\ne0).
\tag{12}
\end{equation}
這種表示也可以看作由逐次分解為兩個組成部分而產生。如果這種過程可以無限繼續，並且所有 $\ma_i\ne0$，則稱為無限多個組成部分之和。

% unit: noether.p17.s05.par.003
在 (12) 中令 $\mx=\ma_i$，得到
% unit: noether.p17.s05.eq.007
\begin{equation}
 \ma_i^2=\ma_i,
 \qquad \ma_i\ma_j=0\quad(i\ne j),
 \qquad i,j=1,\ldots,k,
\tag{13}
\end{equation}
且 $\ma_i\ne0$。反過來，設 $\mG$ 中有一個剩餘類系統 $\ma_1\ldots\ma_k$，它滿足正交條件 (13)，並且滿足
% unit: noether.p17.s05.eq.008
\begin{equation}
 \me=\ma_1+\cdots+\ma_k.
\tag{14}
\end{equation}
我們將證明，這給出了把群表示為 $k$ 個組成部分之和。首先，由假設，對於來自 $\ma_i$ 的每個多項式 $A_i$，因為 $A_i\ma_i=(A_i+M)\ma_i$，所以推出 $M\ma_i=0$；因而 $\ma_i$ 可以同每個類相乘。於是 (14) 蘊含表示 (12)，而這個表示是唯一的。因為若 $0=\mx_1\ma_1+\cdots+\mx_k\ma_k$，則右乘 $\ma_i$ 並用 (13) 得 $0=\mx_i\ma_i$，這就證明了唯一性，從而證明了斷言。單位類分解中出現的剩餘類 $\ma_1\ldots\ma_k$ 稱為屬於該分解的單位；等價地，它們由正交關係 (13) 和條件 (14) 來刻畫。後面的發展本質上將歸結為用單位進行計算。

% unit: noether.p17.s05.par.004
還需要建立群分解為 $k$ 個組成部分的加法分解與相應的模分解之間的聯繫。和 $k=2$ 的情形一樣，通過歸納將得到把模表示為 $k$ 個互素模的最小公倍。

% unit: noether.p17.s05.par.005
設給定分解 (12)，我們把它寫成
% unit: noether.p17.s05.eq.009
\begin{align}
 \mx\me&=\mx\mb+\mx\ma_k,\tag{15}\\
 \mx\mb&=\mx\ma_1+\cdots+\mx\ma_{k-1}.\tag{16}
\end{align}
由 (15) 和定理 I 得知，剩餘類 $\mx\mb$ 與模
% unit: noether.p17.s05.eq.010
\begin{equation}
 \mN=(\mM,A_k)
\tag{17}
\end{equation}
的剩餘群同構。因此，對於這些剩餘類 $\mx\overline{\mb}$，與 (16) 相應的表示也成立：
% unit: noether.p17.s05.eq.011
\begin{equation}
 \mx\overline{\mb}=\mx\overline{\ma}_1+\cdots+\mx\overline{\ma}_{n-1}.
\tag{18}
\end{equation}
現在假設已經證明
% unit: noether.p17.s05.eq.012
\begin{equation}
 \mN=[\mM_1,\ldots,\mM_{k-1}],
\tag{19}
\end{equation}
其中
% unit: noether.p17.s05.eq.013
\begin{equation}
\left\{
\begin{aligned}
 \mM_1&=(\mN,A_2+\cdots+A_{k-1}),\\
 \mM_2&=(\mN,A_1+A_3+\cdots+A_{k-1}),\\
 &\hspace{1.8cm}\vdots\\
 \mM_{k-1}&=(\mN,A_1+\cdots+A_{k-2}).
\end{aligned}\right.
\tag{20}
\end{equation}
對於兩個加項，這個斷言與我們先前的結果一致。此外，這裡的多項式 $A_i$ 可以特別選為來自 $\overline{\ma}_i$ 且同時包含在 $\ma_i$ 中的多項式，因為按 \S~4，後一剩餘類的多項式包含在前一剩餘類中。於是，由 (15)，
% unit: noether.p17.s05.eq.014
\begin{equation}
 \mM=[\mN,\mM_k],
\tag{21}
\end{equation}
其中
% unit: noether.p17.s05.eq.015
\[
 \mM_k=(\mM,A_1+\cdots+A_{k-1}).
\]
由 (19) 和 (21) 按最小公倍的定義得到
% unit: noether.p17.s05.eq.016
\[
 \mM=[[\mM_1,
 \ldots,\mM_{k-1}],\mM_k]=[\mM_1,\ldots,\mM_{k-1},\mM_k],
\]
並且由 (17) 和 (20) 得到
% unit: noether.p17.s05.eq.017
\[
 \mM_1=(\mM,A_k,A_2+\cdots+A_{k-1}).
\]
於是
% unit: noether.p17.s05.eq.018
\[
 \mM_1=(\mM,A_2+\cdots+A_k),
\]
因為後一模由 (13) 還包含多項式
% unit: noether.p17.s05.eq.019
\[
 A_k\equiv A_k(A_2+\cdots+A_k)(\mM).
\]
對於 $\mM_2,\ldots,\mM_{k-1}$ 得到相應的表達式。並且，模 $\mM_1,\ldots,\mM_k$ 是互素的，因為由 (14)
% unit: noether.p17.s05.eq.020
\[
 \frac{1}{k-1}\Bigl((A_2+\cdots+A_k)+(A_1+A_3+\cdots+A_k)+\cdots+(A_1+\cdots+A_{k-1})\Bigr)\equiv1\pmod{\mM}.
\]
它們還形成一個完全互素系統，也就是說，其中任意若干個的最小公倍與其餘各個的最小公倍互素。\footnote{這是 Loewy 引入的一個概念的推廣。} 這可直接從公式 (12) 中把加項分成兩個適當的組而得到。

% unit: noether.p17.s05.par.006
這樣已經證明：剩餘群分解為 $k$ 個加項，對應於把模表示為 $k$ 個模的最小公倍，而這些模形成一個完全互素系統。現在我們也用歸納證明逆命題。

% unit: noether.p17.s05.par.007
設
% unit: noether.p17.s05.eq.021
\[
 \mM=[\mM_1,\ldots,\mM_k]=[[\mM_1,\ldots,\mM_{k-1}],\mM_k]=[\mN,\mM_k],
\]
其中 $\mM_1,\ldots,\mM_k$ 形成完全互素系統。結合定理 I，這給出剩餘群的表示 (15)。現在 $\mM_1,\ldots,\mM_{k-1}$ 也形成完全互素系統。因為若對這個系統的模作某種分組時，一組的最小公倍與另一組的最小公倍有一個公共因子，那麼在把模 $\mM_k$ 加到兩組中的某一組並再取最小公倍時，這個因子仍會保留。\footnote{繼續這個推理可知，在 $\mM_1,\ldots,\mM_k$ 中任意兩個都是互素的；其逆命題，即由此推出 $\mM_1,\ldots,\mM_k$ 形成完全互素系統，一般只在可交換情形中成立；參見 \S~12 中的反例 1。} 對於 $k=2$，“完全互素”與“互素”含義相同；因此我們的斷言可以假定已對 $k-1$ 成立。於是對於剩餘類 $\mx\overline{\mb}$，唯一表示 (18) 成立，從而由於同構，(16) 也成立，於是斷言得證。

% unit: noether.p17.s05.thm.002
\textbf{定理 II.} \emph{一個模 $\mM$ 能表示為 $k$ 個模 $\mM_1,\ldots,\mM_k$ 的最小公倍，且這些模形成完全互素系統，當且僅當它的剩餘群 $\mG$ 可約為 $k$ 個子群 $\mU_1,\ldots,\mU_k$。在適當記號下，$\mU_i$ 同構於 $\mM_i$ 的剩餘群。}

% unit: noether.p17.s06.title
\subsection*{\S~6. 有限性定理.}

% unit: noether.p17.s06.par.001
本節將證明，每個剩餘群都可以表示為有限多個不可約組成部分之和。為此，我們先轉移 Hilbert 的模基定理。

% unit: noether.p17.s06.par.002
從現在起，我們使用 \S~2 中表述的較特殊的假設：$\xi_1,\ldots,\xi_n$ 彼此之間的乘法是可交換的，並且多項式的乘法法則具有形式
% unit: noether.p17.s06.eq.001
\begin{equation}
 \xi_i a=a_i\xi_i+b_i,
 \qquad (a_i\ne0)\quad(i=1,\ldots,n).
\tag{6}
\end{equation}
只需對模證明即可。因為若 $\mS$ 是任意多項式系統，我們從它導出模 $\mM$，方法是把所有可由 $\mS$ 中有限多個多項式 $G_1,\ldots,G_k$ 以形式 $A_1G_1+\cdots+A_kG_k$ 組合出來的多項式都加入 $\mS$。如果 $F_1,\ldots,F_\varrho$ 是 $\mM$ 的一個模基，並且
% unit: noether.p17.s06.eq.002
\[
 F_i=\sum_j A_{ij}G_j\qquad(i=1,\ldots,l),
\]
那麼在這些表示中出現的有限多個多項式 $G_j$ 構成 $\mS$ 的一個基。

% unit: noether.p17.s06.par.003
我們按照 Gordan（\emph{Göttinger Nachrichten} 1899）證明模基定理。證明分為兩部分。第一部分涉及由無限多個冪積組成的系統 $\mP$，由於 $\xi_1,\ldots,\xi_n$ 的可交換性，它基本上保持不變。冪積
% unit: noether.p17.s06.eq.003
\[
 P=\xi_1^{a_1}\cdots\xi_n^{a_n},
\]
其中每一個只在系統中寫一次，要按次數遞增排列，並且在次數相同時按某個固定方式排列。現在考慮 $\mP$ 的子系統 $\mQ$，它由所有不被任何前面的 $P$ 整除的冪積 $P=Q$ 組成。於是沒有任何 $Q$ 被後面的 $Q$ 整除，因為後者或者次數更高，或者在同一次數下不同於前者。另一方面，$\mP$ 的每個 $P$ 至少被某個 $Q$ 整除。否則，設 $P_m$ 是第一個不被任何 $Q$ 整除的 $P$。那麼 $P_m$ 本身不能是 $Q$，所以它被至少一個前面的 $P$ 整除；按假設，這個前面的 $P$ 又被某個 $Q$ 整除，於是 $P_m$ 本身也被一個 $Q$ 整除，矛盾。

% unit: noether.p17.s06.par.004
現在斷言 $Q$ 的個數有限。事實上，令
% unit: noether.p17.s06.eq.004
\[
 Q_0=\xi_1^{h_1}\cdots\xi_n^{h_n}
\]
為第一個 $Q$。由於任意 $Q_\rho=\xi_1^{h_{1\lambda}}\cdots\xi_n^{h_{n\lambda}}$ 不被 $Q_0$ 整除，對於每個 $Q$ 必有至少一個 $h_{\rho\lambda}<h_\lambda$；令 $\lambda$ 表示滿足這一點的最小指標。我們把所有使 $h_{\rho\lambda}$ 具有同一固定值 $c_\lambda$ 的 $Q$ 合成一個類 $K(c_\lambda)$。因為 $c_\lambda<h_\lambda$，這樣的類只有有限多個。但每個類只含有限多個元素。因為如果把類 $K(c_\lambda)$ 中的一個 $Q$ 寫成 $\xi_\lambda^{c_\lambda}Q'$，那麼這些 $Q'$ 也構成一個冪積系統，其中沒有一個被另一個整除，而且它們只含有 $n-1$ 個變量。對於 $n=2$，$Q'$ 只含一個變量；其個數於是為 $=1$，因而有限，所以也可以假定 $n-1$ 個變量的情形中個數有限。斷言得證。

% unit: noether.p17.s06.par.005
第二，令 $\mM$ 是一個多項式模，其項按字典序排列。令 $\mP$ 為出現在最高項中的冪積系統，每個冪積只寫一次。於是對系統 $\mathfrak{Q}=(Q_1,\ldots,Q_\varrho)$ 對應一個由有限多個多項式 $F_1,\ldots,F_\varrho$ 組成的系統：給每個 $Q_i$ 指派一個來自 $\mM$ 的多項式，其最高項等於 $b_iQ_i$，所以 $F_i=b_iQ_i+$，後面是較低項。

% unit: noether.p17.s06.par.006
現在設 $F=bP+\cdots$ 是 $\mM$ 中的任意多項式，其中 $P=\xi_1^{\beta_1}\cdots\xi_n^{\beta_n}Q_i$。形成乘積
% unit: noether.p17.s06.eq.005
\[
 \xi_1^{\beta_1}\cdots\xi_n^{\beta_n}F_i
 =\xi_1^{\beta_1}\cdots\xi_n^{\beta_n}(b_iQ_i+\cdots)
 =c\,\xi_1^{\beta_1}\cdots\xi_n^{\beta_n}Q_i+\cdots
 =cP+\cdots.
\]
這裡由 (6) 有 $c\ne0$，而所有後續項都是字典序中的較低項，因為就指數而言，乘法在略去較低項後是可交換的。\footnote{由此可見，代替法則 (6)，也可以使用稍弱的限制性法則
\[
 \xi_i a=a_i\xi_i+a_{i,i+1}\xi_{i+1}+\cdots+a_{in}\xi_n+b_i\qquad(a_i\ne0)
\]
（在某個固定的 $\xi$ 編號下）。} 差
% unit: noether.p17.s06.eq.006
\[
 F-\frac bc\,\xi_1^{\beta_1}\cdots\xi_n^{\beta_n}F_i
\]
於是只含較低項，並且仍屬於 $\mM$。這個過程經過有限步即終止，因此 $F_1,\ldots,F_\varrho$ 被認出為一個模基。

% unit: noether.p17.s06.thm.003
\textbf{定理 III.} \emph{任意由無限多個多項式組成的系統 $\mS$ 都具有一個有限模基 $G_1,\ldots,G_k$，使得 $\mS$ 中每個多項式 $G$ 都有表示 $G=A_1G_1+\cdots+A_kG_k$。}

% unit: noether.p17.s06.par.007
由此立即得到

% unit: noether.p17.s06.cor.001
\textbf{推論.} \emph{任意由無限多個模 $\mM$ 的剩餘類組成的系統 $\mS$ 都有一個基 $\mx_1,\ldots,\mx_k$，使得 $\mS$ 中每個 $\mx$ 都能表示為 $\mx=A_1\mx_1+\cdots+A_k\mx_k$。}

% unit: noether.p17.s06.par.008
實際上，只需給每個剩餘類指定一個代表，並考慮這些代表同 $\mM$ 中所有多項式構成的系統。

% unit: noether.p17.s06.par.009
現在轉向證明以下定理。

% unit: noether.p17.s06.thm.004
\textbf{定理 IV.} \emph{每個剩餘群都可表示為有限多個不可約組成部分之和；換言之，由定理 II，每個模都可表示為有限多個不可約模的最小公倍，而這些模形成一個完全互素系統。}

% unit: noether.p17.s06.par.010
反設 $\mG$ 是無限多個組成部分 $\mU_1,\mU_2,\ldots$ 之和（參見 \S~5），並令 $\ma_1,\ma_2,\ldots$ 為相應的單位，它們按定義都非零。我們形成由所有這些單位組成的系統 $\mS$；由定理 III 的推論，它有有限基 $\ma_{i_1},\ldots,\ma_{i_v}$，其中 $i_1<\cdots<i_v$。令 $\mu>i_v$。則
% unit: noether.p17.s06.eq.007
\[
 \ma_\mu=\mx_1\ma_{i_1}+\cdots+\mx_v\ma_{i_v}.
\]
這是單位之間的一個線性關係，而這被定義排除。因此，在剩餘群的每個加法分解中，經過有限步後就會出現一個不可約組成部分；由於分解的交換性，我們可以把它放在開頭。所以不失一般性，可以假設 $\mU_i$ 本身就是不可約的。但由上述論證，只可能出現有限多個組成部分 $\mU_i$。斷言得證。

% unit: noether.p17.s07.title
\subsection*{\S~7. 分解唯一的一種情形.}

% unit: noether.p17.s07.par.001
設
% unit: noether.p17.s07.eq.001
\[
 \mG=\mU_1+\cdots+\mU_k=\mB_1+\cdots+\mB_l
\]
是一個剩餘群，其組成部分 $\mU_1,\ldots,\mU_k$、$\mB_1,\ldots,\mB_l$ 都不可約。於是有
% unit: noether.p17.s07.eq.002
\begin{equation}
 \mx\me=\mx\ma_1+\cdots+\mx\ma_k
 =\mx\mb_1+\cdots+\mx\mb_l,
\tag{22}
\end{equation}
其中 $\ma_\chi$ 是與 $\mU_\chi$ 的單位類同構對應的 $\mG$ 的剩餘類，$\mb_\lambda$ 是與 $\mB_\lambda$ 的單位類同構對應的剩餘類。現在把 $\mx$ 換成 $\mx\mb_\lambda$，其中 $\lambda$ 是 $1,\ldots,l$ 中的一個固定指標，就得到
% unit: noether.p17.s07.eq.003
\begin{equation}
 \mx\mb_\lambda=\mx\mb_\lambda\ma_1+\cdots+\mx\mb_\lambda\ma_k.
\tag{23}
\end{equation}
這是群 $\mB_\lambda$ 的唯一表示，但它不是 $\mB_\lambda$ 的加法分解，因為右邊各個剩餘類未必屬於 $\mB_\lambda$。由於 $\mb_\lambda$ 不為零，所有乘積 $\mb_\lambda\ma_\chi$ 不可能都為零。於是，不失一般性，設
% unit: noether.p17.s07.eq.004
\[
 \mb_\lambda\ma_1\ne0,
 \ldots,
 \mb_\lambda\ma_\rho\ne0,
 \qquad
 \mb_\lambda\ma_{\rho+1}=0,
 \ldots,
 \mb_\lambda\ma_k=0.
\]
則
% unit: noether.p17.s07.eq.005
\[
 \mx\mb_\lambda=\mx\mb_\lambda\ma_1+\cdots+\mx\mb_\lambda\ma_\rho.
\]

% unit: noether.p17.s07.par.002
對 $1,\ldots,\rho$ 中固定的 $\chi$，考慮那些非零的並且還滿足 $\mx\mb_\lambda\ma_\chi\ne0$ 的剩餘類 $\mx\mb_\lambda$ 的全體。兩個系統 $\mx\mb_\lambda$ 和 $\mx\mb_\lambda\ma_\chi$ 是同構的。首先，這個對應是一一的：由 $\mx\mb_\lambda\ma_\chi=0$ 按假設推出 $\mx\mb_\lambda=0$，而由後者按表示 (23) 的唯一性又推出 $\mx\mb_\lambda\ma_\chi=0$。此外，和對應於和，同任意多項式的乘積對應於相應的乘積。固定 $\chi$ 時，對於每個滿足 $\mb_\lambda\ma_\chi\ne0$ 的 $\lambda$，都可以作同樣的考察。對於每個 $\ma_\chi$，至少必須有一個這樣的 $\mb_\lambda$，因為 $\ma_\chi=(\mb_1+\cdots+\mb_l)\ma_\chi\ne0$。

% unit: noether.p17.s07.par.003
我們先假定，對於每個 $\chi$ 只有一個 $\lambda$，而對於每個 $\lambda$ 也只有一個 $\chi$，使得 $\mb_\lambda\ma_\chi\ne0$。

% unit: noether.p17.s07.par.004
由這種對應得到 $k=l$；並且可以選擇記號，使得 $\mb_\chi\ma_\chi\ne0$ 且當 $\lambda\ne\chi$ 時 $\mb_\lambda\ma_\chi=0$。這時我們說乘積 $\mb_\lambda\ma_\chi$ 形成一個對角圖式。在這種情形下，由 (23) 得單位為 $\mb_\chi=\mb_\chi\ma_\chi$，並且因為
% unit: noether.p17.s07.eq.006
\begin{equation}
 \ma_\chi=\mb_1\ma_\chi+\cdots+\mb_l\ma_\chi=\mb_\chi\ma_\chi=\mb_\chi,
\tag{24}
\end{equation}
一般有 $\mU_\chi=\mB_\chi$。因此分解是唯一的，而且乘積 $\ma_\chi\mb_\lambda$ 形成一個對角圖式。

% unit: noether.p17.s07.par.005
如果，例如，交換律專門對剩餘類 $\ma_\chi\mb_\lambda$ 成立，那麼上面的假設總是滿足。因為在 (23) 中，每個組成部分 $\mx\mb_\lambda\ma_\chi=\mx\ma_\chi\mb_\lambda$ 此時都包含在 $\mB_\lambda$ 中，所以 (23) 給出群 $\mB_\lambda$ 的一個加法分解。由於 $\mB_\lambda$ 不可約，所以除一個外所有 $\mb_\lambda\ma_\chi$ 都為零。又由 (24)，相應地，對固定的 $\chi$ 只有一個 $\mb_\lambda\ma_\chi$ 不同於零；否則 (24) 由於 $\ma_\chi$ 和 $\mb_\lambda$ 的可交換性，會給出 $\mU_\chi$ 的一個加法分解，而 $\mU_\chi$ 應當是不可約的。\footnote{在 Schmeidler 前引處處理的非可交換雙側情形中，同樣成立；在那裡，一個子群的每個剩餘類同另一子群的每個剩餘類的乘積都消失。於是由於 $\mb_\lambda\ma_\chi=\mb_\lambda\ma_\chi\mb_1+\cdots+\mb_\lambda\ma_\chi\mb_l$，且當 $\lambda\ne\mu$ 時 $\mb_\lambda(\ma_\chi\mb_\mu)=0$，也有 $\mb_\lambda\ma_\chi=\mb_\lambda\ma_\chi\mb_\lambda$，所以這個類包含在 $\mB_\lambda$ 中；由 $\mB_\lambda$ 的不可約性推出，只有一個、因而正好一個 $\chi$ 使得乘積 $\mb_\lambda\ma_\chi\ne0$。同樣可證，對每個 $\ma_\chi$，存在且只存在一個 $\mb_\mu$ 使 $\ma_\chi\mb_\mu\ne0$，從而 $k=l$。於是 $\mb_\lambda\ma_\chi$ 形成對角圖式，這證明了那裡所得的唯一性。}

% unit: noether.p17.s07.par.006
在以下情形中也有同樣的結論：條件 $\mb_\lambda\ma_\lambda\ne0$、當 $\lambda\ne\chi$ 時 $\mb_\lambda\ma_\chi=0$、$\chi=1,\ldots,k$，以及同樣地，對每個 $\mu=1,\ldots,l$，$\mb_\mu\ma_\chi=0$、$\mb_\mu\ma_\mu\ne0$、$\chi=1,\ldots,\sigma$，並非對所有 $\mb_\lambda$，而只對其中一部分（$\lambda=1,\ldots,\sigma$）滿足。可得 $\mU_\chi=\mB_\chi$（$\chi=1,\ldots,\sigma$）；並且若設
% unit: noether.p17.s07.eq.007
\[
 \mU_{\sigma+1}+\cdots+\mU_k=\mU,
 \qquad
 \mB_{\sigma+1}+\cdots+\mB_l=\mB,
\]
則對於 $\mU$ 和 $\mB$ 的單位 $\ma$ 和 $\mb$，也滿足條件 $\mb\ma_\chi=0$、$\mb_\chi\ma=0$（$\chi=1,\ldots,\sigma$）以及 $\mb\ma\ne0$，由此推出 $\mU=\mB$。
%
%
% unit: noether.p17.sec.008
\subsection*{\S~8. 完全可約群加法分解為素群。同構定理.}

% unit: noether.p17.par.173
現在我們考察另一種情形：在其中雖然沒有唯一性，卻有不同分解之間的同構性。這是 Loewy 所考察情形的推廣。

% unit: noether.p17.def.008
\textbf{定義.} \emph{如果一個模除自身和包含所有多項式的單位模之外沒有別的因子，則稱它為素模。素模的剩餘群稱為素群。}

% unit: noether.p17.par.174
因此每一個素群 \emph{a fortiori} 是不可約的。也存在無限素群（參見 \S~11 的定義和 \S~12 的例 2）。

% unit: noether.p17.par.175
現在設給定的剩餘群 $\mG$ 可以以兩種方式表示為有限多個素群 $\mU_1,\ldots,\mU_k$ 以及 $\mB_1,\ldots,\mB_l$ 之和。一個剩餘群只要至少允許一種這樣的素群和表示，並且相應的模 $\mM$，就按照 Loewy 的術語稱為完全可約的。

% unit: noether.p17.par.176
我們考察乘積 $\ma_\varkappa\mb_\lambda$（$\varkappa=1,\ldots,k$，$\lambda=1,\ldots,l$），並要證明：或者 $\ma_\chi\mb_\lambda=0$，或者當 $\mx$ 遍歷 $\mG$ 的所有剩餘類時，$\mx\ma_\chi\mb_\lambda$ 遍歷整個群 $\mB_\lambda$。事實上，屬於 $\mB_\lambda$ 的模是
% unit: noether.p17.eq.025a
\[
 (\mM,\mb_1+\cdots+\mb_{\lambda-1}+\mb_{\lambda+1}+\cdots+\mb_l)=\mM_{\mb_\lambda};
\]
它由 $\mM$ 加入剩餘類 $\mb_1+\cdots+\mb_{\lambda-1}+\mb_{\lambda+1}+\cdots+\mb_l$ 的所有多項式而得到。這個模有因子
% unit: noether.p17.eq.025b
\[
 (\mM,\mb_1+\cdots+\mb_{\lambda-1}+\mb_{\lambda+1}+\cdots+\mb_l,
 \ma_\varkappa\mb_\lambda).
\]
但由於 $\mM_{\mb_\lambda}$ 按假設是素模，這個因子或者等於 $\mM_{\mb_\lambda}$，或者等於單位模。在第一種情形下，$\ma_\varkappa\mb_\lambda$ 屬於 $\mM_{\mb_\lambda}$；於是存在關係
% unit: noether.p17.eq.025c
\[
 \ma_\chi\mb_\lambda=\mx(\mb_1+\cdots+\mb_{\lambda-1}+\mb_{\lambda+1}+\cdots+\mb_l),
\]
由此得出 $\ma_\chi\mb_\lambda=0$。在另一種情形下，存在兩個類 $\mx_0$ 和 $\mx_1$，使得
% unit: noether.p17.eq.025d
\[
 \mx_0\ma_\varkappa\mb_\lambda
 +\mx_1(\mb_1+\cdots+\mb_{\lambda-1}+\mb_{\lambda+1}+\cdots+\mb_l)
 =\me=\mb_1+\cdots+\mb_{\lambda-1}+\mb_{\lambda+1}+\cdots+\mb_l+\mb_\lambda,
\]
於是由於唯一性，$\mx_0\ma_\chi\mb_\lambda=\mb_\lambda$。因此，當 $F$ 或 $\mx$ 遍歷所有多項式或 $\mG$ 的所有剩餘類時，$F\mx_0\ma_\chi\mb_\lambda$，因而 $\mx\ma_\chi\mb_\lambda$ 也遍歷整個群 $\mB_\lambda$。斷言得證。

% unit: noether.p17.par.177
若 $\ma_\chi\mb_\lambda\ne0$，則 $\mU_\chi$ 還同構於 $\mB_\lambda$。因為由 \S~7，所有滿足 $\mx\ma_\chi\ne0$ 且 $\mx\ma_\chi\mb_\lambda\ne0$ 的類 $\mx\ma_\chi$ 的全體，與類 $\mx\ma_\chi\mb_\lambda$ 的系統同構；而這些類窮盡子群 $\mB_\lambda$。還只需證明所指出的類 $\mx\ma_\chi$ 窮盡 $\mU_\chi$。為此，考察所有滿足 $\mx\ma_\chi\mb_\lambda=0$ 的剩餘類 $\mx\ma_\chi$。這個系統具有如下性質：其中兩個類的和以及與任意多項式的乘積仍屬於這個系統。若把所有這些剩餘類加入模 $\mM_{\ma_\chi}$，則又得到一個作為 $\mM_{\ma_\chi}$ 因子的模；因此它或者是 $\mM_{\ma_\chi}$，或者是單位模。但在後一種情形中，$\ma_\chi$ 本身也包含在其中，從而 $\ma_\chi\mb_\lambda=0$，這與假設矛盾。因此不存在滿足 $\mx\ma_\chi\ne0$ 而 $\mx\ma_\chi\mb_\lambda=0$ 的類，於是 $\mU_\chi$ 與 $\mB_\lambda$ 的同構性得證。

% unit: noether.p17.par.178
現在設對固定的 $\chi$，有 $\ma_\chi\mb_{\lambda_1}\ne0,\ldots,\ma_\chi\mb_{\lambda_i}\ne0$，其中 $i>1$。那麼 $\mU_\chi$ 同構於群 $\mB_{\lambda_1},\ldots,\mB_{\lambda_i}$；因此這些群彼此同構。

% unit: noether.p17.par.179
現在我們證明，於是至少兩個 $\mU$ 也必須同構。因為如果在乘積 $\mb_\lambda\ma_\chi$ 中，對每個 $\mb_\lambda$ 只有一個 $\ma_\chi$ 滿足 $\mb_\lambda\ma_\chi\ne0$，從而 $\mb_\lambda=\mb_\lambda\ma_\chi$，那麼 $\mB_\lambda=\mU_\chi$，因為 $\mU_\chi$ 作為素群不能含有真子群。於是 $k=l$，並且在適當編號後 $\ma_\chi=\mb_\chi$；但這樣一來，圖式 $\ma_\chi\mb_\lambda$ 將是對角圖式，而這並非現在的情形。於是以下定理得證。

% unit: noether.p17.thm.005
\textbf{定理 V.} \emph{若一個完全可約群 $\mG$ 可以用兩種方式表示為素群之和，
\[
 \mG=\mU_1+\cdots+\mU_k=\mB_1+\cdots+\mB_l,
\]
則每個群 $\mU$ 至少同構於一個群 $\mB$，反之亦然。若每個 $\mU$ 恰好同構於一個 $\mB$，則兩個分解相同。在另一種情形中，至少有兩個群 $\mU$ 彼此同構，同樣至少有兩個群 $\mB$ 彼此同構。}

% unit: noether.p17.sec.009
\subsection*{\S~9. “同構”與“同類”概念之間的聯繫.}

% unit: noether.p17.par.180
本節說明：一元微分表達式中已知的“同類”概念，如果按其意義自然推廣，就會導致相應剩餘群的同構性。

% unit: noether.p17.def.009
\textbf{定義.}\footnote{對於一元微分表達式，參見 Blumberg 前引處第 25 頁的形式定義。} \emph{兩個模 $\mM$ 與 $\mN$ 稱為同類，如果存在兩個多項式 $P$ 與 $Q$，其中 $P$ 與 $\mM$ 互素，$Q$ 與 $\mN$ 互素，使得對每個 $M\equiv0(\mM)$ 和 $N\equiv0(\mN)$ 都有}
% unit: noether.p17.eq.025
\begin{align}
 MQ&\equiv0(\mN),\tag{25}\\
 NP&\equiv0(\mM)\tag{26}
\end{align}
\emph{成立。這裡的互素性意味着存在兩個多項式 $X$ 和 $Y$，使得}
% unit: noether.p17.eq.027
\begin{align}
 YQ&\equiv1(\mN),\tag{27}\\
 XP&\equiv1(\mM)\tag{28}
\end{align}
\emph{成立。}

% unit: noether.p17.par.181
於是有以下定理。

% unit: noether.p17.thm.006
\textbf{定理 VI.} \emph{兩個模同類，當且僅當它們的剩餘群同構。}

% unit: noether.p17.par.182
1. 設 $\mA$ 是 $\mM$ 的剩餘群，$\mB$ 是 $\mN$ 的剩餘群，並且 $\mA$ 同構於 $\mB$；設 $\ma$ 和 $\mb$ 分別為 $\mA$ 與 $\mB$ 的單位類。再設 $Q\mb$ 是 $\mB$ 中與 $\mA$ 中的單位 $\ma$ 對應的類，同樣 $P\ma$ 是 $\mA$ 中與 $\mB$ 中的單位 $\mb$ 對應的類。我們把這種對應寫成
% unit: noether.p17.eq.029
\begin{align}
 \ma&\sim Q\mb,\tag{29}\\
 P\ma&\sim\mb.\tag{30}
\end{align}
由此得到
% unit: noether.p17.eq.031a
\[
 0=M\ma\sim MQ\mb,
 \qquad\text{也就是說}\qquad MQ\equiv0(\mN),
\]
這證明了 (25)。同樣有
% unit: noether.p17.eq.031b
\[
 NP\ma\sim N\mb=0,
 \qquad\text{也就是說}\qquad NP\equiv0(\mM),
\]
從而得到關係 (26)。此外，由 (29) 有 $P\ma\sim PQ\mb$，而按 (30) 已假定 $P\ma\sim\mb$，故由對應的唯一性得 $PQ\mb=\mb$，於是 $PQ\equiv1(\mN)$，同樣 $QP\equiv1(\mM)$。因此 (27) 和 (28) 在 $X=Q$、$Y=P$ 時得證。

% unit: noether.p17.par.183
2. 反過來，設 $\mM$ 與 $\mN$ 同類，因而關係 (25) 到 (28) 都成立。我們把 $\mA$ 的任意類 $F\ma$ 指派到 $\mB$ 的類 $FQ\mb$。這樣，$\mA$ 中的每個類只對應 $\mB$ 中的一個類；因為由 $F\ma=0$，即 $F=M$，根據 (25) 得 $FQ\mb=0$。此外，這個對應窮盡整個群 $\mB$，因為特殊類 $Y\ma$ 對應於 $YQ\mb$，而按 (27) 它等於 $\mb$。因此我們得到一個從 $\mA$ 到 $\mB$ 上的單值映射 $\Phi$，但尚未證明它是一一的。同樣，方程 (26) 和 (28) 給出一個從 $\mB$ 到 $\mA$ 上的單值映射 $\Psi$，它也尚未證明是一一的。

% unit: noether.p17.par.184
若這兩個映射之一是一一的，則同構性已經證明，因為關於和與乘積的要求都滿足。可是若兩個映射 $\Phi$ 與 $\Psi$ 都不是一一的，那麼在 $\mA$ 中就會有一些類 $G\ma$，在 $\mB$ 中對應於類 $GQ\mb=0$。與所有其他剩餘類 $F\ma$ 相對應的類 $FQ\mb$ 那時已經窮盡群 $\mB$，因此通過映射 $\Psi$，整個群 $\mA$ 會在類 $FQP\ma$ 中返回。如我們所知，$\mB$ 中還存在非零類，它們在 $\Psi$ 下對應於 $\mA$ 中的零類；我們用 $G_1\ma$ 表示所有滿足 $G_1\ma\ne0$ 但 $G_1QP\ma=0$ 的那些類。所有其餘多項式 $F$ 具有這樣的性質：$FQP\ma$ 窮盡群 $\mA$。那些滿足 $FQP\ma=G_1\ma$ 的多項式 $F$ 記作 $G_2$；只要有多項式 $G_1$ 存在，這樣的多項式就總存在。對於每個 $G_2$，有 $G_2QP\ma\ne0$，但 $G_2(QP)^2\ma=0$。相應地，對於每個整數 $\nu$，都有類 $G_\nu\ma$，使得 $G_\nu(QP)^{\nu-1}\ma\ne0$ 但 $G_\nu(QP)^\nu\ma=0$。

% unit: noether.p17.par.185
現在考察由所有多項式 $G_1,G_2,\ldots$ 組成的系統 $\mS$。它具有一個有限模基 $G_{\lambda_1},\ldots,G_{\lambda_e}$。取 $\nu$ 大於指標 $\lambda_1,\ldots,\lambda_e$ 中最大的那個，並把後者記為 $\lambda$。於是
% unit: noether.p17.eq.031c
\[
 G_\nu=A_1G_{\lambda_1}+\cdots+A_eG_{\lambda_e}.
\]
乘以 $(QP)^\lambda$ 後得到
% unit: noether.p17.eq.031d
\[
 G_\nu(QP)^\lambda\equiv0(\mM),
\]
這與我們的假設矛盾。因此同構性成立，並且由 1 還得到更精確的關係 (27)、(28) 在 $X=Q$、$Y=P$ 時成立。

% unit: noether.p17.par.186
藉助定理 VI，定理 V 也可以表述如下。

% unit: noether.p17.thm.007
\textbf{定理 VII.}\footnote{參見 Loewy, pp. 100--101.} \emph{如果一個完全可約模 $\mM$ 能夠以兩種方式表示為素模 $\mP_1,\ldots,\mP_k$ 和 $\mathfrak{Q}_1,\ldots,\mathfrak{Q}_l$ 的最小公倍，而這兩族各自都形成一個完全互素系統，那麼每個 $\mP$ 都與至少一個 $\mathfrak{Q}$ 同類，反之亦然。若每個 $\mP$ 恰好只與一個 $\mathfrak{Q}$ 同類，則兩個分解相同。在另一種情形中，至少有兩個模 $\mP$ 彼此同類，同樣至少有兩個模 $\mathfrak{Q}$ 彼此同類。}
%
%
% unit: noether.p17.s10.title
\subsection*{\S~10. 一個群的無窮多個分解的存在性。}

% unit: noether.p17.s10.par.001
如果一個完全可約群的和表示不是唯一的，那麼按定理 V，在每一種表示中至少存在兩個彼此同構的子群。現在考慮由這樣一個同構子群系所構成的子群
% unit: noether.p17.s10.eq.001
\[
  \mathfrak H=\mA_1+\cdots+\mA_\alpha,
\]
其中 $\mA_1,\mA_2,\ldots,\mA_\alpha$ 全都同構。我們斷言，只要有理性域 $P$ 含有無窮多個“常數”，即含有滿足 $\xi_i a=a\xi_i$ 的量 $a$，則 $\mathfrak H$ 允許無窮多個這樣的不同分解。這個定理甚至在對子群 $\mA_i$ 不作任何進一步假設、甚至不假設不可約性的情況下也成立。因此有：

% unit: noether.p17.s10.thm.008
\textbf{定理 VIII.} \emph{若 $\mG$ 可以表示為有限多個彼此同構的子群之和，那麼只要有理性域 $P$ 含有無窮多個常數，這種表示就能以無窮多個不同方式實現。這裡所謂常數，是指該域中滿足 $\xi_i a=a\xi_i$ 的那些量 $a$。}

% unit: noether.p17.s10.par.002
\textbf{證明。} 設
% unit: noether.p17.s10.eq.002
\[
  \mG=\mA_1+\cdots+\mA_\alpha,
\]
並且對於 $i,\varkappa=1,\ldots,\alpha$，$\mA_i$ 與 $\mA_\varkappa$ 同構（$\alpha\ge2$）。為了給出一個確定的同構對應，先指定一個子群，例如 $\mA_1$。令 $\mA_1$ 與 $\mA_i$ 之間的同構由下式固定：
% unit: noether.p17.s10.eq.031
\begin{equation}
 a_1\sim t_{1i}a_i,
 \qquad
 a_i\sim t_{i1}a_1,
 \qquad
 t_{11}a_1=a_1
 \quad(i,\varkappa=1,\ldots,\alpha),
\tag{31}
\end{equation}
也就是說，在群 $\mA_1$ 中每一個剩餘類都對應於自身。由 $M\mathfrak{a}_1=0$、$M\mathfrak{a}_i=0$ 便得 $M\mathfrak{t}_{1i}\mathfrak{a}_i=0$、$M\mathfrak{t}_{i1}\mathfrak{a}_1=0$；所以類 $\mathfrak{t}_{1i}\mathfrak{a}_i$、$\mathfrak{t}_{i1}\mathfrak{a}_1$ 像單位一樣，允許與類相乘，而不只允許與多項式相乘。由同構的定義可知，對於兩個對應的類 $\my$ 和 $\bar{\my}$，只要它們允許這種類乘法，則 $\mathfrak x\mathfrak y$ 與 $\mathfrak x\bar{\mathfrak y}$ 仍然是對應的類。因此關係 $\mathfrak{a}_r\mathfrak{a}_s=0$ $(r\ne s)$（$r=1,\ldots,\alpha$；$s=1,\ldots,\alpha$）給出，在 $s=1$ 時，
% unit: noether.p17.s10.eq.032
\begin{equation}
\left\{
\begin{aligned}
 \mathfrak{a}_r \mathfrak{t}_{1\varkappa}\mathfrak{a}_\varkappa&=0 &&(\varkappa=1,\ldots,\alpha,\ r\ne1),\\
 \text{且當 }s>1\mathfrak{\text{ 時}\qquad a}_r\mathfrak{ t}_{s1}\mathfrak{a}_1&=0 &&(r\ne s),\\
 \text{因而也有}\qquad t_{1i}a_\varkappa&=a_1t_{1i}a_i,
 & \mathfrak{t}_{i1}\mathfrak{a}_1&=\mathfrak{a}_i\mathfrak{t}_{i1}\mathfrak{a}_1.
\end{aligned}\right.
\tag{32}
\end{equation}
此外，由 (31)，作為 (27) 和 (28) 的推廣，再由對應的唯一性得到
% unit: noether.p17.s10.eq.033
\begin{equation}
 \mathfrak{t}_{1i}\mathfrak{t}_{i1}\mathfrak{a}_1=\mathfrak{a}_1,
 \qquad
 \mathfrak{t}_{\varkappa1}\mathfrak{t}_{1\varkappa}\mathfrak{a}_\varkappa=\mathfrak{a}_\varkappa.
\tag{33}
\end{equation}

% unit: noether.p17.s10.par.003
現在把 $\mA_i$ 與 $\mA_\varkappa$ 之間的同構固定為由 (31) 複合而成的關係；必須這樣固定，因為各個群可能允許自身到自身的同構。於是由 (31)、(32) 和 (33) 得到
% unit: noether.p17.s10.eq.034
\begin{equation}
 \mathfrak{a}_i\sim \mathfrak{t}_{i1}\mathfrak{t}_{1\varkappa}\mathfrak{a}_\varkappa=\mathfrak{t}_{i\varkappa}\mathfrak{a}_\varkappa,
 \qquad
 \mathfrak{t}_{ii}\mathfrak{a}_i=\mathfrak{a}_i,
 \qquad
 \mathfrak{a}_r\mathfrak{t}_{s\varkappa}\mathfrak{a}_\varkappa=0\quad (r\ne s),
\tag{34}
\end{equation}
其中置 $\mathfrak{t}_{i1}\mathfrak{t}_{1\varkappa}=\mathfrak{t}_{i\varkappa}$。由 (33) 和 (34) 還得
% unit: noether.p17.s10.eq.034a
\[
 \mathfrak{t}_{i\varkappa}\mathfrak{t}_{\varkappa l}\mathfrak{a}_l=\mathfrak{t}_{i1}\mathfrak{t}_{1\varkappa}\mathfrak{t}_{\varkappa1}\mathfrak{t}_{1l}\mathfrak{a}_l
 =\mathfrak{t}_{i1}\mathfrak{t}_{1l}\mathfrak{a}_l=\mathfrak{t}_{il}\mathfrak{a}_l;
\]
由此，先前對群 $\mA_1$ 的特殊指定又被消除了。$\mA_i$ 與 $\mA_\varkappa$ 之間的同構關係保持不變，不論在定義中用哪個群 $\mA_\varkappa$ 來代替 $\mA_1$。因此，$\mathfrak{t}_{i\varkappa}$ 是一些類，概括起來滿足關係
% unit: noether.p17.s10.eq.035
\begin{equation}
 M\mathfrak{t}_{i\varkappa}\mathfrak{a}_\varkappa=0;
 \quad \mathfrak{a}_r\mathfrak{t}_{i\varkappa}\mathfrak{a}_\varkappa=0\ (r\ne i);
 \quad \mathfrak{t}_{i\varkappa}\mathfrak{t}_{\varkappa l}\mathfrak{a}_l=\mathfrak{t}_{il}\mathfrak{a}_l;
 \quad \mathfrak{t}_{\varkappa\varkappa}\mathfrak{a}_\varkappa=\mathfrak{t}_{\varkappa i}\mathfrak{t}_{i\varkappa}\mathfrak{a}_\varkappa=\mathfrak{a}_\varkappa.
\tag{35}
\end{equation}
這些關係正是說，屬於 $\mA_i$ 與 $\mA_\varkappa$ 的模 $\mM_i$ 與 $\mM_\varkappa$ 是同種的（參見 \S~9 的定義）；在那裡出現的 $P$ 和 $Q$，在這裡分別由類 $\mathfrak{t}_{i\varkappa}$ 和 $\mathfrak{t}_{\varkappa i}$ 代替，而 (35) 的前兩個關係合起來分別對應於公式 (25) 和 (26)。反過來，由定理 VI，這些關係推出同構。\footnote{而且無需應用有限性定理，因為這些公式本身已經分別給出了一個映射 $\varphi_{i\varkappa}$ 及其逆映射 $\varphi_{\varkappa i}$。事實上，由 $\mathfrak{x}\mathfrak{a}_i=\mathfrak{x}\mathfrak{a}_i \mathfrak{t}_{i\varkappa}\mathfrak{a}_\varkappa \mathfrak{t}_{\varkappa i}\mathfrak{a}_i$ 可知，$\mathfrak{x}\mathfrak{a}_i\ne0$ 也推出 $\mathfrak{x}\mathfrak{a}_i \mathfrak{t}_{i\varkappa}\mathfrak{a}_\varkappa\ne0$。}

% unit: noether.p17.s10.par.004
現在可以由類 $\mathfrak{t}_{i\varkappa}\mathfrak{a}_\varkappa$ 組合出類 $\mathfrak{b}_\rho$，它們將再次成為某些群 $\mB_\rho$ 的單位，於是得到表示
% unit: noether.p17.s10.eq.035a
\[
 \mG=\mB_1+\cdots+\mB_\alpha.
\]
設 $\lambda_{i\varkappa}$ 是常數，其行列式 $|\lambda_{i\varkappa}|\ne0$；令 $\lambda^{i\varkappa}$ 為相應的餘子式除以 $|\lambda_{i\varkappa}|$，於是有關係
% unit: noether.p17.s10.eq.036
\begin{equation}
\left\{
\begin{aligned}
 \sum_{r=1}^{\alpha}\lambda_{ri}\lambda^{r\varkappa}&=\delta_{i\varkappa}=\begin{cases}0,&i\ne\varkappa,\\1,&i=\varkappa,
 \end{cases}\\
 \sum_{r=1}^{\alpha}\lambda_{ir}\lambda^{\varkappa r}&=\delta_{i\varkappa}.
\end{aligned}\right.
\tag{36}
\end{equation}
於是定義
% unit: noether.p17.s10.eq.037
\begin{equation}
 \mathfrak{b}_\rho=\sum_{i,\varkappa}\lambda_{\rho i}\lambda^{\rho\varkappa}\mathfrak{t}_{i\varkappa}\mathfrak{a}_\varkappa
 \qquad(\rho=1,\ldots,\alpha).
\tag{37}
\end{equation}
由 (35)、(36) 和 (37) 得 $M\mathfrak{b}_\rho=0$，並且
% unit: noether.p17.s10.eq.038
\begin{equation}
 \sum_{\rho=1}^{\alpha}\mathfrak{b}_\rho
 =\sum_{i,\varkappa,\rho}\lambda_{\rho i}\lambda^{\rho\varkappa}\mathfrak{t}_{i\varkappa}\mathfrak{a}_\varkappa
 =\sum_{i,\varkappa}\delta_{i\varkappa}\mathfrak{t}_{i\varkappa}\mathfrak{a}_\varkappa
 =\sum_i \mathfrak{t}_{ii}\mathfrak{a}_i=\sum_i \mathfrak{a}_i=\mathfrak{e},
\tag{38}
\end{equation}
此外
% unit: noether.p17.s10.eq.038a
\[
\begin{aligned}
 \mathfrak{b}_\rho \mathfrak{b}_\sigma
 &=\sum_{\substack{i,\varkappa\\ \mu,\nu}}
  \lambda_{\rho i}\lambda^{\rho\varkappa}\lambda_{\sigma\mu}\lambda^{\sigma\nu}
  \mathfrak{t}_{i\varkappa}\mathfrak{a}_\varkappa \mathfrak{t}_{\mu\nu}\mathfrak{a}_\nu \\
 &=\sum_{i,\mu,\nu}\lambda_{\rho i}\lambda^{\rho\mu}\lambda_{\sigma\mu}\lambda^{\sigma\nu}\mathfrak{t}_{i\nu}\mathfrak{a}_\nu
 =\delta_{\sigma\rho}\sum_{i,\nu}\lambda_{\rho i}\lambda^{\sigma\nu}\mathfrak{t}_{i\nu}\mathfrak{a}_\nu
 =\begin{cases}0&(\rho\ne\sigma),\\ \mathfrak{b}_\rho&(\rho=\sigma),\end{cases}
\end{aligned}
\]
$(\rho,\sigma=1,\ldots,\alpha)$。所以表示 $\mathfrak r e=\mathfrak r b_1+\cdots+\mathfrak r b_\alpha$ 是群 $\mG$ 的一個加法分解。只要 $\lambda_{i\varkappa}$ 不僅僅組成一個對角矩陣，群 $\mB_\rho$ 就不同於群 $\mA_\sigma$。因為若形成 $a_\sigma b_\rho$，按 (35) 得
% unit: noether.p17.s10.eq.038b
\[
 a_\sigma b_\rho
 =\sum_{i,\varkappa}\lambda_{\rho i}\lambda^{\rho\varkappa}a_\sigma t_{i\varkappa}a_\varkappa
 =\lambda_{\rho\sigma}
 \sum_{\varkappa}\lambda^{\rho\varkappa}\mathfrak{t}_{\sigma\varkappa}\mathfrak{a}_\varkappa\ne0
 \quad\text{當 }\lambda_{\rho\sigma}\ne0,
\]
因為否則最後一個和式中的每一個項都必須消失，而事實並非如此。因此，若對固定的 $\sigma$ 至少存在兩個 $\rho$ 使得 $\lambda_{\rho\sigma}\ne0$，則 $\mA_\sigma$ 不可能與任何一個 $\mB_\rho$ 相同。

% unit: noether.p17.s10.par.005
然而，群 $\mB_\rho$ 彼此之間、並且與群 $\mA_\sigma$ 也同構。為此，我們首先證明存在類 $\mathfrak{r}_{\rho\sigma}\mathfrak{a}_\sigma$ 和 $\mathfrak{r}^{\sigma\rho}\mathfrak{b}_\rho$，它們實現 $\mA_\sigma$ 與 $\mB_\rho$ 的同構（$\rho,\sigma=1,\ldots,\alpha$）。確切地說，我們證明，對於類
% unit: noether.p17.s10.eq.038c
\[
 \mathfrak{r}_{\rho\sigma}\mathfrak{a}_\sigma=\sum_i\lambda_{\rho i}\mathfrak{t}_{i\sigma}\mathfrak{a}_\sigma,
 \qquad
 \mathfrak{r}^{\sigma\rho}=\sum_\varkappa\lambda^{\rho\varkappa}\mathfrak{t}_{\sigma\varkappa}\mathfrak{a}_\varkappa
 =\mathfrak{r}^{\sigma\rho}\mathfrak{b}_\rho\footnotemark
\]
\footnotetext{由 (37) 可得
\[
 r^{\sigma\rho}b_\rho=\sum_{\varkappa,i,\mu}\lambda^{\rho\varkappa}t_{\sigma\varkappa}a_\varkappa\lambda_{\rho i}\lambda^{\rho\mu}t_{i\mu}a_\mu
 =\sum_{\varkappa,\mu}\lambda^{\rho\varkappa}\lambda_{\rho\varkappa}\lambda^{\rho\mu}t_{\sigma\mu}a_\mu=\sum_\mu\lambda^{\rho\mu}t_{\sigma\mu}a_\mu=r^{\sigma\rho}.
\]}
在模
% unit: noether.p17.s10.eq.038d
\[
 (\mathfrak M,\mathfrak{b}_1+\cdots+\mathfrak{b}_{\rho-1}+\mathfrak{b}_{\rho+1}+\cdots+\mathfrak{b}_\alpha)
\]
與
% unit: noether.p17.s10.eq.038e
\[
 (\mathfrak M,\mathfrak{a}_1+\cdots+\mathfrak{a}_{\sigma-1}+\mathfrak{a}_{\sigma+1}+\cdots+\alpha_\sigma)
\]
之間滿足“同種”的關係：
% unit: noether.p17.s10.eq.039
\begin{equation}
\left\{
\begin{aligned}
 \mathfrak{b}_\varkappa \mathfrak{r}_{\rho\sigma}\mathfrak{a}_\sigma&=0 &&(\varkappa\ne\rho);&
 \mathfrak{r}^{\sigma\rho}\mathfrak{r}_{\rho\sigma}\mathfrak{a}_\sigma&=\mathfrak{a}_\sigma\footnotemark,\\
 \mathfrak{a}_\varkappa \mathfrak{r}^{\sigma\rho}\mathfrak{b}_\rho&=0 &&(\varkappa\ne\sigma);&
 \mathfrak{r}_{\rho\sigma}\mathfrak{r}^{\sigma\rho}\mathfrak{b}_\rho&=\mathfrak{b}_\rho,
\end{aligned}\right.
\tag{39}
\end{equation}
\footnotetext{與 (35) 相比，公式的數目加倍了，因為這裡的逆關係不像那裡那樣通過交換指標得到。}
由此，如 (35) 那樣，藉助對應
% unit: noether.p17.s10.eq.040
\begin{equation}
 \mathfrak{b}_\rho\sim \mathfrak{r}_{\rho\sigma}\mathfrak{a}_\sigma;\qquad,
 \qquad
 \mathfrak{a}_\sigma\sim \mathfrak{r}^{\sigma\rho}\mathfrak{b}_\rho
\tag{40}
\end{equation}
便得同構。事實上，
% unit: noether.p17.s10.eq.040a
\[
\begin{aligned}
 \mathfrak{b}_\varkappa \mathfrak{r}_{\rho\sigma}\mathfrak{a}_\sigma
 &=\sum_{\mu,\nu,i}\lambda_{\varkappa\mu}\lambda^{\varkappa\nu}\mathfrak{t}_{\mu\nu}\mathfrak{a}_\nu\cdot\lambda_{\rho i}\mathfrak{t}_{i\sigma}\mathfrak{a}_\sigma\\
 &=\sum_{\mu,i}\lambda_{\varkappa\mu}\lambda^{\varkappa i}\lambda_{\rho i}\mathfrak{t}_{\mu\sigma}\mathfrak{a}_\sigma
 =\delta_{\rho\varkappa}\sum_\mu\lambda_{\varkappa\mu}\mathfrak{t}_{\mu\sigma}\mathfrak{a}_\sigma
 =\delta_{\rho\varkappa}\mathfrak{r}_{\varkappa\sigma}\mathfrak{a}_\sigma=0\quad(\varkappa\ne\rho),
\end{aligned}
\]
而且
% unit: noether.p17.s10.eq.040b
\[
 \mathfrak{r}^{\sigma\rho}\mathfrak{r}_{\rho\sigma}\mathfrak{a}_\sigma
 =\sum_{\varkappa,i}\lambda^{\rho\varkappa}\mathfrak{t}_{\sigma\varkappa}\mathfrak{a}_\varkappa\cdot\lambda_{\rho i}\mathfrak{t}_{i\sigma}\mathfrak{a}_\sigma
 =\sum_i\lambda^{\rho i}\lambda_{\rho i}\mathfrak{t}_{\sigma\sigma}\mathfrak{a}_\sigma=\mathfrak{a}_\sigma.
\]
反過來同樣得到
% unit: noether.p17.s10.eq.040c
\[
 \mathfrak{a}_\varkappa \mathfrak{r}^{\sigma\rho}\mathfrak{b}_\rho
 =\mathfrak{a}_\varkappa\sum_\mu\lambda^{\rho\mu}\mathfrak{t}_{\sigma\mu}\mathfrak{a}_\mu=0
 \quad(\varkappa\ne\sigma),
\]
以及
% unit: noether.p17.s10.eq.040d
\[
 \mathfrak{r}_{\rho\sigma}\mathfrak{r}^{\sigma\rho}\mathfrak{b}_\rho
 =\sum_{i,\mu}\lambda_{\rho i}\mathfrak{t}_{i\sigma}\mathfrak{a}_\sigma\lambda^{\rho\mu}\mathfrak{t}_{\sigma\mu}\mathfrak{a}_\mu
 =\sum_{i,\mu}\lambda_{\rho i}\lambda^{\rho\mu}\mathfrak{t}_{i\mu}\mathfrak{a}_\mu=\mathfrak{b}_\rho.
\]
由已經證明的同種關係，$\mA_\sigma$ 與 $\mB_\rho$ 的同構隨之得到（同樣不使用有限性定理）。$\mB_\rho$ 與 $\mB_\tau$ 之間的同構由此複合得到；我們把公式也寫出。置
% unit: noether.p17.s10.eq.040e
\[
 \mathfrak{s}_{\rho\tau}\mathfrak{b}_\tau=\mathfrak{r}_{\rho\sigma}\mathfrak{r}^{\sigma\tau}\mathfrak{b}_\tau
 =\sum_{i,j}\lambda_{\rho i}\mathfrak{t}_{i\sigma}\mathfrak{a}_\sigma\cdot\lambda^{\tau j}\mathfrak{t}_{\sigma j}\mathfrak{a}_j
 =\sum_{i,j}\lambda_{\rho i}\lambda^{\tau j}\mathfrak{t}_{ij}\mathfrak{a}_j;
\]
則對應為 $\mathfrak{b}_\rho\sim\mathfrak{s}_{\rho\tau}\mathfrak{b}_\tau$。這裡也可以再次形式地驗證同種關係；若把 $\mathfrak{s}_{\rho\tau}$ 放在 $\mathfrak{t}_{i\varkappa}$ 的位置，它們滿足 (35)。事實上，
% unit: noether.p17.s10.eq.040f
\[
 b_\sigma\ell_{\rho\tau}b_\tau=b_\sigma r_{\rho\sigma}r^{\sigma\tau}b_\tau=0\quad(\sigma\ne\rho),
\]
% unit: noether.p17.s10.eq.040g
\[
 \mathfrak{s}_{\rho\sigma}\mathfrak{s}_{\sigma\tau}\mathfrak{b}_\tau=\mathfrak{r}_{\rho\varkappa}\mathfrak{r}^{\varkappa\sigma}\mathfrak{r}_{\sigma\varkappa}\mathfrak{r}^{\varkappa\tau}\mathfrak{b}_\tau
 =\mathfrak{r}_{\rho\varkappa}\mathfrak{a}_\varkappa \mathfrak{r}^{\varkappa\tau}\mathfrak{b}_\tau=\mathfrak{s}_{\rho\tau}\mathfrak{b}_\tau.
\]

% unit: noether.p17.s10.par.006
於是定理 VIII 得證，因為只須以無窮多種方式選擇 $\lambda_{i\varkappa}$，使它們不構成對角矩陣，並且也不能通過對角矩陣相互變換為同一種選擇。

% unit: noether.p17.s10.par.007
定理 VIII 與定理 V 或 VII 結合，給出如下結果：

% unit: noether.p17.s10.thm.009
\textbf{定理 IX.}\footnote{參見 Loewy, pp. 106--107.} \emph{設模 $\mM$ 是完全可約的，並且是構成一個全然互素系統的素模 $\mP_1,\ldots,\mP_k$ 的最小公倍模；又設有理性域 $P$ 含有無窮多個常數。那麼，$\mM$ 能夠以無窮多個不同方式表示為素模的最小公倍模的充要條件，是在模 $\mP_1,\ldots,\mP_k$ 中至少有兩個同種。}

% unit: noether.p17.s10.par.008
現在我們要導出一個模 $\mM$ 可被無窮多個素模整除的一般判別準則，並且首先證明下面的結論：

% unit: noether.p17.s10.lemma.001
\textbf{輔助定理。}\footnote{參見 Loewy, p. 96.} \emph{若完全可約模 $\mM=[\mP_1,\ldots,\mP_k]$ 可被一個不同於所有 $\mP_i$ 的素模 $\mP'$ 整除，那麼至少有兩個 $\mP_i$ 同種。}

% unit: noether.p17.s10.par.009
我們首先從 $\mP_1,\ldots,\mP_k$ 中選出一個全然互素系統；做法是依次刪去那些包含在尚未刪去的其餘素模的最小公倍模中的素模。因此現在可以假定 $\mP_1,\ldots,\mP_k$ 本身構成一個全然互素系統。設 $a'$ 是 $\mM$ 中對應於模 $\mP'$ 下單位類的一個類。則對於每一個多項式 $F$，有
% unit: noether.p17.s10.eq.041
\[
 F\mathfrak{a}'=F\mathfrak{a}'\mathfrak{a}_1+\cdots+F\mathfrak{a}'\mathfrak{a}_k.
\]
在乘積 $\mathfrak{a}'\mathfrak{a}_i$ 中至少有兩個不為零，譬如 $\mathfrak{a}'\mathfrak{a}_1$ 和 $\mathfrak{a}'\mathfrak{a}_2$；因為如果例如只有 $\mathfrak{a}'\mathfrak{a}_1\ne0$，那麼 $F\mathfrak{a}'=F\mathfrak{a}'\mathfrak{a}_1$，剩餘類 $F\mathfrak{a}'$ 就會形成 $\mA_1$ 的一個子群，從而由於 $\mA_1$ 是素群，它就與 $\mA_1$ 相同，於是模 $\mP'$ 和 $\mP_1$ 也會重合。按 \S~8 的推理可知，$\mA'$ 與 $\mA_1$ 和 $\mA_2$ 同構。

% unit: noether.p17.s10.par.010
現在令 $\mM$ 為任意模。我們考慮所有整除 $\mM$ 的素模的最小公倍模 $\mN$。\footnote{如果這樣的素模有無窮多個，那麼按定理 IV，剩餘群的任何加法分解都不能對應於這種表示。亦參見例 \S~12, 4。} 於是有兩種情形。或者 $\mN$ 不能表示為有限多個素模的最小公倍模；在這種情形下，$\mN$，因而 $\mM$，有無窮多個素模作為因子。或者 $\mN$ 可以表示為有限多個素模的最小公倍模；這時 $\mN$ 稱為 $\mM$ 的最大完全可約模（依照 Loewy 的一個相應構造）。如果在這有限多個素模中有兩個同種，那麼由定理 IX，模 $\mN$，因而 $\mM$，可被無窮多個素模整除。反過來，如果後一種情況發生，那麼 $\mN$ 至少有一個素因子不同於那有限多個其最小公倍模表示 $\mN$ 的模。由輔助定理，$\mN$，因而 $\mM$，有兩個同種的素因子。我們由此證明了如下定理：

% unit: noether.p17.s10.thm.010
\textbf{定理 X.} \emph{若有理性域 $\mathsf{P}$ 含有無窮多個常數，那麼一個模可被無窮多個不同素模整除的充要條件是：要麼不存在最大完全可約模；要麼在存在這樣的模時，它可被至少兩個不同而同種的素模整除。}
%
%

% unit: noether.p17.s11.title
\subsection*{\S~11. 與微分方程組及其積分的關係.}

% unit: noether.p17.s11.par.001
由有限性定理，對於每一個微分式的加群，都可以取有限多個微分式組成的一個系統作為基，使得把該加群中任一微分式置零所得的所有微分方程的共同積分全體，恰好與把基微分式置零所得的有限微分方程組的積分全體相同。由此便給出了與線性偏微分方程組理論的聯繫；在加群的剰餘群為“有限”的情形下，我們還要把這種聯繫再推進一點。

% unit: noether.p17.s11.def.001
\textbf{定義.} \emph{剰餘群稱為有限的，如果存在一個數 $\mu$，即該群的階，使得任意 $\mu+1$ 個剰餘類之間都有一個係數屬於 $\mathsf{P}$ 的線性關係，而同時至少存在一個由 $\mu$ 個剰餘類組成的系統，其間沒有任何這樣的關係；我們稱這樣的系統為剰餘群的一個線性基。在另一種情形下，剰餘群稱為無限的。}

% unit: noether.p17.s11.par.002
兩個有限剰餘群的和具有有限階，並且其階就是兩個階之和。兩個無限群的和，或者一個有限群與一個無限群的和，是無限群。有限群的例子由一變量加群的所有剰餘群給出；但對於多個變量，階也可能是有限的，例如以 $\xi_1^2$, $\xi_2^2$ 為基的加群，其中所有剰餘類都可以由 $1,\xi_1,\xi_2,\xi_1\xi_2$ 的剰餘類線性組成。無限群的例子見 \S~12。

% unit: noether.p17.s11.par.003
現在有如下定理。

% unit: noether.p17.s11.thm.011
\textbf{定理 XI.} \emph{若一個加群具有有限剰餘群，則把該加群的一個基置零而得到的任意微分方程組，恰有與其剰餘群的階相等數目的線性無關積分。反過來，對於一個給定的、由有限多個線性偏微分方程組成的系統，若它只擁有有限個線性無關積分，則這個數等於由這些微分式生成的加群的剰餘群的階。}

% unit: noether.p17.s11.par.004
這裡的有理性領域 $\mathsf{P}$ 由 $n$ 個變量 $x_1,\ldots,x_n$ 的解析函數組成；在某個 $n$ 維區域內，所出現的奇點至多形成較低維數的流形，因而存在任意多的點具有正則的 $n$ 維鄰域。

% unit: noether.p17.s11.par.005
設 $\mM$ 是一個具有有限剰餘群的加群。我們證明存在一個由冪積組成的線性基，具有如下性質：\emph{當任意剰餘類通過這個基來表示時，分母中只會出現 $P$ 中有限多個不同量的冪}。

% unit: noether.p17.s11.par.006
為此，我們把所有冪積 $\xi_1^{a_1}\cdots\xi_n^{a_n}$ 按字典序排列，並取那些模 $\mM$ 與前面諸項線性無關者作為剰餘類的代表 $Q_1,\ldots,Q_m$。按假設，這樣的冪積只有有限多個。其餘的都可以用這些代表線性表示。

% unit: noether.p17.s11.par.007
令 $\mathfrak S$ 為所有未收入線性基的冪積的系統。於是，根據定理 III 證明的第一部分，$\mathfrak S$ 有一個由這類冪積組成的有限子系統 $\mathfrak T$，使得 $\mathfrak S$ 中每個冪積都至少被 $\mathfrak T$ 中的一個冪積整除。令 $P_1,\ldots,P_\rho$ 為 $\mathfrak T$ 的冪積，並令 $M_1=0,\ldots,M_\rho=0\;(\mM)$ 為這樣的關係：藉助它們，$P_1,\ldots,P_\rho$ 各自都可模 $\mM$ 表示為較低冪積的線性組合；例如
% unit: noether.p17.s11.eq.001
\[
 M_i=b_iP_i+\text{較低項}.
\]
若 $P$ 是 $\mathfrak S$ 中任一冪積，因而具有形式
% unit: noether.p17.s11.eq.002
\[
 \xi_1^{h_1}\cdots\xi_n^{h_n}P_i,
\]
則有
% unit: noether.p17.s11.eq.003
\[
\begin{aligned}
 \xi_1^{h_1}\cdots\xi_n^{h_n}M_i
 &=b_i\xi_1^{h_1}\cdots\xi_n^{h_n}P_i+\text{較低項}\\
 &=b_iP+\text{較低項}.
\end{aligned}
\]
現在假定，對所有低於 $P$ 的冪積，已經證明分母中只出現 $b_1,\ldots,b_\rho$ 的冪，而分子由 $M_i$ 的係數通過加法、乘法和微分形成；於是這對 $P$ 本身也成立。這個假定是允許的，因為它對於 $P_1$，即 $\mathfrak S$ 和 $\mathfrak T$ 的第一個元素，已經成立。

% unit: noether.p17.s11.par.008
此外，根據定理 III，$M_1,\ldots,M_\rho$ 是 $\mM$ 的一個加群基。現在考慮 $x_1,\ldots,x_n$ 的一個取值系統 $\beta_1,\ldots,\beta_n$，使得 $b_1\ne0,\ldots,b_\rho\ne0$，並且 $M_i$ 的其餘係數是正則的；也就是說，這是微分方程組 $M_1=0,\ldots,M_\rho=0$ 的一個正則點。冪積 $Q_i=\xi_1^{h_{i1}}\cdots\xi_n^{h_{in}}$ $(i=1,\ldots,m)$ 對應於偏導數
% unit: noether.p17.s11.eq.004
\[
 \frac{\partial^{h_{i1}+\cdots+h_{in}}}{\partial x_1^{h_{i1}}\cdots\partial x_n^{h_{in}}}y,
\]
我們在點 $(\beta_1,\ldots,\beta_n)$ 任意指定這些偏導數的常數值 $(\gamma_1,\ldots,\gamma_m)$。於是微分方程 $M_i=0$ 唯一地把 $y$ 的所有高階導數在該點的值確定為有限值，因為分母中只出現 $b_i$；眾所周知，這便推出存在 $m$ 個線性無關的解析積分。

% unit: noether.p17.s11.par.009
另一方面，在我們的假設下，不可能存在超過 $m$ 個線性無關積分；否則在一個正則點上可以任意指定 $m+1$ 個適當選擇的導數的值，這與任意 $(m+1)$ 個冪積之間必須存在模 $\mM$ 的線性相關相矛盾。

% unit: noether.p17.s11.par.010
反過來，如果一個加群具有無限剰餘群，則用同樣的方法可證明存在無窮多個線性無關積分；所以，在只有 $m$ 個線性無關積分的情形下，剰餘群是有限的，並且按上面所證，其階等於 $m$。

% unit: noether.p17.s11.par.011
在本段最後，我們說明，對於兩個微分式加群，同構的概念，或者等價地說“同種”的概念，對於相應微分方程組的積分來說，與 Poincaré 在一變量情形使用的種類概念具有相同意義（參見 Blumberg, loc. cit., 附錄 I）。

% unit: noether.p17.s11.thm.012
\textbf{定理 XII.} \emph{若兩個微分式加群 $\mM$ 與 $\mN$ 是同種的，則存在兩個微分式 $P$ 和 $Q$，使得當 $y$ 遍歷 $\mM$ 的所有積分、$z$ 遍歷 $\mN$ 的所有積分時，$P(y)$ 和 $Q(z)$ 分別遍歷 $\mN$ 和 $\mM$ 的所有積分。}

% unit: noether.p17.s11.par.012
事實上，按假設（由 \S~9），對於適當選擇的兩個多項式 $P$ 和 $Q$，有關係
% unit: noether.p17.s11.eq.005
\[
 MQ=0(\mN)\qquad PQ=1(\mN)
\]
% unit: noether.p17.s11.eq.006
\[
 NP=0(\mM)\qquad QP=1(\mM).
\]
因此，由於 $NP(y)=0$，函數 $P(y)$ 是 $\mN$ 的一個積分。同樣，$Q(z)$ 是 $\mM$ 的一個積分。由於 $PQ(z)=z$ 且 $QP(y)=y$，$P(y)$ 遍歷 $\mN$ 的所有積分，$Q(z)$ 遍歷 $\mM$ 的所有積分。這樣，每一個非零的 $y$ 都對應一個非零的 $z$，反之亦然。

% unit: noether.p17.s12.title
\subsection*{\S~12. 例子.}

% unit: noether.p17.s12.par.001
1. 作為第一個例子，我們考察一個二階常線性微分方程，其係數被表示為積分的對稱函數。這裡的有理性領域 $\mathsf{P}$ 由積分及其導數的所有有理函數組成；我們限制在該微分方程的一個正則點的足夠小鄰域內。我們把加群 $\mM$ 定義為所有以 $y_1$ 和 $y_2$ 為積分的微分式的全體，並且仍把它們寫成一個不定元 $\xi$ 的多項式。加群基於是隻有一個元素，由多項式
% unit: noether.p17.s12.eq.001
\[
 M=
 \begin{vmatrix} y_1&y_2\\ y_1'&y_2'\end{vmatrix}\xi^2
 -\begin{vmatrix} y_1&y_2\\ y_1''&y_2''\end{vmatrix}\xi
 +\begin{vmatrix} y_1'&y_2'\\ y_1''&y_2''\end{vmatrix}
\]
給出；它在乘以 $\mathsf{P}$ 中因子的意義下唯一確定。現在
% unit: noether.p17.s12.eq.041
\begin{equation}
 \mM=[\mM_1,\mM_2],
\tag{41}
\end{equation}
其中 $\mM_i$ 由所有具有積分 $y_i$ 的微分式組成，其基在乘以 $\mathsf{P}$ 中因子的意義下由
% unit: noether.p17.s12.eq.002
\[
 M_i=y_i\xi-y_i'\qquad(i=1,2)
\]
確定。確實，$\mM$ 被 $\mM_1$ 和 $\mM_2$ 整除，因為它在積分 $y_1$ 和 $y_2$ 上為零，所以被 $[\mM_1,\mM_2]$ 整除；又由於 $[\mM_1,\mM_2]$ 的基多項式的階至少必須為 $2$，故 $\mM=[\mM_1,\mM_2]$。此外，$\mM_1$ 與 $\mM_2$ 互素，因為
% unit: noether.p17.s12.eq.042
\begin{equation}
 \frac{y_2}{D}(y_1\xi-y_1')-\frac{y_1}{D}(y_2\xi-y_2')=1,
 \qquad
 D=D(y)=\begin{vmatrix}y_1&y_2\\ y_1'&y_2'\end{vmatrix}.
\tag{42}
\end{equation}
它們又是素加群，因為其剰餘群的階為 1。

% unit: noether.p17.s12.par.002
不過，除 $y_1$ 和 $y_2$ 外，$\mM$ 還擁有所有積分
% unit: noether.p17.s12.eq.003
\[
 z_1=\alpha_{11}y_1+\alpha_{12}y_2,
 \qquad
 z_2=\alpha_{21}y_1+\alpha_{22}y_2,
 \qquad
 A=|\alpha_{ik}|\ne0,
\]
因而它也等同於兩個互素加群 $\mN_1$ 和 $\mN_2$ 的最小公倍元；它們分別具有積分 $z_1$ 和 $z_2$，其基多項式為
% unit: noether.p17.s12.eq.004
\[
 N_i=z_i\xi-z_i'=(\alpha_{i1}y_1+\alpha_{i2}y_2)\xi-(\alpha_{i1}y_1'+\alpha_{i2}y_2').
\]
因此，$\mM$ 允許無窮多個不同的表示，作為互素素加群的最小公倍元；由定理 VII，兩個加群 $\mN_1$ 和 $\mN_2$ 因而彼此同種，並且與所有 $\mM_1$ 和 $\mM_2$ 同種。事實上，所有這些加群的剰餘群的階都是 1，所以同構。

% unit: noether.p17.s12.par.003
現在我們寫出與表示 (41) 相應的剰餘群加法分解 $\mG=\mA_1+\mA_2$。為此，根據 \S~4，我們從關係 (42) 出發；它表示 $\mM_1$ 與 $\mM_2$ 互素。對於單位 $\mathfrak{a}_1$ 和 $\mathfrak{a}_2$，我們得到由多項式
% unit: noether.p17.s12.eq.005
\[
 -\frac{y_1}{D}(y_2\xi-y_2')
 \quad\text{和}\quad
 \frac{y_2}{D}(y_1\xi-y_1')
\]
生成的剰餘類。因此
% unit: noether.p17.s12.eq.006
\[
 \mM_1=\left(\mM,\frac{y_2}{D}(y_1\xi-y_1')\right),
 \qquad
 \mM_2=\left(\mM,-\frac{y_1}{D}(y_2\xi-y_2')\right).
\]
現在同樣形成與分解 $\mM=[\mN_1,\mN_2]$ 相應的單位 $\mathfrak{b}_1$ 和 $\mathfrak{b}_2$，即用 $z_i$ 代替 $y_i$，得到
% unit: noether.p17.s12.eq.007
\[
 -\frac{z_1}{D(z)}(z_2\xi-z_2')
 =-\frac{\alpha_{11}y_1+\alpha_{12}y_2}{AD(y)}
 ((\alpha_{21}y_1+\alpha_{22}y_2)\xi-(\alpha_{21}y_1'+\alpha_{22}y_2')),
\]
% unit: noether.p17.s12.eq.008
\[
 \frac{z_2}{D(z)}(z_1\xi-z_1')
 =\frac{\alpha_{21}y_1+\alpha_{22}y_2}{AD(z)}
 ((\alpha_{11}y_1+\alpha_{12}y_2)\xi-(\alpha_{11}y_1'+\alpha_{12}y_2')),
\]
因而
% unit: noether.p17.s12.eq.043
\begin{equation}
\left\{
\begin{aligned}
 \mathfrak{b}_1&=\frac{\alpha_{11}y_1+\alpha_{12}y_2}{y_1}
 \left(\frac{\alpha_{22}}{A}\right)\mathfrak{a}_1
 +\frac{\alpha_{11}y_1+\alpha_{12}y_2}{y_2}
 \left(-\frac{\alpha_{21}}{A}\right)\mathfrak{a}_2
 =\frac{z_1}{y_1}\alpha^{11}\cdot \mathfrak{a}_1+\frac{z_1}{y_2}\alpha^{12}\cdot \mathfrak{a}_2,\\
 \mathfrak{b}_2&=\frac{\alpha_{21}y_1+\alpha_{22}y_2}{y_1}
 \left(-\frac{\alpha_{12}}{A}\right)\mathfrak{a}_1
 +\frac{\alpha_{21}y_1+\alpha_{22}y_2}{y_2}
 \left(\frac{\alpha_{11}}{A}\right)\mathfrak{a}_2
 =\frac{z_2}{y_1}\alpha^{21}\mathfrak{a}_1+\frac{z_2}{y_2}\alpha^{22}\mathfrak{a}_2.
\end{aligned}\right.
\tag{43}
\end{equation}
這裡 $(\alpha^{ik})$ 是矩陣 $(\alpha_{ik})$ 的逆矩陣。

% unit: noether.p17.s12.eq.044
反過來，把 $z_i$ 與 $y_i$ 互換，得到
\begin{equation}
\left\{
\begin{aligned}
 \mathfrak{a}_1&=\frac{y_1}{z_1}\alpha_{11}\mathfrak{b}_1+\frac{y_1}{z_2}\alpha_{12}\mathfrak{b}_2,\\
 \mathfrak{a}_2&=\frac{y_2}{z_1}\alpha_{21}\mathfrak{b}_1+\frac{y_2}{z_2}\alpha_{22}\mathfrak{b}_2.
\end{aligned}\right.
\tag{44}
\end{equation}
由這些公式可以看出群 $\mA$ 與 $\mB$ 的同構。按照我們的一般考察，若 $\mathfrak{a}_1\mathfrak{b}_\sigma\ne0$，相應的對應是 $\mathfrak{a}_1\sim \mathfrak{a}_1\mathfrak{b}_\sigma$；因此由 (44)
% unit: noether.p17.s12.eq.009
\[
 \mathfrak{a}_1\sim\frac{y_1}{z_\sigma}\alpha_{1\sigma}\mathfrak{b}_\sigma,
 \qquad\text{若 }\alpha_{1\sigma}\ne0,
\]
而由 (43) 同樣有
% unit: noether.p17.s12.eq.010
\[
 \mathfrak{b}_\sigma\sim\frac{z_\sigma}{y_2}\alpha^{\sigma2}\mathfrak{a}_2,
 \qquad\text{若 }\alpha^{\sigma2}\ne0,
\]
於是
% unit: noether.p17.s12.eq.011
\[
 \frac{y_1}{z_\sigma}\alpha_{1\sigma}\mathfrak{b}_\sigma
 \sim
 \frac{y_1}{z_\sigma}\alpha_{1\sigma}\frac{z_\sigma}{y_2}\alpha^{\sigma2}\mathfrak{a}_2
 =\frac{y_1}{y_2}\alpha_{1\sigma}\alpha^{\sigma2}\mathfrak{a}_2.
\]
因此
% unit: noether.p17.s12.eq.012
\[
 \mathfrak{a}_1\sim\alpha_{1\sigma}\alpha^{\sigma2}\frac{y_1}{y_2}\mathfrak{a}_2
 \quad\text{當 }\alpha_{1\sigma}\alpha^{\sigma2}\ne0,
 \qquad
 \frac{1}{\alpha_{1\sigma}\alpha^{\sigma2}}\frac{y_2}{y_1}\mathfrak{a}_1\sim \mathfrak{a}_2,
\]
這給出了逆對應。若 $(\alpha_{i\varkappa})$ 不是對角矩陣，則此條件 $\alpha_{1\sigma}\alpha^{\sigma2}\ne0$ 對某個 $\sigma$ 總能滿足。對角矩陣的情形必須排除，因為那時加群相同。例如，若只有 $\alpha_{12}=0$，則 $\mA_1=\mB_1$，但並非 $\mA_2=\mB_2$。所以由 $\mA_1+\mA_2=\mB_1+\mB_2$ 和 $\mA_1=\mB_1$ 並不能推出 $\mA_2=\mB_2$（參見注 12）。於是三個加群 $\mM_1$, $\mM_2$ 和 $\mN_2$ 兩兩互素，但並不形成一個完全互素系統，因為 $[\mM_1,\mM_2]$ 被 $\mN_2$ 整除（參見注 15）。

% unit: noether.p17.s12.par.004
此外，若令
% unit: noether.p17.s12.eq.013
\[
 \mathfrak{t}_{12}=\alpha_{1\sigma}\alpha^{\sigma2}\frac{y_1}{y_2}\mathfrak{a}_2,
 \qquad
 \mathfrak{t}_{21}=\frac{1}{\alpha_{1\sigma}\alpha^{\sigma2}}\frac{y_2}{y_1}\mathfrak{a}_1,
\]
則通過計算可確認關係 (35) 成立。

% unit: noether.p17.s12.par.005
2. 任何一個含有多於一個變量、而其基只由一個多項式組成的加群，都給出一個無限群的例子。下面這個稍複雜的例子是為了說明：無限群也可能出現在多元加群中，也就是出現在其加群基包含多於一個多項式的加群中。\footnote{參見 Blumberg 轉述的 Landau 例子（loc. cit., pp. 51--52）。這裡以及下文，我們用 $x$ 和 $y$ 表示獨立變量。}
% unit: noether.p17.s12.eq.014
\[
 \mM=[(\xi+x\eta),(\xi+1)]=[(P),(Q)]
 \qquad\left(\xi=\frac{\partial}{\partial x},\ \eta=\frac{\partial}{\partial y}\right).
\]
乘法規則就是微分式的乘法規則。這裡剰餘群是無限的，因為 $(\xi+x\eta)$ 和 $(\xi+1)$ 的剰餘群都是無限的，而加群 $(\xi+x\eta)$ 與 $(\xi+1)$ 互素，這是由於
% unit: noether.p17.s12.eq.015
\[
 1=(-x\xi-x^2\eta+1)Q+(x\xi+x-1)P.
\]
現在很容易證明 $\mM$ 不含二次多項式；令
% unit: noether.p17.s12.eq.016
\[
 (u_1\xi+u_2\eta+u_3)(\xi+x\eta)=(v_1\xi+v_2\eta+v_3)(\xi+1),
\]
便得 $u_1=u_2=u_3=v_1=v_2=v_3=0$。再在 $\mM$ 中尋找三次多項式時，會發現兩個線性無關的，例如
% unit: noether.p17.s12.eq.017
\[
 (\xi^2+x\xi\eta+\xi+(2+x)\eta)Q=(\xi^2+2\xi+1)P
\]
以及
% unit: noether.p17.s12.eq.018
\[
 (\xi^2+2x\xi\eta+x^2\eta^2+\eta)Q
 =(\xi^2+x\xi\eta+\xi+(x-1)\eta)P.
\]
若 $\mM$ 是單元加群，則二者都必須被其基多項式整除；又因 $\mM$ 不含二次多項式，所以必須被一個三次多項式整除，於是它們應當成比例，但事實並非如此。$\mM$ 的多元性，是 Landau 在一變量情形成立的乘積定理失效的內在原因。

% unit: noether.p17.s12.par.006
3. 另一個例子表明，無限群也會出現在素加群中。\footnote{我們對“素加群”的定義不同於可換情形中所謂素加群；因為在那裡總存在較低維數的因子，除非維數為 $0$，即剰餘群是有限的。} 令 $P$ 為二變量 $x,y$ 的所有有理函數的領域。以
% unit: noether.p17.s12.eq.019
\[
 \xi-\frac{y}{x}=\xi-a
\]
為基的加群 $\mM$ 具有一個無限群，因為它包含 $y$ 的每個冪的剰餘類；另一方面，$\mM$ 是素加群。事實上，我們證明每一個因子
% unit: noether.p17.s12.eq.020
\[
 \mM_1=(\xi-a,F(\xi,\eta))
\]
都等於單位加群。任一多項式 $F(\xi,\eta)$ 模 $(\xi-a)$ 都可表示為一個 $\eta$ 的多項式 $G(\eta)$：
% unit: noether.p17.s12.eq.021
\[
 F(\xi,\eta)\equiv c_0\eta^\nu+c_1\eta^{\nu-1}+\cdots+c_\nu\quad(\mathfrak M).
\]
不失一般性，令 $c_0=1$。按乘法法則，有 $\xi\eta^\nu-\eta^\nu\xi=0$；因此，由於
% unit: noether.p17.s12.eq.022
\[
 \xi\equiv a(\mM_1),
 \qquad
 \eta^\nu\equiv F_1(\mM_1),
\]
其中 $F_1=-(c_1\eta^{\nu-1}+\cdots+c_\nu)$，於是有
% unit: noether.p17.s12.eq.023
\[
 \xi F_1-\eta^\nu a\equiv0(\mM_1).
\]
展開
% unit: noether.p17.s12.eq.024
\[
 \eta^\nu a=a\eta^\nu+\nu\frac{\partial a}{\partial y}\eta^{\nu-1}+\cdots,
\]
並在其中再次用 $F_1$ 代替 $\eta^\nu$，便得到
% unit: noether.p17.s12.eq.025
\[
 \xi(c_1\eta^{\nu-1}+\cdots+c_\nu)-a(c_1\eta^{\nu-1}+\cdots+c_\nu)
 +\nu\frac{\partial a}{\partial y}\eta^{\nu-1}+\cdots\equiv0(\mM_1),
\]
或
% unit: noether.p17.s12.eq.026
\[
 c_1\eta^{\nu-1}a+\frac{\partial c_1}{\partial x}\eta^{\nu-1}+\cdots
 -ac_1\eta^{\nu-1}+\cdots+
 \nu\frac{\partial a}{\partial y}\eta^{\nu-1}+\cdots
 \equiv0(\mM_1),
\]
最後得到
% unit: noether.p17.s12.eq.027
\[
 \left(\frac{\partial c_1}{\partial x}+\nu\frac{\partial a}{\partial y}\right)\eta^{\nu-1}+\cdots
 \equiv0(\mM_1).
\]
現在左邊是一個 $\eta$ 的多項式，精確次數為 $\nu-1$，並且不恆等為零；因為最高係數
% unit: noether.p17.s12.eq.028
\[
 \frac{\partial c_1}{\partial x}+\nu\frac{\partial a}{\partial y}
 =\frac{\partial c_1}{\partial x}+\nu\frac{1}{x}
 \quad\text{不等於 }0,
\]
否則 $c_1=-\nu\log x+\varphi(y)$ 就不會是有理函數。於是此過程可以繼續，並最終得到一個不恆等為零的零次多項式，從而證明單位也包含在 $\mM_1$ 中。\footnote{這個過程歸結為證明：微分方程組所必需的可積性條件未被滿足。}

% unit: noether.p17.s12.par.007
4. 另一方面，若再把函數 $\log x$ 添入有理性領域，則 $\mM$ 有因子
% unit: noether.p17.s12.eq.029
\[
 \left(\xi-\frac{y}{x},\ \eta-\log x+\varphi(y)\right),
\]
其中 $\varphi(y)$ 遍歷 $y$ 的所有有理函數。這些因子彼此全都不同；對於每一個，可積性條件都被滿足，相應的積分為
% unit: noether.p17.s12.eq.030
\[
 y\log x-\int\varphi(y)\,dy+\alpha.
\]
因此單位不能包含在該加群中，否則零函數就會是唯一可能的積分。另一方面，相對於該加群，$\xi$ 與 $\eta$ 都同餘於領域中的某個量；所以剰餘群是有限的，且階為 1。此外，無限多個因子中的任意兩個都互素，而給定加群 $\mM$ 等於所有這些因子的最小公倍元 $\mN$。事實上，它被所有因子整除，因而也被 $\mN$ 整除。若 $\mN$ 是 $\mM$ 的真因子，則 $\mN$ 還必須至少包含一個只依賴於 $\eta$ 的多項式 $F$，設其次數為 $\rho$。令 $\varphi(y)$ 遍歷函數 $y,\ldots,y^{\rho+1}$；在相應積分
% unit: noether.p17.s12.eq.031
\[
 y\log x-\frac{y^{i+1}}{i+1}+\alpha_i\qquad(i=1,\ldots,\rho+1)
\]
之間，不存在係數與 $y$ 無關的線性關係。但是 $\rho$ 階微分方程 $F=0$ 不可能具有 $\rho+1$ 個線性無關積分。

% unit: noether.p17.s12.par.008
然而，模塊 $\mM=\mN$ 並非完全可約。\footnote{這樣我們就得到了定理 X 第一種情形的一個例子。} 否則，它就是這些素加群中有限多個的最小公倍元，其剰餘群因而會有有限階；但是 $\xi-y/x$ 的剰餘群卻包含 $\eta$ 的所有冪的剰餘類，所以是無限的。

% unit: noether.p17.close.001
哥廷根，1919 年 8 月 1 日。

\begin{center}
（1919 年 8 月 4 日收到。）
\end{center}
%
\clearpage
%
% unit: noether.p18.title.001
\begin{center}
\pdfbookmark[0]{18. 關於陣亡於戰爭的 K. Hentzelt 關於消去理論的一項工作}{paper-18}
{\Large\bfseries 18. 關於陣亡於戰爭的 K. Hentzelt 關於消去理論的一項工作}\par
\vspace{1.2em}
\emph{J. Ber. d. DMV} 30 (1921), p. 101
\end{center}

\begin{center}
\textbf{1921 年 9 月 23 日星期五上午 11 時第 11 次會議。}\\
（Loewy 主持。）
\end{center}

% unit: noether.p18.par.001
\textbf{1. E. Noether: 關於陣亡於戰爭的 K. Hentzelt 關於消去理論的一項工作。}

% unit: noether.p18.par.002
Kurt Hentzelt 自 1914 年 10 月在 Dixmuiden 前線失蹤以來，必須被列入死者之中。他的 Erlangen 博士論文的最重要部分，是他留下的、結構上沒有缺口但表述幾乎難以理解的部分；我不久將以自由整理的形式在 \emph{Mathematische Annalen} 上發表。

% unit: noether.p18.par.003
這裡討論的是\textbf{消去理論}的一般問題：確定由多項式組成的一個理想 $\mathfrak M$ 中所有多項式的公共零點。如果預先對變量施以一個帶不定係數的線性變換，那麼主要結果可以表述如下：

% unit: noether.p18.par.004
可以唯一地給理想 $\mathfrak M$ 指派一個結式形式：
% unit: noether.p18.eq.001
\[
 R_{\mathfrak M}=R^{(1)}(x_1,\ldots,x_n)\cdots R^{(i)}(x_i,\ldots,x_n)\cdots R^{(n)}(x_n)=0(\mathfrak M),
\]
% unit: noether.p18.par.005
使得 $R_{\mathfrak M}$ 對 $\mathfrak M$ 的所有零點並且僅對這些零點為零。如果 $\mathfrak N$ 是 $\mathfrak M$ 的一個因子，並且 $R_{\mathfrak N}$ 與 $R_{\mathfrak M}$ 相同，那麼 $\mathfrak N$ 與 $\mathfrak M$ 本身也相同。

% unit: noether.p18.par.006
該定理的第二部分表明，結式形式給出的零點帶有特徵性的重數。它基於以下事實：可以把 $\mathfrak M$ 看作關於 $x_1,\ldots x_{i-1}$ 的冪乘積中的線性形式模；於是 $R_{\mathfrak M}$ 的各個因子 $R^{(i)}(x_i,\ldots,x_n)$，在 $x_{i+1},\ldots x_n$ 被添入係數域時，就成為相應“基本模”關於此模的範。如果引入抽象理想論，則可得到如下解釋：若
% unit: noether.p18.eq.002
\[
 \mathfrak M=[\mathfrak Q,\mathfrak Q_1,\ldots,\mathfrak Q_r]=[\mathfrak Q,\mathfrak L]
\]
% unit: noether.p18.par.007
是 $\mathfrak M$ 分解為初級理想 $\mathfrak Q_\lambda$ 的一個分解，而 $\mathfrak L$ 是 $\mathfrak Q$ 的補，那麼對應於理想 $\mathfrak Q$，有 $R^{(i)}(x_i,\ldots,x_n)$ 的一個初級因子 $Q^{(i)}(x_i,\ldots,x_n)$；它在上面所指出的意義下表示 $\mathfrak L$ 關於 $\mathfrak Q$ 的範。
%
\clearpage
%
% unit: noether.p19.title.001
\pdfbookmark[0]{19. 環域中的理想論}{paper-19}
\section*{19. 環域中的理想論}
\begin{center}
Math. Ann. 83 (1921), pp. 24--66
\end{center}

% unit: noether.p19.toc.001
\begin{center}\textbf{目錄。}\end{center}
\noindent 引言。\\
\S~1. 環域，理想，有限性條件。\\
\S~2. 把一個理想表示為有限多個不可約理想的最小公倍。\\
\S~3. 兩個不同的不可約理想分解中組分個數的相等性。\\
\S~4. 準素理想。兩個不同的不可約理想分解中關聯素理想的唯一性。\\
\S~5. 把一個理想表示為最大準素理想的最小公倍。關聯素理想的唯一性。\\
\S~6. 把一個理想唯一地表示為相對素的不可約理想的最小公倍。\\
\S~7. 孤立理想的唯一性。\\
\S~8. 把一個理想唯一地表示為兩兩互素的不可約理想的乘積。\\
\S~9. 將研究擴展到模。不可約模分解中組分個數的相等性。\\
\S~10. 多項式域的特殊情形。\\
\S~11. 數論和微分表達式理論中的例子。\\
\S~12. 初等因子理論中的例子。

% unit: noether.p19.sec.intro
\subsection*{引言。}

% unit: noether.p19.par.001
本文的內容是把有理整數的分解定理、以及代數數域中理想的分解定理，\emph{轉移到任意整域中的理想，並更一般地轉移到環域中的理想}。為了理解這種轉移，先把有理整數的分解定理用一種稍微不同於通常表述的形式寫出。

% unit: noether.p19.par.002
若在
% unit: noether.p19.eq.001
\[
 a=p_1^{\varrho_1}p_2^{\varrho_2}\cdots p_\sigma^{\varrho_\sigma}=q_1q_2\cdots q_\sigma
\]
中，把素數冪因子 $q_i$ 看作分解的組分，那麼這些組分具有以下特徵性質。

% unit: noether.p19.par.003
1. 它們\emph{兩兩互素}；但沒有一個 $q$ 可以表示為兩兩互素的數的乘積，因此在這個意義下有不可約性。由兩兩互素還推出，乘積 $q_1\cdots q_\sigma$ 等於最小公倍 $[q_1\cdots q_\sigma]$。

% unit: noether.p19.par.004
2. 任意兩個組分 $q_i$ 和 $q_k$ 是\emph{相對素的}；也就是說，若 $b\cdot q_i$ 可被 $q_k$ 整除，則 $b$ 可被 $q_k$ 整除。在這個意義下也有不可約性。

% unit: noether.p19.par.005
3. 每個 $q$ 都是\emph{準素的}；也就是說，若乘積 $b\cdot c$ 可被 $q$ 整除，而 $b$ 不可被整除，則 $c$ 的某個冪\footnote{如果這個冪總是一次冪，那麼眾所周知，所涉及的就是素數。}可被整除。此外，該表示是由\emph{最大準素組分}給出的表示，因為兩個不同的 $q$ 的乘積不再是準素的。就分解為最大準素組分而言，$q$ 也是不可約的。

% unit: noether.p19.par.006
4. 每個 $q$ 都是\emph{不可約的}，其意義是它不能表示為兩個真因子的最小公倍。

% unit: noether.p19.par.007
這些準素數 $q$ 與素數 $p$ 的聯繫在於：對每個 $q$，存在一個且除符號外唯一的 $p$，它是 $q$ 的因子，並且它的某個冪可被 $q$ 整除；這就是關聯素數。若 $p^\varrho$ 是這種冪中最低的一個——$\varrho$ 是 $q$ 的指數——那麼在這裡特別有 $p^\varrho=q$。唯一性定理於是可表述如下：

% unit: noether.p19.par.008
\emph{把一個有理整數分解為不可約的最大準素組分 $q$ 時，在兩個不同分解中，組分的個數、關聯素數（差一個符號不計）以及指數都相同。由於 $p^\varrho=q$，由此還推出 $q$ 本身也相同（差一個符號不計）。}

% unit: noether.p19.par.009
由符號帶來的不定性，眾所周知，在不考慮數本身而考慮由它們導出的理想（所有可被 $a$ 整除的數）時便被消除；於是這個表述同樣適用於（有限）代數數域的理想分解為素理想冪的唯一分解。

% unit: noether.p19.par.010
下面將以一個一般環域為基礎（§~1），它只須滿足如下有限性條件：該域的每個理想都有有限理想基。若無這樣的有限性條件，就未必存在不可約理想和素理想；全體代數整數所成的域就說明了這一點，在其中不存在分解為素理想的分解。

% unit: noether.p19.par.011
結果表明，與組分 $q$ 的四個特徵性質相對應，一般存在四種分開的分解，它們每次都通過進一步分裂而相互產生。其中，分解為兩兩互素的不可約理想時涉及的是乘積表示；其餘三種分解則是作為最小公倍的約化表示（§~2）。準素理想——不可約理想也都是準素的——與關聯素理想之間的聯繫也保持不變：對每個準素理想 $\ideal{Q}$，唯一確定一個關聯素理想 $\ideal{P}$，它是 $\ideal{Q}$ 的因子，並且它的某個冪可被 $\ideal{Q}$ 整除。若 $\ideal{P}^{\varrho}$ 是這種冪中最低的一個——$\varrho$ 是 $\ideal{Q}$ 的指數——那麼在這裡 $\ideal{P}^{\varrho}$ 不必與 $\ideal{Q}$ 相同。唯一性定理於是表述如下。

% unit: noether.p19.par.012
分解 1 和 2 是唯一的；在兩個不同的分解 3 或 4 中，組分個數和關聯素理想相同；\footnote{大概除此以外，指數也相同，更一般地，對應組分還同構。} 在組分中出現的孤立理想（§~7）唯一確定。

% unit: noether.p19.par.013
為了證明分解定理，首先從有限性條件推出“有限鏈定理”，這個定理最早由 Dedekind 對有限數模提出；由此導出每個理想作為有限多個不可約理想的最小公倍的表示 4。通過改造組分可約性的概念，可以由此得到分解 4 為不可約理想時的基本唯一性定理。再通過每次合併有限多個組分，便上升到其餘分解；這些分解的唯一性定理於是作為唯一性定理 4 的推論得到。

% unit: noether.p19.par.014
最後還將證明（§~9），在更弱假設下，表示為有限多個不可約組分仍然成立；環域的交換性並非必要，而且只需考慮相對於該域的一個模來代替理想。在這個更一般的情形下，兩個不同分解中組分個數相等仍然成立；而素與準素的概念則依賴於交換性和理想概念。相反，互素概念對非交換域中的理想仍然有效。

% unit: noether.p19.par.015
四種分開的分解確實出現的最簡單環域，是所有以任意複數為係數的 $n$ 個變量多項式的域。在這裡，各個分解可以藉助代數圖形的行為作非有理的解釋；而關聯素理想的唯一性定理，對應於消去理論的基本定理，即代數圖形可唯一分解為不可約圖形。其他例子由所有由多項式構成的有限整域給出（§~10）。不過即使是全體偶數的簡單域，更一般地說全體可被一個固定數整除的數的域，也已經給出了部分分開的分解的例子（§~11）。非交換域中的理想論的一個例子由初等因子理論提供（§~12），在那裡存在分解為不可約理想、相應地分解為不可約類的唯一分解。這些不可約類完全刻畫初等因子的不可約組成部分；在通常初等因子理論失效的域中，它們也許可以看作初等因子的等價物。

% unit: noether.p19.par.016
關於已有文獻，需要說明如下。對於具有任意復係數或整數係數的多項式域，分解為最大準素理想是由 Lasker 給出的，並由 Macaulay 在若干方面進一步推進。\footnote{E. Lasker, \emph{Zur Theorie der Moduln und Ideale}. Math. Ann. 60 (1905), p. 20, 定理 VII 和 XIII。--- F. S. Macaulay, \emph{On the Resolution of a given Modular System into Primary Systems including some Properties of Hilbert Numbers}. Math. Ann. 74 (1913), p. 66.} 二者都依賴消去理論，因此使用了一個多項式能唯一表示為不可約多項式乘積這一事實。實際上，理想的分解定理獨立於這個假設；代數數域中的理想論已經使人預料到這一點，而本文則證明了它。在 Lasker 和 Macaulay 那裡，準素理想也是以消去理論中的概念為基礎來定義的。

% unit: noether.p19.par.017
分解為不可約理想以及分解為相對素的不可約理想，在文獻中似乎即使對於多項式域也未被注意到；只有 Macaulay 處有一條關於孤立準素理想唯一性的說明。

% unit: noether.p19.par.018
分解為兩兩互素的不可約理想，是 Schmeidler\footnote{W. Schmeidler, \emph{Über Moduln und Gruppen hyperkomplexer Größen}. Math. Zeitschr. 3 (1919), p. 29.} 對多項式域給出的，其中使用了消去理論來證明有限性。不過這裡的唯一性定理只對理想類表述，而不是對理想本身表述。後一種唯一性定理出現在一篇合著論文中，\footnote{E. Noether--W. Schmeidler, \emph{Moduln in nichtkommutativen Bereichen, insbesondere aus Differential- und Differenzenausdrücken}. Math. Zeitschr. 8 (1920), p. 1.} 其中討論的是非交換多項式域中的理想。那裡只使用有限理想基，因此定理和方法對一般環域仍然成立；藉助本文，可以在組分個數相等方面把它們加強（§~11）。

% unit: noether.p19.par.019
這些研究是對上述兩篇論文所依據的概念形成的強烈推廣和進一步發展。兩篇論文的本質是從表示為最小公倍過渡到剩餘類系統的加法分解。這裡為了表述簡單起見，仍然停留在最小公倍上；於是與加法分解相對應的，是把可約性的概念改造成補理想的一個性質（§~3）。不過，根據本質上對應於那篇合著論文的考慮，所有列出的定理也都可以看成剩餘類系統及某些子系統的加法分解定理。這個剩餘類系統形成一個與最初作為基礎的環同樣一般的環；事實上，任意環都可以看成某個理想的剩餘類系統，這個理想對應於環元素之間全部恆等關係；或者，也可以看成這些關係的一個子系統的剩餘類系統，只要把其餘關係也看作在該域中已經滿足。

% unit: noether.p19.par.020
這個說明也給出了 Fraenkel 論文的位置。\footnote{A. Fraenkel, \emph{Über die Teiler der Null und die Zerlegung von Ringen}. J. f. M. 145 (1914), p. 139. \emph{Über gewisse Teilbereiche und Erweiterungen von Ringen}. Habilitationsschrift, Leipzig, Teubner, 1916. \emph{Über einfache Erweiterungen zerlegbarer Ringe}. J. f. M. 151 (1920), p. 121.} Fraenkel 考察滿足若干限制條件的環的加法分解（正則元素的存在、除以這些元素、可分解性條件）；這些條件使得對於相應的理想，四種分解重合。由於這種重合，他的有限性條件——理想只應有有限多個真因子，而他有時又部分地放棄這一條件——也不比我們的條件更強。Fraenkel 的出發點，是由他的論文那些不同的、本質上代數的目標所決定；然後他通過代數擴張達到限制條件較少的更一般的環。

% unit: noether.p19.sec.s01
\subsection*{\S~1. 環域，理想，有限性條件。}

% unit: noether.p19.par.021
\textbf{1.} 作為基礎的域 $\Sigma$ 應當是在抽象定義下的（交換）環；\footnote{這個定義取自 Fraenkel 的 habilitation 論文，省去了他的限制條件 6、I 和 II；為此必須把加法交換律包括進來。因此，這些就是定義一個域的各條定律，只是省去了乘法的可逆性。} 也就是說，$\Sigma$ 由一個元素系統 $a,b,c,\ldots,f,g,h,\ldots$ 組成，其中把一個滿足通常條件的關係定義為相等；並且在其中，由兩個運算（結合方式）——加法和乘法——總能從任意兩個環元素 $a$ 和 $b$ 唯一地得到第三個元素，分別作為和 $a+b$ 以及積 $a\cdot b$。這個環以及除此以外完全任意的運算，必須滿足以下定律：
\begin{enumerate}
\item 加法結合律：$(a+b)+c=a+(b+c)$；
\item 加法交換律：$a+b=b+a$；
\item 乘法結合律：$(a\cdot b)\cdot c=a\cdot(b\cdot c)$；
\item 乘法交換律：$a\cdot b=b\cdot a$；
\item 分配律：$a\cdot(b+c)=a\cdot b+a\cdot c$；
\item 無限制且唯一的減法定律。在 $\Sigma$ 中存在唯一的元素 $x$ 滿足方程 $a+x=b$。（記作 $x=b-a$。）
\end{enumerate}

% unit: noether.p19.par.022
由這些性質推出零元的存在；但環不必有單位元；並且兩個元素的乘積可以為零，而任一因子都不為零。若一個環中乘積為零總推出某個因子為零，並且此外還具有單位元，則稱這樣的環為真正的整域。對於有限和 $a+a\ldots+a$，我們引入通常的縮寫記號 $na$；這裡整數 $n$ 只應看作縮寫符號，而不是環元素，並且由 $a=1\cdot a$、$na+a=(n+1)a$ 遞歸地定義。

% unit: noether.p19.par.023
\textbf{2.} 在 $\Sigma$ 中，所謂一個理想 $\ideal{M}$，\footnote{理想用大德文字母表示。$\ideal{M}$ 是為了讓人想起多項式中的理想這一例子；這個理想通常稱為“模”或型模。順便說，§§~1--3 只使用模的性質，而不使用理想性質；參見 §~9。} 是指 $\Sigma$ 中滿足以下兩個條件的元素系統：
\begin{enumerate}
\item $\ideal{M}$ 與 $f$ 一起還包含 $a\cdot f$，其中 $a$ 是 $\Sigma$ 的任意元素；
\item $\ideal{M}$ 與 $f$ 和 $g$ 一起還包含差 $f-g$；因此與 $f$ 一起還包含每個整數 $n$ 對應的 $nf$。
\end{enumerate}

% unit: noether.p19.par.024
若 $f$ 是 $\ideal{M}$ 的元素，我們照常寫作 $f\equiv0(\ideal{M})$，並說 $f$ 可被 $\ideal{M}$ 整除。若 $\ideal{N}$ 的每個元素同時也是 $\ideal{M}$ 的元素，因而可被 $\ideal{M}$ 整除，我們說：$\ideal{N}$ 可被 $\ideal{M}$ 整除；用符號寫作 $\ideal{N}\equiv0(\ideal{M})$。當 $\ideal{M}$ 含有不同於 $\ideal{N}$ 的元素，因而反過來不被 $\ideal{N}$ 整除時，稱 $\ideal{M}$ 為 $\ideal{N}$ 的真因子。由 $\ideal{N}\equiv0(\ideal{M})$ 和 $\ideal{M}\equiv0(\ideal{N})$ 推出 $\ideal{N}=\ideal{M}$。

% unit: noether.p19.par.025
其他熟知的概念也逐字保留。兩個理想 $\ideal{A}$ 和 $\ideal{B}$ 的最大公因子——$\ideal{D}=(\ideal{A},\ideal{B})$——指的是所有可表示為 $a+b$ 形式的元素的全體，其中 $a$ 遍歷 $\ideal{A}$ 的所有元素，$b$ 遍歷 $\ideal{B}$ 的所有元素；$\ideal{D}$ 又是一個理想。同樣，無窮多個理想的最大公因子——$\ideal{D}=(\ideal{A}_1,\ideal{A}_2,\ldots,\ideal{A}_\nu,\ldots)$——定義為所有元素 $d$ 的全體，這些 $d$ 可表示為每次從有限多個理想中取元素所得的和：
$d=a_{i_1}+a_{i_2}+\cdots+a_{i_n}$；這裡 $\ideal{D}$ 也仍然是一個理想。

% unit: noether.p19.par.026
特別地，若理想 $\ideal{M}$ 含有有限多個元素 $f_1,f_2,\ldots,f_\varrho$，使得
% unit: noether.p19.eq.002
\[
 \ideal{M}=(f_1,\ldots,f_\varrho),\qquad
 f=a_1f_1+\cdots+a_\varrho f_\varrho+n_1f_1+\cdots+n_\varrho f_\varrho
\]
對於每個 $f\equiv0(\ideal{M})$ 都成立，其中 $a_i$ 是環域中的量，$n_i$ 是整數，那麼 $\ideal{M}$ 稱為有限理想；$f_1,\ldots,f_\varrho$ 稱為一個理想基。

% unit: noether.p19.par.027
下面我們只以滿足有限性條件的環 $\Sigma$ 為基礎：\emph{$\Sigma$ 中的每個理想都是有限的，因此具有一個理想基。}

% unit: noether.p19.par.028
\textbf{3.} 從有限性條件直接推出支撐以下所有考察的結果。

% unit: noether.p19.par.029
\textbf{定理 I（有限鏈定理）。}\footnote{Dedekind 首先對數模提出：\emph{Zahlentheorie}, Suppl. XI, \S~172, 定理 VIII（第 4 版）；我們的證明和“鏈”這一名稱取自那裡。關於多項式中的理想，見 Lasker 前引處，第 56 頁（引理）。不過在兩種情形中，這個定理都只是零散地得到應用。我們的應用則始終依賴於選擇公設。} \emph{設 $\ideal{M},\ideal{M}_1,\ideal{M}_2,\ldots,\ideal{M}_\nu,\ldots$ 是 $\Sigma$ 中一個可數無窮的理想系統，其中每個理想都可被下一個理想整除。那麼從某個有限指標 $n$ 起，所有理想都相同，$\ideal{M}_n=\ideal{M}_{n+1}=\cdots$。換言之：若 $\ideal M,\ideal M_1,\ideal M_2,\ldots,\ideal M_\nu,\ldots$ 構成一個單純有序的理想鏈，並且每個理想都是緊前一個理想的真因子，則這個鏈在有限步後中止。}

% unit: noether.p19.par.030
事實上，令 $\ideal{D}=(\ideal{M},\ideal{M}_1,\ideal{M}_2,\ldots,\ideal{M}_\nu,\ldots)$ 為該系統的最大公因子，並令 $f_1\ldots f_k$ 為 $\ideal{D}$ 的一個基；由有限性條件這樣的基總是存在。由整除假設可知，$\ideal{D}$ 的每個元素同時也是鏈中某一個理想的元素；因為從
% unit: noether.p19.eq.003
\[
 f=g+h,\qquad g\equiv0(\ideal M_r),\qquad h\equiv0(\ideal M_s),\qquad (r\le s)
\]
得到 $g\equiv0(\ideal{M}_s)$，因此 $f\equiv0(\ideal{M}_s)$。當 $f$ 是若干組成部分之和時，同樣成立。因此存在有限指標 $n$，使得
% unit: noether.p19.eq.004
\[
 f_1\equiv0(\ideal{M}_n),\quad f_2\equiv0(\ideal{M}_n),\quad\ldots,\quad f_\varrho\equiv0(\ideal{M}_n).
\]
反過來由於 $\ideal{M}_n\equiv0(\ideal{D})$，得到 $\ideal{M}_n=\ideal{D}$；又由於 $\ideal{M}_n\equiv0(\ideal{D})$ 且 $\ideal{D}=\ideal{M}_n\equiv0(\ideal{M}_r)$，對每個 $r>n$ 也有 $\ideal{M}_r=\ideal{D}=\ideal{M}_n$，定理得證。

% unit: noether.p19.par.031
還應指出，反過來，理想基的存在又可由這個定理推出，因此有限性條件也本可以用這種不依賴基的形式來表述。
%
%
% unit: noether.p19.sec.002
\subsection*{\S~2. 把一個理想表示為有限多個不可約理想的最小公倍。}

% unit: noether.p19.par.039
理想 $\ideal B_1,\ideal B_2,\ldots,\ideal B_k$ 的最小公倍 $[\ideal B_1,\ideal B_2,\ldots,\ideal B_k]$，照通常的方式，定義為同時能被 $\ideal B_1$、$\ideal B_2$、……、$\ideal B_k$ 整除的元素全體；用符號寫為：
% unit: noether.p19.eq.018
\[
 \text{由 } f\equiv0(\ideal B_i),\quad (i=1,2,\ldots,k),
 \quad\text{推出}\quad
 f\equiv0([\ideal B_1,\ideal B_2,\ldots,\ideal B_k]),
\]
並且反過來也成立。最小公倍又是一個理想；理想 $\ideal B_i$ 也稱為該分解的組分。

% unit: noether.p19.par.040
\textbf{定義 I.} 表示 $\ideal M=[\ideal B_1\cdots\ideal B_k]$ 稱為約化表示，如果沒有一個 $\ideal B_i$ 包含在其餘理想的最小公倍 $\ideal A_i$ 中，並且沒有一個 $\ideal B_i$ 能由一個真因子替換。\footnote{一個非約化表示的例子是
\[
 (x^2,xy)=[(x),(x^2,xy,y^\lambda)]
\]
對於每個指數 $\lambda\ge2$ 都成立；對應於 $\lambda=1$ 的表示 $[(x),(x^2,y)]$ 是相應的一個約化表示。這個表示是戰死的 K. Hentzelt 告訴我的，作為準素理想分解不唯一性的最簡單例子。} 如果這些條件只對理想 $\ideal B_i$ 滿足，則稱該表示關於 $\ideal B_i$ 是約化的。最小公倍 $\ideal A_i=[\ideal B_1,\ldots,\ideal B_{i-1},\ideal B_{i+1},\ldots,\ideal B_k]$ 稱為 $\ideal B_i$ 的補理想。只滿足第一條件的表示稱為最短表示。

% unit: noether.p19.par.041
因此，在把一個理想表示為最小公倍時，依照下面的結果，只限於約化表示就足夠了。

% unit: noether.p19.thm.001
\textbf{引理 I.} \emph{一個理想作為有限多個理想的最小公倍的任一表示，至少能以一種方式由一個約化表示代替；這樣的表示尤其可以通過逐次分解達到。}

% unit: noether.p19.par.042
事實上，設 $\ideal M=[\ideal B_1^*\cdots\ideal B_l^*]$ 是 $\ideal M$ 的任一表示。我們依次刪去那些包含在保留下來的理想的最小公倍中的 $\ideal B_i$。由於其餘理想仍然給出 $\ideal M$，所以在這樣得到的表示
% unit: noether.p19.eq.019
\[
 \ideal M=[\ideal B_1\cdots\ideal B_k]=[\ideal A_i,\ideal B_i]
\]
中第一條件得到滿足，因此它是一個最短表示；並且如果把任意 $\ideal B_i$ 換成一個真因子，這個條件仍然得到滿足。第二條件則總是可以由有限鏈定理（定理 I）達到。若設
% unit: noether.p19.eq.020
\[
 \ideal B_i\supset \ideal B_i'\supset \ideal B_i''\supset\cdots\supset \ideal B_i^{(\nu)}\supset\cdots,
 \qquad
 \ideal M=[\ideal U_i,\ideal B_i]=[\ideal U_i,\ideal B_i']=\cdots=[\ideal U_i,\ideal B_i^{(\nu)}]=\cdots,
\]
其中每個 $\ideal B_i^{(\nu)}$ 都是其緊前一個理想的真因子，那麼由該定理，鏈 $\ideal B_i,\ideal B_i',\ldots,\ideal B_i^{(\nu)},\ldots$ 必須在有限步後中止；因此在表示 $\ideal M=[\ideal A_i,\ideal B_i^{(n)}]$ 中，理想 $\ideal B_i^{(\nu)}$ 不再能由真因子替換；如果 $\ideal A_i$ 被一個真因子替換，這一點當然仍然成立。把這個過程依次應用於每個 $\ideal B_i$，每次都用已經約化的 $\ideal B$ 來形成補理想，就得到一個約化表示。\footnote{這種表示並不由給定的表示唯一決定，這一點由前面的例子說明。對於 $(x^2,xy)=[(x),(x^2,xy,y^\lambda)]$，其中 $\lambda\ge2$，除了 $[(x),(x^2,y)]$ 之外，對任意 $\mu$，還有 $[(x),(x^2,\mu x+y)]$ 也是一個約化表示。}

% unit: noether.p19.par.043
為了逐次得到這樣的表示，還需要說明：由各個約化表示
% unit: noether.p19.eq.021
\[
 \ideal M=[\ideal B_1,\ideal C_1],\qquad
 \ideal C_1=[\ideal B_2,\ideal C_2],\quad\ldots,\quad
 \ideal C_{l-1}=[\ideal B_l,\ideal C_l]
\]
可以推出所得表示 $\ideal M=[\ideal B_1\cdots\ideal B_k,\ideal C_l]$ 也是約化的。為此，只需證明：由約化表示
% unit: noether.p19.eq.022
\[
 \ideal M=[\ideal B,\ideal C],\qquad \ideal C=[\ideal C_1,\ideal C_2],
\]
可得約化表示 $\ideal M=[\ideal B,\ideal C_1,\ideal C_2]$。實際上，按假設，沒有 $\ideal B$ 包含在它的補理想中；若這種情形對某個 $\ideal C_i$ 發生，則與第一個表示中的假設相反，$\ideal C$ 可以由一個真因子替換，因為由第二個約化表示，$\ideal C_1$ 和 $\ideal C_2$ 都是 $\ideal C$ 的真因子；所以這個表示是最短表示。又按假設，沒有 $\ideal B$ 能由一個真因子替換；若這對某個 $\ideal C_i$ 是可能的，由於 $\ideal C$ 的表示是約化的，這就會與假設相反地對應於把 $\ideal C$ 換成一個真因子。於是引理得證。

% unit: noether.p19.par.044
\textbf{定義 II.} 如果理想 $\ideal M$ 能表示為兩個真因子的最小公倍，則稱 $\ideal M$ 為可約；在相反情形下稱 $\ideal M$ 為不可約。

% unit: noether.p19.par.045
現在我們藉助有限鏈定理 I，並使用約化表示，證明：

% unit: noether.p19.thm.002
\textbf{定理 II.} \emph{每個理想都可以表示為有限多個不可約理想的最小公倍。}\footnote{這種表示一般並不唯一，這一點由前面的例子說明：對任意 $\mu$，$(x^2,xy)=[(x),(x^2,\mu x+y)]$。兩個組分對任意 $\mu$ 都是不可約的。事實上，$(x)$ 的所有因子都形如 $(x,g(y))$，其中 $g(y)$ 表示 $y$ 的一個多項式；所以其中任意兩個的最小公倍也具有這種形式，因而是 $(x)$ 的一個真因子。理想 $(x^2,\mu x+y)$ 只具有唯一的因子 $(x,y)$，所以也必然是不可約的。}

% unit: noether.p19.par.046
因為任意理想 $\ideal M$ 或者是不可約的；這時 $\ideal M=[\ideal M]$ 就是定理 II 所要求的表示；或者有 $\ideal M=[\ideal B_1,\ideal C_1]$，其中 $\ideal B_1$ 和 $\ideal C_1$ 是 $\ideal M$ 的真因子，並且由引理 I 可以假定該表示為約化表示。對於 $\ideal C_1$，同樣的二擇一情形成立：它或者是不可約的，或者存在一個約化表示 $\ideal C_1=[\ideal B_2,\ideal C_2]$。

% unit: noether.p19.par.047
這樣繼續下去，就得到一列約化表示
% unit: noether.p19.eq.023
\begin{equation}
 \ideal M=[\ideal B_1,\ideal C_1]=[\ideal B_1,\ideal B_2,\ideal C_2]=\cdots=[\ideal B_1,\ldots,\ideal B_n,\ideal C_n]=\cdots .
\tag{1}
\end{equation}
在鏈 $\ideal C_1,\ideal C_2,\ldots,\ideal C_\nu,\ldots$ 中，每個 $\ideal C_i$ 都是其緊前一個理想的真因子，因此該鏈在有限步後終止；存在一個指標 $n$，使得 $\ideal C_n$ 不可約。再由引理 I，表示 $\ideal M=[\ideal B_1\cdots\ideal B_n\ideal C_n]$ 是約化的；因此 $\ideal C_n$ 不能包含在它的補理想 $\ideal A_n$ 中，並且在表示 $\ideal M=[\ideal A_n,\ideal C_n]$ 中，理想 $\ideal C_n$ 不能由一個真因子替換。必要時把 $\ideal A_n$ 換成一個真因子，\footnote{事實上，這個表示關於 $\ideal A_n$ 也約化；這一點將在 \S~3（引理 IV）中作為引理 I 的逆命題得到證明。}從而使表示成為約化表示，由此可見，每個可約理想都容許作為一個不可約理想和一個與之互補的理想的最小公倍的約化表示。因此，在列 (1) 中，可以不失一般性地假定所有 $\ideal B_i$ 都不可約；重複上述論證，就得到一個不可約的 $\ideal C_n$ 的存在，從而證明了定理 II。
%
%
% unit: noether.p19.sec.003
\subsection*{\S~3. 兩種不同的既約理想分解中成分數的相等。}

% unit: noether.p19.par.050
為了證明成分數相等，首先必須用補元的性質來表達一個理想的可約性或既約性，如下述定理所示。

% unit: noether.p19.thm.003
\textbf{定理 III。}\footnote{定理 III 對應於 Schmeidler 以及 Noether--Schmeidler 的工作中從模到剩餘群的過渡（見序論）。$\ideal A$ 對應於剩餘群；$\ideal N_1$ 與 $\ideal N_2$ 對應於剩餘群所分解成的子群。}\emph{設最短表示 $\ideal M=[\ideal A,\ideal C]$ 關於 $\ideal C$ 是約化的。那麼 $\ideal C$ 可約的充要條件是，存在兩個理想 $\ideal N_1$ 與 $\ideal N_2$，它們是 $\ideal M$ 的真因子，並且滿足}
% unit: noether.p19.eq.024
\begin{equation}
 \ideal N_1\equiv0(\ideal A);\qquad
 \ideal N_2\equiv0(\ideal A);\qquad
 [\ideal N_1,\ideal N_2]=\ideal M.
\tag{2}
\end{equation}
\emph{由此還可推出：如果條件 (2) 滿足，而 $\ideal C$ 是既約的，那麼至少有一個 $\ideal N_i$ 不是 $\ideal M$ 的真因子；即 $\ideal N_i=\ideal M$。}

% unit: noether.p19.par.051
設 $\ideal C=[\ideal C_1,\ideal C_2]$，其中 $\ideal C_1,\ideal C_2$ 是 $\ideal C$ 的真因子。於是
% unit: noether.p19.eq.025
\[
 \ideal M=[\ideal A,\ideal C]=[\ideal A,\ideal C_1,\ideal C_2]
       =[[\ideal A,\ideal C_1],[\ideal A,\ideal C_2]].
\]
這裡理想 $[\ideal A,\ideal C_i]$ 是 $\ideal M$ 的真因子；否則 $[\ideal A,\ideal C]$ 關於 $\ideal C$ 就不是約化的。由於被 $\ideal A$ 整除的條件也滿足，條件 (2) 的必要性得證。（表示 (2) 不是約化的，因為某個 $[\ideal A,\ideal C_i]$ 可以由 $\ideal C_i$ 代替。）

% unit: noether.p19.par.052
反過來，設 (2) 滿足。構造理想
% unit: noether.p19.eq.026
\[
 \ideal C_1=(\ideal C,\ideal N_1),\qquad
 \ideal C_2=(\ideal C,\ideal N_2),\qquad
 \ideal C^*=[\ideal C_1,\ideal C_2].
\]
於是 $\ideal C$ 同時被 $\ideal C_1$ 和 $\ideal C_2$ 整除，因而也被它們的最小公倍數 $\ideal C^*$ 整除。為了證明 $\ideal C^*$ 被 $\ideal C$ 整除，設
% unit: noether.p19.eq.027
\[
 f\equiv0(\ideal C^*);\quad f\equiv0(\ideal C_1);\quad f\equiv0(\ideal C_2),
 \quad\text{也即}\quad f=c+n_1;\quad f=\bar c+n_2,
\]
其中 $c,\bar c,n_1,n_2$ 分別是 $\ideal C,\ideal N_1,\ideal N_2$ 中的量；特別地，$n_i$ 被 $\ideal A$ 整除。於是差
% unit: noether.p19.eq.028
\[
 g=c-\bar c=n_2-n_1
\]
既被 $\ideal C$ 整除，又被 $\ideal A$ 整除，因而被 $\ideal M$ 整除。又因為 $n_1=n_2+m$，所以 $n_1$（同樣 $n_2$）既被 $\ideal N_1$ 整除，又被 $\ideal N_2$ 整除，因而被 $\ideal M$ 整除。因此
% unit: noether.p19.eq.029
\[
 f=c+m;\qquad f\equiv0(\ideal C);\qquad \ideal C^*=\ideal C.
\]
此時 $\ideal C_1$ 和 $\ideal C_2$ 是 $\ideal C$ 的真因子；因為若 $\ideal C_i=(\ideal C,\ideal N_i)=\ideal C$，則 $\ideal N_i$ 被 $\ideal C$ 整除，又由於它被 $\ideal U$ 整除，便等於 $\ideal M$，這與假設矛盾。因此 $\ideal C=[\ideal C_1,\ideal C_2]$ 已被認作可約；定理 III 得證。

% unit: noether.p19.par.053
注意，幾乎同樣的思考還給出下述結果。

% unit: noether.p19.lem.002
\textbf{引理 II。}\emph{在一個最短表示 $\ideal M=[\ideal A,\ideal C]$ 中，如果理想 $\ideal C$ 能用一個真因子替換，那麼 $\ideal C$ 是可約的。}

% unit: noether.p19.par.054
也就是說，設
% unit: noether.p19.eq.030
\[
 \ideal M=[\ideal A,\ideal C]=[\ideal A,\ideal C_1],
\]
並置
% unit: noether.p19.eq.031
\[
 \ideal C^*=[\ideal C_1,(\ideal A,\ideal C)].
\]
那麼 $\ideal C$ 又被 $\ideal C^*$ 整除。由 $f\equiv0(\ideal C^*)$ 還得到
% unit: noether.p19.eq.032
\[
 f=c_1=a+c.
\]
差 $a=c_1-c$ 因而既被 $\ideal A$ 整除，又被 $\ideal C_1$ 整除，所以被 $\ideal M$ 整除；於是 $c_1=c+m$，$f\equiv0(\ideal C)$，$\ideal C^*=\ideal C$。由於按照假設，$\ideal C_1$ 和 $(\ideal A,\ideal C)$ 都是 $\ideal C$ 的真因子，故 $\ideal C=\ideal C^*$ 已被認作可約。

% unit: noether.p19.par.055
因此，既約的 $\ideal C$ 不能被一個真因子替換。

% unit: noether.p19.par.056
現在給定 $\ideal M$ 作為有限多個既約理想的最小公倍數的兩個不同最短表示：
% unit: noether.p19.eq.033
\[
 \ideal M=[\ideal B_1\cdots\ideal B_k]=[\ideal D_1\cdots\ideal D_l].
\]
根據引理 II 後面的註記，這兩個表示同時也是約化的。我們先證明：

% unit: noether.p19.lem.003
\textbf{引理 III。}\emph{對於每一個補元 $\ideal A_i=[\ideal B_1\cdots\ideal B_{i-1}\ideal B_{i+1}\cdots\ideal B_k]$，存在一個理想 $\ideal D_j$，使得 $\ideal M=[\ideal A_i,\ideal D_j]$。}

% unit: noether.p19.par.057
事實上，若置 $\ideal M=[\ideal D_1,\ideal C_1]$，$\ideal C_1=[\ideal D_2,\ideal C_{12}]$ 等等，則得到
% unit: noether.p19.eq.034
\[
 \ideal M=[\ideal A_i,\ideal M]=[\ideal A_i,\ideal D_1,\ideal C_1]
       =[[\ideal A_i\ideal D_1],[\ideal A_i\ideal C_1]].
\]
這裡對 $\ideal N_1=[\ideal A_i,\ideal D_1]$、$\ideal N_2=[\ideal A_i,\ideal C_1]$，定理 III 的條件 (2) 被滿足，因為 $\ideal M=[\ideal A_i,\ideal B_i]$ 關於 $\ideal B_i$ 是約化的，並且該表示是最短的。既然 $\ideal B_i$ 按假設是既約的，必有一個 $\ideal N_i$ 等於 $\ideal M$。

% unit: noether.p19.par.058
如果 $[\ideal U_i,\ideal D_1]=\ideal M$，引理即得證；如果 $[\ideal U_i,\ideal C_1]=\ideal M$，則相應地有 $\ideal M=[[\ideal U_i\ideal D_2],[\ideal U_i\ideal C_2]]$，按同樣的推理又必須有一個成分等於 $\ideal M$。如此繼續下去，或者得到 $\ideal M=[\ideal A_i,\ideal D_j]$，其中 $j<l$；或者得到 $\ideal M=[\ideal A_i,\ideal C_{1\ldots l-1}]$。由於 $\ideal C_{1\ldots l-1}=\ideal D_l$，引理便得證。

% unit: noether.p19.par.059
由此得到下述定理。

% unit: noether.p19.thm.004
\textbf{定理 IV。}\emph{一個理想作為既約理想的最小公倍數的兩個不同最短表示中，成分的個數相同。}

% unit: noether.p19.par.060
事實上，由引理對 $i=1$ 得
% unit: noether.p19.eq.035
\[
 \ideal M=[\ideal A_1,\ideal B_1]=[\ideal A_1,\ideal D_{j_1}]
          =[\ideal D_{j_1},\ideal B_2,\ldots,\ideal B_k].
\]
現在考慮兩個分解
% unit: noether.p19.eq.036
\[
 \ideal M=[\ideal D_{j_1},\ideal B_2,\ldots,\ideal B_k]
          =[\ideal D_1,\ideal D_2,\ldots,\ideal D_l],
\]
並對 $\ideal B_2$ 的補元 $\ideal U_2=[\ideal D_{j_1},\ideal B_3,\ldots,\ideal B_k]$ 重複上述推理，得到
% unit: noether.p19.eq.037
\[
 \ideal M=[\ideal U_2,\ideal B_2]=[\ideal U_2,\ideal D_{j_2}]
          =[\ideal D_{j_1},\ideal D_{j_2}\ideal B_3,\ldots,\ideal B_k],
\]
並且繼續這個過程得到
% unit: noether.p19.eq.038
\[
 \ideal M=[\ideal D_{j_1},\ideal D_{j_2},\ldots,\ideal D_{j_k}].
\]
由於假設中表示 $\ideal M=[\ideal D_1\cdots\ideal D_l]$ 是最短的，因此沒有任何一個 $\ideal D$ 可以省去；所以 $\ideal D_{j_i}$ 中彼此不同的理想必須窮盡所有 $\ideal D$，從而 $k\ge l$。在引理和隨後的推理中把 $\ideal B$ 與 $\ideal D$ 全部互換，就相應得到 $l\ge k$，因此 $k=l$，成分數相等得證。由此還推出，理想 $\ideal D_{j_i}$ 兩兩不同；否則用這些 $\ideal D_{j_i}$ 所作的最短表示會出現少於 $k$ 個成分。因此可以選擇記號，使 $\ideal D_{j_i}=\ideal D_i$。同理，所有中間表示 $\ideal M=[\ideal D_{j_1}\cdots\ideal D_{j_i},\ideal B_{i+1}\cdots\ideal B_k]$ 也都是最短的，並且按引理 II 後的註記是約化的。

% unit: noether.p19.par.061
成分數相等導出引理 I 的一個逆命題。

% unit: noether.p19.lem.004
\textbf{引理 IV。}\emph{在一個約化表示中，如果把成分合併成若干組，並形成每組的最小公倍數，那麼所得表示仍然是約化的。換言之，由一個約化表示}
% unit: noether.p19.eq.039
\[
 \ideal M=[\ideal C_{11}\cdots\ideal C_{1\mu_1};\;\ldots;\;,
           \ideal C_{\sigma1}\cdots\ideal C_{\sigma\mu_\sigma}]
\]
\emph{可推出，如果置 $\ideal N_i=[\ideal C_{i1}\cdots\ideal C_{i\mu_i}]$，則 $\ideal M=[\ideal N_1\cdots\ideal N_\sigma]=[\ideal N_i,\ideal L_i]$ 也是約化的。}

% unit: noether.p19.par.062
首先要注意，$\ideal N_i$ 不能包含在它的補元 $\ideal L_i$ 中，因為這對它的任何因子 $\ideal C_{ij}$ 都不會發生；所以該表示是最短的。為了證明 $\ideal N_i$ 不能被任何真因子替換，我們把各個 $\ideal C$ 分解為既約理想 $\ideal B$，\footnote{這裡始終指用 $\ideal B$ 作的最短、因而約化的表示。}於是按照引理 I 得到約化表示
% unit: noether.p19.eq.040
\[
 \ideal M=[\ideal B_{11}\cdots\ideal B_{1\lambda_1};\;\ldots;\;,
           \ideal B_{\sigma1}\cdots\ideal B_{\sigma\lambda_\sigma}],
 \qquad
 \ideal N_i=[\ideal B_{i1}\cdots\ideal B_{i\lambda_i}].
\]

% unit: noether.p19.par.063
現在設 $\ideal M=[\ideal N_i^*,\ideal L_i]$ 關於 $\ideal N_i^*$ 是約化的，其中 $\ideal N_i^*$ 是 $\ideal N_i$ 的一個真因子。由引理 II 得
% unit: noether.p19.eq.041
\[
 \ideal N_i=[\ideal N_i^*, (\ideal N_i,\ideal L_i)];
\]
而這個表示必然關於 $\ideal N_i^*$ 是約化的，否則在 $\ideal M$ 中 $\ideal N_i$ 也可由一個真因子替換。必要時再把 $(\ideal N_i,\ideal L_i)$ 也替換為一個真因子，從而得到 $\ideal N_i$ 的一個約化表示。現在把 $\ideal N_i$ 的兩個成分分解為既約理想，則 $\ideal N_i$ 的既約理想個數 $\lambda_i$ 按成分的個數相加而成；對應於真因子 $\ideal N_i^*$ 的既約理想個數必然小於 $\lambda_i$。可是這樣一來，把 $\ideal M=[\ideal N_i^*,\ideal L]$ 分解為既約理想，也會得到少於 $\sum_i\lambda_i$ 個理想，這與成分數相等矛盾。作為特殊情形 $\sigma=2$，還可推出表示 $\ideal M=[\ideal N_i,\ideal L_i]$ 關於補元 $\ideal L_i$ 也是約化的。
%
%
% unit: noether.p19.sec.004
\subsection*{\S~4. 準素理想。兩個不同不可約理想分解中相應素理想的唯一性。}

% unit: noether.p19.par.080
以下討論的是準素理想和不可約理想之間的聯繫。

% unit: noether.p19.def.003
\textbf{定義 III。} 若由 $a\cdot b\equiv0(\ideal Q)$ 且 $a\not\equiv0(\ideal Q)$ 必然推出 $b^\varkappa\equiv0(\ideal Q)$，其中指數 $\varkappa$ 是一個有限數，則理想 $\ideal Q$ 稱為準素的。

% unit: noether.p19.par.081
這個定義也可以這樣表述：若乘積 $a\cdot b$ 可被 $\ideal Q$ 整除，則或者某一個因子可被它整除，或者每一個因子的某個冪可被它整除。特別地，如果 $\varkappa$ 總可取為 $1$，則該理想稱為素理想。

% unit: noether.p19.par.082
由準素理想（相應地，素理想）的定義，藉助基的存在，得到只涉及理想乘積的定義\footnote{兩個理想 $\ideal A\cdot\ideal B$ 的乘積照通常意義理解為由所有量 $a\cdot b$ 及其有限和構成的理想。}：

% unit: noether.p19.def.003a
\textbf{定義 IIIa。} 若由 $\ideal A\cdot\ideal B\equiv0(\ideal Q)$ 且 $\ideal A\not\equiv0(\ideal Q)$ 必然推出 $\ideal B^\lambda\equiv0(\ideal Q)$，則理想 $\ideal Q$ 稱為準素的。如果 $\lambda$ 總是等於 $1$，則該理想稱為素理想。因此，對素理想 $\ideal P$，由 $\ideal A\cdot\ideal B\equiv0(\ideal P)$ 且 $\ideal A\not\equiv0(\ideal P)$ 總可推出 $\ideal B\equiv0(\ideal P)$。

% unit: noether.p19.par.083
因為在 IIIa 中取 $\ideal A=(a)$、$\ideal B=(b)$ 時，定義 III 作為特殊情形被包含在內，所以按 IIIa 為準素的每一個理想也按 III 為準素。反過來，設 $\ideal Q$ 按 III 為準素，並設 IIIa 的假設 $\ideal A\cdot\ideal B\equiv0(\ideal Q)$ 成立。於是或者推出 $\ideal A\equiv0(\ideal Q)$，或者至少存在一個量 $a\not\equiv0(\ideal Q)$，使得 $a\cdot\ideal B\equiv0(\ideal Q)$ 且 $a\not\equiv0(\ideal Q)$。若 $b_1\ldots b_r$ 是 $\ideal B$ 的一個理想基，則因 $a\cdot b_i\equiv0(\ideal Q)$，由定義 III 得到
% unit: noether.p19.eq.040
\[
 b_1^{\varkappa_1}\equiv0(\ideal Q),\quad \ldots,\quad b_r^{\varkappa_r}\equiv0(\ideal Q).
\]
由於
% unit: noether.p19.eq.041
\[
 b=f_1b_1+\cdots+f_rb_r+n_1b_1+\cdots+n_rb_r,
\]
於是對於 $\lambda=\varkappa_1+\cdots+\varkappa_r$，任意 $\lambda$ 個量 $b$ 的乘積都可被 $\ideal Q$ 整除；由此證明了按 III 為準素的理想滿足定義 IIIa。特別對素理想 $\ideal P$，由 $a\cdot\ideal B\equiv0(\ideal P)$，從而對每個 $b\equiv0(\ideal B)$ 有 $a\cdot b\equiv0(\ideal P)$，並且 $a\not\equiv0(\ideal P)$，可推出 $b\equiv0(\ideal P)$，因而 $\ideal B\equiv0(\ideal P)$。這樣就證明了兩個定義的等價性。

% unit: noether.p19.par.084
準素理想和素理想之間的聯繫由如下觀察給出：所有具有“某個冪可被 $\ideal Q$ 整除”這一性質的元素 $p$ 構成的整體 $\ideal P$ 是一個素理想。首先顯然 $\ideal P$ 是理想；因為與 $p_1$ 和 $p_2$ 一起，$ap_1$ 以及 $(p_1-p_2)$ 也具有上述性質。再按定義 IIIa 中所用的基論證，還得到存在一個數 $\lambda$，使得 $\ideal P^\lambda\equiv0(\ideal Q)$。

% unit: noether.p19.par.085
現在設
% unit: noether.p19.eq.042
\[
 a\cdot b\equiv0(\ideal P);\qquad a\not\equiv0(\ideal P).
\]
於是由 $\ideal P$ 的定義得到
% unit: noether.p19.eq.043
\[
 a^\lambda\cdot b^\lambda\equiv0(\ideal Q);\qquad a^\lambda\not\equiv0(\ideal Q);
\]
因此由 $\ideal Q$ 的定義得到
% unit: noether.p19.eq.044
\[
 b^{\lambda\varkappa}\equiv0(\ideal Q)\qquad\text{從而}\qquad b\equiv0(\ideal P),
\]
由此證明 $\ideal P$ 是素理想。$\ideal P$ 也可定義為所有具有“某個冪可被 $\ideal Q$ 整除”這一性質的理想 $\ideal B$ 的最大公因子。因為每個這樣的 $\ideal B$ 按定義都可被 $\ideal P$ 整除；所以這些 $\ideal B$ 的最大公因子 $\ideal D$ 也可被 $\ideal P$ 整除。反過來，$\ideal P$ 本身就是這樣的一個理想 $\ideal B$，因而可被 $\ideal D$ 整除，由此得出 $\ideal P=\ideal D$。

% unit: noether.p19.par.086
所以 $\ideal P$ 是一個素理想，它是 $\ideal Q$ 的因子，並且它的某個冪可被 $\ideal Q$ 整除；這一性質唯一地定義了它。事實上，由
% unit: noether.p19.eq.045
\[
 \ideal Q\equiv0(\ideal P);\qquad \ideal P^\lambda\equiv0(\ideal Q)
 \ideal Q\equiv0(\bar{\ideal P});\qquad \bar{\ideal P}^{\mu}\equiv0(\ideal Q)
\]
推出
% unit: noether.p19.eq.046
\[
 \ideal P^\lambda\equiv0(\bar{\ideal P});\qquad
 \bar{\ideal P}^{\mu}\equiv0(\ideal P),
\]
於是由素理想的性質推出
% unit: noether.p19.eq.047
\[
 \ideal P\equiv0(\bar{\ideal P});\qquad
 \bar{\ideal P}\equiv0(\ideal P);\qquad
 \ideal P=\bar{\ideal P}.
\]

% unit: noether.p19.thm.005
綜上可得：

\textbf{定理 V。} \emph{對於每一個準素理想 $\ideal Q$，存在且只存在一個素理想 $\ideal P$，它是 $\ideal Q$ 的因子，並且它的某個冪可被 $\ideal Q$ 整除；$\ideal P$ 稱為 $\ideal Q$ 的“相應素理想”\footnote{反命題不成立可由例子 $\ideal M=(x^2,xy)$ 看出。素理想 $(x)$ 滿足全部條件，但 $\ideal M$ 不是準素的。}。$\ideal P$ 定義為所有具有“某個冪可被 $\ideal Q$ 整除”這一性質的理想 $\ideal B$ 的最大公因子。若 $\varrho$ 是滿足 $\ideal P^\varrho\equiv0(\ideal Q)$ 的最小數，則稱 $\varrho$ 為 $\ideal Q$ 的指數\footnote{一般並不像在有理整數或代數整數的領域中那樣有 $\ideal P^\varrho=\ideal Q$；例子
\[
 \ideal Q=(x^2,y),\qquad \ideal P=(x,y),\qquad
 \ideal P^2=(x^2,xy,y^2)\equiv0(\ideal Q),
 \quad \text{但 }\ideal Q\not\equiv0(\ideal P^2);
\]
說明了這一點。因此 $\ideal Q$ 與 $\ideal P^2$ 不同。}。}

% unit: noether.p19.par.087
現在我們證明準素性與不可約性之間的聯繫：

% unit: noether.p19.thm.006
\textbf{定理 VI。} \emph{每一個非準素理想都是可約的；換言之，每一個不可約理想都是準素的\footnote{這裡反命題不成立，例如由
\[
 \ideal Q=(x^2,xy,y^\lambda)=[(x^2,y),(x,y^\lambda)],
\]
其中 $\lambda\ge2$，可以看出。這裡 $\ideal Q$ 是準素的，但可約。$\ideal Q$ 是準素的，是因為它包含 $x,y$ 的所有 $\lambda$ 維冪乘積；因此每個無常數項的多項式都有某個冪可被 $\ideal Q$ 整除。但若在 $a\cdot b\equiv0(\ideal Q)$ 中，多項式 $b$ 含有常數項——於是對每個 $\varkappa$ 都有 $b^\varkappa\not\equiv0(\ideal Q)$——那麼由於 $\ideal Q$ 的基多項式的齊次性，$a\cdot b$ 的每一個齊次組成部分都可被 $\ideal Q$ 整除，因而 $a$ 必須可被 $\ideal Q$ 整除。}。}

% unit: noether.p19.par.088
設 $\ideal K$ 是一個非準素理想。於是按定義 III 至少存在一對量 $a,b$，使得
% unit: noether.p19.eq.048
\begin{equation}
 a\cdot b\equiv0(\ideal K),\qquad
 a\not\equiv0(\ideal K),\qquad
 b^\varkappa\not\equiv0(\ideal K)\quad \text{對每個 }\varkappa.
\tag{3}
\end{equation}

% unit: noether.p19.par.089
現在構造兩個理想
% unit: noether.p19.eq.049
\[
 \ideal L_0=(\ideal K,a);\qquad \ideal N_0=(\ideal K,b),
\]
它們按 (3) 是 $\ideal K$ 的真因子，並且由 (3) 有
% unit: noether.p19.eq.050
\begin{equation}
 \ideal L_0\cdot\ideal N_0\equiv0(\ideal K).
\tag{4}
\end{equation}

% unit: noether.p19.par.090
對於最小公倍數 $\ideal K_0=[\ideal L_0,\ideal N_0]$ 的元素 $f$，現在有如下二擇一情形。或者從
% unit: noether.p19.eq.051
\[
 f\equiv0(\ideal L_0),\qquad f\equiv0(\ideal N_0),\qquad \text{即}\qquad f\equiv a_1\cdot b\pmod{\ideal K}
\]
總可推出一個表示
% unit: noether.p19.eq.052
\[
 f\equiv l_0\cdot b\pmod{\ideal K},\qquad l_0\equiv0(\ideal L_0);
\]
於是按 (4) 有 $f\equiv0(\ideal K)$，從而 $\ideal K_0\equiv0(\ideal K)$；又因 $\ideal K\equiv0(\ideal K_0)$，也有 $\ideal K=\ideal K_0$，這就證明了 $\ideal K$ 可約。

% unit: noether.p19.par.091
或者至少存在一個 $f\equiv0(\ideal K_0)$，對它不存在這樣的 $l_0$。於是以屬於這個 $f$ 的 $a_1$ 構造
% unit: noether.p19.eq.053
\[
 \ideal L_1=(\ideal L_0,a_1)=(\ideal K,a,a_1),\qquad
 \ideal N_1=(\ideal K,b^2).
\]
那麼由於 $a_1\cdot b\equiv0(\ideal L_0)$，由 (4) 也有
% unit: noether.p19.eq.054
\begin{equation}
 \ideal L_1\cdot\ideal N_1\equiv0(\ideal K),
\tag{4'}
\end{equation}
並且 $\ideal L_1$ 是 $\ideal L_0$ 的真因子。

% unit: noether.p19.par.092
對於 $\ideal K_1=[\ideal L_1,\ideal N_1]$ 的元素 $f$，同樣的二擇一情形成立。或者從
% unit: noether.p19.eq.055
\[
 f\equiv0(\ideal L_1),\qquad f\equiv0(\ideal N_1),
\]
即從 $f\equiv a_2\cdot b^2\pmod{\ideal K}$，總可推出
% unit: noether.p19.eq.056
\[
 f\equiv l_1\cdot b^2,
 \qquad l_1\equiv0(\ideal L_1),
\]
於是按 (4') 有
% unit: noether.p19.eq.057
\[
 \ideal K=\ideal K_1.
\]

% unit: noether.p19.par.093
或者至少有一個 $f$ 不存在這樣的 $l_1$；這就導致構造 $\ideal L_2=(\ideal L_1,a_2)$、$\ideal N_2=(\ideal K,b^{2^2})$，並且有 $\ideal L_2\cdot\ideal N_2\equiv0(\ideal K)$，其中 $\ideal L_2$ 是 $\ideal L_1$ 的真因子。如此繼續，一般地定義
% unit: noether.p19.eq.058
\[
 \ideal L_0=(\ideal K,a),\quad
 \ideal L_1=(\ideal L_0,a_1);\quad\ldots;\quad
 \ideal L_\nu=(\ideal L_{\nu-1},a_\nu);\quad\ldots;
\]
% unit: noether.p19.eq.059
\[
 \ideal N_0=(\ideal K,b),\quad
 \ideal N_1=(\ideal K,b^2),\quad \ldots,\quad
 \ideal N_\nu=(\ideal K,b^{2^\nu}),\quad \ldots,
\]
其中 $a_i$ 由存在一個 $f$ 來定義，使得
% unit: noether.p19.eq.060
\[
 f\equiv0(\ideal L_{i-1}),\qquad f\equiv0(\ideal N_{i-1}),
\]
即
% unit: noether.p19.eq.061
\[
 f\equiv a_i\cdot b^{2^{i-1}}\pmod{\ideal K},\qquad
 a_i\not\equiv0(\ideal L_{i-1}).
\]
這樣一般有 $\ideal L_i\cdot\ideal N_i\equiv0(\ideal K)$；$\ideal N_i$ 按 (3) 是 $\ideal K$ 的真因子，而 $\ideal L_i$ 是 $\ideal L_{i-1}$ 的真因子。由有限鏈定理 I，$\ideal L$ 的鏈必須在有限步後終止，設終止於 $\ideal L_n$。於是對每個 $f\equiv0(\ideal L_n)$、$f\equiv0(\ideal N_n)$，都有 $f\equiv l_n\cdot b^{2^n}\;(\ideal K)$ 且 $l_n\equiv0(\ideal L_n)$；因而由上面的論證得到 $\ideal K=[\ideal L_n,\ideal N_n]$，這就證明了 $\ideal K$ 可約。

% unit: noether.p19.par.094
由剛才證明的結果，相應素理想的唯一性如下推出。設
% unit: noether.p19.eq.062
\[
 \ideal M=[\ideal B_1\cdots\ideal B_k]=[\ideal D_1\cdots\ideal D_k]
\]
是 $\ideal M$ 作為不可約理想的最小公倍數的兩個最短、因而約化的表示；它們的項數按定理 IV 一致。於是按該定理，其中出現的中間表示（如那裡所指出，可取指標 $j_i=i$）也是最短表示：
% unit: noether.p19.eq.063
\[
\begin{aligned}
 \ideal M&=[\ideal D_1\cdots\ideal D_{i-1}\ideal B_i\ideal B_{i+1}\cdots\ideal B_k]
       =[\ideal D_1\cdots\ideal D_{i-1}\ideal D_i\ideal B_{i+1}\cdots\ideal B_k]\\
       &=[\bar{\ideal A}_i,\ideal B_i]=[\bar{\ideal A}_i,\ideal D_i].
\end{aligned}
\]
於是得到
% unit: noether.p19.eq.064
\[
 \overline{\ideal A}_i\cdot\ideal B_i\equiv0(\ideal D_i),\qquad
 \overline{\ideal A}_i\not\equiv0(\ideal D_i),
 \qquad
 \overline{\ideal A}_i\cdot\ideal D_i\equiv0(\ideal B_i),\qquad
 \overline{\ideal A}_i\not\equiv0(\ideal B_i).
\]

% unit: noether.p19.par.095
現在，由定理 VI，不可約理想 $\ideal B_i$ 和 $\ideal D_i$ 是準素的；於是存在兩個數 $\lambda_i$ 和 $\mu_i$，使得
% unit: noether.p19.eq.065
\begin{equation}
 \ideal B_i^{\lambda_i}\equiv0(\ideal D_i),
 \qquad
 \ideal D_i^{\mu_i}\equiv0(\ideal B_i).
\tag{5}
\end{equation}
分別以 $\ideal P_i$ 和 $\bar{\ideal P}_i$ 表示 $\ideal B_i$ 和 $\ideal D_i$ 的相應素理想；也就是說，$\ideal P_i^{\varrho_i}\equiv0(\ideal B_i)$，$\bar{\ideal P}_i^{\sigma_i}\equiv0(\ideal D_i)$。由 (5) 得到
% unit: noether.p19.eq.066
\[
 \ideal P_i^{\lambda_i\varrho_i}\equiv0(\bar{\ideal P}_i),
 \qquad
 \bar{\ideal P}_i^{\mu_i\sigma_i}\equiv0(\ideal P_i),
\]
於是由素理想的性質得到
% unit: noether.p19.eq.067
\[
 \ideal P_i\equiv0(\bar{\ideal P}_i),
 \qquad
 \bar{\ideal P}_i\equiv0(\ideal P_i),
 \qquad
 \ideal P_i=\bar{\ideal P}_i.
\]
這就證明了：

% unit: noether.p19.thm.007
\textbf{定理 VII。} \emph{在一個理想作為不可約理想的最小公倍數的兩個不同最短表示中，相應素理想一致；其中也可以有相同的素理想出現\footnote{例如注 19 的例子說明了這一點：
\[
 (x^2,xy,y^\lambda)=[(x^2,y),(x,y^\lambda)],
\]
其中 $\lambda\ge2$。這裡兩個相應素理想都是 $(x,y)$。}，並且在每一個分解中出現的次數相同。因此，理想本身至少可以以一種方式兩兩對應，使得每次一個理想 $\ideal B_i$ 的某個冪可被與它對應的 $\ideal D_i$ 整除，反過來也成立。它們的個數按定理 IV 一致\footnote{關於不可約理想中所含“孤立”理想的唯一性，參見 \S~7。}。}
%
%
% unit: noether.p19.sec.005
\subsection*{\S~5. 把一個理想表示為最大準素理想的最小公倍數。相應素理想的唯一性。}

% unit: noether.p19.def.004
\textbf{定義 IV。} \emph{如果一個最短表示 $\ideal M=[\ideal Q_1\cdots\ideal Q_\alpha]$ 中所有 $\ideal Q$ 都是準素的，但任意兩個 $\ideal Q$ 的最小公倍數不再是準素的，那麼這個表示稱為最大準素理想的最小公倍數表示。}

% unit: noether.p19.par.096
至少存在一個這樣的表示，這是由把 $\ideal M$ 表為不可約理想的最小公倍數所得出的。因為這些理想都是準素的；或者此時已經有了由最大準素理想給出的表示，或者任意兩個理想的最小公倍數又成為準素的。由於這裡理想的個數減少了一個，重複這個過程便在有限步之後得到所求的表示。

% unit: noether.p19.par.097
按引理 IV，這個表示是既約的。反過來，每一個由最大準素理想給出的既約表示都以這種方式產生，正如把 $\ideal Q$ 分解為不可約理想所顯示的那樣。

% unit: noether.p19.par.098
為了在這裡由定理 VII 推出相應的唯一性定理，必須按照下面的定理研究它同相應素理想的聯繫。

% unit: noether.p19.thm.008
\textbf{定理 VIII。} \emph{若準素理想 $\ideal N_1,\ideal N_2,\ldots,\ideal N_\lambda$ 都有同一個相應素理想 $\ideal P$，則它們的最小公倍數 $\ideal Q=[\ideal N_1,\ideal N_2,\ldots,\ideal N_\lambda]$ 也是準素的，並且以 $\ideal P$ 為相應素理想。反過來，若 $\ideal Q=[\ideal N_1\cdots\ideal N_\lambda]$ 是準素理想 $\ideal Q$ 的一個既約表示，則所有 $\ideal N_i$ 都是準素的，並且以 $\ideal Q$ 的相應素理想 $\ideal P$ 作為它們的相應素理想。}

% unit: noether.p19.par.099
為證明斷言的第一部分，首先注意：若對每個 $i$ 有 $\ideal P^{\varrho_i}\equiv0(\ideal N_i)$，則也有 $\ideal P^\tau\equiv0(\ideal Q)$，其中 $\tau$ 表示這些指數 $\varrho_i$ 中最大的一個。又因為 $\ideal P$ 也是 $\ideal Q$ 的因子，所以若 $\ideal Q$ 是準素的，則 $\ideal P$ 必然就是相應素理想。由
% unit: noether.p19.eq.068
\[
 \ideal A\cdot\ideal B\equiv0(\ideal Q),\qquad
 \ideal B^k\not\equiv0(\ideal Q)\quad\text{（對每個 }k\text{）}
\]
因此推出
% unit: noether.p19.eq.069
\[
 \ideal B\not\equiv0(\ideal P);
 \quad\text{所以}\quad
 \ideal B^k\not\equiv(\ideal N_i)\quad\text{（對每個 }k\text{）};
\]
從而 $\ideal A\equiv0(\ideal N_i)$；於是 $\ideal A\equiv0(\ideal Q)$，這就證明了 $\ideal Q$ 是準素的，並且 $\ideal P$ 是相應素理想。

% unit: noether.p19.par.100
反過來，先設
% unit: noether.p19.eq.070
\[
 \ideal Q=[\ideal N_1\cdots\ideal N_\lambda]=[\ideal N_i,\ideal L_i]
\]
是 $\ideal Q$ 由準素理想 $\ideal N_i$ 給出的一個最短表示，並且分別以 $\ideal P_i$ 為相應素理想。由
% unit: noether.p19.eq.071
\[
 \ideal L_i\cdot\ideal N_i\equiv0(\ideal Q);\qquad
 \ideal L_i\not\equiv0(\ideal Q)
\]
（由於表示是最短的）於是得到
% unit: noether.p19.eq.072
\[
 \ideal N_i^{\sigma_i}\equiv0(\ideal Q);
 \qquad\text{或者}\qquad
 \ideal P_i^{\varrho_i\sigma_i}\equiv0(\ideal Q).
\]
由於 $\ideal P_i$ 同時也是 $\ideal Q$ 的因子，所以對每個 $i$，$\ideal P_i$ 都與 $\ideal Q$ 的相應素理想 $\ideal P$ 相同。

% unit: noether.p19.par.101
還要證明，在每一個既約表示 $\ideal Q=[\ideal N_1\cdots\ideal N_\lambda]$ 中，各 $\ideal N_i$ 都是準素的。\footnote{這裡既約表示是本質的，這由非既約表示的例子說明：
\[
 \ideal Q=[x^2,xy,y^\lambda]=[(x^2,xy,y^2,yz),(x,y^\lambda)],
 \qquad \lambda\ge2.
\]
這裡 $(x^2,xy,y^2,yz)=[(x^2,y),(x,y^2,z)]$ 按上面所證並不是準素的；因為後一個表示是由準素理想給出的最短表示，但相應素理想 $(x,y)$ 和 $(x,y,z)$ 不同。（按注 19，$\ideal Q$ 是準素的。）} 為此，把各 $\ideal N_i$ 分解為它們的不可約理想 $\ideal B$；在這樣得到的最短表示
% unit: noether.p19.eq.073
\[
 \ideal Q=[\ideal B_1\cdots\ideal B_\mu]
\]
中，每個理想都是準素的，並且按剛剛證明的結果以 $\ideal P$ 為相應素理想。於是按上面已經證明的定理第一部分，對每個 $\ideal N_i$ 也有同樣的結論，從而定理 VIII 完全得到證明。

% unit: noether.p19.par.102
\textbf{附加說明。} 由此還推出，一個素理想必然是不可約的。因為由既約表示 $\ideal P=[\ideal N,\ideal L]$，按定理 VIII 得到 $\ideal N\equiv0(\ideal P)$；所以由於 $\ideal P\equiv0(\ideal N)$，也有 $\ideal P=\ideal N$，同樣有 $\ideal P=\ideal L$。$\ideal P$ 的不可約性也可直接得到：因為由 $\ideal P=[\ideal N,\ideal L]$ 可推出 $\ideal N\cdot\ideal L\equiv0(\ideal P)$，$\ideal N\not\equiv(\ideal P)$，$\ideal L\not\equiv0(\ideal P)$，這同素理想的定義性質相矛盾。

% unit: noether.p19.par.103
現在設
% unit: noether.p19.eq.074
\[
 \ideal M=[\ideal Q_1\cdots\ideal Q_\alpha]=[\bar{\ideal Q}_1\cdots\bar{\ideal Q}_\beta]
\]
是 $\ideal M$ 作為最大準素理想的最小公倍數的兩個既約表示。把各 $\ideal Q$ 分解為它們的不可約理想 $\ideal B$，把各 $\bar{\ideal Q}$ 相應地分解為不可約理想 $\ideal D$，便得到 $\ideal M$ 作為不可約理想的最小公倍數的兩個既約表示；在這些表示中，按定理 VII，分量的數目以及相應素理想都一致。按定理 VIII，出現在一個固定 $\ideal Q_i$ 中的所有不可約理想 $\ideal B$ 都有同一個相應素理想 $\ideal P_i$；而屬於 $\ideal Q_k$ 的 $\ideal P_k$ 必然與它不同，因為否則按定理 VIII 就不會有一個由最大準素理想給出的表示。於是 $\ideal Q$ 的個數 $\alpha$ 等於 $\ideal B$ 的不同相應素理想 $\ideal P$ 的個數；這些不同的 $\ideal P$ 構成 $\ideal Q$ 的相應素理想。對於 $\bar{\ideal Q}$ 相對於它們分解成 $\ideal D$ 的情形，同樣成立。於是由定理 VII 推出 $\ideal Q$ 和 $\bar{\ideal Q}$ 的數目相等，並且它們的相應素理想一致。同時還看出，在本節開頭指出的把不可約理想合併為最大準素理想的過程，正是把所有具有同一個相應素理想的理想，並且只把這些理想，合併在一起。定理 VIII 進一步顯示了最大準素理想的不可約性性質：它們不允許作為最大準素理想的最小公倍數的既約表示。

% unit: noether.p19.par.104
綜上，已經證明：

% unit: noether.p19.thm.009
\textbf{定理 IX。} \emph{在一個理想作為最大準素理想的最小公倍數的兩個既約表示中，分量的數目以及相應素理想一致，而這些相應素理想彼此都不同。換句話說：可以把每個 $\ideal Q$ 與唯一的一個 $\bar{\ideal Q}$ 對應起來，使得 $\ideal Q$ 的某個冪可被 $\ideal Q'$ 整除，反過來也成立。}\footnote{不同表示的一個例子是注 12 在定理 II 中給出的：
\[
 (x^2,xy)=[(x),(x^2,\mu x+y)]
\]
其中 $\mu$ 任意。由於相應素理想 $\ideal P_1=(x)$、$\ideal P_2=(x,y)$ 彼此不同，所以這裡涉及的是最大準素理想。--- 關於“孤立的”最大準素理想的唯一性，參見 \S~7。} \emph{--- $\ideal Q$ 和 $\ideal Q'$ 關於分解為最大準素理想具有不可約性性質。}

% unit: noether.p19.par.105
\textbf{附加說明。} 還要指出，如果不假設表示既約而只假設它最短，定理 IX 在本質上仍保持成立。若例如
% unit: noether.p19.eq.075
\[
 \ideal M=[\ideal Q_1\cdots\ideal Q_i^*\cdots\ideal Q_\alpha]
\]
關於 $\ideal Q_i^*$ 是既約的，且 $\ideal Q_i^*$ 是 $\ideal Q_i$ 的真因子，$\ideal L_i$ 是 $\ideal Q_i$ 的補，則按引理 IV 得到
% unit: noether.p19.eq.076
\[
 \ideal Q_i=[\ideal Q_i^*,(\ideal U_i,\ideal Q_i)],
\]
而這個表示關於 $\ideal Q_i^*$ 是既約的。按定理 VIII，在應用它時必要時要把 $(\ideal L_i,\ideal Q_i)$ 替換為一個真因子，於是 $\ideal Q_i^*$ 是準素的，並且具有同 $\ideal Q_i$ 相同的相應素理想 $\ideal P_i$。繼續這個過程表明，對每一個這樣的表示，都可以指定一個由最大準素理想給出的既約表示，使得分量的數目和相應素理想一致。因此，即使在最短表示的情形下，兩個不同表示中的分量數目和相應素理想也一致。

% unit: noether.p19.par.106
這樣唯一確定的、彼此不同的相應素理想，簡稱為“$\ideal M$ 的相應素理想”。
%
%
% unit: noether.p19.sec.006
\subsection*{\S~6. 把一個理想唯一表示為相對素-不可約理想的最小公倍數。}

% unit: noether.p19.def.005
\textbf{定義 V。} \emph{若由
% unit: noether.p19.eq.077
\[
 \ideal T\ideal R\equiv0(\ideal S)
\]
必然推出 $\ideal T\equiv0(\ideal S)$，則稱理想 $\ideal R$ 對 $\ideal S$ 是相對素的。若 $\ideal R$ 對 $\ideal S$ 相對素，同時 $\ideal S$ 對 $\ideal R$ 相對素，則稱 $\ideal R$ 和 $\ideal S$ 互相相對素。}\footnote{相對素條件並不對稱。例如 $\ideal R=(x^2,y)$ 對 $\ideal S=(x)$ 是相對素的；但是 $\ideal S$ 對 $\ideal R$ 不是相對素的，因為 $\ideal S^2\equiv0(\ideal R)$，而 $\ideal S\not\equiv0(\ideal R)$。}

% unit: noether.p19.par.107
\emph{如果一個理想不能表示為互相相對素的真因子的最小公倍數，則稱它為相對素-不可約的。}

% unit: noether.p19.par.108
特別地，若不用 $\ideal T$，而取所有滿足 $\ideal T\ideal R\equiv0(\ideal S)$ 的 $\ideal T$ 的最大公因子 $\ideal T_0$，則也有 $\ideal T_0\ideal R\equiv0(\ideal S)$ 且 $\ideal S\equiv0(\ideal T_0)$。因此，當 $\ideal R$ 對 $\ideal S$ 相對素時，$\ideal T_0=\ideal S$；當 $\ideal R$ 對 $\ideal S$ 不是相對素時，$\ideal T_0$ 是 $\ideal S$ 的真因子。\footnote{這樣定義的 $\ideal T_0$ 與 Lasker 的“剩餘模”相同；而在把它推廣到數模而非理想時，它與 Dedekind 的兩個模的“商”相同。Lasker，前引處，第 49 頁；Dedekind（\emph{Zahlentheorie}），第 504 頁。}

% unit: noether.p19.par.109
唯一性證明基於下列定理。

% unit: noether.p19.thm.010
\textbf{定理 X。} \emph{1. 若 $\ideal R$ 對理想 $\ideal S_1\ldots\ideal S_\lambda$ 都是相對素的，則 $\ideal R$ 對它們的最小公倍數 $\ideal S$ 也是相對素的。}

% unit: noether.p19.thm.010b
\emph{2. 若理想 $\ideal S_1\ldots\ideal S_\lambda$ 對 $\ideal R$ 都是相對素的，則它們的最小公倍數 $\ideal S$ 對 $\ideal R$ 也是相對素的。}

% unit: noether.p19.thm.010c
\emph{3. 若 $\ideal R$ 對 $\ideal S$ 是相對素的，而 $\ideal S=[\ideal S_1\cdots\ideal S_\lambda]$ 是 $\ideal S$ 的一個既約表示，則 $\ideal R$ 對每個 $\ideal S_i$ 也是相對素的。}

% unit: noether.p19.thm.010d
\emph{4. 若 $\ideal S$ 對 $\ideal R$ 是相對素的，則 $\ideal S$ 的每個因子 $\ideal S_i$ 對 $\ideal R$ 也是相對素的。}

% unit: noether.p19.par.110
1. 因為由 $\ideal T\cdot\ideal R\equiv0(\ideal S)$ 必然推出 $\ideal T\cdot\ideal R\equiv0(\ideal S_i)$，所以按假設得到 $\ideal T\equiv0(\ideal S_i)$，從而 $\ideal T\equiv0(\ideal S)$。

% unit: noether.p19.par.111
2. 置
% unit: noether.p19.eq.078
\[
 \ideal C_1=[\ideal S_2\cdots\ideal S_\lambda]
 \ideal C_{12}=[\ideal S_3\cdots\ideal S_\lambda];\ \ldots
 \ideal C_{12\ldots\lambda-1}=\ideal S_\lambda.
\]
由 $\ideal T\cdot\ideal S\equiv0(\ideal R)$ 得到 $\ideal T\cdot\ideal C_1\cdot\ideal S_1\equiv0(\ideal R)$；所以按假設有 $\ideal T\cdot\ideal C_1\equiv0(\ideal R)$。再由此得到 $\ideal T\cdot\ideal C_{12}\cdot\ideal S_2\equiv0(\ideal R)$，所以 $\ideal T\cdot\ideal C_{12}\equiv0(\ideal R)$，最後得到 $\ideal T\cdot\ideal C_{12\ldots\lambda-1}=\ideal T\cdot\ideal S_\lambda\equiv0(\ideal R)$；因此 $\ideal T\equiv0(\ideal R)$。

% unit: noether.p19.par.112
3. 令 $\widehat{\ideal S}_i$ 表示 $\ideal S_i$ 的補；由 $\ideal T_0\cdot\ideal R\equiv0(\ideal S_i)$，再結合 $\ideal C_i\cdot\ideal R\equiv0(\ideal C_i)$，得到 $[\ideal T_0,\ideal C_i]\cdot\ideal R\equiv0(\ideal S)$。若 $\ideal T_0$ 是 $\ideal S_i$ 的真因子，則由於 $\ideal S=[\ideal S_i,\widehat{\ideal S}_i]$ 是既約表示，$[\ideal T_0,\widehat{\ideal S}_i]$ 也是 $\ideal S$ 的真因子，這同假設矛盾。

% unit: noether.p19.par.113
4. 由 $\ideal T_0\ideal S_i\equiv0(\ideal R)$，其中 $\ideal T_0$ 是 $\ideal R$ 的真因子，也推出 $\ideal T_0\cdot\ideal S\equiv0(\ideal R)$；因此 $\ideal T_0$ 將是 $\ideal R$ 的真因子，同假設矛盾。

% unit: noether.p19.par.114
定義 V 特別表明，每個理想都對由 $\Sigma$ 的所有元素組成的單位理想 $\ideal O$ 是相對素的；\footnote{$\ideal O$ 只在關於整除性和最小公倍數的意義上起單位的作用，而不在乘積形成的意義上起單位的作用。例如，對於所有沒有常數項的 $x$ 的整係數多項式所成的領域，有 $\ideal O=(x)$；對於所有偶數所成的領域，有 $\ideal O=(2)$。} 但是下面的定理只適用於不同於 $\ideal O$ 的理想。

% unit: noether.p19.par.115
由剛剛證明的定理 X 首先得到：

% unit: noether.p19.lem.005
\textbf{引理 V。} \emph{一個理想作為互相相對素且不同於 $\ideal O$ 的理想的最小公倍數的每個表示，都是既約的。}

% unit: noether.p19.par.116
設
% unit: noether.p19.eq.079
\[
 \ideal M=[\ideal R_1\ideal R_2\cdots\ideal R_\sigma]=[\ideal R_i,\ideal L_i]
\]
是這樣的一個表示。按定理 X，2，$\ideal L_i$ 也對 $\ideal R_i$ 是相對素的；因此 $\ideal R_i$ 不能包含在 $\ideal L_i$ 中，表示就是最短的。因為如果 $\ideal L_i\equiv0(\ideal R_i)$，且 $\ideal L_i$ 對 $\ideal R_i$ 相對素，則由於 $\ideal O\cdot\ideal L_i\equiv0(\ideal R_i)$，也會推出 $\ideal O\equiv0(\ideal R_i)$；於是 $\ideal R_i=\ideal O$，這按假設被排除。

% unit: noether.p19.par.117
現在設 $\ideal R_i$ 可由真因子 $\ideal R_i^*$ 替換；則按引理 II，$\ideal R_i$ 可約，並且得到
% unit: noether.p19.eq.080
\[
 \ideal R_i=[\ideal R_i^*,(\ideal R_i,\ideal L_i)].
\]
必要時把 $(\ideal R_i,\ideal L_i)$ 替換為一個真因子 $(\ideal R_i,\ideal L_i)^*$，同樣把 $\ideal R_i^*$ 替換為一個真因子，使得產生一個 $\ideal R_i$ 的既約表示，則定理 X，3 表明 $\ideal L_i$ 對 $(\ideal R_i,\ideal L_i)^*$ 也是相對素的；而由於表示最短，後者不同於 $\ideal O$。但是 $(\ideal R_i,\ideal L_i)^*$ 包含在 $(\ideal R_i,\ideal L_i)$ 中，且 $(\ideal R_i,\ideal L_i)$ 包含在 $\ideal L_i$ 中，由此得出矛盾。

% unit: noether.p19.par.118
由定理 X 還得到把這一點同相應素理想聯繫起來的定理。

% unit: noether.p19.thm.011
\textbf{定理 XI。} \emph{若 $\ideal R$ 對 $\ideal S$ 是相對素的，並且 $\ideal S$ 不同於 $\ideal O$，則 $\ideal R$ 的任何相應素理想\footnote{$\ideal R$ 的相應素理想的定義見 \S~5 末尾。}都不能被 $\ideal S$ 的一個相應素理想整除。反過來，若不存在這樣的整除關係，則 $\ideal R$ 對 $\ideal S$ 是相對素的，當然 $\ideal S$ 也不同於 $\ideal O$。}

% unit: noether.p19.par.119
設
% unit: noether.p19.eq.081
\[
 \ideal R=[\ideal Q_1\cdots\ideal Q_\alpha],\qquad
 \ideal S=[\ideal Q_1^*\cdots\ideal Q_\beta^*]
\]
是 $\ideal R$ 和 $\ideal S$ 由最大準素理想給出的既約表示；$\ideal P_1,\ldots,\ideal P_\alpha$ 以及 $\ideal P_1^*,\ldots,\ideal P_\beta^*$ 為相應素理想。我們用下列形式證明斷言：若某個 $\ideal P$ 可被某個 $\ideal P^*$ 整除，則 $\ideal R$ 不可能對 $\ideal S$ 相對素，反過來也如此。

% unit: noether.p19.par.120
於是設 $\ideal P_\mu\equiv0(\ideal P_\nu^*)$，因而也有 $\ideal Q_\mu\equiv0(\ideal P_\nu^*)$；由 $\ideal P_\nu^*$ 的定義推出 $\ideal Q_\mu^{\sigma_\nu}\equiv0(\ideal Q_\nu^*)$，所以也有 $\ideal R^{\sigma_\nu}\equiv0(\ideal Q_\nu^*)$。現在令 $\ideal R^\tau$ 為可被 $\ideal Q_\nu^*$ 整除的 $\ideal R$ 的最低次冪。若 $\tau=1$，則 $\ideal R$ 可被 $\ideal Q_\nu^*$ 整除；但由於按假設 $\ideal S$ 不同於 $\ideal O$，所以 $\ideal R$ 對 $\ideal Q_\nu^*$ 不是相對素的。若 $\tau\ge2$，則 $\ideal R^{\tau-1}\cdot\ideal R\equiv0(\ideal Q_\nu^*)$，而 $\ideal R^{\tau-1}\not\equiv0(\ideal Q_\nu^*)$，所以 $\ideal R$ 對 $\ideal Q_\nu^*$ 不是相對素的；\footnote{$\ideal R^0$ 未定義，因為 $\Sigma$ 不必含有單位；所以必須預先處理 $\tau=1$ 的情形。即使在 $\Sigma$ 有單位的情形中，$\tau=0$ 也被關於 $\ideal S$ 的假設排除。} 因此在兩種情形中，按定理 X，3，$\ideal R$ 對 $\ideal S$ 也不是相對素的。

% unit: noether.p19.par.121
反過來，若 $\ideal R$ 對 $\ideal S$ 不是相對素的，則按定理 X，1，$\ideal R$ 至少對一個 $\ideal Q_\nu^*$ 不是相對素的。因此有
% unit: noether.p19.eq.082
\[
 \ideal T_0\cdot\ideal R\equiv0(\ideal Q_\nu^*)
 \ideal T_0\not\equiv0(\ideal Q_\nu^*),
\]
並且由於 $\ideal Q_\nu^*$ 是準素的，得到
% unit: noether.p19.eq.083
\[
 \ideal R^\tau\equiv0(\ideal Q_\nu^*),\qquad
 \text{所以也有}\quad
 \ideal Q_1^{\tau}\cdots\ideal Q_\alpha^{\tau}\equiv0(\ideal Q_\nu^*).
\]
於是對相應素理想推出
% unit: noether.p19.eq.084
\[
 \ideal P_1^{\tau\varrho_1}\cdots\ideal P_\alpha^{\tau\varrho_\alpha}\equiv0(\ideal P_\nu^*),
\]
並且按素理想的性質，$\ideal P_\nu^*$ 因而包含在至少一個 $\ideal P$ 中，定理 XI 得證。

% unit: noether.p19.par.122
由定理 X 和 XI，現在如下得到分解為相對素-不可約理想的存在性\footnote{這種分解的存在性也可以直接證明，完全類似於 \S~2 中所證明的分解為有限多個不可約理想的存在性。}和唯一性。

% unit: noether.p19.par.123
設 $\ideal M=[\ideal Q_1\cdots\ideal Q_\alpha]$ 是 $\ideal M$ 由最大準素理想給出的一個既約（或至少最短）表示；$\ideal P_1,\ldots,\ideal P_\alpha$ 為相應素理想。我們把這些 $\ideal P$ 分成若干組，使得某一組中的任何理想都不能被另一不同組中的理想整除，而每一個單獨的組又不能分裂成兩個都具有這個性質的子組。

% unit: noether.p19.par.124
為了構造這樣的分組，要注意按定義，每個組 $G$ 連同任一理想 $\ideal P$ 也必須包含它在 $\ideal P_1,\ldots,\ideal P_\alpha$ 中出現的所有因子和倍數（即可被 $\ideal P$ 整除的那些）。例如，令 $\ideal P_j^{(i)}$ 為 $\ideal P$ 的所有倍數，$\ideal P_{j_1}^{(i_1)}$ 為 $\ideal P_j^{(i)}$ 的所有因子，$\ideal P_{j_1}^{(i_1i_2)}$ 為 $\ideal P_{j_1}^{(i_1)}$ 的所有倍數，依此類推；一般地，$\ideal P_{j_1\ldots j_{\lambda-1}}^{(i_1\ldots i_\lambda)}$ 表示 $\ideal P_{j_1\ldots j_{\lambda-1}}^{(i_1\ldots i_{\lambda-1})}$ 的所有倍數，而 $\ideal P_{j_1\ldots j_\lambda}^{(i_1\ldots i_\lambda)}$ 表示 $\ideal P_{j_1\ldots j_{\lambda-1}}^{(i_1\ldots i_\lambda)}$ 的所有因子。由於總共只涉及有限多個理想 $\ideal P$，這個過程必然在有限步之後中止，也就是說，不再產生與所有先前理想都不同的理想。這樣得到的理想 $\ideal P$ 的全體，事實上構成了一個具有所需性質的組 $G$。

% unit: noether.p19.par.125
因為按定義，$G$ 連同每個 $\ideal P_{j_1\ldots j_{\lambda-1}}^{(i_1\ldots i_{\lambda-1})}$ 也包含所有倍數 $\ideal P_{j_1\ldots j_{\lambda-1}}^{(i_1\ldots i_\lambda)}$；連同每個 $\ideal P_{j_1\ldots j_{\lambda-1}}^{(i_1\ldots i_\lambda)}$ 也包含所有因子 $\ideal P_{j_1\ldots j_\lambda}^{(i_1\ldots i_\lambda)}$。但是其中也包含 $\ideal P_{j_1\ldots j_{\lambda-1}}^{(i_1\ldots i_{\lambda-1})}$ 的所有因子，而 $\ideal P_{j_1\ldots j_{\lambda-1}}^{(i_1\ldots i_\lambda)}$ 的倍數本身又仍然是 $\ideal P_{j_1\ldots j_{\lambda-1}}^{(i_1\ldots i_\lambda)}$。因此，不在 $G$ 中的理想既不能是 $G$ 中理想的因子，也不能是它們的倍數。

% unit: noether.p19.par.126
但 $G$ 也滿足不可約性條件。若存在分成兩個子組 $G^{(1)}$ 和 $G^{(2)}$ 的分法，並且例如 $G^{(1)}$ 含有 $\ideal P_{j_1\ldots j_\lambda}^{(i_1\ldots i_\lambda)}$（相應地，含有 $\ideal P_{j_1\ldots j_{\lambda-1}}^{(i_1\ldots i_\lambda)}$），則它也含有所有先前的理想，因為這些理想交替地是倍數和因子（相應地，是因子和倍數）；所以它也含有 $\ideal P$，從而含有整個組 $G$。對不在 $G$ 中的理想作相同處理，就得到所有 $\ideal P$ 的一個分組 $G_1,\ldots,G_\sigma$，並且該分組具有所需性質。這樣的分組是唯一的；因為若 $G'_1,\ldots,G'_\tau$ 是第二個分組，而 $\ideal P_{j_1\ldots j_\lambda}^{(i_1\ldots i_\lambda)}$（相應地，$\ideal P_{j_1\ldots j_{\lambda-1}}^{(i_1\ldots i_\lambda)}$）是 $G'_i$ 的一個元素，則按上面的論證，$G'_i$ 含有整個組 $G$，並且由於不可約性性質，它同 $G$ 相同。

% unit: noether.p19.par.127
現在令 $\ideal P_{i\mu}$（其中 $\mu$ 從 $1$ 到 $\lambda_i$）表示合併在一個組 $G_i$ 中的那些理想 $\ideal P$；由於所有 $\ideal P$ 都彼此不同，對應於一個最短表示的準素理想 $\ideal Q_{i\mu}$ 就唯一確定。置
% unit: noether.p19.eq.085
\[
 \ideal R_i=[\ideal Q_{i1}\cdots\ideal Q_{i\lambda_i}],
 \qquad
 \text{於是有}\quad
 \ideal M=[\ideal R_1\cdots\ideal R_\sigma].
\]
我們證明，由此已經得到 $\ideal M$ 到相對素-不可約理想的分解。首先要注意，即使 $\ideal R_i=[\ideal Q_{i_1}\cdots\ideal Q_{i\lambda_i}]$ 只是最短表示，定理 XI 仍可應用，因為按定理 IX 的附加說明，相應素理想已由此唯一確定。現在，$\ideal R_i$ 的任何相應素理想 $\ideal P_{i\mu}$ 都不能被 $\ideal R_j$ 的相應素理想 $\ideal P_{j\nu}$ 整除，反過來也不能；因此按定理 XI，$\ideal R_i$ 和 $\ideal R_j$ 互相相對素。並且按定理 XI，由於組 $G_i$ 的不可約性性質，每個單獨的 $\ideal R_i$ 都是相對素-不可約的，因為在可約情形中所涉及的總是不同於 $\ideal O$ 的理想。再按引理 V，這還是一個既約表示，即使最初只是從由 $\ideal Q$ 給出的最短表示出發。反過來，按定理 XI，每個分解為相對素-不可約理想的分解都導向上面給出的 $\ideal P_{i\mu}$ 的分組。

% unit: noether.p19.par.128
現在設
% unit: noether.p19.eq.086
\[
 \ideal M=[\bar{\ideal R}_1\cdots\bar{\ideal R}_\tau]
\]
是 $\ideal M$ 由相對素-不可約理想給出的第二個表示；它按引理 V 是既約的。於是，正如把這些 $\bar{\ideal R}$ 分解為最大準素理想所表明的那樣，相應素理想一致；因此，這些素理想的分組也一致，而其唯一性已經在上面證明。於是 $\tau=\sigma$；並且可以選擇記號，使得 $\ideal R_i$ 和 $\bar{\ideal R}_i$ 屬於同一個組。若
% unit: noether.p19.eq.087
\[
 \ideal M=[\ideal R_i,\ideal L_i]=[\bar{\ideal R}_i,\bar{\ideal L}_i]
\]
是由理想和補給出的表示，則由於 $\bar{\ideal R}_i$ 被歸入同 $\ideal R_i$ 相同的組，按定理 XI，$\ideal L_i$ 對 $\bar{\ideal R}_i$ 也是相對素的，而 $\bar{\ideal L}_i$ 對 $\ideal R_i$ 是相對素的。由於
% unit: noether.p19.eq.088
\[
 \ideal R_i\cdot\ideal L_i\equiv0(\bar{\ideal R}_i),
 \qquad
 \bar{\ideal R}_i\cdot\bar{\ideal L}_i\equiv0(\ideal R_i)
\]
因此得到
% unit: noether.p19.eq.089
\[
 \ideal R_i\equiv0(\bar{\ideal R}_i),
 \qquad
 \bar{\ideal R}_i\equiv0(\ideal R_i),
 \qquad
 \ideal R_i=\bar{\ideal R}_i.
\]
這樣證明了：

% unit: noether.p19.thm.012
\textbf{定理 XII。} \emph{每個理想都能唯一地表示為有限多個互相相對素且相對素-不可約的理想的最小公倍數。}
%
\clearpage
%

%
%
%
%

\setcounter{footnote}{28}

\subsection*{\S~7. 孤立理想的唯一性。}

\textbf{定義 VI。} \emph{若最短表示 $\ideal M=[\ideal R,\ideal L]$ 關於 $\ideal L$ 是約化的，則當 $\ideal R$ 的任何相應素理想都不包含於 $\ideal L$ 的相應素理想中時，稱 $\ideal R$ 為孤立理想；換言之，當 $\ideal L$ 對 $\ideal R$ 是相對素的時，稱 $\ideal R$ 為孤立理想。}

於是表示 $\ideal M=[\ideal R,\ideal L]$ 滿足引理 V 中表示 $\ideal M=[\ideal R_i,\ideal L_i]$ 的條件，而引理 V 證明：

\textbf{引理 VI。} \emph{若 $\ideal R$ 是最短表示 $\ideal M=[\ideal R,\ideal L]$ 的一個孤立理想，且該表示關於 $\ideal L$ 是約化的，則該表示關於 $\ideal R$ 也是約化的。}

因此，對孤立理想而言，總是處在約化表示中。在把 $\ideal R$ 和 $\ideal L$ 分解為不可約理想時出現的相應素理想，互相補足為在 $\ideal M$ 的相應分解中出現的、唯一確定的相應素理想。於是，屬於 $\ideal R$ 的不可約理想分解的任何相應素理想，都不包含於 $\ideal M$ 的其餘相應素理想中。反過來，若這一條件成立，並且 $\ideal R$ 至少出現在一個關於 $\ideal R$ 約化的表示 $\ideal M=[\ideal R,\ideal L]$ 中，因而也出現在一個約化表示 $\ideal M=[\ideal R,\ideal L^*]$ 中，則按定義 VI，$\ideal R$ 是孤立的。由此得到下列不依賴於特殊補 $\ideal L$ 的定義。

\textbf{定義 VIa。} \emph{若屬於 $\ideal R$ 的不可約理想分解的素理想不包含於屬於 $\ideal M$ 的相應分解的其餘相應素理想中，並且 $\ideal R$ 至少出現在一個關於 $\ideal R$ 約化的表示 $\ideal M=[\ideal R,\ideal L]$ 中，則稱 $\ideal R$ 為孤立理想。}\footnote{如果引入通常的相應素理想（\S~5 末尾），則還必須加入一個特殊條件，即 $\ideal L$ 的相應素理想全都不同於 $\ideal R$ 的相應素理想。按照定義 VIa，就不再需要預先假定表示關於補是約化的。}

設
\[
 \ideal M=[\ideal R,\ideal L]=[\bar{\ideal R},\bar{\ideal L}]
\]
是 $\ideal M$ 由孤立理想 $\ideal R$ 和 $\bar{\ideal R}$ 以及補 $\ideal L$ 和 $\bar{\ideal L}$ 給出的兩個表示，並且 $\ideal R$ 和 $\bar{\ideal R}$ 的相應素理想一致。把 $\ideal L$ 和 $\bar{\ideal L}$ 分別替換為除子 $\ideal L^*$ 和 $\bar{\ideal L}^{*}$，使得表示成為約化的；則 $\ideal L^*$ 和 $\bar{\ideal L}^{*}$ 的相應素理想也一致。由定理 XI，$\ideal L^*$ 對 $\bar{\ideal R}$ 是相對素的，而 $\bar{\ideal L}^{*}$ 對 $\ideal R$ 是相對素的。由於
\[
 \ideal R\cdot\ideal L^*\equiv0(\bar{\ideal R}),
 \qquad
 \bar{\ideal R}\bar{\ideal L}^{*}\equiv0(\ideal R),
\]
於是得到
\[
 \ideal R\equiv0(\bar{\ideal R}),
 \qquad
 \bar{\ideal R}\equiv0(\ideal R),
 \qquad
 \ideal R=\bar{\ideal R}.
\]
因此，孤立理想由其相應素理想唯一確定。這特別加強了關於分解為不可約理想、相應地分解為最大準素理想的定理 VII 和 IX；按照定理 IX 的附加說明，這裡只需假定最短表示。綜上，得到：

\textbf{定理 XIII。} \emph{在一個理想作為不可約理想、相應地作為最大準素理想的最小公倍數的每個最短表示中，孤立的不可約理想、相應地孤立的最大準素理想，是唯一確定的；不唯一性只涉及非孤立的不可約理想、相應地非孤立的最大準素理想。}\footnote{對於多項式域中的理想，這一定理在分解為最大準素理想的情形下已由 Macaulay 不加證明地給出；他關於孤立和非孤立（imbedded）準素理想的定義，可以看作下面定義的一種非有理化表述。} \emph{一般地說，孤立理想由其相應素理想唯一確定。}

在這樣一個由不可約理想、相應地由最大準素理想給出的最短表示中，若理想 $\ideal B_i$、相應地 $\ideal Q_j$ 是非孤立的，則按定義，它們的補 $\ideal A_i$、相應地 $\ideal L_j$ 可被 $\ideal P_i$、相應地 $\ideal P_j$ 整除。因此
\[
 \ideal A_i^{\varrho_i}\equiv0(\ideal B_i)
 \qquad \text{相應地} \qquad
 \ideal L_j^{\sigma_j}\equiv0(\ideal Q_j).
\]
反過來，若這些關係成立，則 $\ideal P_i$、相應地 $\ideal P_j$ 包含在補的至少一個相應素理想中，因而按照定理 IX 的附加說明，也包含在 $\ideal L_j$ 的那個除子 $\ideal L_j^*$ 的一個相應素理想中，後者給出關於 $\ideal L_j^*$ 約化的表示。因此 $\ideal B_i$、相應地 $\ideal Q_j$ 是非孤立的。特別地，若某個不可約理想 $\ideal B_i$ 的相應素理想在 $\ideal M$ 的分解中出現不止一次，則 $\ideal B_i$ 總是非孤立的。\emph{所以，非孤立準素理想還可由如下性質刻畫：每個補的某個冪會被它整除；孤立準素理想則由這種情形不可能發生來刻畫。}

\subsection*{\S~8. 把一個理想唯一表示為互素-不可約理想的乘積。}

如果作為基礎的環領域 $\Sigma$ 有單位元，也就是說有一個元素 $\varepsilon$，使得對 $\Sigma$ 的每個元素 $a$ 都有 $\varepsilon a=a$，\footnote{由於乘法是交換的，$\Sigma$ 當然至多隻有一個單位元：若 $\varepsilon_1$ 和 $\varepsilon_2$ 是兩個單位元，則 $\varepsilon_1\varepsilon_2=\varepsilon_2=\varepsilon_1$。} 那麼互素理想可以定義如下。

\textbf{定義 VIII。} \emph{若兩個理想 $\ideal R$ 和 $\ideal S$ 的最大公因子是由 $\Sigma$ 的所有元素組成的單位理想 $\ideal O=(\varepsilon)$，則稱它們為互素的。若一個理想不能表示為兩兩互素理想的最小公倍數，則稱它為互素-不可約的。}

注意，兩個互素理想總是互相相對素的。按定義，存在元素 $r\equiv0(\ideal R)$ 和 $s\equiv0(\ideal S)$ 使得 $\varepsilon=r+s$。由 $\ideal T\cdot\ideal R\equiv0(\ideal S)$ 得到 $\ideal T\cdot r\equiv0(\ideal S)$，因而 $\ideal T\cdot\varepsilon=\ideal T\equiv0(\ideal S)$；相應地，由 $\bar{\ideal T}\cdot\ideal S\equiv0(\ideal R)$ 得到 $\bar{\ideal T}\equiv0(\ideal R)$。\footnote{反命題不成立：例如 $\ideal R=(x)$ 和 $\ideal S=(y)$ 互相相對素，但並不互素。} 因此引理 V 表明，每個由兩兩互素理想給出的表示也都是約化的。

唯一性證明依賴於定理 X 的類似物：

\textbf{定理 XIV。} \emph{若 $\ideal R$ 與理想 $\ideal S_1,\ldots,\ideal S_\lambda$ 中每一個都互素，則 $\ideal R$ 也與 $\ideal S=[\ideal S_1\cdots\ideal S_\lambda]$ 互素。反過來，若 $\ideal R$ 和 $\ideal S$ 互素，則 $\ideal R$ 與每個 $\ideal S_j$ 都互素。若 $\ideal R=[\ideal R_1\cdots\ideal R_\mu]$，並且每個 $\ideal R_i$ 都同每個 $\ideal S_j$ 互素，則 $\ideal R$ 和 $\ideal S$ 互素；這裡反命題也成立。}

事實上，若 $\ideal R$ 與每個 $\ideal S_j$ 互素，則存在元素 $s_j$，使得
\[
 s_j\equiv0(\ideal S_j),
 \qquad
 s_j\equiv\varepsilon(\ideal R).
\]
因此
\[
 s_1\cdot s_2\cdots s_\lambda\equiv0(\ideal S),
 \qquad
 s_1\cdot s_2\cdots s_\lambda\equiv\varepsilon(\ideal R),
 \qquad
 (\ideal R,\ideal S)=(\varepsilon).
\]
由於 $(\ideal R,\ideal S)$ 可被每個 $(\ideal R,\ideal S_j)$ 整除，反命題隨之成立。同一論證反覆應用，即得第二部分。因為若在固定的 $i$ 下 $\ideal R_i$ 與 $\ideal S_1,\ldots,\ideal S_\lambda$ 互素，則 $\ideal R_i$ 與 $\ideal S$ 互素；若這對每個 $i$ 都成立，則由於互素性是相互的，$\ideal S$ 與 $\ideal R$ 互素。反過來，$\ideal R$ 與 $\ideal S$ 的互素性給出 $\ideal S$ 與 $\ideal R_i$ 的互素性，因而也給出 $\ideal R_i$ 與 $\ideal S$ 的互素性。

分解為互素-不可約理想的存在性和唯一性的證明，\footnote{分解的存在性仍可直接仿照 \S~2 來證明；唯一性證明也可以直接進行（參見導言中關於 Schmeidler 以及 Noether--Schmeidler 的說明）。這裡給出的證明同時說明了互素-不可約理想的結構。} 與相對素理想的相應證明一樣，歸結為一個唯一的分組。不過，由定理 XII，$\ideal M$ 的相對素-不可約理想 $\ideal R_1,\ldots,\ideal R_\sigma$ 已經唯一確定，所以這裡不需要再回到相應素理想。

我們把 $\ideal M$ 的唯一確定的相對素-不可約理想 $\ideal R_1,\ldots,\ideal R_\sigma$ 分成若干組，使得\emph{一組中的每個理想都同任何另一組中的每個理想互素，而單獨一組不能再分成兩個子組，使得一個子組的每個理想都同另一個子組的每個理想互素}。這樣的分組如下得到。按定義，一個組 $G$ 必須連同每個理想 $\ideal R$ 一起包含所有不與 $\ideal R$ 互素的理想，記這些為 $\ideal R_{i_1}$；再令 $\ideal R_{i_1i_2}$ 為所有不與這些理想互素的理想；一般地，令 $\ideal R_{i_1\ldots i_\lambda}$ 表示不與 $\ideal R_{i_1\ldots i_{\lambda-1}}$ 互素的理想。由於總共只出現有限多個理想，這個過程經過有限步後終止。這樣得到的理想全體形成一個具有所需性質的組 $G$：由構造可知，$G$ 外的一切理想都同 $G$ 中的一切理想互素。若 $G$ 被分成兩個子組 $G^{(1)}$ 和 $G^{(2)}$，並且 $\ideal R_{i_1\ldots i_\lambda}$ 屬於 $G^{(1)}$，則由於互素性是相互的，$G^{(1)}$ 也必須包含 $\ideal R_{i_1\ldots i_{\lambda-1}},\ldots,\ideal R_{i_1}$，因而也包含 $\ideal R$，從而包含整個組 $G$。這就證明了不可約性。對其餘理想重複這一構造，得到所有 $\ideal R$ 的分組 $G_1,\ldots,G_\tau$。

該分組是唯一的：若 $G'_1,\ldots,G'_\tau$ 是第二個分組，並且 $\ideal R_{i_1\ldots i_\lambda}$ 是 $G'_i$ 的一個元素，則 $G'_i$ 含有 $\ideal R$，從而含有 $G_1$；又由不可約性條件，它不能含有不同於 $G_1$ 的元素。因此 $G'_i=G_1$。

令 $\ideal T_i$ 為合併在組 $G_i$ 中的那些理想 $\ideal R$ 的最小公倍數。我們證明
\[
 \ideal M=[\ideal T_1\cdots\ideal T_\tau]
\]
是 $\ideal M$ 由互素-不可約理想給出的一個表示。先把各 $\ideal T$ 分解為 $\ideal R$，即可看出 $[\ideal T_1\cdots\ideal T_\tau]$ 的確表示 $\ideal M$。由定理 XIV，各 $\ideal T$ 兩兩互素，並且每個 $\ideal T$ 都是互素-不可約的。這證明了這種表示的存在性。

為證唯一性，設 $\ideal M=[\bar{\ideal T}_1\cdots\bar{\ideal T}_{\bar\tau}]$ 是同類的第二個表示。把各 $\bar{\ideal T}$ 分解為相對素-不可約理想 $\ideal R$，定理 XIV 表明，出現在不同 $\bar{\ideal T}$ 中的理想 $\ideal R$ 互相互素，因而也互相相對素；所以它們同 $\ideal M$ 的唯一確定的相對素-不可約理想 $\ideal R$ 一致。各 $\bar{\ideal T}_i$ 還按照定理 XIV 誘導出這些 $\ideal R$ 的一個具有上述性質的分組 $G'_i$。由於該分組唯一，並且每個 $\bar{\ideal T}_i$ 又由組 $G'_i=G_i$ 唯一確定，故有 $\bar{\ideal T}_i=\ideal T_i$；唯一性得證。

對於兩兩互素的理想，最小公倍數還等於乘積。因為由定理 XIV，若 $\ideal M=[\ideal T_1\cdots\ideal T_\tau]$，則補 $\ideal L_i$ 也同 $\ideal T_i$ 互素；於是存在元素
\[
 t_i\equiv0(\ideal T_i),
 \qquad
 l_i\equiv0(\ideal L_i),
 \qquad
 \varepsilon=t_i+l_i.
\]
由 $f\equiv0(\ideal T_i)$ 和 $f\equiv0(\ideal L_i)$ 得到
\[
 f=f\varepsilon=ft_i+fl_i,
 \qquad
 f\equiv0(\ideal L_i\cdot\ideal T_i).
\]
反過來，$\ideal L_i\cdot\ideal T_i$ 可被 $[\ideal L_i,\ideal T_i]$ 整除，於是 $[\ideal L_i,\ideal T_i]=\ideal L_i\cdot\ideal T_i$；繼續對 $\ideal L_i$ 作同樣處理，得到
\[
 \ideal M=[\ideal T_1\cdots\ideal T_\tau]
       =\ideal T_1\cdot\ideal T_2\cdots\ideal T_\tau.
\]
因此：

\textbf{定理 XV。} \emph{每個理想都能唯一地表示為有限多個兩兩互素且互素-不可約的理想的乘積。}

\setcounter{footnote}{34}
\subsection*{\S~9. 把研究推廣到模。分解為不可約模時分量數目的相等。}

現在我們說明，前三節中關於不可約理想而非準素理想和素理想的內容，在較弱的假設下仍然有效。那些部分沒有使用乘法交換律，而只涉及理想作為模的性質。因此它們對於相對於非交換領域的模仍然有效；這類模現在需要定義。模的定義以一個具有下列性質的雙重領域 $(\Sigma,T)$ 為基礎。

$\Sigma$ 是一個抽象定義的非交換環，也就是說，$\Sigma$ 是由元素 $a,b,c,\ldots$ 組成的系統，對它定義了兩種合成：環加法 $(+)$ 和環乘法 $(\times)$；它們滿足 \S~1 中列出的法則，除了環乘法的交換律 4。

$T$ 是由元素 $\alpha,\beta,\gamma,\ldots$ 組成的系統；它同 $\Sigma$ 相關聯，也定義了兩種合成：加法，它由任意兩個元素 $\alpha,\beta$ 產生唯一的第三個元素 $\alpha+\beta$；以及用 $\Sigma$ 的元素 $c$ 乘 $T$ 的元素 $\alpha$，唯一地產生 $T$ 的元素 $c\alpha$。\footnote{這裡處理的是一種“右側”乘法、一個“右側”領域 $T$，因而也處理“右側”模和理想。若改以左側乘法 $\alpha c$ 為基礎，則得到相應的左側模和理想理論；$M$ 將連同 $\alpha$ 一起包含 $\alpha\cdot c$。結合律則取形式 $(\gamma\cdot b)\cdot a=\gamma\cdot(b\times a)$。}

這些合成滿足下列法則：
\begin{enumerate}
\item 加法結合律：$(\alpha+\beta)+\gamma=\alpha+(\beta+\gamma)$。
\item 加法交換律：$\alpha+\beta=\beta+\alpha$。
\item 無限制且唯一的減法法則：在 $T$ 中有且只有一個元素 $\xi$ 滿足方程 $\alpha+\xi=\beta$（記作 $\xi=\beta-\alpha$）。
\item 乘法結合律：$a\cdot(b\cdot\gamma)=(a\times b)\cdot\gamma$。
\item 分配律：
\[
 (a+b)\cdot\gamma=a\cdot\gamma+b\cdot\gamma,
 \qquad
 c\cdot(\alpha+\beta)=c\cdot\alpha+c\cdot\beta.
\]
\end{enumerate}
由這些條件，如所周知，可推出零元的存在，以及減法同乘法結合時分配律也成立：
\[
 (a-b)\cdot\gamma=a\cdot\gamma-b\cdot\gamma,
 \qquad
 c\cdot(\alpha-\beta)=c\cdot\alpha-c\cdot\beta,
\]
其中 $(\dot{-})$ 表示 $\Sigma$ 中的減法。若 $\Sigma$ 含有單位元 $\varepsilon$，則還要求對 $T$ 的每個元素 $\alpha$ 都有 $\varepsilon\alpha=\alpha$。

所謂 $(\Sigma,T)$ 中的模 $M$，指 $T$ 的元素系統，滿足兩個條件：
\begin{enumerate}
\item 連同 $\alpha$，$M$ 也含有 $c\alpha$，其中 $c$ 是 $\Sigma$ 的任意元素；
\item 連同 $\alpha$ 和 $\beta$，$M$ 也含有差 $\alpha-\beta$，因而也含有每個整數 $n$ 的 $n\alpha$。\footnote{整數在這裡仍理解為縮寫符號，而不是環元素。}
\end{enumerate}
按此定義，$T$ 本身構成 $(\Sigma,T)$ 中的一個模。若特殊地，領域 $T$ 及其上所規定的合成同領域 $\Sigma$ 及其上有效的合成一致，則模 $M$ 成為 $\Sigma$ 中的右理想 $M$。若再假設 $\Sigma$ 交換，則得到通常的理想概念；因此它是模概念的一個特殊情形。\footnote{最簡單的模例子是整線性形式的模：這裡 $\Sigma$ 由所有有理整數構成，而 $T$ 由所有整線性形式構成。若在 $\Sigma$ 和 $T$ 中取代數整數而非有理整數，或例如取所有偶數，則得到稍一般的模。若不考慮線性形式，而把每個係數組合看作一個元素，則 $\Sigma$ 和 $T$ 中的合成實際上不同。非交換多項式環領域中的理想，是 Noether--Schmeidler 合作論文的主題。其他特殊非交換領域中的理想，在 Hurwitz 關於四元數數論的講義（Berlin, Springer 1919）以及那裡引用的 Du Pasquier 的著作中處理。}

對於模，\S~1 的所有定義仍然有效。因此 $\alpha\equiv0(M)$，相應地 $N\equiv0(M)$，表示 $\alpha$，相應地 $N$ 的每個元素，都是 $M$ 的元素；換言之，$\alpha$，相應地 $N$，可被 $M$ 整除。若 $M$ 含有不同於 $N$ 的元素，則 $M$ 是 $N$ 的真因子；由 $N\equiv0(M)$ 和 $M\equiv0(N)$ 推出 $M=N$。最大公因子和最小公倍數的定義同樣逐字保持。若模 $M$ 特別含有有限多個元素 $\alpha_1,\ldots,\alpha_\rho$，使得
\[
 M=(\alpha_1,\ldots,\alpha_\rho),
\]
也就是說，對每個 $\alpha\equiv0(M)$ 有
\[
 \alpha=c_1\alpha_1+\cdots+c_\rho\alpha_\rho+n_1\alpha_1+\cdots+n_\rho\alpha_\rho,
\]
其中 $c_i$ 是 $\Sigma$ 的元素，而 $n_i$ 是整數，則稱 $M$ 為有限模，並稱 $\alpha_1\ldots\alpha_\varrho$ 為模基。

以下我們同 \S~1 一樣，只假定領域 $(\Sigma,T)$ 滿足有限性條件：$(\Sigma,T)$ 中的每個模都是有限的，因而有一個模基。

對於這樣的領域 $(\Sigma,T)$，定理 I 關於有限鏈的結論也對模成立，正如那裡給出的證明所表明的；於是 \S\S~2 和 3 的所有前提都滿足。定義 I 以及關於最短表示和約化表示的引理 I 可直接轉移；關於每個模可表示為有限多個不可約模的最小公倍數的定理 II 也同樣轉移，其中引理 II 表明，每個這樣的最短表示同時也是約化的。定理 III 也仍然有效，它通過補的性質表達一個模的可約性；由此推出引理 III，最後推出定理 IV，即兩個不同的、把一個模表示為不可約模最小公倍數的最短表示，其分量數目相等。定理 IV 還給出引理 IV，作為引理 I 中關於約化表示命題的逆命題。

\emph{附加說明。} 同一論證表明，如果對於雙側領域 $T$，也就是同時為右側和左側的領域，在各處都把模理解為雙側模，即連同 $\alpha$ 也含有 $c\cdot\alpha$ 和 $\alpha\cdot c$，並且連同 $\alpha$ 和 $\beta$ 也含有 $\alpha-\beta$ 的模，那麼所有這些定理和定義仍然有效。

這些定理全都只依賴於整除性和最小公倍數概念，而進一步的唯一性定理則本質上依賴於乘積概念，因此不能直接轉移。於是，準素理想和素理想的定義不能轉移到模，因為 $T$ 的兩個元素的乘積沒有定義。即使形式上可轉移到非交換環，這一轉移也會失去意義，因為在這裡相應素理想的存在不能證明；\footnote{因為由 $p_1^\rho\equiv0(\ideal Q)$ 和 $p_2^\sigma\equiv0(\ideal Q)$，這裡不能推出 $(p_1-p_2)^{\rho+\sigma}\equiv0(\ideal Q)$；同樣，由 $(ab)^\rho\equiv0(\ideal Q)$ 也不能推出 $a^\rho b^\rho\equiv0(\ideal Q)$。因此，$\ideal P$ 既不能被證明為一個理想，也不能被證明具有素理想的性質。對於雙側理想 $\ideal P$，確實可以證明它是一個理想，但不能證明它是素理想。} 證明一個不可約理想是準素理想也同樣失敗。下面的唯一性定理正以這兩個事實為基礎。相反，若非交換環有單位元，則可以定義互素和互素-不可約理想；而在定理 II 的證明中所用的論證表明，每個理想都可以表示為有限多個兩兩互素且互素-不可約理想的最小公倍數。\footnote{對於特殊的“完全可約”理想，這裡也可以建立唯一性定理；參見 Noether--Schmeidler 合作論文。}

最後，我們給出一個足以保證 $(\Sigma,T)$ 中有限性條件成立的判據：若 $\Sigma$ 含有單位元並滿足有限性條件，並且 $T$ 本身是 $(\Sigma,T)$ 中的一個有限模，則 $(\Sigma,T)$ 中的每個模都是有限的。

事實上，由 $\Sigma$ 中單位元的存在，對於
\[
 \ideal M=(f_1,\ldots,f_e),
 \qquad
 f=\bar b_1f_1+\cdots+\bar b_\varrho f_\varrho+n_1f_1+\cdots+n_\varrho f_\varrho,
\]
其中 $\bar b_i$ 是 $\Sigma$ 的元素而 $n_i$ 是整數，還得到一個表示
\[
 f=b_1f_1+\cdots+b_\varrho f_\varrho
\]
其中 $b_i$ 屬於 $\Sigma$。因為 $f_i=\varepsilon f_i$，所以 $n_if_i=n_i\varepsilon f_i$；而 $n_i\varepsilon=(\varepsilon+\cdots+\varepsilon)$ 屬於 $\Sigma$，故 $b_i=\bar b_i+n_i\varepsilon$ 是 $\Sigma$ 的元素。關於 $T$ 的假設說明，$T$ 的每個元素 $\alpha$ 都有表示
\[
 \alpha=\bar a_1\tau_1+\cdots+\bar a_k\tau_k+n_1\tau_1+\cdots+n_k\tau_k,
\]
因而也有
\[
 \alpha=a_1\tau_1+\cdots+a_k\tau_k,
\]
第二個表示如上同樣由第一個表示推出，因為 $\varepsilon\tau_i=\tau_i$。

當 $\alpha$ 遍歷 $(\Sigma,T)$ 中模 $M$ 的元素時，$\tau_k$ 的係數 $a_k$ 遍歷 $\Sigma$ 的一個理想 $\ideal M_k$；由前述，對每個 $a_k\equiv0(\ideal M_k)$ 有
\[
 a_k=b_1a_k^{(1)}+\cdots+b_\varrho a_k^{(\varrho)}.
\]
令 $\alpha^{(i)}$ 為 $M$ 的一個元素，其 $\tau_k$ 的係數為 $a_k^{(i)}$。則 $\alpha-b_1\alpha^{(1)}\ldots-b_\varrho\alpha^{(\varrho)}$ 是屬於 $M$ 的元素，並且只依賴於 $\tau_1,\ldots,\tau_{k-1}$。對這些不再含 $\tau_k$ 且構成一個模 $M'$ 的元素全體重複這一過程；經過有限次重複即證明斷言。

\subsection*{\S~10. 多項式領域的特殊情形。}

1. 設作為基礎的環領域 $\Sigma$ 由 $x_1,\ldots,x_n$ 中具有任意復係數的所有多項式組成。由 Hilbert 基定理（Ann. 36），有限性條件對它成立。現在的問題是我們的定理同消去理論和模理論中的已知定理的關係。

這種關係由 Hilbert 一個已知定理的下列特殊情形建立：\footnote{Über die vollen Invariantensysteme. Math. Ann. 42 (1893), \S~3, p.~313.}

若 $f$ 在 $x_1,\ldots,x_n$ 的每個有限值系統上都消失，而這些值系統是一個素理想 $\ideal P$ 的所有多項式的零點（$\ideal P$ 的零點），則 $f$ 可被 $\ideal P$ 整除。換言之，一個素理想 $\ideal P$ 由所有在其零點上消失的多項式組成。\footnote{正如 Lasker 所示 [Math. Ann. 60 (1905), p. 607]，這個特殊情形對於齊次形式也可以直接證明；反過來，Hilbert 定理在齊次和非齊次情形中也可由它重新推出。我們可把該定理表述如下：若理想 $\ideal R$ 在 $M$ 的所有零點上消失，則 $\ideal R$ 的某個冪可被 $M$ 整除。這個定理以及上述特殊情形，只有當 $x$ 的取值範圍代數閉時才成立；所以它不能僅從我們的定理推出，而必須使用根的存在，例如在利用如下事實之處：若一個理想的零點只由 $x_1=0,\ldots,x_n=0$ 組成，則它從某個維度起含有 $x$ 的冪的所有乘積。剩餘證明可藉助我們的定理，相比 Lasker 的證明有所簡化。因為 Lasker 還必須訴諸 Hilbert 定理來證明一個理想分解為最大準素理想。}

因此，若乘積 $fg$ 在 $\ideal P$ 的所有零點上消失，則至少有一個因子消失；這些零點形成一個不可約代數簇。反過來，如果採用這個不可約代數簇的定義，則在一個不可約代數簇上消失的所有多項式的全體形成一個素理想。因此，素理想和不可約代數簇一一對應。由於 $\ideal P^\rho\equiv0(\ideal Q)$ 且 $\ideal Q\equiv0(\ideal P)$，準素理想的零點同其相應素理想的零點一致。\footnote{Macaulay（參見導言）把準素理想的這一性質，即擁有一個不可約代數簇，作為定義；Lasker 在定義中只納入簇的流形概念，除此之外則以抽象方式定義。只在 $x_1=0,\ldots,x_n=0$ 處消失的準素理想，在 Lasker 那裡具有特殊地位。} 因此，一個理想作為最大準素理想的最小公倍數的表示，給出了該理想所有零點到不可約代數簇的分解；並且，正如 Lasker 所示，反命題也成立。相應素理想唯一性的證明，在這裡對應於消去理論的基本定理，即一個代數簇唯一分解為不可約代數簇；而對於某些特殊多項式領域，在那裡多項式不能按該領域的不可約多項式唯一表示為乘積，因而也不存在消去理論，這一證明可作為該消去理論定理的等價物。

同孤立準素理想對應的不可約代數簇，正是那些出現在“最小預解式”中的代數簇；\footnote{例如參見 J. König, Einleitung in die allgemeine Theorie der algebraischen Größen (Leipzig, Teubner, 1903), p.~235.} 因為每個非孤立最大準素理想的零點，同時也是至少一個孤立最大準素理想的零點，即其相應素理想可被非孤立理想的相應素理想整除的那一個。孤立準素理想的唯一性由此在指數中給出新的不變量重數。關於把準素理想分解為不可約理想的唯一性定理，也可視為在重數意義上對消去理論的補充。

經過這些說明，各種不同分解可按它們關於代數簇的行為來解釋。兩兩互素理想對應於沒有共同零點的代數簇；對於互相相對素的理想，一個理想的任何不可約代數簇都不同時是公共零點；最大準素理想只在彼此全都不同的不可約代數簇上消失；而在分解為不可約理想時，相同的不可約代數簇可以帶重數出現。

還應注意，除了通常的多項式領域，也可把所有齊次形式的領域作為基礎；因為容易驗證，在那裡有效的合成下，一般定理仍然成立，其中加法只對同一維度的形式定義。\footnote{在不唯一情形中，非齊次分解可以同齊次分解並存，這由例子 $(x^3,xy,y^3)=[(x^3,y),(y^3,x)]=[(xy,x^3,y^3,x+y^2),(xy,x^3,y^3,y+x^2)]$ 表明。}

按照上述，一個最簡單的顯示四種不同分解的例子--其公式列在下面--由一條直線和兩條同它相交且同它異面的直線給出，其中後一類的一條直線還含有一個不同於交點且具有較高重數的點。分解為互素-不可約理想，對應於分解為這條直線和同它異面的簇；在分解為相對素-不可約理想時，這個簇分裂為兩條直線；分解為最大準素理想對應於分離出那個較高重數點；而分解為不可約理想則迫使該點進一步分解。

取這個點為原點，取所考察的直線為 $y$ 軸，取與它相交的直線平行於 $x$ 軸，取同它異面的直線平行於 $z$ 軸，則這樣的構型例如由下列不可約理想表示：\footnote{其中前三個作為素理想是不可約的；$\ideal B_4$ 不可約，因為它只有除子 $(x^2,y,z)$ 和 $(x,y,z)$；$\ideal B_5$ 不可約，因為每個除子都含有多項式 $xy$。}
\[
\begin{gathered}
 \ideal B_1=(x-1,y),\quad
 \ideal B_2=(y-1,z),\quad
 \ideal B_3=(x,z),\quad
 \ideal B_4=(x^3,y,z),\quad
 \ideal B_5=(x^2,y^2,z).
\end{gathered}
\]
相應素理想為
\[
 \ideal P_1=\ideal B_1,
 \quad \ideal P_2=\ideal B_2,
 \quad \ideal P_3=\ideal B_3,
 \quad \ideal P_4=\ideal P_5=(x,y,z).
\]
最大準素理想為
\[
 \ideal Q_1=\ideal B_1,
 \quad \ideal Q_2=\ideal B_2,
 \quad \ideal Q_3=\ideal B_3,
 \quad \ideal Q_4=[\ideal B_4,\ideal B_5]=(x^3,y^2,x^2y,z).
\]
相對素-不可約理想為
\[
 \ideal R_1=\ideal Q_1,
 \quad \ideal R_2=\ideal Q_2,
 \quad \ideal R_3=[\ideal Q_3,\ideal Q_4]=(x^3,x^2y,xy^2,z).
\]
互素-不可約理想為
\[
 \ideal S_1=\ideal R_1,
 \quad
 \ideal S_2=[\ideal R_2,\ideal R_3]=((y-1)x^3,(y-1)x^2y,(y-1)xy^2,z).
\]
於是總理想為
\[
\begin{aligned}
 \ideal M&=[\ideal S_1,\ideal S_2]=\ideal S_1\ideal S_2 \\
 &=((x-1)(y-1)x^3,(y-1)x^2y,(y-1)xy^2,(x-1)z,y(y-1)x^3,yz),
\end{aligned}
\]
其中
\[
 1=-(y-1)x^3+(y-1)(x^3-1)+y,
\]
\[
 (y-1)(x^3-1)+y\equiv0(\ideal S_1),
 \qquad -(y-1)x^3\equiv0(\ideal S_2).
\]
這裡 $\ideal B_1,\ideal B_2,\ideal B_3$ 是孤立的，因而在分解為不可約理想以及最大準素理想時也唯一確定。理想 $\ideal B_4$ 和 $\ideal B_5$，相應地 $\ideal Q_4$，是非孤立的；它們並不唯一確定，例如可以替換為
\[
 \ideal D_4=(x^3,y+\lambda x^2,z),
 \qquad
 \ideal D_5=(x^2+\mu xy,y^2,z),
\]
而 $\ideal Q_4$ 同樣可以替換為
\[
 \overline{\ideal Q}_4=(x^3,x^2y,y^2+\lambda xy,z).
\]

2. 正如一般多項式領域和整多項式領域滿足有限性條件一樣，每個有限整多項式領域也由 Hilbert 關於模基的定理滿足有限性條件；\footnote{反過來，如果一個多項式領域滿足有限性條件，並且每個多項式至少有一個表示，其中乘子關於 $x$ 的次數較低，則該領域是一個有限整領域。} 係數可以屬於任意域。我們再給出一個例子，即所有偶多項式構成的領域；這是最簡單的領域，在其中由於 $x^2y^2=(xy)^2$，多項式不能按該領域的不可約多項式唯一表示為乘積。這同前一個例子是同一構型，現在由下列不可約理想給出：
\[
\begin{gathered}
 \ideal B_1=(x^2-1,xy,y^2,yz),\quad
 \ideal B_2=(y^2-1,xz,yz,z^2),\\
 \ideal B_3=(x^2,xy,xz,yz,z^2),\quad
 \ideal B_4=(x^4,x^3y,y^2,xz,yz,z^2),\\
 \ideal B_5=(x^2,y^2,xz,yz,z^2).
\end{gathered}
\]
相應素理想為
\[
 \ideal P_1=\ideal B_1,
 \quad \ideal P_2=\ideal B_2,
 \quad \ideal P_3=\ideal B_3,
 \quad \ideal P_4=\ideal P_5=(x^2,xy,y^2,xz,yz,z^2).
\]
若要證明 $\ideal P_1$ 是素理想，只需注意該領域的每個多項式模 $\ideal P_1$ 都具有形式
\[
 f=\varphi(z^2)+xz\,\psi(z^2)\pmod{\ideal P_1}.
\]
因此若
\[
 f_1=\varphi_1(z^2)+xz\,\psi_1(z^2),
 \qquad
 f_2=\varphi_2(z^2)+xz\,\psi_2(z^2),
\]
則由 $f_1f_2\equiv0(\ideal P_1)$ 得到方程
\[
 \varphi_1\varphi_2+z^2\psi_1\psi_2=0,
 \qquad
 \varphi_2\psi_1+\varphi_1\psi_2=0,
\]
因而 $f_1\equiv0(\ideal P_1)$ 或 $f_2\equiv0(\ideal P_1)$。完全同樣可見 $\ideal P_2$ 是素理想；$\ideal P_3$ 是素理想，因為該領域的每個多項式模 $\ideal P_3$ 都同餘於一個 $y^2$ 中的多項式；$\ideal P_4$ 由該領域的所有多項式組成。因此 $\ideal B_1,\ideal B_2,\ideal B_3$ 作為素理想是不可約的，而 $\ideal B_4$ 和 $\ideal B_5$ 各自只有唯一的真因子 $\ideal P_4$，因而必然不可約。

由這些不可約理想得到最大準素理想
\[
\begin{gathered}
 \ideal Q_1=\ideal B_1,
 \quad \ideal Q_2=\ideal B_2,
 \quad \ideal Q_3=\ideal B_3,
 \quad \ideal Q_4=[\ideal B_4,\ideal B_5]=(x^4,x^3y,y^2,xz,yz,z^2),
\end{gathered}
\]
相對素-不可約理想
\[
 \ideal R_1=\ideal Q_1,
 \quad \ideal R_2=\ideal Q_2,
 \quad
 \ideal R_3=[\ideal Q_3,\ideal Q_4]=(x^4,x^3y,x^2y^2,xy^3,xz,yz,z^2),
\]
以及互素-不可約理想
\[
 \ideal S_1=\ideal R_1,
 \quad \ideal S_2=[\ideal R_2,\ideal R_3]
 =((y^2-1)x^4,(y^2-1)x^3y,(y^2-1)x^2y^2,
 (y^2-1)xy^3,xz,yz,z^2).
\]
同第一個例子一樣，
\[
 [\ideal B_4,\ideal B_5]=[\ideal D_4,\ideal D_5],
 \qquad
 [\ideal Q_3,\ideal Q_4]=[\ideal Q_3,\overline{\ideal Q}_4],
\]
其中
\[
 \ideal D_4=(x^4,xy+\lambda x^2,\ldots),
 \qquad
 \ideal D_5=(x^2+\mu xy,\ldots),
 \qquad \lambda\mu\ne1,
\]
並且
\[
 \overline{\ideal Q}_4=(x^4,x^3y,y^2+\lambda xy,\ldots).
\]
其餘不可約理想、相應地最大準素理想 $\ideal B_1,\ideal B_2,\ideal B_3$，作為孤立理想是唯一確定的。

\subsection*{\S~11. 數論以及微分表達式理論中的例子。}

1. 設基礎領域 $\Sigma$ 由所有偶有理整數構成。這樣，$\Sigma$ 可以同所有有理整數構成的領域一一對應，方法是把 $\Sigma$ 中的數 $2a$ 對應到數 $a$。於是立刻得到，$\Sigma$ 中每個理想都是主理想 $(2a)$；在基表示 $2c=n\cdot2a$ 中，對理想的每個元素 $2c$，奇整數 $n$ 只是不定有限和的縮寫。

該領域的素理想是 $\ideal P_0=\ideal O=(2)$ 以及 $\ideal P=(2p)$，其中 $p$ 表示奇素數；因此每個素理想都可被 $\ideal P_0$ 整除，但不可被其他素理想整除。準素理想為
\[
 \ideal Q_{\varrho_0}=(2\cdot2^{\varrho_0})
 \qquad\text{和}\qquad
 \ideal Q_\varrho=(2p^\varrho).
\]
它們同時也是不可約理想；並且根據關於素理想的說明，屬於不同奇素數的兩個理想互相相對素，但沒有任何 $\ideal Q$ 同任何 $\ideal Q_{\varrho_0}$ 相對素。\footnote{事實上，對奇素數 $p_1\ne p_2$，由 $2b\cdot2p_1^{\varrho_1}\equiv0(2p_2^{\varrho_2})$ 總得到 $2b\equiv0(2p_2^{\varrho_2})$；但由 $2b\cdot2p^\varrho\equiv0(2\cdot2^{\varrho_0})$ 只能得到 $2b=2\cdot2^{\varrho_0-1}\pmod{2\cdot2^{\varrho_0}}$。} 因而 $\ideal Q_0$ 是唯一的非孤立準素理想。$a$ 到素數冪的唯一分解，對應於理想 $(2a)$ 由最大準素、同時不可約的理想給出的唯一表示：
\[
 (2a)=[(2\cdot2^{e_0}),(2p_1^{e_1}),\ldots,(2p_\mu^{e_\mu})].
\]
因此在這裡，同多項式領域中的例子不同，即使非孤立最大準素理想也是唯一確定的。若 $\varrho_0=0$，這同時也是由互相相對素的理想給出的表示；而若 $\varrho_0>0$，則該理想是相對素-不可約的。

所以，在所有整數構成的領域中四種不同分解彼此重合，而在這裡，一方面只有分解為最大準素理想和分解為不可約理想這兩種分解重合；另一方面，當 $\varrho_0>0$ 時，互素-不可約分解（每個理想都是互素-不可約的，因為該領域沒有單位元）和相對素-不可約分解重合；而當 $\varrho_0=0$ 時，互素-不可約分解和相對素-不可約分解變得不同。同時，在這裡已經得到一個例子，其中一個素理想可被另一個素理想整除而不與它相同；更一般地可以看出，乘積表示並不從整除性推出。後一事實--這是該領域沒有單位元的結果--也說明，儘管每個理想都是主理想，該領域中的數並不存在按該領域不可約數的唯一乘積表示；因此，引入最小公倍數就成為必要。還應注意，如果不用所有偶整數，而用所有可被一個固定素數或素數冪整除的整數，情況也完全相同。

不過，如果 $\Sigma$ 由所有可被一個合數 $g=\mathfrak p_1^{\nu_1}\ldots\mathfrak p_s^{\nu_s}$ 整除的數構成，則不可約理想和準素理想也會變得不同。每個理想仍然是主理想 $(g\cdot a)$，素理想仍然是 $\ideal P_0=\ideal O=(g)$ 和 $\ideal P=(g\cdot p)$，其中 $p$ 是不同於出現在 $g$ 中的素數的一個素數。但是不可約理想為 $\ideal B_{\lambda_i}=(g\cdot p_i^{\lambda_i})$ 和 $\ideal B_e=(g\cdot p^e)$；準素理想為 $\ideal Q_e=\ideal B_e$ 以及不同於不可約理想的那些理想
\[
 \ideal Q_{\lambda_1\ldots\lambda_s}=(g\cdot p_1^{\lambda_1}\cdots p_s^{\lambda_s}),
\]
其中 $\ideal B_{\lambda_i}$ 和 $\ideal Q_{\lambda_1\ldots\lambda_\nu}$ 全都有同一個相應素理想 $\ideal P_0=(g)$。這裡分解為不可約理想仍然唯一，因而分解為最大準素理想也唯一；所以在這裡非孤立理想也同樣唯一確定。

2. 非交換環領域的一個例子由 Noether--Schmeidler 論文中處理的非交換多項式領域的理想理論給出。特別地，那裡有“完全可約”的理想，也就是說，這些理想的分解分量兩兩互素並且沒有真因子；這些分量因而當然是不可約的。因此，除了那裡按照 \S~9 證明的同構性之外，還得到兩個不同分解中分量數目的相等。對於偏或常線性微分表達式系統的分解--它作為該論文的一個特殊情形出現--由此得到一個結果，這個結果即使在單個常微分表達式的已知情形中也似乎未曾被注意到。

同時，固定理想 $M$ 的所有剩餘類所成的系統 $T$，連同非交換多項式領域 $\Sigma$，給出一個雙重領域 $(\Sigma,T)$，其中 $T$ 對 $\Sigma$ 具有模性質。因為兩個剩餘類之差仍是一個剩餘類，而一個剩餘類同任意多項式的乘積也仍如此；然而兩個剩餘類的乘積並不存在（前引文獻，\S~3）。因此，那裡稱為“子群”的剩餘類系統，給出了雙重領域 $(\Sigma,T)$ 中的模的例子，其中基礎環領域 $\Sigma$ 是非交換的。

\subsection*{\S~12. 初等因子理論中的例子。}

這裡所說的是由一般發展所決定的初等因子理論的理解方式，不過該理論本身被假定為已知。

令 $\Sigma$ 為所有含 $n^2$ 個元素的整矩陣構成的領域，加法和乘法按矩陣的通常意義定義。於是 $\Sigma$ 是一個非交換環領域；因此理想一般是單側的，雙側理想只是特殊情形。\footnote{這些領域的理想理論是 Du Pasquier 的著作主題：Zahlentheorie der Tettarionen, dissertation, Zürich, Vierteljahrsschr. d. Naturf. Ges. Zürich, 51 (1906); Zur Theorie der Tettarionenideale, ibid., 52 (1907)。第二篇的內容是證明每個理想都是主理想。}

我們首先證明每個理想都是主理想。為此，在右理想的情形中，對每個矩陣 $A=(a_{ik})$ 配以整線性形式的模
\[
 \mathsf{A}=(a_{11}\xi_1+\cdots+a_{1n}\xi_n,
 \ldots,
 a_{n1}\xi_1+\cdots+a_{nn}\xi_n).
\]
反過來，對這個模，對應於每個給出 $\mathsf{A}$ 的一個基的矩陣，因而除了 $A$ 之外，也對應於 $UA$，其中 $U$ 是幺模的。更一般地，對乘積 $PA$，對應一個模 $\mathsf{B}$，它是 $\mathsf{A}$ 的倍數。$\mathsf{A}$ 中的單個線性形式由這樣的矩陣 $P$ 給出，即這些矩陣只有一行非零。

現在令 $A_1,A_2,\ldots,A_\nu,\ldots$ 為理想 $\ideal M$ 的所有元素；令 $\mathsf{A}_1,\mathsf{A}_2,\ldots,\mathsf{A}_\nu,\ldots$ 為配給它們的模；令 $\mathsf{A}$ 為這些模的最大公因子，而 $UA$ 為配給這個模 $\mathsf{A}$ 的最一般矩陣。按最大公因子的定義，對 $\mathsf{A}$ 中的每個單個線性形式，都對應一個矩陣 $P_1A_{i_1}+\cdots+P_\sigma A_{i_\sigma}$，其中如上所述，各 $P$ 只有一行非零。由此可知，對應於 $\mathsf{A}$ 的一個基的矩陣 $A$，也有這樣的表示，而且現在用一般的 $P$，因而成為 $\ideal M$ 的元素。又由於每個模 $\mathsf{A}_i$ 都可被 $\mathsf{A}$ 整除，每個矩陣 $A_i$ 也可被 $A$ 整除；於是 $A$，以及一般地 $UA$，構成 $\ideal M$ 的一個基。若處理左理想，則相應地要把每個矩陣的列看作一個模的基；每個理想都是主理想，對它而言，除 $A$ 以外，$AV$ 也是一個基，其中 $V$ 是任意幺模矩陣。

為了同初等因子理論相聯繫，下面我們以雙側理想為基礎；因此特別地，只要 $A$ 屬於理想，則 $PAQ$ 也屬於理想。由前述，這種理想的最一般基為 $UAV$，其中 $U$ 和 $V$ 都是幺模的。於是基元素窮盡一個等價矩陣類，\footnote{單側理想的基元素對應於右類或左類。} 並且理想和類之間存在一一對應，因而理想同該類的初等因子系統 $(a_1\mid a_2\mid\cdots\mid a_n)$ 之間也存在一一對應；這裡 $a_i$ 如所周知是非負整數，每一個都整除下一個。因此，由初等因子確定的標準形中出現的該類矩陣，可看作理想的一個特殊基；理想、相應地類的整除性，推出初等因子的整除性，反過來也成立。

但 Du Pasquier 已在前引文獻 \S~11 中證明，對於每個雙側理想，秩為 $n$，並且所有初等因子一致。為了能夠考察一般初等因子的情形，我們必須不是從理想出發，而是直接從雙側類出發（這些類同樣用大德文字母表示）。一個類 $\mathfrak A=UAV$ 可被另一個類 $\mathfrak B$ 整除，是指 $A=PBQ$。

一般地說：\emph{兩個類的最小公倍數（最大公因子）通過形成相應初等因子系統的最小公倍數（最大公因子）而得到}。\footnote{在論文 Zur Theorie der Moduln, Math. Ann. 52 (1899), p.~1 中，E. Steinitz 通過初等系統的最小公倍數和最大公因子定義類的最小公倍數（最大公因子）。獨立地，類的最小公倍數作為“Kongruenzkomposition”出現在 H. Brandt, Komposition der binären quadratischen Formen relativ einer Grundform, J. f. M. 150 (1919), p.~1 中。} 設
\[
 (a_1\mid a_2\mid\cdots\mid a_n),\quad
 (b_1\mid b_2\mid\cdots\mid b_n),\quad
 (c_1\mid c_2\mid\cdots\mid c_n),
\]
其中 $c_i=[a_i,b_i]$，分別是 $\mathfrak A$、$\mathfrak B$ 和 $\mathfrak C^*$ 的初等因子系統，並令 $\mathfrak C=[\mathfrak A,\mathfrak B]$。於是 $\mathfrak C^*$ 可被 $\mathfrak A$ 和 $\mathfrak B$ 整除，因而可被 $\mathfrak C$ 整除；反過來，$\mathfrak C$ 的初等因子可被 $\mathfrak C^*$ 的初等因子整除，所以 $\mathfrak C$ 可被 $\mathfrak C^*$ 整除，從而 $\mathfrak C=\mathfrak C^*$。最大公因子的證明相應得到。

因此，初等因子 $a_i$ 唯一表示為素數冪的最小公倍數，對應於把 $\mathfrak A$ 表示為類 $\mathfrak Q$ 的最小公倍數，其中 $\mathfrak Q$ 的初等因子由一個素數的冪給出，記號為
\[
 \mathfrak Q\sim(p^{r_1}\mid p^{r_2}\mid\cdots\mid p^{r_\varrho}\mid0\mid\cdots\mid0),
 \qquad r_1\le r_2\le\cdots\le r_\varrho.
\]
特別地，若秩 $\varrho$ 等於 $n$，則這裡得到一個由互素且互素-不可約類給出的分解；儘管領域非交換，該分解仍然唯一。

類 $\mathfrak Q$ 可以進一步分解為對應於下列初等因子系統的類：
\[
\begin{gathered}
 \mathfrak B_1\sim(p^{r_1}\mid\cdots\mid p^{r_1}),\quad
 \mathfrak B_2\sim(1\mid p^{r_2}\mid\cdots\mid p^{r_2});\quad \ldots,\\
 \mathfrak B_\varrho\sim(1\mid\cdots\mid1\mid p^{r_\varrho}\mid\cdots\mid p^{r_\varrho});\quad
 \mathfrak B_{\varrho+1}\sim(1\mid\cdots\mid1\mid0\mid\cdots\mid0),
\end{gathered}
\]
其中在 $\mathfrak B_\nu$ 中數 $1$ 出現 $(\nu-1)$ 次。例如，如果 $r_1=r_2=\cdots=r_\mu=0$，則 $\mathfrak B_1,\ldots,\mathfrak B_\mu$ 等於單位類，應從分解中省去；若 $\varrho=n$，對 $\mathfrak B_{\varrho+1}$ 也同樣如此。再若例如 $r_\nu=r_{\nu+1}=\cdots=r_{\nu+\lambda}$，則 $\mathfrak B_{\nu+1},\ldots \mathfrak B_{\nu+\lambda}$ 是 $\mathfrak B_\nu$ 的真因子，也應刪去。把剩下的那些記作 $\mathfrak B_{i_1},\ldots,\mathfrak B_{i_k}$，它們現在給出一個最短表示。於是
\[
 \mathfrak Q=[\mathfrak B_{i_1},\ldots,\mathfrak B_{i_k}]
\]
是 $\mathfrak Q$ 分解為不可約類的唯一分解。事實上，設
\[
 \mathfrak B_\nu\sim(1\mid\cdots\mid1\mid p^{r_\nu}\mid\cdots\mid p^{r_\nu})
\]
可表示為下列兩個類的最小公倍數：
\[
 \mathfrak C\sim(1\mid\cdots\mid1\mid p^{s_1}\mid\cdots\mid p^{s_\lambda})
 \quad\text{和}\quad
 \mathfrak D\sim(1\mid\cdots\mid1\mid p^{t_1}\mid\cdots\mid p^{t_\mu}).
\]
那麼，在 $\mathfrak C$ 和 $\mathfrak D$ 的初等因子系統中，數 $1$ 必須位於前 $(\nu-1)$ 個位置；在第 $\nu$ 個位置，指數 $s_\nu$ 或 $t_\nu$ 中的一個，例如 $s_\nu$，必須等於 $r_\nu$。但是由於
\[
 r_\nu=s_\nu\le s_{\nu+1}\le\cdots\le s_i\le r_\nu,
\]
得到 $\mathfrak C=\mathfrak B_\nu$，所以 $\mathfrak B_\nu$ 是不可約的。\footnote{不過，$\mathfrak B$ 仍然可分解為單側類；這裡初等因子和類之間不再有唯一關係，因而也不再有到不可約單側類的唯一分解。下面這個由 H. Brandt 告知我的例子，展示了分解為右側類的情形（其中類由基矩陣、或由相應模表示）：
\[
\begin{gathered}
 \mathfrak B=[\mathfrak C_1,\mathfrak C_2]=[\mathfrak D_1,\mathfrak D_2];\quad
 \mathfrak B\sim\begin{pmatrix}p&0\\0&p\end{pmatrix};\quad
 \mathfrak C_1\sim\begin{pmatrix}1&0\\0&p\end{pmatrix};\quad
 \mathfrak C_2\sim\begin{pmatrix}p&0\\0&1\end{pmatrix};\\
 \mathfrak D_1\sim\begin{pmatrix}1&0\\0&p\end{pmatrix};
 \quad
 \mathfrak D_2\sim\begin{pmatrix}p&0\\(p-1)&1\end{pmatrix}.
\end{gathered}
\]
事實上，模 $(\xi,p\eta)$、$(p\xi,\eta)$ 和 $(p\xi,(p-1)\xi+\eta)$ 各自唯一的真因子都是模 $(\xi,\eta)$，因而它們不可約並且彼此不同。} 對 $\mathfrak B_{\varrho+1}$ 也同樣成立，只需把 $p$ 替換為 $0$。但是，正如形成最小公倍數所表明的，每個這樣的不可約類給出一個確定指數，以及該指數在 $\mathfrak Q$、相應地在 $\mathfrak A$ 的初等因子系統中第一次出現的位置；而 $\mathfrak B_{\varrho+1}$ 給出秩。由於這些數由 $\mathfrak Q$、相應地 $\mathfrak A$ 的初等因子系統唯一確定，並且由於初等因子和類之間的關係是一一對應的，所以 $\mathfrak Q$ 以及任意類到不可約類的分解都是唯一的。

綜上：\emph{每個由有界元素數的整矩陣給出的雙側類 $\mathfrak A$，都可以唯一地表示為有限多個不可約雙側類的最小公倍數。每個不可約類表示 $\mathfrak A$ 的初等因子系統中的一個固定素因子、一個相應指數，以及該指數第一次出現的位置。與因子 $0$ 對應的不可約類給出 $\mathfrak A$ 的秩。}

\begin{flushleft}
Erlangen, 1920 年 10 月。
\end{flushleft}

%
\clearpage
%

%
%
%

\setcounter{footnote}{0}
\pdfbookmark[0]{20. 絕對不可約性的一個代數判據}{paper-20}
\section*{20. 絕對不可約性的一個代數判據}
\begin{center}
Math. Ann. 85 (1922), pp. 26--33
\end{center}

一個多項式--其係數取自任意的、抽象定義的域--稱為\emph{絕對不可約}，如果它在可把係數域擴張到的代數閉域中仍然不可約。\footnote{每個域都可以，而且本質上唯一地，可以擴張為代數閉域，也就是說擴張為一個每個一元多項式都分解為一次因子的域；這是 E. Steinitz 在其 \emph{Algebraische Theorie der K\"orper} (J. f. M. 137 (1910), p.~167) 中證明的。由此他給出了代數學基本定理的有理等價物。} 下面我將證明，絕對不可約性，與相對於一個預先給定的域的不可約性不同，可以用代數方式把握。更確切地說，有如下定理：

\emph{對每個含 $n\geqq2$ 個變元、具有不定係數的多項式，可以指定一個關於這些係數以及另外一些不定元的、具有整係數的整有理函數，即可約性形式，使得對於每個同次數的特化多項式，絕對不可約的充要條件由該可約性形式不為零給出。換言之：對於每一組使可約性形式在不定元中恆等消失、且不導致次數降低的係數特化值，相應的特化多項式是可約的；反過來也成立。}

如果不考慮多項式，而考慮多一個變元的齊次形式，則次數條件被更簡單的條件所代替，即係數系統不恆等為零。

作為該定理的直接推論，還得到最早由 A. Ostrowski 提出的命題\footnote{Zur arithmetischen Theorie der algebraischen Gr\"o{\ss}en. G\"ott. Nachr. 1919, p.~279（引理 p.~296）。對於由通常數構成的域的情形，據我所知，Ostrowski 也得到了上述主定理；他將很快發表有關研究。在我的證明安排中，主定理對任意的、抽象定義的域成立，這依賴於這樣一個事實：為不定元構造的可約性形式具有與所取基域無關的整數係數。}：

\emph{每個以代數數為係數的絕對不可約多項式，在包含其係數域的有限代數數域 $\Omega$ 的每個素理想模下仍然不可約，至多有限多個 $\Omega$ 中的素理想例外。特別地，$\Omega$ 也可以是有理數域。}

關於對相對整函數理論的進一步應用\footnote{D. Hilbert, Mathematische Probleme. G\"ott. Nachr. 1900, p.~253, Problem 14.}，我打算在別處討論。

1. 我們先證明，\emph{每個含 $n\geqq2$ 個變元並具有不定係數的多項式都是絕對不可約的}。設
\begin{equation}
\begin{aligned}
F(x)&=\sum A_{i_1\ldots i_n}x_1^{i_1}\ldots x_n^{i_n} &&(i_1+\ldots+i_n\leqq l),\\
G(x)&=\sum B_{j_1\ldots j_n}x_1^{j_1}\ldots x_n^{j_n} &&(j_1+\ldots+j_n\leqq m),\\
H(x)&=F(x)G(x)=\sum C_{k_1\ldots k_n}(A,B)x_1^{k_1}\ldots x_n^{k_n} &&(k_1+\ldots+k_n\leqq l+m),
\end{aligned}
\tag{1}
\end{equation}
其中 $A$ 和 $B$ 表示不定元，而 $C(A,B)$ 是 $A$ 和 $B$ 的整雙線性組合。因此商
\[
D=C(A,B)/C_{0\ldots0}(A,B)
\]
是商 $A/A_{0\ldots0}$ 和 $B/B_{0\ldots0}$ 的有理函數；但是這些函數的個數大於其自變量的個數，因為相應的個數分別為
\[
\binom{l+m+n}{n}-1,
\qquad
\binom{l+n}{n}-1,
\qquad
\binom{m+n}{n}-1,
\]
並且
\[
\binom{l+m+n}{n}+1>
\binom{l+n}{n}+\binom{m+n}{n}
\quad\text{當 } n\geqq2,\ l>0,\ m>0.
\]
於是 $D(A,B)$ 之間至少存在\emph{一個}代數關係，從而 $C(A,B)$ 之間也至少存在一個代數關係：
\begin{equation}
 \varPhi(Z)\ne0\quad[\text{在 }Z_{k_1\ldots k_n}\text{ 中恆等}],
 \qquad
 \varPhi(C(A,B))=0\quad[\text{在 }A,B\text{ 中恆等}],
\tag{2}
\end{equation}
其中 $\varPhi$ 表示一個整係數多項式。

因此，以不定元 $Z$ 為係數的多項式
\begin{equation}
 E(x)=\sum Z_{k_1\ldots k_n}x_1^{k_1}\ldots x_n^{k_n}
 \qquad(k_1+\ldots+k_n\leqq l+m)
\tag{3}
\end{equation}
在 $n\geqq2$ 時必定是絕對不可約的。\footnote{因為按照第 2 節的說法，不定元 $Z$ 不是 $\mathfrak P$ 的零點。}

2. 在代入 $Z=C(A,B)$ 後恆等地在 $A,B$ 中消失的所有多項式 $\varPhi(Z)$ 構成一個\emph{多項式理想}，\footnote{這樣的整體通常稱為模；但事實上，刻畫它的性質--它與 $\varPhi_1$ 和 $\varPhi_2$ 一起也含有 $\varPhi_1\mp\varPhi_2$，與 $\varPhi$ 一起也含有 $\Psi\varPhi$，其中 $\Psi$ 表示 $Z$ 中任意的整係數多項式--正是理想性質。}更確切地說，是一個素理想 $\mathfrak P$；因為如果在此代入下一個乘積 $\varPhi_1\varPhi_2$ 消失，那麼至少有一個因子消失。由於 (2) 在 $A,B$ 中恆等成立，它對每個特化值系 $\mathsf{A},\mathsf{B}$ 仍然成立，其中 $\mathsf{A},\mathsf{B}$ 以及相應的 $\mathsf{\Gamma}=C(\mathsf{A},\mathsf{B})$ 表示某個域中的量。\emph{因此，每個在擴張域中可約的多項式，其係數 $\mathsf{\Gamma}$ 滿足條件 $\varPhi(\mathsf{\Gamma})=0$，也就是 $\mathfrak P$ 的零點，}或者說是 $\mathfrak P$ 的有限多個基多項式的零點。現在還要證明反命題。

為此要注意，當通過引入一個新變元 $y$ 把多項式 (1) 變成次數為 $l,m,l+m$ 的齊次形式 $F(x,y)$、$G(x,y)$、$H(x,y)$ 時，$A,B,C$ 之間的所有關係都保持不變。於是，例如當 $F(x,y)$\footnote{由把 $F,G,\ldots,f,g,\ldots$ 中的不定元替換成特化值系而得到的多項式和形式，在下文一律用上劃線表示：$\bar F,\bar G,\ldots,\bar f,\bar g,\ldots$。}變成 $y$ 的冪時，相應的係數 $\Gamma$ 也成為 $\mathfrak P$ 的零點；對它們而言，轉到非齊次多項式會導致 $\bar H(x)$ 的次數降低，而不是分解。下面將證明，除了這種次數降低之外，反命題總是成立；特別地，\emph{對齊次形式而言，只要並非所有 $\mathsf{\Gamma}$ 都消失，反命題就無例外地成立。換言之，$\mathfrak P$ 的每個零點都由參數表示 $\mathsf{\Gamma}=C(\mathsf{A},\mathsf{B})$ 給出，其中 $\mathsf{A},\mathsf{B}$ 是某個擴張域中的量}。\footnote{“一般地說”$\mathfrak P$ 的零點由參數表示給出，是由素理想 $\mathfrak P$ 定義不可約代數簇這一性質推出的。但眾所周知，這並不意味着--這是我最初忽略、且 A. Ostrowski 很早以前提醒我的一點--參數表示在任何特化值系下都不會失效。這裡給出的證明依賴較為初等的概念。}

我通過直接構造有限多個函數 $\varPhi(Z)$ 來證明這一點：藉助 Kronecker 代換
\[
 x_\lambda=\xi^{d^{\,n-\lambda}}
\]
我把 (1) 中的多項式指定為\emph{一}個變元的多項式；證明這裡在指數中出現空缺，並通過形成相應係數的範數而得到函數 $\varPhi(Z)$ 和所需的判據。

3. 置\footnote{Festschrift, p.~11.}
\begin{equation}
 d=l+m+1,\qquad
 i=i_1d^{n-1}+i_2d^{n-2}+\ldots+i_n,
\quad
 j=j_1d^{n-1}+\ldots+j_n,
\quad
 k=k_1d^{n-1}+\ldots+k_n.
\tag{4}
\end{equation}
於是，在 Kronecker 代換下，(1) 中的多項式 $F,G,H$ 分別對應於多項式
\begin{equation}
\begin{aligned}
 f(\xi)&=\sum A_{i_1\ldots i_n}\xi^i &&(i_1+\ldots+i_n\leqq l),\\
 g(\xi)&=\sum B_{j_1\ldots j_n}\xi^j &&(j_1+\ldots+j_n\leqq m),\\
 h(\xi)&=\sum C_{k_1\ldots k_n}(A,B)\xi^k &&(k_1+\ldots+k_n\leqq l+m).
\end{aligned}
\tag{5}
\end{equation}
這個對應是眾所周知的一一對應，因為對 (5) 中出現的一切誘導指數，由 $k=0$ 可推出 $k_1=0,\ldots,k_n=0$；因而由 $i+j=i'+j'$ 也推出 $i_1+j_1=i'_1+j'_1;\ldots;i_n+j_n=i'_n+j'_n$。所以，不同的乘積 $\xi^i\xi^j$ 只有在相應乘積 $x_1^{i_1}\ldots x_n^{i_n}x_1^{j_1}\ldots x_n^{j_n}$ 也給出同一指數時，才可能給出同一個指數 $\xi^k$。於是得到
\begin{equation}
 h(\xi)=f(\xi)g(\xi),
\tag{6}
\end{equation}
並且，如果關係 (6) 成立，且 $f(\xi)$ 和 $g(\xi)$ 中沒有出現誘導指數以外的指數，則反過來得到
\begin{equation}
 H(x)=F(x)G(x).
\tag{7}
\end{equation}
這裡，\emph{在 $f(\xi)$ 和 $g(\xi)$ 的誘導指數中確實出現空缺}。因為 $f(\xi)$ 的最高指數等於 $l\cdot d^{n-1}$，次高指數等於 $(l-1)d^{n-1}+d^{n-2}$；它們的差為 $d^{n-2}\cdot(d-1)>1$，這是因為 $d=l+m+1\geqq3$ 且 $n\geqq2$；對 $g(\xi)$ 也同樣如此。

對不定元 $A,B$ 成立的關係 (6) 和 (7)，對每個特化值系 $\mathsf{A},\mathsf{B}$ 仍然成立。為了保持次數，我藉助變元 $\eta$ 和 $y$ 將 (6) 和 (7) 齊次化。\emph{因此，一個齊次形式 $\bar{H}(x,y)$ 在某個代數擴張域中分裂，當且僅當在該域中相應的二元形式 $\bar{h}(\xi,\eta)$ 可以分裂為兩個因子，且在這兩個因子中不出現誘導指數以外的指數。}

4. 為了實行第 2 節末尾提到的範數構造，令
\[
 a_1,b_1;\ldots; a_p,b_p;\quad a_{p+1},b_{p+1};\ldots; a_{p+q},b_{p+q}
\]
為新的不定元。置
\begin{equation}
\begin{aligned}
 \varphi(\xi,\eta)&=(a_1\xi+b_1\eta)\ldots(a_p\xi+b_p\eta)
     =r_p\xi^p+r_{p-1}\xi^{p-1}\eta+\ldots+r_0\eta^p,\\
 \psi(\xi,\eta)&=(a_{p+1}\xi+b_{p+1}\eta)\ldots(a_{p+q}\xi+b_{p+q}\eta)
     =s_q\xi^q+s_{q-1}\xi^{q-1}\eta+\ldots+s_0\eta^q,\\
 \chi(\xi,\eta)&=\varphi(\xi,\eta)\psi(\xi,\eta)
     =t_{p+q}\xi^{p+q}+t_{p+q-1}\xi^{p+q-1}\eta+\ldots+t_0\eta^{p+q},
\end{aligned}
\tag{8}
\end{equation}
其中 $r,s,t$ 分別表示齊次化的基本對稱函數；也就是對兩列 $a,b$ 中每一列均為線性的對稱函數，所有在每一列中齊次且同維數的整對稱函數，都能且只能以一種方式由它們整地、並以整數係數組合而成。

現在用另外的不定元 $u_{\mu\nu}$ 形成線性形式
\begin{equation}
 \zeta(u,a,b)=\sum u_{\mu\nu}r_\mu s_\nu,
\tag{9}
\end{equation}
其中求和遍及預先給定的指標 $\mu$ 和 $\nu$。把 $\zeta(u)$ 與所有由置換 $a,b$ 這列而得、且與 $\zeta(u)$ 以及彼此不同的共軛相乘--也就是當 $\zeta(u)$ 被看作 $r_\mu(a,b)$ 和 $s_\nu(a,b)$ 所成域的函數時，$\zeta(u)$ 關於 $t(a,b)$ 所成域的範數 $N(\zeta(u))$--就得到所有列 $a,b$ 的一個整對稱函數，它在每一列中齊次且同維數。因此按上面所述，它是 $t(a,b)$ 和 $u_{\mu\nu}$ 的唯一確定的、具有整數係數的整函數：
\begin{equation}
 N(\zeta(u,a,b))=T(t(a,b),u)\qquad[\text{在 }a,b,u\text{ 中恆等}].
\tag{10}
\end{equation}
同樣的理由表明，也有
\[
N(z-\zeta(u))=z^\delta+T_{\delta-1}(t,u)z^{\delta-1}+\ldots+T(t,u)
\qquad[\text{在 }a,b,u,z\text{ 中恆等}],
\]
其中所有 $T$ 都表示 $t$ 和 $u$ 的整係數整函數。這個恆等式的兩邊在 $z=\zeta(u)$ 時消失；並且由於基本對稱函數 $r(a,b)$、$s(a,b)$ 的代數獨立性，當按 (8) 把 $t(a,b)$ 置為 $r$ 和 $s$ 的整雙線性組合 $w(r,s)$，並把 $\zeta(u)$ 看作 $r$ 和 $s$ 的函數時，它們實際上在 $r,s$ 中恆等消失。因此得到
\begin{equation}
\zeta(u)^\delta+T_{\delta-1}(w(r,s),u)\zeta(u)^{\delta-1}+\ldots+T(w(r,s),u)=0
\quad[\text{在 }r,s,u\text{ 中恆等}].
\tag{11}
\end{equation}\footnote{如果 (9) 中的求和遍及所有指標，那麼 (11) 就給出 Kronecker 對 Gauss 定理的推廣，並且本質上採用 Kronecker 的證明安排（Zur Theorie der Formen h\"oherer Stufen, Berliner Ber. 1883, Werke, 2, p.~417）。}

恆等式 (10) 和 (11) 對 $a,b$ 的所有特化值系 $\alpha,\beta$，以及 $r,s$ 的所有特化值系 $\varrho,\sigma$，仍然成立。特別地，如果選擇 $\varrho,\sigma$ 使得所有乘積 $\varrho_\mu\sigma_\nu$ 都消失--這裡 $\mu,\nu$ 表示 (9) 中出現的指標--從而 $\zeta(u)$ 在代入後消失，那麼令 $w(\varrho,\sigma)=\tau$，由 (11) 得到
\begin{equation}
 T(\tau,u)=0\qquad[\text{在 }u\text{ 中恆等}].
\tag{12}
\end{equation}

反過來，設 $\tau_0,\ldots,\tau_{p+q}$ 是不全為零的值系，並且滿足 (12)。由於在一個代數擴張域中存在分解
\begin{equation}
 \bar\chi(\xi,\eta)=\tau_{p+q}\xi^{p+q}+\ldots+\tau_0\eta^{p+q}
=(\alpha_1\xi+\beta_1\eta)\ldots(\alpha_{p+q}\xi+\beta_{p+q}\eta),
\tag{12'}
\end{equation}
其中任何一個因子中 $\alpha$ 和 $\beta$ 都不同時為零，\footnote{這個（有理的）“代數學基本定理”是存在性定理，第 2 節末尾所說的反命題無例外成立正是建立在它之上。}雖然某些 $\alpha$ 或 $\beta$ 可以為零，於是有 $\tau=t(\alpha,\beta)$。把這些值代入 (10)，並利用 (12)，可知 $\zeta(u,\alpha,\beta)$ 或其某個共軛因子在 $u$ 中恆等消失。因此，在適當重新編號 $\alpha,\beta$ 後，若置 $r(\alpha,\beta)=\varrho$、$s(\alpha,\beta)=\sigma$，這正是說 $\bar\chi(\xi,\eta)$ 至少允許一個分解：
\begin{equation}
\bar\chi(\xi,\eta)=\bar\varphi(\xi,\eta)\bar\psi(\xi,\eta)
=\sum \varrho_\mu\xi^\mu\eta^{p-\mu}\cdot\sum \sigma_\nu\xi^\nu\eta^{q-\nu},
\tag{13}
\end{equation}
使得 $\zeta(u)$ 中出現的所有乘積 $\varrho_\mu\sigma_\nu$ 都消失，但並非所有 $\varrho$ 或所有 $\sigma$ 都消失。

5. 現在令 (8) 中的次數 $p,q$ 分別等於 $ld^{n-1}$ 和 $md^{n-1}$，也就是 (5) 中 $f(\xi)$ 和 $g(\xi)$ 的次數。把指標 $\mu,\nu$ 分成 $\mu^{(1)},\nu^{(1)},\mu^{(2)},\nu^{(2)}$；其中 $\mu^{(1)}$ 遍歷 $f(\xi)$ 中缺失的所有指數，$\nu^{(1)}$ 遍歷 $\psi(\xi,\eta)$ 中 $\xi$ 的所有指數；相應地，$\nu^{(2)}$ 遍歷 $g(\xi)$ 中缺失的所有指數，$\mu^{(2)}$ 遍歷 $\varphi(\xi,\eta)$ 中 $\xi$ 的所有指數。又令
\[
 e(\xi)=\sum Z_{k_1\ldots k_n}\xi^k\qquad(k_1+\ldots+k_n\leqq l+m)
\]
表示與多項式 (3) 對應的一元多項式；令
\[
 S(Z,u)
\]
表示 $Z$ 和 $u$ 的整係數多項式，它由 $T(t,u)$--也就是 (10) 中唯一確定的右端，此處 (9) 中的求和遍及上述指標 $\mu^{(1)},\nu^{(1)},\mu^{(2)},\nu^{(2)}$--通過把 $t$ 替換成係數 $Z$ 以及 $e(\xi)$ 的相應零點而得到。

現在把 $r,s$ 替換成 $f(\xi)$ 和 $g(\xi)$ 的係數 $A,B$ 以及相應零點，則在指定的求和下，$\zeta(u)$ 通過該代入而消失。此外，$t=w(r,s)$ 變為 $h(\xi)$ 的係數 $C(A,B)$ 及其相應零點；所以 $T(t,u)$ 變為 $S(C(A,B),u)$。由恆等式 (11) 於是得到
\begin{equation}
 S(C(A,B),u)=0\qquad[\text{在 }A,B,u\text{ 中恆等}].
\tag{14}
\end{equation}
由於 (14) 對每個特化值系仍然成立，所以那些在擴張域中分裂為兩個因子的特化 $H(x)$ 或更一般的 $H(x,y)$，其係數 $\mathsf{\Gamma}=C(\mathsf{A},\mathsf{B})$ 滿足條件
\begin{equation}
 S(\Gamma,u)=0\qquad[\text{在 }u\text{ 中恆等}].
\tag{15}
\end{equation}

反過來，假設對於多項式 $\bar H(x)$、相應地對於形式 $\bar H(x,y)$ 的不全為零的係數 $\mathsf{\Gamma}$，條件 (15) 成立。按上面所述，這等價於 $T(\tau,u)=0$，其中 $\tau$ 表示 $\bar h(\xi,\eta)$ 的全部係數。因此由 (13)，在 $\mathsf{\Gamma}$ 的一個擴張域中存在分解
\[
 \bar h(\xi,\eta)=\bar f(\xi,\eta)\cdot\bar g(\xi,\eta),
\]
並且由於指定的求和，$\bar f$ 和 $\bar g$ 中不出現非誘導指數；於是存在關係
\[
 \bar H(x,y)=\bar F(x,y)\cdot\bar G(x,y).
\]
由此，只要在 $y=1$ 時沒有因子變為零次--特別是當 $\mathsf{\Gamma}$ 不導致 $H(x)$ 的次數降低時--還得到關係
\[
 \bar H(x)=\bar F(x)\cdot\bar G(x).
\]

\emph{因此，條件 (15) 是 $\bar H(x,y)$ 在擴張域中分裂為兩個因子的充要條件；在排除次數降低時，它也是 $\bar H(x)$ 分裂為兩個因子的充要條件。}

此外還推出，$S(Z,u)$ 不可能在 $Z$ 和 $u$ 中恆等消失，因為第 1 節已證明，當 $n\geqq2$ 時，以不定元為係數的多項式，因而相應的多一個變元的齊次形式，都是絕對不可約的。於是，若用 $U$ 表示 $u$ 的冪乘積，則有
\[
 S(Z,u)=\sum \varPhi_\lambda(Z)U_\lambda,
\]
其中 $\varPhi_\lambda(Z)$ 依 (14) 是屬於 $\mathfrak P$ 的多項式。由此證明，$\mathfrak P$ 沒有除參數表示 $\mathsf{\Gamma}=C(\mathsf{A},\mathsf{B})$ 所給出的零點之外的其他零點，其中 $\mathsf{A}$ 和 $\mathsf{B}$ 屬於 $\mathsf{\Gamma}$ 的某個代數擴張域。第 2 節末尾所說的反命題由此得到證明。

6. 現在可以直接表述\emph{絕對不可約性的條件}。只需把 (3) 中 $E(x)$ 的次數按一切可能方式分解為兩個正整數之和 $l+m$，並對每個這樣的分解形成形式 $S(Z,u)$。這些形式的乘積
\begin{equation}
 R(Z,u)=\prod S(Z,u)
\tag{16}
\end{equation}
就構成 $E(x)$ 的可約性形式；因為按上面所述，$E(x)$、相應地 $E(x,y)$ 的每個分解都對應於 (16) 中一個因子的消失，反過來也成立。這樣便證明了引言中提出的主定理，同時也說明了如何用有限多個步驟實際構造可約性形式。

7. 引言中提到的 Ostrowski 定理由可約性形式 (16) 的存在性立即推出。設
\[
 \bar E(x)=\sum \mathsf{\Gamma}_{k_1\ldots k_n}x_1^{k_1}\ldots x_n^{k_n}
 \qquad(k_1+\ldots+k_n\leqq v)
\]
是一個以代數整數 $\mathsf{\Gamma}$ 為係數的絕對不可約多項式；設 $\Omega$ 表示包含所有 $\mathsf{\Gamma}$ 的有限代數數域，$\mathfrak p$ 是 $\Omega$ 的一個素理想。

於是，為 $\bar E(x)$ 的確切次數構造的可約性形式 $R(Z,u)$ 在代入 $Z=\mathsf{\Gamma}$ 後仍然不為零；也就是說
\[
 R(\mathsf{\Gamma},u)=\sum\mathsf{ P}_\lambda U_\lambda\ne0\qquad[\text{在 }u\text{ 中恆等}].
\]
因此理想 $\mathfrak r=(\mathsf{P}_1,\mathsf{P}_2,\ldots)$ 只被 $\Omega$ 中有限多個素理想整除。所以，如果 $\mathfrak p$ 不同於這些有限多個素理想，那麼 $R(\mathsf{\Gamma},u)$ 模 $\mathfrak p$ 不恆等消失；又因為模 $\mathfrak p$ 的剩餘類仍構成一個域，便推出 $\bar E(x)$ 模 $\mathfrak p$ 不可約，更確切地說是絕對不可約。對 $\bar E(x,y)$ 也是如此，因此模 $\mathfrak p$ 時也不會發生次數降低。

如果 $\bar E(x)$ 的係數是代數分數，那麼還必須要求 $\mathfrak p$ 不屬於整除共同分母 $\Lambda$ 的有限多個素理想；在其餘所有素理想模下，可以用 $\Lambda\bar E(x)$ 代替 $E(x)$，從而滿足上述前提。定理得證。

\begin{flushleft}
Erlangen, August 1921.
\end{flushleft}
\begin{center}
（1921 年 8 月 28 日收到。）
\end{center}

\clearpage

%
\clearpage
%

%
%
%

\pdfbookmark[0]{21. 形式變分法與微分不變量}{paper-21}
\section*{21. 形式變分法與微分不變量}
\begin{center}
《數學科學百科全書》III, 3 (1922), pp. 68--71（載 R. Weitzenböck, \emph{Differentialinvarianten}）
\end{center}

\setcounter{footnote}{148}
\subsection*{28. 形式變分法與微分不變量。\footnote{這一編號出自 E. Noether。}}

所有微分不變量的生成以及主要定理的推出，最統一地可以用變分法的形式方法來完成。這個原則可追溯到 Riemann\textsuperscript{72)}，並且主要由 Lipschitz\textsuperscript{73)} 得到了形式論意義上的展開。

從屬於一個齊次形式 $f(dx)$ 的變分問題出發--其中假定
\[
\left|\frac{\partial^2 f}{\partial dx_i\,\partial dx_k}\right|\ne0
\]
--即
\begin{equation}
 \delta\int_{t_0}^{t_1} f(x')\,dt=0,
 \qquad \left(x_i'=\frac{dx_i}{dt}\right),
\tag{140}
\end{equation}
經過分部積分，人們被引向如下不變方程組：
\[
 2\psi_i=\frac{d}{dt}\frac{\partial f}{\partial x_i'}-\frac{\partial f}{\partial x_i}=0,
\]
其中拉格朗日表達式 $\psi_i$ 對於 $dx$ 是逆變的。從形式上看，若把 $\psi_i$ 通過與分部積分相應的、無積分的形式恆等式來定義（拉格朗日中心方程\footnote{對於特殊的 $f$，這一名稱出自 Heun；參見他的論文 IV, 1 11, p. 447。}），這一點最容易看出：
\[
 \delta f-df_\delta=-2\sum\psi_i(d,d)\,\delta x_i,
 \qquad
 \left(f_\delta=\sum\frac{\partial f}{\partial dx_i}\,\delta x_i\right),
\]
這裡 $f$ 和 $df_\delta$ 都是不變量，因此右端也是不變量。特別地，對於一個二次形式得到
\[
 \delta\sum g_{ik}dx_i dx_k-2d\sum g_{ik}dx_i\delta x_k
 =-2\sum\psi_\mu(d,d)\,\delta x_\mu.
\]
在這裡用協變的 $(d+\lambda\delta)$ 替代 $d$，並比較 $\lambda$ 的各次冪，就得到相應的恆等式組：
\begin{equation}
\left\{
\begin{aligned}
D\sum g_{ik}dx_i dx_k-2d\sum g_{ik}dx_iDx_k
   &=-2\sum\psi_\mu(d,d)\,Dx_\mu,\\
D\sum g_{ik}dx_i\delta x_k-\delta\sum g_{ik}dx_iDx_k
   -d\sum g_{ik}\delta x_iDx_k
   &=-2\sum\psi_\mu(d,\delta)\,Dx_\mu,\\
D\sum g_{ik}\delta x_i\delta x_k-2\delta\sum g_{ik}\delta x_iDx_k
   &=-2\sum\psi_\mu(\delta,\delta)\,Dx_\mu.
\end{aligned}
\right.
\tag{141}
\end{equation}
由此可知 $\psi_\mu(d,\delta)$，特別也包括 $\psi_\mu(d,d)$ 與 $\psi_\mu(\delta,\delta)$，都是逆變向量。由 (141) 可算出
\[
 \psi_\mu(d,\delta)=\sum g_{i\mu}\,d\delta x_i+
 \sum\left[\begin{smallmatrix}ik\\ \mu\end{smallmatrix}\right]dx_i\delta x_k,
\]
並由此得到協變向量
\begin{equation}
 p^\sigma(d,\delta)=\sum g^{\sigma\mu}\psi_\mu
 =d\delta x_\sigma+
 \sum\left\{\begin{smallmatrix}ik\\ \sigma\end{smallmatrix}\right\}dx_i\delta x_k.
\tag{142}
\end{equation}
令 $p^\sigma(d,d)=0$，顯然就給出測地線方程。

知道協變向量 $p^\sigma(d,\delta)$，就直接導向協變導數（參見第 19 號）。例如，若 $h(d,\delta)$ 是兩列 $dx$ 與 $\delta x$ 的形式，則表達式
\begin{equation}
 h^{(1)}(dx,\delta x)=\delta h-
 \sum\frac{\partial h}{\partial dx_\sigma}p^\sigma(d,\delta)-
 \sum\frac{\partial h}{\partial \delta x_\sigma}p^\sigma(\delta,\delta)
\tag{143}
\end{equation}
作為不變量之差本身仍是不變量，並且其構造方式使二階微分相互抵消。因此 $h^{(1)}(dx,\delta x)$ 表示一個只含一階微分的形式，即 $h$ 的協變導數。若 $h$ 含有更多列 $d^{(1)}x,\ldots,d^{(\lambda)}x$，相應的說法同樣成立。\footnote{Lipschitz, J. f. Math. 72 (1870), p. 17。}

Riemann 與 Lipschitz 對曲率形式的構造也基於同一原則。人們從 $\delta^2 f$ 過渡到“第二變分的正規形式”
\begin{equation}
 \Omega=\delta^2\sum g_{ik}dx_i dx_k
 -2d\delta\sum g_{ik}dx_i\delta x_k
 +d^2\sum g_{ik}\delta x_i\delta x_k,
\tag{144}
\end{equation}
它作為若干不變量的合成，本身又是不變量；並且在這裡，作為對拉格朗日中心方程的推廣，三階微分已經被消去。二階微分的消去，類比於協變求導，通過形成差
\begin{equation}
 K=\Omega-2\sum\{p^\sigma(d,\delta)\psi_\sigma(d,\delta)-p^\sigma(\delta,\delta)\psi_\sigma(d,d)\},
\tag{145}
\end{equation}
來完成。它也可寫成\footnote{Lipschitz, J. f. Math. 82 (1876), § 1。}
\begin{equation}
\begin{aligned}
 \frac12 K={}&d\sum\psi_\sigma(d,\delta)\delta x_\sigma
 -\delta\sum\psi_\sigma(d,d)\delta x_\sigma\\
 &-\sum\{p^\sigma(d,\delta)\psi_\sigma(d,\delta)-p^\sigma(\delta,\delta)\psi_\sigma(d,d)\}.
\end{aligned}
\tag{146}
\end{equation}
$K$ 的形成法則也可以這樣表述：在 $\Omega$ 中，通過令 $p(d,\delta)$ 與 $p(\delta,\delta)$--相應地，即 (141) 的右端--為零來消去二階微分。這就是 Riemann 對曲率形式的定義（參見第 21 號）；但必須明確指出，正如 (146) 特別清楚地顯示的那樣，只有在完成微分之後才可實行這種置零。Lipschitz 的形式構造方式避開了這一困難。相應地，協變導數 $h^{(1)}$ (143) 也可以這樣定義：在 $\delta h$ 中，通過令 $p(d,\delta),\ldots$ 為零來消去二階微分。

由曲率形式通過逐次形成協變導數，可得到一個形式序列 $[\Phi_4,\Phi_5,\ldots$（Christoffel 處），按 Christoffel 與 Ricci 的說法，在添入 $f$ 之後，這些形式的射影不變量窮盡了二次形式的全部微分不變量；而且它們在不再需要其他可積性條件的情況下，對於線性變換的等價性，表達了兩個二次微分形式的等價性（參見第 22 號）。這一約化定理的內在理由在於由變分問題 (140) 產生的 Riemann 正規座標的存在（參見第 20 號）：這些座標把通過一點的測地線變為直線，因此在 $x$ 的任意變換下作線性變換。於是所有不變量的形成以及等價性問題，都被歸結為 $f$ 按正規座標展開時各個齊次成分的相應問題；並且可以證明，這些成分可由 $f,\Phi_4,\Phi_5,\ldots$ 線性組合而成。對於二次形式，Vermeil\textsuperscript{92)} 已在 E. Noether\textsuperscript{93)} 的一則札記之後詳細完成了這一點。按照這則札記，當以任意微分表達式為基礎時，定理和方法仍然保持有效（參見第 22、23 號）；於是，與純粹消去論的方法相比，形式變分法的方法的影響範圍便顯現出來。\footnote{關於這些說明以及進一步的幾何解釋，可參見“文獻”中列出的 F. Klein 的討論班講演。}

Levi-Civita\textsuperscript{134)} 曾把置 $p(d,\delta)=0$ 從幾何上解釋為“向量的平行移動”（也參見 Hessenberg\textsuperscript{133)}）。這個由 Weyl 以公理方式把握的概念，本質上侷限於二次形式，但它容許另一種推廣。當群被擴大時（即把 $g_{ik}$ 乘以任意函數），只要除二次形式以外還以一個線性形式為基礎，這一概念仍然保持有效。此時 $p(d,\delta)$ 被唯一確定，可以進行與上述相應的不變量構造，並且可以定義測地座標；但 $p(d,d)=0$ 此時不再來自一個變分問題！\footnote{H. Weyl, \emph{Reine Infinitesimalgeometrie}, Math. Ztschr. 2 (1918), p. 384--411，以及 \emph{Raum, Zeit, Materie}（第 3、4 版）。} 關於約化定理與等價性的問題，在這裡一般地還沒有被處理；只有對於最低階不變量，R. Weitzenböck 給出了一個約化定理。\footnote{Wien. Ber. 129 (1920), p. 683 與 697。}

最後還要指出一般的“不變變分問題”，也就是那些使（單重或多重）積分 $J$ 對 Lie 意義下某個群保持不變的問題。Hamel\footnote{Math. Ann. 59 (1904), p. 416--434; Ztschr. Math. Phys. 50 (1904), p. 1--54。}、Herglotz\footnote{Ann. d. Phys. (4) 36 (1911), p. 511，尤其 § 9。}、Lorentz\footnote{Verslag der Amsterd. Akad., 1916 年 4 月、9 月、10 月，1917 年 5 月。}、Fokker\footnote{同刊，1917 年 1 月。}、Weyl\footnote{Ann. d. Phys. (4) 54 (1917), p. 117, § 2。} 以及 Klein\footnote{Gött. Nachr. 1917, p. 469--482；同刊 1918, p. 171--189。} 的工作研究了一些特殊的、物理上有趣的情形。E. Noether\footnote{Gött. Nachr. 1918, p. 235--257。} 的原則性表述表明：若 $J$ 對一個 $G_\varrho$（具有 $\varrho$ 個本質參數的有限群）不變，則對應有 $\varrho$ 個線性無關的散度；若它對一個包含 $\varrho$ 個任意函數直到 $\sigma$ 階導數的無限群不變，則對應有 $\varrho$ 個拉格朗日表達式及其直到 $\sigma$ 階導數之間的依賴關係。在兩種情形下，逆命題也成立。由於拉格朗日表達式成為群的（相對）不變量，人們同時就得到一個生成不變量的過程。

\begin{center}
（1921 年 3 月完成。）
\end{center}


\clearpage
%
\clearpage
%

%
%
%

\setcounter{footnote}{0}
\pdfbookmark[0]{22. K. Hentzelt 的改寫稿： 關於多項式理想與結式的理論}{paper-22}
\section*{22. K. Hentzelt 的改寫稿：\\ 關於多項式理想與結式的理論}
\begin{center}
Math. Ann. 88 (1923), pp. 53--79
\end{center}

以下討論的是\emph{消去論的一般問題，即確定一個多項式理想中所有多項式的公共零點}；等價地說，是確定構成該理想的任意有限個基多項式的零點。同時，文中還會給出所出現重數的概念性解釋。多項式的係數可以取自任意一個域；於是零點屬於由該域導出的代數閉域。此外，不失一般性，可以假定理想已經是變換後的理想（§~3）。這時，核心結果如下：

\emph{對每個多項式理想 $\mideal$，都可以配屬一個結式形式：}
\[
 R_{\mideal}=R^{(1)}(x_1\ldots x_n)\cdot\ldots\cdot
 R^{(i)}(x_i\ldots x_n)\cdot\ldots\cdot R^{(n)}(x_n)\zero{\mideal},
\]
\emph{使得 $R_{\mideal}$ 在 $\mideal$ 的所有零點且僅在這些零點處消失。若 $\nideal$ 是 $\mideal$ 的一個約數，並且結式形式 $R_{\nideal}$ 與 $R_{\mideal}$ 相同，則理想 $\nideal$ 與 $\mideal$ 也相同。}\footnote{Kurt Hentzelt 自 1914 年 10 月起在 Dixmuiden 前線失蹤，必須列入陣亡者之中。這裡給出的論文，是對他的博士論文《Zur Theorie der Formenmoduln und Resultanten》中最本質部分的完全自由改寫；他於 1914 年夏在 Erlangen 的 E. Fischer 指導下以該論文取得博士學位。這篇完全基於他自己的思想寫成的論文結構無缺；但是輔助命題接着輔助命題，所有概念都用帶有四個、五個指標的公式描述，正文幾乎完全缺失，因此理解起來極其困難。他本人已經來不及進行原計劃的改寫。我在這裡以純粹概念性的形式重述這項工作，由此大大簡化了那些基本思想處處可追溯到 Hentzelt 的證明，並且希望這項工作的美能夠顯現出來。--論文中那些在給定基的情況下涉及用有限多個步驟構造所出現函數的問題的部分，將留待另行發表。（E. N.）}

主定理的第二部分說明，結式形式以特徵性的重數給出零點。這個第二部分類似於代數數域中兩個理想的事實：若其中一個是另一個的約數，那麼它們當且僅當其\emph{範數}相同時相同；這兩個定理的內在理由也是相同的。

事實上，可以把 $\mideal$ 看作線性形式模 $\mathfrak M_{i-1}$：即把每個多項式看成關於 $x_1\ldots x_{i-1}$ 的冪積的線性形式，並把 $x_{i+1}\ldots x_n$ 添入係數域。於是結式因子 $R^{(i)}$ 等於相應基本模（§~1）$\Gmod_{i-1}$ 相對於 $\mathfrak M_{i-1}$ 的\emph{範數}；由此也立即得到\emph{重數--即 $R^{(i)}$ 的一個因子的次數與指數之積--通過線性無關剩餘類個數來解釋}\footnote{一旦考慮理想分解為準素理想，這些後一結果便顯示出其真正意義；與一個準素理想相應的重數可定義為其補元按該理想所得的線性無關剩餘類個數。關於這一點將在別處討論。（E. N.）}。這個事實迄今只在一個特殊情形下為人所知，即結式可對不定係數加以定義的情形\footnote{參見 Macaulay, \emph{The algebraic theory of modular systems}, Cambridge Tracts 19 (Cambridge University Press, 1916), No. 67；又參見 Lasker, Zur Theorie der Moduln und Ideale, Math. Ann. 60 (1905), pp. 20--116, p. 98。上述情形中結式與結式形式一致，是 Hentzelt 在注 1 所提及的尚未發表的定理中證明的。（E. N.）}。

為推出這些結果，§~1 和 §~2 給出關於線性形式模的若干定理，其係數是有理整數或一個不定元的多項式。§~4 的同構定理建立了這些模與 §~3 中引入的多項式理想之間的聯繫。§~5 給出上文提到的主定理第二部分；§~7 給出關於零點的定理，而 §~6 先行給出一般消去論。那裡還在引入新不定元時證明了解析式分解為這些不定元的線性形式；由此，對這裡給出的消去法，Kronecker 的一個斷言獲得了無瑕疵的證明\footnote{König, \emph{Algebraische Größen} (Leipzig 1903, Teubner), V, §~4 中為 Kronecker 消去論所作的證明嘗試，眾所周知並未成功。Macaulay, loc. cit., p. 23 還指出，對於 König 引入 Kronecker 消去論的重數，Hentzelt 主定理的第二部分同樣不成立；那裡還進一步說明，Kronecker 消去論並不對應於準素理想分解。（E. N.）}。

\section*{§ 1. 線性形式模}

在前兩個段落中，所謂整量 $a,b,c,\ldots$ 指有理整數，或指一個不定元的多項式，其係數取自一個任意的、抽象定義的域 $\mathsf{P}$；所謂線性形式 $a(\xi),b(\xi),\ldots$，指有限多個不定元 $\xi_1\ldots\xi_k$ 中的線性形式，其係數為整量。

\emph{一個由線性形式構成的模 $\Amod$} 定義為滿足如下條件：若 $a(\xi)$ 與 $b(\xi)$ 屬於 $\Amod$，則差 $a(\xi)-b(\xi)$ 也屬於 $\Amod$；若 $a(\xi)$ 屬於 $\Amod$，則 $c\cdot a(\xi)$ 也屬於 $\Amod$，其中 $c$ 是任意整量。

所謂 $\Amod$ 的\emph{秩}，照通常意義，是指 $\Amod$ 中線性無關線性形式的最大個數；所謂 $\Amod$ 的第 $\lambda$ 個\emph{行列式因子} $d_\lambda$，是指 $\Amod$ 中所有 $\lambda$ 階行列式的最大公因子（也就是從 $\Amod$ 中任意 $\lambda$ 個線性形式的係數矩陣可形成的所有 $\lambda$ 階行列式）。若 $\rho$ 是 $\Amod$ 的秩，從而 $d_1,\ldots,d_\varrho$ 是非零的行列式因子，則 $\Amod$ 的\emph{初等因子}定義為 $e_1=d_1$、$e_2=d_2/d_1,\ldots,e_\varrho=d_\varrho/d_{\varrho-1}$，以及 $e_{\varrho+i}=0$。眾所周知，$\Amod$ 是一個有限模，它有一個恰由 $\varrho$ 個線性形式 $a_1(\xi),\ldots,a_\varrho(\xi)$ 組成的模基；$\Amod$ 中的每個線性形式都可以用這些形式以整量為係數線性表示：
\[
\Amod=(a_1(\xi),\ldots,a_\varrho(\xi)).
\]
若 $A$ 表示這些線性形式的係數矩陣，則 $\Amod$ 的行列式因子和初等因子與 $A$ 的相應因子一致。由此首先推出：\emph{以下每個初等因子 $e_\lambda$ 都整除 $e_{\lambda+1}$}。初等因子理論給出的矩陣等式
\[
 P A Q=\begin{pmatrix}
 e_1&&\\
 &\ddots&\\
 &&e_\varrho
 %
 \end{pmatrix}
\]
進一步說明：若通過幺模變換 $\xi=Q(\eta)$ 引入新的不定元 $\eta_1\ldots\eta_k$，則有如下模等式：
\begin{equation}
 \Amod=(e_1\eta_1,e_2\eta_2,\ldots,e_\varrho\eta_\varrho).
\tag{1}
\end{equation}

對下文最重要的概念，是如下的基本模概念\footnote{參見 Steinitz, Rechteckige Systeme und Moduln in algebraischen Zahlkörpern II, Math. Ann. 72 (1912), pp. 297--345, No. 36。不過，Hentzelt 至少看起來是在獨立於 Steinitz 的情況下--儘管他在其他方面也與 Steinitz 有接觸點--在他的較簡單情形中以公式 (4) 的形式得到了這個概念。（E. N.）}：

\medskip
\noindent\textbf{定義 I。} \emph{一個模 $\Gmod$ 稱為基本模，如果它沒有相同秩的真約數}\footnote{這裡“約數”按模的意義理解：若 $\Amod$ 的每個元素都包含在 $\Bmod$ 中，則說 $\Amod$ 可被 $\Bmod$ 整除，記作 $\Amod\zero{\Bmod}$；若 $\Bmod$ 含有不同於 $\Amod$ 的元素，則稱 $\Bmod$ 為真約數。}。

因此，$\Gmod$ 是基本模，當且僅當由
\begin{equation}
 c\cdot g(\xi)\zero{\Gmod};\quad c\ne0 \quad\text{總能推出：}\quad g(\xi)\zero{\Gmod}.
\tag{2}
\end{equation}

任意模與基本模之間的聯繫由如下命題給出：

\medskip
\noindent\textbf{定理 I。} \emph{每個模 $\Amod$ 都可被一個且僅一個與它同秩的基本模 $\Gmod$ 整除；$\Gmod$ 被定義為 $\Amod$ 的同秩最小約數，並表示為}
\begin{equation}
 \Gmod=(\eta_1,\ldots,\eta_\varrho);
\tag{3}
\end{equation}
\emph{$\Gmod$ 將稱為 $\Amod$ 的基本模。}

$\Gmod$ 應當是同秩最小約數這一條件，可表述為：$\Gmod$ 含有且僅含有所有滿足某個關係
\begin{equation}
 b\cdot g(\xi)\zero{\Amod};\quad b\ne0
\tag{4}
\end{equation}
的線性形式 $g(\xi)$。由此首先可知 $\Gmod$ 是一個模；因為若
\[
 b_1g_1(\xi)\zero{\Amod};\quad b_1\ne0;
 \qquad
 b_2g_2(\xi)\zero{\Amod};\quad b_2\ne0,
\]
則也有
\[
 b_1b_2\bigl(g_1(\xi)-g_2(\xi)\bigr)\zero{\Amod};\quad b_1b_2\ne0,
\]
當然也有 $b_1\cdot c g_1(\xi)\zero{\Amod}$。而 $\Gmod$ 也是基本模；因為由
\[
 c\cdot g(\xi)\zero{\Gmod};\quad c\ne0
\]
按 (4) 推出
\[
 b c g(\xi)\zero{\Amod};\quad bc\ne0,
\]
於是再由 (4) 得 $g(\xi)\zero{\Gmod}$。

此外，不存在一個不同於該同秩最小約數 $\Gmod$ 的基本模 $\Gmod_1$ 滿足這些條件。因為 $\Gmod_1$ 必須可被 $\Gmod$ 整除，但它作為基本模沒有同秩真約數，因而與 $\Gmod$ 相同。最後，(3) 右端所表示的模是基本模，因為任何真約數都必須還含有一個新的不定元 $\eta$，從而具有更高秩。於是 (3) 給出了 $\Amod$ 的唯一確定的基本模，\emph{定理 I 的所有部分均得證}。

另一個對下文重要的概念，是 Dedekind 的兩個模之商的概念\footnote{Dedekind 討論的是數模，對它們還定義有乘法；因此 Dedekind 的商與這裡定義的商並不相同。（E. N.）}：

\medskip
\noindent\textbf{定義 II。} \emph{兩個線性形式模的商 $\cc=\Amod/\Bmod$，定義為滿足條件}
\[
 c\cdot\Bmod\zero{\Amod}
\]
\emph{的全體整量 $c$。$\cc$ 構成一個整量理想，因而構成一個主理想。相應地，一個線性形式模 $\Amod$ 對一個整量理想 $\bb$ 的商 $\mathfrak C=\Amod/\bb$，定義為滿足條件}
\[
 c(\xi)\cdot\bb\zero{\Amod}
\]
\emph{的全體線性形式 $c(\xi)$。$\mathfrak C$ 又構成一個線性形式模；並且只要 $\bb$ 不同於零理想，它就是與 $\Amod$ 同秩的模。}

這裡只需指出，按定義顯然每個商都具有模性質；而具有模性質的整量系統就是理想。

商的概念導向基本模與最高初等因子之間的一個聯繫：

\medskip
\noindent\textbf{定理 II。} \emph{在 $\Amod$ 的基本模與最高初等因子之間有如下互反關係}
\begin{equation}
 (e_\varrho)=\Amod/\Gmod;\qquad \Gmod=\Amod/(e_\varrho),
\tag{5}
\end{equation}
\emph{其中 $(e_\varrho)$ 表示由 $e_\varrho$ 生成的主理想。}

因為每個初等因子都整除 $e_\varrho$，所以由 (1) 和 (3) 得
\begin{equation}
 e_\varrho\Gmod\zero{\Amod}\quad \text{並因而}\quad (e_\varrho)\cdot\Gmod\zero{\Amod};
\tag{6}
\end{equation}
於是 $(e_\varrho)$ 可被商 $\Amod/\Gmod$ 整除。反過來也成立；因為由
\[
 b\Gmod\zero{\Amod}
 \quad\text{推出}\quad b\eta_\varrho\zero{\Amod}
 \quad\text{從而}\quad b\equiv 0(e_\varrho),
\]
這就證明了 (5) 的第一個公式\footnote{若置 $\Amod_0=\Amod,\ldots,\Amod_\lambda=(\Amod,\eta_\varrho,\ldots,\eta_{\varrho-\lambda+1})$，其中每個 $\Amod_\lambda$ 的最高初等因子為 $e_{\varrho-\lambda}$，則由同樣推理得到：
\[
 (e_\varrho)=\Amod_0/\Amod_1,\ldots,(e_{\varrho-\lambda})=\Amod_\lambda/\Amod_{\lambda+1},\ldots,(e_1)=\Amod_{\varrho-1}/\Amod_\varrho.
\]
更一般地，設 $\zeta_\lambda=b_{\lambda1}\eta_1+\cdots+b_{\lambda\lambda}\eta_\lambda$，設 $(b_{\lambda\lambda},e_\lambda)=1$，並令
\[
 \Bmod_0=\Amod,\ldots,\Bmod_\lambda=(\Amod,\zeta_\varrho,\ldots,\zeta_{\varrho-\lambda+1}),
\]
則 $\Bmod_\lambda$ 也具有最高初等因子 $e_{\rho-\lambda}$，因而有
\[
 (e_{\rho-\lambda})=\Bmod_\lambda/\Bmod_{\lambda+1}.
\]}。(6) 還說明 $\Gmod$ 可被商 $\Amod/(e_\rho)$ 整除；這個商是 $\Amod$ 的一個同秩約數，所以按 $\Amod$ 的基本模定義，它反過來也可被 $\Gmod$ 整除。由此 (5) 的第二個公式也得證。

設 $d$ 是 $\Amod$ 中任意一個不恆等為零的 $\rho$ 階行列式。則 $d$ 可被第 $\rho$ 個行列式因子 $d_\rho$ 整除，因而也可被 $e_\rho$ 整除；於是 $d\Gmod\zero{\Amod}$，由上面的推理可得 $\Gmod=\Amod/(d)$。但對後文來說重要的是，這後一種關於 $\Gmod$ 的商表示也可以不回到初等因子而得到證明，因此具有更一般的有效性：

\medskip
\noindent\textbf{定理 III。} \emph{若把整量理解為多個不定元的多項式}\footnote{實際上只用到如下事實：整量在加法、減法和乘法下封閉，並服從通常的運算法則，因而構成一個環；還用到在這個環中，一個乘積只有在某個因子為零時才會為零，因此該環可通過商構造（添入元素對）擴張為一個域。相反，並未使用模的基表示；所以 (7) 例如也適用於所有代數整數作為係數的範圍。（E. N.）}, \emph{則仍有商表示}
\begin{equation}
 \Gmod=\Amod/(d),
\tag{7}
\end{equation}
\emph{其中 $d$ 是 $\Amod$ 中任意一個不恆等為零的 $\rho$ 階行列式。}

首先顯然，在這一擴張下，秩、基本模和模商這些概念--只是現在模商不再一定是主理想--仍然保持；定理 I 除了基表示以外也保持。

現在設
\[
 a_1(\xi)=a_{11}\xi_1+\cdots+a_{1k}\xi_k;\ \ldots;\
 a_\varrho(\xi)=a_{\varrho1}\xi_1+\cdots+a_{\varrho k}\xi_k
\]
是 $\Amod$ 中的 $\rho$ 個線性無關線性形式；它們一般並不構成模基。把不定元 $\xi$ 這樣編號，使得
\[
 d=|a_{ij}|\ne0;\qquad i,j=1,2,\ldots,\varrho.
\]
由此可知，$\varrho$ 個特殊形狀的線性形式屬於 $\Amod$：
\begin{equation}
 d\xi_\lambda+b_{\lambda,\varrho+1}\xi_{\varrho+1}+\cdots+b_{\lambda,k}\xi_k\zero{\Amod}
 \quad(\lambda=1,2,\ldots,\varrho).
\tag{8}
\end{equation}
但是 $\Amod$ 不可能含有一個只依賴於 $\xi_{\varrho+1}\ldots\xi_k$ 的線性形式
$c_{\varrho+1}\xi_{\varrho+1}+\cdots+c_k\xi_k$；否則至少會有一個 $(\varrho+1)$ 階行列式 $d\cdot c_\sigma$ 非零。

現在，商 $\Amod/(d)$ 作為 $\Amod$ 的同秩約數，可被 $\Gmod$ 整除；反過來也成立。因為對每個 $g(\xi)\zero{\Gmod}$，按定義有
\[
 c\cdot g(\xi)\zero{\Amod};\quad c\ne0
 \quad\text{並因而}\quad d\cdot c\cdot g(\xi)\zero{\Amod}.
\]
但由 (8) 可把
\[
 d\cdot g(\xi)\equiv g_{\varrho+1}\xi_{\varrho+1}+\cdots+g_k\xi_k\,(\Amod),
\]
於是得到
\[
 c g_{\varrho+1}\xi_{\varrho+1}+\cdots+c g_k\xi_k\zero{\Amod}.
\]
由剛才的說明推出 $c\cdot g_{\varrho+1}=0,\ldots,c\cdot g_k=0$；又因 $c\ne0$，也有 $g_{\varrho+1}=0,\ldots,g_k=0$，從而 $d\cdot g(\xi)\zero{\Amod}$。這證明了 (7)。

還要指出，$d\cdot g(\xi)\zero{\Amod}$ 也可直接由兩個模
$(a_1(\xi),\ldots,a_\varrho(\xi))$ 與 $(g(\xi),a_1(\xi),\ldots,a_\varrho(\xi))$ 具有同一秩 $\varrho$ 推出；也就是說，若置 $g(\xi)=c_1\xi_1+\cdots+c_k\xi_k$，則 $(\varrho+1)$ 階行列式
\[
\begin{vmatrix}
 a_1(\xi)&a_{11}&\cdots&a_{1\varrho}\\
 \vdots&\vdots&&\vdots\\
 a_\varrho(\xi)&a_{\varrho1}&\cdots&a_{\varrho\varrho}\\
 g(\xi)&c_1&\cdots&c_\varrho
\end{vmatrix}
\]
為零。上面給出的證明將在 §~4 中接在公式 (30) 之後再次出現。

\section*{§ 2. 範數與模的分解定理}

現在我們重新回到線性形式模；其係數仍取為有理整數，或取為一個不定元的多項式，而這些多項式的係數來自 $P$。

\medskip
\noindent\textbf{定義 III。} \emph{像通常一樣，把所有模 $\Amod$ 同餘的 $\xi$ 中線性形式合併為一個剩餘類；於是記號 $\Gmod\mid\Amod$ 表示 $\Gmod$ 關於 $\Amod$ 的剩餘類系統，也就是關於 $\Amod$ 的、由 $\Gmod$ 的元素組成的全部剩餘類系統。$\Gmod$ 關於 $\Amod$ 的範數--記作 $\Norm(\Gmod\mid\Amod)$--定義為從 $\Gmod$ 到 $\Amod$ 的過渡代換的行列式；也就是說，它是把 $\Amod$ 的任一線性無關模基用 $\Amod$ 的基本模 $\Gmod$ 的相應模基表示出來的那個代換的行列式。}

由 (1)、(3)，以及定理 II 後的附註，得到
\begin{equation}
 \Norm(\Gmod\mid\Amod)=e_1e_2\cdots e_\varrho=d_\varrho;
 \qquad
 \Gmod=\Amod/\bigl(\Norm(\Gmod\mid\Amod)\bigr).
\tag{9}
\end{equation}

由 (1)、(3) 和 (9) 可按熟知方式推出：在有理整數的情形中，$\Norm(\Gmod\mid\Amod)$ 表示 $\Gmod$ 關於 $\Amod$ 的剩餘類個數；而在一個不定元多項式的情形中--此時剩餘類的個數一般為無限--$\Norm(\Gmod\mid\Amod)$ 的次數等於\emph{關於 $\mathsf{P}$ 線性無關的剩餘類個數}。若 $\Bmod$ 是 $\Amod$ 的一個同秩約數，則 $\Bmod$ 除了含有任一這種剩餘類的一個代表以外，也同時含有該剩餘類的全部元素；因而 $\Bmod$ 是由 $\Amod$ 添入有限多個剩餘類，或它們的線性組合而得到的。於是立即得到：

\medskip
\noindent\textbf{定理 IV。} \emph{若 $\Bmod$ 是 $\Amod$ 的同秩約數，因而二者的基本模相同，並且還有 $\Norm(\Gmod\mid\Bmod)=\Norm(\Gmod\mid\Amod)$，則 $\Bmod=\Amod$。}

關於範數分解與模分解之間的聯繫，有如下結果。

\medskip
\noindent\textbf{定理 V。} \emph{設}
\begin{equation}
 \Norm(\Gmod\mid\Amod)=r\cdot s,\qquad (r,s)=1;
\tag{10}
\end{equation}
\emph{則存在唯一的模 $\Rmod$，也存在唯一的模 $\Smod$，二者都是 $\Amod$ 的同秩約數，並且滿足}
\[
 r=\Norm(\Gmod\ \Rmod),\qquad s=\Norm(\Gmod\mid\Smod).
\]
\emph{$\Rmod$ 與 $\Smod$ 定義為}
\[
 \Rmod=\Amod/(s);\qquad \Smod=\Amod/(r),
\]
\emph{並且 $\Amod$ 等於 $\Rmod$ 與 $\Smod$ 的最小公倍數。}\footnote{最小公倍數 $[\Rmod,\Smod]$ 按模的意義理解：即同時可被 $\Rmod$ 和 $\Smod$ 整除的線性形式全體。相應地，最大公約數 $(\Rmod,\Smod)$ 也按模的意義理解：即可以表示為一個 $\Rmod$ 中元素與一個 $\Smod$ 中元素之和的線性形式全體。}
\emph{若置 $r_\varrho=(r,e_\varrho)$，$s_\varrho=(s,e_\varrho)$，則有互反關係}
\begin{equation}
 (s_\varrho)=\Amod/(\Rmod);\quad \Rmod=\Amod/(s_\varrho);
 \qquad
 (r_\varrho)=\Amod/\Smod;\quad \Smod=\Amod/(r_\varrho).
\tag{11}
\end{equation}

事實上，若一般地置 $r_i=(r,e_i)$、$s_i=(s,e_i)$，則由 (9) 與 (10) 得到
\[
 (r_i,s_i)=1;\qquad e_i=r_is_i;\qquad r=r_1\cdots r_\varrho;\qquad s=s_1\cdots s_\varrho,
\]
並且每個 $r_i$ 分別整除後一個 $r_{i+1}$，每個 $s_i$ 分別整除後一個 $s_{i+1}$。於是若置
\[
 \Rmod=(r_1\eta_1,\ldots,r_\varrho\eta_\varrho);\qquad
 \Smod=(s_1\eta_1,\ldots,s_\varrho\eta_\varrho),
\]
則 $\Rmod$ 與 $\Smod$ 都是 $\Amod$ 的同秩約數；由於 $r_i$ 與 $s_i$ 互素，$\Amod$ 成為 $\Rmod$ 與 $\Smod$ 的最小公倍數，並且還得到
\[
 \Norm(\Gmod\mid\Rmod)=r_1\cdots r_\varrho=r;\qquad
 \Norm(\Gmod\mid\Smod)=s_1\cdots s_\varrho=s.
\]
又有
\[
 s_\varrho\Rmod\zero{\Amod};\qquad r_\varrho\Smod\zero{\Amod};
\]
而由 $c\Rmod\zero{\Amod}$ 可推出 $cr_\varrho\eta_\varrho\zero{\Amod}$，從而 $c\zero{(s_\varrho)}$。這就證明了 $(s_\varrho)=\Amod/\Rmod$，同理證明 $(r_\varrho)=\Amod/\Smod$。再由
\[
 b(\eta)=b_1\eta_1+\cdots+b_\varrho\eta_\varrho\zero{\Amod/(s_\varrho)}
\]
並利用所有 $r_i$ 都與 $s_\varrho$ 互素，可推出 $b_i\zero{(r_i)}$；於是 $\Rmod=\Amod/(s_\varrho)$，同樣 $\Smod=\Amod/(r_\varrho)$。這就證明了互反關係 (11)；在特殊情形 $r=1$ 時，它們化為 (5)。

還須證明 $\Rmod$ 與 $\Smod$ 的\emph{唯一確定性}。設 $\overline{\Rmod}$ 與 $\overline{\Smod}$ 是滿足定理 V 條件的模，並設 $\bar r_i$、$\bar s_i$ 為它們的初等因子。由於 $\bar r_i$ 與 $\bar s_i$ 都是 $e_i$ 的約數，而
\[
 r=\bar r_1\cdots\bar r_\varrho,
 \qquad
 s=\bar s_1\cdots\bar s_\varrho,
 \qquad
 (\bar r_i,\bar s_i)=1,
\]
於是也有
\[
 e_i=\bar r_i\bar s_i,
 \qquad
 \bar r_i=(r,e_i)=r_i,
 \qquad
 \bar s_i=(s,e_i)=s_i.
\]
因此 $\overline{\Rmod}$ 與 $\overline{\Smod}$ 在初等因子和基本模上分別同 $\Rmod$ 與 $\Smod$ 一致。於是若通過一個幺模代換 $\eta=Q(\zeta)$ 引入新的不定元 $\xi$，便有
\[
 \overline{\Rmod}=(r_1\zeta_1,\ldots,r_\rho\zeta_\rho);
\]
\[
 \Amod=\bigl(r_1s_1(q_{11}\zeta_1+\cdots+q_{1\rho}\zeta_\rho),\ldots,
 r_\rho s_\rho(q_{\rho1}\zeta_1+\cdots+q_{\rho\rho}\zeta_\rho)\bigr).
\]
由 $\Amod\zero{\overline{\Rmod}}$ 因而得到
\[
 r_i s_i(q_{i1}\zeta_1+\cdots+q_{i\varrho}\zeta_\varrho)\zero{\overline{\Rmod}};
\]
故 $r_is_iq_{ij}\zero{(r_j)}$。由 $(r,s)=1$ 又得到 $r_iq_{ij}\zero{(r_j)}$，也就是 $\Rmod\zero{\overline{\Rmod}}$。但因 $\Norm(\Gmod\mid\Rmod)=\Norm(\Gmod\mid\overline{\Rmod})$，由定理 IV 得 $\Rmod=\overline{\Rmod}$；同樣得到 $\Smod=\overline{\Smod}$。定理 V 的全部部分由此證畢。

\section*{§ 3. 多項式理想}

設 $\overline{\mideal}$ 是 $n$ 個不定元 $y_1\ldots y_n$ 中的一個\emph{多項式理想}，其係數來自一個任意的、抽象定義的域 $\mathsf{P}$；也就是說，$\overline{\mideal}$ 在含有 $\Phi_1(y)$ 與 $\Phi_2(y)$ 的同時，也含有差 $\Phi_1(y)-\Phi_2(y)$，並且在含有 $\Phi(y)$ 的同時，也含有 $A(y)\cdot\Phi(y)$，其中 $A(y)$ 表示任意一個係數在 $\mathsf{P}$ 中的多項式\footnote{這裡用概念上自然得到的名稱“理想”，而不用早先通常使用的“模”或“形式模”，已經是必要的；這樣才能把這些對象同線性形式模區別開來。（E. N.）}。

令 $y$ 經受一個以不定元為係數的變換：
\begin{equation}
 y=U(x);\quad \text{例如}\quad
 \begin{cases}
 y_1=x_1,\\
 y_2=u_{21}x_1+x_2,\\
 \vdots\\
 y_n=u_{n1}x_1+\cdots+u_{n,n-1}x_{n-1}+x_n,
 \end{cases}
\tag{12}
\end{equation}
並把不定元 $u_{\mu\nu}$ 添入這些多項式的係數域，使該係數域變為一個域 $\mathsf{P}(u)$。通過這種添加，$\overline{\mideal}$ 變為 $\overline{\mideal}_{(u)}$；因此，除 $\overline{\mideal}$ 中的多項式 $\Phi_\lambda(y)$ 外，$\overline{\mideal}_{(u)}$ 還含有有限多個 $\Phi_\lambda$ 的所有線性組合 $\sum \alpha_\lambda(u)\Phi_\lambda(y)$，其中 $\alpha_\lambda(u)$ 表示 $u_{\mu\nu}$ 的有理函數，係數來自 $\mathsf{P}$。通過乘以一個適當的 $u$ 中多項式，可以不失一般性地假定 $\alpha_\lambda(u)$ 是 $u$ 中冪積。於是有
\begin{equation}
 \sum \alpha_\lambda(u)\Phi_\lambda(y)\zero{\mideal_{(u)}};
 \qquad
 \Phi_\lambda(y)\zero{\overline{\mideal}}\quad\text{並且反過來也成立}.
\tag{13}
\end{equation}
藉助 (12)，$\overline{\mideal}_{(u)}$ 變為 $x_1\ldots x_n$ 中、係數在 $\mathsf{P}(u)$ 中的一個多項式理想 $\mideal$：
\begin{equation}
 \overline{\mideal}_{(u)}=\mideal\quad \text{通過}\quad y=U(x).
\tag{14}
\end{equation}
關係 (14) 與 (13) 合起來說明如下：
\begin{equation}
 \text{由 } F(x)\zero{\mideal};\quad
 F(x)=\Phi(y)=\sum \alpha_\lambda(u)\Phi_\lambda(y)=\sum \alpha_\lambda(u)F_\lambda(x)
 \text{ 推出 } F_\lambda(x)\zero{\mideal};
\tag{15}
\end{equation}
而反過來，由 (14) 與 (15) 又可重新得到 (13)。

\medskip
\noindent\textbf{定義 IV。} \emph{若 $x_1\ldots x_n$ 中、係數在 $\mathsf{P}(u)$ 中的多項式理想 $\mideal$ 滿足關係 (15)，因而通過 $y=U(x)$ 來自 $\overline{\mideal}_{(u)}$，則稱這樣的理想為變換理想。}

變換理想的特徵在於存在正則多項式；\emph{若一個 $r$ 次多項式含有係數非零的項 $x_i^r$，則稱它關於 $x_i$ 正則}。可以看出，多項式 $F_\lambda(x)$ 關於 $x_1$ 是正則的。下面的關鍵，是把這些變換理想看作某些 $x$ 的冪積中的線性形式模。為此必須確定基本理想的概念；稍後它會過渡為基本模。我們先給出一個從現在起始終保持的記號：

\medskip
\noindent\textbf{記號。} \emph{凡寫作 $F^{(i)},f^{(i)},a^{(i)},\ldots$ 的，都始終理解為 $x$ 中不含 $x_1,\ldots,x_{i-1}$ 的多項式。}

\medskip
\noindent\textbf{定義 V。} \emph{對每個理想 $\mideal$，定義 $n$ 個基本理想 $\gideal_0,\gideal_1,\ldots,\gideal_{n-1}$ 如下：第 $(i-1)$ 級的基本理想 $\gideal_{i-1}$ 含有且只含有所有滿足下述條件的多項式 $G(x)$：存在一個多項式 $b^{(i)}$--一般隨 $G(x)$ 而變--使得}
\begin{equation}
 b^{(i)}G(x)\zero{\mideal};\qquad b^{(i)}\ne0.
\tag{16}
\end{equation}

$\gideal$ 確實是理想，這一點與公式 (4) 後證明 $\Gmod$ 的模性質完全對應，其中 (16) 起着相應的作用。理想 $\gideal_i$ 可被 $\gideal_{i-1}$ 整除，因為每個多項式 $b^{(i+1)}$ 同時也是一個 $b^{(i)}$。

對於基本理想，有：

\medskip
\noindent\textbf{定理 VI。} \emph{變換理想的基本理想仍是變換理想}\footnote{定理 VI 只在 §~7 中使用。}。

證明依賴於 Dedekind--Mertens 的一個熟知定理\footnote{Dedekind: Über einen arithmetischen Satz von Gauß, Mitt. d. deutsch. math. Ges. zu Prag 1892. F. Mertens: Über einen algebraischen Satz, Ber. d. Ak. d. Wissensch. Wien 101 (1892), pp. 1560--1566. -- Kronecker 對 Gauss 定理到帶不定係數多項式的推廣，是這個定理的直接推論。}：設 $a$ 與 $b$ 為不定元，並置
\[
 \sum a_{i_1\ldots i_s}z_1^{i_1}\cdots z_s^{i_s}\cdot
 \sum b_{j_1\ldots j_s}z_1^{j_1}\cdots z_s^{j_s}
 =
 \sum c(a,b)_{k_1\ldots k_s}z_1^{k_1}\ldots z_s^{k_s},
\]
\[
 \sum i\leqq \nu,
 \qquad
 \sum j\leqq \mu,
 \qquad
 \sum k\leqq \nu+\mu.
\]
若 $\Amod,\Bmod,\mathfrak C$ 分別表示由 $a,b,c$ 通過有理整數導出的模，則存在一個指數 $q$，使恆等式成立：
\begin{equation}
 \Amod^q\Bmod=\Amod^{q-1}\mathfrak C.
\tag{17}
\end{equation}
這裡 $\Amod^q$ 由 $a$ 的 $q$ 維冪積及其整數線性組合組成。

現在為證明定理 VI，置
\begin{equation}
\begin{cases}
 G(x)=\Gamma(y)=\sum \alpha_\lambda(u)\Gamma_\lambda(y)=\sum \alpha_\lambda(u)G_\lambda(x),\\
 b^{(i)}(x)=\varphi(y)=\sum \beta_\mu(u)\varphi_\mu(y),
\end{cases}
\tag{18}
\end{equation}
其中 $\alpha_\lambda(u)$ 與 $\beta_\mu(u)$ 表示 $u$ 的冪積。於是由 (16)、(14) 和 (18) 得
\begin{equation}
 \sum\beta_\mu(u)\varphi_\mu(y)\cdot
 \sum\alpha_\lambda(u)\Gamma_\lambda(y)\zero{\overline{\mideal}_{(u)}};
 \qquad
 \sum\beta_\mu(u)\varphi_\mu(y)\ne0.
\tag{19}
\end{equation}
左端是兩個 $u$ 中多項式的乘積，其係數分別為多項式 $\varphi_\mu(y)$ 與 $\Gamma_\lambda(y)$；按 (19) 右端，所得 $u$ 中乘積多項式的係數是 $\overline{\mideal}$ 中的多項式。因此，若在 (17) 中把 $a,b,c$ 替換為這些 $y$ 中多項式，就得到
\[
 \left\{\sum\beta_\mu(u)\varphi_\lambda(y)\right\}^q\cdot \Gamma_\lambda(u)\zero{\overline{\mideal}_{(u)}},
\]
這由 (18) 等價於
\begin{equation}
 b^{(i)}(x)^q\cdot G_\lambda(x)\zero{\mideal};
 \qquad
 G_\lambda(x)\zero{\gideal_{i-1}}.
\tag{20}
\end{equation}
關係 (16)、(18) 和 (20) 表明，\emph{基本理想 $\gideal_{i-1}$ 的確是變換理想。}

\section*{22. K. Hentzelt 的改寫稿：\\ 關於多項式理想與結式的理論\quad（續）}
\begin{center}
Math. Ann. 88 (1923), pp. 53--79
\end{center}

\section*{§ 4. 多項式理想與線性形式模之間的聯繫}

為了把一個始終假定為變換理想的多項式理想 $\mideal$ 看作線性形式模 $\Mmod_{i-1}$ $(i=1\ldots n)$，要把各個多項式 $F(x)$ 看作 $x_1,\ldots,x_{i-1}$ 的冪積 $\xi_\lambda$ 中的線性形式：
\[
 F(x)=\sum a^{(i)}_\lambda\xi_\lambda,
\]
其中，按 §~3 的記號，$a^{(i)}_\lambda$ 是 $x_i\ldots x_n$ 中的多項式。由理想的定義立刻可知，$\Mmod_{i-1}$ 對這些整量 $a^{(i)},b^{(i)},\ldots$ 具有模性質；並且，由於 $x_1,\ldots,x_{i-1}$ 的無限多個冪積 $\xi$ 只線性地出現，憑藉它們的線性無關性，它們在 $\Mmod_{i-1}$ 中起不定元的作用。每個模 $\Mmod_{i-1}$ 含有無限多個不定元 $\xi$，但在每個單獨的線性形式中只出現有限多個不定元。同樣，每個基本理想 $\gideal_{i-1}$ 也要看作線性形式模 $\Gmod_{i-1}$，這裡的不定元仍是 $x_1\ldots x_{i-1}$ 的冪積。對於 $i=1$，只有一個對應於單位元的不定元 $\xi_0$。

現在可以--這是以下一切的基礎--儘管有無限多個不定元 $\xi$，仍按下列定理把問題限制到有限多個不定元中的線性形式模。

\medskip
\noindent\textbf{定理 VII。} \emph{剩餘類系統 $\Gmod_{i-1}\mid\Mmod_{i-1}$ 同某個剩餘類系統 $\Gmod^*_{i-1}\mid\Mmod^*_{i-1}$ 同構，其中 $\Mmod^*_{i-1}$ 是有限多個不定元中的模，$\Gmod^*_{i-1}$ 是它的基本模。把 $x_{i+1},\ldots,x_n$ 添入 $\mathsf{P}(u)$ 之後，模 $\Mmod_{i-1}$ 只有有限多個不等於 $1$ 的初等因子；這些初等因子正是不等於 $1$ 的 $\Mmod^*_{i-1}$ 的初等因子。}

$\Mmod_{i-1}$ 的初等因子按注 8) 中的方式定義。定理 VII 中出現的同構依賴於如下定義。

\medskip
\noindent\textbf{定義 V。} \emph{若兩個剩餘類系統之間可以建立一一對應，使得兩個類之差對應於相應兩個類之差，並且一個類乘以一個整量對應於相應類乘以同一個整量，則稱這兩個剩餘類系統同構。}

定理 VII 的證明通過完全歸納得到。對於 $i=1$，定理是顯然的：$\Mmod_0$ 是一個不定元 $\xi_0$ 中的線性形式模，$\xi_0$ 對應於 $x$ 的零維冪積，也就是單位元；整量是 $x_1\ldots x_n$ 中的所有多項式。基本理想 $\gideal_0$ 是由 $x_1\ldots x_n$ 中所有多項式組成的單位理想，因而 $\Gmod_0$ 是模 $(\xi_0)$，即 $\Mmod_0$ 的基本模；所以這裡 $\Gmod^*_0\mid\Mmod^*_0$ 與 $\Gmod_0\mid\Mmod_0$ 相同。最後，因為 $\Mmod_0$ 的秩為 $1$，它至多隻有一個不等於 $1$ 的初等因子。

我們先從 $i=1$ 過渡到 $i=2$，以便把由此得到的關於 $\Gmod_1\mid\Mmod_1$ 結構的認識用於一般歸納步驟。為此先說明：對於 $\Gmod_0\mid\Mmod_0$ 可以選擇一個關於 $x_1$ 有界的代表系統；由於 $\gideal_1$ 可被 $\gideal_0$ 整除，這隨後將導向 $\Gmod^*_1$ 的有限多個不定元。

作為變換理想，$\mideal$ 按 §~3 至少有一個關於 $x_1$ 正則的多項式 $C^{(1)}(x)$，設其次數為 $k$；因此 $\Mmod_0$ 含有線性形式 $C^{(1)}\xi_0$。由於 $C^{(1)}(x)$ 的正則性，每個多項式 $H(x)$ 都有表示
\begin{equation}
 H(x)\equiv b^{(2)}_0+b^{(2)}_1x_1+\cdots+b^{(2)}_{k-1}x_1^{k-1}\quad (C^{(1)}(x));
\tag{21}
\end{equation}
所以，因為 $C^{(1)}\xi_0$ 屬於 $\Mmod_0$，確實可以為 $\Gmod_0\mid\Mmod_0$ 選擇一個在 $x_1$ 中達不到第 $k$ 維的代表系統。

現在特別對 $\gideal_1$ 中的多項式 $G(x)$ 和 $\mideal$ 中的多項式 $F(x)$ 寫成
\begin{equation}
 \left\{\begin{aligned}
 G(x)&\equiv g^{(2)}_0+g^{(2)}_1x_1+\cdots+g^{(2)}_{k-1}x_1^{k-1} &&(C^{(1)}(x)),\\
 F(x)&\equiv a^{(2)}_0+a^{(2)}_1x_1+\cdots+a^{(2)}_{k-1}x_1^{k-1} &&(C^{(1)}(x)),
 \end{aligned}\right.
\tag{22}
\end{equation}
並令 $\Gmod^*_1$ 表示所有線性形式 $g(\xi)=g^{(2)}_0\xi_0+\cdots+g^{(2)}_{k-1}\xi_{k-1}$ 的系統，相應地令 $\Mmod^*_1$ 表示所有線性形式 $a(\xi)=a^{(2)}_0\xi_0+\cdots+a^{(2)}_{k-1}\xi_{k-1}$ 的系統；那麼 $\gideal_1$ 和 $\mideal$ 的理想性質表明，$\Gmod^*_1$ 與 $\Mmod^*_1$ 是 $\xi_0\ldots\xi_{k-1}$ 中的線性形式模，整量為 $b^{(2)}, c^{(2)}\ldots$。同樣，由 $C^{(1)}(x)$ 導出的主理想給出關於這些整量的一個模：
\[
 \mathfrak C_1=(\zeta_0,\zeta_1,\ldots,\zeta_\nu,\ldots),
\]
其中 $\zeta_\nu=x_1^{\nu}C^{(1)}(x)$。這些 $\zeta_\nu$ 關於整量 $b^{(2)},c^{(2)},\ldots$ 彼此線性無關，並且也同有限多個 $\xi$ 線性無關。由 $\mathfrak C_1$ 可被 $\Mmod_1$ 整除，並且因而可被 $\Gmod_1$ 整除，推出 $\Gmod^*_1$ 可被 $\Gmod_1$ 整除，$\Mmod^*_1$ 可被 $\Mmod_1$ 整除。考慮 (22)，得到
\begin{equation}
 \Gmod_1=(\Gmod^*_1,\mathfrak C_1),\qquad
 \Mmod_1=(\Mmod^*_1,\mathfrak C_1).
\tag{23}
\end{equation}

由 (23) 以及 $\xi$ 與 $\zeta$ 的線性無關性，定理 VII 在 $i=2$ 的情形如下推出。因為 $\mathfrak C_1$ 可被 $\Mmod_1$ 整除，首先有 $\Gmod_1\mid\Mmod_1=\Gmod^*_1\mid\Mmod_1$。為證明它同 $\Gmod^*_1\mid\Mmod^*_1$ 同構，讓 $\Gmod^*_1$ 中若干線性形式 $g^*(\xi)$ 給出 $\Gmod^*_1\mid\Mmod^*_1$ 的一個代表系統。由於 $\xi$ 與 $\zeta$ 的線性無關性，由 $g^*(\xi)\equiv0(\Mmod_1)$ 也推出 $g^*(\xi)\equiv0(\Mmod^*_1)$；因此這些 $g^*(\xi)$ 模 $\Mmod_1$ 仍然不同餘，構成 $\Gmod^*_1\mid\Mmod_1$ 的一個代表系統，並且由這個代表系統給出的從 $\Gmod^*_1\mid\Mmod_1$ 到 $\Gmod^*_1\mid\Mmod^*_1$ 的對應，在定義 V 的意義下是同構的。

完全同樣地得到：由 $b^{(2)}g(\xi)\equiv0(\Mmod_1)$、$b^{(2)}\ne0$，可推出 $b^{(2)}g(\xi)\equiv0(\Mmod^*_1)$、$b^{(2)}\ne0$。既然這對 $\Gmod^*_1$ 中所有 $g(\xi)$ 成立，而且只對這些 $g(\xi)$ 成立--因為按 (23)，$\Gmod^*_1$ 由基本理想 $\gideal_1$ 中同時又是 $\xi$ 中線性形式的全部多項式組成--由此 $\Gmod^*_1$ 被刻畫為 $\Mmod^*_1$ 的基本模。

最後把 $x_3\ldots x_n$ 添入 $\mathsf{P}(u)$，得到
\begin{equation}
\left\{\begin{aligned}
 \Gmod^*_1&=(\eta_1\ldots\eta_\rho);&
 \Mmod^*_1&=(e_1\eta_1,\ldots,e_\rho\eta_\rho);\\
 \Gmod_1&=(\eta_\rho\ldots\eta_1,\zeta_0\ldots\zeta_\nu\ldots);&
 \Mmod_1&=(e_\rho\eta_\rho,\ldots,e_1\eta_1,\zeta_0\ldots\zeta_\nu\ldots);
\end{aligned}\right.
\tag{24}
\end{equation}
因而
\[
 (e_\rho)=\Mmod_1/\Gmod_1=\Mmod_1/(\Mmod_1,\eta_\rho),\qquad
 (e_{\rho-1})=(\Mmod_1,\eta_\rho)/\Gmod_1\quad\hbox{等等。}
\]
於是，考慮注 8)，$e_\rho,e_{\rho-1},\ldots,e_1,1,\ldots,1,\ldots$ 被認出為 $\Mmod_1$ 的初等因子。這證明了定理 VII 對 $i=2$ 的全部部分。

應注意，(22) 與 (23) 只用到了 $C^{(1)}(x)$ 關於 $x_1$ 正則這一點。因此，如果不用 (12)，而採用特殊變換
\begin{equation}
 y_1=x_1;\quad y_2=u_{21}x_1+z_2;\quad\ldots;\quad y_n=u_{n1}x_1+z_n
\tag{25}
\end{equation}
作為基礎，則這些公式以及附屬於它們並導向定理 VII 的論證仍然有效；再同
\begin{equation}
\left\{\begin{array}{l}
 z=U_1(x)\quad\hbox{或}\\[2mm]
 z_2=x_2;\quad z_3=u_{32}x_2+x_3;\quad\ldots;\quad z_n=u_{n2}x_2+\cdots+u_{n,n-1}x_{n-1}+x_n
\end{array}\right.
\tag{26}
\end{equation}
複合，又得到 (12)。若 $\overline{\mideal}_{(u)}$ 通過 (25) 變為 $\widetilde{\mideal}$，而 $\widetilde{\gideal}_1,\widetilde{\Gmod}_1\ldots$ 對 $\widetilde{\mideal}$ 的意義同 $\gideal_1,\Gmod_1\ldots$ 對 $\mideal$ 的意義相同，則
\begin{equation}
 \widetilde{\Gmod}_1=(\widetilde{\Gmod}^{*}_1,\widetilde{\mathfrak C}_1),\qquad
 \widetilde{\Mmod}_1=(\widetilde{\Mmod}^{*}_1,\widetilde{\mathfrak C}_1).
\tag{27}
\end{equation}

這裡 (27) 只是由 (23) 通過 (26) 的反演而來。對 $\mideal$ 與 $C^{(1)}(x)$ 而言這是顯然的，因而對 $\widetilde{\mathfrak C}_1$ 與 $\widetilde\Mmod_1$ 也顯然。對 $\gideal_1$ 也成立，因為
\[
 b^{(2)}(x)G(x)\equiv0(\mideal),\quad b^{(2)}\ne0,
 \qquad
 \widetilde b^{(2)}(z)\widetilde G(x,z)\equiv0(\widetilde\mideal),\quad \widetilde b^{(2)}\ne0
\]
二者通過 $z=U_1(x)$ 互相決定。於是通過 (26) 有 $\widetilde\gideal_1=\gideal_1$；所以也有 $\widetilde\Gmod_1=\Gmod_1$。因此，由 (22) 給出的定義，對 $\Gmod_1^*$ 與 $\Mmod_1^*$ 同對 $\widetilde\Gmod_1^*$ 與 $\widetilde\Mmod_1^*$ 一樣成立，從而也對整個表示 (27) 成立。

對一般歸納步驟，假定下列命題成立：
\begin{equation}
\left\{\begin{aligned}
 \Gmod_{i-1}&=(\Gmod_{i-1}^*,\mathfrak C_{i-1}),&
 \Mmod_{i-1}&=(\Mmod_{i-1}^*,\mathfrak C_{i-1}),\\
 \mathfrak C_{i-1}&=(\zeta_0,\zeta_1,\ldots,\zeta_\nu,\ldots),
\end{aligned}\right.
\tag{28}
\end{equation}
其中 $\Gmod^*_{i-1}$ 與 $\Mmod^*_{i-1}$ 是關於整量 $b^{(i)},c^{(i)},\ldots$ 的線性形式模，它們含有有限多個不定元 $\xi$，即 $x_1,\ldots,x_{i-1}$ 的某些冪積；而這些 $\zeta$ 關於這些整量彼此線性無關，並且也同 $\xi$ 線性無關。對應於 (28) 的表示以及 $\xi$ 與 $\zeta$ 的線性無關性，還應在不用 (12) 而採用特殊變換
\begin{equation}
\begin{aligned}
 y_1&=x_1;\quad \ldots;\quad y_{i-1}=u_{i-1,1}x_1+\cdots+x_{i-1};\\
 y_i&=u_{i,1}x_1+\cdots+u_{i,i-1}x_{i-1}+z_i;\quad\ldots;\quad
 y_n=u_{n,1}x_1+\cdots+u_{n,i-1}x_{i-1}+z_n,
\end{aligned}
\tag{29}
\end{equation}
作為基礎時同樣成立；該特殊變換可通過
\[
 z=U_{i-1}(x)\quad\hbox{或}\quad
 z_i=x_i;\quad z_{i+1}=u_{i+1,i}x_i+x_{i+1};\quad\ldots;\quad
 z_n=u_{n,i}x_i+\cdots+u_{n,n-1}x_{n-1}+x_n
\]
複合成為 (12)。並且，這種特殊表示應當只是由 (28) 通過 $z=U_{i-1}(x)$ 的反演而來。

由假設 (28) 以及 $\xi$ 與 $\zeta$ 的線性無關性，按照與上文接在 (23) 後完全相同的論證，通過指定同一個代表系統，可推出 $\Gmod_{i-1}\mid\Mmod_{i-1}$ 與 $\Gmod^*_{i-1}\mid\Mmod^*_{i-1}$ 的同構；同樣還推出，$\Gmod^*_{i-1}$ 成為 $\Mmod^*_{i-1}$ 的基本模，並且 $\Mmod_{i-1}$ 只有有限多個不等於 $1$ 的初等因子，即 $\Mmod^*_{i-1}$ 中不等於 $1$ 的初等因子。

為完成歸納步驟，首先仍須證明：對於 $\Gmod^*_{i-1}\mid\Mmod^*_{i-1}$，因而也對於 $\Gmod_{i-1}\mid\Mmod_{i-1}$，可以取一個關於 $x_i$ 有界的代表系統。這由定理 III 中在 (7) 證明的商表示得到：
\begin{equation}
 \Gmod^*_{i-1}=\Mmod^*_{i-1}/(C^{(i)}),
\tag{30}
\end{equation}
其中 $C^{(i)}$ 表示 $\Mmod^*_{i-1}$ 中任意一個不恆等於零的 $\rho$ 次行列式，$\rho$ 是 $\Mmod^*_{i-1}$ 的秩。由假設中關於特殊變換 (29) 的部分可知，總可以假定 $C^{(i)}$ 關於 $x_i$ 正則。事實上，與 $\Mmod^*_{i-1}$ 相應的模 $\widetilde\Mmod^*_{i-1}$ 有同樣的秩；若 $\widetilde C^{(i)}$ 表示 $\widetilde\Mmod^*_{i-1}$ 中任意一個不恆等於零的 $\rho$ 次行列式，並且它通過 $z=U_{i-1}(x)$ 變為正則的 $C^{(i)}$，那麼按假設，$C^{(i)}$ 就是 $\Mmod^*_{i-1}$ 中的一個 $\rho$ 次行列式。

由 $C^{(i)}$ 的存在可知，適當編號 $\xi$ 以後，與 (8) 相應，有 $\rho$ 個如下特殊形狀的線性形式屬於 $\Mmod^*_{i-1}$：
\[
 L_\lambda(\xi)=C^{(i)}\xi_\lambda+c^{(i)}_{\lambda,1}\xi_{\rho+1}+\cdots+c^{(i)}_{\lambda,s-\rho}\xi_s
 \qquad(\lambda=1\ldots\rho).
\]
現在，每個 $\xi$ 中的線性形式都可以化為
\begin{equation}
 H(\xi)\equiv a^{(i)}_1\xi_1+\cdots+a^{(i)}_\rho\xi_\rho+b^{(i)}_1\xi_{\rho+1}+\cdots+b^{(i)}_{s-\rho}\xi_s\quad (L_1(\xi),\ldots,L_\rho(\xi)),
\tag{31}
\end{equation}
其中 $a^{(i)}_1\ldots a^{(i)}_\rho$ 關於 $x_i$ 有界，並且達不到 $C^{(i)}$ 的次數 $k$。這個表示是唯一的；因為由
\[
 a^{(i)}_1\xi_1+\cdots+a^{(i)}_\rho\xi_\rho+b^{(i)}_1\xi_{\rho+1}+\cdots+b^{(i)}_{s-\rho}\xi_s\equiv0\quad (L_1(\xi),\ldots,L_\rho(\xi))
\]
通過比較 $\xi_1\ldots\xi_\rho$ 的係數並考慮 $x_i$ 中的次數，可推出所有 $a^{(i)}$ 與 $b^{(i)}$ 都恆等消失。

現在特別設
\[
 H(\xi)\equiv0(G^*_{i-1});\qquad\hbox{于是}\qquad C^{(i)}H(\xi)\equiv0(\Mmod^*_{i-1})\quad\hbox{由 (30)。}
\]
同定理 III 的證明一樣，得到
\[
 C^{(i)}H(\xi)\equiv\sum_{\tau=1}^{s-\rho}\{C^{(i)}b^{(i)}_\tau-\sum_\lambda a^{(i)}_\lambda c^{(i)}_{\lambda,\tau}\}\xi_{\rho+\tau}
 \equiv0(\Mmod^*_{i-1}),
\]
因而，由於如定理 III 中那樣 $\xi_{\rho+1}\ldots\xi_s$ 的係數必須消失，便有
\[
 C^{(i)}b^{(i)}_\tau=\sum_\lambda a^{(i)}_\lambda c^{(i)}_{\lambda,\tau}\qquad(\tau=1\ldots s-\rho).
\]
因此 $b^{(i)}_\tau$ 也有界，並且達不到 $c^{(i)}_{\lambda,\tau}$ 的最大次數。於是，對於 $\Gmod^*_{i-1}\mid\Mmod^*_{i-1}$，也就是對於 $\Gmod_{i-1}\mid\Mmod_{i-1}$，存在一個關於 $x_i$ 有界且達不到某個固定次數 $r$ 的代表系統，這一點已經證明。

現在用 $\vartheta_1\ldots\vartheta_t$ 表示有限多個冪積
\[
 x_i^\alpha\xi_\lambda\quad(\alpha=0,1,\ldots,k-1;\ \lambda=1,2,\ldots,\rho),\qquad
 x_i^\beta\xi_{\rho+\tau}\quad(\beta=0,1,\ldots,r-1;\ \tau=1,2,\ldots,s-\rho)
\]
並把可數無限多個表達式
\[
 x_i^\gamma L_\lambda\quad(\gamma=0,1\ldots\text{ 無限}\ldots;\ \lambda=1,2,\ldots,\rho),
 \qquad
 x_i^\gamma\zeta_\nu\quad(\gamma,\nu=0,1\ldots\text{ 無限}\ldots)
\]
排成一個單無限序列 $\omega_0,\omega_1,\ldots,\omega_\mu\ldots$。那麼 $\vartheta$ 與 $\omega$ 關於整量 $b^{(i+1)},c^{(i+1)},\ldots$ 彼此線性無關，也各自內部線性無關。因為從
\[
 \sum c^{(i+1)}_{\alpha\lambda} x_i^\alpha\xi_\lambda+
 \sum c^{(i+1)}_{\beta\tau} x_i^\beta\xi_{\rho+\tau}+
 \sum d^{(i+1)}_{\gamma\lambda} x_i^\gamma L_\lambda(\xi)+
 \sum e^{(i+1)}_{\gamma\nu}x_i^\gamma\zeta_\nu=0,
\]
每個求和都只含有限多項，而由所假定的 $\xi$ 與 $\zeta$ 的線性無關性，以及 $\xi$ 關於整量 $b^{(i)},\ldots$ 的彼此線性無關性，可推出所有 $e^{(i+1)}_{\gamma\nu}$ 都消失。表示 (31) 的唯一性又給出 $c^{(i+1)}_{\alpha\lambda}$ 和 $c^{(i+1)}_{\beta\tau}$ 的消失；最後，由 $L_\lambda(\xi)$ 的線性無關性得到 $d^{(i+1)}_{\gamma\lambda}$ 的消失。

現在把模 $(\mathfrak C_{i-1},L_1(\xi),\ldots,L_\rho(\xi))$ 看作關於整量 $b^{(i+1)}$ 的模 $\mathfrak C_i$，則 $\mathfrak C_i$ 可被 $\Mmod_i$ 整除，因為 $\mathfrak C_{i-1},L_1(\xi)\ldots L_\rho(\xi)$ 可被 $\Mmod_{i-1}$ 整除，並且
\[
 \mathfrak C_i=(\omega_0,\omega_1,\ldots,\omega_\mu\ldots).
\]
對每個 $H(x)\equiv0(\gideal_{i-1})$，由 (28)、(31) 和剛才所證得到
\[
 H(x)\equiv b^{(i+1)}\vartheta_1+\cdots+b^{(i+1)}_t\vartheta_t\quad(\mathfrak C_i)
\]
這是關於整量 $b^{(i+1)},c^{(i+1)}$ 的一個模方程。又因為 $\gideal_i$ 可被 $\gideal_{i-1}$ 整除，$\mideal$ 可被 $\gideal_i$ 整除，所以現在特別對 $\mathfrak C_i$ 中的每個 $G(x)$ 和 $\Mmod_i$ 中的每個 $F(x)$ 置
\[
 G(x)\equiv g^{(i+1)}_1\vartheta_1+\cdots+g^{(i+1)}_t\vartheta_t\quad(\mathfrak C_i),
 \qquad
 F(x)\equiv a^{(i+1)}\vartheta_1+\cdots+a^{(i+1)}\vartheta_t\quad(\mathfrak C_i),
\]
並用 $\Gmod_i^*$ 表示由所有 $g(\vartheta)$ 組成的模，用 $\Mmod_i^*$ 表示由所有 $a(\vartheta)$ 組成的模，則完全由導向 (23) 的同一論證得到
\begin{equation}
 \Gmod_i=(\Gmod_i^*,\mathfrak C_i),\qquad
 \Mmod_i=(\Mmod_i^*,\mathfrak C_i),\qquad
 \mathfrak C_i=(\omega_0,\omega_1,\ldots,\omega_\mu\ldots).
\tag{32}
\end{equation}
在推出 (32) 時，也只用到了 $C^{(i)}$ 關於 $x_i$ 的正則性，因此若把特殊變換 (29) 的指標提高一位，整個論證仍然有效。於是，與 (27) 相連的同樣考慮表明，這個特殊表示只是由 (32) 通過 $z=U_i(x)$ 的反演得到。

這樣，一般歸納步驟中 (28) 等全部假設都已對下一個指標得到證明；而如那裡所示，定理 VII 可直接由這些假設推出，所以定理 VII 已一般地證明。同時也已說明，由任意選擇 $C^{(i)}$ 所造成的 $\Gmod^*_{i-1},\Mmod^*_{i-1}$ 的任意性，又由剩餘類系統 $\Gmod^*_{i-1}\mid\Mmod^*_{i-1}$ 的同構消除。

特別地，若 §~3 中作為係數域的域 $\mathsf{P}$ 具有特徵零--也就是說，若 $\mathsf{P}$ 中由單位元導出的素域是有理數型--則可以把不定元 $u_{\mu\nu}$ 特化為 $P$ 中的量 $\bar u_{\mu\nu}$，使得對 $i=1,\ldots,n$，各個 $C^{(i)}$ 都仍關於 $x_i$ 正則。上述論證表明，定理 VII 在這種特化下仍然成立。\footnote{Hentzelt 不取任意的 $\mathsf{P}$，而只取全體複數的域，並且主要考察這一後一情形；不定元只在真正的消去論中作為基礎。他隨後還基本按照這裡給出的推理說明，為了形成結式形式，只須在每個 $\Mmod_{i-1}$ 中把 $x$ 的次數推進到某個超過固定界限的有限次數即可。缺少的是定理 VII 給出的 $\Gmod_{i-1}\mid\Mmod_{i-1}$ 與 $\Gmod^*_{i-1}\mid\Mmod^*_{i-1}$ 的同構；正是這個同構說明了為何它獨立於次數數目。（E. N.）} 然而，基本模與初等因子不必一定由一般情形通過特化 $u_{\mu\nu}=\bar u_{\mu\nu}$ 得到。\footnote{對 $\mideal=(x^2,ux+y)$，有 $\gideal_0=(1)$、$\gideal_1=(1)$。若取 $C^{(1)}=x^2$，則 $\Gmod_1^*=(1,x)$、$\Mmod_1^*=(ux+y,xy)$；這個 $\Mmod_1^*$ 的初等因子為 $y^2$ 和 $1$，並且可取 $C^{(2)}=y^2$。當 $u=0$ 時，$C^{(1)}$ 與 $C^{(2)}$ 分別仍關於 $x$ 和 $y$ 正則；但 $\widetilde{\Mmod}_1^*=(y,xy)$ 的初等因子為 $y,y$。所以 $\widetilde{\Gmod}_1\mid\widetilde{\Mmod}_1$ 仍同一個有限模的剩餘類系統同構，但不同構於 $\Gmod_1\mid\Mmod_1$。若反而選擇 $C^{(1)}=ux+y$，並且以保持 $C^{(1)}$ 正則的方式特化，則同構也會保持。（E. N.）}

\section*{§ 5. 多項式理想的結式形式與初等因子形式}

把 $x_{i+1}\ldots x_n$ 添入 $\mathsf{P}(u)$，於是 $\Gmod_{i-1}$ 與 $\Mmod_{i-1}$，以及因而 $\Gmod^*_{i-1}$ 與 $\Mmod^*_{i-1}$，都變成線性形式模，其中整量可以看作一個不定元 $x_i$ 的多項式。由定理 VII，在這種情形下，$\Mmod_{i-1}$ 中不等於 $1$ 的初等因子，連同至多有限多個等於 $1$ 的初等因子，同 $\Mmod^*_{i-1}$ 的那些初等因子相同。這些初等因子的乘積，即 $\Mmod^*_{i-1}$ 的第 $\rho$ 個行列式因子，等於從 $\Gmod^*_{i-1}$ 到 $\Mmod^*_{i-1}$ 的過渡代換的行列式，因而按定義 III 等於 $N(\Gmod^*_{i-1}\mid\Mmod^*_{i-1})$。公式 (24)，或一般 $i$ 時由 (28) 得到的類似公式，表明這些初等因子的乘積也等於從 $\Gmod_{i-1}$ 到 $\Mmod_{i-1}$ 的過渡代換的行列式，因而應記為 $N(\Gmod_{i-1}\cdot\Mmod_{i-1})$。於是
\[
 N(\Gmod_{i-1}\mid\Mmod_{i-1})=N(\Gmod^*_{i-1}\mid\Mmod^*_{i-1})=R^{(i)}(x),
\]
其中多項式 $R^{(i)}(x)$，也就是 $\Mmod^*_{i-1}$ 中所有 $\rho$ 行行列式在多項式意義下的最大公因子，可以假定在 $x_{i+1},\ldots,x_n$ 中為整的並且本原；於是，作為 $C^{(i)}$ 的一個約數，它關於 $x_i$ 正則。同樣，$\Mmod_{i-1}$ 或 $\Mmod^*_{i-1}$ 的最高初等因子 $E^{(i)}(x)$ 也可以假定在 $x_{i+1},\ldots,x_n$ 中為整的並且本原，因而關於 $x_i$ 正則。這導向如下定義。

\medskip
\noindent\textbf{定義 VI。} \emph{多項式 $R^{(i)}(x)$ 稱為 $\mideal$ 的第 $i$ 級結式，$E^{(i)}(x)$ 稱為第 $i$ 級初等因子。於是，第 $i$ 級結式就是基本理想 $\gideal_{i-1}$ 關於 $\mideal$ 的範數；它被解釋為模 $\Gmod_{i-1}$ 關於 $\Mmod_{i-1}$ 的範數，其中 $x_{i+1}\ldots x_n$ 已添入 $\mathsf{P}(u)$。同樣，在同一添加下，$(E^{(i)}(x))$ 等於商 $\Mmod_{i-1}/\Gmod_{i-1}$。$\mideal$ 的結式形式和初等因子形式分別定義為乘積}
\[
 R_\mideal=R^{(1)}\cdot R^{(2)}\cdot\ldots\cdot R^{(n)};\qquad
 E_\mideal=E^{(1)}\cdot E^{(2)}\cdot\ldots\cdot E^{(n)}.
\]

對於結式形式和初等因子形式，有：

\medskip
\noindent\textbf{定理 VIII。} \emph{結式形式可被初等形式整除；反過來，$E_\mideal$ 的某個冪可被 $R_\mideal$ 整除。形式 $E_\mideal$ 與 $R_\mideal$ 都可被 $\mideal$ 整除；更一般地，$E^{(i)}\cdots E^{(n)}\gideal_{i-1}\equiv0(\mideal)$，因而 $R^{(i)}\cdots R^{(n)}\gideal_{i-1}\equiv0(\mideal)$。}

事實上，因為範數可被最高初等因子整除，反過來這個最高初等因子的某個冪又可被範數整除，所以在添入 $x_{i+1},\ldots,x_n$ 之後，這種可整除性對 $R^{(i)}(x)$ 與 $E^{(i)}(x)$ 成立。但 $R^{(i)}(x)$ 與 $E^{(i)}(x)$ 被假定為關於 $x_{i+1}\ldots x_n$ 的本原多項式；因此，這種可整除性對所有不定元 $x_i,x_{i+1},\ldots,x_n$ 成立。於是，按這些表達式在全部不定元 $x$ 中的定義，$R_\mideal$ 可被 $E_\mideal$ 整除，並且 $E_\mideal$ 的某個冪可被 $R_\mideal$ 整除。

進一步，由定理 VII 和定理 II，在添入 $x_{i+1},\ldots,x_n$ 之後，最高初等因子 $E^{(i)}(x)$ 被定義為商 $\Mmod_{i-1}/\Gmod_{i-1}$。因此，回到全部不定元 $x$ 中的多項式，對每個
\begin{equation}
 G(x)\equiv0(\gideal_{i-1})
 \quad\hbox{都存在}\quad b^{(i+1)}\ne0
 \quad\hbox{使得}\quad b^{(i+1)}E^{(i)}(x)G(x)\equiv0(\mideal).
\tag{33}
\end{equation}
特別地，$\gideal_0$ 是單位理想，所以
\[
 b^{(2)}E^{(1)}(x)\equiv0(\mideal);\quad b^{(2)}\ne0,
 \qquad\hbox{从而}\qquad E^{(1)}(x)\equiv0(\gideal_1).
\]
一般地，把 (33) 應用於
\[
 E^{(1)}(x)\cdots E^{(i-1)}(x)\equiv0(\gideal_{i-1}),
\]
便存在 $b^{(i+1)}\ne0$，使得 $b^{(i+1)}E^{(1)}\cdots E^{(i)}\equiv0(\mideal)$；於是 $E^{(1)}\cdots E^{(i)}\equiv0(\gideal_i)$。由於 $\gideal_n=\mideal$，得到
\begin{equation}
 E_\mideal=E^{(1)}\cdots E^{(n)}\equiv0(\mideal)
 \quad\hbox{並且因而}\quad
 R_\mideal=R^{(1)}\cdots R^{(n)}\equiv0(\mideal).
\tag{34}
\end{equation}

如果在 (33) 中讓 $G(x)$ 遍歷 $\gideal_{i-1}$ 的一個理想基中的有限多個多項式，並令 $c^{(i+1)}(x)$ 表示對應的 $b^{(i+1)}(x)$ 的乘積，則
\[
 c^{(i+1)}E^{(i)}(x)\gideal_{i-1}\equiv0(\mideal);
 \qquad c^{(i+1)}\ne0,
 \quad\hbox{因而}\quad E^{(i)}(x)\gideal_{i-1}\equiv0(\gideal_i),
\]
由此如上經有限次重複，得到
\[
 E^{(n)}(x)\cdots E^{(i)}(x)\gideal_{i-1}\equiv0(\mideal),
\]
並且因而也有
\[
 R^{(n)}(x)\cdots R^{(i)}(x)\gideal_{i-1}\equiv0(\mideal).
\]
這證明了定理 VIII 的全部部分。

定理 VIII 本質上依賴於初等因子形式的性質；結式形式對於 $\mideal$ 的特徵意義由下述定理說明。

\medskip
\noindent\textbf{定理 IX。} \emph{若 $\nideal$ 是 $\mideal$ 的一個約數--這裡約數按理想意義理解--並且結式形式 $R_\nideal$ 與 $R_\mideal$ 相同，則理想 $\nideal$ 與 $\mideal$ 也相同。}

令 $\gideal_0\ldots\gideal_{n-1},\gideal_n=\mideal$ 與 $\hideal_0\ldots\hideal_{n-1},\hideal_n=\nideal$ 分別表示 $\mideal$ 與 $\nideal$ 的基本理想。首先 $\gideal_0$ 與 $\hideal_0$ 都是單位理想，所以 $\gideal_0=\hideal_0$。現在假設 $\gideal_{i-1}=\hideal_{i-1}$，於是 $\Mmod_{i-1}$ 與 $\Nmod_{i-1}$ 有同一個基本模 $\Gmod_{i-1}$。假設 $R_\nideal=R_\mideal$，在添入 $x_{i+1}\ldots x_n$ 後可以寫為
\[
 N(\Gmod_{i-1}\mid\Nmod_{i-1})=N(\Gmod_{i-1}\mid\Mmod_{i-1})
 \quad\hbox{或}\quad
 N(\Gmod^*_{i-1}\mid\Nmod^*_{i-1})=N(\Gmod^*_{i-1}\mid\Mmod^*_{i-1}).
\]
這裡，按照 (32)，$\Nmod_{i-1}=(\Nmod^*_{i-1},\mathfrak C_{i-1})$，所以 $\Nmod^*_{i-1}$ 也是 $\xi$ 中的線性形式模，並以 $\Gmod^*_{i-1}$ 為其基本模。由於 $\Mmod_{i-1}$ 可被 $\Nmod_{i-1}$ 整除，這導致 $\Mmod^*_{i-1}$ 可被 $\Nmod^*_{i-1}$ 整除；定理 IV 表明，在這種添加下，兩個模 $\Nmod^*_{i-1}$ 與 $\Mmod^*_{i-1}$ 相同，因而 $\Nmod_{i-1}$ 與 $\Mmod_{i-1}$ 也相同。但添入 $x_{i+1}\ldots x_n$ 意味着向 $\Mmod_{i-1}$，相應地向 $\mideal$，加入所有滿足
\[
 b^{(i+1)}(x)G(x)\equiv0(\mideal),\qquad b^{(i+1)}\ne0
\]
的多項式 $G(x)$。因此，通過這種添加，$\Mmod_{i-1}$，相應地 $\mideal$，變為 $\gideal_i$；同樣，$\nideal$ 變為 $\hideal_i$。所以 $\gideal_i=\hideal_i$。最終得到 $\gideal_n=\hideal_n$，即 $\mideal=\nideal$。

最後，把同一轉移原則應用於定理 V，得到如下結果。

\medskip
\noindent\textbf{定理 X。} \emph{設 $R^{(i)}=S^{(i)}T^{(i)}$ 是 $R^{(i)}$ 在多項式意義下分解為互素多項式 $S^{(i)}$ 與 $T^{(i)}$。則存在唯一的理想 $\jideal$，也存在唯一的理想 $\tideal$，二者都是 $\mideal$ 的約數，並且有同一個第 $(i-1)$ 級基本理想 $\gideal_{i-1}$，使得 $S^{(i)}=N(\Gmod_{i-1}\mid\Smod_{i-1})$ 且 $T^{(i)}=N(\Gmod_{i-1}\mid\mathfrak T_{i-1})$。理想 $\jideal$ 與 $\tideal$ 分別定義為所有滿足下列條件的多項式 $S(x)$ 與 $T(x)$ 的全體：}
\[
 s^{(i+1)}(x)T^{(i)}(x)S(x)\equiv0(\mideal);\quad s^{(i+1)}\ne0,
\]
\[
 t^{(i+1)}(x)S^{(i)}(x)T(x)\equiv0(\mideal);\quad t^{(i+1)}\ne0,
\]
\emph{並且 $\gideal_i$ 成為 $\jideal$ 與 $\tideal$ 的最小公倍數，}
\[
 \gideal_i=[\jideal,\tideal].
\]

最後只需指出：$s^{(i+1)}$ 與 $t^{(i+1)}$ 可以分別對 $\jideal$ 與 $\tideal$ 固定選取，即取為對應於 $\jideal$ 與 $\tideal$ 的基元素的那些 $s^{(i+1)}$ 與 $t^{(i+1)}$ 的乘積；並且，如上所述，通過添入 $x_{i+1},\ldots,x_n$，理想 $\mideal$ 被變為 $\gideal_i$。特別地，$R^{(i)}$ 的分解可以繼續直到不可約多項式的冪--即準素因子--這帶來 $\gideal_i$ 作為最小公倍數的相應表示。

\section*{§ 6. 真正的消去論}

由定理 VIII，或由公式 (34)，結式形式和初等因子形式都可被 $\mideal$ 整除，因而在 $\mideal$ 的所有零點處消失。這裡所謂 $\mideal$ 的“零點”，是指 $x_1,\ldots,x_i$ 的值系，它們屬於由 $\mathsf{P}(u;x_{i+1},\ldots,x_n)$ 導出的代數閉域\footnote{Steinitz 在他的 Algebraische Theorie der Körper (J. f. M. 137 (1910), pp. 167--309) 中指出，每個域都可以本質唯一地擴張為代數閉域，從而給出代數基本定理的有理等價形式。事實上，對每個 $i=1,\ldots,n$，所涉及的是 $\mathsf{P}(u,x_{i+1},\ldots,x_n)$ 的一個有限擴域。（E. N.）}，使得 $\mideal$ 中的每個多項式都在這些值繫上消失，而 $x_{i+1}\ldots x_n$ 被看作已添入係數域 $(i=1,\ldots,n)$。這些添入的 $x_{i+1}\ldots x_n$，如將要說明的，也可以由 $\mathsf{P}(u)$ 中的量代替。反過來，消去論的內容，是從結式形式的零點推出 $\mideal$ 的零點。這個逆向結論依賴於本節要發展的逐次消去；由於下一節將證明其中出現的形式 $D^{(i)}(x)$ 可被 $E^{(i)}(x)$ 整除，所需的逆向結論便由 $E^{(i)}(x)$ 的某個冪可被 $R^{(i)}(x)$ 整除而得出。

這裡給出的逐次消去理論依賴於下列結果。

\medskip
\noindent\textbf{定理 XI。} \emph{設 $\bideal$ 是一個 $t$ 中多項式理想，係數在 $\mathsf{P}$ 中，因而是一個主理想；設 $h(t)$ 是 $\bideal$ 中一個次數為 $k$ 的多項式，並設 $\Bmod$ 是 $\bideal$ 模 $h(t)$ 後所得的 $1,t,\ldots,t^{k-1}$ 中線性形式模。則 $\Bmod$ 的秩為 $k-p$，其中 $p$ 是 $\bideal$ 的基多項式 $f(t)$ 的次數。}

按假設 $h(t)=h_1(t)\cdot f(t)$，其中 $h_1(t)$ 的次數為 $k-p$。因此，由 $f(t),tf(t),\ldots,t^{k-p-1}f(t)$ 表示的 $\bideal$ 模 $h(t)$ 的剩餘類線性無關；並且 $\bideal$ 模 $h(t)$ 的每個剩餘類都可以用這 $k-p$ 個特殊類以 $\mathsf{P}$ 中係數線性表示。但模 $h(t)$ 後，理想 $\bideal$ 變為 $1,t,\ldots,t^{k-1}$ 中的線性形式模，其秩因而正好是 $k-p$。

現在設 $\aideal=\aideal_1$ 是一個變換多項式理想，設 $A^{(1)}(x)$ 是 $\aideal_1$ 中關於 $x_1$ 正則、次數為 $k_1$ 的一個多項式。令 $\mathfrak A_1$ 為 $1,x_1,\ldots,x_1^{k_1-1}$ 中的線性形式模，係數為多項式 $a^{(2)}(x)$，它是 $\aideal_1$ 模 $A^{(1)}(x)$ 後所得的模。令 $\aideal_2$ 表示由 $\mathfrak A_1$ 中所有 $k_1$ 行行列式導出的 $x_2,\ldots,x_n$ 中理想。由定理 XI，$\aideal_2$ 是零理想當且僅當 $\aideal_1$ 中的多項式在多項式意義下有一個關於 $x_1$ 次數 $p_1>0$ 的公約數。若 $\aideal_2$ 不同於零理想，那麼由於 $\aideal$ 假定為變換理想，它含有一個關於 $x_2$ 正則、次數為 $k_2$ 的多項式 $A^{(2)}(x)$；這一點可由把變換 (12) 從特殊變換 (25) 與 (26) 複合出來而直接看出。再令 $\mathfrak A_2$ 表示 $1,x_2,\ldots,x_2^{k_2-1}$ 中的線性形式模，即 $\aideal_2$ 模 $A^{(2)}(x)$ 後所得的模；令 $\aideal_3$ 表示由 $\mathfrak A_2$ 中所有 $k_2$ 行行列式導出的 $x_3,\ldots,x_n$ 中理想。這個理想必須含有一個關於 $x_3$ 正則的多項式 $A^{(3)}(x)$。

如此繼續這個過程，直到達到一個零理想，或者直到 $\aideal_{n+1}$ 成為單位理想；因為 $\aideal_{n+1}$ 不含 $x$，它只能是零理想或單位理想。可以看出，理想 $\aideal_1,\aideal_2,\ldots$ 由 $\aideal$ 唯一確定，與所選擇的多項式 $A^{(\lambda)}(x)$ 無關。對 $\aideal_1=\aideal$ 這是顯然的，因此可假定它對 $\aideal_1,\ldots,\aideal_\lambda$ 已成立。若 $\aideal_\lambda$ 模 $A^{(\lambda)}(x)$ 後變為線性形式模 $\mathfrak A_\lambda$，則 $\mathfrak A_\lambda$ 也可以刻畫為 $\aideal_\lambda$ 中所有在 $x_\lambda$ 中未達到次數 $k_\lambda$ 的多項式全體；因此 $\mathfrak A_\lambda$ 只依賴於 $A^{(\lambda)}(x)$ 的次數 $k_\lambda$。設 $\bar A^{(\lambda)}(x)$ 是 $\aideal_\lambda$ 中關於 $x_\lambda$ 正則、次數 $\bar k_\lambda>k_\lambda$ 的一個多項式，設 $\overline{\mathfrak A}_\lambda$ 為相應的線性形式模，並設 $\bar\aideal_{\lambda+1}$ 為由 $\overline{\mathfrak A}_\lambda$ 中 $\bar k_\lambda$ 行行列式導出的理想。於是有模方程
\[
 \overline{\mathfrak A}_\lambda=(\mathfrak A_\lambda,A^{(\lambda)},x_\lambda A^{(\lambda)},\ldots,x^{\bar k_\lambda-k_\lambda-1}A^{(\lambda)}).
\]
由於 $A^{(i)}(x)$ 關於 $x_i$ 正則，多項式 $A^{(i)},x_iA^{(i)},\ldots$ 可以作為線性形式的新不定元引入；由此可知，$\mathfrak A_i$ 的 $k_i$ 行行列式同 $\overline{\mathfrak A}_i$ 的 $\bar k_i$ 行行列式相同，所以 $\bar\aideal_{i+1}=\aideal_{i+1}$。\footnote{這在原則上同 $\Gmod^*_{i-1}\mid\Mmod^*_{i-1}$ 與 $\Gmod_{i-1}\mid\Mmod_{i-1}$ 的同構所依據的推理相同。}

此外，$\aideal_{i+1}$ 總可被 $\aideal_i$ 整除。若 $\aideal_{i+1}$ 是零理想，這是顯然的；在相反情形中，$\mathfrak A_i$ 的秩為 $k_i$，所以 $\aideal_{i+1}$ 按定義由 $\mathfrak A_i$ 中的 $k_i$ 行行列式組成，且可被 $\mathfrak A_i$ 整除，因而可被 $\aideal_i$ 整除。於是每個 $\aideal_{i+1}$ 都可被 $\aideal$ 整除；所以若 $\aideal_{n+1}$ 成為單位理想，則 $\aideal$ 本身也是單位理想。

現在設 $\aideal$ 不同於單位理想，並且 $\aideal_1\ne0,\ldots,\aideal_i\ne0$、$\aideal_{i+1}=0$。令 $D^{(i)}(x)$ 為 $\aideal_i$ 中所有多項式的最大公約數；這個公約數由定理 XI 存在，且“約數”按多項式意義理解。作為 $A^{(i)}(x)$ 的一個約數，$D^{(i)}(x)$ 本身關於 $x_i$ 正則，次數為 $p_i>0$。把 $x_{i+1},\ldots,x_n$ 添入 $\mathsf{P}(u)$；於是 $D^{(i)}(x)$ 在 $\mathsf{P}(u;x_{i+1}\ldots x_n)$ 的一個代數擴域中分解為 $p_i$ 個線性因子。令 $x_i-\bar x_i$ 為其中一個線性因子，於是 $\aideal_i$ 中所有多項式在 $x_i=\bar x_i$ 時都消失；再由定理 XI，$\aideal_{i-1}$ 中的多項式在此特化下獲得一個最大公約數 $D^{(i-1)}(x)$。作為 $A^{(i-1)}(x)$ 的一個約數，$D^{(i-1)}(x)$ 又關於 $x_{i-1}$ 正則，次數為 $p_{i-1}>0$；特別地，它不可能恆等消失，並且在 $\mathsf{P}(u,\bar x_i,x_{i+1},\ldots,x_n)$ 的一個擴域中分解為 $p_{i-1}$ 個線性因子。若 $x_{i-1}-\bar x_{i-1}$ 是這樣的一個線性因子，則有一個來自 $\aideal_{i-2}$ 的最大公約數與之對應。如此繼續，就得到 $\aideal$ 的有限多個公共零點，它們位於 $\mathsf{P}(u;x_{i+1}\ldots x_n)$ 的一個代數擴域中。反過來，因為 $\aideal_\lambda$ 可被 $\aideal$ 整除，$\aideal$ 的每個零點也是 $\aideal_\lambda$ 的零點。特別地，若 $x_1=\bar x_1,\ldots,x_i=\bar x_i,x_{i+1}=x_{i+1},\ldots,x_n=x_n$ 是 $\aideal$ 的一個零點，則推出 $\aideal_{i+1}=0$，並且 $\aideal_i$ 中的多項式獲得一個帶有線性因子 $x_i-\bar x_i$ 的公約數。綜上：

\medskip
\noindent\textbf{定理 XII。} \emph{設 $\aideal$ 不同於單位理想，且 $\aideal_1\ne0,\ldots,\aideal_i\ne0$、$\aideal_{i+1}=0$。則在把 $x_{i+1},\ldots,x_n$ 添入 $\mathsf{P}(u)$ 之後，存在有限多個關聯值系 $\bar x_1,\ldots,\bar x_i$，它們位於 $\mathsf{P}(u,x_{i+1},\ldots,x_n)$ 的一個代數擴域中，使得 $\aideal$ 中的所有多項式都在 $x_1=\bar x_1,\ldots,x_i=\bar x_i,x_{i+1}=x_{i+1},\ldots,x_n=x_n$ 時消失。反過來，若 $\aideal$ 有這樣的一個零點，則推出 $\aideal_{i+1}=0$，並且 $\aideal_i$ 中所有多項式有一個公約數 $D^{(i)}(x)$，它含有線性因子 $x_i-\bar x_i$。}

定理 XII 特別表明，每個沒有零點的理想 $\aideal$ 都是單位理想。

為進行一個也明確顯示這些值系關聯關係的分解，引入新的不定元 $v$，它們可以添入 $\mathsf{P}(u,x_{i+1},\ldots,x_n)$。置
\begin{equation}
 x_i=-v_1x_1-\cdots-v_{i-1}x_{i-1}+z_i.
\tag{35}
\end{equation}
則 (12) 與 (35) 的複合仍是類型 (12) 的代換；現在只不過把 $u_{i\kappa}$ 替換為 $w_{i\kappa}=u_{i\kappa}-v_\kappa$，更一般地把 $u_{i+\lambda,\kappa}$ 替換為 $w_{i+\lambda,\kappa}=u_{i+\lambda,\kappa}-u_{i+\lambda,i}v_\kappa$。若代換 (35) 把按假設已變換的理想 $\aideal$ 變為 $\bideal$，則 $\bideal$ 只是由 $\aideal$ 通過把 $u_{\mu\nu}$ 替換為 $w_{\mu\nu}$ 並添入 $v$ 而得到；由同一過程，由 $\bideal$ 定義的理想 $\bideal_1,\bideal_2,\ldots$ 由 $\aideal_1,\aideal_2,\ldots$ 得到。反過來，$\aideal_\lambda$ 由 $\bideal_\lambda$ 通過 $v=0$ 得到；因此理想 $\aideal_\lambda$ 與 $\bideal_\lambda$ 同時是零理想，或同時不同於零理想。

現在仍設 $\aideal_1\ne0;\ldots;\aideal_i\ne0;\aideal_{i+1}=0$，因此也有 $\bideal_1\ne0;\ldots;\bideal_i\ne0;\bideal_{i+1}=0$。設 $\bar x_1\ldots\bar x_i,x_{i+1}\ldots x_n$ 是 $\aideal$ 的一個關聯零點系。定義 $\bar z_i=\bar x_i+v_1\bar x_1+\cdots+v_{i-1}\bar x_{i-1}$。於是由 (35)，$\bar x_1\ldots\bar x_{i-1},\bar z_i,x_{i+1}\ldots x_n$ 是 $\bideal$ 的一個零點。因此若 $H^{(i)}(z,x)$ 表示 $\bideal_i$ 中所有多項式的最大公約數，並且可假定它在 $v$ 中為整的且本原，則由定理 XII 的最後部分，$H^{(i)}(z,x)$ 有因子
\[
 z_i-\bar z_i=z_i-(\bar x_i+v_1\bar x_1+\cdots+v_{i-1}\bar x_{i-1}),
\]
並且在添入 $x_{i+1},\ldots,x_n$ 之後，$\aideal$ 的每個零點都對應於這樣的一個因子。還必須說明，這已經窮盡了 $H^{(i)}(z,x)$ 的因子分解，也就是說，不能從 $H^{(i)}$ 分出一個並不分解為 $z$ 與 $v$ 中線性因子的因子 $H^{(i)}_1$。若不是這樣，則由於假定它在 $v$ 中本原，$H^{(i)}$ 至少會含有一個線性因子 $z_i-\bar{\bar z}_i$，其中 $\bar{\bar z}_i$ 屬於 $\mathsf{P}(u,v,x_{i+1}\ldots x_n)$ 的一個代數擴域。若 $\bar{\bar x}_1,\ldots,\bar{\bar x}_{i-1},\bar{\bar x}_i,x_{i+1},\ldots,x_n$ 是 $\aideal$ 的相應零點，則它必須同添入 $x_{i+1},\ldots,x_n$ 後出現的 $\aideal$ 的有限多個零點之一相同，因而必須與 $v$ 無關。因此 $\bar{\bar z}_i$ 必然線性地含有 $v$，而不是代數地或有理地非線性地含有 $v$；這與假設相反，因為還可以從 $H^{(i)}_1$ 分出另一個 $z$ 與 $v$ 中的線性因子 $z_i-(\bar{\bar x}_i+v_1\bar x_1+\cdots+v_{i-1}\bar x_{i-1})$。於是顯式分解為
\[
 H^{(i)}(z,x)=\prod \{z_i-(\bar x_i+v_1\bar x_1+\cdots+v_{i-1}\bar x_{i-1})\}^{\alpha_\lambda}.
\]
由於 $H^{(i)}$ 的次數和各個指數 $\alpha_\lambda$ 必須同 $D^{(i)}(x)$ 中相應的次數和指數一致--因為 $H^{(i)}$ 由 $D^{(i)}$ 通過特化 $u_{\mu\nu}=w_{\mu\nu}$ 得到，而 $D^{(i)}$ 由 $H^{(i)}$ 通過特化 $v=0$ 得到--這個分解還說明，由於不定元 $u_{\mu\nu}$ 已添入 $\mathsf{P}$，從 $\aideal$ 出發的消去中出現的 $D^{(i)}$ 的每個零點 $\bar x_i$，分別對應於來自 $\aideal_{i-1},\ldots,\aideal_1$ 的一個線性最大公約數，或這種公約數的一個冪；因而每次只得到 $\aideal$ 的一個關聯零點系 $\bar x_1,\ldots,\bar x_i,x_{i+1},\ldots,x_n$。\footnote{應指出，這裡給出的逐次消去不同於 Kronecker 的方法，它只依賴於給定的理想 $\aideal$，而與任何理想基無關；特別地，這一點也適用於出現的指數 $\alpha_\lambda$。按 Hentzelt 的說法，若要用這種方法完整執行消去，就需要在商 $\aideal_i/D^{(i)}$ 上繼續這個過程；考慮到下一節，這一步可以省略。（E. N.）}

\section*{§ 7. 結式形式與消去論}

為了建立結式形式與消去論之間的聯繫，把 §~6 的逐次消去專門應用於理想除以第 $(i-1)$ 級基本理想所得的商 $\aideal=\mideal/\gideal_{i-1}$。首先必須證明這個商是變換理想。現在 $\mideal$ 是變換理想，因此由 §~3 的定理 VI，$\gideal_{i-1}$ 也是變換理想。由
\[
 H(x)\equiv0(\mideal/\gideal_{i-1});\qquad
 H(x)=\overline H(y)=\sum\alpha_\lambda(u)\overline H_\lambda(y)=\sum\alpha_\lambda(u)H_\lambda(x)
\]
得到
\[
 \overline H(y)\cdot\overline\gideal_{i-1(u)}\equiv0(\overline\mideal_{(u)}),
 \quad\hbox{因而也有}\quad
 \overline H(y)\overline\gideal_{i-1}\equiv0(\overline\mideal_{(u)}),
\]
或者
\[
 \overline H_\lambda(y)\overline\gideal_{i-1}\equiv0(\overline\mideal),
 \quad\hbox{因而}\quad H_\lambda(x)\gideal_{i-1}\equiv0(\mideal).
\]
這證明了 $H_\lambda(x)$ 可被 $\mideal/\gideal_{i-1}$ 整除，因而這個理想是變換理想，所以 §~6 的消去方法可以應用。

另一方面，(30) 在考慮 (28) 後--見 §~4--表明 $C^{(i)}(x)$ 可被 $\mideal/\gideal_{i-1}$ 整除；這裡 $C^{(i)}(x)$ 表示 $\Mmod^*_{i-1}$ 中任意一個不恆等於零的 $\rho$ 次行列式。由於當 $i>1$ 時這個 $C^{(i)}$ 關於 $x_1$ 的次數為零，對應於理想 $\aideal=\aideal_1$ 的模 $\mathfrak A_1$ 含有多項式 $C^{(i)},x_1C^{(i)},\ldots,x_1^{k_1-1}C^{(i)}$，因而 $(C^{(i)})^{k_1}$ 可被 $\aideal_2$ 整除；同樣，$(C^{(i)})^{k_1k_2}$ 可被 $\aideal_3$ 整除，最後 $(C^{(i)})^{k_1k_2\cdots k_{i-1}}$ 可被 $\aideal_i$ 整除。因此 $\aideal_1,\ldots,\aideal_i$ 都不同於零理想，$i=1$ 時也如此。現在令 $D^{(i)}(x)$ 為 $\aideal_i$ 中所有多項式的最大公約數；只有當 $\aideal_{i+1}$ 是零理想時，它才具有次數 $p_i>0$。由 $D^{(i)}(x)$ 的定義，並考慮 $\aideal_i$ 可被 $\aideal$ 整除，得到
\[
 d^{(i+1)}D^{(i)}(x)\equiv0(\mideal/\gideal_{i-1});\qquad d^{(i+1)}\ne0.
\]
於是 $d^{(i+1)}D^{(i)}$ 是一個不含 $x_1\ldots x_{i-1}$ 的多項式，並屬於 $\mideal/\gideal_{i-1}$。但 $E^{(i)}(x)$ 按定義（定理 II 與 §~4）是 $\Mmod_{i-1}$ 或 $\Mmod^*_{i-1}$ 的最高初等因子，也就是所有屬於 $\mideal/\gideal_{i-1}$ 且不含 $x_1,\ldots,x_{i-1}$ 的多項式在多項式意義下的最大公約數。因此 $d^{(i+1)}D^{(i)}$ 可被 $E^{(i)}(x)$ 整除；又因 $E^{(i)}(x)$ 關於 $x_i$ 正則，$D^{(i)}(x)$ 本身也可被 $E^{(i)}(x)$ 整除。

如果變換 (12) 從一開始就同 (35) 複合，也就是說，如果不定元 $u_{\mu\nu}$ 被 $w_{\mu\nu}$ 替代，同一論證仍可進行；這給出 $H^{(i)}(z,x)$ 可被相應的 $\bar E^{(i)}(z,x)$ 整除。正如 $D^{(i)}(x)$ 與 $H^{(i)}(z,x)$ 通過特化彼此轉化一樣，$E^{(i)}(x)$ 與 $\bar E^{(i)}(z,x)$、以及 $R^{(i)}(x)$ 與 $\bar R^{(i)}(z,x)$ 也通過特化彼此轉化，所以分解為線性因子時的重數也必須彼此一致。再考慮到 $E^{(i)}(x)$ 的某個冪可被 $R^{(i)}(x)$ 整除，得到如下總結。

\medskip
\noindent\textbf{定理 XIII。} \emph{若 $R^{(i)}(x)$ 在 $\mathsf{P}(u,x_{i+1}\ldots x_n)$ 的一個代數擴域中分解為線性因子 $x_i-\bar x_i$，則 $\bar x_i$ 可以且只能以一種方式補全為一個關聯值系 $\bar x_1\ldots\bar x_i,x_{i+1}\ldots x_n$，它是 $\mideal/\gideal_{i-1}$ 的零點，因而也是 $\mideal$ 的零點。由 $R^{(i)}$ 的線性因子以這種方式得到的 $\mideal$ 的有限多個零點--在添入 $x_{i+1}\ldots x_n$ 後為有限多個--由顯式分解}
\begin{equation}
 \bar R^{(i)}(z,x)=\prod\{z_i-(\bar x_{i\nu}+v_1\bar x_{1\nu}+\cdots+v_{i-1}\bar x_{i-1,\nu})\}^{\lambda_\nu}.
\tag{36}
\end{equation}
\emph{彙集起來。由於 $\bar R^{(i)}(z,x)$ 關於 $z_i$ 正則，若把添入 $\mathsf{P}(u)$ 的 $x_{i+1},\ldots,x_n$ 替換為來自 $\mathsf{P}(u)$ 或來自其導出的代數閉域的任意特殊值系 $\bar x_{i+1},\ldots,\bar x_n$，這個分解仍保持不變。}

這也完成了按照結式形式的各個因子對 $\mideal$ 的零點進行分離；添入 $\mathsf{P}(u)$ 的不定元 $x$ 的個數也稱為相應代數對象的維數。

若在 $\bar R^{(i)}(z,x)$ 或相應的 $R^{(i)}(x)$ 的分解中，把在 $\mathsf{P}(u,x_{i+1},\ldots,x_n)$ 上共軛的因子合併--這些因子在一般的 $\mathsf{P}$ 下可能全同或部分相同--就得到 $R^{(i)}(x)$ 分解為在 $\mathsf{P}(u,x_{i+1}\ldots x_n)$ 上不可約的多項式冪：
\[
 R^{(i)}(x)=S^{(i)}_1(x)^{\beta_1}\cdots S^{(i)}_\varkappa(x)^{\beta_\varkappa}.
\]
由定理 X，這個分解對應於一個表示 $\gideal_i=[\jideal_1,\ldots,\jideal_\varkappa]$，使得 $S^{(i)}_\mu(x)^{\beta_\mu}$ 等於 $\gideal_{i-1}$ 關於理想 $\jideal_\mu$ 的範數，或者說等於模 $\Gmod_{i-1}$ 關於與 $\jideal_\mu$ 相應的模的範數。這闡明瞭指數 $\beta_\mu$ 的意義：特別地，若 $\gamma_\mu$ 是 $S^{(i)}_\mu$ 的次數，則重數 $\beta_\mu\gamma_\mu$ 等於 $\gideal_{i-1}$ 模 $\jideal_\mu$ 的、在 $P(u,x_{i+1},\ldots,x_n)$ 上線性無關的剩餘類個數。

最後簡要考慮 $\mathsf{P}$ 具有特徵零的特殊情形。若把 $u_{\mu\nu}$ 特化為 $\mathsf{P}$ 中的量 $\bar u_{\mu\nu}$，並且使 $n$ 個行列式 $C^{(i)}(x)$ $(i=1,2,\ldots,n)$ 仍關於 $x_i$ 正則\footnote{Hentzelt 從一開始就假定了對 $u_{\mu\nu}$ 的這種特化，因此必須調用複雜得多的考慮來證明 $H^{(i)}(z,x)$ 與 $\bar R^{(i)}(z.x)$ 可以分解為 $z$ 與 $v$ 中的線性因子。那裡出現的指數不一定必須同上面的指數一致。（E. N.）}，則 $R^{(i)}$ 的正則性也保持。在這種特化下，$R^{(i)}$ 以及 $R^{(i)}(z,x)$ 都成為 $\Mmod^*_{i-1}$ 中所有 $\rho$ 行行列式的一個公約數，雖然未必是最大公約數。但總可以選擇特化，使後一個性質也保持。因為如模基的存在所表明的，$R^{(i)}(z,x)$ 也可以刻畫為 $\Mmod^*_{i-1}$ 中有限多個 $\rho$ 行行列式的最大公約數。這樣特化後的 $\bar R^{(i)}(z,x)$ 的分解，只需在 (36) 中把 $u_{\mu\nu}$ 替換為 $\bar u_{\mu\nu}$ 即可得到，所以整個消去論都保持不變。

\begin{center}
（1922 年 3 月 17 日收到。）
\end{center}

%
\clearpage
%
%
\begingroup
\def\D{\mathit D}
\def\srcemph#1{\textbf{#1}}
\def\srcname#1{\textsc{#1}}
\def\srcnumdisplay#1#2{\begin{equation}\tag{#1}#2\end{equation}}
\def\PXXIIIitem#1#2{\par\noindent\hangindent=2.2em\hangafter=1\makebox[2.2em][l]{#1.}#2\par}
\pdfbookmark[0]{23. 代數不變量與微分不變量}{paper-23}
\section*{23. 代數不變量與微分不變量}
\begin{center}
Jahresber. d. Deutschen Mathematiker-Vereinigung 32 (1923), S. 177--184
\end{center}

\begin{center}
{\Large\bfseries 代數不變量與微分不變量}\par
\vspace{0.45em}
據 1922 年 9 月 18 日萊比錫報告整理。\par
\vspace{1.0em}
作者：\textsc{Emmy Noether}，哥廷根。
\end{center}

\begingroup
\renewcommand{\thefootnote}{\arabic{footnote})}
\setcounter{footnote}{0}

如果要我報告代數不變量與微分不變量的發展，那麼在兩個領域中，我都願意把討論限制在 \srcname{Hilbert} 所說的、接在樸素時期和形式時期之後的批判時期。對於代數不變量，這一批判時期以 \srcname{Hilbert} 本人的名字為標誌；對於微分不變量，則以 \srcname{Riemann} 的名字為標誌；或者從內容上說：對於代數不變量，它以代數的算術方法為標誌，這些方法本質上正是圍繞這些問題發展起來，並且在這裡顯示了它們的全部鋒芒；對於微分不變量，它以形式變分法的方法為標誌。因此，我願意報告這些算術方法及其超出本特殊主題的意義；而在第二部分，我可以只限於說明微分不變量如何歸屬於代數不變量，這種歸屬正是在變分法方法的基礎上完成的。

\textbf{1.} 眾所周知，\srcemph{代數不變量}以 $x_1,\ldots,x_n$ 中的\srcemph{一般線性群}為基礎：
\srcnumdisplay{1}{x_i=\sum s_{ik}x'_k
\qquad\text{或簡寫為：}\qquad x=S(x'),}
其中 $s_{ik}$ 表示不定元，所以 $\D=|s_{ik}|$ 非零。現在若 $f(a,x),g(b,x),\ldots$ 是 $x$ 中的形式，次數分別為 $\alpha,\beta,\ldots$，並以不定元 $a,b,\ldots$ 為係數，則通過 (1) 得到對 $a,b,\ldots$ 的\srcemph{誘導變換}；因為由
\[
\begin{aligned}
f(a,x)&=\sum a_{i_1\ldots i_n}x_1^{i_1}\cdots x_n^{i_n}
      =\sum a'_{j_1\ldots i_n}{x'_1}^{j_1}\cdots {x'_n}^{j_n}
      =f(a',x'),\\
g(b,x)&=\sum b_{r_1\ldots r_n}x_1^{r_1}\cdots x_n^{r_n}
      =\sum b'_{s_1\ldots s_n}{x'_1}^{s_1}\cdots {x'_n}^{s_n}
      =g(b',x')
\end{aligned}
\]
可推出 $a',b',\ldots$ 成為 $a,b,\ldots$ 的線性齊次函數，其係數關於 $s_{ik}$ 的次數分別為 $\alpha,\beta,\ldots$：
\srcnumdisplay{2}{a'=A_\alpha(a);\qquad b'=B_\beta(b)\ldots .}
現在所謂\srcemph{不變量}，是指相對於變換 (1) 與 (2) 的不變量；也就是說，它是不定元 $a,b,\ldots,x,y,\ldots$ 的整有理函數（即多項式）——其中也有 $y=S(y')$——在施行變換時只差一個因子，即代換行列式的某個冪，而再現自身：
\srcnumdisplay{3}{I(a',b',\ldots,x',y',\ldots)=\D^\rho I(a,b,\ldots,x,y,\ldots).}
$I$ 的係數可以假定為有理數，甚至可以假定為有理整數；因為係數在任意數域 $P$ 中的每個不變量，都可以用 $P$ 中係數線性表示為有限多個這樣的特殊不變量。同樣，假定 $I$ 在每一組變量中分別齊次也不是限制；因為每個不變量也可以表示為有限多個這種特殊不變量的和。

還應指出，反過來，(1) 與 (2) 又由要求 (3) 所決定。正如 \srcname{Ostrowski} 和 \srcname{Schur} 最近指出的那樣，\footnote{A. \textsc{Ostrowski} und I. \textsc{Schur}, Über eine fundamentale Eigenschaft der Invarianten einer binären Form, Math. Zeitschr. 15 (1922)，以及 \textsc{Ostrowski} 尚未發表的相關工作。}誘導變換也可以刻畫為所有在每一組變量中分別作線性變換的全體；除去某些可明確列出的例外，這些變換允許任意一個不含 $x,y,\ldots$ 的不變量 $I$。

兩個不變量的乘積仍是不變量；但它們的和是不變量，當且僅當權——即 (3) 中的指數 $\rho$——相同。不過，如果只限於行列式為一的變換，則任意和仍是不變量，並且窮盡了關於這個受限群的所有不變量。因此得到一個整環，而且是一個\srcemph{齊次整環}；也就是說，在這個整環中，一個多項式屬於該環，就蘊含它的各個齊次分量也分別屬於該環——在這裡這又正好回到按 (3) 在每一組變量中分別齊次的不變量。

這裡出現了不變量論以及一般算術代數發展所圍繞的基本問題：\srcemph{這個整環是否有限生成；也就是說，它是否有一個有限整性基（即有限代數生成組） $I_1,\ldots,I_k$，使每個不變量都能通過 $I_1,\ldots,I_k$ 作多項式表示，$I=G(I_1,\ldots,I_k)$？} \srcname{Hilbert} 對這個問題的解答——\srcname{Gordan} 對二元情形的證明由於其難以駕馭的符號計算不能推廣，而 \srcname{Mertens--Hilbert} 的證明由於使用二元形式分解為線性形式也不能推廣——建立在把整性基問題化歸為理想生成基問題之上。這裡我願意簡要勾勒其思路，強調其中的算術基本思想，然後轉向三個相連的問題圈：用有限步驟實際構造這個基；子群的有限生成性問題；以及考慮整係數條件時的有限生成性問題。

\textbf{2.} 我先回憶（有限）代數數域中\srcemph{理想的定義}。域中的一個數系 $\ma$ 稱為理想，如果
\PXXIIIitem{1}{$\ma$ 屬於該域全體代數整數構成的環 $\mo$，}
\PXXIIIitem{2}{隨 $\alpha$ 與 $\beta$ 一起，差 $\alpha-\beta$ 也屬於 $\ma$，}
\PXXIIIitem{3}{隨 $\alpha$ 一起，$\lambda\alpha$ 也屬於 $\ma$，其中 $\lambda$ 是 $\mo$ 的任意元素。}
而且，眾所周知，每個理想都有一個\srcemph{理想生成基}：$\ma=(\alpha_1,\ldots,\alpha_k)$；也就是說，存在 $\ma$ 中有限多個元素 $\alpha_1,\ldots,\alpha_k$，使得藉助 2. 和 3.，$\ma$ 可由 $\alpha_1,\ldots,\alpha_k$ 導出；於是 $\alpha=\lambda_1\alpha_1+\cdots+\lambda_k\alpha_k$ 窮盡 $\ma$ 的全部元素。\footnote{$k$ 可以減為二這一點，對下文無關緊要。}

這個理想定義只依賴於 $\mo$ 構成一個整環這一事實；因此，當把 $\mo$ 換為任意整環時，理想的概念以及理想生成基的概念都逐字保持不變。

\srcname{Hilbert} 的有限生成性證明現在分解為下列三個定理：
\PXXIIIitem{1}{在以一個有理性域為係數、含 $n$ 個不定元的全體多項式環中，每個理想都滿足理想基定理——因而任意系統 $\mS$ 也滿足。}
\PXXIIIitem{2}{一個由多項式組成的齊次整環 $\mJ$ 是有限生成的，當且僅當理想基定理在 $\mJ$ 中成立。}
\PXXIIIitem{3}{對於不變量整環 $\mJ$，理想基定理以下述加強形式成立：每個相對於全體多項式環形成的理想生成基，同時也是 $\mJ$ 中的理想生成基；因此除}
\[
 I=A_1I_1+\cdots+A_kI_k
 \qquad\text{之外，總還有}\qquad
 I=i_1I_1+\cdots+i_kI_k,
\]
其中 $A$ 表示 $a,b,\ldots,x,y,\ldots$ 中的多項式，而 $i$ 表示不變量。

1. 的證明在原則上依靠同代數數域中相同的論證；在兩種情形中，理想生成基的存在都直接來自一個關於線性形式模的一般定理。\footnote{參見我關於代數函數算術理論的比較報告。Jahresber. 28 (1920), S. 187 und Anm. 2), S. 188.}由 1. 直接推出，每個有限生成的多項式整環中的理想都存在理想生成基；對於齊次整環，反過來也容易看出。只有在 3. 中才使用了不變量的特殊性質；也就是說，在 \srcname{Hilbert} 那裡，3. 來自一個關於 $\Omega$ 過程的熟知定理；而最近 E. \srcname{Fischer} 已經說明，3. 可以——以其他一般定理為基礎——僅由一般線性群的性質推出，即它隨每個變換也含有與之逆轉置（反變）的變換，或含有後者的共軛復變換。\footnote{E. \textsc{Fischer}, Über die Endlichkeit der Invarianten. Göttinger Nachrichten 1915, und Leipziger Vortrag.}

\textbf{3.} \srcemph{基的實際構造}依靠引用\srcemph{函數域}理論，對有限生成性概念作進一步的算術改寫。以下命題成立：
\PXXIIIitem{4}{一個多項式整環 $\mJ$ 是有限生成的，當且僅當 $\mJ$ 中含有一個有限生成子環 $\mJ'$，使 $\mJ$ 在 $\mJ'$ 上整；此時 $\mJ'$ 的每個整性基，都可由這些整元素的一個基本系補充為 $\mJ$ 的一個整性基。\footnote{\textsc{Hilbert} 證明（Über die vollen Invariantensysteme, Math. Annalen 42 (1893), § 1 und 2），這些事實可由所假定的有限生成性推出；反過來，有限生成性可由整依賴推出，則來自有理基的存在（E. Noether, Körper und Systeme rationaler Funktionen, Math. Annalen 76 (1915), § 12）。}}
特別地，像不變量的情形那樣，若處理的是齊次整環 $\mJ$，並且理想基定理在加強形式 3. 下成立，那麼這樣的整環 $\mJ'$ 可以由某些總是存在的有限多個多項式 $I_1,\ldots,I_s$ 作為整性基導出，並且這些多項式一旦全都為零，$\mJ$ 中所有多項式都隨之為零。這個事實又建立在一個一般理想論定理上：
\PXXIIIitem{5}{對每個多項式理想 $\ma$，都存在一個通常的有理整數 $r$，使得對每個在 $\ma$ 的所有零點上消失的理想 $\mb$，$\mb^r$ 可被 $\ma$ 整除。}

在不變量情形中，可以確定 $I_1,\ldots,I_s$ 的\srcemph{次數上界}，因而也可以確定它們的個數 $s$ 的上界；這個上界只依賴於所考察的基本形式的次數。在這一點上，又使用了不變量的特殊性質；相關定理說，一個作了數值特化的基本形式系統有一個非零不變量，當且僅當 $\D=|s_{ik}$ 在變換後的係數上整。

從這個上界出發，K. \srcname{Hentzelt}\footnote{我將很快從其遺稿中發表這項工作。} 又確定了\srcemph{整性基中不變量次數的上界}，這個上界同樣只依賴於基本形式的次數。他做到這一點，是通過給出 5. 中\srcemph{指數} $r$ 的一個\srcemph{上界}；這個上界只依賴於 $\ma$ 的一個理想生成基中多項式的個數和最高次數，與這些多項式的係數無關，並且可由給定的數在有限步內計算出來。原則上，所提出的問題由此完全解決；誠然，給出的界相當高。例如，對於一個三元二次形式，判別式——也就是一個三次不變量——已經構成所有不含 $x$ 的不變量的基，但按 \srcname{Hilbert} 和 \srcname{Hentzelt} 的估計，人們會得到高達數百萬的界。

\textbf{4.} 在一般線性群的\srcemph{子群}的不變量有限生成性問題上，到目前為止只得到了一些部分結果。方法幾乎完全建立在按照 \srcname{Study} 為正交群採用的辦法，把子群的不變量化歸為一般線性群的同時不變量。\srcname{Weitzenböck} 對運動不變量和仿射不變量、\srcname{Deruyts} 對半不變量和切變不變量都實行了這種化歸；\footnote{文獻參見 \textsc{Weitzenböck} 關於不變量論的百科全書條目，III E 1。}這個方法的適用範圍仍未決定。2. 末尾提到的 \srcname{Fischer} 的意見同樣解決了一系列子群的有限生成性問題；特別地，它由此給出了正交群情形的最簡單證明。但半不變量的例子表明——正如 \srcname{Fischer} 最近指出的那樣\footnote{Leipziger Vortrag und Math. Zeitschrift.}——對於子群，儘管有限生成性成立，理想基定理的加強形式 3. 不一定成立。

按域論解釋，這裡出現的有限生成性問題導向 \srcname{Hilbert} 關於\srcemph{相對整函數}的問題；也就是說，問題是：由一個有理函數域中的全體多項式組成的這類整環是否總是有限生成的。沿這個方向，有一個初等可證的事實：有限群的不變量可以直接給出一個優良的整性基——即伽羅瓦預解式的係數系統。\footnote{E. Noether, Der Endlichkeitssatz der Invarianten endlicher Gruppen. Math. Ann. 77 (1916).}

\textbf{5.} 不變量有限生成性定理是否能在\srcemph{整係數條件}方面加強，在一般情形下也尚未解決。誠然，如 \srcname{Hilbert} 所示，理想基定理也適用於全體整係數多項式構成的環；同樣，2. 中理想生成基與整性基之間的聯繫仍然有效；\footnote{特別地，由這個聯繫可推出一般理想論的分解定理適用於所有有限生成整環，正如我在 Math. Ann. 83 (1921) 中所發展的一樣。}但對於——其實並非必要的——加強形式 3.，證明失效了。對於\srcemph{二元不變量}，可以證明\srcemph{整係數情形的有限生成性}；\footnote{E. Noether, Die Endlichkeit des Systems der ganzzahligen Invarianten binärer Formen. Göttinger Nachrichten 1919.}這個證明是二元 \srcname{Mertens--Hilbert} 有限性證明在整係數條件下的加強；不過它除使用整係數理想基定理外，還使用二元形式分解為線性形式，因而不能推廣到更多不定元。另一方面，在類似基礎上，也可得到\srcemph{整係數情形的有限生成性}，這一次是對\srcemph{有限群不變量}而言，作為上述證明的加強；不過這裡得到的又只是存在性證明，而不是基的實際構造。\footnote{見腳註 11 所引工作的 § 4。}在這裡，恐怕也只有函數域理論與一般理想論的進一步發展才能達到目標。

\textbf{6.} 現在我轉向\srcemph{微分不變量}，以及引言中提出的、把它們歸約到代數不變量理論的問題。眾所周知，基礎群通過 $x_i=x_i(y)$ 定義如下：
\srcnumdisplay{4}{d x_i=\sum \frac{\partial x_i}{\partial y_k}d y_k;\qquad
 d\delta x_i=\sum \frac{\partial x_i}{\partial y_k}d\delta y_k+
\sum\frac{\partial^2 x_i}{\partial y_k\partial y_l}d y_k\,\delta y_l;\ \ldots .}
為了轉到 (4) 所誘導的變換，可以從任意函數 $f(x,d x)=g(y,d y)$ 出發；通過要求 $\delta f,\delta^2f,\ldots$ 相對於 (4) 不變，就得到 $f$ 關於 $x$ 與 $d x$ 的導數的變換。例如，若限於關於 $d x$ 齊次的形式
\[
 f(x,d x)=\sum a_{i_1\ldots i_n}(x)d x_1^{i_1}\cdots d x_n^{i_n}
 =\sum a'_{j_1\ldots j_n}(y)d y_1^{j_1}\cdots d y_n^{j_n}
 =g(y,d y),
\]
則 $a'$ 關於 $a$ 線性，$\frac{\partial a'}{\partial y}$ 關於 $a$ 與 $\frac{\partial a}{\partial x}$ 線性，依此類推：
\srcnumdisplay{5}{a'=A_0(a),\qquad
\frac{\partial a'}{\partial y}=A_1\!\left(a,\frac{\partial a}{\partial x}\right),\qquad
\frac{\partial 2a'}{\partial y^2}=A_2\!\left(a,\frac{\partial a}{\partial x},\frac{\partial 2a}{\partial x^2}\right),\ldots}
問題就是相對於變換 (4) 與 (5) 的不變性。這裡出現的所有量
\[
a,\ \frac{\partial a}{\partial x},\ldots,\ d x,\ \delta x,\ d\delta x,\ldots;\ 
\frac{\partial x}{\partial y},\ \frac{\partial^2 x}{\partial y^2},\ldots
\]
都要看作不定元，所以每個單獨變換的行列式當然非零。只有當形式上完成的理論應用於特殊函數時，才會出現限制到這些函數的問題：即要求變換在某個區域中保持唯一可逆，並且存在足夠多的微分商；這正如代數情形中，在數值特化下必須假定行列式 $|s_{ik}|$ 非零一樣。

在把所有出現的論元都看作不定元的這種處理中，人們面對的是一個不能直接用通常代數方法處理的特殊線性群；在這個群中，$d x$ 的變換以及 $a'=A_0(a)$ 與 (1) 和誘導變換 (2) 一致。於是問題產生：是否可以相當一般地通過\srcemph{引入新的論元}，把變換 (4) 和 (5) \srcemph{歸約}為一般線性群 (1) 以及由它\srcemph{誘導的變換} (2)？

事實確是如此。高階微分 $d^2x,\ d\delta x\ldots$ 要由屬於 $f(x,d x)$ 的變分問題所產生的 Lagrange 導數取代，並由通過極化形成的進一步組合取代（Weyl 和 Schouten 意義下的聯絡）；所有這些量都同 $d x$ 同變，所以它們的變換由一般線性群 (1) 給出。另一方面，$\frac{\partial a}{\partial x},\frac{\partial^2a}{\partial x^2},\ldots$（或者更一般地，假定齊次的 $f$ 關於 $x$ 與 $d x$ 的導數）要由某些形式的係數取代；這些形式來自“高階變分的正規形”，其中已經按上述方式消去了 $d^2x,\ d\delta x\ldots$，以及它們的“協變導數”
\[
 [\Omega_i],\quad [\Omega_i^{(1)}],\quad [\Omega_i^{(2)}],\ldots\qquad (i=1,2,\ldots);
\]
而由於 $[\Omega_i^{(k)}]$ 是關於 (4) 與 (5) 的不變量，並且只依賴於 $d x$ 與 $\delta x$，這些係數滿足代數誘導的變換 (2)。\footnote{E. Noether, Invarianten beliebiger Differentialausdrücke. Göttinger Nachrichten 1918.}通過這個\srcemph{歸約定理}，微分不變量理論被置於代數不變量理論之下。

當 $f(x,d x)$ 是關於 $d x$ 的 $\alpha$ 次齊次形式時，$[\Omega_i]$ 的序列在 $[\Omega_\alpha]$ 處終止。對於二次形式，$[\Omega_2]$ 與 \srcname{Riemann} 的曲率形式相同；這裡指出的構造過程，正是 \srcname{Riemann} 在拉丁文獲獎論文中給出的過程，而 \srcname{Lipschitz} 獨立地、從 \srcname{Riemann} 的任教資格論文出發，也發展了這個過程。\footnote{也可參見我在 \textsc{Weitzenböck} 百科全書條目第二部分第 28 號中關於 \textsc{Riemann} 與 \textsc{Lipschitz} 的敘述，見上文腳註 7。}\srcname{Christoffel} 與 \srcname{Ricci} 已經對二次微分形式通過計算完成了歸約定理的證明；這個證明依賴於引入“\srcname{Riemann} 正規座標”，它把從一點發出的極值曲線變成直線。由於這樣一來，$x$ 的任意變換對應於正規座標的一個線性變換，這些座標就實現了向一般線性群的過渡。

現在的問題是：如果按照 \srcname{Weyl} 和 \srcname{Schouten}，定義這個群時不以變分問題為基礎，而只\srcemph{以一個聯絡為基礎}，這樣的歸約定理是否仍然存在？也就是說，對應於 Lagrange 導數，設立表達式 $\psi_i(d\delta x,d x,\delta x)$，通過它們消去高階微分，並且類比二次微分形式再加上一個本身並不必要的限制，即 $\psi_i$ 關於 $d x$ 應當是線性的（對任意齊次的 $f(x,d x)$，極化後的 Lagrange 導數關於 $d x$ 只是一階齊次的）。通過要求 $\psi$ 同 $d x$ 同變，其係數的變換被誘導出來，因而群也像 (5) 中那樣被確定。對於最低階的不變量，\srcname{Weitzenböck} 已經在 \srcname{Weyl} 情形下通過計算確認了歸約定理；一般構造仍然缺失。\footnote{參見：\textsc{Weyl}, Raum, Zeit, Materie; \textsc{Schouten}, Math. Zeitschr. 13 (1922); \textsc{Weitzenböck}, Sitzungsber. d. Wiener Akademie, 1920 u. 1921.}

\begin{center}
（1922 年 10 月 26 日收到。）
\end{center}
\endgroup
\endgroup

%
\clearpage
%

%
%
%

\setcounter{footnote}{0}
\pdfbookmark[0]{24. 消去論與一般理想論}{paper-24}
\section*{24. 消去論與一般理想論}
\begin{center}
Math. Ann. 90 (1923), pp. 229--261
\end{center}

以下所討論的是把消去論--即我賦予 Hentzelt 的表述\footnote{Kurt Hentzelt, Zur Theorie der Polynomideale und Resultanten. 由 E. Noether 改寫。Math. Ann. 88 (1922), pp. 53--79；引用作 H.-N.。}的那種算術形式，並且可以稱為多項式域中的理想論的形式--納入一般理想論\footnote{E. Noether, Idealtheorie in Ringbereichen, Math. Ann. 83 (1921), pp. 24--66；引用作 Ideal Theory。}之中；同時，也要為涉及零點的部分重新奠基。兩個理論的基本概念，以及域論的基本概念\footnote{E. Steinitz, Algebraische Theorie der Körper, J. f. M. 137 (1910), pp. 167--309；引用作 Steinitz。}，都在 § 1 中簡要彙集，因此這裡可以引用它們。

這種納入和重新奠基圍繞四個問題展開：維數概念的算術表述；零點理論；範數與初等因子的分解定理同一般理想論中的分解定理之間的聯繫；用範數和初等因子形式刻畫素理想與準素理想。

\emph{維數概念}的算術表述（§ 4）由素理想鏈給出；這就是“曲面上有曲線，曲線上有點”這一事實的算術等價物。這個算術表述將被證明同消去論中的參數定義一致；它只需作輕微修改便可轉移到任意環域上，並且同若干其他問題的轉移一起，在那裡給出對理想結構的精確理解；這一點我將在別處說明。

在 H.-N. 中，\emph{零點}問題照常是直接針對給定理想處理的，也就是回到逐次消去；這樣仍不能完全避免檢驗性的特徵。與此相對的是一元多項式中通常採用的道路：把給定多項式表示為素函數冪（準素函數）的乘積，並且對每一個素函數分別構造零點域，使它同剩餘類域同構。素函數的次數給出該域的次數；而與素函數有相同零點的準素函數的次數，則給出模該準素函數的剩餘類環的次數。若對每一個素函數再轉到 Galois 域，便得到線性因子的分解；這樣，對原給定多項式來說，零點、分解為線性因子以及重數的問題都得到解決。這裡完全相應地（§ 3），把模素理想的剩餘類系通過形成商擴張為剩餘類域，並構造同它同構的零點域。這個零點域含有一個由附加不定元--例如 $x_{i+1},\ldots,x_n$--生成的子域；範數的次數給出相對於這個子域的次數；同時，通過轉到 Galois 域，得到初等因子形式、因而也得到範數分解為線性因子。對於相應的準素理想，零點則相同；分解也一併給出，因為範數成為素理想的初等因子形式的一個冪（§ 2）；範數的次數給出模該準素理想的剩餘類環相對於由不定元給出的子域的次數。

這樣，任意理想的零點問題就化歸為範數與初等因子的分解定理同一般理想論的分解定理之間的聯繫（§ 5）。一個理想的零點由它各個關聯素理想的零點組成；與此相應，範數分解為最大的準素因子，而這些因子等於各個關聯素理想的初等因子形式的冪；\footnote{這種消去論同一般理想論之間的平行，在 Kronecker 的消去論中並不成立；比較 H.-N., note *）。}由此一般地完成了範數分解為線性因子。在 H.-N. 中，重數則通過基本理想及其分量來解釋：範數的一個最大準素因子的次數，被等同於--在附加 $x_{i+1},\ldots,x_n$ 之後--$(i-1)$ 級基本理想模 $i$ 級基本理想某一分量的線性無關剩餘類的個數。隨後證明，$(i-1)$ 級基本理想同由所有維數高於 $(n-i)$ 的素理想唯一確定的分解的孤立分量相同；而 $i$ 級基本理想的一個分量，同由對應於範數因子的那個素理想以及所有較高維數的素理想所確定的孤立分量相同。相應範數因子的次數--在附加 $x_{i+1},\ldots,x_n$ 之後--等於該分量的剩餘類中只由 $(i-1)$ 級基本理想的類構成的那個子環的次數。

素理想和真準素理想的刻畫（§ 6）作為考察剩餘類域的直接結果而得到。若原來的係數域是完美域，則素理想對應於作為範數和初等因子形式的素函數，真準素理想對應於真準素函數，反之亦然。若係數域是非完美域，則只能給出必要條件或充分條件；例子表明，一個素理想的範數可以成為真準素函數，而一個真準素理想的初等因子形式可以成為素函數。更確切地說，在特徵 $p$ 中，素理想的範數成為其初等因子形式的 $p^f$ 次冪（$f>0$）；而對於初等因子形式為素函數的真準素理想，範數可以成為任意冪，例子同樣說明了這一點。最後考察絕對素理想（§ 7），即當係數域擴張到代數閉域時仍保持為素理想的那些理想。對於係數屬於（有限）代數數域的絕對素理想，這一刻畫導致 A. Ostrowski 首先為絕對素函數陳述的定理的轉移：模數域的每個素理想時，素理想性質仍保持，至多有有限多個例外。

還要指出，最後一個問題同代數數域中理想論的類比。這樣的理想的初等因子形式應看作可被該理想整除的最小整數有理數（注 10），所以對於素理想總是成為素數；反過來，只要其範數為素數，理想便成為素理想；證明也完全平行。因此，上述例子說明，在非完美域上會出現高於一次的素理想和分歧理想的類似物；而在完美域上，只有一次素理想而沒有分歧理想。

\subsection*{§ 1. 消去論、一般理想論與域論的基本概念}

\textbf{1. 域。} 令 $\bar{\mathfrak{J}}$ 表示所有以抽象定義的域 $\mathsf{P}$ 中元素為係數、關於 $y_1,\ldots,y_n$ 的多項式的整環域。令 $y$ 經受一個以不定元 $u_{\mu\nu}$ 為係數的變換：\footnote{考慮到 § 3 及其後各處，這個變換比 H.-N. 中的稍一般；其中所有結果因而更仍然成立。}
\begin{align}
 y_1&=u_{11}x_1+\cdots+u_{1n}x_n,\quad \ldots,\quad
 y_n=u_{n1}x_1+\cdots+u_{nn}x_n, \tag{1}
\end{align}
其逆為
\begin{align}
 x_1&=t_{11}y_1+\cdots+t_{n1}y_n,\quad \ldots,\quad
 x_n=t_{1n}y_1+\cdots+t_{nn}y_n. \tag{1'}
\end{align}
把 $u_{\mu\nu}$，或者等價地由於彼此有理可表而把 $t_{\mu\nu}$，附加到 $\mathsf{P}$ 上，從而得到係數域 $\mathsf{P}(u)=\mathsf{P}(t)$；令 $\frakS'$ 表示所有以 $\mathsf{P}(u)$ 中元素為係數、關於 $x_1,\ldots,x_n$ 的多項式的整環域。消去論以域 $\mathfrak{J}$ 和 $\frakS'$ 為基礎；考察 $\mathfrak{J}$ 中的理想 $\fraka,\frakb,\ldots$ 以及 $\frakS'$ 中的理想 $\fraka',\frakb',\ldots$。任意環中的理想在這裡按通常方式定義：與 $\alpha,\beta$ 同在時也含有 $\alpha-\beta$，並且與 $\alpha$ 同在時也含有 $\lambda\alpha$，其中 $\lambda$ 是任意環元素；$\frako$ 總表示由所有元素組成的單位理想。

\textbf{2. 變換理想。} $\frakS'$ 中的理想 $\frakm$ 稱為變換的，如果它由 $\mathfrak{J}$ 中的一個理想 $\barfrakm$ 通過 (1) 並在附加 $u_{\mu\nu}$ 或 $t_{\mu\nu}$ 後得到。因此，當 $\bar f(y)$ 或 $\bar g(y)$ 遍歷 $\barfrakm$ 中所有多項式時，$\frakm$ 由所有線性組合
\[
 f(x)=\sum U_\lambda(u)\bar f_\lambda(y)=\sum U_\lambda(u)f_\lambda(x)
\]
以及
\[
 g(x)=\sum T_\lambda(t)\bar g_\lambda(y)=\sum T_\lambda(t)g_\lambda(x)
\]
組成。特別地，所有 $f_\lambda(x)=\bar f_\lambda(y)$ 都在其中；或者等價地，所有 $g_\lambda(x)=\bar g_\lambda(y)$ 都在其中。由於 $f(x)$ 或 $g(x)$ 只在相差 $\mathsf{P}(u)=\mathsf{P}(t)$ 中量的意義下確定，不失一般性可把 $U_i,T_i$ 取為 $u$ 或 $t$ 的冪積。若把 $U,T$ 替換為任意其他冪積，特別是把 (1) 的列、相應地 (1') 的行作置換，則多項式 $f(x),g(x)$ 仍變為 $\frakm$ 中的多項式；因此 $\frakm$ 與任一多項式同含所有由置換 $x$ 得到的多項式，只要在依賴於 $u$ 或 $t$ 的係數中同時實行相應置換。以下出現的所有理想，除非明確說明相反，都應視為變換理想。非變換理想通過把 (1) 同
\[
 x_i=v_{i1}z_1+\cdots+v_{in}z_n
\]
複合，變成以 $\mathsf{P}(u,v)$ 為係數的 $z$ 的多項式域中的變換理想；而對於變換理想，這個複合只相當於把 $u_{\mu\nu}$ 替換為 $u,v$ 的雙線性組合。

\textbf{3. 記號。} 以下始終把 $f^{(i)},g^{(i)},\ldots,a^{(i)},b^{(i)},\ldots$ 理解為 $\frakS'$ 中不含 $x_i,\ldots,x_{i-1}$ 的多項式。

\textbf{4. 基本理想。} 對每個理想 $\frakm$ 指派 $n$ 個基本理想 $\frakg_0,\frakg_1,\ldots,\frakg_{n-1}$，規定如下：$(i-1)$ 級基本理想 $\frakg_{i-1}$ 含有且只含有那些多項式 $G(x)$，對它們存在一個多項式 $b^{(i)}\ne0$--一般隨 $G(x)$ 而變--使得 $b^{(i)}G(x)\equiv0\pmod{\frakm}$。由理想基的存在還推出，至少存在一個對所有 $G(x)$ 固定的 $B^{(i)}$，使
\begin{equation}
 B^{(i)}\frakg_{i-1}\equiv0(\frakm),\qquad B^{(i)}\ne0\tag{2}
\end{equation}
變換理想的基本理想本身也成為變換理想（H.-N., Theorem VI）。$i$ 級基本理想 $\frakg_i$ 也可定義為這樣的理想：在把 $x_{i+1},\ldots,x_n$ 附加到 $\mathsf{P}(u)$ 上時，$\frakm$ 變入其中；這裡可限制於在 $x_{i+1},\ldots,x_n$ 中為整的多項式，這時並不失一般性。$\frakg_0$ 總等於單位理想 $\frako$。

\textbf{5. 理想的模表示、初等因子和個別範數。} 如果把每個多項式 $f(x)$ 看作 $x_1,\ldots,x_{i-1}$ 的冪積 $\xi$ 中的線性形式，其係數為多項式 $a^{(i)}$，那麼理想 $\frakm$ 轉化為相對於 $a^{(i)}$ 的域的線性形式模 $\frakM_{i-1}$；也就是說，$\frakM_{i-1}$ 與任意兩個線性形式同含它們的差，並且與線性形式 $l(\xi)$ 同含 $a^{(i)}l(\xi)$，其中 $a^{(i)}$ 任意。這樣的模 $\frakM$ 的基本模，是所有滿足 $b^{(i)}g(\xi)\equiv0\pmod{\frakM}$、$b^{(i)}\ne0$ 的線性形式 $g(\xi)$ 的全體；因此，基本理想 $\frakg_{i-1}$ 在模表示下變成 $\frakM_{i-1}$ 的基本模 $\frakG_{i-1}$。$\frakM_{i-1}$ 與 $\frakG_{i-1}$ 依賴於無窮多個不定元 $\xi$，但每一個單獨線性形式中只出現有限多個這樣的不定元。$\frakM_{i-1}$ 具有如下特徵性質：把 $x_{i+1},\ldots,x_n$ 附加到 $\mathsf{P}(u)$ 後，它只有有限多個不同於一的初等因子；也就是說，經過這種附加後，$\frakG_{i-1}$ 與 $\frakM_{i-1}$ 允許一個基表示--其中 $\zeta$ 表示新的不定元，它們由在 $x_i$ 中為整的可逆線性變換同 $\xi$ 相連，並且每個 $\zeta_i=r_i(\xi)$ 中每次只出現有限多個這樣的不定元--
\begin{align}
\frakG_{i-1}&=(\zeta_0,\zeta_1,\ldots,\zeta_\sigma,\zeta_{\sigma+1},\ldots,\zeta_{\sigma+\nu},\ldots), \notag\\
\frakM_{i-1}&=(E^{(i)}\zeta_0,E_1^{(i)}\zeta_1,\ldots,E_\sigma^{(i)}\zeta_\sigma,
\zeta_{\sigma+1},\ldots,\zeta_{\sigma+\nu},\ldots), \tag{3}
\end{align}
其中每個 $E_\mu^{(i)}$ 整除前一個（H.-N. (24) 與 (28)）。

最高初等因子 $E^{(i)}$ 也定義為 (2) 中出現的所有 $B^{(i)}$ 的最大公因子--在多項式意義下；它可假定在 $x_{i+1},\ldots,x_n$ 中為整且本原，於是它關於 $x_i$ 正則，即若它關於 $x$ 的次數為 $\lambda$，則 $E^{(i)}$ 含有非零係數的項 $x_i^\lambda$。

(3) 中出現的所有初等因子的乘積 $R^{(i)}$ 稱為 $\frakG_{i-1}$ 相對於 $\frakM_{i-1}$ 的範數，也稱為理想 $\frakm$ 的個別範數；用符號表示，並考慮到當把 $x_{i+1},\ldots,x_n$ 附加時 $\frakm$ 變入 $\frakg_i$，有
\begin{equation}
 R^{(i)}=E^{(i)}E_1^{(i)}\cdots E_\sigma^{(i)}
 =N(\frakG_{i-1}\mid\frakM_{i-1})=N(\frakg_{i-1}\mid\frakg_i).
\end{equation}
如 (3) 所示，在附加 $x_{i+1},\ldots,x_n$ 後，$R^{(i)}$ 等於從 $\frakG_{i-1}$ 到 $\frakM_{i-1}$ 的過渡代換的行列式，而 $R^{(i)}$ 的次數給出 $\frakG_{i-1}$ 模 $\frakM_{i-1}$、因而 $\frakg_{i-1}$ 模 $\frakg_i$ 中相對於 $\mathsf{P}(u,x_{i+1},\ldots,x_n)$ 線性無關的剩餘類個數；$R^{(i)}$ 可假定在 $x_{i+1},\ldots,x_n$ 中為整且本原，於是關於 $x_i$ 正則。$R^{(i)}$ 可算作某個總是存在的秩為 $\varrho$ 的模 $\frakM^*_{i-1}$ 的所有 $\varrho$ 階行列式的最大公因子--在多項式意義下；該模具有性質：剩餘類系 $\frakG^*_{i-1}\mid\frakM^*_{i-1}$ 同 $\frakG_{i-1}\mid\frakM_{i-1}$ 同構，其中 $\frakG^*_{i-1}$ 理解為 $\frakM^*_{i-1}$ 的基本模（H.-N., Theorem VII and § 5）。

\textbf{6. 理想的初等因子形式與範數（結式形式）。} 所有最高初等因子的乘積 $E_\frakm=E^{(1)}E^{(2)}\cdots E^{(n)}$ 稱為初等因子形式，所有個別範數的乘積 $R_\frakm=R^{(1)}R^{(2)}\cdots R^{(n)}$ 稱為 $\frakm$ 的範數（結式形式）。\footnote{在 H.-N. 中只使用結式形式這一名稱；範數這一名稱對應於算術性質。}初等因子形式和範數都被理想整除；範數被初等因子形式整除，而後者的某個冪被範數整除。更一般地還有：
{\csname tagsleft@true\endcsname\begin{align}
 E^{(i)}\frakg_{i-1}&\equiv0(\frakg_i),
& E^{(n)}\cdots E^{(i)}\frakg_{i-1}&\equiv0(\frakm)\quad\text{從而}\notag\\
 R^{(i)}\frakg_{i-1}&\equiv0(\frakg_i),
& R^{(n)}\cdots R^{(i)}\frakg_{i-1}&\equiv0(\frakm),\quad\text{(H.-N., Theorem VIII),} \tag*{\raisebox{8.5pt}[0pt][0pt]{(4)}}
\end{align}}
並且還成立：若 $\frakm$ 被 $\frakn$ 整除，且範數 $R_\frakm$ 與 $R_\frakn$ 相同，則理想 $\frakm$ 與 $\frakn$ 也相同（H.-N., Theorem IX）。考慮到，對於 $\frakm$ 的一個約理想，$R^{(1)},R^{(2)},\ldots$ 的相同導致基本理想 $\frakg_1,\frakg_2,\ldots$ 相同，並且總有 $\frakg_0=\frako$，相應地得到：若 $\frakn$ 是 $\frakm$ 的真約理想，則 $\frakn$ 的第一個與 $\frakm$ 相應個別範數不同的個別範數成為真因子。\footnote{另一方面，$R_\frakn$ 後面的因子可以成為 $R_\frakm$ 相應因子的倍數；例如當 $\frakm$ 是素理想而 $\frakn$ 不是單位理想時總是如此。}若 $\frakm$ 含有一個多項式 $G^{(i)}\ne0$，則按定義 $\frakg_0,\frakg_1,\ldots,\frakg_{i-1}$ 都等於單位理想；於是由第 5 點所述個別範數的性質，即決定線性無關剩餘類的個數，可知 $R^{(1)},\ldots,R^{(i-1)}$ 都等於一；因而 $E^{(1)},\ldots,E^{(i-1)}$ 也都等於一。

\textbf{7. 個別範數和初等因子的分解；基本理想的分量。} 素函數是指相對於 $\mathsf{P}(u)$ 不可約的多項式；準素函數是素函數的冪。真準素函數是同時不是素函數的準素函數，也就是涉及高於一次的冪。把一個多項式分解為最大準素因子，是指每個因子都是準素函數，而任意兩個因子的乘積不再是準素的分解。\footnote{素函數和準素函數的定義也可完全類比 11 來表述，只需把“理想”替換為“多項式”。最大準素因子是 12 中分解的最大準素分量的類似物。}與個別範數 $R^{(i)}=N(\frakg_{i-1}\mid\frakg_i)$ 分解為最大準素因子 $Q_\lambda^{(i)}$ 相對應，有 $\frakg_i$ 作為最小公倍數的表示
\[
 \frakg_i=[\frakr_1,\frakr_2,\ldots,\frakr_s],
\]
使得 $Q_\nu^{(i)}\frakg_{i-1}\equiv0\pmod{\frakr_\nu}$；這裡 $\frakr_\nu$ 以 $\frakg_{i-1}$ 為其 $(i-1)$ 級基本理想，並且同其 $i$ 級基本理想相同；這些 $\frakr_\nu$ 稱為基本理想的分量。$\frakr_\nu$ 的最高初等因子等於 $Q_\nu^{(i)}$ 與 $E^{(i)}$ 的最大公因子--在多項式意義下--因而等於 $E^{(i)}$ 的一個最大準素因子。$\frakr_\nu$ 由它們關於基本理想的上述行為，並由“個別範數或最高初等因子應分別為 $R^{(i)}$ 或 $E^{(i)}$ 的最大準素因子”這一要求唯一確定（H.-N., Theorem X）。

\textbf{8. 維數概念。} 若 $R^{(i)}$ 是第一個不同於一的個別範數--或按 6 等價地，若 $\frakg_i$ 是第一個不同於單位理想的基本理想--則給理想 $\frakm$ 指派最高維數 $n-i$；給單位理想 $\frako$ 指派維數 $-1$。若只有一個不同於一的個別範數 $R^{(i)}$，也就是說 $\frakg_i=\frakg_{i+1}=\cdots=\frakm$，則最高維數也簡稱為 $\frakm$ 的維數。若 $\frakm$ 含有一個多項式 $G^{(i)}\ne0$，則由 6 末尾的說明可知，它的最高維數至多為 $n-i$。若 $\frakm$ 的最高維數為 $n-i$，且 $\frakn$ 是 $\frakm$ 的一個約理想，則 $\frakn$ 的最高維數也至多為 $n-i$。

\textbf{9. 範數（結式形式）與零點。} 理想 $\frakm$ 的零點，是指所有屬於 $\mathsf{P}(u;x_{i-1},\ldots,x_n)$ 的適當代數擴張域的 $x_1,\ldots,x_i$ 的值系（$i=1,2,\ldots,n$），使得--在把 $x_{i+1},\ldots,x_n$ 附加到係數域後--$\frakm$ 中所有多項式都消失。在這種附加下，範數（結式形式）的每個因子 $R^{(i)}$ 對應於 $\frakm$ 的若干零點，其個數由 $R^{(i)}$ 的次數給出；並且當 $R^{(i)}$ 寫成 $u_{\mu\nu}$ 與進一步不定元 $v_{\mu\nu}$ 的雙線性組合 $w_{\mu\nu}$ 時，允許顯式分解
\begin{equation}
 R^{(i)}=\prod\{z-(v_1\bar{x}_{1\nu}+\cdots-v_{i-1}\bar{x}_{i-1,\nu}-v_i\bar{x}_{i,\nu})\}^{\lambda_\nu},\notag
\end{equation}
其中 $\bar{x}_{1\nu},\ldots,\bar{x}_{i\nu}$ 表示一個相應的零點系（H.-N., Theorem XIII；新的奠基將在 § 3 給出）。因此，在最高維數為 $n-i$ 的理想的零點中，存在依賴於 $n-i$ 個參數的零點，但不存在依賴更多參數的零點；這證明了第 8 點給出的維數概念同通常表述一致。

\textbf{10. 一般理想論的環域。} 基本域應是一個交換環，其中有限鏈定理成立：每一條理想鏈 $\fraka,\fraka_1,\ldots,\fraka_\nu,\ldots$，其中每個 $\fraka_\lambda$ 是前一個理想的真約理想--約理想按理想意義理解--都在有限步後終止。這個要求等價於另一要求，即每個理想都有理想基（Ideal Theory, Theorem I），因此特別在多項式域中成立。在以下 11、12、13 中，都假定這個一般環域。

\textbf{11. 素理想與準素理想。} 如果一個乘積被 $\frakp$ 整除便推出至少一個因子被 $\frakp$ 整除，則稱理想 $\frakp$ 為素理想；如果一個乘積被 $\frakq$ 整除便推出一個因子被 $\frakq$ 整除，或每一個因子的某個冪被 $\frakq$ 整除，則稱 $\frakq$ 為準素理想--或簡稱準素。用符號表示：
\[
 a\not\equiv0\pmod{\frakp},\quad b\not\equiv0\pmod{\frakp}\quad\text{推出}\quad ab\not\equiv0\pmod{\frakp};
\]
\[
 a\not\equiv0\pmod{\frakq},\quad b^x\not\equiv0\pmod{\frakq}\ \text{對每個冪 }x\quad\text{推出}\quad ab\not\equiv0\pmod{\frakq}.
\]
每個準素理想 $\frakq$ 有且僅有一個關聯素理想 $\frakp$，它是 $\frakq$ 的一個約理想，並且它的某個冪被 $\frakq$ 整除，$\frakp\ne\frako$，$\frakp^\rho\equiv0\pmod{\frakq}$；這裡 $\frakp$ 由所有其某個冪被 $\frakq$ 整除的環元素組成。由定義可見，每個素理想同時也是準素理想；真準素理想指那些同時不是素理想的準素理想，因此 $\rho>1$。

\textbf{12. 分解定理。} 把理想表示為最小公倍數 $\fraka=[\frakq_1,\ldots,\frakq_\alpha]$；若沒有任何 $\frakq$ 被其他理想的最小公倍數--它的補--所包含，則稱該表示最短；若每個 $\frakq$ 都準素，但任意兩個 $\frakq$ 的最小公倍數不準素，則稱它是由最大準素分量給出的表示。每個理想都允許一個由有限多個最大準素分量給出的最短表示；兩個不同的這種表示中，分量個數和相應的素理想都相同，且這些素理想兩兩不同。這樣唯一確定的素理想稱為該理想的素理想，或關聯素理想（Ideal Theory, Theorem IX）。

\textbf{13. 分解的孤立分量。} 若理想 $\frakr=[\frakq_1,\ldots,\frakq_i]$ 中的 $\frakq_i$ 在 $\fraka$ 的某個由最大準素分量給出的最短表示中出現，並滿足條件：它們的關聯素理想沒有一個包含在 $\fraka$ 的其他素理想中，則稱 $\frakr$ 為 $\fraka$ 的分解的孤立分量。分解的孤立分量由其素理想唯一確定；特別地，孤立的最大準素分量唯一確定（Ideal Theory, Theorem XIII）。

\textbf{14. 特徵、完美域與非完美域、第一類擴張。} 若域 $\Omega$ 中由單位導出的素域是有理數域類型，則稱 $\Omega$ 具有特徵零；若這個素域是模素數 $p$ 的剩餘類系類型，則稱其具有特徵 $p$。若 $\Omega$ 上係數的每個素函數都在適當擴張域中分解為互異線性因子，則稱 $\Omega$ 為完美域；否則稱為非完美域。每個特徵零域都是完美的；特徵 $p$ 的域完美當且僅當它與每個元素同含其 $p$ 次根。非完美域中的素函數是關於 $x^{p^f}$ 的 $r$ 次整函數，並在適當擴張域中分解為 $r$ 個互異線性因子的 $p^f$ 次冪；$f$ 稱為指數。若素函數在適當擴張域中分解為互異線性因子，即指數 $f$ 為零，則稱其為第一類；若擴張中的每個元素都是第一類元素，即某個第一類素函數的零點，則稱該擴張為第一類擴張。由附加一個第一類元素得到的每個擴張都是第一類；在完美域上只有第一類擴張（Steinitz, § 11 and § 13）。

\textbf{15. 有限擴張的約化次數與指數。} 若 $\mathcal K$ 是非完美域 $\Omega$ 的有限擴張，則 $\mathcal K$ 含有一個次數為 $r$ 的子域 $\mathcal K_0$，它相對於 $\Omega$ 為第一類，並且由 $\mathcal K$ 中所有且僅有的第一類元素組成。$\mathcal K$ 中每個元素的 $p^f$ 次冪都屬於 $\mathcal K_0$，並且 $\mathcal K$ 中存在恰好次數為 $rp^f$ 的元素。$r$ 稱為約化次數，$f$ 稱為該有限擴張的指數（Steinitz, § 14, 1 and § 14, 5）。

\subsection*{§ 2. 素理想與準素理想的初等因子形式和範數；素理想的約理想}

從現在起，我們始終考察 § 1, 1 中定義的多項式域。首先要證明，對於素理想和準素理想，也可以限制於變換理想，即

\textbf{引理 I.} $\mathfrak{J}$ 中素理想 $\bar{\frakp}$ 的變換理想 $\frakp$ 仍成為素理想；$\mathfrak{J}$ 中準素理想 $\bar{\frakq}$ 的變換理想 $\frakq$ 仍成為準素理想。換言之，在把 $u_{\mu\nu}$ 附加到 $\mathsf{P}$ 時，素理想或準素理想這一性質得以保持。

事實上，設 $\fraka\not\equiv0\pmod{\frakp}$，$\frakb\not\equiv0\pmod{\frakp}$，其中 $\fraka$ 和 $\frakb$ 不必是變換理想；比如取 $\fraka$ 和 $\frakb$ 中的多項式 $a(x)\not\equiv0\pmod{\frakp}$ 與 $b(x)\not\equiv0\pmod{\frakp}$。按 (1') 逆轉 (1)，並不失一般性把 $T$ 理解為 $T_{\mu\nu}$ 的冪積，可得
\[
 a(x)=\sum T_i a_i(y);\qquad b(x)=\sum T_i' b_i(y),
\]
且至少有一個 $a_i(y)$ 和一個 $b_i(y)$ 不被 $\bar{\frakp}$ 整除。設例如在某個 $T$ 的詞典序中，$a_h(y)\not\equiv0\pmod{\bar{\frakp}}$ 與 $b_k(y)\not\equiv0\pmod{\bar{\frakp}}$ 分別為最高項；那麼也有 $a_h(y)b_k(y)\not\equiv0\pmod{\bar{\frakp}}$，從而 $a(x)b(x)\not\equiv0\pmod{\frakp}$，因而 $\fraka\frakb\not\equiv0\pmod{\frakp}$，這證明 $\frakp$ 是素的。這個證明顯然依賴於這樣一個事實：模素理想的剩餘類系形成無零因子的環，而附加不定元保持這一性質。

相應地，設 $\fraka\not\equiv0\pmod{\frakq}$ 且對每個冪 $x$ 都有 $\frakb^x\not\equiv0\pmod{\frakq}$；則由理想基的存在可知，必存在 $\fraka$ 中的 $a(x)$ 和 $\frakb$ 中的 $b(x)$，使 $a(x)\not\equiv0\pmod{\frakq}$ 且對每個冪 $x$ 都有 $b(x)^\varkappa\not\equiv0\pmod{\frakq}$；否則，$\frakb$ 的某個冪將被 $\frakq$ 整除。

令 $a(x)b(x)=c(x)$，並令 $\fraka^*,\frakb^*,\frakc^*$ 分別表示由 $a(x),b(x),c(x)$ 的係數 $a_i(y),b_j(y),c_k(y)$ 導出的理想。則對每個冪 $x$ 都有 $\frakb^{*x}\not\equiv0\pmod{\bar{\frakq}}$，否則 $b(x)^\varkappa$ 將被 $\frakq$ 整除；由於 $\fraka^*\not\equiv0\pmod{\bar{\frakq}}$，還必有對每個冪 $\lambda$ 都有 $\fraka^*\frakb^{*\lambda}\not\equiv0\pmod{\bar{\frakq}}$。但由 Dedekind--Mertens 定理\footnote{比較 H.-N., p. 63 and note 10）。}存在指數 $\lambda$，使 $\fraka^*\frakb^{*\lambda}\equiv \frakc^*\frakb^{*\lambda-1}$；於是 $\frakc^*\not\equiv0\pmod{\bar{\frakq}}$。這給出 $c(x)\not\equiv0\pmod{\frakq}$，從而 $\fraka\frakb\not\equiv0\pmod{\frakq}$，證明 $\frakq$ 是準素的。

在把理想表示為最小公倍數時，也可以限制於變換理想，即

\textbf{引理 II.} 由 $\mathfrak{J}$ 中的表示 $\barfrakm=[\frakc_1,\ldots,\frakc_s]$ 可推出 $\frakS'$ 中的表示 $\frakm=[\frakc_1',\ldots,\frakc_s']$，其中 $\frakc_i'$ 表示 $\frakc_\lambda$ 的變換理想。特別地，在每個由最大準素分量給出的最短表示中，這些分量可假定為變換的。

事實上，$\frakm$ 被 $[\frakc_1',\ldots,\frakc_s']$ 整除，因為 $\frakm$ 中的每個多項式 $f(x)$--必要時乘以 $U(u)$ 的適當冪之後--具有形式 $\sum U_i(u)f_i(y)$，其中每個 $f_i(y)$ 作為 $\barfrakm$ 中的多項式按假設被所有 $\frakc_\lambda$ 整除。反過來，若 $f(x)$ 被每個 $\frakc_i'$ 整除，則它具有上述形式，因而被 $\frakm$ 整除。因此，若從 $\mathfrak{J}$ 中最大準素分量的分解出發，就得到 $\frakS'$ 中的分解 $\frakm=[\frakq_1',\ldots,\frakq_s']$；這裡 $\frakq_i'$ 作為準素理想的變換理想，由引理 I 是準素的，並且它們是最大準素分量，因為 $\mathfrak{J}$ 中不同的關聯素理想 $\bar{\frakp}$ 對應於 $\frakS'$ 中不同的 $\frakp$。

現在需要用初等因子形式和範數對素理想與準素理想作第一個刻畫，尚不完全分離素情形和準素情形。完全類比於代數數域中的理想論，\footnote{代數數域中理想的最高初等因子，應看作被該理想整除的最小整數有理數，因而特別地，對於素理想是被它整除的素數，對於準素理想（素理想的冪）是素數冪。事實上，每個理想 $\fraka^*$ 也是場基元素 $\omega_1,\ldots,\omega_n$ 中的線性形式模 $\mathfrak{A}^*$，其中 $(\omega_1,\ldots,\omega_n)$ 是 $\mathfrak{A}^*$ 的基本模；因此，被 $\fraka^*$ 整除的最小整數有理數成為該模 $\mathfrak{A}^*$ 的最高初等因子，而 $\fraka^*$ 的範數成為 $\mathfrak{A}^*$ 的所有初等因子的乘積。}有

\textbf{定理 I.} 素理想的初等因子形式等於一個素函數，其範數等於該素函數的一個冪。對於準素理想，初等因子形式與範數成為準素函數，即同一個素函數的冪；這個素函數本身是相應關聯素理想的初等因子形式。對於素理想和準素理想，最高維數同純粹維數一致；並且準素理想的這個維數同其關聯素理想的維數一致。

定理 I 對單位理想顯然成立，其初等因子形式和範數都等於一。現在令素理想 $\frakp$ 的最高維數為 $n-i$；由 § 1, 6 的 (4) 得到
\[
 E^{(i)}\frakg_{i-1}\not\equiv0\pmod{\frakp},\quad\text{但}\quad
 E^{(i)}\frakg_i\equiv0\pmod{\frakp},
\]
否則由 § 1, 8 可知 $\frakp$ 的最高維數至多為 $n-i-1$。因此 $\frakg_i\equiv0(\frakp)$，所以 $\frakp$ 同它的 $i$ 級基本理想相同，因而也同 $\frakg_{i+1},\frakg_{i+2},\ldots$ 相同。於是 $R^{(i+1)},\ldots,R^{(n)}$ 等於一，而作為這些範數的因子，$E^{(i+1)},E^{(i+2)},\ldots$ 也等於一；所以初等因子形式等於 $E^{(i)}$，而它必須等於一個素函數 $P^{(i)}$。因為由 § 1, 5 並注意到 $\frakg_{i-1}$ 成為單位理想，$E^{(i)}$ 被定義為所有被 $\frakp$ 整除的 $B^{(i)}$ 的最大公因子--在多項式意義下；若 $E^{(i)}=C_1^{(i)}C_2^{(i)}\equiv0(\frakp)$，則 $E^{(i)}$ 必須包含在因子 $C_1^{(i)}$ 或 $C_2^{(i)}$ 之一中。又因為 $E^{(i)}$ 的某個冪被 $R^{(i)}$ 整除，$R^{(i)}$ 成為這個素函數 $P^{(i)}$ 的一個冪--可能就是一次冪；同時 $R^{(i)}$ 成為 $\frakp$ 的範數。由於只有一個不同於一的個別範數出現，$\frakp$ 的最高維數最終同其純粹維數一致。

相應地，令準素理想 $\frakq$ 的最高維數為 $n-i$；則如上有
\[
 (E^{(i)}\frakg_{i-1})^x\not\equiv0\pmod{\frakq},\quad\text{但}\quad
 (E^{(i)}\frakg_i)^x\equiv0\pmod{\frakq}
\]
對每個冪 $x$ 成立；於是同樣得到 $\frakq=\frakg_i$，其初等因子形式為 $E^{(i)}$，其範數為 $R^{(i)}$，且最高維數為純粹維數。這個維數同關聯素理想的維數一致；因為由 $\frakp^\rho\equiv0\pmod{\frakq}$、$\frakq\equiv0\pmod{\frakp}$ 可知，$\frakp$ 的維數 $n-i$ 至多等於 $\frakq$ 的維數，而 $\frakq$ 的維數至多等於 $\frakp^\rho$ 的維數；由 $(P^{(i)})^\rho\equiv0(\frakp^\rho)$，其中 $P^{(i)}$ 表示 $\frakp$ 的初等因子形式，可知後者最高維數至多為 $n-i$，於是 $n-i$ 就是 $\frakp$ 與 $\frakq$ 的維數。又由 $(P^{(i)})^\rho\equiv0(\frakq)$ 推出 $E^{(i)}$--作為 $\frakq$ 中所有 $B^{(i)}$ 的最大公因子--成為 $P^{(i)}$ 的一個冪，也可能是一次冪；$R^{(i)}$ 也是如此。由此證明定理 I 的全部部分。

用維數概念刻畫素理想由下述定理給出：

\textbf{定理 II.} 素理想沒有同一最高維數的真約理想；反過來，每個具有這個性質的理想都是素理想。

證明依賴於

\textbf{引理 III.} 如果以域 $\Omega$ 中元素為係數的多項式素理想 $\frakp$ 只有有限多個相對於 $\Omega$ 線性無關的剩餘類，則 $\frakp$ 沒有不同於單位理想的真約理想；換言之，模 $\frakp$ 的剩餘類系形成一個域。

假設說明存在有限數 $k$，使任意 $(k+1)$ 個剩餘類之間都有以 $\Omega$ 中元素--更確切地說，以由 $\Omega$ 中元素表示的剩餘類域 $(\Omega)$ 中元素--為係數的線性關係；但至少存在一個由 $k$ 個相對於 $\Omega$ 線性無關的剩餘類組成的系。立即推出，所有剩餘類都能由任一選定的 $k$ 個線性無關剩餘類線性表示。現在令 $\fraka$ 為 $\frakp$ 的真約理想，並令 $a(z)\not\equiv0(\frakp)$ 是 $\fraka$ 中的多項式；由 $c_1f_1(z)+\cdots+c_kf_k(z)\equiv0\pmod{\frakp}$ 得
\[
 c_1a(z)f_1(z)+\cdots+c_ka(z)f_k(z)\equiv0\pmod{\frakp}.
\]
這意味着，與 $f_1,\ldots,f_k$ 的剩餘類一樣，$af_1,\ldots,af_k$ 的剩餘類也構成一個 $k$ 個線性無關元素的系，通過它們特別可以線性表示單位類。因此 $\fraka$ 是單位理想；換言之，剩餘類系形成域，因為對 $a(z)\not\equiv0(\frakp)$ 總存在 $f(z)$ 使 $1\equiv a(z)f(z)(\frakp)$。

為證明定理 II 的第一部分，只需再注意：每個維數為 $(n-i)$ 的素理想--更一般地，每個最高維數為 $(n-i)$ 的理想--相對於 $\mathsf{P}(u;x_{i+1},\ldots,x_n)$ 只有有限多個線性無關剩餘類，即由 § 1, 5 可知，其數目為 $R^{(i)}$ 的次數，因為這裡 $\frakg_{i-1}$ 成為單位理想。因此，$\frakp$ 的每個真約理想 $\fraka$，在把 $x_{i+1},\ldots,x_n$ 附加到 $\mathsf{P}(u)$ 時成為單位理想；換言之，$\fraka$ 的 $i$ 級基本理想成為單位理想，這意味着 $\fraka$ 的最高維數至多為 $n-i-1$。為了使斷言也適用於作為約理想的非變換理想 $\fraka$，只需按 § 1, 2 轉到新的變換域。

反過來，設 $\frakp$ 沒有同一最高維數 $n-i$ 的真約理想，並令 $\fraka\not\equiv0(\frakp)$ 且 $\frakb\not\equiv0(\frakp)$，其中 $\fraka$ 與 $\frakb$ 仍假定為變換理想。由假設，$(\fraka,\frakp)$ 與 $(\frakb,\frakp)$ 作為最大公因子--按理想意義--成為 $\frakp$ 的真約理想，因此有
\[
 E^{(n)}\cdots E^{(i+1)}\frakg_i\equiv0\pmod{(\fraka,\frakp)},
\qquad
 E^{(n)}\cdots E^{(i+1)}\frakg_i\equiv0\pmod{(\frakb,\frakp)}.
\]
這給出
\[
 E^{(n)}\cdots E^{(i+1)}\frakg_i\equiv0\pmod{(\fraka\frakb,\frakp)},
\]
因而必有 $\fraka\frakb\not\equiv0\pmod{\frakp}$，否則 $\frakp$ 的最高維數會小於 $n-i$。由於 $\fraka,\frakb$ 是變換理想，$\frakp$，以及由引理 I 得到的 $\bar{\frakp}$，都是素理想。定理 II 得證。

\subsection*{§ 3. 素理想與準素理想的零點}

首先對於素理想，通往零點理論的直接奠基（§ 1, 9）的過渡由下述引理構成；它推廣引理 III，並適用於任意環中的素理想。

\textbf{引理 IV.} 模每個不同於單位理想的素理想的剩餘類系，都可以通過形成商擴張為一個域，即該素理想的剩餘類域。

模任意理想的剩餘類系形成一個環，因為和、差和積仍屬於該系，並且結合律、交換律和分配律在轉到剩餘類時保持。若該理想特別是素理想，則這個環沒有零因子；因為素理想的定義（§ 1, 11）說明，非零剩餘類的乘積也不為零。若該素理想不同於單位理想，則這個環至少含有一個非零元素，因此可以通過形成商--附加元素對--擴張為一個域；\footnote{Steinitz, § 3。原整環域中存在非零元素的假設已經足夠，因為這時單位的存在由形成商保證，從而原域成為該域的子域；Steinitz 假設整環域中單位存在。}這樣剩餘類域被定義。

首先固定以下一直保留的記號。

\textbf{記號。} 一般用把任一代表元置於括號中的方式表示剩餘類，$(x),\ldots,(a(x)),\ldots$；同樣，剩餘類的域也用括號表示。特別地，$(P)$、$(\mathsf{P}(u))$、$(\mathsf{P}(u,x_{i+1},\ldots,x_n))$ 表示所有且僅有由 $P$、$\mathsf{P}(u)$、$P(u,x_{i+i},\ldots,x_n)$ 中元素表示的剩餘類所組成的剩餘類域。

還回憶如下內容。若域 $\mathcal R$ 相對於基域 $\Omega$ 含有至少一個由 $k$ 個在 $\Omega$ 上代數無關的量--超越量--組成的系，而 $\mathcal R$ 中任意 $k+1$ 個量在 $\Omega$ 上代數相關，則稱 $\mathcal R$ 相對於 $\Omega$ 的代數秩（超越次數）為 $k$。若 $\mathcal R$ 可表示為某個子域 $\Omega(z_1,\ldots,z_s)$ 的代數擴張域，其中 $z$ 在 $\Omega$ 上代數無關--超越--則必有 $s=k$；\footnote{關於這一熟知事實在任意特徵下成立的證明，比較 Steinitz, § 22, 9。--附加 $u$ 時，代數秩不能增加，因為只涉及原元素以 $\Omega(u)$ 中係數作有理組合；也不能降低，因為元素之間以 $\Omega(u)$ 中係數給出的每個關係，都分解為有限多個以 $\Omega$ 中係數給出的關係。}同樣，當 $u$ 表示不含於 $\mathcal R$ 的不定元時，$\mathfrak{K}(u)$ 相對於 $\Omega(u)$ 的代數秩也等於 $k$。

零點域的構造現在完全類比於一元素函數的通常道路，\footnote{Steinitz, § 6。}即

\textbf{定理 III.} 對維數為 $n-i$ 的素理想 $\frakp$ 的剩餘類域 $(\mathcal R)$，可以同構地指派一個 $\frakp$ 的零點域 $R$；它的代數秩為 $n-i$--依賴於 $n-i$ 個參數--並且特別可看作 $\mathsf{P}(u,x_{i-1},\ldots,x_n)$ 的有限擴張域。$(\mathcal R)$ 與 $R$ 之間的同構包含 $(\mathsf{P})$ 與 $\mathsf{P}$、$(\mathsf{P}(u))$ 與 $\mathsf{P}(u)$ 之間的同構，並且當 $x_{i+1},\ldots,x_n$ 被取作參數時，也包含 $(\mathsf{P}(u,x_{i+1},\ldots,x_n))$ 與 $\mathsf{P}(u,x_{i+1},\ldots,x_n)$ 之間的同構。反過來，每個代數秩為 $n-i$ 的零點域--依賴於 $n-i$ 個參數--都同剩餘類域 $(\mathcal R)$ 同構。

定理 III 作為推論給出熟知定理：若理想 $\fraka$ 在素理想 $\frakp$ 的所有零點中消失，則 $\fraka$ 被 $\frakp$ 整除。

為證明這一點，首先注意剩餘類域 $(\mathcal R)$ 相對於 $(\mathsf{P}(u))$ 的代數秩為 $n-i$。事實上，類 $(x_{i+1}),\ldots,(x_n)$ 相對於 $(\mathsf{P}(u))$ 代數無關，因為 $\frakp$ 的維數為 $n-i$，不含不含 $x_1,\ldots,x_i$ 的多項式；而任一其他類都依賴於這些類。因為由初等因子形式 $P^{(i)}$ 屬於 $\frakp$ 可知，按 § 1, 2，對 $\lambda=1,2,\ldots,i$ 也有屬於 $\frakp$ 的多項式，分別依賴於 $x_\lambda,x_{i+1},\ldots,x_n$。因此 $(\mathcal R)$ 成為 $(\mathsf{P}(u,x_{i+1},\ldots,x_n))$ 的有限擴張域；由 § 1, 5--並注意到 $\frakg_{i-1}$ 等於單位理想而 $\frakg_i$ 由定理 I 等於 $\frakp$--知它在該子域上的次數等於範數的次數。

為了轉到同構的零點域，還要注意：$\mathsf{P}(u,x_{i+1},\ldots,x_n)$ 與 $(\mathsf{P}(u,x_{i+1},\ldots,x_n))$ 通過把 $\mathsf{P}(u,x_{i+1},\ldots,x_n)$ 的每個元素指派給由該元素表示的剩餘類而同構地相關。首先，在這個指派下，由 $x_{i+1},\ldots,x_n$ 的多項式組成的整環域同由 $(x_{i+1}),\ldots,(x_n)$ 的多項式組成的整環域分別同構；因為 $\frakp$ 不含不含 $x_1,\ldots,x_i$ 的多項式，所以來自 $\mathsf{P}(u,x_{i+1},\ldots,x_n)$ 的不同多項式對應於不同類；而從同構整環域通過形成商得到同構域。$\mathsf{P}(u,x_{i+1},\ldots,x_n)$ 與 $(\mathsf{P}(u,x_{i+1},\ldots,x_n))$ 之間的同構關係包含 $\mathsf{P}$ 與 $(\mathsf{P})$、$\mathsf{P}(u)$ 與 $(\mathsf{P}(u))$ 之間的同構。為把此同構擴展為 $(\mathcal R)$ 與零點域 $R$ 之間的同構，把某些新元素 $\xi_1,\ldots,\xi_i$ 同構地指派給類 $(x_1),\ldots,(x_i)$；也就是說，每個多項式 $(a(x))$ 對應於多項式 $a(\bar{x}_1,\ldots,\bar{x}_i,x_{i+1},\ldots,x_n)$，並且和與積對應於和與積。通過形成商--其中可限制於來自子域 $(\mathsf{P}(u,x_{i+1},\ldots,x_n))$ 或 $\mathsf{P}(u,x_{i+1},\ldots,x_n)$ 的分母--剩餘類域 $(\mathcal R)$ 於是對應於一個域 $R$，它是 $\mathsf{P}(u,x_{i+1},\ldots,x_n)$ 的有限擴張，次數為範數的次數；它可稱為零點域。這個指派說明，只要 $f(x)\equiv0\pmod{\frakp}$，就有 $f(\xi_1,\ldots,\xi_i,x_{i+1},\ldots,x_n)$ 消失。在整個論證中，$(\mathsf{P}(u,x_{i+1},\ldots,x_n))$ 也可以替換為 $(\mathcal R)$ 中由附加 $n-i$ 個代數無關量所得的任意其他子域。

反過來，很容易證明每個代數秩為 $n-i$ 的零點域 $R$ 都同剩餘類域同構。由於 $R$ 可由量 $\xi_1,\ldots,\xi_i$ 以 $\mathsf{P}(u)$ 中係數有理導出，因此這些量中必須有 $n-i$ 個相對於 $\mathsf{P}(u)$ 代數無關、因而超越的元素。特別地，這必須對 $\xi_{i+1},\ldots,\xi_n$ 成立；因為由 $P^{(i)}\equiv0\pmod{\frakp}$，如上可知 $R$ 是 $\mathsf{P}(u,\xi_{i+1},\ldots,\xi_n)$ 的代數擴張域。由於按假設，對 $\frakp$ 中每個多項式 $f(x)$ 都有 $f(\xi)$ 消失，$(\mathcal R)$ 中所含的 $(x)$ 多項式域的每個類，對應於且只對應於 $R$ 中一個關於 $\xi$ 的多項式元素；和與積對應於和與積。但不同的類對應於 $R$ 中不同的元素；否則一個非零類，並且經適當乘法後一個由 $(\mathsf{P}(u,x_{i+1},\ldots,x_n))$ 中多項式表示的類，會對應於 $R$ 中的零元素，因而同假設相反，量 $\xi_{i+1},\ldots,\xi_n$ 將滿足一個代數相關關係。這個指派窮盡 $R$ 中所有關於 $\xi$ 的多項式；通過形成商--仍可限制於來自 $(\mathsf{P}(u,x_{i+1},\ldots,x_n))$ 或 $\mathsf{P}(u,\bar{x}_{i+1},\ldots,\bar{x}_n)$ 的分母--從這些同構整環域得到同構域。

現在令理想 $\fraka$ 在 $\frakp$ 的所有零點中消失，特別也在所有依賴於 $n-i$ 個參數的零點中消失。若 $a(x)$ 是 $\fraka$ 中任一多項式，則 $a(\xi)$ 消失；由於 $R$ 與 $(\mathcal R)$ 同構，類 $(a(x))$ 就是零類，因而 $a(x)$，於是 $\fraka$，被 $\frakp$ 整除。這證明了定理 III 及其推論。

為了從定理 III 過渡到 § 1, 9 中給出的分解，首先針對素理想，併為了能對指數作更精確的陳述，需要進一步研究初等因子形式，即

\textbf{引理 V.} 若在特徵 $p$ 中，某個素理想或準素理想--更一般地，某個最高維數與純粹維數一致的理想 $\fraka$--的初等因子形式 $E^{(i)}$ 關於 $x_i$ 具有指數 $f$，因而是 $x_i^{p^f}$ 的整函數，則它關於 $x_{i+1},\ldots,x_n$ 以及 $E^{(i)}$ 中出現的不定元 $u_{\mu\nu}$ 或 $t_{\mu\nu}$ 也具有指數 $f$。特別地，若係數域是完美域，則素理想的初等因子形式 $P^{(i)}$ 關於 $x_i$ 總具有指數零，因此是第一類素函數。

為此要注意：當 $\mathsf{P}$ 是特徵 $p$ 的完美域時，$\mathsf{P}(u,x_{i+1},\ldots,x_n)$ 成為非完美域，所以不能從 § 1, 14 直接推出 $P^{(i)}$ 為第一類。設 $E^{(i)}$ 關於 $x_i$ 的指數為 $f$，但關於另一個不定元 $x_{i+\lambda}$ 的指數為 $f'<f$。那麼--通過置換 $u_{\mu\nu}$ 或 $t_{\mu\nu}$--由 § 1, 2 也存在一個非零且被該理想整除的多項式 $G^{(i)}$，它是 $x_i^{p^{f'}}$ 的整函數，並且按初等因子形式的定義被 $E^{(i)}$ 整除。但由於 $G^{(i)}$ 與 $E^{(i)}$ 次數相同，它只能同 $E^{(i)}$ 相差一個來自 $\mathsf{P}(u)$ 的因子；因此必有 $f'=f$。於是 $E^{(i)}$--可以假定它關於 $t_{\mu\nu}$ 為整且本原--具有形式
\[
 E^{(i)}=\sum c_\lambda(t)\,x_i^{\lambda_i p^f}\cdots x_n^{\lambda_n p^f}
       =\sum c_\lambda(t)\,g_\lambda(y^{p^f},t^{p^f})
       =\sum T_\mu(t)\,\bar h_\mu(y^{p^f}),
\]
其中 $T_\mu$ 表示 $t$ 的冪積，而 $\bar h$ 被 $\bar{\fraka}$ 整除。因此，若在 $T_\mu$ 中把每個單獨的 $t$ 的指數替換為其中所含的 $p^f$ 的最大倍數，就得到一個同樣屬於 $\fraka$ 的多項式；由上述表示可見，這等價於在 $c_\lambda(t)$ 中把每個 $t$ 的指數替換為其中所含的 $p^f$ 的最大倍數。經此替換，$E^{(i)}$ 變為一個關於 $t,x_i,\ldots,x_n$ 的多項式；按定義，它被 $E^{(i)}$ 整除；由於所假定的本原性，這種整除性關於 $t$ 也成立。因為它關於任何自變量的次數都沒有增加，所以它必與 $E^{(i)}$ 恆等。特別地，若 $\mathsf{P}$ 是完美域，則 $E^{(i)}$ 作為 $x^{p^f},t^{p^f}$ 的整函數，是一個 $p^f$ 次冪；因此若所論為素理想，必有 $f$；初等因子形式 $P^{(i)}$ 成為第一類素函數。

分解現在由下述定理給出：

\textbf{定理 IV.} 素理想的初等因子形式 $P^{(i)}$ 在由零點域導出的 Galois 域中允許顯式分解
\[
 P^{(i)}=\prod_\nu (x_i-t_{i+1}x_{i+1}^{(\nu)}-\cdots-t_nx_n^{(\nu)})^e
 =\prod_\nu (t_i(y_i-y_i^{(\nu)})+\cdots+t_n(y_n-y_n^{(\nu)}))^e,
\]
其中 $y_1^{(\nu)},\ldots,y_n^{(\nu)}$ 表示原不定元 $y$ 的一個關聯零點系，與 $t_i,\ldots,t_n$ 無關；不同的 $\nu$ 對應於不同的線性因子；指數 $e$ 在 $P$ 完美時等於一，在 $P$ 非完美時等於 $p^f$（$f\geqq0$）。由此，在附加新的不定元 $v$ 後，把 $u_{\mu\nu}$ 換為 $u$ 與 $v$ 的雙線性組合而寫成的初等因子形式的分解也被確定：
\[
 P^{(i)}(w)=\prod_\nu (z_i-z_i^{(\nu)})^e
 =\prod_\nu (v_{i1}(y_1-y_1^{(\nu)})+\cdots+v_{in}(y_n-y_n^{(\nu)}))^e,
\]
其中 $e$ 的意義同上。同樣，素理想的範數，或以 $\frakp$ 為關聯素理想的準素理想的範數--作為 $P^{(i)}$ 的冪--的分解也同時給出。

作為定理 IV 的推論，還得到：若兩個素理想的初等因子形式 $P^{(i)}$ 相同，則它們相同。

證明定理 IV 歸結為證明初等因子形式 $P^{(i)}$ 同擴張域 $(\mathsf{P}(u,x_{i+1},\ldots,x_n),(x_i))$ 在 $(\mathsf{P}(u,x_{i+1},\ldots,x_n))$ 上的基本方程相同。首先，為了解對不定元 $u_{\mu\nu}$ 或 $t_{\mu\nu}$ 的依賴，轉到 $\frakp$ 的剩餘類域 $(\mathcal S)$；按代數秩定義中所說，它相對於 $(P)$ 必有代數秩 $n-i$，這裡 $(P)$ 表示由 $\mathsf{P}$ 中元素表示的子域。因為 $(\mathcal S)$ 由把不定元 $u$ 或 $t$ 附加到一個同構於 $(\mathcal R)$ 的子域而得，該子域從所有且僅有由 $\mathfrak{J}$ 中多項式表示的類通過形成商得到，並且 $(P)$ 與 $(P)$ 彼此對應。因而 $(\mathcal S)$ 相對於 $(\mathsf{P}(u))$ 的代數秩等於這個子域相對於 $(P)$ 的代數秩，進而由同構等於 $(\mathcal R)$ 相對於 $(P)$ 的代數秩。由於 $(\mathcal S)$ 可由類 $(y_1),\ldots,(y_n)$ 以 $(P)$ 中係數有理導出，所以這些類中必須有 $n-i$ 個代數無關的類，並且 $(\mathcal S)$ 作為把這些 $n-i$ 個類附加到 $(P)$ 後生成的子域的有限代數擴張域出現。\footnote{這個說明表明，定理 III 及其推論本可直接為 $\frakp$ 的零點域 $R$ 陳述並證明。只有通過轉到變換理想，才能使任意理想中等維數的最大準素分量在零點中出現相同參數；也只有這樣，任意理想的範數和初等因子形式之間的聯繫才產生。}

若只附加 $x_{i+1},\ldots,x_n$ 中出現的那些不定元 $t_{\mu\nu}$，則再次得到一個剩餘類域，它含有類 $(y_1),\ldots,(y_n)$ 與 $(x_{i+1}),\ldots,(x_n)$，併成為 $(\mathsf{P}(t_{\mu\nu},x_{i+1},\ldots,x_n))$ 的有限代數擴張域；這些域同 $(\mathcal S)$ 的子域同構，但不與它們相同，因為對於不同的不定元，剩餘類雖由同樣的元素表示，卻並不相同，因為它們不含同樣的元素。因此，在類 $(y)$ 對 $(P(t_{\mu\nu},x_{i+1},\ldots,x_n))$ 的依賴中，只進入不定元 $t_{\mu\nu}$。現在通過附加 $t_{i1},\ldots,t_{in}$ 形成類 $(x_i)=(t_{i1}y_1+\cdots+t_{in}y_n)$；令 $(x_{i\nu})=(t_{i1}y_{1\nu}+\cdots+t_{in}y_{n\nu})$ 為相對於 $(P(t_{\mu\nu},x_{i+1},\ldots,x_n))$ 位於 Galois 域中的 $r$ 個互異共軛值。於是--視是否處理第一類擴張而定--所有因子 $x_i-(x_{i\nu})$ 的乘積，或其 $p^f$ 次冪的乘積，是關於 $(P(t_{\mu\nu},x_{i+1},\ldots,x_n))$ 的一個素函數 $(G)$；因為既然存在次數為 $r\cdot p^f$ 的元素（§ 1, 15），$(x_i)$ 必定就是這樣的一個元素，因為每個元素都由對 $t_{\mu\nu}$ 作特殊化從 $(x_i)$ 得出；而 $(x_i)$ 是 $(G)$ 的零點。轉到 $\frakp$ 的同構零點域以及由它導出的 Galois 域後，$G$ 因而允許分解為線性因子 $t-t_{1i}\bar y_{1\nu}-\cdots-t_{ni}\bar y_{n\nu}$，或這些因子的 $p^f$ 次冪。現在，由於 $(G)$ 在 $t=(x_i)$ 時消失，$G(x_i)$ 被 $\frakp$ 整除，因而被 $P^{(i)}$ 整除；而 $G$ 作為相對於 $P(t_{\mu\nu},x_{i+1},\ldots,x_n)$ 的素函數並且關於 $x_i$ 本原，必定同 $P^{(i)}$ 相同。這樣證明了定理 IV 的第一個分解。

為了轉到第二個分解，注意擴張是否為第一類這一性質在附加不定元時保持。因此，在附加 $v_1,\ldots,v_i$ 和所有 $t_{\mu\nu}$ 後，形成類 $(z)=(v_1x_1+\cdots+v_ix_i)$ 以及共軛類 $(z_\nu)=(v_1x_{1\nu}+\cdots+v_ix_{i\nu})$。於是所有 $t-(z_\nu)$ 的乘積，或其 $p^f$ 次冪的乘積，又是一個素函數 $(H)$，它在 $t=(z)$ 時消失。因此 $H$ 如上可表示為線性因子的乘積，並且同由把變換 (1) 與 $z=v_1x_1+\cdots+v_nx_n$ 複合而從 $\frakp$ 得到的素理想的初等因子形式相同；這個初等因子形式又--由於相互特殊化--由把 $P^{(i)}$ 中的 $u_{\mu\nu}$ 替換為 $u$ 與 $v$ 的某些雙線性組合得到。\footnote{H.-N., § 7。}由這些分解，也給出了範數作為 $P^{(i)}$ 的冪的分解，\footnote{對第一類擴張，特別是在完美域上，涉及一次冪；在非完美域上涉及 $p^g$ 次冪，這將在 § 6 的定理 XIII 與 XIV 中證明。}並同樣給出了以 $\frakp$ 為關聯素理想的準素理想的範數與初等因子形式的分解。

最後，若兩個素理想在初等因子形式 $P^{(i)}$ 上相同，則由上述分解可知它們的零點相同；因此它們彼此整除，從而相同。

\subsection*{§ 4. 維數概念的算術表述}

一個不依賴引入變換理想的維數概念的第一表述由下述定理給出：

\textbf{定理 V.} \emph{素理想的維數由其剩餘類域的代數秩給出。一個理想的最高維數等於其關聯素理想的維數數目中的最高者。}

定理 V 的第一部分已經在定理 IV 的第一條說明中證明了，那裡指出 $\frakp$ 的剩餘類域 $(\mathcal R)$ 的代數秩等於 $\frakp$ 的剩餘類域 $(\mathcal R)$ 的代數秩，因而等於 $\frakp$ 的維數，同 § 1, 8 一致。為證明第二部分，令 $\frakm=[\frakq_1,\ldots,\frakq_\alpha]$ 是由最大準素分量給出的最短表示。首先，作為 $\frakm$ 的約理想，沒有一個 $\frakq$ 能具有比 $\frakm$ 更高的維數。但所有 $\frakq$ 的乘積被 $\frakm$ 整除，因而所有單個 $\frakq$ 的初等因子形式的乘積也被 $\frakm$ 整除。若 $n-i$ 是這些 $\frakq$ 的維數數目中的最高者，則 $\frakm$ 含有一個多項式 $G^{(i)}\ne0$，因此不可能有高於 $n-i$ 的維數。由定理 I，$\frakq$ 的維數數目同其關聯素理想的維數數目一致。因此，如果不引入變換理想，而把素理想的維數定義為其剩餘類域的代數秩，則定理 V 的第二部分給出任意理想的最高維數的定義。

另一方面，為了通過素理想鏈把握維數概念--同樣不依賴引入變換理想--需要在定理 II 的補充下證明：

\textbf{定理 VI.} \emph{每個不同於單位理想的素理想--在必要時附加不定元之後--至少有一個素理想作為約理想，其維數恰好降低一個單位。}

事實上，若 $\frakp$ 的維數為 $(n-i)>0$，且 $v$ 表示新的不定元，則
\[
 \frakc=(\frakp,x_i-v)
\]
是具有所需性質的理想。這由一個同構對應推出。可寫 $\frakc=(\frakp^*,x_i-v)$，其中 $\frakp^*$ 由把 $\frakp$ 中的 $x_i$ 替換為不定元 $v$ 並把 $v$ 附加到係數域而得；剩餘類也同樣由把類 $(x_i)$ 替換為 $(v)$ 得到。但當把 $x_i$ 附加到 $\mathsf{P}(u)$ 時，$\frakp$ 變入自身；因為由
\[
 h(x_i)f(x)\equiv0\pmod{\frakp},\qquad h(x_i)\ne0
\]
必然得到 $f(x)\equiv0\pmod{\frakp}$。否則 $\frakp$ 將含有一個只依賴於 $x_i$ 的多項式，從而也含有一個只依賴於 $x_n$ 的多項式；這同假設矛盾，因為這會使它的維數至多為零。因此，在指派 $v\sim x_i$ 下，模 $\frakc$ 的零類必然對應於模 $\frakp$ 的零類，反之當然也如此。於是模 $\frakp$ 的剩餘類系與模 $\frakc$ 的剩餘類系之間存在同構；後者沒有零因子，所以 $\frakc$ 是素理想。在此同構下，$(x_i)$ 與 $(\mathsf{P}(u,x_{i+1},\ldots,x_n))$ 之間的代數相關對應於 $(x_{i+1}),\ldots,(x_n)$ 相對於 $(\mathsf{P}(u,v))$ 的一個關係；而對 $\lambda=1,\ldots,i-1$，$(x_\lambda)$ 與 $(\mathsf{P}(u,x_{i+1},\ldots,x_n))$ 之間的相關仍保持為這樣的相關，因此不再進一步降低代數秩。代數秩恰好降低一個單位。若 $\frakp$ 的維數為零，則 $\frakc$ 成為單位理想；上述結果仍然有效。因此在原域 $\bar{\mathfrak{J}}$ 中，
\[
 \bar{\fraka}=(\bar{\frakp},v+v_1y_1+\cdots+v_ny_n),
\]
其中以 $v_\lambda$ 代替 $-t_{\lambda i}$，就是所需種類的理想。特別地，若 $\mathsf{P}$ 含有無窮多個元素，則總可以把 $v,v_1,\ldots,v_n$ 特殊化為 $\mathsf{P}$ 中元素，使範數和初等因子形式、因而 $\fraka$ 的維數，等於特殊化理想 $\frakb$ 的相應量；\footnote{比較 H.-N., § 7 的最後一段。}不過初等因子形式不必仍為素函數。由定理 V，$\frakb$ 仍至少有一個所需維數的關聯素理想；再回到 $\bar{\mathfrak{J}}$，定理 VI 因而在那裡無需附加任何不定元也成立。

定理 II 和定理 VI 立即給出所需的鏈式表述：

\textbf{定理 VII.} \emph{素理想 $\frakp$ 的維數為 $n-i$，當且僅當--在必要時附加不定元后--至少存在一條由 $1+(n-i)+1$ 個素理想組成的鏈}
\[
 \frakp_0=\frakp,\quad \frakp_1,\ldots,\frakp_{n-i},\quad \frakp_{n-i+1},
\]
\emph{其中每個都是前一個的真約理想，而不存在成員數更多的這種鏈。當 $\mathsf{P}$ 含有無窮多個元素時，這個定義在原域中、且不引入任何不定元時也成立。}\footnote{這個維數的鏈式定義，在代數數域的情形中，給通常的素理想維數零，給單位理想維數 $-1$，給零理想維數一。事實上，後者也滿足素理想的定義：在域中，非零量的乘積總不等於零。按這個維數定義，定理 II 在代數數域中仍成立。}

首先顯然，鏈的最後一個成員必須等於單位理想--它滿足素理想的定義--否則可通過附加 $\frako$ 使鏈加長；對於 $\frako$ 本身，定理 VII 給出維數 $-1$，同 § 1, 8 一致。現在令 $\frakp$ 不同於 $\frako$ 且維數為 $n-i$。鏈中成員數至多為 $1+(n-i)+1$，因為由定理 II，兩個同維數的素理想不能出現。反過來，由定理 VI，在必要時附加新不定元后，至少存在一條這樣的長度的鏈，因為在鏈的每個成員處維數恰好降低一個單位。若把定理 VII 作為維數的定義--這個定義顯然可以在原域中完全同樣地陳述，並且當 $\mathsf{P}$ 含有無窮多個元素時由定理 VI 可知不再需要引入不定元--則任意理想的最高維數定義又由定理 V 的第二部分給出。

\subsection*{§ 5. 基本理想及其分量在一般理想論中的歸類；理想的零點}

問題在於用一般理想論意義下分解的孤立分量（§ 1, 13）來刻畫基本理想及其分量（§ 1, 7）。這種刻畫將顯示範數分解定理同一般理想論分解定理之間的平行，並使得可以為任意理想陳述 § 3 中為素理想和準素理想給出的零點理論。對於基本理想有：

\textbf{定理 VIII.} \emph{$i$ 級基本理想 $\frakg_i$ 是分解的孤立分量，它由維數為 $n-1,n-2,\ldots,n-i$ 的所有關聯素理想唯一確定。若所有關聯素理想具有相同維數，則理想本身也具有這個維數；也就是說，範數中只出現一個因子 $R^{(i)}$。一般地，範數中不同於一的因子 $R^{(i)},R^{(i+\sigma)},R^{(i+\tau)},\ldots$ 的個數，等於關聯素理想中出現的不同維數數目的個數。}

證明之前先給出：

\textbf{引理 VI.} \emph{若理想表示為最小公倍數 $\frakm=[\fraka_1,\ldots,\fraka_\alpha]$，則它的基本理想 $\frakg_i$ 是 $\fraka$ 的 $i$ 級基本理想 $\frakh_i$ 的最小公倍數。}

由
\[
 b^{(i+1)}f\equiv0\pmod{\frakm},\qquad b^{(i+1)}\ne0
\]
可得
\[
 b^{(i+1)}f\equiv0\pmod{\fraka_\lambda},\qquad b^{(i+1)}\ne0;
\]
因此 $\frakg_i$ 被 $[\frakh_{i1},\ldots,\frakh_{i\alpha}]$ 整除。反過來，由
\[
 c_\lambda^{(i+1)}f\equiv0(\fraka_\lambda),\qquad c^{(i+1)}\ne0
\]
可得
\[
 c_1^{(i+1)}\cdots c_\alpha^{(i+1)}f\equiv0(\frakm),\qquad
 c_1^{(i+1)}\cdots c_\alpha^{(i+1)}\ne0,
\]
因而得到反向整除，引理得證。

現在把 $\frakm$ 的關聯素理想按維數排列（某些維數當然可以缺席，$j_\lambda=0$）：
\[
 \frakp_{1,1},\ldots,\frakp_{1,j_1};\quad
 \frakp_{2,1},\ldots,\frakp_{2,j_2};\quad \ldots\quad
 \frakp_{n,1},\ldots,\frakp_{n,j_n},
\]
其中一般 $\frakp_{\lambda,\kappa}$ 表示維數為 $n-\lambda$ 的素理想；令 $\frakq_{\lambda,\kappa}$ 為 $\frakm$ 的某個固定給定分解中出現的相應準素理想。於是按定義，對於 $\lambda>i$ 的所有 $\frakq_{\lambda,\kappa}$，其 $i$ 級基本理想等於單位理想；而對於 $\lambda\leqq i$，由定理 I，它等於 $\frakq_{\lambda,\kappa}$。因此由引理 VI，
\[
 \tag{5}
 \frakg_i=[\frakq_{1,1},\ldots,\frakq_{1,j_1};\ldots;\frakq_{i,1},\ldots,\frakq_{i,j_i}],
\]
這個表示把 $\frakg_i$ 識別為分解的孤立分量，因為屬於 $\frakg_i$ 的素理想沒有一個能包含在其餘素理想之一中，而那些素理想都具有較低維數。

特別地，若只出現維數為 $n-i$ 的關聯素理想，則由 (5) 有 $\frakg_0,\ldots,\frakg_{i-1}$ 等於單位理想，而 $\frakg_i=\frakg_{i+1}=\cdots=\frakm$；於是 $R_\frakm=R^{(i)}$。若出現不同維數的素理想，比如 $j_i,j_{i+\sigma},j_{i+\tau},\ldots$ 非零，則由 (5)
\[
 \frakg_0\ne \frakg_i\ne \frakg_{i+\sigma}\ne \frakg_{i+\tau}\ne\cdots,
\]
而
\[
 \frakg_0=\frakg_1=\cdots=\frakg_{i-1}=\frako,
 \qquad
 \frakg_i=\frakg_{i+1}=\cdots=\frakg_{i+\sigma-1},\ldots;
\]
因此且只有因此，因子 $R^{(i)},R^{(i+\sigma)},R^{(i+\tau)}$ 不同於一。\footnote{考慮引理 II 後，定理 VIII 又推出基本理想--作為變換理想的最小公倍數--成為變換理想。由於定理 VIII 只依賴於 H.-N. 中不使用這一事實的定理，這同時給出了一個揭示內在原因的新證明。在 H.-N. 中，該定理只用於這裡重新發展的零點理論。}

藉助定理 VIII，定理 I 加強為：

\textbf{定理 IX.} \emph{一個理想為準素，當且僅當其初等因子形式--因而其範數--是準素函數。}

條件對準素理想成立已經在定理 I 中說明。反過來，設 $\frakm$ 的初等因子形式為 $Q^{(i)}=(P^{(i)})^\rho$，其中 $P^{(i)}$ 表示素函數。由定理 VIII，$\frakm$ 的所有素理想都有同一維數 $n-i$。由 $\frakm=[\frakq_1,\ldots,\frakq_\alpha]$ 得到
\[
 Q^{(i)}\equiv0\pmod{\frakq_\lambda};
 \quad\text{故}\quad
 (P^{(i)})^\rho\equiv0\pmod{\frakp_\lambda},
 \quad\text{從而}\quad P^{(i)}\equiv0\pmod{\frakp_\lambda}.
\]
由於所有 $\frakp_\lambda$ 都有維數 $n-i$，且 $P^{(i)}$ 是素函數，$P^{(i)}$ 成為每個 $\frakp_\lambda$ 的初等因子形式。因此由定理 IV 的推論，所有 $\frakp_\lambda$ 都相同。於是，由於涉及的是最大準素分量，只能出現一個 $\frakp$；$\frakm$ 是準素的，其中 $\rho=1$ 的特殊素情形不被排除。定理 VIII 和 IX 導出基本理想分量的如下歸類。

\textbf{定理 X.} \emph{$\frakg_i$ 的分量 $\frakr_\lambda$，它對應於 $R^{(i)}$ 的最大準素因子 $Q_\lambda^{(i)}=(P_\lambda^{(i)})^{\rho_\lambda}$，由這樣的分解孤立分量給出：該孤立分量由維數為 $n-1,\ldots,n-i+1$ 的關聯素理想以及那個維數為 $n-i$ 且初等因子形式等於 $P_\lambda^{(i)}$ 的素理想唯一確定。關聯素理想與範數的最大準素因子一一對應。}

由 § 1, 7，$\frakr_\lambda$ 同它們的 $i$ 級基本理想一致，並以 $\frakg_{i-1}$ 為 $(i-1)$ 級基本理想。因此由定理 VIII 或 (5)，它們允許一個最短表示
\[
 \frakr_\lambda=[\frakg_{i-1},\frakq_{\lambda,1},\ldots,\frakq_{\lambda,t_\lambda}],
\]
其中所有 $\frakq$ 具有維數 $n-i$，並屬於不同的素理想。由 $Q_\lambda^{(i)}$ 的定義--並注意到 $\frakr_\lambda$ 同其 $i$ 級基本理想相同--有
\[
 Q_\lambda^{(i)}\frakg_{i-1}\equiv0\pmod{\frakq_{\lambda,\kappa}},
 \quad(\kappa=1,\ldots,t_\lambda),
\]
因而
\[
 (Q_\lambda^{(i)})^{\rho_\lambda}\frakr_\lambda\equiv0\pmod{\frakq_{\lambda,\kappa}};
 \quad\text{故}\quad
 (P_\lambda^{(i)})^{\rho_\lambda}\sigma_\lambda\equiv0\pmod{\frakp_{\lambda,\kappa}},
\]
所以 $P_\lambda^{(i)}\equiv0\pmod{\frakp_{\lambda,\kappa}}$。如在定理 IX 的證明中一樣，推出 $\frakr_\lambda$ 的表示中只能出現一個維數為 $n-i$ 的 $\frakq_\lambda$ 和相應 $\frakp_\lambda$，且 $\frakp_\lambda$ 的初等因子形式由 $P_\lambda^{(i)}$ 給出。剩下要證明 $\frakp_\lambda$ 是 $\frakm$ 的關聯素理想，並且 $\frakq_\lambda$ 的確出現在 $\frakm$ 的某個由最大準素分量給出的表示中。於是 $\frakr_\lambda$ 也被識別為分解的孤立分量，因為沒有素理想能夠被同維或較低維的素理想吸收；也就是說，它正是由 $\frakp_\lambda$ 以及所有較高維數的關聯素理想確定的分量。

為此只需證明 $\frakq_\lambda$ 出現在 $\frakg_i$ 的一個最短表示中，因為由定理 VIII，這樣的表示總可補全為 $\frakm$ 的最短表示。因此只需證明
\[
 \frakg_i=[\frakr_1,\ldots,\frakr_\alpha]=[\frakg_{i-1},\frakq_1,\ldots,\frakq_\alpha]
\]
是最短表示。若例如 $\frakq_\lambda$ 被其補吸收，使得 $\frakg_{i-1}\frakq_1\cdots\frakq_{\lambda-1}\frakq_{\lambda+1}\cdots\frakq_\alpha$ 被 $\frakq_\lambda$ 整除，那麼由於 $\frakg_{i-1}\not\equiv0\pmod{\frakq_\lambda}$，便推出某個 $\mu\ne\lambda$ 的 $\frakq_\mu$ 的某個冪被 $\frakq_\lambda$ 整除。於是同維數的關聯素理想 $\frakp_\mu$ 與 $\frakp_\lambda$ 必須由定理 II 相同；但這是不可能的，因為它們的初等因子形式 $P_\mu^{(i)}$ 與 $P_\lambda^{(i)}$ 按定義不同。$\frakg_i$ 的表示還表明，由定理 VIII，維數為 $n-i$ 的所有關聯素理想都對應於該表示，從而每一個這樣的素理想確實對應於範數的一個最大準素因子。一一對應因此得證。

定理 X 立即給出定理 IV 的轉移：

\textbf{定理 XI.} \emph{範數的各個最大準素因子 $Q_\lambda^{(i)}=(P_\lambda^{(i)})^{\rho_\lambda}$ 允許乘積表示}
\[
 Q_\lambda^{(i)}=\prod \bigl(x_i-t_{1i}\bar y_{1\nu}-\cdots-t_{ni}\bar y_{n\nu}\bigr)^{\delta_\lambda\rho_\lambda}
 =\prod \bigl(t_{1i}(y_1-\bar y_{1\nu})+\cdots+t_{ni}(y_n-\bar y_{n\nu})\bigr)^{\delta_\lambda\rho_\lambda},
\]
\emph{其中 $\bar y_{1\nu},\ldots,\bar y_{n\nu}$ 表示關聯素理想 $\frakp_\lambda$ 在原不定元 $y$ 中的一個關聯零點系，與 $t_{1i},\ldots,t_{ni}$ 無關，並且 $\delta_\lambda$ 的值為一或 $p^{f_\lambda}$。於是範數分解為線性因子就完全給出。$Q_\lambda^{(i)}$ 的次數給出由形成商得到的剩餘類子域 $(\mathsf{P}(u,x_{i+1},\ldots,x_n))$ 上的次數；所取對象是模孤立分量 $\frakr_\lambda$ 的剩餘類子環，該子環由來自 $\frakg_{i-1}$ 的剩餘類表示；對於準素理想，這因而就是所有剩餘類的環。當附加新的不定元 $v$ 時，還得到分解}
\[
 Q_\lambda^{(i)}(u,v)=\prod (z-v_1\bar x_{1\nu}-\cdots-v_i\bar x_{i\nu})^{\delta_\lambda\rho_\lambda}.
\]

由於 $P_\lambda^{(i)}$ 已由定理 X 識別為 $\frakp_\lambda$ 的初等因子形式，這實際上只是代入定理 IV 的分解。關於次數的說明，不過是 § 1, 5 和 7 中所說內容的另一種表述。\footnote{同樣，H.-N. 注 2 中提到的重數表述也是定理 X 的直接結果。因為，在附加 $x_{i+1},\ldots,x_n$ 後，剩餘類系 $\frakg_{i-1}\mid\frakr_\lambda$ 同 $\frakt_\lambda\mid\frakr_\lambda$ 同構，這裡置 $\frakt_\lambda=[\frakg_{i-1},\frakr_1,\ldots,\frakr_{\lambda-1},\frakr_{\lambda+1},\ldots,\frakr_\alpha]$；鑑於因子 $Q_\mu^{(i)}$ 的相對互素性，這由回到線性形式模（H.-N., Theorem V）推出。直接可以證明 $\frakt_\lambda\mid\frakr_\lambda$ 同 $\frakt_\lambda\mid\frakq_\lambda$ 同構；這裡 $\frakt_\lambda$ 由 $\frakq_\lambda$ 的補在附加 $x_{i+1},\ldots,x_n$ 後得到。}

\subsection*{§ 6. 用初等因子形式和範數刻畫素理想與真準素理想}

在定理 IX 中，準素理想由初等因子形式和範數刻畫，但方式是把素情形看作準素的特殊情形。現在需要把素理想與真準素理想分開；這些考察將同 § 3 的考察平行。必要條件或充分條件很容易陳述：

\textbf{定理 XII.} \emph{素理想的一個必要條件是其初等因子形式為素函數；一個充分條件是其範數為素函數。換言之，素理想的初等因子形式總是素函數，而真準素理想的範數總是真準素函數。}

素理想的初等因子形式為素函數已經在定理 I 中證明。現在設範數 $R_\frakq$ 等於一個素函數。要證明在此假設下，$\frakq$ 的每個真約理想 $\fraka$ 都有較低最高維數；於是由定理 II，$\frakq$ 被識別為素理想。事實上，由 § 1, 6，範數 $R_\fraka$ 的第一個與 $R_\frakq$ 的相應因子不同的因子是真因子。由於 $R_\frakq=P^{(i)}$，$R_\fraka$ 的因子 $R^{(i)}$ 必須等於 $P^{(i)}$ 的一個真因子，因而等於一；所以 $\fraka$ 有較低最高維數。

第二個簡單證明依賴於下述引理，下一定理也依賴於它。

\textbf{引理 VII.} \emph{由多項式 $H^{(i)}$ 所表示的、模某個準素理想的初等因子形式 $E^{(i)}$ 的剩餘類系--更一般地，模某個只有維數為 $n-i$ 的關聯素理想的理想的初等因子形式的剩餘類系--同模該理想的剩餘類的一個子系同構。初等因子形式的次數給出這個剩餘類系相對於商域 $(\mathsf{P}(u,x_{i+1},\ldots,x_n))$ 的次數。}

事實上，把由同一多項式 $H^{(i)}$ 表示的類互相對應，便定義了一個對應；和與積對應於和與積。這個指派也是一一的。因為所有屬於模該理想零類的多項式 $H^{(i)}$，按定義都被 $E^{(i)}$ 整除；所以零類對應於模 $E^{(i)}$ 的零類，反之亦然。

現在令一個理想的範數為素函數 $P^{(i)}$，因而同初等因子形式相同。於是模初等因子形式 $P^{(i)}$ 與模 $\frakq$ 的剩餘類系的次數相對於 $(\mathsf{P}(u,x_{i+1},\ldots,x_n))$ 相同；由引理 VII 同構於模 $P^{(i)}$ 的剩餘類系的子系因而窮盡模 $\frakq$ 的剩餘類。由於這個系沒有零因子，$\frakq$ 是素理想。

一般地，定理 XII 不能加強為必要且充分條件，這將由下面的例子說明。不過，當係數域 $\mathsf{P}$ 是完美域（§ 1, 14）時，這是可能的：

\textbf{定理 XIII.} \emph{若係數域 $\mathsf{P}$ 是完美域，則素理想的範數和初等因子形式都成為素函數，因而相同；真準素理想的範數和初等因子形式都成為真準素函數。因此，在完美域上，初等因子形式為素函數這一條件是素理想的必要且充分條件；在該條件中也可以用範數代替初等因子形式。}

考慮定理 XII 後，為證明定理 XIII，只需證明在完美域上，只要理想的初等因子形式為素函數，該理想就是素理想，並且這時範數等於該素函數。在 § 3 的引理 V 中已經證明，在完美域上，只要初等因子形式是素函數，它就是第一類素函數。因此令 $(\mathcal R)$ 表示通過把 $(x_i)$ 附加到 $(\mathsf{P}(u,x_{i+1},\ldots,x_n))$ 而得到的剩餘類域，也就是 $(\mathsf{P}(u,x_i,x_{i+1},\ldots,x_n))$ 模 $P^{(i)}$。則 $(x_i)$ 是第一類素函數 $(P^{(i)})$ 的零點，於是例如由 Lagrange 公式可知，$(y_1),\ldots,(y_n)$ 都包含在 $(\mathcal R)$ 中。按引理 VII，與 $(\mathcal R)$ 同構的、所考察理想 $\frakq$ 的模剩餘類子系--其中分母只來自 $(\mathsf{P}(u,x_{i+1},\ldots,x_n))$--因此含有剩餘類 $(y_1),\ldots,(y_n)$，從而窮盡模 $\frakq$ 的剩餘類系。由於它同構於 $(\mathcal R)$，沒有零因子，$\frakq$ 是素理想。$\frakq$ 的這個剩餘類域 $(\mathcal R)$ 相對於 $(\mathsf{P}(u,x_{i+1},\ldots,x_n))$ 的次數等於 $\frakq$ 的範數次數；另一方面，由於同構於 $(\mathcal R)$，它等於初等因子形式 $P^{(i)}$ 的次數。因此範數和初等因子形式必須相同。

對非完美系數域，定理 XII 仍允許如下加強：

\textbf{定理 XIV.} \emph{若係數域 $\mathsf{P}$ 是特徵 $p$ 的非完美域，則素理想的範數是初等因子形式的 $p^g$ 次冪；且 $0\leqq g\leqq(i-1)f$，其中 $f$ 表示初等因子形式的指數，$n-i$ 表示素理想的維數。}

定理 XIV 等價於斷言，$\frakp$ 的剩餘類域 $(\mathcal R)$ 相對於同構於 $(\mathcal R)$ 的子域 $(\mathcal R')=(\mathsf{P}(u,x_i,x_{i+1},\ldots,x_n))$ 具有次數 $p^g$（$g\geqq0$）。這直接由 § 1, 15 中陳述的 Steinitz 定理推出，更確切地由

\textbf{引理 VIII.} \emph{令 $\Lambda$ 是非完美域 $\Omega$ 的有限擴張，其約化次數為 $r$、指數為 $f$；令 $\Lambda_0$ 是 $\Lambda$ 中所含的第一類域，並令 $M$ 是介於 $\Lambda_0$ 與 $\Lambda$ 之間的中間域。則 $\Lambda$ 中每個元素相對於 $M$ 的次數都是 $p$ 的冪，其中 $p$ 表示 $\Omega$ 的特徵。}

因為 $\Lambda$ 中每個元素 $z$ 的 $p^{f'}$ 次冪（$f'\leqq f$）都屬於 $\Lambda_0$；所以 $z$ 是以 $\Lambda_0$ 中元素為係數的素函數 $G(t)=(t-z)^{p^{f'}}$ 的零點。令 $H(t)=(t-z)^\delta$ 為 $M$ 中以 $z$ 為零點的素函數。則 $H(t)$ 也以 $\Lambda_0(z)$ 中元素為係數，因此 $H(t)$ 的係數屬於 $M$ 與 $\Lambda_0(z)$ 的交域，也就是介於 $\Lambda_0$ 與 $\Lambda_0(z)$ 之間的中間域。因此 $z$ 相對於 $M$ 的次數等於相對於這個中間域的次數，並且作為 $p^{f'}$ 的因子，是 $p$ 的冪。

為證明定理 XIV，只需再注意：若 $r$ 表示約化次數，$f$ 表示 $(\mathcal R)$ 的指數，則域 $(\mathcal R')=(\mathsf{P}(u,x_i,x_{i+1},\ldots,x_n))$ 相對於 $(\mathsf{P}(u,x_{i+1},\ldots,x_n))$ 的次數必為 $r\cdot p^f$；因而 $(\mathcal R)$ 中所含的第一類域 $(\mathcal R_0)$ 必是 $(\mathcal R')$ 的子域。由於 $(\mathcal R)$ 中存在次數為 $r\cdot p^f$ 的元素（§ 1, 15），$(x_i)$ 必然具有這個次數，因為 $(\mathcal R)$ 的所有元素都通過對 $(x_i)$ 中的 $t_{\mu\nu}$ 作特殊化而得。$(\mathcal R)$ 的指數同初等因子形式的指數一致，而 $(x_i)^{p^f}$ 具有次數 $r$，並且是 $(\mathcal R_0)$ 的元素，因而是 $(\mathcal R_0)$ 的本原元。由於 $(\mathcal R)$ 由把有限多個元素--例如 $(x_1),\ldots,(x_{i-1})$--附加到 $(\mathcal R')$ 而得到，對引理 VIII 重複有限次應用便給出所需證明，同時給出對 $g$ 的估計。\footnote{比較 A. Ostrowski, Zur arithmetischen Theorie der algebraischen Größen, Gött. Nachr. 1919, pp. 279--298, 特別是 p. 288，其中相應定理沒有證明地陳述。}

現在補充例子，說明定理 XII 一般不能在定理 XIV 之外再加強。它們涉及一個範數等於初等因子形式平方（$g>0$）的素理想，以及初等因子形式為素函數的真準素理想。後一個例子還說明，此處沒有定理 XIV 的類似物，而範數可以成為初等因子形式的任意冪。

令係數域 $\mathsf{P}$ 由把兩個不定元 $\lambda,\mu$ 附加到模 $2$ 的剩餘類域而得到，因而 $\mathsf{P}$ 是特徵二的非完美域。考慮理想
\[
 \bar{\frakp}=(y_1^2+\lambda,\,y_2^2+\mu),
\]
它必為素理想，因為剩餘類系 $(\mathcal R)$ 同域 $\mathsf{P}(\sqrt\lambda,\sqrt\mu)$ 同構。域 $(\mathcal R)$ 相對於 $(\bar{\mathsf{P}})$ 的次數為四，約化次數 $r=1$，指數 $f=2$。因此初等因子形式 $P^{(2)}$ 是二次的，而範數是四次的，從而是 $P^{(2)}$ 的平方；這也可直接由模表示 (3), § 1, 5 檢驗。得到
\[
 P^{(2)}=(u_{11}u_{22}+u_{12}u_{21})^2x_2^2+(u_{11}^2\mu+u_{21}^2\lambda),
\]
或經適當乘法後，
\[
 x_2^2+(t_{12}^2\lambda+t_{22}^2\mu).
\]

現在令 $\mathsf{P}$ 由把一個不定元 $\lambda$ 附加到模 $2$ 的剩餘類域而得到，並以理想
\[
 \bar{\frakq}=(y_1^2+\lambda,\,y_2^2+\lambda)=(y_1^2+\lambda,\,(y_1+y_2)^2)
\]
為基礎。作為初等因子形式，由在剛才給出的式子中令 $\mu=\lambda$，得到
\[
 P^{(2)}=x_2^2+\lambda(t_{12}^2+t_{22}^2),
\]
所以這是素函數，因為 $t_{12}=(0),\ t_{22}=(1)$ 導致素函數 $x_2^2+\lambda$。但 $\bar{\frakq}$ 是真準素的，因為它含有 $(y_1+y_2)^2$ 而不含 $(y_1+y_2)$。範數必等於平方，也就是 $P^{(2)}$ 的 $p$ 次冪，因為模 $\bar{\frakq}$ 的線性無關剩餘類數為四。

下列例子說明，定理 XIV 的這個最後類似物並不總成立。令 $\mathsf{P}$ 由把不定元 $\lambda$ 附加到模 $3$ 的剩餘類域而得到，並置
\[
 \frakq=(y_1^3+\lambda,\,(y_1-y_2)^2).
\]
則如上，$\frakq$ 是真準素理想。考慮到 $y_2^3+\lambda$ 也被 $\frakq$ 整除，初等因子形式成為
\[
 P^{(2)}=x_2^3+\lambda(t_{12}+t_{22})^3,
\]
這是一個素函數；但範數等於 $P^{(2)}$ 的平方，因為模 $\bar{\frakq}$ 有六個線性無關剩餘類。因此這裡不出現初等因子形式的 $p^g$ 次冪。

第一個例子表明，在以非完美域為係數域時，如同在代數數域中一樣，可能出現高於一次的素理想，其中 $p^g$ 應稱為該素理想的次數。後面的例子應理解為分歧理想的類似物：一個素數可以被某個素理想的較高冪，即被某個準素理想整除。因為如注 10 所示，初等因子形式對應於代數數域中被一個理想整除的最小整數有理數。

\subsection*{§ 7. 絕對素理想}

若以 $\mathsf{P}$ 中元素為係數的素理想在 $P$ 可擴張到的代數閉域 $A$ 中仍為素理想，則稱它為\emph{絕對素理想}。\footnote{眾所周知，代數閉域是其中每個一元多項式都分解為線性因子的域。關於屬於任意域的代數閉域的存在性和本質唯一性，見 Steinitz。}相應地，若以 $\mathsf{P}$ 中元素為係數的素函數 $P$ 在 $\mathsf{A}$ 中仍為素函數，則稱它為\emph{絕對素函數}--絕對不可約多項式。絕對素理想的刻畫由下述定理給出：

\textbf{定理 XV.} \emph{一個絕對素理想的必要且充分條件是，它的初等因子形式--可假定在 $u$ 中為整且本原--作為 $x$ 與 $u$ 的多項式成為絕對素函數。初等因子形式同範數相同，因此該條件也可用範數來表述。}

證明時首先注意，一般地，範數和初等因子形式在係數域的代數擴張下保持。事實上，個別因子 $R^{(i)}$ 與 $E^{(i)}$ 分別定義為多項式意義下的最大公因子，而這一性質在係數域的代數擴張下保持。因此，當從 $P$ 轉到 $\mathsf{A}$，也就是轉到完美域作為係數域時，$P^{(i)}$ 仍為 $\frakp$ 的初等因子形式，因為每個代數閉域由定義立即可見是完美的。由定理 XIII，$\frakp$ 成為絕對素理想的必要且充分條件是，$P^{(i)}$ 相對於 $\mathsf{A}(u)$ 成為素函數；也就是說，若 $P^{(i)}$ 假定在 $u$ 中為整且本原，則它作為 $x$ 與 $u$ 的多項式為絕對素函數。同時，由定理 XIII，一旦取 $\mathsf{A}$ 為係數域，$P^{(i)}$ 就同 $\frakp$ 的範數相同。因此由剛才所指出的事實，在 $\mathsf{P}$ 上同樣成立。定理 XV 得證。

對於以（有限）代數數域 $\mathfrak K$ 中元素為係數的絕對素函數 $P$，有如下定理\footnote{A. Ostrowski, loc. cit. (note 21), lemma p. 296。一個更簡單的證明見 E. Noether, Ein algebraisches Kriterium für absolute Irreduzibilität, Math. Ann. 85 (1922), pp. 26--33, no. 7。}：模 $\mathfrak K$ 的每個素理想時，$P$ 仍為素函數，更確切地說仍為絕對素函數，至多有有限多個 $\mathfrak K$ 的素理想例外。把這個定理轉移到絕對素理想，依賴於

\textbf{定理 XVI.} \emph{若作為係數域不用域，而用（有限）代數數域 $\mathfrak K$ 的所有代數整數構成的環 $\mathfrak o^*$，則多項式理想的範數和初等因子形式，在模 $\mathfrak K$ 的每個素理想時仍分別保持為範數和初等因子形式，至多有有限多個 $\mathfrak K$ 的素理想例外。}

為證明，需要相對於 § 1, 1 修改作為基礎的多項式域。令 $\bar{\mathfrak{J}}$ 由所有以 $\mathfrak o^*$ 中元素、即 $\mathfrak K$ 的代數整數為係數的 $y_1,\ldots,y_n$ 中多項式組成；令 $\mathfrak{J}$ 的係數域為 $\mathfrak o^*[u]$，也就是所有以 $\mathfrak o^*$ 中元素為係數的 $u$ 中多項式。若 $\bar{\frakm}$ 是 $\bar{\mathfrak{J}}$ 中的理想，它通過 (1) 變成 $\mathfrak{J}$ 中的一個變換理想 $\frakm$。要證明在形成範數時，分母中只出現 $\mathfrak o^*[u]$ 中有限多個多項式；於是定理 XVI 直接推出。

首先注意，§ 1, 5 中為了形成個別範數而使用的模表示 (3)，相當於模一個關於 $x_{i-1}$ 正則的多項式來約化 $x_{i-1}$ 的冪，
\[
 C^{(i-1)}=U_{i-1}(u)x_{i-1}^{e_{i-1}}+\text{較低項},
\]
\footnote{比較 H.-N., § 4。}其中 $C^{(i-1)}$ 的所有係數，特別是 $U_{i-1}$，可假定為 $\mathfrak o^*[u]$ 中元素。由於這是 $x_1,\ldots,x_{i-1}$ 的逐次約化，關於 $x$ 為整的模表示除分母中的冪積
\[
 U_1^{\lambda_1}\cdots U_{i-1}^{\lambda_{i-1}}
\]
外，也關於 $\mathfrak o^*[u]$ 為整。乘積 $U_1U_2\cdots U_n$ 作為 $u$ 中多項式，由其係數確定 $\mathfrak K$ 的一個理想 $\mathfrak r^*$，因而只被 $\mathfrak K$ 的有限多個素理想整除。若選擇 $\frakp^*$ 不同於這些有限多個素理想，並使用模 $\frakp^*$ 的剩餘類域作為係數域，則通過把 $\mathfrak o^*$ 的每個數替換為其模 $\frakp^*$ 的剩餘類，原來的模表示得以保持。

此外，由於範數和初等因子形式只在相差 $\mathfrak o^*[u]$ 中因子的意義下確定，它們也可假定為相對於 $\mathfrak o^*[u]$ 被 $\frakm$ 整除。設在這個約定下例如
\[
 R^{(i)}=V_i(u)x_i^{e_i}+\text{較低項},
\]
那麼在秩為 $\rho$、由 $C^{(i-1)}$ 決定的模 $\mathfrak M_{i-1}^*$ 的所有 $\rho$ 階行列式的整除中，分母除 $U_1^{\lambda_1}\cdots U_{i-1}^{\lambda_{i-1}}$ 外只出現 $V_i$ 的冪。進一步選擇 $\frakp^*$ 不同於由 $V_1V_2\cdots V_n$ 確定的理想中出現的有限多個 $\mathfrak K$ 的素理想。於是當取模 $\frakp^*$ 的剩餘類域為係數域時，$R^{(i)}$ 仍為這些行列式的公因子，並且 $\mathfrak M_{i-1}^*$ 的秩得以保持，因為 $C^{(i)}$ 是來自 $\mathfrak M_{i-1}^*$ 的一個 $\rho$ 階行列式。

最後還要選擇 $\frakp^*$ 使 $R^{(i)}$ 在每種情況下仍為最大公因子。令 $D^{(i)}$ 遍歷 $\mathfrak M_{i-1}^*$ 的所有 $\rho$ 階行列式；按假設有
\[
 U_1^{\lambda_1}\cdots U_{i-1}^{\lambda_{i-1}}V_i^{\sigma}D^{(i)}=R^{(i)}T^{(i)}
\]
且這些 $T^{(i)}$ 在多項式意義下沒有含有 $x$ 的公因子。因此，由所有 $T^{(i)}$ 導出的理想含有一個多項式
\[
 G^{(i+1)}=W_i(u)x_{i+1}^{m_i}+\text{較低項},\qquad W_i\ne0.
\]
$G^{(i+1)}$ 可假定關於 $x_{i+1}$ 正則，這是由把變換 (1) 同以新不定元為係數的 $x_{i+1},\ldots,x_n$ 中某一個作複合推出的。若再選擇 $\frakp^*$ 不同於由 $W_1W_2\cdots W_n$ 確定的理想中所含的有限多個素理想，則當使用模 $\frakp^*$ 的剩餘類域作為係數域時，$R_\frakm$ 的確保持為範數。完全相應地，可再進一步選擇 $\frakp^*$，使 $E_\frakm$ 也保持為初等因子形式。定理 XVI 得證。

由定理 XV 和 XVI，並利用絕對素函數的定理，得到該定理的如下推廣：

\textbf{定理 XVII.} \emph{若 $\frakp$ 是一個以（有限）代數數域 $\mathfrak K$ 中代數整數為係數的絕對素理想，則模 $\mathfrak K$ 的每個素理想時，$\frakp$ 仍保持為素理想，更確切地說保持為絕對素理想，至多有有限多個 $\mathfrak K$ 的素理想例外。}

為證明，按定理 XVI 選取 $\frakp$ 的範數 $R_\frakp$，使它相對於 $\mathfrak o^*[u]$ 也被 $\frakp$ 整除。於是
\[
 R_\frakp=T(u)P^{(i)}(x),
\]
其中 $P^{(i)}(x)$ 可假定在 $u$ 中為整且本原。由定理 XV，$P^{(i)}$ 作為以 $\mathfrak K$ 中元素為係數的 $u$ 與 $x$ 中多項式，是絕對素函數。由絕對素函數定理，模 $\mathfrak K$ 的每個素理想時，它都保持這個性質，至多有有限多個 $\mathfrak K$ 的素理想例外。於是選取 $\frakp^*$ 不同於這些有限多個素理想，並由定理 XVI 也不同於 $\mathfrak K$ 的有限多個進一步素理想。則 $P^{(i)}$--其來自 $\mathfrak o^*$ 的係數被替換為模 $\frakp^*$ 的相應剩餘類--在取模 $\frakp^*$ 的剩餘類域為係數域 $\mathsf{P}$ 時，成為 $\frakp$ 的範數；而這個範數是絕對素函數。於是以 $\mathsf{P}$ 為係數域，$\frakp$ 由定理 XV 仍為絕對素理想。定理 XVII 得證。

\begin{center}
Göttingen, May 8, 1923.\\[1em]
（1923 年 3 月 9 日收到。）
\end{center}

%
\clearpage
%
%
\begingroup
\def\srcfn#1#2{\begingroup\renewcommand{\thefootnote}{#1}\footnote{#2}\addtocounter{footnote}{-1}\endgroup}
\def\srcemph#1{\textbf{#1}}
\pdfbookmark[0]{25. 消去論與理想論}{paper-25}
\section*{25. 消去論與理想論}

\begin{center}
J. Ber. d. DMV 33 (1924), S. 116--120
\end{center}

\setcounter{footnote}{0}

\begin{center}
據 1923 年 9 月 25 日在馬爾堡所作報告整理。\srcfn{1)}{詳述及文獻說明見 Math. Ann. 90，並見一篇即將發表的論文。}\\[0.5em]
作者：\textsc{Emmy Noether}，哥廷根。
\end{center}

在下文中，我想說明消去論怎樣能夠完全平行於一元多項式的零點理論來建立；這個建立依託於 \emph{Hentzelt} 所給出的消去論同域論方法--如 \emph{Steinitz} 所發展者--以及同我自己所發展出的理想論方法之間的聯繫。

我先在一元多項式的情形中勾畫這個問題。作為係數域，始終取一個固定域 $P$ 為基礎；素函數等概念都相對於它來理解。這裡可以把 $P$ 假定為 Steinitz 意義下抽象定義的域；不過在特殊情形中，它當然也可以同有理數域，或同所有複數的域--按照消去論中從前通常採用的假設--相重合。一元多項式的零點理論於是分成四個步驟：

\noindent 1. \srcemph{分解定理}，它允許把問題化歸為關於素函數的問題--即化歸為關於在 $P$ 中不可約的多項式的問題。因為由
\[
 a(x)=p_1(x)^{\varrho_1}\cdots p_r(x)^{\varrho_r}=q_1(x)\cdots q_r(x)
\]
--其中 $p(x)$ 表示不同的素函數，$q(x)$ 因而表示“最大準素函數”--可知，$a(x)$ 消失於所有且僅有的素函數 $p(x)$ 的零點。

\noindent 2. \srcemph{為一個素函數 $p(x)$ 構造零點域。} 這樣的零點域 $\mathfrak K=P(\alpha)$--其中 $\alpha$ 表示 $p(x)$ 的一個零點--按 \emph{Steinitz} 總可以構造為 $P$ 的一個擴張域，它同 $p(x)$ 的剩餘類域 $(\mathfrak K)$ 同構。反過來，位於 $P$ 的某個擴張域中的每個零點域都同剩餘類域 $(\mathfrak K)$ 同構；因此，零點域和剩餘類域在抽象意義下是同一的。

\noindent 3. \srcemph{在由 $p(x)$ 確定的伽羅瓦擴張域中把 $p(x)$ 分解為線性因子。} 通過重複第 2 點所給出的構造，就到達一個本質上唯一確定的伽羅瓦擴張域 $P(\alpha_1,\ldots,\alpha_m)$；這個域正好足以使 $p(x)$ 分裂為線性因子。

\noindent 4. \srcemph{重數對於原多項式 $a(x)$ 的意義。} 根據 1、2、3，對於任意多項式 $a(x)$，零點問題以及相應的線性因子分解問題都已經解決。這時可以把各個準素函數 $q(x)$ 的次數--也就是模 $q(x)$ 後相對於 $P$ 線性無關的剩餘類的個數--看作重數；因為在代數閉係數域 $P$ 的情形中，$p(x)$ 都成為線性的，而這個個數恰好給出零點的重數。

現在我轉向一般消去論，也就是說，以所有以 $P$ 中元素為係數、關於 $x_1,\ldots,x_n$ 的多項式組成的多項式環 $\mathfrak{J}$ 為基礎；於是消去論可以定義為這個多項式環中\srcemph{理想的零點理論}。這裡我按通常方式把理想 $\fraka$ 理解為這樣一個多項式系：與 $f$ 和 $g$ 同在時也含有 $f-g$，並且與 $f$ 同在時也含有 $Af$，其中 $A$ 表示 $\mathfrak{J}$ 中任意多項式。所謂 $\fraka$ 的零點，我指屬於 $P$ 的某個擴張域的值系 $\alpha_1,\ldots,\alpha_n$，使得對每個屬於該理想的多項式 $f$，都有 $f(\alpha)$ 消失。若 $f_1,\ldots,f_t$ 是 $\fraka$ 的一組有限生成元（總存在這樣一組）--也就是說，對 $\fraka$ 中每個 $f$ 都有 $f=A_1f_1+\cdots+A_tf_t$--那麼 $\fraka$ 的零點顯然同 $f_1,\ldots,f_t$ 的零點相同；反過來，每個有限多項式系都可以看作由它導出的理想的一組生成元。因此，上面為消去論所給出的零點定義同通常定義一致，但避免了任意指定一個基所帶來的困難。在一元多項式的情形中，只涉及主理想，這個困難並不出現，因為生成多項式在相差一個單位--即來自 $P$ 的一個非零常數因子--的意義下唯一確定；不過即便在那裡，理想觀點也會給出更精確的理解。

相應於一元情形，消去論現在分為四個步驟：

\noindent I. \srcemph{分解定理}--把一個理想表示為最大準素分量--允許把問題化歸為關於素理想的問題。這裡，若 $ab\in\frakp$ 必有 $a\in\frakp$ 或 $b\in\frakp$，則 $\frakp$ 稱為素理想；若 $ab\in\frakq$ 且 $a\notin\frakq$ 時，總有某個 $m\geq 1$ 使 $b^m\in\frakq$，則 $\frakq$ 稱為準素理想。每個準素理想 $\frakq$ 都有唯一的關聯素理想 $\frakp=\sqrt{\frakq}$；因而 $\frakq\subseteq\frakp$，並且存在某個 $m\geq 1$ 使 $\frakp^m\subseteq\frakq$。存在作為最小公倍數的表示：
\[
 \mathfrak a=[\mathfrak q_1,\cdots,\mathfrak q_r];\qquad
 \frakp_1^{\sigma_1}\cdots \frakp_r^{\sigma_r}\equiv 0\,(\fraka).
\]
這裡 $\frakq$ 是最大準素分量；也就是說，每個 $\frakq$ 都是準素的，但任意兩個 $\frakq$ 的最小公倍數不是準素的。此外，可以不失一般性地假定沒有任何 $\frakq$ 整除其餘 $\frakq$ 的最小公倍數，也就是沒有任何 $\frakq$ 可以簡單地刪去；$\frakp$ 是各個 $\frakq$ 的關聯素理想。現在，雖說 $\frakq$ 本身並非唯一，但--這正是一元情形的類似物所在--它們的關聯素理想 $\frakp$ 由 $\fraka$ 唯一確定；並且由於每個 $\frakp$ 都包含 $\fraka$，而反過來各 $\frakp$ 的某個冪積又包含於 $\fraka$，所以理想的零點同其關聯素理想的零點相一致。

\noindent II. \srcemph{為一個素理想 $\frakp$ 構造零點域。} 按定義，模素理想的剩餘類組成一個沒有零因子的環；只要 $\frakp$ 不同於單位理想，這個環就至少含有一個不同於零元素的元素。因此，通過取分式，可以把這個環擴張為 $\frakp$ 的剩餘類域 $(\mathfrak K)$。總可以構造 $P$ 的一個擴張域 $\mathfrak K$，它同 $(\mathfrak K)$ 同構，並且成為 $\frakp$ 的零點域；也就是說，$\mathfrak K=P(\alpha_1,\ldots,\alpha_n)$，其中 $\alpha_1,\ldots,\alpha_n$ 表示 $\frakp$ 的一個零點。這裡 $\mathfrak K$ 相對於 $P$ 的超越次數為 $k$（$0\leq k<n$），因而含有 $k$ 個相對於 $P$ 代數無關的元素；而任意 $(k+1)$ 個元素相對於 $P$ 都成為代數相關。反過來，每個位於 $P$ 的某個擴張域中、超越次數為 $k$ 的零點域，都同剩餘類域 $(\mathfrak K)$ 同構；因此，超越次數為 $k$ 的零點域和剩餘類域在抽象意義下是同一的。

\noindent III. \srcemph{在由 $\frakp$ 確定的伽羅瓦擴張域中把 $\frakp$ 的範數分解為線性因子。} 若把不定元 $x_1,\ldots,x_n$ 看作由原不定元 $y_1,\ldots,y_n$ 通過以不定元為係數的線性變換生成，則可把超越次數為 $k$ 的零點域 $\mathfrak K$ 看作 $P(x_{i+1},\ldots,x_n)$ 的有限代數擴張域（$(k=n-i)$）。於是可以轉到 $\overline{\mathfrak K}$ 關於 $P(x_{i+1},\ldots,x_n)$ 的伽羅瓦擴張域；在其中，帶有零點 $\alpha_i$ 的素函數 $P^{(i)}(x_i,x_{i+1},\ldots,x_n)$ 分裂為線性因子。當回到原不定元 $y$ 時，這個素函數 $P^{(i)}$ 起基本方程的作用；因為 $\alpha_i$ 等於基本形式 $u_1\beta_1+\cdots+u_n\beta_n$，其中 $\beta$ 表示 $y$ 中的一個零點系，它代數地依賴於參數 $x_{i+1},\ldots,x_n$；於是線性因子的分解給出對所有零點的乘積。但素函數 $P^{(i)}$ 或其某個冪，也可以看作 $\frakp$ 的範數 $N(\frakp)$：這是通過把 $\frakp$ 看成 $x_1,\ldots,x_{i-1}$ 的冪積中的線性形式模，把 $P(x_{i+1},\ldots,x_n)$ 取作係數域以代替 $P$，並把範數定義為所有初等因子的乘積來實現的--其中只有有限多個初等因子不同於單位。若係數域 $P$ 是完全域，則總有 $N(\frakp)=P^{(i)}$；若係數域 $P$ 非完全，則範數可能成為 $P^{(i)}$ 的一個冪。

\noindent IV. \srcemph{原來理想的範數以及重數的意義。} 相應地，若把理想 $\fraka$ 逐次看作 $x_1,\ldots,x_{i-1}$ 的冪積中的線性形式模（$(i=1,2,\ldots n)$），則可以證明，在以 $P(x_{i+1},\ldots,x_n)$ 為基礎時，這個模每次只有有限多個不同於單位的初等因子；這是一個有限性條件，它由理想具有有限生成元這一事實推出。每次這些有限多個初等因子的乘積給出各次範數；而範數 $N(\fraka)$ 定義為所有各次範數的乘積。$N(\fraka)$ 的最大準素因子 $Q^{(i)}(x_i,x_{i+1},\ldots,x_n)$ 同 $\fraka$ 的關聯素理想 $\frakp$ 一一對應，並且對應方式為：$Q^{(i)}$ 成為 $P^{(i)}$ 的一個冪，其中 $P^{(i)}$（或在必要時 $P^{(i)}$ 的一個冪）表示 $\frakp$ 的範數 $N(\frakp)$。於是得到一個分解
\[
 N(\fraka)=N(\frakp_1)^{\varrho_1}\cdots N(\frakp_r)^{\varrho_r}=Q_1\cdots Q_r,
\]
--其中在必要時，要用最高初等因子 $E(\frakp)=P^{(i)}$ 代替 $N(\frakp)$--這樣，按照 I、II、III，就為 $\fraka$ 的所有零點完成了範數分解為線性因子。對應於 $\fraka$ 的單個素理想 $\frakp$ 的零點，其重數可以看作 $Q^{(i)}$ 的次數；這個次數可以解釋為相對於 $P(x_{i+1},\ldots,x_n)$ 線性無關的剩餘類個數。更確切地說，這裡涉及的是某些由 $\frakp$ 和更高維數的素理想唯一確定的分解孤立分量的剩餘類：即 $(i-1)$ 級基本理想模由 $\frakp$ 所確定的 $i$ 級基本理想分量所得的剩餘類；後一分量以 $\frakp$ 以及所有更高維數的素理想為關聯素理想，而前一分量只以更高維數的素理想為關聯素理想。

若在這裡也轉到主理想，就得到\srcemph{一元情形的完整歸入}。這時，第 1 點所給出的分解--按照把 $a(x)$ 看作生成多項式還是看作範數--分裂為兩個分開的分解：其一，是按 I 把主理想表示為最大準素分量的最小公倍數，在這裡這等於把它表示為素理想冪的乘積；其二，是按 IV 把範數分解為各個素理想範數的冪的乘積。按第 2 點構造零點域時，實際上涉及的是由 $p(x)$ 定義的素理想的零點域；而在第 3 點中被分解為線性因子的是範數。按第 4 點給出的重數定義被納入 IV 中給出的定義之下，因為零級基本理想等於單位理想，而一元情形中的一級基本理想已經等於理想本身。

\srcemph{歸入通常的消去論}則由如下事實給出：理想的範數具有結式意義。在這種理解下，消去論的舊問題因此成為多項式環的理想論的一部分，而且正是這樣一部分：除了理想的分解定理之外，它還處理這些理想的範數分解以及兩種分解之間的聯繫。

\begin{center}
（1923 年 10 月 19 日收到。）
\end{center}
\endgroup

%
\clearpage
%
\setcounter{footnote}{0}

\pdfbookmark[0]{26. 代數數域中理想論的抽象構造}{paper-26}
\section*{26. 代數數域中理想論的抽象構造}

\begin{center}
J. Ber. d. DMV 33 (1924), p. 102
\end{center}

\textbf{4. E. Noether，哥廷根：代數數域中理想論的抽象構造。}

這裡給出了若干抽象性質；它們同代數數域中的理想論完全等價，因而刻畫出所有擁有這種理想論的環。這些環是具有單位元且無零因子的環，在其中雙鏈條件成立--每一條由理想組成的除子鏈都在有限步後中止，同樣，模每個不同於零理想的理想的每一條倍數鏈也都在有限步後中止--並且這些環相對於它們的商域是\emph{整}閉的。若去掉後一假設，則得到關於有限階中理想論的若干定理，其中一部分已見於 Dedekind（第 3 版），但未給出證明。

詳述將發表於 \emph{Mathematische Annalen}。

%
\clearpage
%
\setcounter{footnote}{0}

\pdfbookmark[0]{27. 理想論中的希爾伯特數}{paper-27}
\section*{27. 理想論中的希爾伯特數}

\begin{center}
J. Ber. d. DMV 34 (1925), p. 101
\end{center}

--- \textbf{E. Noether}：理想論中的希爾伯特數。所謂希爾伯特數，通常指同理想的剩餘類計數有關的那些數。Hilbert 本人在他的“特徵函數”中給出了一個函數；對於多項式整環的理想，這個函數從某個界限起給出線性無關剩餘類的精確數目。對於較低的次數數，Macaulay 用一個生成函數補充了這個函數，而 Ostrowski 能夠將該生成函數顯式求值。不過，Macaulay 還給出了另一類數；事實證明，由於這類數可以抽象地加以把握，它們對於一般理想論最為本質。按照一般分解定理，只需對準素理想 $\mathfrak q$ 給出定義。這裡涉及的是相對於所屬素理想 $\mathfrak p$ 的剩餘類域而言線性無關的剩餘類數。這些數分解為若干子系統，每個子系統分別由 $\mathfrak q,\mathfrak q/\mathfrak p,\mathfrak q/\mathfrak p^2,\ldots$ 的不可約分量數給出。第 $i$ 個子系統的數目，也可由一條從 $\mathfrak q/\mathfrak p^i$ 走向 $\mathfrak q/\mathfrak p^{i-1}$ 的合成列的項數來刻畫。總合成列及其希爾伯特數，構成了在一般情形中--例如在一個代數數域的有限階中--不再有效的那種把準素理想表示為素理想冪的等價替代物。

%
\clearpage
%
\setcounter{footnote}{0}

\pdfbookmark[0]{28. 群特徵標與理想論}{paper-28}
\section*{28. 群特徵標與理想論}

\begin{center}
J. Ber. d. DMV 34 (1925), p. 144
\end{center}

\textbf{3. E. Noether，哥廷根：群特徵標與理想論。}

Frobenius 的群特徵標理論--也就是有限群表示的理論--被理解為一個完全可約環，即群環，的理想論。因此，一般地發展了完全可約環在加法分解下的同構定理；把它表示為直接不可分解的簡單理想的直和時，這種表示在同構意義下唯一，並且屬於同一個雙邊理想的不可分解單邊理想彼此同構。因此，若把同構的理想合併為一類，則不同的不可分解理想類的數目，正好等於雙邊不可分解理想的數目。

專門化到超複量的完全可約系統--特別是群環--時，不可分解理想類和不可約表示類彼此一一對應。進一步，雙邊不可分解理想和特徵標相互對應；而一個雙邊理想的加法不可分解分量的數目，等於這些不可分解理想的秩。這樣，Frobenius 理論就被納入了這一框架。

詳述將發表於 \emph{Math. Ann.}

%
\clearpage
%
%
\begingroup
\def\srcemph#1{\textbf{#1}}
\pdfbookmark[0]{29. 特徵為 p 的有限線性群不變量的有限生成性定理}{paper-29}
\section*{29. 特徵為 $p$ 的有限線性群不變量的有限生成性定理}

\begin{center}
Nachr. v. d. Ges. d. Wiss. zu Göttingen 1926, S. 28--35
\end{center}

\setcounter{footnote}{0}

\begin{center}
作者\\[0.5em]
\textbf{Emmy Noether}，哥廷根。\\[1em]
由 R. Courant 於 1926 年 7 月 30 日會議上報告。
\end{center}

下面證明：有限線性代換群不變量的著名有限生成性定理\footnote{關於一個初等的有限性證明，參見 E. Noether, Der Endlichkeitssatz der Invarianten endlicher Gruppen. Math. Ann. 77 (1916), S. 89--92.} 在把通常的數域改為任意域作為係數域時仍然成立——因而尤其適用於特徵為 $p$ 的域。\footnote{關於域論概念，參見 E. Steinitz, Algebraische Theorie der Körper. J. f. M., 137 (1910), S. 167--309。特徵見 \S\ 4，根域見 \S\ 12，分式域見 \S\ 3。} 不過，一旦群的階可被 $p$ 整除，已知的有限性證明便失效；必須引入更深的算術事實，也就是下面這個此前只在特徵零情形已知的

\srcemph{有限性判別準則}\footnote{參見 E. Noether, Algebraische und Differentialinvarianten. Jahresber. d. d. Math.-Ver. 32 (1923), S. 177--184, Nr. 3, Satz 4 以及那裡列出的文獻。在該判別準則的最後一句中，誤寫為“關於 $\mathfrak R$ 的模生成組”（整元素的基本系），而不是“關於 $\mathfrak R$ 的某個子環”。}：一個由多項式——更一般地，由有理函數——構成、係數來自一個域的整環 $\mathfrak J$，當且僅當 $\mathfrak J$ 中含有一個有限生成子環 $\mathfrak R$，使得 $\mathfrak J$ 在 $\mathfrak R$ 上整時，$\mathfrak S$ 是有限生成的。於是 $\mathfrak R$ 的每一個整性基（即有限代數生成組），都可由 $\mathfrak J$ 關於 $\mathfrak R$ 的某個子環的一個模生成組補充為 $\mathfrak J$ 的一個整性基。

這個判別準則——它本身也有獨立的意義——依賴於有理基的存在，以及這樣一個事實：從一個域過渡到其有限擴張中的整元素時，因子鏈定理仍然成立。第二個定理已由 Artin 和 van der Waerden 在前一篇札記中對任意特徵證明，不過只是在原來的整環滿足某種有限性前提時；在有限群情形中這個前提成立（根環給出一個有限擴張）。對於有理基的存在，我給出一個以任意域為係數域都有效的證明；由此便能在上述限制下證明有限性判別準則。

對於有限群的不變量，有限性判別準則之所以滿足，依據的是對稱函數原理。然而，在特徵零情形，伽羅瓦預解式的係數已經給出整性基；而在一般情形，它們只生成有限性判別準則所要求的有限生成子環。由這個有限生成子環的存在，按同樣的推理可推出“相對”不變量和“模”不變量的有限生成性。

這裡考察的不變量，作為特殊情形，包括 Dickson 及其學派\footnote{參見 Dickson, On Invariants and the theory of Numbers. The Madison Colloquium 1913。較新的文獻見此後出版的 Transactions of the American Mathematical Society 各卷。} 所處理的“模 $p$ 射影不變量”。在這裡，有限生成性定理此前雖已對模不變量成立，卻尚未對一般的“形式”不變量得到證明。

\subsection*{\S\ 1. 有限性判別準則}

1. \srcemph{有理基存在定理。第一種表述：}由 $n$ 個不定元的有理函數組成、係數來自一個域 $P$ 的每個系統 $S$，都有一個\srcemph{有限有理基}；也就是說，系統中存在有限多個函數，使得系統中的每個函數都能用這有限多個函數\srcemph{以 $\mathsf{P}$ 中的係數作有理表示}。

\srcemph{第二種表述：}設 $\mathfrak K$ 是 $\mathsf{P}$ 與 $\mathsf{P}(x_1,\ldots,x_n)$ 之間的一個中間域——其中 $x_1,\ldots,x_n$ 表示不定元——且超越次數為 $t\leqq n$；設 $y_1,\ldots,y_t$ 是 $\mathfrak K$ 中的一個不可約組\footnote{按 Steinitz 的說法，一個組稱為不可約，是指它由代數無關函數組成。一個組稱為超越次數為 $t$，是指它至少含有一個由 $t$ 個元素組成的不可約組，但不含有元素更多的不可約組。關於下面使用的不可約組的簡單定理，參見 Steinitz, \S\ 22。}。那麼 $\mathfrak K$ 是 $\mathsf{P}(y_1,\ldots,y_t)$ 的\srcemph{有限代數擴張}。

這兩種表述事實上等價。只需對所有中間域證明有理基的存在即可；因為由一個系統 $S$ 導出的域是這樣一個中間域 $\mathfrak K$，而 $\mathfrak K$ 的有理基立即給出 $S$ 的有理基，這是因為 $\mathfrak K$ 的每個元素都是 $S$ 中有限多個函數的有理組合。另一方面，由第二種表述可知，$\mathfrak K$ 中任意一個不可約組 $y_1,\ldots,y_t$，再加上表徵 $\mathfrak K$ 為 $\mathsf{P}(y_1,\ldots,y_t)$ 有限擴張的有限多個函數，就給出這樣一個有理基。反過來，如果按第一種表述給出了 $\mathfrak K$ 的一個有理基，那麼按超越次數的定義，每個基元素都必定代數依賴於 $\mathsf{P}(y_1,\ldots,y_t)$；於是 $\mathfrak K$ 就由此被表徵為 $\mathsf{P}(y_1,\ldots,y_t)$ 的有限代數擴張。

下面在\srcemph{第二種表述}下證明該定理。先設 $\mathfrak K$ 與 $\mathsf{P}(x_1,\ldots,x_n)$ 具有相同的超越次數 $n$；設 $y_1,\ldots,y_n$ 是 $\mathfrak K$ 中的一個不可約組，因而也是 $\mathsf{P}(x_1,\ldots,x_n)$ 中的一個不可約組。由於 $y_1,\ldots,y_n$ 的超越次數為 $n$，每個 $x_i$ 都在 $\mathsf{P}(y_1,\ldots,y_n)$ 上代數；所以 $\mathsf{P}(x_1,\ldots,x_n)$，從而其中間域 $\mathfrak K$，成為 $\mathsf{P}(y_1,\ldots,y_n)$ 的有限代數擴張。

現在設 $t<n$；將 $x$ 的編號選定，使得 $y_1,\ldots,y_t,x_{t+1},\ldots,x_n$ 構成 $\mathsf{P}(x_1,\ldots,x_n)$ 中的一個不可約組。於是 $x_{t+1},\ldots,x_n$ 關於 $\mathsf{P}(y_1,\ldots,y_t)$ 代數無關，並且由代數依賴的傳遞律，也關於 $\mathfrak K$ 以及 $\mathsf{P}(y_1,\ldots,y_t)$ 與 $\mathfrak K$ 之間的每個中間域 $\mathfrak M$ 代數無關。這就是說，$\overline{\mathfrak K}=\mathfrak K(x_{t+1},\ldots,x_n)$ 以及 $\overline{\mathfrak M}=\mathfrak M(x_{t+1},\ldots,x_n)$，分別成為 $\mathfrak K$ 以及 $\mathfrak M$ 的純超越擴張；因而 $\overline{\mathfrak K}$ 中每個不屬於 $\mathfrak K$ 的元素都關於 $\mathfrak K$ 超越，同樣，$\overline{\mathfrak M}$ 中每個不屬於 $\mathfrak M$ 的元素都關於 $\mathfrak M$ 超越。再令 $\overline{\mathsf{P}}$ 等於純超越擴張 $\mathsf{P}(x_{t+1},\ldots,x_n)$；於是 $y_1,\ldots,y_t$ 關於 $\overline{\mathsf{P}}$ 代數無關。

現在，把 $\overline{\mathsf{P}}$ 作為係數域，就可以把證明歸結到上面已經處理過的情形。由
\[
 P<\mathfrak K<P(x_1,\ldots,x_t;\,x_{t+1},\ldots,x_n)\text{\footnotemark}
\]
\footnotetext{$\mathsf{P}<\mathfrak K$ 表示：$\mathsf{P}$ 包含在 $\mathfrak K$ 中，即 $P$ 是 $\mathfrak K$ 的子域。}可推出
\[
 \overline P<\overline{\mathfrak K}<\overline P(x_1,\ldots,x_t),
\]
其中 $\overline{\mathfrak K}$ 關於 $\overline{\mathsf{P}}$ 的超越次數為 $t$，與 $\overline{\mathsf{P}}(x_1,\ldots,x_t)$ 相同。因此，$\overline{\mathfrak K}$ 以 $\overline{\mathsf{P}}$ 為係數域，有一個\srcemph{有限有理基}，記為 $y_1,\ldots,y_t; z_1,\ldots,z_s$。這裡 $z$ 可以選自 $\mathfrak K$，因為 $\overline{\mathfrak K}$ 是 $\mathfrak K$ 與 $\overline{\mathsf{P}}$ 的合成域。還需證明 $\mathfrak K$ 等於 $\mathfrak M=\mathsf{P}(y_1,\ldots,y_t;z_1,\ldots,z_s)$。$\mathfrak M$ 是 $\mathsf{P}(y_1,\ldots,y_t)$ 與 $\mathfrak K$ 之間的一個中間域；所以 $\mathfrak K$ 的每個元素都關於 $\mathfrak M$ 代數。另一方面，$\overline{\mathfrak M}$ 等於 $\overline{\mathfrak K}$，因而 $\mathfrak K$ 包含在 $\overline{\mathfrak M}$ 中。於是 $\mathfrak K$ 成為 $\mathfrak M$ 與 $\overline{\mathfrak M}$ 之間的中間域；又由於 $\overline{\mathfrak M}$ 中每個不屬於 $\mathfrak M$ 的元素都關於 $\mathfrak M$ 超越，所以必然有 $\mathfrak K=\mathfrak M$。

\srcemph{推論：}若 $\mathsf{P}<\mathfrak K<\mathfrak L<\mathsf{P}(x_1,\ldots,x_n)$，並且 $\mathfrak L$ 關於 $\mathfrak K$ 代數，則 $\mathfrak L$ 是\srcemph{$\mathfrak K$ 的有限代數擴張}。因為 $\mathfrak L$ 的一個有理基中的每個元素於是都關於 $\mathfrak K$ 代數；從而 $\mathfrak L$ 被表徵為 $\mathfrak K$ 的有限擴張。

2. \srcemph{有限性判別準則的證明。}這個條件顯然是必要的；因為如果 $\mathfrak J$ 是有限生成整環，那麼 $\mathfrak J$ 本身就可看作判別準則中出現的有限生成子環 $\mathfrak R$。

若係數域具有特徵 $p$，這個條件在如下\srcemph{前提——其 $p$ 次根域 $\mathsf{P}^{1/p}$ 是係數域的有限擴張——}\footnote{見第 28 頁腳註 2。}下是充分的；而在特徵零情形，不需要任何限制性前提。為了證明這一點，需要證明 $\mathfrak J$ 關於 $\mathfrak R$ 的某個子環有一個\srcemph{有限}模生成組。

令 $\mathfrak K$ 表示 $\mathfrak R$ 的分式域\footnotemark[\value{footnote}]，令 $\mathfrak L$ 表示 $\mathfrak J$ 的分式域。這些分式域存在，因為這裡所涉及的是無零因子環；並且有
\[
 P<\mathfrak K<\mathfrak L<P(x_1,\ldots,x_n).
\]
因此，由 1. 下的推論，$\mathfrak L$ 是 $\mathfrak K$ 的\srcemph{有限代數擴張域}；又按假設，$\mathfrak J$ 中每個元素都在 $\mathfrak R$ 上整，於是 $\mathfrak J$ 是由 $\mathfrak L$ 中所有在 $\mathfrak R$ 上整的元素組成的環 $\mathfrak S$ 的子環。但是，由於並未假設 $\mathfrak R$ 在 $\mathfrak K$ 中整閉，有限擴張下因子鏈定理仍成立這一事實不能直接引用。必須先過渡到 $\mathfrak R$ 的一個同樣有限生成的子環 $\mathfrak T$，使得這種整閉性成立。

先假設係數域 $\mathsf{P}$ 含有無窮多個元素。那麼可從 $\mathfrak R$ 中選出一個不可約組 $g_1(x),\ldots,g_t(x)$，使得 $\mathfrak R$，從而 $\mathfrak J$，在
\[
 \mathfrak T=P[g_1(x),\ldots,g_t(x)]
\]
上整。事實上，設 $f_1(x),\ldots,f_k(x)$ 是按假設存在的 $\mathfrak R$ 的整性基；若它們不構成不可約組，則在必要地重新編號之後，設 $F(f_k;f_1,\ldots,f_{k-1})=0$ 是 $f_k$ 所滿足的一個關係。若把 $f_1,\ldots,f_{k-1}$ 替換為
\[
 f_1+t_1f_k,\ldots,f_{k-1}+t_{k-1}f_k;
\]
則當 $t_i$ 取為不定元時，$f_k$，從而也包括 $f_1,\ldots,f_{k-1}$，都在這些新函數上整；由於 $\mathsf{P}$ 有無窮多個元素，可把 $t_i$ 特化為 $\mathsf{P}$ 中的元素，使這種整依賴保持。於是經過有限多步，並利用整依賴的傳遞律，就得到所求的環 $\mathfrak T=\mathsf{P}[g_1(x),\ldots,g_t(x)]$\footnote{參見 Hilbert, Über die vollen Invariantensysteme, \S\ 1. Math. Ann. 42 (1893), S. 313--373.}。此時 $\mathfrak L$ 是 $\mathfrak T$ 的分式域 $\mathsf{P}(g_1(x),\ldots,g_t(x))$ 的有限擴張域，而由 $\mathfrak L$ 中所有在 $\mathfrak R$ 上整的元素組成的環 $\mathfrak S$，與由 $\mathfrak L$ 中所有在 $\mathfrak T$ 上整的元素組成的環相同。

在 $\mathfrak T$ 中，所有理想滿足因子鏈定理，並且 $\mathfrak T$ 在其分式域中整閉；二者都因為 $\mathfrak T$ 同構於 $t$ 個不定元的多項式環。由 $\mathsf{P}^{1/p}$ 關於 $\mathsf{P}$ 有限這一假設，還可推出 $\mathfrak T^{1/p}$ 關於 $\mathfrak T$ 有限\footnote{參見引言中引用的 Artin--van der Waerden 札記。}。因此，在 $\mathfrak S$ 中，所有 $\mathfrak T$-模都滿足因子鏈定理；特別地，$\mathfrak J$ 作為 $\mathfrak S$ 中包含 $\mathfrak T$ 的一個子環，有一個\srcemph{有限 $\mathfrak T$-模生成組}\footnotemark[\value{footnote}]。此外，如果 $\mathfrak L$ 是 $\mathsf{P}(g_1(x),\ldots,g_t(x))$ 的第一類擴張，則 $\mathfrak S$ 的有限 $\mathfrak T$-模生成組的存在不再需要對係數域 $P$ 作任何限制性前提；因此特別是在 $P$ 具有特徵零時總是成立。\footnote{例如參見 E. Noether, Abstrakter Aufbau der Idealtheorie in algebraischen Zahl- und Funktionenkörpern, \S\ 3. Math. Ann. 96 (1926), S. 26--61.}

設 $h_1(x),\ldots,h_s(x)$ 是 $\mathfrak S$ 的一個 $\mathfrak T$-模生成組；那麼對 $\mathfrak J$ 中任意 $f(x)$ 的表示
\[
 f(x)=A_1(g)h_1(x)+\cdots+A_s(g)h_s(x)
\]
表明 $g_1(x),\ldots,g_t(x),h_1(x),\ldots,h_s(x)$ 構成所求的整性基；這裡 $A$ 作為 $\mathfrak T$ 的元素，表示 $g$ 中的多項式。

若域 $\mathsf{P}$ 只含有有限多個元素，那麼通過添加一個不定元 $u$——也就是一個關於 $\mathfrak J$ 超越的元素——所得的域 $\mathsf{P}(u)$ 仍含有無窮多個元素。因此，環 $\mathfrak J(u)$ 以 $\mathsf{P}(u)$ 為係數域有一個有限整性基；並且由於 $\mathfrak J(u)$ 是 $\mathfrak J$ 與 $\mathsf{P}(u)$ 的合成環，這個基可以通過按 $u$ 的冪分解而從 $\mathfrak J$ 中選出，不過 $\mathfrak S$ 的基元素 $g_i$ 現在將不再構成不可約組。這樣，$\mathfrak T$ 中每個多項式 $f(x)$ 都有一個表示
\[
 C(u)f(x)=G(u,g_i,h_i),
\]
其中 $C(u)$ 表示 $\mathsf{P}[u]$ 中一個非零多項式，而 $G$ 是 $u,g_i,h_i$ 的多項式。比較 $u$ 的某個適當冪的係數，就表明 $g_i(x)$ 與 $h_i(x)$ 給出 $\mathfrak J$ 的一個整性基。於是\srcemph{有限性判別準則已在一般情形中得到證明}。

\subsection*{\S\ 2. 不變量的有限生成性}

1. 所謂\srcemph{特徵為 $p$ 的有限線性群}，是指由 $h$ 個線性代換 $A_1,\ldots,A_h$ 組成的一個群 $\mathfrak H$，這些代換的行列式非零，且係數來自特徵為 $p$ 的一個域 $\mathsf{P}$。這裡 $A_k$ 表示線性代換
\[
\begin{aligned}
 x_1^{(k)}&=a_{11}^{(k)}x_1+\cdots+a_{1n}^{(k)}x_n,\\[-0.2em]
 &\ldots,\\[-0.2em]
 x_n^{(k)}&=a_{n1}^{(k)}x_1+\cdots+a_{nn}^{(k)}x_n,
\end{aligned}
\]
或者簡記為 $(x^{(k)})=A_k(x)$。群 $\mathfrak H$ 因而把帶有元素 $x_1,\ldots,x_n$ 的組 $(x)$ 變為帶有元素 $x_1^{(k)},\ldots,x_n^{(k)}$ 的各組 $(x^{(k)})$。由於 $A_1,\ldots,A_h$ 中必須包含恆等變換，各組 $(x^{(k)})$ 中也包含組 $(x)$。

所謂該群的一個\srcemph{整有理（絕對）不變量}，是指 $x_1,\ldots,x_n$ 的一個這樣的整有理函數（即多項式）——係數來自 $\mathsf{P}$\footnote{這並不限制一般性，因為任意一個係數來自某個特徵為 $p$ 的域的不變量——為了定義這些不變量，該域必須與 $\mathsf{P}$ 有一個共同的包域——都可以線性地分解為有限多個係數來自 $\mathsf{P}$ 的不變量。一般情形中，只需把係數表示為關於 $\mathsf{P}$ 線性無關的那些係數；於是，不變性就在這些獨立係數中恆等地成立。}——它在應用 $A_1,\ldots,A_h$ 時恆等地保持不變。因此，各組 $(x^{(k)})$ 的所有對稱函數——係數來自 $P$——都屬於這些不變量。反過來，對每一個這樣的不變量都有
\[
 f(x)=f(x^{(1)})=\cdots=f(x^{(h)})
 \quad\text{或}\quad
 h\cdot f(x)=f(x^{(1)})+\cdots+f(x^{(h)}).
\]
由此可知，在群的階 $h$ 不被特徵整除的情形——因而特別是在特徵零情形——各組 $(x^{(k)})$ 的、係數來自 $\mathsf{P}$ 的對稱函數已經窮盡所有不變量。由於這些對稱函數有一個有限整性基，在這種情形下有限生成性證明就完成了；同時，整性基給出為伽羅瓦預解式\footnote{見第 28 頁腳註 1。}的係數系統 $G_{\alpha\alpha_1\ldots\alpha_n}(x)$：
\[
\begin{aligned}
\Phi(z,u)&=\prod\bigl(z-u_1x_1^{(k)}-\cdots-u_nx_n^{(k)}\bigr)\\
&=z^h+\sum G_{\alpha\alpha_1\ldots\alpha_n}(x)
 z^{\alpha}u_1^{\alpha_1}\ldots u_n^{\alpha_n}
 \quad(\alpha\ne h;\ \alpha+\alpha_1+\cdots+\alpha_n=h).
\end{aligned}
\]

2. 在 $h$ 能被 $p$ 整除的情形，並不必然每個群不變量都是各組 $(x^{(k)})$ 的對稱函數，因為現在不能再由 $h\cdot f(x)$ 推出 $f(x)$。但是，由有限多個不變量 $G_{\alpha\alpha_1\ldots\alpha_n}(x)$——即伽羅瓦預解式的係數——導出的環 $\mathfrak R$，構成所有不變量的系統 $\mathfrak J$ 的一個有限生成子環；還需證明，$\mathfrak J$ 關於 $\mathfrak R$ 滿足有限性判別準則的前提。

首先，不失一般性\footnote{見第 33 頁腳註。}，可把由全部代換系數 $a_{\mu\nu}^{(k)}$ 導出的 $\mathsf{P}$ 的子域 $\overline{\mathsf{P}}$ 作為係數域。這個由有限多個元素導出的域，是由向素域添加有限多個代數元素或超越元素而得到的；並且由於素域是完美的，遂有 $\overline{\mathsf{P}}^{1/p}$ 關於 $\overline{\mathsf{P}}$ 有限。\footnote{見第 32 頁腳註 2。}

還要證明 $\mathfrak J$ 在 $\mathfrak R$ 上整。按 1.，在 $h$ 個組 $x^{(k)}$ 中也出現原來的 $x_1,\ldots,x_n$；因此，當
\[
 z=u_1x_1+\cdots+u_nx_n
\]
時，伽羅瓦預解式 $\Phi(z,u)$ 關於 $u$ 恆等為零。每次把除 $u_i$ 外的所有 $u$ 特化為零，並令 $u_i$ 等於單位元，就得到一個關係
\[
 x_i^h+\sum x_i^\alpha G_{\alpha00\ldots\alpha_i\ldots00}(x)=0
 \quad (\alpha\ne h;\ \alpha+\alpha_i=h),
\]
它表明每個 $x_i$ 都在 $\mathfrak R$ 上整。於是對於 $x$ 中任意多項式也同樣成立；特別地，$\mathfrak J$ 在 $\mathfrak R$ 上整。因此有限性判別準則的前提滿足，\srcemph{不變量的有限生成性得證}。

3. 還應指出，同樣的推理給出\srcemph{相對不變量與模不變量的有限生成性}。若一個多項式 $f(x)$ 的所有 $f(x^{(k)})$ 都只與 $f(x)$ 相差一個來自 $\overline{\mathsf{P}}$ 的非零因子，則稱 $f(x)$ 為該群的一個\srcemph{相對不變量}。絕對不變量的系統——特別是由 $G_{\alpha\alpha_1\ldots\alpha_n}(x)$ 導出的環 $\mathfrak R$——因而構成由所有相對不變量導出的環的一個有限生成子環；並且如上所述，關於 $\mathfrak R$ 的有限性判別準則前提滿足。

若對於來自 $\overline{\mathsf{P}}$ 的每個特殊取值組 $x_1=\alpha_1,\ldots,x_n=\alpha_n$，都有 $f(x)=f(x^{(k)})$，則稱一個多項式 $f(x)$ 為該群的一個\srcemph{模不變量}。每個絕對不變量同時也是模不變量；但如果 $\overline{\mathsf{P}}$ 只含有有限多個元素，反過來則未必成立。在這裡，關於由 $G_{\alpha\alpha_1\ldots\alpha_n}(x)$ 導出的子環 $\mathfrak R$，有限性判別準則的前提同樣滿足。
\endgroup

%
\clearpage
%
%
\pdfbookmark[0]{30. 代數數域和函數域中理想論的抽象構造}{paper-30}
\section*{30. 代數數域和函數域中理想論的抽象構造}

\noindent\emph{Math. Ann. 96 (1927), S. 26--61.}
\setcounter{footnote}{0}

\begin{center}
Emmy Noether（哥廷根）撰。
\end{center}

下面給出所有這些環的一個抽象刻畫：它們的理想論同代數數域中全體整元的理想論一致；也就是說，它們的理想都能唯一地表示為素理想的冪的乘積。屬於這些環的最著名例子，還包括一個不定元的代數函數域中全體整元的環；更一般地，當有多個不定元時，只要同 Kronecker 一樣限制在最高維數的理想上，也就是過渡到某個商環，即函數域範圍，也屬於這種情形。

對於作為基礎的交換環，與這個理想論完全等價的公理如下：\footnote{關於理想論基本概念的定義，可參見 E. Noether, \emph{Idealtheorie in Ringbereichen}, Math. Ann. 83 (1921), S. 24--66（下文引用為 \emph{Idealtheorie}）。順便說，最本質的概念也會在本文中，特別是在 \S\S\ 1, 4, 5 中，作簡短表述。為求完整，其中也收入了一些對本文直接目的並非必要的內容。}

\begin{description}
\item[I. 除子鏈條件。] 每一條由理想組成的鏈，只要其中每個理想都是前一理想的真除子，就在有限步後中止；換言之，對每條由理想組成的除子鏈，都有一個指標，從該指標起所有理想都相等。
\item[II. 模每個非零理想的倍數鏈條件。] 每一條由理想組成的鏈--其所有理想都是某個固定的、不同於零理想的理想的除子--只要其中每個理想都是前一理想的真倍數，就在有限步後中止。
\item[III.] 乘法單位元的存在。
\item[IV.] 無零因子的環。
\item[V. 商域中的整閉性。] 商域中每個關於該環為整的元素都屬於該環。
\end{description}

從這些公理出發，我一步一步發展出限制越來越強的理想論，直到得到所要求的那個理想論；反過來，我也證明，在那裡所有這些公理事實上都滿足。

由除子鏈條件 I 可推出，正如我在 \emph{Idealtheorie} \S\ 4 中已經證明的，一個理想可以表示為有限多個屬於不同素理想的準素理想的最小公倍數。我在這裡簡短重複這個證明（\S\ 6），因為使用理想商的概念可以使它簡化，並且因為我要更正一個疏漏：原來的證明在某處使用了並未假設的乘法單位元的存在。\footnote{P. Urysohn 提醒我注意到這個疏漏；不過我對證明所作的修改比這裡給出的版本稍繁。這裡的版本出自 E. Artin。他早先已經自行補上了這個缺口，並且也獨立於我想到了用理想商作簡化。--另一個簡化避免了引入“約化表示”；我取自 W. Krull 的一個相關問題：\emph{Algebraische Theorie der Ringe III}, Math. Ann. 92 (1924), S. 181--213, Satz I。}

倍數鏈條件的假設導致（\S\ 7）這樣的結果：在模每個不同於零理想的理想所得的剩餘類環中，除包含所有元素的單位理想外，一個素理想不能有真除子。由此，這些理想的準素分量就唯一確定，但屬於單位理想的那些除外。這個結果以稍弱的形式已見於 Masazo Sono；\footnote{Masazo Sono, \emph{On the Reduction of Ideals}, Memoirs of the College of Science, Kyoto University, Series A, 7 (1924), S. 191--204.} 他假設剩餘類環中存在合成列，這同該環中的“雙鏈條件”等價；我在最後（\S\ 10）會簡短說明這一點。所謂雙鏈條件，我指的是除子鏈條件和倍數鏈條件同時成立，並且不再限制於不同於零理想的理想。

若再滿足公理 III--單位元的存在--則單位理想不會作為關聯素理想出現；因此準素分量唯一確定；它們兩兩互素，而且它們的最小公倍數等於它們的乘積。公理 I, II, III 對代數數域的所有有限階都成立；在這裡，Dedekind（Dirichlet--Dedekind, \emph{Zahlentheorie}, 第 3 版，第 11 補篇，\S\ 172）雖未給出完全展開的證明，卻已經陳述了把一個理想唯一表示為兩兩互素的同類理想（準素分量）的乘積。

由假設 IV--無零因子出現--可知，一個不同於零理想和單位理想的理想的各個不同冪全都彼此不同，並且商域存在。若最後還滿足公理 V，即商域中的整閉性，那麼準素理想--由於 IV 已唯一確定--成為素理想的冪，由此得到通常的分解定理（\S\ 8）。後一個證明依賴於一個由整閉性推出的結論；在數域情形，這個結論已由 Dedekind（第 3 版，\S\ 172）得出：可以把一個真分式元素表示為整元之商，使得分子的平方也不能被分母整除。在後來的敘述中，取代這個結論的--同樣作為整閉性的後果--是遠不如它透明的一般化 Gauss 定理，或相應的、斷言更多一些的模定理；這些定理在應用時都預設基元素，而這裡則直接用理想本身來工作。

在反向證明中 \hbox{\(\S 9\)}，不同於公理 V 的各公理幾乎直接得出；而 V 則來自如下證明：真分式元素總能表示成一種形式，使得分子的任何冪都不能被分母整除，這因而是對原先由 V 得出的結論的加強。不過，這個證明只有在從理想表示推出通常結論之後才完成：也就是模某個固定理想的主理想表示，以及分式理想的理論，即域中的理想商理論。由素理想冪表示推出整閉性這一事實，也已經為 Dedekind 所知，正如第 3 版 \S\ 172 第 522 頁下方的一條註記所顯示的。

公理 I 至 V 表明，通常彼此不同的四種分解--分解為最小的兩兩互素理想、相互素理想、最大準素分量以及不可約理想--在這裡相互重合；反過來，由這種重合以及公理 I--IV 可推出整閉性。對於帶有零因子的環，即使其餘公理滿足--其中 V 應定義為商環中的整閉性--也未必推出這些分解相互重合；有限階按某些理想所得剩餘類環的例子說明了這一點 \hbox{\(\S 9\)}。

在 \S\ 1 中，我概述關於一個環為整的整元理論，一直到 Dedekind 從整閉性得出的結論。在 \S\ 2 中，我說明鏈條件如何從環轉移到由線性形式構成的模，其中這個環作為係數域和乘子範圍。\footnote{這種模定理的形式出自 Prüfer（Neue Begründung der algebraischen Zahlentheorie, \S\ 2, Math. Ann. 94 (1924), S. 198--243），並且比通常的基定理轉移更透明。} 由此我得出（\S\ 3）這樣的結論：當過渡到第一類代數擴張域的整元時，公理 I 至 IV 對每個階都保持有效；而對於由全體整元組成的階，V 也成立。這個事實--對於代數數域當然已知--使開頭提到的函數域可以歸入這裡；特別是，只要簡單證明由多個不定元的整係數多項式導出的函數域範圍滿足這些公理，Kronecker 的理想論也就被完全納入。\footnote{差別只在判別式定理處出現，因為函數域範圍按一個素數所得的剩餘類環會成為一個不完美域，從而可能出現第二類擴張。}

隨後發展的理想論並不預設這些僅用於歸類的 \S\S\ 2 和 3。在 \S\ 4 和 \S\ 5 中，將給出獨立於公理 I 至 V 的基礎；這樣就比原來的奠基更清楚地顯示出，理論的哪些部分不依賴有限性假設。

\subsection*{\S\ 1. 整元理論}

設給定一個交換環\footnote{眾所周知，一個環是在一個集合中給出一個等號關係，並再給出兩個滿足通常法則的運算--加法和乘法--時定義的：加法是結合的、交換的並且可唯一逆轉；乘法是結合的--若所論為交換環則還交換--並且對加法滿足分配律。如果作為基礎的等號定義不是集合論意義上的同一，那麼，為了在每個子系統中保持同一個等號定義，這樣的子系統（子環、模、理想）除了含有某個元素外，還必須含有所有同它相等的元素。同樣，在每個擴張中，必須假設新的等號定義包含舊的等號定義。諸如“有限多個元素”“唯一對應”等概念，都按這個等號定義來理解；例如，“有且只有一個元素”就是指在相等意義下只有一個元素。順便說，始終可以從等號過渡到集合論意義上的同一：把相等的元素歸入一個類，並把這些類作為新元素。這裡關於等號定義的說明出自 R. Hölzer。} $\mT$，它\emph{無零因子}，並具有\emph{乘法單位元 $e$}（公理 III 和 IV）；在 $\mT$ 中選定一個固定的、含有單位元的\emph{子環} $\mR$。模、階、整這些概念以下始終相對於這個固定環 $\mR$；為簡短起見，後面在記號中將省略這一點。\footnote{如果允許零因子，並且代之以在 $\mR$ 中假設除子鏈條件，整元理論仍然成立；按 \S\ 2，這也使下面引理的證明成為可能。不過，整元理論的這種形式後面並不使用，所以只在腳註中指出。所討論的一切定義在零因子存在時仍然保持；如所熟知且容易看出，多數定義還需要更弱的前提。不能保持的是 Dedekind 的表述：“一個數若有一個包絡，就是整的。”因為在零因子出現時，兩個有限模的商未必仍然有限；所以這裡階是直接定義的，而不是作為模商定義的。} 屬於 $\mR$ 的元素用拉丁字母表示，$\mT$ 中一般元素用希臘字母表示。

1. $\mT$ 中元素的一個系統 $\mM$，照通常說法稱為 \emph{$\mR$-模}--簡稱\emph{模}--如果它隨兩個元素 $\alpha$ 與 $\beta$ 一起也含有它們的差，並且隨 $\alpha$ 一起也含有 $r\alpha$，這裡 $r$ 表示 $\mR$ 中任意元素。若 $\mM$ 的每個元素也包含在 $\mN$ 中，則稱 $\mM$ 可被 $\mN$ 整除，記作
\[
        \mM\equiv 0 \pmod \mN,
\]
並稱 $\mM$ 為倍數，$\mN$ 為除子。其元素全都屬於 $\mR$ 的 $\mR$-模，稱為 $\mR$ 中的\emph{理想}。

顯然，$\mT$ 中任意多個模的交--最小公倍數 $[\ldots\mM_i\ldots]$--仍是一個模；同樣，$\mT$ 本身也是模。因此，由 $\mT$ 中任意元素系統 $\Sigma$ \emph{導出的模}唯一存在，它就是所有包含 $\Sigma$ 的模的交。任意一個模系統的和--最大公約數 $(\ldots\mM_i\ldots)$--是由它們的並集導出的模。兩個模的積 $\mM\mN$ 定義為由所有元素 $\mu\nu$ 導出的模，其中 $\mu$ 遍歷 $\mM$ 的全部元素，$\nu$ 遍歷 $\mN$ 的全部元素。因此積滿足結合律和交換律，因為環元素滿足這些律。由 $\mT$ 中的分配律，還可得到模的分配律
\[
        (\ldots,\mM_i,\ldots)\mN=(\ldots,\mM_i\mN,\ldots),
\]
因為正是按 $\mT$ 中的這個分配律，左右兩邊都是由所有元素 $\mu_i\alpha$ 導出的模。

由於 $\mR$ 本身是一個模，$\mR$ 是乘子範圍這一條件可表示為 $\mR\mM\equiv0(\mM)$。又由於按假設 $\mR$ 含有單位元，也有 $\mM\equiv0(\mR\mM)$，從而 $\mM=\mR\mM$。

若至少存在一個由\emph{有限多個}元素組成的系統 $\Sigma$，使得 $\mM$ 由它導出，則稱模 $\mM$ 為\emph{有限}的；$\Sigma$ 的元素 $\mu_1,\ldots,\mu_r$ 稱為
\[
        \mM=(\mu_1,\ldots,\mu_r)
\]
的一個模基。兩個有限模--因而有限多個有限模--的和與積仍然有限，因為它們分別由基元素的並集和基元素的乘積導出。後一事實來自如下表示：$\mM$ 中所有元素 $\mu$ 都可用基表示為 $r_1\mu_1+\ldots+r_k\mu_k$\footnote{若 $\mR$ 不含單位元，則基表示具有 $r_1\mu_1+\ldots+r_k\mu_k+n_1\mu_1+\ldots+n_k\mu_k$ 的形式；其中整數 $n$ 不是環元素，而是重複加法或減法的符號。}，並且來自 $\mT$ 中乘法的交換律。這兩點在這裡都是第一次使用。

2. $\mT$ 中一個含有 $\mR$ 的擴張環--也就是說，包含 $\mR$ 的所有元素--稱為一個 \emph{$\mR$-階}，簡稱\emph{階}。$\mT$ 中任意多個階的交仍是一個階，$\mT$ 也同樣是階；因此，由任意元素系統 $\Sigma$ 導出的階唯一存在。階的和與積可相應地像模那樣定義；特別地，對每個階都有 $\mO^2=\mO$。

每個階同時也是 $\mR$-模；若一個階是有限模，則稱它為有限階。由於和與積通過形成模和、模積而產生，由 1. 知有限階的和與積仍是有限階；並且有\emph{傳遞律}：若 $\mA$ 是有限 $\mR$-階，$\mB$ 是有限 $\mA$-階，則 $\mB$ 也是有限 $\mR$-階。因為 $\mB$ 同時成為 $\mR$-階，從而成為 $\mR$-模。若 $\alpha_1,\ldots,\alpha_r$ 構成 $\mA$ 關於 $\mR$ 的一個模基，$\beta_1,\ldots,\beta_s$ 構成 $\mB$ 關於 $\mA$ 的一個模基，那麼顯然所有乘積 $\alpha_i\beta_k$ 構成 $\mB$ 關於 $\mR$ 的一個模基。

3. $\mT$ 中一個元素 $\alpha$，若\emph{由 $\alpha$ 導出的階 $\mR_\alpha$ 是有限的}，則稱它\emph{關於 $\mR$ 為整，或簡稱整}；換言之，若模
\[
        \mA_\mu=(\alpha^0,\alpha,\ldots,\alpha^{\mu-1})
\]
的除子鏈在有限步後中止，即 $\mA_n=\mA_{n+1}=\ldots=\mA_{n+\nu}=\ldots$，也稱 $\alpha$ 為整。

這個定義還允許在 4. 中使用的第二種表述：$\mT$ 中元素 $\alpha$ 稱為整，如果至少存在一個不同於零模的主模 $\mC$--即 $\mC$ 由某個元素 $\gamma\ne0$ 導出--使得階 $\mR_\alpha$ 與 $\mC$ 的積成為一個有限模；換言之，使得模鏈 $\mC\mA_\mu$ 在有限步後中止。\footnote{如果 $\mR$ 中允許零因子，則必須要求存在一個正則主模；也就是說，必須假設 $\gamma$ 為非零因子，即正則元素；對於無零因子的環，這就歸結為 $\gamma\ne0$。}

最後，這個定義也同通常的表述相同：$\mT$ 中元素 $\alpha$ 稱為整，如果它至少滿足一個方程
\[
        \alpha^n+r_1\alpha^{n-1}+\ldots+r_n=0,
\]
其中 $r$ 為 $\mR$ 中元素。

這些不同定義事實上相同。因為 $\mR_\alpha$ 同由 $\alpha$ 的所有冪導出的模一致，所以在 $\mR_\alpha$ 有限時，總可從這些冪中選取一個模基。如果這樣一個基例如由 $e=\alpha^0,\alpha,\ldots,\alpha^{n-1}$ 給出，那麼這意味着 $\mA_\mu$ 的鏈在 $\mA_n$ 處中止；反過來，由 $\mA_n=\mA_{n+\nu}\ldots$ 可得到這個基。而 $\mA_\mu$ 的鏈中止，也同方程
\[
        \alpha^n+r_1\alpha^{n-1}+\ldots+r_n=0
\]
的成立相同。若 $\mR_\alpha$ 有限，則 $\mC\mR_\alpha$ 也有限，因而 $\mC\mA_\mu$ 的鏈中止；另一方面，由 $\mC\mR_\alpha$ 的有限性，即由 $\mC\mA_\mu$ 的中止，也可推出 $\mA_\mu$ 的中止，從而推出 $\mR_\alpha$ 的有限性。因為 $\mC$ 被假設為不同於零模的主模，所以 $\mC\mA_n=\mC\mA_{n+\nu}$ 也給出 $\mA_n=\mA_{n+\nu}$。

整元的環性質和整依賴的傳遞律建立在下面的

\medskip
\noindent\textbf{引理。} 設 $\mS$ 是 $\mT$ 中的一個有限階，$\alpha$ 是 $\mS$ 中任意元素。那麼由 $\alpha$ 導出的階 $\mR_\alpha$ 是有限的；換言之，有限階中的所有元素都是整的。

證明按通常推理得到。若 $\sigma_1,\ldots,\sigma_n$ 是 $\mS$ 的一個模基，則由於 $\mS$ 具有環性質，$\alpha\sigma_i$ 也屬於 $\mS$。於是有
\[
        \alpha\sigma_i=r_{i1}\sigma_1+\ldots+r_{in}\sigma_n,
\]
由此得到
\[
        \left|r_{ik}-\alpha e_{ik}\right|=0,
\]
其中 $e_{ik}=0$ 當 $i\ne k$，$e_{ii}=e$；這就證明了 $\alpha$ 為整。為了得到行列式消失，要用到 $\mT$，從而 $\mS$，被假設為無零因子。\footnote{如果在 $\mR$ 中假設除子鏈條件，但允許零因子，那麼該引理是 \S\ 2 中模定理的直接後果；按該定理，鏈條件可轉移到由 $\mS$ 構成的模上。對於數域而言，這個證明也出自 Dedekind（第 4 版，\S\ 173, III），並且是文獻中對鏈條件的第一次應用。在 4. 下將給出的結論中，Dedekind 取而代之使用了一個更特殊的事實，即代數數域中模一個理想的剩餘類個數是有限的。}

下列系統 $\mS$，即 $\mT$ 中全體整元所成的系統，構成一個階。$\mS$ 包含 $\mR$；因為 $\mR$ 中的元素都是整的，這是由於 $\mR$ 以單位元為基構成一個有限 $\mR$-模。$\mS$ 也具有環性質；因為由任意兩個整元 $\alpha,\beta$ 導出的階 $\mR_{\alpha,\beta}$ 等於積 $\mR_\alpha\cdot\mR_\beta$，因而是有限的。因此，按輔助定理，$\alpha\pm\beta$ 和 $\alpha\beta$ 作為有限階中的元素都是整的。

\medskip
\noindent\textbf{整依賴的傳遞律。} 若 $\mS$ 是由整元組成的階，那麼 $\mT$ 中每個關於 $\mS$ 為整的元素，也關於 $\mR$ 為整。按定義，$\alpha$ 滿足一個方程
\[
        \alpha^m+\sigma_1\alpha^{m-1}+\ldots+\sigma_m=0,
\]
因而也關於由 $\sigma_1,\ldots,\sigma_m$ 導出的階 $\overline{\mS}=\mR_{\sigma_1,\ldots,\sigma_m}$ 為整；而這個階作為積 $\mR_{\sigma_1}\ldots\mR_{\sigma_m}$ 是有限的。按 2. 中給出的傳遞律，由 $\alpha$ 導出的有限 $\overline{\mS}$-階 $\overline{\mS}_\alpha$ 於是也是一個有限 $\mR$-階，並且 $\alpha$ 作為有限階中的元素按引理為整。

4. 若 $\mT$ 中每個關於 $\mR$ 為整的元素都屬於 $\mR$，則稱環 $\mR$ \emph{在 $\mT$ 中整閉}。按 3. 中關於整元定義的第二種表述，$\mT$ 中一個元素 $\alpha$ 當且僅當至少存在一個不同於零模的主模 $\mC$，使得 $\mC\mR_\alpha$ 構成一個有限模時，屬於 $\mR$。\footnote{如果在 $\mR$ 中允許零因子，則 $\mC$ 又必須被假設為正則主模；同樣，推論 I 中的元素 $c$ 也必須被假設為正則。}

Dedekind（第 3 版，\S\ 172）從整閉環的概念得出兩個重要結論，使這個概念可以用於理想論：

\medskip
\noindent\textbf{Dedekind 推論 I。}\emph{設 $\mR$ 在 $\mT$ 中整閉，並且進一步假定 $\mR$ 中的理想滿足除子鏈條件（公理 I）。若 $\beta$ 是 $\mT$ 中的一個非整元素，$c\ne0$ 是 $\mR$ 中的一個元素，則存在一個指數 $\sigma\geqq0$，使元素列
\[
        c,\;c\beta,\ldots,c\beta^r,\ldots
\]
中的 $c,\ldots,c\beta^\sigma$ 都是整的，而其餘所有元素都不是整的。}

若所有元素 $c\beta^r$ 都是整的，則由於 $\mR$ 的整閉性，它們都屬於 $\mR$。於是模鏈 $\mC\mB_1,\mC\mB_2,\ldots,\mC\mB_r,\ldots$——其中 $\mC$ 表示由 $c$ 導出的模，$\mB_r$ 表示由 $\beta^0,\ldots,\beta^{r-1}$ 導出的模——便成為理想鏈 $\mc_1,\ldots,\mc_r,\ldots$，其中 $\mc_r=(c,c\beta,\ldots,c\beta^{r-1})$，這條鏈依假設在有限步後終止。但這樣一來，$\mR_\beta$ 就會是有限的，$\beta$ 就會是整的，與假設矛盾；所以各 $c\beta^r$ 中必有非整元素。此外，只要其中一個元素不是整的，後面所有元素就都不是整的；因為若 $c\beta^\tau$ 是整的，因而屬於 $\mR$，則對 $\rho<\tau$，$c\beta^\rho$ 滿足方程
\[
        (c\beta^\rho)^\tau-c^{\,\tau-\rho}(c\beta^\tau)^\rho=0,
\]
因而也是整的。因此，若 $c\beta^{\sigma+1}$ 是第一個非整元素，則因為 $c$ 是整的，必有 $\sigma\geqq0$，這就是所求指數。

把擴張環 $\mT$ 專門化為 $\mR$ 的商域--也就是通過添入所有元素對而得到的域\footnote{按 Steinitz 所熟知的表述，J. f. M. 137。如果在 $\mR$ 中允許零因子，則應以商環代替商域，即添入所有分母為 $\mR$ 中正則元素的商對。這裡 Dedekind 推論 II 不再成立；因為 $n$ 一般不會是正則元素。}--就得到

\medskip
\noindent\textbf{Dedekind 推論 II。}\emph{設 $\mR$ 在其商域 $\mK$ 中整閉，並且 $\mR$ 中的理想滿足除子鏈條件。若 $\beta$ 是 $\mK$ 中的一個非整元素，則 $\beta$ 可表示為 $\mR$ 中元素之商：
\[
        \beta=m/n,
\]
並且還可使 $m^2/n$ 也不是整的。}

置 $\beta=b/c$，其中 $c\ne0$。在列 $c,c\beta=b,c\beta^2\ldots$ 中，前兩個元素必定都是整的，所以推論 I 中的指數 $\sigma\geqq1$。置 $m=c\beta^\sigma,\qquad n=c\beta^{\sigma-1}$；由於 $\sigma\geqq1$，它們都是該列中的整元；並且 $m/n=\beta$，而 $m^2/n=c\beta^{\sigma+1}$ 不是整的。

\emph{整閉性的傳遞律}具有如下形式：若 $\mR$ 是 $\mT$ 中任意環，並令 $\mS$ 表示 $\mT$ 中所有關於 $\mR$ 為整的整元組成的系統，那麼 $\mS$ 也由 $\mT$ 中所有關於 $\mS$ 為整的整元組成。因為每個關於 $\mS$ 為整的整元，按 3. 中的傳遞律，也關於 $\mR$ 為整，因此屬於 $\mS$。

\subsection*{\S\ 2. 有限模範圍中的鏈條件}

使鏈條件得以轉移的模定理，在這裡要給出的應用中，只對一個擴張環的模出現。不過，在一般模範圍的情形中證明仍然相同，所以以下就以這樣一個模範圍為基礎。

元素 $\alpha,\beta\ldots$ 的一個系統 $M$ 稱為\emph{關於環 $\mR$ 的模範圍}，如果在 $\mathsf{M}$ 中給定了兩個運算，並且這兩個運算都唯一地給出 $\mathsf{M}$ 中的元素：一個是元素之間的加法結合，另一個是用 $\mR$ 中元素乘 $\mathsf{M}$ 中元素；如果 $\mathsf{M}$ 對加法構成一個阿貝爾群，而乘法滿足結合律；並且如果分配律以兩種形式成立。\footnote{參見 \emph{Idealtheorie}, \S\ 9。}

對於一個模範圍，\S\ 1, 1. 中給出的所有不涉及乘法的模定義顯然都保持有效。特別地，若存在 $M$ 中有限多個元素 $\xi_1,\ldots,\xi_k$，使得 $M$ 等於由 $\xi_1,\ldots,\xi_k$ 導出的模，也就是當 $\mR$ 中存在單位元時，等於所有線性形式 $r_1\xi_1+\ldots+r_k\xi_k$ 組成的系統；而在沒有單位元時，還要加上附加項 $n_i\xi_i$（參見注 8），則稱 $M$ 是一個\emph{有限模範圍}。

\medskip
\noindent\textbf{模定理。}\emph{設 $\mathsf{M}$ 是關於一個帶單位元的交換環 $\mR$ 的有限模範圍，並且 $\mR$ 中理想滿足除子鏈條件，或者相應地滿足倍數鏈條件。則在 $\mathsf{M}$ 中，模滿足相應的除子鏈條件，或者相應地滿足倍數鏈條件。}

證明基於這樣一點：可以把 $\mathsf{M}$ 中的每一個模唯一地配給 $\mR$ 中有限多個理想所組成的一個系統，使得模的一個真除子，或者相應地一個真倍數，每次都對應於至少一個真除子理想，或者相應地一個真倍數理想。

$\mathsf{M}$ 中的一個元素稱為\emph{具有長度 $i$}，如果它至少能以一種方式表示為 $\xi_1,\ldots,\xi_i$ 的線性形式，而不能表示為 $\xi_1,\ldots,\xi_{i-1}$ 的線性形式。現在給 $\mathsf{M}$ 中任意一個模 $\mA$ 配給 $\mR$ 中的 $k$ 個理想 $\mathfrak a_1,\ldots,\mathfrak a_k$：令 $\mA_i$ 表示 $\mA$ 中所有長度 $\leqq i$ 的元素所成的模，並令 $\mathfrak a_i$ 表示 $\mA_i$ 中 $\xi_i$ 的所有係數組成的系統。於是首先得到：

若 $\mB$ 是 $\mA$ 的一個\emph{真除子}，而 $\mathfrak b_1\ldots,\mathfrak b_k$ 是配給 $\mB$ 的理想，則這些 $\mathfrak b$ 中至少有一個是相應 $\mathfrak a$ 的\emph{真除子}。事實上，由 $\mA\equiv0(\mB)$ 得到 $\mA_\lambda\equiv0(\mB_\lambda)$，從而對 $\lambda=1,2,\ldots,k$ 有 $\mathfrak a_\lambda\equiv0(\mathfrak b_\lambda)$。依假設，$\mB$ 中必有不屬於 $\mA$ 的元素，因而也有其中長度最小的元素，設其長度為 $i$。此時 $\mathfrak b_i$ 是 $\mathfrak a_i$ 的\emph{真除子}：若
\[
        \beta=b_1\xi_1+\ldots+b_i\xi_i
\]
是 $\mB$ 中這樣的最小長度元素，則 $b_i\equiv0(\mathfrak b_i)$，但 $b_i\not\equiv0(\mathfrak a_i)$。因為若 $b_i\equiv0(\mathfrak a_i)$，便存在元素
\[
        \alpha=a_1\xi_1+\ldots+b_i\xi_i\equiv0(\mA),
\]
於是存在 $\beta-\alpha\equiv0(\mB)$、$\not\equiv0(\mA)$，而其長度小於 $\beta$，這與 $i$ 的選擇矛盾。

現在給定一條除子鏈，或者相應地一條倍數鏈
\[
        \mA_1,\mA_2,\ldots,\mA_r,\ldots
\]
並假設在 $\mR$ 中相應的鏈條件成立。若 $\mathfrak a_{r1},\ldots,\mathfrak a_{rk}$ 是配給 $\mA_r$ 的理想，則序列
\[
        \mathfrak a_{1\lambda},\mathfrak a_{2\lambda},\ldots,\mathfrak a_{r\lambda},\ldots
\]
對 $\lambda=1,\ldots,k$ 構成除子鏈，或者相應地構成倍數鏈。因此按假設，存在一個指標 $\mu$，使得對每個 $\lambda$，理想 $\mathfrak a_{\mu\lambda}$ 都等於所有後繼理想。於是按上面已經證明的結果，模 $\mA_\mu$ 也等於所有後繼模；由此\emph{鏈條件便轉移過去}。

立即得到如下結論：

\medskip
\noindent\textbf{模定理的推論。}\emph{若在 $\mR$ 中，模每個不同於零理想的理想的倍數鏈條件成立，則在 $\mathsf{M}$ 中，模每個模 $\mC$ 的倍數鏈條件成立，只要配給 $\mC$ 的理想 $\mathfrak c_1,\ldots,\mathfrak c_k$ 全都不同於零理想。}在這些假設下，倍數鏈 $\mathfrak a_{1\lambda},\ldots,\mathfrak a_{r\lambda},\ldots$ 都在有限步後中止。

\medskip
\noindent\textbf{附加說明。} 若 $\mR$ 是沒有單位元的環，則除子鏈條件以及模每個非零理想的倍數鏈條件仍可轉移。為此，對任一模 $\mA$，不考察 $\mA$ 本身，而考察所有下列形式的元素所成的系統 $\overline{\mA}$：
\[
 a_1\xi_1+\ldots+a_k\xi_k+n_1\xi_{k+1}+\ldots+n_k\xi_{2k};
\]
把 $\xi_{k+\lambda}$ 換成 $\xi_\lambda$ 後，它又過渡為 $\mA$。同上面一樣，可以給 $\overline{\mA}$ 配給 $2k$ 個理想 $\mathfrak a_1,\ldots,\mathfrak a_k,\mathfrak n_1,\ldots,\mathfrak n_k$；其中 $\mathfrak a_\lambda$ 是 $\mR$ 中的理想，$\mathfrak n_\lambda$ 是整數環中的理想。由於對 $\mathfrak n_\lambda$ 而言，除子鏈條件以及模每個非零理想的倍數鏈條件都成立，所以對 $\overline{\mA}$，從而對 $\mA$，除子鏈條件的證明仍然有效；對於倍數鏈條件，則必須要求配給 $\mC$ 或相應地配給 $\overline{\mC}$ 的 $2k$ 個理想全都不是零理想。

\subsection*{\S\ 3. 向有限擴張域的過渡}

為了把數域和函數域納入這個框架，還需要說明導言中表述的公理 I 至 V 在過渡到有限擴張域時如何轉移。

1. \emph{設在 $\mathfrak R$ 中，除公理 II--倍數鏈條件--外，公理 I 至 V 全部成立。令 $\mathfrak K$ 表示 $\mathfrak R$ 的商域，令 $\mathfrak L$ 表示 $\mathfrak K$ 的一個第一類有限擴張域。\footnote{按照 Steinitz，一個代數擴張域稱為``第一類''，如果其中每個元素都是某個素函數的零點，而該素函數在適當的擴張域中分解為兩兩不同的線性因子。} 那麼，在 $\mathfrak L$ 中所有關於 $\mathfrak R$ 為整的量所組成的系統 $\mathfrak S$ 中，同樣的公理成立；並且在 $\mathfrak S$ 中的每一個階中，除整閉性外，所有這些公理也仍成立。}

按 \S\ 1, 3.，$\mathfrak S$ 構成一個環，且公理 III 和 IV--單位元的存在以及無零因子環--在其中成立。按 \S\ 1, 4. 的結尾，$\mathfrak S$ 在 $\mathfrak L$ 中整閉；同時 $\mathfrak L$ 成為 $\mathfrak S$ 的商域，因為 $\mathfrak L$ 中每個元素乘以 $\mathfrak R$ 中某個適當元素後都會變成 $\mathfrak S$ 中的元素。因此公理 V 成立。

公理 I 的證明由熟知的推理給出。作為第一類擴張，$\mathfrak L$ 可由在 $\mathfrak K$ 上添入一個單一元素 $\alpha$ 得到；按上面的結果，可以假設 $\alpha$ 關於 $\mathfrak R$ 為整。除 $\alpha$ 外，$\mathfrak L$ 關於 $\mathfrak K$ 的 Galois 域中所含的共軛量 $\alpha',\alpha'',\ldots$ 也為整；因此它們的差積為整，這個差積的平方 $D$ 也為整。由於 $\alpha,\alpha',\ldots$ 全都彼此不同，$D$ 就是 $\mathfrak K$ 中一個不同於零的量，於是由於 $\mathfrak R$ 在 $\mathfrak K$ 中整閉，$D$ 是 $\mathfrak R$ 中的一個量。按通常推理，由 $\mathfrak R$ 的整閉性可知，$\mathfrak S$ 包含於關於 $\mathfrak R$ 的有限模範圍 $M$ 中；這個 $M$ 由元素 $\xi_i=\alpha^i/D$ 生成。為此，只須在把 $\mathfrak S$ 中元素用 $\xi_i$ 表示時過渡到共軛元素，並解出表示係數所滿足的線性方程組。按 \S\ 2 中的模定理，除子鏈條件對 $M$ 中所有 $\mathfrak R$-模成立，因而特別對 $\mathfrak S$ 中所有理想成立。同樣，除子鏈條件也對 $\mathfrak S$ 中任意階的理想成立，因為這裡同樣是在處理 $M$ 中的 $\mathfrak R$-模；而且對每個階，公理 III 和 IV 也成立。

2. \emph{若在 $\mathfrak R$ 中公理 I 至 V 全部成立，則同樣的結論也對 $\mathfrak S$ 成立。在 $\mathfrak S$ 中的每一個階中，除整閉性外，所有公理仍成立。}

考慮到 1.，只剩下證明公理 II 成立。按 \S\ 2 中模定理的推論，只需證明如下事實：若把 $\mathfrak S$ 中一個不同於零理想的理想看作 $\mathsf{M}$ 中的 $\mathfrak R$-模 $\mathfrak C$，則配給 $\mathfrak C$ 的 $\mathfrak R$ 中理想 $\mathfrak c_i$ 全都不同於零理想。現在，$\mathfrak C$ 關於 $\mathfrak K$ 具有同 $\mathsf{M}$ 相同的有限線性秩；因為 $\mathfrak C$ 除含有任意一個元素 $\gamma$ 外，也含有 $\gamma\alpha^i$。因此，對每個 $\xi_i$，存在 $\mathfrak R$ 中一個元素 $c_i\ne0$，使得 $c_i\xi_i$ 屬於 $\mathfrak C$；由於由 $\xi_i$ 給出的表示具有唯一性--這裡指標只取小於域次數的那些值--$c_i$ 屬於 $\mathfrak c_i$，所以 $\mathfrak c_i$ 不同於零理想。這個考慮對所有同 $\mathsf{M}$ 具有相同秩的階仍然有效；對於較低秩的階，則過渡到它的商域，這個商域作為 $\mathfrak K$ 與 $\mathfrak L$ 之間的中間域也屬於第一類，並且它關於 $\mathfrak K$ 的次數同該階的線性秩一致；因此上述考慮仍然保持有效。最後還要指出，$\mathfrak S$ 中不同於 $\mathfrak S$ 的階決不會是整閉的；因為每個階也包含 $\mathfrak R$，若它整閉，就必須也包含 $\mathfrak S$。

因此，為了把數域和函數域歸入這裡，只需證明這些基本範圍--有理整數、一個不定元的多項式以及多個不定元多項式的函數域範圍--滿足公理 I 至 V。

3. \emph{對於有理整數環 $\mathfrak R$，以及對於以一個域為係數的一個不定元多項式環 $\mathfrak R$，公理 I 至 V 成立。} 這裡，I 和 II 或者來自這樣一個事實：模每個不同於零的數，或者模這樣一個不定元多項式，只有有限多個，或者相應地只有有限多個線性無關的剩餘類；也可來自每個理想都是主理想這一事實。因為由此得到理想唯一分解為素理想冪的乘積，於是一個不同於零理想的理想只可被有限多個理想整除。公理 III 和 IV 顯然成立；後者是等號定義的一個後果。整閉性 V 來自如下事實：由於 $\mathfrak R$ 中元素到單位--一的因子--為止可唯一分解為不可分解元，商域中不屬於 $\mathfrak R$ 的每個元素都可表示為 $\mathfrak R$ 中兩個互素元素的商；因此分子的任何冪都不能被分母整除。這樣的元素因而不能滿足任何把它刻畫為關於 $\mathfrak R$ 為整的方程。

4. \emph{現在令 $\mathfrak R$ 表示以一個域為係數，或者以有理整數為係數的多個不定元全體多項式所成的環。} 那麼，公理 III、IV、V 如 3. 中一樣成立；因為這裡元素到單位為止分解為不可分解元的分解也唯一。除子鏈條件 I 成立，是 Hilbert 理想基存在定理的一個直接後果；而倍數鏈條件 II 在這裡並不成立。

不過，通過過渡到\emph{函數域範圍}，也就是相當於限制在最高維數的理想上，也還能得到公理 II 的成立；這時 I 的成立可像 3. 中那樣證明，而無需援引 Hilbert 定理。事實上，向 $\mathfrak R$ 添入一個不定元 $u$，也就是考察以 $\mathfrak R$ 中元素為係數的 $u$ 的全體多項式所成的環 $\mathfrak R^*$，其中等號按係數相等來定義。若像通常那樣，把 $\mathfrak R^*$ 中一個多項式稱為關於 $\mathfrak R$ 本原的，是指它的係數的最大公因子--這裡的因子是 $\mathfrak R$ 中多項式意義下的因子，而不是理想因子--成為 $\mathfrak R$ 的一個單位，那麼 $\mathfrak R^*$ 中每個多項式都是 $\mathfrak R$ 中一個元素同 $\mathfrak R^*$ 中一個本原多項式的乘積，並且 $\mathfrak R^*$ 中本原多項式的乘積仍然本原。

\emph{$\mathfrak R^*$ 的商域中所有分母為 $\mathfrak R^*$ 中本原多項式的元素的全體，因此構成一個環，即 $\mathfrak R$ 的函數域範圍 $\mathfrak F$。} 在這個函數域範圍 $\mathfrak F$ 中，單位恰好是那些分子和分母都是 $\mathfrak R^*$ 中本原多項式的元素；因此 $\mathfrak F$ 中每個元素都能唯一地表示為 $\mathfrak R$ 中一個元素同 $\mathfrak F$ 中一個單位的乘積。由此，$\mathfrak R$ 中元素的分解定理轉移到 $\mathfrak F$ 中元素上；因此，對 $\mathfrak F$ 而言，除公理 III 和 IV 外，V 也如 3. 中一樣成立。

此外，在 $\mathfrak F$ 中每個理想都是主理想。事實上，一個理想除了包含 $\mathfrak F$ 中任意一個元素外，也包含由它乘以一個單位而得到的 $\mathfrak R$ 中元素；並且除了包含 $\mathfrak R$ 中任意兩個元素外，也包含它們的最大公因子。因為除 $f(x)=t(x)\overline f(x)$ 和 $g(x)=t(x)\overline g(x)$ 外，理想中還含有 $t(x)(\overline f(x)+u\overline g(x))$；若 $\overline f$ 與 $\overline g$ 互素，並且它們的這個線性組合因而成為一個單位，則理想中也含有 $t(x)$。因此，從理想中任意一個元素 $f(x)$ 出發，若理想中有一個不被它整除的 $g(x)$，則最大公因子 $t(x)$ 也屬於該理想，併成為 $f(x)$ 的一個真除子；有限次重複這一過程便得到理想的一個基元素，從而該理想被識別為主理想。於是，在 $\mathfrak F$ 中理想作為素理想冪乘積的唯一分解成立，它對應於 $\mathfrak R$ 中元素的分解；並且由此像 3. 中一樣得到公理 I 和 II 的成立。這裡所考慮的 $\mathfrak F$ 的商域的擴張域，只涉及那些可通過添入 $\mathfrak S$ 中一個元素而生成的擴張域。\footnote{這些擴張域中關於 $\mathfrak F$ 為整的量組成的系統，也可通過從 $\mathfrak S$ 過渡到一個相應的函數域範圍來得到，正如從 $\mathfrak R$ 過渡到 $\mathfrak F$ 一樣。參見 J. König, \emph{Algebraische Größen}, Leipzig 1903, S. 468/69。} 因此，結合 1. 和 2.，就已經認明數域和函數域滿足公理 I 至 V。

\subsection*{\S\ 4. 同構定理。直和}

以下只假定環性質或模性質，而不預設其他公理。

1. 若 $\mathsf{M}$ 和 $\overline{\mathsf{M}}$ 是關於 $\mR$ 的模範圍（\S\ 2），則稱 $M$ 同態於 $\overline M$（更準確地說，模同態），記作 $\mathsf{M}\sim\overline{\overline{\mathsf{M}}}$，如果 $M$ 中每個元素都對應於 $\overline M$ 中一個且只有一個元素，並且這種對應使 $\overline M$ 被窮盡；\footnote{均按照 $\overline M$ 中採用的等號定義；參見注 6。} 同時，在這種對應下，差以及同 $\mR$ 中同一元素相乘都彼此對應；也就是說，若 $\beta\sim\overline\beta$ 且 $\gamma\sim\overline\gamma$，則總有 $(\beta-\gamma)\sim(\overline\beta-\overline\gamma)$ 和 $r\beta\sim r\overline\beta$。

於是，$\mathsf{M}$ 中一個 $\mR$-模 $\mB$ 同態地對應於 $\overline M$ 中一個 $\mR$-模 $\overline{\mB}$；而 $\mathsf{M}$ 中所有同 $\overline{\mathsf{M}}$ 中某個 $\mR$-模 $\overline{\mC}$ 的元素相對應的元素全體 $\mC^*$，在 $\mathsf{M}$ 中構成一個由 $\overline{\mC}$ 唯一確定的 $\mR$-模，即同 $\overline{\mC}$ 相配的模 $\mC^*$，並且 $\mC^*\sim\overline{\mC}$。特別地，若 $\mA$ 是同 $\overline{\mathsf{M}}$ 中零元相配的模，則 $\mC^*$ 成為 $\mA$ 的一個除子，並分解為 $\mathsf{M}$ 中模 $\mA$ 同餘的元素類，使這些類同 $\overline{\mC}$ 中元素一一對應。若從 $\mathsf{M}$ 中的 $\mB$ 過渡到 $\overline M$ 中的 $\overline{\mB}$，再從那裡返回到 $\mB^*$，則得到 $\mB^*=(\mB,\mA)$，也就是 $\mB$ 與 $\mA$ 的最大公約數；因此，當 $\mB$ 是 $\mA$ 的除子時，有 $\mB^*=\mB$。

若 $\mathsf{M}$ 與 $\overline{\mathsf{M}}$ 的元素之間的對應可逆且一一，則稱這些範圍同構（更準確地說，模同構），記作 $M\cong\overline M$。

若 $\mA$ 是 $M$ 中任意一個 $\mR$-模，則通過把模 $\mA$ 的同餘作為新的等號關係，可得到一個同 $\mathsf{M}$ 同態的模範圍 $\overline{\mathsf{M}}$，即剩餘類模 $\mathsf{M}\mid\mA$。這時，$M$ 中每個元素所對應的，恰好是 $\overline M$ 中同它相等的全部元素。若在 $\overline M$ 中從等號定義過渡到同一性，這就意味着把 $\overline M$ 中相等的所有元素合併成一個類--剩餘類--並把這些剩餘類看作 $\overline M$ 的新元素\footnote{參見注 9。}。

\emph{每個同態都由過渡到剩餘類模而產生；} 因為若 $\mathsf{M}\sim\overline{\overline{\mathsf{M}}}$，並且 $\mA$ 是同 $\overline{\overline{\mathsf{M}}}$ 中零元相配的模，則如上所示，$\overline{\overline{\mathsf{M}}}$ 同構於剩餘類模 $\mathsf{M}\mid\mA$。

2. \textbf{第一同構定理。} \emph{設 $\overline{\mathsf{M}}$ 表示剩餘類模 $\mathsf{M}\mid\mA$，並且 $\mC$ 是 $\mA$ 的除子，則有同構 $\overline{\mathsf{M}}\mid\overline{\mC}\simeq \mathsf{M}\mid\mC$。} 事實上，模 $\mA$ 的同餘同時也是模 $\mC$ 的同餘；因此，模 $\mA$ 相等的元素在模 $\mC$ 下仍相等。於是可以這樣形成剩餘類模 $\mathsf{M}\mid\mC$--也就是模 $\mC$ 的等號--先把元素按 $\mA$ 等同，即過渡到 $\overline{\mathsf{M}}$，再在其中把模 $\mC$ 相等的元素合併；這在 $\overline{\mathsf{M}}$ 中對應於按 $\overline{\mC}$ 等同，也就是形成 $\overline{\mathsf{M}}\mid\overline{\mC}$。

\textbf{第二同構定理。} \emph{若 $\mB$ 和 $\mA$ 是 $\mathsf{M}$ 中的模，則有同構 $(\mB,\mA)\mid\mA\cong \mB\mid[\mB,\mA]$。} 因為按 1.，當把 $(\mB,\mA)\mid\mA$ 置為 $\overline{\mB}$ 時，$\mB$ 同態於 $\overline{\mB}$；而在這裡，同 $\overline{\mB}$ 中零元相對應的恰好是 $[\mB,\mA]$ 中所有元素，所以仍按 1. 得到上述同構。

3. 若 $\mR$ 和 $\overline{\mR}$ 是（交換）環，則稱 $\mR$ 同態於 $\overline{\mR}$（更準確地說，環同態），記作 $\mR\sim\overline{\mR}$，如果 $\mR$ 中每個元素都對應於 $\overline{\mR}$ 中一個且只有一個元素，並且這種對應使 $\overline{\mR}$ 被窮盡；同時在這種對應下差與積彼此對應。若元素之間的對應可逆且一一，則稱這些環同構（更準確地說，環同構），記作 $\mR\simeq\overline{\mR}$。

若把 1. 和 2. 中的模替換為 $\mR$ 中的理想，\footnote{對於非交換環，必須以雙邊理想為基礎。} 並把模同態和模同構的概念替換為環同態和環同構，則所有考慮仍然保持有效。特別地，若 $\mc$ 是 $\ma$ 的除子，則通過把模 $\ma$ 的同餘引入為新的等號關係，就得到剩餘類環 $\mc\mid\ma$；這裡必須注意，現在剩餘類的乘積也已定義。

於是 $\mc$ 同態於 $\mc\mid\ma$；每個同態都以這種方式產生，並且如 2. 中一樣，同構定理成立：

\textbf{第一同構定理。} 若 $\overline{\mR}$ 表示剩餘類環 $\mR\mid\ma$，並且 $\mc$ 是 $\ma$ 的除子，則在環同構意義下有 $\overline{\mR}\mid\overline{\mc}\simeq\mR\mid\mc$。

\textbf{第二同構定理。} 若 $\mb$ 和 $\ma$ 是 $\mR$ 中的理想，則有環同構 $(\mb,\ma)\mid\ma\simeq\mb\mid[\mb,\ma]$。
\footnote{這些從群論中熟知的同構定理，對於模而言，首先以稍為特殊的形式見於 Dedekind（參見第 3 版，S. 484）。上面關於理想的形式見 Masazo Sono: \emph{On Congruences. I, II, III, IV}, Memoirs of the College of Science, Kyoto Imperial University, 2 (1917), 3 (1918, 1919)，參見 2, S. 215。每次明白說出的都只是第二同構定理。}

4. 下面彙集關於互素理想的計算規則，\footnote{確切地說，是採用可追溯到 Dedekind 的形式（第 4 版，\S\ 178, III 至 VIII）。近來敘述中普遍出現的藉助單位元的計算要繁瑣得多。} 後面 5. 中關於直和的定理將由這些規則推出。在交換環 $\mR$ 中，假定乘法單位元 $e$ 存在。因此，對 $\mR$ 中每個理想都有 $\ma=\mo\ma$，其中 $\mo$ 表示單位理想。兩個理想 $\ma,\mb$ 稱為互素，如果它們的最大公約數 $(\ma,\mb)$ 等於 $\mo$。

\emph{4$\alpha$. 若 $(\ma,\mb)=\mo$，則 $(\ma,\mb\mc)=(\ma,\mc)$。} 因為 $(\ma,\mc)=(\ma,\mb)\cdot(\ma,\mc)=(\ma^2,\ma\mb,\ma\mc,\mb\mc)$ 可被 $(\ma,\mb\mc)$ 整除，反過來也一樣。由此得出：

\emph{4$\beta$. 由 $\mc\mb\equiv0(\ma)$ 和 $(\ma,\mb)=\mo$ 可推出 $\mc\equiv0(\ma)$。} 事實上，$\mc\mb\equiv0(\ma)$ 等價於 $(\ma,\mb\mc)=\ma$；因此由於 $(\ma,\mb)=\mo$，也有 $(\ma,\mc)=\ma$，或者 $\mc\equiv0(\ma)$。又得到：

\emph{4$\gamma$. 若理想 $\ma_1,\ldots,\ma_r$ 中每一個都同理想 $\mb_1,\ldots,\mb_s$ 中每一個互素，則 $\ma_i$ 的乘積同 $\mb_i$ 的乘積互素。} 因為由 $(\ma,\mb)=\mo$ 和 $(\ma,\mc)=\mo$ 得到 $(\ma,\mb\mc)=\mo$，經過有限次重複便得出斷言。

由 4$\alpha$ 和 4$\gamma$ 得到：

\emph{4$\delta$. 若理想 $\mb_1,\ldots,\mb_r$ 兩兩互素，也就是當 $i\ne k$ 時 $(\mb_i,\mb_k)=\mo$，則它們的最小公倍數等於它們的乘積。} 令 $(\ma,\mb)=\mo$，$\mv=[\ma,\mb]$；則 $\mv=\mo\mv=(\ma\mv,\mb\mv)\equiv0(\ma,\mb)$；另一方面 $\ma\mb\equiv0(\mv)$，所以 $\mv=\ma\mb$。結合 4$\gamma$ 將此論證有限次重複，便得到斷言。

\emph{4$\varepsilon$. 若理想 $\mb_1,\ldots,\mb_r$ 兩兩互素，並置 $\ma_i=[\mb_1,\ldots,\mb_{i-1},\mb_{i+1},\ldots,\mb_r]=\mb_1\ldots\mb_{i-1}\mb_{i+1}\ldots\mb_r$，則 $(\ma_1,\ldots,\ma_{i-1},\ma_{i+1},\ldots,\ma_r)=\mb_i$，並且 $(\ma_1,\ldots,\ma_r)=\mo$。} 因為按分配律，$(\ma_1,\ma_2)=\mb_3\ldots\mb_r(\mb_2,\mb_1)=\mb_3\ldots\mb_r$；於是 $(\ma_1,\ma_2,\ma_3)=\mb_4\ldots\mb_r(\mb_3,\mb_1\mb_2)=\mb_4\ldots\mb_r$，再經適當編號下的有限次重複便得到斷言。

$\ma_i$ 稱為表示 $\mm=[\mb_1,\ldots,\mb_r]$ 中 $\mb_i$ 的補理想。

5. 在交換環 $\mR$ 中仍假定乘法單位元存在。若 $\mR$ 中每個元素 $c$ 都能以一種且只有一種方式表示為 $c=a_1+a_2+\ldots+a_r$，其中每個 $a_i$ 是 $\ma_i$ 中元素，則稱環 $\mR$ 是理想 $\ma_1,\ldots,\ma_r$ 的\emph{直和}--記作 $\mR=\ma_1+\ma_2+\ldots+\ma_r$。

\emph{若零理想是兩兩互素的理想 $\mb_1,\ldots,\mb_r$ 的最小公倍數，並且 $\ma_i$ 是 $\mb_i$ 的補理想，則 $\mR$ 是理想 $\ma_1,\ldots,\ma_r$ 的直和。} 因為由 $\mo=(\ma_1,\ldots,\ma_r)$ 可知，$\mR$ 中每個元素至少有一個通過 $\ma_i$ 給出的加法表示；又由於 $[\ma_i,\mb_i]=(0)$，這個表示是唯一的。事實上，由 $0=a_1+a_2+\ldots+a_r=a_i+b_i$ 可得 $a_i=0$。按第二同構定理，$\ma_i$ \emph{同構於剩餘類環 $\mR\mid\mb_i$}。並且有``正交關係'' $\ma_i\ma_k=(0)$（當 $i\ne k$）以及 $\ma_i^2=\ma_i$；後者來自 $\ma_i=\mo\ma_i$。

\emph{反過來，若存在把 $\mR$ 表示為直和的表示 $\mR=\ma_1+\ma_2+\ldots+\ma_r$，並置 $\mb_i=(\ma_1,\ldots,\ma_{i-1},\ma_{i+1},\ldots,\ma_r)=\ma_1+\ldots+\ma_{i-1}+\ma_{i+1}+\ldots+\ma_r$，則 $\mb_i$ 兩兩互素，並且它們的最小公倍數等於零理想。} 因為由 $\mo=(\ma_i,\mb_i)$ 得到，當 $k\ne i$ 時有 $\mo=(\mb_k,\mb_i)$，這是由於 $\mb_k$ 是 $\ma_i$ 的除子。進一步，若 $c$ 可被所有 $\mb_i$ 整除，則有
\[
 c=a_2^{(1)}+\ldots+a_r^{(1)}=a_1^{(2)}+a_3^{(2)}+\ldots+a_r^{(2)}=\ldots=a_1^{(r)}+\ldots+a_{r-1}^{(r)},
\]
其中每個 $a_i^{(\lambda)}\equiv0\pmod{\ma_i}$。由表示的唯一性可知，由於每次都有一個分量為零，對每個 $\lambda$ 都有 $0=a_1^{(\lambda)},\ldots,0=a_r^{(\lambda)}$，從而 $c=0$。\footnote{5. 中給出的定理是關於直和與直接交之間關係的一般定理的特殊情形。參見 H. Prüfer, Theorie der Abelschen Gruppen I, Math. Zeitschr. 20 (1924), S. 165--187, \S\ 6。那裡給出的考慮，即使在群概念作如此一般的表述以至模和理想都作為特殊情形歸入其中時，也仍保持有效。}

\subsection*{\S\ 5. 素理想和準素理想}

仍設 $\mR$ 是一個交換環，不預設任何進一步的公理--甚至不預設單位元存在。下面彙集關於素理想和準素理想的基本事實。\footnote{在 \emph{Idealtheorie}, \S\ 4 中，預設了公理 I，即除子鏈條件；因此在那裡並沒有清楚凸顯出素理想和準素理想的哪些性質獨立於這個公理。}

1. $\mR$ 中的理想 $\mpideal$ 稱為\emph{弱素理想}，如果剩餘類環 $\mR\mid\mpideal$ 是無零因子環；也就是說，由 $a\not\equiv0(\mpideal)$ 和 $b\not\equiv0(\mpideal)$ 總能推出 $ab\not\equiv0(\mpideal)$。$\mR$ 中的理想 $\mpideal^*$ 稱為\emph{強素理想}，如果剩餘類環 $\mR\mid\mpideal^*$ 沒有\emph{零因子理想}；也就是說，由 $\ma\not\equiv0(\mpideal^*)$ 和 $\mb\not\equiv0(\mpideal^*)$ 總能推出 $\ma\mb\not\equiv0(\mpideal^*)$。

\emph{每個強素理想同時也是弱素理想}，這是把 $\ma,\mb$ 特化為主理想即可看出的。反過來，\emph{每個弱素理想也同時是強素理想}；因此可以徑稱\emph{素理想}。事實上，設 $\mpideal$ 是弱素理想；設 $\ma\not\equiv0\pmod{\mpideal}$ 且 $\mb\not\equiv0\pmod{\mpideal}$；再取 $\ma$ 與 $\mb$ 中的元素 $a$ 與 $b$，使得 $a\not\equiv0\pmod{\mpideal}$ 且 $b\not\equiv0\pmod{\mpideal}$。於是 $ab\not\equiv0\pmod{\mpideal}$，從而 $\ma\mb\not\equiv0\pmod{\mpideal}$；所以 $\mpideal$ 也是強素理想。

2. $\mR$ 中一個理想 $\mq$ 稱為\emph{弱準素理想}，如果在剩餘類環 $\mR\mid\mq$ 中每個零因子的某個冪都消失；也就是說，若由 $a\not\equiv0\pmod{\mq}$ 以及對每個 $\varkappa$ 都有 $b^\varkappa\not\equiv0\pmod{\mq}$，總能推出 $ab\not\equiv0\pmod{\mq}$。$\mR$ 中所有在 $\mR\mid\mq$ 中成為零因子的元素所組成的系統 $\mpideal$ 構成一個理想，而且是一個素理想；它成為 $\mq$ 的除子，稱為\emph{相伴素理想}。因為隨 $a$ 一起，$ra$ 也是零因子；隨 $a$ 和 $b$ 一起，$a-b$ 也是零因子，後者是因為，只要 $a^\varkappa$ 和 $b^\lambda$ 都可被 $\mq$ 整除，$(a-b)^{\varkappa+\lambda}$ 也可被 $\mq$ 整除。若再設 $a$ 不是零因子，則按定義 $a$ 的任何冪也不是零因子；因此，由 $a\not\equiv0(\mpideal)$ 和 $b\not\equiv0(\mpideal)$ 得到 $a^\varkappa b^\varkappa=(ab)^\varkappa\not\equiv0(\mq)$，從而 $ab\not\equiv0(\mpideal)$。

$\mR$ 中一個理想 $\mq^*$ 稱為\emph{強準素理想}，如果在剩餘類環 $\mR\mid\mq^*$ 中每個零因子理想的某個冪都消失；也就是說，若由 $\ma\not\equiv0\pmod{\mq^*}$ 以及對每個 $\varkappa$ 都有 $\mb^\varkappa\not\equiv0\pmod{\mq^*}$，總能推出 $\ma\mb\not\equiv0\pmod{\mq^*}$。同弱準素理想一樣可以證明：$\mR$ 中所有在 $\mR\mid\mq^*$ 中成為零因子理想的理想的最大公約數，構成一個素理想，它成為 $\mq^*$ 的除子，即\emph{相伴素理想}。反過來，所有具有同一個相伴素理想 $\mpideal$ 的準素理想，都說成``屬於 $\mpideal$''。

\emph{每個強準素理想同時也是弱準素理想}，這是把 $\ma,\mb$ 特化為主理想即可看出的。\emph{但一般說來，反命題並不成立}。\footnote{下面這個由 R. Hölzer 告知我的例子說明了這一點。令 $\mR$ 為以一個域為係數的可數多個不定元 $x_i$ 的多項式範圍；令 $\mq=(x_1^2,x_2^3,\ldots,x_\nu^{\nu+1},\ldots,x_i x_k\,[i\ne k]\ldots)$。剩餘類環中的零因子恰好是所有可被 $\mpideal=(x_1,x_2,\ldots,x_\nu,\ldots)$ 整除的多項式，並且每一個這樣的零因子的某個冪都可被 $\mq$ 整除；因此 $\mq$ 是以 $\mpideal$ 為相伴素理想的弱準素理想。可是 $\mq$ 不是強準素理想；因為若置 $\ma=(x_1,x_3,x_5,\ldots,x_{2\nu+1},\ldots)$ 且 $\mb=(x_2,x_4,\ldots,x_{2\nu},\ldots)$，則 $\ma\mb\equiv0(\mq)$，但 $\ma$ 或 $\mb$ 的任何冪都不能被 $\mq$ 整除。} \emph{然而，若在 $\mR$ 中預設公理 I，即除子鏈條件，則每個弱準素理想同時也是強準素理想；} 因此可以徑稱\emph{準素理想}。於是設 $\mq$ 是弱準素理想；設 $\ma\not\equiv0\pmod{\mq}$，並且對每個 $\varkappa$ 都有 $\mb^\varkappa\not\equiv0\pmod{\mq}$。那麼存在 $\ma$ 中元素 $a$ 和 $\mb$ 中元素 $b$，使得 $a\not\equiv0\pmod{\mq}$，且對每個 $\varkappa$ 都有 $b^\varkappa\not\equiv0\pmod{\mq}$；於是 $ab\not\equiv0\pmod{\mq}$，從而 $\ma\mb\not\equiv0\pmod{\mq}$。因為如果對 $\mb$ 中每個 $b$ 都存在一個指數 $\lambda$，使得 $b^\lambda\equiv0(\mq)$，則有 $\mb^{\lambda_1+\ldots+\lambda_r}\equiv0(\mq)$，其中 $\lambda_1,\ldots,\lambda_r$ 表示有限多個基元素的指數。\footnote{為了能從公理 I 推出有限理想基，必須在 $\mR$ 中預設一個固定的良序（參見 \S\ 6）。} 後一個考慮還表明：存在一個（最小）指數 $\varrho$，使得 $\mpideal^\varrho\equiv0(\mq)$，其中 $\mpideal$ 表示相伴素理想；$\varrho$ 稱為 $\mq$ 的\emph{指數}。

3. 以下再次\emph{不預設}除子鏈條件。有限多個理想的最小公倍數 $[\ma_1,\ldots,\ma_r]$ 稱為一個\emph{最短表示}，如果沒有任何 $\ma_i$ 進入其餘理想的最小公倍數，也就是沒有任何 $\ma_i$ 可以刪去。

\emph{有限多個屬於同一素理想 $\mpideal$ 的弱準素理想，或者相應地強準素理想，其最小公倍數仍然是以 $\mpideal$ 為相伴素理想的弱準素理想。一個由有限多個屬於不同素理想的弱準素理想，或者相應地強準素理想，所給出的最短表示，不是弱準素理想，或者相應地不是強準素理想。} 設 $\mf=[\mq_1,\ldots,\mq_r]$，並設 $\mpideal$ 是弱準素理想 $\mq_i$ 的相伴素理想；那麼 $\mpideal$ 也由所有某個冪可被 $\mf$ 整除的元素全體組成。由 $a\not\equiv0\pmod{\mf}$ 以及對每個 $\varkappa$ 都有 $b^\varkappa\not\equiv0\pmod{\mf}$，可推出 $b\not\equiv0\pmod{\mpideal}$，並且對至少一個指標 $i$ 有 $a\not\equiv0\pmod{\mq_i}$；於是 $ab\not\equiv0\pmod{\mq_i}$，從而 $ab\not\equiv0\pmod{\mf}$。若把所有元素都替換為理想，就不能再推出 $\mb\not\equiv0\pmod{\mpideal}$。

現在設 $\mm=[\mq_1,\ldots,\mq_r]$ 是一個由弱準素理想給出的最短表示，並且當 $i\ne k$ 時 $\mpideal_i\ne\mpideal_k$，又設 $r\geqq2$。那麼至少有一個相伴素理想，例如 $\mpideal_1$，不包含於其餘任何一個之中；因為 $\mpideal$ 的每一條真倍數鏈都必須至多經過 $r$ 步而中止。因此，存在 $\mpideal_i$ 中元素 $a_i$，使得 $a_i^{\varrho_i}\equiv0\pmod{\mq_i}$，並且 $a_2\not\equiv0\pmod{\mpideal_1},\ldots,a_r\not\equiv0\pmod{\mpideal_1}$。若進一步 $q_1\not\equiv0\pmod{\mm}$ 是 $\mq_1$ 中的一個元素，則 $q_1 a_2^{\varrho_2}\ldots a_r^{\varrho_r}\equiv0\pmod{\mm}$，但 $(a_2^{\varrho_2}\ldots a_r^{\varrho_r})$ 的任何冪都不能被 $\mm$ 整除；否則會推出可被 $\mpideal_1$ 整除。因此 $\mm$ 不是弱準素理想。若把所有元素都替換為理想，也就是把 $q_1$ 替換為 $\mq_1$，並把 $a_i$ 替換為 $\ma_i$，其中 $\ma_i\not\equiv0\pmod{\mpideal_1}$ 且 $\ma_i^{\varrho_i}\equiv0\pmod{\mq_i}$，就得到強準素理想情形的證明。\footnote{這個使除子鏈條件不再需要的、對我原來證明的修改，要感謝 B. L. van der Waerden。}

4. \emph{模商和理想商。} 設 $\mT$ 是 $\mR$ 的一個擴張環，設 $\mA,\mB,\mC,\ldots$ 是 $\mT$ 中的 $\mR$-模。\emph{商} $\mC=\mA:\mB$ 定義為滿足下列條件的一個模：$\mC\mB\equiv0(\mA)$，並且由 $\mD\mB\equiv0(\mA)$ 可推出 $\mD\equiv0(\mC)$；也就是說，$\mC$ 是使 $\mC\mB\equiv0\pmod{\mA}$ 成立的``最小''模。由此，商被唯一地確定為所有滿足 $\mD\mB\equiv0\pmod{\mA}$ 的模 $\mD$ 的最大公約數；這樣的 $\mD$ 存在，因為零模當然就是這樣一個模。

由定義可得：若 $\mB$ 是 $\overline{\mB}$ 的倍數，則 $\mA:\overline{\mB}$ 是 $\mA:\mB$ 的倍數。若 $\mA$ 是 $\overline{\mA}$ 的倍數，則 $\mA:\mB$ 是 $\overline{\mA}:\mB$ 的倍數。並且有
\[
        \mA:\mB\mC=(\mA:\mB):\mC=(\mA:\mC):\mB .
\]
若擴張環 $\mT$ 同 $\mR$ 重合，則\emph{模商}過渡為\emph{理想商} $\ma:\mb$，所以上面的計算規則也適用於它。由於 $\ma\mb\equiv0(\ma)$，這裡還可推出 $\ma\equiv0(\ma:\mb)$。

5. \emph{若 $\ma:\mb=\ma$，則稱 $\mb$ 對 $\ma$ 為素。} 因此，一個理想 $\mb$ 當且僅當由 $\md\mb\equiv0\pmod{\ma}$ 總能推出 $\md\equiv0\pmod{\ma}$ 時，才對 $\ma$ 為素。由 4. 中的計算規則可得：\emph{若 $\mb$ 和 $\mc$ 都對 $\ma$ 為素，則它們的乘積和最小公倍數也對 $\ma$ 為素。}若 $\mb$ 對 $\ma$ 為素，同時 $\ma$ 也對 $\mb$ 為素，則稱 $\ma$ 和 $\mb$ 相互為素。由 \S\ 4, 4$\beta$ 可得：若 $\ma$ 和 $\mb$ 互素，則它們也相互為素。

\setcounter{footnote}{25}

\subsection*{\S\ 6. 假定除子鏈條件下的理想論}

作為基礎，設有一個（交換）環 $\mR$，其中\emph{公理 I，即除子鏈條件成立}；並且在以下總是假定 $\mR$ 的\emph{所有元素帶有一個固定的良序}。這樣也同時給出了 $\mR$ 中\emph{所有理想的一個良序}。事實上，這兩個假設立即推出 $\mR$ 中每個理想都有一個由有限多個元素組成的理想基。於是從 $\mR$ 中元素的良序出發，轉到 $\mR$ 的所有有限子集的一個準字典序排列，\footnote{這應理解如下：有限子集 $\mathfrak A_n$ 的元素 $a_1,\ldots,a_n$ 被如此標記，使指標的次序同 $\mR$ 中給定的良序一致。每個子集 $\mathfrak A_n$ 都先於每個含 $m>n$ 個元素的子集 $\mathfrak A_m$。若 $\mathfrak A_n$ 與 $\mathfrak B_n$ 是具有同樣多元素的子集，則稱 $\mathfrak A_n$ 早於 $\mathfrak B_n$，如果第一個與相應的 $b_i$ 不同的元素 $a_i$ 在 $\mR$ 的良序中先於 $b_i$。這樣，每個有限子集系統都有一個第一元素 $\mathfrak A_n$。在給定 $\mR$ 的良序時，因此不需要進一步的選擇公設；特別地，若 $\mR$ 可數，如代數數域的情形，就完全不涉及選擇公設。關於選擇公設在理想論中的作用，\emph{Idealtheorie} 注 9 曾有所指出，但未作詳論。} 並給每個理想指定一個卓異基，即在各種可能基中按有限子集良序排在第一位的基；於是，隨着有限子集被良序化，理想也被良序化。

\textbf{定理 I。} \emph{在假定除子鏈條件時，$\mR$ 中每個理想都允許表示為有限多個不可約理想的最小公倍數；這裡不可約理想是指不能表示為兩個真除子的最小公倍數的理想。}

為證明這一點，只須說明：若定理 I 對某個理想 $\mm$ 不成立，則 $\mm$ 有一個真除子，定理 I 對這個真除子同樣不成立；由此便可構造一條不在有限步後中止的除子鏈，同所假定的除子鏈條件相矛盾。事實上，$\mm$ 必須是可約的，因為否則 $\mm=[\mm]$ 就給出了所需的表示。若 $\mm=[\ma,\mb]$，則定理 I 不可能同時對 $\ma$ 和 $\mb$ 成立；否則 $\mm$ 的相應表示也會隨之得到。因此，存在 $\mm$ 的真除子，使定理 I 對它們不成立；令 $\ma_1$ 為這些理想中按理想良序排在第一位者。再按同樣方式從 $\ma_1$ 構造一個真除子 $\ma_2$，一般地從 $\ma_i$ 構造一個真除子 $\ma_{i+1}$，就得到一條定義良好且不在有限步後中止的除子鏈，違背假設。\footnote{顯然，若對全體模的系統假定除子鏈條件，則定理 I 對模範圍也同樣成立。定理 I 具有純粹集合論性質；同時，它也是證明中唯一使用理想良序的定理。撇開一切運算，它所涉及的是如下集合論概念：在集合 $\mathfrak M$ 中，指定其冪集的一個子集 $\Sigma$，也就是一個子集系統；假定 $\Sigma$ 已被良序化，並把 $\Sigma$ 的元素稱為 $\Sigma$-集合。在 $\Sigma$ 中假定鏈條件：每條 $\Sigma$-集合鏈 $\mathfrak A_1,\mathfrak A_2,\ldots,\mathfrak A_\nu\ldots$，只要 $\mathfrak A_\nu$ 是 $\mathfrak A_{\nu-1}$ 的真上集，就在有限步後終止。若一個 $\Sigma$-集合 $\mathfrak A$ 是兩個都嚴格包含 $A$ 的 $\Sigma$-集合的交，則稱 $\mathfrak A$ 可約；否則稱它不可約。上述論證於是表明：每個 $\Sigma$-集合都可表示為有限多個不可約 $\Sigma$-集合的交。}

由於預設了除子鏈條件，按 \S\ 5, 2 可以徑稱準素理想。準素性和不可約性之間的聯繫由下述定理給出：

\textbf{定理 II。} \emph{在假定除子鏈條件時，每個不可約理想都是準素理想；換言之，每個非準素理想都是可約的。}

過渡到剩餘類環 $\mR\mid\mm$ 時，由於同態性，除子鏈條件保持成立；$\mm$ 作為最小公倍數的表示對應於零理想的表示；並且由於第一同構定理（\S\ 4, 3），$\mm$ 的準素除子又對應於 $\mR\mid\mm$ 中的準素理想，反過來也一樣。因此可以只限於剩餘類環中零理想的分解。由於這個環仍是同樣一般性的環，可以從一開始就假定待分解的理想是 $\mR$ 的零理想。

於是，設 $\mR$ 的零理想不是準素理想；則至少存在一對理想 $\ma,\mb$，其中還可以假定 $\mb$ 為主理想，使得 $a\ne(0)$，並且對每個 $\varkappa$ 都有 $\mb^\varkappa\ne(0)$，而 $\ma\mb=(0)$。形成按假設在有限步後中止的理想商除子鏈 $\ma,\ \ma:\mb,\ldots,\ma:\mb^\nu,\ldots$；設例如 $\mt=\ma:\mb^{m-1}=\ma:\mb^m\ldots$ 等於這條鏈中所有後繼理想；也就是說，$\mt=\mt:\mb$，或者 $\mb$ 對 $\mt$ 為素。進一步，$\mt$ 作為 $\ma$ 的除子不同於零理想，並且對每個指數 $\varkappa$，$\mb^\varkappa$ 也不同於零。為了證明定理 II，因此只須證明一個表示
\[
        (0)=[\mt,\mb^{m+1}].
\]
按定義有
\[
        \mb^{m-1}\mt\equiv0\pmod{\ma},\qquad\text{因而}\qquad \mb^m\mt=(0),
\]
所以只剩下證明：
\[
        \mc=[\mt,\mb^{m+1}]\equiv0\pmod{\mb^m\mt}.
\]
這由 $\mb$ 是主理想並且對 $\mt$ 為素這兩個假設推出。$\mc$ 中每個元素 $c$，由於可被 $\mb^{m+1}$ 整除，都允許一個表示；這裡 $n$ 表示一個整數符號：
\[
        c=k b^{m+1}+n b^{m+1}=r b^m,
\]
其中 $r$ 又表示 $\mR$ 中一個元素。由 $c\equiv0\pmod{\mt}$ 得到 $r b^m\equiv0\pmod{\mt}$，並由 $\mt:\mb=\mt$ 得到 $r\equiv0\pmod{\mt}$。這樣便有 $c\equiv0\pmod{\mt\mb^m}$，定理 II 得證。

\textbf{定理 III。}\emph{在假定除子鏈條件時，每個理想都可作一個最短表示，即表示為有限多個屬於不同素理想的、因而為最大準素分量的最小公倍數\footnote{其中相伴素理想和孤立分量唯一確定。參見 \emph{Idealtheorie}，或 W. Krull，\emph{Ein neuer Beweis für die Hauptsätze der allgemeinen Idealtheorie}，Math. Ann. 90 (1923), S. 55--64。第 \S\ 7 節將對那裡出現的簡單特殊情形直接證明唯一性；在那裡被認明為唯一的準素理想正是孤立分量。}。}

為得到這種表示，只須把依定理 I 存在的不可約理想表示——從而依定理 II 是準素理想表示——換成一個最短表示，這是始終可以做到的。把屬於同一素理想的準素理想合併，便依 \S\ 5, 3 得到所求表示。準素分量可稱為最大的，因為其中任意幾個的最小公倍數已不再是準素理想。

\subsection*{\S\ 7. 假定雙鏈條件下的理想論}

同目前已經發展出的理論相比，這裡得到的簡化建立在 1. 和 2. 中將給出的輔助定理上。

1. \emph{若在一個無零因子的交換環中倍數鏈條件成立，則該環同時就是一個域。} 需要證明的是，對 $a\ne0$，方程 $ax=b$ 在環中總有一個解。它不可能有多於一個解，則由環被假定無零因子推出。

令 $\ma$ 為由 $a$ 導出的主理想。按假設，冪的序列在有限步後中止；設 $\ma^m$ 等於所有後繼冪。若 $\mb$ 表示由 $b$ 導出的主理想，則有
\[
        \ma^m\mb=\ma^{m+1}\mb,
        \qquad \ma^{m+1}\mb=(a^{m+1}b).
\]
這特別給出元素 $a^m b$ 的一個表示，其中 $n$ 表示一個整數符號：
\[
        a^m b=k a^{m+1}b+n a^{m+1}b=a^{m+1}c,
\]
其中 $c$ 又是 $\mR$ 中的元素。由於按假設沒有零因子，由此得到 $b=ac$，從而證明了域性質。

2. 由 1. 立即得到輔助定理：\emph{若在一個交換環中倍數鏈條件成立，則一個素理想沒有任何不同於 $\mo$ 的真除子。}

同樣得到附加結論：\emph{若只滿足公理 II，即模每個非零理想的倍數鏈條件，則每個非零素理想沒有任何不同於 $\mo$ 的真除子。}事實上，模素理想 $\mpideal$ 的剩餘類環按定義是無零因子環；又因同態性，倍數鏈條件在其中也成立，所以依 1. 它是一個域。由於 $\mpideal$ 的除子同該剩餘類環中的理想一一對應，$\mpideal$ 沒有任何不同於 $\mo$ 的真除子\footnote{這裡並不要求 $\mR$，從而也不要求 $\mo$——由 $\mR$ 全體元素組成的單位理想——具有單位元，儘管依 1. 剩餘類環中存在單位類。最簡單的例子是全體偶數組成的系統；它屬於附加結論的較弱情形。}。

\textbf{定理 IV。} \emph{若在一個交換環中雙鏈條件成立，則在一個理想的每個最短表示中，那些不屬於單位理想 $\mo$ 的準素分量都是唯一確定的。}

\emph{附加結論。} 在假定公理 I 和 II 時，定理 IV 對每個不同於零理想的理想成立。

設
\[
 \mm=[\mq,\mq_1,\ldots,\mq_r]
    =[\overline{\mq},\overline{\mq}_1,\ldots,
      \overline{\mq}_{\overline r}]
\]
是 $\mm$ 由最大準素分量給出的兩個最短表示，並設 $\mo,\mpideal_1,\ldots,\mpideal_r$ 以及 $\mo,\overline{\mpideal}_1,\ldots,\overline{\mpideal}_{\overline r}$ 分別是相伴素理想。若沒有出現屬於 $\mo$ 的準素理想，就省去 $\mq$ 或 $\overline{\mq}$。由於
\[
 \mo^\varrho\mpideal_1^{\varrho_1}\ldots\mpideal_r^{\varrho_r}
 \equiv0(\mm),
\]
每個 $\overline{\mpideal}_i$ 必進入某個 $\mpideal_i$，所以依輔助定理與它相同。反過來，每個 $\mpideal_k$ 也必與某個 $\overline{\mpideal}_k$ 相同；因此，不同於 $\mo$ 的素理想彼此一致，可編號為 $\mpideal_\lambda=\overline{\mpideal}_\lambda$。又有 $\mq\mq_1\ldots\mq_r\equiv0(\overline{\mq}_\lambda)$，從而 $\mq_\lambda\equiv0(\overline{\mq}_\lambda)$，因為其餘分量不被 $\mpideal_\lambda$ 整除，因而對 $\overline{\mq}_\lambda$ 為素。反向同理有 $\overline{\mq}_\lambda\equiv0(\mq_\lambda)$，所以不屬於 $\mo$ 的準素分量唯一確定。最後，若 $\mq$ 確實出現，則由表示的最短性，$\overline{\mq}$ 也必出現；這樣也證明了相伴素理想的唯一性。這個唯一性證明還表明，任何不含屬於 $\mo$ 的準素分量的表示本身就是最短的。

3. 若在雙鏈條件的假設之外再加上單位元的存在，則定理 IV 加強為定理 V。

\textbf{定理 V。} \emph{若在一個帶單位元的交換環中雙鏈條件成立，則每個理想都可唯一地表示為有限多個兩兩互素的準素理想的乘積。}

\emph{附加結論。} 在假定公理 I、II、III 時，定理 V 對每個不同於零理想的理想成立。

由於單位元存在，有 $\mo^2=\mo$，於是每個屬於 $\mo$ 的準素理想都等於 $\mo$，因而不能出現在一個不同於 $\mo$ 的理想的最短表示中。因此，按定理 IV，準素分量唯一確定。同時，任何表示在刪去 $\mo$ 後都會成為一個最短表示。按 2. 中的輔助定理，兩個不同的素理想總是互素的；於是其準素分量也同樣互素，最小公倍數就成為乘積。

由這些假設還得到下列輔助定理。

\textbf{輔助定理。}\emph{在一個帶單位元並滿足雙鏈條件的交換環中，除單位理想外，對理想 $\mm$ 為素的素理想，恰好就是那些不進入 $\mm$ 的素理想。為素與互素這兩個概念重合。}

\emph{附加結論。}在假定公理 I、II、III 時，須再假定 $\mm$ 不是零理想。

若 $\mpideal$ 不進入 $\mm$，則它不同於所有屬於 $\mm$ 的素理想；依第 2 點的輔助定理，它不被其中任何一個不同於 $\mo$ 的素理想整除。所以 $\mpideal$ 對每個準素分量為素，從而對 $\mm$ 為素；同時它同每個準素分量互素，從而同 $\mm$ 互素。

若 $\mpideal\ne\mo$ 進入 $\mm$，則它必與一個相伴素理想相同。例如設它屬於分量 $\mq=\mq_1$；於是 $\mq:\mpideal$ 是 $\mq$ 的真除子，因為 $\mpideal^{\varrho-1}\equiv0(\mq:\mpideal)$，但 $\mpideal^{\varrho-1}\not\equiv0(\mq)$。同時，$\mq:\mpideal$ 作為 $\mpideal^{\varrho-1}$ 的除子，只能被一個不同於 $\ne\mo$ 的素理想整除，因而仍是準素理想。由分解的唯一性，可被 $\mm:\mpideal$ 整除的乘積 $(\mq:\mpideal)\cdot\mq_2\ldots\mq_r$ 也成為 $\mm$ 的真除子；所以 $\mpideal$ 不對 $\mm$ 為素。

因此，一個理想 $\mt\ne\mo$ 當且僅當屬於 $\mt$ 的素理想沒有一個進入 $\mm$ 時，才對 $\mm$ 為素；在這種情形下，$\mt$ 也同 $\mm$ 互素。反過來，兩個互素理想總是相互為素（\S\ 5, 5）；所以在上述假設下，這兩個概念重合。

\subsection*{\S\ 8. 在商域中整閉條件下的理想論}

到目前為止所發展的理想論，現在要通過加入最後兩個公理而加強為通常的形式。

1. \emph{輔助定理。}設 $\mR$ 是帶單位元且無零因子的交換環，並且在 $\mR$ 中假定除子鏈條件成立（公理 I、III、IV）。那麼，由 $\ma=\ma\mb$ 和 $\ma\ne(0)$ 總能推出 $\mb=\mo$。所以，對每個既不同於零理想又不同於單位理想的理想，都有
\[
        \mc^\varrho\ne\mc^{\varrho+1}.
\]
\footnote{相反，在同樣的假設下，不能從 $\ma\mb=\ma\mc$ 推出 $\mb=\mc$。例如，令 $\mR$ 為一個域上關於 $x,y$ 的多項式範圍，並置 $\ma=(x,y)$、$\mb=(x^2,xy,y^2)$、$\mc=(x^2,y^2)$。事實上 $\ma\mb=\ma\mc=\ma^3$，但 $\mb\ne\mc$。}
證明與 \S\ 1, 3 的輔助定理相同，因為可把 $\ma\mb$ 看成關於 $\mb$ 的有限模。令 $\alpha_1,\ldots,\alpha_n$ 是由公理 I 保證存在的 $\ma$ 的一個理想基；由 $\ma=\ma\mb$ 得到方程組
\[
 \alpha_i=b_{i1}\alpha_1+\ldots+b_{in}\alpha_n,
 \qquad b_{ik}\equiv0(\mb),\quad i=1,\ldots,n.
\]
因為假定 $\mR$ 沒有零因子，由此得到 $|b_{ik}-e_{ik}|=0$，其中 $i\ne k$ 時 $e_{ik}=0$，而 $e_{ii}=e$。由方程
\[
        e^n+b_1e^{n-1}+\ldots+b_n=0,
\]
其中 $b_i\equiv0(\mb)$，便得到 $e\equiv0(\mb)$，從而 $\mb=\mo$。

2. 現在只需證明：在公理 I 至 IV 之外再加上商域中整閉性（公理 V）的假定，就可推出每一個不同於零理想的準素理想都是它的相伴素理想的一個冪。由此得到的表示
\[
        \mm=\mpideal_1^{\varrho_1}\ldots\mpideal_r^{\varrho_r}
\]
對所有不同於零理想的理想的唯一性，已經由定理 V 和 1. 中的輔助定理證明；同時也已經證明，這些素理想沒有任何不同於單位理想的真除子。證明先借助 Dedekind 推論 II（\S\ 1, 4）對指數二給出；一般情形則用完全歸納法。

\emph{輔助定理。}\emph{在假定公理 I 至 V 時，指數為二的準素理想除 $\mq=\mpideal^2$ 外別無其他。}

為證明這一點，須說明：由
\[
 \mpideal^2\equiv0(\mq),\qquad
 \mq\equiv0(\mpideal),\qquad
 \text{但 }\mq\not\equiv0(\mpideal^2)
\]
必然推出 $\mq=\mpideal$。於是取
\[
 c\equiv0(\mq),\qquad
 \text{因而 }c\equiv0(\mpideal),\qquad
 \text{但 }c\not\equiv0(\mpideal^2).
\]
由 \S\ 7, 3 的輔助定理，$\mo c:\mpideal$ 是 $\mo c$ 的真除子；所以存在一個元素 $b$，使得
\[
 b\mpideal\equiv0(\mo c)\equiv0(\mq),
 \qquad \text{但 }b\not\equiv0(\mo c).
\]
因此 $\gamma=b/c$ 不是整的，即它屬於商域而不屬於 $\mR$。

現在須證明 $b\not\equiv0(\mpideal)$；由於 $\mq$ 是屬於 $\mpideal$ 的準素理想，由此便得到 $\mpideal\equiv0(\mq)$。

依 Dedekind 推論 II，存在 $\mR$ 中的元素 $m,n$，使得
\[
        \gamma=b/c=m/n
\]
並且 $m^2/n$ 也不是整的；也就是說，
\[
 bn=mc,\qquad
 m\not\equiv0(\mo n),\qquad
 m^2\not\equiv0(\mo n).
\]
把 $b\mpideal$ 乘以 $m$，得到
\[
 mb\mpideal\equiv0(\mo mc),
 \qquad\text{因而}\qquad
 \equiv0(\mo nb),
 \qquad\text{從而}\qquad
 m\mpideal\equiv0(\mo n),
\]
最後一步成立是因為所論的環沒有零因子。由此得到
\[
        m\not\equiv0(\mpideal),
\]
因為 $m^2\not\equiv0(\mo n)$ 而 $n\equiv0(\mpideal)$；後一關係來自 \S\ 7, 3 的輔助定理：由於 $m\not\equiv0(\mo n)$，$\mo n:\mpideal$ 是 $\mo n$ 的真除子。於是四個關係
\[
 m\not\equiv0(\mpideal),\quad
 n\equiv0(\mpideal),\quad
 c\equiv0(\mpideal),\quad
 c\not\equiv0(\mpideal^2),\quad bn=mc
\]
給出 $b\not\equiv0(\mpideal)$。事實上，因為 $\mpideal^2$ 是準素理想，先有 $mc\not\equiv0(\mpideal^2)$；再由 $bn=mc$ 得到 $b\not\equiv0(\mpideal)$。輔助定理得證。

3. \emph{輔助定理。}\emph{在假定公理 I 至 V 時，指數為 $\varrho$ 的準素理想除 $\mq=\mpideal^\varrho$ 外別無其他。}

該輔助定理可由第 2 點的輔助定理推出，無須再使用這些公理\footnote{從輔助定理 2 推出輔助定理 3 的證明，見 Masazo Sono，\emph{On Congruences II}，\S\S\ 9 和 10；參見注 19。}。首先須證明
\[
 \text{若 }c\equiv0(\mpideal),\quad c\not\equiv0(\mpideal^2),
 \quad\text{則}\quad
 \mpideal^\sigma=(\mo c^\sigma,\mpideal^{\sigma+1})
 \quad\text{對每個 }\sigma\text{ 都成立。}
\]
由第 2 點的輔助定理，有 $\mpideal=(\mo c,\mpideal^2)$。所以，若假定
\[
        \mpideal^{\sigma-1}=(\mo c^{\sigma-1},\mpideal^\sigma),
\]
則乘以 $\mpideal$ 並應用分配律，得到
\[
\begin{aligned}
 \mpideal^\sigma
  &=(\mpideal c^{\sigma-1},\mpideal^{\sigma+1})\\
  &=(\mo c^\sigma,\mpideal^2c^{\sigma-1},\mpideal^{\sigma+1})\\
  &=(\mo c^\sigma,\mpideal^{\sigma+1}).
\end{aligned}
\]
待證的輔助定理可化為如下第二種表述：
\[
 \text{若 }\mq\equiv0(\mpideal^\sigma),\quad
 \mq\not\equiv0(\mpideal^{\sigma+1}),\quad
 \mpideal^{\sigma+\lambda}\equiv0(\mq),\quad \lambda\geqq1,
\]
則也有
\[
        \mpideal^{\sigma+\lambda-1}\equiv0(\mq).
\]
事實上，對每個準素理想，由 $\mq\equiv0(\mpideal)$ 而 $\not\equiv0(\mpideal^{\varrho+1})$，總存在這樣的指數 $\sigma\geqq1$；而由 $\mpideal^\varrho\equiv0(\mq)$ 以及第 1 點輔助定理所給出的 $\mpideal^\varrho\ne\mpideal^{\varrho+1}$，可知 $\mq\not\equiv0(\mpideal^{\varrho+1})$。有限次應用第二種表述，便得到 $\mpideal^\sigma\equiv0(\mq)$，從而 $\mq=\mpideal^\sigma$。

在第二種表述的假設下，利用 $\mpideal^\sigma$ 的上述表示，對 $\mq$ 中每個元素 $q$ 都有
\[
 q\equiv bc^\sigma(\mpideal^{\sigma+1})
 \qquad\text{其中 } b\not\equiv0(\mpideal),
 \qquad\text{或者}\qquad
 bc^\sigma\equiv0(\mq,\mpideal^{\sigma+1}).
\]
由於 $(\mq,\mpideal^{\sigma+1})$ 是準素理想，由此得到
\[
 c^\sigma\equiv0(\mq,\mpideal^{\sigma+1})
 \quad\text{或者}\quad
 c^{\sigma+\lambda-1}\equiv0(\mq,\mpideal^{\sigma+\lambda}),
\]
從而 $\mpideal^{\sigma+\lambda-1}\equiv0(\mq)$。

4. 總結起來得到：

\textbf{定理 VI。} \emph{若在一個交換環中公理 I 至 V 成立，則每個不同於零理想和單位理想的理想，都可唯一地表示為有限多個不同於零理想和單位理想的素理想的冪的乘積；這些素理想沒有任何不同於單位理想的真除子。}

\subsection*{\S\ 9. 由所假定的分解推出公理}

為了也能反過來從通常理想分解的存在推出這些公理，而這種通常理想分解按定理 VI 又由公理 I 至 V 推出，必須以下列形式來假定分解。

\textbf{假設。}\emph{在交換環 $\mR$ 中，每個不是由零元素或單位元素生成的理想，都可唯一地表示為素理想冪的乘積。這些素理想是並非由單位元素生成的單純理想；也就是說，它們沒有不同於 $\mo$ 的真除子。反過來，所有單純理想都是素理想。這裡的唯一性以如下嚴格形式成立：}
\[
        \ma^\varrho\ne\ma^{\varrho+1}
\]
\emph{對每個不是由零元素或單位元素生成的理想都成立。}

這裡並不預先假定單位元素存在。若沒有單位元素，則``不是由單位元素生成''這一條件不構成任何限制。

\textbf{1. 證明公理 III（乘法單位元的存在）。}假定公理 III 不成立；須證明 $\mo^2$ 是單純理想卻不是素理想，從而與假設矛盾。由嚴格唯一性假設，$\mo\ne\mo^2$；所以 $\mo\not\equiv0(\mo^2)$，而 $\mo^2\equiv0(\mo^2)$，因而 $\mo^2$ 不是素理想。但 $\mo^2$ 是單純的；因為 $\mo^2$ 不可能被任何不同於 $\mo$ 的素理想整除，所以依所假定的乘積表示，它沒有任何不同於 $\mo$ 或 $\mo^2$ 的除子。

\textbf{2. 證明公理 IV（環沒有零因子）。}若 $a$ 是零因子，即 $ab=0$，但 $a\ne0$ 且 $b\ne0$，則把由 $a$ 和 $b$ 導出的主理想的表示相乘，得到
\[
        (0)=\mpideal_1^{\varrho_1}\ldots\mpideal_k^{\varrho_k},
\]
其中各 $\mpideal$ 都不同於零理想和單位理想；後一點是因為單位元已經存在。把這個表示取成最短的，意即刪去任何一個因子都會給出零理想的一個真除子；由於關於零理想並未假定唯一表示，這個條件並非自動成立。若該表示中只出現一個素理想，則
\[
        (0)=\mpideal^\varrho=\mpideal^{\varrho+1},
\]
與嚴格唯一性要求矛盾。若出現一個以上的素理想，則
\[
\begin{aligned}
 \mpideal_1^{\varrho_1+1}
  &=(\mpideal_1^{\varrho_1+1},
     \mpideal_1^{\varrho_1}\ldots\mpideal_k^{\varrho_k})\\
  &=\mpideal_1^{\varrho_1}
    (\mpideal_1,\mpideal_2^{\varrho_2}\ldots\mpideal_k^{\varrho_k})\\
  &=\mpideal_1^{\varrho_1};
\end{aligned}
\]
因為單位元存在，$\mpideal_1$ 同其餘各素理想互素。這裡同樣得到與嚴格唯一性要求的矛盾。

\textbf{3. 證明公理 I 和 II（鏈條件）。}這些假設對每個非零理想給出：\emph{整除關係蘊含乘積表示}；也就是說，由 $\mm\equiv0(\ma)$ 和 $\ma\ne\mo$ 得到 $\mm=\ma\mb$，其中 $\mb\not\equiv0(\mm)$。設
\[
 \mm\equiv0(\ma),\qquad
 \mm=\mpideal_1^{\varrho_1}\ldots\mpideal_r^{\varrho_r},
 \qquad
 \ma=\overline{\mpideal}_1^{\sigma_1}\ldots
     \overline{\mpideal}_{\overline r}^{\sigma_{\overline r}}.
\]
於是每個 $\overline{\mpideal}$ 都同某個 $\mpideal$ 相同，並且 $\sigma_i\leqq\varrho_i$；所以斷言成立，可取 $\mb=\mpideal_1^{\varrho_1-\sigma_1}\ldots\mpideal_r^{\varrho_r-\sigma_r}$，其中某些 $\sigma_i$ 當然可以為零。因而，$\mm$ 只有有限多個除子，恰好對應於有限多個組合 $0\leqq\sigma_i\leqq\varrho_i$；所以模 $\mm$ 的剩餘類環滿足雙鏈條件。這樣便證明了公理 I 和 II：除子鏈條件在 $\mR$ 本身中也成立，因為零理想的每個真除子都不是零理想。

\emph{公理 V，即商域中的整閉性的證明，要用到理想分解的通常推論；所以須先推導這些推論。}

4. \emph{模每個非零理想的剩餘類環都是主理想環。}可假定該理想不同於單位理想，因為對只含零元的剩餘類環，斷言顯然成立。

先設該理想是準素理想 $\mq=\mpideal^\varrho$，並且 $c\equiv0(\mpideal)$、$\not\equiv0(\mpideal^2)$。由第 3 點得到 $\mo c=\mpideal\ma$，其中 $(\ma,\mpideal)=\mo$。所以
\[
        \mpideal^\sigma=(\mo c^\sigma,\mpideal^\varrho)
        \qquad\text{對每個 }\sigma\leqq\varrho;
\]
因而，在模一個準素理想的剩餘類環中，每個理想都是主理想。一般地，設
\[
 \mm=\mq_1\mq_2\ldots\mq_r,
 \qquad \mq_i=\mpideal_i^{\varrho_i},
\]
並相應地把剩餘類環表示為直和
\[
        \mR\mid\mm=\overline{\ma}_1+\ldots+\overline{\ma}_r.
\]
則 $\overline{\ma}_i$ 同構於 $\mR\mid\mq_i$（\S\ 4, 5），所以 $\overline{\ma}_i$ 中每個理想都是主理想。若 $\overline{\mc}$ 是剩餘類環中的任意理想，則 $\overline{\ma}_i\overline{\mc}$ 是主理想 $\overline{\ma}_i\overline c_i$；於是
\[
\begin{aligned}
 \overline{\mc}
  &=\overline{\mo}\,\overline{\mc}
   =\overline{\ma}_1\overline{\mc}+\ldots+
    \overline{\ma}_r\overline{\mc}\\
  &=\overline{\ma}_1\overline c_1+\ldots+
    \overline{\ma}_r\overline c_r\\
  &=\overline{\ma}_1\overline c+\ldots+
    \overline{\ma}_r\overline c
   =\overline{\mo}\,\overline c,
\end{aligned}
\]
其中置 $\overline c=\overline c_1+\ldots+\overline c_r$。因為 $\overline{\ma}_i\overline{\ma}_k=(0)$，並且 $\overline c_i\equiv0(\overline{\ma}_i)$，所以 $\overline{\ma}_i$ 中的主理想 $\overline{\ma}_i\overline c_i$ 與 $\overline{\ma}_i\overline c$ 相同。

主理想環定理還可表述如下：

\emph{若 $\mc$ 是任意理想，則可把它乘以一個同給定理想 $\mb$ 互素的理想而化為主理想。}

事實上，置 $\mm=\mb\mc$；則 $\mc$ 模 $\mm$ 後成為主理想，即
\[
        \mc=(\mo c,\mm)=(\mc\ma,\mc\mb)=\mc(\ma,\mb);
\]
所以 $\mo c=\mc\ma$ 且 $(\ma,\mb)=\mo$。

\textbf{5. 分式理想理論。}商域 $\mK$ 中的每個有限 $\mR$-模稱為分式理想。

\textbf{定理。}\emph{$\mK$ 中不同於零模的有限 $\mR$-模關於乘法構成 Abel 群。}

兩個有限 $\mR$-模的積仍是有限 $\mR$-模；乘法滿足結合律和交換律；由於 $\mA=\mo\mA$，單位理想是這個系統的單位元。因此，只須證明方程 $\mathfrak A\mX=\mo$ 在有限 $\mR$-模的系統中總有解。

\textbf{預備說明。}若方程 $\mathfrak A\mT=\mathfrak B$ 有且只有一個解，則 $\mT$ 等於模商 $\mathfrak B:\mathfrak A$（\S\ 5, 4）。因為
\[
 \mT\equiv0(\mathfrak B:\mathfrak A);
 \quad\text{所以}\quad
 \mathfrak B=\mathfrak A\mT
 \equiv0\bigl(\mathfrak A\cdot(\mathfrak B:\mathfrak A)\bigr)
 \equiv0(\mathfrak B),
 \quad\text{也就是}\quad
 \mathfrak A\cdot(\mathfrak B:\mathfrak A)=\mathfrak B.
\]
現在先設 $\mathfrak A$ 是\emph{主模}，$\mathfrak A=(\alpha)$；則 $\mX=(e/\alpha)$ 是 $\mathfrak A\mX=\mo$ 的解，並且 $\mX$ 仍是有限 $\mR$-模。由此得到：\emph{每個有限 $\mR$-模都是 $\mR$ 中兩個理想的模商}，這也說明了“分式理想”這一名稱。事實上，設 $\tau_1,\ldots,\tau_r$ 是 $\mT$ 的一個模基，置 $\tau_i=t_i/a$，並令 $\mt$ 是 $\mR$ 中由 $t_i$ 導出的理想；則 $\mo a\cdot\mT=\mt$，所以依預備說明 $\mT=(\mt:\mo a)$，這裡的商取在 $\mR$ 中。

現在可以一般地給出 $\mathfrak A\mX=\mo$ 的解。置 $\mathfrak A=a:\mo c$，並取 $\mb$ 使 $\ma\mb$ 等於主理想 $\mo a$。於是
\[
 \mathfrak A c=\ma,
 \qquad \mathfrak A\mb c=\mo a,
 \qquad \mathfrak A(\mb c:\mo a)=\mo.
\]
這裡 $\mX$ 是一個理想除以一個主理想所得的商，所以仍是有限 $\mR$-模。由這樣\emph{已經證明的群性質}還可得到：\emph{任意兩個理想的模商 $\ma:\mb$}，作為方程 $\mb\mX=\ma$ 的解，\emph{是有限 $\mR$-模}。

6. 由群性質得到以下進一步推論。

6$\alpha$. \emph{可以約去}；也就是說，$\mathfrak T=\mathfrak A\mathfrak M:\mathfrak B\mathfrak M$ 當且僅當 $\mathfrak T=\mathfrak A:\mathfrak B$。事實上，從 $\mathfrak T\mathfrak B\mathfrak M=\mathfrak T\mathfrak A\mathfrak M$ 得到 $\mathfrak T\mathfrak B=\mathfrak T\mathfrak A$，反過來也一樣。

6$\beta$. \emph{每個分式理想（有限 $\mR$-模）都可表示為 $\mathfrak C=\ma:\mb$，其中 $\ma$ 和 $\mb$ 互素；這個條件唯一地確定 $\ma$ 和 $\mb$。}表示的存在性由第 5 點和 6$\alpha$ 得到；而從 $\mathfrak C=\ma:\mb=\overline{\ma}:\overline{\mb}$ 得到 $\ma\overline{\mb}=\mb\overline{\ma}$，若 $\ma,\mb$ 以及 $\overline{\ma},\overline{\mb}$ 各自互素，便得到唯一性。

6$\gamma$. \emph{每個由非整元素（即屬於 $\mK$ 而不屬於 $\mR$ 的元素）導出的主模 $\mo\eta$，都可表示為主理想之商，並且分子的任何冪都不被分母整除。}取約化表示 $\mo\eta=\mb:\mc$，其中 $(\mb,\mc)=\mo$；再依第 4 點取 $\mo b=\mb\ma$ 且 $(\ma,\mc)=\mo$。於是
\[
        \mo\eta=\ma\mb:\ma\mc=\mo b:\ma\mc.
\]
因為 $\ma\mc=\mo b:\mo\eta$ 也成為主理想，所以
\[
        \mo\eta=\mo b:\mo c.
\]
由於 $(\ma\,\mc)=\mo$ 且 $(\mb,\mc)=\mo$，$\ma\mb$ 的任何冪都不被 $\mc$ 整除。

7. \emph{證明公理 V，即商域中的整閉性。}由 6$\gamma$，\emph{$\mK$ 中每個非整元素都有一個商表示}
\[
        \eta=a/c,
\]
\emph{使得在 $\mR$ 中分子的任何冪都不被分母整除。}因為由 $\mo\eta=\mo b:\mo c$ 可知，$\eta$ 與商 $b/c$ 只相差 $\mR$ 中的一個單位。這樣的元素 $\eta$ 不可能滿足任何把它刻畫為關於 $\mR$ 為整的方程；因為由
\[
        \eta^n+r_1\eta^{n-1}+\ldots+r_n=0
\]
會得到 $\mR$ 中的 $a^n\equiv0(\mo c)$。

8. \emph{推論：若交換環 $\mR$ 滿足公理 I 至 IV，並且每個準素理想都不可約，則公理 V，即商域中的整閉性，也成立。}由公理 I 至 III，每個素理想冪 $\mpideal^\varrho$ 同時是準素理想；特別是 $\mpideal^2$ 也是準素理想，因而依假設不可約。由此須證明
\[
        \mpideal=(\mo c,\mpideal^2);
\]
因為這樣便由 \S\ 8, 3 的輔助定理得到每個準素理想都是素理想冪，而這同其餘假設一起依第 7 點推出整閉性。

由除子鏈條件，有
\[
        \mpideal=(\mo c_1,\ldots,\mo c_k,\mpideal^2),
\]
其中作為 $c$ 的乘子的，只須考慮模 $\mpideal$ 的剩餘類，而且可假定 $c$ 關於模 $\mpideal$ 的剩餘類域線性無關。若 $k>1$，這種線性無關性便給出表示
\[
 \mpideal^2=[\mq_1,\mq_2],
 \qquad \mq_1=(\mo c_1,\mpideal^2),
 \qquad \mq_2=(\mo c_2,\ldots,\mo c_k,\mpideal^2),
\]
使 $\mq_1$ 和 $\mq_2$ 都成為 $\mpideal^2$ 的真除子；這會使 $\mpideal^2$ 可約，與假設矛盾。因此 $\mpideal=(\mo c,\mpideal^2)$ 以及由此推出的結論得證。

反過來，由公理 I 至 V 立即可知每個準素理想都不可約，因為素理想冪 $\mpideal^\varrho$ 不可能表示為形如 $\mpideal^\sigma$ 和 $\mpideal^\tau$ 的真除子的最小公倍數。

9. \emph{若允許環中有零因子，則由公理 I 至 III 以及商環中的整閉性，並不能推出每個準素理想都不可約}\footnote{參見注 19。}。令 $\mT$ 是一個滿足除整閉性外公理 I 至 IV 的環；這種環確實存在。例如，若 $\mR$ 滿足公理 I 至 V，則在 $\mR$ 的商域的有限擴張域中，除全體關於 $\mR$ 為整的元素所成系統外，各個有限階都是這樣的環（\S\ 3, 2）。由推論 8，$\mT$ 中存在可約的準素理想 $\mq$。令 $\mpideal^\sigma$ 是 $\mq$ 的一個真倍數，並令 $\overline{\mR}$ 為剩餘類環 $\mT\mid\mpideal^\sigma$；於是與 $\mq$ 對應的理想 $\overline{\mq}$ 是一個不同於零理想的可約準素理想。在 $\overline{\mR}$ 中，公理 I 至 III 成立，而且商環中的整閉性也成立。事實上，因為 $\mT$ 中每個不被 $\mpideal$ 整除的元素都同 $\mpideal^\sigma$ 互素，所以 $\overline{\mR}$ 中所有正則元素都是單位；因而 $\overline{\mR}$ 與它的商環相同。

\begin{center}
{\bfseries\large \S\ 10.\par}
\vspace{0.35em}
{\bfseries\large 雙鏈條件與合成列\par}
\end{center}

下文將證明，對於任意假定為良序的模範圍（\S\ 2），\emph{雙鏈條件成立這一假設}——即每條模的除子鏈和每條模的倍數鏈都在有限步後終止——與\emph{存在合成列這一假設}等價\footnote{該範圍中的模關於加法構成 Abel 群；確切地說，因為乘子範圍由 $\mR$ 中的元素組成，所以它們是\emph{廣義} Abel 群。由於所論的是 Abel 群，合成列的概念與主列的概念重合。若把“模”理解為任意一個非交換的、也可以是廣義的群的全體正規因子，本節的定理仍然成立。這裡採用的術語與前文略有不同，正是為了使人聯想到群論。}。把模範圍特化為交換環，便得到該環全體理想系統的相應結論；此時在第 2 點中須處處把模同構換成環同構。

若在模範圍 $\mG$ 中，除 $\mG$ 本身和零模 $\mE$ 外再無其他 $\mR$-模，則稱 $\mG$ \emph{單純}。若 $\mG|\mA$ 單純，則稱模 $\mA$ \emph{在 $\mG$ 中單純}（極大正規子模）。

\emph{倍數鏈}
\[
        \mG,\mA_1,\ldots,\mA_r,\mE
\]
\emph{稱為 $\mG$ 的長度為 $r$ 的合成列}，如果鏈中所有模互不相同，而且每一個模都\emph{在前一個模中單純}，也就是說，剩餘類模（商群）$\mA_{i-1}|\mA_i$ 以及 $\mG|\mA_1$ 和 $\mA_r$ 都是單純模範圍。通過形成剩餘類模 $\mG|\mA$，$\mG$ 中存在一條止於 $\mA\ne\mE$ 的合成列這一情形，可以歸結為通常的合成列情形。

1. \emph{若假定 $\mG$ 中雙鏈條件成立，則 $\mG$ 中存在合成列}——這裡假定 $\mG\ne\mE$。由 $\mG$ 的良序性和除子鏈條件，可以像在 \S\ 6 中那樣，藉助準字典序得到全體模所成系統的一個良序。憑藉這個良序，除子鏈條件蘊含至少存在一個在 $\mG$ 中單純的模。事實上，令 $\mE_1$ 為該良序中 $\mE$ 的第一個真除子，一般地令 $\mE_i$ 為 $\mE_{i-1}$ 的第一個真除子，則完全確定的鏈
\[
        \mG,\mG_1,\ldots,\mG_i,\ldots
\]
在有限步後終止，因而必然導向一個在 $\mG$ 中單純的模 $\mA_1$。若 $\mA_1\ne\mE$，則用同一方法可得一個在 $\mA_1$ 中單純的模 $\mA_2$；一般地，若 $\mA_{i-1}\ne\mE$，則存在一個在 $\mA_{i-1}$ 中單純的模 $\mA_i$。這樣便得到一條完全確定的倍數鏈
\[
        \mG,\mA_1,\ldots,\mA_i,\ldots,
\]
它依倍數鏈條件在有限步後終止，因而給出 $\mG$ 的一條合成列。

2. \emph{若假定 $\mG$ 中存在合成列，則 $\mG$ 中雙鏈條件成立。}證明採用完全歸納法，而且無須假定 $\mG$ 具有良序。在歸納步驟中，同時只以存在合成列為假設而證明 Jordan--Hölder 定理\footnote{關於環的情形，參見 Masazo Sono（注 19），\emph{On Congruences I}，\S\S\ 11--13。該文還討論了非交換群中合成列的對應物，即列 $\mR,\mA_1,\ldots,\mA_r,\mE$，其中每個 $\mA_i$ 都是 $\mA_{i-1}$ 中的理想且在其中單純，但不必是 $\mR$ 中的理想。事實上，只要把倍數鏈理解為 $\mG,\mA_1,\ldots,\mA_i,\ldots$ 這種列，其中每個 $\mA_i$ 都是 $\mA_{i-1}$ 的正規因子，並相應地定義除子鏈，上述全部論證對於非交換群仍然成立。另參見 Dedekind，\emph{Über die von drei Moduln erzeugte Dualgruppe}（Math. Ann. 53 (1900), S. 371--403）。該文對於更一般的範圍（模群），也只在存在合成列這一假設下證明了 Jordan--Hölder 定理（\S\ 6, XVI）。在那裡，第二同構定理由一個稍弱的對應關係取代，因而定理的表述也稍弱；但其證明在原理上與這裡給出的證明相同。}。

先設 $\mG$ 單純，即 $r=0$，於是只存在唯一的合成列 $\mG,\mE$。此時以下命題成立：
\[
\begin{array}{ll}
\alpha) & \text{通過每個在 $\mG$ 中單純的模都可以作出一條合成列。}\\
\beta) & \text{Jordan--Hölder 定理：每條合成列的長度都相同，並且商群}\\
        & \text{（剩餘類模）的系統除排列次序外都相同；亦即相應的商群彼此同構。}\\
\gamma) & \text{通過任意一個不同於 $\mG$ 的模都可以作出一條合成列。}\\
\delta) & \text{倍數鏈條件成立。}\\
\varepsilon) & \text{除子鏈條件成立。}
\end{array}
\]
因此，對於每個具有長度 $\bar r<r+1$ 的合成列的模範圍，都可以假定命題 $\alpha)$ 至 $\varepsilon)$ 已經成立。現在要在如下\emph{假設下證明這些命題成立：$\mG$ 中存在一條合成列}
\[
        \mG,\mA,\mA_1,\ldots,\mA_r,\mE
\]
\emph{其長度為 $r+1$。}

\textbf{2$\alpha$.} 若不存在不同於 $\mA$ 而又在 $\mG$ 中單純的模，則無須證明。於是設 $\mB$ 是一個在 $\mG$ 中單純且不同於 $\mA$ 的模；須證明 $\mG$ 中存在一條通過 $\mB$ 的合成列。因為 $\mA$ 和 $\mB$ 都在 $\mG$ 中單純而又彼此不同，所以必有 $(\mA,\mB)=\mG$；再置 $\mM=[\mA,\mB]$，由第二同構定理（\S\ 4, 2）得到
\[
        \mG|\mB\simeq\mA|\mM
        \qquad\text{以及}\qquad
        \mG|\mA\simeq\mB|\mM.
\]
因此 $\mM$ 在 $\mA$ 和 $\mB$ 中都單純。由於 $\mA$ 有一條長度為 $r<r+1$ 的合成列，依命題 $\alpha)$ 和 $\beta)$，$\mA$ 中便有一條長度為 $r$ 且通過 $\mM$ 的合成列，例如
\[
        \mA,\mM,\mM_1,\ldots,\mM_{r-1},\mE;
\]
於是
\[
        \mG,\mB,\mM,\mM_1,\ldots,\mM_{r-1},\mE
\]
就是 $\mG$ 中一條通過 $\mB$ 的合成列。

\textbf{2$\beta$.} 為證明 Jordan--Hölder 定理，設
\[
 (1)\quad \mG,\mA,\mA_1,\ldots,\mA_r,\mE
 \qquad\text{以及}\qquad
 (2)\quad \mG,\mB,\mB_1,\ldots,\mB_s,\mE
 \quad\text{其中 } s\geqq r
\]
是 $\mG$ 的兩條合成列。若這裡 $\mA=\mB$，則由對 $r<r+1$ 已成立的歸納假設，結論立即得到。在另一情形下，同 2$\alpha$ 中構造的列
\[
 (3)\quad \mG,\mA,\mM,\mM_1,\ldots,\mM_{r-1},\mE
 \qquad\text{以及}\qquad
 (4)\quad \mG,\mB,\mM,\mM_1,\ldots,\mM_{r-1},\mE
\]
比較，便得到如下結果：依對 $r<r+1$ 所作的歸納假設，(1) 與 (3) 的商群系統相同；(2) 與 (4) 的商群系統也相同，因為 (4) 中通過 $\mB$ 的合成列長度為 $r$；所以 $s=r$。再由 2$\alpha$ 中給出的同構，(3) 與 (4) 的商群系統也相同，從而 (1) 與 (2) 的商群系統相同，Jordan--Hölder 定理得證。

\textbf{2$\gamma$.} 預備說明。\emph{設 $\mA,\mB,\mH$ 是 $\mG$ 中的模，並且 $\mB$ 在 $\mA$ 中單純；那麼模 $(\mB,\mH)$ 與 $(\mA,\mH)$ 或者相同，或者 $(\mB,\mH)$ 在 $(\mA,\mH)$ 中單純。}由第二同構定理（\S\ 4, 2）得到
\[
        (\mA,\mB,\mH)|(\mB,\mH)\simeq
        \mA|[\mA,(\mB,\mH)]
\]
或者——由於 $\mB\equiv0(\mA)$——
\[
        (\mA,\mH)|(\mB,\mH)\simeq \mA|\mC
        \quad\text{其中}\quad
        \mC=[\mA,(\mB,\mH)].
\]
這裡 $\mC\equiv0(\mA)$；另一方面，由 $\mB\equiv0(\mA)$ 以及 $\mB\equiv0((\mB,\mH))$，又有 $\mB\equiv0(\mC)$；所以 $\mC$ 位於 $\mA$ 與 $\mB$ 之間，依假設必與 $\mA$ 或 $\mB$ 相同，預備說明由此得證。

現在再令 $\mG,\mA,\mA_1,\ldots,\mA_r,\mE$ 是 $\mG$ 的一條合成列，並令 $\mH$ 是 $\mG$ 中任意一個不同於 $\mG$ 的模。為了證明有一條合成列通過 $\mH$，形成模的列
\[
        \mG,(\mA,\mH),(\mA_1,\mH),\ldots,(\mA_r,\mH),\mH.
\]
依預備說明，其中彼此不同的模 $\mG,\mC,\mC_1,\ldots,\mC_s,\mH$ 構成一條從 $\mG$ 到 $\mH$ 的合成列；所以 $\mC$——在 2$\alpha$ 的特殊情形中它與 $\mH$ 重合——是一個在 $\mG$ 中單純的模。因此，由 2$\alpha$，$\mC$ 中存在一條長度為 $r<r+1$ 的合成列；從而依歸納假設，可以在 $\mC$ 中作出一條通過 $\mH$ 的合成列。再添上 $\mG$，便得到 $\mG$ 中一條通過 $\mH$ 的合成列。反覆應用這一論證還表明，尤其可以把這樣的合成列取為從 $\mG$ 到 $\mH$ 所構造之列的延長。

\textbf{2$\delta$.} 證明倍數鏈條件。仍假定 $\mG$ 中存在一條長度為 $r+1$ 的合成列，並令 $\mG,\mB,\mB_1,\ldots,\mB_\nu,\ldots$ 是一條倍數鏈；這裡假定每一個模都是前一個模的真倍數。假定該鏈從 $\mG$ 開始並不限制一般性，因為總可以把 $\mG$ 添加為首項。由 2$\gamma$，$\mG$ 中存在一條通過 $\mB$ 的合成列，因而 $\mB$ 有一條長度更短的合成列。於是依歸納假設，鏈 $\mB,\mB_1,\ldots,\mB_\nu,\ldots$ 在有限步後終止；添加 $\mG$ 延長而成的鏈也同樣終止。

\textbf{2$\varepsilon$.} 證明除子鏈條件。仍假定 $\mG$ 中存在一條長度為 $r+1$ 的合成列，並令 $\mE,\mC,\mC_1,\ldots,\mC_\nu,\ldots$ 是一條除子鏈；這裡假定每一個模都是前一個模的真除子。假定該鏈從 $\mE$ 開始同樣不限制一般性。由 2$\gamma$，$\mG$ 中存在一條通過 $\mC$ 的合成列，所以 $\mG|\mC$ 中存在一條長度更短的合成列。於是依歸納假設，$\mG|\mC$ 中的除子鏈 $\mC|\mC,\mC_1|\mC,\ldots,\mC_\nu|\mC,\ldots$ 在有限步後終止。由一一對應——依第一同構定理（\S\ 4, 2）——原來以 $\mC$ 開始的鏈也在有限步後終止，因此添加 $\mE$ 延長而成的鏈也同樣終止。

\emph{至此已經證明：存在合成列與雙鏈條件成立這兩個假設彼此等價。}

\begin{center}
（1925 年 8 月 13 日收稿。）
\end{center}
%
\clearpage
%
\setcounter{footnote}{0}

%
\pdfbookmark[0]{31. 代數數域或函數域的階的判別式定理}{paper-31}
\section*{31. 代數數域或函數域的階的判別式定理}
\begin{center}
Journal f. d. reine u. angew. Math. 157 (1927), pp. 82--104

哥廷根的 \emph{埃米·諾特}\\[0.4em]
\rule{9em}{0.4pt}
\end{center}

眾所周知，Dedekind 的判別式定理說：一個素數進入域判別式，當且僅當它至少被某個素理想的平方整除。

我將在下文證明，這一定理對於有限代數數域的所有有限階仍然成立；也就是說，對於所有含有全體有理整數之環的代數整數環，定理可表述如下：一個素數 $p$ 進入某個階的判別式，當且僅當在該階中主理想 $(p)$ 分解為最大準素分量時，至少出現一個真正的準素理想，即一個並非素理想的準素理想。\footnote{參見 E. Noether, \emph{Abstrakter Aufbau der Idealtheorie in algebraischen Zahl- und Funktionenkörpern}, Math. Ann. 96 (1926), pp. 26--61。所有定義均參照該文；下文稱之為“理想論”。以下“階”始終意指“有限階”，“環”意指交換環。}

因為對於整數環（極大階），即域中全體整數的系統，除素理想的冪以外不再有別的準素理想，所以 Dedekind 的定理由此作為特例得出。

證明通過過渡到該階模 $(p)$ 的剩餘類環，化歸為如下環的分解理論：這些環相對於其內含的一個域具有有限秩；換言之，它們構成一個帶有主單位、以同一個域為係數域的交換超複量系統。在這裡，係數域特別是有理整數模 $p$ 的剩餘類域，而 $(p)$ 的分解在剩餘類環中對應於零理想的分解。

因此，若把這種有限秩環稱為完全可約的，只要它的零理想可以表示為有限多個素理想的最小公倍數，那麼問題就在於尋找完全可約性的判據。首先得到的判據是：當係數域為完美域時，環完全可約，當且僅當把係數域擴張為代數閉域後所得的擴張環完全可約。而對於這個擴張環，判別式不為零極易看出正是完全可約性的充要判據；\footnote{Hilbert 在一個特殊情形中也作類似論證：Zur Theorie der aus $n$ Haupteinheiten gebildeten komplexen Größen, Gött. Nachr. 1896。} 又因為階模每個素數所得的判別式化為剩餘類環的判別式，判別式定理於是得到證明。\footnote{相關研究也包括非交換情形，但限於整數環（極大階）並假定有理整數係數，可見 A. Speiser 近作 \emph{Allgemeine Zahlentheorie}, Naturf. Gesellschaft Zürich 71 (1926), pp. 8--48。亦參看該文所引美國文獻。}

若考慮函數域而非數域，則階現在必須包含全體有理整函數的基整環；對於多個不定元的代數函數或整代數函數，則必須包含具有有理整數係數的全體整泛函的基整環。在這種情形下，剩餘類環中也可能出現非完美系數域，例如在泛函整環情形下按一個素數取模時。對於非完美域作為係數域的情形，須按模素理想後的剩餘域是第一類擴張域還是第二類擴張域，來區分第一類和第二類素理想，\footnote{採用 Steinitz 的通用說法，見 J. f. M. 137。參看本文 \S 2 第 4 節的注。} 並相應地區分第一類和第二類完全可約性。上面所述判據始終指第一類完全可約性；在完美情形中也只會出現這一種。因此，在最一般情形下，判別式定理說：基整環中的一個素元進入某個階的判別式，當且僅當在 $(p)$ 的理想分解中，至少出現一個真正準素分量，或出現一個第二類素理想。

對於極大階，這一事實最早由 Kronecker 在例子中注意到；此前的紀念文仍錯誤地給出了僅對數域極大階有效的表述，並且沒有完整證明。對於極大階，Ostrowski\footnote{A. Ostrowski, Zur arithmetischen Theorie der algebraischen Größen, Nachr. Gött. Ges. d. Wissensch. 1919。} 給出了上述判別式定理以及與之相連的進一步分歧定理的證明；同時該文也給出詳盡的文獻綜述。最後，對於所有相對域的階，判別式定理通過對基整環的每個素理想過渡到局部化環（原文稱 Quotientenring）而直接推出。這種過渡的思想在特殊情形中已知，W. Krull\footnote{參見其 Danzig 講演，Jhrbr. d. D. Math. Ver. 34, p. 121。} 從原則上加以強調；在更一般研究框架下，H. Grell\footnote{\emph{Beziehungen zwischen den Idealen verschiedener Ringe}（將發表於 Math. Annalen）。} 給出了局部化環的系統理論。

這裡對所有階以及對完美和非完美系數域所作的統一處理，除使用理想論外，還基於一開始就以有限秩環的理論、因而以有限擴張的一般概念為基礎，而不是以由一個本原元生成的擴張概念為基礎。因此，一個階的模基以及其元素之間的關係，取代了域的本原元冪及其定義方程，或者取代了基本型的冪及基本方程。\footnote{這一觀點同 Dedekind 關於高等複數的看法密切相關：Zur Theorie der aus $n$ Haupteinheiten gebildeten komplexen Größen, Gött. Nachr. 1885。}

這樣的模基在對基整環的每個素元取模後，仍化為剩餘類環的一組基；而即便對極大階而言，若基由本原元的冪組成，這一點也未必成立；一旦涉及多個不定元的整代數函數，由基本型的冪組成的基也未必如此。\footnote{參看 Ostrowski，前引文，“第一類不規則性”。} 換句話說：以基整環為係數的一元多項式環，按任意定義方程，甚至按基本方程取模後所得的剩餘類環，模 $p$ 未必同極大階模 $(p)$ 的剩餘類環同構。通常表述中的主要困難就在於這種同構的消失。由於這裡並不引入這個輔助多項式環，這些困難也就自動避免。同時，剩餘類環的加法分解理論可以理解為方程乘法分解的對應物；在這裡同樣，分解為互素因子對應於多項式剩餘類環中的加法分解。

須指出，判別式定理只能代表階的一般分歧理論中的第一個定理，正如在極大階情形中，判別式定理還伴隨着更精細的分歧理想理論，後者在那裡允許更精確地決定指數。這個指數的決定，依賴於方程準素因子的冪表示，可解釋為從 $\mathfrak p$ 到 $\mathfrak q=\mathfrak p^\varrho$ 的一個合成列長度的決定，並且大概會以這種形式出現在一般分歧理論中。

\subsection*{\S 1. 域的擴張環}

\textbf{1.} 我們先列出關於任意環的若干說明。若 $T$ 是一個環，$S$ 是 $T$ 中的一個元素系統，則所謂“由 $S$ 在 $T$ 中導出的環”，是指 $T$ 中所有含有 $S$ 的環的交。相應地，當 $P$ 是 $T$ 的一個子環時，也可定義“由 $S$ 在 $T$ 中導出的理想”，或“由 $S$ 在 $T$ 中導出的 $P$-模”。\footnote{參見“理想論”，\S 1。}

若 $R$ 是 $T$ 的一個子環，而 $S$ 是 $T$ 的一個元素系統，則由 $R$ 與 $S$ 在 $T$ 中導出的環，記作由把 $S$ 環附加到 $R$ 而得到的環 $R[S]$。這個環由 $S$ 中元素的所有多項式構成，係數取自 $R$，其中相等關係和運算關係由 $T$ 中既定的相等關係和運算關係規定。由於這些相等關係和運算關係已經固定，在特殊情形中只需形成所有多項式的一個子系統就夠了，例如只形成以 $R$ 為係數、關於 $S$ 的所有線性形式。這樣的特殊情形例如在第 2 節出現：當 $R$ 是帶單位元的環，而 $S$ 是同一單位元下的一個域時，所有線性組合 $\sum c_i\gamma_i$，其中 $c_i$ 屬於 $S$、$\gamma_i$ 屬於 $R$，構成環 $R[S]$。

但是，當只給定環 $R$，另外給定一個不含於 $R$ 的元素系統 $S$，並給定其相等關係，或者還給定相互運算關係時，也可以構造由把 $S$ 附加到 $R$ 而得的擴張環 $R[S]$。可以證明，總存在一個含有 $R$ 和 $S$ 的環，即形式上由 $S$ 中元素、以 $R$ 為係數的所有多項式構成，並按通常運算法則規定運算的環。對於相等的定義，除運算法則外仍有完全的自由，唯一限制是必須包含在 $R$ 與 $S$ 中已經有效的相等定義。因此只須把那些由於 $R$ 與 $S$ 中的相等，以及由於運算，可以從 $R$ 與 $S$ 中相等元素以相同方式形成的多項式判為相等。這樣構造出的、含有 $R$ 與 $S$ 的環 $T$，便同 $T$ 中由 $R$ 與 $S$ 導出的環一致。\footnote{參看 Steinitz 對商域的構造。}

\textbf{2. 域的擴張環。} 從現在起，設所論環 $\mR$ 具有單位元，且含有一個域 $P$ 作為子環，因而是域的上環。於是 $\mR$ 的單位元 $e$ 也是 $\mathsf{P}$ 的單位元，而該環是一個 $\mathsf{P}$-模。

若 $\overline{\mathsf{P}}$ 是 $\mathsf{P}$ 的一個上域，則由把 $\overline{\mathsf{P}}$ 環附加到 $\mR$ 而得到的擴張環 $\bR=\mR[\overline{\mathsf{P}}]$ 說明如下。

假定 $\mR$ 與 $\overline{\mathsf{P}}$ 的集合論交 $[\mathfrak R,\overline{\mathsf{P}}]$ 為 $P$，並且不屬於 $P$ 的 $\mR$ 元素與 $\overline{\mathsf{P}}$ 元素之間沒有運算關係。這個假定總可通過把 $\overline{\mathsf{P}}$ 替換為 $\mathsf{P}$ 的一個等價擴張域，或把 $\mR$ 替換為一個等價擴張環來實現。\footnote{依 Steinitz，兩 個 $\mathsf{P}$ 的擴張域若能互相同構且使 $\mathsf{P}$ 逐元素對應於自身，稱為“等價”。$\mathsf{P}$ 的等價擴張環也相應定義。通過引入新符號，總能造成等價擴張；在下文中，這種使用等價擴張的自由是本質性的。}

把所有形式形成的線性組合聲明為 $\bR$ 的元素：
\[
   \bar c_{i_1}\gamma_{i_1}+\cdots+\bar c_{i_s}\gamma_{i_s},
\]
其中 $\bar c_i$ 遍歷 $\overline{\mathsf{P}}$ 的全體元素，$\gamma_i$ 遍歷 $\mR$ 的全體元素。\footnote{通常，$\mR$ 或 $\bR$ 的元素用希臘字母表示，$\mathsf{P}$ 或 $\overline{\mathsf{P}}$ 的元素用拉丁字母表示。} 這個定義應是唯一的；也就是說，若把 $\gamma$ 替換為在 $\mR$ 中與之相等的元素，把 $\bar c$ 替換為在 $\overline{\mathsf{P}}$ 中與之相等的元素，則 $\bR$ 中的元素也聲明為相等。

和定義為
\[
  \sum \bar c_i\gamma_i+\sum \bar d_i\gamma_i
  =\sum(\bar c_i+\bar d_i)\gamma_i .
\]
這裡在須固定的相等定義意義下是唯一的：把一個加項替換為與其相等的加項，所得結果仍相同。其中 $\bar c$ 與 $\bar d$ 中可以出現零，從而形式上相同的 $\gamma_i$ 都出現。

積同樣按此意義唯一地定義為
\[
  \sum \bar c_i\gamma_i\cdot \sum \bar d_k\gamma_k
  =\sum \bar c_i\bar d_k\gamma_i\gamma_k .
\]
對於相等定義，聲明如下：若 $\gamma_{i_1},\ldots,\gamma_{i_s}$ 關於 $P$ 線性無關，則兩個元素
\[
   \bar c_{i_1}\gamma_{i_1}+\cdots+\bar c_{i_s}\gamma_{i_s}
   \quad\text{与}\quad
   \bar d_{i_1}\gamma_{i_1}+\cdots+\bar d_{i_s}\gamma_{i_s}
\]
相等，當且僅當每個 $\bar c_{i_k}$ 在 $\overline{\mathsf{P}}$ 中等於 $\bar d_{i_k}$。換言之：把關於 $P$ 線性無關的 $\mR$ 元素，聲明為關於 $\overline{\mathsf{P}}$ 仍線性無關。

由和與積的定義，所有 $\bR$ 元素的相等定義也由此給出。事實上，它們均可由 $\mR$ 中線性無關元素表示。因為從 $\mR$ 中的關係
\[
  \mu=c_1\gamma_1+\cdots+c_s\gamma_s
\]
出發，按積的定義乘以 $\bar b=be$，得到
\[
  \bar b\mu=\bar b e\cdot(c_1\gamma_1+\cdots+c_s\gamma_s)
  =\bar b c_1\gamma_1+\cdots+\bar b c_s\gamma_s,
\]
於是由和的定義，得到所有元素均可由 $\mR$ 中線性無關元素表示。

因此，若所有元素已經表示為 $\mR$ 中線性無關元素的線性組合，則相等定義就是 $\overline{\mathsf{P}}$ 中係數相等。於是，對同時屬於 $\mR$ 與 $\bR$，或同時屬於 $\overline{\mathsf{P}}$ 與 $\bR$ 的元素，這一定義同原處有效的相等定義一致；而對 $\bR$ 中其餘元素，由於交的假定和進一步條件，$\mR$ 與 $\overline{\mathsf{P}}$ 中有效的相等除了定義中要求的唯一性外不再給出任何關係。\footnote{若沒有這個假定，則至少會有一個屬於 $\overline{\mathsf{P}}$ 但不屬於 $P$ 的元素 $\bar c$，使得 $\bar c=\bar c e=\gamma$。此時，按 $\mR$ 中的相等定義，$e$ 與 $\gamma$ 關於 $\mathsf{P}$ 確為線性無關，但關於 $\overline{\mathsf{P}}$ 卻相關。由於 $\overline{\mathsf{P}}$ 與 $\mR$ 元素之間存在關係，上述相等定義便不可能成立；例如，當把基域替換為中間域時，有限擴張域的次數會降低。相反，上述擴張的特徵正在於可能出現新的零因子。例如，有理係數域 $\mathsf{P}$ 上多項式環按 $x^2+1$ 取模所得的剩餘類環沒有零因子，並同域 $\mathsf{P}(\sqrt{-1})$ 同構；而過渡到以 $\mathsf{P}(\sqrt{-1})$ 為係數的多項式環時，因 $(x+i)(x-i)=(x^2+1)$ 而出現零因子，儘管沒有一個因子被 $x^2+1$ 整除。相應說法也適用於每一個有限數域；這正是 Dedekind 關於高等複數首先表達的觀念。參看 \S 6 第 4 節末尾。} 同樣，和與積唯一性的要求也不引入進一步關係，因為相等元素可以假定為逐係數相等，因而形式上給出相同的和與積表達式。

現在可立即驗證，環運算的所有公理均滿足；因此由環附加生成的擴張環 $\mR[\overline{\mathsf{P}}]$ 總是存在。

\textbf{3. $\mR$ 與 $\bR$ 中的理想。} 若 $\ma$ 是 $\mR$ 中的一個理想，則模積 $\bR\ma$ 是 $\bR$ 中的一個理想 $\ba$，即 $\ma$ 的擴張理想。同時 $\ba$ 是由 $\ma$ 在 $\bR$ 中導出的理想；它由所有線性組合 $\sum \bar c_i a_i$ 構成，每個和都只含有限多項，其中 $a_i$ 遍歷 $\ma$ 的元素，$\bar c_i$ 遍歷 $\overline{\mathsf{P}}$ 的元素。

若 $\overline{\mb}$ 是 $\bR$ 中的一個理想，則交 $[\overline{\mb},\mR]$ 是 $\mR$ 中的一個理想 $\mb$，即 $\overline{\mb}$ 的收縮理想。\footnote{按 H. Grell 在原文第 83 頁注 6 所引工作中的記號。} $\mR$ 的每個理想 $\ma$ 都是收縮理想，即其擴張理想 $\ba$ 的收縮理想。事實上，$[\mR,\ba]$ 的每個元素 $a$ 都可寫成 $\sum\bar c_i a_i$，其中 $a_i$ 可假定關於 $\mathsf{P}$ 或 $\overline{\mathsf{P}}$ 線性無關。於是元素 $a,a_i$ 在 $\bR$ 中線性相關，因此按相等定義在 $\mR$ 中也線性相關；故 $a$ 屬於 $\ma$。反包含顯然也成立，所以 $\ma=[\mR,\ba]$。

若 $\bp$ 是 $\bR$ 中的素理想，則 $\mpideal=[\mR,\bp]$ 是 $\mR$ 中的素理想。因為由第二同構定理，\footnote{“理想論”，\S 4，第 3 節。那裡定理只對同一環的理想陳述；同樣考察表明，它以上述形式對 $\bR$ 的任意理想 $\overline{\ma}$ 也有效。環同態 $\bR\to\bR/\overline{\ma}$ 對子環 $\mR$ 誘導出同態 $\mathfrak R\sim(\mathfrak R,\overline{\mathfrak a})\,|\,\overline{\mathfrak a}$。在此同態下，恰好 $[\mR,\overline{\ma}]$ 中的元素對應於零元，因此得到環同構 $(\mathfrak R,\overline{\mathfrak a})\,|\,\overline{\mathfrak a}\simeq \mathfrak R\,|\,[\mathfrak R,\overline{\mathfrak a}]$。左右兩邊的類可由 $\mR$ 中同一元素表示，因為證明每一步都把元素配給它所代表的類。} 剩餘類環 $\mR|\mpideal$ 同構於 $\bR/\bp$ 的一個子環，因而無零因子。也就是說，有環同構
\[
   (\mR,\bp)/\bp
   \simeq \mR/[\mR,\bp],
\]
其中 $(\mR,\bp)$ 表示模 $\mR$ 與 $\bp$ 的和，即最大公約數。

\subsection*{\S 2. 有限秩環；作為直和的表示}

\textbf{1. 有限秩環。} 若 $\mathfrak A$ 是有限 $P$-模，則 $\mathfrak A$ 關於 $P$ 具有有限秩，設為 $k$：也就是說，$\mathfrak A$ 中有且至多有 $k$ 個關於 $P$ 線性無關的元素，而 $\mathfrak A$ 中任意 $k$ 個關於 $P$ 線性無關的元素因而構成 $\mathfrak A$ 的一個模基。若一個有限秩 $\mathsf{P}$-模 $\mathfrak B$ 被同一有限秩的另一個 $\mathfrak A$ 整除，$\mathfrak B=O(\mathfrak A)$，則二者相同。

從現在起，設 $\mR$ 關於 $P$ 的有限秩為 $n$。於是 $\mR$ 中每個 $P$-模都有有限秩。因此，在 $\mR$ 中任意一條由 $P$-模組成的真除子鏈或真倍數鏈中，每個模的秩都必須分別大於或小於緊前一項的秩；鏈在有限步後中止。特別地，$\mR$ 中所有理想的系統滿足“雙鏈定理”。

由此可知，\footnote{“理想論”，\S 7，第 2 節。} $\mR$ 的一個素理想 $\mpideal$ 除單位理想外沒有真除子；因此，模一個非單位素理想所得的剩餘類環 $\mR/\mpideal$ 是一個域。若 $\mq$ 是屬於 $\mpideal$ 的一個準素理想，即 $\mq=O(\mpideal)$ 且 $\mpideal^\rho=O(\mq)$，則剩餘類環 $\mR/\mq$ 是準素環，也就是其中每個零因子的某個冪、在這裡無論如何第 $\rho$ 次冪、為零的環；並且 $\mR/\mq$ 只含零因子和單位。後一事實立即來自：每個不被 $\mpideal$ 整除的理想都與 $\mpideal$ 互素，因而也與 $\mq$ 互素。

由雙鏈定理以及單位元的存在還可推出，\footnote{“理想論”，\S 7，第 3 節，定理 V。} $\mR$ 的每個理想 $\ma$ 都可唯一表示為有限多個兩兩互素準素理想的乘積。由於互素性，這個乘積等於最小公倍數，並對應於把剩餘類環 $\mR/\ma$ 表示為準素環的直和。\footnote{“理想論”，\S 4，第 5 節。現在彙編中補充的內容散見於許多文獻。} 也就是說，$\mR/\ma$ 是這些環的最大公約數，而 $\mR/\ma$ 每個元素的和表示是唯一的。

\textbf{2. 有限秩環表示為準素環直和。} 設 $\mR$ 的零理想表示為
\[
   (0)=\mq_1\mq_2\cdots\mq_r
       =[\mq_1,\mq_2,\ldots,\mq_r],
\]
相應的直和表示為
\[
   \mR=\mR_1+\mR_2+\cdots+\mR_r,
   \qquad
   \mR_i\mR_k=(0)\ (i\ne k),
   \qquad
   \mR_i^2=\mR_i,
   \qquad
   \mR\mR_i=\mR_i .
\]
則 $\mR_i$ 是 $\mR$ 的一個理想，且
\[
   \mR_i
   =\mq_1\cdots\mq_{i-1}\mq_{i+1}\cdots\mq_r
   =[\mq_1,\ldots,\mq_{i-1},\mq_{i+1},\ldots,\mq_r],
\]
\[
   \mathfrak q_i=\mathfrak R_1+\cdots+\mathfrak R_{i-1}+\mathfrak R_{i+1}+\cdots\mathfrak R_r,
\]
並且 $\mR_i$ 同剩餘類環 $\mR/\mq_i$ 同構：$\mR_i$ 中每個元素 $\gamma_i$ 對應於 $\mR/\mq_i$ 中由 $\gamma_i$ 表示的類。\footnote{“理想論”，\S 4，第 5 節。} 若 $\ma$ 是 $\mR$ 的任意理想，則分配律給出
\[
   \ma=\mR\ma
   =\mR_1\ma+\cdots+\mR_r\ma
   =\ma_1+\cdots+\ma_r;
\]
這裡，由於 $\mathfrak a_i=\mathfrak R_i\mathfrak a_i=\mathfrak R\mathfrak a_i$，各分量 $\ma_i$ 既是 $\mR_i$ 中的理想，也是 $\mR$ 中的理想。作為理想，$\ma_i$ 和 $\mR_i$ 本身一樣都是有限 $P$-模；由於和是直和，它們的秩相加等於 $\ma$ 的秩，分別也等於 $\mR$ 的秩。

若單位元的和表示為
\[
   e=\varepsilon_1+\cdots+\varepsilon_r,
   \qquad
   \varepsilon_i\varepsilon_k=0\ (i\ne k),
   \qquad
   \varepsilon_i^2=\varepsilon_i,
   \qquad
   \varepsilon_i\equiv e\pmod{\mq_i},
\]
則 $\mR$ 中每個元素 $\gamma$ 的表示唯一地由
\[
   \gamma=\gamma e=\gamma\varepsilon_1+\cdots+\gamma\varepsilon_r
   =\gamma_1+\cdots+\gamma_r
\]
給出，其中分量 $\gamma_i=\gamma\varepsilon_i$ 屬於 $\mR_i$；同時 $\mR_i=\mR\varepsilon_i$。由“正交關係”得到
\[
   \gamma\delta=(\gamma_1\varepsilon_1+\cdots\gamma_r\varepsilon_r)
      (\delta_1\varepsilon_1+\cdots+\delta_r\varepsilon_r)
   =\gamma_1\delta_1\varepsilon_1+\cdots+\gamma_r\delta_r\varepsilon_r.
\]
因此 $\gamma_i\delta_i$ 是積的第 $i$ 個分量，類似地 $\gamma_i+\delta_i$ 是和的第 $i$ 個分量；這些也可由從 $\mR$ 到 $\mR_i$ 的同態推出。

環 $\mR_i$ 含有一個同構於 $\mathsf{P}$ 的子域 $\mathsf{P}_i=\mathsf{P}\varepsilon_i$，並且關於 $\mathsf{P}_i$ 有有限秩；這個秩等於 $\mR_i$ 作為 $P$-模所具有的秩。因為對 $P$ 的每個元素 $c$，有 $c\varepsilon_i\equiv ce\equiv c\pmod{\mq_i}$，所以只要 $c\ne0$ 就有 $c\varepsilon_i\ne0$；$P$ 與 $\mathsf{P}_i$ 之間的同態成為同構。其次，$\mR_i$ 元素關於 $P$ 的每個線性相關關係，乘以 $\varepsilon_i$ 後都化為關於 $\mathsf{P}_i$ 的相關關係；反過來，因為 $\mathsf{P}_i=\mathsf{P}\varepsilon_i$，關於 $\mathsf{P}_i$ 的每個關係同時也是關於 $\mathsf{P}$ 的關係。因此關於 $\mathsf{P}$ 和 $\mathsf{P}_i$ 的秩相同。

\textbf{3. 擴張環的直和表示。} 若
\[
   \mR=\mR_1+\cdots+\mR_r,
\]
則對 $\bR$ 有表示
\[
   \bR=\bR_1+\cdots+\bR_r,
\]
其中 $\bR_i$ 表示 $\mR_i$ 在 $\bR$ 中的擴張理想，同時也是由 $\mR_i$ 在 $\bR$ 中導出的環。事實上，由 $(0)=\mq_1\cdots\mq_r$ 乘以 $\bR$ 得到乘積表示 $(0)=\overline{\mq}_1\cdots\overline{\mq}_r$，其中各 $\overline{\mq}_i$ 仍兩兩互素。因此可得 $\bR$ 的直和表示；而該表示的第 $i$ 個分量由
\[
   \overline{\mq}_1\cdots \overline{\mq}_{i-1}\,
   \overline{\mq}_{i+1}\cdots \overline{\mq}_r
\]
給出，故等於 $\mR_i$ 的擴張理想。按 \S 1 第 3 節，這個擴張理想又同由 $\mR_i$ 關於 $\overline{\mathsf{P}}_i=\overline{\mathsf{P}}\varepsilon_i$ 形成的擴張環 $\mR_i[\overline{\mathsf{P}}_i]$ 相同；因此在擴張時只需擴張各個分量。一般而言，$\bR_i$ 於是未必仍是準素環。

\textbf{4. 第一類和第二類素理想。} 由第 1 節，模每個非單位素理想所得的剩餘類環 $\mR/\mpideal$ 是一個域。這個域含有一個同構於 $P$ 的子域 $(P)=(P,\mpideal)/\mpideal=P/[P,\mpideal]\simeq P$，因為 $\mpideal$ 不能含有 $\mathsf{P}$ 中任何非零元素；所以 $(\mathsf{P})$ 恰由那些可由 $\mathsf{P}$ 元素表示的類組成。\footnote{參看 \S 1 第 3 節關於第二同構定理的注。} 因而 $\mR/\mpideal$ 關於 $(\mathsf{P})$ 有有限秩，所以是一個有限代數擴張域。按這個擴張為第一類還是第二類，\footnote{採用 Steinitz 的通用術語：一個代數擴張稱為第一類，如果每個元素都是某個素函數的零點，而該素函數在適當的擴張域中可分解為互異線性因子；否則稱為第二類。} $\mpideal$ 稱為第一類或第二類素理想。

\subsection*{\S 3. 完全可約環}

\textbf{1.} 若在 $\mR$ 的零理想乘積表示中（\S 2，第 2 節），所有準素分量都成為素理想，或者等價地按 \S 2 第 1、2 節，$\mR$ 可表示為有限多個域的直和，則稱 $\mR$ 為\emph{完全可約}。若所有素理想都是第一類（\S 2，第 4 節），則稱該完全可約環為第一類完全可約；反之則為第二類完全可約。

若 $\mathsf{P}$ 是完美域，則只會出現第一類完全可約性；特別地，對代數閉域總是如此。下文中 $\overline{\mathsf{P}}$ 專指由 $\mathsf{P}$ 導出的代數閉包，而 $\bR$ 等於 $\mR[\overline{\mathsf{P}}]$。

\textbf{2. 定理。} \emph{環 $\mR$ 為第一類完全可約，當且僅當帶有代數閉域 $\overline{\mathsf{P}}$ 的擴張環 $\bR=\mR[\overline{\mathsf{P}}]$ 完全可約。}

證明分三步進行。

\textbf{2a.} 由 $\bR$ 完全可約，更一般地由任一擴張環完全可約，可推出 $\mR$ 完全可約（第一類或第二類）。事實上，按假設 $\bR$ 的零理想有表示
\[
        (0)=[\bp_1,\ldots,\bp_s]
\]
為素理想的最小公倍數。收縮到 $\mR$ 得到零理想的表示
\[
        (0)=[[\mR,\bp_1],\ldots,[\mR,\bp_s]],
\]
而按 \S 1 第 3 節，理想 $[\mR,\bp_s]$ 都是素理想。

\textbf{2b.} 由 $\mR$ 的第一類完全可約性，可推出 $\bR$ 的第一類完全可約性；更一般地，可推出介於 $\mR$ 與 $\bR$ 之間每個中間環的第一類完全可約性。

由於按 \S 2 第 3 節，$\mR$ 的直和表示給出 $\bR$ 的相應表示，使得每個分量 $\bR_i$ 成為擴張環 $\mR_i[\overline{\mathsf{P}}_i]$，其中 $\overline{\mathsf{P}}_i$ 是同構於 $\overline{\mathsf{P}}$ 的、$\mR_i$ 中與 $P$ 同構的子域 $\mathsf{P}_i$ 的一個擴張域，所以只須考察單個分量 $\mR_i$。因此可不失一般性地假定 $\mR$ 是一個域，且是其子域 $\mathsf{P}$ 上的第一類有限擴張域。於是 $\mR$ 由給 $\mathsf{P}$ 附加一個第一類本原元生成；等價地，$\mR$ 同一元多項式環 $\mathsf{P}[x]$ 按第一類素函數 $f(x)$ 取模的剩餘類環 $P[x]/f(x)$ 同構。若 $Q$ 是 $P$ 的一個擴張域，則 $Q[x]/f(x)$ 同構於 $\mR[Q]$；因為從 $\mathsf{P}[x]$ 過渡到 $Q[x]$ 時，剩餘類環的秩保持不變，等於 $f(x)$ 的次數，所以涉及的正是 \S 1 第 2 節給出的擴張環構造。在 $Q[x]$ 中，多項式 $f(x)$ 分解為兩兩互素的第一類素函數，每個都在 $Q[x]$ 中生成一個素理想；因此 $Q[x]/f(x)$，從而按同構 $\mR[Q]$，也是第一類完全可約。由於這對每個擴張域 $Q$ 都成立，特別對代數閉域 $\overline{\mathsf{P}}$ 成立，2b 得證。

\textbf{2c.} 若 $\mR$ 是第二類完全可約，則 $\bR$ 不完全可約。

同 2b 一樣，只需設 $\mR$ 是一個域，並且是 $\mathsf{P}$ 上的第二類有限擴張域。須證明，在 $\bR$ 作為準素環直和的表示中至少有一個真正的準素分量。為此，只需給出 $\bR$ 的一個非零元素，其某個冪為零。因為若 $\bar\beta\ne0$，則至少有一個分量 $\bar\beta_i$ 非零；若 $\bar\beta^\rho=0$，則 $\bar\beta_i^\rho=0$（\S 2，第 2 節，從 $\bR$ 到 $\bR_i$ 的同態），於是 $\bR_i$ 不是域。

令 $\gamma$ 為 $\mR$ 中一個第二類元素，其指數為 $f\geqq1$，並令
\[
        F(\gamma)=\gamma^{kp^f}+c_1\gamma^{(k-1)p^f}+\cdots+c_k=0
\]
為 $\gamma$ 在 $\mathsf{P}$ 上滿足的最低次數方程，其中 $p$ 表示 $\mathsf{P}$ 的特徵。於是 $\mR$ 中的元素
\[
        e,\gamma,\gamma^2,\ldots,\gamma^{kp^f-1}
\]
關於 $\mathsf{P}$ 線性無關，因而按相等定義在 $\bR$ 中關於 $\overline{\mathsf{P}}$ 仍線性無關。若現在置
\[
        G(\gamma)=\gamma^{kp^{f-1}}+\sqrt[p]{c_1}\,\gamma^{(k-1)p^{f-1}}+
        \cdots+\sqrt[p]{c_k},
\]
則 $G(\gamma)$ 是 $\bR$ 中的非零元素；因為 $p>1$ 且 $f\geqq1$，指數 $kp^{f-1}$ 總小於 $kp^f$，所以 $G(\gamma)$ 中出現的 $\gamma$ 的冪之間不可能有相關關係。但 $G(\gamma)^p=F(\gamma)$，故為零。2c 得證。

由 2a 與 2c 可知，$\bR$ 的第一類完全可約性推出 $\mR$ 的第一類完全可約性；2b 給出反向。定理得證。

\textbf{3. 推論。} \emph{若 $\mR$ 是第一類完全可約，則 $\bR$ 是秩一環的直和，因而是同構於 $\overline{\mathsf{P}}$ 的域的直和。這樣的秩一環分解已經存在於某個環 $\mR[Q]$ 中，其中 $Q$ 是 $\mathsf{P}$ 的一個有限擴張域。}

事實上，由 2b，環 $\bR$ 成為若干域的直和；這些域作為同構於 $\overline{\mathsf{P}}$ 的子域上的有限擴張，因 $\overline{\mathsf{P}}$ 代數閉而同這些子域相同。因此得到
\[
        \bR=\overline{\mathsf{P}}\bar\varepsilon_1+\cdots+\overline{\mathsf{P}}\bar\varepsilon_n,
        \qquad e=\bar\varepsilon_1+\cdots+\bar\varepsilon_n;
\]
$e$ 的各分量 $\bar e_i$ 同時構成 $\bR$ 的一個線性無關模基。

若按 \S 1 第 2 節，把 $\bar e_i$ 表示為 $\mR$ 中元素 $\gamma$ 的線性組合，係數在 $\overline{\mathsf{P}}$ 中，則這些係數的全體確定 $P$ 的一個有限擴張域 $Q$。於是 $\bar e_i$ 已屬於 $\mR[Q]$，由其模基性質得到
\[
        \mR[Q]=Q\bar e_1+\cdots+Q\bar e_n.
\]
因此，在 $\mR[Q]$ 中已經有秩一環分解；而在 $\mR[Q]$ 與 $\bR$ 之間的每個中間環中，不可分分量由 $Q\bar e_i$ 的擴張環產生。\footnote{這樣的最小擴張域 $Q$ 可刻畫為由 $\mR$ 的分量 $\mR_i=\mR\varepsilon_i$ 所確定的伽羅瓦域的合成域。具體說，按 2b 選擇 $f_i(x)$，使 $\mathsf{P}[x]/f_i(x)$ 同構於分量 $\mR_i$，並令 $Q^{(i)}$ 為位於 $\overline{\mathsf{P}}$ 中且恰足以把 $f_i(x)$ 分解為線性因子的伽羅瓦擴張域。則 $Q$ 是所有 $Q^{(i)}$ 的合成域。$f_i(x)$ 分解為線性因子，對應於把 $Q^{(i)}[x]/f_i(x)$ 表示為秩一域的直和；$Q^{(i)}$ 是使此發生的最小擴張域。按同構，對 $\mR_i[Q^{(i)}\bar e_i]$ 也同樣成立，因而對 $\mR[Q]$ 成立。}

\subsection*{\S 4. 有限秩環的矩陣表示、跡和判別式}

在本節和下一節中，對基礎環的假設比前面略為一般，以便同時包括數域和函數域中的階。這裡所要處理的是對本質上已知事實的一般而統一的表述。

\emph{假設。} $\mR$ 是整環 $\mZ$ 的一個擴張環；$\mR$ 的單位元已經屬於 $\mZ$。對於 $\mR$ 中的每一個 $\mZ$-模，特別是對於 $\mR$ 本身，至少存在一個由關於 $\mZ$ 線性無關元素組成的有限模基。除 $\alpha_1,\ldots,\alpha_s$ 外，元素 $\beta_1,\ldots,\beta_s$ 構成同一 $\mZ$-模的線性無關模基，當且僅當變換行列式是 $\mZ$ 的一個單位。

這個假設顯然對迄今考慮的有限秩環成立，此時 $\mZ$ 是域 $P$。它也對數域和函數域中的階成立，此時 $\mZ$ 分別為有理整數環或泛函整環。\footnote{參見“理想論”，\S 3。}

\textbf{1. 同態於 $\mR$ 的矩陣環。} 令 $\alpha_1,\ldots,\alpha_s$ 為理想 $\ma$ 的一個線性無關 $\mZ$-模基，令 $\gamma$ 為 $\mR$ 的任意元素。則 $\gamma\alpha_i$ 也屬於 $\ma$，並有係數在 $\mZ$ 中的方程\footnote{$\mZ$ 的元素用拉丁字母表示。}
\[
        \gamma\alpha_i=c_{1i}\alpha_1+\cdots+c_{si}\alpha_s,
        \qquad (i=1,\ldots,s).
\]
若用 $(\tau_1,\ldots,\tau_s)$ 表示由有關元素組成的一行矩陣，則矩陣形式為
\[
(\gamma\alpha_1,\ldots,\gamma\alpha_s)
=(\alpha_1,\ldots,\alpha_s)
\begin{pmatrix}
 c_{11}&\cdots&c_{1s}\\
 \vdots&&\vdots\\
 c_{s1}&\cdots&c_{ss}
\end{pmatrix}
=(\alpha_1,\ldots,\alpha_s)C.
\]
\emph{當 $\gamma$ 遍歷 $\mR$ 的全體元素時，由基 $\alpha_i$ 指定出的矩陣 $C$ 的全體形成一個同態於 $\mR$ 的環 $R_\alpha$；$R_\alpha$ 同剩餘類環 $\mR/((0):\ma)$ 同構，其中 $(0):\ma$ 表示零理想對 $\ma$ 的理想商。}因為對差 $\beta-\gamma$ 指定的是差 $B-C$，積的情形也同樣成立，因為
\[
 (\beta\gamma\alpha_1,\ldots,\beta\gamma\alpha_s)
  = (\beta\alpha_1,\ldots,\beta\alpha_s)C
  = (\alpha_1,\ldots,\alpha_s)BC.
\tag*{\footnotemark}
\]
\footnotetext{\label{fn:p31-noncomm}這種記法特意安排成在非交換環 $\mR$ 的情形下仍然成立；在這種情形中，假定 $\mZ$ 與所有元素可交換。}
特別地，$\mZ$ 的元素 $c$ 對應於對角矩陣 $cE$，其中 $E$ 是與單位元 $e$ 對應的矩陣 $\left(\begin{smallmatrix}e&&\\&\ddots&\\&&e\end{smallmatrix}\right)$。恰好那些使 $\gamma\ma$ 成為零理想的元素 $\gamma$ 被指定為零矩陣；由此得到上述同構。

\emph{同一理想的另一組基 $\bar\alpha_1,\ldots,\bar\alpha_s$ 生成一個同構於 $R_\alpha$、事實上與之等價的矩陣環：}
\[
        R_{\bar a}=P^{-1}R_\ma P.
\]
事實上，若 $\bar\alpha_1,\ldots,\bar\alpha_s$ 是另一組線性無關 $\mZ$-模基，則
\[
(\bar\alpha_1,\ldots,\bar\alpha_s)=(\alpha_1,\ldots,\alpha_s)P,
\]
其中 $P$ 的行列式是 $\mZ$ 的一個單位。因此
\[
(\gamma\bar\alpha_1,\ldots,\gamma\bar\alpha_s)
=(\alpha_1,\ldots,\alpha_s)CP
=(\bar\alpha_1,\ldots,\bar\alpha_s)P^{-1}CP,
\tag*{\textsuperscript{\ref{fn:p31-noncomm}}}
\]
所以相應矩陣 $C'$ 與 $C$ 滿足 $C'=P^{-1}CP$，並有 $R_{\bar\alpha}=P^{-1}R_\alpha P$。

\textbf{2. 理想類和表示類。} 兩個理想 $\ma$ 與 $\mb$ 屬於同一理想類，若它們作為 $\mR$-模同構；也就是說，若可在它們的元素之間建立一一對應，使差以及被 $\mR$ 同一元素所乘都相對應，因而由 $\alpha\simeq\beta$ 對任意 $\gamma\in\mR$ 總有 $\gamma\alpha\simeq\gamma\beta$。按第 1 節矩陣表示產生的兩個環 $R_\alpha$ 與 $R_{\bar\alpha}$ 屬於同一表示類，若它們等價；也就是說，若存在一個幺模矩陣 $P$，使
\[
        R_{\bar \alpha}=P^{-1}R_\alpha P.
\tag*{\footnotemark}
\]
\footnotetext{這裡不討論這樣一個在若干特殊情形中已經解決的問題：每一個同態矩陣環是否都能按第 1 節由 $\mR$ 的一個理想生成；即使允許使用線性相關的基，也不在這裡討論。按定義，表示類只由按第 1 節從 $\mR$ 的理想生成的矩陣環組成。}

\emph{理想類與表示類一一對應。}第 1 節已經證明，同一理想的不同模基生成屬於同一表示類的環。屬於同一理想類的不同理想 $\ma,\mb$ 也生成屬於同一表示類的環：若 $\beta_1,\ldots,\beta_s$ 是在同構下與 $\ma$ 的基 $\alpha_1,\ldots,\alpha_s$ 對應的 $\mb$ 的基，則 $\gamma\alpha_i$ 與 $\gamma\beta_i$ 對應，而 $c_{1i}\alpha_1+\cdots+c_{si}\alpha_s$ 與 $c_{1i}\beta_1+\cdots+c_{si}\beta_s$ 也對應；所以矩陣環 $R_\alpha$ 與 $R_\beta$ 相同。

反過來，若 $R_\alpha$ 與 $R_{\bar\beta}$ 屬於同一表示類，即 $R_{\bar\beta}=P^{-1}R_\alpha P$，其中 $\alpha_i$ 是 $\ma$ 的一組基而 $\bar\beta_i$ 是 $\mb$ 的一組基，則
\[
(\beta_1,\ldots,\beta_s)=(\bar\beta_1,\ldots,\bar\beta_s)P^{-1}
\]
也是 $\mb$ 的一組基，並且 $R_\beta=PR_{\bar\beta}P^{-1}=R_\alpha$。把 $\ma$ 的元素 $a=c_1\alpha_1+\cdots+c_s\alpha_s$ 對應到 $\mb$ 的元素 $\beta=c_1\beta_1+\cdots+c_s\beta_s$，便得到一一對應，且和與差相互對應。由 $R_\alpha=R_\beta$ 還可知，$\gamma\alpha_i$ 與 $\gamma\beta_i$ 由各自基的同一線性形式表示，所以一般地 $\gamma\alpha$ 與 $\gamma\beta$ 對應；因此 $\ma$ 與 $\mb$ 屬於同一理想類。單位理想的類稱為\emph{主類}；$\mR$ 的每一組基因而都給出主表示類。

\textbf{3. 一個元素關於一個類的跡。} 令 $C$ 為 $R_\alpha$ 中指定給 $\mR$ 元素 $\gamma$ 的矩陣。從 $R_\alpha$ 過渡到等價環可知，$|C|$，更一般地 $|tE-C|$，其中 $t$ 為不定元，是表示類的不變量。因此按第 2 節，它同時也是理想類的不變量，簡稱為\emph{類不變量}。於是 $|tE-C|$ 關於 $t$ 的各個係數都是類不變量；其中 $(-t^{s-1})$ 的係數稱為 $\gamma$ 關於該類的跡：
\[
        S_{(\ma)}(\gamma)=c_{11}+c_{22}+\cdots+c_{ss}.
\]
常數項 $\pm |C|$ 稱為 $\gamma$ 關於該類的範數。

\emph{對於固定類 $(\ma)$，跡 $S_{(\ma)}(\gamma)$ 構成一個 $\mZ$-模，而把 $\mR$ 看作 $\mZ$-模時，$\mR$ 同態到這個模。}事實上，由第 1 節所證明的從 $\mR$ 到 $R_\alpha$ 的環同態，有 $(\beta-\gamma)\sim(B-C)$ 以及 $c\beta\sim cEB=cB$；所以
\[
        S_{(\ma)}(\beta-\gamma)=S_{(\ma)}(\beta)-S_{(\ma)}(\gamma),
        \qquad
        S_{(\ma)}(c\beta)=cS_{(\ma)}(\beta).
\]
\emph{若 $\mR$ 是無零因子的環，則一個元素的跡與所取類無關；可直接稱為跡。範數也一樣。}因為此時每個非零理想都具有環的秩，而一個理想的每組基同時也是商域中單位理想的一組基。於是所有跡和範數在商域中同時由單位理想生成，因而在 $\mR$ 中與類無關。

\textbf{4. 理想關於一個類的判別式。} 成為理想與類的不變量的並不是判別式本身，而只是它在 $\mZ$ 中導出的主理想。令 $\omega_1,\ldots,\omega_s$ 與 $\bar\omega_1,\ldots,\bar\omega_s$ 為 $\mR$ 的一個理想 $\mt$ 的兩組線性無關 $\mZ$-模基。則行列式
\[
        |S_{(\ma)}(\omega_i\omega_j)|,
        \qquad
        |S_{(\ma)}(\bar\omega_i\bar\omega_j)|
\]
\emph{只相差 $\mZ$ 中一個單位的平方。}基元素乘積的兩個行列式也同樣如此；而由於模同態，跡之間有同產品之間相同的線性關係，所以在同變換下行列式取得同一個二次因子。

每個行列式 $|S_{(\ma)}(\omega_i\omega_j)|$ 稱為 $\mt$ 關於類 $(\ma)$ 的一個判別式，它只在一個單位平方因子的意義下確定；由它在 $\mZ$ 中導出的主理想稱為判別式理想。\emph{若 $\mR$ 是無零因子的環，則理想的判別式同所用的類無關。}因為按第 3 節，行列式的每一個元素都已經具有這種獨立性。

\subsection*{\S 5. 類的直和}

\textbf{1. 理想類或表示類的直和。} 若理想 $\ma$ 是直和 $\ma=\mb+\mc$，則由同構性可知，該類中的每個理想也都是直和 $\bar\ma=\bar\mb+\bar\mc$。這裡 $\bar\mb$ 分別同構於 $\mb$，$\bar\mc$ 同構於 $\mc$，把這些理想看作 $\mR$-模；因此可以談論理想類的直和及其分量類。

若表示類中的一個環 $R_\alpha$ 是若干子環的直和，且這些子環同時是 $R_\alpha$ 中的理想，則由同構性可知，該表示類中每個環 $R_{\bar\alpha}$ 也有同樣性質，並且各分量是等價環。因此可以談論表示類的直和及其分量類。\emph{理想類的直和對應於表示類的直和；在這個意義上說類的直和。}\footnote{並不在這裡討論所有表示類直和在多大程度上都能由理想類直和產生。}

令 $\beta_1,\ldots,\beta_l$ 為 $\mb$ 的一組基，$\gamma_1,\ldots,\gamma_m$ 為 $\mc$ 的一組基；由於和為直和，所有 $\beta$ 與 $\gamma$ 合在一起構成 $\ma$ 的一組線性無關基。若 $R_\alpha$ 是由這組基生成的矩陣環，則 $R_\alpha$ 中指定給 $\mR$ 任意元素 $\delta$ 的矩陣具有形式
\[
        \begin{pmatrix}D^{(\beta)}&0\\0&D^{(\gamma)}\end{pmatrix}.
\]
定義 $R^{(\beta)}$ 為全體形如 $\left(\begin{smallmatrix}D^{(\beta)}&0\\0&0\end{smallmatrix}\right)$ 的矩陣所成的系統，並類似地定義 $R^{(\gamma)}$。$R^{(\beta)}$ 在同態下對應於一切使 $\delta\mc$ 消失的 $\delta\in\mR$，也就是商理想 $(0):\mc$；所以 $R^{(\beta)}$ 是 $R_\alpha$ 的一個理想，$R^{(\gamma)}$ 也一樣。$R_\alpha$ 是這兩個理想的和，而且是直和，因為乘積 $R^{(\beta)}R^{(\gamma)}$ 為零。最後，把 $D^{(\beta)}$ 對應到 $R^{(\beta)}$ 中含有 $D^{(\beta)}$ 與零塊的矩陣，給出 $R^{(\beta)}$ 與由 $\mb$ 生成的環 $\mR_\beta$ 之間的同構；對 $\mc$ 也一樣。因此，理想類 $\ma$ 的分量類 $(\mb),(\mc)$ 在這個意義下對應於 $R_\alpha$ 的分量類。

\textbf{2. 跡和判別式在類直和下的行為。} 若 $\ma=\mb+\mc$，並用塊矩陣來決定跡 $S_{(\ma)}(\delta)$，即可讀出：\emph{一個元素關於一個類所取的跡，等於它關於各分量類所取跡的和。}又因乘積 $\mb\mc$ 為零，還可讀出：\emph{若 $\beta$ 是 $\mb$ 的元素，則 $S_{(\ma)}(\beta)=S_{(\mb)}(\beta)$；$\beta$ 關於類 $(\ma)$ 的跡等於它關於分量類 $(\mb)$ 的跡。}

\emph{由此可知，分量理想 $\mb$ 關於類 $(\ma)$ 的判別式理想，等於 $\mb$ 關於分量類 $(\mb)$ 的判別式理想。}\emph{最後，$\ma$ 關於 $(\ma)$ 的判別式理想，等於各分量關於分量類的判別式理想的乘積。}因為在構造判別式和跡時都採用與直和相適應的 $\beta,\gamma$ 基，則 $\ma$ 關於 $(\ma)$ 的判別式理想的一組基為
\[
\left|
\begin{array}{cc}
S_{(\ma)}(\beta_i\beta_k)&0\\
0&S_{(\ma)}(\gamma_\mu\gamma_\nu)
\end{array}
\right|
=
\left|
\begin{array}{cc}
S_{(\mb)}(\beta_i\beta_k)&0\\
0&S_{(\mc)}(\gamma_\mu\gamma_\nu)
\end{array}
\right|
\quad
\left(
\begin{matrix}
 i,k=1,\ldots,l\\
\mu,\nu=1,\ldots,m
\end{matrix}
\right).
\]
右端分塊行列式正是 $|S_{(\mb)}(\beta_i\beta_k)|\,|S_{(\mc)}(\gamma_\mu\gamma_\nu)|$。

\textbf{3. 過渡到擴張環。} 令 $\mX$ 為 $\mZ$ 的一個無零因子擴張環，且 $\mR$ 與 $\mX$ 的交為 $\mZ$；令由附加 $\mX$ 得到的擴張環 $\bR=\mR[\mX]$ 定義為同秩的環。若 $\ma$ 與 $\mt$ 是 $\mR$ 中的理想，$\ba$ 與 $\bt$ 是它們在 $\bR$ 中的擴張理想，則 $\bt$ 關於 $(\ba)$ 的判別式理想，也就是 $\mt$ 關於 $(\ma)$ 的判別式理想在 $\mX$ 中的擴張理想。因為 $\ma$ 或 $\mt$ 的每組基，也同時是 $\ba$ 或 $\bt$ 的 $\mX$-模基。特別地，$\bR$ 關於主類的判別式理想，是 $\mR$ 關於主類的判別式理想的擴張理想；以此方式在 $\mZ$ 或 $\mX$ 中定義的理想，以下簡稱為相應環的\emph{判別式理想}。

\subsection*{\S 6. 第一類完全可約性的判別式準則}

從現在起，重新假定 $\mR$ 是域 $\mathsf{P}$ 的一個擴張環，而 $\overline{\mathsf{P}}$ 是 $\mathsf{P}$ 的一個擴張域。判別式不為零將成為第一類完全可約性的充要條件。按 \S 5 第 3 節，$\mR$ 的判別式理想可藉助擴張環 $\bR=\mR[\overline{\mathsf{P}}]$ 來決定；按 \S 5 第 2 節，只需限制到 $\bR$ 加法分解所含的準素環。因此先考察準素環的判別式理想。

\textbf{1. 特殊 $P$-模基。} 若 $\ma$ 與 $\mt$ 是 $\mR$ 的理想，且 $\ma$ 被 $\mt$ 整除，則 $\ma$ 的每組基 $a_1,\ldots,a_s$ 都可通過添入元素 $\tau_{s+1},\ldots,\tau_t$，補成 $\mt$ 的一組線性無關模基。\footnote{由這組基給出的 $\mR$ 中任意元素 $\gamma$ 的矩陣表示具有形狀 $\left(\begin{smallmatrix}C_{\ma}&0\\C_1&C_2\end{smallmatrix}\right)$；反過來，由 $\mt$ 給出的這種矩陣表示又推出 $\mt$ 有一個倍理想 $\ma$。在完全可約情形中，不可分分量除零理想和單位理想外沒有別的理想，所以這種表示不可能出現在這些分量理想的表示類中。再結合 \S 5 第 1 節所給出的類的直和之間的聯繫，便得到文獻中通常採用的完全可約性定義。} 這是關於域 $P$ 的有限秩這一假設的直接推論。若 $\mR$ 是準素環，$\mpideal$ 是屬於零理想的素理想，$\rho$ 為指數，使 $\mpideal^\rho=(0)$ 且 $\mpideal^{\rho-1}\ne(0)$，則可通過逐次把 $\mpideal^{\rho-1}$ 的一組基補成 $\mpideal^{\rho-2}$ 的一組基，並一般地把 $\mpideal^i$ 的一組基補成 $\mpideal^{i-1}$ 的一組基（把 $\mR$ 計作 $\mpideal^0$），得到 $\mR$ 的一組特殊 $P$-模基。

藉助這樣的基可證明：\emph{若 $\mR$ 是準素環，則每個被相應素理想 $\mpideal$ 整除的元素，其關於主類所取的跡為零。}事實上，令
\[
\pi_{0,1},\ldots,\pi_{0,i_0};\quad
\pi_{1,1},\ldots,\pi_{1,i_1};\quad\ldots;\quad
\pi_{\varrho-1,1},\ldots,\pi_{\varrho-1,i_{\varrho-1}}
\]
為這樣一組特殊 $P$-模基，其中
\[
\pi_{\lambda,j}\equiv0\pmod{\mpideal^\lambda},
\qquad
\pi_{\lambda,j}\not\equiv0\pmod{\mpideal^{\lambda+1}}.
\]
若 $\pi\equiv0\pmod{\mpideal}$，則 $\pi\pi_{\lambda,j}\equiv0\pmod{\mpideal^{\lambda+1}}$；因此由這組基指定給 $\pi$ 的矩陣在對角線上全為零，$\pi$ 的跡也為零。

\textbf{2. 帶代數閉子域 $\mathsf{P}$ 的準素環的判別式理想。} 若 $\mR$ 的零理想是素理想且 $\mathsf{P}$ 代數閉，則 $\mR$ 秩為一，因而同 $\mathsf{P}$ 相同（\S 3，第 3 節）；單位元 $e$ 可取為模基。於是
\[
        S_{(\mR)}(e^2)=S_{(\mR)}(e)=e,
\]
判別式理想為單位理想。

帶代數閉子域 $\mathsf{P}$ 的真正準素環的判別式理想為零理想。若使用第 1 節的特殊模基，或只把 $\mpideal$ 的任意一組基補成 $\mR$ 的一組基，則因為剩餘類環 $\mR/\mpideal$ 的秩為一，第一基塊只有一個元素，可取為單位元 $e$。所有不同於 $e^2$ 的基元素乘積都被 $\mpideal$ 整除；它們的跡為零。又因 $\mR$ 被假定為真正準素環，它至少還有一個不同於 $e$ 的基元素。因此，判別式作為至少二階的行列式消失，其中除 $S_{(\mR)}(e^2)$ 外所有項均為零。

\textbf{3. 定理。} \emph{有限秩環 $\mR$ 在子域 $\mathsf{P}$ 上第一類完全可約，當且僅當其判別式理想是 $\mathsf{P}$ 中的單位理想；換言之，當且僅當在單位平方因子意義下確定的判別式非零。}

按 \S 3 第 2 節，第一類完全可約性等價於斷言擴張環 $\bR=\mR[\overline{\mathsf{P}}]$ 完全可約，也就是它的不可分分量都是同構於 $\overline{\mathsf{P}}$ 的域。按 \S 5 第 3 節，$\mR$ 的判別式理想、同時也是 $\bR$ 的判別式理想，不是零理想就是單位理想；按 \S 5 第 2 節，$\bR$ 的判別式理想是各分量關於這些分量自身所取判別式理想的乘積。各個判別式理想從第 2 節所考察者得到，只需每次把同構於 $\overline{\mathsf{P}}$ 的域 $\overline{\mathsf{P}}_i=\overline{\mathsf{P}} e_i$ 替換為 $\overline{\mathsf{P}}$ 本身。因此，若 $\bR$ 完全可約，則其判別式理想作為單位理想的乘積仍為單位理想；若 $\bR$ 不完全可約，則至少一個分量的判別式理想等於零理想，因而 $\bR$ 的判別式理想為零。定理得證。

\textbf{4. 在第一類完全可約情形中，由共軛元素表示跡、範數和判別式。} 設在第一類完全可約情形中，按 \S 3 第 3 節有
\[
        \mR[Q]=Qe_1+\cdots+Qe_n
\]
表示為秩一環的直和，並令 $Q$ 為能作出這種表示的最小擴張域。於是 $\mR$ 的 $n$ 個分量由 $K_i e_i$ 給出，其中每個 $K_i$ 是 $Q$ 的子域，而 $\mR$ 同態到 $K_i$。

取 $e_1,\ldots,e_n$ 為模基，則每個元素
\[
        \gamma=c_1e_1+\cdots+c_ne_n
\]
被指定為對角矩陣 $\operatorname{diag}(c_1,\ldots,c_n)$，其對角元 $c_i$ 在 $\gamma\in\mR$ 時屬於 $K_i$。因此
\[
        S(\gamma)=c_1+\cdots+c_n,
        \qquad
        N(\gamma)=c_1\cdots c_n,
\]
其中 $S(\gamma)$ 是關於主類的跡，$N(\gamma)$ 是範數。若 $\alpha_1,\ldots,\alpha_n$ 是由 $\mR$ 元素組成的 $\mR$ 的一組模基，並且
\[
        \alpha_j=\alpha_j^{(1)}e_1+\cdots+\alpha_j^{(n)}e_n,
\]
則 $\alpha_i\alpha_j=\sum_{\lambda=1}^n\alpha_i^{(\lambda)}\alpha_j^{(\lambda)}e_\lambda$，從而
\[
\begin{aligned}
|S(\alpha_i\alpha_j)|
&=\left|\sum_{\lambda=1}^{n}\alpha_i^{(\lambda)}\alpha_j^{(\lambda)}\right|\\
&=\left|
\begin{matrix}
\alpha_1^{(1)}&\cdots&\alpha_n^{(1)}\\
\vdots&&\vdots\\
\alpha_n^{(1)}&\cdots&\alpha_n^{(n)}
\end{matrix}
\right|^2.
\end{aligned}
\]
也就是說，判別式表示為共軛基元素行列式的平方。特別地，若 $\mR$ 本身是一個域，因而是子域 $\mathsf{P}$ 的第一類擴張域，則從 $\mR$ 到 $K_i$ 的同態成為同構，並且 $\mathsf{P}$ 中的元素逐一對應於自身，因為 $\mathsf{P}$ 的相應分量為 $\mathsf{P}\varepsilon_i$。由於各 $K_i$ 都位於同一個包絡域 $Q$ 中，它們便是關於 $P$ 的共軛域；而且這 $n$ 個同構彼此不同，這可由上面非零判別式 $|S(\alpha_i\alpha_k)|$ 的行列式表示看出。

\emph{因此，若 $\mR$ 是域，則它在 $K_1,\ldots,K_n$ 中的直和表示，給出位於伽羅瓦擴張域內、與 $\mR$ 同構的 $n$ 個關於 $\mathsf{P}$ 的共軛 $n$ 次擴張域。跡、範數和判別式的表示於是成為熟知的由共軛元素給出的表示。}還須再次強調，$\mR$ 只被假定為與各 $K_i$ 等價的一個域，而不是已經位於伽羅瓦擴張域 $Q$ 中的域。\footnote{參見 \S 1 第 2 節的最後一條注。}

\subsection*{\S 7. 代數數域或函數域的階的判別式定理}

\textbf{1.} 所考慮的數域和函數域都納入如下帶有固定整元概念的域型；在多個不定元的代數函數和整代數函數情形中，這裡所指的是泛函整環的擴張。\footnote{參見“理想論”，\S 3 第 3、4 節。}

令 $\mo$ 為\emph{無零因子且有單位元的主理想整環}（有理整數環、以一個域為係數的一元多項式環，或泛函整環），令 $\Omega$ 為其商域。令 $K$ 為 $\Omega$ 的一個\emph{第一類有限擴張}，令 $\mO$ 為 $\mathfrak{K}$ 中\emph{全體在 $\mo$ 上整的元素組成的環}。\footnote{關於相對於一個環為整的元素之定義，參見“理想論”，\S 1 第 3 節。}\emph{$\mO$ 的每個含有 $\mo$ 的子環都稱為一個階；$\mO$ 本身稱為整數環（極大階；原文 „Hauptordnung“）。}

$\mO$ 中每個 $\mo$-模，特別是一個階的每個理想以及該階本身，都有\emph{一組關於 $\mo$ 線性無關的模基}。\footnote{“理想論”，\S 3 第 1 節。}若 $K$ 在 $\Omega$ 上次數為 $n$，則只需限制於關於 $\mo$ 秩為 $n$ 的階；因為任何秩較低的階通過形成商域便生成 $\Omega$ 與 $K$ 之間一個次數等於該秩的中間域，而全部假設對它仍然成立。\emph{因此每個階 $\mT$ 都是滿足 \S 4 與 \S 5 假設的環，並且是無零因子環；於是階的判別式 $D_\mT$ 只在 $\mo$ 的一個單位平方因子的意義下唯一確定。}這個判別式也在 $\Omega$ 的一個單位因子意義下同把 $K$ 看作 $\Omega$-模時的域判別式相同，因此\emph{非零}，因為 $K$ 被假定為第一類擴張（\S 6，第 3 節）。同時，該判別式具有 \S 6 第 4 節中由共軛基元素行列式平方給出的表示。

$\mT$ 中的理想論由如下定理給出：\footnote{“理想論”，\S 7 第 3 節。}\emph{在 $\mT$ 中，每個非零理想都可唯一表示為有限多個兩兩互素的準素理想的乘積，而相應的素理想除單位理想外沒有真除子。}

\textbf{2. 從階過渡到域上的有限秩環。} 令 $p$ 為 $\mo$ 的一個素元，使由 $p$ 在 $\mo$ 中導出的主理想成為 $\mo$ 的一個不同於零理想和單位理想的素理想 $\mo p$；相應地，令 $\mT p$ 表示由 $p$ 在 $\mT$ 中導出的主理想。\footnote{從這裡起採用這種主理想記法，以便同時標明所考察的環。}\emph{剩餘類環}
\[
        \mR=\mT/\mT p
\]
\emph{成為一個在子域 $P$ 上秩為 $n$ 的環，其中 $P$ 同構於域 $\mo/\mo p$；$\mT$ 同態到 $\mT\,|\,\mT p$。}

同態 $\mT\sim\mT\,|\,\mT p$ 對任意剩餘類環都成立。$\mR$ 含有一個同構於 $\mo/\mo p$ 的子域這一點來自第二同構定理：\footnote{參見 \S 1 的最後一條注。}子環 $(\mo,\mT p)|\mT p$ 同構於 $\mo|[\mo,\mT p]=\mo|\mo p$。令 $\alpha_1,\ldots,\alpha_n$ 為 $\mT$ 關於 $\mo$ 的一組線性無關模基，令 $(\alpha_i)$ 為這些元素模 $\mT p$ 決定的剩餘類。若 $(\alpha_i)$ 關於 $P$ 有任何線性相關關係，則它來自一個同餘
\[
        c_1\alpha_1+\cdots+c_n\alpha_n\equiv0\pmod{\mT p},
\]
也就是來自關係 $c_1\alpha_1+\cdots+c_n\alpha_n=py$，其中 $y\in\mT$。由 $\gamma$ 在基 $\alpha_i$ 中的表示及該表示的唯一性，可知每個 $c_i$ 都被 $p$ 整除；因此 $P$ 中元素 $(c_i)$ 消失。於是 $(\alpha_i)$ 構成 $\mR$ 在 $P$ 上的一組基。

\emph{在同態下，$\mT$ 的判別式 $D_\mT$ 化為 $\mR$ 的判別式，而 $\mT p$ 的理想分解化為 $\mR$ 中零理想的分解。}事實上，剛剛證明了剩餘類 $(\alpha_i)$ 關於 $P$ 線性無關，故可用這組基把 $\mR$ 的判別式寫成 $|S((\alpha_i)(\alpha_k))|$；它正是由 $\mT$ 關於基 $\alpha_i$ 的判別式 $|S(\alpha_i\alpha_k)|$ 經同態所得。這裡跡的定義以矩陣表示為基礎，而不採用只對 $\mT$ 成立、對 $\mR$ 未必成立的共軛元素表示；因此跡以及判別式確實只由環元素通過環運算構成。

若
\[
        \mT p=\mq_1\cdots\mq_r=[\mq_1,\ldots,\mq_r]
\]
是 $\mT p$ 在 $\mT$ 中的表示，則在同態下它化為 $\mR$ 中的
\[
        (0)=[(\mq_1),\ldots,(\mq_s)].
\]
由第一同構定理，\footnote{參見“理想論”，\S 4 第 3 節及 \S 6（定理 II 之後）。}$\mR\,|\,(\mq_i)$ 同構於 $\mT\,|\,\mq_i$；因此 $\mq_i$ 與 $(\mq_i)$ 同時是真正準素理想或素理想，而 $(\mq_i)$ 通過同態仍兩兩互素。

\textbf{3. 判別式定理。} \emph{$\mo$ 的一個素元 $p$ 整除階 $\mT$ 的判別式 $D_\mT$，當且僅當在 $\mT$ 中主理想 $\mT p$ 分解為兩兩互素準素分量時，至少出現一個真正準素分量，或至少出現一個第二類素理想。}\footnotemark

\emph{若剩餘域 $\mo/\mo p$ 是完美域，則每個整除 $D_\mT$ 的素元 $p$，且只有這樣的 $p$，具有至少一個真正準素分量。若 $\mo/\mo p$ 是完美域且 $\mT$ 取為極大階 $\mO$，則 $p$ 整除判別式，當且僅當它至少被 $\mO$ 的某個素理想平方整除。}

這是因為，由 $\mT p$ 在 $\mT$ 中的理想分解與零理想在 $\mR=\mT\,|\,\mT p$ 中的分解之間的對應關係可知，$\mT p$ 中出現真正準素分量或第二類素理想，正意味着 $\mR$ 不是第一類完全可約。按 \S 6 第 3 節，這當且僅當 $\mR$ 的判別式為零；而按第 2 節的同態，這恰恰說明 $D_\mT$ 被 $p$ 整除。\footnotetext{第二類素理想出現的一個例子如下。令 $\mo$ 為由 $x$ 的所有有理整數係數多項式生成的泛函整環，令 $\Omega$ 為其商域，令 $K$ 為附加 $\sqrt{x}$ 後得到的擴張域。考慮關於 $\mo$ 以 $1,\sqrt{x}$ 為模基的階 $\mT$；代換 $x=y^2$ 表明它是極大階。其判別式為
$\left|\begin{smallmatrix}1&\sqrt{x}\\1&-\sqrt{x}\end{smallmatrix}\right|^2=2^2x$。主理想 $\mT2$ 在 $\mT$ 中仍為素理想，但為第二類。事實上，由於這組模基，$\mT$ 同構於剩餘類環 $\mo[t]|t^2-x$，其中 $t$ 是一個新的不定元；所以 $\mT\,|\,\mT2=\mR$ 同構於 $\mathsf{P}[t]\,|\,t^2-x$，這裡 $\mathsf{P}$ 是由 $x$ 以及泛函整環中出現的不定元 $u$ 所成、係數模 $2$ 的有理函數域。由於 $t^2-x$ 是 $\mathsf{P}[t]$ 中的素函數，而且是第二類素函數，$\mR$ 是一個域，並且是 $\mathsf{P}$ 的第二類擴張域；故 $\mT2$ 是第二類素理想。相反，$\mT x=(\mT\sqrt{x})^2$ 是真正準素理想。一般說來，$\mR$ 並不一定由對 $\mathsf{P}$ 附加一個元素生成；只要把一個不定元 $x$ 換成 $x,y$，便在由 $1,\sqrt{x},\sqrt{y},\sqrt{xy}$ 導出的階中得到一個例子。這個階仍是極大階，而 $\mT2$ 是第二類素理想，且相應的 $\mR$ 只有對 $\mathsf{P}$ 附加兩個元素後才生成。關於這些例子，參見 Ostrowski 前引文獻。}

\subsection*{\S 8. 相對域的階的判別式定理}

若不以主理想整環 $\mo$ 為起始環，而已經以一個數域或函數域的整數環（極大階）為起始環，更一般地說，以一個乘法環\footnote{“乘法環”這一簡短稱呼取自 W. Krull：《論乘法環》，Heidelberger Berichte，1925，第 5 篇，第 3 號。不過 Krull 把這裡出現的無零因子環稱為“正則乘法環”。}為起始環，即一個其中極大階的理想論成立的環，則代替判別式 $D_\mT$ 出現的是 $\mT$ 關於 $\mo$ 的判別式理想\footnote{在“相對數域”的文獻中，只出現極大階的判別式理想，也就是“相對判別式”。參見 Hilbert《數論報告》\S 14，以及 Hecke《代數數理論》\S 38。}；判別式定理通過以下考察完全轉移。

\textbf{1.} 令 $\mo$ 為一個乘法環，即一個無零因子且有單位元的環，在其中每個不同於零理想和單位理想的理想都可唯一表示為素理想的冪乘積，且這些素理想除單位理想外沒有真除子。\footnote{關於乘法環的公理體系，參見“理想論”，例如導論。}令 $\Omega,K,\mO,\mT$ 如 \S 7 第 1 節定義；則在 $\mT$ 中有那裡所述的理想論。\footnote{“理想論”，\S 3 第 2 節及 \S 7 第 3 節。}

特別地，若一個乘法環 $\mo$ 除零理想和單位理想外只有一個素理想 $\mpideal$，則 $\mo$ 是主理想整環。事實上，取 $p\equiv0\pmod{\mpideal}$ 且 $p\not\equiv0\pmod{\mpideal^2}$。按假設，每個不同於零理想和單位理想的理想都是 $\mpideal$ 的一個冪，特別是主理想 $\mo p$ 也是；由 $p\not\equiv0\pmod{\mpideal^2}$ 必有 $\mo p=\mpideal$。所以 $\mpideal$ 的每個冪都是主理想，零理想和單位理想也同樣是主理想。

\textbf{2. 跡和判別式理想。} $\mT$ 的每個元素 $\gamma$ 的跡，\emph{定義為 $\gamma$ 在 $K$ 中的跡，其中 $K$ 被看作 $\Omega$-模}。$K$ 中任意 $n$ 個關於 $\Omega$ 線性無關的元素都可選作一組基，以構造定義跡的矩陣表示。若 $\alpha_1,\ldots,\alpha_n$ 是這樣的系統，則由於 $K$ 被假定為第一類擴張，
\[
        |S(\alpha_i\alpha_j)|
\]
是 $\Omega$ 的一個非零元素，具體而言，是共軛基元素行列式的平方（\S 6，第 4 項）。由 $\mo$ 在 $\Omega$ 中的整閉性\footnote{“理想論”，\S 9 第 7 節。}可知，只要所有 $\alpha_i$ 屬於 $\mO$，這個行列式就是 $\mo$ 的元素。

令 $\tau_1,\ldots,\tau_n$ 遍歷 $\mT$ 中所有 $n$ 個元素組成的系統，並對每個這樣的系統形成行列式 $|S(\tau_i\tau_j)|$。這些行列式屬於 $\mo$，且當 $\tau$ 線性無關時非零。\emph{由所有這些行列式在 $\mo$ 中導出的理想，稱為階 $\mT$ 關於 $\mo$ 的判別式理想；}所以它不是零理想。

再令 $\beta_1,\ldots,\beta_m$（$m\geqq n$）為 $\mT$ 的一組總是存在的 $\mo$-模基。\footnote{“理想論”，\S 9 及 \S 3 第 1 節。}從中每次取 $n$ 個 $\beta$ 元素所形成的有限多個行列式，構成判別式理想的一組理想基。這由 $\mT$ 到 $\mT$ 中各跡之系統的模同態性（\S 4 第 3 節）推出：由乘積 $\tau_i\tau_j$ 構成的行列式 $|(\tau_i\tau_j)|$，都可用 $\mo$ 中係數線性表示為從每 $n$ 個 $\beta$ 元素形成的有限多個行列式 $|(\beta_i\beta_j)|$；因而 $|S(\tau_i\tau_j)|$ 對 $|S(\beta_i\beta_j)|$ 也有同樣關係。特別地，若該階有一組線性無關模基 $\alpha_1,\ldots,\alpha_n$，則行列式 $|S(\alpha_i\alpha_k)|$ 可取為判別式理想的基；這個定義同 \S 7 第 1 節的定義一致。

\textbf{3. 過渡到局部化環（原文 Quotientenring）。}\footnote{本節所引各定理的證明，參見 H. Grell 在原文第 83 頁注 6 所引的論文。} 若 $\ma$ 是 $\mT$ 的任意理想，則局部化環 $\mT_\ma$ 定義為 $\mT$ 的商域中那些分母為與 $\ma$ 互素的 $\mT$ 元素的全體。\footnote{由於 $\mT$ 中的素理想除單位理想外沒有真除子，這裡的“素於”與理想意義下的“互素”相同。}除零理想和單位理想外，環 $\mT_\ma$ 沒有別的素理想，只有屬於 $\ma$ 的素理想的擴張理想。$\mT_\ma$ 的非平凡理想，與 $\mT$ 中那些相應素理想包含在屬於 $\ma$ 的素理想之中的理想之間存在一一對應，且對應的剩餘類環同構；此外零理想和單位理想也一一對應。

同樣的考察適用於乘法環 $\mo$。特別地，若 $\mpideal$ 是一個素理想，則 $\mo_\mpideal$ 除零理想和單位理想外只有一個素理想；由於剩餘類環同構，$\mo_\mpideal$ 同 $\mo$ 一樣是乘法環，故由第 1 節它是主理想整環。由 $\mpideal$ 導出的素理想的基成為 $\mo_\mpideal$ 的一個素元 $p$。$\mo$ 的任意理想 $\mc$ 的擴張理想，由 $\mc$ 中屬於 $\mpideal$ 的準素分量生成；因此它等於 $(p)^\rho$ 或單位理想，視 $\mpideal$ 是否以第 $\rho$ 次冪出現在 $\mc$ 中而定。若兩個理想在所有局部化環 $\mo_\mpideal$ 中的擴張一致，則這兩個理想相同，因為它們恰被同樣的素理想 $\mpideal$ 以同樣的冪次整除。

\textbf{4. 判別式理想在過渡到局部化環下的行為。} 令 $\mpideal$ 為 $\mo$ 的一個素理想，並置 $\mP=\mT\mpideal$ 為 $\mT$ 中的擴張理想。\emph{局部化環 $\mT_\mP$ 具有由關於主理想整環 $\mo_\mpideal$ 線性無關元素組成的一組模基；同時 $\mT_\mP$ 是 $\mo_\mpideal$ 的商域的擴張域 $K$ 中的一個有限 $\mo_\mpideal$-階。}\footnote{證明見 H. Grell。}\emph{因此在 $\mT_\mP$ 中，關於 $\mo_\mpideal$ 的判別式定理（\S 7，第 3 節）成立。}

於是判別式 $D_{\mT_\mP}$ 是 $\mT$ 關於 $\mo$ 的判別式理想在 $\mo_\mpideal$ 中的擴張理想的一組基。事實上，選擇（這總是可能的） $\mT$ 中元素 $\alpha_1,\ldots,\alpha_n$ 作為 $\mT_\mP$ 關於 $\mo_\mpideal$ 的一組模基；則 $|S(\alpha_i\alpha_j)|$ 屬於判別式理想，而根據第 1 項，局部化環中的每個行列式 $|S(\tau_i\tau_j)|$ 都被 $|S(\alpha_i\alpha_j)|$ 整除。\footnote{當 $\mo$ 是一個代數數域的極大階時，由此得到同 Hilbert 所定義的相對判別式相一致；該定義又已知同 Hecke 的定義相一致。設 $\omega_1,\ldots,\omega_m$ 是相對域極大階的一組 $\mo$-模基；Hilbert 特別選擇關於有理整數環的模基。按 Hilbert 的定義，從每 $n$ 個基元素及其共軛元素形成所有行列式；“相對判別式”就是由每兩個這種行列式（可以相同，也可以不同）的乘積在 $\mo$ 中導出的理想。因此在 $\mo_\mpideal$ 中，相對判別式化為基 $\alpha_1,\ldots,\alpha_n$ 及其共軛元素所成行列式的平方，也就是 $|S(\alpha_i\alpha_j)|$。所以，相對判別式與相對極大階的判別式理想在每一個局部化環 $\mo_\mpideal$ 中的擴張都相同；由第 3 節最後的結論，這兩個理想本身相同。}

\textbf{5. 一般判別式定理。} \emph{乘法環 $\mo$ 的一個素理想 $\mpideal$ 出現在某個階關於 $\mo$ 的判別式理想中，當且僅當在 $\mT\mpideal$ 於 $\mT$ 中分解為兩兩互素準素分量時，至少出現一個真正準素分量，或至少出現一個第二類素理想。}

由第 3 和第 4 節可知，$\mo$ 的一個素理想 $\mpideal$ 出現在 $\mT$ 的判別式理想中，當且僅當 $\mo_\mpideal$ 的素元 $p$ 整除判別式 $D_{\mT_\mP}$。這又等價於說剩餘類環 $\mT_\mP\,|\,\mT_\mP p$ 不是第一類完全可約；而由於這個剩餘類環同構於 $\mT\,|\,\mT\mpideal$，一般判別式定理得證。

作為數域相對判別式的特例，得到：\emph{數域極大階的一個素理想 $\mpideal$ 出現在擴張域的相對判別式中，當且僅當在擴張域的極大階中，$\mpideal$ 被某個素理想的平方整除。}

哥廷根，1926 年 3 月 30 日。

%
\clearpage
%
\setcounter{footnote}{0}

%
\pdfbookmark[0]{32. 與 R. Brauer 合作：不可約表示的極小分裂域}{paper-32}
\section*{32. 與 R. Brauer 合作：不可約表示的極小分裂域}
\emph{Sitz. Ber. d. Preuß. Akad. d. Wiss. 1927, pp. 221--228.}

\begin{center}\emph{（Schur 先生提交。）}\end{center}

關於如下數域的問題，最早由 I. Schur 處理：在一個基域 $P$ 上，一個有限群在 $\mathrm{P}$ 上不可約的表示，在哪些相對於 $\mathrm{P}$ 次數最小的數域中分解為絕對不可約組成部分。\footnote{Arithmetische Untersuchungen über endliche Gruppen, Sitzungsber. d. Berl. Ak. d. Wiss. 1906, p. 164. Beiträge zur Theorie der Gruppen linearer Substitutionen, Transact. of the Am. Math. Soc. Ser. 2, Vol. XV, 1909, p. 159。在第二篇文章中，結果被轉移到完全可約超複系統。一種表示同它的所有變換（相似表示）合併為一個類。} 例如設 $\mathrm{P}$ 含有這樣一個組成部分的特徵標。Schur 於是證明，這些域關於 $\mathrm{P}$ 的次數，即指數，等於那些全都屬於同一類的絕對不可約組成部分的個數；因而，同一個域已經由分裂出一個絕對不可約組成部分這一要求所確定；並且，使這種分裂發生的任意域關於 $\mathrm{P}$ 的次數都是該指數的倍數。所有這些域均稱為分裂域，即便在 $\mathrm{P}$ 不含該特徵標的情形下也如此。Schur 通過一個例子說明，次數最小的分裂域未必同構。若 $\mathrm{P}$ 不含相應的單特徵標，則分裂域必須包含這樣一個特徵標的域以及它關於 $\mathrm{P}$ 的共軛域；屬於不同共軛特徵標的絕對不可約表示固然彼此共軛，但顯然屬於不同的類。為了分裂出同一類中的一個絕對不可約表示系統，在這裡也只需取相應特徵標的域以及相應分裂域之一。

下文將進一步考察分裂域問題。次數最小的分裂域可以由所配屬的非交換除環來刻畫；一般分裂域則由與這些非交換除環不變地相連的雙側單環來刻畫。進一步，通過四元數除環的例子，將說明極小分裂域的次數沒有上界；這裡，一個分裂域稱為極小的，如果它的任一真子域都不是分裂域。這種無界性本質上來自一個數論存在定理，其證明由 H. Hasse 在緊接着的短文中給出。不過這個例子還會以初等計算方式處理，更像是一種驗證；由此，不借助一般理論和數論存在定理，也可通過初等地證明一個不那麼深入但已經足夠的存在定理，來說明次數的無界性。\footnote{這個例子出自 E. Noether；它基於對 R. Brauer 一個例子的修改，在 Brauer 的例子中一般地顯示了次數超過指數的極小分裂域的存在。該例子的初等處理出自 R. Brauer。分裂域的刻畫追溯到兩位作者的研究，而這些研究在其他方面追求不同目標。對 E. Noether 而言，問題是在一般有限性假設下，以模論和理想論為基礎建立表示論；對 R. Brauer 而言，則是在給定交換完美基域上，藉助因子系構造有限次數的非交換除環。兩位作者獨立得到次數最小分裂域的刻畫；R. Brauer 是在完美基域情形下，E. Noether 則沒有這個限制。R. Brauer 是在同 E. Noether 關於非交換嵌入可能性的一個問題相聯繫時，發現一般分裂域的刻畫；E. Noether 在這裡也能去掉對完美基域的限制。} 還須指出，在這個例子中，我們再次處理的是有限群的分裂域，即四元數群的分裂域，而 Schur 的例子也以它為基礎。那裡給出的表示同時也是四元數除環的（同構）表示。

\subsection*{1. 分裂域的刻畫}

先簡短彙集基本概念。令 $\mS$ 表示關於域 $\mathrm{P}_0$ 的一個超複系統。所謂 $\mS$ 的一個表示 $\Gamma$，是在 $\mathrm{P}$ 中，指一個元素取自 $\mathrm{P}$ 的矩陣系統，使 $\Gamma$ 成為 $\mS$ 的同態像，並且 $\mathrm{P}_0$ 中的元素 $\rho_0$ 被指定為對角矩陣 $\rho_0E$；因此 $\mathrm{P}$ 必須包含 $\mathrm{P}_0$。一個表示 $\Gamma$ 連同它的所有變換 $\mathit{P}^{-1}\Gamma \mathit{P}$ 構成一個表示類。有限群 $\mathfrak G$ 的表示包含在這樣定義的表示之中；所配屬的超複系統 $\mS$，即“群環”，由所有“群數”組成，也就是由群元素的全部線性組合組成，係數取自 $\mathrm{P}_0$，原來的群元素被視為線性無關，乘法由群乘法和算術法則定義。所謂“可約”“不可約”“完全可約”和“絕對不可約”表示，按通常方式定義；這些術語也推廣到表示類。若 $\mS$ 在 $\mathrm{P}$ 中的所有表示類都完全可約，則稱系統 $\mS$ “在 $P$ 中完全可約”。

現在假定 $\mS$ 在 $\mathrm{P}_0$ 中完全可約；又假定 $\mathrm{P}_0$ 是完美域，從而對 $\mathrm{P}_0$ 的每個代數擴張 $\mathrm{P}$ 也同樣如此。在 $\mathrm{P}_0$ 的每個擴張 $\mathrm{P}$ 下，表示類仍保持完全可約。\footnote{若 $\mathrm{P}_0$ 不是完美的，則完全可約性保持，當且僅當中心的分量 $\mathfrak{K}$ 成為 $\mathrm{P}_0$ 的“第一類擴張”。在這個假設下，正文中的所有進一步推論仍然有效。} $\mS$ 的中心成為有限多個域的直和；$\mathrm{P}_0$ 的每個同構於這樣的域 $K$ 的擴張 $P$，都是某個（單）特徵標的域。若把 $\mathrm{P}_0$ 擴張到同構於 $\mathfrak{K}$ 的 $\mathrm{P}$，則從 $\mS$ 中分裂出一個雙側單環，配屬於該特徵標的絕對不可約表示類。按 Wedderburn 定理，這個環是 $\mathrm{P}$ 上一個矩陣環與一個非交換除環 $\mA$ 的直積；$\mA$ 在同構意義下唯一確定，其中心為 $P$，關於 $P$ 的次數為 $m^2$，其中 $m$ 表示屬於該特徵標的 Schur 指數。$\mS$ 的每個分裂域，就屬於該特徵標的表示而言，同時也是 $\mA$ 關於 $\mathrm{P}$ 中表示的分裂域，反之亦然。$\mS$ 的共軛表示可從 $\mathrm{P}_0$ 上相應的雙側單環出發得到；其所配屬的除環 $\mA_0$ 成為同構於 $\mA$ 的除環，並以 $\mathfrak{K}$ 為中心。因此，分裂域問題完全化歸為所配屬的非交換除環 $\mA$ 及其分裂的問題。\footnote{這對應於 Schur 的化歸：考慮 $\mathrm{P}$ 上同 $\mS$ 在 $\mathrm{P}$ 中不可約表示可交換的矩陣系統。這些矩陣的轉置給出 $\mA$ 在 $\mathrm{P}$ 中的一個表示。} 特別地，有如下定理：

$\mA$ 的所有極大交換子域關於 $P$ 都有同一次數 $m$，其中 $m$ 表示 Schur 指數。所有這些極大交換子域都是次數最小的分裂域，而每個次數最小的分裂域都包含在它們之中。每個分裂域的次數都是 $m$ 的倍數。環 $\mA$ 關於 $\mathrm{P}$ 只有一個絕對不可約表示類，同樣也只有一個不可約表示類。

為了刻畫一般分裂域，令 $\mA_r$ 表示 $\mA$ 與 $P$ 上次數為 $r^2$ 的矩陣環的直積（即元素取自 $\mathrm{P}$ 的所有 $r$ 行矩陣系統）。於是 $\mA_r$ 是雙側單環，並以 $\mA$ 為所配屬的非交換除環；而且每個關於 $\mathrm{P}$ 為超複的、配屬於 $\mA$ 的雙側單環，都以這種方式得到。$\mA_r$ 也可刻畫為元素取自 $\mA$ 的所有 $r$ 行矩陣系統。

$\mA$ 的次數為 $mr$ 的每個分裂域，只要這樣的域存在，都同構於 $\mA_r$ 的一個極大交換子域。反過來，$\mA_r$ 的所有極大交換子域都是分裂域，各自次數為 $ms$，其中 $s$ 是 $r$ 的一個除數。在正則情形中，$\mA_r$ 的所有極大交換子域次數均為 $mr$。這裡所謂正則情形，是指基域 $\mathrm{P}$ 具有如下性質：對於每個在 $P$ 上交換且在 $P$ 上有限的域 $\Sigma$，仍存在 $\Sigma$ 上任意次數的交換擴張。\footnote{因此，例如當 $\mathrm{P}$ 為全體實數的域時，就不是正則情形。}

在 $r>1$ 時，這些分裂域中是否也出現極小分裂域，還不能僅由嵌入 $\mA_r$ 來決定。下面的例子說明，這種情形確實會對無窮多個 $r$ 發生。

\subsection*{2. 極小分裂域次數的無界性}

為了說明這種無界性，現在以四元數除環為例，把它視為有理數域 $\mathrm{P}$ 上的一個超複系統；除單位 $1$ 外，其基元素為 $i,j,k$，並有熟知關係：
\[
 i^2=j^2=k^2=-1;\qquad ij=k;\quad jk=i;\quad ki=j;\quad ji=-k;\quad kj=-i;\quad ik=-j.
\]
將證明，恰好那些代數數域是分裂域：在附加它們之後，存在一個“冪等”元素 $r$，使 $r^2=r$，且 $r$ 不同於零元和單位元。這些數域也可刻畫為使 $(-1)$ 可表示為三個平方和、因而也可表示為兩個平方和的數域。\footnote{$(-1)$ 可表示為三個平方和總蘊含可表示為兩個平方和，這可由恆等式
\[
(c^2+d^2)(a^2+b^2+c^2+d^2)=(ac+bd)^2+(ad-bc)^2+(c^2+d^2)^2
\]
看出；該恆等式例如來自四元數範數的乘積定理。對於分裂而言，不要求存在冪等元而只要求存在範數為零的元素，也已經足夠。} 事實上，置
\[
        r=\alpha+\frac12\beta i+\frac12\gamma j+\frac12\delta k;
\]
則
\[
        r^2=\left(\alpha^2-\frac14\beta^2-\frac14\gamma^2-\frac14\delta^2\right)
          +\alpha\beta i+\alpha\gamma j+\alpha\delta k,
\]
所以 $r^2=r$ 等價於
\[
4\alpha^2-4\alpha-\beta^2-\gamma^2-\delta^2=0;\quad
2\alpha\beta-\beta=0;\quad 2\alpha\gamma-\gamma=0;\quad 2\alpha\delta-\delta=0,
\]
從而，由於 $\beta,\gamma,\delta$ 不應同時為零，又等價於
\[
        (-1)=\beta^2+\gamma^2+\delta^2.
\]

每個這樣的冪等元素都產生一個分裂，而這種冪等元素的存在又是分裂的必要條件；這一點來自如下事實：完全可約系統的每個不可約表示由一個單側單理想、更準確地說由相應理想類生成，並且每個這樣的理想類都生成該表示。這些單側單理想全部出現在一個單側直和分解中，而這個分解決定了成為單位各分量的“本原”冪等元素。反過來，每個冪等元素都給出直和分解 $\mS=r\mS+(e-r)\mS$，其中 $e$ 表示單位；“本原”的冪等元分裂出單理想。\footnote{本原冪等元素與不可約表示之間的聯繫已見於 M. Herzberger 的博士論文：Über Systeme hyperkomplexer Größen, Chapter 4, Berlin 1923。單理想生成不可約表示這一點，在交換情形中例如由 E. Noether, Der Diskriminantensatz für Ordnungen, Crelle 157 (1926) 證明；在非交換情形中由 W. Krull, Theorie und Anwendung der verallgemeinerten Abelschen Gruppen, Heidelberger Berichte 1926 證明。另參 Speiser 群論末尾的說明。若 $a_1,\cdots,a_m$ 表示關於分裂域的理想基，$c$ 是任意環元素，並且 $a_i c=\gamma^i_1 a_1+\cdots+\gamma_{im} a_m$，則 $(\gamma_{ik})$ 是指定給 $c$ 的矩陣；所有不可約表示均由理想生成這一事實還要更深。} 一個非交換除環，作為關於其中心的超複系統，在擴張到關於分裂域的超複系統時，成為 $m$ 個絕對單的單側理想的直和，其中 $m$ 是 Schur 指數。

在四元數除環的情形中，$m$ 為二；因此每個 $\ne0$ 且 $\ne1$ 的冪等元素已經產生分解為絕對單理想的分解，而這正對應於分裂為絕對不可約表示類。由此刻畫了上述域為分裂域。

進一步得到：所有次數為 $2^n$ 的循環域 $\Omega_n$，只要是在有理數域上且其中 $(-1)$ 可表示為兩個平方和的，都是四元數除環的極小分裂域。因為分裂域不可能是實的，這是由 $(-1)$ 作為平方和的表示決定的；另一方面，$\Omega_n$ 中次數為 $2^{n-1}$ 的子域，因而 $\Omega_n$ 的每個真子域，都必然是實的，因為它是 $\Omega_n$ 與全體實代數數的域的交。也就是說，$\Omega_n$ 在這個交之上是次數二的伽羅瓦域，所以這個交必須同次數為 $2^{n-1}$ 的子域一致。

按 Hasse 的存在定理，\footnote{$\Omega_2$ 的存在已可從次數 4 循環域的參數表示讀出，例如見 Weber（Kleines Lehrbuch der Algebra, § 93）。} 對每個 $n$ 都有這樣的域 $\Omega_n$；極小分裂域的次數沒有上界。事實上，這裡指數的每個冪都作為一個極小分裂域的次數出現。\footnote{後來 R. Brauer 還能進一步證明，四元數除環甚至有每個偶數次數的極小分裂域，因而有所有作為分裂域次數可能出現的次數。對於每個 $n\geqq2$，這樣的域 $P(\Theta)$ 可由方程
\[
 f(\Theta)=\Theta^{2n}+a^2p^2\Theta^2+b^2p^2\beta^2=0;
 \quad \text{因此}\quad -1=\left(\frac{ap}{\Theta^{n-1}}\right)^2+
       \left(\frac{bp\beta}{\Theta^n}\right)^2
\]
定義；通過適當選擇 $a,b,p,\beta$，可做到：1. 對每個 $n$，$f(x)$ 不可約；2. $\mathrm{P}(\vartheta)$ 除 $\mathrm{P}$ 外只有一個子域 $\mathrm{P}(\vartheta^2)$；3. $\mathrm{P}(\vartheta^2)$ 的共軛域中也有實域，因而它不能是分裂域。一個可能的 $a,b,p,\beta$ 選擇例如為
\[
        f(x)=x^{2n}+(25\cdot9)4x^2+9\cdot64.
\]
證明基於 Furtwängler 構造本原方程的方法的一個修改（Math. Ann. 85, p. 34, § 3），但細節執行相當繁瑣。}

\subsection*{3. 四元數除環分裂域的初等處理}

令 $\mathfrak{O}$ 表示由以下八個矩陣構成的四元數群表示：\footnote{參見 Sylvester, Math. Papers Vol. III p. 647, Vol. IV p. 122。}
\[
\pm E=\pm\mat{1&0\\0&1},\qquad
\pm I=\pm\mat{i&0\\0&-i},\qquad
\pm J=\pm\mat{0&1\\-1&0},\qquad
\pm K=\pm\mat{0&i\\ i&0}.
\]
由 $\mathfrak{O}$ 生成的群環，正是把四元數除環視為超複系統時的一個表示；我們將初等地證明，$\mathfrak{O}$ 可在域 $\mathrm{K}$ 中表示，當且僅當 $-1$ 可在 $\mathrm{K}$ 中表示為兩個平方和。

\emph{a)} 令 $\mathfrak{O}^*=T^{-1}\mathfrak{O} T$ 是一個在域 $\mathrm{K}$ 上有理的表示，並置 $I^*=T^{-1}IT$，$J^*=T^{-1}JT$。由於 $J$ 與 $J^*$ 是兩個在 $\mathrm{K}$ 上有理且相似的矩陣，也存在一個在 $\mathrm{K}$ 上有理的變換 $R$，把 $J^*$ 變為 $J$：
\[
        R^{-1}J^*R=J.
\]
若一開始就把 $\mathfrak{O}^*$ 替換為同樣 $\mathrm{K}$-有理的表示 $R^{-1}\mathfrak{O}^*R$，便可看出不失一般性可假定 $J^*=J$。令
\[
        I^*=\mat{a&b\\ c&d}\qquad (a,b,c,d\text{ 為 }\mathrm{K}\text{ 中的數}).
\]
由 $IJ=-JI$ 得 $I^*J^*=-J^*I^*$，又因 $J^*=J$，於是
\[
        a=-d,
        \qquad b=c.
\]
由 $I^2=-E$ 得 $I^{*2}=-E$，故
\[
        \mat{-1&0\\0&-1}=\mat{a&b\\ b&-a}^2
        =\mat{a^2+b^2&0\\0&a^2+b^2},
\]
所以 $a^2+b^2=-1$；於是 $-1$ 確實可在 $\mathrm{K}$ 中表示為兩個平方和。

\emph{b)} 反過來，設 $-1$ 可在域 $\mathrm{K}$ 中表示為兩個平方和。若例如 $-1=a^2+b^2$，置
\[
        I^*=\mat{a&b\\ b&-a},\qquad
        J^*=\mat{0&1\\-1&0},\qquad
        \mathit{K}^*=I^*J^*.
\]
直接計算表明，$\pm E$，$\pm I^*$，$\pm J^*$，$\pm \mathit{K}^*$ 構成一個在 $K$ 上有理且同 $\mathfrak{O}$ 相似的群。\footnote{分裂域的刻畫也可藉助因子系容易地得到。}

還需證明：對無窮多個 $n$，存在次數為 $2^n$ 的循環域，它們是四元數群的極小分裂域。

令 $p$ 為一個有理素數，使得對某個合適的 $r>0$，
\[
        (*)\qquad 2^r+1\equiv0\pmod p
\]
並令 $2^n$ 為整除 $p-1$ 的最高的 2 的冪。

首先證明，素數 $p$ 對無窮多個 $n$ 出現。令 $k$ 為任意正有理整數，令 $p$ 為 $2^{2^k}+1$ 的一個素因子。則所需關於 $p$ 的同餘在 $r=2^k$ 時成立。又 $2\pmod p$ 的指數為 $2^{k+1}$，因此，對於整除 $p-1$ 的最高 2 冪 $2^n$，數 $n$ 大於 $k$。於是無論如何，都有任意大的 $n$ 與某些素數 $p$ 相對應。

令 $\varepsilon$ 為一個本原 $p$ 次單位根。現在證明，在 $\mathrm{P}(\varepsilon)$ 中，數 $-1$ 可表示為兩個平方和：
\[
        -1=a^2+b^2=(a+ib)(a-ib)\qquad (a,b\text{ 為 }\mathrm{P}(\varepsilon)\text{ 中的數}).
\]
這等價於說，在 $(4p)$ 次單位根的域 $\mathrm{P}(\varepsilon,i)=\mathrm{P}(\varepsilon\cdot i)$ 中，存在關於 $\mathrm{P}(\varepsilon)$ 的相對範數恰為 $-1$ 的數。

數 $1+i\varepsilon^\nu$（$(\nu=1,2,\ldots,p-1)$）的相對範數為 $(1+i\varepsilon^\nu)(1-i\varepsilon^\nu)=1+\varepsilon^{2\nu}$；因而所有數 $1+\varepsilon^\nu$（$(\nu=1,2,\ldots,p-1)$）都是從 $\mathrm{P}(\varepsilon i)$ 到基域 $\mathrm{P}(\varepsilon)$ 的相對範數。因此，同樣結論適用於
\[
\Pi=\prod_{\nu=0}^{r-1}(1+\varepsilon^{2^\nu})
=(1+\varepsilon)(1+\varepsilon^2)(1+\varepsilon^4)\ldots(1+\varepsilon^{2^{r-1}})
=\sum_{\lambda=0}^{2^r-1}\varepsilon^\lambda.
\]
若在 $\Pi$ 中添上項 $\varepsilon^{2^r}=\varepsilon^{-1}=\varepsilon^{p-1}$（由 $(*)$），則 $\Pi+\varepsilon^{p-1}$ 成為若干段 $1+\varepsilon+\varepsilon^2+\ldots+\varepsilon^{p-1}=0$ 的和，因此
\[
        \Pi=-\varepsilon^{-1}.
\]
現在 $\varepsilon$ 本身也是從 $\mathrm{P}(\varepsilon i)$ 中某個數到 $\mathrm{P}(\varepsilon)$ 的相對範數，即 $\varepsilon^{\frac{p+1}{2}}$ 的相對範數。因此類似斷言也適用於
\[
        \Pi\cdot\varepsilon=-1.
\]
所以 $-1$ 可在 $\mathrm{P}(\varepsilon)$ 中表示為兩個平方和；\footnote{對於形如 $n8\pm3$ 的素數，這一論證已見於 E. Landau, Über die Darstellung definiter Funktionen durch Quadrate, Math. Ann. vol. 62, pp. 272--285。} $\mathrm{P}(\varepsilon)$ 是四元數群的一個分裂域。

令 $\mathrm{K}$ 為 $\mathrm{P}(\varepsilon)$ 中次數為 $2^n$ 的循環域。我們斷言 $\mathrm{K}$ 也是一個分裂域。\footnote{把次數為 $2^n$ 的循環分裂域選為分圓域的子域，是仿照 Hasse 的例子。} 令 $m$ 為四元數除環關於以 $\mathrm{K}$ 為基域時的指數。則 $m=1$ 或 $m=2$。但是，由於存在關於 $K$ 相對次數為奇數 $\frac{p-1}{2^n}$ 的分裂域 $\mathrm{P}(\varepsilon)$，第二種情形不可能；故 $m=1$。這正是說 $\mathrm{K}$ 是一個分裂域。

因此，對無窮多個 $n$，存在次數為 $2^n$ 的循環域，它們是極小分裂域。

%
\clearpage
%
%
\begingroup
\def\srcemph#1{\textbf{#1}}
\def\srctextsc#1{\textbf{#1}}
\setcounter{footnote}{0}
\pdfbookmark[0]{33. 算術觀點下的超複量與表示論}{paper-33}
\section*{33. 算術觀點下的超複量與表示論}
\srcemph{Atti Congresso Bologna 2 (1928), S. 71--73.}

\begingroup
\renewcommand{\thefootnote}{(\arabic{footnote})}

\begin{center}
\textsc{Emmy Noether} (Göttingen -- Germania)\\[-0.1em]
\rule{3em}{0.4pt}\\[0.8em]
\srctextsc{超複量與表示論}\\
\srctextsc{（算術觀點下）}\footnote{詳細論述將發表於 Math. Zeitschr.；在那裡的一篇續文還將討論非交換伽羅瓦理論。關於“分裂域”，亦參見 R. Brauer 與 E. Noether, Ber. Berl. Ak. 1927；關於一般結構定理，參見 G. Köthe 在本卷 Atti 中的報告。}
\end{center}

我想說明，怎樣把表示論——尤其是群表示論——理解為模的算子同構類與理想的算子同構類的理論，以及怎樣把後兩者納入一個推廣的群概念，即帶算子的群。\footnote{帶算子的群可追溯到 W. Krull (Math. Zeitschr. 23) 和 O. Schmidt (Math. Zeitschr. 29)。本文對給定表示之約化問題的處理，在表示除環為交換的限制下，也來自 W. Krull, Theorie und Anwendung der verallgemeinerten Abelschen Gruppen, Heidelberger Ber. 1926。}

所謂\srcemph{一個群的表示}——在預先給定的除環 $P$ 中——眾所周知，是指如下情形：對元素 $a,b$（它們屬於群 $\mathfrak{G}$）分別配以矩陣 $A,B$，矩陣的係數取自 $P$，使乘積 $ab$ 對應於乘積 $AB$；於是這個矩陣系統成為該群的一個同態像。與此同時，也給出了“群環”的一個表示；群環由所有“群數”（即群環元素）組成，也就是所有線性組合 $a\alpha+b\beta+....+c\gamma$，其中 $a,b,....,c$ 每次取有限多個群元素，而 $\alpha,\beta,....,\gamma$ 相應地取 $P$ 中有限多個元素；這裡把 $a,b,....,c$ 視為線性無關，乘法由群乘法和運算法則定義。群環的表示同時關於加法同態；因此，群的表示被納入\srcemph{環的表示}這一一般問題之下，在其中要求同時關於加法和乘法同態。

這裡有兩個本質問題。第一個是\srcemph{給定表示的約化}問題，以及在約化中出現的不可約組成部分是否唯一的問題；第二個則是詢問一個給定環在一個給定的、不必交換的除環中的全部可能表示，例如只考慮其中的不可約表示。

第一個問題遠為簡單；這一點已經表現在，對於待表示的環以及通常非交換的表示除環，都不需要任何特殊假設。表示由固定階數的矩陣給出這一事實，已經提供了一切必要的有限性假設。藉助帶算子的群的理論，這個問題可以完全處理；因此我先簡短說明這一理論。

稱群 $\mathfrak{B}$ 為一個\srcemph{帶算子系的群}，如果同時給定一個符號系統 $\varTheta,H$——即算子系——並且表達式 $\varTheta(a),H(a)$ 對每個 $\mathrm{a}$（取自 $\mathfrak{B}$）都唯一確定，且給出 $\mathfrak{B}$ 中的一個元素，而且滿足分配律 $\varTheta(ab)=\varTheta(a)\varTheta(b)$。兩個帶有同一算子系的群稱為算子同構，如果它們在通常意義下同構，並且由 $a$ 與 $\bar a$ 的對應還推出 $\varTheta(a)$ 與 $\varTheta(\bar a)$、$H(a)$ 與 $H(\bar a)$ 的對應。若只把自身也容許該算子系的子群算作可容許子群（因而所謂單群，是指除單位元外沒有可容許的真正規子群），那麼關於合成列的若爾當--赫爾德定理仍以如下形式成立：如果一個帶算子系的群確有合成列，那麼其因子群（合成因子）的系統在算子同構意義下唯一確定，至多差一個排列。在同一假設下，分解為直接不可分解因子的定理也成立，並且這種分解在算子同構意義下唯一。

它同表示問題的聯繫在於，理想和模尤其屬於帶算子的群：它們相對於加法被看作阿貝爾群，而算子由同環元素相乘給出。每一個表示都由一個表示模產生，也就是由一個關於該環和表示除環的模產生。更準確地說，每一個模類，即彼此算子同構的表示模所成的類，產生同一個表示。這樣，表示約化的唯一性問題便由合成列定理和直積定理回答。把合成列應用於表示模，就對每個矩陣給出如下形式的約化：
\[
\begin{pmatrix}
B_1&0&0\\
&B_2&0\\
A_{ik}&....&B_r
\end{pmatrix}
\]
其中對應於因子群的對角矩陣產生一個“不可約”表示，也就是說，它本身不再容許這樣的約化。

既然這些因子群的類唯一確定，對角表示也同樣唯一，這裡的唯一性是在表示等價的意義下。類似地，直積定理給出“分解”為“不可分解”組成部分的唯一性：
\[
\begin{pmatrix}
C_1&&&\\
&C_2&&\\
&&....&\\
&&&C_s
\end{pmatrix},
\]
其中每個 $C_i$ 隨後還可以按合成列定理唯一地繼續約化。

我想在超複系統的表示情形下，更一般地在滿足“雙鏈條件”的環的情形下，勾畫\srcemph{第二個問題}。其結果是，每一個不可約（單）模的算子同構類都是某個理想的算子同構類；因而，由單側理想組成的合成列已經通過其合成因子給出全部不可約表示。更準確地說，只須考察模去根基（在此即 Jacobson 根）所得商環中的合成列，而該商環成為一個“完全可約”環（現代稱半單環）。問題由此化歸為完全可約環的結構研究；因而，把自同態除環中的表示作為出發點就顯得與問題相適應。這裡所說的是如下事實：這樣一個環的一個雙側單分量中的所有單右理想都屬於同一個類，因此有同一個自同態環；由於這些理想是單的，這個自同態環成為除環，同時也是單左理想的自同態環。這個單環本身同自同態除環上的全體矩陣系統同構；而從這一表示出發，便產生關於一個子除環的一切表示，特別是關於該超複系統的係數域的一切表示，只要把自同態除環的元素本身替換為相應的矩陣。於是，在超複系統情形下已知的 Wedderburn 結構定理，在這裡通過來自一般群論的自同態除環概念得到新的解釋，並且同時被識別為超複系統表示論的來源。由此產生非交換除環的伽羅瓦理論中的新問題，它們同“分裂域”問題緊密相連。同時，通過不可分解理想的自同態環，也展望到一般的、並非完全可約的環的結構定理。群論，在它擴展到最一般的帶算子的群之後，在這裡也再次顯現為一種整理原則。

\endgroup
\endgroup

%
\clearpage
%
\setcounter{footnote}{0}

%
%
%
%
\pdfbookmark[0]{34. 超複量與表示論}{paper-34}
\section*{34. 超複量與表示論}
\emph{Math. Zs. 30 (1929), pp. 641--692.}

\subsection*{引言}

關於超複系統的最重要的一般定理可追溯到 Molien（\emph{Math. Annalen} 第 41 卷以及 1897 年 Dorpat 報告）。基本上與此獨立，Frobenius 不久以後把超複系統理論及其表示理論，特別是有限群的表示論，發展為一個統一理論。其基礎是 Dedekind 提出的群行列式概念，更一般地說，是任意超複系統的行列式概念。Frobenius 證明，這個行列式的各個不可約因子（其中係數域始終取全體複數域）對應於各個不可約表示類，並且所有不可約表示類都以這種方式窮盡。這樣的一個表示類由其“特徵標系”完全刻畫；在有限群的情形中，這個特徵標系來自由群的共軛元素類導出的那個交換超複系統的行列式分解。該行列式分解為線性因子，特徵標作為係數出現；這是 Dedekind 首先得到的結果的直接推廣：有限阿貝爾群的群行列式分解為線性因子，其係數就是這個阿貝爾群的各個特徵標（與 Frobenius 的通信）。\footnote{以下內容是 B. L. van der Waerden 根據我 1927/28 年冬季學期的講課所作的自由整理。我們共同完成了付印稿的加工。我還要感謝 B. L. van der Waerden 提出的若干批評性意見。}

然而，這些在概念上簡單而透明的結果，在 Frobenius 那裡是通過繁難計算得到的。後來的發展旨在簡化這些結果的推導，同時也旨在把它們推廣到任意域作為係數域的情形。這樣的後續發展在超複系統和表示論兩方面完全沿着分離的道路進行。超複系統由 Wedderburn 得到算術化處理：每個超複系統唯一地成為雙側直接不可分解環的直和；對於無根基的系統，這些不可分解環同構於全體 $n$ 階矩陣的系統，其元素取自一個相應的、不必交換的除環，而這個除環在本質上唯一確定。\footnote{參見 Dickson, \emph{Algebras and their Arithmetics}（德文版：\emph{Algebren und ihre Zahlentheorie}, Zürich 1927）中的文獻。} 表示論則由 Burnside 和 I. Schur 給出初等基礎；他們不依賴超複系統，而是直接從一個給定表示出發，並使用矩陣定理。Schur 特別處理瞭如下問題：一個在給定域上不可約的表示，在其中絕對分裂的最小次數的數域。那些絕對不可約組成部分屬於有限多個共軛表示類。所求數域的次數等於這些類的數目與指標（每個共軛類中組成部分的數目）的乘積。Schur 理論的主要輔助工具，是所有與一個不可約表示可交換的矩陣所組成的系統。\footnote{參見 Speiser 群論第 10 章中的文獻。}

在下面的純算術基礎中，超複量與表示論再次表現為一個統一整體，即作為非交換環的一般理論的特例；這些環只滿足某些有限性條件。更準確地說，這是一種關於這些環的模類與理想類的理論，其主要結果是：不可約模類已經由相應的理想類窮盡；特別地，對於無根基的環，所有模類都分解為不可約模類（變為完全可約）。這正是 Frobenius 結果的算術等價物：正則群行列式（相應地，系統行列式）的不可約因子已經窮盡不可約表示類的全體，並且對於無根基系統，表示具有完全可約性。也就是說，考察系統行列式及其分解可以理解為過渡到範數；而在這裡，則直接處理理想的直和分解及其合成列。表示論藉助表示模這一概念，變成模類理論。

模類與理想類的引入有純粹群論的基礎。模和理想被看作關於加法的阿貝爾群，同時受到如下條件的限制：它們容許由環元素給出的某些乘法；也就是說，它們構成“帶算子的群”（§1）。對於這樣的群，通常的同構被“算子同構”（§2）取代；彼此算子同構的群（在一個固定域內）被合併為一個類。這樣的類概念，例如在數域理想的情形中，就與通常的類概念一致。帶算子群的理論可追溯到 W. Krull 和 O. Schmidt（參見注 6）；它將在第一章中系統展開。這裡仍然成立的關於合成列、直積（相應地直和）以及完全可約群的定理，也就是在上述分類意義下的唯一性定理，構成後文一切內容的基礎。

這些定理首先給出任意表示中不可約對角組成部分的一般唯一性定理（第三章 §16）；其根據是表示類與表示模類，也就是帶算子群的類，之間的一一對應。進一步關於所有表示的總體問題，則要求更精細地研究待表示的環的結構；在特殊情形中，就是超複系統的結構。第二章在群論觀點的基礎上重新得到並推進 Wedderburn 的結果；在這種觀點下，單側理想的考察位於前景。事實表明，對於右理想的“多重鏈定理”，或者與之相同的“極小條件”（在任意右理想集合中，集合內至少有一個極小成員），已經足以作為有限性條件。\footnote{如果假定“雙鏈定理”，Wedderburn 的證明方法可以轉移。參見 E. Artin, Zur Theorie der hyperkomplexen Zahlen, \emph{Hamb. Abh.} 1927。另參見 A. Suschkewitsch, Über die endlichen Gruppen ohne das Gesetz der eindeutigen Umkehrbarkeit, \emph{Math. Annalen} 99 (1928), pp. 30--50；其中對只有一個結合的、不可逆運算的（有限）域建立了平行的結構定理。Suschkewitsch 將“核”分裂為右群和左群，對應於把有單位元的雙側單環表示為右理想與左理想的和。} 由此得到：滿足極小條件且無根基的環，正同有單位元的（右）完全可約環相同。在這樣的環中，一個雙側不可分解環的所有單右理想屬於同一類，因而作為帶算子群具有同一個自同構環；由於這些理想是單的，這個自同構環成為除環。這個類的自同構除環正是 Wedderburn 所發現的矩陣表示的中介。

只要要求關於自同構除環或關於其子域的表示，特別是關於一個超複系統的係數域的表示，那麼在完全可約情形中，自同構除環上的這種矩陣表示也同時解決了表示問題。作為把模類歸約為理想類的定理的直接推論，不存在 Wedderburn 表示以及由它通過把自同構除環的元素以自然方式替換為子域上的矩陣而產生的表示以外的其他不可約表示。

因此在完全可約情形中，不同的不可約表示類的數目與理想類一樣多；這個數目又與不同的雙側不可分解、因而單的環的數目一致。最後，這個數目又等同於中心的不可分解組成部分的數目，而這些組成部分成為交換域。在係數域代數閉的超複系統情形中，後面這些域完全由中心的各個環同態（不可約表示）決定；除了一個數值因子外，它們正是特徵標。

所得到的結果同時也給出了帶根基系統的不可約表示的概觀；這是就這些表示可歸約為模去根基後的剩餘類環的表示而言的，而該剩餘類環成為一個無根基系統。

後文將說明，對於無根基的超複系統，“分裂域”理論也從自同構除環推出；也就是說，係數域的那些交換擴張的理論，在這些擴張中表示分解為絕對不可約表示，換言之，發生到絕對單右理想的直和分解。分裂域與自同構除環的分裂域相同，並且所有次數最小的分裂域都同構於自同構除環的極大交換子域。它與上面提到的 Schur 研究之間的聯繫在於：與表示可交換的那些矩陣的轉置，正好給出自同構除環的一個表示。\footnote{參見 R. Brauer--E. Noether, Über minimale Zerfällungskörper irreduzibler Darstellungen, Sitz.-Ber. Preuß. Akad. Wiss. 1927, p. 221。} 這些定理將在非交換除環的伽羅瓦理論框架內系統展開。

作為基礎的概念，即算子同態和自同構環，也給出只滿足極小條件的帶根基一般環的結構；這將從另一個方面加以展開。

\section*{第一章. 群論基礎}

\subsection*{§ 1. 帶算子的群}

設 $\mG$ 是一個群（有限或無限），元素為 $a,b,\ldots$。所謂群 $\mG$ 的一個算子域 $\Omega$，是指一組新的符號 $\mathsf{H},\Theta,\ldots$，使得對 $\mG$ 中每個 $a$ 和 $\Omega$ 中每個 $\Theta$，都有一個 $\mG$ 中的 $\Theta a$ 被唯一確定，並且滿足分配律：
\[
        \Theta(ab)=\Theta a\cdot\Theta b.
\]
因此，每個算子都定義了該群到自身的一個同態；這裡群被映到自身或映到一個子群上。單位元映為單位元，逆元映為逆元。

如果不同的算子也定義不同的同態，則稱一個算子域為\emph{絕對的}。從任意算子域都可以得到一個絕對算子域：只要把產生同一同態的所有元素等同起來。絕對算子域是群到自身的全體同態集合的某個子集的一一像。

相對於一個固定算子域 $\Omega$，群 $\mG$ 的一個\emph{容許子群} $\mH$，是指一個容許該算子域的子群；也就是說，對 $\mH$ 中每個 $a$ 和 $\Omega$ 中每個 $\Theta$，都有 $\Theta a$ 屬於 $\mH$。

\emph{例.} 1. 令算子為內自同構：$\Theta a=c^{-1}ac$。容許子群就是正規子群。

2. 令算子為所有自同構。容許子群就是在所有自同構下映到自身的“特徵子群”。\footnote{這裡解釋的概念來自 W. Krull, Über verallgemeinerte endliche Abelsche Gruppen, \emph{Math. Zeitschr.} 23 (1925), pp. 161--196, 以及 O. Schmidt, Über unendliche Gruppen mit endlicher Kette, \emph{Math. Zeitschr.} 29 (1928), pp. 34--41。}

3. 設 $\mG$ 是一個環，也就是關於加法的阿貝爾群，並且其中還定義了乘法，滿足
\[
        r(a+b)=ra+rb,\qquad (a+b)r=ar+br,\qquad ab\cdot c=a\cdot bc.
\]
每個元素 $r$ 同時定義兩個算子：算子 $rx$ 和 $xr$。容許子群就是“理想” $\mathfrak a$：左理想容許運算 $rx$，即 $r\mathfrak a\subseteqq\mathfrak a$；右理想容許運算 $xr$，即 $\mathfrak a r\subseteqq\mathfrak a$；雙側理想則容許這兩種運算。當環 $\mG$ 是一個除環時，也就是當 $\mG-\{0\}$ 關於乘法成群時，所有理想都成為平凡的 $(=\{0\}\text{ 或 }=\mG)$。\footnote{因此，無論對於環還是對於除環，都不假定乘法交換律。}

4. \emph{關於一個環 $\mo$ 的模.} 設 $\mo$ 是一個環，$\mM$ 是一個按加法書寫的阿貝爾群。假設給定了 $\mo$ 的元素與 $\mM$ 的元素之間的乘法 $r\cdot a$；乘積仍應屬於 $\mM$。要求：
\[
\left.
\begin{aligned}
 r(a+b)&=ra+rb,\\
 rs\cdot a&=r\cdot sa,\\
 (r+s)a&=ra+sa,
\end{aligned}
\right\}\qquad r,s\in\mo,\quad a,b\in\mM.
\]
於是 $\mM$ 稱為一個 $\mo$-模；更準確地說，由於來自 $\mo$ 的乘子寫在左邊，它稱為一個左 $\mo$-模。環 $\mo$ 同時也是算子域。如果它是絕對的，也就是說，如果不同的環元素也給出不同的運算，就稱它為模 $\mM$ 的一個\emph{絕對乘子域}。

如同上面一樣，可以從任意乘子域過渡到絕對乘子域：把產生同一運算的元素等同起來。絕對乘子域仍然是一個環。模 $\mM$ 的容許子群稱為\emph{子模}。特別地，如果 $\mo=\mM$，就又得到左理想。\footnote{對於右模，同樣要求：
\[
(a+b)r=ar+br,\qquad a\cdot rs=ar\cdot s,\qquad a(r+s)=ar+as.
\]}

5. \emph{雙模.} 如果 $\mM$ 同時是左 $\mo$-模和右 $\mo'$-模，並且對於 $\mo$ 中的 $a$、$\mM$ 中的 $\alpha$ 以及 $\mo'$ 中的 $a'$，恆有
\[
        a\cdot\alpha a'=a\alpha\cdot a',
\]
則稱 $\mM$ 為一個雙模。

6. \emph{阿貝爾群的自同構環.}\footnote{參見 A. Châtelet, \emph{Les Groupes Abéliens finis}, Paris, Gauthier-Villars, 1925, p. 99。} 對一個按加法書寫的阿貝爾群的自同態，可以用公式
\[
        (\mathsf{H}+\Theta)a=\mathsf{H}a+\Theta a,\qquad (\mathsf{H}\Theta)a=\mathsf{H}(\Theta a)
\]
定義加法和乘法。如此定義的運算顯然仍表示同態，因而仍屬於該系統。這個系統構成一個環，而該阿貝爾群可按 4 的方式看作關於這個環的模。

下文中，“子群”一詞總是指容許子群。兩個容許子群的交 $\mA\cap\mB$ 仍是容許子群。如果兩個子群 $\mA,\mB$ 中至少有一個是正規子群，則積 $\mA\mB$ 也是容許子群。

\subsection*{§ 2. 同構定理}

一個群 $\mG$ 到一個群 $\overline{\mG}$ 上的映射，若 $\mG$ 與 $\overline{\mG}$ 應具有同一個算子域，則稱為算子同態，如果第一，它是通常意義下的同態（$ab\rightarrow \bar a\bar b$）；第二，每當 $a$ 映到 $\bar a$ 時，元素 $\Theta a$ 也映到 $\Theta\bar a$。\footnote{可以把算子同態的概念加以推廣，允許群 $\mG,\overline{\mG}$ 具有各自的算子域；此時同態不僅把 $\mG$ 映到 $\overline{\mG}$ 上，也把算子域彼此映到一起，其方式是：如果 $a$ 映到 $\bar a$ 且 $\Theta$ 映到 $\overline\Theta$，則 $\overline\Theta\bar a$ 是 $\Theta a$ 的像。在這種一般形式中，算子同態的概念也包含環同態的概念；參見 §8。} 記號為 $\mG\sim\overline{\mG}$。如果對應是一一的，就稱為算子同構。記號為 $\mG\simeq\overline{\mG}$。與普通群的情形一樣，可以證明同態定理：

\emph{如果 $\overline{\mG}$ 是 $\mG$ 的一個算子同態像，那麼 $\overline{\mG}$ 算子同構於一個商群 $\mG/\mN$；這裡 $\mN$ 是一個容許正規子群，由 $\mG$ 中所有對應於 $\overline{\mG}$ 的單位元的元素組成。反過來，每個容許正規子群 $\mN$ 定義一個商群 $\mG/\mN$，它容許 $\mG$ 的算子域，並且是 $\mG$ 的一個算子同態像。}

如果一個群除了自身和單位群以外沒有正規子群，就稱它為\emph{單的}。單群的每個同態或者是同構，或者把每個元素都對應到單位元。

\emph{第一同構定理.} \emph{設 $\overline{\mG}$ 是 $\mG$ 的同態像，$\overline{\mA}$ 是 $\overline{\mG}$ 的正規子群，而 $\mA$ 是 $\mG$ 中那些對應到 $\overline{\mA}$ 的元素的全體。則 $\mA$ 仍是正規子群，並且}
\[
        \overline{\mG}/\overline{\mA}\simeq \mG/\mA. \tag{1}
\]
\emph{證明.} $\mG\sim\overline{\mG}$，$\overline{\mG}\sim\overline{\mG}/\overline{\mA}$，因此 $\mG\sim\overline{\mG}/\overline{\mA}$。所以 $\overline{\mG}/\overline{\mA}$ 同構於 $\mG$ 中的一個商群；相應的正規子群正是那些在 $\overline{\mG}/\overline{\mA}$ 中對應於單位元的元素全體；這正是 $\mA$。

\emph{補充.} 若按照同態定理取 $\overline{\mG}=\mG/\mN$，則 $\mA$ 必定包含 $\mN$。由 $\mA$ 可以恢復 $\overline{\mA}$：$\overline{\mA}=\mA/\mN$。對應 $\mA\leftrightarrow\overline{\mA}$ 是一一的。公式 (1) 也可寫為
\[
        (\mG/\mN)/(\mA/\mN)\simeq \mG/\mA.
\]

\emph{第二同構定理.} \emph{設 $\mA$ 是子群，$\mB$ 是 $\mG$ 的正規子群。則 $\mA\cap\mB$ 是 $\mA$ 中的正規子群，並且}
\[
        \mA\mB/\mB\simeq \mA/\mA\cap\mB.
\]
\emph{證明.} 在同態 $\mG\sim\mG/\mB$ 下，特別地，$\mA$ 的元素 $a$ 被對應到某些剩餘類 $a\mB$，這些剩餘類合在一起構成群 $\mA\mB/\mB$。於是 $\mA\sim\mA\mB/\mB$，由同態定理即得斷言。

一個群 $\mG$ 到自身或到一個子群上的算子同態映射，稱為 $\mG$ 的自同態。如果 $\mG$ 特別是一個按加法書寫的阿貝爾群（帶算子），並且如上（§1，例 6）定義同態的和與積，那麼該群到自身的算子同態構成一個環，即自同構環。

\emph{一個單阿貝爾群（帶算子）的自同構環是一個除環。}

\emph{證明.} 每個同態或者把該群映到自身上，或者映到零群上（因為不存在其他容許子群）。那些並非把一切元素都映到零的同態，根據上面關於單群的說明，都是同構；而一個群到自身上的同構映射構成一個群。因此，在去掉零算子以後，該環成為一個群；所以該環本身是一個除環。

特別地，如果 $\mG$ 是一個左 $\mo$-模，並且把算子同態寫作右算子，那麼 $\mG$ 成為一個\emph{關於左側的 $\mo$ 和右側的自同構環的雙模}。因為新的算子 $\Gamma$ 被選為同態這一事實，正由公式
\[
        (a+b)\Gamma=a\Gamma+b\Gamma,
        \qquad ra\cdot\Gamma=r\cdot a\Gamma,
\]
表示；這些公式（連同前面已經解釋的其餘公式）刻畫了雙模。

反過來，如果給定一個雙模，那麼這些同一公式表示：右乘子域的每個元素都誘導一個關於左算子的算子同態，反之亦然。

\subsection*{§ 3. 合成列}

如果在 $\mG$ 中存在一列有限的容許子群
\[
        \mG\supset \mA_1\supset \mA_2\supset \ldots\supset \mA_r=\mE\tag{2}
\]
（$\mE$ 是隻由單位元 $e$ 組成的群），使得每個 $\mA_i$ 都是前一項中的正規子群，並且不可能在該列相鄰兩項之間再插入一個具有同樣性質的子群，那麼這個列稱為一條合成列。

因子群 $\mG/\mA_1,\ \mA_1/\mA_2,\ldots,\ \mA_{r-1}/\mE$ 稱為合成因子。它們是單群，也就是說，除了自身和單位群以外沒有容許正規子群。

如果在算子中尤其納入所有內自同構，那麼列 (2) 完全由正規子群組成，稱為主列。如果納入所有自同構，則稱為特徵列。在 §1 後面的例子中，這分別就是 $\mo$ 中理想的合成列，或者模與雙模的合成列。

\emph{Jordan--Hölder 定理.} \emph{如果對於一個群 $\mG$ 存在兩條不同的合成列：}
\[
        \mG\supset \mA\supset \mA_2\supset \ldots\supset \mA_r=\mE,
        \qquad
        \mG\supset \mB\supset \mB_2\supset \ldots\supset \mB_s=\mE,
\]
\emph{那麼它們具有相同的長度，$r=s$，並且合成因子}
\[
        \mG/\mA,\ \mA/\mA_2,\ldots,\mA_{r-1}/\mE
\]
\emph{在某種排列下同構於}
\[
        \mG/\mB,\ \mB/\mB_2,\ldots,\mB_{s-1}/\mE.
\]
證明可參見例如 E. Noether, Abstrakter Aufbau usw., \emph{Math. Annalen} 96 (1926), §10, p. 57。同時那裡還證明了：

\emph{如果在 $\mG$ 中存在一條合成列，那麼通過 $\mG$ 的每一個正規子群 $\mH$ 都可以作出一條合成列。}

合成列的存在對於有限群是顯然的；更一般地，對於滿足極大條件和極小條件的群也成立。

\emph{極大條件}是說，在任一子群集合中都有一個極大子群，即集合中沒有另一個子群再包含它。等價地說，每一條升鏈
$\mA_1\subset\mA_2\subset\mA_3\ldots$
在有限多項之後終止。

\emph{極小條件}是說，在任一子群集合中都有一個極小子群，即它不再包含該集合中的其他子群；等價地說，每一條降鏈
$\mA_1\supset\mA_2\supset\mA_3\ldots$
在有限多項之後終止。

\emph{如果只允許正規子群作為子群}（例如在阿貝爾群的情形），\emph{那麼反過來，極大條件和極小條件也從合成列的存在推出。}

\emph{證明.} 每個正規子群都有一條合成列，其長度稱為該正規子群的長度。若 $\mA\subset\mB$，則 $\mA$ 的長度小於 $\mB$ 的長度，因為可以通過 $\mA$ 作出 $\mB$ 的一條合成列。因此，每個長度最短的子群同時是極小的，每個長度最大的子群同時是極大的。

\subsection*{§ 4. 直積與交}

如果
\[
\begin{array}{ll}
1.& \mA,\mB\text{ 是 }\mG\text{ 中的正規子群},\\
2.& \mA\mB=\mG,\\
3.& \mA\cap\mB=\mE,
\end{array}
\]
則稱群 $\mG$ 為兩個因子的直積，$\mG=\mA\times\mB$。

\emph{這個定義顯然等價於：} $\mG$ 的每個元素 $g$ 都能唯一表示為 $g=ab$，其中 $a$ 屬於 $\mA$，$b$ 屬於 $\mB$，並且 $\mA$ 的元素與 $\mB$ 的元素可交換：$ab=ba$。

如果令
\[
\mB_i=\mA_1\cdots\mA_{i-1}\mA_{i+1}\cdots\mA_n,
\]
並且對每個 $i$ 都有 $\mG=\mA_i\times\mB_i$ 為直積，則稱群 $\mG$ 為 $n$ 個因子的直積，
\[
        \mG=\mA_1\times\cdots\times\mA_n.
\]

\emph{這個定義又等價於：} 每個 $\mG$ 中的 $g$ 都能唯一表示為
\[
        g=a_1a_2\ldots a_n,\qquad a_i\in\mA_i,
\]
並且 $\mA_i$ 的元素與 $\mA_k$ 的元素可交換。

以下事實常被使用：
\begin{enumerate}
\item 若 $\mA_1\times\cdots\times\mA_n=\mR$，且 $\mR\times\mH=\mG$，則 $\mA_1\times\cdots\times\mA_n\times\mH=\mG$。
\item 若 $\mG=\mA_1\times\ldots\times\mA_n=\mC_1\times\ldots\times\mC_n$，且 $\mC_i\subseteqq\mA_i$，則 $\mC_i=\mA_i$。（這由 $\mA_i$ 的元素藉助 $\mC_j$ 的表示推出。）
\item 若 $\mG=\mA\times\mB$，且 $\mR\supseteq\mA$，則 $\mR=\mA\times(\mR\cap\mB)$。（這由 $\mR$ 的元素表示為 $ab$ 推出，其中第二個因子同時屬於 $\mR$ 和 $\mB$。）
\item 若 $\mG=\mA\times\mB$，則 $\mA\simeq\mG/\mB$ 且 $\mB\simeq\mG/\mA$（第二同構定理）。由此還推出：
\item 若 $\mG=\mA\times\mB$，並且 $\mA$ 與 $\mB$ 分別有長度為 $m$ 和 $n$ 的合成列，則 $\mG$ 有一條長度為 $m+n$ 的合成列。
\item 若 $\mA\simeq\overline{\mA}$ 且 $\mB\simeq\overline{\mB}$，則 $\mA\times\mB\simeq\overline{\mA}\times\overline{\mB}$。
\end{enumerate}

如果一個群不能表示為若干個不等於 $\ne\mE$ 的因子的直積，則稱它為直接不可分解的。

顯然：\emph{每個滿足極小條件的群都是有限多個直接不可分解群的直積。}\footnote{如 W. Krull 在交換情形中以及 O. Schmidt 在一般情形中所證明的（參見注 6），在極大條件與極小條件的假設下，這種表示在同構意義下唯一。因為可以把所有內自同構加入算子域而不改變直接不可分解性的概念，所以要求正規子群的有限性條件，或要求主列存在，就已經足夠。}

如果把有關的群按加法書寫，那麼“積”和“直積”的概念就轉化為“和”和“直和”。記號為：群的和寫作 $(\mA_1,\mA_2,\ldots)$，直和寫作 $\mA_1+\mA_2+\ldots$。

“直接交”的概念雖然後文不再使用，但是下面關於直接交與直積之間關係的定理說明，後面幾章的理想理論也可以用交而不是和來表述，由此恢復同熟悉的交換理論之間的聯繫。

如果
\[
\begin{array}{ll}
1.& \mR,\mS\text{ 是 }\mG\text{ 中的正規子群},\\
2.& \mR\cap\mS=\mD,\\
3.& \mR\mS=\mG,
\end{array}
\]
則稱群 $\mD$ 為群 $\mR$ 與 $\mS$ 關於 $\mG$ 的直接交。

同樣地，若令
\[
\mR_i=\mS_1\cap\cdots\cap\mS_{i-1}\cap\mS_{i+1}\cap\cdots\cap\mS_n,
\]
並且對每個 $i$，交 $\mD=\mR_i\cap\mS_i$ 都是直接的，則稱 $\mD$ 為 $n$ 個群 $\mS_1,\ldots,\mS_n$ 的直接交。

由定義可知，$\mD$ 必須是 $\mG$ 中的正規子群。若 $\mD=\mS_1\cap\ldots\cap\mS_n$ 是直接的，並且在同態 $\mG\sim\mG/\mD$ 下，群 $\mG,\mS,\mD$ 分別過渡到 $\overline{\mG},\overline{\mS},\overline{\mE}$，則 $\overline{\mE}=\overline{\mS}_1\cap\ldots\cap\overline{\mS}_n$ 成為直接交。反過來，從每個這樣的表示 $\overline{\mE}=\overline{\mS}_1\cap\ldots\cap\overline{\mS}_n$，也可返回到一個直接表示 $\mD=\mS_1\cap\ldots\cap\mS_n$。藉助這種對應，關於任意 $\mD$ 的直接交定理被歸約為特殊情形 $\mD=\mE$。

在 $\mD=\mE$ 且只有兩個因子的情形中，直積與直接交的條件相同。對於 $n$ 個因子，直積與交之間有如下的一一關係：

\emph{I. 若 $\mG=\mA_1\times\ldots\times\mA_n=\mA_i\times\mB_i$，則 $\mE=\mB_1\cap\ldots\cap\mB_n=\mB_i\cap\mC_i$ 是直接交，並且 $\mC_i=\mA_i$。}

\emph{II. 若 $\mE=\mS_1\cap\ldots\cap\mS_n=\mR_i\cap\mS_i$ 是直接交，則 $\mG=\mR_1\times\ldots\times\mR_n=\mR_i\times\mT_i$，並且 $\mT_i=\mS_i$。}

\emph{I 的證明.} 設 $\mB=\mB_1\cap\ldots\cap\mB_n=\mB_i\cap\mC_i$，則 $\mC_i\supseteqq\mA_i$；我們證明 $\mC_i\subseteqq\mA_i$。例如設 $c$ 屬於 $\mC_1$，則 $c$ 屬於 $\mB_2,\ldots,\mB_n$，所以關於 $\mA$ 的分量表示給出
\[
        c=a_1a_2\ldots a_n=a_1ea_3\ldots a_n=a_1a_2e\ldots a_n=\ldots=a_1\ldots a_{n-1}e.
\]
由於所有 $\mA_i$ 的積是直積，這些表示彼此相同；除第一位外，每一位都曾出現一次 $e$，於是 $c=a_1$，或 $\mC_1\subseteqq\mA_1$，即 $\mC_1=\mA_1$，相應地 $\mC_i=\mA_i$。由此得 $\mB=\mB_i\cap\mC_i=\mB_i\cap\mA_i=\mE$；且 $\mB_i\mC_i=\mB_i\mA_i=\mG$，所以交是直接的。

\emph{II 的證明.} 設 $\mH=\mR_1\ldots\mR_n=\mR_i\mT_i$，則 $\mT_i\subseteqq\mS_i$；我們證明 $\mS_i\subseteqq\mT_i$。例如設 $s$ 屬於 $\mS_1$。藉助 $\mG=\mS_i\times\mR_i$，有
\[
        s=s_1e=s_2r_2=\cdots=s_nr_n.
\]
取 $t_1=er_2\ldots r_n=t_ir_i$，並注意 $\mR_i$ 與 $\mS_i$ 的元素逐個可交換，則得到
\[
        st_1^{-1}=s_it_i^{-1},
\]
因而這是所有 $\mS_i$ 的元素，所以等於 $e$。因此 $\mS_1\subseteqq\mT_1$，或 $\mT_1=\mS_1$，相應地 $\mT_i=\mS_i$。由此得 $\mH=\mR_i\mT_i=\mR_i\mS_i=\mG$，並且 $\mR_i\cap\mT_i=\mR_i\cap\mS_i=\mE$；所以 $\mG$ 是 $\mR_i$ 的直積。

\clearpage

\section*{34. 超複量與表示論}
\emph{續：第一章，§§5--7，以及第二章，§§8--14.}

\subsection*{§ 5. 完全可約群}

如果一個群是有限多個單群的直積，則稱它為\emph{完全可約}：
\[
        \mG=\mA_1\times\ldots\times\mA_n.
\]
在這種情形中，群
\[
        \mA_1\times\ldots\times\mA_n
        \supset \mA_1\times\ldots\times\mA_{n-1}
        \supset\ldots\supset \mA_1\supset \mE
\]
構成一條合成列。其長度為 $n$，合成因子為
\[
        \sim \mA_n,\ \mA_{n-1},\ldots,\mA_1.
\]

\emph{若 $\mG$ 完全可約，則每個正規子群都是直因子，而另一個因子可以選為那些出現在 $\mG$ 的一個給定直積分解中的單群的乘積。正規子群 $\mH$ 本身也完全可約。}

\emph{證明.} 有
\begin{equation}
        \mG=\mH\mA_1\mA_2\ldots\mA_n.\tag{3}
\end{equation}
現在令 $\mH_i=\mH\mA_1\mA_2\ldots\mA_i$，則 $\mH_{i+1}=\mH_i\mA_{i+1}$。或者 $\mA_{i+1}\subseteqq\mH_i$，因而 $\mH_{i+1}=\mH_i$；這時我們在 (3) 中刪去因子 $\mA_{i+1}$。或者 $\mH_i\cap\mA_{i+1}$ 是 $\mA_{i+1}$ 的一個真正規子群，因而等於 $=\mE$，於是 $\mH_i\times\mA_{i+1}$ 為直積。刪去多餘因子以後，(3) 變為
\[
        \mG=\mH\times\mA_{i_1}\times\ldots\times\mA_{i_r},
\]
所以 $\mH$ 是直因子。進一步可得
\[
        \mH\simeq \mG/\mA_{i_1}\times\ldots\times\mA_{i_r}
        \simeq \mA_{i_1}\times\ldots\times\mA_{j_\mu},
\]
其中 $\mA_{j_1},\ldots,\mA_{j_\mu}$ 是其餘的 $\mA_i$；因此 $\mH$ 完全可約。

\emph{推論.} \emph{完全可約群的每一個分解為直接不可分解因子的分解，都是分解為單因子的分解，因而給出一條合成列；由此推出這些因子在同構意義下唯一確定。}

\subsection*{§ 6. 關於域的模。超複系統}

關於一個（不必交換的）域而具有有限基的模，給出帶算子的完全可約群的一個幾乎平凡的例子。

設 $\mG$ 是一個右 $\mathsf{K}$-模，並且 $K$ 的單位元 $e$ 同時也是單位算子：對於 $\mG$ 中的 $a$ 有 $ae=a$。由一個 $a$ 生成的每個子模 $a\mathsf{K}$ 都是單的；也就是說，它等於由其中任一非零元素生成的模。因此，若 $\mH$ 是由某些元素生成的子模，而 $a$ 是一個不屬於 $\mH$ 的元素，則 $\mH\cap a\mathsf{K}=\mE$，所以 $\mH+a\mathsf{K}$ 是直和。這樣，從任一基元素 $a_1$ 出發，可以不斷添入新的基元素 $a_2,a_3,\ldots$，並得到 $\mG$ 作為直和
\[
        \mG=a_1K+a_2K+\cdots+a_nK.
\]
\emph{所以 $\mG$ 是完全可約的。}

數 $n$，也就是合成列的長度，稱為 $\mG$ 關於 $\mathsf{K}$ 的\emph{秩}。由於 $\mG$ 中元素表示的唯一性（並且由於已假定 $e$ 為單位算子），基 $a_1,\ldots,a_n$ 是線性無關的。

每個子模 $\mH$ 都是直和項：
\[
        \mG=\mH+\mR=(h_1\mathsf{K}+\ldots+h_s\mathsf{K})+(r_1\mathsf{K}+\ldots+r_{n-s}\mathsf{K}),
\]
或者說：\emph{$\mH$ 的每個線性無關基都可以擴充為 $\mG$ 的一個線性無關基。} 這些 $r_i$ 甚至可以從原來的基元素 $a_i$ 中選取。

設 $(a_1,\ldots,a_n)$ 與 $(c_1,\ldots,c_n)$ 是線性無關基。於是
\[
        c_k=\sum a_i\pi_{ik},
\]
或者寫成矩陣形式：
\[
        (c_1,\ldots,c_n)=(a_1,\ldots,a_n)P;
        \qquad
        P=\begin{pmatrix}
        \pi_{11}&\ldots&\pi_{1n}\\
        \vdots&&\vdots\\
        \pi_{n1}&\ldots&\pi_{nn}
        \end{pmatrix}.
\]
反過來，
\[
        (a_1,\ldots,a_n)=(c_1,\ldots,c_n)Q.
\]
由此可得
\[
        (a_1,\ldots,a_n)=(a_1,\ldots,a_n)PQ,
        \qquad
        (c_1,\ldots,c_n)=(c_1,\ldots,c_n)QP,
\]
所以，由於 $a_i$ 與 $c_i$ 都線性無關，
\[
        PQ=QP=E.
\]

同時又是環的特殊 $\mathsf{K}$-模，就是“超複系統”。

一個環 $\mo$ 若同時是關於一個交換域 $\mathsf{K}$ 的右模，則稱為\emph{關於 $\mathsf{K}$ 的超複系統}（文獻中也稱為“$\mathsf{K}$ 上的代數”），如果：
\begin{enumerate}
\item 它的秩是有限的（設 $u_1,\ldots,u_n$ 為一個線性無關基）；
\item $ab\cdot\varkappa=a\cdot b\varkappa=a\varkappa\cdot b$。這可表述為：$\mathsf{K}$ 與 $\mo$ 以可交換方式相聯結；
\item $K$ 的單位元 $e$ 同時也是單位算子：對於 $\mo$ 中的 $a$ 有 $ae=a$。
\end{enumerate}
由於
\[
        \Bigl(\sum_i u_i\alpha_i\Bigr)\Bigl(\sum_k u_k\beta_k\Bigr)
        =\sum_{i,k}(u_i u_k)(\alpha_i\beta_k),
\]
只要除 $\mathsf{K}$ 外還給出\emph{乘法表}，即給出每個乘積 $u_i u_k$ 怎樣通過 $u_j$ 表示，
\[
        u_i u_k=\sum_j u_j\gamma^j_{ik},
\]
這個系統就被唯一確定。$\gamma^j_{ik}$ 必須滿足由結合律推出的熟知關係。

如果 $\mo$ 像我們以後常常假定的那樣，具有一個單位元 $e$（對所有 $a$ 有 $ea=ae=a$），那麼可以把 $K$ 中的元素 $x$ 與 $e\varkappa$ 等同，從而把 $\mathsf{K}$ 看作 $\mo$ 的一個子域。於是這個超複系統也可以描述為：\emph{關於一個位於中心中的域而具有有限秩的環。}

超複系統的一個例子是一個有限群的\emph{群環}：取該有限群的元素作為基元素 $u_i$，並取群乘法作為乘法。這裡的域 $\mathsf{K}$ 是任意的。

\footnotetext{條件 2 說明，當把 $\mo$ 同時看作 $\mo$-左模和 $\mo$-右模時，用 $K$ 中一個元素相乘意味着 $\mo$ 的一個同態。藉助條件 3，這些同態甚至是同構。}

\subsection*{§ 7. 矩陣}

方陣 $(\pi_{ik})$，其中 $\pi_{ik}$ 屬於環 $\mo$，在通常的矩陣乘法
$\sum \pi_{ij}\varrho_{jk}=\sigma_{ik}$ 和加法
$\pi_{ik}+\varrho_{ik}=\sigma_{ik}$ 下構成一個環。

一個以域 $\mathsf{K}$ 中元素為元素的矩陣，如果它同時有右逆和左逆，就稱為\emph{正則的}。

我們在 §6 已見：\emph{一個 $\mathsf{K}$-模的兩個線性無關基之間的過渡矩陣是正則的。}

\emph{判定正則性時，有一個右逆就足夠。} 的確，取一個具有線性無關基
$(a_1,\ldots,a_n)$ 的右 $K$-模（由未定元 $a_1,\ldots a_n$ 的線性形式構成的模），並令
\[
        (c_1,\ldots,c_n)=(a_1,\ldots,a_n)P.
\]
於是
\[
        (c_1,\ldots,c_n)P^{-1}=(a_1,\ldots,a_n)PP^{-1}=(a_1,\ldots,a_n),
\]
所以 $(c_1,\ldots,c_n)$ 又是一個基，因而仍然線性無關，於是矩陣 $P$ 是正則的。
此外，右逆等於左逆。

同樣地，有一個左逆也足夠。

\emph{如果 $P$ 有左逆，那麼 $P$ 不是左零因子。}
\emph{證明。} 從 $PA=0$ 得出 $A=P^{-1}PA=0$。

因此，$P$ 也只能有一個右逆；按 §6，此右逆與左逆重合。

按照通常記法，令 $P_\varkappa$ 表示矩陣 $(\pi_{ik}\varkappa)$，令 $C_{ik}$ 表示在 $ik$
位置為單位元、其餘位置全為零的矩陣，則
\[
        (\pi_{ik})=\sum C_{ik}\pi_{ik};
\]
因此，$\mathsf{K}$ 上的矩陣構成一個秩為 $n^2$ 的 $K$-模。$C_{ik}$ 滿足關係
\[
        C_{ij}C_{jk}=C_{ik},
        \qquad
        C_{ij}C_{\bar j k}=0,\quad j\ne\bar j.
\]
於是，$\mathsf{K}$ 上的矩陣構成一個有限 $\mathsf{K}$-模；如果 $\mathsf{K}$ 是交換的，則構成一個超複系統。

\section*{第二章. 非交換理想論}

\subsection*{§ 8. 環的同態定理}

若 $\ma$ 是 $\mo$ 中的雙邊理想，則剩餘類環 $\mo/\ma$ 不僅是一個
$\mo$-模，而且也顯然是一個環。$\mo$ 的每一個同態像仍是一個環，並且同構於一個商環
$\mo/\ma$，其中 $\ma$ 是雙邊理想。證明與群的情形相同。

特別地，如果 $\mo$ 是一個除環，則除了 $(0)$ 和 $\mo$ 以外沒有其他理想；
因此這裡的同態或者是同構，或者把每個元素都映到零。

再特別地，設環 $\mo$ 是一個右 $\mathsf{K}$-模，其中 $\mathsf{K}$ 應當是逐元素與 $\mo$ 可交換的環：
\[
        ab\cdot x=a\cdot bx=ax\cdot b
\]
（例如，關於交換域 $\mathsf{K}$ 的超複系統就是這種情形）。這時，作為容許的右理想、
左理想或雙邊理想，只考慮那些同時也是 $\mathsf{K}$-模的理想；並且除環同態外，還要求關於
$\mathsf{K}$ 的乘法的算子同態。容許理想由此成為關於 $\mo$ 與 $\mathsf{K}$ 的雙模。
（對於超複系統，若不加說明而說“理想”，總是隻指關於作為基礎的係數域 $\mathsf{K}$ 而容許的理想。）
在這種擴大的意義下，上述同態定理仍然成立。

環同態的一個例子，是從一個乘子域 $\mo$ 過渡到絕對乘子域（§1，例 4）。因此，絕對乘子域
同構於商環 $\mo/\md$，其中 $\md$ 是 $\mo$ 中那些零化 $\mM$ 的元素所組成的雙邊理想。

由同態定理的成立，可以像 §2 中那樣推出環同構和算子同構的第一、第二同構定理。

$\mo$ 的子集的所有乘積 $\mA\mB$ 都應理解為模乘積：$\mA\mB$ 是所有和式
$\sum ab$ 的全體，其中 $a\in\mA$、$b\in\mB$；而 $a\mB$ 是所有和式
$\sum ab$ 的全體，其中 $b\in\mB$。如果 $\mB$ 是關於某個作為算子域的環的右模，
那麼乘積 $\mA\mB$，相應地 $a\mB$，也仍是右模。（特別地：右理想、關於係數域
$\mathsf{K}$ 的容許理想，等等。）若用 $(\mA,\mB)$ 表示模 $\mA,\mB$ 的和，則有計算規則
\[
        \mA\mB\cdot\mC=\mA\cdot\mB\mC,
\]
\[
        \mA(\mB,\mC)=(\mA\mB,\mA\mC),
        \qquad
        (\mA,\mB)\mC=(\mA\mC,\mB\mC).
\]
直和記作 $\mA+\mB$（§4）。

以下我們常常考察那些對右理想滿足極大條件或極小條件（§3）的環。超複系統尤其屬於這一類；
因為按照剛才的約定，只有 $\mathsf{K}$-模才被容許作為理想，而它們的秩因而保持有界（§6）。

\subsection*{§ 9. 冪等元素. 右理想的直和分解}

設 $\mo$ 是一個有單位元的環。

若 $\mr$ 是右理想，則 $\mr\mo\subseteqq\mr$；由於有單位元，甚至有 $\mr\mo=\mr$。
相應地，對於左理想有 $\mo\ml=\ml$。

元素 $c$ 稱為冪等的，如果 $c^2=c$。

\emph{若 $\mo=\mr_1+\ldots+\mr_n$ 是一個分解為右理想的直和分解，並且}
\[
        e=e_1+\ldots+e_n,
        \qquad e_i\in\mr_i,
\]
\emph{則有}
\[
        e_i^2=e_i,
        \qquad e_i e_k=0\quad(i\ne k)
        \quad\hbox{（正交关系）},
\]
\[
        \mr_i=e_i\mo.
\]

\emph{證明.} 設 $r$ 是 $\mr_1$ 的一個元素。於是
\[
        r=er=e_1r+\ldots+e_nr,
        \qquad e_ir\in\mr_i,
\]
另一方面也有
\[
        r=r+0+\ldots+0,
\]
所以
\[
        e_1r=r,
        \qquad e_ir=0\quad(i\ne1).
\]
由第一個公式推出 $\mr_1\subseteqq e_1\mo$；另一方面 $e_1\mo\subseteqq\mr_1$，因此 $e_1\mo=\mr_1$。
令 $r=e_1$ 特化，得到
\[
        e_1^2=e_1,
        \qquad e_i e_1=0\quad(i\ne1).
\]
把指標 $1$ 換成 $k$ 時，同樣成立。

\emph{反過來，若 $e=\sum e_i$，$e_i e_k=0$ $(i\ne k)$，$e_i^2=e_i$，並且置 $\mr_i=e_i\mo$，則 $\mo=\mr_1+\ldots+\mr_n$。}

\emph{證明.} $\mo$ 的每個元素 $r$ 都有
\[
        r=er=e_1r+\ldots+e_nr,
\]
因而
\[
        \mo=(e_1\mo,\ldots,e_n\mo).
\]
但是這種表示是唯一的；因為由
\[
        0=e_1a_1+\ldots+e_na_n
\]
用 $e_i$ 左乘得到
\[
        e_i^2a_i=e_i a_i=0.
\]

由一個右側表示
\[
        \mo=\mr_1+\ldots+\mr_n=e_1\mo+\ldots+e_n\mo
\]
於是得到一個左側表示
\[
        \mo=\ml_1+\ldots+\ml_n=\mo e_1+\ldots+\mo e_n.
\]

若 $\mr_i$ 是直接不可分解的，則 $\ml_i$ 也直接不可分解；因為例如 $\ml_1$ 的一個分解將意味着
\[
        \ml_1=\mo e_1=\mo e_1'+\mo e_1'',
        \qquad
        \mo=\mo e_1'+\mo e_1''+\mo e_2+\ldots+\mo e_n.
\]
由此得到右側分解
\[
        \mo=e_1'\mo+e_1''\mo+e_2\mo+\ldots+e_n\mo
        =\mr_1'+\mr_1''+\mr_2+\ldots+\mr_n,
\]
\[
        \mr_1\cong \mo/(\mr_2+\ldots+\mr_n)=\mr_1'+\mr_2',
\]
於是 $\mr_1$ 將是可分解的。

\footnotetext{另有一個逆命題：若 $e_i e_k=0$，$e_i^2=e_i$，$c=\sum e_i$，則 $c$ 是環 $e_1\mo+\ldots+e_n\mo$ 的左單位元。該和為直和這一點，完全像上面一樣看出。}

\subsection*{§ 10. 分解為雙邊理想}

再設 $\mo$ 是一個有單位元的環，並且
\begin{equation}
        \mo=\ma_1+\ldots+\ma_n\tag{1}
\end{equation}
是把 $\mo$ 分解為雙邊理想的一個分解，則正交關係對於這些理想本身也成立：
\begin{equation}
        \ma_i\ma_k=0\quad(i\ne k),
        \qquad
        \ma_i^2=\ma_i.                                             \tag{2}
\end{equation}

\emph{證明.} 令 $\mb_1=\ma_2+\ldots+\ma_n$，$\mo=\ma_1+\mb_1$；於是 $\ma_1\mb_1\subseteqq \ma_1\cap\mb_1=0$，故 $\ma_1\mb_1=0$，因此更有
\[
        \ma_1\ma_i=0\quad(i\ne1);
\]
\[
        \ma_1=\ma_1\mo=\ma_1(\ma_1+\mb_1)=(\ma_1^2,\ma_1\mb_1)=\ma_1^2.
\]

反過來：\emph{由 (1) 和 (2) 可知，模 $\ma_i$ 是雙邊理想。}
\emph{證明.}
\[
        \mo\ma_i=(\ma_1+\cdots+\ma_n)\ma_i=\ma_i^2=\ma_i,
        \qquad
        \ma_i\mo=\ma_i.
\]

\emph{環 $\ma_i$ 中的右理想也是 $\mo$ 中的右理想。}
\emph{證明.}
\[
        \mr\mo=\mr(\ma_i+\mb_i)=(\mr\ma_i,\mr\mb_i)=(\mr,0)=\mr.
\]

\emph{若 $\mr$ 是 $\mo$ 中的右理想，則 $\mr$ 是右理想 $\mr\ma_i\subseteqq\ma_i$ 的直和。}
\emph{證明.} $\mr=\mr\mo=(\mr\ma_1,\ldots,\mr\ma_n)$，並且 $\mr\ma_i\mo=\mr\ma_i$，故 $\mr\ma_i$ 是包含在 $\mr$ 中的右理想。再者 $\mr\ma_i\subseteqq\ma_i$，所以該和是直和：
\[
        \mr=\mr\ma_1+\ldots+\mr\ma_n.
\]

\emph{特別地，若 $\mr$ 是直接不可分解的，則 $\mr$ 只能有一個分量；因而 $\mr\subseteqq$ 必須包含在某個 $\ma_i$ 中。}

根據這些定理，一旦知道各個 $\ma_i$ 的理想理論，也就掌握了 $\mo$ 中的理想理論。因此我們通常限於研究雙邊不可分解的環。

\emph{包含在不同 $\ma_i$ 中的兩個右理想絕不可能算子同構。}
\emph{證明.} 包含在 $\ma_1$ 中的理想被所有其他 $\ma_k$ 零化，但不被 $\ma_1$ 零化。相反，包含在 $\ma_2$ 中的理想被 $\ma_1$ 零化。

所以，如果能夠把環 $\mo$ 分解為雙邊不可分解理想
\[
        \ma_1+\ldots+\ma_n
\]
而每個 $\ma_i$ 又分解為單邊不可分解右理想
\[
        \mr_i'+\mr_i''+\cdots,
\]
那麼要尋找彼此算子同構的 $\mr$，只能在屬於同一個 $\ma_i$ 的那些右理想中尋找。但是也可能在同一個 $\ma_i$ 中還存在不同的 $\mr$ 的類，也就是說存在彼此並非算子同構的類，如下面的例子所示。

設 $\mathsf{K}$ 是有理數域，
\[
        \mo=e_1\mathsf{K}+e_2\mathsf{K}+u\mathsf{K}
\]
是一個超複系統，其中 $e_i$ 和 $u$ 都同 $\mathsf{K}$ 中的數可交換，並且乘法表為
\[
\begin{array}{c|ccc}
     & e_1&e_2&u\\ \hline
 e_1 & e_1&0&u\\
 e_2 & 0&e_2&0\\
 u   & 0&u&0
\end{array}
\]
單位元是 $e=e_1+e_2$。這些 $e_i$ 滿足正交關係。因此右側分解為
\[
        \mo=e_1\mo+e_2\mo=(e_1,u)+(e_2),
\]
左側分解為
\[
        \mo=\mo e_1+\mo e_2=(e_1)+(e_2,u).
\]
由於 $e_2\mo$ 和 $\mo e_1$ 顯然不可分解，$\mo e_2$ 和 $e_1\mo$ 也必然不可分解。

雙邊分解是不可能的；因為那樣一個分量必須包含 $e_1\mo$，另一個分量必須包含 $e_2\mo$；於是第一個分量會包含 $u$，但第二個分量也會包含 $ue_2=u$。

但是 $e_1\mo$ 和 $e_2\mo$ 並不算子同構，因為它們關於 $\mathsf{K}$ 的秩不同。

\emph{設 $\mo=\ma_1+\ldots+\ma_n$，並且 $\ma_i$ 是雙邊不可分解的。那麼它們被唯一地確定。}

\emph{證明.} 又設 $\mo=\mc_1+\ldots+\mc_n$，於是
\[
        \ma_i=\ma_i\mo=(\ma_i\mc_1,\ldots,\ma_i\mc_n).
\]
後一和是直和，因為 $\ma_i\mc_k\subseteqq\mc_k$ 且 $\ma_i\mc_k\subseteqq\ma_i$；然而由於 $\ma_i$ 不可分解，所有 $\ma_i\mc_k$ 都必須為零，只有一個例外，即 $\ma_i\mc_{j_i}$。因此
\[
        \ma_i=\ma_i\mc_{j_i}.
\]
同樣有
\[
        \mc_{j_i}=\mo\mc_{j_i}=\ma_1\mc_{j_i}+\cdots+\ma_n\mc_{j_i},
\]
且 $\ma_i\mc_{j_i}\ne0$，所以其餘項全為 $=0$，
\[
        \mc_{j_i}=\ma_i\mc_{j_i}=\ma_i,
\]
證畢。

\subsection*{§ 11. 中心}

環 $\mo$ 的中心 $\mZ$ 是所有同每個環元素可交換的 $z$ 的全體（即對所有 $a$ 都有 $za=az$）。$\mZ$ 是一個交換環。

$\mo$ 中的每個單邊理想 $\ma$ 在 $\mZ$ 中都有一個“收縮理想” $\ma\cap\mZ$。因為 $\ma$ 和 $\mZ$ 都是 $\mZ$-模，所以 $\ma\cap\mZ$ 也是。

$\mZ$ 中的每個理想 $\mA$ 都有一個擴張理想：即 $\mA$ 在 $\mo$ 中生成的理想
\[
        \ma=(\mA,\mo\mA).
\]
這是雙邊理想，因為
\[
        \ma\mo=(\mo\mA,\mA)\mo=(\mo\mA\mo,\mA\mo)
        = (\mo^2\mA,\mo\mA)\subseteqq \mo\mA\subseteqq\ma.
\]
若假定 $\mo$ 有單位元，則可以寫成 $\ma=\mo\mA$。在這種情況下，有下列定理：

\emph{$\mo$ 的每個雙邊分解都同中心的一個分解一一對應：由}
\begin{equation}
        \mo=\ma_1+\ldots+\ma_n,
        \qquad \mA_i=\ma_i\cap\mZ                         \tag{1}
\end{equation}
\emph{推出}
\begin{equation}
        \mZ=\mA_1+\ldots+\mA_n,
        \qquad \mo\mA_i=\ma_i,                              \tag{2}
\end{equation}
\emph{反之亦然。}

\emph{證明.} 設 $z$ 是中心元素，並且
\begin{equation}
        z=z_1+\ldots+z_n\tag{3}
\end{equation}
是它按 (1) 給出的分解。於是 $z_i$ 也是中心元素；因為
\[
        za=z_1a+\ldots+z_na=az=az_1+\ldots+az_n;
\]
由和的直和性，並且由於 $\ma_i$ 是雙邊的，得到
\[
        z_i a=az_i.
\]
因此 $z_i$ 屬於 $\ma_i\cap\mZ=\mA_i$；所以 (3) 表示的正是所斷言的和分解。特別地，$e=\sum e_i$，所以 $e_i$ 都是中心元素。最後有
\[
        z=ze_1+\ldots+ze_n,
\]
從而
\[
        \mA_i=\mZ e_i,
        \qquad
        \mo\mA_i=\mo\mZ e_i=\mo e_i=\ma_i.
\]

反過來，若 $\mZ=\mA_1+\ldots+\mA_n$ 是中心的一個分解，則由 §9 得到
\[
        \mA_i=\mZ e_i,
        \qquad e_i^2=e_i,
        \qquad e_i e_k=0\quad(i\ne k).
\]
因此，藉助正交關係 \hbox{\(\S 9\)}，有
\[
\begin{aligned}
        \mo&=\mo e=\mo e_1+\ldots+\mo e_n
        =\mo\mZ e_1+\ldots+\mo\mZ e_n\\
        &=\mo\mA_1+\ldots+\mo\mA_n
        =\ma_1+\ldots+\ma_n .
\end{aligned}
\]
現令 $\ma_i\cap\mZ=\mA_1'$，則由斷言的第一部分知道：
\[
        \mZ=\mA_1'+\ldots+\mA_n'\quad\hbox{为直和}.
\]
但 $\mA_i'\supseteqq\mA_i$，所以 $\mA_i'=\mA_i$（§4，2）。

\emph{推論.} $\ma_i$ 和 $\mA_i$ 同時是雙邊直接可分解的，或同時不是。

\subsection*{§ 12. 冪零理想}

理想 $\mc$ 稱為\emph{冪零的}，如果 $\mc^\varrho=0$。

\emph{例.} §10 中的理想 $(u)$。

\emph{如果在 $\mo$ 中存在一個冪零右理想 $\mr$，那麼也存在一個冪零左理想，甚至雙邊理想。}

\emph{證明.} 設 $\mr^\varrho=0$。那麼 $\mo\mr$；或者，如果 $\mo$ 不含單位元並且 $\mo\mr$ 可能為零，則取 $(\mo\mr,\mr)$；這是一個左理想，並且
\[
        (\mo\mr)^\varrho=\mo\mr\mo\mr\ldots\mo\mr
        \subseteqq\mo\mr\mr\ldots\mr=\mo\mr^\varrho=0.
\]

\emph{兩個冪零右理想的和仍然是一個冪零右理想。}

\emph{證明.} 設 $\mc^\varrho=0$，$\md^\sigma=0$。在
\[
        (\mc,\md)^{\varrho+\sigma-1}
        =(\ldots,\md\mc\ldots\md\ldots\mc,\ldots)
\]
中，每一項要麼至少含有 $\varrho$ 個因子 $\mc$，要麼至少含有 $\sigma$ 個因子 $\md$。在第一種情形下有
\[
        \md\mc\cdots\md\cdots\mc\cdots
        \subseteq\md\mc\mc\cdots\mc=\md\mc^\varrho=0;
\]
在後一種情形下，對因子 $\md$ 同理。因此
\[
        (\mc,\md)^{\varrho+\sigma-1}=0.
\]

現在若 $\mo$ 中的右理想滿足極大條件，則存在一個極大的冪零右理想。它包含所有其他冪零右理想；否則可以同這樣的一個理想取和。設它為 $\mc$。由於 $\mo\mc$ 也是冪零右理想，所以 $\mo\mc\subseteqq\mc$，因而 $\mc$ 是左理想。每個冪零左理想 $\md$ 也包含在 $\mc$ 中；因為 $(\md,\md\mo)$ 是冪零右理想，所以 $\md\subseteqq\mc$。 \emph{於是存在一個極大的（雙邊）冪零理想 $\mc$，或稱“根基”，它包含所有其他右側和左側的冪零理想。}

在 §10 的例子中，$(u)$ 是極大的冪零理想。

如果根基是零理想，則稱之為一個“無根基環”（“半單環”，“Dedekind 系統”）。按根基取模所得的剩餘類環總是一個無根基環。

\emph{如果 $\mZ$ 是 $\mo$ 的中心，$\mc$ 是 $\mo$ 的根基，那麼 $\mC=\mc\cap\mZ$ 是 $\mZ$ 的根基。}

\emph{證明.} 顯然 $\mC$ 是冪零的。若在 $\mZ$ 中還存在一個更大的冪零理想 $\overline{\mC}$，那麼它可以擴張成一個不全包含在 $\mc$ 中的冪零理想 $\overline{\mc}$，因為
\[
        (\overline{\mC}\mo)^\varrho=\overline{\mC}^{\varrho}\mo^\varrho=0.
\]

\emph{推論.} 如果 $\mo$ 是無根基環，那麼 $\mZ$ 也是。反過來並不成立，§10 中的例子即說明這一點。$\mo/\mc$ 的中心包含一個同構於 $\mZ/\mC$ 的環；但它可以是真子環（§10 中的例子）。

\subsection*{§ 13. 完全可約環}

按 §5，如果一個環是有限多個單右理想的直和，則稱它為右完全可約的。

\emph{一個帶單位元的右完全可約環沒有 $\ne(0)$ 的冪零理想，因而沒有根基。}

\emph{證明.} 每個右理想 $\mr$ 都是直和項（§5）：
\[
        \mo=\mt+\mr=e_1\mo+e_2\mo.
\]
由於 $e_2^2=e_2$，也有 $e_2^\varrho=e_2$。如果 $\mr$ 是冪零的，則 $e_2^\varrho=0$，從而 $e_2=0$，於是 $\mr=0$，證畢。

現在證明兩個逆命題。第一：\emph{滿足右理想極小條件的無根基環，關於右理想是完全可約的。}

\emph{證明.} 設 $\mt$ 是一個 $\ne0$ 的極小右理想。我們要證明，$\mt$ 含有一個冪等元素。

因為 $\mt^2\subseteqq\mt$ 且 $\mt^2\ne0$，所以 $\mt^2=\mt$；因為 $\mt^2$ 是右理想。因此存在 $\mt$ 中的一個 $a$，使得 $a\mt\ne0$。於是必有 $a\mt=\mt$，因為 $a\mt$ 是右理想。

$\mt$ 中所有被 $a$ 零化的 $b$（即 $ab=0$）的全體是一個右理想，並且不等於 $\mt$，所以等於 $0$。因此，由 $ab=0$ 可推出 $b=0$。

由於 $\mt=a\mt$，$a$ 必須能表示成 $ac$ 的形式：$a=ac$，且 $c\ne0$。於是得到
\[
        ac=ac^2,
        \qquad a(c-c^2)=0,
        \qquad c-c^2=0,
        \qquad c^2=c.
\]
接着證明 $\mt$ 是直和項。$c\mo$ 是一個包含於 $\mt$ 的非零右理想，因為 $c^2=c$ 屬於 $c\mo$；所以 $c\mo=\mt$。$\mo$ 中每個元素 $r$ 都有一個表示：
\[
        r=cr+(r-cr) \qquad \hbox{（单侧 Peirce 分解）}.
\]
元素 $cr$ 構成理想 $\mt$；元素 $r-cr$ 構成另一個右理想 $\mr$，它被 $c$ 零化：$c\mr=0$。因此
\[
        \mo=(\mt,\mr).
\]
由於 $\mt$ 的元素不被 $c$ 零化，這個和是直和：
\begin{equation}
        \mo=\mt+\mr.                                             \tag{1}
\end{equation}

現在按照極小條件，把 $\mo$ 表示為直接不可分解理想的直和：
\[
        \mo=\mr_1+\ldots+\mr_n.
\]
這些 $\mr_i$ 必須是單的（即極小的）。因為如果例如 $\mr_1$ 不是單的，而 $\mt$ 是 $\mr_1$ 中的一個極小理想，則由 (1) 得到
\[
        \mr_1=\mt+\mr\cap\mr_1 \qquad \hbox{\(\S 4,3\)};
\]
於是 $\mr_1$ 可分解，矛盾。因此 $\mo$ 是完全可約的。

第二：\emph{滿足極小條件的無根基環有單位元。}

為了先證明左單位元的存在，只需證明以下事實：\emph{如果一個右理想 $\ma=c_1\mo\ne\mo$ 有左單位元 $c_1$，那麼存在一個長度更大的右理想 $c\mo$，它同樣有左單位元 $c$。} 因為完全可約性，也就是所有長度的有界性，已經證明；於是經過有限多步就達到 $\mo$ 本身的一個左單位元。

於是設 $\ma=c_1\mo$ 且 $c_1^2=c_1$。由公式
\[
        r=c_1r+(r-c_1r)
\]
如上得到 Peirce 分解 $\mo=c_1\mo+\mr$。如上令 $c_2$ 是 $\mr$ 中的一個冪等元素，即 $c_2^2=c_2$，並且 $c_1c_2=0$（因為 $c_1$ 零化 $\mr$ 的所有元素）。現在置 $e_1=c_1$，$e_2=c_2-c_2c_1$，則有
\[
        e_1^2=e_1,
        \quad e_1e_2=0,
        \quad e_2e_1=0,
        \quad e_2c_2=c_2,
        \quad e_2\ne0,
        \quad e_2^2=e_2.
\]
根據注中所作的說明，$c=e_1+e_2$ 因而成為環
\[
        c\mo=e_1\mo+e_2\mo
\]
的左單位元；由於 $e_2^2=e_2\ne0$，這個環作為右理想比 $\ma$ 有更大的長度。

為證明所構造的左單位元 $e$ 也是右單位元，我們把 Peirce 分解用於左理想：
\[
        \mo=\mo e+\ml.
\]
有 $\ml e=0$，$\ml=e\ml$，所以 $\ml^2=\ml e\ml=0$。又按假設不存在非零冪零左理想，故 $\ml=0$；於是 $e$ 也是右單位元，從而就是單位元。

\emph{定理. 由右側完全可約性和單位元的存在，推出雙邊完全可約性：}
\[
        \mo=\md_1+\cdots+\md_s,
        \qquad \md_i \hbox{ 双边单}.
\]

同前一個證明一樣，只需證明每個極小的（雙邊）理想 $\ma$ 都是直和項。作為右理想，$\ma$ 是直和項：
\[
        \mo=\ma+\mr=e_1\mo+e_2\mo.
\]
相應的左側分解為
\[
        \mo=\mc+\ml=\mo e_1+\mo e_2.
\]
於是得到
\[
        \mr\mc\subseteqq \mr\ma\subseteqq \mr\cap\ma=0,
        \qquad
        \ml\ma=\mo e_2e_1\mo=0,
\]
\[
        (\ma\ml)^2=\ma\ml\ma\ml=0,
        \qquad \hbox{所以 } \ma\ml=0;
\]
\[
        \mr=\mr\mo=(\mr\mc,\mr\ml)=\mr\ml,
        \qquad
        \ml=\mo\ml=(\ma\ml,\mr\ml)=\mr\ml.
\]
因此 $\mr=\ml$，所以 $\mr$ 是雙邊的；於是 $\ma$ 是雙邊直和項。從這裡以後，證明同單側完全可約性的證明相同。

這些 $\md_i$ 作為環也都是雙邊單的，因為 $\md_i$ 中的所有理想同時也是 $\mo$ 中的理想。下面研究它們的結構。

\subsection*{§ 14. 帶單位元的雙邊單完全可約環}

設 $\mo$ 是這樣一個環：
\[
        \mo=\mr_1+\ldots+\mr_n=e_1\mo+\ldots+e_n\mo,
        \qquad \mr_i\hbox{ 单},
\]
於是得到
\[
        \mo=\ml_1+\ldots+\ml_n=\mo e_1+\ldots+\mo e_n,
        \qquad \ml_i\hbox{ 直接不可分解 }\hbox{\(\S 9\)}.
\]
$\ml_i\mr_i$ 是非零的雙邊理想（因為 $e_i^2$ 屬於 $\ml_i\mr_i$），所以等於 $=\mo$。

置 $\mA_{ik}=\mr_i\ml_k$，於是
\[
        \mo=(\mr_1,\ldots,\mr_n)\cdot(\ml_1,\ldots,\ml_n)
        =(\ldots,\mA_{ij},\ldots)
\]
是直和；因為若 $0=\sum_{i,j}a_{ij}$，由於 $a_{ij}$ 屬於 $\mr_i$ 且 $\sum\mr_i$ 是直和，便有
\[
        0=\sum_j a_{ij},
\]
再由 $a_{ij}$ 屬於 $\ml_j$，得到
\[
        0=a_{ij}.
\]

進一步，$\mA_{ij}\mr_j=\mr_i\ml_j\mr_j=\mr_i\mo=\mr_i$，所以 $\mA_{ij}\ne0$。取 $a_{ij}\ne0$ 屬於 $\mA_{ij}$。於是
\[
        a_{ij}\mr_j\subseteqq\mr_i,
        \qquad
        a_{ij}\mr_j\ne0 \hbox{ 因为 } a_{ij}e_j=a_{ij},
\]
所以
\[
        a_{ij}\mr_j=\mr_i.
\]
若對 $\mr_j$ 中的每個 $r_j$ 置
\[
        a_{ij}r_j=r_i,
\]
則映射 $r_j\to r_i$ 是從 $\mr_j$ 到 $\mr_i$ 的算子同態。被映到零的元素構成一個包含於 $\mr_j$ 的理想，所以是零理想；因此這是一個\emph{同構}。由此：

\emph{在 $\mo$ 的一個表示中出現的任意兩個單右理想 $\mr_i$ 和 $\mr_j$ 都是算子同構的；這些同構由 $\mA_{ij}$ 的元素來中介。} 由於任意兩個不同的單右理想 $\mr,\mr'$ 總是在 $\mo$ 的至少一個表示中共同出現，所以任意兩個這樣的理想也同構。

\emph{從 $\mr_j$ 到 $\mr_i$ 的所有同態都由 $\mA_{ij}$ 的元素來中介。}

\emph{證明.} 如果 $e_j\to a_{ij}$，則 $e_j^2\to a_{ij}e_j$，所以 $a_{ij}=a_{ij}e_j$ 屬於 $\mr_i\ml_j=\mA_{ij}$；進一步有
\[
        \mr_j=e_j\mr_j\to a_{ij}\mr_j=\mr_i.
\]

\emph{特別地，$\mr_i$ 到自身的同態由 $\mA_{ii}$ 來中介。}

$\mA_{ii}$ 中兩個元素的積對應於同態的積，和對應於同態的和（定義見 §1, 6）。不同的 $a_{ii}$ 給出不同的自同構，因為 $e_i$ 被映到 $a_{ii}e_i=a_{ii}$。所以環 $\mA_{ii}$ 同構於 $\mr_i$ 的自同構環。

但單理想的自同構環是除環（§2），所以 $\mA_{ii}$ 是除環。又因為所有 $\mr_i$ 都算子同構，它們的自同構環也環同構；因此 $\mA_{ii}$ 在環同構意義下由 $\mo$ 唯一確定。可能出現的不同 $\mA_{ii}$ 只是單右理想的抽象自同構環的不同具體實現。

\emph{主定理. 每個帶單位元的完全可約雙邊單環 $\mo$，都同構於除環 $K$ 上的 $n$ 階矩陣環。（除環 $\mathsf{K}$ 同構於 $\mo$ 的右理想的自同構除環。）}

\emph{證明. 矩陣單位 $c_{ik}$ 的構造。} 令 $\Gamma_{11}$ 為 $\mr_1$ 的恆等自同構，令 $\Gamma_{i1}$ 為一個任意同構 $\Gamma_{i1}\mr_1=\mr_i$，最後置
\[
        \Gamma_{ik}=\Gamma_{i1}\Gamma_{k1}^{-1}.
\]
於是一般地有
\[
        \Gamma_{ik}\Gamma_{kl}=\Gamma_{il}.
\]
若 $\Gamma_{ik}$ 由 $\mA_{ik}$ 中的元素 $c_{ik}$ 中介，則進一步得到：
\[
        c_{ik}=e_i c_{ik}=c_{ik}e_k,
        \qquad
        c_{ik}c_{kl}e_l=c_{il}e_l,
\]
因而
\[
        c_{ik}c_{kl}=c_{il},
        \qquad
        c_{ij}c_{kl}=c_{ij}e_j e_k c_{kl}=0\quad(j\ne k).
\]
所以這些 $c_{ik}$ 是矩陣單位（§7）。

\footnotetext{近來 B. L. van der Waerden 找到了一個更簡單的證明，它直接用抽象自同構除環運算；該證明將收入他的代數學著作（Grundlehren d. Math. Wiss., Berlin: Julius Springer）。[1929 年 6 月 11 日。]}

\emph{$\mathsf{K}$ 的構造。} 對應
\[
        a_{ii}=c_{i1}a_{11}c_{1i}
\]
把 $\mA_{11}$ 中的每個 $a_{11}$ 分配給 $\mA_{ii}$ 中的一個“共軛元素” $a_{ii}$。

這個對應是環同構，因為同 $a$ 和 $c$ 相應的同構 $\Delta,\Gamma$ 由
\[
        \Delta_{ii}=\Gamma_{i1}\Delta_{11}\Gamma_{i1}^{-1}
\]
聯繫起來；這是一個眾所周知給出環同構的關係。現在構造元素
\[
        \alpha=a_{11}+\ldots+a_{nn}
        =a_{11}+c_{21}a_{11}c_{12}+\ldots+c_{n1}a_{11}c_{1n}.
\]
每個 $a_{11}$ 都對應於一個 $\alpha$（由於直和性，這是一一對應）。全體 $\alpha$ 構成一個除環 $K\sim\mA_{11}$；因為由
\[
        \alpha=a_{11}+\ldots+a_{nn},
        \qquad
        \beta=b_{11}+\ldots+b_{nn},
\]
得到
\[
        \alpha+\beta=(a_{11}+b_{11})+\ldots+(a_{nn}+b_{nn}),
        \qquad
        \alpha\cdot\beta=a_{11}b_{11}+\ldots+a_{nn}b_{nn};
\]
所以對應 $a_{11}\to\alpha$ 是同構。

進一步，
\[
        c_{ik}\alpha=c_{ik}a_{kk}=c_{ik}c_{k1}a_{11}c_{1k}=c_{i1}a_{11}c_{1k},
\]
\[
        \alpha c_{ik}=a_{ii}c_{ik}=c_{i1}a_{11}c_{1i}c_{ik}=c_{i1}a_{11}c_{1k},
\]
所以每個 $\alpha$ 都同所有 $c_{ik}$ 可交換。

最後，由同樣的公式得到
\[
        c_{ik}\mathsf{K}=c_{ik}\mA_{kk}=c_{ik}\mr_k\ml_k=\mr_i\ml_k=\mA_{ik},
\]
因此
\[
        \mo=\sum c_{ik}\mathsf{K}.
\]

\emph{對應。} 若把 $\mo$ 的每個元素寫成 $\sum c_{ik}\alpha_{ik}$ 的形式，並給該元素對應矩陣 $(\alpha_{ik})$，則由矩陣乘法的定義立刻可知，這個對應是一個同構。主定理由此得證。

\emph{逆命題. 除環 $A$ 上的 $n$ 階矩陣環是雙邊單且完全可約的；$A$ 同構於右理想的自同構環，而 $Ae$ 就是前一定理中構造出的 $\mathsf{K}$。}

\emph{證明.} 設矩陣單位（§7）為 $c_{ik}$。置
\[
        \mr_i=\sum_k c_{ik}A.
\]
顯然 $\mo=\mr_1+\ldots+\mr_n$。這些 $\mr_i$ 是單右理想；因為每個非零元素 $\sum c_{ik}\alpha_k\ne0$ 都生成整個 $\mr_i$。事實上，若 $\alpha_j\ne0$，則
\[
        \Bigl(\sum c_{ik}\alpha_k\Bigr)\cdot(c_{jl}\alpha_j^{-1})=c_{il},
\]
而 $\sum c_{il}\beta_l$ 遍歷 $\mr_i$ 的所有元素。

這些 $\mr_i$ 算子同構：$\mr_i=c_{ik}\mr_k$。

由此推出：$\mo$ 是雙邊不可分解的；所以，依據 §13 的定理、已經證明的完全可約性以及單位元的存在，它是雙邊單的。也可以這樣看出這一點：任意非零的 $\sum\sum c_{ik}\alpha_{ik}\ne0$ 都生成整個雙邊理想 $\mo$。

進一步有
\[
        \mr_i=c_{ii}\mo,
        \qquad
        \ml_i=\mo c_{ii},
\]
\[
        e=\sum c_{ii},
        \qquad c_{ii}=e_i,
\]
\[
        \mA_{ii}=\mr_i\cap\ml_i=c_{ii}\Lambda\simeq\Lambda.
\]
最後置 $a_{11}=c_{11}\lambda$，則得到
\[
        \alpha=a_{11}+c_{21}a_{11}c_{12}+\ldots
        =c_{11}\lambda+c_{22}\lambda+\ldots=e\lambda,
\]
從而全部命題都得證。

\emph{$\mo$ 的中心就是 $\mathsf{K}$ 的中心，因此是一個交換域。}

\emph{證明.} $\mZ(\mathsf{K})$ 由所有同一切 $x$ 可交換的 $\zeta$ 組成。但這些元素也同所有 $c_{ik}$ 可交換，因此也同所有和 $\sum c_{ik}a_{ik}$ 可交換。所以
\[
        \mZ(\mathsf{K})\subseteqq\mZ(\mo).
\]
反過來，設 $z$ 屬於 $\mZ(\mo)$：
\[
        z=\sum_\mu\sum_\nu c_{\mu\nu}\gamma_{\mu\nu}.
\]
於是
\[
        zc_{ik}=c_{ik}z,
        \qquad
        \sum_\mu c_{\mu k}\gamma_{\mu i}=\sum_\nu c_{i\nu}\gamma_{k\nu},
\]
所以
\[
        \gamma_{\mu i}=0\quad(\mu\ne i),
        \qquad
        c_{ik}\gamma_{ii}=c_{ik}\gamma_{kk},
        \qquad
        \gamma_{ii}=\gamma_{kk}=\gamma;
\]
\[
        z=\sum c_{ii}\gamma=e\gamma=\gamma.
\]
所以 $z$ 屬於 $\mathsf{K}$，並且同所有 $x$ 可交換；故 $z$ 屬於 $\mZ(\mathsf{K})$。

由於矩陣環當然也是左側完全可約的，而每個帶單位元的右完全可約環都是矩陣環的直和，於是得到：

\emph{每個帶單位元的右完全可約環也是左完全可約的，反之亦然。}

並且：\emph{帶單位元的完全可約環的中心，是對應於雙邊單矩陣環的若干交換域的直和。}

還剩下以下定理：

\emph{任意兩個不同的分解 $\mo=\sum c_{ik}\mathsf{K}$，$\mo=\sum c'_{ik}\mathsf{K}'$，可以由內自同構 $\alpha'\to x^{-1}\alpha'x$，$c_{ik}=x^{-1}c_{ik}x$ 互相轉化。}

\emph{證明.} 設
\[
        \mo=\mr_1+\ldots+\mr_n=\mr_1'+\cdots+\mr_n',
\]
並設 $a,b$ 是中介理想 $\mr_1,\mr_1'$ 的兩個互逆同構的元素：
\[
        \mr_1=a\mr_1',
        \qquad
        \mr_1'=b\mr_1,
        \qquad
        ab=c_{11},
        \qquad
        ba=c'_{11}.
\]
置
\[
        x=\sum_i c_{i1}ac'_{1i},
        \qquad
        y=\sum_i c'_{i1}bc_{1i},
\]
於是
\[
        xy=e,
        \qquad
        y=x^{-1},
        \qquad
        x^{-1}c_{ik}x=c'_{ik},
        \qquad
        x^{-1}\alpha x=\alpha'.
\]

\section*{34. 超複量與表示論}
\emph{續篇與結束：第三章，§§15--19；第四章，§§20--26。}

\subsection*{第三章. 模與表示論}

\subsection*{§ 15. 表示與表示模}

設 $\mo$ 是一個環，$\mathsf{K}$ 是一個帶單位元的環。（在後面的應用中，$\mathsf{K}$ 總是除環。）

$\mo$ 在 $K$ 中的一個 $n$ 次表示，是一個同態
\[
        \mo\sim\mD,
\]
其中 $\mD$ 是 $K$ 上 $n$ 階矩陣所成的一個環。

所謂 $\mo$ 關於 $\mathsf{K}$ 的一個\emph{表示模}，是指一個雙模 $\mM$（§1, 5），它是左 $\mo$-模並且是右 $\mathsf{K}$-模：
\[
        \mo\mM\subseteqq\mM,
        \qquad
        \mM \mathsf{K}\subseteqq\mM,
\]
並且它還是有限多個單生成元 $\mathsf{K}$-模的直和：
\[
        \mM=x_1\mathsf{K}+\ldots+x_n\mathsf{K};
\]
其中 $\mathsf{K}$ 的單位元作為單位算子。

\emph{每個表示模給出一個表示。} 設 $c$ 屬於 $\mo$，並且
\[
        cx_k=\sum x_i\gamma_{ik},
\]
或者寫成矩陣形式：
\[
        c(x_1,\ldots,x_n)=(x_1,\ldots,x_n)C,
        \qquad C=(\gamma_{ik}).
\]
於是矩陣 $C$ 對 $c$ 構成一個表示，因為
\[
        (b+c)x_k=\sum x_i(\beta_{ik}+\gamma_{ik}),
\]
並且
\[
%
        bcx_k
        =b\sum x_j\gamma_{jk}
          =\sum b x_j\gamma_{jk}
          =\sum\sum x_i\beta_{ij}\gamma_{jk}
        =\sum x_i(\sum\beta_{ij}\gamma_{jk})
%
\]
或者
\[
        bc(x_1,\ldots,x_n)=(x_1,\ldots,x_n)BC.
\]

反過來：$\mo$ 的每個表示 $\mD$ 都屬於一個表示模，並且屬於該模的一組確定基。也就是說，把 $\mM$ 理解為 $x_1,\ldots,x_n$ 中所有形式線性形式的全體：
\[
        y=\sum x_i\alpha_i,
        \qquad \alpha_i\in K.
\]
那麼 $\mM$ 是一個右 $\mathsf{K}$-模。進一步，如果元素 $c$ 對應於矩陣 $C=(\gamma_{ik})$，定義
\[
        cx_k=\sum x_i\gamma_{ik},
        \qquad
        c\biggl(\sum_k x_k\alpha_k\biggr)
        =\sum_i x_i\sum_k\gamma_{ik}\alpha_k.
\]
由同態關係
\[
        c+d\to C+D,
        \qquad cd\to CD
\]
可得
\[
        (c+d)x_i=cx_i+dx_i,
        \qquad cdx_i=c(dx_i),
\]
而對和 $\sum x_i\alpha_i$ 也同樣成立。其餘雙模性質
\[
        c(y+z)=cy+cz,
        \qquad c(y\alpha)=(cy)\alpha
\]
是顯然的。因此 $\mM$ 是表示模；按構造，它正好屬於表示 $\mo\sim\mD$。

現在設 $(y_1,\ldots,y_n)$ 是同一個表示模的另一組基：
\[
        (y_1,\ldots,y_n)=(x_1,\ldots,x_n)\cdot P,
        \qquad
        (x_1,\ldots,x_n)=(y_1,\ldots,y_n)\cdot P^{-1}.
\]
若相對於基 $x$ 有 $b\sim B$，則
\[
        b(y_1,\ldots,y_n)=b(x_1,\ldots,x_n)P
        =(x_1,\ldots,x_n)BP
        =(y_1,\ldots,y_n)P^{-1}BP.
\]
稱表示 $b\sim B$ 和 $b\sim P^{-1}BP$ 等價，並把它們歸入同一個表示類。由於每個正則矩陣 $P$ 都把 $\mM$ 的一組基變為另一組基，故已證明：

\emph{每個表示模唯一地給出一個確定的表示類。}\textsuperscript{15a)}

{\footnotesize\noindent
15) 關於交換域的表示模最早見於 W. Krull, \emph{Theorie und Anwendung der verallgemeinerten Abelschen Gruppen}, Sitzungsber. Heidelberger Ak. 1926, 第一論文。\par
15a) 校樣補記（1929 年 4 月 14 日）。據 B. L. van der Waerden 告知，可以通過區分“線性變換”和“矩陣”這兩個概念，直接得到一個不依賴特殊基選擇、因而不變的聯繫。線性變換是兩個線性形式模之間的同態；矩陣則是這個同態在一組特定基下的表達（表示）。I. 每個表示模都把乘子域 $\mo$ 唯一地分配給該模到自身的一套同態的線性變換系統。因為按照 §2 末尾，雙模 $\mM$ 的左乘子關於右算子（這裡即乘以 $\mathsf{K}$）生成 $\mM$ 到自身的算子同態，而且這個對應是同態的。II. 每一個與 $\mo$ 同態的、線性形式模到自身的線性變換系統，都給出一個表示模。\par}

顯然：兩個算子同構的表示模對應於同一個表示類。反過來也成立：如果兩個表示模 $(x_1,\ldots,x_n)$ 和 $(\bar x_1,\ldots,\bar x_n)$ 產生同一個表示，那麼對應
\[
        x_i\to\bar x_i,
        \qquad
        \sum_i x_i\alpha_i\to\sum_i\bar x_i\alpha_i
\]
給出一個同構。因為乘法規則完全由矩陣 $C$ 決定。

特殊情形：如果 $\mM$ 的兩組基產生同一個表示，那麼它們由 $\mM$ 到自身上的一個算子同構互相轉化。

或者說：若對所有 $C$ 和一個固定的 $P$ 都有 $C=PCP^{-1}$，則對應
\[
        y_i\mapsto x_i,
        \qquad y_i=\sum_kx_k\pi_{ik},
\]
是 $\mM$ 到自身上的一個同構。反過來：雙模 $\mM$ 到自身上的每個同構 $y_i\mapsto x_i$，都由一個同所有 $C$ 可交換的矩陣 $P$ 中介。

還可以把表示概念推廣到這樣的環：它還同一個交換乘子域 $P$ 交換地聯繫在一起（§8）。在這種情形下，只取那些在中心中含有 $P$ 的表示環 $K$；並且要求表示不僅是環同態 $\mo\sim\mD$，甚至還應當是算子同態：由 $a\sim A$ 應推出
\[
        a\varrho\sim A\varrho\qquad(\varrho\in P).
\]
對於表示模，這一要求意味着附加的計算規則
\[
        am\cdot\varrho=a(m\varrho)=a\varrho\cdot m.
\]
這時就說，模與環“同 $P$ 交換地聯繫”。

\subsection*{§ 16. 可約表示}

如果 $\mU$ 是表示模 $\mM$ 的一個子模，並且可以為 $\mM$ 選取一組基，它由 $\mU$ 的一組基 $z_1,\ldots,z_t$ 補以 $y_1,\ldots,y_r$ 而成，也就是說
\[
        \mM=y_1\mathsf{K}+\ldots+y_r\mathsf{K}+z_1\mathsf{K}+\ldots+z_t\mathsf{K}\textsuperscript{16)}
\]
那麼這些表示具有形式
\[
        C=\begin{pmatrix}R&0\\ S&T\end{pmatrix},
\]
其中矩陣 $T$ 自身構成一個由 $\mU$ 中介的 $t$ 次表示，矩陣 $R$ 構成一個由 $\mM/\mU$ 中介的 $r$ 次表示。

\emph{證明.} 有
\[
        (cy_1,\ldots,cy_r,cz_1,\ldots,cz_t)
        =(y_1,\ldots,y_r,z_1,\ldots,z_t)\cdot C.
\]
由於 $cz_i$ 作為 $\mU$ 的元素，只能用 $z_i$ 表出，所以 $C$ 的右上角塊為零。若把 $C$ 的其餘部分按上面的排列稱為 $R,S,T$，則特別有
\[
        (cz_1,\ldots,cz_t)=(z_1,\ldots,z_t)\cdot T;
\]
因此這些 $T$ 構成由 $\mU$ 中介的表示。進一步有
\[
        (cy_1,\ldots,cy_r)\equiv (y_1,\ldots,y_r)R\pmod{\mU},
\]
而 $y_i$ 構成模 $\mU$ 意義下線性無關的一組基。所以這些 $R$ 構成由 $\mM/\mU$ 生成的表示。

反過來，如果給定一個“可約的”表示
\[
        C=\begin{pmatrix}R&0\\ S&T\end{pmatrix},
\]
其中 $R$ 和 $T$ 是方陣，那麼在相應的表示模中，每個 $c$ 同最後 $t$ 個基元素 $z_1,\ldots,z_t$ 的積都只用這些元素表出；也就是說，$\mU=(z_1,\ldots,z_t)$ 是一個子模。

{\footnotesize\noindent 16) 換句話說：$\mU$ 作為 $\mathsf{K}$-模是一個直和項。如果 $\mathsf{K}$ 是除環，那麼這個前提總是滿足。\par}

\emph{推論.} 設
\[
        \mM\supset \mU_1\supset \mU_2\supset\cdots\supset \mU_e=0
\]
是 $\mM$ 的一條合成列，並且每個 $\mU_i$ 在前一項中作為 $K$-模是直和項，從而可以選擇一組基
\[
        z_{11},\ldots,z_{1r_1},\ z_{21},\ldots,z_{2r_2},\ldots,
        z_{e1},\ldots,z_{er_e}
\]
使得所有滿足 $i>\nu$ 的 $z_{ik}$ 構成 $\mU_\nu$ 的一組基。於是由 $\mM$ 中介的表示具有形式
\[
 C=\begin{pmatrix}
 R_{11}&0&\ldots&0\\
 R_{12}&R_{22}&\ldots&0\\
 \vdots&\vdots&\ddots&\vdots\\
 R_{1e}&R_{2e}&\ldots&R_{ee}
 \end{pmatrix},                                      \tag{2}
\]
並且 $R_{\nu\nu}$ 構成由合成因子 $\mU_{\nu-1}/\mU_\nu$ 中介的表示。

各個表示 $R_{\nu\nu}$ 是不可約的，因為合成因子 $\mU_{\nu-1}/\mU_\nu$ 是單的。若假定 $K$ 是除環，則反過來，每個已經儘可能化為形式 (2) 的表示也給出一條合成列。按 Jordan--Hölder 定理，合成因子在算子同構意義下唯一確定；因此這些 $R_{\nu\nu}$ 在等價表示意義下並且在排列順序以外唯一確定。

如果 $\mM=\mA+\mB$ 是兩個表示模 $\mA=(y_1,\ldots,y_r)$、$\mB=(z_1,\ldots,z_s)$ 的直和，那麼由基 $(y_1,\ldots,y_r,z_1,\ldots,z_s)$ 中介的表示顯然具有形式
\[
        C=\begin{pmatrix}R&0\\0&S\end{pmatrix},
\]
其中 $R$ 表示由 $\mA$ 中介的元素 $c$ 的表示，$S$ 表示由 $\mB$ 中介的表示；反過來也成立。若先把模 $\mM$ 表示為直接不可分解成分的直和，並且如上在這些成分中尋求合成列，則在適當選取基後，表示相應地呈塊狀形式。若仍假定 $\mathsf{K}$ 是除環，則反過來，每個儘可能化約的表示也給出模由直接不可分解和項以及其中合成因子所作的表示。按 Krull 與 Otto Schmidt 定理，直接不可分解的表示成分的類在排列順序以外唯一確定。

如果模 $\mM$ 是完全可約的，那麼所有方塊都只由一個不可約表示構成，並稱這個表示為完全可約的。

\subsection*{§ 17. 帶單位元的環的表示模的直和分解}

設 $\mM$ 是一個 $\mo$-模，$\mo$ 是帶單位元的環。則有
\[
        \mM=\mM_e+\mM_0,
\]
其中單位元在 $\mM_e$ 上是單位算子，在 $\mM_0$ 上是零算子。（如果 $\mM$ 同時是右 $K$-模，那麼 $\mM_e$ 和 $\mM_0$ 也都是右 $\mathsf{K}$-模。）

\emph{證明.} $\mM$ 的每個 $m$ 都可以表示為
\[
        m=em+(m-em).
\]
令所有 $em$ 的全體為 $\mM_e$，所有 $m-em$ 的全體為 $\mM_0$。那麼 $\mM_e$ 和 $\mM_0$ 顯然是加法群；如果 $\mM$ 是右 $K$-模，它們也都是右 $K$-模。進一步有
\[
        r em=erm,
\]
所以 $\mM_e$ 也是 $\mo$-模；同樣
\[
        r(m-em)=rm-rem=rm-rm,
\]
所以 $\mM_0$ 也是 $\mo$-模。對於 $em$，$e$ 是單位算子；對於 $m-em$，$e$ 是零算子。因此只有零元同時屬於 $\mM_e$ 和 $\mM_0$，故和 $\mM_e+\mM_0$ 是直和。

如果討論的是表示模，那麼 $\mM_0$ 產生平凡表示，即每個元素都被賦予零。把直和項 $\mM_0$ 分離出去以後，剩下由 $\mM_e$ 中介的表示，在這個表示中單位元對應於單位矩陣。因此，我們從現在起可以並且將只限於這樣的表示。

設
\[
        \mo=\ma_1+\ldots+\ma_s
\]
是雙邊理想的直和，
\[
        e=e_1+\ldots+e_s
\]
是 $e$ 的相應分解，並且 $\mM$ 是一個 $\mo$-模，其中 $e$ 是單位算子。則
\[
        \mM=\ma_1\mM+\ldots+\ma_s\mM
\]
是直和。顯然，$\mM=\mo\mM=(\ma_1\mM,\ldots,\ma_s\mM)$。置
\[
        \mb_i=\ma_1+\cdots+\widehat{\ma_i}+\cdots+\ma_s,
\]
則相應地 $\mb_i$ 是一個雙邊理想，並且
\[
        \mM=(\ma_i\mM,\mb_i\mM),
\]
而這個和是直和，因為 $e_i$ 在 $\ma_i\mM$ 上是單位算子，在 $\mb_i\mM$ 上是零算子。

根據這個定理，我們通常限制於雙邊不可分解環的表示。


\subsection*{§ 18. 完全可約環的模論與表示論}

設 $\mo$ 是完全可約且雙邊單的環，$\mM$ 是一個有限的 $\mo$-模，並且 $\mo$ 的單位元在 $\mM$ 上是單位算子。則 $\mM$ 完全可約，並且它的簡單組成成分與簡單左理想 $\ml_i$ 算子同構。

\emph{證明.} 設
\[
        \mo=\ml_1+\ldots+\ml_n,
        \qquad
        \mM=(\mo m_1,\ldots,\mo m_k),
\]
於是
\[
        \mM=(\ldots,\ml_i m_k,\ldots).                         \tag{4}
\]
那些非零的 $\ml_i m_k$ 同構於 $\ml_i$，因為映射 $a\mapsto am_k$ 是算子同構。因此這些 $\ml_i m_k$ 是簡單的。從表示 (4) 中刪去那些已經包含在前面各項之和中的 $\ml_i m_k$，這個和就成為直和。

\emph{推論.} 如果此外 $\mM$ 是直接不可分解的，那麼 $\mM$ 是簡單的，並且同構於某個 $\ml_i$。

現在設 $\Delta$ 是這個簡單 $\mM$ 的自同構除環，$\Lambda$ 是 $\ml_i$ 的自同構除環。由於 $\mM$ 與 $\ml_i$ 算子同構，得到環同構 $\Delta\simeq\Lambda$。按照 §1，可以把 $\mM$ 和 $\ml$ 分別看作以 $\Delta$ 和 $\Lambda$ 為右作用域的雙模；這些雙模同時也是表示模。由它們中介的表示因而經由環同構 $\Delta\simeq\Lambda$ 相互轉化。因此，為研究這一表示類，我們可以限於由 $\ml_1$ 在其自同構除環中中介的表示。

特別地，為 $\ml_1$ 選取一組矩陣單位作為基
\[
        \ml_1=(c_{11},\ldots,c_{n1})
\]
並把先前（§14）中 $\mo$ 的反除環（當時記作 $K$）作為自同構除環 $\Lambda$ 的實現；在先前的討論中，應以左理想代替右理想，並把自同構寫在右邊。若
\[
        c=\sum_{i,k}c_{ik}\alpha_{ik}
\]
是 $\mo$ 的一個元素，則
\[
        (cc_{11},\ldots,cc_{n1})=(c_{11},\ldots,c_{n1})(\alpha_{ik}).
\]
因此，由 $\ml_1$ 中介的表示由矩陣 $(\alpha_{ik})$ 給出。對於 $\ml_\nu=(c_{1\nu},\ldots,c_{n\nu})$，得到完全相同的表示。

如果 $\mo$ 是雙邊單成分的直和，那麼這些成分依次同構地被表示，而其餘成分作用為零。

進一步，若把 $\mo$ 的每個元素 $a$ 寫成
\[
        a=\sum c_{ik}\alpha_{ik},
\]
則 $a\mapsto(\alpha_{ik})$ 就是在 $K$ 中由 $\ml_1$ 中介的表示。

設 $\Gamma$ 是 $K$ 的一個子除環，使 $K$ 關於 $\Gamma$ 有有限秩 $t$：
\[
        K=\chi_1\Gamma+\cdots+\chi_t\Gamma.
\]
那麼 $K$ 關於 $\Gamma$ 成為它自身的表示模。$K$ 的元素 $\beta$ 由 $\Gamma$ 中的矩陣 $B$ 表示，其取得方式為
\[
        \beta\varkappa_j=\sum_i\chi_i\beta_{ij},\qquad B=(\beta_{ij}).
\]
現在把 $\ml_1$ 看作關於子除環 $\Gamma$ 的表示模，就得到如下結論：

\emph{由 $\ml_1$ 中介的 $\mo$ 在 $\Gamma$ 中的表示，是把上面給出的 $\mo$ 在 $\mathsf{K}$ 中的表示
\[
        a\to(\alpha_{ik})
\]
裡的元素 $\alpha_{ik}$ 替換為矩陣 $A_{ik}$；這些矩陣是在 $\mathsf{K}$ 關於 $\Gamma$ 的表示中同 $\alpha_{ik}$ 對應的矩陣。}

\emph{證明.}
\[
        \ml_1=\sum_i c_{i1}\mathsf{K}=\sum_i\sum_j c_{i1}\varkappa_j\Gamma.
\]
關於基 $(\ldots,c_{i1}\varkappa_j,\ldots)$，元素 $c_{ik}$ 由一個塊矩陣表示：在第 $(i,k)$ 塊中是單位矩陣 $E_t$，其餘塊是 $t$ 階零矩陣 $0$。同樣，$K$ 中的元素 $\alpha$ 由帶有相應 $t$ 階矩陣 $A$ 的塊對角矩陣表示。因此
\[
        a=\sum c_{ik}\alpha_{ik}
\]
對應於塊矩陣 $(A_{ik})$。從這裡過渡到任意的 $\mM$，對於表示仍然意味着一個環同構。

\emph{注.} 這裡所研究的左理想的自同構除環及其有限子除環中的表示，在同構意義下，正是那些能夠由簡單 $\mo$-模或左理想中介的表示。因為按照前面（§2）的一個說明，如果把 $\mo$-模 $\mM$ 看作關於某個除環 $\mathsf{K}$ 的雙模，那麼 $\mathsf{K}$ 的元素產生模同態。因此這個除環 $\mathsf{K}$ 必然以同態方式對應到該 $\mo$-模的自同構除環的一個子除環。由於單位元是單位算子，這個同態是同構；並且整個自同構除環關於相應的子除環必須具有有限次數，否則表示模也會具有無限秩，這是不可能的。

最後還要指出，這裡所研究的不可約表示還不是全部，而只是那些由簡單 $\mo$-模中介的表示。一般還存在別的表示模，例如把一個除環 $\Omega$ 看作關於子除環 $\Omega_0$ 的左作用域、關於它自身的右作用域的雙模；這樣的模作為 $\mo$-模是可約的，但作為雙模不可約，因此仍然導致不可約表示。

\subsection*{§ 19. 模與表示模中的簡單合成因子}

設 $\mo$ 是一個關於左理想滿足極大條件與極小條件的環，$\mc$ 是最大冪零理想。則有：

\emph{每個簡單 $\mo$-模 $\mM$ 要麼被 $\mo$ 零化，要麼同構於無根基環 $\mo/\mc$ 中的一個簡單左理想 $\ml$。在後一種情形下，該模被 $\mo/\mc$ 中所有不包含 $\ml$ 的雙邊理想零化，並且絕對乘子域（§1）同構於包含 $\ml$ 的那個雙邊單環。}

\emph{證明.} 設 $\mc^{\varrho}=0$。必須有 $\mc\mM=0$；因為否則 $\mc\mM=\mM$，於是 $\mc^\varrho\mM=\mM$，這同 $\mc^\varrho=0$ 所推出的矛盾 $\mM=0$ 相沖突。因此，也可以把 $\mM$ 看作 $\mo/\mc$-模：同一個模元素相乘時，模 $\mc$ 的每個剩餘類中的所有元素給出同樣的結果。

作為無根基環，$\mo/\mc$ 有單位元。因此 $\mM$ 分解為兩個模的直和：一個被 $\mo/\mc$ 零化，另一個上單位元是單位算子（§17）。由於 $\mM$ 是簡單的，兩個直和項中只能出現一個。如果單位元是單位算子，則同 §18 一樣推出它同構於一個簡單左理想。事實上，若 $m\ne0$ 是 $\mM$ 中的一個元素，則 $(\mo/\mc)m\ne0$；因此對於 $\mo/\mc$ 中至少一個簡單左理想 $\ml$，有 $\ml m\ne0$，而它必須窮盡 $\mM$。其餘斷言對於這樣的左理想是清楚的，並立即轉移到同它們同構的所有模上。

\emph{推論.} 如果 $\mM$ 是一個具有合成列的 $\mo$-模，那麼它的簡單合成因子要麼被 $\mo$ 零化，要麼同構於 $\mo/\mc$ 中的簡單左理想。

若 $\mathsf{P}$ 是一個同 $\mo$ 交換地聯繫在一起的環，則所有結果仍然成立。此時只需把模替換為關於右側 $\mathsf{P}$ 與左側 $\mo$ 的雙模，並要求它們滿足條件
\[
        (am)\varrho=a(m\varrho)=(a\varrho)m \qquad(\S 15).
\]
此外，$\mathsf{P}$ 的單位元也應當是單位算子。所謂子模，始終只考慮那些允許同 $\mathsf{P}$ 相乘的子模，如同理想的情形一樣。我們在證明中構造的子模 $\mc\mM,\mc^2\mM,\ldots$ 都滿足這個條件。

特別地，若 $\mathsf{P}$ 是一個域，而 $\mM$ 是一個有限秩的 $\mathsf{P}$-模，則 $\mM$ 同時是表示模。由 $\mM$ 中介的表示不僅是環同態，而且關於 $\mathsf{P}$ 也是算子同態（§15）。$\mathsf{P}$ 中具有這個性質的所有表示都由這樣的模給出。不可約表示屬於簡單模，反過來也一樣。因此，本節各定理可以立即解釋為關於 $\mo$ 在 $\mathsf{P}$ 中的表示的定理：除了零表示以外，所有不可約表示都由 $\mo/\mc$ 的簡單左理想中介。於是，如果按照 §16 藉助合成列來化約任意一個表示，則在主對角線上，除零以外，只會出現等價於由 $\mo/\mc$ 的簡單左理想中介的那些表示。

\emph{注.} 在關於 $\mo$ 和 $\mathsf{P}$ 所作的上述假設下，不需要對 $\mo$ 另加有限性條件。因為所有表示同時也是其絕對乘子域的同構表示；而這個絕對乘子域作為一個有限階矩陣環的同構原像，本身是一個有限秩的 $\mathsf{P}$-模。又因為按定義只把 $\mathsf{P}$-模看作可容許理想，所以這個絕對乘子域滿足極大條件和極小條件。

如果把 $\mathsf{P}$ 假定為交換域，那麼這個絕對乘子域成為關於 $\mathsf{P}$ 的一個超複系統。只要主對角線上出現的不全是零，這個假定就總是滿足的。因為這時，按照 §18 的注，$\mathsf{P}$ 同構於 $\ml$ 的自同構除環的一個子除環；在通過 $\mo/\mc$ 的子除環 $K$（§14）來實現時，由於交換聯繫性，它必須位於 $K$ 的中心中。

\subsection*{第四章. 群與超複系統的表示}

\subsection*{§ 20. 超複系統的定位}

我們考察關於一個交換域 $P$ 的超複系統 $\mo$。按照約定，在 $\mo$ 中作為理想來考慮的只限於 $\mathsf{P}$-模，因此左理想和右理想都滿足極大條件與極小條件（參見 §8）。所以第二章所闡明的全部環論都適用。

以下所謂超複系統 $\mo$ 的表示，始終指在域 $\mathsf{P}$ 中的表示，並且是這樣的表示：由 $a\sim A$ 可推出
\[
        a\varrho\sim A\varrho \qquad(\varrho\in P)
\]
（即關於 $P$ 的模同態性）。對於表示模，這意味着
\[
        (a\varrho)m=a(m\varrho)
\]
（參見 §15 結尾）。因此，§19 的定理適用於這些表示模；由此可知，所有不可約表示都由
\[
        \mo_0=\mo/\mc
\]
的簡單左理想中介，並且不同等價類的不可約表示的個數，正等於分解
\[
        \mo_0=\ma_1+\cdots+\ma_s
\]
中雙邊單和項 $\ma_\nu$ 的個數。

現在，為了實際確定由這樣一個左理想 $\ml_\nu$ 所定義的 $\mo$ 的表示，首先從 $\mo$ 的元素過渡到它們在 $\mo_0$ 中的相應剩餘類。由於乘以 $P$ 的一個元素給出 $\mo_0$ 中所有理想的同構，域 $P$ 同態地對應於每個簡單左理想的自同構除環 $\mathsf{K}_\nu$ 的一個子除環 $e^{(\nu)}\mathsf{P}$；又由於單位元是單位算子，這個同態實際上是同構。由簡單左理想 $\ml_\nu$ 在 $\mathsf{P}$ 中中介的表示，可以這樣得到：先研究 $\ml_\nu$ 在這個子除環 $e^{(\nu)}\mathsf{P}$ 中中介的表示，再借助同構 $e^{(\nu)}\mathsf{P}\simeq \mathsf{P}$ 過渡到 $\mathsf{P}$。自同構除環 $\mathsf{K}_\nu$ 關於 $\mathsf{P}$ 有有限秩；因此問題是關於 $\ml_\nu$ 的自同構除環的一個有限次子除環中的表示。

於是，為了得到所求表示，先把 $\mo_0$ 的每個元素 $c$ 按其雙邊成分表示為
\[
        c=c_1+\ldots+c_s.
\]
現在只需尋找 $c_\nu$ 的表示矩陣，因為其餘的 $c_i$ 零化理想 $\ml_\nu$，因而由零矩陣表示。按 §18 可如下得到這個矩陣：把 $c_\nu$ 用 $\ma_\nu$ 的矩陣單位 $c_{ik}^{(\nu)}$ 表示為
\[
        c_\nu=\sum c_{ik}^{(\nu)}\alpha_{ik}^{(\nu)}
\]
並在矩陣 $(\alpha_{ik}^{(\nu)})$ 中，把每個元素 $\alpha_{ik}^{(\nu)}$ 替換為矩陣 $A_{ik}^{(\nu)}$；這個矩陣是在由 $\mathsf{K}_\nu$ 中介的、$\mathsf{K}_\nu$ 到 $\mathsf{P}$ 的表示中同它對應的矩陣。

\emph{注.} 除環 $\mathsf{K}_\nu$ 關於 $\mathsf{P}$ 有有限秩。每個元素 $x$ 都滿足一個以 $e^{(\nu)}\mathsf{P}$ 中元素為係數的方程 $f(\varkappa)=0$，因為 $x$ 的冪之間必定存在一個線性相關關係。特別地，若 $P$ 是代數閉的，則這個方程分解為一次因子；由於 $\mathsf{K}_\nu$ 是除環，$x$ 已經是某個一次因子的根，因而 $x$ 已經屬於 $e^{(\nu)}\mathsf{P}$，也就是說 $\mathsf{K}_\nu=e^{(\nu)}\mathsf{P}$。在這種情形下，矩陣 $(\alpha_{ik}^{(\nu)})$ 本身就形成不可約表示。

由這個注和 §19 結尾可得“Burnside 定理”：

\emph{設 $\mathsf{P}$ 是代數閉的，並且以交換方式同一個環 $\mo$ 聯繫在一起。那麼 $P$ 中每個 $n$ 階不可約表示恰好含有 $n^2$ 個線性無關的矩陣。}\footnote{因此，並不預先假定 $\mo$ 是超複系統；例如，$\mo$ 可以是由一個無限群中任意有限多個元素的所有線性組合構成的群環。}

由這個表述立刻推出通常的表述：\emph{一個 $P$ 中的 $n$ 階矩陣系統，如果任意兩個矩陣的乘積也屬於該系統，那麼它不可約，當且僅當不能給出一個以 $P$ 中元素為係數的線性齊次關係
\[
        \sum \alpha_{ik}x_{ik}=0
\]
使得該系統中每個矩陣的元素 $x_{ik}$ 都滿足這個關係。} 事實上，從任何這樣的矩陣系統出發，只要添入所有線性組合，就可以導出一個環和一個 $P$-模，並把這個環看作它自身的表示。

\emph{證明.} 根據 §19，絕對乘子域 $\mD$ 同構於一個帶單位元的雙邊單且完全可約的環 $\mo_0$；而這樣的環是關於左理想的自同構除環的全矩陣環（比如全體 $m$ 階矩陣的環），由於代數閉性，這個自同構除環同 $\mathsf{P}$ 重合。$\mD$ 的每個不可約表示，因而給定的表示，也都由一個簡單左理想生成。這個左理想的秩是 $m$，表示的次數是 $n$，所以 $m=n$；於是 $\mo_0$ 中恰好有 $n^2$ 個線性無關元素，$\mD$ 中也同樣如此。

同樣容易推出“廣義 Burnside 定理”：\footnote{G. Frobenius 和 I. Schur, Über die Äquivalenz der Gruppen linearer Substitutionen, Sitzungsber. Berlin 1906。}

\emph{在同樣的假設下，設 $\mD$ 是 $\mo$ 的一個完全可約表示，它分解為彼此不等價的組成部分，其次數為
\[
        n_1,n_2,\ldots,n_s;
\]
則 $\mD$ 關於 $P$ 的秩為
\[
        n_1^2+\cdots+n_s^2.
\]}

\emph{證明.} 仍令 $\mo$ 表示絕對乘子域；它必然是一個超複系統。不同的不等價表示由左理想 $\ml_1,\ldots,\ml_s$ 生成，這些左理想屬於 $\mo/\mc$ 的不同雙邊單成分 $\ma_i$，其中 $\mc$ 是根基。因此，給定表示由 $\ml_1+\cdots+\ml_s$ 生成，並且以 $\ma_1+\cdots+\ma_s$ 為絕對乘子域。於是用上面相同的秩計算即可得到斷言。

現在設 $\mo$ 是關於任意域的無根基超複系統。根據 §13，$\mo$ 完全可約並且有單位元。由 §17 和 §18 的定理推出：$\mo$ 的每個表示都是完全可約的。

反過來，有如下結論：同 $P$ 交換地聯繫在一起的環 $\mo$ 的每個完全可約表示，都以一個無根基的超複系統為絕對乘子域。

\emph{證明.} 絕對乘子域 $\mD$ 同構於一個由 $P$ 中矩陣組成的環和 $P$-模，因此是一個超複系統。根基 $\mc$ 的元素按 §19 在每個不可約表示中都由零表示，因而在整個表示中也由零表示，所以 $\mc=0$。

（如果在表示中只出現等價的組成部分，即它們全都由同一個左理想 $\ml$ 中介，那麼絕對乘子域就是雙邊單的。）

超複系統 $\mo$ 的\emph{正則表示}，指由單位理想 $\mo$ 自身中介的表示。（在群環的情形中，特別選取由群元素 $u_1,\ldots,u_h$ 組成的基；參見 §6。）由於 $\mo$ 的合成因子中包含 $\mo/\mc$ 的所有左理想，如果把正則表示按照 §16 公式 (2) 藉助一條合成列變換到適當的基上，那麼所有不可約表示已經作為對角矩陣 $R_{ii}$ 包含在正則表示中。

一個表示只要知道基元素 $u_1,\ldots,u_h$ 的表示矩陣，就已經確定。特別地，如果所討論的是群環，也就是說 $u_1,\ldots,u_h$ 是一個群的元素，那麼群的每個矩陣同態表示也產生群環的一個表示；因此，群的表示問題是超複系統表示問題的一個特例。

\subsection*{§ 21. 基域的擴張。中心的表示}

若
\[
        \mo=a_1P+\cdots+a_hP
\]
是一個超複系統，而 $\Omega$ 是基域 $\mathsf{P}$ 的一個擴域，則可以形成
\[
        \mo\Omega=a_1\Omega+\ldots+a_h\Omega,
\]
其中基元素 $a_i$ 仍採用舊的乘法規則。\textsuperscript{18a)} 於是 $\mo\Omega$ 仍然是關於 $\Omega$ 的超複系統。

$\mo$ 的每個表示都給出 $\mo\Omega$ 的一個表示，因為一旦知道基元素 $a_i$ 的表示，就知道所有元素的表示。\textsuperscript{19)}

{\footnotesize\noindent 18a) 更準確地說：引入帶有舊乘法規則的新符號 $a_i$；於是 $a_1\Omega+\ldots+a_p\Omega$ 包含一個同構於 $\mo$ 的子環，因此事後可以把這些 $a_i$ 同舊的 $a_i$ 等同起來。\par
19) 當然仍然只考慮那些作為 $P$-模同態的表示；由 $a\sim A$ 應推出 $a\varrho\sim A\varrho$（$\varrho\in P$），對於 $\Omega$ 也相應如此。\par}

如果一個理想或表示模，以及相應的表示，在過渡到代數閉域以後仍保持不可約，就稱為\emph{絕對不可約}。

若 $\mo\Omega$ 沒有根基，則 $\mo$ 也沒有根基；因為 $\mo$ 中的一個冪零理想 $\mc$ 會給出 $\mo\Omega$ 中的擴張理想 $\mc\Omega$。反命題並不成立，如下文將看到的那樣。

$\mo\Omega$ 的中心等於擴張後的中心 $\mZ\Omega$。若 $\mo$ 完全可約，則 $\mZ$ 是若干域的直和：
\[
        \mZ=\mZ_1+\cdots+\mZ_s.
\]
如果 $\mo$ 過渡到 $\mo\Omega$ 後仍保持完全可約（也就是說，沒有新的冪零理想加入），那麼新的中心
\[
        \mZ\Omega=\mZ_1\Omega+\cdots+\mZ_s\Omega
\]
在 $\Omega$ 中仍分解為若干域的直和；當 $\Omega$ 是代數閉的時，按照 §20 的一個說明，這些域關於 $\Omega$ 的秩為 $1$，並且同構於 $\Omega$。

對於中心 $\mZ$ 的表示，有以下定理。

\emph{交換超複系統 $\mZ$ 在代數閉域 $\Omega$ 中的每個不可約表示都是一階的。（或者說：不可約表示同 $\mZ$ 到 $\Omega$ 的同態相同。）}

\emph{證明.} $\mZ$ 的每個不可約表示給出 $\mZ\Omega$ 的一個不可約表示，因而也給出 $\mZ\Omega/\mc$ 的一個不可約表示，其中 $\mc$ 是 $\mZ\Omega$ 的根基。$\mZ\Omega/\mc$ 是一個無根基的交換系統，所以按 §14 結尾，它是若干域的直和。由於假定 $\Omega$ 是代數閉的，這些域都是一階的，並且生成所有不可約表示（§20）。

同時還得到：由於一階的等價表示必然相同（$A^{-1}\alpha A=\alpha$），$\mZ$ 到 $\Omega$ 的不同同態的數目，等於 $\mZ\Omega/\mc$ 的秩。

最後這個說明，在 $\mZ$ 是 $\mathsf{P}$ 上的一個域這一特殊情形中給出如下定理：

\emph{若 $\mZ$ 是一個交換域，則 $\mZ$ 完全可約（無根基），當且僅當 $\mZ$ 是 $\mathsf{P}$ 的第一類擴張。}

這是因為，對於一個域，同態成為同構；它們的數目等於 $\mZ\Omega/\mc$ 的秩，因而當且僅當這個數目等於域次數（即 $\mZ$ 是第一類擴張）時，才有 $\mc=0$。

若 $\mZ$ 完全可約：
\[
        \mZ=\mZ_1+\cdots+\mZ_t,
\]
則在每個表示中，各個域 $\mZ_\nu$ 也被同態地映出；並且在一個不可約表示中，其中一個域被同構地表示，其餘域由零表示（§19）。

把這些定理應用於超複系統，首先對無根基系統得到：

\emph{若 $\mo_\Omega$，因而 $\mo$，是無根基系統，那麼在 $\mo$ 的每個絕對不可約表示中，中心元素 $z$ 都由對角矩陣 $E\xi$ 表示，並且 $z\mapsto\xi$ 是 $\mZ$ 到 $\Omega$ 的一階表示。這樣在 $\mo$ 的絕對不可約表示類與 $\mZ$ 的絕對不可約表示類之間建立的對應是一一的。因此，這些表示類的數目等於中心的秩。}\textsuperscript{20)}

{\footnotesize\noindent 20) 相應的定理不僅適用於代數閉域 $\Omega$ 中的表示，也適用於各個自同構除環中的表示；這一點將在別處說明。\par}

\emph{證明.} 雙邊不可分解理想的分解
\[
        \mo\Omega=\ma_1+\ldots+\ma_s
\]
同中心分解為若干域
\[
        \mZ\Omega=\mZ_1+\ldots+\mZ_s
\]
一一對應，其中 $\mZ_i$ 是 $\ma_i$ 的中心。在 $\mo$ 的一個不可約表示中，一個 $\ma_i$ 被同構地表示，其餘的由零表示，而表示類由 $\ma_i$ 唯一確定。$\ma_i$ 被表示為一個全矩陣環，而這樣的全矩陣環的中心由對角矩陣 $E\xi$ 組成。對應 $z\mapsto\xi$ 是 $\mZ$ 的一個同態；其中一個 $\mZ_\nu$ 被同構地表示，其餘的由零表示。這樣一切都已證明。

關於什麼時候一個無根基的 $\mo$ 對應於同樣無根基的 $\mo_\Omega$，還可推出：

\emph{如果 $\mZ$ 有一個成分 $\mZ_i$ 是 $\mathsf{P}$ 的第二類擴張，那麼 $\mo\Omega$ 必定有根基。}\textsuperscript{21)} 因為在這種情形下，$\mZ_i\Omega$ 已經有根基。

{\footnotesize\noindent 21) 如果 $\mZ$ 只有第一類成分，那麼 $\mo\Omega$ 將沒有根基；這一點同樣將在後面證明。對於特徵為零的情形，I. Schur 已經證明了等價的定理：在 $\mathsf{P}$ 中的不可約表示，在係數域作任意擴張後仍保持完全可約（\emph{Beiträge zur Theorie der Gruppen linearer Substitutionen}, Transact. Am. Math. Soc. 15 (1909), p. 159）。\par}

\subsection*{§ 22. 在阿貝爾群上的應用}

設
\[
        G=G_1\times G_2\times\cdots\times G_r
\]
是一個有限阿貝爾群，分解為階為 $h_i$ 的循環群 $G_i$。其元素為
\[
        a=a_1^{\alpha_1}\ldots a_r^{\alpha_r}.
\]
設 $P$ 是一個域，其特徵不整除群的階
\[
        h=h_1h_2\cdots h_r.
\]

群環 $\mZ$ 由所有和
\[
        \sum_{\alpha_1,\ldots,\alpha_r}
        A_{\alpha_1\ldots\alpha_r}a_1^{\alpha_1}\cdots a_r^{\alpha_r},
        \qquad A_{\alpha_1\ldots\alpha_r}\in P,
\]
組成，並且藉助 $x_i\mapsto a_i$ 是多項式環 $\mathsf{P}[x_1,\ldots,x_r]$ 的一個同態像。因此
\[
        \mZ\simeq P[x_1,\ldots,x_r]/\mathfrak m,
        \qquad
        \mathfrak m=(x_1^{h_1}-e,\ldots,x_r^{h_r}-e).
\]
在擴域 $\Omega$ 中，$z_i^{h_i}-e$ 分解為互不相同的因子 $z_i-\varepsilon_i$；於是 $\mathfrak m$ 成為互不相同的素理想
\[
        (x_1-\xi_1^{(\alpha_1)},\ldots,x_r-\xi_r^{(\alpha_r)})
\]
的乘積（或者交），每個素理想都只有一個零點 $(\xi_1^{(\alpha_1)},\ldots,\xi_r^{(\alpha_r)})$。與取交相對應的是一個直接和分解（§4）：
\[
        \mZ\Omega=\mZ_1+\cdots+\mZ_h,
\]
其中每個 $\mZ_i\sim\Omega$，從而給出表示
\[
        a_i\mapsto \xi_i^{(\alpha_i)},
        \qquad\hbox{或更一般地}\qquad
        a\mapsto \chi^{(\alpha)}(a).
\]
這些表示就是特徵標。它們的個數是
\[
        h_1h_2\cdots h_r=h,
\]
這正如一般理論所要求的那樣。

\subsection*{§ 23. 超複系統的行列式}

設
\[
        \mo=a_1P+\cdots+a_hP
\]
是一個超複系統。我把 $h$ 個不定元 $x_1,\ldots,x_h$ 添加到 $\mathsf{P}$ 中，並形成
\[
        \mo^*=a_1\mathsf{P}(x)+\ldots+a_h\mathsf{P}(x).
\]
這些 $x$ 應當同 $a_i$ 可交換；於是 $\mo^*$ 中的計算規則已經由此確定。

在 $\mo^*$ 中有“$\mo$ 的一般元素”
\[
        w=a_1x_1+\ldots+a_hx_h.
\]
如果在一個表示中 $a_i\to A_i$，那麼同 $w$ 對應的是
\[
        W=A_1x_1+\ldots+A_hx_h.
\]
$W$ 稱為屬於該表示的系統矩陣（或者在特殊情形下，如果 $a_i$ 構成一個群，因而 $\mo$ 是群環，則稱為群矩陣）。若所討論的是正則表示，則得到正則系統矩陣。

$W$ 的元素 $w_{ij}$ 是 $x$ 的線性形式。因此，當表示為 $n$ 階表示時，“系統行列式” $|W|$ 的次數為 $n$。特別地，正則系統行列式的次數為 $h$。

系統行列式在過渡到等價表示時不改變，因為
\[
        |PWP^{-1}|=|P|\,|W|\,|P^{-1}|=|W|.
\]
在從 $(a_1,\ldots,a_h)$ 過渡到新基 $(b_1,\ldots,b_h)$，並從 $w=\sum a_ix_i$ 過渡到 $w=\sum b_jy_j$ 時，可以通過在舊元素中施行一個帶正則替換矩陣的替換
\[
        x_i=\sum_j\varrho_{ij}y_j
\]
來得到 $W$ 的新元素，從而也得到新的行列式。

若表示模有一條合成列，則在適當選取基後，矩陣 $W$ 具有如下形式：
\[
        W=\begin{pmatrix}
        W_1&0&\cdots&0\\
        *&W_2&\cdots&0\\
        \vdots&\vdots&\ddots&\vdots\\
        *&*&\cdots&W_s
        \end{pmatrix}.
\]
任意表示中出現的各個 $|W_i|$，不會超出 $\mo$ 的正則表示中出現的那些；甚至不會超出 $\mo/\mc$ 的正則表示中出現的那些，其中 $\mc$ 是根基。

如果 $\mathsf{P}$ 是代數閉的，那麼屬於一個不可約表示的行列式 $|W_i|$ 是 $x$ 中的一個素函數；而屬於不等價表示的，是不同的素因子。

\emph{證明.} 因為 $\mo$ 的所有不可約表示也都是 $\mo/\mc$ 的表示，所以我們可以限於無根基的環 $\mo/\mc$。在其中把矩陣單位 $c_{ik}$ 作為基引入；於是一般元素為
\[
        w=\sum_{\nu,i,k}c_{ik}^{(\nu)}x_{ik}^{(\nu)}.
\]
不可約表示的矩陣為
\[
        W_\nu=(x_{ik}^{(\nu)}).
\]
函數 $|W_\nu|=|x_{ik}^{(\nu)}|$ 眾所周知是不可約的，並且顯然彼此不同。

為了計算各個 $|W_\nu|$，可以把 $\mo/\mc$ 的正則系統矩陣分解為其素因子。每個素因子出現的次數，同相應不可約表示的階數一樣多，因為相應理想在合成列中也出現這麼多次。

也可以以 $\mo$ 的正則系統矩陣為基礎；但這時每個不可約因子出現得更頻繁，確切地說，出現次數等於相應單理想作為左理想中的合成因子出現的次數。在這個正則表示中，$\mo$ 被看作左理想；若把 $\mo$ 看作右理想，則得到第二個正則系統矩陣（Frobenius 所說的“antistrophe 矩陣”），它含有同樣的不可約因子（即所有不可約表示的系統行列式），但可能帶有不同的指數（見 §10 的例子）。

交換系統的系統行列式分解為線性因子，因為所有不可約表示都是一階的。這些線性因子本身就是不可約表示；因此在阿貝爾群的情形中，它們給出特徵標。這一事實正是 Dedekind 研究非阿貝爾群的群行列式時的出發點。

\subsection*{§ 24. 跡與特徵標}

在超複系統 $\mo$ 的一個表示 $\mo\sim\mathfrak D$ 中，如果元素 $a$ 對應於矩陣 $A$，則定義
\[
        \operatorname{Sp}_{D}(a)=\operatorname{Sp} A.
\]
跡是線性函數：
\[
        \mathit{Sp}(c+d)=\mathit{Sp}(c)+\mathit{Sp}(d),
        \qquad
        \mathit{Sp}(c\alpha)=\alpha\mathit{Sp}(c).
\]
等價表示具有相同的跡。

在一個可約表示中的跡，等於由合成因子中介的各個表示中的跡之和。

\emph{證明.} 在適當選取基時，矩陣具有塊形式
\[
        \begin{pmatrix}
        A_{11}&0&\cdots&0\\
        *&A_{22}&\cdots&0\\
        \vdots&\vdots&\ddots&\vdots\\
        *&*&\cdots&A_{ss}
        \end{pmatrix},
\]
所以它的跡為
\[
        \operatorname{Sp}(A_{11})+\operatorname{Sp}(A_{22})+\cdots+\operatorname{Sp}(A_{ss}).
\]

主跡是指正則表示中的跡。約化跡是指在各個不同不可約表示中的跡之和。

若 $c$ 是最大冪零理想 $\mc$ 的元素，則對於每個表示都有 $\mathit{Sp}(c)=0$。

\emph{證明.} 只需通過 $\mo/\mc$ 的簡單左理想來考察表示；因為在任何其他表示中，只會出現這些合成因子，而跡按加法組合而成。但是 $\mathfrak{c}$ 零化 $\mo/\mc$ 的所有元素；因此 $\mc$ 中的每個 $c$ 都對應於零矩陣，所以 $\mathit{Sp}(c)=0$。

若 $\mo$ 是雙邊單的，並且 $\mathsf{P}$ 是代數閉的，也就是說 $\mo$ 是一個矩陣環
\[
        \mo=\sum c_{ik}P,
\]
則在 $\mo$ 的不可約表示中，元素
\[
        a=\sum c_{ik}\alpha_{ik}
\]
對應於矩陣 $(\alpha_{ik})$。因此
\[
        \operatorname{Sp}_{\rm red}(a)=\sum_i\alpha_{ii},
        \qquad
        \operatorname{Sp}_{\rm pr}(a)=n\,\operatorname{Sp}_{\rm red}(a)=n\sum_i\alpha_{ii}.
\]
特別地，
\[
        \operatorname{Sp}_{\rm red}(c_{ik})=0\quad(i\ne k),\qquad
        \operatorname{Sp}_{\rm red}(c_{ii})=e,
\]
並且
\[
        \operatorname{Sp}_{\rm pr}(c_{ik})=0\quad(i\ne k),\qquad
        \operatorname{Sp}_{\rm pr}(c_{ii})=n e.
\]

從 $\mathsf{P}$ 過渡到代數閉域 $\Omega$ 時，主跡不改變，因為它可以用同一組基來計算。約化跡則並非如此。

無根基系統的元素 $a$ 在絕對不可約表示中的跡稱為特徵標，記作 $\chi(a)$；如果需要指出所指的是哪一個表示，則記作 $\chi^{(\nu)}(a)$。

在一個 $n_\nu$ 階不可約表示中，中心元素按照 §21 由對角矩陣 $E\zeta^{(\nu)}$ 表示，其中 $z\to\zeta^{(\nu)}$ 是中心的一階表示，或者說是中心到 $\Omega$ 的一個同態。由此可知跡的值為 $n_\nu\zeta^{(\nu)}$。因此，中心的同態 $\Theta$ 與特徵標 $\chi$ 由關係
\[
        \chi^{(\nu)}(z)=n_\nu\theta^{(\nu)}(z)                 \tag{1}
\]
聯繫起來。

在交換情形中，$n_\nu=1$，特徵標本身給出同態（參見 §22）。

如果域 $\Omega$ 的特徵為零，而我們在下文總將這樣假定，那麼可以把 (1) 除以 $n$：
\[
        \theta^{(\nu)}(z)=\frac{\chi^{(\nu)}(z)}{n_\nu}.
\]
$\theta^{(\nu)}$ 的同態性質由下列公式表達：
\[
        \frac{\chi^{(\nu)}(z+z')}{n_\nu}
        =
        \frac{\chi^{(\nu)}(z)}{n_\nu}
        +
        \frac{\chi^{(\nu)}(z')}{n_\nu},
\]
\[
        \frac{\chi^{(\nu)}(zz')}{n_\nu}
        =
        \frac{\chi^{(\nu)}(z)}{n_\nu}
        \frac{\chi^{(\nu)}(z')}{n_\nu}.
\]

一個表示類僅由矩陣的跡就已經唯一確定。要知道所有矩陣的跡，當然只需知道給定表示中基元素的跡。

\emph{證明.} 如果知道每個不可約表示模 $\mM_\nu$ 作為直和項在表示模 $R$ 中出現的次數，那麼 $R$ 就已確定。若這個次數為 $\mu_\nu$，則 $e^{(\nu)}$ 在該表示中的跡等於 $p_\nu n_\nu$，其中 $n_\nu$ 是不可約表示的階。因此
\[
        \mu_\nu=\frac{\operatorname{Sp}(e^{(\nu)})}{n_\nu}.
\]

\subsection*{§ 25. 判別式}

設
\[
        \mo=a_1P+\cdots+a_hP
\]
是一個超複系統。所謂判別式矩陣，是指其元素為主跡 $\operatorname{Sp}(a_i a_j)$ 的矩陣。所謂約化判別式矩陣，是指在代數閉域中由約化跡構成的相同矩陣。該矩陣的行列式分別稱為判別式和約化判別式。

在擴張基域時，判別式保持不變。過渡到另一組基時，它乘以變換行列式的平方。

\emph{證明.} 設
\[
        (b_1,\ldots,b_h)=(a_1,\ldots,a_h)P.
\]
則有
\[
        (\operatorname{Sp}(a_i a_j))=(\operatorname{Sp}(a_i b_j))P.             \tag{1}
\]
同樣，若 $P^t$ 為轉置矩陣，則
\[
        (\operatorname{Sp}(a_i b_j))=P^t(\operatorname{Sp}(b_i b_j)).           \tag{1a}
\]
由 (1) 和 (1a) 得
\[
        (\operatorname{Sp}(a_i a_j))=P^t(\operatorname{Sp}(b_i b_j))P.
\]
取行列式後即得所述結論。

因此，判別式只是在差一個來自 $P$ 的平方因子意義下確定；但是它消失或不消失是一種不變性質。由 (1) 還可推出：判別式也可以在差一個非零因子意義下由兩組不同的基來確定，即作為 $|\operatorname{Sp}(a_i b_j)|$ 來確定。

如果 $\mo$ 含有一個冪零理想 $\mc$，則判別式消失。

\emph{證明.} 選取
\[
        c_1,\ldots,c_r,d_1,\ldots,d_s
\]
作為 $\mo$ 的基元素，其中 $c_1,\ldots,c_r$ 構成 $\mc$ 的一組基。於是判別式矩陣具有形式
\[
        \begin{pmatrix}
        \operatorname{Sp}(c_i c_j)&\operatorname{Sp}(c_i d_j)\\
        \operatorname{Sp}(d_i c_j)&\operatorname{Sp}(d_i d_j)
        \end{pmatrix}
        =
        \begin{pmatrix}
        0&0\\
        0&\mathit{Sp}(d_i d_k)
        \end{pmatrix}.
\]
（這對於約化判別式同樣成立。）所以行列式消失。

推論. 如果過渡到代數閉擴域之後出現一個冪零理想，那麼判別式也消失。

設
\[
        \mo=\ma_1+\ldots+\ma_s
\]
是雙邊直和，並令 $M,M_1,\ldots,M_s$ 分別為環 $\mo,\ma_1,\ldots,\ma_s$ 的判別式矩陣。於是
\[
\begin{pmatrix}
he&&&0\\
&he&&\\
&&\ddots&\\
0&&&he
\end{pmatrix}.
\]
因而 $|M|=|M_1|\cdots|M_s|$。

\emph{證明.} 選取一組由 $\ma_1,\ldots,\ma_s$ 的各組基合成的 $\mo$ 的基。若 $a\in\ma_i$，則
\[
        \mathit{Sp}_{\mo}(\ma)=\mathit{Sp}_{\ma_i}(a),
\]
這裡兩邊所說的都是主跡，但分別在環 $\mo$ 與 $\ma_i$ 中取。再者，當因子位於不同的雙邊組成部分中時，$a_i a a_j=0$；因此混合跡也消失。由此得到斷言；同一結論也適用於約化判別式。

為了形成矩陣環 $\sum c_{ik}\mathsf{P}$ 的判別式，選取兩組不同的基，即把 $c_{ik}$ 一次按第一指標排列，一次按第二指標排列。乘法表如下：
\[
\begin{array}{c|c|c|c}
& \overbrace{c_{11}\;\ldots\;c_{1n}}^{\mathfrak r_1}
& \overbrace{c_{21}\;\ldots\;c_{2n}}^{\mathfrak r_2}
& \ldots \\ \hline
\mathfrak l_1\left\{\begin{array}{c}c_{11}\\[-1mm]\vdots\\[-1mm]c_{n1}\end{array}\right.
& (c_{ik})&0&0\\[1.3em]\hline
\mathfrak l_2\left\{\begin{array}{c}c_{12}\\[-1mm]\vdots\\[-1mm]c_{n2}\end{array}\right.
&0&(c_{ik})&0\\[1.3em]\hline
\vdots&0&0&(c_{ik})
\end{array}
\]
判別式為
\[
        D=\left|\begin{array}{cccc}
        \mathit{Sp}(c_{11})& & &0\\
        &\mathit{Sp}(c_{22})& & \\
        & &\ddots& \\
        0& & &\mathit{Sp}(c_{nn})
        \end{array}\right|^{n}
        =\begin{cases}
        e, & \hbox{对于约化迹},\\
        n^{n^2}\cdot e, & \hbox{对于主迹}.
        \end{cases}
\]

推論. 若 $\mo$ 在代數閉擴域中分解為次數為 $n_1,\ldots,n_s$ 的矩陣環，則判別式為
\[
        D=n_1^{n_1^2}n_2^{n_2^2}\ldots n_s^{n_s^2}\cdot e
\]
而約化判別式為
\[
        D_{\rm red}=e.
\]
因此，對於沒有根基的系統，約化判別式總是 $\ne0$；而通常的判別式只有在沒有任何 $n_i$ 被域的特徵整除時才如此。於是：

\emph{約化判別式不消失，是系統無根基（在域 $\mathsf{P}$ 的代數閉包之後）的必要且充分條件。}

\emph{在特徵為零的情形下，判別式不消失，是不出現根基（在域 $\mathsf{P}$ 的代數閉包之後）的必要且充分條件。}\textsuperscript{22)}

{\footnotesize\noindent 22) 對於交換系統，參見 E. Noether, \emph{Diskriminantensatz für Ordnungen ...}, J. f. M. 157 (1927), pp. 82--104, §§4--6。那裡的證明方法相同，但細節上更為複雜；從一組基到另一組基的過渡（那裡 §4, 4.）應由本段開頭給出的過渡取代，因為那裡出現的行列式會消失。\par}

\subsection*{§ 26. 群環的定位}

設 $a_1,\ldots,a_h$ 是一個有限群的元素，並在一個特徵不整除 $h$ 的域中形成群環。那麼判別式 $D\ne0$，因而該群環是一個無根基的環。

\emph{證明.} 我們先證明如下事實，其中 $\mathit{Sp}$ 從現在起表示主跡：
\[
        \operatorname{Sp}(e)=he,
        \qquad
        \operatorname{Sp}(a_i)=0\quad(a_i\ne e).
\]
事實上，在正則表示中有
\[
        e\mapsto
        \begin{pmatrix}
        e&&&0\\
        0&e&\cdots&0\\
        &&\ddots&\\
        0&&&e
        \end{pmatrix},
        \qquad \operatorname{Sp}(e)=he.
\]
若 $a_i\ne e$，則 $a_i a_1,\ldots,a_i a_h$ 是基元素的一個無不動點排列；相應矩陣的主對角線上只有零，因此 $\mathit{Sp}(a_i)=0$。

現在取兩組基 $a_1,\ldots,a_h$ 與 $a_1^{-1},\ldots,a_h^{-1}$。矩陣 $(\mathit{Sp}\operatorname{Sp}(a_i a_j^{-1}))$ 是以 $he$ 為主對角元素的對角矩陣，因而
\[
        D=h^h e\ne0.
\]

所謂一個群元素的類，是指在群環中形成的和
\[
        K_a=\sum s^{-1}as,\tag{1}
\]
其中只對互不相同的 $s^{-1}as$ 求和。這些類 $K_a$ 是中心元素，因為它們同所有群元素可交換。它們生成中心；因為如果群環中的一個元素 $\sum_i\alpha_i a_i$ 同所有 $s$ 可交換，那麼
\[
        \sum_i\alpha_i a_i=s^{-1}\biggl(\sum_i\alpha_i a_i\biggr)s
        =\sum_i\alpha_i(s^{-1}a_i s),
\]
所以每個共軛類上的係數相等。

因此，中心的秩等於共軛群元素的類數，從而絕對不可約表示的數目也等於這個類數。矩陣 $A$ 與 $S^{-1}AS$ 具有相同的跡；所以群元素 $a$ 與 $s^{-1}as$ 具有相同的跡。對 (1) 兩邊取跡，即得
\[
        \chi(K_a)=h_a\chi(a),
\]
其中 $h_a$ 是類 $K_a$ 中元素的數目。用文字來說：一個群元素的特徵標等於該類的特徵標除以這個類中元素的數目。

\begin{center}
（1928 年 8 月 12 日收稿。）
\end{center}

%
\clearpage
%
\setcounter{footnote}{0}

%
%
%
%
%
\pdfbookmark[0]{35. 整係數函數的極大域}{paper-35}
\section*{35. 整係數函數的極大域}

\begin{center}
\textbf{埃米·諾特（哥廷根）}\\[0.3em]
\emph{Rec. Soc. Math. Moscou \textbf{36} (1929), S. 65--72}
\end{center}

以下研究與 Hilbert 關於相對整函數的問題\footnote{D.\ Hilbert, Mathematische Probleme (Gott.\ Nachr.\ 1900, S.\ 253--297), Problem 14.}密切相接；它們源自 E.\ Hecke 偶然提出的一個問題。他需要其中一個特殊情形作為輔助引理，並且能夠直接用計算處理它\footnote{E.\ Hecke, Über die Konstruktion relativ-Abelscher Zahlkörper durch Modulfunktionen von zwei Variabeln (Math.\ Ann.\ 74 (1913), S.\ 465--510), \S\,2. 在 Hecke 那裡，有限整域 $\mathfrak{I}$ 由兩個變量中三個冪級數的商生成，並且這些商之間存在一個代數關係。}。

問題的表述如下。所謂給定區域 $\mathfrak{Z}$ 內、關於 $x_1,\ldots,x_n$ 的整係數函數（即多項式、冪級數或它們的商）所成的一個\textit{極大域}，我指的是這樣一個區域：它不能再通過添入整數分母而擴大；也就是說，若 $a\cdot g(x)$ 屬於該區域，其中 $g(x)$ 屬於 $\mathfrak{Z}$，則總推出 $g(x)$ 也屬於該區域；這裡 $a$ 是一個非零有理整數。由 $\mathfrak{Z}$ 中每一個整域 $\mathfrak{I}$ 都唯一地生成一個極大域 $M(\mathfrak{I})$；它由 $\mathfrak{Z}$ 中所有函數 $g(x)$ 組成，條件是至少存在一個 $a\neq0$，使得 $a\cdot g(x)$ 屬於 $\mathfrak{I}$。

現在要問的是：如果原來的整域 $\mathfrak{I}$ 是有限的，也就是說，它由 $\mathfrak{Z}$ 中有限多個函數 $f_1(x),\ldots,f_r(x)$ 的所有整係數多項式組成，那麼關於這個極大域 $M(\mathfrak{I})$ 的分母 $a$ 能說些什麼。這裡 $\mathfrak{Z}$ 應取為 $x_1,\ldots,x_n$ 中的所有整係數多項式或冪級數，或者取為這樣的商：它們除了 $f(x)$ 中已經出現的分母及其冪以外不含其他分母。

我證明：在冪級數或其商的情形中，附加假設 $f(x)$ 之間只有代數依賴關係出現，則這些 $a$ 總可以選成只含有限多個不同素數作為因子；不過，一般說來，這些素數可以以任意高的冪出現。

後一點立即可見：如果 $a\cdot g(x)$ 屬於 $\mathfrak{I}$，但對於任何指數 $\varrho$，$a^\varrho g(x)$ 都不屬於 $\mathfrak{I}$，那麼所有冪 $a^\varrho$ 都會出現（例如 $\mathfrak{I}=[2x]$；$x=(2x)/2$，$x^2=(2x)^2/2^2$）。

關於有界性的問題只能這樣表述：能否給出 $M(\mathfrak{I})$ 的一個通常為無限的生成元系，使得相應分母中各素數的指數保持有界？由於這裡只涉及有限多個素數，按照已知推理\footnote{參見 Hilbert 前引文獻，或 E.\ Noether, Die Endlichkeit des Systems der ganzzahligen Invarianten binärer Formen (Gött.\ Nachr.\ 1919, S.\ 138--156), \S\,5 中更詳細的敘述。}，這個問題等同於 $M(\mathfrak{I})$ 的有限性問題，也就是整係數形式下 Hilbert 關於相對整函數的問題；用本文采用的方法，我不能判定這個問題\footnote{因此，例如在一個基本型系統的整係數（射影）不變量情形中，這裡證明的是：它們能夠由通常意義下的一個完整不變量系統表示，只是在分母中只出現有限多個不同素數；這個結果並不由通常的 $\varOmega$-過程方法推出。但是，關於整係數有限性本身並沒有得到任何結論；此前只在二元基本型的情形中知道這一點。參見注 3 所引的短文。}。

證明只有有限多個不同素數 $p$ 會在分母 $a$ 中出現，其根據是：每一個這樣的 $p$ 至少對應於由 $\mathfrak{I}$ 的生成函數 $f(x)$ 之間的一個模 $p$ 關係，而這個關係並不是以整係數在 $x$ 中恆等成立的關係。換言之，若 $\mathfrak{o}$ 表示不定元 $u_1,\ldots,u_r$ 中的整係數多項式環，則藉助對應 $u_i\mapsto f_i(x)$，$\mathfrak{I}$ 成為 $\mathfrak{o}$ 的同態像；這個同態由 $\mathfrak{o}$ 中的一個素理想 $\mathfrak{a}$ 給出（也就是說，$\mathfrak{I}$ 同構於剩餘類環 $\mathfrak{o}/\mathfrak{a}$）。相應地，$\mathfrak{I}_p$ 成為 $\mathfrak{o}_p$ 的同態像；這裡下標表示轉到模 $p$ 剩餘類，這在排除 $\mathfrak{I}$ 的分母中出現的有限多個素數後可以進行。這個同態又由 $\mathfrak{o}_p$ 中的一個素理想 $\mathfrak{c}_p$ 給出，且 $\mathfrak{c}_p$ 包含 $\mathfrak{a}_p$。於是，當且僅當 $\mathfrak{c}_p$ 成為 $\mathfrak{a}_p$ 的真上理想（理想意義下的真因子）時，$p$ 才會在 $M(\mathfrak{I})$ 的分母中出現。這只能在兩種情況下發生：或者 $\mathfrak{I}_p$ 的超越度相對於 $\mathfrak{I}$ 降低，或者 $\mathfrak{a}_p$ 失去其素理想性質\footnote{所謂一個系統的“超越度”，通常指該系統代數依賴於其上的代數獨立函數的個數；由代數獨立函數組成的系統稱為不可約系統。所謂一個素理想的超越度（維數），指的是它的剩餘類域的超越度。一個素理想沒有同一超越度的真素因子；參見例如 B.\,L.\,v.\,d.\ Waerden, Zur Nullstellentheorie der Polynomideale, Math.\ Ann.\ 96 (1926), S.\ 183--208。}。這些情形只會對有限多個 $p$ 發生，其理由是：即使模 $p$，一個代數依賴關係也導致函數行列式消失（\S\,1）；並且進一步可知 $\mathfrak{a}$ 是絕對素理想，所以它只會模有限多個 $p$ 失去素理想性質\footnote{E.\ Noether, Eliminationstheorie und allgemeine Idealtheorie, Math.\ Ann.\ 90 (1923), S.\ 229--261, \S\,7。該定理通過消去理論\textsuperscript{*)}由較簡單的定理推出：一個絕對不可約多項式只會模有限多個素數失去這個性質；見 A.\ Ostrowski, Zur arithmetischen Theorie der algebraischen Größen, Gött.\ Nachr.\ 1919, S.\ 279--298, 特別是 S.\ 296。較簡單的證明見 E.\ Noether, Ein algebraisches Kriterium für absolute Irreduzibilität, Math.\ Ann.\ 85 (1922), S.\ 26--33, \S\,7。如果 $\mathfrak{I}$ 有一個整基，其包含的函數個數恰好比其超越度多一個，正如 Hecke（注 2）所考察的情形，那麼 $\mathfrak{a}$ 成為主理想，此時只需引用這個較簡單的定理，而不必使用消去理論。到處都可以用一個（有限）代數數域的素理想代替素數。顯然，當 $\mathfrak{a}$ 是零理想時，該定理也成立。}。

還要指出，如果用一個（有限）代數數域的整數及其素理想來代替有理整數與素數，那麼全部考慮只需作很小的修改仍然成立（\S\,4）。

% -----------------------------------------------------------
\bigskip
\S~\textbf{1.}
\textbf{特徵為 $p$ 的係數域中的函數行列式。}

\smallskip

\noindent\textbf{1.}
設 $f_1(x_1,\ldots,x_n),\ldots,f_t(x_1,\ldots,x_n)$ 是有理函數，或更一般地，是形式定義的冪級數或這些冪級數的商，其係數屬於一個特徵為 $p$ 的域 $P$。導數 $\partial f_i/\partial x_k$ 也按形式定義，並且所有通常的運算法則保持有效。用 $\varOmega$ 表示 $\mathrm{P}$ 的一個代數閉擴域。

於是與特徵為零時一樣，有：

\noindent\textbf{定理。}
\textit{如果 $f(x)$ 之間存在一個係數屬於 $\varOmega$ 的代數依賴關係，那麼函數矩陣 $(\partial f_i/\partial x_k)$（$i=1,\ldots,t$；$k=1,\ldots,n$）的秩小於 $t$}\footnote{在原來的表述中，我曾假定係數域 $P$ 是完美域；通過從 $P$ 過渡到 $\varOmega$，把它推廣到任意域、包括非完美域，應歸功於 B.\,L.\,v.\,d.\ Waerden。在特徵 $p$ 時，即使 $P$ 是完美域，逆命題也不成立；例子是 $f=x^p$，此時 $\partial f/\partial x=0$。}。

\noindent\textit{證明。}
在多項式環 $\varOmega[u_1,\ldots,u_t]$ 中，令 $\mathfrak{m}$ 表示所有滿足 $G(f)=0$（在 $x$ 中恆等）的多項式 $G(u)$ 所成的素理想。按假設，$\mathfrak{m}$ 不是零理想；取 $\mathfrak{m}$ 中一個次數最低的多項式 $G(u)\neq0$（因而也不是常數），由此特別推出 $G(u)$ 不可分解。於是像通常一樣得到：
\[
\left(\frac{\partial G}{\partial u_1}\right)_f\cdot \frac{\partial f_1}{\partial x_k}+\cdots+
\left(\partial G/\partial u_t\right)_f\partial f_t/\partial x_1=0
\quad (k=1,\ldots,n),
\]
由此可知：或者函數矩陣的秩小於 $t$，或者所有 $(\partial G/\partial u_i)_f$ 都消失。但是，由於 $G(u)$ 已被假定為次數最低的多項式，後一種情況推出所有 $\partial G/\partial u_i$ 都消失。

因此有
\[
        G(u)=H(u_1^p,\ldots,u_t^p).
\]
又因為係數域 $\varOmega$ 被假定為代數閉，因而是完美域，進一步有
\[
        G(u)=H(u_i^p)=(T(u))^p,
\]
這與 $G(u)$ 的選擇矛盾。這個證明當然也適用於特徵為零的情形，只是最後一步不再需要。

\medskip\noindent\textbf{2.} \textbf{推論。}

設 $F_1(x_1,\ldots,x_n),\ldots,F_t(x_1,\ldots,x_n)$ 是整係數函數（多項式、冪級數或這些函數的商），其函數矩陣的秩恰好為 $t$。素數 $p$ 選得使它既不出現在 $F_i(x)$ 的分母中，也不整除函數矩陣的全部 $t$ 階行列式。

那麼 $F_1(x),\ldots,F_t(x)$ 在模 $p$ 後仍然代數獨立。

\noindent\textit{證明。}
令 $\bar{P}$ 表示模 $p$ 的剩餘類域，並令 $f_1(x),\ldots,f_t(x)$ 表示由 $F_1(x),\ldots,F_t(x)$ 把整係數替換為它們模 $p$ 的剩餘類而得到的函數。由於 $p$ 不在 $F_i(x)$ 的分母中出現，這些 $f_i(x)$ 有定義；它們的函數矩陣 $(\partial f_i/\partial x_k)$ 也是由 $(\partial F_i/\partial x_k)$ 把整係數替換為模 $p$ 剩餘類而得到的，因為在 $(\partial F_i/\partial x_k)$ 中不出現 $F_i(x)$ 之外的其他分母。由於 $p$ 的選擇，$(\partial f_i/\partial x_k)$ 保持秩 $t$；因此 $f_i(x)$ 在 $\varOmega$ 中代數獨立，從而更在 $\bar{P}$ 中代數獨立。

% -----------------------------------------------------------
\bigskip
\S~\textbf{2.}
\textbf{定理的表述。}

\smallskip

\noindent\textbf{1.}
如引言中那樣，令 $\mathfrak{I}$ 表示一個關於 $x_1,\ldots,x_n$ 的整係數函數所成的整域（無零因子環），這些函數可以是有理整數係數的多項式或冪級數，也可以是這些函數的商。在冪級數的情形中，它們可以是形式定義的，也可以是在某個區域內收斂的；所用到的只是這些函數組成一個整域。

用 $\mathfrak{Q}$ 表示 $\mathfrak{I}$ 的商域；$\mathfrak{Z}$ 表示介於 $\mathfrak{I}$ 與 $\mathfrak{Q}$ 之間的一個整域。如引言中所說，若從 $a\cdot g(x)$ 屬於 $M$（其中 $a$ 是非零整數，$g(x)$ 屬於 $\mathfrak{Z}$）總能推出 $g(x)$ 屬於 $M$，則稱 $M$ 在 $\mathfrak{Z}$ 中是\textit{極大的}。對於 $\mathfrak{I}$，存在唯一的由 $\mathfrak{I}$ 在 $\mathfrak{Z}$ 中生成的極大域 $\mathfrak{M}(\mathfrak{I})$，它是所有包含 $\mathfrak{I}$ 的 $\mathfrak{Z}$ 中極大域的交；因為 $\mathfrak{Z}$ 本身就是這樣的一個域，並且任意多個這種極大域的交仍然極大且包含 $\mathfrak{I}$。$M(\mathfrak{I})$ 由 $\mathfrak{Z}$ 中所有函數 $g(x)$ 組成，條件是至少存在一個 $a\neq0$，使得 $a\cdot g(x)$ 屬於 $\mathfrak{I}$。

事實上，這些 $g(x)$ 組成 $\mathfrak{Z}$ 中一個包含 $\mathfrak{I}$ 的極大域 $N$，因為若 $b\cdot h(x)$ 屬於 $\mathfrak{N}$，其中 $b\neq0$ 且 $h(x)$ 屬於 $\mathfrak{Z}$，則 $ab\cdot h(x)$ 屬於 $\mathfrak{I}$ 且 $ab\neq0$，所以 $h(x)$ 屬於 $N$。另一方面，$\mathfrak{Z}$ 中任一包含 $\mathfrak{I}$ 的極大域 $M$ 都包含 $N$；因為若 $a\cdot g(x)$ 屬於 $\mathfrak{I}$，則由於 $\mathfrak{I}$ 包含於 $M$，$a\cdot g(x)$ 屬於 $M$，於是 $g(x)$ 屬於 $M$。

\medskip\noindent\textbf{2.}
現在假定 $\mathfrak{I}$ 是帶單位元的有限整域，以 $f_1(x),\ldots,f_r(x)$ 為整基；也就是說，$\mathfrak{I}$ 由 $f_i(x)$ 的所有整係數多項式組成。再令 $\mathfrak{Z}$ 由 $x$ 中所有整係數多項式或冪級數組成，或由這樣的商組成：它們不含除有限多個 $f_i(x)$ 中已經出現的分母以外的其他分母；這些分母當然可以任意次冪出現。

還假定至少在一個 $f_i(x)$ 中，通常是在所有 $f_i(x)$ 中，$x$ 真正出現；否則 $\mathfrak{I}$ 與 $\mathfrak{Z}$ 只是某個分數 $1/c$ 的所有整係數多項式所成的環，沒有什麼需要證明。

如果所涉及的是多項式或有理函數，那麼商域 $\mathfrak{Q}$ 關於有理數域 $\mathbf{K}$ 的超越度為 $t$，其中 $1 \leqq t \leqq r$；而如果例如 $f_1(x),\ldots,f_t(x)$ 組成一個不可約系統\footnote{所謂“不可約系統”，通常指由代數獨立函數組成的系統。}，則 $\mathfrak{Q}$ 是純超越擴張 $\mathrm{K}(f_1(x),\ldots,f_t(x))$ 的有限代數擴張。同時，$f_1(x),\ldots,f_r(x)$ 的函數矩陣以及 $f_1(x),\ldots,f_t(x)$ 的函數矩陣都恰好具有秩 $t$\footnote{這個附加前提在如下情形中滿足：正如 Hecke 所考察的情形（注 2）那樣，若所涉及的是解析函數，且 $f_1(x),\ldots,f_t(x)$ 解析獨立，而 $f_{t+1}(x),\ldots,f_r(x)$ 代數依賴於它們。這裡有意把前提表述為形式化的“函數矩陣的秩”，以便它在形式定義的冪級數情形中也有清楚意義。這樣就使全部考慮保持在純形式代數範圍內；解析函數的定理只在特殊情形中用來證明該前提確實滿足。}。

為了在冪級數或其商的情形中也保持 $\mathfrak{Q}$ 的這種結構，還要作一個附加前提：若 $t$ 是 $f_i(x)$ 的函數矩陣的秩，並且例如同時也是 $f_1(x),\ldots,f_t(x)$ 的函數矩陣的秩，從而按 \S\,1, 1.\ 它們組成一個不可約系統，則 $\mathfrak{Q}$ 仍應是純超越擴張 $\mathbf{K}(f_1(x),\ldots,f_t(x))$ 的有限代數擴張；也就是說，$f_{t+1}(x),\ldots,f_r(x)$ 應相對於 $\mathbf{K}(f_1(x),\ldots,f_t(x))$ 代數依賴\footnote{如果所涉及的是解析函數，並且如 Hecke 所考察的情形（注 2）那樣，$f_1(x),\ldots,f_t(x)$ 解析獨立，而 $f_{t+1}(x),\ldots,f_r(x)$ 代數依賴於它們，則這個附加前提滿足。}。這裡仍有 $t \geqq 1$，因為在特徵為零的情形中，不可能所有 $\partial f_i/\partial x_k$ 都消失。

同時，這個條件也對整基中任意 $t$ 個函數所組成的系統成立，只要它們的函數矩陣恰好秩為 $t$；因為這些函數組成一個不可約系統，其餘函數因而代數依賴於它們。

\medskip\noindent\textbf{3.}
現在令 $\mathfrak{o}$ 表示不定元 $u_1,\ldots,u_r$ 中的整係數多項式環。由對應 $u_i\mapsto f_i(x)$ 立即可見，$\mathfrak{I}$ 是 $\mathfrak{o}$ 的環同態像。因此 $\mathfrak{I}$ 環同構於 $\mathfrak{o}/\mathfrak{a}$，其中 $\mathfrak{a}$ 是素理想，因為 $\mathfrak{I}$ 是無零因子環；$\mathfrak{a}$ 由所有滿足 $G(f)$ 在 $x$ 中恆等為零的多項式 $G(u)$ 組成。

如引言已經提到，所要證明的是如下定理：

\medskip\noindent\textbf{定理。}
\textit{在 2.\ 的附加前提下，令 $\mathfrak{o}_p,\mathfrak{a}_p,\mathfrak{I}_p$ 表示分別由 $\mathfrak{o},\mathfrak{a},\mathfrak{I}$ 把整係數替換為其模 $p$ 剩餘類而得到的對象。那麼，除至多有限多個素數外，仍有：從 $\mathfrak{o}_p$ 到 $\mathfrak{I}_p$ 的同態由 $\mathfrak{a}_p$ 給出；這些例外素數包括 $f_i(x)$ 的分母中出現的素數，因為要使 $\mathfrak{I}_p$ 有定義。於是 $\mathfrak{I}_p$ 環同構於 $\mathfrak{o}_p/\mathfrak{a}_p$，也即同構於 $\mathfrak{o}/(\mathfrak{a},p\mathfrak{o})$。}

由這個定理，通過下面的推論便得到 $M(\mathfrak{I})$ 的分母 $a$ 中只出現有限多個素數這一事實。

\medskip\noindent\textbf{推論。}
\textit{若 $p$ 選得使 $\mathfrak{I}_p$ 同構於 $\mathfrak{o}_p/\mathfrak{a}_p$，那麼只要 $p\cdot g(x)$ 屬於 $\mathfrak{I}$ 且 $g(x)$ 屬於 $\mathfrak{Z}$，就必有 $g(x)$ 屬於 $\mathfrak{I}$。}

換言之，對於所有不同於有限多個例外素數的素數，$\mathfrak{I}$ 在 $\mathfrak{Z}$ 中是極大的。

有 $\mathfrak{o}_p/\mathfrak{a}_p$ 同構於 $\mathfrak{o}/(\mathfrak{a},p\mathfrak{o})$。因為按定義 $\mathfrak{o}_p$ 等於剩餘類環 $\mathfrak{o}/p\mathfrak{o}$，相應地，$\mathfrak{a}_p$ 等於 $(\mathfrak{a},p\mathfrak{o})/p\mathfrak{o}$。所以由第一同構定理，$\mathfrak{o}_p/\mathfrak{a}_p$，並且按假設 $\mathfrak{I}_p$，都同構於 $\mathfrak{o}/(\mathfrak{a},p\mathfrak{o})$。

這意味着：若在 $\mathfrak{I}$ 中 $H(f_i(x))\equiv0\pmod p$，則在 $\mathfrak{o}$ 中
\[
        H(u)\equiv0 \pmod{(\mathfrak{a},p\mathfrak{o})}.
\]

也就是說：若 $H(f(x))=p g(x)$，其中 $H(f(x))$ 屬於 $\mathfrak{I}$ 且 $g(x)$ 屬於 $\mathfrak{Z}$，則有
\[
        H(u)-pF(u)\equiv0\pmod{\mathfrak{a}},
\]
因而也有 $H(f_i(x))=pF(f_i(x))$，所以 $g(x)=F(f_i(x))$；於是 $g(x)$ 屬於 $\mathfrak{I}$。

最後，有限次重複給出：若 $a$ 不被任何例外素數整除，則從 $a\cdot g(x)$ 屬於 $\mathfrak{I}$ 且 $g(x)$ 屬於 $\mathfrak{Z}$ 總推出 $g(x)$ 屬於 $\mathfrak{I}$。推論由此得證。

% -----------------------------------------------------------
\bigskip
\S~\textbf{3.}
\textbf{定理的證明。}

\smallskip

\noindent\textbf{1.} \textbf{輔助引理：}
\textit{給出 $\mathfrak{I}$ 與 $\mathfrak{o}$ 之間同態的素理想 $\mathfrak{a}$（\S\,2, 3.）是絕對素理想。}

令 $\mathfrak{o}^*$ 表示 $u_1,\ldots,u_r$ 中、以任意代數數（也允許分數代數數）為係數的多項式環，令 $\mathfrak{a}^*$ 表示在 $\mathfrak{o}^*$ 中由 $\mathfrak{a}$ 生成的理想；也就是說，$\mathfrak{a}^*$ 由所有線性組合 $a_1G_1(u)+\cdots+a_sG_s(u)$ 組成，其中 $a_i$ 為任意代數數，而 $G_i(u)$ 是 $\mathfrak{a}$ 中的多項式。要證明 $a^*$ 是素理想。

只要證明 $a^*$ 給出了從 $\mathfrak{o}^*$ 到 $\mathfrak{I}^*$ 的同態，就完成了證明；這裡 $\mathfrak{I}^*$ 是由 $f_i(x)$ 的所有以任意代數數為係數的多項式組成的整域。換言之，要證明：若 $H(u)$ 屬於 $\mathfrak{o}^*$ 且 $H(f_i(x))$ 消失，則 $H(u)$ 屬於 $\mathfrak{a}^*$。

令 $\omega_1,\ldots,\omega_s$ 表示由 $H(u)$ 的有限多個係數生成的有限數域的一組線性無關基。於是有
\[
        a\cdot H(u)=\omega_1H_1(u)+\cdots+\omega_sH_s(u),
\]
其中 $H_1(u),\ldots,H_s(u)$ 屬於 $\mathfrak{o}$，而 $a\neq0$ 是一個有理整數。由 $H(f)=0$ 可得
\[
        \omega_1H_1(f)+\cdots+\omega_sH_s(f)=0,
\]
從而由線性無關性得到 $H_1(f)=0,\ldots,H_s(f)=0$。所以 $H_1(u),\ldots,H_s(u)$ 都屬於 $\mathfrak{a}$，因而 $H(u)$ 屬於 $\mathfrak{a}^*$。

\medskip\noindent\textbf{2.}
現在可以把定理證明寫成如下形式。

若 $p$ 選得滿足
\begin{itemize}
\item[I.] $\mathfrak{I}_p$ 有定義，
\item[II.] $\mathfrak{a}_p$ 保持為素理想（而且保持為絕對素理想），
\item[III.] $\mathfrak{I}_p$ 的超越度相對於 $\mathfrak{I}$ 的超越度沒有降低，
\end{itemize}
則從 $\mathfrak{o}_p$ 到 $\mathfrak{I}_p$ 的同態由 $\mathfrak{a}_p$ 給出。除至多有限多個例外素數外，所有素數都滿足這些條件。

關於條件 I、II、III 的例外素數只有有限多個，這一點由前文立即推出。

I 是清楚的：例外素數只是 $f_i(x)$ 的分母中出現的有限多個素數。II 由輔助引理推出：$\mathfrak{a}$ 是絕對素理想，所以除至多有限多個素數外，它模素數後仍保持這個性質（參見注 6）。III 最後由 \S\,1, 2.\ 的關於函數行列式的推論給出，因為按照 \S\,2, 2.，$f_1(x),\ldots,f_t(x)$ 的函數矩陣恰好具有秩 $t$。

現在選取一個不同於這些例外素數的 $p$。從 $\mathfrak{o}_p$ 到 $\mathfrak{I}_p$ 的同態由一個理想 $\mathfrak{c}_p$ 給出。它是素理想，因為 $\mathfrak{I}_p$ 也是無零因子環；它的超越度由 $\mathfrak{I}_p$ 的超越度給出，因而至少等於 $t$。這是因為 $\mathfrak{I}_p$ 同構於 $\mathfrak{o}_p/\mathfrak{c}_p$，所以 $\mathfrak{c}_p$ 的剩餘類域同構於 $\mathfrak{I}_p$ 的商域。

又因為 $\mathfrak{a}_p$ 也是素理想，並且是 $\mathfrak{c}_p$ 的倍理想（下理想），所以 $\mathfrak{a}_p$ 與 $\mathfrak{c}_p$ 當且僅當它們的超越度相同時相等\footnote{參見例如 B.\,L.\,v.\,d.\ Waerden, Zur Nullstellentheorie der Polynomideale, Math.\ Ann.\ 96 (1926), S.\ 183--208, \S\,3。}。作為 $\mathfrak{c}_p$ 的倍理想，$\mathfrak{a}_p$ 的超越度也至少等於 $t$；並且 $u_1,\ldots,u_t$ 的模 $\mathfrak{a}_p$ 類組成一個不可約系統。現在只需證明這個超越度恰好等於 $t$，於是全部證明就完成了。

這裡，在冪級數或其商的情形中，要用到關於 $\mathfrak{I}$ 的商域 $\mathfrak{Q}$ 的結構的附加前提（\S\,2, 2.）。按該前提，素理想 $\mathfrak{a}$ 在每種情形中都有超越度 $t$；並且編號已經選得使 $u_1,\ldots,u_t$ 的模 $\mathfrak{a}$ 類組成一個不可約系統，而 $u_{t+1},\ldots u_r$ 的類代數依賴於它們。因此，$\mathfrak{a}$ 中包含整係數且本原的多項式
\[
G_1(u_{t+1};u_1,\ldots,u_t),\ldots,G_{r-t}(u_r;u_1,\ldots,u_t) \tag{1}
\]
它們模 $p$ 後變為 $\mathfrak{a}_p$ 中的非零多項式
\[
G^*_1(u_{t+1};u_1,\ldots,u_t),\ldots,G^*_{r-t}(u_r;u_1,\ldots,u_t). \tag{2}
\]
非零性來自本原性。這裡，在每個 $G^*_j$ 中，相應的 $u_{t+j}$ 必須真正出現；否則，$u_1,\ldots,u_t$ 的模 $\mathfrak{a}_p$ 類之間就會存在一個代數依賴關係，這與 $p$ 的選擇相矛盾。於是已經證明 $\mathfrak{a}_p$ 的超越度為 $t$；因此它與 $\mathfrak{c}_p$ 相同，並給出該同態。證明至此完成。

% -----------------------------------------------------------
\bigskip
\S~\textbf{4.}
\textbf{推廣到整代數係數。}

\smallskip

\noindent
如引言末尾已經指出，如果不用有理整數，而以一個（有限）代數數域 $\mathbf{K}$ 的整數為基礎，並且不用素數 $p$，而用這個數域的素理想 $\mathfrak{p}$ 為基礎，那麼全部考慮經過很小修改仍然成立。其中 \S\,1 完全保持不變，導向 \S\,2 中定理表述的那些考慮也保持不變；只是在證明（\S\,3, 2.）以及 \S\,2, 3.\ 的推論中需要作若干補充。

\medskip\noindent\textbf{1.}
在 \S\,3, 2.\ 中，還要把有限多個素理想加入例外素理想之列：這些素理想出現在多項式 $G_1(u),\ldots,G_{r-t}(u)$ 中（見 \S\,3, 2., (1)），因為現在不能再預先假定這些多項式是本原的。這樣就保證了多項式 $G^*_1(u),\ldots,G^*_{r-t}(u)$（見 \S\,3, 2., (2)）不消失；這相當於對條件 III 的一個加強。條件 II 保持不變，因為關於絕對不可約性的定理沒有改變；當然條件 I 也保持不變。

\medskip\noindent\textbf{2.}
作為最終結果的 \S\,2, 3.\ 的推論，應依照那裡最後的形式改寫如下：

\noindent\textbf{推論。}
\textit{若 $\mathbf{K}$ 中的整數 $a$ 不被有限多個例外素理想中的任何一個整除，則從 $a\cdot g(x)$ 屬於 $\mathfrak{I}$ 且 $g(x)$ 屬於 $\mathfrak{Z}$，總推出 $g(x)$ 屬於 $\mathfrak{I}$。}

若把 $\mathbf{K}$ 中的整數 $\lambda_1,\lambda_2,\ldots$ 選得使每一個 $\lambda$ 都被一個例外素理想整除，但不被它的平方整除，並且不被其餘例外素理想整除，那麼作為分母 $a$，這些 $\lambda$ 的冪乘積就足夠了。

\noindent\textit{證明。}
設
\[
        (a)=\mathfrak{p}_1^{\alpha_1}\cdots\mathfrak{p}_g^{\alpha_g}
\]
是 $a$ 的素理想分解，並且 $\mathfrak{p}_1,\ldots,\mathfrak{p}_g$ 都不是例外素理想；於是對於每一個 $\mathfrak{p}_i$，$\mathfrak{I}_p$ 同構於 $\mathfrak{o}_p/\mathfrak{a}_p$。從 $\mathbf{K}$ 的全體整數環過渡到由 $(a)$ 定義的商環 $R_a$：也就是把且僅把那些與 $(a)$ 互素的整數加入分母，換言之，這些整數不被 $\mathfrak{p}_1,\ldots,\mathfrak{p}_g$ 中任何一個整除。在 $R_a$ 中，$\mathfrak{p}_1,\ldots,\mathfrak{p}_g$ 仍然是素理想，但成為主理想；所有其餘素理想都變為單位理想；$(a)$ 的分解保持不變\footnote{參見例如 H.\ Grell, Zur Theorie der Ordnungen in algebraischen Zahl- und Funktionenkörpern, Math.\ Ann.\ 97 (1927), S.\ 524--558, \S\,2, 1。}。因此，當過渡到以 $R_a$ 為係數域時，推論及其證明完全保持 \S\,2, 3.\ 中的形式。

這就是說：若 $a\cdot g(x)$ 屬於 $\mathfrak{I}$，則 $g(x)=F(f_i(x))$，其中多項式 $F(u)$ 的係數屬於 $R_a$。若 $d$ 是這些係數的一個公分母，因而 $d$ 與 $a$ 互素，那麼也可以表述為：從 $a\cdot g(x)$ 屬於 $\mathfrak{I}$ 推出 $\delta \cdot g(x)$ 屬於 $\mathfrak{I}$，其中 $d$ 與 $a$ 互素。於是 $(d,a)$ 等於單位理想，從而得到：由 $a\cdot g(x)$ 屬於 $\mathfrak{I}$ 推出 $g(x)$ 屬於 $\mathfrak{I}$。這證明了上述結論的第一部分。

為證明第二部分，令 $R^*$ 是由有限多個例外素理想定義的 $\mathrm{K}$ 中的商環，並令 $\lambda_1,\ldots,\lambda_s$ 為這些素理想在 $R^*$ 中的基元素；這些基元素可以取為整數。在 $R^*$ 中，$a$ 除去一個 $R^*$ 的單位外，成為這些 $\lambda$ 的一個冪乘積。回到整數，就得到
\[
        \gamma a=\beta\lambda_1^{\alpha_1}\cdots\lambda_s^{\alpha_s},
\]
其中 $\gamma$ 與 $\beta$ 都不被任何例外素理想整除。若 $\alpha \cdot g(x)$ 屬於 $\mathfrak{I}$，則
\[
        \beta\cdot\lambda_1^{\alpha_1}\cdots\lambda_s^{\alpha_s}g(x)
\]
也屬於 $\mathfrak{I}$；於是按推論第一部分，
\[
        \lambda_1^{\alpha_1}\cdots\lambda_s^{\alpha_s}g(x)
\]
屬於 $\mathfrak{I}$。

\medskip
\hrule width 2.7cm
\smallskip
\noindent\textsuperscript{*)} 更確切地說，由於 $\mathfrak{a}$ 具有特徵為零的係數，這一結論與該文的輔助定理 V 無關。事實上，該輔助定理的證明並不正確（在所給出的代換下，一切都可能恆等消失）；正確的證明直到 B.\,L.\,v.\,d. Waerden 的《Bézout 定理的一種推廣》，\emph{Math. Ann.} 99 (1928)，第 497--541 頁，才得以給出。參見該文注 37。

該文第 257 頁還需要作一個容易看出的修正（定理 XIV 在 \S\,7 中同樣沒有用到）。原文說，$(\mathfrak{R})$ 中的所有元素都是通過把 $(x_i)$ 中的 $t_{\mu,\nu}$ 特殊化為 $(\mathfrak{R}')$ 中的元素而產生的；實際上，它們是由一個更一般的線性形式
\[
\sum s_{\lambda_1\ldots\lambda_n}(y_1)^{\lambda_1}\cdots(y_n)^{\lambda_n}
\]
經特殊化得到的，而該線性形式的指數與 $(x_i)$ 的指數相同。

\bigskip
\noindent\rule{\textwidth}{0.4pt}

\medskip
\noindent\textbf{俄文摘要（漢譯）}

\begin{center}
\textbf{論整係數函數的極大域}\\[0.4em]
\textbf{埃米·諾特（哥廷根）}\\[0.3em]
\emph{（摘要）}
\end{center}

\smallskip
\textit{每一個整域 $\mathfrak{I}$（Integritätsbereich，即無零因子環）定義了同一類型的某個極大域 $\mathfrak{M}(\mathfrak{I})$；它由所有函數 $g(x)$ 組成，條件是至少存在一個 $a\ne0$，使得 $a\cdot g(x)$ 包含在 $\mathfrak{I}$ 中。}

\textit{本文研究這個極大域 $\mathfrak{M}(\mathfrak{I})$ 的“分母” $a$ 的問題，並且假定域 $\mathfrak{I}$ 是有限的。本文證明，這些 $a$ 總可以選得使它們的因子中（總體而言）只出現有限多個素數。}

\textit{同時，在 $g(x)$ 為冪級數或這些級數的商的情形中，假定定義域 $\mathfrak{I}$ 的函數 $f_1(x),\ldots,f_r(x)$ 彼此之間只能由代數關係聯繫起來。}

\textit{證明的基礎是把這個問題歸約為一般理想論中的某些問題。這個歸約通過如下方式完成：把素數 $p$ 在分母 $a$ 中的出現看作函數 $f_i$ 之間的一個模 $p$ 關係。}

\begin{center}
\emph{（Rec. Math. XXXVI:1；1929）}
\end{center}

\medskip
\noindent\rule{\textwidth}{0.4pt}

%
\clearpage
%
\setcounter{footnote}{0}

%
%
%
%
%
\pdfbookmark[0]{36. 理想微分與不同}{paper-36}
\section*{36. 理想微分與不同}

\begin{center}
\emph{J. Ber. d. DMV39 (1930), S. 17}
\end{center}

E. Noether，哥廷根：理想微分與不同。

本文說明，代數數域的不同如何可被理解為一個定義\emph{理想}的\emph{微分商}。這種定義與通常定義相一致，是由一些本身也很有意義的結構定理推出的；這些定理在類似意義下表示 Lagrange 插值公式的推廣。

詳細論述將發表於 \emph{Math. Ann.}

%
\clearpage
%
%
\begingroup
\def\srcemph#1{\textbf{#1}}
\setcounter{footnote}{0}
\pdfbookmark[0]{37. 無高階分歧的域中的正規基}{paper-37}
\section*{37. 無高階分歧的域中的正規基}

\begin{center}
\emph{Journal f. d. reine u. angew. Math. 167 (1932), S. 147--152}
\end{center}

\begin{center}
作者：\emph{Emmy Noether}，哥廷根。
\end{center}

所謂伽羅瓦數域 \(K/k\) 的正規基，是指 \(K\) 的整數環（極大階）\(\frakO\) 作為 \(k\) 的整數環 \(\frako\) 上的模，有一組由共軛元素組成的基。眾所周知\footnote{參見 A. Speiser, Gruppendeterminante und Körperdiskriminante, § 6. Math. Ann. 77 (1916), S. 546--562。}，這樣一個基存在的必要條件是，在 \(K/k\) 中不出現高階分歧域。即使只限制在一個素位上，也就是說，過渡到 \(k\) 的 \(\frakp\)-進完備化 \(k_\frakp\) 以及 \(K\) 的相應標量擴張 \(K_\frakp\)，這一點仍然成立。下面我證明其逆命題：

\srcemph{對於所有素位 \(\frakp\)，只要 \(\frakp\) 不整除 \(K/k\) 的次數，就存在正規基。特別地，如果 \(K/k\) 沒有高階分歧，則在每個分歧素位上都存在正規基；在那裡，判別式成為一個群行列式的平方。}

對於最後這個斷言還要補充說，與 E. Artin 的一個猜想相反，過渡到阿廷導子\footnote{E. Artin, Über die gruppentheoretische Struktur der Diskriminante. J. f. Math. 164 (1931), S. 1--11。}還需要進一步的考慮：按不可約特徵標分解群行列式，會得到廣義拉格朗日預解因子（原文稱 Wurzelzahlen）；適當合併之後，得到在由特徵標擴大的基域中的一個分解。在最簡單的情形中，藉助這些預解因子的理想論分解，我能夠把這些新因子與阿廷導子等同起來。\footnote{[1931 年 10 月 6 日。] 一般情形中仍需要理想論考慮，這一點由 M. Deuring 告訴我的下列例子說明；這個例子可以任意變化。它給出一個非極大階，其判別式可以表示為群行列式的平方，而共軛特徵標卻對應於不同的理想因子：令 \(k\) 為五次單位根域，\(K=k(\sqrt[5]{2})\)。考察階
\([1,2\sqrt[5]{2},\sqrt[5]{2^2},\sqrt[5]{2^3},\sqrt[5]{2^4}]\)；並且是在素位 \((2)\) 上考察，按本短文的考慮，它在那裡有正規基。按特徵標作分解時，相應群行列式的因子恰好就是上述基元素。因此“導子”除了 \(1\) 以外成為：\(2\sqrt[5]{2}\cdot\sqrt[5]{2^4};\;\sqrt[5]{2^2}\cdot\sqrt[5]{2^3};\;\sqrt[5]{2^3}\cdot\sqrt[5]{2^2};\;\sqrt[5]{2^4}\cdot 2\sqrt[5]{2}\)，也就是 \(4,2,2,4\)；所以對於共軛特徵標它們是不同的。} 我還要指出，即使在一般情形中沒有正規基，藉助一個“按伽羅瓦模的合成列”，仍然總可以分解為廣義拉格朗日預解因子；並且由此也像上面一樣得到在擴大的基域中的一個分解；這裡又開始出現理想論問題。

在上述前提下，正規基存在性的證明基於這樣一個事實：把 \(\frakO/\frako\) 看作伽羅瓦模（即帶伽羅瓦群作用的模），也就是看作一個由伽羅瓦群代換作用的 \(\frako\)-模時，它作為 \([\Gg]_{\frako}\)-模，與用 \(\frako\) 作標量擴張所得整係數群環中的一個單側理想同構。這個群環在所有溫分歧素位上給出一個極大階；而 \(\frakp\)-進半單代數的極大階中的理想都是主理想\footnote{參見 H. Hasse, Über \(\frakp\)-adische Schiefkörper und ihre Bedeutung für die Arithmetik hyperkomplexer Zahlsysteme, Sätze 46 und 47. Math. Ann. 104 (1931), S. 495--534。在交換代數中，只需過渡到 \(k\) 在 \(\frakp\) 處的局部化環，就已經有主理想性質；因此當處理阿貝爾群和阿貝爾域時，可以用這種較弱的局部化來定義該素位。}，所以這些理想在這裡由一個元素通過群或群環的代換產生。回到伽羅瓦模，這恰好意味着正規基的存在。至於高階分歧素位，整係數群環變為一個非極大階，與 \(\frakO/\frako\) 對應的理想變為非主理想。

顯然，當所涉及的是一變量函數域，並且常數域的特徵不整除 \(K/k\) 的次數時，所有這些考慮仍然保持有效。

\subsection*{§1. 經 \(\frakp\)-進標量擴張的整係數群環}

令 \(\frako\) 以及 \(\frako_\frakp\) 分別表示代數數域 \(k\) 及其 \(\frakp\)-進完備化 \(k_\frakp\) 的整數環，其中 \(\frakp\) 是 \(k\) 中的一個素理想。再令 \((\Gg)\) 表示群 \(\Gg\) 的有理係數群環，\([\Gg]\) 表示其整係數群環，群 \(\Gg\) 的元素個數為 \(n\)。所謂 \((\Gg)_k\) 或 \((\Gg)_{k_\frakp}\)，是指係數分別取自 \(k\) 或 \(k_\frakp\) 的群環；相應地，\([\Gg]_\frako\) 或 \([\Gg]_{\frako_\frakp}\) 是把整係數群環擴展為係數分別取自 \(\frako\) 或 \(\frako_\frakp\) 的群環。

於是 \((\Gg)_k\) 和 \((\Gg)_{k_\frakp}\) 是根基為零的代數（即半單代數），而 \([\Gg]_\frako\) 以及 \([\Gg]_{\frako_\frakp}\) 分別是 \((\Gg)_k\) 和 \((\Gg)_{k_\frakp}\) 中的階。

\paragraph{定理 1.}
\srcemph{如果素理想 \(\frakp\) 不整除 \(\Gg\) 的元素個數 \(n\)，則 \([\Gg]_{\frako_\frakp}\) 是 \((\Gg)_{k_\frakp}\) 的一個極大階；如果 \(\frakp\) 整除 \(n\)，則 \([\Gg]_{\frako_\frakp}\) 不是極大的。}

整係數群環的判別式，因而 \([\Gg]_\frako\) 以及 \([\Gg]_{\frako_\frakp}\) 的判別式，等於元素個數 \(n\) 的一個冪。\footnote{參見例如 E. Noether, Hyperkomplexe Größen und Darstellungstheorie, § 26. Math. Ztschr. 30 (1929), S. 641--692。} 因此，如果 \(\frakp\) 不整除 \(n\)，則 \([\Gg]_{\frako_\frakp}\) 的判別式成為一個單位。這就是說，在固定 \(\frako_\frakp\) 的情況下，\([\Gg]_{\frako_\frakp}\) 不能擴大為一個更大的階；因為這樣的擴大對應於從判別式中分離出一個非單位的平方因子。而 \(\frako_\frakp\) 已經被假定為整數環，也就是 \(k_\frakp\) 中的極大階。於是定理 1 的第一部分，也就是下文唯一要用到的部分，得到證明。

反之，如果 \(\frakp\) 整除 \(n\)，則與平凡表示相應的直和
\[
        \frakC=E^{(1)}[\Gg]_{\frako_\frakp}+(1-E^{(1)})[\Gg]_{\frako_\frakp},
        \qquad
        E^{(1)}=\frac1n\sum S,
\]
給出 \([\Gg]_{\frako_\frakp}\) 的一個真擴張。因為由 \(E^{(1)}=\frac1n\sum S\) 可知 \(\frakC\) 是真擴張；\(\frakC\) 是一個階，其相關生成元為 \(E^{(1)}\) 和 \((1-E^{(1)})S\)，其中 \(S\in\Gg\)（且 \(E^{(1)}S=E^{(1)}\)）。

\paragraph{定理 2.}
\srcemph{如果 \(\frakp\) 不整除 \(n\)，則關於 \([\Gg]_{\frako_\frakp}\) 的每一個整理想或分式理想都是主理想。}

這是因為，按定理 1，\([\Gg_{\frako_\frakp}]\) 是一個 \(\frakp\)-進半單代數的極大階；因此其中的理想都是主理想。\footnote{Hasse，前引。那裡只對單代數證明了主理想性質；但由已知推理可從中推出半單情形。因為一個極大階的各個分量仍然極大；所考察理想的各個分量於是成為主理想，例如由 \(a_i\) 生成。於是 \(a=\sum a_i\) 成為原理想的生成元。對於阿貝爾群還可參見注 3。}

\subsection*{§2. 伽羅瓦模、模同構與正規基}

設 \(K/k\) 是伽羅瓦擴張，\(\Gg\) 為其群；\(z^S\) 表示由 \(\Gg\) 中代換 \(S\) 作用於 \(z\) 得到的 \(K\) 中元素。把 \(K/k\) 看作 \(k\)-模，也就是看作以 \(k\) 為標量環的加法阿貝爾群時，\(\Gg\) 的代換在其上產生模自同構。因此可以通過如下規定，使 \(K/k\) 成為關於 \((\Gg)_k\) 的模：
\[
        z\Bigl(\sum_i S_i c_i\Bigr)=\sum_i z^{S_i}c_i,
        \qquad z\in K,\; S_i\in\Gg,\; c_i\in k.
\]

\paragraph{定義。}
\srcemph{\(K/k\) 中的每一個 \((\Gg)_k\)-模稱為有理伽羅瓦模。同樣，\(K/k\) 中的 \([\Gg]_\frako\)-模稱為整伽羅瓦模；特別地，\(\frakO/\frako\) 是一個整伽羅瓦模。}

\paragraph{定理 3.}
\srcemph{作為 \((\Gg)_k\)-模，\(K/k\) 與 \((\Gg)_k\) 同構。}

這是因為 \(K/k\) 總有一個域正規基，也就是說，有一個由共軛元素 \(z^S\) 組成的 \(k\)-基。\footnote{事實上，若 \(a_1,\ldots,a_n\) 是 \(K/k\) 的一個基，而 \(u_i\) 表示不定元，則當 \(S\) 遍歷 \(\Gg\) 的元素時，\(\sum u_i a_i^S\) 線性無關；因此也可以把 \(u_i\) 專化為 \(k\) 中元素，並保持線性無關。} 對應
\[
        S\to z^S,
        \qquad
        \sum S_i c_i\to\sum z^{S_i}c_i
        \quad(c_i\in k),
        \quad\hbox{也即 } ST\to z^{ST},
\]
給出一個 \((\Gg)_k\)-模同態；由於兩邊關於 \(k\) 的秩相同，因而成為同構。

\paragraph{定理 4.}
\srcemph{作為 \([\Gg]_\frako\)-模，\(\frakO/\frako\) 與 \((\Gg)_k\) 中一個最高秩 \(n\) 的 \([\Gg]_\frako\)-模同構。同樣，\(\frakO_\frakp/\frako_\frakp\) 作為 \([\Gg]_{\frako_\frakp}\)-模，與 \((\Gg)_{k_\frakp}\) 中一個秩為 \(n\) 的 \([\Gg]_{\frako_\frakp}\)-模同構。}

定理 3 中給出的 \((\Gg)_k\)-同構 \(K/k\to(\Gg)_k\)，把 \([\Gg]_\frako\)-模 \(\frakO/\frako\) 對應到 \((\Gg)_k\) 中一個秩為 \(n\) 的 \([\Gg]_\frako\)-模。進一步，只要在定理 3 的對應中令 \(c_i\) 遍歷 \(k_\frakp\) 的所有元素，就可把這個 \((\Gg)_k\)-同構擴張為從 \(K_\frakp/k_\frakp\) 到 \((\Gg)_{k_\frakp}\) 的一個 \((\Gg)_{k_\frakp}\)-同構。這個同構於是誘導出 \(\frakO_\frakp/\frako_\frakp\) 的 \([\Gg]_{\frako_\frakp}\)-同構。這裡 \(K_\frakp/k_\frakp\) 要看作 \(k_\frakp\) 上的有限維代數（原文“超複系統”）。它由 \(K/k\) 作標量擴張而得；一般說來會有零因子，並且是若干同構域的直和，因而是半單代數。

\paragraph{定理 5.}
\srcemph{在每一個不整除 \(K/k\) 的次數 \(n\) 的素位 \(\frakp\) 上，\(\frakO_\frakp/\frako_\frakp\) 有一個正規基。}\footnote{在阿貝爾域的情形中，按注 3 和注 5，也可以用在該素位處的局部化環來定義素位。}

事實上，按定理 1，\([\Gg]_{\frako_\frakp}\) 在那裡成為極大階；於是 \((\Gg)_{k_\frakp}\) 中秩為 \(n\) 的 \([\Gg]_{\frako_\frakp}\)-模成為整理想或分式理想，並且按定理 2 是主理想。因此，如果與 \(\frakO_\frakp/\frako_\frakp\) 對應的理想等於 \(W[\Gg]_{\frako_\frakp}\)，則它有 \(\frako_\frakp\)-基 \(WS_i\)；上述模同構說明，如果 \(w\) 表示與 \(W\) 對應的 \(\frakO_\frakp/\frako_\frakp\) 中元素，那麼 \(w^{S_i}\) 構成 \(\frakO_\frakp/\frako_\frakp\) 的一個 \(\frako_\frakp\)-基。所以這 \(n\) 個共軛元素 \(w^{S_i}\) 構成 \(\frakO_\frakp/\frako_\frakp\) 的一個正規基。

如果 \(\frakp\) 整除 \(n\)，則與 \(\frakO_\frakp/\frako_\frakp\) 對應的模不可能是主模，因為在這種情況下不可能存在正規基。

\paragraph{對定理 3 的補充說明。}
由定理 3 所證明的 \((\Gg)_k\) 與 \(K/k\) 的模同構\footnote{這個證明顯然只要求 \(K\) 是 \(k\) 上第一類的伽羅瓦擴張，並且 \(k\) 有無限多個元素。後一限制是不必要的：M. Deuring 想出了一個對有限域也有效的定理 3 的證明，由此反過來得到域正規基的存在。同時還得到原始元存在性的一個證明，它對有限域和無限域統一適用。}可以推出 Speiser 關於 Klein 型問題的結果\footnote{同前。伽羅瓦模這一概念就是從那裡抽象出來的。如果 \(k\) 的特徵整除 \(n\)，則要處理的是按伽羅瓦模的合成列及其不可約因子。}，其表述如下：

\srcemph{\(K/k\) 的伽羅瓦群 \(\Gg\) 的每個在 \(k\) 上不可約的表示，都由 \(K/k\) 中一個不可約有理伽羅瓦模產生；每個絕對不可約表示都由 \(K_{\mathsf{Z}}/\mathsf{Z}\) 中一個伽羅瓦模產生。}

這裡 \(\mathsf{Z}\) 表示包含 \(k\) 的一個分裂域，它分裂群環的相應分量；\(K_{\mathsf{Z}}/\mathsf{Z}\) 要看作 \(\mathsf{Z}\) 上的有限維代數，由 \(K/k\) 作標量擴張而得。特別地，如果可以選擇 \(\mathsf{Z}\)，使它與 \(K\) 在 \(k\) 上線性無交，也就是使 \(K_{\mathsf{Z}}/\mathsf{Z}\) 成為輔助擴張（原文 akzessorische Erweiterung；Speiser 所處理的情形），則 \(K_{\mathsf{Z}}\) 仍然是域。

這個斷言本身藉助模同構，歸結為對 \((\Gg)_k\) 或 \((\Gg)_{\mathsf{Z}}\) 的相應斷言；在那裡，相對於 \((\Gg)_k\) 或 \((\Gg)_{\mathsf{Z}}\) 的不可約理想產生表示。

如果 \(v_1,\ldots,v_t\) 是這樣一個伽羅瓦模的基，而 \(S\to \bar S\) 是相應表示，則有
\[
        \begin{pmatrix}\vdots\\ v_i^S\\ \vdots\end{pmatrix}
        =\bar S\begin{pmatrix}\vdots\\ v_i\\ \vdots\end{pmatrix}.
\]
這就是 Klein 和 Speiser 的表述。

\subsection*{§3. 無高階分歧的域中作為群行列式的判別式}

為了分解判別式，眾所周知只需考察它在各個分歧素位上的分解；對於沒有高階分歧的域，所涉及的就只是有正規基的素位。

\paragraph{定理 6.}
\srcemph{對於沒有高階分歧的伽羅瓦域 \(K/k\)，在每個分歧素位上，判別式都等於伽羅瓦群的群行列式的平方，其中不定元由正規基代入。按照不可約表示，群行列式分解為廣義拉格朗日預解因子；判別式分解為用特徵標擴大的基域中的理想，並且複共軛特徵標對應於同一理想。}

事實上，若 \(w^S\) 是正規基，則判別式 \(\Delta\) 等於
\[
        D=\bigl|w^{ST^{-1}}\bigr|,\qquad S,T\in\Gg
\]
的平方。這裡 \(D\) 正是以 \(w^S\) 代替不定元 \(x_S\) 得到的群行列式。

分解為預解因子由 Speiser 前引文獻中的推理得到。設
\[
        D=D_1^{f_1}\cdots D_t^{f_t},
        \qquad
        M=f_1M_1+\cdots+f_tM_t
\]
分別為群行列式和群矩陣按照絕對不可約表示的分解。若 \(S\to\bar S\) 表示與 \(M_\lambda\) 相應的表示，其中 \(\bar S\) 位於擴大的基域中，則有
\[
        M_\lambda=\sum_S w^S\bar S,
        \qquad
        M_\lambda^{T^{-1}}=\sum_S w^{ST^{-1}}\bar S
          =\sum_R w^R\bar R\bar T=M_\lambda\bar T,
\]
進而取行列式得到
\[
        D_\lambda^{T^{-1}}=D_\lambda|\bar T|=D_\lambda\varepsilon_T.
\]
這就把 \(D_\lambda\) 識別為廣義拉格朗日預解因子；\(\varepsilon_T\) 作為 \(\bar T\) 的行列式，成為一個單位根（即 \(\Gg\) 模其導出子群所得的商群（交換化）的不可約表示）。在循環群的情形中，\(D_\lambda\) 簡化為拉格朗日預解和 \(\sum w^S\chi_\lambda(S)\)。

一般說來，\(D_\lambda\) 位於由 \(K_\frakp/k_\frakp\) 把係數域 \(k_\frakp\) 按相應特徵標作標量擴張所得的有限維代數中；並且它們是這個代數中的\srcemph{整元素}。因為用不定元 \(x_S\) 構成的行列式，其係數是特徵標域中的整數；而 \(w^S\) 是 \(K_\frakp\) 中的整元素。

為了從群行列式的分解過渡到判別式的分解，需要把每個表示同其複共軛表示（在此等價於對偶表示）配對。若 \(\lambda,\bar\lambda\) 是這樣一對錶示，則同時有
\[
        D_\lambda^{T^{-1}}=D_\lambda\varepsilon_T,
        \qquad
        D_{\bar\lambda}^{T^{-1}}=D_{\bar\lambda}\varepsilon_T^{-1};
\]
同時，對於不定元 \(x_S\) 形成的、屬於 \(\lambda,\bar\lambda\) 的行列式是複共軛的；一般地說，不定元 \(x_S\) 下的共軛特徵標對應於共軛行列式。

因此若令 \(\Delta_\lambda=D_\lambda D_{\bar\lambda}\)，就只需添入第 \(\lambda\) 個特徵標的實子域，並且有
\[
        \Delta_\lambda^T=\Delta_\lambda
        \qquad\hbox{对所有 }T\hbox{ 属于 }\Gg;
\]
於是\footnote{由於這裡是在有限維代數意義下對一個域作標量擴張，通常伽羅瓦理論中的推理必須作修改。設 \(\mathsf{P}\) 為所考慮的 \(k_\frakp\) 的擴域，並設
\[
        (K_\frakp)_{\mathsf{P}}=(K_\frakp)_{\mathsf{P}} e^{S_1}+\cdots+(K_\frakp)_{\mathsf{P}} e^{S_r}
\]
為分解為域的直和。於是 \( (K_\frakp)_{\mathsf{P}} e^{S_i}/\mathsf{P}e^{S_i}\) 是伽羅瓦的（其群是 \(\frakp\) 的一個素因子之分解群的子群），而 \(e^{S_i}\) 關於陪集互為共軛。由 \(\Delta_\lambda^T=\Delta_\lambda\) 先推出 \(\Delta_\lambda\) 的第 \(i\) 個分量位於 \(\mathsf{P}e^{S_i}\) 中，所以比如等於 \(\gamma_i e^{S_i}\)，其中 \(\gamma_i\) 在 \(\mathsf{P}\) 中。再由 \(e^{S_i}\) 的共軛性推出所有 \(\gamma_i\) 相等；所以實際上 \(\Delta_\lambda=\gamma\sum e^{S_i}=\gamma\) 位於 \(\mathsf{P}\) 中。} \(\Delta_\lambda\) 位於用第 \(\lambda\) 個特徵標的實子域擴大的基域中。

\srcemph{判別式分解}
\[
        \Delta=\Delta_1^{f_1}\cdots\Delta_t^{f_t}
\]
\srcemph{因此是在由各特徵標的實子域擴大的基域中的分解，並且複共軛特徵標對應於同一因子。對於其餘共軛特徵標，一般不能直接作出結論（參見 2a）。}

這樣定理 6 的各部分都得到了證明。如果能夠證明由 \(\Delta_\lambda\) 生成的理想實際上已經位於基域 \(k\) 中，並且對於共軛特徵標相等，那麼只要把相應的 \(\Delta_\lambda\) 的冪乘積分配給合成特徵標，它們就與阿廷導子一致。因為函數方程成立，而且還成立如下事實：由正規基可直接推出（參見 Speiser 前引文獻），子群的平凡表示對應於相應子域的判別式。在每一個素位上相等，就推出完整理想相等。若限制在有理不可約特徵標上，則上述事實給出相等。對於絕對不可約特徵標，下面還要在最簡單的情形中直接確認這種相等。

\paragraph{定理 7.}
\srcemph{設 \(k\) 是有理數域，且 \(K/k\) 是次數為素數 \(l\)、與判別式互素的循環擴張，因而沒有高階分歧。則對於非平凡表示，由 \(\Delta_\lambda\) 生成的理想彼此相等，並且與 \(K/k\) 的導子一致。}

這由已知的拉格朗日預解因子分解推出。\footnote{Hilbert, Zahlbericht, § 108。由於按定理 5 正規基的存在已經確定，而此處的素位可按注 7 直接用局部化環定義，所以不再需要引用《數論報告》中所用的事實，即絕對循環域是分圓域。} 在 \(K/k\) 的一個分歧素位 \(p\) 上，對於 \(k_\varepsilon\)，即 \(l\) 次單位根域，以及 \(K_\varepsilon\)（\(K_\varepsilon\) 仍然是域，因為 \(k_\varepsilon\) 與 \(K\) 在 \(k\) 上線性無交，並且按注 7 這裡不需要 \(p\)-進完備化），有
\[
        (p)=\frakp_1\cdots\frakp_{l-1},
        \qquad
        \frakp_i=\frakP_i^{\,l},
\]
其中 \(\frakp\) 表示 \(k_\varepsilon\) 中的素理想，\(\frakP\) 表示 \(K_\varepsilon\) 中的素理想。進一步，對於預解因子有
\[
        (\Omega_\lambda)=\frakP_1^{r_1}\cdots\frakP_{l-1}^{r_{l-1}},
        \qquad
        (\Omega_{\bar\lambda})=
        \frakP_1^{l-r_1}\cdots\frakP_{l-1}^{l-r_{l-1}},
\]
其中 \(r_1,\ldots,r_{l-1}\) 是數字 \(1,\ldots,l-1\) 的一個排列。這裡，預解因子 \(\Omega_\lambda\) 與群行列式的因子 \(D_\lambda\) 相同；因此得到
\[
        (\Delta_\lambda)=(D_\lambda D_{\bar\lambda})
        =(\Omega_\lambda\Omega_{\bar\lambda})
        =\frakP_1^{l}\cdots\frakP_{l-1}^{l}
        =\frakp_1\cdots\frakp_{l-1}=(p);
\]
也就是說，它與 \(K/k\) 的導子一致。

\begin{center}
1931 年 8 月 24 日收稿。
\end{center}
\endgroup

%
\clearpage
%
\setcounter{footnote}{0}

%
%
%
%
\pdfbookmark[0]{38. 與 R. Brauer 和 H. Hasse 合作：代數理論中一個主定理的證明}{paper-38}
\section*{38. 與 R. Brauer 和 H. Hasse 合作：代數理論中一個主定理的證明}

\begin{center}
\emph{Journal f. d. reine u. angew. Math. 167 (1932), S. 399--404}
\end{center}

\begin{center}
\textbf{代數理論中一個主定理的證明}\\[0.5em]
由柯尼斯堡的 \emph{R. Brauer}、馬堡的 \emph{H. Hasse} 與哥廷根的 \emph{E. Noether} 合作完成\footnote{本文的撰寫由 H. Hasse 承擔。}。
\end{center}

經過我們共同的努力，終於證明了下列定理；它對於代數數域上線性代數的結構理論，並且也在更廣的範圍內，都具有基本意義：

\paragraph{主定理.}
\emph{代數數域上的每個正規除代數都是循環的；或者也可說，是 Dickson 型的。}

這個結果本質上應歸功於 \(p\)-進方法。能把它獻給這種方法的奠基人 Kurt Hensel 先生，以慶祝他的七十歲生日，對我們來說是一件特別高興的事。

我們的證明由三個歸約組成；我們三人各貢獻了其中一個。\footnote{這裡按照這些歸約產生的先後順序敘述；這同系統順序正好相反。}

\paragraph{1.}
第一歸約由 H. Hasse 給出，它基於他最近建立的代數數域上循環代數理論。\footnote{H. Hasse, \emph{Theorie der zyklischen Algebren über einem algebraischen Zahlkörper}, Gött. Nachr. 1931。一個較詳細的證明敘述即將以 \emph{Theory of cyclic algebras over an algebraic number field} 為題發表於 Trans. Amer. Math. Soc.；在那裡特別也發展了 E. Noether 在講課中提出的、在該處和本文中都起基礎作用的分裂域與交叉積理論。下文以 H 引用這篇工作。H, 1--6 除語言不同外即構成先前那篇短文。}

\paragraph{歸約 1.}
主定理已經證明，只要證明：
\[
\mathrm{I}.\quad A\text{ 是 }\Om\text{ 上處處分裂的代數 }\Longrightarrow A\sim\Om .
\]
為簡化措辭，以下所謂“\(\Om\) 上的代數 \(A\)”始終指“代數數域 \(\Om\) 上的正規單代數 \(A\)”，也就是中心為 \(\Om\) 的單超複系統。記號 \(A\sim\Om\) 表示 \(A\) 是 \(\Om\) 上的全矩陣代數，也就是 \(A\) 屬於由 \(\Om\) 自身作為 \(\Om\) 上除代數所確定的那個特殊相似類（即相應的 \(\Om\) 上除代數相同）[H, 5]。最後，所謂 \(\Om\) 上處處分裂的代數，是指對 \(\Om\) 的每個素位 \(\frakp\)，由把係數域 \(\Om\) 擴張到相應的 \(\frakp\)-進域 \(\Om_{\frakp}\) 所得到的 \(\frakp\)-進擴張代數都分裂：
\[
A_{\frakp}\sim \Om_{\frakp}.
\]

\paragraph{證明.}
設 \(D\) 是 \(\Om\) 上的一個除代數。若 \(Z\) 是 \(\Om\) 上的一個循環域，且對 \(\Om\) 的每個素位 \(\frakp\)，\(Z\) 的 \(\frakp\)-次數都是 \(D\) 的 \(\frakp\)-指標的倍數 [參見 H, 17 Bb 的注]，那麼對於 \(Z\) 中相應的素因子 \(\frakP\)，\(\frakP\)-進域 \(\mathsf{Z}_{\mathfrak P}\) 都是 \(D_{\frakp}\) 的分裂域 [H, 18.1]。於是，把 \(D\) 的係數域擴張到 \(\mathsf{Z}\) 得到的代數 \(D_{\mathsf{Z}}\)，作為 \(\mathsf{Z}\) 上的代數，處處分裂。如果 I 已對每個 \(\Om\) 成立，就有 \(D_{\mathsf{Z}}\sim \mathsf{Z}\)。這就是說，\(D\) 有循環分裂域 \(\mathsf{Z}\)，因而可循環表示 [H, 5]。於是 \(D\) 還具有一個循環分裂域，其次數與 \(D\) 自身的次數相同；這就是說，\(D\) 是循環的 [H, 6, Satz 6]。所以在假設 I 下，主定理已經證明。

\paragraph{2.}
第二歸約的決定性啓發來自 R. Brauer。他曾寫信告訴 H. Hasse：與此相關的、關於除代數指數精確值的問題（這個問題後來也由主定理一併解決，見下）可藉助 Sylow 群定理歸約到具有可解分裂域的情形。這個想法也可用於循環性問題的相應進一步歸約。

\paragraph{歸約 2.}
定理 I 已經證明，只要證明：
\[
\mathrm{II}.\quad A\text{ 處處分裂且在 }\Om\text{ 上有可解分裂域 }\Longrightarrow A\sim\Om .
\]

\paragraph{證明.}
設 \(A\) 是 \(\Om\) 上處處分裂的代數。令 \(\mathsf{K}\) 是 \(A\) 的一個伽羅瓦分裂域，再令 \(p\) 為任意素數，並令 \(\Sigma\) 是 \(K/\Om\) 中關於 \(p\) 的一個 Sylow 域（即相應於一個 \(p\)-Sylow 子群的伽羅瓦群不變域）。那麼 \(A_{\Sigma}\) 是 \(\Sigma\) 上的一個同樣處處分裂的代數，並且有可解分裂域 \(\mathsf{K}\)。如果 II 已經成立（對每個 \(\Om\)），則 \(A_{\Sigma}\sim\Sigma\)。這就是說，\(A\) 有一個次數與 \(p\) 互素的分裂域 \(\Sigma\)。因此 \(A\) 的指標作為這個次數的因子，也與 \(p\) 互素。因為這對每個素數 \(p\) 都成立，所以指標等於 \(1\)，即 \(A\sim\Om\)。在假設 II 下，I 已經證明。\footnote{藉助 Sylow 群定理歸約到可解分裂域的想法，R. Brauer 早已用過；他用它證明指標的每個素因子也出現在指數中（\emph{Über den Zusammenhang von arithmetischen und invariantentheoretischen Eigenschaften von Gruppen linearer Substitutionen}, Berl. Akad.-Ber. 1926）。近來 A. A. Albert 對這個想法以及 R. Brauer 與 E. Noether 理論中的一系列一般定理，給出了不依賴表示論的簡單證明：1. \emph{On direct products, cyclic division algebras, and pure Riemann matrices}; 2. \emph{On direct products}; 兩篇均見 Trans. Amer. Math. Soc. 33 (1931)；關於此處涉及的歸約，特別見第二篇中的 Satz 23。\par \emph{校樣期間補記.} 另外，A. A. Albert 在得知 H. Hasse 私信告知主定理已由他對阿貝爾代數證明（見下文）之後，立即獨立地由此推出以下事實：\par a) 對 \(2^e\) 型次數的主定理；\par b) 下文的定理 1（指數 = 指標）；\par c) 除了歸約 2 的基本思想外，還包括下列歸約 3 的思想；當然沒有同歸約 1 聯繫起來，因此得出的結果是：若 \(D\) 是 \(\Om\) 上素數冪次數 \(p^e\) 的除代數，則存在一個關於 \(\Om\) 次數與 \(p\) 互素的擴張域 \(\Om'\)，使得 \(D_{\Om'}\) 是循環的。\par 這三個結果自然都已由我們後來完成的主定理證明所涵蓋。它們仍表明，A. A. Albert 對主定理的證明也有獨立貢獻。最後，A. A. Albert（在知道我們的主定理證明之後）還指出，我們的核心定理 I 可由他一篇正在印刷中的論文（Bull. Amer. Math. Soc. 37 (1931)）中的 Satz 13, 10, 9 幾行推出。這些定理的證明本質上基於同我們的歸約 2 和 3 相同的推理。}

\paragraph{3.}
第三歸約由 E. Noether 給出；這源於 H. Hasse 的一個通知，其中他用一般歸約 1 和相應因子系的公式化歸約，證明了主定理在阿貝爾分裂域特殊情形中的成立。E. Noether 從中提煉出了第三歸約的核心。另一方面，R. Brauer 也曾獨立地想到過這個第三歸約。H. Hasse 在已經掌握前兩個歸約之後，認識到它將導向最終證明。

\paragraph{歸約 3.}
定理 II 已經證明，只要證明：
\[
\mathrm{III}.\quad A\text{ 處處分裂且有素數次循環分裂域 }\Longrightarrow A\sim\Om .
\]

\paragraph{證明.}
設 \(A\) 是 \(\Om\) 上處處分裂的代數，並且有一個關於 \(\Om\) 的可解分裂域 \(\mathsf{K}\)。取域鏈
\[
K=\Lambda_0>\Lambda_1>\cdots>\Lambda_r=\Om
\]
位於 \(\mathsf{K}\) 與 \(\Om\) 之間，使得每一步 \(\Lambda_i/\Lambda_{i+1}\) 都是素數次數的循環擴張。於是 \(A_{\Lambda_1}\) 是 \(\Lambda_1\) 上處處分裂的代數，並且有素數次數循環分裂域 \(\mathsf{K}=\Lambda_0\)。如果 III 已經成立（對每個 \(\Om\)），則 \(A_{\Lambda_1}\sim\Lambda_1\)。這就是說，\(\Lambda_1\) 已是 \(A\) 的分裂域。現在從 \(\Lambda_1\) 而不是從 \(\Lambda_0\) 出發作完全相同的證明，得到 \(A_{\Lambda_2}\sim\Lambda_2\)，如此繼續，直到最後 \(A_{\Lambda_r}\sim\Lambda_r\)，也就是 \(A\sim\Om\)。在假設 III 下，II 已經證明。

\paragraph{4.}
現在，III 已由 H. Hasse 的結果 [H, 24] 確立。因此，通過歸約 3、2、1 逆向推出主定理成立。

E. Noether 還指出如下事實：III 的正確性，更一般地說，對任意次數循環分裂域的情形，歸結為 Hasse 範數定理\footnote{\label{fn:p38-hasse-normensatz}見 H, 3.11；證明見 H. Hasse, \emph{Beweis eines Satzes und Widerlegung einer Vermutung über das allgemeine Normenrestsymbol}, Gött. Nachr. 1931。}。但對於歸約 3，只需要素數次數的特殊情形，也就是原來的 Hilbert--Furtwängler 範數定理\footnote{見 H. Hasse, \emph{Bericht über neuere Untersuchungen und Probleme in der Theorie der algebraischen Zahlkörper II}, Jahresber. der D. M.-V., Erg.-Bd. 6 (1930), § 8 及其中所引文獻。}。因此，歸約 III 給出了 Hasse 範數定理的一個新的簡單證明。證明如下。

設 \(Z\) 是 \(\Om\) 上的循環域，\(\alpha\) 是 \(\Om\) 中一個數，並且對 \(\Om\) 的每個素位 \(\frakp\)，範數剩餘符號都滿足
\[
\left(\frac{\alpha,Z}{\frakp}\right)=1 .
\]
那麼每個循環代數
\[
A=(\alpha,Z)
\]
\footnote{\label{fn:p38-hasse-notation}採用 H, 1 的記號。這裡略去了同 \(\alpha\) 相聯繫的 \(Z\) 的自同構 \(S\)，因為它對當前論證無關緊要。} 處處分裂 [H, 17.7]。由歸約 3 得到--只使用 Hilbert--Furtwängler 範數定理--\(A\sim\Om\)。於是 \(\alpha\) 是 \(Z\) 中某個元素的範數 [H, 15.4]。

關於範數定理，我們還要指出：定理 I 應看作該定理向較高情形（也包括非阿貝爾情形）的正確推廣；而逐字推廣如 H. Hasse 所示\textsuperscript{\ref{fn:p38-hasse-normensatz}}，並不一般成立。

\begin{center}
\textbf{推論（H. Hasse）.}
\end{center}

\paragraph{5.}
作為主定理最直接的一項推論，可以指出：H. Hasse 所發展的代數數域上循環代數理論，現在已經具有代數數域上正規單代數（特別是正規除代數）的一般結構理論與不變量理論的意義。特別是，現在已經一般地證明了：

\paragraph{定理 1.}
代數數域上的一個正規單代數的指數等於它的指標。

此外，由定理 I 得到如下值得注意的事實：

\paragraph{定理 2.}
\(\Om\) 上一個真正正規除代數的基本理想至少被 \(\Om\) 的一個素位整除，而且至少被兩個素位整除。

這裡基本理想可定義為（約化）不同的（約化）範數。\footnote{在有理四元數代數的特殊情形中，這歸結為 H. Brandt 引入的“基本數”概念（\emph{Idealtheorie in Quaternionenalgebren}, Math. Ann. 99 (1928)）。這裡討論的不是一個數，而是一個理想，所以必須稱為“基本理想”；不要把它同 Dedekind 用於不同與判別式的“基本理想”和“基本數”混淆，後者正由於上述理由在相對域中不可用。-- 關於（約化）不同的定義，見下一條腳註。一個理想的（約化）範數，E. Noether 也逐素位地定義；確切地說，按照每個素位處任意理想都是主理想這一事實，取各素位處主理想生成元的（約化）數範數。} 另外，一個無限素位當且僅當該代數在那裡化為四元數代數（而不是實數域或複數域）時，才被納入基本理想。

\paragraph{證明.}
由 H. Hasse 的結果\footnote{\label{fn:p38-hasse-padisch}H. Hasse, \emph{Über \(\mathfrak p\)-adische Schiefkörper und ihre Bedeutung für die Arithmetik hyperkomplexer Zahlsysteme}, Math. Ann. 104 (1931)；見其中 Satz 42, 59。}可知，\(\Om\) 的一個素位 \(\frakp\) 出現在基本理想中，當且僅當 \(\frakp\)-指標不同於 \(1\)。若所有 \(\frakp\)-指標都等於 \(1\)，那麼代數處處分裂，由定理 I 可知它就不是真正的除代數。（也不可能只有一個 \(\frakp\)-指標不同於 \(1\)，因為範數剩餘符號的階正是 \(\frakp\)-指標 [H, 17.7]，而這些符號滿足乘積定理，即互反律 [H, 3.8]。）

\paragraph{6.}
定理 I 還使得可以確定（作算術刻畫）\(\Om\) 上一個代數 \(A\) 的所有分裂域 \(\mathsf{K}\)。這精確推廣了 H. Hasse 在 \(A,\mathsf{K}\) 均為循環的特殊情形中已經得到的事實 [H, 6, Satz 2]：

\paragraph{定理 3.}
對 \(\Om\) 上的一個代數 \(A\)，一個關於 \(\Om\) 的代數數域 \(\mathsf{K}\) 是分裂域，當且僅當對於 \(\Om\) 的每個素位 \(\frakp\)，以及 \(\mathsf{K}\) 中位於其上的所有素因子 \(\frakP_i\)，\(\mathsf{K}\) 的相應 \(\frakP_i\)-次數 \(n_{\frakP_i}\) 都是 \(A\) 的 \(\frakp\)-指標 \(m_{\frakp}\) 的倍數。

這裡，\(K\) 的 \(\frakP_i\)-次數是 \(\mathsf{K}_{\mathfrak P_i}\) 關於 \(\Om_{\frakp}\) 的次數。也就是說，如果
\[
\frakp=\prod_i \frakP_i^{\,e_{\frakP_i}},
        \qquad \mathsf{N}_{\mathsf{K}\Omega}(\mathfrak P_i)=\mathfrak p^{f_{\mathfrak P_i}},
\]
是 \(\frakp\) 在 \(K\) 中的分解，那麼
\[
n_{\frakP_i}=f_{\frakP_i}e_{\frakP_i}
\]
是剩餘次數 \(f_{\frakP_i}\) 與 \(\frakP_i\) 關於 \(\frakp\) 的分歧指數 \(e_{\frakP_i}\) 的乘積 [H, 4]。而 \(A\) 的 \(\frakp\)-指標是 \(\frakp\)-進擴張 \(A_{\frakp}\) 的指標 \(m_{\frakp}\) [H, 6]。

\paragraph{證明.}
\(\mathsf{K}\) 是 \(A\) 的分裂域，也就是 \(A_{\mathsf{K}}\sim \mathsf{K}\)，按定理 I 等價於對每個 \(\frakP_i\) 都有
\[
A_{K_{\frakP_i}}\sim K_{\frakP_i}.
\]

對一個有限素位 \(\frakp\)，設
\[
        A_{\frakp}\sim(\pi,W_{\frakp})
\]
\textsuperscript{\ref{fn:p38-hasse-notation}}為 \(A_{\frakp}\) 的算術上特出的循環表示 [H, 16；另見前引文獻\textsuperscript{\ref{fn:p38-hasse-padisch}}中的 Satz 38]。這裡 \(W_{\frakp}\) 是 \(\Om_{\frakp}\) 上次數為 \(m_{\frakp}\) 的未分歧域，\(\pi\) 是 \(\Om_{\frakp}\) 中恰被 \(\frakp^1\) 整除的一個數。又令 \(\mathit{K}_{\mathfrak P_i}\) 為複合域，並令
\[
        \Delta_{\frakP_i}=K_{\frakP_i}\cap W_{\frakp}
\]
為 \(W_{\frakp}\) 與 \(\mathsf{K}_{\mathfrak P_i}\) 的交。由於 \(\mathsf{K}_{\mathfrak P_i}\) 中所含的最大未分歧子域關於 \(\Om_{\frakp}\) 的次數是 \(f_{\frakP_i}\)，而在 \(\Om_{\frakp}\) 上每個次數的未分歧域唯一，所以 \(\Delta_{\frakP_i}\) 關於 \(\Om_{\frakp}\) 的次數為
\[
        d_{\frakP_i}=(m_{\frakp},f_{\frakP_i}).
\]
因此，\(\mathit{K}_{\mathfrak P_i}\) 關於 \(\mathsf{K}_{\mathfrak P_i}\) 的次數是 \(\frac{m_{\mathfrak p}}{d_{\mathfrak P_i}}\)，並且仍未分歧。

現在
\[
        A_{K_{\frakP_i}}=(A_{\frakp})_{K_{\frakP_i}}
        \sim(\pi,K_{\frakP_i}W_{\frakp})
\]
[H, 15.5]。於是 \(A_{\mathsf{K}_{\mathfrak P_i}}\sim \mathsf{K}_{\mathfrak P_i}\) 等價於 \(\pi\) 是 \(\mathit{K}_{\mathfrak P_i}\) 關於 \(\mathsf{K}_{\mathfrak P_i}\) 的範數。由於 \(K_{\frakP_i}W_{\frakp}/K_{\frakP_i}\) 未分歧，這又當且僅當 \(\pi\) 在 \(\frakP_i\) 處的階數 \(e_{\frakP_i}\) 是 \(K_{\frakP_i}W_{\frakp}/K_{\frakP_i}\) 的次數 \(\frac{m_{\mathfrak p}}{d_{\mathfrak P_i}}\) 的倍數。等價地，注意到
\[
\left(\frac{m_{\frakp}}{d_{\frakP_i}},\frac{f_{\frakP_i}}{d_{\frakP_i}}\right)=1,
\]
這同
\[
\frac{f_{\frakP_i}}{d_{\frakP_i}}e_{\frakP_i}
  =\frac{n_{\frakP_i}}{d_{\frakP_i}}
\]
是 \(\frac{m_{\mathfrak p}}{d_{\mathfrak P_i}}\) 的倍數相同，也就是同 \(n_{\frakP_i}\) 是 \(m_{\frakp}\) 的倍數相同，正如所斷言。

如果 \(\frakp\) 是一個無限素位，並且不是平凡情形 \(m_{\frakp}=1\)，則 \(m_{\frakp}=2\)，而 \(A_{\frakp}\) 是實數域 \(\Om_{\frakp}\) 上的四元數代數。此時 \(A_{\mathsf{K}_{\mathfrak P_i}}=(A_{\mathfrak p})_{\mathsf{K}_{\mathfrak P_i}}\sim \mathsf{K}_{\mathfrak P_i}\)，等價於 \(\mathsf{K}_{\mathfrak P_i}\) 是複數域，也就是 \(n_{\frakP_i}=f_{\frakP_i}=2=m_{\frakp}\)（這裡 \(e_{\frakP_i}\) 不出現）；這也再次給出所斷言的條件（此時 \(n_{\frakP_i}\) 與 \(m_{\frakp}\) 都只能取 \(1\) 或 \(2\)）。

\paragraph{7.}
如果把定理 3 所回答的問題反過來，即固定 \(\Om\) 上的代數域 \(K\)，考察所有由 \(K\) 分裂的 \(\Om\) 上代數 \(A\)，就會得到重要結果。這些代數 \(A\) 組成由 \(K\) 確定的 R. Brauer 群 \(\frakA\) 的一個子群 \(\frakK\)；這裡 \(\frakA\) 是 \(\Om\) 上所有代數（更準確說，所有相似代數類）的群。\footnote{R. Brauer, \emph{Über Systeme hyperkomplexer Größen}, Jahresber. d. D. M.-V. 38 (1929), S. 47/48。另見 H, 13.1。} 若假定 \(K/\Om\) 是伽羅瓦擴張，則由此得到的定理可看作類域論的若干主定理（相對阿貝爾數域理論）的推廣，推廣到一般相對伽羅瓦數域。

\paragraph{定理 4.}
\emph{（分解定理.）} 設 \(K/\Om\) 為伽羅瓦擴張，\(\frakp\) 是 \(\Om\) 中一個不整除 \(K\) 的相對判別式的素理想。則 \(K\) 中位於 \(\frakp\) 上的素因子的相對次數 \(f\)，等於對所有屬於 \(K\) 所對應群 \(\frakK\) 的代數 \(A\)，使
\[
        A_{\frakp}^{\,f}\sim\Om_{\frakp}
\]
成立的最早指數。

\paragraph{證明.}
\(f\) 是 \(\frakp\) 在 \(K\) 中的（對所有素因子相同的）\(\frakp\)-次數；另一方面，\(A\) 的 \(\frakp\)-指標等於 \(A_{\frakp}\) 的指標，也即指數。因此由定理 3 可知，\(f\) 至少是上述最早指數的倍數。為了完成證明，只需說明在群 \(\frakK\) 中存在 \(\frakp\)-指標恰為 \(f\) 的代數 \(A\)。

由 Frobenius 密度定理，必定還存在 \(\Om\) 中另一個不整除 \(K\) 的相對判別式的素理想 \(\frakp'\)，使得它在 \(\mathsf{K}\) 中的素因子也有相對次數 \(f\)。令 \(Z\) 是 \(\Om\) 上的一個循環域，使得它的 \(\frakp\)-次數和 \(\frakp'\)-次數都等於 \(f\)。再令 \(\alpha\) 是 \(\Om\) 中的一個數，使得範數剩餘符號
\[
\left(\frac{\alpha,Z}{\frakp}\right)\quad\text{与}\quad
\left(\frac{\alpha,Z}{\frakp'}\right)
\]
具有互為倒數的、階為 \(f\) 的值，而對 \(Z\) 的導子的所有其餘素因子 \(\frakq\)，都有
\[
\left(\frac{\alpha,Z}{\frakq}\right)=1.
\]
按照算術級數定理的推廣，還可把 \(\alpha\) 選得使它除可能含有 \(\frakp,\frakp'\) 及這些 \(\frakq\) 之外，只含有 \(\Om\) 中唯一一個一次冪出現的素理想 \(\frakr\)。於是由範數剩餘符號的乘積定理（互反律）又有
\[
\left(\frac{\alpha,Z}{\frakr}\right)=1
\]
[H, 3.8--10]。那麼每個代數
\[
        (\alpha,Z)=A
\]
的 \(\frakp\)-指標與 \(\frakp'\)-指標都等於 \(f\)，而在 \(\Om\) 的所有其他素位處指標為 \(1\) [H, 17.7]。由定理 3，並結合對 \(\frakp\) 的假設以及 \(\frakp'\) 的選擇，可知 \(\mathsf{K}\) 是 \(A\) 的分裂域。這樣就確實在群 \(\frakK\) 中找到了 \(\frakp\)-指標為 \(f\) 的代數 \(A\)。

\paragraph{定理 5.}
\emph{（唯一性與序定理.）} 若 \(\mathsf{K}\leqq \mathsf{K}'\)，則 \(\frakK\leqq\frakK'\)，反之亦然。

特別地，伽羅瓦域 \(K\) 與代數群 \(\frakK\) 的對應是可逆唯一的。

\paragraph{證明.}
a) 由 \(\mathsf{K}\leqq \mathsf{K}'\) 顯然推出 \(\frakK\leqq\frakK'\)；因為每個被 \(\mathsf{K}\) 分裂的代數 \(A\)，當然也被 \(\mathsf{K}'\) 分裂。

b) 反過來，設 \(\frakK\leqq\frakK'\)。那麼由分解定理，對相對判別式的每個非因子 \(\frakp\)，都有 \(f_{\frakp}\mid f'_{\frakp}\)。特別地，對幾乎所有滿足 \(f'_{\frakp}=1\) 的 \(\frakp\)，也有 \(f_{\frakp}=1\)。由通常的解析推理方法\footnote{參見 H. Hasse [前引第 6 注]，§ 25, III。}即可推出 \(\mathsf{K}\leqq \mathsf{K}'\)。

\paragraph{8.}
最後還要說明，主定理對 I. Schur\footnote{I. Schur, \emph{Arithmetische Untersuchungen über endliche Gruppen linearer Substitutionen}, Berl. Akad.-Ber. 1906。} 所研究的問題有實質推進；該問題是：一個有限群的絕對不可約表示可在哪些數域中實現。

\paragraph{定理 6.}
有限群 \(\frakG\) 的所有絕對不可約表示都可在圓分域中實現；例如，當 \(n\) 是 \(\frakG\) 的階且 \(h\) 足夠大時，總可在 \(n^h\) 次單位根域中實現。

\paragraph{證明.}
轉到 \(\frakG\) 的有理群環 \(G\) 後，\(\frakG\) 的絕對不可約表示 \(\Gamma_i\) 就成為半單代數 \(G\) 的簡單組成部分 \(G_i\) 的絕對不可約表示；它們的中心分別是相應特徵標的域 \(\Om_i\)。\footnote{參見：a) R. Brauer 與 E. Noether, \emph{Über minimale Zerfällungskörper irreduzibler Darstellungen}, Berl. Akad.-Ber. 1927, § 1。b) R. Brauer, \emph{Über Systeme hyperkomplexer Zahlen}, Math. Zeitschr. 30 (1929), Satz 3。c) E. Noether, \emph{Hyperkomplexe Größen und Darstellungstheorie}, Math. Zeitschr. 30 (1929), §§ 21, 24, 26。} 由於 \(\frakG\) 有限，這些中心當然都是圓分域，並且確切說都是 \(n\) 次單位根域的子域。

由定理 3 可知，若對 \(\Om_i\) 的每個素位 \(\frakp\)，一個關於 \(\Om_i\) 的循環域 \(Z_i\) 的 \(\frakp\)-次數 \(n_{i\frakp}\) 都是 \(G_i\) 的 \(\frakp\)-指標 \(m_{i\frakp}\) 的倍數，那麼 \(Z_i\) 是 \(G_i\) 的分裂域。\(m_{i\frakp}\) 只有在 \(G_i\) 關於 \(\Om_i\) 的基本理想的素因子處才可能不同於 \(1\)，因此當然只可能出現在 \(G\) 的絕對判別式的素因子處；而這個判別式整除 \(n^n\)--這裡 \(n^n\) 是 \(G\) 中一個（非極大）階的判別式\footnote{見 E. Noether [前引第 12c 注], § 26。}--所以至多隻涉及 \(n\) 的素因子 \(\frakp\)。為了使這些 \(\frakp\) 處的 \(n_{i\frakp}\) 成為 \(m_{i\frakp}\) 的倍數，只需規定這些 \(\frakp\) 在 \(Z_i\) 中以分別可被 \(m_{i\frakp}\) 整除的階數分歧。很容易看出，當 \(h\) 足夠大時，\(n^h\) 次單位根域就能做到這一點。

I. Schur 在上述論文的最後一個定理中指出，在所有當時已知情形中，只用 \(n\) 次單位根域就夠了。它是否總是成立，以及本文方法是否足以判定這個問題，仍有待進一步研究。

\begin{center}
1931 年 11 月 11 日收稿。
\end{center}

%
\clearpage
%
\setcounter{footnote}{0}

%
%
%
%
\pdfbookmark[0]{39. 超複系統及其與交換代數和數論的關係}{paper-39}
\section*{39. 超複系統及其與交換代數和數論的關係}

\begin{center}
\emph{Verhandl. Intern. Math.-Kongreß Zürich 1 (1932), S. 189--194}
\end{center}

\begin{center}
由 \textsc{Emmy Noether} 撰，哥廷根
\end{center}

\paragraph{1.}
超複系統，也就是代數的理論，在近幾年有了強勁發展；但是直到最近，這一理論對於交換性問題的意義才變得清楚。今天我要報告非交換理論對於交換理論的這種意義；具體說，我將沿着兩個可追溯到 Gauss 的經典問題來說明，即主屬定理以及與它密切相關的範數定理。這些問題在歷史進程中不斷改變表述形式：在 Gauss 那裡，它們作為二次型理論的收尾出現；隨後，它們在通過類域論刻畫相對循環數域和阿貝爾數域時起到本質作用；最後，它們可以表述為關於自同構以及關於代數分裂的定理，而這一最後表述同時把這些定理推廣到任意相對伽羅瓦數域。

通過這幅稍後還要展開的草圖，我同時想說明把非交換理論應用於交換理論的\emph{原則}：\emph{藉助代數理論，人們力圖為關於二次型或循環域的已知事實取得不變而簡單的表述，也就是說，取得只依賴於代數結構性質的表述。一旦這些不變表述已經得到證明，而且上述例子正是如此，那麼這些事實向任意伽羅瓦域的轉移也就自動獲得。}

\paragraph{2.}
在展開細節之前，我還要先概述各種方法和進一步結果。首先必須指出，主要困難在於為一般伽羅瓦域取得相應表述；若沒有超複方法，原本甚至沒有著手點。在上面列舉的例子中，向非交換理論的根本過渡，是通過\emph{同時考察域與群}而取得的，具體說，是藉助“交叉積”及其乘法常數，也就是“因子系”（見 3）。這樣得到一個關於基域的單純正規代數，而每一個這樣的代數本質上也都能這樣生成。這樣的交叉積最早由 Dickson 考察，\footnote{參見他的著作 \emph{Algebren und ihre Zahlentheorie}, Zürich 1927, § 34。}而因子系理論則由 Speiser、Schur 和 R. Brauer\footnote{例如參見 R. Brauer, \emph{Untersuchungen über die arithmetischen Eigenschaften von Gruppen linearer Substitutionen}, 以及該文注 2 所列文獻，Math. Z. 28 (1928)。}從完全不同的出發點發展出來，即從絕對不可約表示的問題出發。只有通過這兩種理論的融合，才可能得到一個足夠簡單且影響範圍足夠廣的結構，從而能夠用它處理交換性問題。\footnote{這一結構最初由我在 1929/30 年冬季學期的一門課程中發展出來，後來收入 H. Hasse, \emph{Theory of cyclic algebras}, Trans. Amer. Math. Soc. 34 (1932), 第 2 章。關於本報告涉及的整個領域，M. Deuring 關於超複數及其數論應用的一篇報告將提供綜述，該報告擬收入 \emph{Ergebnisse der Mathematik} 叢書。}

與此同時，對一些已知事實也得到新的、透明的證明：這裡我想指出 H. Hasse 給出的循環域互反律的超複證明，這一證明即將發表於 Math. Ann.，\footnote{H. Hasse, \emph{Die Struktur der R. Brauerschen Algebrenklassengruppe über einem algebraischen Zahlkörper (insbesondere Normenrestsymbol und Reziprozitätsgesetz)}. Math. Ann. 107 (1932/33)。}它依靠的是基於交叉積理論的範數剩餘符號的不變表述。還要指出 C. Chevalley 最近在同一基礎上給出的局部類域論的超複奠基；不過為此還必須發展若干關於因子系的新代數定理。\footnote{C. Chevalley, \emph{Sur la théorie du symbole de restes normiques}; 將發表於 J. f. Math. 169。}同時我還必須作一個限制性說明：只靠交叉積方法，看來並不能得到伽羅瓦數域的完整理論。這一點來自 Artin 尚未發表的新結果；這些結果按照上述原則接續 Hasse 的證明，但只能給出計數相等，而不是完整的同構定理。

能夠給出完整同構，準確說是算子同構的方法，已經存在於代數方面。這裡涉及的是 A. Speiser\footnote{A. Speiser, \emph{Gruppendeterminante und Körperdiskriminante}. Math. Ann. 77 (1916)。}若干思路的延續，即把伽羅瓦域理解為“伽羅瓦模”：也就是說，把它看成基域上的一個模，其中伽羅瓦群的代換作為算子出現。於是域與群環（群代數）之間存在算子同構，其意義是：元素之間有一一對應，基域上的線性形式相互對應，並且域中伽羅瓦群的代換對應於群環中的乘法。我提出的這一定理\footnote{E. Noether, \emph{Normalbasis bei Körpern ohne höhere Verzweigung}, Satz 3, J. f. Math. 167 (1932)（證明中含有一處缺口）。}已經由 M. Deuring 證明，\footnote{M. Deuring, \emph{Galoissche Theorie und Darstellungstheorie}. Math. Ann. 107 (1932)。}他在此基礎上給出了伽羅瓦理論的一個證明，其中算子同構實現了群與域的對應。進一步的結構定理，同樣由 Deuring 給出，與 Artin \(L\)-級數的形式事實相平行，並給出通向 Artin 導子的結構性途徑。這些 Artin \(L\)-級數和導子\footnote{E. Artin, \emph{Über eine neue Art von \(L\)-Reihen}, Math. Sem. Hamburg 3 (1924). \emph{Zur Theorie der \(L\)-Reihen mit allgemeinen Gruppencharakteren}, 同刊, 8 (1931). \emph{Die gruppentheoretische Struktur der Diskriminanten algebraischer Zahlkörper}, J. f. M. 164 (1931).}由一般群特徵標構成，和 Speiser\textsuperscript{6)} 的工作一起，是數論與表示論之間的第一次聯繫，是越過阿貝爾域範圍的第一次推進。它們給整個發展以強烈推動；特別是伽羅瓦模理論正是以此為取向發展起來的。

\paragraph{3.}
現在我要具體追蹤置於開頭的兩個問題，即範數定理和主屬定理。首先是\emph{交叉積的定義}：設 \(K/k\) 是一個 \(n\) 次伽羅瓦域，\(\Gg\) 為其群。所謂交叉積，是指把 \(K\) 和 \(\Gg\) 同時嵌入一個代數 \(A\)，使得 \(K\) 的自同構成為內自同構。令 \(u_{S_1},\ldots,u_{S_n}\) 表示與這 \(n\) 個群元素相應的符號；於是先把 \(A\) 設為關於 \(K\) 的秩 \(n\) 線性形式模：
\[
\tag{1}
A=u_{S_1}K+\ldots+u_{S_n}K .
\]
也就是說，\(A\) 由所有線性形式 \(u_{S_1}a_1+\ldots+u_{S_n}a_n\) 組成，其中 \(a_i\) 在 \(K\) 中任意取值。通過要求由 \(u_S\)，更一般地由 \(u_SK^\ast\)\footnote{\(K^\ast\) 由 \(K\) 去掉零元而得；這一記號通用。}生成內自同構，\(A\) 便成為一個環，也就是一個關於 \(k\) 的秩 \(n^2\) 代數。這個要求表示為
\[
\tag{2}
u_S^{-1}zu_S=z^S,\qquad \text{或}\qquad zu_S=u_Sz^S
\]
對每個 \(z\in K\) 成立，\footnote{\(z^S\) 照常表示由代換 \(S\) 作用於 \(z\) 後所得的元素。}並且乘法滿足
\[
\tag{3}
u_Su_T=u_{ST}a_{S,T}\quad \text{其中 }a_{S,T}\text{ 在 }K^\ast\text{ 中},
\]
\[
\tag{4}
\begin{gathered}
a_{ST,R}\,a_{S,T}^{\,R}=a_{S,TR}\,a_{T,R}\\
\text{（由结合律 \([u_Su_T]u_R=u_S[u_Tu_R]\) 得到）}.
\end{gathered}
\]
\(A\) 稱為 \(K\) 與 \(\Gg\) 在因子系 \(a_{S,T}\) 下的交叉積；可以證明，\(A\) 成為 \(k\) 上的一個單純正規代數，也就是某個全矩陣環 \(D_r\)，其係數取自在相應除代數 \(D\) 中，並且 \(K\) 成為極大交換子域，因而成為分裂域（即把係數域 \(k\) 擴張到一個與 \(K\) 同構的域後，所得代數分裂為以中心為係數的全矩陣環）。反過來，對預先給定的除代數 \(D\)，總存在矩陣環 \(D_r\)，可按上述方式生成作為交叉積。

如果從 \(u_S\) 過渡到 \(v_S=u_Sc_S\)，其中 \(c_S\) 在 \(K^\ast\) 中，這仍生成同一個自同構，但會產生“相伴”的因子系：
\[
\tag{5}
\bar a_{S,T}=a_{S,T}\frac{c_S^{\,T}c_T}{c_{ST}} .
\]
相伴的因子系合併為一個類 \((a)\)；同樣，把所有與 \(A\) 相似的代數，也就是所有 \(r=1,2,\ldots\) 的 \(D_r\)，合併為一個類 \(\mathfrak A\)。\emph{類 \(\mathfrak A\) 與類 \((a)\) 一一對應：具有固定分裂域 \(K\) 的這些類，在直接乘法下構成一個阿貝爾群，它同因子系類的逐項乘積群同構。單位元是分裂代數類，或者說是所有變換量所構成的系統}
\[
        \frac{c_S^{\,T}c_T}{c_{ST}} .
\]
這裡涉及的是 R. Brauer 所指出的代數類群。

\paragraph{4.}
現在我通過\emph{特殊化}到\emph{循環分裂域}來說明同\emph{範數概念}的聯繫，並由此按照開頭說明的原則得到推廣\emph{範數定理}的表述。若 \(Z\) 是循環域，\(S\) 是其群的一個生成代換，那麼可以讓 \(S\) 的冪對應於 \(u\) 的冪，即
\[
\begin{aligned}
(1')\qquad &A=Z+uZ+\cdots+u^{n-1}Z\\
(2')\qquad &zu=uz^S\\
(3')\qquad &u^n=a\\
(4')\qquad &a\text{ 位於基域 }k^\ast\text{ 中}\\
(5')\qquad &\bar a=a\cdot \N(c),\quad \text{當置 }v=uc\text{ 時}.
\end{aligned}
\]
因此這裡每一個因子系只由一個位於基域中的元素 \(a\) 組成，記作 \(A=(\alpha,Z)\)；因子系的單位類由來自 \(Z^\ast\) 的範數給出，而代數類群同商群
\[
        k^\ast/N(Z^\ast)
\]
同構。於是循環代數 \((\alpha,Z)\) 當且僅當 \(a\) 是 \(Z\) 中某個元素的範數時分裂。

範數與分裂之間的這一聯繫給出了“範數定理”的表述，即\emph{關於分裂代數的定理：若一個代數在每個素位處分裂，則它本身分裂。}這裡的“素位”按數論中的通常方式定義：把基域 \(k\) 替換為它的 \(p\)-進擴張 \(k_{\mathfrak p}\)，其中 \(\mathfrak p\) 是 \(k\) 的一個素理想；對有限多個無限素位，則相應地把 \(k\) 及其共軛域用實數域來擴張。

事實上，循環域的範數定理包含在這一定理之中。因為對循環代數 \((\alpha,Z)\) 而言，根據上面的說明，這一定理等價於如下斷言：若 \(a\) 在每個有限或無限素位處都是 \(p\)-進範數，則 \(a\) 是 \(Z\) 中某個數的範數；或者不經過 \(p\)-進過渡地說：\emph{若 \(a\) 按 \(k\) 的每個素理想 \(\mathfrak p\) 都是範數剩餘（並滿足某些符號條件），則 \(a\) 是 \(Z\) 中某個數的範數。}後一種表述正是類域論中利用熟知的解析工具所證明的範數定理。而關於分裂代數的一般定理，則可由這一循環特殊情形通過純代數-算術的考察得到。\footnote{R. Brauer, H. Hasse, E. Noether, \emph{Beweis eines Hauptsatzes in der Theorie der Algebren}. J. f. M. 167 (1932)。}Hasse 已由此得出第一個重要推論：\emph{代數數域上的每個單純正規代數都是循環的。}正是在尋找這個長期被猜測事實的證明時，才產生了上述一般表述。

\paragraph{5.}
從關於分裂代數的定理還可推出第二個後果，即開頭提出的主屬定理，仍然只需純代數-算術的推理。\footnote{該證明擬發表於 Math. Ann.。}它的不變表述基於這樣一個事實：定義交叉積的關係 (2) 到 (5) 完全是乘法性的；因此，如果把 \(K^\ast\) 替換為一個阿貝爾群 \(\mathfrak J\)，只要該群的自同構群含有一個同 \(\Gg\) 同構的子群，這些關係仍有意義。於是取代 (1) 的，是群論意義下“用 \(\mathfrak J\) 擴張 \(\Gg\)”。若把 \(\mathfrak J\) 取為 \(K\) 中所有理想構成的群，因子系便成為理想系統；\(\mathfrak J\) 中的一個類劃分誘導出因子系的一個類劃分，而且一般會得到比原來的理想類劃分更細的劃分。因為要求 \(u_S\) 的乘法對絕對理想類是單值的，只說變換量 \(c_S^{\,T}c_T/c_{ST}\)，也就是元素因子系的單位類，落在理想因子系的單位類中。事實上，已知特殊情形表明，稍微不那麼細的劃分已經足夠。我作如下定義：\emph{把那些在 \(K\) 的所有有限和無限分歧素位處生成分裂代數的元素 \(a_{S,T}\)，計入因子系的單位類。}

這樣產生的 \(\Gg\) 的擴張記作 \(\Gg^\ast\)；\((\mathfrak c)\) 表示 \(\mathfrak c\) 的絕對理想類。於是：

\emph{主屬定理的不變表述為：在這樣定義的類劃分下，若通過代換 \(v_S=u_S(\mathfrak c_S)\) 產生 \(\Gg^\ast\) 的一個自同構，其中所有這些 \((\mathfrak c_S)\) 構成主屬，那麼這個自同構是內的，並且由一個理想類 \((\mathfrak b)\) 生成。}

已知特殊情形可由一個等價而稍顯明白的表述推出：\emph{若由理想類 \((\mathfrak c_S)\) 形成的變換量 \((\mathfrak c_S^{\,T})(\mathfrak c_T)/(\mathfrak c_{ST})\) 屬於因子系的單位理想類，則存在一個理想類 \((\mathfrak b)\)，使得 \((\mathfrak c_S)\) 成為符號意義下的 \((1-S)\) 次冪：對 \(\Gg\) 中所有 \(S\)，都有 \((\mathfrak c_S)=(\mathfrak b)/(\mathfrak b^{S})\)。}因為前提恰好表達了自同構性質；而這個自同構是內的這一事實，則由
\[
        v_S=(\mathfrak b)^{-1}u_S(\mathfrak b)
        =u_S(\mathfrak b)^{1-S}=u_S(c_S)
\]
來表達。

特殊化到循環情形就得到（注意歸一化）：若 \(N([\mathfrak c])\) 位於因子系的單位理想類中，則 \((\mathfrak c)\) 成為符號意義下的 \((1-S)\) 次冪：
\[
        (c)=(\mathfrak b)^{1-S}.
\]

\paragraph{6.}
為了由此轉到循環域和二次型的熟知表述，我首先指出，定理雖然是針對完整理想類說出的，但如果像通常那樣只限於同分歧素位互素的理想，它的內容仍保持不變。

這樣，這裡定義的二次域主屬就轉化為 Gauss 的主屬；因為理想類對應於二次型，而類的範數對應於可由這些型表示的數。單位類由在 \(K\) 的分歧素位處生成分裂代數的那些數給出，這就是說，這些可表示的數在分歧素位處成為二次剩餘；相應的型因而具有主型的全特徵，所以構成 Gauss 的主屬。符號意義下的 \((1-S)\) 次冪轉化為平方化，這是熟知的。

對於循環域，表述轉化為如下形式：主屬由所有理想類組成，其範數在有限和無限分歧素位處都成為範數剩餘。但這又等價於熟知的“按導子的範數剩餘”；通常的定理便這樣作為特殊化得到。

此外，即使在任意伽羅瓦域的一般情形中，也可以引入一個只由分歧素位組成的導子，使得在對因子系作歸一化時，對每個這樣的素位，射線只含有單位類中的元素。

這裡產生的問題是：這同概述 2 中提到的 Artin 導子有什麼聯繫？那些導子畢竟也是由同一些素理想組成的；因此也產生了它同伽羅瓦模理論，即第二種超複方法之間有什麼聯繫的問題。這兩種方法究竟能推進到何種程度，還有待將來說明。

%
\clearpage
%
\setcounter{footnote}{0}

%
%
%
%
\pdfbookmark[0]{40. 非交換代數}{paper-40}
\section*{40. 非交換代數}

\begin{center}
\emph{Math. Zs. 37 (1933), 514--541}
\end{center}

\begin{center}
\textbf{非交換代數}\\[0.5em]
由\\[0.25em]
哥廷根的 Emmy Noether 撰
\end{center}

眾所周知，交換代數的主要定理都包含在伽羅瓦理論之中；在伽羅瓦理論之前，是關於分解域和分裂域的理論，也就是關於那些足以使一個給定多項式分出一個一次因子，或者完全分解為一次因子的域的理論。

我在這裡把代數的相應部分發展到非交換情形，特別是超複情形。具體說，我原則上使用非交換方法，即在非交換域中的表示。最後我會說明，上面提到的交換代數定理，如何能夠與此完全平行地藉助在交換域中的表示來奠基。

這裡所依據的表示論，在一段簡短的自同構理論之後（§1），是對那種基於表示模理論的表示論的推進。\footnote{E. Noether, Hyperkomplexe Größen und Darstellungstheorie, Math. Zeitschr. 30 (1929), S. 641--692，下文引用為“表示論”。另參見 van der Waerden, \emph{Moderne Algebra II} 中的轉述。}尤其是，除直接表示外，我還考察互反表示；兩者可以相互歸約，而現在互反表示由互反表示模產生。這樣做的優點是：互反表示模，也就是一個雙模，也可以理解為按某個擴張環的一側模（§2）。由此，在非交換域中的表示被歸約為擴張環中的理想論；不可約表示類對應於擴張環的不可約理想類，這同超複系統在其交換系數域中表示時對系統本身成立的事實完全類比（§3 與 §4）。

由此，通過一個觀察（§5）可以得到非交換域上矩陣環的結構定理：\emph{每個子環都意味着由這個域給出的一個表示，或者由互反同構域給出的一個互反表示。}這個互反同構域構成了交換情形中極小分解域和分裂域的第一個類似物，因為它給出所有一次互反表示；並且當按中心的秩有限時，雖然此前並未假定中心有限，還給出分解為秩 \(1\) 直和項的完全分裂。這給出了域的伽羅瓦理論（§6），同時也表達瞭如下事實（§7）：互反同構的除代數，即按中心有限秩的域，在 R. Brauer 的代數類群中產生互逆類。

伽羅瓦理論的第二種奠基，對簡單系統也更一般有效，在正文中會先出現；它幾乎是關於可交換子環的一個定理的直接後果，而這個定理本身又幾乎直接接在前置的自同構考察之後。

到這裡為止，考察完全是非交換的。不過，對於\emph{交換}分裂域的問題，也通過 §5 中預先給出的觀察，即通過由除代數表示分裂域，而以非交換方式處理（§7）。最後是前面提到的向交換情形的轉移（§8），以及任意系統的分裂域理論（§9）；這同時給出了同通常交換情形表示論的聯繫。

R. Brauer 正是在交換情形中的這種表示論，以及“無理”的因子系之上，建立了分裂域理論。\footnote{參見我們的合著短文：\emph{Über minimale Zerfällungskörper irreduzibler Darstellungen}. Sitz. Ber. d. Preuß. Ak. d. Wiss. 1927, S. 221--228。（第 222 頁的注含有一份“說明報告”。主要定理是獨立且幾乎同時產生的。）另參見 R. Brauer, \emph{Über Systeme hyperkomplexer Zahlen}, § 3. Math. Zeitschr. 30 (1929), S. 79--107。A. A. Albert 後來獨立重新發現了這些定理。}由於這種交換的奠基，他和後來的 Albert 一樣，不得不假定中心是完全域；按非交換的奠基來看，這是一個不必要的限制。在同一限制下，同樣藉助交換情形中的表示論，在知道我的非交換域伽羅瓦理論以後，R. Brauer 和 K. Shoda 又進一步發展了這一理論：R. Brauer 給出了上面提到的可交換子環定理，K. Shoda 則獨立給出了也適用於半單系統的完整理論。\footnote{R. Brauer, \emph{Über die algebraische Struktur von Schiefkörpern}, § 2. Journ. f. Math. 166 (1932), S. 241--252。K. Shoda, \emph{Über die galoissche Theorie der halbeinfachen hyperkomplexen Systeme}. Math. Ann. 107 (1932), S. 252--258。J. Levitzki 也曾發展出這些結果。}

最後還要提到，對於域的情形，van der Waerden II 中有一段簡短敘述（§128），接續的是我 1928 年夏季學期的課程；我在那裡第一次講授這些問題。van der Waerden 的敘述相對於這門課帶來了一系列簡化，我有些是在第二門課程中（1929/30 年冬季學期），\footnote{第一門課由 G. Köthe 整理，第二門課由 M. Deuring 整理。}有些直到本文才採納並推進。特別是，把伽羅瓦理論的超複證明方法轉移到交換情形（§8），最早來自第二門課；在那裡，非交換證明（§6, 3.）相對於第一門課已經大為簡化。可交換子環定理（§5），以及以它為基礎的非交換伽羅瓦理論的第二種證明方法（§6, 1. 和 2.），是在第二門課之後，在知道 R. Brauer（注 3）和 van der Waerden §128 的工作後才加入的；不過在 §5, 3. 中，所需的有限性假設比那裡弱得多。

第一門課中的某些推理方式，即差分理想取交的方法，雖然較複雜，卻有一個優點：它能轉移到無限秩系統和整系統上。\footnote{這一方法實質上見於 G. Köthe, \emph{Schiefkörper unendlichen Ranges über dem Zentrum} § 5, Math. Ann. 105 (1931), S. 15--39。又見注 10）。這種取交方法現在已被一個幾乎平凡的事實（§4, 1.）替代，即雙側理想屬於不變模；這是 van der Waerden 首先注意到的。這樣，一切都可藉助更簡單的直和分解來奠基；但在無限秩和整情形中，這種方法失效。}對於交換的整系統，由此可達到理想微分與不同之間的聯繫，\footnote{參見布拉格報告，Jahresber. d. Deutsch. Math. Ver. 39 (1930), S. 17（斜分頁）。}我以後還會更詳細地討論這一點。

分裂域理論的主要應用，在於交叉積及其因子系的理論；該理論本身又構成數論應用的基礎。本文將不再進入這一方向。\footnote{交叉積理論是在第二門課中發展出來的；它經過少量改動，以便不預設本文給出的理論，收入 H. Hasse: \emph{Theory of cyclic algebras over an algebraic numberfield}, Kap. II. Transactions of the Amer. Math. Soc. 34 (1932), S. 171--214。M. Deuring 將在 \emph{Ergebnisse der Mathematik} 的一篇報告中給出一種更接近課程講稿的敘述。}

\subsection*{§1. 自同構、模與雙模}

基於表示模的表示論，依託於帶算子或不帶算子的阿貝爾群自同構環理論。為避免重複，下面把其中隱含使用的推理，也就是映射與運算法則之間的關係，表述為幾個簡單命題。

\paragraph{1. 乘法性映射。結合律。}
先設 \(\frG\) 是一個不帶算子的群，\(\frA\) 是它的絕對自同構域，也就是 \(\frG\) 到自身的所有同態組成的系統。因為乘積 \(\sigma\tau\) 由
\[
        g(\sigma\tau)=(g\sigma)\tau,
        \qquad g\in\frG,\quad \sigma,\tau\in\frA                         \tag{1}
\]
定義，並且顯然滿足結合律，所以 \(\frA\) 對乘法封閉。

眾所周知，當給定一組符號 \(\Theta,H,\ldots\) 構成的集合 \(\frB\)，使得聯結 \(g\Theta,gH,\ldots\) 都給出 \(\frG\) 中唯一確定的元素，併產生 \(\frG\) 到自身的同態：
\[
        (gh)\Theta=g\Theta\cdot h\Theta ,
\]
則 \(\frG\) 成為一個帶算子的群。如果算子域 \(\frB\) 對乘法封閉，那麼 \(\frB\) 到相應自同構域部分的唯一映射，當且僅當滿足結合關係
\[
        g(\Theta H)=(g\Theta)H, \qquad g\in\frG,\ \Theta,H\in\frB,                       \tag{1a}
\]
時，才是乘法同態。相應地，當處理左算子時，定義的是互反同態。

事實上，\(\frB\) 到 \(\overline{\frB}\) 的映射由對 \(\frG\) 中所有 \(g\) 滿足 \(g\Theta=g\sigma\) 給出，因而也由
\[
        (g\Theta)H=(g\sigma)\tau=g(\sigma\tau)
\]
給出；所以 (1a) 成為乘法同態的充要條件。左算子產生互反同態，是因為寫在左邊的自同構 \(\sigma^\ast g,\tau^\ast\sigma^\ast g\) 必須從右向左讀：\(\tau^\ast\sigma^\ast g=g\sigma\tau\)，而算子始終從左向右讀。

\paragraph{2. 算子同態性映射。}
若 \(\frG\) 是帶算子的群，則算子同態性由
\[
        (g\Theta)\sigma=(g\sigma)\Theta,                                          \tag{2}
\]
定義；對左算子則由
\[
        (\Theta g)\sigma=\Theta(g\sigma)                                           \tag{2*}
\]
定義，也就是說，由交換相聯性或者貫穿結合律來定義。由 1. 的推理方式，也就是映射到自同構域，立即得到對兩個不同算子域的如下命題：

若給定兩個算子域 \(\frB\) 與 \(\frC\)，其元素分別為 \(\Theta,H,\ldots\) 與 \(\bar\Theta,\bar H,\ldots\)，則 \(\frC\) 的元素產生 \(\frB\)-自同構，同時 \(\frB\) 的元素產生 \(\frC\)-自同構，當且僅當這兩個域與 \(\frG\) 交換相聯，或者說滿足貫穿結合律：
\[
        (g\Theta)\bar H=(g\bar H)\Theta,                                                    \tag{2a}
\]
相應地，對左算子滿足
\[
        (\Theta g)\bar H=\Theta(g\bar H).                                                    \tag{2a*}
\]

\paragraph{3. 按環的模與雙模。}
環 \(\frR\) 上的右模 \(\frM\)，是一個加法阿貝爾群，以 \(\frR\) 的元素作為右算子，並且除結合關係 (1a) 外還滿足分配關係
\[
        (g+h)\rho=g\rho+h\rho,\qquad
        g(\rho+\sigma)=g\rho+g\sigma.                                  \tag{3a}
\]
左模的情形相應定義。

雙模要分為兩種。若一個群同時是按兩個環 \(\frR\) 與 \(\frS\) 的右模，並且除分別對 \(\frR\) 與 \(\frS\) 成立的聯結外，還滿足
\[
        (m\rho)\sigma=(m\sigma)\rho,
\]
則稱為雙模。若它是 \(\frR\)-左模、\(\frS\)-右模，並且還滿足貫穿結合律
\[
        \rho(m\sigma)=(\rho m)\sigma,
\]
也稱為雙模。

這些聯結的意義來自如下事實：在阿貝爾群的情形中，絕對自同構域構成一個環，即絕對自同構環。\footnote{對於非阿貝爾群，涉及的是一種“廣義”環；參見 H. Fitting 一篇將發表於 Math. Annalen 的論文。[Math. Annalen 107, S. 514--542.] }具體地，由映射推理得到：

若一個（加法記號的）阿貝爾群 \(\frM\) 的算子域 \(\frR\) 是一個環，那麼 \(\frR\) 到絕對自同構環的一個子集 \(\overline{\frR}\) 的唯一映射，當且僅當 \(\frM\) 是按 \(\frR\) 的右模時，也就是當 (1a) 與 (3a) 成立時，才是環同態；對左模，則是互反環同態。

若 \(\frM\) 是按環 \(\frR\) 與 \(\frS\) 的右模，那麼 \(\frM\) 當且僅當成為雙模，即滿足 (2a)，才使得 \(\frR\) 可以環同態地映入 \(\frS\)-自同構環的一個子環，同時 \(\frS\) 可以環同態地映入 \(\frR\)-自同構環的一個子環。若 \(\frM\) 是按 \(\frS\) 的右模、按 \(\frR\) 的左模，則雙模條件，也就是 (2a*)，意味着 \(\frS\) 到 \(\frR\)-自同構環一個子環的同態映射，以及 \(\frR\) 到 \(\frS\)-自同構環一個子環的互反同態映射。

由此特別得到：

1) 若 \(\frR\) 是有單位元的環，則作為 \(\frR\)-左模，\(\frR\) 與它的自同構環直接同構；作為 \(R\)-右模，則與它的自同構環互反同構。

2) 若 \(\frR\) 是一個一般非交換域 \(\mathsf{A}\) 上的全矩陣環，
\[
        \frR=\sum c_{ik}\mathsf{A}=\sum \mathsf{A} c_{ik},
\]
則 \(\mathsf{A}\) 與單右理想的自同構域互反同構，並與單左理想的自同構域直接同構。

\paragraph{4. 從右雙模過渡到按乘積環的一側模。}
右雙模的重要性本質上基於這一過渡。

\paragraph{定理。}
若 \(\frM\) 是按一個環 \(\frT\) 的右模，而 \(\frT\) 含有兩個元素兩兩可交換的子環 \(\frR\) 與 \(\frS\)，則 \(\frM\) 也可理解為按 \(\frR\) 與 \(\frS\) 的雙模。反過來，若已經有按 \(\frR\) 與 \(\frS\) 的雙模，並且存在一個乘積環 \(\frT\)，其中 \(\frR\) 與 \(\frS\) 的元素兩兩可交換，則 \(\frM\) 也可理解為按 \(\frT\) 的右模；這裡 \(\frT\) 的聯結延拓給定的 \(\frR\) 與 \(\frS\) 的聯結。

第一條斷言來自如下事實：按定義，\(\frT\) 可同態映入絕對自同構環的一個子環 \(\overline{\frT}\)，並且 \(\frR,\frS\) 分別變為元素兩兩可交換的子環 \(\overline{\frR},\overline{\frS}\)。這正是刻畫雙模的關係 (2a)。

反過來，若給定 \(\frM\) 作為 \(\frR,\frS\)-雙模（右模），則 \(\frR,\frS\) 又可同態映入絕對自同構環中元素兩兩可交換的子環 \(\overline{\frR},\overline{\frS}\)。在絕對自同構環中，對於 \(\overline{\frR},\overline{\frS}\) 存在乘積環\footnote{所謂乘積環，僅指由兩個具有公共上環的環所生成的環；這個乘積不必是直積。-- 給定交時，乘積環並不總是存在，下例說明這一點：設 \(\frako\) 為整數環，並令 \(\frR=\frako[x]\)、\(\frS=\frako[y]\)，其中 \(2x-1=0\)、\(2y-1=0\)；如果不規定 \(x\) 與 \(y\) 之間的聯結，那麼 \(\frako\) 是 \(\frR\) 與 \(\frS\) 的交。但若 \(\frR\) 與 \(\frS\) 位於一個公共上環中，就有 \((2x-1)y-x(2y-1)=0\)，因而 \(x=y\)。所以不存在以 \(\frako\) 為交的乘積環。} \(\overline{\frT}\)，由於元素兩兩可交換，它由所有如下形式的元素組成：
\[
        \sum \bar r_i\bar s_i+\bar r+\bar s
\]
（若 \(\overline{\frR},\overline{\frS}\) 有單位元，則附加項 \(\bar r,\bar s\) 可省去）。若現在也存在一個乘積環 \(\frT\)，其中 \(\frR\) 與 \(\frS\) 元素兩兩可交換，那麼它相應地由所有 \(\sum r_i s_i+r+s\) 形式的元素組成。於是同態
\[
        \frR\to\overline{\frR},\qquad \frS\to\overline{\frS}
\]
可由對應
\[
        \sum r_i s_i+r+s
        \to
        \sum \bar r_i\bar s_i+\bar r+\bar s
\]
補成一個同態 \(\frT\to\overline{\frT}\)；因此聯結
\[
        m\Bigl(\sum_i r_i s_i+r+s\Bigr)
        =\sum_i (m r_i)s_i+m r+m s
\]
是唯一的，並且按照 §1, 3. 定義了一個 \(\frT\)-模，包含給定的 \(\frR,\frS\)-模。

\subsection*{§2. 互反表示與直接表示}

\paragraph{1. 線性形式模。}
從 §1 所發展的理論過渡到表示論，基於把那裡的模專門化為線性形式模，並考察它們的自同構環。

一個 \(\frS\)-右模 \(\frM\) -- 其中 \(\frS\) 是有單位元的環 -- 若 \(\frM\) 是 \(n\) 個單項 \(\frS\)-模的直和：
\[
        \frM=m_1\frS+\ldots+m_n\frS,
\]
並且 \(m_i\frS\) 與 \(\frS\) 算子同構，也就是說由 \(m_i s=0\) 總推出 \(s=0\)，則稱為 \(\frS\) 中的線性形式模。由此可見，\(\frS\) 不僅可以同態地，而且可以同構地映入絕對自同構環的一個子環。

\paragraph{定理 1。}
若 \(\frM\) 是 \(\frS\) 中的線性形式模，則 \(\frM\) 的 \(\frS\)-自同構環 \(\mA\) 與 \(\frS\) 中所有 \(n\) 階矩陣組成的環 \(\mbarA\) 互反同構，而不僅僅是同態。

\(\frM\) 的任意 \(\frS\)-自同構 \(\alpha\)，完全由如下對應描述：
\[
        m_i\to \overline m_i=m_i\alpha;
\]
因為由此按定義得到
\[
        \sum m_i s_i\to
        \sum \overline m_i s_i
        =\sum (m_i\alpha)s_i
        =\sum (m_i s_i)\alpha.
\]
若現在用矩陣記法寫成
\[
        (m_1\alpha,\ldots,m_n\alpha)=(m_1,\ldots,m_n)A,
\]
那麼，當 \(\alpha\) 遍歷所有自同構時，對應 \(\alpha\to A\) 給出 \(\mA\) 與 \(\mbarA\) 之間的一一對應。因為既已假定 \(\frM\) 為線性形式模，\(A\) 就由 \(\alpha\) 唯一確定；並且 \(\frS\) 中每個 \(n\) 階矩陣 \(A\) 都產生一個自同構。進一步有：
\[
        \alpha+\beta\to A+B,\qquad
        \alpha\beta\to BA.
\]
後一條由
\[
        (m_1\alpha\beta,\ldots,m_n\alpha\beta)
        =(m_1,\ldots,m_n)AB
        =(m_1,\ldots,m_n)BA
\]
得到，其中用到 \(ms\beta=m\beta s\)。

\paragraph{2. 表示與表示模。}
若把環 \(\frR\) 在 \(\frS\) 中的 \(n\) 階互反表示或直接表示，理解為 \(\frR\) 到 \(\frS\) 中所有 \(n\) 階矩陣組成的環的某個子環中的互反環同態或直接環同態，那麼定理 1 也可以表述為：

\paragraph{定理 1'。}
\(\frS\) 中 \(n\) 階線性形式模的自同構環 \(\mA\)，允許由 \(\frS\) 中所有 \(n\) 階矩陣組成的全矩陣環 \(\mbarA\) 給出一個同構的（忠實的）互反表示。

從這裡出發，直接表示的通常定理也同樣得到互反表示的表述。

\paragraph{定義。}
\(\frS\) 中的線性形式模 \(\frM\)，如果是按 \(\frR\) 與 \(\frS\) 的雙模，並且同時是 \(\frR\)-右模與 \(\frS\)-右模，則稱為 \(\frR\) 在 \(\frS\) 中的互反表示模。如果 \(\frM\) 是雙模，並且同時是 \(\frR\)-左模、\(\frS\)-右模，則稱為直接表示模。

\paragraph{定理 2。}
每個互反表示模或直接表示模，都產生 \(\frR\) 在 \(\frS\) 中的一類等價的互反表示或直接表示；並且所有互反表示或直接表示都是這樣產生的。

先設 \(\frM\) 是互反表示模。按 §1, 3.，\(\frR\) 可直接同態地映到 \(\frM\) 的 \(\frS\)-自同構環的一個子環 \(\overline{\frR}\) 上；按 §2 的定理 1'，\(\overline{\frR}\) 允許一個互反的（忠實的）表示，因此這也成為 \(\frR\) 的互反表示。反過來，若給定一個互反表示 \(\frR\to\frR^*\)，則把它同 \(\frR^*\) 到 \(\frS\)-自同構環的一個子環 \(\overline{\frR}\) 的互反同構複合，就得到 \(\frR\) 到 \(\overline{\frR}\) 上的直接同態；於是 \(\frM\) 成為一個 \(\frR\)-模，並且確切地說成為雙模，也就是按 §1, 3. 的表示模。過渡到其他 \(\frS\)-基底給出等價表示的類。

若討論的是直接表示模，則 \(\frR\) 被互反同態地映到 \(\overline{\frR}\) 上。於是同 \(\overline{\frR}\) 的互反表示複合，給出 \(\frR\) 的一個直接表示。反過來，由一個直接表示得到到 \(\overline{\frR}\) 上的互反同態映射；隨後 \(\frM\) 成為按 §1, 3. 的直接表示模。

\paragraph{說明。}
當然，兩種表示可以相互歸約：\(\frR\) 的直接表示，就是與 \(\frR\) 互反的環的互反表示，反過來也一樣。另一種歸約是在從 \(\frS\) 過渡到一個互反環時，把矩陣替換為它們的轉置。這對應於從一個 \(\frR\)-左、\(\frS\)-右模過渡到一個 \(\frR\)-和 \(\frS\)-左模。

\paragraph{3. 從互反表示模過渡到按擴張環的模。}
互反表示模的優點，首先在於 §1, 4. 所允許的向擴張環的過渡。對於這裡的特殊情形，即 \(\frM\) 是按 \(\frS\) 的線性形式模，得到的正是表示模的過渡定理：

\emph{若 \(\frM\) 是某個環 \(\frT\) 上的右模，而 \(\frT\) 含有兩個元素兩兩可交換的子環 \(\frR\) 與 \(\frS\)，並且 \(\frM\) 成為按 \(\frS\) 的線性形式模，則 \(\frM\) 也可理解為 \(\frR\) 到 \(\frS\) 的互反表示模。}

\emph{反過來，若給定 \(\frR\) 到 \(\frS\) 的一個互反表示模，並且存在一個乘積環 \(\frT\)，其中 \(\frR\) 與 \(\frS\) 元素兩兩可交換，則 \(\frM\) 也可理解為 \(\frT\)-模，並且其 \(\frT\)-聯結是給定的 \(\frR\)-、\(\frS\)-聯結的延拓。}

表示論中本質使用的是如下

\paragraph{表示模過渡定理的推論。}
若 \(\frM\) 是 \(\frR\) 在 \(\frS\) 中的互反表示模，同時又是以 \(\frT\) 為乘積環的 \(\frT\)-模，則 \(\frM\) 到某個模 \(\frN\) 的 \(\frR\)-、\(\frS\)-同構，等價於 \(\frT\)-同構。\emph{因此，每一類同構的 \(\frT\)-模，都對應於 \(\frR\) 在 \(\frS\) 中的一類互反表示；反過來也一樣。}

這些定理同 §1, 2. 和 3. 結合，給出了下面要用的基本聯繫：與一個表示可交換的矩陣，同 \(\frR\)-、\(\frS\)-自同構之間的聯繫。

\paragraph{關於可交換矩陣的定理。}
\emph{若 \(\frR\to\frR^*\) 是 \(\frR\) 在 \(\frS\) 中的一個 \(n\) 階互反表示，並且 \(\frB^*\) 表示所有與 \(\frR^*\) 元素兩兩可交換的矩陣（即 \(\frS\) 中的 \(n\) 階矩陣）組成的環，則 \(\frB^*\) 給出生成的表示模的 \(\frR\)-、\(\frS\)-自同構的一個互反同構表示；這裡所謂 \(\frR\)-、\(\frS\)-自同構，也就是那些同時為 \(\frR\)-自同構的 \(\frS\)-自同構；若乘積環 \(\frT\) 存在，則也就是 \(\frT\)-自同構。}

設 \(\frB\) 是 \(\frM\) 的所有 \(\frR\)-、\(\frS\)-自同構組成的環。作為 \(\frS\)-自同構環 \(\mA\) 的子環，\(\frB\) 允許一個在 \(\frS\) 中的互反同構表示。按 §1, 2. 和 3.，\(\frB\) 由 \(\mA\) 中所有同 \(\frR\) 交換相聯、也就是滿足關係 (2a) 的自同構組成。因此，\(\frB\) 由所有與 \(\overline{\frR}\) 中那些自同構元素兩兩可交換的自同構組成；這裡 \(\overline{\frR}\) 表示 \(\frR\) 在 \(\mA\) 中的像。由於 \(\frB\) 與 \(\overline{\frR}\) 的表示涉及的是互反同構，而不僅僅是同態，這就證明了所要證明的事實。

\paragraph{說明。}
\(\frR\)-、\(\frS\)-自同構意味着一種對應 \(m_i\to m_i'\)，使得藉助這些 \(m_i\) 得到同一個表示 \(\frR^*\)；這些 \(m_i'\) 不必一定給出 \(\frM\) 的一個 \(\frS\)-基底，這對應於 \(\frB^*\) 中的矩陣不必一定可逆這一事實。若和 \(\sum m_i'\frS\) 不再是直和，則“表示”要在轉義中理解：
\[
        (m_1',m_2',\ldots m_n')a=(m_1',\ldots m_n')A
\]
仍然成立，但 \(A\) 不再由 \(a\) 唯一確定。

關於過渡定理還有下面的說明：對角矩陣 \(E\cdot s\) 給出 \(\frS\) 的一個\emph{直接同構}表示，也就是 \(\frS\) 的“恆等”表示。這裡之所以是\emph{直接}同構，是因為 \(\frM\) 的 \(\frS\)-自同構環含有一個與 \(\frS\) 互反同構的子環，而這裡所說的正是這個子環的互反表示。事實上，\(\frM\) 的每個單項直和項 \(m_i\frS\) 作為右模都與 \(\frS\) 算子同構，因此它的 \(\frS\)-自同構環也與 \(\frS\) 互反同構（§1, 3.，推論 1）。

因此，由 \(\frR\) 與 \(\frS\) 唯一定義的、從 \(\frT\) 到 \(\frS\) 中矩陣的對應，\emph{並不是環 \(\frT\) 的一個表示，而是 \(\frR\) 的互反表示到一個按 \(\frS\) 算子同態的表示的延拓：}
\[
        \sum r_i s_i\to \sum R_i s_i.
\]
特別是，\(\frR\) 的互反表示也就對於位於 \(\frR\) 中心內的 \(\frR\) 與 \(\frS\) 的交 \([\frR,\frS]\)，具有算子同態性。

\subsection*{§3. 按域的模}

以下只涉及在一般非交換域中的表示。因此先彙集若干關於按域的線性形式模的簡單定理。

\paragraph{1. 子模相對於全模給定基的正規基。}
設 \(A\) 為域，並且
\[
        \frN=x_1A+\ldots+x_nA
\]
是秩 \(n\) 的線性形式模；再設
\[
        \frL=z_1A+\ldots+z_\ell A
\]
是秩 \(\ell\le n\) 的子模。如果 \(z_i\) 具有如下形式
\[
        z_i=x_i-(x_{\ell+1}\alpha_{i,\ell+1}+\ldots+x_n\alpha_{i,n}),
        \qquad i=1,\ldots,\ell,
\]
則稱這些 \(z_i\) 是相對於 \(x\) 的正規基。每個模 \(\frL\)，在適當重排 \(x_i\) 後，都有一個正規基。事實上，在適當重排後有
\[
        \frN=\frL+x_{l+1}\mathsf{A}+\ldots+x_n\mathsf{A},
\]
也就是
\[
        x_i\equiv x_{l+1}\alpha_{i,l+1}+\ldots+x_n\alpha_{i,n}
        \pmod{\frL},\qquad i=1,\ldots,\ell,
\]
因此相應的
\[
        z_i=x_i-(x_{\ell+1}\alpha_{i,\ell+1}+\ldots+x_n\alpha_{i,n})
\]
位於 \(\frL\) 中，並且窮盡 \(\frL\)。

\paragraph{2. 擴張模。}
仍設 \(A\) 為域，\(P\) 為子域。秩 \(n\) 的 \(A\)-模 \(N\)，如果 \(M\) 是 \(N\) 中同秩的子模，則稱為 \(P\)-模 \(M\) 的擴張模，記作 \(N=M_{\mathsf{A}}\)。若
\[
        M=y_1P+\ldots+y_nP,
\]
則有
\[
        N=M_A=y_1A+\ldots+y_nA.
\]
反過來，對於每個 \(P\)-模 \(M=y_1P+\ldots+y_nP\)，在模同構意義下唯一存在擴張模 \(M_{\mathsf{A}}\)。因為模 \(y_1\mathsf{A}+\ldots+y_n\mathsf{A}\) 含有一個同 \(M\) 同構的子模，可以把它與 \(M\) 等同起來。

\paragraph{推論 a)。}
若 \(M\) 中的元素 \(z_1,\ldots,z_s\) 在 \(P\) 上線性無關，則它們在 \(N=M_A\) 中也在 \(A\) 上線性無關（因為它們可以補成 \(M\) 的一個基，從而也補成 \(N\) 的一個基）。因此，\(M\) 中每個 \(P\)-模
\[
        T=z_1P+\ldots+z_sP
\]
都生成一個擴張模
\[
        T_A=z_1A+\ldots+z_sA\subset M_A.
\]

\paragraph{推論 b)。}
\(M\) 中每個 \(P\)-模 \(T=z_1P+\ldots+z_sP\) 都是限制模，也就是其擴張模 \(T_{\mathsf{A}}\) 與 \(M\) 的交：
\[
        T=T_A\cap M .
\]
因為 \(T\subseteq T_A\cap M\)；反過來，若 \(a\) 位於 \(T_A\cap M\)，則有
\[
        a=\sum z_i\alpha_i,
\]
於是 \(M\) 中的元素 \(a,z_1,\ldots,z_t\) 在 \(\mathsf{A}\) 上相關，因此按推論 a) 也在 \(\mathsf{P}\) 上相關；所以 \(a\) 位於 \(T\) 中。

\paragraph{3. 關於不變模的定理。}
仍設 \(N=M_A\) 是秩 \(n\) 的 \(P\)-模 \(M\) 的擴張模；再設 \(P\) 是對於 \(A\) 的一群環自同構 \(\frG\) 的全不變量域。群 \(\frG\) 作為 \(N=M_A\) 的算子域，由如下規定定義：
\[
        Gm=m\quad\text{對來自 }M\text{ 的 }m\text{ 以及來自 }\frG\text{ 的 }G
\]
並且
\[
        G\sum m_i\alpha_i=\sum m_i\,G(\alpha_i)
        \quad\text{對來自 }M\text{ 的 }m_i\text{ 以及來自 }A\text{ 的 }\alpha_i.
\]

\paragraph{引理。}
在這些規定下，\(M\) 由 \(N\) 中在 \(\frG\) 下保持不動的全部元素組成。事實上，若 \(x_1,\ldots,x_n\) 是 \(M\) 的一個 \(P\)-基，且 \(w=x_1\alpha_1+\ldots+x_n\alpha_n\)，那麼對於每個來自 \(\frG\) 的 \(G\)，由 \(G(w)=w\) 可推出也有 \(G(\alpha_i)=\alpha_i\)；於是按假設 \(\alpha_i\) 位於 \(\mathsf{P}\) 中。

\paragraph{定理。}
模 \(M\) 的子模 \(L\) 的擴張模 \(L_A\)，並且只有這些擴張模，是相對於作為算子域的 \(\frG\) 的容許子模。

前一半是清楚的。反過來，設 \(L\) 相對於 \(\frG\) 是容許的，並設 \(z_1,\ldots,z_l\) 是 \(L\) 相對於 \(x_1,\ldots,x_n\) 的一個正規基（見 1.）；這裡 \(x_1,\ldots,x_n\) 被選為 \(M\) 的 \(P\)-基，所以 \(\frG\) 逐個保持它們不動。按假設
\[
        G(z_i)=x_i-\bigl(x_{l+1}G(\alpha_{i,l+1})
              +\ldots+x_nG(\alpha_{i,n})\bigr)
\]
對於每個來自 \(\frG\) 的 \(G\) 都是 \(L\) 的元素；於是，當 \(G\) 固定時，它可由這些 \(z\) 線性表示。因此，如同比較 \(x_1,\ldots,x_l\) 的係數所顯示的，對每個來自 \(\frG\) 的 \(G\) 都有 \(G(z_i)=z_i\)。於是，按引理，這些 \(z\) 位於 \(M\) 中，而 \(L\) 成為 \(M\) 中 \(P\)-模
\[
        T=z_1P+\ldots+z_\ell P
\]
的擴張，並且 \(T=L\cap M\)（按 2）。\footnote{藉助簡單的良序論證，本段各定理即使在 \(M\) 關於 \(\mathsf{P}\) 的秩變為無限時仍然成立；參見 G. Köthe, \emph{Ein Beitrag zur Theorie der kommutativen Ringe ohne Endlichkeitsvoraussetzung} § 1, Gött. Nachr. 1931, S. 195--207。}

\subsection*{§4. 超複系統及其表示類}

\paragraph{1. 擴張定理。}
現在把 §2 的表示定理同 §3 的模定理結合起來。我們不再考察一般的環 \(\frR\)，而考察帶單位元的超複系統 \(S,T\ldots\)，它們定義在交換域 \(\mathsf{P}\) 上，也就是
\[
        S=x_1\mathsf{P}+\ldots+x_n\mathsf{P};
\]
並且，不再取任意的表示環 \(\frS\)，而只取域 \(\mathsf{A},\mathsf{B}\ldots\)，如果沒有另行說明，就取中心為 \(\mathsf{P}\) 的域。在這種情形下，總存在一個乘積環，使 \(S\) 與 \(\mathsf{A}\) 元素兩兩可交換，並且交為 \(\mathsf{P}\)；更確切地說，這裡涉及的是直積 \(S\times \mathsf{A}\)，它同時成為 §3 意義下 \(S\) 的擴張模 \(S_{\mathsf{A}}\)，
\[
        S_{\mathsf{A}}=x_1\mathsf{A}+\ldots+x_n\mathsf{A},
\]
因此也稱為擴張環。

在這種特化下，不變模定理（§3, 3.）變為如下

\paragraph{擴張定理。}
\(S_{\mathsf{A}}\) 中的每個雙側理想，都是 \(S\) 中某個理想的擴張理想；確切地說，是它同 \(S\) 的交的擴張。

事實上，若所取的群 \(\frG\) 是 \(A\) 的內自同構群，則 \(P\) 作為 \(A\) 的中心成為全不變量域；同時，\(A\) 與 \(S\) 的元素兩兩可交換，說明 \(\frG\) 按 §3, 3. 擴張到 \(S_{\mathsf{A}}\) 上，使 \(S\) 成為全不變量範圍。\(S_{\mathsf{A}}\) 中每個雙側理想 \(\mathfrak a\) 都是容許子群，因為 \(\alpha^{-1}\mathfrak a\alpha\leqq\mathfrak a\)；所以按 §3, 3.，它是其交模 \(\mathfrak a\cap S\) 的擴張，而這個交模成為 \(S\) 中的理想。

附加說明：若 \(S\) 是交換的，則 \(S\) 成為 \(S_{\mathsf{A}}\) 的中心。因為 \(S\) 位於中心中，並且相對於由 \(\mathsf{A}\) 生成的內自同構，是全不變量範圍。更一般地，\(S\) 的中心也成為 \(S_{\mathsf{A}}\) 的中心。

記號：\(\mathsf{A}_r,\mathsf{B}_n\ldots\) 以下始終表示在 \(\mathsf{A},\mathsf{B}\ldots\) 上指定階數的矩陣環，也就是
\[
        A_r=\sum A c_{ik}=\sum c_{ik}A.
\]

下面將本質地使用一個推論，即簡單系統的擴張定理：若 \(S\) 是簡單系統\footnote{簡單系統照通常含義，是指“雙側簡單”。}，則 \(S_{\mathsf{A}}\) 也簡單：
\[
        S_A=\sum B c_{ik}=\sum c_{ik}B=B_t
\]
其中 \(\mathsf{B}\) 是相應的域，其中心包含 \(\mathsf{P}\)。這裡 \(\mathsf{A}\)，因而 \(\mathsf{B}\)，可以在 \(\mathsf{P}\) 上有無限秩；只假定 \(S\) 是超複的。若 \(S\) 與 \(\mathsf{A}\) 有共同中心 \(\mathsf{P}\)，則 \(\mathsf{P}\) 也成為 \(S_{\mathsf{A}}\) 的中心。

更一般地還成立：若 \(S\) 是簡單超複系統，則直積 \(S\times \mathsf{A}_r\) 也簡單。因為 \(S\times \mathsf{A}_r\) 與 \(S_{\mathsf{A}}\times \mathsf{P}_r\) 一致；若 \(S_{\mathsf{A}}=\mathsf{B}_t\)，它就等於 \(\mathsf{B}_{tr}\)。

\paragraph{2. 表示類。}
所謂超複系統的互反表示或直接表示，照通常理解，是指關於係數域 \(\mathsf{P}\) 算子同態（§2 末尾）並且不同於零表示的表示。把過渡定理及其推論（§2, 3.）同 1. 中的擴張定理結合，得到如下

\paragraph{關於表示類的定理。}
設 \(S\) 是帶單位元的超複系統，以 \(\mathsf{P}\) 為係數域，且 \(\mathsf{A}\) 是中心為 \(\mathsf{P}\) 的域；那麼 \(S\) 在 \(\mathsf{A}\) 中有多少個不同的不可約互反表示類，正等於 \(S_{\mathsf{A}}\) 按其根基取的商環中有多少個簡單右理想類。-- 若 \(S\) 是簡單系統，則它在 \(\mathsf{A}\) 中恰有一個不可約互反表示類和一個不可約直接表示類。\footnote{參見“表示論”注 20），其中處理的是 \(\mathsf{A}\) 為 \(S\) 的相應自同構域的情形。}

這是因為，按過渡定理的推論（§2, 3.），簡單互反表示類同 \(S_{\mathsf{A}}\) 上的簡單模類一一對應。由於 \(S_{\mathsf{A}}\) 關於域 \(\mathsf{A}\) 具有有限秩，它滿足單側理想的極大條件和極小條件；於是斷言的第一部分直接來自“表示論”§19。若 \(S\) 簡單，則按擴張定理 \(S_{\mathsf{A}}\) 也簡單。由於極大條件和極小條件，它同時完全可約，因此只有一個簡單單側理想類；也就是說，\(S\) 在 \(\mathsf{A}\) 中只有一個不可約互反表示類。與 \(S\) 互反同構的系統也隨 \(S\) 一起是簡單的，並且只有一個不可約互反表示類；因此 \(S\) 也只有一個不可約直接表示類。

\paragraph{3. 秩關係。}
當 \(S\) 簡單時，
\[
        S_A=\sum Bc_{ik}
\]
無論關於 \(A\) 還是關於 \(B\) 都具有有限秩；這一事實給出秩關係：
\[
        n=rt\quad\text{其中}\quad n=(S:P),\quad t^2=(S_A:B),
        \quad r=\text{不可約互反表示的階}.                              \tag{1}
\]

因為 \(S_{\mathsf{A}}\) 本身就是 \(S\) 在 \(\mathsf{A}\) 中的互反表示模（該表示已經在 \(\mathsf{P}\) 中）。作為 \(\mathsf{B}\) 上的矩陣環，\(S_{\mathsf{A}}\) 分解為 \(t\) 個簡單右理想，也就是 \(S_{\mathsf{A}}\)-模；這些模關於 \(\mathsf{A}\) 的秩等於不可約互反表示的階 \(r\)。由於 \(n\) 同時也是 \(S_{\mathsf{A}}\) 關於 \(\mathsf{A}\) 的秩，秩關係隨之得出。

\paragraph{4. 當 \(\mathsf{A}\) 關於 \(\mathsf{P}\) 有有限秩時，秩關係的加強。}
在這種情形下有三個關係：
\[
        (S:P)=n=rt,                                                     \tag{1}
\]
\[
        (B:P)\,t=(A:P)\,r,                                               \tag{2}
\]
\[
        (S:P)(B:P)=(A:P)\,r^2.                                           \tag{3}
\]
其中 (2) 表示 \(S_{\mathsf{A}}\) 中簡單右理想關於 \(P\) 的秩；按 3.，它由關於 \(\mathsf{B}\) 的秩或關於 \(\mathsf{A}\) 的秩來表示。而 (3) 表示矩陣環 \(\mathsf{B}_n\) 中一個簡單右理想的秩，並且由 (2) 按照 (1) 乘以 \(r\) 得到。

\subsection*{§5. 表示定理在矩陣環上的應用}

§4 中考察的 \(S\) 在 \(\mathsf{A}\) 中的互反表示，也可以理解為 \(S\) 到 \(A_r\) 的某個子環中的直接同態；這裡 \(A\) 表示與 \(\mathsf{A}\) 互反同構的域。以下各段的原則反過來是：把矩陣環 \(A_r\) 的研究歸結為在互反同構域 \(\mathsf{A}\) 中的表示論。

\paragraph{1. 在 \(A_r\) 中的可約與不可約嵌入。}
設 \(\mathsf{A}\) 與 \(A\) 是兩個在其中心 \(\mathsf{P}\) 上具有有限或無限秩的互反同構域。\footnote{一般說來，互反的域或系統將用相應的希臘字母和拉丁字母表示。}若 \(P\) 上的簡單系統 \(S\) 在 \(A\) 中具有 \(r\) 階的不可約或可約互反表示，則稱 \(S\) 在 \(A_r\) 中可作不可約嵌入，或者可作可約嵌入。若 \(S\) 在 \(A_r\) 中可作不可約嵌入，則它可約地嵌入所有矩陣環 \(\mathsf{A}_{rs}\)（\(s>1\)）中，因為它恰有這些階數且只有這些階數的可約表示（§4, 2.）。因此，\(S\) 直接同構於 \(A_r\) 和 \(A_{rs}\) 的某些子環。該同構是 \(\mathsf{P}\) 上恆等映射的延拓；\footnote{如同在表示中一樣，下文出現的同構始終是延拓 \(\mathsf{P}\) 上恆等映射的同構，即按 \(\mathsf{P}\) 的同構；以後不再每次另行說明。}按 §4, 3.，這裡 \(r\) 是 \(S\) 關於 \(P\) 的秩 \(n\) 的一個因子。

用這種方式，可以得到各個 \(A_f\)（\(f=1,2\ldots\)）中所有包含 \(P\)、並且按 \(P\) 具有有限秩的簡單子環；因為每個這樣的子環都與 \(\mathsf{A}_f\) 的某個子環互反同構，從而容許在 \(\mathsf{A}\) 中有一個可約或不可約互反表示。

\paragraph{2. 關於內自同構的定理。}
若 \(S^{(1)}\) 與 \(S^{(2)}\) 是 \(A_f\) 中兩個包含 \(P\)、並且按 \(\mathsf{P}\) 具有有限秩的簡單子環，且 \(S^{(1)}\) 與 \(S^{(2)}\) 按 \(\mathsf{P}\) 同構，則這個同構由 \(A_f\) 的一個內自同構產生。

因為 \(S^{(1)}\) 與 \(S^{(2)}\) 表示同一個簡單系統 \(S\) 在 \(A_f\) 中的嵌入（一般是可約的），從而表示在 \(A\) 中的 \(f\) 階直接表示；所以它們屬於 \(A\) 中同一個可約直接表示類；完成過渡的矩陣產生所需的自同構。

若 \(A\) 本身按 \(P\) 具有有限秩，即為超複的，則該定理化為熟知命題：\(A_f\) 中兩個包含中心 \(P\)、且按 \(\mathsf{P}\) 同構的簡單子系統，通過 \(A_f\) 的一個內自同構互相過渡。特別地，\(A_f\) 的每個自同構都是內自同構。\footnote{當 \(A\) 按 \(\mathsf{P}\) 具有無限秩時，這一點不再成立（參見注 5 所引 Köthe 的工作）。}

\paragraph{3. 關於元素逐個可交換子環的定理。}
這個定理仍按本段開頭提到的原則，由 §2, 3. 中關於可交換矩陣的定理推出：

若 \(S\) 是來自 \(A_f\) 的一個包含 \(P\) 的簡單系統，\(R\) 表示 \(A_f\) 中同 \(S\) 的元素逐個可交換的全體元素，則 \(R\) 也成為矩陣環：\(R=B_s\)。\(S_{\mathsf{A}}\) 的對應域 \(B\) 與 \(R\) 的對應域 \(\overline B\) 互反同構。\(R\) 與 \(S\) 的交成為 \(S\) 的中心。當且僅當 \(S\) 不可約地嵌入 \(A_f\) 時，\(R\) 成為域。

事實上，按 §2, 3.，\(R\) 直接同構於 \(S\) 在 \(A\) 中的互反表示模 \(\frM\) 的自同構環；但這正是把 \(\frM\) 看作 \(S_{\mathsf{A}}\)-模時的自同構環。若 \(\frM\) 分解為 \(s\) 個簡單 \(S_{\mathsf{A}}\)-模，並且如 §4, 1. 那樣寫成 \(S_{\mathsf{A}}=\sum \mathsf{B}c_{ik}\)，則 \(R\) 同構於 \(\sum Bc_{ik}\)，其中 \(\mathsf{B}\) 與 \(B\) 互反同構（因為按 §1, 3. 的結尾，\(\mathsf{B}\) 與 \(S_{\mathsf{A}}\) 中簡單右理想，即簡單 \(S_{\mathsf{A}}\)-模的自同構域，按 \(P\) 互反同構）。當且僅當 \(s\) 等於一，也就是 \(S\) 被不可約嵌入時，\(R\) 成為域，即 \(R=B\)。\(R\) 與 \(S\) 的交就是 \(S\) 的中心，這一點直接來自 \(R\) 的定義。

\paragraph{4. 超複矩陣環情形下的加強交換定理。}
在這種情形中，§4, 4. 的秩關係給出如下加強：

來自 \(A_f\) 的、包含 \(P\) 的簡單子系統（其中 \(A\) 按其中心 \(\mathsf{P}\) 具有有限秩）分解為成對的 \(S,\overline S\)，使得 \(\overline S\) 由同 \(S\) 的元素逐個可交換的全體元素組成，反之亦然。\(S_{\mathsf{A}}\) 的對應域與 \(\overline S\) 的對應域互反同構，\(S\) 的對應域與 \(\overline S_{\mathsf{A}}\) 的對應域也互反同構。\(S\) 與 \(\overline S\) 的交等於共同中心。當且僅當一個子環被不可約嵌入時，另一個子環成為域。\(S\) 與 \(\overline S\) 按 \(P\) 的秩數之積，等於 \(A_f\) 的秩。若
\[
        f=rs=\overline r\,\overline s
\]
其中 \(S\) 在 \(A_r\) 中可作不可約嵌入，而 \(\overline S\) 在 \(A_{\overline r}\) 中可作不可約嵌入，則 \(S\) 與 \(\overline S\) 的對應域同構於 \(A_g\) 中一對可交換子環，其中 \(f=g s\overline s\)。

本質上只需證明關於秩數之積的斷言。先設 \(S\) 被不可約嵌入；於是 \(\overline S\) 是域，並且與 \(\mathsf{B}\) 互反同構，其中 \(S_{\mathsf{A}}=\mathsf{B}_t\)。這時 §4, 4. 的秩關係 (3) 給出乘積斷言，因為現在 \(f=r\)，並且 \(\overline S\) 的秩等於 \(\mathsf{B}\) 的秩。若 \(S\) 被可約嵌入，即 \(\overline S=B_s\)，則 \(f=rs\)，把 (3) 乘以 \(s^2\) 就得到該斷言。由此立刻推出，\(\overline S\) 由同 \(S\) 的元素逐個可交換的全體元素組成；於是除最後一個斷言外，其餘全部都由 3. 得出。最後一個斷言則由兩次過渡到不可約嵌入得到。在 \(A_r\) 中（其中 \(f=rs\)），\(S\) 被不可約嵌入，\(S\) 與 \(R\) 成為成對可交換系統，而 \(R\) 同構於域 \(B\)。另一方面，由於 \(S\) 與 \(\overline S\) 的互反性，\(S\) 等於 \(\overline B_{\overline s}\)，這裡 \(\overline B\) 對於 \(\overline S\) 的意義，正如 \(B\) 對於 \(S\) 的意義；因此 \(R\) 可約地嵌入 \(A_r\)，並且不可約地嵌入 \(A_g\)，其中 \(r=g\overline s\)。在這裡，那對逐個可交換的系統便同構於 \(B\) 與 \(\overline B\)。

\subsection*{§6. 簡單系統的伽羅瓦理論}

§5, 4. 的加強交換定理，表達了以中心作為基礎域時簡單系統的伽羅瓦理論。我們先考察域，然後考察一般簡單系統。

\paragraph{1. 按中心具有有限秩的域的伽羅瓦理論。}
\(A\) 的伽羅瓦群 \(\Gcal\) 定義為一切延拓中心 \(P\) 上恆等映射的自同構所成的群；也就是說，按 §5, 2.，它是所有內自同構所成的群。因此有
\[
        \Gcal \simeq A^*/P^*,
\]
其中 \(A^*\) 和 \(\mathsf{P}^*\) 分別表示 \(A\) 和 \(\mathsf{P}\) 中非零元素的乘法群。於是，\(\Gcal\) 的子群 \(\mathfrak H\) 同 \(A^*\) 中包含 \(\mathsf{P}^*\) 的子群 \(H^*\) 一一對應，關係為
\[
        \mathfrak H \sim H^*/P^* .
\]
\(\Gcal\) 的一個子群 \(\mathfrak H\) 稱為閉的，如果 \(H^*\) 添入零以後成為一個域 \(H\)。\(C\) 容許群 \(\mathfrak H\) 這一事實，等價於 \(C\) 同 \(H\) 元素逐個可交換。由此，§5, 4. 的交換定理化為如下

\paragraph{非交換域伽羅瓦理論的主定理。}
位於 \(\mathsf{P}\) 與 \(A\) 之間的域 \(C\)，同 \(\Gcal\) 的閉子群 \(\mathfrak H\)，可以這樣一一對應：\(C\) 成為 \(\mathfrak H\) 的全不變量域，而 \(\mathfrak H\) 成為 \(C\) 的全不變量群。

\paragraph{2. 簡單系統的伽羅瓦理論。}
若 \(A_f\) 是中心為 \(\mathsf{P}\) 的簡單系統，則伽羅瓦群 \(\Gcal\) 仍為內自同構群，因此同
\[
        A_f^*/P^*
\]
同構；這裡 \(A_f^*\) 表示 \(A_f\) 中正則元素的乘法群。\(\Gcal\) 的子群 \(\mathfrak H\simeq H^*/\mathsf{P}^*\) 稱為簡單閉的，如果由 \(H^*\) 生成的子環 \(H\) 是簡單系統，並且 \(H^*\) 由 \(H\) 中全體正則元素組成。這裡從交換定理過渡到伽羅瓦理論，依靠如下

\paragraph{輔助定理。}
\(A_f\) 中每個包含 \(P\) 的簡單系統，都由其正則元素所成的群 \(H^*\) 生成。若 \(\mathsf{P}\) 含有無限多個元素，這是清楚的；因為用不定元形成的“通有元素”是正則的，並且通過適當特化，可以在 \(\mathsf{P}\) 上形成正則的基元素。不過按 Shoda，此輔助定理一般成立。\footnote{Shoda，注 3）所引的工作。那裡避免了不定元的引入。}

藉助輔助定理，\(S\) 同一個簡單系統 \(\mathfrak S\) 可交換，又等價於 \(S\) 容許由 \(\mathfrak S\) 的正則元素誘導的自同構。因此，§5, 4. 的交換定理在這裡也化為如下

\paragraph{簡單系統的伽羅瓦理論主定理。}
\(A_f\) 中包含 \(\mathsf{P}\) 的簡單系統 \(S\)，同 \(\Gcal\) 的簡單閉子群 \(\mathfrak H\)，可以這樣一一對應：\(S\) 成為 \(\mathfrak H\) 的全不變量範圍，而 \(\mathfrak H\) 成為 \(S\) 的全不變量群。

\paragraph{3. 藉助延拓原則的證明。}
我還要在域的情形下給出主定理的第二個證明；這個證明基於同構，也就是一階表示，的延拓原則，並且由於不使用秩數計數，需要較少的有限性假設。特別地，這樣可以看出，向無限秩的域 \(S\) 轉移將可能沿着什麼方向進行。這個證明也不使用“只有一個不可約互反表示類”這一事實；因此它在交換情形中同樣成立，參見 §8, 2. 我先列出若干輔助定理。

\paragraph{前提。}
設 \(P\) 上的簡單系統 \(S\) 容許在 \(\mathsf{A}\) 中有一個一階互反表示，因而是域；這裡 \(\mathsf{A}\) 關於其中心 \(\mathsf{P}\) 可以具有有限秩或無限秩。按 §4, 3.，給定 \(S_{\mathsf{A}}\) 分解為 \(n\) 個簡單且算子同構的右理想：
\[
        S_{\mathsf{A}}=\mathfrak r_1+\ldots+\mathfrak r_n
            =e_1S_{\mathsf{A}}+\ldots+e_nS_{\mathsf{A}}
            =e_1\mathsf{A}+\ldots+e_n\mathsf{A} ;
\]
於是，由 \(e_i\) 生成的互反表示定義為
\[
        e_i s=e_i\sigma .
\]

\paragraph{輔助定理 1。}
由 \(e_1,\ldots,e_n\) 生成的這 \(n\) 個互反表示彼此不同，因而是同一個表示類中的不同表示。因此，\(S\) 至少擁有同它的秩一樣多的不同表示。

證明時，我們按 §2 末尾，把來自 \(S\) 的表示擴張為 \(S_{\mathsf{A}}\) 的對應，使之按 \(\mathsf{A}\) 算子同態。這裡這些對應由
\[
        e_i w=e_i\omega
\]
定義，其中 \(\omega\in A\)，而 \(w\) 是 \(S_{\mathsf{A}}\) 的任意元素。若 \(e_i\) 和 \(e_j\)（\((i\ne j)\)）生成同一個表示，則對 \(S_{\mathsf{A}}\) 中每個 \(w\) 都應有 \(e_iw=e_jw\)。但這同 \(e_ie_i=e_i\) 以及 \(e_je_i=0\) 矛盾。

\paragraph{輔助定理 2。}
若 \(T\) 是 \(\mathsf{P}\) 與 \(S\) 之間的域，且 \(s\) 是 \(S\) 關於 \(T\) 的秩，則 \(T\) 到 \(\mathsf{A}\) 中的每個互反同構至少容許 \(s\) 個不同延拓。

設這個互反同構，也就是 \(T\) 在 \(\mathsf{A}\) 中的互反表示，由一個分解給出：
\[
        T_{\mathsf{A}}=E_1T_{\mathsf{A}}+\ldots+E_hT_{\mathsf{A}}
            =E_1A+\ldots+E_hA,
        \qquad E_it=E_i\tau ,
\]
其中 \(E_i\) 是冪等元。這並不限制一般性，因為一個由基元素 \(E_i\alpha\) 生成的表示，也可以由 \(\alpha^{-1}\cdot E_i\alpha\)，也就是一個冪等元，生成。用 \(S\) 作右乘，得到
\[
        S_A=E_1S_A+\ldots+E_hS_A .
\]
這裡各個 \(E_iS_{\mathsf{A}}\) 關於 \(\mathsf{A}\) 的秩全都等於 \(s\)。事實上，如藉助 \(S\) 與 \(\mathsf{A}\) 的可交換性把 \(S\) 的一個 \(T\)-基代入所顯示的，\(E_iS_{\mathsf{A}}\) 的秩至多為 \(s\)：
\[
        S=Tw_1+\ldots+Tw_s,
\]
\[
        E_iS_{\mathsf{A}}=E_iT_{\mathsf{A}}w_1+\ldots+E_iT_{\mathsf{A}}w_s
              =E_iw_1\mathsf{A}+\ldots+E_iw_s\mathsf{A} .
\]
由於總和 \(S_{\mathsf{A}}\) 的秩為 \(n=hs\)，每個 \(E_iS_{\mathsf{A}}\) 的秩恰為 \(s\)。於是有
\[
        E_iS_A=r_{i1}+\ldots+r_{is}
              =e_{i1}A+\ldots+e_{is}A,
\]
其中，由 \(e_{i1},\ldots,e_{is}\) 生成的這 \(s\) 個表示，按輔助定理 1 全都不同。它們全都是由 \(E_i\) 生成的 \(T\) 的表示的延拓；因為由 \(E_it=E_i\tau\) 可得
\[
        (e_{i1}+\ldots+e_{is})t=(e_{i1}+\ldots+e_{is})\tau
\]
從而由於直和分解，有 \(e_{il}t=e_{il}\tau\)。

\paragraph{定義。}
由 \(T\) 誘導的、\(S\) 到 \(\mathsf{A}\) 的互反同構 \(\mathfrak I\) 的類劃分 \(\mathfrak I_T\)，定義如下：\(\mathfrak I\) 中全部且僅有那些在 \(T\) 上誘導同一同構的同構，被看作等價的。

\paragraph{主定理。}
若 \(\mathfrak I_T\) 是由 \(T\) 誘導的類劃分，則 \(T\) 是相對於這個類劃分的極大子域。

事實上，若 \(Z\) 是位於 \(T\) 與 \(S\) 之間的中間域，並且相對於 \(T\) 的次數為 \(\ell\)，則按輔助定理 2，\(\mathfrak I_T\) 中的每個類在 \(\mathfrak I_Z\) 中至少分裂為 \(\ell\) 個類。從主定理的這個表述過渡到通常表述，是通過把同構集合 \(\mathfrak I\) 的各元素同某個固定元素的逆相乘，使之進入自同構群；這時 \(\mathfrak I_T\) 的一個類轉入 \(T\) 的不變量群。閉子群是全不變量群這一斷言也包含在其中；對 \(\mathfrak H\simeq H^*/\mathsf{P}^*\)，只要把域 \(H\) 作為中間域 \(T\) 代入即可。

\subsection*{§7. 相似代數的類。分裂域}

從現在起，只討論按其中心 \(P\) 具有有限秩的域 \(\mathsf{A},A\) 以及矩陣環 \(A_f\)；也就是說，只討論 \(\mathsf{P}\) 上的簡單正規代數，以下簡稱為 \(\mathsf{P}\) 上的代數，或者僅稱代數。於是 \(\mathsf{A},A\) 是除代數。一個相似代數類 \(A,A_f\) 中，彙集所有具有同一對應 \(A\) 的 \(A_f\)。同構的 \(A\) 在這裡加以認同。

\paragraph{1. 代數類群。}
若 \(A,B\) 是 \(P\) 上的代數，則它們的乘積，也就是按 \(P\) 的直積，仍是同樣的對象；因為按 §3 末尾有
\[
        A_r\times_P B_s=C_t,
\]
其中 \(C\) 是一個除代數，中心為 \(\mathsf{P}\)，並且在同構意義下唯一確定。把 \(A_r,B_s\) 乘以 \(\mathsf{P}\) 上的矩陣環可見，這種乘法對各個類是唯一確定的；因此這些類構成一個乘法封閉的交換系統。不僅如此，還成立如下

\paragraph{R. Brauer 定理。}
相似代數的類在直積構成下形成一個阿貝爾群。

事實上，該系統有一個單位類，由 \(\mathsf{P}\) 上的矩陣環，也就是同 \(\mathsf{P}\) 相似的代數組成。並且，每個類 \((A)\) 都有逆類 \((A)^{-1}\)，由同 \(\mathsf{A}\) 相似的代數組成；這裡 \(\mathsf{A}\) 與 \(A\) 互反同構。這由 §5, 3. 的交換定理在特化 \(S=A\) 時推出。因為 \(A\) 在自身中可作不可約嵌入；於是得到 \(B=\mathsf{P}\) 和 \(A_{\mathsf{A}}=\mathsf{P}_t\)。\footnote{關於代數類群的更精細定理使用因子系統理論；此處不進入這一點。應當指出，到這裡為止沒有出現關於交換域的考察；例如，秩 \((A:\mathsf{P})\) 是平方數這一事實既未使用也未證明。}

\paragraph{2. 一個代數類的分裂域。}
分裂域理論仍然完全基於 §5 開頭所說的原則；交換擴張域在 \(\mathsf{A}\) 或 \(A\) 中表示。先給出如下

\paragraph{輔助定理。}
若 \(A\) 是中心為 \(\mathsf{P}\) 的域，\(\mathsf{Z}\) 是 \(\mathsf{P}\) 的有限交換擴張域，則 \(A_{\mathsf{Z}}\) 是中心為 \(\mathsf{Z}\) 的簡單系統。因為這按 §4, 1. 對同 \(A_{\mathsf{Z}}\) 互反同構的系統 \(Z_{\overline A}\) 成立。因此，對於 \(P\) 上同 \(A\) 相似的每個代數類 \((A)\)，\(Z\) 產生一個擴張類 \((A)_{\mathsf{Z}}\)，由 \(Z\) 上同 \(A_{\mathsf{Z}}\) 相似的簡單代數組成。

一個擴張類的對應除代數 \(D\)，可直接從 §5, 3. 的交換定理得到：

\paragraph{關於擴張類的對應除代數的定理。}
設 \((A)_{\mathsf{Z}}\) 是擴張類，且 \(Z\) 表示 \(\mathsf{Z}\) 在 \(A_r\) 中的一個不可約嵌入，也就是表示。則
\[
        (A)_Z=(D),
\]
其中 \(D\) 由 \(A_f\) 中同 \(Z\) 的元素逐個可交換的全體元素組成。只需在 §5, 3. 中把那裡的簡單系統 \(S\) 替換為 \(Z\)，同時注意 \(Z_{\mathsf{A}}\) 與 \(A_{\mathsf{Z}}\) 互反同構。若涉及的是 \(\mathsf{Z}\) 的可約嵌入，則同 \(Z\) 的元素逐個可交換的全體元素給出 \(D\) 上的一個矩陣環。

眾所周知，\(P\) 的交換擴張域 \(Z\) 稱為類 \((A)\) 的分裂域，如果擴張類 \((A)_{\mathsf{Z}}\) 完全分裂，也就是說，等於 \(\mathsf{Z}\) 上的單位類，即 \(\mathsf{Z}\) 上所有矩陣環的類。

\paragraph{關於分裂域刻畫的定理。}
\(P\) 的一個有限交換擴張域 \(Z\)，當且僅當它在 \(A_f\) 中的不可約嵌入給出 \(A_f\) 的一個極大交換子域時，是類 \((A)\) 的分裂域。

因為 \(\mathsf{Z}\) 是 \((A)\) 的分裂域這一性質，等價於 \((D)=(\mathsf{Z})\)。因此，按關於對應除代數的定理，不可約嵌入 \(Z\) 必須是極大交換子域。若這一條件滿足，則必有 \(D=Z\)，因為 \(D\) 中每個 \(a\) 添接到 \(Z\) 上都會生成一個交換擴張域。

\paragraph{說明。}
1. 同時可得：一個極大交換子域 \(Z\) 在不可約嵌入時生成一個極大交換子環。因為同 \(Z\) 的元素逐個可交換的全體元素形成一個域。

2. 該定理同時給出分裂域的存在性；因為對應除代數 \(A\) 必定具有極大交換子域。

\paragraph{3. 秩關係、指數、無限交換擴張。}
§5, 4. 的秩斷言首先給出熟知事實：一個簡單代數關於其中心的秩是平方數。事實上，若 \(Z\) 在 \(A\) 中極大交換，則由於在 \(A\) 中每個嵌入都是不可約的，有
\[
        (Z:P)\cdot(Z:P)=(A:P),
\]
因而
\[
        (A:P)=m^2,\qquad (Z:P)=m,
\]
其中 \(m\) 是 Schur 指數；它同時表示所有可嵌入除代數 \(A\) 本身的分裂域的次數。進一步，對一般分裂域的次數 \(n\)，有
\[
        n=m\cdot r,
\]
這可由 \((A_r:\mathsf{P})=m^2\cdot r^2\) 得到，也可由 §4, 3. 的秩關係得到；後者是因為指數 \(m\) 也等於 \(A_{\mathsf{Z}}\) 的簡單分量數，而 \(A_{\mathsf{Z}}\) 成為 \(\mathsf{Z}\) 上秩為 \(m^2\) 的矩陣環。更一般地，若 \(A_\Lambda\sim D\)，\(L\) 是 \(A\) 在 \(A_r\) 中的不可約嵌入，且 \(d^2\) 是除代數 \(D\) 關於其中心 \(L\) 的秩 \((D:L)\)，則有
\[
        l\cdot\,d=m\cdot r .
\]
因為按 §5, 4. 以及 2. 中關於對應除代數的定理，有
\[
        (A_r:P)=(L:P)\cdot(D:P),
\]
所以
\[
        m^2r^2=l^2d^2 .
\]

由分裂域的存在性，對 \(P\) 的無限交換擴張域 \(\Omega\)，無論代數的還是超越的，可以推出如下事實：擴張類 \((A)_\Omega\) 是中心為 \(\Omega\) 的簡單代數類。特別是，若 \(\Omega\) 是代數閉的，則 \(A_\Omega\) 成為次數為 \(m\) 的矩陣環。因此，指數等於“絕對分量數”。事實上，若 \(\mathsf{Z}\) 是 \(A\) 的分裂域，\(\bar\Omega\) 是 \(\Omega\) 與 \(Z\) 的合成域，則 \(A_{\Omega'}\) 是 \(\Omega'\) 上的矩陣環，因而沒有根基，並且中心為 \(\Omega'\)。於是 \(A_\Omega\) 也沒有根基且中心為 \(\Omega\)，因為 \(A_\Omega\) 中不同於 \(\Omega\) 的中心元素，也會位於 \(A_{\Omega'}\) 的中心，因而成為域。於是 \(A_\Omega\) 是 \(\Omega\) 上的簡單代數；當 \(\Omega\) 代數閉時，它必定是 \(\Omega\) 上的矩陣環。\footnote{一般說來，也存在各種次數的“極小”分裂域，也就是說，沒有真子域仍為分裂域的分裂域。參見 Brauer--Noether，注 2）所引的工作。}

\paragraph{4. 可分分裂域的存在性。}
每個類 \((A)\) 都具有可分分裂域，甚至具有可嵌入 \(A\) 本身的可分分裂域。\footnote{參見 G. Köthe, Über Schiefkörper mit Unterkörpern zweiter Art über dem Zentrum, Journ. f. Math. 166 (1932), S. 182--184。這裡給出的簡單得多的證明來自 M. Zorn 的一個觀察。}

設基礎域 \(\mathsf{P}\) 的特徵為 \(p\)，指數為
\[
        m=s p^f
\]
其中 \(s\) 與 \(p\) 互素。若 \(\mathsf{Z}\) 是 \(A\) 的一個次數為 \(m\) 的分裂域，且 \(\mathsf{Z}_0\) 是其中所含的第一類擴張，也就是
\[
        (Z:Z_0)=p^f,
\]
則 \(A_{\mathsf{Z}_0}\sim D\) 的指數等於 \(p^f\)。現在要證明存在 \(\mathsf{Z}_0\) 的一個可分擴張域 \(\mathsf{Z}_1\)，使得 \(D_{\mathsf{Z}_1}\) 的指數成為 \(p^f\) 的真因子。由於當 \(\Omega\) 代數閉時，\(D_\Omega\) 按 3. 沒有根基，所以 \(D\) 的約化判別式不為零。於是至少存在一個 \(D\) 中的元素 \(d\)，其約化跡不為零；此時 \(d\) 不能已經位於基礎域 \(\mathsf{Z}_0\) 中，因為 \(\mathsf{Z}_0\) 中每個元素 \(\alpha\) 的跡等於 \(p^f\alpha\)，因而為零。因此 \(\mathsf{Z}_1=\mathsf{Z}_0(d)\) 是真擴張，而且是可分擴張；\(D_{\mathsf{Z}_1}\) 的指數成為 \(p^f\) 的真因子。有限次重複這一過程，便得到類 \((A)\) 的一個可分分裂域，其次數 \(m\) 同 \(A\) 的指數一致。

\subsection*{§8. 交換情形中的分裂域、分解域與伽羅瓦理論}

為了從關於中心上的簡單系統的分裂域，過渡到任意簡單系統的分裂域，還必須發展中心的分裂理論；並且還要添入一種同 §6, 3. 平行的伽羅瓦理論。本段中出現的一切域和系統都是交換的。

\paragraph{1. 交換簡單系統的分裂域與分解域。}
設交換系統 \(Z\) 在 \(\mathsf{P}\) 上簡單，也就是一個域；又設 \(\Omega\) 是 \(\mathsf{P}\) 的代數閉擴張域。於是眾所周知：\(Z_\Omega\) 當且僅當 \(Z\) 在 \(\mathsf{P}\) 上可分、也就是第一類擴張時，仍然是沒有根基的系統。因為當且僅當如此時，它才擁有同它關於 \(P\) 的秩一樣多的一階同構映射到 \(\Omega\) 中，因而也擁有同樣多的不同表示；在這裡，表示與表示類相重合。

\(Z_\Omega\) 中可能出現根基，因而並不必然有分解為絕對簡單分量的直和分解；這一事實引出如下定義：\(P\) 的一個擴張域 \(\Lambda\) 稱為分裂域，如果 \(Z_\Lambda\) 分裂為一階合成因子；換句話說，如果所有絕對不可約表示已經都位於 \(\Lambda\) 中。\(P\) 的一個擴張域 \(T\) 稱為分解域，如果 \(Z_{\mathsf{T}}\) 至少分出一個一階合成因子；換句話說，如果 \(Z\) 在 \(\mathsf{T}\) 中至少存在一個絕對不可約表示。

分裂域和分解域可由下述命題刻畫：

\paragraph{定理。}
一個域 \(\mathsf{T}\) 當且僅當含有一個同 \(Z\) 同構的子域 \(Z'\) 時，是分解域；一個域 \(\Lambda\) 當且僅當含有一個子域 \(\Gamma\)，且 \(\Gamma\) 與屬於 \(Z\) 的伽羅瓦域同構時，是分裂域。

這直接由藉助絕對不可約表示給出的分裂域與分解域定義得到。特別地，類似於非交換情形，可以推出：一個域 \(Z'\) 當且僅當同 \(Z\) 同構時，才是極小分解域，也就是說，沒有真子域仍為分解域。一個極小分解域，當且僅當 \(Z\) 在 \(\mathsf{P}\) 上正規、也就是為伽羅瓦域時，同時也是分裂域。

這正是把一個簡單代數稱為在其中心上正規的原因。在固定的代數閉 \(\Omega\) 內，在 \(\mathsf{P}\) 上只有唯一一個極小分裂域，即屬於 \(\mathsf{Z}\) 的伽羅瓦域；這同非交換情形中一般存在無窮多個非同構的極小分裂域\footnote{R. Brauer--E. Noether，同前引。}形成對照。其原因在於，交換情形中直和分解是唯一的，而非交換情形中的唯一性只到算子同構為止；或者等價地說，在交換情形中，只有有限多個不同的絕對不可約表示，而在非交換情形中，有無窮多個不同的絕對不可約表示，雖然它們屬於同一個類。

\paragraph{2. 可分域 \(Z/\mathsf{P}\) 情形中伽羅瓦理論的超複奠基。}
同 §6, 3. 完全類似，這裡對任意在 \(\mathsf{P}\) 上可分的 \(Z\)，有同構理論；當 \(Z\) 是伽羅瓦域時，便可以過渡到通常的自同構群。

\paragraph{前提。}
設簡單系統 \(Z\) 在 \(\mathsf{P}\) 上是有限可分擴張域；\(\Omega\) 表示一個在 \(\mathsf{P}\) 上有限或無限的分裂域；\(\mathsf{Z}'=\mathsf{Z}^{(1)}\) 是同 \(Z\) 同構的子域；\(\Gamma\) 是相應的伽羅瓦域。按 1. 因而有直和分解
\[
        Z_\Omega=\mathfrak r^{(1)}+\ldots+\mathfrak r^{(n)}
        =e^{(1)}Z_\Omega+\ldots+e^{(n)}Z_\Omega
        =e^{(1)}\Omega+\ldots+e^{(n)}\Omega .
\]
由 \(e^{(i)}\) 生成的表示 \(Z\to\mathsf{Z}^{(i)}\)，其中 \(\mathsf{Z}^{(i)}\subseteq\Gamma\)，由下式定義：
\[
        e^{(i)}z=e^{(i)}z^{(i)}.
\]

\paragraph{輔助定理 1。}
由 \(e^{(1)},\ldots,e^{(n)}\) 生成的這 \(n\) 個表示全都不同。證明同 §6, 3. 一樣，也可由 1. 得出；因為它們屬於不同的表示類。

\paragraph{輔助定理 2。}
若 \(T\) 是 \(\mathsf{P}\) 與 \(Z\) 之間的中間域，且 \(s\) 是 \(Z\) 關於 \(T\) 的秩，那麼 \(T\) 到 \(\Omega\) 中的每個同構至少、因而也恰好，容許 \(s\) 個延拓。

證明同 §6, 3. 一樣。這裡把“至少”加強為“恰好”，來自 \(Z_\Omega\) 的唯一分解；按此分解，\(Z\) 在 \(\Omega\) 中不是至少，而是恰好擁有 \(n\) 個同構。

現在如 §6, 3. 那樣，定義由 \(T\) 誘導的有限同構集合 \(\mathfrak I\) 的類劃分 \(\mathfrak I_T\)：\(\mathfrak I\) 中全部且僅有那些在 \(T\) 上誘導同一同構的同構，在 \(\mathfrak I_T\) 中被視為等價。於是證明如下

\paragraph{主定理。}
若 \(\mathfrak I_T\) 是由 \(T\) 誘導的類劃分，則 \(T\) 是相對於這個類劃分的極大子域。

從這裡出發，在 \(Z\) 為伽羅瓦域的情形中，正如 §6, 3. 一樣，可以通過同某個同構的逆相複合，過渡到自同構群以及通常表述。關於子群的相應斷言，則由冪等元的考察得到；這個考察本身也有意義。

\paragraph{3. \(Z_\Omega\) 的冪等元。}
令 \(S\) 遍歷把 \(Z\) 帶到各個 \(\mathsf{Z}^{(i)}\) 中的 \(n\) 個代換，即同構 \(\mathsf{Z}\to Z\) 與 \(Z\to\mathsf{Z}^{(i)}\) 的複合，其中前者表示與 \(Z\to\mathsf{Z}\) 互逆的同構。不必假定 \(Z\) 是伽羅瓦域。對 \(Z_{\mathsf{Z}}\) 中的元素 \(a\)，把代換 \(S\) 定義為算子，規定它們在 \(Z\) 上誘導恆等映射；於是對每個 \(a\)，都定義了位於 \(Z_{\mathsf{Z}^{(i)}}\) 中的共軛元 \(a^S\)。在這些規定下，有

\paragraph{定理。}
\(Z_\Omega\) 的 \(n\) 個冪等元 \(e^{(i)}\) 互為共軛：
\[
        e^{(i)}=e^S,\qquad e=e^{(1)};
\]
它們位於 \(n\) 個共軛的擴張系統 \(Z_{\mathsf{Z}^{(i)}}\) 中。

事實上，應用 \(S\) 會把定義關係
\[
        ez=e\zeta^{(1)},\qquad e^2=e
\]
分別變為
\[
        e^S z=e^S\zeta^{(i)},\qquad (e^S)^2=e^S .
\]
由於分解的唯一性，冪等元也唯一確定；所以 \(e^{(i)}=e^S\)。又因為 \(\mathsf{Z}\) 是分解域，因而已經實現表示 \(ez=e\zeta\)，所以 \(e\) 已經位於 \(Z_{\mathsf{Z}}\) 中，從而 \(e^{(i)}\) 位於 \(Z_{\mathsf{Z}^{(i)}}\) 中。

若特別地 \(Z\) 是伽羅瓦域，則關於子群的伽羅瓦理論主定理部分立即推出；因為若 \(\mathfrak H\) 是伽羅瓦群 \(\Gcal\) 的子群，則由於直和各分量唯一確定，元素
\[
        \sum_{H\in\mathfrak H} e^H
\]
容許來自 \(\mathfrak H\) 的全部替換，而不容許除此以外的替換。由於群 \(\Gcal\) 是為 \(\mathsf{Z}\) 定義的，這裡討論的是 \(\mathsf{Z}\) 的伽羅瓦理論；藉助同構，\(Z'\) 的相應理論也同時得到處理。

\paragraph{4. 冪等元同互補基的聯繫。}
這裡仍不必假定 \(Z\) 為伽羅瓦域。

\paragraph{定理。}
若 \(a_1,\ldots,a_n\) 是 \(Z\) 在 \(\mathsf{P}\) 上的一組基，並置冪等元
\[
        e=a_1\beta_1+\ldots+a_n\beta_n,
\]
也就是
\[
        e^S=a_1\beta_1^S+\ldots+a_n\beta_n^S,
\]
則
\[
        \beta_1^S,\ldots,\beta_n^S
\]
表示在由 \(e^S\) 給出的映射下，\(a_1,\ldots,a_n\) 的互補基 \(b_1,\ldots,b_n\) 的像。

事實上，對由 \(e^S\) 給出的映射 \(e^Sz=e^S\zeta^S\)，特別是 \(e^Sa_i=e^S\alpha_i^S\)，由於
\[
        e^S\cdot e^S=e^S,\qquad e^S\cdot e^R=0\quad(S\ne R)
\]
成立，便有對應 \(e^S\to1\)、\(e^R\to0\)。於是得到
\[
        1=\alpha_1^S\beta_1^S+\ldots+\alpha_n^S\beta_n^S;\qquad
        0=\alpha_1^S\beta_1^R+\ldots+\alpha_n^S\beta_n^R\quad(S\ne R).
\]
寫成矩陣形式，這正是互補基的定義方程：
\[
\begin{pmatrix}
\alpha_1 & \ldots & \alpha_n\\
\ldots & \ldots & \ldots\\
\alpha_1^S & \ldots & \alpha_n^S\\
\ldots & \ldots & \ldots
\end{pmatrix}
\cdot
\begin{pmatrix}
\beta_1 & \ldots & \beta_1^R & \ldots\\
\ldots & \ldots & \ldots & \ldots\\ \ldots & \ldots & \ldots & \ldots\\
\beta_n & \ldots & \beta_n^R & \ldots
\end{pmatrix}
=
E.
\]
過渡到跡定義，是因為對
\[
        a_i=\sum_S e^S\alpha_i^S,\qquad b_i=\sum_S e^S\beta_i^S
\]
而言，該矩陣關係在交換行之後變為
\[
        \mathit{Sp}(a_i b_i)=1;\qquad \mathit{Sp}(a_i b_k)=0\quad\text{對於 }i\ne k,
\]
其中 \(a_i\) 與 \(b_i\) 位於 \(Z\) 中，因為它們容許全部替換 \(S\)。\footnote{參見“表示論”§25。}\footnote{互補基與冪等元之間的關係最早見於 Dedekind, \emph{Zur Theorie der aus \(n\) Haupteinheiten gebildeten komplexen Größen}; 參見《全集》第二卷，第 4--5 頁。}

互補基的反變性，在這裡來自這樣一個事實：冪等元 \(e^S\) 是 \(Z\) 的不變量，且與所取基無關；也就是說，在過渡到 \(\mathsf{Z}/\mathsf{P}\) 的另一組基
\[
        a_1',\ldots,a_n'
\]
時，所有
\[
        a_1\beta_1^S+\ldots+a_n\beta_n^S
\]
都轉入自身。

\subsection*{§9. 任意系統的分裂域與分解域}

\paragraph{1. 定義以及向簡單情形的歸約。}
在一般情形中，同 §8 所給定義相應，有如下定義：設 \(S\) 是 \(P\) 上的超複系統。\(P\) 的一個交換擴張域 \(\Lambda\) 稱為 \(S\) 的分裂域，如果 \(S_\Lambda\) 在由單側理想構成的一條合成列中分裂為絕對簡單合成因子；換句話說，如果 \(S\) 在 \(\Lambda\) 中的全部不可約表示已經都是絕對不可約的。\(P\) 的一個擴張域 \(T\) 稱為分解域，如果 \(S_{\mathsf{T}}\) 至少分出一個絕對簡單合成因子，也就是如果至少有一個在 \(\mathsf{T}\) 中不可約的表示已經是絕對不可約的。

這些定義顯然包含 §8 中對交換情形給出的定義，也包含 §7 中對簡單代數給出的定義；後者是因為這裡 \(A_\Lambda\) 在中心的任一交換擴張下都是完全可約的，因而合成因子同直和分解的分量一致。但是，只有中心上的全矩陣環給出絕對簡單分量，從而也給出絕對不可約表示。又因為 \(A_\Lambda\) 是雙側簡單的，所以它分解為算子同構的分量；於是，分出一個絕對簡單分量就已經意味着完全分裂：分解域與分裂域在這裡重合。

由這些定義立即得到如下事實：\(S\) 的分裂域是由每個在 \(\mathsf{P}\) 中不可約的表示各取一個分裂域而成的合成域；\(S\) 的分解域則是任一在 \(\mathsf{P}\) 中不可約表示的分解域。由於簡單系統、按根基取的剩餘類環的分量、以及在 \(\mathsf{P}\) 中不可約表示之間有一一對應，所以只需考察簡單系統。在這裡，結果將由 §7 與 §8 組合得到。

\paragraph{2. 簡單系統的分裂域與分解域。}
先考察簡單系統 \(A\) 的中心 \(Z\) 在 \(\mathsf{P}\) 上可分的情形。由 §8, 1. 與 §8, 3. 得到：若 \(\Omega\) 表示 \(Z\) 的一個分裂域，那麼同 \(Z_\Omega\) 的直和分解相應的 \(A_\Omega\) 的雙側分解，給出 \(A_\Omega\) 的 \(n\) 個互為共軛且同 \(A\) 同構的分量
\[
        e^S A_\Omega ,
\]
它們成為其中心
\[
        e^S Z_\Omega=e^S Z
\]
上的簡單代數。這裡，把替換 \(S\) 定義為 \(A_\Omega\) 的算子，規定它們在 \(A\) 上誘導恆等映射。

事實上，按 §8, 3.，冪等元是共軛的；因而按對 \(A\) 的替換定義，各分量 \(e^S A_\Omega\) 也同樣共軛。由於 \(A\) 是雙側簡單的，\(e^S A_\Omega\) 同 \(A\) 環同構。又由於 \(e^S Z_\Omega=e^S Z\)，有 \(e^S A_\Omega=e^S A_{Z^S}\)；所以中心同系數域重合，這些分量於是為正規代數。

§7, 2. 的結果進一步給出：若 \(A\) 是 \(\mathsf{P}\) 上的簡單系統，且其中心 \(Z\) 在 \(\mathsf{P}\) 上可分，則 \(A_\Omega\) 也沒有根基，因而完全可約；這裡 \(\Omega\) 可以表示任一有限或無限擴張域。因為按 §7, 2.，\(e^S A_\Omega\) 在每次係數擴張下仍是簡單的，因而沒有根基。所以直和也同樣如此；而這個直和就是 \(A_\Omega\)，或者當 \(\Omega\) 不是 \(Z\) 的分裂域時，它由 \(A_\Omega\) 再作係數擴張得到。

進一步，由 §7, 2. 得到如下

\paragraph{分解域與分裂域的刻畫。}
若 \(A\) 是中心 \(Z\) 可分的除代數，則分解域的不可約嵌入，恰由 \(A_f\) 中全部且僅有的包含 \(Z\) 的極大交換子域給出；分裂域全部且僅是一個分解域同中心的一個分裂域的合成域，特別是同該中心所屬伽羅瓦域的合成域。即使中心不是伽羅瓦域，一個極小分解域也可能是分裂域。

因為一個分解域必須包含一個同 \(Z\) 同構的域；關於分解域的嵌入斷言，現在由前述事實和 §7, 2. 推出。一個分裂域還必須包含一個屬於 \(Z\) 的分裂域 \(\Omega\)。但是，\(\Omega\) 上的每個分解域同時也是分裂域，因為各分量 \(e^{(i)}A_\Omega\) 是簡單的，並且具有一個同 \(\Omega\) 同構的中心。若 \(Z\) 是伽羅瓦域，則極小分解域與極小分裂域重合。但即使在非伽羅瓦情形中，也可能與交換情形不同，存在同時也是分裂域的極小分解域；這正發生在這樣一個極小分解域已經包含屬於 \(Z\) 的伽羅瓦域時。例如：中心同 \(\mathsf{P}(\sqrt[3]{2})\) 同構的四元數域，其中 \(\mathsf{P}\) 是有理數域；屬於 \(\mathsf{P}(\sqrt[3]{2})\) 的伽羅瓦域就是一個極小分解域和分裂域。

最後，若中心不可分，則 \(Z_\Omega\)，從而 \(A_\Omega\)，成為帶根基的系統。設 \(\mathfrak C\) 是 \(Z_\Omega\) 的根基，而 \(\mathfrak c\) 是 \(A_\Omega\) 中的擴張理想。於是，對 \(Z_\Omega/\mathfrak C\)，有一種同 §8, 2. 完全類似的分解為共軛分量的直和分解；這又如上面一樣給出 \(A_\Omega/\mathfrak c\) 的直和分解，使得各分量 \(\overline e^{(i)}A\) 同 \(A\) 同構，並以 \(\overline e^{(i)}\mathsf{Z}^{(i)}\) 為中心。因此，關於分解域與分裂域的定理仍然保持成立。進一步，秩數計數表明，與 \(\mathfrak C\) 的一條合成列相應的合成因子，對於 \(A_\Omega/\mathfrak c\) 成為算子同構，而不僅僅是同態。

用不可約表示來表述，則得到如下概括：若給定一個超複系統的、在 \(\mathsf{P}\) 中不可約的表示，那麼當且僅當相應中心在 \(\mathsf{P}\) 上可分時，該表示才成為絕對完全可約的。在任何情形中，該表示的分裂域與分解域都可通過如上嵌入到相應簡單系統中來刻畫。特別地，一個分解域的最低次數，等於中心的次數同相應代數類在該中心上的指數之積；每個分解域的次數都是這個數的倍數。該表示按合成列意義分裂為若干不同但共軛的絕對不可約表示類，其數目等於中心中所含最大可分子域的次數。每個類出現的重數，等於指數同中心關於這個可分擴張的次數之積。

在完全基礎域上，由於只有可分擴張域，我們便回到 I. Schur 的已知結果，同時添入一個已見於 R. Brauer 的嵌入刻畫。

\begin{center}
（1932 年 6 月 8 日收稿。）
\end{center}

%
\clearpage
%
%
\begingroup
\newcommand{\srcnumdisplay}[3][3.4em]{%
  \par\smallskip\noindent
  \begin{minipage}[t]{#1}\vspace*{0pt}#2\end{minipage}%
  \begin{minipage}[t]{\dimexpr\linewidth-#1\relax}\vspace*{0pt}\centering
  \(\displaystyle #3\)
  \end{minipage}\par\smallskip}
\pdfbookmark[0]{41. 相對伽羅瓦數域的主屬定理}{paper-41}
\section*{41. 相對伽羅瓦數域的主屬定理}

\begin{center}
著者\\[0.4ex]
\textbf{Emmy Noether（哥廷根）}\\[0.8ex]
\emph{Math. Ann. 108 (1933), S. 411--419}
\end{center}

近來人們發現，非交換方法，特別是代數理論，使我們能夠把關於相對循環數域和相對阿貝爾數域的已知定理表述並證明為適用於任意相對伽羅瓦域的定理。\footnote{參見我在蘇黎世所作的報告：\emph{Hyperkomplexe Systeme in ihren Beziehungen zur kommutativen Algebra und Zahlentheorie}, Verhandl. des internat. Math. Kongr. Zürich, 1932, Bd. I；其中尤其詳述了範數定理和主屬定理的超複表述，並列有相關文獻。關於分裂代數的定理，另參 H. Hasse, \emph{Die Struktur der R. Brauerschen Algebrenklassengruppe}. (6.1). Math. Annalen 107 (1932)。}我首先要提到“範數定理”；一般說來，它可以表述為一個關於分裂代數的定理。下文將說明，與此相應，主屬定理也可以獲得一般的超複表述，即表述為關於內自同構的定理，或者表述為關於“交叉表示”的定理。這個證明由上述關於分裂代數的定理經純粹的代數—算術推理得出；而在上述定理本身中，循環情形的範數定理——只須考慮素數次數即已足夠——構成其超越性核心。為了便於理解這一表述，我先陳述下文並不使用的“最小情形中的主屬定理”\footnote{我稱之為“最小情形”，因為它是任意域中主屬定理的代數類似物；而所謂“小情形”已經包含了過渡到 $p$-進基域的意義。}。這是一個初等代數定理\footnote{該定理見 A. Speiser, \emph{Zahlentheoretische Sätze aus der Gruppentheorie}, S. 3, Math. Zeitschr. 5 (1919), S. 1--6。該文已經把它表述為關於交叉表示的定理。}；在循環特例中，它化為 Hilbert《數論報告》中熟知的定理 90：$N(a)=1$，當且僅當 $a=b^{1-S}$。

我也把這個定理直接表述為關於內自同構或交叉表示的定理，並藉助關於代數的內自同構及交叉積表示的一般定理來證明它。在真正的主屬定理中，取代後一構造的是理想類群與伽羅瓦群的交叉積，但其中要對因子系採用一種更細的誘導理想類劃分，從而把類域論中射線類劃分的類似物納入進來。由於這裡尚缺少關於自同構和表示的一般定理——主屬定理可以看作朝這個方向邁出的第一步——所以我如上所述，藉助分裂代數定理作一化歸。由此，定理被化歸為關於理想而不是理想類的相應定理；在這裡，Brandt 關於一個代數的極大階之間關係的定理\footnote{Brandt 的極大階理論見 H. Hasse, \emph{Über \(\mathfrak p\)-adische Schiefkörper und ihre Bedeutung für die Arithmetik hyperkomplexer Zahlsysteme}, Math. Ann. 104 (1931)；該文對這一理論作了系統表述，並詳盡地重新奠定了基礎。}可以看作自同構定理的等價形式。事實上，把這些定理同對交叉積之極大階的考察\footnote{C. Chevalley（Hamb. Ber.）、H. Hasse 和 E. Noether 將各發表一篇短文；後兩篇將收入一卷 Herbrand 紀念文集。}結合起來，便會得到一個與最小情形中的證明本質上平行的證明，儘管由於必須逐一回到各個位，證明要複雜一些。因此，我寧願採用 E. Artin 指出的一個幾乎平凡的證明；在這裡，它構成 Speiser 證明\textsuperscript{3)} 的一個甚至更簡單的類似物。

\subsection*{§ 1.}
\begin{center}
\textbf{最小情形中的主屬定理}
\end{center}

本節中，\(k\) 表示一個任意的交換域（不必是數域），\(K/k\) 是一個 \(n\) 次可分伽羅瓦擴張，\(\mathfrak G\) 是相應的伽羅瓦群。

\noindent\textbf{1.} 先把關於交叉積的已知事實簡要彙集如下。\footnote{這一理論由 H. Hasse 根據我的一次課程（1929/30）寫成，見 \emph{Theory of cyclic algebras}, Kap. 2, Am. Transact. 34 (1932)。另參 M. Deuring 關於超複數的一篇綜述，該文將收入 \emph{Ergebnisse der Mathematik}。這裡的概述取自我在注 1 所引的蘇黎世報告。}所謂交叉積，是指把 \(K\) 與 \(\mathfrak G\) 同時嵌入一個代數 \(A\)，使 \(K\) 的自同構成為內自同構。令 \(u_{S_1},\ldots,u_{S_n}\) 表示與這 \(n\) 個群元素相應的符號；先把 \(A\) 取為 \(K\) 上秩為 \(n\) 的線性形式模：

\srcnumdisplay{(1)}{A=u_{S_1}K+\ldots+u_{S_n}K.}

要求由 \(u_S\)，更一般地由 \(u_SK^*\)\footnote{\(K^*\) 由 \(K\) 去掉零元而得；下文一律使用這一記號。}，生成內自同構；在這一要求下，\(A\) 成為一個環，也就是 \(k\) 上秩為 \(n^2\) 的代數。這個要求具體表示為

\srcnumdisplay{(2)}{u_S^{-1}zu_S=z^S\quad\text{或}\quad zu_S=u_Sz^S}

其中 \(z\) 是 \(K\) 中的任意元素，並且

\srcnumdisplay{(3)}{u_Su_T=u_{ST}a_{S,T}\quad\text{其中 }a_{S,T}\text{ 屬於 }K^*,}

\srcnumdisplay{(4)}{a_{ST,R}a_{S,T}^{R}=a_{S,TR}a_{T,R}}

（結合律等價於 \([u_Su_T]u_R=u_S[u_Tu_R]\)）。

再用
\[
\sum u_Sb_S\cdot\sum u_Tc_T=\sum u_Su_Tb_S^Tc_T
\]
定義任意兩個元素的乘積，\(A\) 就成為 \(K\) 與 \(\mathfrak G\) 關於因子系 \(a_{S,T}\) 的交叉積，記作 \(A=(a_{S,T},K,\mathfrak G)\)。可以證明，\(A\) 是 \(k\) 上的中心單代數，因而是相應除代數 \(D\) 上的某個 \(r\) 階矩陣環 \(D_r\)；並且 \(K\) 是極大交換子域，因而是分裂域。反過來，對一個預先給定的除代數 \(D\)，總有一些矩陣環 \(D_r\) 能按上述方式生成為交叉積。

若從 \(u_S\) 過渡到 \(v_S=u_Sc_S\)，其中 \(c_S\) 屬於 \(K^*\)，所生成的自同構不變，而得到“相伴”的因子系

\srcnumdisplay{(5)}{\bar a_{S,T}=a_{S,T}c_S^Tc_T/c_{ST}.}

特別地，由 (5) 可見，因子系 \(\bar a_{S,T}=c_S^Tc_T/c_{ST}\) 與單位因子系相伴；這些量 \(c_S^Tc_T/c_{ST}\) 稱為變換量。相伴的因子系歸入同一個類 (\((a)\))，因而所有變換量 \(c_S^Tc_T/c_{ST}\) 都歸入類 ((1))。同樣，把所有與 \(A\) 相似的代數，即 \(r=1,2,\ldots\) 時的所有 \(D_r\)，歸入一個類 \(\mathfrak A\)。類 \(\mathfrak A\) 與 (\((a)\)) 一一對應。\emph{具有固定分裂域 \(K\) 的這些類，在直接乘積運算下構成一個阿貝爾群；這個群與因子系各類在逐項乘積下構成的群同構。單位元分別是分裂代數類，以及由全部變換量 \(c_S^Tc_T/c_{ST}\) 構成的系統。}這就是 R. Brauer 的代數類群。

再簡要給出循環代數這一特例。若 \(Z/k\) 是循環擴張，\(S\) 是其群的一個生成代換，則可令 \(S\) 的各次冪同 \(u\) 的各次冪對應，於是有

\srcnumdisplay{(1')}{A=Z+uZ+\ldots+u^{n-1}Z.}
\srcnumdisplay{(2')}{zu=uz^S.}
\srcnumdisplay{(3')}{u^n=\alpha.}
\srcnumdisplay{(4')}{\alpha\text{ 屬於基域 }k^*.}
\srcnumdisplay{(5')}{\bar\alpha=\alpha\cdot N(c),\quad\text{若令 }v=uc.}

因此，這裡的每個因子系都只由基域中的一個元素 \(\alpha\) 構成——記作 \(A=(\alpha,Z,S)\)；因子系的單位類由 \(Z^*\) 中各元素的範數給出，而代數類群與商群 \(k^*/N(Z^*)\) 同構。

\noindent\textbf{2.} 最小情形中主屬定理的表述利用這樣一個事實：關係 (2) 至 (5) 純粹是乘法關係，因而同時定義了 \(\mathfrak G\) 的一個擴張群 \(\mathfrak G^*\)，它由 \(n\) 個陪集 \(u_SK^*\) 中的元素組成：

\srcnumdisplay{(6)}{\mathfrak G^*=\{u_{S_1}K^*,\ldots,u_{S_n}K^*\}.}

\(\mathfrak G^*\) 以 \(K^*\) 為阿貝爾正規子群，而 \(\mathfrak G^*/K^*\) 與 \(\mathfrak G\) 同構。\footnote{因此，關係 (2) 至 (5) 無非就是群擴張的定義關係。眾所周知，只要 \(K^*\) 的全自同構群中有一個子群——這裡僅把 \(K^*\) 看作乘法群——是 \(\mathfrak G\) 的同態像，這樣的擴張就存在。交叉積的特殊之處在於：這個同態成為同構；而且，上述子群恰取為同時保持乘法和加法的全部自同構，也就是 \(K\) 的伽羅瓦群；這樣，在得到群擴張的同時也得到環擴張。}\(\mathfrak G^*\) 由所有把 \(K\) 變換到自身的（可逆）元素 \(g^*\) 組成，也就是說，它們滿足 \(g^{*-1}Kg^*=K\)；因而 \(K\) 的環自同構由且僅由 \(\mathfrak G^*\) 中的元素生成。事實上，每個這樣的 \(g^*\) 都生成 \(K\) 的一個環自同構。反過來，每個這樣的自同構都由某個元素 \(v=u_Sw\) 的共軛變換生成，其中 \(w\) 在 \(K\) 上誘導恆等映射。於是 \(w\) 與 \(K\) 的所有元素可交換，故 \(w\) 屬於 \(K\)；又因 \(K\) 是極大交換子環，所以 \(w\in K^*\)。

\emph{最小情形中的主屬定理。第一種表述：\(\mathfrak G^*\) 的每個群自同構，只要它是 \(K^*\) 上恆等自同構的延拓，就是內自同構，並且由 \(K^*\) 中的一個元素生成。第二種表述：若所有變換量 \(c_S^Tc_T/c_{ST}\) 都等於一，則 \(c_S\) 是符號意義下的 \((1-S)\) 次冪；也就是說，存在 \(b\in K^*\)，使得對每個 \(S\in\mathfrak G\)，都有 \(c_S=b^{1-S}=b/b^S\)。第三種表述：群 \(\mathfrak G\) 在 \(K^*\) 中只有一個屬於單位因子系的一次交叉表示類。}

這裡，如果 \(C_S^TC_T=C_{ST}a_{S,T}\)，就稱表示 \(u_S\to C_S\) 是關於因子系 \(a_{S,T}\) 的交叉表示。若 \(C_S=B^{-S}D_SB\)，則兩個表示屬於同一個類。

我先證明第一種表述，再說明它與第二、第三種表述的一致性。這個證明依據的是：所述類型的每個群自同構都可以延拓為 \(A\) 的一個環自同構，而眾所周知，後者是內自同構。設在這個自同構下，\(u_S\) 的像為 \(v_S\)；那麼 \(\sum v_SK\) 又是 \(\mathfrak G\) 與 \(K\) 關於同一因子系構成的交叉積，因為按假設，所有關係 (2) 至 (4) 都保持不變。對應關係
\[
\sum u_Sb_S\to\sum v_Sb_S
\]
給出一個環自同構，並且 \(K\) 中的每個元素都對應於自身。因而，根據上面對 \(\mathfrak G^*\) 的討論，這個自同構由 \(K^*\) 中的一個元素生成。

過渡到第二種表述依據的是：由 (2) 可知，\(v_S\) 必須具有形式 \(v_S=u_Sc_S\)；由於因子系也保持不變，根據 (5)，相應的變換量必須等於一。反過來，若變換量等於一，則由 (2) 至 (5) 可見，代換 \(v_S=u_Sc_S\) 給出 \(\mathfrak G^*\) 的一個所述類型的自同構。根據第一種表述，於是有
\[
v_S=b^{-1}u_Sb=u_Sb/b^S,
\qquad\text{因而}\qquad c_S=b^{1-S}.
\]
在循環情形中，第二種表述的假設根據 (5') 特別化為 \(N(c)=1\)；而 \(c=b^{1-S}\) 正是熟知的定理。\footnote{這裡給出的證明尤其也適用於 \(k\) 只有有限多個元素的情形；Hilbert 的證明在這種情形下失效，但 Speiser 的證明仍然有效。}

第三種表述與第二種表述直接等價，因為後者的假設恰好說 \(u_S\to c_S\) 構成一個關於單位因子系的一次交叉表示；結論 \(c_S=b^{1-S}\) 則說這個表示等價於單位表示：\(c_S\sim1\)。

不過還應指出，第三種表述只是下面這個更一般定理的一個特例：對一個給定的因子系 \(a_{S,T}\)，只有一個不可約交叉表示類；它的次數為 \(m\)，這裡 \(m\) 是相應代數的 Schur 指數\footnote{I. Schur, \emph{Einige Bemerkungen zu der vorstehenden Arbeit des Herrn A. Speiser}\textsuperscript{3)}, Math. Zeitschr. 5 (1919), S. 7--10。關於一個以交叉表示模為基礎的證明［出自我在注 6 所引的課程］，參見注 6 所引 M. Deuring 的綜述。}，也就是相應除代數的次數。

\subsection*{§ 2.}
\begin{center}
\textbf{主屬定理}
\end{center}

\noindent\textbf{1. 關於類劃分的預備說明。} 在類域論中關於循環相對域的主屬定理裡，眾所周知會出現兩種不同的類劃分：上域中的絕對理想類和下域中的射線類。在一般情形中，與此相應的是：上域理想的一種類劃分會為因子系誘導出一種更細的理想類劃分。

首先令 \(\mathfrak I\) 表示 \(K\) 中全部理想組成的群。由於 \(\mathfrak G\) 在 \(\mathfrak I\) 上誘導自同構，所以按 (2) 至 (5) 定義的、由 \(\mathfrak G\) 作出的擴張，由 \(n\) 個陪集 \(\{\ldots u_S\mathfrak I\ldots\}\) 中的元素組成。在 \(\mathfrak I\) 中規定等價關係，過渡到絕對理想類群 \(\overline{\mathfrak I}\) 後，\(n\) 個陪集 \(u_S\overline{\mathfrak I}\) 仍然有定義；這是因為 \(\mathfrak G\) 也在 \(\overline{\mathfrak I}\) 上誘導自同構，主類在 \(\mathfrak G\) 的作用下保持不變。\footnote{一般地，只要選作主類的是 \(\mathfrak I\) 的一個容許 \(\mathfrak G\) 作用的子群，\(\{\ldots u_S\overline{\mathfrak I}\ldots\}\) 就有定義；這時，由該子群中的理想得到的變換量取代主理想的變換量。}這就是說，從 \(\mathfrak I\) 到 \(\overline{\mathfrak I}\) 的等價關係可以由 \(\{\ldots u_S\mathfrak I\ldots\}\) 傳遞到 \(\{\ldots u_S\overline{\mathfrak I}\ldots\}\)：可以把等價理想 \(\mathfrak a\) 和 \(\mathfrak b\) 分別對應到等價的乘積 \(u_S\mathfrak a\) 和 \(u_S\mathfrak b\)。特別地，\(u_S\) 與所有乘積 \(u_S(c_S)\) 等價，其中 \((c_S)\) 遍歷 \(\mathfrak I\) 中的全部主理想。

現在要求乘法關係 (3) 在這種等價意義下一意確定，就等於要求由主理想得到的全部變換量 \((c_S^T)(c_T)/(c_{ST})\) 都等價於單位因子系。換句話說：

\emph{在由 \(n\) 個陪集組成的系統 \(\{\ldots u_S\overline{\mathfrak I}\ldots\}\) 中，\(\overline{\mathfrak I}\) 表示 \(K\) 的絕對理想類群，\(H\) 特別表示主類；對 \(u_SH\) 而言，乘法關係 (3) 當且僅當滿足下列條件時才一意確定：在因子系的誘導理想類劃分中，主類包含由主理想得到的全部變換量，也就是由 \(H\) 中的生成元 \((c_S)\) 得到的全部變換量。}

根據這一事實，我給出以下定義，作為射線類劃分的推廣：

\emph{定義：因子系的誘導理想類劃分規定如下：主類由所有這樣的主理想 \((a_{S,T})\) 組成——對每個這樣的主理想系統，至少存在一組基元素 \(a_{S,T}\)，使代數 \((a_{S,T},K,\mathfrak G)\) 在 \(K\) 的所有分歧位上都分裂。}

這些主理想 \((a_{S,T})\) 按因子系的乘法組成一個群；這個群包含變換量 \((c_S^T)(c_T)/(c_{ST})\)，因為因子系 \(c_S^Tc_T/c_{ST}\) 生成處處分裂的代數。

\noindent\textbf{2. 主屬定理的表述。} 與最小情形中的定理相應，我把這個定理表述成三個彼此等價的形式。

\emph{第一種表述：以因子系的誘導理想類劃分為基礎，如果代換 \(v_S=u_S\bar{\mathfrak c}_S\)\footnote{\(\bar{\mathfrak c}_S\) 表示 \(\mathfrak c_S\) 的理想類。}產生 \(n\) 個陪集 \(\{\ldots u_S\overline{\mathfrak I}\ldots\}\) 的一個自同構——也就是說，\(v_S\) 在這種類劃分的意義下滿足同樣的關係 (3)——，那麼這個自同構是內自同構，並且由一個理想類 \(\bar{\mathfrak b}\) 生成。}

\emph{第二種表述：若由理想類 \(\bar{\mathfrak c}_S\) 構成的變換量 \(\bar{\mathfrak c}_S^T\bar{\mathfrak c}_T/\bar{\mathfrak c}_{ST}\) 全都屬於因子系誘導類劃分中的主類——由理想類構成的全部這些“向量” \(\{\ldots\bar{\mathfrak c}_S\ldots\}\) 組成主屬——，則類 \(\bar{\mathfrak c}_S\) 是符號意義下的 \((1-S)\) 次冪；也就是說，存在一個理想類 \(\bar{\mathfrak b}\)，使得對所有 \(S\in\mathfrak G\)，都有 \(\bar{\mathfrak c}_S=\bar{\mathfrak b}^{1-S}=\bar{\mathfrak b}/\bar{\mathfrak b}^S\)。}

\emph{第三種表述：以因子系的誘導理想類劃分為基礎，群 \(\mathfrak G\) 在 \(\overline{\mathfrak I}\) 中只有一個屬於因子系單位類的一次交叉表示類。}

這三個表述的等價性同 §~1 一樣得到；從第一種表述到第二種表述的過渡甚至更簡單，因為這裡已經假定自同構由 \(v_S=u_S\bar{\mathfrak c}_S\) 生成。這個假設是必要的，因為 \(\overline{\mathfrak I}\) 現在未必還是最大的交換子群（\(\overline{\mathfrak I}\) 的所有類都可能是不變類）。關於第三種表述，應當指出：由於 \(\overline{\mathfrak I}\) 中只定義了乘法運算，所以這裡作為表示矩陣出現的，只能是每一行和每一列都恰有一個非零元素的矩陣；對一次表示而言，這並不是限制。

\noindent\textbf{3. 主屬定理的證明。} 我證明第二種表述。證明基於兩個輔助引理。

\emph{輔助引理 1（理想的主屬定理）：若由理想 \(\mathfrak c_S\) 構成的變換量 \(\mathfrak c_S^T\mathfrak c_T/\mathfrak c_{ST}\) 等於單位理想，則 \(\mathfrak c_S\) 是符號意義下的 \((1-S)\) 次冪；對所有 \(S\in\mathfrak G\)，都有 \(\mathfrak c_S=\mathfrak b^{1-S}\)。與此等價的仍有第一種和第三種表述；在第一種表述中，必須像理想類的情形一樣，預先假定自同構由 \(v_S=u_S\mathfrak c_S\) 生成。}

\emph{證明。}\footnote{參見引言的末尾。}\(\mathfrak b=\sum\mathfrak c_S\) 即滿足要求；這裡的和指模之和，也就是最大公因子。事實上，
\[
\mathfrak b=\sum\mathfrak c_S;\quad
\mathfrak b^T=\sum\mathfrak c_S^T;\quad
\mathfrak b^T\mathfrak c_T=\sum\mathfrak c_S^T\mathfrak c_T
=\sum\mathfrak c_{ST}=\mathfrak b;\quad
\text{所以}\quad\mathfrak c_T=\mathfrak b^{1-T}.
\]

\emph{輔助引理 2（分裂代數定理的另一種表述）：若代數 \(A=(a_{S,T},K,\mathfrak G)\) 在 \(K\) 的所有（有限和無限）分歧位上都分裂，並且主理想 \((a_{S,T})\) 是理想的變換量——\((a_{S,T})=\mathfrak c_S^T\mathfrak c_T/\mathfrak c_{ST}\)——，那麼 \(A\) 處處分裂，也就是整體分裂。於是，\((a_{S,T})\) 是主理想的變換量：\((a_{S,T})=(d_S^T)(d_T)/(d_{ST})\)。\footnote{更精確地說，有 \(a_{S,T}=d_S^Td_T/d_{ST}\)。不過在主屬定理的證明中，只用到關於主理想的這個較弱結論。}反過來，每個處處分裂的代數都滿足這個條件。}

證明依據 \(A\) 在 \(\mathfrak p\)-進擴張下的熟知行為。令 \(k_{\mathfrak p}\) 表示 \(k\) 的 \(\mathfrak p\)-進擴張，相應地令 \(K_{\mathfrak P}\) 表示 \(K\) 的 \(\mathfrak P\)-進擴張，其中 \(\mathfrak P\) 是位於 \(\mathfrak p\) 上方的一個素理想；令 \(A_{\mathfrak p}\) 表示由 \(k\) 擴張到 \(k_{\mathfrak p}\) 後得到的 \(A\) 的擴張。再令 \(\mathfrak Z\) 表示與 \(\mathfrak P\) 相應的分解群，\(a_{\mathfrak Z}\) 表示因子系的相應截段（也就是說，\(a_{\mathfrak Z}\) 由那些 \(S,T\in\mathfrak Z\) 的 \(a_{S,T}\) 組成）。於是有
\[
A_{\mathfrak p}\sim(a_{\mathfrak Z},K_{\mathfrak P}/k_{\mathfrak p},\mathfrak Z).\footnotemark
\]
\footnotetext{參見 Hasse\textsuperscript{6)}, 14。一個稍簡單的證明如下。令 \(Z\) 表示與 \(\mathfrak P\) 相應的分解域；眾所周知，\(k_{\mathfrak p}\) 包含一個與它同構的 \(Z\)。於是 \(A_Z\sim(a_{\mathfrak Z},K/Z,\mathfrak Z)\)。這是因為，\(A_Z\) 與 \(A\) 中所有逐元素與 \(Z\) 可交換的元素所成的代數相似（例如參見 v. d. Waerden, Bd. II, S. 210，或我即將在 Math. Zeitschr. 發表的關於非交換代數的論文）。再把 \(A_Z\) 擴張到 \(A_{\mathfrak p}\)，就得到上述結果。選擇哪一個共軛的 \(\mathfrak P\) 並無影響；當 \(\mathfrak G=\sum\mathfrak ZT\) 時，不同選擇所得代數彼此同構，並通過 \(u_T\) 的共軛變換相互過渡。}

若 \(\mathfrak p\) 特別是\emph{不分歧的}，那麼 \(\mathfrak Z\) 是循環群；此時，上式是藉助不分歧慣性域 \(K_{\mathfrak P}/k_{\mathfrak p}\) 對 \(A_{\mathfrak p}\) 給出的標準表示。\footnote{例如參見注 4 所引 H. Hasse 的論文，§ 3。}

現在要證明，在輔助引理 2 的假設下，\(A_{\mathfrak p}\) 在這個不分歧情形中也分裂。對無限位 \(\mathfrak p\)，這裡 \(K_{\mathfrak P}=k_{\mathfrak p}\)，所以 \(A_{\mathfrak p}\) 分裂。對有限位 \(\mathfrak p\)，首先可知，因子系 \(a_{\mathfrak Z}\) 與一個全部由 \(K_{\mathfrak P}\) 中的單位組成的因子系相伴。事實上，在 \(K_{\mathfrak P}\) 中，所有理想 \(\mathfrak c_S\) 都成為主理想 \((c_S)\)；因而關於 \((a_{S,T})\) 的假設就是說：
\[
a_{S,T}=e_{S,T}\frac{c_S^Tc_T}{c_{ST}},
\]
其中 \(e_{S,T}\) 是單位。於是因子系 \(a_{\mathfrak Z}\) 與由單位 \(e_{S,T}\) 組成的因子系相伴。

進一步，在過渡到\emph{規範化的循環表示} (1') 至 (5') 時，這一性質仍然保持；也就是說，在 \(A_{\mathfrak p}=(\alpha,K_{\mathfrak P}/k_{\mathfrak p},R)\) 中，\(\alpha\) 是單位，其中 \(R\) 是 \(\mathfrak Z\) 的一個生成元。事實上，\(e_{S,T}\) 是因子系這一事實可表為關係
\[
u_{R^i}u_{R^j}=u_{R^{i+j}}e_{R^i,R^j}.
\]
令 \(u_R=u\)，便得到
\[
u^i=u_{R^i}e_{R,R^{i-1}}e_{R,R^{i-2}}\ldots e_{R,R};
\]
特別地，若 \(f\) 是 \(\mathfrak Z\) 的階，並注意 \(u_E=e_{E,E}\)，則
\[
u^f=\alpha=e_{E,E}e_{R,R^{f-1}}\ldots e_{R,R}.
\]
而在 \(K_{\mathfrak P}/k_{\mathfrak p}\) 不分歧時，所有單位都是範數\textsuperscript{16)}，所以 \(A_{\mathfrak p}\) 分裂。於是 \(A\) 處處分裂，從而根據分裂代數定理整體分裂。換句話說，從假設可推出 \(a_{S,T}\) 是元素的變換量：
\[
a_{S,T}=d_S^Td_T/d_{ST},\qquad d_S\in K^*;
\]
因而，較弱地說，\((a_{S,T})\) 是主理想的變換量。反過來，若 \(A\) 處處分裂，則主理想 \((a_{S,T})\) 也成為理想的變換量，因為這一點在每個位上都成立；而且 \(A\) 尤其在各分歧位上分裂。這確實只是分裂代數定理之假設的另一種表述。

由這兩個輔助引理立即得到\emph{主屬定理的證明}。對每個 \(S\)，令 \(\mathfrak c_S\) 表示類 \(\bar{\mathfrak c}_S\) 中的一個理想。根據因子系之誘導類劃分的定義，關於這些類的假設恰好說：\(\mathfrak c_S\) 的變換量是滿足輔助引理 2 之假設的主理想 \((a_{S,T})\)。因此，存在主理想 \((d_S)\)，使得對仍屬於類 \(\bar{\mathfrak c}_S\) 的理想 \(\mathfrak d_S=\mathfrak c_S/(d_S)\)，輔助引理 1 的假設成立：\(\mathfrak d_S\) 的變換量等於單位理想。於是，根據輔助引理 1，存在一個理想 \(\mathfrak b\)，使得 \(\mathfrak d_S=\mathfrak b^{1-S}\)。這就是說，類 \(\bar{\mathfrak c}_S\) 是類 \(\bar{\mathfrak b}\) 在符號意義下的 \((1-S)\) 次冪。

在循環情形中，主屬定理化為熟知的定理；這是因為，在採用規範化表示時，因子系誘導類劃分中的主類化為這樣一組主理想 \((\alpha)\)：\(\alpha\) 在所有分歧位上都是範數剩餘。

\begin{center}
（1932 年 10 月 27 日收稿。）
\end{center}
\endgroup

%
\clearpage
%
%
\begingroup
\newcommand{\srcnumdisplay}[3][3.4em]{%
  \par\smallskip\noindent
  \begin{minipage}[t]{#1}\vspace*{0pt}#2\end{minipage}%
  \begin{minipage}[t]{\dimexpr\linewidth-#1\relax}\vspace*{0pt}\centering
  \(\displaystyle #3\)
  \end{minipage}\par\smallskip}
\setcounter{footnote}{0}

\pdfbookmark[0]{42. 分裂交叉積及其極大階}{paper-42}
\section*{42. 分裂交叉積及其極大階}

\begin{center}
\emph{Actualités scientifiques et industrielles 148 (1934), S. 5--15}
\end{center}

在分裂交叉積最簡單的情形中\footnote{交叉積理論以我的一次課程為基礎，載於 H. Hasse, \emph{Theory of cyclic algebras over an algebraic numberfield}, Transact. Amer. Math. Soc. 34 (1932) 第 II 節。這裡我只使用最初的基本概念。Brandt 的極大階理論，連同更廣泛的結果，在 H. Hasse, \emph{Ueber \(p\)-adische Schiefkörper und ihre Bedeutung für die Arithmetik hyperkomplexer Zahlsysteme}, Math. Ann. 104 (1931) 中有系統敘述並得到詳盡的重新奠基（下文引作 H.）。}——也就是在產生分裂代數、因而具有與單位因子系相伴之因子系的那些交叉積中——各種關係在很大程度上都可以顯式追蹤；即使分裂域（極大交換子域）不是伽羅瓦的，這一點仍然成立。與此有關的簡單代數事實在 § 1 中展開。在 § 2 中，當基域為代數數域時，我藉助基礎分裂域 \(k\) 的模和互補模，給出所有極大階及其理想的顯式表示\footnote{這些顯式表示我早已知道；後來繼續作這些考察的契機，是 H. Hasse 告知我注 (4) 所引的定理。}。我還給出一種“區域”劃分：凡與 \(k\) 有相同交集的極大階，都歸入同一個區域。這樣，諸區域便與 \(k\) 中的諸階一一對應，而主階對應於“主區域”。同一區域中的極大階，通過由相應階的模構成的理想彼此變換；特別地，在主區域中，所涉及的是 \(k\) 中理想的擴張（即關於主階的模的擴張）。在 § 3 中，我轉向任意簡單代數。把區域劃分加以精細化——除了交集相同外，還須要求極大階在分歧位處的 \(p\)-進分量相同——便可以把分裂情形中所得的定理移植過來；特別地，此時一個主區域中的極大階仍通過 \(k\) 中理想的擴張彼此變換。

關於主區域的定理可見 Chevalley\footnote{\emph{Sur certains idéaux d'une algèbre simple}, Abh. math. Sem. Hambg., 1934。} 和 Hasse\footnote{\emph{Ueber gewisse Ideale in einer einfachen Algebra}，本文前一篇；本輯第 I 篇。}；Chevalley 還說明，如果沒有這種涉及分歧位的精細劃分，定理一般就不再成立。

\subsection*{§ I.}
\begin{center}
\textbf{分裂交叉積的矩陣單位與分裂域的互補基}
\end{center}

\noindent\textbf{1.——} 本節中，\(\Omega\) 表示任意基域。設 \(k/\Omega\) 是 \(n\) 次可分伽羅瓦擴張，\(\Gg\) 是含有元素 \(\mathrm{S},\mathrm{T},\ldots\) 的伽羅瓦群，\(u_s,u_t,\ldots\) 是相應的算子。我們考察\emph{因子係為一的交叉積}，也就是分裂代數
\[
K=\Gg\ast k=u_{S_1}k+\ldots+u_{S_n}k,
\qquad u_Su_T=u_{ST},
\]
\[
zu_s=u_sz^s\quad\text{（對每個 \(k\) 中的 \(z\)）}.
\]
因子係為一還給出進一步的關係。令
\[
E_\Gg=\sum_Su_S,
\]
則有
\srcnumdisplay{(1)}{\mathrm{E}_{\mathfrak g}u_s=u_s\mathrm{E}_{\mathfrak g}=\mathrm{E}_{\mathfrak g}}
對每個 \(\Gg\) 中的 \(\mathrm{S}\) 都成立（也就是說，\(\mathrm{E}_{\mathfrak g}\) 生成 \(\Gg\) 的單位表示\footnote{若 \(k\) 的特徵不整除 \(n\)，則 \(\frac1n\mathrm{E}_{\mathfrak g}\) 是單位表示的冪等元。不過，即使特徵整除 \(n\)，\(\mathrm{K}\) 的正規形式仍然成立；唯一的假設是 \(k/\Omega\) 可分。}）；並且
\srcnumdisplay{(2a)}{z\mathrm{E}_{\mathfrak g}=\sum_su_sz^s,}
\srcnumdisplay{(2b)}{E_\Gg zE_\Gg=E_\Gg\cdot\sum_Su_Sz^S
  =E_\Gg\cdot\Sp(z),}
其中 \(\Sp(z)\) 表示 \(z\) 作為 \(k/\Omega\) 的元素時的跡。由這些關係可得下述定理。

\medskip
\noindent\textbf{定理。——} \emph{若 \(a_1,\ldots,a_n\) 是 \(k/\Omega\) 的一個基，\(\widetilde a_1,\ldots,\widetilde a_n\) 是與之互補的基，則
\(c_{ik}=\widetilde a_i\mathrm{E}_{\mathfrak g} a_k\) 給出 \(\mathrm{K}\) 的一個矩陣單位系。因此，\(\mathrm{K}\) 可以寫成如下模乘積：}
\[
\mathrm{K}=k\mathrm{E}_{\mathfrak g}k=k\cdot\sum_su_s\cdot k
 =\sum_{i,k}(\widetilde a_i\mathrm{E}_{\mathfrak g}a_k)\,\Omega.
\]
事實上，由 (2b) 得
\[
\begin{aligned}
c_{ij}c_{lk}
 &=\widetilde a_i\mathrm{E}_{\mathfrak g} a_j\cdot\widetilde a_l\mathrm{E}_{\mathfrak g} a_k\\
 &=\widetilde a_i\mathrm{E}_{\mathfrak g} a_k\cdot\Sp(a_j\widetilde a_l)\\
 &=c_{ik}\cdot\Sp(a_j\widetilde a_l).
\end{aligned}
\]
而根據互補基的定義式 \(\Sp(a_j\widetilde a_l)=\delta_{jl}\)，這恰好就是矩陣單位的性質。

\emph{注 1。} 若 \(\mathrm{Z}\) 是 \(\Omega\) 的任意交換擴域，那麼當 \(a_i\) 表示 \(k_{\mathrm Z}/\mathrm{Z}\) 的一個基、而 \(k_{\mathrm Z}\) 成為諸域的直和時，上述一切仍然成立。這裡群的代換被定義為在 \(\mathrm{Z}\) 上誘導恆等映射；這與 \(u_s\) 在 \(\mathrm{K}_{\mathrm Z}\) 中同 \(\mathrm{Z}\) 的每個元素可交換相一致。

\emph{注 2。} 等式 \(\mathrm{K}=k\mathrm{E}_{\mathfrak g}k\) 也可直接從右端是簡單系 \(\mathrm{K}\) 的一個非零理想這一事實輕易推出。

\medskip
\noindent\textbf{2.——} 現在假定基域 \(\Omega\) 的\emph{特徵為零}。\emph{即使 \(k\) 是 \(\mathrm{K}/\Omega\) 的非伽羅瓦分裂域，1 中給出的分裂代數表示 \(\mathrm{K}=k\mathrm{E}_{\mathfrak g}k\) 仍然有效。} 這可以通過過渡到伽羅瓦分裂域來證明：設 \(\bar k\) 是與 \(k\) 相應的伽羅瓦域，\(\bar\Gg\) 是它的群，\(\bar u_s\) 是相應的算子，而 \(\bar{\mathrm K}=\bar k\mathrm{E}_{\bar{\mathfrak g}}\bar k\) 是一個分裂交叉積。子域 \(k\) 對應於階為 \(\mathrm{h}\) 的子群 \(\Hh\)。令
\[
\bar\Gg=\Hh\mathrm{S}_1+\ldots+\Hh\mathrm{S}_n
 =\mathrm{S}_1^{-1}\Hh+\ldots+\mathrm{S}_n^{-1}\Hh;
\qquad
\mathrm{E}_{\mathfrak h}=\frac1h\sum\bar u_{\mathfrak h},
\]
\[
\mathrm{E}_{\mathfrak g}=\frac1h\mathrm{E}_{\bar{\mathfrak g}}=\sum_i u_{s_i},
\qquad\text{其中}\qquad
u_{s_i}=\mathrm{E}_{\mathfrak h}\bar u_{s_i}.
\]
於是，冪等元 \(\mathrm{E}_{\mathfrak h}\) 尤其生成 \(\Hh\) 的單位表示，而 \(\mathrm{E}_{\mathfrak g}\) 同 \(\mathrm{E}_{\bar{\mathfrak g}}\) 一樣生成 \(\bar\Gg\) 的單位表示。\emph{對於這樣定義的 \(\mathrm{E}_{\mathfrak g}\) 和 \(u_s\)，關係 (1) 在一側成立，而關係 (2) 完全成立。} 這裡 \(z\) 遍歷 \(k\) 的所有元素，因而 \(z^s\) 遍歷共軛域 \(k^s\) 的元素；作為 \(\bar k\) 的子域，\(k^s\) 也是 \(\bar{\mathrm K}\) 的子域。首先有
\srcnumdisplay{(3)}{\mathrm{E}_{\mathfrak g}=\mathrm{E}_{\mathfrak g}\mathrm{E}_{\mathfrak h}=\mathrm{E}_{\mathfrak h}\mathrm{E}_{\mathfrak g}=\mathrm{E}_{\mathfrak h}\mathrm{E}_{\mathfrak g}\mathrm{E}_{\mathfrak h},}
\srcnumdisplay{(4)}{z\mathrm{E}_{\mathfrak h}=\mathrm{E}_{\mathfrak h}z,\quad\text{並因而}}
\srcnumdisplay{(5)}{zu_s=u_sz^s.}
事實上，由 \(\mathrm{E}_{\bar{\mathfrak g}}\) 的定義和關係 (1)，(3) 對 \(\mathrm{E}_{\bar{\mathfrak g}}\) 成立，因而也對 \(\mathrm{E}_{\mathfrak g}\) 成立；而 (4) 所說的是，作為 \(k\) 中的元素，\(z\) 容許 \(\Hh\) 的各個代換，即 \(zu_h=u_hz\)。把 (4) 右乘 \(\bar u_{s_i}\)，便得到 (5)。現在得到關係 (1) 的一側形式：
\srcnumdisplay{(1')}{E_\Gg u_{S_i}=E_\Gg,
\quad\text{因為 }E_\Gg u_{S_i}=E_\Gg E_\Hh\bar u_{S_i}
=E_\Gg\bar u_{S_i}=E_\Gg;}
對 (5) 求和便得到 (2a)，再由 (1') 得到 (2b)。

現在，\(\mathrm{K}=k\mathrm{E}_{\mathfrak g}k\) 的證明與 1 中完全一樣：要證明 \(\widetilde a_i\mathrm{E}_{\mathfrak g}a_k\) 是矩陣單位，只使用了 (2b)。作為 \(\Omega\) 上的全矩陣環，\(\mathrm{K}\) 含有每一個 \(n\) 次域作為子域，特別也含有 \(k\)。

\medskip
\noindent\textbf{3.——} 為了 § 3，我還要指出，即使 \(\mathrm{K}\) 不分裂，也可以從 \(K\) 過渡到 \(\bar{\mathrm K}\)\footnote{由此可以得到非伽羅瓦情形下與交叉積表示相似的表示；這大概還將在一篇學位論文中加以闡述。}。因為除 \(k\) 外，\(\bar k\) 也是分裂域，所以 \(\mathrm{K}\) 與由 \(\bar\Gg\) 和 \(\bar k\) 構成的交叉積 \(\bar{\mathrm K}\) 相似。又因為 \(k\) 是分裂域，與 \(k\) 對應的子群 \(\Hh\) 所對應的因子係為一；所以冪等元 \(\mathrm{E}_{\mathfrak h}\) 如 2 中定義。反過來，由此得到 \(\mathrm{K}\) 同構於 \(\mathrm{E}_{\mathfrak h}\bar{\mathrm K}\mathrm{E}_{\mathfrak h}\)。事實上，這個代數同左理想 \(\bar{\mathrm K}\mathrm{E}_{\mathfrak h}\) 的自同態環同構，因而與 \(\bar{\mathrm K}\) 相似。它的秩與 \(\mathrm{K}\) 的秩一致；因為 \(\bar{\mathrm K}\mathrm{E}_{\mathfrak h}\) 在 \(\bar k\) 上的秩為 \(n\)（\(\bar u_{s_i}^{-1}\mathrm{E}_{\mathfrak h}\) 在 \(\bar k\) 上線性無關），而 \(\bar{\mathrm K}\) 的秩為 \(nh\)。這就是說，\(\bar{\mathrm K}\) 分解為 \(h\) 個與 \(\bar{\mathrm K}\mathrm{E}_{\mathfrak h}\) 同構的理想；對自同態環而言，這意味着再乘上一個次數為 \(h\) 的全矩陣代數。

當 \(\mathrm{K}\)，因而 \(\bar{\mathrm K}\)，分裂時，這樣又得到 2 的結果；因為
\[
\begin{aligned}
\mathrm{E}_{\mathfrak h}\bar{\mathrm K}\mathrm{E}_{\mathfrak h}
 =\mathrm{E}_{\mathfrak h}\bar k\mathrm{E}_{\bar{\mathfrak g}}\bar k\mathrm{E}_{\mathfrak h}
 =(\mathrm{E}_{\mathfrak h}\bar k\mathrm{E}_{\mathfrak h})\mathrm{E}_{\bar{\mathfrak g}}(\mathrm{E}_{\mathfrak h}\bar k\mathrm{E}_{\mathfrak h})
 &=\Sp(\bar k/k)\mathrm{E}_{\mathfrak h}\mathrm{E}_{\bar{\mathfrak g}}\mathrm{E}_{\mathfrak h}\Sp(\bar k/k)\\
 &=k\mathrm{E}_{\mathfrak g}k.
\end{aligned}
\]

\subsection*{§ II.}
\begin{center}
\textbf{分裂交叉積的極大階及其區域劃分}
\end{center}

從現在起，設 \(\Omega\) 是一個（有限）代數數域\footnote{當 \(\Omega\) 是一個單變量函數域時，本節的考察仍然成立。}，\(\mathrm{K}\) 是 \(\Omega\) 上次數為 \(n\) 的分裂代數，而 \(k\) 是一個伽羅瓦或非伽羅瓦的極大交換子域；因而根據 § 1，\(\mathrm{K}=k\mathrm{E}_{\mathfrak g}k\)。再設 \(\mathfrak a,\widetilde{\mathfrak a}\) 是 \(k\) 中一對互補的有限 \(\oo\)-模（\(o\) 是 \(\Omega\) 的主階，\(\oo\) 是 \(k\) 的主階）。現在要考察 \(\mathrm{K}\) 中相對於 \(k\) 的極大階的結構；這一結構由下列定理描述。

\medskip
\noindent\textbf{定理 1。——} \emph{模乘積 \(\mathfrak O_{\mathfrak a}=\widetilde{\mathfrak a}\mathrm{E}_{\mathfrak g}\mathfrak a\) 是 \(\mathrm{K}\) 中的一個極大階；特別地，\(\mathfrak O=\mathfrak O_{\mathfrak o}=\widetilde{\mathfrak o}\mathrm{E}_{\mathfrak g}\mathfrak o\) 是極大階。}

\medskip
\noindent\textbf{定理 2。——} \emph{從 \(\mathfrak O\) 到 \(\mathfrak O_{\mathfrak a}\) 的變換，由相對於 \(\mathfrak O\) 的左理想 \(\mathfrak L=\widetilde{\mathfrak o}\mathrm{E}_{\mathfrak g}\mathfrak a\)，以及與它互逆的右理想 \(\mathfrak R=\widetilde{\mathfrak a}\mathrm{E}_{\mathfrak g}\mathfrak o\) 實現。}

\emph{注。} 在乘積 \(\mathfrak L\mathfrak R\) 中，兩個互補模 \(\mathfrak a,\widetilde{\mathfrak a}\) 的相遇確實導致互逆關係；這與單純的乘積 \(\mathfrak a\widetilde{\mathfrak a}\) 不同，是由 \(\mathrm{E}_{\mathfrak g}\) 所引起的取跡作用造成的。

\medskip
\noindent\textbf{定理 3。——} \emph{相對於 \(\mathfrak O\) 的所有左理想和右理想，都由定理 2 中給出的理想窮盡；因此，諸極大階 \(\mathfrak O_{\mathfrak a}\) 窮盡 \(\mathrm{K}\) 中的全部極大階。任意兩個極大階 \(\mathfrak O_{\mathfrak a}\) 和 \(\mathfrak O_{\mathfrak b}\)，通過 \(\mathfrak L_{\mathfrak a}=\widetilde{\mathfrak a}\mathrm{E}_{\mathfrak g}\mathfrak b\)、\(\mathfrak R_{\mathfrak a}=\widetilde{\mathfrak b}\mathrm{E}_{\mathfrak g}\mathfrak a\) 相互變換；\(\mathfrak L_{\mathfrak a}\) 和 \(\mathfrak R_{\mathfrak a}\) 窮盡相對於 \(\mathfrak O_{\mathfrak a}\) 的左理想和右理想。}

\medskip
\noindent\textbf{定理 4。——} \emph{當 \(\mathfrak O_{\mathfrak a}\) 遍歷 \(\mathrm{K}\) 中的全部極大階時，交集 \([\mathfrak O_{\mathfrak a},k]\) 遍歷 \(k\) 中的全部階；每一次，這個交集都等於 \(\mathfrak a\) 的階。把屬於 \(k\) 中同一個階的全部極大階合併為一個區域後，一個區域內部的變換由相對於 \(\mathfrak O_{\mathfrak a}\) 的左理想 \(\widetilde{\mathfrak a}\mathrm{E}_{\mathfrak g}\mathfrak b\) 及其互逆右理想實現，其中 \(\mathfrak a\) 和 \(\mathfrak b\) 是 \(k\) 中相應階的模。}

\medskip
\noindent\textbf{定理 4a。——} \emph{特別地，主區域中的極大階——也就是含有 \(k\) 的主階 \(o\) 的全部極大階——通過理想 \(\widetilde{\mathfrak a}\mathrm{E}_{\mathfrak g}\mathfrak b\) 相互變換，其中 \(\mathfrak a\) 和 \(\mathfrak b\) 是 \(k\) 中的理想。換句話說，若 \(\mathfrak L\) 是主區域中一個極大階的左理想，並且 \(\mathfrak L\) 的右階含有 \(o\)，那麼 \(\mathfrak L\) 是一個 \(\mathfrak o\)-理想的擴張。}

證明基於向各個位的過渡，而且每次只須過渡到 \(p\) 處的局部化環，其中 \(p\) 遍歷 \(\Omega\) 中的全部素理想。令 \(o\) 表示 \(\Omega\) 的主階，\(o_p\) 表示 \(p\) 處的局部化環（原文稱 Quotientenring），也就是所有 \(p\)-整元素——即能用與 \(p\) 互素的分母表示的 \(\Omega\) 中元素——組成的主理想環。若 \(\mathfrak C\) 是 \(\mathrm{K}\) 中的任意 \(o\)-模，則稱擴張模 \(\mathfrak C_p=\mathfrak C\oo_p\) 為它的分量。於是有下述輔助引理。

\medskip
\noindent\emph{輔助引理。\(\mathfrak C\) 是全部擴張模 \(\mathfrak C_p\) 的交。} 事實上，若元素 \(w\) 屬於 \(\mathfrak C_p\)，則存在一個不被 \(p\) 整除的 \(\sigma\in\oo\)，使得 \(\sigma w\in\mathfrak C\)。現在設 \(w\) 屬於全部 \(\mathfrak C_p\) 的交，並令 \(g\) 為所有滿足 \(\sigma w\in\mathfrak C\) 的 \(\sigma\in\oo\) 所成的理想。由剛才的論證，\(g\) 不可能被任何 \(p\) 整除，因而 \(g=o\)\footnote{若通過 \(p\)-進擴張來定義各個位，相應的輔助引理表述為：所有 \(p\)-進分量 \(\mathfrak C^{(p)}\) 與 \(\mathrm{K}\) 的交等於 \(\mathfrak C\)。這是因為 \([\mathfrak C^{(p)},\mathrm{K}]=\mathfrak C_p\)；用 \(\mathfrak C_p\) 的一個 \(o_p\)-基即可看出這一點，該基同時也是 \(\mathfrak C^{(p)}\) 的基。［在 H. 定理 66 中，輔助引理只對（最高秩的）理想作了表述。］}。應當注意，這個輔助引理對模 \(\mathfrak C\) 的秩沒有任何假設。

\emph{推論。} 因而，每個模都由它的各個分量確定。若 \(\mathfrak D\) 是 \(\mathfrak C\) 的真上模，則至少有一個 \(\mathfrak D_p\) 是 \(\mathfrak C_p\) 的真上模。特別地，若 \(\mathfrak M\) 是 \(\mathrm{K}\) 中的一個階，而所有分量 \(\mathfrak M_p\) 都是相對於 \(o_p\) 的極大階，那麼 \(\mathfrak M\) 本身也是極大階。

\medskip
\noindent\emph{定理 1 的證明。} \(\mathfrak O_{\mathfrak a}\) 是一個階，因為 \(a\) 和 \(\widetilde{\mathfrak a}\) 都是有限 \(\oo\)-模，所以 \(\mathfrak O_{\mathfrak a}\) 也是有限 \(\oo\)-模；而且 \(\mathfrak O_{\mathfrak a}\) 具有最高秩。進一步，
\[
\mathfrak O_a^2<\mathfrak O_a;
\quad\text{因為}\quad
\mathfrak O_a^2
 =\widetilde aE_\Gg a\cdot\widetilde aE_\Gg a
 =\widetilde aE_\Gg a\cdot\Sp(\widetilde a a)\footnotemark.
\]
\footnotetext{\(\Sp(\mathfrak c)\) 一般表示由 \(\mathfrak c\) 的全部元素的跡組成的 \(o\) 中的理想。}根據 § 1，並且由於 \(\Sp(\widetilde{\mathfrak a}\mathfrak a)<o\)，便有 \(\mathfrak O_{\mathfrak a}^2<\mathfrak O_{\mathfrak a}\)。

為了證明 \(\mathfrak O_{\mathfrak a}\) 是極大的（從而包含 \(o\)），根據上述推論，只須證明每個分量都是極大的。由於 \(o_p\) 是主理想環，\(\mathfrak a_p\) 和 \(\widetilde{\mathfrak a}_p\) 分別具有相互獨立的互補基 \(a_1,\ldots,a_n\) 和 \(\widetilde a_1,\ldots,\widetilde a_n\)。根據 § 1，乘積 \(\widetilde a_i\mathrm{E}_{\mathfrak g}a_k\) 構成一個矩陣單位系；因此
\[
\widetilde a_pE_\Gg a_p=\sum_{i,k}c_{ik}o_p
\]
是極大階。

\emph{注。} 由於 \(\mathfrak O_{\mathfrak a}\) 是極大階，所以 \(\mathfrak O_{\mathfrak a}^2=\mathfrak O_{\mathfrak a}\)，還可推出 \(\Sp(\widetilde{\mathfrak a}\mathfrak a)=o\)。這一點也可直接推出：互補基表明該關係在每個位都成立，因而由輔助引理可知它無條件成立。

\medskip
\noindent\emph{定理 2 的證明。} 由剛剛確立的取跡規則可得
\[
\begin{aligned}
\mathfrak O\mathfrak L
 &=\widetilde{\mathfrak o}\mathrm{E}_{\mathfrak g}\mathfrak o\,\widetilde{\mathfrak o}\mathrm{E}_{\mathfrak g}\mathfrak a
 =\widetilde{\mathfrak o}\mathrm{E}_{\mathfrak g}\mathfrak a=\mathfrak L,
&\qquad
\mathfrak R\mathfrak O
 &=\widetilde{\mathfrak a}\mathrm{E}_{\mathfrak g}\mathfrak o\,\widetilde{\mathfrak o}\mathrm{E}_{\mathfrak g}\mathfrak o
 =\widetilde{\mathfrak a}\mathrm{E}_{\mathfrak g}\mathfrak o=\mathfrak R,\\
\mathfrak L\mathfrak R
 &=\widetilde{\mathfrak o}\mathrm{E}_{\mathfrak g}\mathfrak a\,\widetilde{\mathfrak a}\mathrm{E}_{\mathfrak g}\mathfrak o
 =\widetilde{\mathfrak o}\mathrm{E}_{\mathfrak g}\mathfrak o=\mathfrak O,
&
\mathfrak R\mathfrak L
 &=\widetilde{\mathfrak a}\mathrm{E}_{\mathfrak g}\mathfrak o\,\widetilde{\mathfrak o}\mathrm{E}_{\mathfrak g}\mathfrak a
 =\widetilde{\mathfrak a}\mathrm{E}_{\mathfrak g}\mathfrak a=\mathfrak O\,.
\end{aligned}
\]

\medskip
\noindent\emph{定理 3 的證明。} 設 \(\mathfrak L\) 是相對於 \(\mathfrak O\) 的一個左理想，則特別有
\[
\mathfrak L=\widetilde oE_\Gg\mathfrak L
 =(\widetilde oE_\Gg l_1,\ldots,\widetilde oE_\Gg l_t),
\]
其中 \(l_1,\ldots,l_t\) 例如是 \(\mathfrak L\) 的一個 \(o\)-基。令 \(l_\mu=\sum\limits_i u_{s_i}z_{\mu,i}\)，其中 \(z_{\mu,i}\in k\)\footnote{當 \(k\) 不是伽羅瓦的時，必須把 \(z\) 換成 \(\bar z\mathrm{E}_{\mathfrak h}\)，因而把 \(a\) 換成 \(\bar a\mathrm{E}_{\mathfrak h}\)；於是得到 \(E_\Gg l=E_\Gg\bar aE_\Hh=E_\Gg a\)，其中 \(a\in k\)。事實上，由 (3)、(4)、(5) 得 \(kE_\Gg k=kE_\Gg kE_\Hh=\sum_i u_{S_i}\{k^{S_i},k\}E_\Hh\)。}，則
\[
E_\Gg l_\mu=E_\Gg\sum_i z_{\mu,i}=E_\Gg a_\mu,
\]
所以 \(\mathfrak L=\widetilde{\mathfrak o}\mathrm{E}_{\mathfrak g}\mathfrak a\)。當 \(\mathfrak L\) 是正則理想時，\(a\) 必須具有最高秩 \(n\)，於是 \(\widetilde aE_\Gg o\) 是與之互逆的右理想。通過合成可知，\(\widetilde{\mathfrak a}\mathrm{E}_{\mathfrak g}\mathfrak b\) 遍歷相對於 \(\mathfrak O_{\mathfrak a}\) 的全部左理想，而 \(\widetilde bE_\Gg a\) 遍歷與它們互逆的右理想。

\medskip
\noindent\emph{定理 4 的證明。} 交集 \(t=[k,\mathfrak O_a]\) 是一個階，因為 \(\mathfrak O_{\mathfrak a}\) 只含整元素。作為 \(\mathfrak O_{\mathfrak a}\) 的子集，\(t\) 是一個右乘子環，而且是 \(o\) 中滿足這一條件的最大階；所以它至少包含 \(a\) 的階 \(o_a\)。為了證明 \(t=o_a\)，仍只須證明各個分量相等。設 \(x\in t_p\)；那麼，對每一個 \(c_{ik}=\widetilde a_i\mathrm{E}_{\mathfrak g} a_k\)，\(c_{ik}x\) 都是係數為整元素（即屬於 \(o_p\)）的諸 \(c_{ik}\) 的線性組合。另一方面，\(a_kx\) 是係數屬於 \(\Omega_p\) 的諸 \(a_k\) 的線性組合。由於通過 \(c_{ik}\) 的表示是唯一的，這些係數必定屬於 \(o_p\)；也就是說，\(x\) 屬於 \(\mathfrak a_p\) 的階。因而，\(t\) 等於 \(a\) 的階。每個 \(k\) 中的階都可作為這種交集出現，是因為 \(k\) 中的模 \(a\) 尤其也遍歷所有階。現在，關於變換的斷言直接由定理 3 得出。

\medskip
\noindent\emph{定理 4a 的證明。} 這個定理只是定理 4 對主區域的特殊化。先考慮出發極大階 \(\mathfrak O=\widetilde oE_\Gg o\)。若 \(o\) 也包含在 \(\mathfrak L=\widetilde{\mathfrak o}\mathrm{E}_{\mathfrak g}\mathfrak a\) 的右階中，則 \(\widetilde aE_\Gg a\) 也屬於主區域，因而 \(a\) 是 \(k\) 中的理想，並且 \(\mathfrak L=\mathfrak Oa\)。若從 \(\mathfrak O_{\mathfrak a}\) 出發，則相應地有 \(\mathfrak L=\widetilde aE_\Gg b\)，而 \(b\)，因而 \(a^{-1}b\)，是 \(k\) 中的理想；於是
\(\mathfrak L=\widetilde aE_\Gg aa^{-1}b=\mathfrak O_a a^{-1}b\)。

\subsection*{§ III.}
\begin{center}
\textbf{任意交叉積的極大階}
\end{center}

只要把區域劃分加以精細化，為分裂交叉積推導出的定理在一般情形中也仍然成立。關於主區域的事實這時原樣成立；其餘事實則必須表述為關於各個位的命題。這裡所謂位，現在應理解為 \(p\)-進擴張。

\emph{定義。} 設 \(K\) 是 \(\Omega\) 上的中心單代數，\(k\) 是一個極大交換子域。相對於 \(k\) 的一個區域，是由 \(K\) 中所有滿足下列條件的極大階組成的：它們與 \(k\) 有相同的交集，並且在 \(K\) 的分歧位處具有相同的 \(p\)-進分量。若與 \(k\) 的交集是主階，則稱為主區域。

對於主區域，有下列定理。

\medskip
\noindent\textbf{定理 1。——} \emph{一個主區域中的極大階，通過以 \(k\) 中理想的擴張理想作變換而彼此變換。換句話說，主區域中極大階的那些與不同理想互素、並且其右階含有 \(o\) 的左理想，都是 \(k\) 中理想的擴張。}

\medskip
\noindent\textbf{定理 2。——} \emph{若 \(K\) 的次數為素數，則只有一個主區域；這個區域通過 \(k\) 中理想的擴張理想變換到自身。相對於一個固定極大階 \(\mathfrak O\) 的相應理想，是把 \(k\) 中的理想同 \(\mathfrak O\) 中的雙邊理想一起擴張所得的理想。}

\emph{證明。} 同一個主區域中的兩個極大階只在有限多個位處彼此不同；根據定義，這些都是 \(K\) 分裂的位。在這些位處，\(K\) 具有 § 1 和 § 2 中所考察的形式；當 \(k\) 是伽羅瓦的時，這一點很清楚，因為此時因子系與單位因子系相伴，因而可以用單位因子系來替換；根據 § 1 的 2 和 3，這一點在非伽羅瓦情形中也隨之成立。因此，在每一個這樣的位處，因而整體上\footnote{參見 H. § 8，或 Hasse 的前一篇論文，即注 (4)。}，所涉及的都是擴張理想。

若 \(K\) 的次數為素數，則它或者分裂（§ 2），或者是除代數；在後一種情形中，它在全部分歧位處仍然是除代數，因而只有一個主區域，所以仍由擴張理想實現變換。最一般的變換，是把擴張理想同那些把極大階變換到自身的理想合成，也就是同雙邊理想合成；眾所周知，雙邊理想只在分歧位處才與中心理想、因而與擴張理想不同。由此得到其餘斷言。

此外還有：

\medskip
\noindent\textbf{定理 3。——} \emph{任意區域中的極大階，通過與不同理想互素的理想彼此變換；這些理想在每個位處都由相應階的模按 § 2 定理 3 所述的方式組成。對於一個固定理想，這些模的分量只在有限多個位處不同於該階。}

這直接由 § 2 得出；只須注意，在各個位處，算子 \(u_s\) 必須換成相應的等價算子。

至於對 \(k\) 中的每一個階，這裡是否也存在以該階為交集的極大階，則需要在分歧位處作更細緻的考察。主區域總是存在的；這一點可從交叉積表示——必要時是廣義交叉積表示——看出：只要因子系被取為整的，它就導出一個含有 \(o\) 的階\footnote{另一個證明見 Hasse 和 Chevalley，即注 (4) 和 (3)。}。

另一個問題是，算術上的區域劃分在多大程度上可以刻畫代數分裂域。換句話說：不同的分裂域是否總會產生不同的區域劃分？它們是否已經產生不同的主區域？或者，全部分裂域的主區域是否就已經窮盡全部極大階？

僅憑一個極大階肯定不能作這種刻畫，最簡單的例子已經說明這一點。例如，四元數體的極大階
\[
\left[1,i,j,\frac{1+i+j+k}{2}\right]
\]
除含有主階 \([1,i]\) 外，還含有三次單位根域的主階
\[
\left[1,\frac{1+i+j+k}{2}\right].
\]

\begin{flushleft}
哥廷根，1932 年 8 月。
\end{flushleft}
\endgroup

%
\clearpage
%
%
\begingroup
\def\srcfn#1#2{\begingroup\renewcommand{\thefootnote}{#1}\footnote{#2}\addtocounter{footnote}{-1}\endgroup}
\def\srcfnmark#1{\begingroup\renewcommand{\thefootnote}{#1}\footnotemark\addtocounter{footnote}{-1}\endgroup}
\def\srcfntext#1#2{\begingroup\renewcommand{\thefootnote}{#1}\footnotetext{#2}\endgroup}
\setcounter{footnote}{0}

\pdfbookmark[0]{43. 理想微分與不同}{paper-43}
\section*{43. 理想微分與不同。}

\begin{center}
\emph{Journal f. d. reine u. angew. Math. 188 (1950), S. 1--21}\\[0.75ex]
\emph{Emmy Noether}（\(\dagger\)）著\srcfnmark{1)}。
\end{center}

\srcfntext{1)}{本文是 E. Noether 的遺稿，寫於 1927/28 年冬，由 H. Grell 惠交編者。文中詳盡闡述了作者 1929 年在布拉格德國數學家協會（D. M. V.）大會上所作的同名報告；見 Jahresbericht D. M. V. 39 (1930), S. 17 的斜體部分。從 § 6, 3 起，手稿顯然原擬再作更細緻的修訂。不過，內行讀者不難理解這裡保留下來的提綱式論述；除少量潤飾外，原稿未作改動。}

\subsection*{引言。}

眾所周知，代數數域分歧理論的主定理斷言：設一個素理想以第 \(\varrho\) 次冪出現在它所位於的素數 \(p\) 中，則不同（即分歧理想）至少能被該素理想的第 \((\varrho-1)\) 次冪整除；更確切地說，當 \(\varrho\) 不能被 \(p\) 整除時，恰能被第 \((\varrho-1)\) 次冪整除，而當 \(\varrho\) 能被 \(p\) 整除時，則能被更高次冪整除。

這個定理類似於如下事實：若一個線性因子以第 \(\varrho\) 次冪出現在多項式 \(f(x)\) 中，則 \(f(x)\) 的微分商 \(f'(x)\) 至少能被該線性因子的第 \((\varrho-1)\) 次冪整除；更確切地說，當 \(\varrho\) 不能被係數環的特徵整除時，恰能被第 \((\varrho-1)\) 次冪整除，而當 \(\varrho\) 能被該特徵整除時，則能被更高次冪整除。

下面我將說明，這裡不只是形式上的類比。\emph{不同可以看作數域 \(\mathsf{K}\) 的一個定義理想在一個相應的、以 \(x_1,\ldots,x_n\) 為不定元的整係數多項式環中的微分商；該微分商取在 \(x=\omega\) 處，也就是取 \(x_1=\omega_1,\ldots,x_n=\omega_n\)，其中 \(\omega_1,\ldots,\omega_n\) 表示 \(\mathsf{K}\) 中全體整數所成系統---主階 \(\mathfrak o\)---的一組模基。}這裡，定義理想 \(\mathfrak M\) 對應於各 \(\omega\) 之間的全體關係，因而由所有滿足 \(f(\omega)=0\) 的整係數多項式 \(f(x)\) 組成。微分商 \(\mathfrak M'[x\to\omega]\) 則定義為取在 \(x=\omega\) 處的差分商。

事實上，由於 \(f(\omega)=0\)，可將理想 \(\mathfrak M\) 看成由所有差 \(f(x)-f(\omega)\) 組成；再構造差分理想 \(\mathfrak P=(x_1-\omega_1,\ldots,x_n-\omega_n)\) 以及差分商 \(\mathfrak A=\mathfrak M:\mathfrak P\)。這裡取商是指通常的理想商，並且是在以 \(\mathsf{K}\) 中的整數、亦即 \(\mathfrak o\) 中的數為係數的 \(x\) 的多項式環內進行的。於是，不同 \(\mathfrak D\) 定義為
\[
\mathfrak D=\mathfrak A[x\to\omega].
\]
係數環擴張是必要的，否則差分理想和差分商根本無法定義。因此，這正是單不定元 \(x\) 的多項式微分商
\[
f'(\xi)=(f(x)-f(\xi)):(x-\xi),
\]
在 \(x=\xi\) 處取值的直接推廣；這裡也必須添入 \(\xi\) 來擴張多項式環，才能使各差有定義。事實上，只要各 \(\omega\) 的基由單個元素的各次冪組成，理想微分商也就過渡為多項式微分商。

這裡給出的不同定義因而直接銜接於已知概念，但它引入了非不變的中間對象，因為理想 \(\mathfrak M\) 依賴於所選的基。若不用整係數多項式環，而引入一個與主階 \(\mathfrak o\) 同構的環 \(\mathfrak O\)---它將與模 \(\mathfrak M\) 的剩餘類環同構---並通過 \(\mathfrak o\) 把它的係數環擴張為 \(\mathfrak O_{\mathfrak o}\)，便可得到一個等價的、\emph{完全不變的定義}\srcfn{2)}{係數環擴張---即直積形成---在 § 1 中也針對非交換環加以處理，其一般性超過本文所需。對於本文的目的，§ 2 在交換環和有限基這一特殊化之下便已足夠。}。這裡，差分理想 \(\mathfrak P\) 由所有差 \((x-\xi)\) 導出，其中 \(x\) 和 \(\xi\) 在 \(\mathfrak O\) 與 \(\mathfrak o\) 的同構之下相互對應；差分商 \(\mathfrak A\) 等於零理想關於 \(\mathfrak P\) 的商；把 \(\mathfrak A\) 中的每個 \(x\) 都換成與它同構對應的 \(\xi\)，便得到不同。

下文將把這個不變定義置於首位，並由此直接定義任意交換環相對於一個子環的不同；這裡只須滿足關於係數擴張可能性的若干假設，例如在存在模基時這些假設總能滿足，必要時可先過渡到某些商環。這樣一來，特別是數域任意階的不同便有了定義；一個階模某個素數冪的剩餘類環的不同也同樣得到定義；此外還有相對不同，以及多不定元代數函數情形中的不同。不同的微分定義與通常定義的一致性，來自對數域通過直積形成而對應的伽羅瓦擴張環所作的一項精細結構考察；這項考察以直和分解為基礎\srcfn{3)}{這裡是對我關於判別式的論文所依據的思想的進一步發展，但不以讀者熟悉該文為前提。參見 E. Noether, Der Diskriminantensatz für die Ordnungen eines algebraischen Zahl- und Funktionenkörpers, J. Math. 157 (1927), 82--104.}。在單不定元多項式的剩餘類環這一特殊情形中，這項考察歸結為對 Lagrange 插值公式的分析。更確切地說，我的意思如下：

引入一個與數域 \(\mathsf{K}\) 同構並包含 \(\mathfrak O\) 的環 \(\mathfrak K\)，再用 \(K\) 的伽羅瓦域 \(F\) 擴張它的係數域---有理數。這個擴張後的系統 \(\mathfrak K_\Gamma\) 成為 \(n\) 個簡單分量的直和；它們對應於 \(K\) 的 \(n\) 個同構（即向各共軛域的過渡），並且分別是差分商 \(\mathfrak A\) 及其各共軛的擴張。同時，\(\mathfrak K_\Gamma\) 的零理想成為由差分理想及其各共軛導出的 \(n\) 個理想的交。各簡單分量具有 \(\Gamma e^{(i)}\) 的形式，其中 \(e^{(i)}\) 表示單位元的一個分量；由於 \(e^{(1)}\) 在 \(x=\xi\) 時過渡為單位元，差分商等於 \(\mathfrak d e^{(1)}\)，這裡 \(\mathfrak d\) 表示 \(\mathfrak o\) 的不同。不同 \(\mathfrak d\) 由體中所有滿足如下條件的元素組成：該元素乘以 \(e^{(1)}\) 後屬於 \(\mathfrak O_{\mathfrak o}\)，即該階與自身的直積。對各共軛也有相應的結論。

為了由此得到通過互補模給出的定義，應注意 \(e^{(i)}\) 成為關於 \(\mathfrak O\) 或 \(\mathfrak K\) 中與各 \(\omega\) 相對應的基 \(x_1,\ldots,x_n\) 的線性形式；而各 \(e^{(i)}\) 中 \(x\) 的係數，恰好給出 \(\mathfrak o\)（及其各共軛）的互補模的一組基。由差分商 \(\mathfrak A\) 的係數屬於 \(\mathfrak o\) 這一事實立即可知，不同等於 \(\mathfrak o\) 關於其互補模的商，從而得到 Dedekind 的定義。上述考察並不限於主階；它們把每個階的不同刻畫為該階關於其互補模的商。

在主階的情形中，微分定義與通過基本方程給出的定義的一致性，來自上面提到的事實：一旦基由單個元素---這裡即基本形式---的各次冪組成，理想微分商便過渡為多項式微分商。由此也就給出了後一種定義與 Dedekind 定義之間的概念聯繫。

在主階的情形中，以基本方程為出發點，眾所周知，可直接推出開頭提到的主定理。為了精確確定指數，需要考察主階模 \(p^t\) 的剩餘類環的不同 \(\mathfrak d[p^t]\)，其中 \(t\) 足夠大---取 \(t=n\) 即已足夠，而對於二次域這也是必要的---；關鍵在於，把不同 \(\mathfrak d\) 模 \(p^t\) 考察時，它恰好過渡為 \(\mathfrak d[p^t]\)。這裡仍可把考察限制於模 \(\mathfrak p^{t\varrho}\) 的剩餘類環，因為環的直和對應於不同的直和。後一個環又與某個單不定元多項式的剩餘類環同構，其係數模 \(p^t\) 取值；正如 Ö. Ore 已經證明的那樣---他從另一途徑得到這樣一個多項式\srcfn{4)}{編者注：手稿中沒有補全這一引文。參見本文僅具提綱性質的 § 7 中所列的提示。}---對這個多項式求微分，藉助補充數便能精確統觀一切可能的指數。

最後，這項結構考察還給出瞭如下著名定理：\(\mathsf{K}\) 的一個上域的不同，等於相對不同與 \(\mathsf{K}\) 的不同的乘積。這是因為同樣的乘法性質對直和分解中的各單位元成立，因而也對各互補模成立。

通過過渡到某些商環，以上所述的全部事實對相對不同也完全成立，正如我在關於判別式的論文中為相對判別式所闡明的那樣。再過渡到函數域之後，對於係數體特徵為零的代數函數體，一切結論仍然成立；而在特徵為 \(p\) 以及整代數函數的情形中，結構性質固然仍然保留（實質上甚至無須過渡到函數域），但分歧理論會因不可分擴張體的出現而發生變化，這裡不再展開。我要指出，最重要的結果是上面概述的、對多不定元代數函數也成立的完全不變結構定理（§ 5 和 § 6）；過渡到互補模已經意味着把論證分解為係數恆等式，這一步只是為了把已知結果納入其中。
% 原文：第43篇，印刷頁 4--6 （期刊掃描頁 693--695）
\subsection*{§ 1. 環的係數環擴張或直積形成。定義理想\srcfnmark{5)}}
\srcfntext{5)}{對於他對本節提出的批評意見，我謹向 B. L. van der Waerden 致謝。}

本節以一般環為基礎，並不預設乘法交換律。另一方面，為了簡便起見，假定單位元存在；除非明確註明相反情形，這一假設在全文中均予採用。

\paragraph{1. 直積的定義。}
如果滿足下列條件，環 \(\mathfrak O\times\mathfrak o\) 或 \(\mathfrak O_{\mathfrak o}\) 稱為環 \(\mathfrak O\) 與 \(\mathfrak o\) 關於 \(\mathfrak h\) 的\emph{直積}，也稱為\emph{由 \(\mathfrak O\) 通過把係數環 \(\mathfrak h\) 擴張為 \(\mathfrak o\) 而得到的環}\srcfn{6)}{下文大多采用第二種寫法，用以提示係數環的擴張；第一種則是直積的通常名稱。}：

\begin{enumerate}
\item[(1)] \(\mathfrak O_{\mathfrak o}\) 包含 \(\mathfrak O\) 和 \(\mathfrak o\) 作為子環，並且等於這兩個環的乘積（也就是說，它等於 \(\mathfrak O_{\mathfrak o}\) 中所有同時包含 \(\mathfrak O\) 和 \(\mathfrak o\) 的環的交）。

\item[(2)] \(\mathfrak O\) 與 \(\mathfrak o\) 的交等於 \(\mathfrak h\)，其中 \(\mathfrak h\) 包含 \(\mathfrak O_{\mathfrak o}\) 的單位元。

\item[(3)] \(\mathfrak O\) 與 \(\mathfrak o\) 的元素兩兩可交換；因此，\(\mathfrak O_{\mathfrak o}\) 由全體雙線性組合
\[
x_{i_1}y_{i_1}+\cdots+x_{i_r}y_{i_r}
\]
組成，其中各組 \(x\) 和 \(y\) 分別遍歷 \(\mathfrak O\) 和 \(\mathfrak o\) 中任意有限多個元素所成的全部系統；而 \(\mathfrak h\) 位於 \(\mathfrak O\) 和 \(\mathfrak o\) 的中心，因而也位於 \(\mathfrak O_{\mathfrak o}\) 的中心。

\item[(4)] \(\mathfrak O_{\mathfrak o}\) 中的兩個元素相等，當且僅當（當然，只要滿足這一條件就一定相等）它們至少能以一種方式，藉助一方面在 \(\mathfrak O\) 中、另一方面在 \(\mathfrak o\) 中成立的關係而形式地變為相等。其含義如下：若
\[
x_1y_1+\cdots+x_ny_n=0
\quad\hbox{在 }\mathfrak O_{\mathfrak o}\hbox{ 中},
\]
則可添入若干形式消失的組合 \((\gamma h)\delta-\gamma(h\delta)\) 之和，其中 \(h\) 屬於 \(\mathfrak h\)，\(\gamma\) 屬於 \(\mathfrak O\)，\(\delta\) 屬於 \(\mathfrak o\)，從而把該式改寫為
\[
(z_1\alpha_1+\cdots+z_r\alpha_r)+(t_1\beta_1+\cdots+t_s\beta_s)
\]
並滿足 \(z_1=0,\ldots,z_r=0\)；\(\beta_1=0,\ldots,\beta_s=0\)。這裡 \(\alpha,\beta\) 屬於 \(\mathfrak o\)，\(z,t\) 屬於 \(\mathfrak O\)，因而 \(z=0\) 表示 \(\mathfrak O\) 中各 \(x,\gamma,h\) 之間的關係，而 \(\beta=0\) 表示 \(\mathfrak o\) 中各 \(y,\delta,h\) 之間的關係。\srcfn{7)}{若不假定單位元存在，則 (3) 和 (4) 中還會出現附加的線性項。}
\end{enumerate}

\(\mathfrak O_{\mathfrak o}\) 中的兩個元素相等，當且僅當它們的差以上述方式消失。

因此，由直積的要求，等式關係與運算關係（加法和乘法）都由各因子中的相應關係唯一確定。由此可知：若 \(\overline{\mathfrak O},\overline{\mathfrak o}\) 及其交 \(\overline{\mathfrak h}\) 分別同構於 \(\mathfrak O,\mathfrak o\) 和 \(\mathfrak h\)，並且直積 \(\mathfrak O\times\mathfrak o\) 存在，那麼 \(\overline{\mathfrak O}\times\overline{\mathfrak o}\) 也存在，並且與 \(\mathfrak O\times\mathfrak o\) 同構。

設一個環 \(\mathfrak S\) 等於 \(\overline{\mathfrak O}\) 與 \(\overline{\mathfrak o}\) 的乘積，並且 \(\overline{\mathfrak O}\) 與 \(\overline{\mathfrak o}\) 的元素兩兩可交換，則 \(\mathfrak S\) 是 \(\mathfrak O\times\mathfrak o\) 的一個同態像，其中從 \(\overline{\mathfrak O}\) 到 \(\mathfrak O\) 以及從 \(\overline{\mathfrak o}\) 到 \(\mathfrak o\) 的映射仍為同構。因為，\(\mathfrak O\times\mathfrak o\) 中的等式蘊含在 \(\mathfrak S\) 中的等式之內，而在 \(\overline{\mathfrak O}\) 或 \(\overline{\mathfrak o}\) 上，它又分別與 \(\mathfrak O\) 或 \(\mathfrak o\) 在 \(\mathfrak O\times\mathfrak o\) 中的等式一致。\srcfn{8)}{另見 § 1, 4。}

\paragraph{2. 關於子環的直積存在的充要條件。}
設 \(\mathfrak O\) 與 \(\mathfrak o\) 是兩個環，它們的交為位於二者中心的環 \(\mathfrak h\)。關於 \(\mathfrak h\) 的直積未必存在，下面的例子說明了這一點：令 \(\mathfrak h\) 為全體有理整數所成的環，令 \(\mathfrak O=\mathfrak h[x]\) 並滿足 \(1-2x=0\)，再令 \(\mathfrak o=\mathfrak h[y]\) 並滿足 \(1-2y=0\)，其中 \(y\) 表示一個新符號。這樣，\(\mathfrak O\) 與 \(\mathfrak o\) 就是 \(\mathfrak h\) 的兩個等價環，它們在集合論意義下的交為 \(\mathfrak h\)，而符號 \(x\) 與 \(y\) 之間沒有給定任何運算關係。不可能存在一個同時包含 \(\mathfrak O\) 和 \(\mathfrak o\) 的環 \(R\)，使它們在 \(R\) 中---此時 \(x\) 與 \(y\) 之間已有運算關係---的交為 \(\mathfrak h\)。因為在 \(R\) 中有：
\[
0=(1-2x)y-x(1-2y)=y-x,
\]
而這裡甚至沒有用到 \(x\) 與 \(y\) 的可交換性。\srcfn{9)}{位於一個固定環內的交換環的商環---這裡即添入 \(\frac12\)---本來就是唯一的。這類似於兩個等價而相異的伽羅瓦域擴張不能嵌入同一個公共擴域。儘管如此，這裡仍可通過直積形成而嵌入一個含零因子的環。}

\emph{直積存在的充要條件}如下：

\emph{至少必須存在一個同時包含 \(\mathfrak O\) 和 \(\mathfrak o\) 的環 \(R\)，使得 \(\mathfrak O\) 與 \(\mathfrak o\) 的交為 \(\mathfrak h\)，並使 \(\mathfrak O\) 與 \(\mathfrak o\) 的元素兩兩可交換。}

事實上，如果直積存在，它本身就是這樣的一個環 \(R\)。反過來，為了構造直積，可以假定---下文始終作此假定---在 \(\mathfrak O\) 和 \(\mathfrak o\) 中不屬於 \(\mathfrak h\) 的元素之間沒有運算關係，因為通過過渡到 \(\mathfrak h\) 的等價擴張---即引入另一些符號---總能達到這一點。\srcfn{9a)}{參見 § 2, 1 及腳註 12)。}再設 \(\mathfrak h\) 包含 \(\mathfrak O\) 和 \(\mathfrak o\) 的單位元 \(e\)（不過也請參見注 7)）。現在把 \(\mathfrak T\) 定義為全體雙線性組合
\[
x_{i_1}y_{i_1}+\cdots+x_{i_r}y_{i_r}
 =y_{i_1}x_{i_1}+\cdots+y_{i_r}x_{i_r}
\]
所成的集合，其中各組 \(x\) 和 \(y\) 分別遍歷 \(\mathfrak O\) 與 \(\mathfrak o\) 中任意有限多個元素所成的全部系統，因而各 \(x\) 與各 \(y\) 在 \(\mathfrak T\) 中可交換。當且僅當兩個這樣的組合能借助 \(\mathfrak O\) 和 \(\mathfrak o\) 中的關係，按照 1, (4) 所精確定義的意義形式地變為相等時，才把它們規定為相等。再把
\[
\alpha_{i_1}x_{i_1}+\cdots+\alpha_{i_r}x_{i_r}
\quad\hbox{与}\quad
\beta_{i_1}x_{i_1}+\cdots+\beta_{i_r}x_{i_r}
\]
的和與積分別定義為
\[
(\alpha_{i_1}+\beta_{i_1})x_{i_1}+\cdots+(\alpha_{i_r}+\beta_{i_r})x_{i_r}
\quad\hbox{与}\quad
\sum\alpha_{i_\lambda}\beta_{i_\mu}x_{i_\lambda}x_{i_\mu}.
\]
這裡，在各 \(\alpha,\beta\) 之中也允許出現零。這個定義在上述等式意義下是唯一的，因為任何零關係（按 1, (4) 的意義）在右乘或左乘 \(\mathfrak o\) 或 \(\mathfrak O\) 中的一個元素之後仍然成為這樣的零關係，而這些零關係的和與差也有同樣性質。系統 \(\mathfrak T\) 構成一個同時包含 \(\mathfrak O\) 與 \(\mathfrak o\) 的環，並且在等式、加法和乘法方面滿足直積的各項條件。作為一般嵌入假設的結果，\(\mathfrak T\) 也滿足交的條件。事實上，令 \(\mathfrak S\) 為 \(R\) 中 \(\mathfrak O\) 與 \(\mathfrak o\) 的乘積；根據 1，\(\mathfrak S\) 是 \(\mathfrak T\) 的一個同態像，而在 \(\mathfrak O\) 與 \(\mathfrak o\) 上該同態為同構。若在 \(\mathfrak T\) 中有關係 \(x=y\)，其中 \(x\) 屬於 \(\mathfrak O\)，\(y\) 屬於 \(\mathfrak o\)，則在 \(\mathfrak S\) 中也有同一關係；因此，按假設，這兩個元素在 \(\mathfrak S\) 中都屬於 \(\mathfrak h\)，而由於映射在 \(\mathfrak O\) 和 \(\mathfrak o\) 上分別為同構，對 \(\mathfrak T\) 也有同樣結論。因而，嵌入假設既必要又充分。

\paragraph{3. 環關於子環的定義理想。}
若 \(\mathfrak O\) 中的元素 \(\ldots,z,\ldots\) 所成的系統 \(S\) 滿足：\(\mathfrak O\) 等於在 \(\mathfrak O\) 中由 \(\mathfrak h\) 與 \(S\) 生成的環---即等於 \(\mathfrak O\) 中所有同時包含 \(\mathfrak h\) 與 \(S\) 的環的交；用符號寫作 \(\mathfrak O=\mathfrak h[S]\)---，則稱 \(S\) 為 \(\mathfrak O\) 關於 \(\mathfrak h\) 的\emph{生成系}。仍假定 \(\mathfrak h\) 是 \(\mathfrak O\) 的中心的一個子環；於是，\(\mathfrak O\) 是非交換多項式環 \(\mathfrak h[..Z..]\) 的一個同態像，其中 \(Z\) 表示不定元，並且不定元系統與 \(S\) 具有相同的基數。這裡所謂非交換多項式環，是指所有有限和 \(\sum h_iP_i=\sum P_i h_i\) 所成的系統，其中 \(P_i\) 是冪乘積，包括與指數 \(0\) 對應的單位元：
\[
P_i=Z_{i_1}^{\varrho_1}Z_{i_2}^{\varrho_2}\cdots Z_{i_n}^{\varrho_n}
\]
（當然，相同的指標可以重複出現），而等式定義為關於各個不同冪乘積的係數相等。由上述同態可知，\(\mathfrak O\) 與剩餘類環 \(\mathfrak h[..Z..]/\mathfrak M\) 同構，其中雙邊理想 \(\mathfrak M\) 由所有滿足如下條件的多項式 \(F(Z)\)（即 \(\mathfrak h[..Z..]\) 的元素）組成：\(F(z)\) 在 \(\mathfrak O\) 中消失；這裡 \(Z\) 與 \(z\) 在該同態之下相互對應。\emph{理想 \(\mathfrak M\) 稱為關於 \(\mathfrak h\) 的一個定義理想，更確切地說，是與生成系 \(z\) 相對應的理想。}由定義可知，若兩個環關於同構對應的子環 \(\mathfrak h,\overline{\mathfrak h}\) 和生成系 \(S,\overline S\) 同構，那麼它們的定義理想也同構；特別地，當 \(\mathfrak h\) 與 \(\overline{\mathfrak h}\) 重合時，還可把這兩個定義理想視為相同。

因為已經假定 \(\mathfrak h\) 位於 \(\mathfrak O\) 的中心，所以只要 \(\mathfrak O\) 有一個僅由單個元素 \(z\) 組成的生成系 \(S\)，\(\mathfrak O\) 就必為交換環。若 \(\mathfrak M\) 還是主理想，則稱由 \(\mathfrak M\) 的一個基多項式 \(G(Z)\) 給出的等式 \(G(z)=0\) 為 \(\mathfrak O\) 關於 \(\mathfrak h\) 的一個定義方程。

\paragraph{4. 直積的定義理想。}
設（如 3 中）\(\mathfrak O\) 與 \(\mathfrak h[..Z..]/\mathfrak M\) 同構，其中 \(\mathfrak M\) 是定義理想。在 \(\mathfrak h[..Z..]\) 之外，再考察多項式環 \(\mathfrak o[..Z..]\)，即全體有限和
\[
\sum \gamma_iP_i=\sum P_i\gamma_i
\qquad\hbox{其中 }\gamma_i\hbox{ 属于 }\mathfrak o
\]
所成的系統，並把等式定義為係數相等。\(\mathfrak o[Z]\) 包含 \(\mathfrak h[Z]\)，並且是 \(\mathfrak h[..Z..]\) 與 \(\mathfrak o\) 關於 \(\mathfrak h\) 的直積；換言之，它由 \(\mathfrak h[..Z..]\) 通過擴張係數環而得到。\srcfn{10)}{因為等式關係與運算關係都得到滿足，並且由 \(\gamma=h+h_1P_1+\cdots+h_rP_r\)，按照 \(\mathfrak o[..Z..]\) 中的等式可知 \(\gamma=h\)。對此可參見 § 2, 2。本項並不使用 \(\mathfrak o[..Z..]\) 是直積這一事實。}以 \(\mathfrak M_{\mathfrak o}\) 表示 \(\mathfrak M\) 關於 \(\mathfrak o[..Z..]\) 的擴張理想---即 \(\mathfrak o[..Z..]\) 中所有包含 \(\mathfrak M\) 的雙邊理想的交---；因而，它是 \(\mathfrak M\) 中各元素以 \(\mathfrak o\) 中元素為係數的全體線性組合所成的系統，並且由於 \(\mathfrak o\) 與各 \(Z\) 可交換，它已經是一個雙邊理想。\emph{如果關於 \(\mathfrak h\) 的直積 \(\mathfrak O_{\mathfrak o}\) 存在，那麼 \(\mathfrak O_{\mathfrak o}\) 與 \(\mathfrak o[..Z..]/\mathfrak M_{\mathfrak o}\) 同構。}事實上，\(\mathfrak O_{\mathfrak o}\) 是 \(\mathfrak o[..Z..]\) 的一個同態像，因而與模 \(\mathfrak O_{\mathfrak o}\) 關於 \(\mathfrak o\) 的定義理想所成的剩餘類環同構，該定義理想對應於生成系 \(z\)。而根據 1, (4) 的等式定義，\(\mathfrak O_{\mathfrak o}\) 中各 \(z\) 之間的關係可以寫成 \(\mathfrak O\) 中已有關係以 \(\mathfrak o\) 中元素為係數的線性組合；所以，\(\mathfrak M_{\mathfrak o}\) 就是定義理想。

\paragraph{注。}
同樣的考察還給出直積及其相應定義理想的一種對稱但較為繁複的表示。設 \(\mathfrak o\) 與 \(\mathfrak h[..Y..]/\mathfrak N\) 同構，其中 \(Y\) 表示不同於 \(Z\) 的符號，從而多項式環 \(\mathfrak h[..Y..Z..]\) 有定義。於是，\(\mathfrak O_{\mathfrak o}\) 同構於
\[
\mathfrak h[..Y..Z..]/(\mathfrak M,\mathfrak N,..YZ-ZY..),
\]
其中右側的定義理想由 \(\mathfrak M,\mathfrak N\) 和所有組合 \(YZ-ZY\) 生成。（後一類組合表示 \(\mathfrak o\) 與 \(\mathfrak O\) 中元素的可交換性，前兩者表示 \(\mathfrak O_{\mathfrak o}\) 中的關係都是 \(\mathfrak O\) 和 \(\mathfrak o\) 各自關係的推論。）
% 原文：第43篇，印刷頁 7--10 （期刊掃描頁 696--699）
\subsection*{§ 2. 具有獨立模基的環。直積的構造。}

\paragraph{1. 直積的構造。}
若 \(\mathfrak O\) 中的元素 \(\ldots,t,\ldots\) 所成的系統 \(T\) 滿足：\(\mathfrak O\) 等於在 \(\mathfrak O\) 中由 \(T\) 導出的 \(\mathfrak h\)-模，則稱 \(T\) 為 \(\mathfrak O\) 的一組 \(\mathfrak h\)-模基。由於單位元存在，\(\mathfrak O\) 此時由全體線性組合 \(h_{i_1}t_{i_1}+\cdots+h_{i_r}t_{i_r}\) 組成。若
\[
h_{i_1}t_{i_1}+\cdots+h_{i_r}t_{i_r}=0
\quad\hbox{总能推出}\quad
h_{i_1}=0,\ldots,h_{i_r}=0,
\]
則稱這組模基是獨立的。\emph{如果 \(\mathfrak O\) 有一組包含單位元的獨立模基，那麼，對於 \(\mathfrak h\) 的任意擴張環 \(\mathfrak o\)，只要 \(\mathfrak h\) 是其中心的一個子環，並且 \(\mathfrak O\) 與 \(\mathfrak o\) 的集合論意義下的交 \([\mathfrak O,\mathfrak o]\) 等於 \(\mathfrak h\)，直積 \(\mathfrak O_{\mathfrak o}\) 就存在。}\srcfn{11)}{與此對稱，也可以對 \(\mathfrak o\) 作出基的假設。根據 § 2, 2，在不失一般性之下---參見 E. Noether, Der Diskriminantensatz für die Ordnungen eines algebraischen Zahl- oder Funktionenkörpers, dies. Journ. 157 (1927), 82--104, §§ 1, 2---假定 \(\mathfrak O\) 與 \(\mathfrak o\) 中不屬於 \(\mathfrak h\) 的元素之間不存在運算關係。}根據 § 1, 2，只須構造一個相應的嵌入環 \(\mathfrak R\)；在這裡，該環隨後就會證明為直積本身。把 \(\mathfrak R\) 定義為全體線性組合
\[
\gamma_{i_1}t_{i_1}+\cdots+\gamma_{i_r}t_{i_r}
=t_{i_1}\gamma_{i_1}+\cdots+t_{i_r}\gamma_{i_r}
\]
所成的系統，其中各 \(t_{i_r}\) 遍歷獨立基 \(T\) 中任意有限多個元素所成的全部系統，相應的 \(\gamma\) 取自 \(\mathfrak o\)。等式定義為關於各 \(t_{i_\nu}\) 的係數相等；和與積則按通常的運算法則定義，如 § 1, 2 所述。由基的性質，任意兩個 \(t_{i_\nu}\) 的積都是有限多個基元素以 \(\mathfrak h\) 中元素為係數的線性組合，所以 \(\mathfrak R\) 成為一個環。\(\mathfrak R\) 包含 \(\mathfrak O\)，因為 \(\mathfrak h\) 是 \(\mathfrak o\) 的子環；它也包含 \(\mathfrak o\)，因為 \(T\) 包含單位元。又因為各 \(\gamma\) 分別與各 \(h_{i_\nu}\) 以及各 \(t_{i_\nu}\) 兩兩可交換，\(\mathfrak O\) 與 \(\mathfrak o\) 在 \(\mathfrak R\) 中的元素兩兩可交換。最後，\(\mathfrak O\) 與 \(\mathfrak o\) 在 \(\mathfrak R\) 中的交為 \(\mathfrak h\)，因為關係
\[
\gamma=\gamma e=h_1t_1+\cdots+h_rt_r
\]
按照等式定義---並利用 \(e\) 按假設是基的一個元素---說明其中一個 \(t\) 等於 \(e\)，其餘項消失，因而 \(\gamma=h\)。所以直積存在；它等於 \(\mathfrak R\)，因為 \(\mathfrak R\) 等於 \(\mathfrak O\) 與 \(\mathfrak o\) 的乘積，並滿足可交換性與等式條件（由於基元素獨立，這個條件歸結為係數相等）。\(\mathfrak O\) 的這組基 \(T\) 同時成為 \(\mathfrak O_{\mathfrak o}\) 的一組 \(\mathfrak o\)-基。

\paragraph{2. 擴張模與收縮模。}
若 \(\mathfrak B\) 是 \(\mathfrak O\) 中的一個 \(\mathfrak h\)-模，則其擴張 \(\mathfrak B_{\mathfrak o}\) 是指 \(\mathfrak O_{\mathfrak o}\) 中的這樣一個 \(\mathfrak o\)-模：它等於 \(\mathfrak O_{\mathfrak o}\) 中所有包含 \(\mathfrak B\) 的 \(\mathfrak o\)-模的交。若 \(\mathfrak C\) 是 \(\mathfrak O_{\mathfrak o}\) 中的一個 \(\mathfrak o\)-模，則它在 \(\mathfrak O\) 中的收縮模 \(\mathfrak C_{\mathfrak h}\) 定義為交 \([\mathfrak C,\mathfrak O]\)。若 \(\mathfrak B\) 同時是 \(\mathfrak O\) 中的理想，則 \(\mathfrak B_{\mathfrak o}\) 是 \(\mathfrak O_{\mathfrak o}\) 中的理想；若 \(\mathfrak C\) 是 \(\mathfrak O_{\mathfrak o}\) 中的理想，則 \([\mathfrak C,\mathfrak O]\) 是 \(\mathfrak O\) 中的理想。

\paragraph{定理。}
\emph{如果 \(\mathfrak O\) 中的一個模 \(\mathfrak B\) 有一組獨立 \(\mathfrak h\)-模基，而且這組基可以補成 \(\mathfrak O\) 的一組獨立 \(\mathfrak h\)-模基，那麼 \(\mathfrak B\) 是它在 \(\mathfrak O_{\mathfrak o}\) 中的擴張模 \(\mathfrak B_{\mathfrak o}\) 的收縮。}

\paragraph{證明。}
以 \(T\) 表示 \(\mathfrak B\) 的一組獨立 \(\mathfrak h\)-模基；按假設，可用 \(\mathfrak O\) 中由元素 \(\overline t\) 組成的系統 \(\overline T\) 把它補成 \(\mathfrak O\) 的一組獨立 \(\mathfrak h\)-模基。於是，\(T\) 同時是 \(\mathfrak B_{\mathfrak o}\) 的一組獨立 \(\mathfrak o\)-基，而 \(T\) 與 \(\overline T\) 合在一起則是 \(\mathfrak O_{\mathfrak o}\) 的這樣一組基。

\(\mathfrak B\) 包含在交 \([\mathfrak B_{\mathfrak o},\mathfrak O]\) 中。反過來，設 \(c\) 是這個收縮模的一個元素。作為 \(\mathfrak B_{\mathfrak o}\) 的元素，\(c\) 可以用 \(T\) 中的元素表示；同時，這也是用 \(\mathfrak O_{\mathfrak o}\) 的 \(\mathfrak o\)-基 \(T,\overline T\) 作出的一個表示：
\[
c=\gamma_1t_1+\cdots+\gamma_rt_r
\qquad\hbox{其中 }\gamma\hbox{ 属于 }\mathfrak o.
\]
作為 \(\mathfrak O\) 的元素，\(c\) 又有一個用 \(\mathfrak O\) 的 \(\mathfrak h\)-基 \(T,\overline T\) 作出的表示，即：
\[
c=h_{i_1}t_{i_1}+\cdots+h_{i_r}t_{i_r}+h_{j_1}\overline t_{j_1}+\cdots+h_{j_s}\overline t_{j_s}.
\]
由於 \(\mathfrak O\) 是 \(\mathfrak O_{\mathfrak o}\) 的子環，而 \(\mathfrak h\) 是 \(\mathfrak o\) 的子環，這同時也是用 \(\mathfrak O_{\mathfrak o}\) 的 \(\mathfrak o\)-基 \(T,\overline T\) 作出的表示。因此，這兩個表示的對應係數必相等；各 \(\overline t\) 的係數消失，並且
\[
c=h_{i_1}t_{i_1}+\cdots+h_{i_r}t_{i_r}
\]
表明 \(c\) 是 \(\mathfrak B\) 的元素。由此，\(\mathfrak B\) 已被識別為收縮模。

\paragraph{附註。}
證明所用的假設比定理中表述的略弱，並且還證明了其餘各項；因而有如下加強的表述：

\emph{如果 \(\mathfrak O\) 有一組獨立 \(\mathfrak h\)-模基 \(T,\overline T\)，其中 \(T\) 包含於 \(\mathfrak B\)，並且 \(T\) 同時是 \(\mathfrak B_{\mathfrak o}\) 的一組（獨立）\(\mathfrak o\)-模基，那麼 \(\mathfrak B\) 是其擴張 \(\mathfrak B_{\mathfrak o}\) 的收縮模。此時還可證明 \(T\) 是 \(\mathfrak B\) 的一組獨立 \(\mathfrak h\)-模基。}

\paragraph{3. 直積存在的一個充分判別準則。}
有時（例如在相對不同的情形中）只有在擴張係數環 \(\mathfrak h\) 之後，獨立模基才存在。這時可由下列充分判別準則推得直積存在：

\emph{若存在 \(\mathfrak h\) 的一個擴張環 \(\mathfrak f\)，使得關於 \(\mathfrak h\) 的直積 \(\mathfrak O_{\mathfrak f}\) 和 \(\mathfrak o_{\mathfrak f}\) 存在，並且關於 \(\mathfrak f\) 的直積 \(\mathfrak O_{\mathfrak f}\times\mathfrak o_{\mathfrak f}\) 也存在，那麼關於 \(\mathfrak h\) 的 \(\mathfrak O\times\mathfrak o\) 存在。}

須證明 \(\mathfrak R=\mathfrak O_{\mathfrak f}\times\mathfrak o_{\mathfrak f}\) 是 § 1, 2 意義下的嵌入環。事實上，\(\mathfrak O\) 與 \(\mathfrak o\) 作為 \(\mathfrak O_{\mathfrak f}\) 與 \(\mathfrak o_{\mathfrak f}\) 的子環，其元素兩兩可交換；又有
\[
[\mathfrak O,\mathfrak o]\subseteqq[\mathfrak O_{\mathfrak f},\mathfrak o_{\mathfrak f}]=\mathfrak f;
\]
同樣，\([\mathfrak O,\mathfrak o]\subseteqq\mathfrak O\)；所以 \([\mathfrak O,\mathfrak o]\subseteqq[\mathfrak O,\mathfrak f]=\mathfrak h\)，從而 \([\mathfrak O,\mathfrak o]=\mathfrak h\)。由此證明了 \(\mathfrak R\) 是一個嵌入環。

\paragraph{注。}
上述證明只對 \(\mathfrak O_{\mathfrak f}\) 使用了關於交的假設，對 \(\mathfrak o_{\mathfrak f}\) 並未使用；所以，後一個環也可以僅僅是 \(\mathfrak o\) 與 \(\mathfrak f\) 的乘積。
\subsection*{§ 3. 環相對於子環的不同。定義理想的微分商。}

從現在起，假定所出現的一切環都是交換環。\srcfn{12)}{若無這一假定，就必須定義右不同與左不同，並且必須考察二者彼此符合到何種程度，以及在半單系統的階的情形下，它們同已知構造符合到何種程度（作為 § 4 中研究的推廣）。}此外，始終假定環 $\mathfrak O$ 和 $\mathfrak o$ 是子環 $\mathfrak h$ 的同構且等價的擴張（亦即從 $\mathfrak O$ 到 $\mathfrak o$ 的映射是 $\mathfrak h$ 上恆等映射的延拓），並且關於 $\mathfrak h$ 的直積 $\mathfrak O_{\mathfrak o}$ 存在。

\paragraph{1. 不同的定義。}
假定關於 $\mathfrak h$ 的直積 $\mathfrak O_{\mathfrak o}$ 存在。令 $x$ 遍歷 $\mathfrak O$ 的所有元素，並令 $\xi$ 每次表示 $\mathfrak o$ 中由該同構所對應的元素；以 $\mathfrak P$ 表示由所有差 $x-\xi$ 生成的 $\mathfrak O_{\mathfrak o}$ 的雙邊\emph{差分理想}，
\[
\mathfrak P=\{..(x-\xi)..\}\srcfnmark{13)}
\]
\srcfntext{13)}{這一記號將始終用於由集合生成的理想。}
\emph{差分商} $\mathfrak A$ 定義為零理想關於 $\mathfrak P$ 的理想商，因而由 $\mathfrak O_{\mathfrak o}$ 中所有使 $a\mathfrak P$ 消失的元素 $a$ 組成，
\[
\mathfrak A=(0):\mathfrak P.
\]
$\mathfrak o$ 關於 $\mathfrak h$ 的\emph{不同} $\mathfrak d$ 定義為 $\mathfrak A[x\to\xi]$，也就是說，它由 $\mathfrak A$ 產生，方法是在 $\mathfrak A$ 的每個 $a=\alpha_{i_1}x_{i_1}+\cdots+\alpha_{i_r}x_{i_r}$ 中，將各 $x$ 換成與之同構對應的 $\xi$。相應地，$\mathfrak O$ 關於 $\mathfrak h$ 的不同 $\mathfrak D$ 定義為 $\mathfrak A[\xi\to x]$；由於 $\mathfrak A$ 的定義關於 $x$ 與 $\xi$ 對稱，$\mathfrak d$ 與 $\mathfrak D$ 是由交換 $\xi$ 與 $x$ 得到的，因此在從 $\mathfrak o$ 到 $\mathfrak O$ 的同構下彼此對應。由於 $\mathfrak A$ 是 $\mathfrak O_{\mathfrak o}$ 中的理想，特別地既是 $\mathfrak o$-模又是 $\mathfrak O$-模，故 $\mathfrak d$ 與 $\mathfrak D$ 分別成為 $\mathfrak o$ 與 $\mathfrak O$ 中的理想。

\paragraph{2. 定義理想的微分商。}
設（§ 1, 3）$\mathfrak M$ 是 $\mathfrak O$ 關於 $\mathfrak h$ 的定義理想，與 $z$ 的生成系相對應。於是（§ 1, 3）$\mathfrak M$ 也是 $\mathfrak o$ 關於 $\mathfrak h$ 的定義理想，與 $\xi$ 的同構生成系相對應；而且（§ 1, 4）$\mathfrak O_{\mathfrak o}$ 同構於 $\mathfrak o[..Z..]/\mathfrak M_{\mathfrak o}$。現在，差分理想與差分商在多項式環 $\mathfrak o[..Z..]$ 中分別定義為
\[
\mathfrak Q=\{..(Z-\xi)..\}
\]
以及
\[
\mathfrak C=\mathfrak M_{\mathfrak o}:\mathfrak Q
=\{..(F(Z)-F(\xi))..\}:\{..(Z-\xi)..\}.
\]
\emph{微分商} $\mathfrak M'$ 定義為 $\mathfrak C[\xi\to Z]$。這裡應當注意，這一定義是唯一的，並不依賴於用生成元 $\xi$ \emph{表示} $\mathfrak o$ 中元素的方式。事實上，由於 $\mathfrak C$ 是理想商，它包含 $\mathfrak M_{\mathfrak o}$；同一個 $\mathfrak o$ 中元素的兩個表示相差一個多項式 $F(\xi)$，因而在代入 $Z$ 後，相差一個屬於 $\mathfrak M$、從而已經屬於 $\mathfrak C$ 的 $F(Z)$。理想 $\mathfrak M'$ 成為 $\mathfrak h[..Z..]$ 中包含 $\mathfrak M$ 的理想。\emph{微分商 $\mathfrak M'$ 在 $Z\to\xi$ 時過渡為 $\mathfrak o$ 關於 $\mathfrak h$ 的不同 $\mathfrak d$。}這是因為在同態映射下，不僅理想彼此對應，商也彼此對應。因此，在 $\mathfrak o[..Z..]$ 與 $\mathfrak O_{\mathfrak o}$ 之間、即使 $Z$ 與 $z$ 對應的同態下，$\mathfrak A$ 成為 $\mathfrak C$ 的同態像，因為相對於 $\mathfrak Q$ 與 $\mathfrak M_{\mathfrak o}$，$\mathfrak P$ 與零理想也有同樣的關係。\srcfn{13a)}{編者注：為了能夠斷定商彼此對應，關鍵在於 $\mathfrak M_{\mathfrak o}$ 表示 $(0)$ 的原像集合。}由此可知，在該同態下，$\mathfrak C$ 中的代換 $\xi\to Z$ 對應於 $\mathfrak A$ 中的代換 $\xi\to z$；因此，在 $Z$ 與 $z$ 對應之下，$\mathfrak O$ 的不同 $\mathfrak D$ 成為 $\mathfrak M'$ 的同態像。所以
\[
\mathfrak M'[Z\to z]=\mathfrak D
\qquad\hbox{且}\qquad
\mathfrak M'[Z\to\xi]=\mathfrak D[z\to\xi]=\mathfrak d.
\]

\paragraph{3. 有一個定義方程時的不同。}
設 $\mathfrak O$ 有一個定義方程 $G(z)=0$（§ 1, 3 末尾），並進一步\emph{假定}最高次項的係數等於單位元，即
\[
G(Z)=Z^n+h_1Z^{n-1}+\cdots+h_n
\]
且係數 $h_i$ 屬於 $\mathfrak h$。那麼，它與多項式微分商之間的聯繫由下述結果給出：

\paragraph{定理。}
\emph{$\mathfrak o$ 或 $\mathfrak O$ 關於 $\mathfrak h$ 的不同有定義（即關於 $\mathfrak h$ 的直積 $\mathfrak O_{\mathfrak o}$ 存在），並且分別成為以 $G'(\xi)$ 或 $G'(z)$ 為生成元的主理想，其中 $G'(Z)$ 表示多項式微分商。理想微分商 $\mathfrak M'$ 等於 $(G'(Z),G(Z))$。}

\paragraph{證明。}
由假定首先可知，元素 $e,z,\ldots,z^{n-1}$ 構成 $\mathfrak O$ 的一個獨立模基。事實上，因為 $G(Z)$ 的最高次係數是單位元，這些冪確實構成一個 $\mathfrak h$-模基；而這個模基是獨立的，因為 $\mathfrak M$ 不可能包含次數低於 $n$ 的多項式。多項式 $K(Z)G(Z)$ 的次數確實等於兩個因子次數之和，這仍然是由關於最高次係數的假定所致。於是，由於該模基存在，關於 $\mathfrak h$ 的直積 $\mathfrak O_{\mathfrak o}$ 存在，從而不同有定義。其餘結論由考察 $\mathfrak o[..Z..]$ 中的差分商 $\mathfrak C$ 得到（參見 2）。這個差分商成為以如下多項式為生成元的主理想：
\[
%
C(Z,\xi)=(Z^{n-1}+Z^{n-2}\xi+\cdots+\xi^{n-1})
+h_1(Z^{n-2}+\cdots+\xi^{n-2})+\cdots+h_{n-1}.
%
\]
事實上，$C$ 由 $G(Z)-G(\xi)$ 除以 $Z-\xi$ 的形式除法得到，因而屬於 $\mathfrak C$，因為 $Z-\xi$ 成為 $\mathfrak Q$ 的生成元。任意多項式模 $C$ 均可約化為一個關於 $Z$ 的次數達不到 $n-1$ 的多項式 $D$；於是乘積 $D(Z-\xi)$ 的次數達不到 $n$，所以只有當它消失時才可能屬於 $\mathfrak M$。但是，由 $D(Z-\xi)=0$ 可推出 $D=0$，因為 $Z-\xi$ 中 $Z$ 的係數是單位元。所以 $C$ 成為 $\mathfrak C$ 的生成元。而由在 $C$ 及其倍式中代換 $\xi\to Z$（亦即 $G(Z)=C(Z-\xi)$ 及其倍式），現在得到
\[
\mathfrak M'=(G'(Z),G(Z)),
\]
進而得到
\[
\mathfrak d=(G'(\xi))
\qquad\hbox{且}\qquad
\mathfrak D=(G'(z)).
\]
% 原文：第43篇，印刷頁 10--13 （期刊掃描頁 699--702）
\subsection*{\S{} 4. 直和的不同。}

對 $\mathfrak O$ 作如下假定：

\begin{enumerate}
\item[(1)] $\mathfrak O$ 是諸理想的直和
\[
\mathfrak O=\mathfrak R_1+\cdots+\mathfrak R_r
=\mathfrak R_i+\mathfrak S_i
=\mathfrak Oe_1+\cdots+\mathfrak Oe_r,
\]
其中 $e_i$ 表示單位元的各個分量。

\item[(2)] 每個 $\mathfrak R_i$ 都有一個獨立的 $\mathfrak h_i$-模基 $T_i$，其中 $\mathfrak h_i$ 是 $\mathfrak h$ 在 $\mathfrak R_i$ 中的分量，即 $\mathfrak h_i=\mathfrak h e_i$。這個基 $T_i$ 包含 $\mathfrak R_i$ 的單位元 $e_i$ 作為一個元素。

\item[(3)] 每個交集 $[\mathfrak h,\mathfrak S_i]$ 都等於零元（$\mathfrak S_i$ 不含常數）。
\end{enumerate}

由於存在同構，同樣的假定對 $\mathfrak o$ 也成立。

\paragraph{定理。}
\emph{在上述假定下，$\mathfrak O$ 關於 $\mathfrak h$ 的不同成為各 $\mathfrak R_i$ 關於 $\mathfrak h_i$ 的不同的直和。}

證明以證實一系列個別事實為基礎。

\paragraph{I.}
首先可知，$\mathfrak O$ 也有一個包含 $e$ 的獨立 $\mathfrak h$-模基，因而 $\mathfrak O$ 或 $\mathfrak o$ 關於 $\mathfrak h$ 的不同有定義；具體而言，這個基由 $\mathfrak R_i$ 的各個基 $T_i$ 的並集 $T$ 構成。事實上，由
\[
x=xe_1+\cdots+xe_r
\]
以及
\[
xe_i=(h_1e_i)t_{1i}e_i+\cdots+(h_\lambda e_i)t_{\lambda i}e_i
=h_1(t_{1i}e_i)+\cdots+h_\lambda(t_{\lambda i}e_i)
\]
可知，$T$ 確實成為 $\mathfrak O$ 的一個 $\mathfrak h$-基。這個基還是獨立的；因為由 $T$ 中元素之間的一個關係，藉助直和可得到每個 $T_i$ 中元素之間的關係，因而對所有係數都有 $h e_i=0$。但根據假定 (3)，這又給出 $h=0$；事實上，由於
\[
h\equiv h e_i\ (\mathfrak S_i)
\]
故從 $h e_i=0$ 可知 $h$ 屬於 $[\mathfrak h,\mathfrak S_i]$。此外，$T$ 包含 $\mathfrak R_i$ 的各個單位元 $e_i$，而它們的和是 $e$，所以可在 $T$ 中以 $e$ 替換一個 $e_i$，由此得到滿足一切條件的 $\mathfrak h$-基 $\overline T$。於是便得到 $\mathfrak O_{\mathfrak o}$ 以及 $\mathfrak O$ 關於 $\mathfrak h$ 的不同的存在性。

\paragraph{II.}
交集條件 (3) 現在蘊含如下更強的條件：
\[
[\mathfrak o,\mathfrak S_{i(\mathfrak o)}]=0,
\]
其中 $\mathfrak S_{i(\mathfrak o)}$ 表示 $\mathfrak O_{\mathfrak o}$ 中由 $\mathfrak S_i$ 導出的理想。設
\[
\alpha=\alpha e=\beta_1t_1+\cdots+\beta_\lambda t_\lambda,
\]
其中 $t_1,\ldots,t_\lambda$ 是從 $T$ 中刪去 $T_i$ 所得的基的元素，而 $\alpha$ 與 $\beta$ 屬於 $\mathfrak o$。若以 $e$ 替換 $e_i$ 而構成 $\overline T$，那麼 $e,t_1,\ldots,t_\lambda$ 就是 $\mathfrak O$ 的 $\mathfrak h$-基 $\overline T$ 中彼此不同的元素。但是，由於 $\mathfrak O_{\mathfrak o}$ 中的等式在每個這種基中都化為係數相等，關係
\[
\alpha=\alpha e=\beta_1t_1+\cdots+\beta_\lambda t_\lambda
\]
便給出 $\alpha=0$，從而\emph{證明了加強的交集性質}。由此特別得到（如上）：由
\[
\alpha\text{ 是 }\mathfrak o\text{ 的元素}\quad\text{且}\quad\alpha\ne0
\]
可推出 $\alpha e_i\ne0$。

\paragraph{III.}
這個交集條件的一個直接推論是：$\mathfrak o$ 與它的、對應於 $\mathfrak O$ 的分解的各個\emph{分量} $\mathfrak o e_i$ \emph{環同構}；此同構包含同構
\[
\mathfrak h\simeq\mathfrak h_i=\mathfrak h e_i.
\]
事實上，由於 $e_i^2=e_i$，並且 $e_i$ 與 $\mathfrak o$ 的所有元素可交換，$\mathfrak o e_i$ 成為 $\mathfrak o$ 的環同態像；這一對應又是一一的，因為由 $\alpha\ne0$ 也可推出 $\alpha e_i\ne0$。在這一對應下，子環 $\mathfrak h$ 對應於分量 $\mathfrak h_i=\mathfrak h e_i$。

\paragraph{IV.}
若進一步採用 $\mathfrak o$ 的相應分解，便得到\emph{另外的同構}：
\[
\mathfrak o=\mathfrak r_1+\cdots+\mathfrak r_r
=\mathfrak r_i+\mathfrak s_i
=\mathfrak o\varepsilon_1+\cdots+\mathfrak o\varepsilon_r.
\]
於是有：
\[
\mathfrak o\simeq\mathfrak o e_i,
\quad \mathfrak O\simeq\mathfrak O\varepsilon_i,
\quad \mathfrak r_\lambda\simeq\mathfrak r_\lambda e_i,
\quad \mathfrak R_\lambda\simeq\mathfrak R_\lambda\varepsilon_i;
\quad \mathfrak h\simeq\mathfrak h e_i\simeq\mathfrak h\varepsilon_i\simeq\mathfrak h e_i\varepsilon_i.
\]
如果在上述考察中從 $\mathfrak o$ 而不是 $\mathfrak O$ 出發，關係 $\mathfrak O\simeq\mathfrak O\varepsilon_i$ 就與 $\mathfrak o\simeq\mathfrak o e_i$ 完全對應。同樣，接下來的兩個關係彼此對應；其中第一個關係來自如下事實：在 $\mathfrak o\simeq\mathfrak o e_i$ 下，子環 $\mathfrak r_\lambda$ 也對應於 $\mathfrak r_\lambda e_i$；同理得到
\[
\mathfrak h\simeq\mathfrak h e_i\simeq\mathfrak h\varepsilon_i.
\]
而由於 $\mathfrak h\varepsilon_i$ 是 $\mathfrak o$（也是 $\mathfrak r_i$）的子環，最後還有
\[
\mathfrak h\varepsilon_i\simeq\mathfrak h\varepsilon_i e_i.
\]

\paragraph{V.}
由 $\mathfrak O$ 的直和表示可知：\emph{若 $\mathfrak C=(0):\mathfrak B$ 是 $\mathfrak O$ 中零理想關於理想 $\mathfrak B$ 的商，則每個分量 $\mathfrak C_i$ 等於 $\mathfrak R_i$ 的零理想關於分量 $\mathfrak B_i$ 的商。}由
\[
\mathfrak B\mathfrak C=(\mathfrak B e_1)(\mathfrak C e_1)+\cdots+(\mathfrak B e_r)(\mathfrak C e_r)=0
\]
根據\emph{直和}可逐一推出
\[
(\mathfrak B e_i)(\mathfrak C e_i)=0.
\]
反過來，若 $\mathfrak R_i$ 中的元素 $r e_i$ 滿足條件 $(\mathfrak B e_i)(r e_i)=0$，那麼對 $i\ne k$ 也有 $(\mathfrak B e_k)(r e_i)=0$。

所以 $r e_i$ 是 $\mathfrak C$ 的元素，同時也是 $\mathfrak R_i$ 的元素，因而屬於分量 $\mathfrak C_i=\mathfrak C e_i$。由此識別出 $\mathfrak C_i=\mathfrak C e_i$ 是零理想關於 $\mathfrak B_i$ 的商；它也就是 $\mathfrak R_i$ 的零理想關於 $\mathfrak B_i$ 的商，因為由
\[
r e_i\mathfrak B_i\equiv0\quad(\mathfrak S_i)
\]
總能推出 $r e_i\mathfrak B_i=0$。\srcfn{14)}{本定理也可直接由第 2 節所用的定理得到，即在同態映射下，不僅理想彼此對應，商也彼此對應。}由於除 $\mathfrak O$ 外，
\[
\mathfrak O_{\mathfrak o}=\mathfrak O_{\mathfrak o}e_1+\cdots+\mathfrak O_{\mathfrak o}e_r
\]
也成為直和，\emph{所以該定理對 $\mathfrak O_{\mathfrak o}$ 中的理想也成立。}

\paragraph{VI. 不同的分量分解。}
按定義（參見 § 3, 1），原有 $\mathfrak D=\mathfrak A[\xi\to x]$；故有
\[
\mathfrak D=\mathfrak D e_1+\cdots+\mathfrak D e_r
=\mathfrak A[\xi\to x]e_1+\cdots+\mathfrak A[\xi\to x]e_r.
\]
但是，由於 $e_i$ 屬於 $\mathfrak O$，因而對應於其自身；又由於 $\varepsilon_i$ 在該同構下過渡為 $e_i$，並且 $e_i^n=e_i$，所以
\[
\mathfrak A[\xi\to x]e_i=\mathfrak A e_i[\xi\to x]
=\mathfrak A e_i\varepsilon_i[\xi\to x].
\]
而根據 V，$\mathfrak A e_i\varepsilon_i$ 等於商 $(0):\mathfrak P e_i\varepsilon_i$，也等於 $\mathfrak R_i\mathfrak r_i$ 的零理想關於 $\mathfrak P e_i\varepsilon_i$ 的商。事實上，$\mathfrak O$ 與 $\mathfrak o$ 的直和分解經乘法產生分解
\[
\mathfrak O_{\mathfrak o}=\sum\mathfrak O_{\mathfrak o}e_i\varepsilon_k
=\sum\mathfrak R_i\mathfrak r_k,
\]
這個分解是直和，因為 $\sum e_i\varepsilon_k=e$，並且這些單位元滿足正交關係。

\paragraph{VII.}
\emph{證明 $\mathfrak D e_i=\mathfrak A e_i\varepsilon_i[\xi\to x]=(0):\mathfrak P e_i\varepsilon_i$ 是 $\mathfrak R_i$ 關於 $\mathfrak h_i$ 的不同。}這一證明基於下列事實：$\mathfrak R_i\mathfrak r_i$ 等於關於 $\mathfrak h e_i\varepsilon_i$ 的直積 $\mathfrak R_i\varepsilon_i\times\mathfrak r_i e_i$；其差分理想由 $\mathfrak P e_i\varepsilon_i$ 給出，其差分商由 $\mathfrak A e_i\varepsilon_i$ 給出，而 $\mathfrak D e_i\varepsilon_i$ 成為 $\mathfrak R_i\varepsilon_i$ 的不同，相應地 $\mathfrak d\varepsilon_i e_i$ 成為 $\mathfrak r_i e_i$ 的不同。於是，同構
\[
\mathfrak R_i\varepsilon_i\simeq\mathfrak R_i
\qquad\text{和}\qquad
\mathfrak r_i\simeq\mathfrak r_i e_i
\]
分別給出 $\mathfrak D e_i$ 與 $\mathfrak d\varepsilon_i$，作為 $\mathfrak R_i$ 關於 $\mathfrak h e_i$ 的不同與 $\mathfrak r_i$ 關於 $\mathfrak h\varepsilon_i$ 的不同。

\paragraph{VIII.}
\emph{$\mathfrak R_i\mathfrak r_i$ 等於關於 $\mathfrak h e_i\varepsilon_i$ 的直積 $\mathfrak R_i\varepsilon_i\times\mathfrak r_i e_i$。}由於 $\mathfrak R_i$ 有 $\mathfrak h_i$-模基 $T_i$，從同構 (IV) $\mathfrak R_i\simeq\mathfrak R_i\varepsilon_i$ 可知，$T_i\varepsilon_i$ 成為 $\mathfrak R_i\varepsilon_i$ 的一個 $\mathfrak h e_i\varepsilon_i$-模基。現在，根據 IV，$\mathfrak r_i e_i$ 同構於 $\mathfrak R_i\varepsilon_i$——因為藉助 $\mathfrak o\simeq\mathfrak O$ 有 $\mathfrak r_i\simeq\mathfrak R_i$——並且是 $\mathfrak h e_i\varepsilon_i$ 的擴張環。因此，該直積由 $T_i\varepsilon_i$ 中元素的一切線性形式組成，係數取自 $\mathfrak r_i e_i$，而等式定義為各系數相等。但藉助 $\mathfrak O\times\mathfrak o$，這恰好成為 $\mathfrak R_i\mathfrak r_i$，因為 $T_i=T_i e_i$、$\mathfrak r_i=\mathfrak r_i\varepsilon_i$，並且 $T_i$ 屬於 $\mathfrak O$ 的一個模基。

\paragraph{IX.}
\emph{$\mathfrak R_i\varepsilon_i\times\mathfrak r_i e_i$ 的差分理想由 $\mathfrak P e_i\varepsilon_i$ 給出，差分商由 $\mathfrak A e_i\varepsilon_i$ 給出。}為看清這一點，把 $\mathfrak P$ 寫成如下形式：
\[
\mathfrak P=\{\ldots,xe_1-\xi\varepsilon_1,\ldots,xe_r-\xi\varepsilon_r,\ldots\}.
\]
於是
\[
\mathfrak P e_i\varepsilon_i
=\{\ldots,xe_i\varepsilon_i-\xi\varepsilon_i e_i,\ldots\},
\]
所以它確實等於 $\mathfrak R_i\varepsilon_i\times\mathfrak r_i e_i$ 的差分理想。而因為（根據 VI）$\mathfrak A e\varepsilon_i$ 等於此環的零理想關於 $\mathfrak P e_i\varepsilon_i$ 的商，所以它被識別為差分商。

\paragraph{X.}
\emph{$\mathfrak R_i\varepsilon_i$ 的不同等於 $\mathfrak D e_i\varepsilon_i$。}$\mathfrak R_i\varepsilon_i$ 的不同通過如下方式產生：在 $\mathfrak A e_i\varepsilon_i$ 中，以每個 $\mathfrak r_i e_i$ 中元素的同構對應元素替換它；也就是每次以 $(xe_i)\varepsilon_i$ 替換 $(\xi\varepsilon_i)e_i$。換言之，由於 $\xi\varepsilon_i$ 也是 $\mathfrak o$ 中的一個 $\xi$，它過渡為 $xe_i$：在 $\mathfrak A e_i\varepsilon_i$ 中作代換 $\xi\to x$——這把 $(\xi\varepsilon_i)e_i$ 變為 $(xe_i)e_i=xe_i$——然後乘以 $\varepsilon_i$。而（根據 VI）原有 $\mathfrak A e_i\varepsilon_i[\xi\to x]=\mathfrak D e_i$，所以 $\mathfrak R_i\varepsilon_i$ 關於 $\mathfrak h e_i\varepsilon_i$ 的不同等於 $\mathfrak D e_i\varepsilon_i$；與此對稱，$\mathfrak r_i e_i$ 的不同等於 $\mathfrak d e_i\varepsilon_i$。

\paragraph{XI.}
\emph{$\mathfrak R_i$ 的不同等於 $\mathfrak D e_i$。}事實上，在該同構下，$\mathfrak R_i$ 的不同通過乘以 $\varepsilon_i$ 過渡為 $\mathfrak R_i\varepsilon_i$ 的不同；而後者等於 $\mathfrak d e_i\varepsilon_i$，所以前者等於 $\mathfrak D e_i$。因此：
\[
\mathfrak D=\mathfrak D e_i+\cdots+\mathfrak D e_r,
\]
\emph{並且 $\mathfrak D e_i$ 是 $\mathfrak R_i$ 的不同；}與此對稱，
\[
\mathfrak d=\mathfrak d\varepsilon_1+\cdots+\mathfrak d\varepsilon_r,
\]
\emph{並且 $\mathfrak d\varepsilon_i$ 是 $\mathfrak r_i$ 的不同，}定理得證。
% 原文：第43篇，印刷頁 13--17 （期刊掃描頁 702--706）
\subsection*{\S{} 5. 域的伽羅瓦擴張環。結構定理。}

\paragraph{1. 基礎環。}
設 $\mathsf{K}$ 是一個\emph{基域 $P$ 上的 $n$ 次域}，並且假定 $\mathsf{K}$ 是可分擴張。設 $\Gamma$ 是相應的伽羅瓦域；它在同構意義下唯一確定，包含 $n$ 個共軛域 $\mathsf{K}^{(i)}$，並且是它們的合成域（複合域），其中規定 $\mathsf{K}=\mathsf{K}^{(1)}$。
用 $\mathfrak K$ 表示 $P$ 的一個擴張，它同 $\mathsf{K}$、因而也同所有 $\mathsf{K}^{(i)}$ 同構，並且相對於 $\mathsf{P}$ 與它們等價；此外還要求其與 $\Gamma$ 的交 $[\mathfrak K,\Gamma]$ 就是 $\mathsf{P}$（亦即 $\Gamma$ 與 $\mathfrak K$ 不存在共同的包絡域）。\srcfn{15)}{這裡當然假定——這並不限制一般性——$K$ 與 $\mathfrak K$ 中不屬於 $P$ 的元素之間不存在任何關係（參見第 1 節第 2 小節、第 2 節第 1 小節以及腳註 11）。}
由於 $\mathfrak K$ 具有一個獨立、有限且包含單位元的 $\mathsf{P}$-模基，根據第 2 節第 1 小節，直積 $\mathfrak K_\Gamma$、$\mathfrak K_{\mathsf{K}}$、$\mathfrak K_{\mathsf{K}^{(i)}}$ 都存在；$\mathfrak K_\Gamma$ 是它的 $n$ 個子環 $\mathfrak K_{\mathsf{K}^{(i)}}$ 的直積。稱 $\mathfrak K_\Gamma$ 為 $\mathfrak K$ 的\emph{伽羅瓦擴張環}。同時，$\mathfrak K_\Gamma$ 還是 $\mathfrak K_{\mathsf{K}}$ 與 $\Gamma$ 關於 $\mathsf{K}$ 的直積，因為 $\mathfrak K$ 的每個獨立 $P$-基也成為 $\mathfrak K_{\mathsf{K}}$ 的獨立 $K$-基；對各共軛域也同樣如此。\srcfn{16)}{對所有共軛域一致成立的定理，通常只對 $\mathfrak K_{\mathsf{K}}$ 表述，也就是省去指標。}這些環的理想之間有如下關係：

\paragraph{定理。}
\emph{$\mathfrak K$ 中的每個理想，都是它在 $\mathfrak K_{\mathsf{K}}$ 或 $\mathfrak K_\Gamma$ 中的擴張的收縮理想；$\mathfrak K_{\mathsf{K}}$ 中的每個理想，都是它在 $\mathfrak K_\Gamma$ 中的擴張的收縮理想。}

\paragraph{證明。}
由於 $\mathsf{P}$ 和 $\mathsf{K}$ 都是域，$\mathfrak K$ 中的每個 $P$-模以及 $\mathfrak K_{\mathsf{K}}$ 中的每個 $K$-模，都分別具有一個獨立有限的 $\mathsf{P}$-模基或 $\mathsf{K}$-模基，而且該基可以分別補充成 $\mathfrak K$ 或 $\mathfrak K_{\mathsf{K}}$ 的這種基。於是該定理由第 2 節第 2 小節中的定理推出。

用 $\mathfrak P_{\mathsf{K}^{(i)}}^{(i)}$ 表示 $\mathfrak K_{\mathsf{K}^{(i)}}$ 的差分理想，即
\[
\mathfrak P_{\mathsf{K}^{(i)}}^{(i)}=\{\ldots,x-\xi^{(i)},\ldots\},
\]
特別地：
\[
\mathfrak P_{\mathsf{K}}=\{\ldots,x-\xi,\ldots\}.
\]
由於 $\xi$ 或 $\xi^{(i)}$ 分別是 $\mathsf{K}$ 或 $\mathsf{K}^{(i)}$ 中通過同構與 $x$ 對應的元素，所以這 $n$ 個差分理想彼此共軛。

用 $\mathfrak A_{\mathsf{K}}$ 或 $\mathfrak A_{\mathsf{K}^{(i)}}$ 表示 $\mathfrak K_{\mathsf{K}}$ 中的差分商 $(0):\mathfrak P_{\mathsf{K}}$ 及其各共軛。它們在 $\mathfrak K_\Gamma$ 中的擴張理想分別記作 $\mathfrak P_\Gamma$ 和 $\mathfrak A_\Gamma$。由本小節關於擴張理想與收縮理想的定理，特別得到：

\emph{差分理想 $\mathfrak P_{\mathsf{K}}$ 和差分商 $\mathfrak A_{\mathsf{K}}$，分別是它們的擴張 $\mathfrak P_\Gamma$ 和 $\mathfrak A_\Gamma$ 的收縮理想。}

\paragraph{2. 伽羅瓦擴張環的直和分解。}
差分理想及差分商同伽羅瓦擴張環結構之間的關係，由下述定理給出。

\paragraph{結構定理。}
\emph{伽羅瓦擴張環 $\mathfrak K_\Gamma$ 的零理想，是 $n$ 個共軛差分理想的擴張之交；$\mathfrak K_\Gamma$ 本身，是 $n$ 個共軛差分商的擴張之直和。$\mathfrak K_{\mathsf{K}^{(i)}}$ 是相應差分理想與差分商的直和，而它的零理想是二者之交。}

\paragraph{證明。}
按每個 $\mathfrak P_\Gamma^{(i)}$ 取的剩餘類環關於 $\Gamma$ 的秩為一，因為每個元素模 $\mathfrak P_\Gamma^{(i)}$ 都同 $\Gamma$ 中的一個元素同餘，而 $\mathfrak P_\Gamma^{(i)}$ 不可能包含 $\Gamma$ 中任何非零元素（因為令 $x\to\xi^{(i)}$ 時 $\mathfrak P_\Gamma^{(i)}$ 消失）。由於可分性，各 $\mathfrak P_\Gamma^{(i)}$ 全都不同，因而它們兩兩互素（任意兩個之和等於 $\mathfrak K_\Gamma$）。因此，理想
\[
\mathfrak P_\Gamma^{(1)},
[\mathfrak P_\Gamma^{(1)},\mathfrak P_\Gamma^{(2)}],\ldots,
[\mathfrak P_\Gamma^{(1)},\ldots,\mathfrak P_\Gamma^{(n)}]
\]
的秩每一步至少降低一個單位；所有這些理想之交成為零理想，而由於兩兩互素，少於全體共軛理想的交必定異於零理想（所以秩在每一步恰好降低一個單位）。由熟知的定理，這便給出 $\mathfrak K_\Gamma$ 如下唯一確定的直和表示：\srcfn{17)}{例如參見 E. Noether，\emph{Abstrakter Aufbau der Idealtheorie in algebraischen Zahl- und Funktionenkörpern}，Math. Ann. 96 (1927)，26--61，第 4 節第 5 小節。}
\begin{align}
\mathfrak K_\Gamma
&=\mathfrak C_\Gamma^{(1)}+\cdots+\mathfrak C_\Gamma^{(n)}
=\Gamma e^{(1)}+\cdots+\Gamma e^{(n)}
=\mathfrak C_\Gamma^{(i)}+\mathfrak P_\Gamma^{(i)}, \tag{1}\\
e&=e^{(1)}+\cdots+e^{(n)},\qquad
\mathfrak C_\Gamma^{(i)}=
[\mathfrak P_\Gamma^{(1)},\ldots,\mathfrak P_\Gamma^{(i-1)},
\mathfrak P_\Gamma^{(i+1)},\ldots].\notag
\end{align}
還須證明 $\mathfrak C_\Gamma$ 等於 $\mathfrak A_\Gamma$。只要證明 $\mathfrak K_{\mathsf{K}}$ 中相應的二者相等，這一結論便隨之成立。在這裡有分解：
\begin{equation}
\mathfrak K_{\mathsf{K}}=\mathfrak C_{\mathsf{K}}+\mathfrak P_{\mathsf{K}}
=\mathfrak K_{\mathsf{K}} e^{(1)}+\mathfrak K_{\mathsf{K}}(e-e^{(1)}).\tag{2}
\end{equation}
事實上，$e^{(1)}$ 屬於 $\mathfrak K_{\mathsf{K}}$；因為 $\mathfrak C_\Gamma^{(1)}$ 在 $\mathfrak P_\Gamma^{(2)},\ldots,\mathfrak P_\Gamma^{(n)}$ 的所有置換下保持不變，所以作為 $\mathfrak C_\Gamma^{(1)}$ 中單位元分量的 $e^{(1)}$ 也同樣保持不變。因此，用 $\mathfrak K_{\mathsf{K}}$ 的一個獨立 $K$-基來表示 $e^{(1)}$——這是可能的，因為該基同時也是 $\mathfrak K_\Gamma$ 的獨立 $\Gamma$-基——則各個係數在 $\mathsf{K}^{(2)},\ldots,\mathsf{K}^{(n)}$ 的所有置換下都保持不變，因而屬於 $\mathsf{K}$。由於 $(e^{(1)})^2=e^{(1)}$，和式表示
\[
\mathfrak K_{\mathsf{K}}=\mathfrak K_{\mathsf{K}} e^{(1)}+\mathfrak K_{\mathsf{K}}(e-e^{(1)})
\]
是直和；它在 $\mathfrak K_\Gamma$ 中的擴張仍然是直和，故由 (1) 得到 $\mathfrak K_\Gamma=\mathfrak C_\Gamma^{(1)}+\mathfrak P_\Gamma^{(1)}$。於是根據第 2 節第 2 小節的末尾部分，有
\[
\mathfrak K_{\mathsf{K}} e^{(1)}=[\mathfrak C_\Gamma,\mathfrak K_{\mathsf{K}}]
=[\Gamma e^{(1)},\mathfrak K_{\mathsf{K}}]=\mathsf{K}e^{(1)}=\mathfrak C_{\mathsf{K}};
\]
\[
\mathfrak K_{\mathsf{K}}(e-e^{(1)})=[\mathfrak P_\Gamma,\mathfrak K_{\mathsf{K}}]=\mathfrak P_{\mathsf{K}},
\]
從而 (2) 得證。

由 (2) 又可推出 $\mathfrak C_{\mathsf{K}}=\mathfrak A_{\mathsf{K}}$，進而 $\mathfrak C_\Gamma=\mathfrak A_\Gamma$。事實上，由 $e^{(1)}(e-e^{(1)})=0$ 可得 $\mathfrak C_{\mathsf{K}}\mathfrak P_{\mathsf{K}}=0$；反過來，$b\mathfrak P_{\mathsf{K}}=0$ 也給出 $b(e-e^{(1)})=0$，所以 $b=be=be^{(1)}$，於是
\[
\mathfrak C_{\mathsf{K}}=(0):\mathfrak P_{\mathsf{K}}=\mathfrak A_{\mathsf{K}},
\]
並且擴張後有 $\mathfrak C_\Gamma=\mathfrak A_\Gamma$。因此：
\begin{align}
\mathfrak K_\Gamma
&=\mathfrak A_\Gamma^{(1)}+\cdots+\mathfrak A_\Gamma^{(n)}
=\Gamma e^{(1)}+\cdots+\Gamma e^{(n)};
&0&=[\mathfrak P_\Gamma^{(1)},\ldots,\mathfrak P_\Gamma^{(n)}]; \tag{3}\\
\mathfrak K_{\mathsf{K}^{(i)}}
&=\mathfrak A_{\mathsf{K}^{(i)}}+\mathfrak P_{\mathsf{K}^{(i)}}^{(i)}
=\mathsf{K}^{(i)}e^{(i)}+\mathfrak P_{\mathsf{K}^{(i)}}^{(i)};
&0&=[\mathfrak A_{\mathsf{K}^{(i)}},\mathfrak P_{\mathsf{K}^{(i)}}^{(i)}].\notag
\end{align}
至此，該定理的各部分全都得證。還應指出，這是一個\emph{完全不變的結構定理}，因為其中出現的所有構造（$\mathfrak K_{\mathsf{K}}$、$\mathfrak K_\Gamma$、$\mathfrak P_{\mathsf{K}}$、$\mathfrak A_{\mathsf{K}}$）都不依賴於基來定義。\srcfn{18)}{這裡也可以藉助定義方程以及相應多項式的線性因式分解來表述，但這樣的表述並非不變的。}

\paragraph{3. 由共軛元素給出的 $\mathfrak K$ 的分量表示。}
與共軛差分理想及差分商相對應，$\mathfrak K$ 的元素可通過共軛分量來表示：

\paragraph{定理。}
\emph{設
\[
x=\xi^{(1)}e^{(1)}+\cdots+\xi^{(n)}e^{(n)}
\]
是 $\mathfrak K$ 中元素 $x$ 的 $\mathfrak K$-分量表示，則這 $n$ 個分量彼此共軛。特別地，$\xi^{(1)},\ldots,\xi^{(n)}$ 正是通過 $\mathfrak K$ 到 $\mathsf{K}^{(1)},\ldots,\mathsf{K}^{(n)}$ 的各同構與 $x$ 對應的 $n$ 個共軛值。}

$e^{(1)},\ldots,e^{(n)}$ 彼此共軛，可由第 2 小節中的關係式 (3) 看出：把 $\mathfrak P_{\mathsf{K}}$ 映到 $\mathfrak P_{\mathsf{K}^{(i)}}^{(i)}$ 的 $\Gamma$ 的自同構，也同時把 $e^{(1)}$ 映到 $e^{(i)}$，而後者確實是 $\mathfrak K_{\mathsf{K}^{(i)}}$ 中的元素。同一組關係式還給出 $xe^{(i)}=\xi^{(i)}e^{(i)}$；這是因為 $\mathfrak P_{\mathsf{K}^{(i)}}^{(i)}e^{(i)}=0$，從而
\[
(x-\xi^{(i)})e^{(i)}=0
\]
（也可由以下事實推出：作為域的 $\mathfrak K$ 同它的分量同構，而該分量就是 $\mathsf{K}^{(i)}e^{(i)}$，其中 $x$ 與 $\xi^{(i)}e^{(i)}$ 同構對應）。

由該定理可得：

\paragraph{推論。}
\emph{設 $\mathfrak B$ 是 $\mathfrak K$ 中的一個絕對模（即它在包含任意兩個元素的同時也包含二者之差），並設 $\mathfrak B^{(1)},\ldots,\mathfrak B^{(n)}$ 是它在 $\mathfrak K_\Gamma$ 中的各分量，則 $\mathfrak B$ 等於交 $[\mathfrak K,\mathfrak B^{(1)}+\cdots+\mathfrak B^{(n)}]$。}

事實上，$\mathfrak B$ 包含在這個交中；反過來，若 $\mathfrak B^{(1)}+\cdots+\mathfrak B^{(n)}$ 中的一個元素還屬於 $\mathfrak K$，則根據該定理，它必為共軛分量之和，因而是 $\mathfrak B$ 中的元素。

\paragraph{4. $\mathfrak K$ 的互補基及其同單位元分量的關係。}
從不變結構轉到基表示時，這一結構把各個基歸併為互補基對：

\paragraph{定理。}
\emph{$\mathfrak K$ 的每個 $\mathsf{P}$-基 $t_1,\ldots,t_n$，都對應着 $\mathfrak K$ 中的一個互補基 $T_1,\ldots,T_n$，而後者的互補基又是原來的 $t$ 基。$\mathsf{K}^{(i)}$ 中同互補基 $T_1,\ldots,T_n$ 同構的基 $\mathsf{A}_1^{(i)},\ldots,\mathsf{A}_n^{(i)}$，被定義為單位元分量 $e^{(i)}$ 在由 $t$ 給出的基表示
\[
e^{(i)}=\mathsf{A}_1^{(i)}t_1+\cdots+\mathsf{A}_n^{(i)}t_n
\]
中的係數。若 $\alpha_1^{(i)},\ldots,\alpha_n^{(i)}$ 是與各 $t_i$ 對應的基，則 $\alpha$ 與 $\mathsf{A}$ 的矩陣互逆。若 $\mathfrak K$ 的兩個基 $(s)$ 與 $(t)$ 通過變換
\[
(s)=(t)P
\]
相聯繫，則相應互補基之間有逆變關係
\[
(S)=P^{-1}(T).
\]}

\paragraph{證明。}
設 $e^{(i)}=A_1^{(i)}t_1+\cdots+A_n^{(i)}t_n$ 是單位元的 $n$ 個分量 $e^{(i)}$ 的基表示，其中取定了 $\mathfrak K$ 的一個 $P$-基 $t_1,\ldots,t_n$；該基於是也成為 $\mathfrak K_\Gamma$ 的 $\Gamma$-基。根據第 2、3 小節，$A_1^{(i)},\ldots,A_n^{(i)}$ 屬於 $K^{(i)}$，而不同的 $e^{(i)}$ 對應於彼此共軛的係數系 $A$。若 $\alpha_1^{(i)},\ldots,\alpha_n^{(i)}$ 表示 $K^{(i)}$ 中與 $t_1,\ldots,t_n$ 同構對應的基，則還有以下關係：
\begin{equation}
e^{(i)}[t\to\alpha^{(i)}]=e;
\qquad e^{(i)}[t\to\alpha^{(j)}]=0\quad\text{當 }i\ne j. \tag{1}
\end{equation}
因為 $\mathfrak P_{\mathsf{K}^{(i)}}^{(i)}$ 包含元素 $t_1-\alpha_1^{(i)},\ldots,t_n-\alpha_n^{(i)}$——它們構成 $\mathfrak P_{\mathsf{K}^{(i)}}^{(i)}$ 的一個非獨立 $P$-模基——所以在代入 $t=\alpha^{(i)}$ 時該理想消失。於是由
\[
e\equiv e^{(i)}\;(\mathfrak P_{\mathsf{K}^{(i)}}^{(i)})
\]
可得 (1) 的第一個關係；第二個關係則由 $i\ne j$ 時 $e^{(i)}$ 屬於所有 $\mathfrak P_{\mathsf{K}^{(j)}}^{(j)}$ 而得到。按係數展開後，關係式 (1) 可寫成矩陣形式：
\begin{equation}
\begin{pmatrix}
\alpha_1^{(1)}&\ldots&\alpha_n^{(1)}\\
\vdots&&\vdots\\
\alpha_1^{(n)}&\ldots&\alpha_n^{(n)}
\end{pmatrix}
\begin{pmatrix}
A_1^{(1)}&\cdots&A_1^{(n)}\\
\vdots&&\vdots\\
A_n^{(1)}&\cdots&A_n^{(n)}
\end{pmatrix}=E_n. \tag{2}
\end{equation}
其中 $E_n$ 表示單位矩陣；由此，$\mathsf{A}$ 已由 $\alpha$ 唯一確定。關係式 (2) 表明兩個矩陣的行列式都不為零。它還表明，$\mathsf{A}_1^{(i)},\ldots,\mathsf{A}_n^{(i)}$ 關於 $\mathsf{P}$ 線性無關；因為若存在一個線性相關關係，它也會對各共軛系成立，從而降低 $\mathsf{A}$ 矩陣的秩，這與 $E_n$ 的秩為 $n$ 相矛盾。因此，每組 $\mathsf{A}_1^{(i)},\ldots,\mathsf{A}_n^{(i)}$ 都構成 $\mathsf{K}^{(i)}$ 的一個基；這些基彼此共軛，所以它們都同 $\mathfrak K$ 中的同一個基 $T_1,\ldots,T_n$ 同構，而根據第 3 小節，$T_\lambda$ 由
\[
T_\lambda=\mathsf{A}_\lambda^{(1)}e^{(1)}+\cdots+\mathsf{A}_\lambda^{(n)}e^{(n)}
\]
確定。稱基 $T_1,\ldots,T_n$ 為 $t_1,\ldots,t_n$ 的互補基，相應地，稱 $\mathsf{A}_1^{(i)},\ldots,\mathsf{A}_n^{(i)}$ 為 $\alpha_1^{(i)},\ldots,\alpha_n^{(i)}$ 的互補基。由於 $\mathsf{A}$ 已通過 (2) 唯一確定，而這一關係——通過轉到轉置矩陣——又從 $T$ 的基、亦即從 $\mathsf{A}$ 確定 $\alpha$，所以原來的 $t$ 基就是 $T$ 基的互補基。因此，單位元各分量所給出的對應把所有基都分成互補基對。若 $s,S$ 是另一對互補基，滿足
\[
e^{(i)}=\mathsf{B}_1^{(i)}s_1+\cdots+\mathsf{B}_n^{(i)}s_n
\qquad\text{且}\qquad (s)=(t)P,
\]
則由
\[
\mathsf{A}_1t_1+\cdots+\mathsf{A}_nt_n=\mathsf{B}_1s_1+\cdots+\mathsf{B}_ns_n
\]
——在代入 $(s)=(t)P$ 後作為關於 $t$ 的恆等式——可知，$\mathsf{B},\mathsf{A}$ 的變換與 $s,t$ 的變換互為逆變。因而由同構可知，$S$ 和 $T$ 也有同樣的關係。至此，該定理的各部分全都得證。

\paragraph{注。}
若 $c=C_1t_1+\cdots+C_nt_n$ 是 $\mathfrak K$ 中一個元素關於某個 $\mathsf{P}$-基的表示，而 $c=c_1T_1+\cdots+c_nT_n$ 是同一元素關於互補基的表示，則
\[
C_\lambda=\operatorname{Sp}(cT_\lambda),
\qquad c_\lambda=\operatorname{Sp}(ct_\lambda),
\]
其中 $\operatorname{Sp}(x)$ 照例表示各共軛之和：
\[
\operatorname{Sp}(x)=\xi^{(1)}+\xi^{(2)}+\cdots+\xi^{(n)}
\]
對 $\mathfrak K$ 中任意 $x$ 都如此。事實上，
\[
c=\gamma^{(1)}e^{(1)}+\cdots+\gamma^{(n)}e^{(n)}
=\sum\limits_\lambda\gamma^{(1)}\mathsf{A}_\lambda^{(1)}t_\lambda+
\cdots+\sum\limits_\lambda\gamma^{(n)}\mathsf{A}_\lambda^{(n)}t_\lambda,
\]
由此比較係數便得到 $C_\lambda=\operatorname{Sp}(cT_\lambda)$。第二個關係可同樣通過交換 $t$ 與 $T$ 得到。\srcfn{19)}{本小節的討論表明，不變結構定理實質上比非不變的基定理簡單。這裡只是為了說明它們同已有結果的聯繫，才使用後者。事實上，伽羅瓦擴張環的分解以及單位元分量的引入，已經完成了互補基（下一節中的互補模）所完成的一切。下一節中，不同與互補模之間的聯繫也只是為了歸入已有理論才給出。對第 5 小節中拉格朗日插值公式的討論同樣如此。}

\paragraph{5. 以定義方程為基礎。拉格朗日插值公式。}
[編者注：本小節在手稿中仍未填寫。]
% 原文：第43篇，印刷頁 17--21 （期刊掃描頁 706--710）
\subsection*{\S{} 6. 一個階的不同。}

現在要在整性的意義下加強那些只涉及域及其伽羅瓦擴張環的結構定理：在 (P) 中指定某些子環，並在 \(\mathfrak K\) 和 (K) 中指定關於這些子環為整的元素，也就是各個階。§ 5 的全部記號繼續沿用。\srcfn{20)}{關於元素相對於一個環為整以及整閉的概念，參見腳註 17 所引論文的 § 1。}

\paragraph{1. 基礎環。}
設 \(\mathfrak h\) 是 (P) 的一個子環，使 (P) 成為 \(\mathfrak h\) 的分式域；再設 \(\mathfrak h\) 在 (P) 中\emph{整閉}（(P) 中每個相對於 \(\mathfrak h\) 為整的元素都屬於 \(\mathfrak h\)）。再令 \(\mathfrak O\) 表示 \(\mathfrak K\) 中的一個 \(\mathfrak h\)-階，也就是 \(\mathfrak K\) 的一個包含 \(\mathfrak h\) 的子環，其分式域為 \(\mathfrak K\)，並且具有一個有限的、不一定獨立的 \(\mathfrak h\)-模基。於是 \(\mathfrak O\) 的所有元素都是（相對於 \(\mathfrak h\)）整的；這個 \(\mathfrak h\)-模基同時成為 \(\mathfrak K\) 的一個——通常相關的——(P)-模基。令 \(\mathfrak o\) 表示 (K) 中與 \(\mathfrak O\) 同構的階，\(\mathfrak o^{(i)}\) 表示 \(\mathsf{K}^{(i)}\) 中的各個共軛階。\emph{令 \(\mathfrak g\) 為 \(\Gamma\) 中由這 (n) 個共軛階生成的環；}那麼 \(\mathfrak g\) 也是一個階。事實上，\(\mathfrak g\) 包含 \(\mathfrak h\)。各個 \(\mathfrak o^{(i)}\) 的 \(\mathfrak h\)-基中所有元素的乘積構成 \(\mathfrak g\) 的一個 \(\mathfrak h\)-基，同時構成 \(\Gamma\) 的一個 (P)-基。

\emph{直積 \(\mathfrak O_{\mathfrak o}\) 和 \(\mathfrak O_{\mathfrak g}\) 存在，}這由 § 2, 3 的充分判據得出（其中的環 \(\mathfrak f\) 要用 (P) 替換）。在 \(\mathfrak K_\Gamma\) 中，關於 \(\mathfrak h\) 的直積 \(\mathfrak O_{\mathsf{P}}=\mathfrak K\) 首先存在。因為 \(\mathfrak h\) 包含在交 \([\mathfrak O,\mathsf{P}]\) 中；反過來，交中的每個元素作為 \(\mathfrak O\) 的元素都是整的，因而由於 \(\mathfrak h\) 在 (P) 中整閉，它作為 (P) 的元素也屬於 \(\mathfrak h\)。取 \(\mathfrak O\) 中任意 (n) 個線性無關元素，它們以 (P) 中元素為係數的線性組合便給出作為直積的 \(\mathfrak K\)。同樣有
\[
\mathfrak o_{\mathsf{P}}=\mathsf{K}
\qquad\text{以及}\qquad
\mathfrak g_{\mathsf{P}}=\Gamma;
\]
由關於 \(\mathfrak h\) 的直積 \(\mathfrak K_{\mathsf{K}}=\mathfrak O_{\mathsf{P}}\times\mathfrak o_{\mathsf{P}}\)，按 § 2, 3 的判據可得關於 \(\mathfrak h\) 的 \(\mathfrak O_{\mathfrak o}\) 存在；同樣，\(\mathfrak K_\Gamma=\mathfrak O_{\mathsf{P}}\times\mathfrak g_{\mathsf{P}}\) 給出 \(\mathfrak O_{\mathfrak g}\) 的存在。隨着 \(\mathfrak O_{\mathfrak o}\) 的存在，各個共軛環 \(\mathfrak O_{\mathfrak o^{(i)}}\) 也存在。稱 \(\mathfrak O_{\mathfrak g}\) 為\emph{階 \(\mathfrak O\) 的伽羅瓦擴張環}。

因此，\(\mathfrak O\) 和 \(\mathfrak o\) 相對於 \(\mathfrak h\) 的不同都有定義；各個共軛不同也同樣有定義。令
\[
\mathfrak P=\{\ldots,y-\eta,\ldots\}
\]
表示 \(\mathfrak O_{\mathfrak o}\) 的差分理想，並令
\[
\mathfrak P^{(i)}=\{\ldots,y-\eta^{(i)},\ldots\}
\]
表示 \(\mathfrak O_{\mathfrak o^{(i)}}\) 中的各個共軛差分理想；再令
\[
\mathfrak A^{(i)}=(0):\mathfrak P^{(i)}
\quad\text{——商在 \(\mathfrak O_{\mathfrak o^{(i)}}\) 中取得——}
\]
表示各個共軛\emph{差分商}，特別地，
\[
\mathfrak A=(0):\mathfrak P\quad\text{在 }\mathfrak O_{\mathfrak o}\text{ 中。}
\]
再令
\[
\mathfrak D=\mathfrak A[\eta\to y]
\]
表示 \(\mathfrak O\) 的不同，並令
\[
\mathfrak d^{(i)}=\mathfrak A[y\to\eta^{(i)}]
\]
表示這 (n) 個共軛階 \(\mathfrak o^{(i)}\) 的不同，特別地，
\[
\mathfrak d=\mathfrak A[y\to\eta]
\]
表示 \(\mathfrak o\) 的不同。\srcfn{21)}{例如，當 \(\mathfrak h\) 是一個數域的主階，特別是 \(\mathfrak h\) 為全體有理整數所成的系統時，\(\mathfrak h\) 在 (P) 中整閉這一假設成立；因此任意階的不同和相對不同，包括相對域中的情形，都有定義。再者，\(\mathfrak h\) 可以是 (n) 個不定元的多項式環，這便給出 (n) 個不定元的代數函數域中的不同。}

\paragraph{2. 結構定理。}
與域的伽羅瓦擴張環的結構定理（§ 5, 2）相對應的是：

\paragraph{階的伽羅瓦擴張環結構定理。}
\emph{相對於階和域的差分理想與差分商具有如下關係：\(\mathfrak P_{\mathsf{K}}\) 是 \(\mathfrak O_{\mathfrak o}\) 的差分理想 \(\mathfrak P\) 的擴張；同樣，\(\mathfrak A_{\mathsf{K}}\) 是 \(\mathfrak O_{\mathfrak o}\) 的差分商 \(\mathfrak A\) 的擴張，而 \(\mathfrak A\) 是 \(\mathfrak A_{\mathsf{K}}\) 的限制。差分商、不同和單位元分量之間具有關係}
\[
\mathfrak A=\mathfrak d e^{(1)},
\qquad
\mathfrak A^{(i)}=\mathfrak d^{(i)}e^{(i)},
\]
\emph{因而}
\[
\mathfrak d=[\mathfrak o,\mathfrak A^{(1)}+\cdots+\mathfrak A^{(n)}].
\]
\emph{所以，\(\mathfrak o\) 的不同 \(\mathfrak d\) 由 (K) 中所有滿足 \(xe^{(1)}\in\mathfrak O_{\mathfrak o}\) 的元素 (x) 組成；它不是零理想。}

\(\mathfrak P_{\mathsf{K}}\) 是 \(\mathfrak P\) 的擴張，是因為 \(\mathfrak K=\mathfrak O_P\)，同樣 \(K=\mathfrak o_P\)，所以 \(\mathfrak P_{\mathsf{K}}\) 中每個 \(x-\xi\) 都能由 \(\mathfrak P\) 中的各個 \(y-\eta\)，以 \(K\) 中元素為係數作線性表示。由此還可得，\(\mathfrak A_{\mathsf{K}}=(0):\mathfrak P_{\mathsf{K}}\) 也可定義為 \(\mathfrak K_{\mathsf{K}}\) 中的商 \((0):\mathfrak P\)。因此，\(\mathfrak O_{\mathfrak o}\) 中的 \(\mathfrak A=(0):\mathfrak P\) 由 \(\mathfrak O_{\mathfrak o}\) 中屬於 \(\mathfrak A_{\mathsf{K}}\) 的全部元素組成，即 \(\mathfrak A=[\mathfrak O_{\mathfrak o},\mathfrak A_{\mathsf{K}}]\)，所以它是 \(\mathfrak A_{\mathsf{K}}\) 的限制。由於 \(\mathfrak A_{\mathsf{K}}=\mathsf{K}e^{(1)}\)，於是 \(\mathfrak A=\mathfrak z e^{(1)}\)，其中 \(\mathfrak z\) 是 \(K\) 中的一個 \(\mathfrak o\)-模，因為 \(\mathfrak A\) 本身也是一個 \(\mathfrak o\)-模。但根據不同的定義，
\[
\mathfrak d=\mathfrak A[x\to\xi]
=(\mathfrak z e^{(1)})[x\to\xi]
=\mathfrak z\bigl(e^{(1)}[x\to\xi]\bigr)=\mathfrak z e=\mathfrak z
\]
（參見 § 5, 4, (1)）。所以 \(\mathfrak z=\mathfrak d\)；差分商 \(\mathfrak A\) 成為 \(\mathfrak o\) 的不同 \(\mathfrak d\) 在 (K) 中的分量，各個共軛情形也相應成立。於是，根據 § 5, 3 的推論，便得到上面關於 \(\mathfrak d\) 的關係。並且
\[
\mathfrak A=\mathfrak d e^{(1)}=[\mathsf{K}e^{(1)},\mathfrak O_{\mathfrak o}]
\]
表明，\(\mathfrak d\) 由 (K) 中所有使 \(\varkappa e^{(1)}\) 屬於 \(\mathfrak O_{\mathfrak o}\) 的元素 (x) 組成。由此 \(\mathfrak d\) 不是零理想；因為 \(e^{(1)}\) 作為 \(\mathfrak K_{\mathsf{K}}=\mathfrak O_{\mathsf{K}}\) 的元素，是 \(\mathfrak O\) 中有限多個元素以 (K) 中係數作成的線性組合，因而是 \(\mathfrak O_{\mathfrak o}\) 中一個元素除以 \(\mathfrak o\) 中一個元素（甚至可以取自 \(\mathfrak h\)）所得的商；這個分母給出 \(\mathfrak d\) 的一個非零元素。於是，\(\mathfrak d\) 關於 (P) 的秩為 (n)，或者說 \(\mathfrak d_{\mathsf{P}}=\mathsf{K}\)，從而 \(\mathfrak A_{\mathsf{K}}\) 是 \(\mathfrak A\) 的擴張。

\begin{center}
\textbf{以下為提綱。}
\end{center}

\paragraph{3. 在基獨立時同互補模的關係。}
若 \(t_1,\ldots,t_n\) 是 \(\mathfrak O\) 的一個獨立 \(\mathfrak h\)-基，那麼，只要變換矩陣 (P) 的元素屬於 \(\mathfrak h\) 且為幺模，\(\mathfrak K\) 的每個基 \((s)=(t)P\) 也是如此。此時，\((S)=P^{-1}(T)\) 表明，除 (T) 外，(S) 也是同一個 \(\mathfrak h\)-模的基；這個模稱為 \(\mathfrak O\) 的互補模。

現在，2 中結構定理的最後一部分說明，\(\mathfrak d\) 等於 (K) 中的商
\[
\mathfrak o:\mathfrak c,
\]
其中 \(\mathfrak e\) 表示互補模。事實上，設
\[
e^{(1)}=\mathsf{A}_1t_1+\cdots+\mathsf{A}_nt_n,
\]
其中 \(\mathsf{A}_\lambda\) 是 \(\mathfrak e\) 的一個基；那麼 \(\mathfrak d\) 由所有使 \(\delta e^{(1)}\) 屬於 \(\mathfrak O_{\mathfrak o}\) 的元素 \(\delta\) 組成，也就是說，要求 \(\delta A_i\in\mathfrak o\)，所以 \(\mathfrak d=\mathfrak o:\mathfrak e\)。

\paragraph{特殊情形。}
正則階以 \(e,z,\ldots,z^{n-1}\) 為模基時：
\[
\mathfrak d=\{f'(z)\}.
\]

\paragraph{注。}
由 § 5, 4 的注還可得到 \(\mathfrak O\) 的互補模 \(\mathfrak E\) 的另一個熟知定義。\emph{\(\mathfrak E\) 由 \(\mathfrak K\) 中所有這樣的元素 (c) 組成：跡 \(\operatorname{Sp}(ct)\)，也就是任取 \(\mathfrak O\) 中的 (o) 時的跡 \(\operatorname{Sp}(co)\)，都是整的，即屬於 \(\mathfrak h\)；}關於 \(\mathfrak o\) 和 (K) 的 \(\mathfrak e\) 也相應如此。

事實上，設 \(c=c_1T_1+\cdots+c_nT_n\)，其中 \(c_i\in\mathfrak h\)，這是 \(\mathfrak E\) 中元素 (c) 關於與 (t) 互補的 \(\mathfrak K\) 的基所作的表示。那麼 \(c_\lambda=\operatorname{Sp}(ct_\lambda)\)，故 \(c_i\in\mathfrak h\)；反過來，由 \(\operatorname{Sp}(ct_\lambda)\in\mathfrak h\) 可知 \(c_i\in\mathfrak h\)，因而 \(c\in\mathfrak C\)。

\paragraph{補充。}
\emph{若階中存在獨立的 \(\mathfrak h\)-模基，則差分商 \(\mathfrak A^{(i)}\) 是除 \(\mathfrak P^{(i)}\) 外所有差分理想的交；更確切地說，}
\[
\mathfrak A^{(i)}=
[\mathfrak O_{\mathfrak o^{(i)}},
\mathfrak P_{\mathfrak g}^{(1)},\ldots,
\mathfrak P_{\mathfrak g}^{(i-1)},
\mathfrak P_{\mathfrak g}^{(i+1)},\ldots,
\mathfrak P_{\mathfrak g}^{(n)}].
\]
因為根據 § 2, 2 的加強形式，此時 \(\mathfrak P\) 也是 \(\mathfrak P_{\mathsf{K}}\) 的限制理想；事實上，若 \(e,t_2,\ldots,t_n\) 是 \(\mathfrak O\) 的一個模基，因而也是 \(\mathfrak O_{\mathfrak o}\) 的一個模基，那麼
\[
e,t_2-\alpha_2,\ldots,t_n-\alpha_n
\]
也是一個模基；而 \(t_2-\alpha_2,\ldots,t_n-\alpha_n\) 同時是 \(\mathfrak P_{\mathsf{K}}\) 的一個 (K)-模基。因此，§ 2, 2（補充）的假設得到滿足。現在，原有
\[
\mathfrak A_\Gamma=[\mathfrak P_\Gamma^{(2)},\ldots,\mathfrak P_\Gamma^{(n)}];
\qquad
\mathfrak A_\mathfrak g=[\mathfrak A_\Gamma,\mathfrak O_\mathfrak g]
=[\ldots,[\mathfrak P_\Gamma^{(\lambda)},\mathfrak O_\mathfrak g],\ldots]
=[\mathfrak P_\mathfrak g^{(2)},\ldots,\mathfrak P_\mathfrak g^{(n)}].
\]
而 \(\mathfrak A=[\mathfrak A_\mathfrak g,\mathfrak O_{\mathfrak o}]\)；所以
\[
\mathfrak A=[\mathfrak O_{\mathfrak o},\mathfrak P_\mathfrak g^{(2)},\ldots,\mathfrak P_\mathfrak g^{(n)}].
\]

\paragraph{4. 在基本方程存在時同其微分商的關係。}
把階 \(\mathfrak O\) 以 (n) 個不定元 \(u_1,\ldots,u_n\) 擴張，也就是形成直積 \(\mathfrak O_{\mathfrak U}\)，其中 (U) 表示多項式環 \(\mathfrak h[u_1,\ldots,u_n]\)，故 \(\mathfrak O_{\mathfrak U}=\mathfrak O[u_1,\ldots,u_n]\)。於是，\(\mathfrak O\) 中每個理想 \(\mathfrak C\) 的擴張 \(\mathfrak C_{\mathfrak U}\)，由 (u) 的所有係數屬於 \(\mathfrak C\) 的多項式組成；反過來，\(\mathfrak C\) 又被刻畫為 \(\mathfrak C_{\mathfrak U}\) 的所有係數所成的系統。

現在假定有一個基本方程；也就是說，當過渡到由 \(\mathfrak h[u]\) 中一個或多個本原多項式 (H(u)) 所確定的商環時，\(e,U,\ldots,U^{n-1}\) 應構成一個 \(\mathfrak h\)-模基。於是，理想 \(\mathfrak C_{\mathfrak U}\) 過渡為由 \(\mathfrak C_{\mathfrak U}\) 乘以單位（(H(u)) 及其正負整數次冪）而得到的系統。因此，為了重新得到 \(\mathfrak C\)，可以先通過乘以這些單位回到 \(\mathfrak C_{\mathfrak U}\)（在 (u) 中為整），再由此回到 \(\mathfrak C\)。特別地，令 \(\mathfrak O\) 的不同 \(\mathfrak D\) 過渡為 \(\mathfrak D_{\mathfrak U}\)，其中 \(\mathfrak D_{\mathfrak U}\) 是 \(\mathfrak O_{\mathfrak U}\) 的不同。但在 \(\mathfrak O_{\mathfrak U}\) 中——通過用本原的 (H(u)) 取商——存在一個定義方程 (G(U)=0)。所以 (G'(U)) 成為不同的基多項式，也就是說，不同由 (G'(U)) 的係數導出。（還必須更精確地證明，\((G'(U))\) 確實只包含 \(\mathfrak D_{\mathfrak U}\) 而不包含更多東西；這也許只在 \(\mathfrak o\) 在 (K) 中整閉時成立。因為乘以 (H(u)) 的因子，還可能從 \(\mathfrak O_{\mathfrak U}\) 中產生更多仍位於 \(\mathfrak D_{\mathfrak U}\) 內的元素。在整閉的情形下〔五條公理〕，這可由係數理想〔內容〕相乘的定理推出。於是，通過考察基本方程模 (p)，便可建立分歧理論。）

\subsection*{\S{} 7. 主階模 \(p^t\) 的分歧理論。}

設給定一個數域，\(\mathfrak h\) 為其整數；或者也可以給定一個相對於主階按某個素理想所得商環的相對域，此時 (p) 表示這個商環的素理想的基元。添入不定元 \(u_1,\ldots,u_n\)；於是，無論主階還是按 \(p^t\) 取的剩餘類環，都具有 \(\mathfrak h\)-模基 \(e,U,\ldots,U^{n-1}\)，因此兩種情形中的不同都由 \(G'(u)\) 給出。\emph{主階的不同模 \(p^t\) 後，過渡為按 \(p^t\) 取的剩餘類環的不同}（這對任意階是否也成立？）。如果
\[
p=\mathfrak p_1^{\varrho_1}\cdots\mathfrak p_r^{\varrho_r},
\]
那麼按 \(p^t\) 取的剩餘類環 \(\mathfrak R\) 是一個直和，其分量同按 \(\mathfrak p_i^{t\varrho_i}\) 取的剩餘類環同構，即
\[
\mathfrak R=\mathfrak R e_1+\cdots+\mathfrak R e_r;
\]
其中，當 \(i\ne k\) 時，\(e_k\) 被 \(\mathfrak p_i^{t\varrho_i}\) 整除。因此，若
\[
\mathfrak d[p^t]=\mathfrak d[p^t]e_1+\cdots+\mathfrak d[p^t]e_r
=\mathfrak d_1+\cdots+\mathfrak d_r
\]
且 (t) 足夠大，那麼 \(\mathfrak p_i\) 在 \(\mathfrak d_i\) 中出現的冪次，同它在 \(\mathfrak d[p^t]\) 中、因而同它在 \(\mathfrak d\) 中出現的冪次相同。但根據 § 4，\(\mathfrak d_i\) 是 \(\mathfrak R_i\) 的不同；所以只須考察按 \(\mathfrak p^{\varrho t}\) 取的剩餘類環的不同。選取 \(\xi\)（《數論報告》定理 29、30），使 \(\xi\) 為整數，並且
\[
\varphi(\xi)\equiv0(\mathfrak p),
\quad \not\equiv0(\mathfrak p^2),
\]
其中 \(\varphi(x)\) 是模 (p) 的一個 (f) 次素函數。對 \(t\geqq2\)，把這個 \(\xi\) 換成 \(\xi^*=\xi e_1\)，其中 \(e_1\) 是剩餘類環 \(\mathfrak R\) 模 \(p^t\) 的分解中出現的第一個冪等元；那麼，除 \(\varphi(\xi^*)\equiv0(\mathfrak p)\)、\(\not\equiv0(\mathfrak p^2)\) 外，甚至還有 \(\mathfrak p_1\equiv(p,\varphi(\xi^*))\)。於是，關於模 \(p^t\) 的整數的一組基由 \(\xi^\mu\varphi(\xi)^\nu\) 給出，其中 \(\mu=0,1,\ldots,f-1\)，\(\nu=0,1,\ldots,t-1\)。一個獨立的 \(\mathfrak h\)-基（模 \(p^t\) 的剩餘類）由
\begin{equation}
\xi^{\mu_0}\varphi(\xi)^{\nu_0}
\quad\text{其中}\quad
\mu_0=0,1,\ldots,f-1;
\quad \nu_0=0,1,\ldots,\varrho-1. \tag{1}
\end{equation}
給出。事實上，\((\varphi(\xi),\mathfrak p^{\varrho t})=\mathfrak p\)；所以，由於剩餘類環中 \((p)=\mathfrak p^\varrho\)，從主理想過渡到元素便有 \(\varphi(\xi)=pM(\xi)\)，其中 \(M(\xi)\) 是剩餘類環中的單位，也就是說 \(M(\xi)\not\equiv0(\mathfrak p)\)。因此（由於 \(p^t=0\)），可以逐次用較低次冪表示 \(\xi^{\varrho t}\)，先在 \(M(\xi)\) 中，再在剩餘類環中表示；所以 (1) 確實是一組基。關於這組基的定義方程為
\[
F(x)=\varphi(x)^\varrho+pM(x),
\]
其中 (M(x)) 模 (p) 不被 \(\varphi(x)\) 整除（Ö. Ore, Math. Ann. 96, S. 348），因而
\[
M(\xi)\not\equiv0(\mathfrak p).
\]
但這組基也是獨立的；因為由
\[
a_0(\xi)+a_1(\xi)\varphi(\xi)+\cdots+a_{\varrho-1}(\xi)\varphi(\xi)^{\varrho-1}
\equiv0(\mathfrak p^{t\varrho})
\]
先可推出模 (p) 同餘於零；從而 \(a_0(\xi)\equiv0(\mathfrak p)\)，於是
\[
a_0(x)\equiv0(p);
\]
約去 (p) 後逐次得到
\[
a_0(x)\equiv0(p),\ldots,a_{\varrho-1}(x)\equiv0(p),
\]
再約去 (p) 後得到
\[
a_0(\alpha)\equiv0(p^2),\ldots,a_{\varrho-1}(\alpha)\equiv0(p^2),
\]
如此繼續，直至
\[
a_0(\alpha)\equiv0(p^t),\ldots,a_{\varrho-1}(\alpha)\equiv0(p^t).
\]
於是，剩餘類環的不同由多項式微分商 \(F'(\xi)\) 給出——全部問題都歸結為 Ore, S. 345, § 3。由 Ore 定理 10（參見 Dedekind--Hensel）可知，對每個 \(\mathfrak p\)，取 \(p^t=p^n\) 已經足夠；並且還可由 Ore 得到輔助數。

通過過渡到商環，所有這些同時也適用於相對不同；從而 Ore, Math. Ann. 97, S. 594, § 4 直接由 Ore, Math. Ann. 96 得出。只須補充關於如下可表示性的斷言：上域的不同等於下域的不同與相對不同的乘積（參見本文引言）。

\begin{center}
1949 年 10 月 25 日收稿。
\end{center}
\endgroup

%
\clearpage
%
%
%
%
\providecommand{\tightlist}{}
\providecommand{\mathscr}[1]{\mathcal{#1}}
\def\srcfn#1#2{\begingroup\renewcommand{\thefootnote}{#1}\footnote{#2}\addtocounter{footnote}{-1}\endgroup}
\addcontentsline{toc}{section}{論文 44：超複量代數（講義）}
\markboth{論文 44}{超複量代數}
\begin{center}
\vspace*{3.5em}
{\LARGE\bfseries 超複量代數\par}
\vspace{3.5em}
{\large E. Noether 教授講授\par}
{\large 1929/30 冬季學期\par}
\vspace{1.5em}
{\large M. Deuring 教授整理\par}
\vspace{4em}
{\Large\bfseries 目錄\par}
\end{center}
\vspace{1em}
\begingroup
\small
\setlength{\parindent}{0pt}
\newcommand{\tocline}[2]{\noindent #1\dotfill #2\par}
\newcommand{\tocsec}[3]{\noindent\hspace*{1.2em}\S\,#1.\quad #2\dotfill #3\par}
\tocline{\textbf{導言}}{3}
\medskip
\tocline{\textbf{第一章：表示}}{3}
\tocsec{1}{直接表示的定義 --- 互反表示的定義}{3}
\tocsec{2}{表示類}{4}
\tocsec{3}{表示模}{4}
\tocsec{4}{表示模與表示之間的聯繫}{5}
\medskip
\tocline{\textbf{第二章：交換域的伽羅瓦理論}}{7}
\tocsec{5}{超複系統的係數擴張}{7}
\tocsec{6}{交換系統的不可約表示}{8}
\tocsec{7}{一個域的同構映射}{8}
\tocsec{8}{分裂域}{9}
\tocsec{9}{從 $\mathsf{Z}_1$ 到各 $\mathsf{Z}_i$ 的同構映射}{9}
\tocsec{10}{伽羅瓦群}{9}
\tocsec{11}{伽羅瓦理論主定理}{9}
\tocsec{12}{分量 $e_i$ 的形式意義}{11}
\tocsec{13}{一般延拓定理}{12}
\medskip
\tocline{\textbf{第三章：阿貝爾群}}{13}
\tocsec{14}{群環}{13}
\tocsec{15}{阿貝爾群的群環}{14}
\tocsec{16}{特徵標關係}{15}
\tocsec{17}{阿貝爾群的伽羅瓦理論}{15}
\medskip
\tocline{\textbf{第四章：雙邊單環}}{18}
\tocsec{18}{一個引理}{18}
\tocsec{19}{雙邊單超複系統在其係數域的非交換擴域中的表示}{19}
\tocsec{20}{非交換域}{22}
\tocsec{21}{非交換域的伽羅瓦理論}{26}
\clearpage
\tocline{\textbf{第五章：因子系}}{29}
\tocsec{22}{具有給定中心的域所成的群}{29}
\tocsec{23}{因子系}{30}
\tocsec{24}{因子系的乘法}{33}
\tocsec{25}{具有伽羅瓦極大交換子域的 $\mathfrak K_r$ 的正規表示}{39}
\tocsec{26}{交叉表示的乘法}{43}
\medskip
\tocline{\textbf{第六章：交叉積理論}}{44}
\tocsec{27}{把 $\mathfrak K_r$ 表示為交叉積}{44}
\tocsec{28}{因子系的乘積定理}{47}
\tocsec{29}{極小意義下的主屬定理}{50}
\tocsec{30}{專門化到循環分裂域}{50}
\tocsec{31}{循環特殊情形的應用}{52}
\endgroup
\clearpage

\section*{導言}\label{einleitung}

本講義討論環的一般表示論；它一方面由有限群的表示論發展而來，另一方面由超複數系的代數理論和算術理論發展而來。與以往的表示論（Frobenius、Schur）相比，這種較為一般的理解有如下優點：

由於所考察的不是群，而是群環，更一般地說，是一個超複系統或任意一個環，便可以運用環論和理想論，從而使表示論擺脫許多繁雜難明的計算。不過，運用一般環論和理想論不僅簡化了表示論，而且還把它向前推進了——例如在研究一個群同它的子群之間的關係時——並且為它開闢了新的應用領域。藉助一般表示論，可以為伽羅瓦理論給出一種新的、十分優美的奠基，並且——這也許更為重要——建立非交換域的伽羅瓦理論。

群論、環論和理想論的基本概念以及表示論中的一般定理，見 Noether 女士的論文《超複量與表示論》，Mathematische Zeitschrift 30，第 641 頁。這篇論文下文以 M. Z. 引用，應當視為本講義的補充。M. Z. 中的相同內容，也見於 Noether 女士 1927/28 冬季學期《超複量與群特徵標》課程的整理稿。

這裡簡要說明表示和表示模的概念，以便引入——只在形式上提供新內容的——互反表示模和互反表示。

\section*{第一章. 表示}\label{kapitel-i.-darstellungen}

設 \(\mathfrak{o}\) 是一個環——即被表示的環——，\(\mathsf T\) 是另一個環——即表示在其中進行的環。

\subsection*{§ 1. 直接表示的定義}\label{definition-der-direkten-darstellung}

一個 \(n\) 階的、\(\mathfrak{o}\) 在 \(\mathsf T\) 中的\emph{直接表示}，是一個子環 \(\mathfrak D\)，它屬於環 \(\mathsf T_n\)，即由 \(n\) 階、元素取自 \(\mathsf T\) 的矩陣所成的環；\(\mathfrak{o}\) 環同態地映到這個子環上。用符號表示為

\[
\mathfrak{o}\to\mathfrak D\subseteqq\mathsf T_n.
\]

\subsection*{互反表示的定義}\label{definition-der-reziproken-darstellung}

一個 \(n\) 階的、\(\mathfrak{o}\) 在 \(\mathsf T^*\) 中的\emph{互反表示}，是一個子環 \(\mathfrak D^*\)，它屬於環 \(\mathsf T_n^*\)，即由 \(n\) 階、元素取自 \(\mathsf T^*\) 的矩陣所成的環；\(\mathfrak{o}\) 互反環同態地映到這個子環上。所謂互反環同態，是指：若 \(\mathfrak{o}\) 的元素 \(c,d\) 分別對應於 \(\mathfrak D^*\) 的元素 \(c^*,d^*\)，則 \(\mathfrak{o}\) 的元素 \(cd\) 應當對應於 \(\mathfrak D^*\) 的元素 \(d^*\cdot c^*\)，而 \(\mathfrak{o}\) 的元素 \(c+d\) 應當對應於 \(\mathfrak D^*\) 的元素 \(c^*+d^*\)。用符號表示為

\[
\mathfrak{o}\nrightarrow\mathfrak D^*\subseteqq T_n^*.
\]

舉一個互反同態的例子：把交換域上的每個 \(n\) 階矩陣對應到它的轉置矩陣，這個對應就是一個互反環同構。

一般地，對任意環 \(\mathsf T\)，可以按如下方式構造一個與它互反同構的環 \(\mathsf T^*\)：

給每個 \(\mathsf T\) 元素 \(c\) 對應一個符號 \(c^*\)，並由 \(c^*+d^*=(c+d)^*\) 定義這些符號的加法，由

\[
c^*\cdot d^*=(dc)^*.
\]

定義它們的乘法。符號 \(c^*\) 的全體構成一個與 \(\mathsf T\) 互反同構的環 \(\mathsf T^*\)。

\subsection*{§ 2. 表示類}\label{darstellungsklassen}

把一個直接表示 \(\mathfrak D\)，即 \(\mathfrak{o}\) 在 \(\mathsf T\) 中的表示，變成一個同構表示 \(\mathfrak D'\) 的方法，是用一個可逆矩陣 \(P\) 逐元素作變換；也就是說，對 \(\mathfrak{o}\) 的元素 \(c\)，不再對應 \(C\subseteqq\mathfrak D\)，而對應 \(P^{-1}CP\subseteqq P^{-1}\mathfrak D P\)。

兩個能夠以這種方式彼此變換的表示，稱為\emph{等價}表示。

與一個表示等價的所有表示構成一個\emph{表示類}（直接表示類，d. D. K.）。

互反表示的等價性也完全同樣定義；若兩個互反表示中的一個可由另一個通過一個可逆矩陣的變換得到，就稱它們等價。與一個固定的互反表示等價的所有互反表示構成一個\emph{互反表示類}。

\subsection*{\texorpdfstring{\(\S\) 3. 表示模}{\textbackslash S 3. 表示模}}\label{s-3.-darstellungsmoduln}

\emph{直接表示模}（d. D. M.）的定義。

\(\mathfrak{o}\) 關於 \(\mathsf T\) 的直接表示模，是每一個 \(\mathfrak{o}\)、\(\mathsf T\)-雙模 \(\mathfrak M\)（也就是說，一個以加法記號書寫的阿貝爾群；它關於 \(\mathfrak{o}\) 是左模，關於 \(\mathsf T\) 是右模，並滿足貫通結合律：\(c\in\mathfrak{o}\)，\(m\in\mathfrak M\)，\(\tau\in\mathsf T\)；\((cm)\tau=c(m\tau)\)），而且它可寫成有限多個單生成元 \(\mathsf T\)-模的直和——其中 \(\tau=0\) 可由 \(x_i\tau=0\) 推出——

\[
\mathfrak M=x_1\cdot\mathsf T+x_2\cdot\mathsf T+\cdots+x_n\cdot\mathsf T.
\]

\emph{互反表示模}（r. D. M.）的定義。

\(\mathfrak{o}\) 關於 \(\mathsf T^*\) 的互反表示模，是每一個同時為 \(\mathfrak{o}\)-左模和 \(\mathsf T^*\)-左模的 \(\mathfrak M^*\)；它滿足交換律（若 \(c\in\mathfrak{o}\)，\(m\in\mathfrak M^*\)，\(\tau\in\mathsf T^*\)，則 \(c(\tau^*m^*)=\tau^*(cm)\)），並可寫成有限多個簡單 \(\mathsf T^*\)-模的直和——其中 \(\tau^*=0\) 可由 \(\tau^*\cdot x_i^*=0\) 推出——

\[
\mathfrak M^*=\mathsf T^*\cdot x_1+\cdots+\mathsf T^*\cdot x_n.
\]

\subsection*{§ 4. 表示模與表示之間的聯繫}\label{der-zusammenhang-zwischen-darstellungsmoduln-und-darstellungen}

\textbf{1.} \emph{每個表示模都以唯一方式導出一個表示類（d. D. M. 導出 d. D. K.，r. D. M. 導出 r. D. K.）。}

具體如下：

直接表示模 \(\mathfrak M\) 乘以一個 \(\mathfrak{o}\) 元素 \(c\)，意味着 \(\mathfrak M\) 到自身的一個同態；由於結合律 \((cm)\tau=c(m\tau)\)，它容許 \(\mathsf T\) 的元素作為算子。因此，它是 \(\mathfrak M\) 到自身的一個\emph{線性變換}；也就是說，對應 \(m\mapsto cm=m'\) 具有如下性質：

\[
\text{若 } y_i\mapsto y_i'(=cy_i) \text{，則也有 }
\sum y_i\tau_i\mapsto \sum y_i'\tau_i\left(=c\sum y_i\tau_i\right).
\]

若把兩個線性變換的和 \(\mathfrak T_1+\mathfrak T_2\)，其中這兩個線性變換是

\[
\mathfrak T_1:\ m\mapsto m'(=c_1m)
\quad\text{和}\quad
\mathfrak T_2:\ m\mapsto m''(=c_2m)
\]

定義為變換 \(m\mapsto m'+m''(=(c_1+c_2)m)\)，把它們的積 \(\mathfrak T_1\cdot\mathfrak T_2\) 定義為複合變換 \(m\mapsto (m'')'=c_1c_2m\)，那麼 \(\mathfrak M\) 的由 \(\mathfrak{o}\) 元素生成的各個不同線性變換的全體，就是一個環同態像 \(\overline{\mathfrak D}\)，它是 \(\mathfrak{o}\) 的環同態像。

現在任取一組基 \(x_1,\ldots,x_n\)，它屬於 \(\mathfrak M\)。由 \(c\) 生成的 \(\mathfrak M\) 的線性變換，只要知道基元素 \(x_i\) 怎樣變換，顯然就已經確定；也就是說，只要矩陣 \((\gamma_{ij})=C\)，即方程組

\[
x'_j=cx_j=\sum x_i\gamma_{ij}
\]

已經知道即可，因為對其他元素 \(x=\sum x_j\tau_j\)，簡單地有 \(x'=\sum x'_j\tau_j\)。這些矩陣 \(C\) 與各個不同的線性變換一一對應（因為它們彼此確定），並且易見構成一個與 \(\overline{\mathfrak D}\) 同構的環，因而構成 \(\mathfrak{o}\) 的一個環同態像，即一個 \(n\) 階的 \(\mathfrak{o}\) 的表示。

\(\overline{\mathfrak D}\) 與這個表示之間的同構體現在如下公式中。若由 \(c\in\mathfrak{o}\) 生成的線性變換對應於矩陣 \(C\)，即

\[
c(x_1,\ldots,x_n)=(x_1,\ldots,x_n)C
\]

成立，並且同樣地，\(d\in\mathfrak{o}\) 與 \(D\in\mathsf T_n\) 相互對應，

\[
d(x_1,\ldots,x_n)=(x_1,\ldots,x_n)D
\]

那麼顯然有

\[
\begin{aligned}
(c+d)(x_1,\ldots,x_n)
&=(cx_1+dx_1,\ldots,cx_n+dx_n)\\
&=(x_1,\ldots,x_n)(C+D)
\end{aligned}
\]

以及

\[
\begin{aligned}
cd(x_1,\ldots,x_n)
&=c(dx_1,\ldots,dx_n)=c((x_1,\ldots,x_n)D)\\
&=(cx_1,\ldots,cx_n)D=(x_1,\ldots,x_n)CD.
\end{aligned}
\]

因此，由 \(\mathfrak{o}\) 生成的 \(\mathfrak M\) 到自身的線性變換，在藉助一組指定基 \(x_\nu\)——它是 \(\mathfrak M\) 的基——實現時，給出一個 \(n\) 階的、\(\mathfrak{o}\) 在 \(\mathsf T\) 中的表示。

若 \((y_1,\ldots,y_n)\) 是 \(\mathfrak M\) 的另一組基，則它由 \(x_1,\ldots,x_n\) 乘一個可逆矩陣得到：

\[(y_1,\ldots,y_n)=(x_1,\ldots,x_n)P\]

若 \(\mathfrak M\) 的一個線性變換在基 \(x_\nu\) 下由矩陣 \(C\) 實現，那麼它在基 \(y_\nu\) 下由 \(P^{-1}CP\) 實現；因為若

\[c(x_1,\ldots,x_n)=(x_1,\ldots,x_n)C\]
成立，則

\[
\begin{aligned}
c(y_1,\ldots,y_n)
&=c(x_1,\ldots,x_n)P=((x,\ldots,x_n)C)P\\
&=(x_1,\ldots,x_n)CP=((y_p,\ldots,y_n)P^{-1})CP\\
&=(y_1,\ldots,y_n)P^{-1}CP.
\end{aligned}
\]

由此立即推出：

\emph{\(\mathfrak M\) 給出一個表示類中的所有表示，只要採用 \(\mathfrak M\) 的所有基來實現由 \(\mathfrak{o}\) 所誘導的 \(\mathfrak M\) 的線性變換。}

\begin{enumerate}
\def\labelenumi{\arabic{enumi}.}
\setcounter{enumi}{1}
\tightlist
\item
  \emph{每個表示類都以 1. 中所述的方式由一個表示模生成。}（我們仍只討論直接表示。）
\end{enumerate}

\emph{證明.} 任取該類中的一個表示，設它由同態對應 \(c\mapsto C\) 給出。若 \(n\) 是這個表示的階數，就用 \(n\) 個不定元 \(x_1,\ldots,x_n\) 構成一個 \(\mathsf T\)-右模

\[\mathfrak{M} = x_1 \mathsf{T} + \dots + x_n \mathsf{T}\]

（採用通常的運算規則 \(\sum x_i\tau_i+\sum x_i\overline{\tau}_i=\sum x_i(\tau_i+\overline{\tau}_i)\)，\(\sum x\tau_i=\sum x_i\overline{\tau}_i\) 當且僅當 \(\tau_i=\overline{\tau}_i\) 對所有 \(i\) 都成立，以及 \((\sum x_i\tau_i)\varrho=\sum x_i\tau_i\varrho\)），並通過如下規定使它成為一個 \(\mathfrak{o}\)-左模：規定 \(c\sum x_i\tau_i=\sum x_i'\tau_i\)，只要 \((x_1',\ldots,x_n')=(x_1,\ldots,x_n)C\)，其中 \(C\) 表示該表示中對應於 \(\mathfrak{o}\) 元素 \(c\) 的矩陣。

這個規定使 \(\mathfrak M\) 成為一個 \(\mathfrak{o}\)、\(\mathsf T\)-雙模；因為第一，由同態關係

\[
c+d\mapsto C+D \quad \text{和} \quad cd\mapsto CD
\]

可得

\[
(c+d)x_i=cx_i+dx_i
\]

以及

\[
cd\cdot x_i=c\cdot dx_i,\quad\text{此外}\quad
c\left(\sum x_i\tau_i+\sum x_i\overline{\tau}_i\right)
=c\sum x_i\tau_i+c\sum x_i\overline{\tau}_i,
\]

這些就是左模性質；第二，顯然有

\[c(\sum x_i \, \tau_i \, \varrho) = (c \, \sum x_i \, \tau_i) \cdot \varrho \, ,\]

這就是貫通結合律。

因此，\(\mathfrak{M}\) 是表示模；若用基 \(x_1, \ldots, x_n\) 來實現它的線性變換，則所採用的表示由這個模按 1. 中所述的方式生成。其他各組基給出該類中的其餘表示。

對互反表示和互反表示模，一切都完全相同；結果是，模

\[\mathfrak{M}^* = \mathsf{T}^* x_1^* + \dots + \mathsf{T}^* x_n^*\]

生成 \(\mathfrak D\) 在 \(\mathsf T^*\) 中的一個互反表示，方法是給 \(\mathfrak{o}\) 的元素 \(c\) 對應矩陣 \(C\)，它出現在方程

\[
c\begin{pmatrix}x_1\\ \vdots\\ x_n\end{pmatrix}
=\begin{pmatrix}cx_1\\ \vdots\\ cx_n\end{pmatrix}
=C\begin{pmatrix}x_1\\ \vdots\\ x_n\end{pmatrix}
\]

中。不同的基仍給出一個類中的各個不同表示。

\section*{第二章. 交換域的伽羅瓦理論}\label{kapitel-ii.-galoissche-theorie-kommutativer-koerper}

迄今所發展的表示論足以用來奠定伽羅瓦理論；與通常的奠基相比，它的優點在於無須訴諸特殊的生成元系和關於不定元多項式的定理。\srcfn{1}{以下內容可參看 M. Z. §§ 15 --- 20。M. Z. § 21 包含下一節關於伽羅瓦理論的最初幾個定理。關於超複系統的定義，見 M. Z. § 6。}

\subsection*{§ 5. 超複系統的係數擴張}\label{koeffizientenerweiterung-hyperkomplexer-systeme}

我們表述如下的一般定理；它的證明幾乎不言自明，因而從略。設 \(\mathfrak{o}=c_1\mathsf P+\cdots+c_n\mathsf P\) 是交換域 \(\mathsf P\) 上的一個超複系統。設 \(\Omega\) 是 \(\mathsf P\) 的一個擴張環，並且其中心包含 \(\mathsf P\)。我們把\emph{擴張系統} \(\mathfrak{o}_{\Omega}\) 定義為所有形式和 \(\sum c_i\omega_i\)，\(\omega_i\in\Omega\)，的全體；先作如下規定：

\[
\begin{array}{l}
\sum c_i\omega_i=\sum c_i\omega_i' \quad\text{當且僅當 } \omega_i=\omega_i' \text{ 對所有 } i;\\
\sum c_i\omega_i+\sum c_i\omega_i'=\sum c_i(\omega_i+\omega_i');\\
\omega\sum c_i\omega_i=\sum c_i\omega\omega_i;\quad
\sum c_i\omega_i\cdot\omega=\sum c_i\omega_i\omega;\\
\sum c_i\omega_i=\sum \omega_i c_i
\end{array}
\]

由此把它變成一個同 \(\Omega\) 交換相聯的 \(\Omega\)-模，再作如下進一步規定：

\[
\sum c_i\omega_i\cdot \sum c_i\omega_i
=\sum_i\sum_j c_ic_j\omega_i\omega_j
\]

——其中 \(c_ic_j\) 作為 \(\mathfrak{o}\) 的元素已經有了定義——從而使它成為 \(\mathfrak{o}\) 的一個擴張環。\(\mathfrak{o}_{\Omega}\) 的秩為 \(n\)（關於 \(\Omega\)），並可寫成

\[
\mathfrak{o}_{\Omega}=c_1\Omega+\cdots+c_n\Omega
\]

的形式。

\emph{\(\mathfrak{o}_{\Omega}\) 的定義與基 \(c_1,\ldots,c_n\) 無關。若 \(\Omega\) 是交換域，則 \(\mathfrak{o}_{\Omega}\) 是 \(\Omega\) 上的一個超複系統。}

\subsection*{§ 6. 交換系統的不可約表示}\label{die-irreduziblen-darstellungen-kommutativer-systeme}

設 \(\mathfrak Z=z_1\mathsf P+\cdots+z_n\mathsf P\) 是 \(\mathsf P\) 上的一個交換系統。我們要建立 \(\mathfrak Z\) 在一個代數閉擴域 \(\Omega\)——它是 \(\mathsf P\) 的擴域——中的所有不可約表示。由於 \(\mathfrak Z\) 在 \(\Omega\) 中的每個表示都導出 \(\mathfrak Z_\Omega=z_1\Omega+\cdots+z_n\Omega\) 的一個表示——只要知道那些已經包含在 \(\mathsf P\) 中的基元素的表示，就足以建立整個表示——，反過來，\(\mathfrak Z_\Omega\) 的每個表示都包含 \(\mathfrak Z\) 的一個表示，所以我們也可以轉而考察 \(\mathfrak Z_\Omega\) 的表示。按照 M. Z. § 20，這些表示包含在 \(\mathfrak Z_\Omega\) 的正則表示中；所謂正則表示，就是 \(\mathfrak Z_\Omega\) 的這樣一種表示：把 \(\mathfrak Z_\Omega\) 本身作為表示模。甚至只須把 \(\mathfrak Z_\Omega/\mathfrak C\) 作為表示模，其中 \(\mathfrak C\) 表示 \(\mathfrak Z_\Omega\) 的根基（見 M. Z. § 19）。

不過，\(\mathfrak Z_\Omega/\mathfrak C\) 是一個完全可約系統，因而是若干個 \(\Omega\) 的有限次擴域的直和。由於 \(\Omega\) 是代數閉域，這些直和項都與 \(\Omega\) 同構。

\[
\mathfrak Z_\Omega/\mathfrak C=e_1\Omega+\cdots+e_t\Omega.
\]

各 \(e_\nu\) 是單位元的分量。因此有：

\emph{\(\mathfrak Z_\Omega\) 恰有 \(t\) 個在 \(\Omega\) 中不同的不可約表示類，它們全都是一階的。這裡 \(t\) 是 \(\mathfrak Z_\Omega/\mathfrak C\) 的秩，因而至多等於 \(n\)。}

這些表示類中的每一個都只包含一個表示，因為在交換域中變換一個一階表示並不改變該表示：\(\pi^{-1}\cdot\alpha\cdot\pi=\alpha\)。

所以，恰有 \(t\) 個不同的同態映射 \(\Theta_1,\ldots,\Theta_t\)，它們把 \(\mathfrak Z_\Omega\) 映到 \(\Omega\) 的子環上。

\subsection*{§ 7. 一個域的同構映射}\label{die-isomorphismen-eines-koerpers}

現在設 \(\mathfrak Z\) 是一個域。每個同態 \(\Theta_i\)，若它把 \(\mathfrak Z_\Omega\) 映到 \(\Omega\) 的一個子環上，就必須把 \(\mathfrak Z\) 同構地映到一個子域 \(\mathsf{Z}_i\)，而這個子域屬於 \(\Omega\)。因為要麼 \(\mathsf{Z}_i=0\)——這裡並非如此，因為 \(\mathfrak Z\) 的單位元對應於 \(\Omega\) 的單位元——，要麼 \(\mathsf{Z}_i\) 是 \(\mathfrak Z\) 的一個同構像。因此：

\emph{若 \(\mathfrak Z\) 是 \(\mathsf P\) 的一個 \(n\) 次有限擴域，則在 \(\Omega\) 中，與 \(\mathfrak Z\) 同構的不同域 \(\mathsf{Z}_i\) 的數目等於秩 \(t\)，即 \(\mathfrak Z_\Omega/\mathfrak C\) 的秩。所以，\(\mathsf{Z}_i\) 的數目至多等於 \(n\)。}

現在稱 \(\mathfrak Z\) 在 \(\mathsf P\) 上為\emph{第一類}，如果各 \(\mathsf{Z}_i\) 的數目等於 \(n\)；如果這個數目小於 \(n\)，則稱為\emph{第二類}。

\emph{\(\mathfrak Z\) 在 \(\mathsf P\) 上為第一類，當且僅當 \(\mathfrak Z_\Omega\) 完全可約。}

這裡可以看出，這種伽羅瓦理論的奠基比通常奠基優越之處。通常的伽羅瓦理論考察某個 \(\mathsf{Z}_i\) 到其他各共軛域的同構映射。這裡避免了這種不對稱性。為了證明一個域的共軛域不多於它的次數，通常要使用本原元。這裡，這個事實簡單地來自：\(\mathfrak Z_\Omega\) 的所有不可約表示都包含在正則表示中。

\subsection*{§ 8. 分裂域}\label{der-zerfaellungskoerper}

各 \(\mathsf{Z}_i\) 是 \(\mathfrak Z\) 在 \(\Omega\) 中的不可約表示。因此，若域 \(\Gamma\) 位於 \(\mathsf P\) 與 \(\Omega\) 之間，即 \(\mathsf P\subseteqq\Gamma\subseteqq\Omega\)，並且包含所有 \(\mathsf{Z}_i\)，那麼 \(\mathfrak Z_\Gamma/\mathfrak C\) 必已分解為 \(t\) 個在 \(\Gamma\) 上為一次的域。反過來，要使 \(\mathfrak Z_\Gamma/\mathfrak C\) 分解為 \(t\) 個域——這些域於是就是 \(\mathfrak Z_\Gamma\) 的一階不可約表示——，各 \(\mathsf{Z}_i\) 必須包含在 \(\Gamma\) 中。若稱 \(\mathfrak Z\) 的分裂域為一個域 \(\Gamma\)，且它具有這樣的性質：\(\mathfrak Z_\Gamma/\mathfrak C\) 成為 \(t\) 個域的直和，那麼我們得到：

\emph{與 \(\mathfrak Z\) 同構的各 \(\mathsf{Z}_i\) 的伽羅瓦域，也就是它們的複合域，是 \(\mathfrak Z\) 的極小分裂域。}
\subsection*{§ 9. 從 \(\mathsf{Z}_1\) 到各 \(\mathsf{Z}_i\) 的同構}\label{die-isomorphismen-von-z1-auf-die-zi}

若只考察 \(\mathfrak Z\)，則 \(\Theta_i\) 是從 \(\mathfrak Z\) 到 \(\mathsf{Z}_i\) 的一個同構映射。因此，\(S_i=\Theta_i\Theta_1^{-1}\) 是從 \(\mathsf{Z}_1\) 到 \(\mathsf{Z}_i\) 的一個同構映射。由於各 \(\Theta_i\) 彼此不同，各 \(S_i\) 也彼此不同。各 \(S_i\) 就是通常在伽羅瓦理論中所考察的、從域 \(\mathsf{Z}_1\) 到其各共軛域的同構。除此之外不存在別的同構，因為不存在別的 \(\Theta_i\)。

\subsection*{§ 10. 伽羅瓦群}\label{die-galoissche-gruppe}

若 \(\mathfrak Z\) 是伽羅瓦域，也就是說，各 \(\mathsf{Z}_i\) 作為元素集合彼此重合，那麼各 \(S_i\) 就組成所有逐點固定 \(\mathsf P\) 的、從 \(\mathsf{Z}\) 到自身的同構所成的群。事實上，每個乘積 \(S_iS_j\) 都是 \(\mathsf{Z}\) 的一個自同構；又因除各 \(S_i\) 外別無其他同構，所以 \(S_iS_j\) 等於某個 \(S_k\)。因此，各 \(S_i\) 組成 \(Z\) 的所有逐點固定 \(\mathsf P\) 的自同構所成的群；這就是\emph{\(\mathsf{Z}\) 關於 \(\mathsf P\) 的伽羅瓦群}。

\subsection*{§ 11. 伽羅瓦理論主定理}\label{der-hauptsatz-der-galoisschen-theorie}

現在設 \(\mathfrak Z\) 是 \(\mathsf P\) 的第一類伽羅瓦擴張。我們來證明伽羅瓦理論主定理：

\emph{介於 \(\mathsf P\) 與 \(\mathfrak Z\) 之間的域 \(\Gamma\)，可以與 \(Z\) 關於 \(\mathsf P\) 的伽羅瓦群 \(\mathfrak G\) 的子群 \(\mathfrak H\) 一一對應，使得與域 \(\Gamma\) 對應的群 \(\mathfrak H\) 由所有逐點固定 \(\Gamma\) 的 \(Z\) 的自同構組成，而 \(\Gamma\) 則是所有在 \(\mathfrak H\) 所含的每個自同構下保持不變的 \(\mathsf{Z}\)-元素的全體。}

我們將它簡要表述如下：

\begin{enumerate}
\item \emph{若 \(\mathfrak H\) 是 \(\Gamma\) 的不變群，則 \(\Gamma\) 是 \(\mathfrak H\) 的不變域。}
\item \emph{若 \(\Gamma\) 是 \(\mathfrak H\) 的不變域，則 \(\mathfrak H\) 是 \(\Gamma\) 的不變群。}
\end{enumerate}

\paragraph{1. 可歸結為如下輔助定理：}\label{laesst-sich-zurueckfuehren-auf-den-hilfssatz}

\emph{若 \(\mathsf P \subseteqq \Sigma \subseteqq \mathsf T \subseteqq \mathsf Z\)，並且 \(\mathsf T\) 在 \(\Sigma\) 上的次數為 \(t\)，那麼從 \(\Sigma\) 到某個共軛域（它必然包含在 \(\mathsf Z\) 中）的每一個同構，都能以 \(t\) 種不同的方式延拓為從 \(\mathsf T\) 到各共軛域的同構。}

事實上，若假定 \(\mathfrak H\) 是 \(\Sigma\) 的不變群，而 \(\mathsf T\) 是 \(\mathfrak H\) 的不變域，則可以作如下推論：\(\mathfrak H\) 的元素，是 \(\Sigma\) 的恆同同構在整個 \(\mathsf Z\) 上的各個延拓。若只把這些延拓限制到 \(\mathsf T\) 上，則由輔助定理可知，它們分成 \(t\) 個類，使得同一類中的同構在 \(\mathsf T\) 上相同，不同類中的同構在 \(\mathsf T\) 上不同。但由於 \(\mathsf T\) 是 \(\mathfrak H\) 的不變域，必有 \(t=1\)、\(\mathsf T=\Sigma\)，證畢。

為證明輔助定理，我們回到介於 \(\mathsf P\) 與 \(\mathsf Z\) 之間、經 \(\Theta_1\) 分別與 \(\Sigma\) 和 \(\mathsf T\) 對應的域 \(\mathfrak S\)、\(\mathfrak T\)。於是必須證明：

\emph{從 \(\mathfrak S\) 到某個 \(\Sigma_i\) 的每一個同構 \(\mathsf{H}_i\)，都恰好能以 \(t\) 種不同的方式延拓為從 \(\mathfrak T\) 到 \(\mathsf Z\) 的各子域 \(\mathsf T_i\) 的同構。}

由於假定 \(\mathfrak Z\) 在 \(\mathsf P\) 上是第一類的，\(\mathfrak Z_\Omega/\mathfrak C\) 與 \(\mathfrak Z_\Omega\) 相同；因此，\(\mathfrak S_\Omega\) 分解為其次數 \(s\) 所指定的那麼多個域：
\[
\mathfrak S_\Omega = E_1\Omega+\cdots+E_s\Omega .
\]
作為表示模，各 \(E_\nu\Omega\) 產生從 \(\mathfrak S\) 到 \(s\) 個子域 \(\Sigma_1,\ldots,\Sigma_s\) 的 \(s\) 個同構。為找出 \(\mathfrak T\) 的同構，我們需要考察 \(\mathfrak T_\Omega\)。由於各 \(E_\nu\Omega\) 是 \(\mathfrak S_\Omega\) 的理想，有 \(E_\nu\Omega=\mathfrak S_\Omega E_\nu\)。此外，顯然
\[
\mathfrak T_\Omega=\mathfrak T\cdot\Omega=\mathfrak T\cdot\mathfrak S_\Omega .
\]
所以
\[
\mathfrak T_\Omega
=E_1\mathfrak S_\Omega\cdot\mathfrak T+\cdots+E_s\mathfrak S_\Omega\cdot\mathfrak T
=E_1\mathfrak T_\Omega+\cdots+E_s\mathfrak T_\Omega .
\]
這是 \(\mathfrak T_\Omega\) 的一個理想分解。\(E_\nu\cdot\mathfrak T_\Omega\) 分解為 \(t\) 個簡單理想。事實上，有次數關係
\[
[\mathfrak T_\Omega\cdot E_\nu:\mathfrak S_\Omega E_\nu]
\leqq [\mathfrak T_\Omega:\mathfrak S_\Omega]=[\mathfrak T:\mathfrak S]=t,
\]
因而，由 \(\mathfrak S_\Omega E_\nu=\Omega E_\nu\cong\Omega\)，有
\[
[\mathfrak T_\Omega\cdot E_\nu:\Omega]
=[\mathfrak T_\Omega E_\nu:\mathfrak S_\Omega E_\nu]\leqq t,
\]
而由於所有次數 \([\mathfrak T_\Omega E_\nu:\Omega]\) 的總和必須等於 \(s\cdot t\)，便有
\[
[\mathfrak T_\Omega\cdot E_\nu:\Omega]=t,\qquad
E_\nu\mathfrak T_\Omega=e_\nu^{(1)}\cdot\Omega+\cdots+e_\nu^{(t)}\Omega .
\]
每個 \(e_\nu^{(i)}\Omega\) 都給出 \(\mathfrak T_\Omega\) 的一個不可約表示。\emph{只要再證明}：把所有這些表示限制到 \(\mathfrak S_\Omega\) 上時，它們都同由 \(E_\nu\cdot\Omega\) 產生的表示重合，證明便完成了。這一點是清楚的，因為由
\[
s\cdot E_\nu=E_\nu\sigma_\nu\quad
(s\in\mathfrak S_\Omega\text{ 由 }\sigma_\nu\text{ 表示})
\]
可得
\[
\sum s\cdot e_\nu^{(i)}=\sum e_\nu^{(i)}\sigma_\nu,
\]
於是，由和表示的唯一性可得
\[
s\cdot e_\nu^{(i)}=e_\nu^{(i)}\sigma_\nu,
\]
所以新的表示也把元素 \(s\) 配給元素 \(\sigma_\nu\)。

\paragraph{2. 設 \(\mathsf T\) 是 \(\mathfrak H\) 的不變域。}
我們要認出 \(\mathfrak H\) 是 \(\mathsf T\) 的不變群；為此，\emph{只須}找出一個 \(\mathsf T\)-元素，使它在某個不屬於 \(\mathfrak H\) 的自同構下不保持不變。

\(\mathsf Z\) 的伽羅瓦群 \(\mathfrak G\) 的元素 \(S_i\)，是 \(Z\) 的、逐點固定 \(\mathsf P\) 的自同構。我們通過如下規定，把它們變成 \(\mathfrak Z_{\mathsf Z}\) 的自同構：
\[
S_i z=z,\quad \text{若 } z \text{ 屬於 } \mathfrak Z.
\]
把 \(S_i\) 作用於
\[
\mathfrak Z_{\mathsf Z}=e_1\mathsf Z+\cdots+e_n\mathsf Z
\]
上，必然會置換各直和分量 \(e_\nu\mathsf Z\)，因為這些分量是唯一確定的；特別地，必有
\[
S_i e_\nu=e_{\nu'},\qquad\text{且}\qquad S_i e_\nu\ne S_k e_\nu
\quad\text{當 } i\ne k
\]
成立。若 \(S_1,\ldots,S_h\) 是 \(\mathfrak H\) 的元素，我們便任取一個 \(e_\nu\)，例如記為 \(e'\)，並考察由各 \(S_i\)（\(i\leqq h\)）從它產生的所有 \(e_\nu\)：
\[
e_1=S_1e',\ldots,e_h=S_he'.
\]
由於當 \(i\leqq h, j\leqq h\) 時有
\[
S_i e_j=S_iS_j e'=S_ke'=e_k \quad \text{其中}\quad k\leqq h
\]
所以 \(\mathfrak H\) 中的每個 \(S_i\) 都在 \(e_1,\ldots,e_h\) 之間作置換；因此
\[
E_1=e_1+\cdots+e_h
\]
是一個在所有 \(S_i\in\mathfrak H\) 下均保持不變的 \(\mathfrak Z\)-元素。若用 \(\mathfrak Z\) 的一組基表示 \(E_1\)，它的係數必定是在 \(\mathfrak H\) 下保持不變的 \(\mathsf Z\)-元素；因而這些係數屬於 \(\mathsf T\)。

\(E_1\) 在每個不屬於 \(\mathfrak H\) 的 \(S_i\) 下都不保持不變。事實上，對於這樣的 \(S_i\)，
\[
S_i e_1=e_i+\cdots
\]
含有未出現於 \(e_1,\ldots,e_h\) 中的分量 \(e_i\)，而 \(E_1\) 的和表示是唯一的。

\subsection*{§ 12. 分量 \(e_i\) 的形式意義}\label{die-formale-bedeutung-der-komponenten-ei}

我們固定 \(\mathfrak Z\) 的一組基 \(z_1,\ldots,z_n\)：
\[
\mathfrak Z=z_1\mathsf P+\cdots+z_n\mathsf P
\]
（現在只假定 \(\mathfrak Z\) 是第一類的）。同構 \(\Theta_i\) 把系統 \(z_1,\ldots,z_n\) 映到 \(\mathsf Z_i\) 的一組基 \(\zeta_1^{(i)},\ldots,\zeta_n^{(i)}\)，其映射由下列關係給出：
\[
z_\nu e_i=e_i\zeta_\nu^{(i)} .
\]

反過來，單位元的分量 \(e_j\) 必須能夠用各 \(z_\nu\) 及取自 \(\mathsf Z\) 的係數表示為
\[
e_j=\sum_\nu \varrho_\nu^{(j)}z_\nu .
\]
因此有
\[
e_j=e_j^2=\sum \varrho_\nu^{(j)}z_\nu e_j
=\sum \varrho_\nu^{(j)}\cdot\zeta_\nu^{(j)}\cdot e_j
\]
以及
\[
0=e_je_k=\sum \varrho_\nu^{(j)}\cdot z_\nu\cdot e_k
=\sum \varrho_\nu^{(j)}\cdot\zeta_\nu^{(k)}\cdot e_k,
\]
所以
\[
\sum_\nu \varrho_\nu^{(j)}\cdot\zeta_\nu^{(j)}=1;
\qquad
\sum_\nu \varrho_\nu^{(j)}\cdot\zeta_\nu^{(k)}=0,\quad
\text{當 } j\ne k.
\]
這就是說，矩陣
\[
\begin{pmatrix}
\zeta_1^{(1)} & \cdots & \zeta_n^{(1)}\\
\cdots & \cdots & \cdots\\
\cdots & \cdots & \cdots\\
\cdots & \cdots & \cdots\\
\zeta_1^{(n)} & \cdots & \zeta_n^{(n)}
\end{pmatrix}
\quad\text{与}\quad
\begin{pmatrix}
\varrho_1^{(1)} & \cdots & \varrho_1^{(n)}\\
\cdots & \cdots & \cdots\\
\cdots & \cdots & \cdots\\
\cdots & \cdots & \cdots\\
\varrho_n^{(1)} & \cdots & \varrho_n^{(n)}
\end{pmatrix}
\]
彼此互逆。也可以這樣說：\(\varrho_1^{(i)},\ldots,\varrho_n^{(i)}\) 是與 \(\zeta_1^{(i)},\ldots,\zeta_n^{(i)}\) 互補的基。它在數域的“不同”（即分歧理想）理論中起作用。

\subsection*{§ 13. 一般延拓定理}\label{der-allgemeine-fortsetzungssatz}

第 10 頁所證明的關於延拓各 \(\mathsf{H}_i\) 的定理可以推廣。

設 \(\mathfrak Z=z_1\mathsf P+\cdots+z_n\mathsf P\) 是一個交換的、完全可約的超複系統。在這種情形下，\(\mathfrak S\) 和 \(\mathfrak T\) 是 \(\mathfrak Z\) 的子系統——並且必然也是完全可約的——且滿足
\[
\mathsf P\subseteqq\mathfrak S\subseteqq\mathfrak T\subseteqq\mathfrak Z .
\]
由 M. Z. § 13，這三個系統各有一個單位元。但這三個單位元完全不必相同。（當 \(\mathfrak S\) 和 \(\mathfrak T\) 都是 \(\mathfrak Z\) 的、不同於 \(\mathfrak Z\) 本身的理想時，它們肯定不同。）設 \(E\) 是 \(\mathfrak S\) 的單位元。系統 \(\mathfrak T\cdot E\) 是一個有限秩的 \(\mathfrak S\)-模：
\[
E\mathfrak T=a_1\mathfrak S+\cdots+a_t\mathfrak S
\]
（一般而言，\(\mathfrak T\) 本身並不是！）。

\emph{證明。} \(\mathfrak S\) 是若干域的直和：\(\mathfrak S=\mathfrak K_1+\cdots+\mathfrak K_\gamma\)。與此相應，\(\mathfrak T E\) 也有一個直和分解。\(\mathfrak T E\) 的每個直和分量必定包含某個域 \(\mathfrak K_i\)，否則它就會被 \(E\) 零化；但 \(E\) 是 \(\mathfrak T E\) 的單位元。因此，\(\mathfrak T E\) 是各 \(\mathfrak K_i\) 的擴張環的直和，也就是說，是如下形式的直和：
\[
\mathfrak T E=a_{11}\mathfrak K_1+\cdots+a_{1p_1}\mathfrak K_1+\cdots+a_{rp_r}\mathfrak K_r,
\]
或者說，\(E\mathfrak T=a_{11}\mathfrak S+\cdots+a_{rp_r}\mathfrak S\)，證畢。

\begin{center}
一般延拓定理。
\end{center}

\emph{\(\mathfrak S\) 在 \(\Omega\) 中的每個不可約表示——由一個同態 \(\mathsf{H}_i\) 給出——都恰好能以 \(t\) 種不同的方式延拓為 \(\mathfrak T\) 的表示。}

證明與第 10--11 頁的證明相應。我們有
\[
\mathfrak S_\Omega=E_1\Omega+\cdots+E_s\Omega,\qquad
E_\nu\Omega=\mathfrak S_\Omega E_\nu,
\]
\[
\mathfrak T_\Omega\cdot E=\mathfrak T\cdot E\Omega
=\mathfrak T\mathfrak S\Omega=\mathfrak T\mathfrak S_\Omega,
\]
\[
\mathfrak T_\Omega\cdot E
=\mathfrak T(E_1\mathfrak S_\Omega+\cdots+E_s\mathfrak S_\Omega)
=E_1\mathfrak T_\Omega+\cdots+E_s\mathfrak T_\Omega,
\]
\[
[\mathfrak T_\Omega E_\nu:\Omega]
=[\mathfrak T_\Omega E_\nu:\Omega E_\nu]
=[\mathfrak T_\Omega E_\nu:E_\nu\mathfrak S_\Omega]
\leqq[\mathfrak T_\Omega\cdot E:\mathfrak S_\Omega]
=[\mathfrak T\cdot E:\mathfrak S]=t,
\]
\[
\sum[\mathfrak T_\Omega\cdot E_\nu:\Omega]=st
\]
所以
\[
[\mathfrak T_\Omega E_\nu:\Omega]=t.
\]
因而，\(\mathfrak T_\Omega E_\nu\) 分解為 \(t\) 個理想；其中每一個作為 \(\mathfrak T_\Omega E\) 的直和分量，從而也作為 \(\mathfrak T_\Omega\) 的直和分量，都給出 \(\mathfrak T_\Omega\) 的一個表示；這個表示如同第 10--11 頁所述，可看出是由 \(\mathsf{H}_\nu\) 給出的表示的延拓。\(\mathfrak T_\Omega\) 的其餘不可約表示——由 \(\mathfrak T_\Omega\) 中不包含在 \(\mathfrak T_\Omega\cdot E_\nu\) 內的那些直和分量產生——不可能是由各 \(\mathsf{H}_\nu\) 產生的 \(\mathfrak S_\Omega\) 的表示的延拓，因為這些理想被 \(\mathfrak T_\Omega\cdot E_\nu\) 零化，因而更被 \(\mathfrak S_\Omega\) 零化。

\section*{第三章. 阿貝爾群}\label{kapitel-iii.-abelsche-gruppen}

\subsection*{§ 14. 群環}\label{der-gruppenring}

設 \(\mathfrak G\) 是一個有限群，其元素為 \(a_1,\ldots,a_h\)，而 \(\mathsf P\) 是一個交換域。以群乘法作為乘法，我們構造一個超複系統
\[
\mathfrak o[\mathfrak G]=a_1\mathsf P+\cdots+a_h\mathsf P
\]
即\emph{\(\mathfrak G\) 在 \(\mathsf P\) 上的群環}。（M. Z. § 6；Speiser, Theorie d. Gruppen end. Ord. § 48。）

我們要研究 \(\mathfrak o[\mathfrak G]\) 在 \(\mathsf P\) 中或在 \(\mathsf P\) 上的代數閉域 \(\Omega\) 中的表示。如前所述，不可約表示由 \(\mathfrak o[\mathfrak G]/\mathfrak C\) 的簡單左理想產生。因此，同伽羅瓦理論中一樣，我們所關心的問題是：\(\mathfrak o[\mathfrak G]\) 何時有根基。對此有如下的一般定理：

\begin{center}
\emph{\(\mathfrak o[\mathfrak G]\) 完全可約，當且僅當 \(\mathfrak G\) 的階數 \(h\) 不被 \(\mathsf P\) 的特徵 \(p\) 整除。}
\end{center}

我們先證明：
\[
\text{由 } h\equiv0(p) \text{ 可推出 } \mathfrak C\ne 0.
\]
\[
\mathfrak a=\mathsf P(a_1+\cdots+a_h)
\]
是 \(\mathfrak o\) 的一個雙側理想。事實上，
\[
\mathsf P\cdot(a_1+\cdots+a_h)a
=\mathsf P\cdot(a_1a+\cdots+a_ha)
=\mathsf P(a_1+\cdots+a_h)=\mathfrak a
\]
並且
\[
a\cdot\mathsf P(a_1+\cdots+a_h)
=\mathsf P(aa_1+\cdots+aa_h)
=\mathsf P(a_1+\cdots+a_h)=\mathfrak a
\]

這來自各 \(a_i\) 的群性質。\(\mathfrak a\) 顯然不為零，但它是冪零的：
\[
\left(\sum_i a_i\right)^2=\sum_i\sum_j a_i a_j
=h\sum_j a_j=0.
\]

反過來的命題
\[
\text{由 } h\not\equiv0(p) \text{ 可推出 } \mathfrak C=0,
\]
我們現在只對阿貝爾群加以證明。（在 M. Z. § 26 中，這部分定理已在一般情形下得到證明。）

\subsection*{§ 15. 阿貝爾群的群環}\label{die-gruppenringe-abelscher-gruppen}

我們用 \(\mathfrak A=\{a_1,\ldots,a_h\}\) 表示一個阿貝爾群。於是 \(\mathfrak o\) 是交換的，所以我們可以應用關於交換超複系統的一般定理。

\emph{若 \(h\) 不被 \(\mathsf P\) 的特徵 \(p\) 整除，則 \(\mathfrak o[\mathfrak A]\) 在 \(\Omega\) 中恰有 \(h\) 個彼此不同的不可約表示。（因為 \(\mathfrak o\) 是交換的，所以這些表示都是一次的。）}

由此又可推出，\(\mathfrak o[\mathfrak A]/\mathfrak C\) 與 \(\mathfrak o\) 有相同的秩，因而 \(\mathfrak C=0\)；這正是剛才尚未證明的定理第二部分在阿貝爾群情形下的結論。

\emph{證明。} \(\mathfrak A\) 是若干階數為 \(h_i\) 的循環群 \(\mathfrak Z_i\) 的直積：
\[
\mathfrak A=\mathfrak Z_1\times\cdots\times\mathfrak Z_t,\qquad
h=h_1h_2\ldots h_t .
\]
若 \(z_i\) 是 \(\mathfrak Z_i\) 的生成元，那麼它在一個（一次）不可約表示下所對應的元素，必定是一個 \(h_i\) 次單位根 \(\varepsilon_\nu^{(h_i)}\)。反過來，若給每個 \(z_i\) 配以一個 \(h_i\) 次單位根 \(\varepsilon_\nu^{(h_i)}\)，便得到 \(\mathfrak o\) 在 \(\Omega\) 中的一個不可約表示。由於 \(h\not\equiv0(p)\)，從而更有 \(h_i\not\equiv0(p)\)，所以 \(\mathfrak o\) 恰有 \(h_1h_2\ldots h_t=h\) 個不同的表示，證畢。

（也很容易看出，若 \(h\equiv0(p)\)，從而至少有一個 \(h_i\) 滿足 \(h_i\equiv0(p)\)，那麼至多有 \(h_i/p\) 個不同的 \(\varepsilon_\nu^{(h_i)}\)，因而表示的總數少於 \(h\)；所以 \(\mathfrak o_\Omega\) 必須含有一個非零根基。但這並不能直接推出任何關於 \(\mathfrak o\) 的根基的結論。）

我們用 \(\Theta_1,\ldots,\Theta_h\) 表示從 \(\mathfrak o_\Omega\) 到 \(\Omega\) 的各子環的 \(h\) 個不同的同態映射；每個 \(\Theta_i\) 都有 \(\mathfrak o_\Omega\) 的一個直和分量 \(e_i\Omega\) 作為其表示模。有時我們只把 \(\Theta_i\) 作用於 \(\mathfrak A\)，此時便稱它為 \(\mathfrak A\) 的一個表示。這些從 \(\mathfrak A\) 到單位根所成的群的映射，就是 \(\mathfrak A\) 的\emph{\(h\) 個特徵標}。其中一個特殊的特徵標是\emph{主特徵標}，它把 \(\mathfrak A\) 的每個元素都映到 \(1\)。設相應的 \(\Theta_i\) 為 \(\Theta_1\)。

由主特徵標給出的表示對每個群都存在，而不只對阿貝爾群存在。由於這個表示已經位於 \(\mathsf P\) 中，所以當 \(\mathfrak o\) 完全可約時，\(\mathfrak o\) 必須分離出一個給出這一表示的簡單直和分量 \(E\cdot\mathfrak o\)。我們可以直接算出 \(E\)。它必有形式 \(E=\sum x_i a_i\)。首先，
\[
aE=E\cdot 1\quad\text{或}\quad \sum x_i a a_i=\sum x_i a_i
\]
由此，對 \(a=a_j a_i^{-1}\)，有
\[
x_1a_ja_i^{-1}a_1+\cdots+x_i a_j+\cdots+x_na_ja_i^{-1}a_n
=x_1a_1+\cdots+x_na_n
\]
因而
\[
x_i=x_j
\]

其次，

\[
E=E^2=x^2\left(\sum a_i\right)^2=x^2h\sum a_i=xhE.
\]

在 \(h\equiv0(p)\) 的情形下，由此得 \(E=0\)；也就是說，\(\mathfrak o\) 不包含主特徵標的表示模，因而不可能是完全可約的。這樣，第 13 頁上的定理又有一半得到了證明。

但若 \(h\not\equiv0(p)\)，從而 \(\mathfrak o\) 完全可約，則立即得到

\[
x=\frac{1}{h},\qquad E=\frac{1}{h}(a_1+\cdots+a_h).
\]

\subsection*{§ 16. 特徵標關係}\label{die-charakterenrelationen}

\(E=(1/h)\cdot\sum a_i\) 給出主特徵標這一事實，立即導出特徵標關係。必有

\[a \cdot e_i = e_i \cdot \Theta_i a\]

因而

\[e_i \cdot \Theta_i e_j = \begin{cases} e_i & i = j \\ 0 & i \neq j \end{cases}\]

或者

\[\Theta_i e_j = \begin{cases} 1 & i = j \\ 0 & i \neq j \end{cases}\]

特別地，

\[\Theta_i e_1 = \begin{cases} 1 & i = 1 \\ 0 & i \neq 1 \end{cases}\]

或者，由於 \(e_1 = (1/h) \sum a_i\)，

\[
\sum_\nu \Theta_i a_\nu=\begin{cases} h & i=1\\ 0 & i\ne 1 \end{cases}
\]

這就是熟知的特徵標關係：

\emph{一個特徵標在所有群元素上的和為零，主特徵標除外；對主特徵標，這個和等於 \(h\)。}

\subsection*{§ 17. 阿貝爾群的伽羅瓦理論}\label{die-galoissche-theorie-abelscher-gruppen}

對於阿貝爾群的群環，可以建立與交換域的伽羅瓦理論中所證明的定理相似的定理。在伽羅瓦理論中，各 \(\Theta_i\Theta_1^{-1}\) 組成一個群 \(\mathfrak{G}\)。主定理表明，\(\mathfrak{G}\) 的子群與 \(\mathsf{Z}\) 的子域之間存在一個以某些不變性為基礎的一一對應。

\(\mathfrak{o}[\mathfrak{A}]\) 的各同態 \(\Theta_i\) 可以合併成一個群 \(\mathsf A\)，它的子群同 \(\mathfrak{A}\) 的子群之間有一種類似的關係。

\emph{若通過下式定義乘積 \(\Theta_i \Theta_k = \Theta\)：}

\[\Theta a = \Theta_i a \cdot \Theta_k a,\]

\emph{則全體 \(\Theta_i\) 組成一個與 \(\mathfrak A\) 同構的群 \(\mathsf A\)。}

\emph{證明。} 由於

\[\begin{aligned}
\Theta a \cdot \Theta b &= \Theta_i a \cdot \Theta_k a \cdot \Theta_i b \cdot \Theta_k b \\
&= \Theta_i a \cdot \Theta_i b \cdot \Theta_k a \cdot \Theta_k b \\
&= \Theta_i a b \cdot \Theta_k a b = \Theta a b
\end{aligned}\]

所以 \(\Theta = \Theta_i \cdot \Theta_k\) 是從 \(\mathfrak A\) 到 \(\mathsf A\) 的一個同態；又因為除了各 \(\Theta_i\) 之外，\(\mathfrak A\) 沒有其他同態，所以 \(\Theta\) 等於某個 \(\Theta_j\)。因此，各 \(\Theta_i\) 的確組成一個含有 h 個元素的群。

現在令 \(\varepsilon_{h_\nu}\) 表示一個本原 \(h_\nu\) 次單位根。於是，\(\Theta_i\) 把 \(\mathfrak A\) 的循環因子 \(\mathfrak Z_\nu\) 的生成元 \(z_\nu\) 映到 \(\varepsilon_{h_\nu}\) 的一個冪：

\[
\Theta_i z_\nu=\varepsilon_{h_\nu}^{\varrho_\nu^{(i)}}.
\]

我們把 \(\mathfrak A\) 的元素 \(\prod_\nu z_\nu^{\varrho_\nu^{(i)}}\) 配給同態 \(\Theta_i\)。這是從 \(\mathsf A\) 到 \(\mathfrak A\) 的一一映射；事實上，由於 \(\Theta_i\) 由各 \(\varrho_\nu^{(i)}\) 確定，不同的 \(\Theta_i\) 對應於不同的 \(\prod_\nu z_\nu^{\varrho_\nu^{(i)}}\)。這個映射是同構。因為一方面，

\[
\Theta_i\Theta_j z_\nu=\Theta_i z_\nu\cdot \Theta_j z_\nu
=\varepsilon_{h_\nu}^{\varrho_\nu^{(i)}+\varrho_\nu^{(j)}},
\]

另一方面，

\[
\prod_\nu z_\nu^{\varrho_\nu^{(i)}}\cdot
\prod_\nu z_\nu^{\varrho_\nu^{(j)}}=
\prod_\nu z_\nu^{\varrho_\nu^{(i)}+\varrho_\nu^{(j)}},
\quad\text{證畢。}
\]

現在我們定義：

\emph{\(\mathsf A\) 的一個子群 \(\mathsf B\) 的不變範圍：}

由 \(\mathfrak A\) 中所有被每個 \(\Theta\in\mathsf B\) 映到 1 的元素組成的子群 \(\mathfrak B\)。

\emph{\(\mathfrak A\) 的一個子群 \(\mathfrak B\) 的不變群：}

由 \(\mathsf A\) 中所有把每個 \(a\in\mathfrak B\) 映到 1 的元素組成的子群 \(\mathsf B\)。

並表述主定理如下：

\emph{\(\mathfrak A\) 的子群 \(\mathfrak B\) 與 \(\mathsf A\) 的子群 \(\mathsf B\) 可以彼此一一對應，使得當 \(\mathfrak B\) 與 \(\mathsf B\) 相互對應時，\(\mathsf B\) 是 \(\mathfrak B\) 的不變群，而 \(\mathfrak B\) 是 \(\mathsf B\) 的不變範圍。}

所以我們必須證明：

\begin{enumerate}
\item \emph{若 \(\mathsf B\) 是 \(\mathfrak B\) 的不變群，則 \(\mathfrak B\) 是 \(\mathsf B\) 的不變範圍。}
\item \emph{若 \(\mathfrak B\) 是 \(\mathsf B\) 的不變範圍，則 \(\mathsf B\) 是 \(\mathfrak B\) 的不變群。}
\end{enumerate}

\paragraph{1.}
如同第 10 頁那樣，這可由一般延拓定理推出。我們假設 \(\mathsf B\) 是 \(\mathfrak B\) 的不變群，\(\mathfrak B'\) 是 \(\mathsf B\) 的不變範圍；於是 \(\mathfrak B'\) 包含 \(\mathfrak B\)。\(\mathsf B\) 的各元素，是 \(\mathfrak B\) 上那個把 \(\mathfrak B\) 映到 1 的同態在整個 \(\mathfrak A\) 上的各個延拓。若只把這些延拓限制到 \(\mathfrak B'\) 上，則由一般延拓定理可知，它們分成 \(t\) 個類，使得同一類中的同態彼此相同，而不同類中的同態在 \(\mathfrak B'\) 上不同；這裡 \(t\) 是 \(\mathfrak o[\mathfrak B']\) 在 \(\mathfrak o[\mathfrak B]\) 上的次數，也就是 \(\mathfrak B\) 在 \(\mathfrak B'\) 中的指數，因為 \(\mathfrak o[\mathfrak B]\) 和 \(\mathfrak o[\mathfrak B']\) 有同一個單位元，即整個 \(\mathfrak A\) 的單位元。由於 \(\mathfrak B'\) 是 \(\mathsf B\) 的不變範圍，所以必有 \(t=1\)，從而 \(\mathfrak B'=\mathfrak B\)，證畢。

\textbf{2.} 我們把它歸結為 1。

我們把 \(\mathfrak A\) 的元素 \(a\) 看作 \(\mathsf A\) 上的同態。具體地說，我們把 \(a\Theta\) 定為單位根 \(\Theta a\)，簡寫為 \(a\Theta=\Theta a\)。這樣，\(a\) 的確被定義為 \(\mathsf A\) 的一個同態，因為
\[
a\Theta\cdot a\Theta'=\Theta a\cdot\Theta' a=\Theta\Theta' a=a\Theta\Theta'.
\]

兩個同態 \(a,b\) 的乘積就是通常的群乘積：
\[
\begin{array}{cccccc}
ab\Theta&=&a\Theta\cdot b\Theta&=&\Theta a\cdot\Theta b=\Theta ab&=ab\Theta\\
&&\text{同態的乘積}&&\text{群的乘積}&
\end{array}
\]

作為群元素彼此不同的 \(a,b\)，作為同態也彼此不同：若對所有 \(\Theta\) 都有 \(a\Theta=b\Theta\)，則對所有 \(\Theta\) 都有 \(\Theta a=\Theta b\)，或者說，對所有 \(\Theta\) 都有 \(\Theta ab^{-1}=1\)。因此，各 \(\Theta\) 是由 \(ab^{-1}\) 生成的 \(\mathfrak A\) 的循環子群 \(\mathfrak Z\) 上把所有元素映到 1 的同態在整個 \(\mathfrak A\) 上的各個延拓。所以，由延拓定理，\(\mathfrak Z\) 在 \(\mathfrak A\) 中的指數等於 \(h\)，也就是說，\(\mathfrak Z=e,\ a=b\)。

現在，我們把在這一新觀點下作出的陳述
\[
\begin{array}{ll}
\text{a)}& \mathfrak B\ \text{是 }\mathsf B\text{ 的不變群。}\\
\text{b)}& \mathsf B\ \text{是 }\mathfrak B\text{ 的不變範圍。}
\end{array}
\]
翻譯回原來的觀點：

對於 a)，\(\mathfrak B\) 由所有滿足對每個 \(\Theta\in\mathsf B\) 都有 \(a\Theta=1\) 的 \(a\in\mathfrak A\) 組成；也就是說，\(\mathfrak B\) 由所有滿足對每個 \(\Theta\in\mathsf B\) 都有 \(\Theta\cdot a=1\) 的 \(a\in\mathfrak A\) 組成。因此，a) 翻譯為
\[
\text{a')} \quad \mathfrak B\ \text{是 }\mathsf B\text{ 的不變範圍。}
\]

相應地，b) 翻譯為
\[
\text{b')} \quad \mathsf B\ \text{是 }\mathfrak B\text{ 的不變群。}
\]

所以，在 1 中已經證明的命題

若 \(\mathfrak B\) 是 \(\mathsf B\) 的不變群，則 \(\mathsf B\) 是 \(\mathfrak B\) 的不變範圍，

翻譯後恰好就是待證命題：

若 \(\mathfrak B\) 是 \(\mathsf B\) 的不變範圍，則 \(\mathsf B\) 是 \(\mathfrak B\) 的不變群。

若把第 15 頁上 \(\mathfrak A\) 的特徵標關係改為對 \(\mathsf A\) 寫出，便得到等式
\[
\sum_{i=1}^{h} \Theta_i a_\nu =
\begin{cases}
h & \nu=1\\
0 & \nu\ne 1.
\end{cases}
\]

所有特徵標在一個群元素上的和為零，單位元除外；對單位元，這個和等於 \(h\)。

\section*{第四章. 雙邊單環}\label{kapitel-iv.-zweiseitig-einfache-ringe}
\subsection*{§ 18. 一個引理}\label{ein-hilfssatz}

我們以 \(\mathsf K\) 表示一個域，以 \(\mathfrak G\) 表示 \(\mathsf K\) 的一組自同構（即它到自身的同構映射）。在屬於 \(\mathfrak G\) 的每一個自同構下都保持不變的全體 \(\mathsf K\)-元素，組成 \(\mathsf K\) 的一個子域 \(\mathsf P\)，稱為 \(\mathfrak G\) 的\emph{不變域}。

若 \(\tau\ne 0\) 是任意一個 \(\mathsf K\)-元素，則把每個 \(\mathsf K\)-元素 \(\alpha\) 映到 \(\tau^{-1}\alpha\tau\)，便得到 \(\mathsf K\) 的一個自同構；仿照群論中的相應名稱，我們稱這種自同構為\emph{內自同構}。域 \(\mathsf K\) 的全體內自同構所成群的不變域 \(\mathsf P\)，正是 \(\mathsf K\) 的中心。

現在給定域 \(\mathsf K\)、它的一個自同構群 \(\mathfrak G\)，以及該群的不變域 \(\mathsf P\)。設

\[
\mathfrak M=x_1\mathsf P+\cdots+x_n\mathsf P
\]

是秩為 \(n\) 的一個 \(\mathsf P\)-模。我們定義 \(\mathfrak G\)-元素 \(\gamma\) 在 \(\mathfrak M\)-元素上的作用為

\[
\gamma\cdot x_i=x_i,\qquad\text{因而}\qquad
\gamma\sum \varrho_i x_i=\sum \varrho_i x_i.
\]

於是，\(\mathfrak G\) 成為擴張模

\[
\mathfrak M_{\mathsf K}=x_1\mathsf K+\cdots+x_n\mathsf K
\]

的一個自同構群；而 \(\mathfrak G\) 在 \(\mathfrak M_{\mathsf K}\) 中的不變區域就是 \(\mathfrak M\)。在這些假設下有如下定理：

\begingroup\itshape
若
\[
\mathfrak C=z_1\mathsf K+\cdots+z_r\mathsf K
\]

是 \(\mathfrak M_{\mathsf K}\) 的一個子模，並且每個 \(\mathfrak G\)-元素 \(\gamma\) 都把它映到自身（不必逐點不變），則 \(\mathfrak C\) 是 \(\mathfrak M\) 的某個子模 \(\mathfrak c\) 以 \(\mathsf K\) 擴張所得的擴張模：

\[
\mathfrak C=\mathfrak c_{\mathsf K}.
\]
\endgroup

\emph{證明。} 1. 先作一點說明：必有 \(\mathfrak c=\mathfrak C\cap\mathfrak M\)。因為一般地都有
\[
\mathfrak c=\mathfrak c_{\mathsf K}\cap\mathfrak M.
\]

\(\mathfrak c\subseteqq\mathfrak c_{\mathsf K}\cap\mathfrak M\) 是顯然的。為證明 \(\mathfrak c_{\mathsf K}\cap\mathfrak M\subseteqq\mathfrak c\)，取 \(\mathfrak c\) 的一組基 \(y_1,\ldots,y_r\)，

\[
\mathfrak c=y_1\mathsf P+\cdots+y_r\mathsf P,
\]

於是

\[
\mathfrak c_{\mathsf K}=y_1\mathsf K+\cdots+y_r\mathsf K.
\]

每個 \(\mathfrak c_{\mathsf K}\cap\mathfrak M\)-元素 \(c\)，作為 \(\mathfrak c_{\mathsf K}\)-元素有基表示
\[
c=\sum_{i=1}^{r}y_i\varkappa_i
\]
而作為 \(\mathfrak M\)-元素有基表示

\[
c=\sum_{i=1}^{r}y_i\varrho_i+\sum_{j=r+1}^{n}y_j\varrho_j
\]

這裡 \(y_{r+1},\ldots,y_n\) 是一組 \(\mathfrak M\)-元素，把 \(y_1,\ldots,y_r\) 補成 \(\mathfrak M\) 的一組基（M. Z. § 6）。

兩種表示都可以看作 \(\mathfrak M_{\mathsf K}\) 中的基表示，所以它們相同：\(\varkappa_i=\varrho_i\ (i=1,\ldots,r)\)，\(\varrho_j=0\ (j=r+1,\ldots,n)\)。因此

\[
c=\sum_{i=1}^{r}y_i\varrho_i
\qquad\text{證畢。}
\]

\textbf{2.} 為證明所述定理，我們構造 \(\mathfrak C\) 的一組基 \(z_1,\ldots,z_r\)，使每個 \(\mathsf{Z}_i\) 都屬於 \(\mathfrak M\)。這樣，若置

\[
\mathfrak c=z_1\mathsf P+\cdots+z_r\mathsf P
\]

就確實有

\[
\mathfrak C=\mathfrak c_{\mathsf K}.
\]

任取 \(\mathfrak C\) 的一組基 \(y_1,\ldots,y_r\)，再從 \(\mathfrak M\) 的一組基 \(x_1,\ldots,x_n\) 中選取適當的元素 \(x_{r+1},\ldots,x_n\)，把它補成 \(\mathfrak M_{\mathsf K}\) 的一組基；後者當然也是 \(\mathfrak M_{\mathsf K}\) 的基。按照 M. Z. § 5，這是可行的。

\[
\mathfrak M_{\mathsf K}=\mathfrak C+x_{r+1}\mathsf K+\cdots+x_n\mathsf K.
\]

由此，對於 \(\mathfrak M\) 的基的前 \(r\) 個元素 \(x_1,\ldots,x_r\)，有等式

\[
x_i=z_i-\sum_{j=r+1}^{n}x_j\varkappa_j^{(i)},\qquad
z_i\in\mathfrak C,\quad \varkappa_j^{(i)}\in\mathsf K.
\]

顯然，\(z_i\) 與 \(x_{r+1},\ldots,x_n\) 一起構成 \(\mathfrak M_{\mathsf K}\) 的一組基，故它們線性無關；因此 \(z_1,\ldots,z_r\) 是 \(\mathfrak C\) 的一組基。\(z_1,\ldots,z_r\) 正是我們所求形式的基：各 \(z_i\) 都屬於 \(\mathfrak M\)。

設 \(\gamma\) 是一個 \(\mathfrak G\)-元素。一方面有

\[
\begin{array}{rcl}
(1)\qquad \gamma z_i
&=&\gamma x_i+\displaystyle\sum_{j=r+1}^{n}\gamma x_j\cdot \gamma\varkappa_j^{(i)}
=x_i+\displaystyle\sum_{j=r+1}^{n}x_j\gamma\varkappa_j^{(i)}.
\end{array}
\]

另一方面，
\[
\begin{array}{rcl}
(2)\qquad \gamma z_i&=&\displaystyle\sum_{\mu=1}^{r}z_\mu\sigma_\mu^{(i)},
\qquad \sigma_\mu^{(i)}\in\mathsf K,
\end{array}
\]

因為 \(\mathfrak C\) 容許所有 \(\gamma\in\mathfrak G\) 的作用。於是還可把 (2) 寫成

\[
\begin{array}{rcl}
(3)\qquad \gamma z_i
&=&\displaystyle\sum_{\mu=1}^{r}x_\mu\sigma_\mu^{(i)}
+\displaystyle\sum_{j=r+1}^{n}x_j\sum_{\mu=1}^{r}\varkappa_j^{(\mu)}\sigma_\mu^{(i)}
\end{array}
\]

比較 (1) 與 (3) 中的係數，得到

\[
\sigma_j^{(i)}=
\begin{cases}
1 & i=j,\\
0 & i\ne j.
\end{cases}
\]

所以

\[
\gamma\cdot z_i=z_i.
\]

因此各 \(z_i\) 屬於 \(\mathfrak G\) 的不變區域，也就是屬於 \(\mathfrak M\)，證畢。

\subsection*{§ 19. 雙邊單超複系統在其係數域的非交換擴域中的表示}\label{darstellungen-zweiseitig-einfacher-hyperkomplexer-systeme-in-nichtkommutativen-erweiterungskoerpern-ihrer-koeffizientenkoerper}

設

\[
\mathfrak S=x_1\mathsf P+\cdots+x_n\mathsf P
\]

是一個雙邊單、從而完全可約的超複系統，而 \(\mathsf K\) 是以 \(\mathsf P\) 為中心的一個域。

\textbf{定理 1.} \emph{\(\mathfrak S_{\mathsf K}\) 是雙邊單的（因而也完全可約）。}

\emph{證明。} 1. \(\mathfrak S_{\mathsf K}\) 的每個雙邊理想 \(\mathfrak a\)，在 \(\mathsf K\) 的內自同構群 \(\mathfrak G\) 下都不變。事實上，若 \(\varkappa\) 是一個非零的 \(\mathsf K\)-元素，並置 \(\varkappa^{-1}\mathfrak a\varkappa=\bar{\mathfrak a}\)，則由 \(\mathfrak a\) 的理想性質有 \(\bar{\mathfrak a}\subseteqq\mathfrak a\)；但把 \(\bar{\mathfrak a}\) 看成 \(\mathsf K\)-模時，它在 \(\mathsf K\) 上的秩與 \(\mathfrak a\) 相同，所以必有 \(\bar{\mathfrak a}=\mathfrak a\)。

\textbf{2.} 按剛剛建立的引理，\(\mathfrak S_{\mathsf K}\) 的每個雙邊理想 \(\mathfrak a\) 必是 \(\mathfrak S\) 的某個子模 \(\mathfrak A\) 的擴張理想，即 \(\mathfrak a=\mathfrak A_{\mathsf K}\)。\(\mathfrak A\) 是 \(\mathfrak S\) 的雙邊理想；因為 \(\mathfrak A=\mathfrak S\cap\mathfrak a\)，故
\[
\mathfrak S\mathfrak A\subseteqq\mathfrak S\mathfrak S\cap\mathfrak S\mathfrak a\subseteqq\mathfrak S\cap\mathfrak a=\mathfrak A
\]
\[
\mathfrak A\mathfrak S\subseteqq\mathfrak S\mathfrak S\cap\mathfrak a\mathfrak S\subseteqq\mathfrak S\cap\mathfrak a=\mathfrak A.
\]

而 \(\mathfrak S\) 只有兩個雙邊理想 0 與 \(\mathfrak S\)，所以 \(\mathfrak S_{\mathsf K}\) 也只有兩個雙邊理想 0 與 \(\mathfrak S_{\mathsf K}\)，證畢。

現在我們來研究 \(\mathfrak S\) 在 \(\mathsf K\) 中的不可約表示。我們將建立互反表示，因為這樣在形式上較為便利。歸根結底，採用直接表示還是互反表示並無區別，因為兩者一一對應。

設 \(\mathfrak M\) 是 \(\mathfrak S\) 在 \(\mathsf K\) 中的一個互反表示模，也就是有限秩的 \(\mathfrak S\)-左模兼 \(\mathsf K\)-左模，
\[
\mathfrak M=\mathsf K\cdot y_1+\cdots+\mathsf K\cdot y_r
\]
並滿足交換律
\[
s(\varkappa m)=\varkappa (sm).
\]

我們證明，\(\mathfrak M\) 也可看作 \(\mathfrak S_{\mathsf K}\)-左模。首先必須說明一個 \(\mathfrak S_{\mathsf K}\)-元素與一個 \(\mathfrak M\)-元素的乘法。若 \(s\in\mathfrak S,\ \varkappa \in\mathsf K,\ m\in\mathfrak M\)，置
\[
s\varkappa \cdot m=s\cdot(\varkappa \cdot m)
\]
相應地，對於一般的 \(\mathfrak S_{\mathsf K}\)-元素 \(\sum s_i\varkappa _i\)，置
\[
\sum s_i\varkappa _i\cdot m=\sum s_i\cdot(\varkappa _i\cdot m).
\]

可以驗證，這個定義是唯一的，不依賴於把一個 \(\mathfrak S_{\mathsf K}\)-元素表示為各乘積 \(s_i\varkappa _i\) 之和的方式。

所以問題在於由
\[
\begin{array}{rcl}
(1)\qquad \displaystyle\sum_{i=1}^{p}s_i\varkappa _i&=&0
\end{array}
\]
推出
\[
\sum_{i=1}^{p}s_i(\varkappa _i m)=0
\]
假設 \(s_1,\ldots,s_q\) 線性無關，而 \(s_{q+1},\ldots,s_p\) 可用 \(s_1,\ldots,s_q\) 表出：
\[
s_j=\sum_{i=1}^{q}s_i\varrho_i^{(j)},\qquad j=q+1,\ldots,p.
\]

各系數必須屬於 \(\mathsf P\)，因為 \(s_j\) 的這種表示是唯一的，而這個表示已經可以在 \(\mathfrak S\) 中作出。(1) 給出

\[
\varkappa_i+\sum_{j=q+1}^{p}\varrho_i^{(j)}\varkappa_j=0,
\qquad i=1,\ldots,q.
\]

因此

\[
\begin{aligned}
\sum_{i=1}^{p}s_i(\varkappa _i m)
&=s_1(\varkappa _1m)+\cdots+s_q(\varkappa _qm)
  +\left(\sum_i s_i\varrho_i^{(q+1)}\right)(\varkappa _{q+1}m)\\
&\quad+\cdots+\left(\sum_i s_i\varrho_i^{(p)}\right)(\varkappa _pm)\\
&=s_1(\varkappa _1m)+\cdots+s_q(\varkappa _qm)\\
&\quad+s_1\varrho_1^{(q+1)}(\varkappa _{q+1}m)+\cdots+s_q\varrho_q^{(q+1)}(\varkappa _{q+1}m)
  \qquad\text{由模的性質}\\
&\quad+\cdots\\
&\quad\cdots\\
&\quad+s_1\varrho_1^{(p)}(\varkappa _pm)+\cdots+s_q\varrho_q^{(p)}(\varkappa _pm)\\
&=s_1(\varkappa _1m)+\cdots+s_q(\varkappa _qm)\\
&\quad+s_1(\varrho_1^{(q+1)}\varkappa _{q+1}m)+\cdots+s_q(\varrho_q^{(q+1)}\varkappa _{q+1}m)\\
&\quad+\cdots\\
&\quad\cdots\\
&\quad+s_1(\varrho_1^{(p)}\varkappa _pm)+\cdots+s_q(\varrho_q^{(p)}\varkappa _pm)
  \qquad\begin{gathered}\text{由於}\\ \text{结合律 }s'(sm)\\ =s's\cdot m\text{ 以及}\\ \varkappa '(\varkappa m)=\varkappa '\varkappa \cdot m\end{gathered}\\
&=s_1\left(\varkappa _1m+\sum\varrho_1^{(q+1)}\varkappa _{q+1}\cdot m\right)+\cdots\\
&=s_1\left(\left(\varkappa _1+\sum\varrho_1^{q+1}\varkappa _{q+1}\right)m\right)+\cdots\\
&=0.
\end{aligned}
\]

在 \(\mathfrak M\) 的各項 \(\mathfrak S_{\mathsf K}\)-模性質中，分配律 \((S+S')m=Sm+S'm,\ S(m+m')=Sm+Sm'\) 是清楚的。還須注意結合律 \(S(S'm)=SS'\cdot m\)。由於分配律成立，我們可以只考察情形

\[
s\varkappa \cdot(s'\varkappa ' m)=s\varkappa (s'\varkappa '\cdot m)
\]

一方面，按照 \(s\varkappa \cdot m\) 的定義，

\[
s\varkappa \cdot(s'\varkappa '\cdot m)=s\cdot(\varkappa \cdot(s'\cdot(\varkappa '\cdot m)))
\]

另一方面

\[
s\varkappa  s'\varkappa '\cdot m=ss'\varkappa \varkappa '\cdot m=ss'\cdot(\varkappa \varkappa '\cdot m)
=s\cdot(s'\cdot(\varkappa \cdot(\varkappa '\cdot m)))
\]

由兩個交換關係，後一式也就是 \(=s\cdot(\varkappa \cdot(s'\cdot(\varkappa '\cdot m)))\)，因為

\[
\varkappa \cdot s'=s'\cdot \varkappa
\]
\[
s'\cdot(\varkappa \cdot m)=\varkappa \cdot(s'\cdot m)
\]

所以 \(\mathfrak M\) 確實是一個 \(\mathfrak S_{\mathsf K}\)-左模，而且當然是有限的。反過來，每個有限 \(\mathfrak S_{\mathsf K}\)-左模都是 \(\mathfrak S\) 在 \(\mathsf K\) 中的一個互反表示模；因為它顯然是有限 \(\mathsf K\)-模，而且由

\[
s(\varkappa \cdot m)=s\varkappa \cdot m=\varkappa s\cdot m=\varkappa (s\cdot m)
\]

也滿足交換條件。（另請比較 § 24 中較為簡單的證明。）

現在，按 M. Z. § 18，每個有限 \(\mathfrak S_{\mathsf K}\)-模都是簡單模之和，而這些簡單模與 \(\mathfrak S_{\mathsf K}\) 的簡單左理想在 \(\mathfrak S_{\mathsf K}\)-算子的意義下同構。反過來，\(\mathfrak S_{\mathsf K}\) 的每個左理想都是 \(\mathfrak S\) 在 \(\mathsf K\) 中的表示模；由於 \(\mathfrak S_{\mathsf K}\) 的所有簡單左理想在算子的意義下同構，故有如下定理：

\textbf{定理 2.} \emph{\(\mathfrak S\) 在 \(\mathsf K\) 中唯一的不可約互反表示類，是由 \(\mathfrak S_{\mathsf K}\) 的一個簡單左理想生成的表示類。}

若這個不可約表示的次數為 \(r\)，則任一表示的次數都是 \(r\) 的倍數；當然，可以把不可約表示並列接合，從而對 \(r\) 的任一給定倍數 \(rs\) 構造一個 \(rs\) 次表示。這對應於構成一個表示模，它是 \(s\) 個簡單模的直和。

我們已經知道 \(\mathfrak S\) 在 \(\mathsf K\) 中的一個表示，即以 \(\mathfrak S\) 本身為表示模而生成的表示（主表示，M. Z. § 20）；這個表示已經位於 \(\mathsf P\) 中。它的次數是 \(n\)，所以 \(n\) 能被 \(r\) 整除：

\[
n=rt.
\]

商數 \(t\) 很容易確定：\(r\) 是一個單獨左理想在 \(\mathsf K\) 上的秩，\(n\) 是整個 \(\mathfrak S_{\mathsf K}\) 的秩，所以 \(t\) 是 \(\mathfrak S_{\mathsf K}\) 分解成簡單左理想時的項數；從而 \(t^2\) 是 \(\mathfrak S_{\mathsf K}\) 在其簡單理想的自同構域上的秩。

現在把我們的考察專門化到 \(\mathfrak S\) 是交換域的情形。

\textbf{定理 3.} \emph{若 \(\mathfrak S\) 是交換域，則 \(\mathfrak S\) 是 \(\mathfrak S_{\mathsf K}\) 的中心。}

\emph{證明。} 每個 \(\mathfrak S\)-元素都與每個 \(\mathfrak S\)-元素和每個 \(\mathsf K\)-元素可交換，因而也與每個 \(\mathfrak S_{\mathsf K}\)-元素可交換，所以 \(\mathfrak S\) 包含在 \(\mathfrak S_{\mathsf K}\) 的中心中。一個 \(\mathfrak S_{\mathsf K}\)-元素只有在至少與每個 \(\mathsf K\)-元素可交換時，才與每個 \(\mathfrak S_{\mathsf K}\)-元素可交換；這就是說，它的係數必須屬於 \(\mathsf K\) 的中心，即 \(\mathsf P\)，亦即它必須屬於 \(\mathfrak S\)，證畢。

對於交換環，互反同構與直接同構沒有區別，所以每個互反表示也可看作直接表示。因而，\(\mathfrak S\) 在 \(\mathsf K\) 中的不可約表示也必須由一個直接表示模給出。這個直接表示模是 \(\mathfrak S_{\mathsf K}\) 的一個簡單右理想。事實上，\(\mathfrak r\) 是秩為 \(r\) 的有限 \(\mathsf K\)-右模，並且由於 \(\mathfrak r\mathfrak S=\mathfrak S\mathfrak r\)，也是 \(\mathfrak S\)-左模（當然滿足貫通的結合律），所以它是秩為 \(r\) 的一個直接表示模；而因為只有一個不可約表示類，它就必須生成給定的表示。

\subsection*{§ 20. 非交換域}\label{nichtkommutative-koerper}

現在以

\[
\mathfrak K=y_1\mathsf P+\cdots+y_m\mathsf P
\]

表示一個以 \(\mathsf P\) 為中心的非交換域。我們將為 \(\mathfrak K\) 建立一系列定理，它們類似於交換有限代數域擴張理論中的相應定理。

若 \(\mathsf{Z}=\xi_1\mathsf P+\cdots+\xi_n\mathsf P\) 表示 \(\mathsf P\) 上的一個超複系統，則兩個擴張系統 \(\mathfrak K_{\mathsf{Z}}\) 與 \(\mathsf{Z}_{\mathfrak K}\) 相同，因為對於 \(\mathfrak K_{\mathsf{Z}}\) 可以寫成

\[
\mathfrak K_{\mathsf{Z}}=y_1\xi_1\mathsf P+\cdots+y_m\xi_n\mathsf P
\]

而對於 \(\mathsf{Z}_{\mathfrak K}\) 可以寫成

\[
\mathsf{Z}_{\mathfrak K}=\xi_1y_1\mathsf P+\cdots+\xi_ny_m\mathsf P
\]

由於 \(y_i\xi_j=\xi_jy_i\)，這兩個表達式相同。

下文中的 \(\mathsf{Z}\) 始終是一個交換域，所以可以應用 § 19 的各定理。按 § 19 的定理 1，\(\mathfrak K_{\mathsf{Z}}\) 總是雙邊單的。

\textbf{定理 1.} \emph{設 \(\Omega\) 是 \(\mathsf P\) 的一個代數閉擴域，則 \(\mathfrak K_\Omega\) 的中心等於 \(\Omega\)。}

\emph{證明。} 我們依據 § 19 的相應定理 3。\(\Omega\) 包含在 \(\mathfrak K_\Omega\) 的中心中，這是清楚的。若 \(\alpha\) 是 \(\mathfrak K_\Omega\) 中心的任一元素，則它的係數屬於 \(\mathsf P\) 的某個有限代數擴域 \(Z\)（即由這些係數生成的擴域）。當然，\(\alpha\) 也必須屬於 \(\mathfrak K_{\mathsf{Z}}\) 的中心；而這個中心就是 \(\mathsf{Z}\)，所以 \(\alpha\) 屬於 \(\mathsf{Z}\)，從而也屬於 \(\Omega\)，證畢。

以完全類似的方式可以證明：

\textbf{定理 2.} \emph{\(\mathfrak K_\Omega\) 是雙邊單的。}

\emph{證明。} 設 \(\mathfrak a=z_1\Omega+\cdots+z_r\Omega\) 是一個雙邊理想。\(z_1,\ldots,z_r\) 的係數屬於 \(\mathsf P\) 的某個有限擴域 \(\mathsf{Z}\)；在 \(\mathfrak K_{\mathsf{Z}}\) 中，\(\mathfrak a'=z_1\mathsf{Z}+\cdots+z_r\mathsf{Z}\) 是雙邊理想。因為 \(\mathfrak K_{\mathsf{Z}}\) 雙邊單，所以 \(\mathfrak a'\) 或為零，或等於 \(\mathfrak K_{\mathsf{Z}}\)；因而 \(\mathfrak a\) 或為 0，或等於 \(\mathfrak K_\Omega\)。也就是說，\(\mathfrak K_\Omega\) 雙邊單，證畢。

按這個定理，\(\mathfrak K_\Omega\) 是其中心 \(\Omega\) 的某個擴域上的矩陣環。這個域作為整個 \(\mathfrak K_\Omega\) 的子環，在 \(\Omega\) 上必須是有限次的，所以它只能等於 \(\Omega\)：

\textbf{定理 3.} \emph{\(\mathfrak K_\Omega\) 是 \(\Omega\) 上的矩陣環，}

\[
\mathfrak K_\Omega=\sum\Omega c_{ik}
\]

\emph{因而 \(\mathfrak K\) 在其中心 \(\mathsf P\) 上的次數是一個平方數 \(m=t^2\)。}

\(\mathfrak K_\Omega\) 所分解成的右（或左）理想的數目 \(t\)，簡稱為 \(\mathfrak K\) 的\emph{絕對分量數}。

\(\Omega\) 所具有的性質，即 \(\mathfrak K_\Omega\) 分解為 \(t\) 個簡單理想，當然也已為 \(\mathsf P\) 的某些有限擴域 \(\mathsf{Z}\) 所具有；例如，把各 \(c_{ik}\) 的係數添加到 \(\mathsf P\) 所得到的那些 \(\mathsf{Z}\)。現在研究這些分裂域。

\textbf{定義。} 若 \(\mathsf P\) 的有限代數擴域 \(Z\) 使 \(\mathfrak K_{\mathsf{Z}}\) 所分解成的簡單右理想數已經等於絕對分量數 \(t\)，則稱 \(Z\) 為 \(\mathfrak K\) 的\emph{分裂域}\srcfn{1}{請比較第 9 頁的分裂域概念}。

\textbf{定理 4.} \emph{若 \(\mathsf{Z}\) 是 \(\mathfrak K\) 的分裂域，則 \(\mathsf{Z}\) 是 \(\mathfrak K_{\mathsf{Z}}\) 的各右理想的自同構域；反過來，若 \(\mathfrak K_{\mathsf{Z}}\) 的各右理想的自同構域等於 \(\mathsf{Z}\)，則 \(\mathsf{Z}\) 是分裂域。}

\emph{證明。} 1. 設 \(\mathsf{Z}\) 是分裂域，\(\mathfrak T\) 是 \(\mathfrak K_{\mathsf{Z}}\) 的各右理想的自同構域，即

\[
\mathfrak K_{\mathsf{Z}}=\sum\mathfrak T c_{ik}
\]

於是

\[
\mathfrak K_\Omega=\sum\mathfrak T_\Omega c_{ik}
\]

所以，\(\mathfrak K_\Omega\) 在 \(\Omega\) 上的秩是 \(t^2\) 與 \(\mathfrak T_\Omega\) 的秩之積；另一方面它等於 \(t^2\)。故 \(\mathfrak T_\Omega\) 的秩、因而也就是 \(\mathfrak T\) 的秩等於 1，於是 \(\mathfrak T=\mathsf{Z}\)，證畢。

\textbf{2.} 若 \(\mathsf{Z}\) 是 \(\mathfrak K_{\mathsf{Z}}\) 的各簡單理想的自同構域，則

\[
\mathfrak K_{\mathsf{Z}}=\sum\mathsf Zc_{ik}
\]

從而

\[
\mathfrak K_\Omega=\sum\Omega c_{ik}
\]

所以 \(c_{ik}\) 的數目必等於 \(t^2\)，證畢。

\textbf{定理 5.} \emph{設 \(\mathfrak Z\) 是 \(\mathfrak K\) 的一個分裂域 \(Z\) 在 \(\mathfrak K\) 中的不可約表示——按第 22 頁的定理 2，只有一個不可約表示類——則這個表示是矩陣環 \(\mathfrak K_r\) 的一個極大交換子域（\(r\) 是 \(\mathsf{Z}\) 的不可約表示的次數）。}

\emph{反過來，若 \(\mathfrak Z\) 是 \(\mathsf P\) 的一個交換擴域 \(Z\) 的不可約表示，並且是 \(\mathfrak K_r\) 的一個極大交換子域，則 \(Z\) 是分裂域。}

\emph{證明。} 1. 設分裂域 \(\mathsf{Z}\) 在 \(\mathsf P\) 上的次數為 \(n\)，則其不可約表示 \(\mathfrak Z\) 的次數 \(r\) 等於 \(n/t\)（第 22 頁）。設 \(\mathfrak Z\subseteqq\mathfrak Z^*\subseteqq\mathfrak K_r\)，其中 \(\mathfrak Z^*\) 是在 \(\mathsf P\) 上次數為 \(n^*\geqq n\) 的一個極大交換域。我們須證明 \(\mathfrak Z^*=\mathfrak Z\)；只須證明 \(n^*=n\) 即可。若 \(r^*\) 表示一個與 \(\mathfrak Z^*\) 同構的、\(\mathsf{Z}\) 的擴域 \(\mathsf{Z}^*\) 的不可約表示次數，則仍有 \(n^*=r^*t\)。另一方面，\(r\) 能被 \(r^*\) 整除，因為在 \(\mathfrak Z^*\) 中存在一個 \(\mathsf{Z}^*\) 的 \(r\) 次表示，故 \(r^*\leqq r\)。由

\[
n\leqq n^*=r^*t\leqq rt=n \qquad\text{可得}\qquad n^*=n.
\]

2. 設 \(\mathsf P\) 的交換擴域 \(\mathsf{Z}\) 在 \(\mathfrak K\) 中的不可約 \(r\) 次表示 \(\mathfrak Z\)，是 \(\mathfrak K_r\) 的一個極大交換子域。設 \(\mathfrak K_{\mathsf{Z}}\) 的各簡單理想的自同構域為 \(\mathfrak T\)，即

\[
\mathfrak K_{\mathsf{Z}}=\sum_{1}^{t'}\mathfrak T c_{ik}.
\]

\(\mathfrak T\) 包含 \(\mathsf{Z}\)。我們須證明 \(\mathfrak T=\mathsf{Z}\)。\(\mathfrak K_r\) 的一個簡單右理想 \(\mathfrak r\)，不僅是 \(\mathsf{Z}\) 在 \(\mathfrak K\) 中的表示模，也是 \(\mathfrak T\) 的表示模，因為

\[
\mathfrak r=\mathfrak T c_{i1}+\cdots+\mathfrak T c_{ik},
\]

所以 \(\mathfrak T\mathfrak r=\mathfrak r\)，即 \(\mathfrak r\) 是 \(\mathfrak T\)-左模，當然同時也是 \(\mathfrak K\)-右模。

因此，在 \(\mathfrak K_r\) 中存在一個與 \(\mathfrak T\) 同構並且包含 \(\mathfrak Z\) 的子域 \(\overline{\mathfrak T}\)。若 \(\mathfrak T\) 不同於 \(\mathsf{Z}\)，從而 \(\overline{\mathfrak T}\) 不同於 \(\mathfrak Z\)，那麼把一個不屬於 \(\mathfrak Z\) 的 \(\mathfrak T\)-元素 \(\alpha\) 添加到 \(\mathfrak Z\)，便會在 \(\mathfrak K_r\) 中得到 \(\mathfrak Z\) 的一個真交換擴域，這與 \(\mathfrak Z\) 的極大性矛盾，證畢。

由定理 5 立即得到：

\textbf{定理 6.} \emph{\(\mathfrak K\) 的各極大交換子域在 \(\mathsf P\) 上的次數為 \(t\)，並且都是分裂域\srcfn{1}{“每個分裂域都含有 \(\mathfrak K\) 的一個極大交換子域”這一猜想是不正確的；可能存在不具有這個性質的分裂域，與之相應的 \(r\) 必定大於 1。參見 R. Brauer 與 E. Noether，Über minimale Zerfällungskörper irreduzibler Darstellungen；Sitzungsberichte d. preuss. Akademie der Wisssch. 1927 XXXII，第 221 頁。}。}

現在建立一個關於 \(\mathfrak K_r\) 的子環同構的定理；對 \(\mathfrak K\) 的伽羅瓦群結構及相關事實的更精確認識，都以它為基礎。

\textbf{定理 7.} \emph{設 \(\mathfrak S_1\) 和 \(\mathfrak S_2\) 是 \(\mathfrak K\) 上的矩陣環 \(\mathfrak K_r\) 的雙邊單子環，都包含 \(\mathsf P\)，並且它們之間存在一個逐點固定 \(\mathsf P\) 的同構 \(\mathfrak S_1\cong\mathfrak S_2\)。那麼，這個同構是由一個正則的 \(\mathfrak K_r\)-元素 \(\tau\) 所給出的變換。}

\emph{若 \(\sigma_1\in\mathfrak S_1\) 與 \(\sigma_2\in\mathfrak S_2\) 相互對應，則}

\[
\tau^{-1}\sigma_1\tau=\sigma_2.
\]

\emph{證明。} 構造一個與 \(\mathfrak S_1\) 和 \(\mathfrak S_2\) 都互反同構的環 \(\Sigma\)，並把它的一個與 \(\mathsf P\) 同構的子域同一為 \(\mathsf P\)。這樣，\(\mathfrak S_1,\mathfrak S_2\) 是同一個 \(\Sigma\) 在 \(\mathfrak K\) 中次數相同為 \(r\) 的兩個互反表示。按第 22 頁的定理 2，\(\Sigma\) 在 \(\mathfrak K\) 中只有一個不可約互反表示類，因而每個給定次數也只有一個表示類（它是不可約類的倍數）；所以 \(\mathfrak S_1\) 與 \(\mathfrak S_2\) 必須屬於同一個類，從而通過一個變換彼此得到，證畢。

由定理 7 立即得到：

\textbf{定理 8.} \emph{逐點固定 \(\mathsf P\) 的 \(\mathfrak K\)-自同構群——也就是 \(\mathfrak K\) 關於其中心的伽羅瓦群 \(\mathfrak G\)——就是 \(\mathfrak K\) 的內自同構群。}

只需在定理 7 中置 \(r=1\) 和 \(\mathfrak S_1=\mathfrak S_2=\mathfrak K\)。

下面兩個定理表明，迄今建立的一般理論對於處理特殊問題是有用的。

\textbf{定理 9.} \emph{以實數域 \(\mathsf{R}\) 為中心的有限次非交換域，只有四元數域。}

\emph{證明。} 設 \(\mathfrak K\) 是以 \(\mathsf{R}\) 為中心的一個非交換域（在 \(\mathsf{R}\) 上有限秩）。按定理 6，\(\mathfrak K\) 在 \(\mathsf{R}\) 上的次數，是 \(\mathfrak K\) 的一個極大交換子域的次數 \(t\) 的平方；而這個子域只能是一個與複數域同構的域，所以 \(t=2\)，\(\mathfrak K\) 在 \(\mathsf{R}\) 上的次數為 4。設 \(\mathfrak Z_1\) 是 \(\mathfrak K\) 的一個極大交換子域，可以置 \(\mathfrak Z_1=\mathsf{R}(i)\)，其中 \(i^2=-1\)。於是 \(\mathfrak Z_2=\mathsf{R}(-i)\) 也是一個與 \(\mathfrak Z_1\) 同構的域；它作為集合與 \(\mathfrak Z_1\) 相同，而同構由 \(i\mapsto-i\) 給出。按定理 7，存在一個元素 \(j^*\in\mathfrak K\)，使

\[
j^{*-1}ij^*=-1.
\]

現在，\(\mathfrak Z_1(j^*)\) 在 \(\mathfrak Z_1\) 上的秩至少為 2，因為 \(j^*\) 不可能屬於 \(\mathfrak Z_1\)。所以 \(\mathfrak Z_1(j^*)\) 在 \(\mathsf{R}\) 上的秩為 4，從而等於 \(\mathfrak K\)，並可寫成

\[
\mathfrak K=\mathsf R+\mathsf Ri+\mathsf Rj^*+\mathsf Rij^*.
\]

由於

\[
\begin{aligned}
j^{*-2}1j^{*2}&=1\\
j^{*-2}ij^{*2}&=i\\
j^{*-2}jj^{*2}&=j^*\\
j^{*-2}ij^*j^{*2}&=ij^*
\end{aligned}
\]

\(j^{*2}\) 與 \(\mathfrak K\) 的所有元素可交換，所以屬於中心 \(\mathsf{R}\)；即 \(j^{*2}\) 是一個實數 \(r\)。\(r\) 不可能為正，否則多項式 \(x^2-r\) 在交換域 \(\mathsf{R}(j^*)\) 中會有四個根 \(+\sqrt r,-\sqrt r,+j^*\) 與 \(-j^*\)。所以 \(j^{*2}=-m^2,\ m\in \mathsf{R}\)。置 \(j=j^*m^{-1}\)，則 \(j^2=-1\)，而 \(1,i,j,k=ij\) 就是四元數單位，證畢。

這個證明顯然也適用於以任一實閉域代替 \(\mathsf{R}\)。

\textbf{定理 10（Wedderburn 定理）。} \emph{每個由有限多個元素組成的域都是交換的。}

\emph{證明。} 設 \(\mathfrak K\) 是一個只有有限多個元素的域，\(\mathsf P\) 是它的中心；設 \(t^2\) 是 \(\mathfrak K\) 在 \(\mathsf P\) 上的次數。\(\mathfrak K\) 的所有極大交換子域在 \(\mathsf P\) 上的次數都是 \(t\)，所以它們具有相同的元素數；按照伽羅瓦域理論中的一個已知定理（Moore 定理），它們彼此同構。又因為每個伽羅瓦域都是其素域的伽羅瓦擴張，可以選擇這些同構，使 \(\mathsf P\) 的每個元素都與自身對應。由定理 7，\(\mathfrak K\) 的所有極大交換子域都由某一個 \(\mathfrak Z\) 通過變換得到。另一方面，\(\mathfrak K\) 正是其所有極大交換子域的並集（只向 \(\mathsf P\) 添加\emph{一個}元素，就生成一個交換域）。現在只考察 \(\mathfrak K\) 的乘法群 \(\mathfrak K^*\)（即 \(\mathfrak K\) 去掉零），便得到：\(\mathfrak K^*\) 是與其子群 \(\mathfrak Z^*\) 共軛的各子群之並。但當 \(t>1\) 時，這會產生矛盾。事實上，若 \(\mathfrak K^*\) 有 \(M\) 個元素，\(\mathfrak Z^*\) 有 \(N\) 個元素，而 \(L\) 是 \(\mathfrak Z^*\) 在 \(\mathfrak K^*\) 中的指數，則 \(M=NL\)，不同共軛子群的數目至多為 \(L\)。即使這個數目恰為 \(L\)，若要使所有共軛子群合起來窮盡整個群，其中任意兩個也不應有共同元素；但它們全都含有單位元。

\subsection*{§ 21. 非交換域的伽羅瓦理論}\label{die-galoissche-theorie-der-nichtkommutativen-koerper}

以 \(\mathsf{G}\) 表示 \(\mathfrak K\) 關於其中心 \(\mathsf P\) 的伽羅瓦群。按定理 8，\(\mathsf{G}\) 與 \(\overline{\mathsf G}=\mathfrak K^*/\mathsf P^*\) 同構；我們把這個同構看作固定給定的。若 \(G\) 的一個子群 \(H\)，在同構 \(\mathsf{G}\cong\mathfrak K^*/\mathsf P^*\) 下所對應的 \(\mathfrak K^*\)-子群 \(\mathfrak H^*\)，添加\emph{零}以後成為一個域，則稱這個子群為\emph{閉子群}。

\textbf{定理 11. 伽羅瓦理論主定理。} \emph{\(G\) 的閉子群 \(H\) 與介於 \(\mathsf P\) 和 \(\mathfrak K\) 之間的域 \(\mathfrak S\)，可以一一對應；若 \(\mathsf{H}\) 與 \(\mathfrak S\) 相互對應，則 \(\mathsf{H}\) 是 \(\mathfrak S\) 的不變群，而 \(\mathfrak S\) 是 \(\mathsf{H}\) 的不變域。}

\emph{證明。} 把定理分成三個分命題：
\begin{enumerate}
\item 一個域 \(\mathfrak S\) 的不變群 \(H\) 是閉的。
\item 若 \(\mathfrak S\) 具有不變群 \(\mathsf{H}\)，則 \(\mathfrak S\) 是 \(\mathsf{H}\) 的不變域。
\item 若 \(\mathsf{H}\) 具有不變域 \(\mathfrak S\)，則 \(\mathsf{H}\) 是 \(\mathfrak S\) 的不變群。
\end{enumerate}

\begin{enumerate}
\item 在同構 \(G\cong \mathfrak K^*/\mathsf P^*\) 下，域 \(\mathfrak S\) 的不變群 \(\mathsf{H}\) 所對應的是全體與 \(\mathfrak S\) 中元素逐點可交換的非零 \(\mathfrak K\)-元素所成的集合 \(\mathfrak H^*\)。這些元素添加 0 後自然構成一個域，所以 \(H\) 是閉的。

\item 我們把第三個分命題歸結為第二個。域 \(\mathfrak S\) 對應於 \(G\) 的一個子群 \(S\)，即屬於 \(\mathfrak S^*/\mathsf P^*\) 的 \(\mathsf{G}\) 的子群。同樣，群 \(\mathsf{H}\) 對應於 \(\mathfrak K\) 的一個子域 \(\mathfrak H\)，並有 \(\mathsf{H}\cong \mathfrak H^*/\mathsf P^*\)。若 \(\mathfrak S\) 是 \(\mathsf{H}\) 的不變域，則 \(\mathfrak S\) 由所有與 \(\mathfrak H\) 逐點可交換的 \(\mathfrak K\)-元素組成；也可以說，\(\mathsf{S}\) 是 \(\mathfrak H\) 的不變群。相應地，若 \(\mathsf{H}\) 是 \(\mathfrak S\) 的不變群，則 \(\mathfrak H\) 是 \(\mathsf{S}\) 的不變域。因此第三個分命題還可表述為：

\item[\(3'\).] 若 \(\mathfrak H\) 具有不變群 \(\mathsf{S}\)，則 \(\mathfrak H\) 是 \(\mathsf{S}\) 的不變域；這正是現在還須證明的第二個分命題。

\item[3.] 除了若干小差別外，這個證明與交換情形相同：
\end{enumerate}

構造一個與 \(\mathfrak K\) 互反同構的域 \(\mathsf{K}\)，並把它的中心同一為 \(\mathsf P\)。我們可以不考察 \(\mathfrak K\) 的自同構，而考察 \(\mathfrak S\) 到 \(\mathsf{K}\) 的互反同構。

我們還要證明一個引理，它類似於交換情形中使用的延拓定理：

設 \(\mathfrak S\) 與 \(\mathfrak T\) 是介於 \(\mathsf P\) 與 \(\mathfrak K\) 之間的域，
\[
\mathsf P \subseteqq \mathfrak S \subseteqq \mathfrak T \subseteqq \mathfrak K.
\]
我們斷言，\(\mathfrak S\) 到 \(\mathsf{K}\) 的一個子域的每個互反同構（即 \(\mathfrak S\) 在 \(\mathsf{K}\) 中的一次互反表示），至少可以按照 \(\mathfrak T\) 在 \(\mathfrak S\) 上的次數 \(s\) 所給出的那麼多種方式，延拓為 \(\mathfrak T\) 的互反同構。

把域 \(\mathsf P\)、\(\mathfrak S\)、\(\mathfrak T\)、\(\mathfrak K\) 都以 \(\mathsf{K}\) 擴張：
\[
\mathsf{K}\subseteqq \mathfrak S_{\mathsf{K}}\subseteqq \mathfrak T_{\mathsf{K}}\subseteqq \mathfrak K_{\mathsf{K}}.
\]
按第 20 頁的定理 1，每個 \(\mathfrak S_{\mathsf{K}}\) 都是雙邊單而且完全可約的；其彼此在算子的意義下同構的各簡單左理想，是 \(\mathfrak S\) 在 \(\mathsf{K}\) 中唯一不可約表示類的表示模。這個表示類的次數為 1——按照 \(\mathsf{K}\) 的構造，\(\mathsf{K}\) 中存在與 \(\mathfrak S\) 互反同構的域——所以各簡單左理想的秩全都為 1；也就是說，它是數目等於它在 \(\mathsf P\) 上次數的左理想之和：
\[
\mathfrak S_{\mathsf{K}}=\mathsf K E_1+\cdots+\mathsf K E_p,\qquad \mathsf K E_\nu=\mathfrak S_{\mathsf{K}} E_\nu.
\]

於是
\[
\begin{aligned}
\mathfrak T_{\mathsf{K}}&=\mathfrak T\cdot \mathsf{K}=\mathfrak T\cdot \mathfrak S_{\mathsf{K}},\\
\mathfrak T_{\mathsf{K}}&=\mathfrak T\mathfrak S_{\mathsf{K}} E_1+\cdots+\mathfrak T\mathfrak S_{\mathsf{K}} E_p
              =\mathfrak T_{\mathsf{K}} E_1+\cdots+\mathfrak T_{\mathsf{K}} E_p.
\end{aligned}
\]
每個理想 \(\mathfrak T_{\mathsf{K}}\cdot E_\nu\) 分解為 \(s\) 個簡單理想，因為
\[
[\mathfrak T_{\mathsf{K}} E_\nu:\mathsf{K}]=[\mathfrak T_{\mathsf{K}} E_\nu:\mathsf K E_\nu]=[\mathfrak T_{\mathsf{K}} E_\nu:\mathfrak S_{\mathsf{K}} E_\nu]\leqq[\mathfrak T_{\mathsf{K}}:\mathfrak S_{\mathsf{K}}]=[\mathfrak T:\mathfrak S]=s
\]
且
\[
\sum [\mathfrak T_{\mathsf{K}} E_\nu:\mathsf{K}]=sp
\]
所以
\[
[\mathfrak T_{\mathsf{K}} E_\nu:\mathsf{K}]=s,
\qquad
\mathfrak T_{\mathsf{K}} E_\nu=\mathsf K e_1^{(\nu)}+\cdots+\mathsf K e_s^{(\nu)}.
\]

由 \(\mathsf K e_i^{(\nu)}\) 給出的 \(\mathfrak T\) 的表示，是由 \(\mathsf K E_\nu\) 給出的 \(\mathfrak S\) 的表示的延拓；因為若 \(\alpha\in \mathfrak S\)，則
\[
\alpha E_\nu=\sigma E_\nu,\qquad
\sum \alpha e_i^{(\nu)}=\sigma\sum e_i^{(\nu)},\qquad
\alpha e_i^{(\nu)}=\sigma e_i^{(\nu)}.
\]

現在只須證明由各 \(\mathsf K e_i^{(\nu)}\) 給出的表示全都不同。在交換情形中這是清楚的，因為它們屬於不同的表示類；在這裡相反，它們全都屬於同一個表示類，所以必須考察表示本身。為此，考察由各 \(\mathsf K e_i^{(\nu)}\) 給出的整個 \(\mathfrak T_k\) 的 \(\mathsf{K}\)-同態映射；一旦知道 \(\mathfrak T\) 的表示，這些映射就已經完全確定，因為構造整個映射只須知道基元素怎樣表示。因此，只須說明由各 \(\mathsf K e_i^{(\nu)}\) 給出的這些 \(\mathfrak T_{\mathsf{K}}\) 映射全都不同。若一次用表示模 \(\mathsf K e_i^{(\nu)}\) 表示 \(e_i^{(\nu)}\)，另一次用 \(\mathsf K e_j^{(\nu)}\) 表示它，則由
\[
e_i^{(\nu)2}=e_i^{(\nu)}
\]
第一次得到 1，而由
\[
e_i^{(\nu)}e_j^{(\nu)}=0
\]
第二次得到 0 作為 \(e_i^{(\nu)}\) 的表示元素。因此兩個表示確實不同。（完全相同的推論也可用於交換情形，只是在那裡是多餘的。）

從現在起，我們完全像第 10 頁和第 16 頁那樣結束證明。設 \(\mathsf{H}\) 是 \(\mathfrak S\) 的不變群，而 \(\mathfrak T\) 是 \(H\) 的不變域。\(H\) 的元素正是 \(\mathfrak S\) 的恆等自同構到整個 \(\mathfrak K\) 的所有延拓。由剛剛證明的引理，其中至少有 \(s\) 個在 \(\mathfrak T\) 上彼此不同，所以必有 \(s=1\)，\(\mathfrak T=\mathfrak S\)，證畢。

\section*{第五章：因子系}\label{kapitel-v.-faktorensysteme}

\subsection*{§ 22. 具有給定中心的域所成的群}\label{die-gruppe-der-koerper-mit-gegebenem-zentrum}

我們的研究以一個固定的交換域 \(\mathsf P\) 為基礎。全體矩陣環 \(\mathfrak K_r\)——其中自同構域 \(\mathfrak K\) 是以 \(\mathsf P\) 自身為中心、在 \(\mathsf P\) 上有限秩的域——關於“直接乘積構造”組成一個乘法封閉系統。兩個這樣的環 \(\mathfrak K_r\) 和 \(\mathfrak L_s\) 的直接積，就是由把 \(\mathfrak K_r\) 的基域擴張到 \(\mathfrak L_s\) 所得到的環 \(\mathfrak K_r\cdot\mathfrak L_s\)：
\[
\mathfrak K_r\times\mathfrak L_s=\mathfrak K_r\cdot\mathfrak L_s.
\]
按第 22 頁的注，這種直接乘積構造是交換的：
\[
\mathfrak K_r\times\mathfrak L_s=\mathfrak L_s\times\mathfrak K_r.
\]
斷言全體 \(\mathfrak K_r\) 所成的系統關於乘法封閉，無非就是說 \(\mathfrak K_r\times\mathfrak L_s=\mathfrak K_r\mathfrak L_s\) 總是又一個矩陣環，其自同構域 \(\mathfrak M\) 以 \(\mathsf P\) 為中心——相對於 \(\mathsf P\) 的有限秩當然成立。按第 20 頁的定理 1，\(\mathfrak K_r\times\mathfrak L_s\) 雙邊單而且完全可約，因而是 \(\mathsf P\) 的某個擴域 \(\mathfrak M\) 上的矩陣環。\(\mathfrak M\) 的中心等於 \(\mathsf P\)，直接來自如下兩點：第一，\(\mathsf P\) 確實包含在中心中；第二，這個中心必須包含在 \(\mathfrak K\)（或 \(\mathfrak L\)）的中心 \(\mathsf P\) 中。乘法 \(\mathfrak K_r\times\mathfrak L_s\) 顯然是結合的。各 \(\mathfrak K_r\) 組成一個乘法封閉系統，但不構成群；因為若 \(\mathfrak K_r\times\mathfrak L_s=\mathfrak M_m\)，則必有 \(m\geqq rs\)，所以當 \(r>m\) 時，不可能找到滿足等式 \(\mathfrak K_r\times\mathfrak L_s=\mathfrak M_m\) 的 \(\mathfrak L_s\)。

現在，把自同構域相同為 \(\mathfrak K\) 的所有 \(\mathfrak K_r\) 合併成一個類 \(\{\mathfrak K\}\)。於是有：若 \(\mathfrak K_r\) 與 \(\mathfrak K_{r'}\) 屬於同一個類 \(\{\mathfrak K\}\)，而 \(\mathfrak L_s\) 與 \(\mathfrak L_{s'}\) 屬於同一個類 \(\{\mathfrak L\}\)，則 \(\mathfrak K_r\times\mathfrak L_s\) 與 \(\mathfrak K_{r'}\times\mathfrak L_{s'}\) 也屬於同一個類。只須證明 \(\mathfrak K_r\times\mathfrak L_s\) 與 \(\mathfrak K\times\mathfrak L\) 屬於同一類；也就是說，若 \(\mathfrak K\times\mathfrak L=\mathfrak M_m\)，則 \(\mathfrak K_r\times\mathfrak L_s\) 也是 \(\mathfrak M\) 上的矩陣環。這很容易看出：
\[
\mathfrak K_r\times\mathfrak L
=\left(\sum_1^r\mathfrak K c_{ik}\right)\mathfrak L
=\sum_1^r\mathfrak K_{\mathfrak L}c_{ik}
=\sum_1^r\mathfrak M_m c_{ik}.
\]
這是 \(\mathfrak M_m\) 上的 \(r\) 次矩陣環，當然也就是 \(\mathfrak M\) 上的 \(mr\) 次矩陣環。隨後以完全相同的方式繼續：
\[
\mathfrak K_r\times\mathfrak L_s
=\left(\sum_1^s\mathfrak L d_{ik}\right)_{\mathfrak K_r}
=\sum_1^s\mathfrak L_{\mathfrak K_r}d_{ik}
=\sum_1^s\mathfrak M_{mr}d_{ik}
=\mathfrak M_{mrs}
\]
從而斷言已證。我們記下：由 \(\mathfrak L\times\mathfrak K=\mathfrak M_m\)，可推出 \(\mathfrak K_r\times\mathfrak L_s=\mathfrak M_{mrs}\)。

由已證內容可知，若把兩個類 \(\{\mathfrak K\}\) 與 \(\{\mathfrak L\}\) 的乘積定義為：從 \(\{\mathfrak K\}\) 中取一個代表、從 \(\{\mathfrak L\}\) 中取一個代表，再取這兩個代表的乘積所屬的類，那麼這個定義與代表的選擇無關。

所以映射 \(\mathfrak K_r\mapsto\{\mathfrak K_r\}\) 是從全體 \(\mathfrak K_r\) 所成系統到全體類 \(\{\mathfrak K\}\) 所成集合的一個同態；因此，全體類 \(\{\mathfrak K\}\) 所成的集合 \(\mathscr K\) 是一個乘法封閉系統。我們證明 \(\mathscr K\) 是一個群。單位元顯然是類 \(\{\mathsf P\}\)，因為 \(\mathfrak K_r\times\mathsf P=\mathfrak K_r\)。

所以，只須為每個類 \(\{\mathfrak K\}\) 證明逆類 \(\{\mathfrak K\}^{-1}\) 的存在。可知必須置 \(\{\mathfrak K\}^{-1}=\{\overline{\mathfrak K}\}\)，其中 \(\{\overline{\mathfrak K}\}\) 表示與 \(\mathfrak K\) 互反同構的域。換言之，必須證明 \(\mathfrak K\times\overline{\mathfrak K}\) 是 \(\mathsf P\) 上的一個矩陣環；或者，若 \(\mathfrak K\) 的秩為 \(t^2\)，因而 \(\mathfrak K\times\overline{\mathfrak K}\) 在 \(\mathsf P\) 上的秩為 \(t^4\)，則必須證明 \(\mathfrak K\times\overline{\mathfrak K}\) 分解為 \(t^2\) 個簡單左理想。後一點來自如下事實：\(\mathfrak K\times\overline{\mathfrak K}\) 的一個簡單左理想，是 \(\overline{\mathfrak K}\) 在 \(\mathfrak K\) 中一個不可約互反表示的表示模；這個表示的次數是 1，故該左理想在 \(\mathfrak K\) 上的秩為 1，也就是在 \(\mathsf P\) 上的秩為 \(t^2\)。

若在 \(\mathsf P\) 與 \(\mathfrak K\) 之間不存在另一個中心同樣為 \(\mathsf P\) 的域，我們就稱 \(\mathfrak K\) 為一個\emph{簡單域}。

\textbf{定理。} \emph{每個域 \(\mathfrak K\) 都是一些簡單子域 \(\mathfrak K^{(i)}\) 的直接積：}
\[
\mathfrak K=\mathfrak K^{(1)}\times\cdots\times\mathfrak K^{(p)}.
\]

\emph{證明。} 只須證明：若 \(\mathfrak K^{(1)}\) 是 \(\mathfrak K\) 的一個以 \(\mathsf P\) 為中心的子域，則 \(\mathfrak K\) 有一個以 \(\mathsf P\) 為中心的子域 \(\mathfrak S\)，滿足
\[
\mathfrak K^{(1)}\times\mathfrak S=\mathfrak K
\]
無論如何，存在一個類 \(\{\mathfrak S\}\) 具有性質
\[
\{\mathfrak K^{(1)}\}\times\{\mathfrak S\}=\{\mathfrak K\}
\]
即
\[
\{\mathfrak S\}=\{\overline{\mathfrak K^{(1)}}\}\times\{\mathfrak K\}.
\qquad
\text{設}\qquad
\overline{\mathfrak K^{(1)}}\times\mathfrak K=\mathfrak S_\lambda.
\]
於是
\[
\mathfrak K^{(1)}\times\mathfrak S_\lambda
=\mathfrak K^{(1)}\times\overline{\mathfrak K^{(1)}}\times\mathfrak K
=\mathsf P_{t^2}\times\mathfrak K
=\mathfrak K_{t^2},
\qquad
t^2=\text{秩 }\mathfrak K^{(1)}.
\]
但 \(\mathfrak K^{(1)}\times\mathfrak S_\lambda\) 是一個次數至少等於 \(\lambda\) 的矩陣環，所以 \(\lambda\leqq t^2\)。\(\mathfrak S_\lambda=\overline{\mathfrak K^{(1)}}\times\mathfrak K\) 必須分解為在 \(\mathfrak K\) 上秩為 1 的簡單理想，因為 \(\overline{\mathfrak K^{(1)}}\) 容許在 \(\mathfrak K\) 中的一次互反表示；故 \(\lambda\geqq t^2\)，從而 \(\lambda=t^2\)，
\[
\mathfrak K_{t^2}=\mathfrak K^{(1)}\times\mathfrak S_{t^2}.
\]
因此，\(\mathfrak K^{(1)}\times\mathfrak S\) 本身必須是域，而且
\[
\mathfrak K^{(1)}\times\mathfrak S_{t^2}
=\bigl(\mathfrak K^{(1)}\times\mathfrak S\bigr)_{t^2}
=\mathfrak K_{t^2}.
\]
\(\mathfrak K^{(1)}\times\mathfrak S\) 是 \(\mathfrak K_{t^2}\) 的一種分解中各簡單左理想的自同構域，而 \(\mathfrak K\) 是 \(\mathfrak K_{t^2}\) 的另一種分解中各簡單左理想的自同構域。因為 \(\mathfrak K_{t^2}\) 的任意兩個不同左理想都能納入某個分解，所以 \(\mathfrak K\) 與 \(\mathfrak K^{(1)}\times\mathfrak S\) 同構，並且可選擇這個同構使 \(\mathsf P\) 逐點不變。因此，按第 25 頁的定理 7，這個同構由 \(\mathfrak K_{t^2}\) 的一個內自同構生成：
\[
\mathfrak K
=\lambda^{-1}\bigl(\mathfrak K^{(1)}\times\mathfrak S\bigr)\lambda
=\lambda^{-1}\mathfrak K^{(1)}\lambda\times\lambda^{-1}\mathfrak S\lambda,
\qquad
\lambda\in\mathfrak K_{t^2}.
\]
\(\lambda^{-1}\mathfrak K^{(1)}\lambda\) 是 \(\mathfrak K\) 的一個與 \(\mathfrak K^{(1)}\) 同構的子域；所以，按第 25 頁的定理 7，存在來自 \(\mathfrak K\) 的一個合適元素 \(\mu\)，使
\[
\mu^{-1}\lambda^{-1}\mathfrak K^{(1)}\lambda\mu=\mathfrak K^{(1)}
\]
從而
\[
\mu^{-1}\mathfrak K\mu
=\mathfrak K
=\mathfrak K^{(1)}\times\mu^{-1}\lambda^{-1}\mathfrak S\lambda\mu,
\]
斷言得證。

\subsection*{§ 23. 因子系}\label{faktorensysteme}

在下面的研究中，假設基域 \(\mathsf P\) 是完美的；也可以改用更一般的假設：只考察那些以 \(\mathsf P\) 為中心、並且每個 \(\mathfrak K_r\) 都含有一個在 \(\mathsf P\) 上第一類的極大交換子域的非交換域 \(\mathfrak K\)。（正如 Köthe 後來已經證明的，這個性質對每個 \(\mathfrak K\) 都成立。）

以 \(Z\) 表示中心為 \(\mathsf P\) 的非交換域 \(\mathfrak K\) 的一個分裂域；設 \(Z\) 在 \(\mathfrak K\) 中的不可約表示 \(\mathfrak Z\) 的次數為 \(r\)，於是它是 \(\mathfrak K_r\) 的一個極大交換子域。設 \(\Gamma\) 是 \(\mathsf{Z}\) 的伽羅瓦域，即
\[
\mathfrak Z_\Gamma=e_1\Gamma+\cdots+e_n\Gamma
\qquad\text{（第 9 頁）}
\]
其中 \(n\) 是 \(\mathsf{Z}\) 在 \(\mathsf P\) 上的次數。對於以 \(\Gamma\) 擴張的環 \(\mathfrak K_r\)，也得到分解
\[
\mathfrak K_{r\Gamma}=\mathfrak K_{r\Gamma}e_1+\cdots+\mathfrak K_{r\Gamma}e_n
\]
\[
\mathfrak K_{r\Gamma}=\mathfrak l_1+\cdots+\mathfrak l_n,\qquad
\mathfrak l_i=\mathfrak K_{r\Gamma}e_i.
\]
各 \(\mathfrak l_i\) 顯然是左理想。它們甚至是簡單左理想；因為 \(\mathsf{Z}\) 是分裂域，\(\mathfrak K_{r\Gamma}\) 是 \(\Gamma\) 上次數為 \(r\cdot t=n\) 的矩陣環（定理 4，第 23 頁；第 22 頁有 \(n=rt\)）。這樣得到的 \(\mathfrak K_{r\Gamma}\) 分解成簡單左理想的分解，特別適合於 \(\mathfrak K\) 和 \(\mathsf{Z}\) 的代數性質。事實上，若像第 18 頁那樣定義 \(\Gamma/\mathsf P\) 的伽羅瓦群 \(\mathfrak G\) 的一個置換 \(S\) 在全體 \(\mathfrak K_{r\Gamma}\)-元素上的作用，即對每個 \(\mathfrak K_r\)-元素 \(z\) 置
\[
S(z)=z
\]
則各理想 \(\mathfrak l_i\) 彼此共軛；也就是說，把一個 \(S\) 作用於某個 \(\mathfrak l_j\)，會得到另一個理想：
\[
S(\mathfrak l_j)=\mathfrak l_i,
\]
因為單位元的每個不可分解分量 \(e_j\) 必須變成另一個 \(e_i\)。（另見第 11 頁。）

通過表示模 \(e_i\Gamma\)，\(\mathfrak Z\) 被映到 \(\Gamma\) 的一個子域 \(\mathsf{Z}_i\)。\(e_i\) 已經位於 \(\mathfrak Z_{\mathsf{Z}_i}\) 中，因為 \(\mathfrak Z\) 的表示 \(\mathsf{Z}_i\) 已經包含在 \(\mathsf{Z}_i\) 中；所以 \(\mathfrak Z_{\mathsf{Z}_i}\) 必須分離出表示模 \(e_i\mathsf{Z}_i\)。設 \(\{\mathsf{Z}_i,\mathsf{Z}_k\}\) 是 \(\mathsf{Z}_i\) 與 \(\mathsf{Z}_k\) 的複合域。\(\mathfrak K_{r\{\mathsf{Z}_i,\mathsf{Z}_k\}}\) 分離出兩個左理想
\[
\bar{\mathfrak l}_i=\mathfrak K_{i\{\mathsf{Z}_i,\mathsf{Z}_k\}}e_i,\qquad
\bar{\mathfrak l}_k=\mathfrak K_{r\{\mathsf{Z}_i,\mathsf{Z}_k\}}e_k
\]
它們是簡單的，因為 \(\bar{\mathfrak l}_i=\mathfrak K_{r\Gamma}\bar{\mathfrak l}_i\)。因此
\[
\bar{\mathfrak l}_i=\sum_{\lambda=1}^m\{\mathsf{Z}_i,\mathsf{Z}_k\}c_{\lambda i}^{(ik)},\qquad
\bar{\mathfrak l}_k=\sum_{\lambda}\{\mathsf{Z}_i,\mathsf{Z}_k\}c_{\lambda k}^{(ik)}.
\]
\(c_{\lambda\varrho}^{(ik)}\) 是 \(\mathfrak K_{r\{\mathsf{Z}_i,\mathsf{Z}_k\}}\) 的一個矩陣單位系。

每個以 \(\mathfrak K_{r\{\mathsf{Z}_i,\mathsf{Z}_k\}}\) 為算子域的、從 \(\bar{\mathfrak l}_i\) 到 \(\bar{\mathfrak l}_k\) 的同構 \(a_i\mapsto a_k\)，都由元素 \(p_{ik}\) 來標識，其中 \(e_i\mapsto p_{ik}\)；因為這時 \(a_i=a_ie_i\mapsto a_ip_{ik}\)，所以 \(\mathfrak l_i\cdot p_{ik}=\mathfrak l_k\)。所有可能的這種 \(p_{ik}\) 顯然為
\[
p_{ik}=c_{ik}^{(ik)}\gamma_{ik},\qquad \gamma_{ik}\ne0\quad\text{來自}\quad \{\mathsf{Z}_i,\mathsf{Z}_k\}.
\]
映射 \(a_i\mapsto a_ip_{ik}\) 同時也給出從 \(\mathfrak l_i\) 到 \(\mathfrak l_k\) 的一個算子同構。現在選取某些這樣的 \(p_{ik}\) 系，方法是給它們施加一個共軛性條件。

定義 \(S\) 在指標 \(i\) 上的作用為
\[
S(i)=k,\qquad\text{若}\qquad S(\mathsf{Z}_i)=\mathsf{Z}_k.
\]
對於每個指標對 \((\nu,\mu)\)，至少存在一個與之共軛的對 \(S(\nu,\mu)\)，其形式為 \(S(\nu,\mu)=(1,k)\)。從每一類共軛指標對中選取一個固定的對 \((1,k)\)，再對這些 \((1,k)\) 置
\[
p_{1k}=c_{1k}^{(1k)}\gamma_{1k},\qquad \gamma_{1k}\ne0\quad\text{來自}\quad \{\mathsf{Z}_1,\mathsf{Z}_k\},
\]
其餘任意；然後，若 \(S(1,k)=(\nu,\mu)\)，就置
\[
p_{\nu\mu}=S(p_{1k}).
\]
於是 \(p_{\nu\mu}\) 確實具有形式
\[
p_{\nu\mu}=c_{\nu\mu}^{(\nu\mu)}\gamma_{\nu\mu},\qquad
\gamma_{\nu\mu}\ne0\quad\text{來自}\quad \{\mathsf{Z}_\nu,\mathsf{Z}_\mu\}.
\]
因為，把 \(S\) 作用於同構
\[
a_1\mapsto a_1p_{1k},\qquad\text{特別地}\qquad e_1\mapsto p_{1k}
\quad\text{從 }\bar{\mathfrak l}_1\text{ 到 }\bar{\mathfrak l}_k
\]
會給出同構
\[
a_\nu\mapsto a_\nu p_{\nu\mu},\qquad\text{特別地}\qquad e_\nu\mapsto p_{\nu\mu}
\quad\text{從 }\bar{\mathfrak l}_\nu\text{ 到 }\bar{\mathfrak l}_\mu.
\]

若以 \(c_{\lambda\varrho}\) 表示 \(\mathfrak K_{r\Gamma}\) 的一個矩陣單位系，則顯然還可把 \(p_{\nu\mu}\) 寫成
\[
p_{\nu\mu}=c_{\nu\mu}\bar\gamma_{\nu\mu}
\]
關係
\[
c_{ij}c_{1k}=0,\quad\text{若}\quad j\ne1,
\qquad
c_{ij}c_{jk}=c_{ik}
\]
因而導致關係
\[
\begin{aligned}
p_{ij}p_{1k}&=0,\quad &&\text{若 }j\ne1,\\
p_{ij}p_{jk}&=\alpha_{ik}^{(j)}p_{ik},\quad
&&\alpha_{ik}^{(j)}\text{ 來自 }\Gamma\text{（更準確地說來自 }\{\mathsf{Z}_i,\mathsf{Z}_j,\mathsf{Z}_k\}\text{）}
\end{aligned}
\]
這 \(n^3\) 個量 \(\alpha\)，組成以 \(\mathsf{Z}\) 為基礎時 \(\mathfrak K\) 的\emph{因子系}，它屬於“\emph{僞矩陣單位}”系 \(p_{ik}\)。

各 \(\alpha\) 彼此共軛：由 \(S(i,k,j)=(\nu,\mu,\tau)\) 可得
\[
S\bigl(\alpha_{ik}^{(j)}\bigr)=\alpha_{\nu\mu}^{\tau}.
\]

每個僞矩陣單位系 \(p_{ik}^*\)，都是從一個 \(p_{ik}\) 系得到的：把初始的各 \(p_{ik}\) 乘以任意來自 \(\{\mathfrak l_1,\mathfrak l_k\}\) 的 \(\delta_{ik}\ne0\)：
\[
p_{1k}^*=p_{1k}\delta_{1k}
\]
然後仍在 \(S(1,k)=(\nu,\mu)\) 時置
\[
p_{\nu\mu}^*=S(p_{1k}^*)
\]
因此
\[
p_{\nu\mu}^*=p_{\nu\mu}\delta_{\nu\mu}
\]
其中 \(\delta_{\nu\mu}\) 來自 \(\{\mathsf{Z}_\nu,\mathsf{Z}_\mu\}\)；各 \(\delta_{\nu\mu}\) 彼此共軛，並且由 \(S(i,k)=(\nu,\mu)\) 可得 \(S(\delta_{ik})=\delta_{\nu\mu}\)。若 \(\alpha^*\) 是屬於 \(p_{ik}^*\) 的因子系，則有
\[
\alpha_{ik}^{(i)*}=\frac{\delta_{ij}\delta_{jk}}{\delta_{ik}}\alpha_{ik}^{(j)}.
\]

因子系 \(\alpha^*\) 與 \(\alpha\) 稱為\emph{相伴的}（\(p_{ik}^*\) 與 \(p_{ik}\) 也同樣稱為相伴的）；以 \(\mathsf{Z}\) 為基礎時 \(\mathfrak K\) 的所有因子系組成一個相伴因子系類 \(\{\alpha\}\)。

類 \(\{\alpha\}\) 不依賴於定義中所採用的 \(\mathsf{Z}\) 在 \(\mathfrak K_r\) 中的表示 \(\mathfrak Z\)。事實上，從一個 \(\mathfrak Z\) 過渡到另一個表示，可以通過用 \(\mathfrak K_r\) 中的一個元素 \(\lambda\) 作變換來實現；這個元素在 \(\mathfrak G\) 下不變。因此，由屬於 \(\mathfrak Z\) 的一個 \(p_{ik}\) 系得到一個系
\[
p'_{ik}=\lambda^{-1}p_{ik}\lambda,
\]
它屬於 \(\mathfrak Z'\)；但 \(p_{ik}\) 與 \(p'_{ik}\) 的因子系完全相同。

因子系無非是 \(\mathfrak K_{n\Gamma}\) 的一張乘法表，只不過它特別適應於 \(\mathfrak K\) 的代數性質。
\subsection*{§ 24. 因子系的乘法}\label{multiplikation-von-faktorensystemen}

除了說非交換域 \(\mathfrak K\) 的分裂域以外，我們也說 \(\{\mathfrak K\}\) 的分裂域；這種說法有意義，因為若 \(\mathfrak K_{\mathsf{Z}}\) 分解成絕對不可約分量，則每個環 \(\mathfrak K_{r\mathsf{Z}}\) 也如此，反之亦然。

任意兩個域 \(\mathfrak K\)、\(\mathfrak S\)，若中心為 \(\mathsf P\)，都有公共分裂域，例如 \(\mathfrak K\) 的一個分裂域與 \(\mathfrak S\) 的一個分裂域的複合域。

\textbf{定理 1.} \emph{那些 \(\{\mathfrak K\}\) 若以域 \(\mathsf{Z}\) 為分裂域，就組成一個子群 \(\mathscr K_{\mathsf{Z}}\)；它位於群 \(\mathscr K\) 中，而後者由所有 \(\{\mathfrak K\}\) 組成。}

\emph{證明.} 若 \(\mathfrak K\) 與 \(\mathfrak S\) 都以 \(\mathsf{Z}\) 為分裂域，則 \(\mathfrak K_{\mathsf{Z}}\) 與 \(\mathfrak S_{\mathsf{Z}}\) 都分解，因而 \((\mathfrak K\times\mathfrak S)_{\mathsf{Z}}\) 也分解成絕對不可約分量；亦即，\(\mathsf{Z}\) 是 \(\{\mathfrak K\times\mathfrak S\}\) 的分裂域（定理 4，第 23 頁）。

互反同構的域具有完全相同的分裂域。

一個輔助考察。
\[
\mathfrak K_r=\mathfrak L_1+\cdots+\mathfrak L_r,\qquad
\mathsf P_m=\mathfrak z_1+\cdots+\mathfrak z_m
\]
是把 \(\mathfrak K_r\)、\(\mathsf P_m\) 分解成簡單左理想的分解，則
\[
\mathfrak K_{rm}=\mathfrak K_r\times\mathsf P_m
=\mathfrak L_1\mathfrak z_1+\cdots+\mathfrak L_i\mathfrak z_j+\cdots+\mathfrak L_r\mathfrak z_m
\]
是 \(\mathfrak K_{rm}\) 的一個分解，包含 \(rm\) 個非零、因而簡單的左理想。因此，每個有 \(r\) 項的左理想 \(\mathsf{A}\)，若屬於 \(\mathfrak K_{rm}\)，都同有 \(r\) 項的左理想
\[
\mathfrak L_1\mathfrak z_1+\cdots+\mathfrak L_r\mathfrak z_1=\mathfrak K_r\mathfrak z_1
\]
通過一個 \(\mathfrak K_r\) 的內自同構相聯繫：
\[
\mathsf{A}=\tau^{-1}\mathfrak K_r\tau\cdot\tau^{-1}\mathfrak z_1\tau
=\mathfrak K'_r\cdot\mathfrak z'_1.
\]

若 \(\mathfrak K_{r\Gamma}=\mathfrak l_1+\cdots+\mathfrak l_n\) 是把 \(\mathfrak K_{r\Gamma}\) 分解成共軛左理想的一個分解，它屬於某個 \(\mathfrak Z\)，而這個域位於 \(\mathfrak K_r\) 中，且 \(\mathfrak l_i=\mathfrak K_re_i\)，則由於
\[
S(\tau^{-1}\mathfrak l_i\tau\cdot\mathfrak z'_1)
=S(\tau^{-1}\mathfrak l_i\tau)\,\mathfrak z'_1
=\tau^{-1}(S\mathfrak l_i)\tau\cdot\mathfrak z'_1
=\tau^{-1}\mathfrak l_k\tau\cdot\mathfrak z'_1
\]
在
\[
\tag{1}
\mathsf{A}_\Gamma=\mathfrak l'_1\mathfrak z'_1+\cdots+\mathfrak l'_n\mathfrak z'_1
\]
中得到把 \(\mathsf{A}_\Gamma\) 分解成共軛左理想的一個分解。這裡
\[
\mathfrak l'_i\mathfrak z'_1=\mathfrak K_{rm\Gamma}\cdot e'_iE',
\]
其中置 \(\tau^{-1}e_i\tau=e'_i\)，而 \(E'\) 表示 \(\mathfrak z_1\) 的單位元。

任何另一個分解
\[
\tag{2}
\mathsf{A}_\Gamma=\mathfrak q_1+\cdots+\mathfrak q_n,\qquad
S\mathfrak q_i=\mathfrak q_k\quad\text{若}\quad
S\mathfrak l'_i\mathfrak z'_1=\mathfrak l'_k\mathfrak z'_1
\]
為共軛左理想的分解，都由 (1) 經一個 \(\mathfrak K_{rm}\)-元素 \(\lambda\) 的變換得到：
\[
\mathfrak q_i=\lambda^{-1}\mathfrak l'_i\mathfrak z'_1\lambda,
\qquad
\text{亦即，若 \(e_i^*\) 是 \(\mathfrak q_i\) 的單位元，則}\quad
e'_i=\lambda^{-1}e'_iE'\lambda.
\]

這個證明是下列引理的直接推論：

\textbf{引理.} \emph{設 \(\mathsf{A}\) 是 \(\mathfrak K_n\) 的一個左理想，並設}
\[
\mathsf{A}=\mathfrak l_1+\cdots+\mathfrak l_n
=\mathfrak K_{n\Gamma}e_1+\cdots+\mathfrak K_{n\Gamma}e_n
\]
\emph{是一個按假設存在的共軛左理想分解；這裡不失一般性，可以假定這些 \(l_i\) 彼此共軛。}

\emph{設 \(\mathfrak z\) 與 \(\mathsf{A}\) 算子同構，並滿足同樣的假設：}
\[
\mathfrak z=\mathfrak q_1+\cdots+\mathfrak q_n
=\mathfrak K_{n\Gamma}e_1^*+\cdots+\mathfrak K_{n\Gamma}e_n^*
\]
\emph{（亦即 \(\mathfrak q_i\) 與 \(e_i^*\) 各自共軛）。則存在一個正則元素 \(\lambda\)，它來自 \(\mathfrak K_n\)，使 \(\mathfrak q_i=\lambda^{-1}\mathfrak l_i\lambda\)。}

\emph{證明. 1. 可以假定這些 \(e_i\) 共軛，理由如下：}若 \(e_1\) 是 \(\mathfrak l_1\) 的任意一個生成冪等元，即 \(\mathfrak l_1=\mathfrak K_{n\Gamma}e_1\)，則對 \(\Gamma\) 施用一個自同構便得 \(\mathfrak l_i=\mathfrak K_{n\Gamma}e_i\)，其中 \(e_i\) 與 \(e_1\) 共軛，並且仍是生成冪等元。

2. \emph{每個 \(\mathfrak l\) 都與每個 \(\mathfrak q\) 同構；}因為這些 \(\mathfrak l\) 彼此共軛，故具有相同的（絕對）長度；又因為 \(\mathfrak l\) 與 \(\mathfrak q\) 的數目相同，而 \(\mathfrak z\) 與 \(\mathsf{A}\) 按假設同構。

3. 用對應
\[
\tag{3}
e_1\mapsto s_1=e_1s_1;\qquad
e_1^*\mapsto t_1=e_1^*t_1;\qquad
s_1t_1=e_1;\qquad
t_1s_1=e_1^*
\]
確定 \(\mathfrak l_1\) 與 \(\mathfrak q_1\) 之間的一個同構及其逆同構。施用所有自同構，(3) 變為
\[
e_i\mapsto s_i=e_is_i;\qquad
e_i^*\mapsto t_i=e_i^*t_i;\qquad
s_it_i=e_i;\qquad
t_is_i=e_i^*,
\]
由此確定同構 \(\mathfrak q_i=\mathfrak l_is_i\) 與 \(\mathfrak l_i=\mathfrak q_it_i\)。

4. 現在置
\[
s=s_1+\cdots+s_n;\qquad t=t_1+\cdots+t_n;
\]
則有
\[
\mathfrak q_i=\mathfrak l_is;\qquad
\mathfrak l_i=\mathfrak q_it;\qquad
st=e_1+\cdots+e_n=\xi;\qquad
ts=e_1^*+\cdots+e_u^*=\xi^*,
\]
其中 \(\xi\) 與 \(\xi^*\) 屬於 \(\mathfrak K_n\)；並且
\[
\mathsf{A}\xi=\mathsf{A},\qquad \mathfrak z\xi^*=\mathfrak z\quad
(\text{因為 }\mathfrak l_i\xi=\mathfrak l_il_i=\mathfrak l_i),
\]
\[
\xi^2=\xi,\qquad \xi^{*2}=\xi^*.
\]
現在置 \(e=\overline E+E=\overline E^*+E^*\)，並因而有
\[
\mathfrak K_n=\overline{\mathsf A}+\mathsf{A}
=\mathfrak K_n\overline\xi+\mathfrak K_n\xi
=\overline{\mathfrak z}+\mathfrak z
=\mathfrak K_n\overline\xi^*+\mathfrak K_n\xi^*,
\]
則 \(\overline{\mathsf A}\) 與 \(\overline{\mathfrak z}\) 也同構；若
\[
\overline E\mapsto\overline s=\overline E\,\overline s
\qquad\text{且}\qquad
\overline E^*\mapsto\overline t=\overline E^*\overline t
\]
是互逆同構，並置
\[
\lambda=s+\overline s,\qquad \mu=t+\overline t,
\]
則
\[
\lambda\mu=\mu\lambda=e,\qquad
\mathfrak l_i\lambda=\mathfrak q_i,\qquad
\lambda^{-1}\mathfrak q_i=\mathfrak q_i
\quad
(\lambda^{-1}\mathfrak q_i\subseteqq\mathfrak q_i
=\lambda\cdot\lambda^{-1}\mathfrak q_i\leqq\lambda^{-1}\mathfrak q_i)
\]
從而
\[
\lambda^{-1}\mathfrak l_i\lambda=\mathfrak q_i,\qquad \text{證畢。}
\]

對於分解 (1)，還可以確定元素
\[
r_{ik}=c_{ik}^{\prime(ik)}\gamma_{ik},\qquad \gamma_{ik}c\{\mathsf{Z}_i\mathsf{Z}_k\},
\]
使得 \(e'_iE'\mapsto r_{ik}\) 給出理想 \(\mathfrak l'_i\mathfrak z'_1\) 與 \(\mathfrak l'_k\mathfrak z'_1\) 之間的算子同構，並且滿足共軛條件 \(S(r_{ik})=r_{\nu\mu}\)，其中 \(S(i,k)=(\nu,\mu)\)。它們由
\[
r_{ij}r_{jk}=\alpha_{ik}^{(j)}r_{ik}
\]
確定一個因子系 \(\alpha\)，它屬於 \(\mathsf{A}_\Gamma\)。現在重要的是，在確定 \(\mathsf{A}_\Gamma\) 的因子系時，也可以使用任意的分解 (2)。事實上，若把屬於 (1) 的一個系 \(r_{ik}\) 替換為 \(r'_{ik}=\lambda^{-1}r_{ik}\lambda\)，則 \(e_i^*\mapsto r'_{ik}\) 給出從 \(\mathfrak q_i\) 到 \(\mathfrak q_k\) 的算子同構；這些 \(r'_{ik}\) 滿足共軛條件，因為 \(\lambda\) 是 \(\mathfrak q\)-不變的；而且 \(r'_{ik}\) 與 \(r_{ik}\) 具有相同的因子系。

取一組僞矩陣單位 \(p_{ik}\)，它們屬於 \(\mathfrak K\)，則
\[
r_{ik}=\tau^{-1}p_{ik}\tau\cdot E'
\]
是一個 \(r_{ik}\)-系，屬於 \(\mathsf{A}_\Gamma\)，且有相同的因子系。於是得到如下引理：

\textbf{定理.} \emph{藉助任意一個分解}
\[
\mathsf{A}_\Gamma=\mathfrak q_1+\cdots+\mathfrak q_n
\]
\emph{把 \(\mathsf{A}_\Gamma\) 分解成共軛左理想而形成的 \(\mathsf{A}_\Gamma\) 的因子系，與 \(\mathfrak K_r\) 的因子系相同。}

下面的定理說明：若 \(\alpha_{ik}^{(j)}\)、\(\beta_{ik}^{(j)}\) 分別是 \(\{\mathfrak K\}\)、\(\{\mathfrak S\}\) 以 \(\mathsf{Z}\) 為基礎的因子系，則 \(\varepsilon_{ik}^{(j)}=\alpha_{ik}^{(j)}\cdot\beta_{ik}^{(j)}\) 是 \(\{\mathfrak K\times\mathfrak S\}\) 以 \(\mathsf{Z}\) 為基礎的因子系。因此，把\emph{乘積 \(\{\alpha\}\times\{\beta\}\)}定義為相伴因子系的兩個類 \(\{\alpha\}\)、\(\{\beta\}\) 所對應的類 \(\{\alpha_{ik}^{(j)}\beta_{ik}^{(j)}\}\)，是有意義的。

\textbf{定理 2.} \emph{\(\{\mathfrak K\}\) 以 \(\mathsf{Z}\) 為基礎時具有相伴因子系的類 \(\{\alpha\}\)。所有這些 \(\{\alpha\}\)，即所有 \(\{\mathfrak K\}\) 以 \(\mathsf{Z}\) 為分裂域時所得的類，組成一個群；映射 \(\{\mathfrak K\}\mapsto\{\alpha\}\) 將這個群同構地映到 \(\mathscr K_{\mathsf{Z}}\) 上。}

\textbf{為證明這一點}，必須證明兩件事：

1. 若 \(\{\mathfrak K\}\) 的因子係為 \(\alpha_{ik}^{(j)}\)，\(\{\mathfrak S\}\) 的因子係為 \(\beta_{ik}^{(j)}\)，則 \(\varepsilon_{ik}^{(j)}=\alpha_{ik}^{(j)}\beta_{ik}^{(j)}\) 是 \(\{\mathfrak K\times\mathfrak S\}\) 的因子系。

2. 若 \(\{\mathfrak K\}\) 的因子係為 \(\alpha_{ik}^{(j)}=1\)，則 \(\mathfrak K=\mathsf P\)。

1. 設 \(\mathsf{Z}\) 在 \(\mathfrak K\)、\(\mathfrak S\) 中的不可約表示次數分別為 \(r,s\)，並設
\[
\mathfrak K_{r\Gamma}=\mathfrak l_1+\cdots+\mathfrak l_n,\qquad
\mathfrak S_{s\Gamma}=\mathfrak m_1+\cdots+\mathfrak m_n
\]
是共軛理想分解，其中 \(\mathfrak l_i\) 與 \(\mathfrak m_i\) 應是相互對應的理想（分別由 \(\mathfrak Z_\Gamma\) 與 \(\mathfrak Z'_\Gamma\) 的對應分量導出）。於是
\[
(\mathfrak K_r\times\mathfrak S_s)_\Gamma
=\mathfrak l_1\mathfrak m_1+\cdots+\mathfrak l_i\mathfrak m_j+\cdots+\mathfrak l_n\mathfrak m_n
\]

是分解成 \(n^2\) 個非零、因而簡單的理想的分解。因此
\[
\mathsf{A}_1=\mathfrak l_1\mathfrak m_1+\mathfrak l_2\mathfrak m_2+\cdots+\mathfrak l_n\mathfrak m_n
\]
是絕對長度為 \(n\) 的一個理想。\(\mathsf{A}_1\) 在所有 \(S\) 下不變，因為 \(\mathfrak m_i\) 與 \(\mathfrak l_i\) 按同樣方式變換；因此存在一個 \(\mathfrak K_r\times\mathfrak S_s\)-理想 \(\mathsf{A}\)，使
\[
\mathsf{A}_1=\mathsf{A}_\Gamma
\]
（引理，第 18 頁）。

若 \(u\) 是 \(\mathfrak Z\) 在自同構域 \(\mathfrak M\) 中的不可約表示次數，而這個域對應於 \(\mathfrak K\times\mathfrak S\) 的簡單理想，則 \(\mathsf{A}\) 分解成 \(u\) 個簡單理想（因為絕對長度為 \(n\)）。因此，我們的引理給出：\(\mathsf{A}\) 的因子系也就是 \(\mathfrak M\) 的因子系。

而 \(\mathsf{A}_\Gamma\) 分解成共軛理想的一個分解，正是定義中使用的分解：
\[
\tag{3}
\mathsf{A}_\Gamma=\mathfrak l_1\mathfrak m_1+\mathfrak l_2\mathfrak m_2+\cdots+\mathfrak l_n\mathfrak m_n.
\]
若 \(\alpha_{ik}^{(j)}\)、\(\beta_{ik}^{(j)}\) 分別屬於僞矩陣單位 \(p_{ik}\)、\(q_{ik}\)，而這些單位分別屬於 \(\mathfrak K\)、\(\mathfrak S\)，則顯然 \(r_{ik}=p_{ik}q_{ik}\) 是屬於分解 (3) 的 \(\mathsf{A}_\Gamma\) 的一組僞矩陣單位；相應的因子系 \(\varepsilon_{ik}^{(j)}=\alpha_{ik}^{(j)}\beta_{ik}^{(j)}\) 因而也是 \(\mathfrak M\) 的一個因子系。

2. 設 \(\{\mathfrak K\}\) 的因子係為 \(\alpha_{ik}^{(j)}=1\)，相應的僞矩陣單位為 \(p_{ik}\)，並有共軛理想分解
\[
\mathfrak K_{r\Gamma}=\mathfrak l_1+\cdots+\mathfrak l_n
\]
我們必須證明 \(\mathfrak K=\mathsf P\)，或 \(\mathfrak K_r=\mathsf P_r\)，或證明存在一個有 \(n\) 項的左理想 \(\mathfrak l\)，它屬於 \(\mathfrak K_r\)；因為 \(\mathfrak K_r\) 的一個簡單左理想的秩為 \(t^2r\)，故必有 \(n=t^2rh\)，其中 \(h\) 為整數，或者必有 \(t=1\)。

\(\mathfrak K_{r\Gamma}\)-理想
\[
\mathfrak L=\mathfrak l_1(p_{11}+\cdots+p_{1n})
\]
非零——\(e_1\sum_i p_{1i}\ne0\)——並且簡單，因而秩為 \(n\)。\(\Omega\) 在 \(\mathfrak G\) 下不變：
\[
S\mathfrak L
=\mathfrak L\mathfrak l_1\sum_\nu S p_{1\nu}
=\mathfrak l_i\sum_\mu p_{i\mu}
=\mathfrak l_i\sum_\mu p_{i1}p_{1\mu}
\]
因為 \(\alpha_{ik}^{(j)}=1\)。所以
\[
S\mathfrak L=\mathfrak l_ip_{i1}\sum_\mu p_{1\mu}
=\mathfrak l_1\sum_\mu p_{1\mu}
=\mathfrak L.
\]
從而 \(\mathfrak L=\mathfrak l_\Gamma\)，\emph{其中 \(\mathfrak l\) 表示 \(\mathfrak K_r\) 的一個秩為 \(n\) 的理想。}

\textbf{定理 3.} \emph{每個元素 \(\{\mathfrak K\}\) 若屬於群 \(\mathscr K\)，都有有限階 \(\lambda\)；這個階是 \(t\) 的一個因子，而後者是 \(\mathfrak K\) 的絕對分量數。}

\emph{第一證明（Brauer）.} 我們取 \(\mathfrak K_{r\Gamma}\) 分解成共軛左理想的兩個分解
\[
\mathfrak K_{r\Gamma}=\mathfrak l_1+\cdots+\mathfrak l_n=\mathfrak z_1+\cdots+\mathfrak z_n,\qquad
\mathfrak l_i=\mathfrak K_{r\Gamma}e_i,\quad
\mathfrak z_i=\mathfrak K_{r\Gamma}e_i^*.
\]
設 \(a_i\mapsto a_ir_{ik}\)，特別是 \(e_i\mapsto r_{ik}\)，給出從 \(\mathfrak l_i\) 到 \(\mathfrak z_k\) 的一個算子同構。在其逆映射下，\(e_k^*\) 變為一個 \(s_{ki}\)，而這個元素屬於 \(\mathfrak l_i\)，且 \(s_{ki}\cdot r_{ik}=e_k^*\)。若將這些 \(r_{ik}\) 選為共軛的，
\[
S(r_{ik})=r_{\nu\mu}\quad\text{若}\quad S(i,k)=(\nu,\mu),
\]
則這些 \(s_{ik}\) 也共軛：
\[
S(s_{ki})S(r_{ik})=S(s_{ki}r_{ik})=S(e_k^*)=e_\mu^*
=s_{\mu\nu}r_{\nu\mu}=s_{\mu\nu}S(r_{ik}),
\]
故 \(S(s_{ki})=s_{\mu\nu}\)。由於 \(e_i\mapsto e_ir_{ij}s_{jk}\) 給出從 \(\mathfrak l_i\) 到 \(\mathfrak l_k\) 的一個算子同構，必有
\[
r_{ij}s_{jk}=p_{ik}\xi_{ik}^{(j)},\qquad
\xi_{ik}^{(j)}\ne0\in\Gamma
\]

由於這些 \(r\)、\(s\)、\(p\) 都是共軛的，\(\xi\) 也如此：由 \(S(i,j,k)=(\nu,\tau,\mu)\) 得
\[
S(\xi_{ik}^{(j)})=\xi_{\nu\mu}^{\tau}.
\]

由
\[
r_{i\lambda}s_{\lambda j}r_{j\lambda}\varepsilon_{\lambda k}
=r_{i\lambda}e_\lambda^*s_{\lambda k}
=r_{i\lambda}s_{\lambda k}
\]
得到
\[
p_{ij}\xi_{ij}^{(\lambda)}\cdot p_{jk}\xi_{ik}^{(\lambda)}
=p_{ik}\xi_{ik}^{(\lambda)}
\]
或者，由於 \(p_{ij}p_{jk}=p_{ik}\alpha_{ik}^{(j)}\)，得到
\[
\alpha_{ik}^{(j)}
=\frac{\xi_{ik}^{(\lambda)}}{\xi_{ij}^{(\lambda)}\xi_{jk}^{(\lambda)}}.
\]
置 \(\delta_{\nu\mu}=\prod_{\lambda}\xi_{\nu\mu}^{(\lambda)}\)，則 \(\delta_{\nu\mu}\) 在所有 \(S\) 下不變，只要該自同構保持 \(\nu\) 與 \(\mu\) 不動，因而屬於 \(\{\mathsf{Z}_\nu,\mathsf{Z}_\mu\}\)；此外，\(\delta_{\nu\mu}\) 也共軛，因為 \(\xi_{\nu\mu}^{(j)}\) 是共軛的。現在既然
\[
\alpha_{ik}^{(j)n}
=\frac{\prod_{\lambda}\xi_{ik}^{(\lambda)}}
       {\prod_{\lambda}\xi_{ij}^{(\lambda)}\prod_{\lambda}\xi_{jk}^{(\lambda)}}
=\frac{\delta_{ik}}{\delta_{ij}\delta_{jk}}
\]
則由相伴因子系的公式（第 32 頁）可知，\(\alpha_{ik}^{(j)n}\) 與系 \(\beta_{ik}^{(j)}=1\) 相伴；又由於 \(\{\mathfrak K\}\) 的階等於 \(\{\alpha\}\) 的階，所以 \(\{\mathfrak K\}\) 的階是 \(n\) 的一個因子。特別把 \(\mathfrak Z\) 選為 \(\mathfrak K\) 的一個極大交換子域時，\(n=t\)，故 \(\lambda\) 是 \(t\) 的一個因子。

\emph{第二證明（Schur）.} 這個證明以如下表示為基礎：把 \(p_{ik}\) 表示為正則的 \(n\) 階矩陣 \(P_{ik}\)，矩陣取自 \(\Gamma\)，並滿足與 \(p_{ik}\) 相同的共軛條件。矩陣 \(P_{ik}\) 屬於 \(\{\mathsf{Z}_i,\mathsf{Z}_k\}\)；若 \(p_{ik}\) 對應於 \(P_{ik}\)，則 \(S(p_{ik})\) 對應於 \(S(P_{ik})\)。這一表示是同態的；亦即，若 \(P_{ik}\) 對應於 \(p_{ik}\)，\(P_{kj}\) 對應於 \(p_{kj}\)，則 \(P_{ik}P_{kj}\) 對應於 \(p_{ik}p_{kj}\)，或者
\[
P_{ik}P_{kj}=\alpha_{ij}^{(k)}P_{ij}
\]
（關係 \((P_{ik_1}P_{k_2j}=0\) 在 \(k_1\neq k_2\) 時不可能成立）。有了這種 \(p_{ik}\) 的表示後，推理如下：轉到行列式便得
\[
|P_{ik}|\cdot |P_{kj}|=\alpha_{ij}^{(k)n}\cdot |P_{ij}|,
\]
所以，因為矩陣 \(P_{ik}\) 正則，\(|P_{\nu\mu}|\ne0\)，
\[
\alpha_{ij}^{(k)n}=\frac{\delta_{ik}\delta_{kj}}{\delta_{ik}},
\]
其中置 \(|P_{\nu\mu}|=\delta_{\nu\mu}\)。此後與 Brauer 的證明相同。

\emph{將 \(p_{ik}\) 表示為矩陣 \(P_{ik}\)：}

分解中的理想 \(\mathfrak l_i\)：
\[
\mathfrak K_{r\Gamma}=\mathfrak l_1+\cdots+\mathfrak l_n
\]
是 \(\mathfrak K_r\) 在 \(\Gamma\) 中的互反表示模（定理 2，第 32 頁）。固定一組基 \(k_1,\ldots,k_n\)，它是 \(\mathfrak K_r\) 關於 \(\mathfrak Z\) 的基，則
\[
\mathfrak K_{r\Gamma}=k_1\mathfrak Z_\Gamma+\cdots+k_n\mathfrak Z_\Gamma,
\]
所以
\[
\mathfrak l_i=\mathfrak K_{r\Gamma}e_i
=k_1\mathfrak Z_\Gamma e_i+\cdots+k_n\mathfrak Z_\Gamma e_i
=k_1e_i\Gamma+\cdots+k_ne_i\Gamma.
\]
由此得到
\[
\mathfrak l_k=\mathfrak l_i p_{ik}
=k_1p_{ik}\Gamma+\cdots+k_np_{ik}\Gamma.
\]
兩組基 \((k_1e_i,\ldots,k_ne_i)\) 與 \((k_1p_{ik},\ldots,k_np_{ik})\) 都屬於 \(\mathfrak l_k\)，可以通過乘以一個正則矩陣 \(P_{ik}\) 相互轉換：
\[
P_{ik}\cdot
\begin{pmatrix}
k_1e_i\\ \vdots\\ k_ne_i
\end{pmatrix}
=
\begin{pmatrix}
k_1p_{ik}\\ \vdots\\ k_np_{ik}
\end{pmatrix},
\qquad
\text{簡記為}\quad
P_{ik}(k_\varrho e_i)=(k_\varrho p_{ik}).
\]
由此得到 \(P_{ik}\) 的共軛性：
\[
S(P_{ik})(k_\varrho S(e_k))=(k_\varrho S(p_{ik}))
\quad\text{或}\quad
S(P_{ik})(k_\varrho e_\mu)=(k_\varrho p_{\nu\mu})
=P_{\nu\mu}(k_\varrho e_\mu),
\]
即
\[
S(P_{ik})=P_{\nu\mu}\quad\text{若}\quad S(i,k)=(\nu,\mu).
\]
由此又可知 \(P_{ik}\) 屬於 \(\{\mathsf{Z}_i,\mathsf{Z}_k\}\)，即屬於那些 \(S\) 的不變域，而這些自同構保持 \(i\) 與 \(k\) 不動。同態性也容易看出：
\[
\begin{split}
P_{ik}P_{kj}(k_\varrho e_j)
&=P_{ik}(k_\varrho p_{kj})
=(k_\varrho p_{ik}p_{kj})\\
&=\alpha_{ij}^{(k)}(k_\varrho p_{ij})
=\alpha_{ij}P_{ij}(k_\varrho e_j),
\quad\text{因而}\\
&\cdot\, P_{ik}P_{kj}=\alpha_{ij}^{(k)}P_{ij},
\quad\text{證畢。}
\end{split}
\]

定理 3 可以加強如下：

\textbf{定理 4.} \emph{若 \(p\) 是 \(t\) 的一個素因子，而後者是 \(\mathfrak K\) 的絕對分量數，則 \(p\) 也整除階 \(\lambda\)，即 \(\{\mathfrak K\}\) 的階。（Brauer）}

\emph{證明.} 設 \(p^\sigma\) 是 \(p\) 的整除次數 \(n\) 的最高冪，其中這個次數屬於分裂域 \(\mathsf{Z}\)；設 \(\mathfrak H\) 的階為 \(p^\sigma\)，它是 \(\mathfrak G\) 的一個子群（由 \emph{Sylow 定理}知其存在，Speiser，第 1 版，§ 19，第 42 頁，定理 43），並設 \(\mathsf{T}\) 是 \(\mathfrak H\) 的不變域。

於是 \([\Gamma:\mathsf{T}]=p^\sigma\)，而 \([\mathsf{T}:\mathsf P]=n/p^\sigma\) 不被 \(p\) 整除，更不被 \(t\) 整除。因此 \(\mathsf{T}\) 不是分裂域；自同構域 \(\mathfrak S\)，即屬於 \(\mathfrak K_{\mathsf{T}}=\mathfrak S_m\) 的那個域，不等於 \(\mathsf P\)。但 \(\mathfrak S\) 以 \(\Gamma\) 為分裂域，而其在中心 \(\mathsf{T}\) 上的次數——這個中心屬於 \(\mathfrak S\)——為 \(p^\sigma\)，所以 \(\mathfrak S\) 的絕對分量數是某個冪 \(p^\beta>1\)；因而由定理 3，\(\{\mathfrak S\}\) 的階是某個冪 \(p^\alpha>1\)，它是 \(p\) 的冪。又由
\[
\{\mathfrak K\}^{\lambda}=\{\mathsf P\},
\qquad
\{\mathfrak K_{\mathsf{T}}\}^{\lambda}=\{\mathsf{T}\}
\]
可知 \(\lambda\equiv0\pmod {p^\alpha}\)，證畢。（關於不一定有 \(\lambda=t\)，Brauer 給出了例子\srcfn{1}{Brauer: Untersuchungen über die arithmetischen Eigenschaften von Gruppen linearer Substitutionen II. Math. Zeitschr. 31, S. 733 § 5, (1930).}。）

\textbf{定理 5.} \emph{設 \(t\) 的素因子分解為 \(t=\prod_i p_i^{\varrho_i}\)，其中 \(\varrho_i\geqq1\)，且 \(p_i\ne p_j\)，如果 \(i\ne j\)。則 \(\mathfrak K=\mathfrak K^{(1)}\times\mathfrak K^{(2)}\times\cdots\)，其中 \(\mathfrak K^{(i)}\) 是絕對分量數為 \(p_i^{\varrho_i}\) 的一個域。}

\emph{證明.} 由定理 3 與 4，階 \(\lambda\)，即 \(\{\mathfrak K\}\) 的階，為
\[
\lambda=\prod_i p_i^{\alpha_i},\qquad 1\leqq\alpha_i\leqq\varrho_i.
\]
由 \(\{\mathfrak K\}\) 生成的循環子群 \(\mathscr \mathsf{Z}\) 位於 \(\mathscr K\) 中，並且是子群 \(\mathscr \mathsf{Z}_i\) 的直積，這些子群的階分別為 \(p_i^{\alpha_i}\)，故
\[
\{\mathfrak K\}=\prod_i\{\mathfrak K^{(i)}\},
\]

由此可知，每個分裂域 \(\mathsf{Z}\) 若屬於 \(\mathfrak K\)，也是 \(\mathfrak K^{(i)}\) 的分裂域；所以絕對分量數 \(p_i^{\sigma_i}\)，即 \(\mathfrak K^{(i)}\) 的絕對分量數，整除 \(t\)，且 \(\sigma_i\leqq\varrho_i\)。設 \(\prod_i \mathfrak K^{(i)}=\mathfrak K_m\)，則
\[
\prod p_i^{\varrho_i}=\prod p_i^{\varrho_i}\cdot m,
\qquad
\sigma_i\geqq\varrho_i,
\]
因而
\[
\sigma_i=\varrho_i,\qquad m=1,\qquad
\mathfrak K=\prod \mathfrak K^{(j)},
\quad\text{證畢。}
\]

\subsection*{§ 25. 具有伽羅瓦極大交換子域的 \(\mathfrak K_r\) 的正規表示\protect\footnotemark}
\label{normaldarstellung-von-r_r-mit-galoisschem-maximalem-kommutativem-teilkoerper-1}
\footnotetext{從 § 27 起有更簡單的論證。}

設 \(\mathfrak Z\) 是 \(\mathfrak K_r\) 的極大交換子域，並且是伽羅瓦的。由定理 7（第 25 頁），這 \(n\) 個自同構 \(\mathsf{H}_i\) 若屬於 \(\mathfrak Z/\mathsf P\)，就是 \(\mathfrak K_r\) 的內自同構；亦即存在元素 \(u_i\)，使
\[
\mathsf{H}_i z=u_i^{-1}zu_i\quad\text{對所有 }z\text{；其取值範圍為 }\mathfrak Z.
\]
每個 \(u_i\) 在相差一個非零因子 \(y\) 的意義下唯一確定，而該因子來自 \(\mathfrak Z\)。事實上，若 \(u^{-1}zu=u'^{-1}zu'\) 對所有 \(z\) 成立，而這些元素來自 \(\mathfrak Z\)，則 \((u^{-1}u')^{-1}z(u^{-1}u')=z\) 對所有 \(z\in\mathfrak Z\) 成立。因此，把 \(u^{-1}u'\) 添入 \(\mathfrak Z\) 所得的是 \(\mathfrak K_r\) 的一個交換子域；由 \(\mathfrak Z\) 的極大性，這個子域與 \(\mathfrak Z\) 重合。所以 \(u^{-1}u'\) 屬於 \(\mathfrak Z\)。

我們證明下列定理：
\[
\mathfrak K_r=u_1\mathfrak Z+\cdots+u_n\mathfrak Z.
\]
顯然，只需證明這些 \(u_i\) 關於 \(\mathfrak Z\) 線性無關。證明通過直接給出 \(u_i\) 來進行。設
\[
\mathfrak Z_{\mathsf{Z}}=e_1\mathsf{Z}+\cdots+e_n\mathsf{Z},\qquad
\mathfrak K_{r\mathsf{Z}}=\mathfrak l_1+\cdots+\mathfrak l_n,\quad
\mathfrak l_i=\mathfrak K_{r\mathsf{Z}}e_i
\]
並設 \(p_{ik}\) 是屬於這個分解的一組僞矩陣單位。第 8 頁中，我們把由表示模 \(e_i\mathsf{Z}\) 給出的、從 \(\mathfrak Z\) 到 \(\mathsf{Z}\) 的同構記作 \(\Theta_i\)，於是
\[
S_i=\Theta_i\Theta_1^{-1}
\]
是這 \(n\) 個 \(\mathsf{Z}\) 的自同構（但這只限於把它們施用於 \(\mathsf{Z}\) 本身；這裡始終作為基礎的、把 \(S_i\) 施用於整個 \(\mathfrak K_{r\mathsf{Z}}\) 的作用則不再如此！）。為簡單起見，對於 \(\Theta_i\) 在 \(S_i=\Theta_i\Theta_1^{-1}\) 中的寫法，我們改寫為 \(\Theta_{S_i}\)；同時也在 \(p_{ik}\) 中把指標 \(i,k\) 替換為相應的自同構 \(S\)：
\[
p_{ik}=p_{S_iS_k}.
\]

這 \(n\) 個 \(\mathfrak Z\) 的自同構為（\(E\) 是恆等代換）
\[
\mathsf{H}=\Theta_E^{-1}\Theta.
\]
現在把這些 \(\mathsf{H}\) 編號如下：
\[
\mathsf{H}_S=\Theta_E^{-1}\Theta_{S^{-1}}.
\]

在這些條件下，可以置

\[
u_T=\sum_S p_{S,ST}=\sum_S S(p_{E,T})
\]

這是因為：

1. \(u_T^{-1}\) 存在：
\[
u_T^{-1}=\sum_S p_{ST,S}
\frac{1}{\alpha_{S,S}^{ST}\overline{\gamma}_{S,ST}\overline{\gamma}_{ST,S}}
\]

（\(\bar{\gamma}\) 的意義見第 32 頁。）

2. \(u_T\) 是 \(\mathfrak K_r\)-元素，因為它在 \(\mathfrak G\) 下不變：
\[
R(u_T)=\sum_S RS(p_{E,T})=\sum_{S'} S'(p_{E,T})=u_T.
\]

3. \(u_T\) 生成 \(\mathsf{H}_T\)；亦即，若 \(z\in\mathfrak Z\)，則 \(u_T^{-1}zu_T=\mathsf{H}_T(z)\)，或者
\[
z u_T=u_T \mathsf{H}_T(z).
\]

事實上

\[
z=\sum_S\Theta_S(z)e_S,\qquad
u_T=\sum_S S(e_Ep_{E,T})=\sum_S e_S S(p_{E,T})
\]

所以

\[
zu_T=\sum_{S,S'}\Theta_S(z)e_Se_{S'}S(p_{E,T})
=\sum_S\Theta_S(z)S(p_{E,T}).
\]

另一方面，

\[
u_T=\sum_S S(p_{E,T}e_T)=\sum_S S(p_{E,T})e_{ST}
\]

並且，對於施用於 \(\mathfrak Z\)-元素，有如下等式：

\[
\Theta_S \mathsf{H}_T=\Theta_S\Theta_E^{-1}\Theta_{T^{-1}}
=S\Theta_{T^{-1}}=ST^{-1}\Theta_E=\Theta_{ST^{-1}}
\]

所以

\[
\mathsf{H}_T(z)=\sum_S e_S\Theta_S \mathsf{H}_T(z)=\sum_S e_S\Theta_{ST^{-1}}(z)
\]

即

\[
\begin{aligned}
u_T\mathsf{H}_T(z)
&=\sum_{S,S'} S(p_{E,T})e_{ST}\cdot e_{S'}\Theta_{S'T^{-1}}(z)\\
&=\sum_S S(p_{E,T})\Theta_S(z)
\end{aligned}
\]

因此

\[
z\cdot u_T=u_T\cdot \mathsf{H}_T(z).
\]

4. 這些 \(u_S\) 關於 \(\mathfrak Z\) 線性無關。

由
\[
\sum_T z_Tu_T=0
\]
並利用 \(z=\sum_S\Theta_S(z)e_S\)，得到

\[
\sum_{T,S,L}\Theta_S(z_T)e_Sp_{L,LT}
=\sum_{S,T}\Theta_S(z_T)p_{S,ST}\quad\text{所以}\quad
\Theta_S(z_T)=0,
\]

從而 \(z_T=0\)。

從現在起，我們把 \(\mathsf{H}_S(z)\) 寫作 \(z^S\)，稱為“符號冪”。這些 \(u_S\) 必須滿足如下形式的關係：

\[
u_Ru_T=u_{RT}a_{R,T},\qquad a_{R,T}\ne0\quad\text{且來自 }\mathfrak Z.
\]

由於

\[
\begin{aligned}
u_Ru_T
&=\sum_{S,S'} S(p_{E,R})S'(p_{E,T})
=\sum_S p_{S,SR}p_{SR,SRT}\\
&=\sum_S p_{S,SRT}\alpha_{S,SRT}^{(SR)}
=\sum_S p_{S,SRT}\alpha_{S,SRT}^{(SR)}e_{SRT}
\end{aligned}
\]

並且

\[a_{R,\,T} = \sum\limits_{S} e_{SRT}\, \Theta_{SRT} \left(a_{R,\,T}\right),\]

所以

\[u_{RT} a_{R,T} = \sum\limits_{S} p_{S,SRT} \Theta_{SRT} (a_{R,T})\]

從而

\[\Theta_{SRT}(a_{R,\,T}) = \alpha_{S,SRT}^{(SR)}, ~~ a_{R,\,T} = \sum\limits_{S} e_{SRT} \alpha_{S,SRT}^{(SR)} = \sum\limits_{L} e_{L} \alpha_{LT^{-1}R^{-1},L}^{(LT^{-1})}.\]

我們把 \(a_{R,\,T}\) 這個系稱為 \(\mathfrak K_r\) 的小因子系。大因子系 \(\alpha_{BC}^{(A)}\) 與小因子系 \(a_{R,\,T}\) 相互確定。兩個大因子系的乘積 \(\varepsilon_{BC}^{(A)} = \alpha_{BC}^{(A)} \beta_{BC}^{(A)}\)，所對應的小因子系乘積為 \(f_{R,\,T} = a_{R,\,T}b_{R,\,T}\)；這裡大因子系 \(\alpha_{BC}^{(A)}\)、\(\beta_{BC}^{(A)}\) 分別對應於小因子系 \(a_{R,\,T}\)、\(b_{R,\,T}\)。這由上面的公式立即可得。大單位因子系 \(\alpha_{BC}^{(A)} = 1\) 對應於小單位因子系 \(a_{R,\,T} = 1\)，反之亦然。

可以把 \(u_S\) 替換為 \(u_S'=u_Sr_S\)，其中 \(r_S\neq 0\) 是從 \(\mathfrak Z\) 任意選取的。屬於這些 \(u_S'\) 的小因子系記為 \(a_{R,T}'\)：

\[
\text{(1)}\qquad u_R'u_T'=u_{RT}'a_{R,T}'
\]

稱 \(a_{R,T}^{\prime}\) 與 \(a_{R,T}\) 相伴。有

\[u_R r_R u_T r_T = u_{RT} r_{RT} a'_{R,T} = u_R u_T r_R^T r_T = u_{RT} a_{R,T} r_R^T r_T\]

從而

\[
\text{(2)}\qquad a'_{R,T}=a_{R,T}\frac{r_R^T r_T}{r_{RT}}
\]

屬於 \(u_S' = u_S r_S\) 的僞矩陣單位系如下：

\[p_{S,ST}^{\prime}=p_{S,ST}r_{T}\]

因此，小因子系 \(a'_{R,\,T}\) 屬於一個大因子系 \(\alpha^{(A)\prime}_{BC}\)，而這個大因子系與屬於 \(a_{R,\,T}\) 的因子系 \(\alpha^{(A)}_{BC}\) 相伴；相伴的小因子系對應於相伴的大因子系，反之亦然。相伴大因子系的類 \(\{\alpha\}\) 到相應的相伴小因子系的類 \(\{a\}\) 的對應，是一個群同構。

現在我們來用矩陣表示 \(u_S\)，這些 n 階矩陣取自 \(\mathfrak{Z}\)；這種表示與迄今所考察的表示有本質區別。

把 \(\mathfrak Z\)-基

\[
\text{(I)}\qquad k_{S_1}, \ldots, k_{S_n}
\]

把 \(\mathfrak K_r\) 的這組基從右邊乘以 \(u_S\)，得到

\[
\text{(II)}\qquad k_{S_1}u_S, \ldots, k_{S_1}u_S.
\]

把 (I) 變換為 (II) 的（正則）矩陣 \(A_S\)，

\[
\text{(III)}\qquad (k_{S_1}u_S, \ldots, k_{S_n}u_S) = (k_{S_1}, \ldots, k_{S_n})A_S
\]

稱為表示 \(u_S\) 的矩陣。這個表示應當適用於所有 \(\mathfrak K_{r}\)-元素，它們像 \(u_S\) 一樣生成 \(\mathfrak Z/\mathsf{P}\) 的一個自同構；也就是適用於所有元素 \(u_Sr_S\)，其中 \(r_S \neq 0\) 從 \(\mathfrak Z\) 任意選取。所有這些元素組成一個群 \(\mathfrak G^*\)；它含有正規子群 \(\mathfrak Z^*\)，即 \(\mathfrak Z\) 的乘法群。商群 \(\mathfrak G^*/\mathfrak Z^*\) 與 \(\mathfrak G\) 同構，而後者是 \(\mathfrak Z/\mathsf{P}\) 的伽羅瓦群：\(\mathfrak G^*\) 關於 \(\mathfrak Z^*\) 的每個剩餘類中的每個元素，都生成與該剩餘類對應的 \(\mathfrak Z\) 的自同構。群 \(\mathfrak G^*\) 按照剛才給出的方式表示為

\(u_S \mapsto A_S\)。

但這個表示既不是同態表示，也不是互反同態表示；由

\(u_S \mapsto A_S\)，\(u_T \mapsto A_T\)

並不能推出 \(u_Su_T \mapsto A_SA_T\)，而由於

\[
\begin{aligned}
(k_{S_1}, \ldots, k_{S_n})u_Su_T
&=(k_{S_1}, \ldots, k_{S_n})A_Su_T
=(k_{S_1}, \ldots, k_{S_n})u_TA_S^T\\
&=(k_{S_1}, \ldots, k_{S_n})A_TA_S^T:\ u_Su_T\mapsto A_TA_S^T.
\end{aligned}
\]

這個 \(\mathfrak{G}^*\) 的表示稱為交叉表示。若 \(h_{S_1}, \ldots, h_{S_n}\) 是 \(\mathfrak{K}_r\) 關於 \(\mathfrak{Z}\) 的另一組基，並且

\[
\text{(IV)}\qquad (h_{S_1}, \ldots, h_{S_n}) = (k_{S_1}, \ldots, k_S)M
\]

則由 (III) 與 (IV) 得

\[(h_{S_1},\ldots,h_{S_n})\,u_S=(h_{S_1},\ldots,h_{S_n})\,M^{-1}A_S\,M^S,\]

即：若 \(u_S\) 在 \(k\)-基下由 \(A_S\) 表示，則在 \(h\)-基下由 \(M^{-1}A_SM^S\) 表示，其中 \(M\) 是把 \(k\)-基變換成 \(h\)-基的矩陣。

用來定義交叉表示的公式 (III) 表明，交叉表示與環 \(\mathfrak K_r\) 的關係，類似於普通的互反表示（但從右邊而不是從左邊書寫）與它的表示模的關係。差別來自：這裡取代交換關係

\[(m^*\tau^*)c = (m^*c)\tau^*\]
（見第 5 頁）

而出現的是

\[
(kz)u_S=(ku_S)z^S\qquad k\text{ 屬於 }\mathfrak K_r,\quad z\text{ 屬於 }\mathfrak Z
\]

也可以構造 \(\mathfrak{G}^*\) 的交叉表示，所藉助的是右理想 \(\mathfrak r\) 的一組基，而這個理想屬於 \(\mathfrak{K}_r\)。\(\mathfrak{K}_r\) 的每個右理想以這種方式生成一個等價表示類；這裡，表示 \(u_S \mapsto A_S, u_S \mapsto B_S\) 稱為等價，如果存在一個矩陣 \(M\) 使 \(B_S = M^{-1}A_SM^S\)。

現在，與其侷限於由 \(\mathfrak K_r\) 的右理想生成的、\(\mathfrak G^*\) 在 \(\mathfrak Z\) 中的表示，我們可以一般地如下定義一個 \(m\) 次交叉表示，即 \(\mathfrak G^*\) 在 \(\mathfrak Z\) 中的表示：

\begin{enumerate}
\item 每個 \(u_S\) 對應一個 \(m\) 階 \(\mathfrak Z\)-矩陣 \(A_S\)。
\item 由 \(u_S \mapsto A_S\)、\(u_T \mapsto A_T\) 推出 \(u_Su_T \mapsto A_TA_S^T\)。
\end{enumerate}

我們取一個 \(\mathfrak Z\)-右模，其秩為 \(m\)：

\[\mathfrak{M}=x_1\mathfrak{Z}+\cdots+x_m\mathfrak{Z}\]

並藉助

\[(x_1u_S,\ldots,x_mu_S)=(x_1,\ldots,x_m)A_S\]

把它變成一個 \(\mathfrak G^*\)-右模，其交換關係為

\[(mz) u_S = (mu_S) z^S\]

這與第 5 頁中由給定表示構造表示模完全相應。現在我們像第 20 頁一樣說明，\(\mathfrak{M}\) 成為一個有限 \(\mathfrak K_r\)-右模，方式是規定

\[m \cdot z u_S = mz \cdot u_S = mu_S \cdot z^S\]

以及一般地規定

\[m \cdot \sum z_S u_S = \sum m z_S \cdot u_S = \sum m u_S \cdot z_S^S.\]

這裡根本不需要類似第 20 頁計算的東西，因為我們的規定使用了一組確定的基。

但根據 M.Z. § 18，每個有限簡單 \(\mathfrak K_r\)-模 \(\mathfrak{M}\) 都與 \(\mathfrak K_r\) 的一個簡單右理想算子同構；因此，與第 22 頁定理 2 相應，我們推得：\(\mathfrak{G}^*\) 在 \(\mathfrak Z\) 中只有\emph{一個}不可約交叉表示類，即屬於簡單右理想的那個類。

這個不可約表示的次數為 \(t\)，其中 \(t\) 是 \(\mathfrak K\) 的指標（絕對指數數）。這由簡單的秩計數得出，因為這個次數等於一個簡單右理想的 \(\mathfrak Z\)-秩 \((\mathfrak r:\mathfrak Z)\)。第 22 頁給出：

\[
(\mathfrak Z:\mathsf{P})=n=rt;\qquad (\mathfrak K_r:\mathsf{P})\equiv n^2;\qquad (\mathfrak r:\mathsf{P})\equiv n\cdot t
\]

最後一個等式是由於 \(\mathfrak K_r = \mathfrak{r}_1 + \cdots + \mathfrak{r}_r\)。

由
\[(\mathfrak{r}: \mathsf{P}) = (\mathfrak{r}: \mathfrak Z) \cdot (\mathfrak Z: \mathsf{P})\]
或者 \(n \cdot t = (\mathfrak{r}: \mathfrak Z) \cdot n\)，得到 \((\mathfrak{r}: \mathfrak Z) = t\)。

\subsection*{§ 26. 交叉表示的乘法}\label{multiplikation-der-verschraenkten-darstellungen}

迄今我們總把交叉表示看作 \(\mathfrak{G}^*\) 的表示。但也可以把它看作 \(\mathfrak{G}\) 的表示，而後者是 \(\mathfrak Z/\mathsf{P}\) 的伽羅瓦群；方式如下：把每個元素 \(S\)，即 \(\mathfrak{G}\) 的元素，配給矩陣 \(A_S\)，而這個矩陣表示一個元素 \(u_S\)；後者來自 \(\mathfrak K_r\)，並通過變換——\(u_S^{-1}zu_S=z^S\)——生成自同構 \(S\)。

為了刻畫這樣一個表示 \(S\mapsto A_S\)\srcfn{1}{這樣一個表示（與超複系無關）
\[
S\mapsto A_S;\quad T\mapsto A_T;\quad ST\mapsto A_T\mathrm{A}_s^T=A_{ST}\cdot a_{S,T^{-1}}
\]
（或者它的轉置）是 Speiser（Math. Zschr. 5）定義因子系 \(a_{S,T}\) 的出發點。因此，這些因子系與這裡從另一基礎引入的因子系相同。}，必須說明它取自哪個 \(\mathfrak K_r\)，使用了哪些 \(u_S\)，以及使用了哪一組 \(k_{S_1},\ldots,k_{S_n}\)，這組基屬於 \(\mathfrak K\) 在 \(\mathfrak Z\) 上的基。更換基得到等價表示，更換 \(u\) 得到相伴因子系 \(a_{S,\,T}\)；因而本質的是取表示的那些 \(\mathfrak K_r\)，亦即相伴因子系的類 \(\{a_{S,\,T}\}\)。

現在給定 \(\mathfrak G\) 的兩個這種表示：

\[
\begin{array}{lll}
S\mapsto u_S\mapsto A_S & \text{相應的小因子系} & a_{R,T},\\
S\mapsto v_S\mapsto B_S & \text{相應的小因子系} & b_{R,T}.
\end{array}
\]

Kronecker 乘積 \(A\times B\)，對於兩個矩陣 \(A=(a_{ik})\)、\(B=(b_{\mu\nu})\)，是指矩陣 \(A\times B=(a_{ik}b_{\mu\nu})\)。

\noindent\textbf{斷言.} \emph{\(S\to A_S\times B_S\) 也是 \(\mathfrak G\) 的一個交叉表示；並且，它屬於乘積 \(\mathfrak K_r\times\mathfrak L_s\)，如果 \(A_S\) 來自 \(\mathfrak K_r\)，而 \(B_S\) 來自 \(\mathfrak L_s\)。}

\noindent\emph{證明.} 表示 \(S\mapsto A_S\times B_S\) 是一個交叉表示。設
\[
S\mapsto A_S\times B_S,\qquad T\mapsto A_T\times B_T
\]
且
\[
ST\mapsto A_{ST}\times B_{ST}
=(A_TA_S^T\times B_TB_S^T)a_{S,T}^{-1}b_{S,T}^{-1},
\]
則必須證明
\[
(A_T\times B_T)(A_S\times B_S)^T a_{S,T}^{-1}b_{S,T}^{-1}
=(A_TA_S^T\times B_TB_S^T)a_{S,T}^{-1}b_{S,T}^{-1}.
\]
這由下列考察得到：

把矩陣 \(A\) 與 \(B\) 分別用矩陣單位 \(c_{ik}\) 與 \(d_{\mu\nu}\) 表示（這些矩陣單位與 \(\mathfrak K_r\) 或 \(\mathfrak L_s\) 的矩陣單位無關）：
\[
A=\sum c_{ik}\alpha_{ik};\qquad
B=\sum d_{\mu\nu}\beta_{\mu\nu},\qquad\text{於是}
\]
\[
A\times B=B\times A
=\sum c_{ik}d_{\mu\nu}\cdot\alpha_{ik}\cdot\beta_{\mu\nu},
\]
也就是等於以新矩陣單位 \(c_{ik}d_{\mu\nu}=d_{\mu\nu}c_{ik}\) 表示的一個矩陣。因此，\(A\)-矩陣與 \(B\)-矩陣的乘法，一一對應於兩個矩陣環的直積中的元素乘法：
\[
\sum c_{ik}\mathfrak Z\times\sum d_{\mu\nu}\mathfrak Z.
\]
所以，每個 \(A\)-矩陣與每個 \(B\)-矩陣可交換（但各個 \(A\)-矩陣彼此之間或各個 \(B\)-矩陣彼此之間並非如此）；而每個 \(A\)-矩陣與每個 \(B\)-矩陣的這種可交換性，在乘以 \(a_{S,T}^{-1}b_{S,T}^{-1}\) 後，恰好給出所要證明的關係。

\section*{第六章. 交叉積理論}\label{kapitel-vi.-theorie-der-verschraenkten-produkte}

\subsection*{§ 27. 把 \(\mathfrak K_r\) 表示為交叉積}

現在要重新獨立地奠定 § 24 中從因子系發展出來的交叉積理論。這些交叉積的意義在於：一個伽羅瓦域及其群在 \(\mathfrak K_r\) 中被\emph{共同}描述。

\noindent\textbf{定理 1.} \emph{若 \(\mathfrak Z\) 是 \(\mathfrak K_r\) 的伽羅瓦分裂域，因而是極大交換子域，而 \(\mathfrak G\) 是它的伽羅瓦群，則 \(\mathfrak K_r\) 等於交叉積 \(\mathfrak Z\cdot\mathfrak G\)。反之，每個交叉積 \(\mathfrak Z\cdot\mathfrak G\) 都生成一個這種 \(\mathfrak K_r\)。}

\noindent\textbf{證明與定義. 定義 1.} 設 \(\mathfrak G=\{E,S,\ldots,T\}\) 是 \(\mathfrak Z\) 的伽羅瓦群；\(\mathfrak K_r\) 等於交叉積，意思是
\[
\mathfrak K_r=\mathfrak Zu_E+\cdots+\mathfrak Zu_T
\]
其中 \(zu_S=u_Sz^S\)，且 \(u_Su_T=a_{S,T}u_{ST}\)（這些 \(a_{S,T}\) 屬於 \(\mathfrak Z\)，並且是小因子系）。

\noindent\textbf{定義 2.} 生成群 \(\mathfrak G^*\) 定義為所有正則元素 \(g\) 的全體，這些元素來自 \(\mathfrak K_r\)，並通過變換把整個 \(\mathfrak Z\) 映入自身。\(g^{-1}\mathfrak Zg=\mathfrak Z\)，所以 \(g^{-1}zg=z^S\)，其中 \(S\) 屬於 \(\mathfrak G\)。因此，\(\mathfrak G\) 是 \(\mathfrak G^*\) 的群同態像，因為 \(\mathfrak Z\) 的每個自同構都由一個 \(g\) 生成，即 \(\mathfrak G^*\to\mathfrak G\)，並且進一步有
\[
\mathfrak G\cong\mathfrak G^*/\mathfrak T^*
\quad\text{且}\quad
\mathfrak T^*\supseteqq\mathfrak Z^*,
\]
（稍後將證明 \(\mathfrak T^*=\mathfrak Z^*\)——尚未證明 \(\mathfrak Z\) 也是極大交換子環；這也可以獨立證明。）

\noindent\textbf{1.} \emph{表示為交叉積的證明.} 設 \(u_E,u_S,\ldots,u_T\) 是 \(\mathfrak K_r\) 中分別生成自同構 \(E,\ldots,T\) 的元素，即 \(zu_S=u_Sz^S\)，特別地，由 2 得 \(u_E=e\)。要證明和
\[
\mathfrak M=\{\mathfrak Zu_E,\ldots,\mathfrak Zu_T\}
\]
是直和 \(\mathfrak Zu_E+\cdots+\mathfrak Zu_T\)，於是由於秩相同，它等於 \(\mathfrak K_r\)。這可如下推出（Schwarz）：

\(\mathfrak M\) 是 \(\mathfrak Z\) 的左模與右模，因而是 \(\mathfrak Z\) 在 \(\mathfrak Z\) 中的表示模；而且 \(\mathfrak M\) 生成這 \(n\) 個不同的表示。所以，因為 \(\mathfrak Z\) 是第一類交換的，\(\mathfrak M\) 的秩至少為 \(n\)；其餘證明完全相同。

\noindent\textbf{2.} \(\mathfrak G\cong\mathfrak G^*/\mathfrak Z^*\)，即 \(\mathfrak T^*=\mathfrak Z^*\)；只有來自 \(\mathfrak Z\) 的元素才通過變換在 \(\mathfrak Z\) 上誘導恆等映射。

事實上，設 \(tz=zt\) 對每個 \(z\) 都成立，而該元素來自 \(\mathfrak Z\)，並置
\[
t=z_0+z_1u_S+\cdots+z_{n-1}u_T,
\]
於是
\[
zt=zz_0+\cdots+zz_{n-1}u_T
\]
而
\[
tz=z_0z+\cdots+z_{n-1}u_Tz
=zz_0+z^{S^{-1}}z_1u_S+\cdots+z^{T^{-1}}z_{n-1}u_T,
\]
因此，由於該和為直和，
\[
(z-z^{S^{-1}})z_1u_S=0,\ldots,
(z-z^{T^{-1}})z_{n-1}u_T=0
\]
並且，因為元素 \(u_S,\ldots\) 正則，得到
\[
(z-z^{S^{-1}})z_1=0,\ldots,
(z-z^{T^{-1}})z_{n-1}=0,
\qquad\text{對每個 }z.
\]
（也可由秩等於 \(n\) 得到。）但按照定義，對 \(S,\ldots\) 等每個不等於 \(E\) 的自同構，至少存在一個 \(z\) 使 \(z-z^{S^{-1}}\neq0\)；也存在一個 \(\bar z\) 使 \(\bar z-\bar z^{T^{-1}}\neq0\)。所以 \(z_i=0\)，\(i=1,\ldots,(n-1)\)，從而 \(t=z_0\) 屬於 \(\mathfrak Z^*\)，證畢。

\noindent\textbf{3.} 這些 \(u\) 滿足乘法律 \(u_Su_T=a_{S,T}u_{ST}\)，其中 \(a_{S,T}\) 滿足交叉結合律
\[
a_{R,S}\cdot a_{RS,T}=a_{S,T}^{R^{-1}}\cdot a_{R,ST}
\]
（這些 \(a_{S,T}\) 是小因子系）。

事實上，由 2 得
\[
\mathfrak Z^*u_S\cdot\mathfrak Z^*u_T=\mathfrak Z^*u_{ST};
\]
因此
\[
u_Su_T=a_{S,T}u_{ST};
\]
其中 \(a_{S,T}\) 來自 \(\mathfrak Z^*\)，因而來自 \(\mathfrak Z\)，且 \(\neq0\)。再由 \(u_R\cdot u_Su_T=u_Ru_S\cdot u_T\)，得到
\[
\begin{aligned}
u_Ra_{S,T}u_{ST}
&=a_{S,T}^{R^{-1}}\cdot u_Ru_{ST}
 =a_{S,T}^{R^{-1}}a_{R,ST}u_{RST}\\
&=a_{R,S}u_{RS}\cdot u_T
 =a_{R,S}a_{RS,T}u_{RST},
\end{aligned}
\]
這就是所述的乘法律。

\noindent\textbf{4. 逆向：}交叉積 \(=\mathfrak K_r\)。給直和
\[
\mathfrak M=\mathfrak Zu_E+\cdots+\mathfrak Zu_T
\]
賦予環結構，所用關係為 \(zu_S=u_Sz^S\)、\(u_Su_T=a_{S,T}u_{ST}\)（其中 \(a_{S,T}\) 滿足交叉結合關係，因而 \(u_Ru_Su_T\) 唯一確定）：
\[
\begin{array}{ll}
1.&z_1\cdot z_2u_S=z_1z_2\cdot u_S;\\[2pt]
2.&zu_S\cdot u_T=z\cdot u_Su_T
\end{array}
\]
由此 \(\mathfrak M\) 成為結合環。

\noindent\textbf{逆定理 2.} \emph{交叉積 \(\mathfrak Z\cdot\mathfrak G\) 是以 \(\mathsf{P}\) 為中心的雙邊單環（依照 R. Brauer 的思路）。}證明基於以下引理：

\noindent\textbf{引理 1.} \emph{每個與 \(\mathfrak Z\) 中所有元素可交換的元素都屬於 \(\mathfrak Z\)。}

\noindent\emph{證明.} 設 \(w=\sum_S c_{\lambda_S}u_S\)（\(c_{\lambda_S}\) 屬於 \(\mathfrak Z\)），並假定 \(wz=zw\) 對每個 \(z\) 成立，而該元素來自 \(\mathfrak Z\)。於是
\[
\sum_S zc_{\lambda_S}u_S
 =\sum_S c_{\lambda_S}u_Sz
 =\sum_S c_{\lambda_S}z^{S^{-1}}u_S
\]
所以有 \(zc_{\lambda_S}=c_{\lambda_S}z^{S^{-1}}\)，或者 \(c_{\lambda_S}(z-z^{S^{-1}})=0\)，對每個 \(z\) 與 \(S\) 都成立（這些 \(u\) 線性無關！）。由此可知 \(c_{\lambda_S}=0\) 當 \(S\neq E\)，因為對每個這樣的 \(S\)，至少存在\emph{一個} \(z\) 使 \(z-z^{S^{-1}}\neq0\)；因此 \(w\) 屬於 \(\mathfrak Z\)。

\noindent\textbf{引理 2.} \emph{每個 \(\mathfrak Z\)-雙模 \(\mathfrak A\)，若與 \(\mathfrak Zu_S=u_S\mathfrak Z\) 雙重同構（\(\mathfrak Z\) 作為左、右算子），都與 \(\mathfrak Zu_S\) 相同。}

\noindent\emph{引理 2 的證明.} \(\mathfrak A\) 作為左模與 \(\mathfrak Zu_S\) 同構；因此，若對應 \(u_S\mapsto a\) 來自 \(\mathfrak A\)，則 \(\mathfrak Zu_S\mapsto\mathfrak Za\)，所以 \(\mathfrak A=\mathfrak Za\)。由雙重同構性得到 (1) \(za=az^S\)（由於 \(u_Sz\mapsto az\)）；並由此得到 \(\mathfrak A=\mathfrak Zu_S\)。事實上，置 \(a=u_Su_S^{-1}a=u_Sw\)，則 \(za=zu_Sw=u_Sz^Sw\)；另一方面，由 (1) 有 \(za=u_Swz^S\)，所以（乘以 \(u_S\)）得到 \(z^Sw=wz^S\)，對每個 \(z\) 都成立，而該元素來自 \(\mathfrak Z\)。於是 \(z^S\) 也跑遍整個 \(\mathfrak Z\)，\(w\) 與 \(\mathfrak Z\) 中所有元素可交換，故由引理 1 屬於 \(\mathfrak Z\)。

\noindent\textbf{引理 3.} \emph{每個 \(\mathfrak Z\)-雙模，若取自 \(\mathfrak Z\cdot\mathfrak G\)，都形如 \(\mathfrak Zu_{S_1}+\cdots+\mathfrak Zu_{S_k}\)，其中 \(u_{S_i}\) 表示任意 \(k\) 個不同者，選自 \(u_S,\ldots,u_T\)。}

\noindent\emph{證明.} 交叉積 \(\mathfrak Z\cdot\mathfrak G\) 作為 \(\mathfrak Z\)-雙模完全可約，因為每個直和項 \(\mathfrak Zu_S\) 的秩為 1。作為雙模，這些直和項又屬於不同的類，因而彼此不雙重同構。所以每個 \(\mathfrak Z\)-雙模 \(\mathfrak M\)，由完全可約性，成為一個直和項，並且本身完全可約。\(\mathfrak M=\bar z\mathfrak A_{S_i}\)，其中 \(\mathfrak A_{S_i}\) 與 \(\mathfrak Zu_{S_i}\) 同構；因而由引理 2，與 \(\mathfrak Zu_{S_i}\) 相同。（所有 \(\mathfrak A_{S_i}\) 屬於不同的類。）

\noindent\textbf{引理 4.} \emph{每個包含 \(\mathfrak Z\) 的環，若取自 \(\mathfrak Z\cdot\mathfrak G\)，特別是 \(\mathfrak Z\cdot\mathfrak G\) 本身，都是雙邊單的。——包含 \(\mathfrak Z\) 的子環與 \(\mathfrak G\) 的子群一一對應。}

\noindent\emph{證明.} 每個包含 \(\mathfrak Z\) 的環 \(\mathfrak o\) 是一個 \(\mathfrak Z\)-雙模，因而由引理 3 形如 \(\sum\mathfrak Z\cdot u_{S_i}\)；因為 \(\mathfrak o\) 是環，兩個 \(u_{S_i}\) 的乘積仍屬於 \(\mathfrak o\)，相應的 \(\mathfrak Z\)-模亦然。由於元素只有有限多個，這些 \(u_{S_i}\) 因而在因子系意義下組成一個群。所以 \(\mathfrak o=\mathfrak Z\cdot\mathfrak H\)，其中 \(\mathfrak H\) 是 \(\mathfrak G\) 的子群。（更準確地說，這些 \(\mathfrak Z\cdot u_{S_i}\) 組成一個乘法系統。）現在若 \(\mathfrak a\) 是一個雙邊理想且 \(\neq0\)，它取自 \(\mathfrak o\)，則是 \(\mathfrak o\)-雙模，因而是 \(\mathfrak Z\)-雙模和 \(\mathfrak o\) 的子模。所以，它除含有某個 \(\mathfrak Zu_{S_i}\) 外還含有整個群 \(\mathfrak H\)，從而成為 \(\mathfrak o=\mathfrak Z\cdot\mathfrak H\)。

\noindent\textbf{引理 5.} \emph{\(\mathsf{P}\) 是 \(\mathfrak Z\cdot\mathfrak G\) 的中心。}

\noindent\emph{證明.} 若 \(w\) 與 \(\mathfrak Z\cdot\mathfrak G\) 中所有元素可交換，則特別與 \(\mathfrak Z\) 中所有元素可交換，所以由引理 1，該元素屬於 \(\mathfrak Z\)。於是，由 \(wu_S=u_Sw^S\)，可交換時得到 \(w=w^S\)，所以 \(w\) 屬於 \(\mathsf{P}\)。

\noindent\emph{注 1.} 證明表明，即使 \(\mathfrak Z\) 是 \(\mathsf{P}\) 上的無限代數伽羅瓦域，定理仍成立，因為關於完全可約性的定理保持成立（Krull，Zahlringe，Math. Ann. 99）。此時，\(\mathfrak Z\cdot\mathfrak G\) 是所有有限和 \(\sum_{S_\lambda}u_{S_\lambda}\cdot c_{S_\lambda}\) 的系統。

只有在引理 4 中，必須對 \(\mathfrak o\) 假定 \(\mathfrak o=\mathfrak Z\cdot\mathfrak H\)；這當然對全交叉積 \(\mathfrak Z\cdot\mathfrak G\) 成立。因為，僅由 \(u_{S_i}u_{S_j}\) 屬於 \(\mathfrak o\)，不能再推出該系統具有群性質，如果該環在 \(\mathfrak Z\) 上為無限。

\noindent\emph{注 2.} 在直接奠基時，應在此處接上 §§ 25/26 中發展的交叉矩陣表示理論。下文 § 28 末尾將用到：不可約交叉表示的次數為 \(t\)，其中 \(t\) 是 \(\mathfrak K\) 的絕對分量數（指標）。再次指出，這個表示理論給出：這裡作為乘法常數引入的因子系，與文獻中通常的概念相同。
\subsection*{§ 28. 因子系的乘積定理}\label{produktsatz-fuer-faktorensysteme}

設
\[
\mathfrak K_r=\mathfrak Z\cdot\mathfrak G
=\mathfrak Zu_E+\cdots+\mathfrak Zu_T,
\qquad u_Su_T=u_{ST}a_{S,T},
\]
\[
\overline{\mathfrak K}_r
=\overline{\mathfrak Z}\,\overline{\mathfrak G}
=\overline{\mathfrak Z}\bar u_E+\cdots+\overline{\mathfrak Z}\bar u_T,
\qquad
\bar u_S\bar u_T=\bar u_{ST}\bar b_{S,T},
\]
（\(\mathfrak Z\) 與 \(\overline{\mathfrak Z}\) 是同構的伽羅瓦域，\(\mathfrak G\)、\(\overline{\mathfrak G}\) 是它們的群），於是
\[
\mathfrak K_r\times\overline{\mathfrak K}_r
=\widetilde{\mathfrak Z}\,\widetilde{\mathfrak G}\times \mathsf{P}_n
=\mathfrak A_f\times \mathsf{P}_n
=\mathfrak A_j
\]
且 \(\tilde u_S\tilde u_T=\tilde u_{ST}\tilde c_{S,T}\)（\(\widetilde{\mathfrak Z}\) 與 \(\mathfrak Z\) 及 \(\overline{\mathfrak Z}\) 同構）。這是由秩的計數以及 \(\mathfrak Z\) 也是分裂域這一事實推出的。

\noindent\textbf{斷言。} \emph{有 \(\tilde c_{S,T}=\tilde a_{S,T}\tilde b_{S,T}\)，其中 \(\tilde a,\tilde b\) 是 \(\widetilde{\mathfrak Z}\) 中在相應同構下分別對應於 \(a\) 和 \(\bar b\) 的元素。}

證明基於如下事實：\(\mathfrak A_j\) 中任一絕對長度為 \(n\) 的左理想的自同態環都與 \(\mathfrak A_f=\widetilde{\mathfrak Z}\widetilde{\mathfrak G}\) 同構，並且對於 \(\mathfrak A_j\) 中一個適當的左理想，這個自同態環可以直接確定。更一般地，下列為方便起見針對 \(\mathfrak K_r\) 表述的定理成立，其中可用 \(s\) 代替 \(n\)：

\noindent\textbf{1. 定理 1.} \emph{\(\mathfrak K_{rs}=\mathfrak K_r\times \mathsf{P}_s\) 中絕對長度為 \(n\) 的左理想 \(\mathfrak l\) 的自同態環與 \(\mathfrak K_r\) 同構。}

證明由以下輔助引理推出。

\noindent\textbf{輔助引理 1.} \emph{\(\mathfrak K_{rs}\) 中所有絕對長度為 \(n\) 的左理想都算子同構。因此，它們的自同態環也環同構。特別地，\(\mathfrak l=\mathfrak K_r\times\mathfrak b\) 具有絕對長度 \(n\)，其中 \(\mathfrak b\) 表示 \(\mathsf{P}_s\) 中的一個單理想。}

\noindent\emph{證明。} 左理想具有絕對長度 \(n\)，即分解成 \(n\) 個絕對不可分解的單理想，這就使它們在 \(\mathfrak K_{rs}\) 中具有相同長度，亦即在 \(\mathfrak K_{rs}\) 中出現同樣多的單理想直和項。由此得到它們的算子同構性，進而得到自同態環的環同構。又 \(\mathfrak K_r=\mathfrak Z\cdot\mathfrak G\) 的絕對長度為 \(n\)，即其秩為 \(n^2\)；所以 \(\mathfrak K_r\times \mathsf{P}_s\) 的絕對長度為 \(n\cdot s\)，從而 \(\mathfrak K_r\times\mathfrak b\) 的絕對長度為 \(n\)。

\noindent\textbf{輔助引理 2.} \emph{若 \(\mathfrak o\) 是帶單位元的環，\(\mathfrak l=\mathfrak o e_1\) 是左理想且為直和項，亦即 \(e_1\) 是冪等元，則 \(\mathfrak l\) 的自同態環同構於 \(e_1\mathfrak o e_1\)。}

\noindent\emph{證明。} 乘以元素 \(e_1re_1\) 給出 \(\mathfrak l\) 到自身的一個算子同態；不同的 \(e_1re_1\) 給出不同的同態，因為它們把 \(e_1\) 對應到不同的元素；而每個同態都以這種方式產生：\(e_1\mapsto re_1\)，並且由 \(e_1e_1=e_1\) 有 \(re_1=e_1re_1\)。而 \(e_1\mathfrak o e_1\) 中的和與積分別對應於所配同態的和與積。

\noindent\textbf{輔助引理 3.} \emph{置 \(\mathsf{P}_s=\sum_1^s c_{ik}\mathsf{P}\)，其中 \(c_{ik}\) 是矩陣單位，並置 \(\mathfrak b=c_{11}\mathsf{P}+\cdots+c_{s1}\mathsf{P}\)，則 \(\mathfrak K_r\times\mathfrak b\) 的自同態環為 \(\mathfrak K_rc_n\)，它與 \(\mathfrak K_r\) 同構。}

\noindent\emph{證明。} 因為 \(c_{11}\) 與 \(\mathfrak K_r\) 的每個元素可交換，並且 \(\mathfrak K_r\) 中沒有元素被 \(c_{11}\) 零化；這是因為 \(c_{11}\) 是 \(\mathfrak K_r\times \mathsf{P}_s\) 的生成冪等元，而且
\[
c_{11}\cdot\mathfrak K_r\times \mathsf{P}_s\cdot c_{11}
=\mathfrak K_r\times c_{11}\mathsf Pc_{11}
=\mathfrak K_rc_{11}.
\]
定理 1 直接由這些輔助引理推出。

\noindent\textbf{2. 確定 \(\mathfrak K_r\times\overline{\mathfrak K}_r\) 並將其分解為絕對長度為 \(n\) 的左理想。} 按定義有
\[
\mathfrak A_f=\mathfrak K_r\times\overline{\mathfrak K}_r
=\mathfrak Z\cdot\mathfrak G\times\overline{\mathfrak Z}\cdot\overline{\mathfrak G},
\]
其中帶橫線與不帶橫線的元素逐一可交換，因而
\[
\mathfrak A_j
=(\mathfrak G\times\overline{\mathfrak G})\cdot(\mathfrak Z\times\overline{\mathfrak Z})
=\mathfrak G\times\overline{\mathfrak G}\cdot(\mathfrak Ze_1+\cdots+\mathfrak Ze_n)
\]
且 \(\mathfrak Ze_i=\overline{\mathfrak Z}e_i\)（一個伽羅瓦域與自身相乘）。由於所有直和項相對於 \(\mathfrak Z\) 都有相同的秩，因而也有相同的絕對長度，而 \(\mathfrak A_f\) 的絕對長度 \(=n^2\)，所以每個直和項的絕對長度都是 \(n\)。下文以
\[
\mathfrak l=(\mathfrak G\times\overline{\mathfrak G})\mathfrak Ze_1,
\]
為基礎，其中用 \(ze_1=e_1\bar z\) 固定 \(\mathfrak Z\) 與 \(\overline{\mathfrak Z}\) 之間的一個特定同構。

\noindent\textbf{3. \(e_1\mathfrak A_fe_1\) 的交叉表示，並由此證明乘積定理。} --- 根據 1 和 2，乘積 \(\mathfrak K_r\times\overline{\mathfrak K}_r\) 的因子系同構於 \(e_1\mathfrak A_je_1\equiv\mathfrak A\) 的一個交叉表示的因子系。

\noindent\textbf{定理 2.} \emph{\(\mathfrak A=\widetilde{\mathfrak Z}\cdot\widetilde{\mathfrak G}\)，其中 \(\widetilde{\mathfrak Z}=\mathfrak Ze_1=\overline{\mathfrak Z}e_1\)，\(\overline{\mathfrak G}=\{-,e_1u_S\bar u_Se_1,-\}\)，並且
\[
e_1u_S\bar u_Se_1\cdot e_1u_T\bar u_Te_1
=e_1u_{ST}\bar u_{ST}e_1\cdot a_{S,T}e_1\bar b_{S,T}e_1,
\]
所以令 \(\tilde a_{S,T}=a_{S,T}e_1\) 及 \(\tilde b_{S,T}=\bar b_{S,T}e_1\)，乘積定理即得證；（\(\tilde a\) 與 \(\tilde b\) 位於同一個域 \(\widetilde{\mathfrak Z}\) 中）。}

\noindent\emph{證明。} \(e_1\) 成為該環的單位元；與 \(\tilde z=ze_1\) 共軛的元素於是由 \(\tilde z^Se_1\) 給出；必須證明元素 \(u_S\bar u_Se_1\) 誘導這些置換，從而證明該交叉表示。

\noindent\textbf{輔助引理 1.} \(e_1u_S\bar u_S=u_S\bar u_Se_1\)，因而也即 \(=e_1u_S\bar u_Se_1\)。

\noindent\emph{證明。} 置 \(u_S^{-1}e_1u_S=e^S\)\srcfn{1}{這一約定使得 \(e_1u_S=u_Se^S\)。}，亦即 \(e_1=e^E\)，則 \(e^S\) 跑遍 \(e_1,e_2,\ldots,e_n\) 的共軛元。於是有 (1)\srcfn{2}{因為 \(e_1\bar u_R=\bar u_Re^{R^{-1}}\)。}
\[
e_1\bar u_R=\bar u_Re^{R^{-1}}.
\]
事實上，\(e_1=\sum a_i\bar A_i=\sum a_i^R\bar A_i^R\)，其中 \(a_i\) 跑遍任一 \(\mathfrak Z\)-基，而 \(\bar A_i\) 跑遍與之互補的 \(\overline{\mathfrak Z}\)-基；所以同一事實也適用於 \(a_i^R\) 和 \(\bar A_i^R\)。因此 \(e_1\bar u_R=\bar u_Re^{R^{-1}}\)，從而 (1) 得證。由 (1) 推出
\[
e_1u_S\bar u_S=u_Se^S\bar u_S=u_S\bar u_S\cdot e^{SS^{-1}}=u_S\bar u_Se_1,
\]
輔助引理 1 得證。

\noindent\textbf{輔助引理 2.} \(ze_1\cdot e_1u_S\bar u_Se_1=e_1u_S\bar u_Se_1\cdot z^Se_1\)。

\noindent\emph{證明。} 因為
\[
ze_1\cdot e_1u_S\bar u_Se_1
=zu_S\bar u_Se_1\quad(\text{輔助引理 1})
=u_Sz^S\bar u_Se_1
=u_S\bar u_Se_1\cdot z^Se_1
\equiv b=e_1b,
\]
（\(z\) 與 \(\bar u\) 及 \(e_1\) 可交換）。--- 令 \(u'=u_S\bar u_Se_1\)。輔助引理 2 說明 \(u'\) 標引 \(\mathfrak Ze_1\) 的置換，因而定理所給的交叉表示成立。

\noindent\textbf{輔助引理 3.} \emph{因子系等於 \(a_{S,T}e_1\cdot\bar b_{S,T}e_1\)。}

\noindent\emph{證明。}
\[
\begin{aligned}
e_1u_S\bar u_Se_1\cdot e_1u_T\bar u_Te_1
&=u_S\bar u_S\cdot u_T\bar u_Te_1
 =u_{ST}a_{S,T}\cdot\bar u_{ST}\bar b_{S,T}e_1\\
&=u_{ST}\bar u_{ST}\cdot a_{S,T}\bar b_{S,T}e_1
 =e_1u_{ST}\bar u_{ST}e_1\cdot a_{S,T}e_1\cdot\bar b_{S,T}e_1.
\end{aligned}
\]
這樣，2 中的定理即告證明，因而因子系的乘積定理也得證。

\noindent\textbf{4. 相伴因子系的類的乘積定理：}
\[
\{a_{S,T}\}\cdot\{b_{S,T}\}=\{a_{S,T}\cdot b_{S,T}\}.
\]
這直接由 3 推出，因為改用新的生成元 \(v\) 和 \(\bar v\)，也意味着改用 \(v\bar ve_1\)；事實上，\(u=vc\) 和 \(\bar u=\bar v\bar d\) 給出
\[
u\bar ue_1=vc\cdot\bar v\bar de_1=v\bar v\cdot c\bar de_1=v\bar ve_1\cdot c\bar de_1.
\]

\noindent\textbf{5.} 若 \(\{a\}\) 是 \(\{\mathfrak K\}\) 的因子系，而 \(t\) 是 \(\{\mathfrak K\}\) 的指標，則 \(\{a\}^t\) 等於單位類。這是因為 \(\mathfrak G\) 的不可約交叉表示的次數為 \(t\)，正如 Schur 在表示 \(p_{ik}\) 之後所作的證明一樣（§ 24，第 37 頁）。

\noindent\textbf{6.} 群 \(\mathscr K_{\mathfrak Z}\)，即與 \(\mathfrak Z\) 對應的 \(\{\mathfrak K\}\) 所成的群（也就是說，以 \(\mathfrak Z\) 為分裂域），與相伴因子系的類 \(\{a_{S,T}\}\) 所成的群同構。根據 § 27 中的構造，\(\{\mathfrak K\}\) 與 \(\{a_{S,T}\}\) 之間存在一一對應；根據 4，這個對應是同態，因而存在同構。特別地，由此推出單位類彼此對應：

\noindent\emph{交叉積等於 \(\mathsf{P}\)，即等於中心上的矩陣環，當且僅當因子系與單位系相伴。}

\subsection*{§ 29. 極小意義下的主屬定理}\label{hauptgeschlechtssatz-im-minimalen}

設 \(\mathfrak K_r=\mathfrak Z\cdot\mathfrak G\) 為交叉積，\(\mathfrak G^*\) 是所有通過內自同構把 \(\mathfrak Z\) 整體映到自身的元素所成的群（寫成陪集即 \(\mathfrak Z^*u_E,\mathfrak Z^*u_S\) 等）。

\noindent\textbf{定義。} “\emph{極小意義下的主屬}” \(\mathscr S\) 由 \(\mathfrak G^*\) 的所有群自同構組成，這些自同構是 \(\mathfrak Z^*\) 上恆等映射的延拓。這個名稱的理由可由專門化到循環域的情形看出。見 § 30。

\noindent\textbf{定理。} \emph{主屬的每個自同構都具有形式 \(u_S\mapsto u_S\cdot b^Sb^{-1}=u_Sb^{S-1}\)，其中 \(b\) 屬於 \(\mathfrak Z\)；反過來，每個這樣的對應都產生主屬的一個自同構。}

\noindent\emph{證明。} 設在該群自同構下 \(u_S\mapsto v_S\)、\(z\mapsto z\)，後者對 \(\mathfrak Z\) 中每個 \(z\) 成立，並且乘積彼此對應；則通過對應
\[
\sum u_Sc_{\lambda_S}\mapsto\sum v_Sc_{\lambda_S},
\]
可將該自同構擴充為 \(\mathfrak K_r\) 的一個環自同構，其中 \(c_{\lambda_S}\) 屬於 \(\mathfrak Z\)。而眾所周知，每個這樣的自同構都是內自同構。

設 \(b\) 是實現它的元素，於是 \(v_S=bu_Sb^{-1}\)、\(z=bzb^{-1}\)；由於 \(b\) 與所有 \(z\) 可交換，它屬於 \(\mathfrak Z\)，從而
\[
v_S=bu_Sb^{-1}=u_Sb^Sb^{-1}.
\]
反過來，置 \(v_S=u_Sb^Sb^{-1}\)，其中 \(b\) 屬於 \(\mathfrak Z\)，則 \(v_S=bu_Sb^{-1}\)；所以這是 \(\mathfrak K_r\) 的一個自同構，它逐點固定 \(\mathfrak Z\)，因而按定義也將 \(\mathfrak G^*\) 整體映到自身；所以它是主屬的一個自同構。

\subsection*{§ 30. 專門化到循環分裂域}\label{spezialisierung-auf-zyklische-zerfaellungskoerper}

若特別用一個循環域 \(\mathfrak Z\) 構造交叉積，就可得到對若干已知範數定理的解釋。首先有：

\noindent\textbf{定理 1.} \emph{以 \(\mathsf{P}\) 上固定的循環分裂域 \(\mathfrak Z\) 為分裂域的 \(\{\mathfrak K\}\) 所成的群 \(\mathscr K_{\mathfrak Z}\)，同構於群 \(\mathsf{P}^*/N\mathfrak Z^*\)，其中 \(N\mathfrak Z^*\) 是由 \(\mathfrak Z\) 中各元素 \(z\neq0\) 的所有範數 \(Nz\) 組成的群。}

\noindent\emph{證明。} 證明基於交叉表示的規範化以及關於這種規範化的若干輔助引理。

\noindent\textbf{定義。} 循環交叉表示稱為\emph{規範化的}，如果它具有形式
\[
\mathfrak K_r=\mathfrak Z+\mathfrak Zu+\cdots+\mathfrak Zu^{n-1},
\]
並滿足 \(zu=uz^S\)、\(u^n=a\)，其中 \(S\) 是 \(\mathfrak G\) 的一個生成元（也就是說，置 \(u_{S^r}=u^r\)，\(r=0,\ldots,n-1\)，而不是更一般地置 \(u_{S^r}=b_ru^r\)，其中 \(b_r\) 屬於 \(\mathfrak Z\)）。

\noindent\textbf{輔助引理 1.} \emph{循環交叉表示總可以假定為規範化的。此時 \(a\) 屬於 \(\mathsf{P}\)，而 \(P\) 中每個 \(a\) 都給出這樣的表示。}

\noindent\emph{證明。} 第一部分很清楚，因為由 \(u^{-1}zu=z^S\) 推出 \(u^{-r}zu^r=z^{S^r}\)，所以 \(u^r\) 產生置換 \(z\mapsto z^{S^r}\)。又由於 \(u^n\) 誘導恆等置換 \(z\mapsto z^{S^n}\)，所以 \(u^n=a\)，其中 \(a\) 屬於 \(\mathfrak Z\)。但 \(u^r\) 與所有 \(u^r\) 可交換，所以 \(a\) 也如此，因而 \(a^{S^r}=a\)；於是 \(a\) 屬於 \(\mathsf{P}\)。

所以 \(a\) 是唯一實質性的乘法常數。這樣一來，結合條件就沒有內容；\(\mathsf{P}\) 中每個 \(a\) 都給出一個結合系統，從而給出一個 \(\mathfrak K_r\)。

\noindent\textbf{輔助引理 2.} \emph{不必考察相伴因子系的整個類 \(\{a_{S,T}\}\)，只須考察該類所含的規範化因子系，即由規範化表示產生的因子系。類的這個子集稱為規範化類。完整類 \(\{a_{S,T}\}\) 所成的群同構於規範化類所成的群。}

\noindent\emph{證明。} 因為根據輔助引理 1，沒有一個規範化類是空的，所以這個對應是一一的。運算仍是原來的運算，因為它是通過任一代表定義的。

\noindent\textbf{輔助引理 3.} \emph{規範化因子系所成的群同構於乘法群 \(\mathsf{P}^*\)，規範化類所成的群同構於 \(\mathsf{P}^*/N\mathfrak Z^*\)。}

\noindent\emph{證明。} 由輔助引理 1 可知，每個規範化因子系都具有形式 \([1,1,\ldots,a,a,\ldots]\)，因為
\[
u^\nu\cdot u^\mu=u^{\nu+\mu}\quad(\nu+\mu<n),
\qquad
u^\nu\cdot u^\mu=au^\varrho\quad(\tau+\mu=n+\varrho,\ 0\leqq\varrho\leqq n-2).
\]
而對應
\[
[1,\ldots,a,a,\ldots]\mapsto a,
\qquad
[1,1,\ldots,b,b,\ldots]\mapsto b,
\]
\[
[1,\ldots,ab,\ldots]\mapsto ab,
\]
是到 \(\mathsf{P}^*\) 上的群同構，因為仍根據輔助引理 1，\(a\) 跑遍 \(\mathsf{P}^*\) 的所有元素。--- 兩個相伴的規範化因子系通過變換 \(v=uc\)（\(c\) 屬於 \(\mathfrak Z\)）彼此產生。於是
\[
v=uc,\qquad v^2=u^2\cdot c^S\cdot c;
\]
\[
v^r=u^r c^{S^{r-1}}c^{S^{r-2}}\ldots c;
\qquad
v^n=u^nN(c).
\]
所以由 \(u^n=a\) 推出 \(v^n=aN(c)\)。因此，一個規範化類由因子系 \([1,\ldots,aN(c),aN(c),\ldots]\) 組成，其中 \(c\) 跑遍 \(\mathfrak Z\) 中所有元素 \(\neq0\)。輔助引理的第二部分得證。

注意到那裡的 \(\{\mathfrak K\}\) 與類 \(\{a_{S,T}\}\) 之間的同構，以上定理 1 即由輔助引理 2 和 3 推出。

\noindent\textbf{定理 2.} \emph{極小意義下的主屬群同構於滿足 \(N(c)=1\) 的那些 \(\mathfrak Z\)-元素 \(c\) 所成的群。因此，由 \(N(c)=1\) 推出存在 \(\mathfrak Z\) 中的 \(b\neq0\)，使得 \(c=b^S\cdot b^{-1}=b^{S-1}\)（符號意義下的 \(S-1\) 次冪）。}

\noindent\emph{證明。} 在輔助引理 3 的證明中已得到：\(N(c)=1\) 與變換 \(v=uc\) 給出主屬的一個自同構是同一回事；不同的 \(c\) 對應於不同的自同構，反之亦然，而乘積對應於乘積。

\noindent\textbf{注。} 定理 2 說明了“主屬”這個名稱的合理性。按其意義，可以把商群
\[
\mathfrak Z^*/\mathscr H
\quad(\mathscr H=\text{主屬})
\cong\mathfrak Z^*/\mathfrak Z^{*S-1}
\]
稱為極小意義下的屬群；所以它同構於 \(N(\mathfrak Z^*)\)，並且結合定理 1 有
\[
\begin{aligned}
\text{極小意義下的屬群}
&=\mathfrak Z^*/\mathscr H
 \cong\mathfrak Z^*/\mathfrak Z^{*S-1}
 \cong N(\mathfrak Z^*)\\
&=\text{以 }\mathfrak Z\text{ 為分裂域的 }\{\mathfrak K\}\text{ 所成的群}
 =\mathscr K_{\mathfrak Z}
 \cong \mathsf{P}^*/N(\mathfrak Z^*).
\end{aligned}
\]

\subsection*{§ 31. 循環特殊情形的應用}\label{anwendungen-des-zyklischen-spezialfalles}

循環交叉表示為 § 20 中已經證明的下列定理給出新的證明：每個有限域都是交換的，並且除四元數域外，實數域不存在真正的非交換域擴張。--- 這些證明大致可以稱為極小意義下類域論的證明。

\noindent\textbf{1. 每個有限域都是交換的（定理 10，§ 20）。}

證明基於關於交換有限域（伽羅瓦域）的以下兩個已知事實。

\noindent\textbf{1.} \emph{伽羅瓦域在其每個子域上都是循環的。}

\noindent\textbf{2.} \emph{若 \(\mathfrak Z\) 是伽羅瓦域，因而 \(\mathfrak Z\) 在 \(\mathsf{P}\) 上循環，則 \(\mathsf{P}\) 中每個元素都是某個 \(\mathfrak Z\)-元素的範數。\(\mathsf{P}^*=N(\mathfrak Z^*)\)。}

現在設 \(\mathfrak K\) 有限，\(\mathsf{P}\) 是它的中心，\(\mathfrak Z\) 是一個循環分裂域；則根據 § 30 末尾和事實 1，以 \(\mathfrak Z\) 為分裂域的 \(\{\mathfrak K\}\) 所成的群 \(\mathscr K_{\mathfrak Z}\) 同構於 \(\mathsf{P}^*/N\mathfrak Z^*\)，所以根據事實 2，它等於單位類 \(\{\mathsf{P}\}\)。這就是說，不存在不同於 \(\mathsf{P}\) 且以 \(\mathsf{P}\) 為中心的域，故 \(\mathfrak K=\mathsf{P}\)。

\noindent\textbf{2. 實數域 \(\mathsf{P}\) 上唯一真正的非交換域擴張是四元數域（定理 9，§ 20）。}

\noindent\emph{證明。} \(\mathsf{P}\) 必須是中心，因為 \(\mathsf{P}\) 僅有複數域 \(\mathfrak Z\) 這一個交換擴域；若以它為中心，則由於 \(\mathfrak Z\) 是代數閉域，不存在非交換域擴張。同時，也只有 \(\mathfrak Z\) 可能作為分裂域，而 \(\mathfrak Z\) 在 \(\mathsf{P}\) 上是循環的。進一步，每個正數都是範數，所以群 \(\mathsf{P}^*/N\mathfrak Z^*\) 的階為 2；因此除 \(\{\mathsf{P}\}\) 外只有一個類 \(\{\mathfrak K\}\)，這就是四元數域。而且，若在相伴類中取 \(-1\) 作為規範化因子系，同時就得到正規形。因為此時 \(\mathfrak K=\mathfrak Z+\mathfrak Zu\)，且 \(u^2=-1\)。所以 \(\mathfrak K=\mathsf P+\mathsf Pi+(\mathsf P+\mathsf Pi)u\)，其中 \(1^2=-1\)、\(u^2=-1\)。

\noindent\textbf{關於 1 的注。} 已知事實 2 也可以根據 § 30 末尾作如下證明：有 \(\mathscr S=\mathfrak Z^{*S-1}\cong\mathfrak Z^*/\mathsf{P}^*\)，因為所有且僅有 \(\mathsf{P}^*\) 中的元素產生主屬的恆等自同構，或者說在施行 \(S-1\) 後變成單位元。所以當 \(\mathfrak Z\) 有限時，\(\mathscr S\) 的元素數等於 \(\mathsf{P}^*\) 在 \(\mathfrak Z^*\) 中的指數。於是 \(\mathfrak Z^*/\mathscr S\) 的元素數等於 \(\mathsf{P}^*\) 的元素數；因此對 \(N(\mathfrak Z^*)=\mathfrak Z^*/\mathscr S\) 也同樣如此，故 \(\mathsf{P}^*=N(\mathfrak Z^*)\)。（參見“屬的數目 = 歧義類的數目”。不過，這種數論表述處處都須加上“極小意義下”。）

\noindent\textbf{3.} \emph{若 \(\mathsf{P}\) 是一個 \(p\)-進基域，\(\mathfrak Z\) 是 \(P\) 上 \(n\) 次循環域，則以 \(\mathfrak Z\) 為分裂域的 \(\{\mathfrak K\}\) 所成的群 \(\mathscr K_{\mathfrak Z}\) 同構於 \(\mathfrak Z\) 的伽羅瓦群。}

根據 § 30 末尾，這是“小類域論”的一個直接推論；該理論斷言
\[
P^*/N\mathfrak Z^*\cong\mathfrak G.
\]

\noindent\textbf{關於 3 的注。} 由此可知，\(\mathscr K_{\mathfrak Z}\) 中元素的階，即 \(\{\mathfrak K\}\) 的指數，等於其指標。事實上，\(\mathscr K_{\mathfrak Z}\) 的一個生成元的指數 \(=n\)，所以指標是 \(n\) 的因子和該指數的倍數；所以 \(\mathscr K_{\mathfrak Z}\) 只含指標為 \(n\) 的因子的那些類，並且 \(\mathfrak Z\) 是它們的分裂域。在此基礎上，H. Hasse 在 Math. Ann. 104 中證明，\(\mathscr K_{\mathfrak Z}\) 由指標為 \(n\) 的因子的所有類 \(\{\mathfrak K\}\) 組成。所以，\(p\)-進域 \(P\) 上每個 \(n\) 次循環分裂域 \(\mathfrak Z\) 都導出同一個域類 \(\{\mathfrak K\}\) 所成的群 \(\mathscr K_{\mathfrak Z}\)。

以下兩個事實：

\noindent\textbf{1.} \(P\) 上指標為 \(n\) 的因子的所有 \(\{\mathfrak K\}\) 所成的群是 \(n\) 階循環群。

\noindent\textbf{2.} \(\mathsf{P}\) 上每個 \(n\) 次循環分裂域都產生 \(\{\mathfrak K\}\) 的整個群。

可以看作小類域論的逆定理的實質內容。


\clearpage
%
%

%
%
% 原刊：Math. Ann. 97 (1927)，印刷頁561--563。
%

\newpage
\part*{第三部\hspace{1em}Kapferer---Noether 論文}
\addcontentsline{toc}{part}{第三部：Kapferer---Noether 論文}

\begin{center}
\Large\textbf{Noether 代數函數基本定理的}\\[0.3em]
\Large\textbf{必要且充分的重數條件}\\[1em]
\normalsize\textit{H. Kapferer}\\[0.5em]
\small\textit{（附一則與 E. Noether 合作的補充）}\\[1em]
\normalsize \foreign{Math.\ Ann.}\ \textbf{97} (1927), S.~559--567
\end{center}

\vspace{1em}

以下內容給出 Noether 基本定理的一個加強形式，即以必要且充分的重數條件來表述\textsuperscript{1)}。
我按照該定理的理想論理解\textsuperscript{2)}，但在此處及下文不使用理想論的概念和定理，把 Noether 基本定理所包含的事實複合體概括為以下兩個分開的命題\textsuperscript{3)}：

\smallskip
\textbf{定理 Ia.} 對於 \(x,y\) 中兩個互素多項式 \(\varphi,\psi\) 的每一個公共零點 \(x=a,\ y=b\)（係數可為任意複數），存在一個最小的正有理整數 \(\varrho\)，使得模
\((\varphi,\psi,\mathfrak p^{\varrho})\) 表示的多項式全體，恰同模
\((\varphi,\psi,\mathfrak p^{\sigma})\) 表示的多項式全體相同；這裡
\(\mathfrak{p}=(x-\alpha,y-\beta)\)，而 \(\sigma\) 表示任意大於 \(\varrho\) 的有理整數\textsuperscript{4)}。

\smallskip
\textbf{定理 Ib.} 若把這個在概念上唯一確定的多項式全體記為 \(\mathfrak q\)，或者在談到 \(s\) 個不同零點中第 \(\nu\) 個零點 \(x=a_{\nu},\ y=b_{\nu}\) 時記為 \(\mathfrak q_{\nu}\)，則有：

一個給定多項式 \(K\)（在 \(x,y\) 中）具有可除性性質
\(K\equiv0(\varphi,\psi)\) 的必要且充分條件，是同時滿足下列 \(s\) 個同餘式
\[
K \equiv 0(\mathfrak{q}_{1}), \quad K \equiv 0(\mathfrak{q}_{2}), \quad\ldots,\quad
K \equiv 0(\mathfrak{q}_{s}).
\]

從歷史上說，參看 Max Noether 關於此題的奠基論文導言\textsuperscript{5)}，基本定理產生於求解以下基本問題的努力：對於每一個給定的 \(x,y\) 中多項式 \(K\)，應當能夠判定它是否可被 \((\varphi,\psi)\) 整除；也就是說，還要判定是否存在 \(x,y\) 中兩個另外的多項式 \(A\) 和 \(B\)，使得方程
\[
K = A\varphi + B\psi
\]
成為關於 \(x,y\) 的恆等式。

Noether 定理並不解決這個問題；毋寧說，它只是把所提出的問題作一個雖然基本的歸約\textsuperscript{6)}，歸約到一個較簡單的問題，即問各個同餘式 \(K\equiv0(\mathfrak q)\) 成立的必要且充分條件是什麼。那麼，何時有 \(K\equiv0(\mathfrak q)\)？只有這個後一問題得到回答，M. Noether 原先的問題才算解決。本文的目標正是如此，並且實際上達到了這個目標：給出有限多個、實踐上也不難檢驗的重數條件，它們對於
\(K\equiv0(\mathfrak q)\) 既必要又充分。一個已知的、文獻中常被使用和引用的充分條件是下面這種形式；但它並不是必要條件：

\smallskip
\textit{為了使同餘式 \(K\equiv0(\varphi,\psi)\) 成立，只要 \(\varphi=0,\ \psi=0\) 的每一個交點都以“足夠”高的重數出現在多項式 \(K\) 中即可。}

\smallskip
這個命題作為特殊情形包含在定理 I 中；因為 \(K\) 中的 \(\varrho\) 重點等價於
\(K\equiv0(\mathfrak p^{\varrho})\)，於是更有
\(K\equiv0(\varphi,\psi,\mathfrak p^{\varrho})\)，也就是說還有
\(K\equiv0(\mathfrak q)\)。但是反命題並不成立。

\vspace{0.5em}
\noindent\rule{\textwidth}{0.4pt}
\vspace{0.3em}

{\footnotesize
\textsuperscript{1)} 詳細敘述可參看我同題的論文，該文載於 A.\ Loewy（Freiburg i.\ Br.）發起的 1927 年 Heidelberg Akademie-Berichte 專號 ``Beiträge zur Algebra''，其中收有多位作者的論文。

\textsuperscript{2)} 例如參看 B.\ L.\ van der Waerden, ``Zur Nullstellentheorie der Polynomideale''. \foreign{Math.\ Annalen} 96 (1926), S.~203.

\textsuperscript{3)} 在所引論文的 \S~1 中，我分別且直接地證明了這兩個命題，即沒有使用理想論概念。

\textsuperscript{4)} 選擇字母 \(\mathfrak p\) 和 \(\mathfrak q\)，是為了提醒讀者：這裡所指的對象在內容上同近代文獻中慣用的素理想 \(\mathfrak p\) 與準素理想 \(\mathfrak q\) 概念相一致。

\textsuperscript{5)} Max Noether, ``Über einen Satz aus der Theorie der algebraischen Funktionen'', \foreign{Math.\ Annalen} 6 (1872).

\textsuperscript{6)} 所說的歸約，大致可同初等數論中人人熟悉的歸約相比：自然數 \(A\) 被另一個自然數 \(B\) 整除這一性質，等價於 \(A\) 分別被
\(p_1^{a_1}\), \(p_2^{a_2}\), \ldots, \(p_s^{a_s}\) 整除；這裡
\(p_1,p_2,\ldots,p_s\) 表示
\(B=p_1^{a_1}p_2^{a_2}\ldots p_s^{a_s}\) 的不同素因子。因此，正是“準素理想”這個概念（正文中記為 \(\mathfrak q\)）應看作素數冪的對應物。[參看 E.\ Noether, ``Idealtheorie in Ringbereichen'', \foreign{Math.\ Annalen} 83 (1921), 導言。]
}

\newpage
\section*{主定理}

對所提出問題的解答可以簡短地表述如下：對於單個 \(\mathfrak q\)，可以配給一個由 \(t\) 個正有理整數組成的系統
\((g_1),(g_2),\ldots,(g_t)\)；同時，對於模 \((K,\mathfrak q)\)，可以配給同樣數目的、可正可負的有理整數
\((K_1),(K_2),\ldots,(K_t)\)，於是有下列\textbf{主定理}：

\begin{satz}[主定理]\label{satz:kap-II}
\(K\equiv0(\mathfrak q)\) 的必要且充分條件，是同時滿足 \(t\) 個條件
\[
(K_1)\geqq(g_1),\quad (K_2)\geqq(g_2),\quad \ldots,\quad (K_t)\geqq(g_t).
\]
其中有
\[
(g_1)+(g_2)+\ldots+(g_t)=\mu,
\]
這裡 \(\mu\) 表示與 \(\mathfrak q\) 相應並由結式定義的交點重數。
\end{satz}

證明時，不妨假定 \(x=0,\ y=0\) 是多項式對 \((\varphi,\psi)\) 中與
\(\mathfrak q\) 相應的零點。進一步，一般約定：若 \(A(x,y)\) 是
\(x,y\) 中任一多項式\textsuperscript{7)}，則符號 \((A)\) 表示
\(y\) 作為因子在 \(A(0,y)\) 中出現的次數；若 \(A(0,y)\) 恆等為零，
則令 \((A)=\infty\)。與 \(\mathfrak q\) 相配的數
\((g_1),(g_2),\ldots,(g_t)\) 由下列方程組定義；這個方程組構成該方法的基礎：

\begin{equation}\tag{1}\label{eq:kap-1}
\begin{aligned}
f_1\cdot g_1(0,y):y^{(g_1)}
 - g_1\cdot f_1(0,y):y^{(g_1)} &= x\cdot h_1,\\
f_2\cdot g_2(0,y):y^{(g_2)}
 - g_2\cdot f_2(0,y):y^{(g_2)} &= x\cdot h_2,\\
&\vdots \qquad\qquad\vdots\qquad\qquad\vdots\\
f_r\cdot g_r(0,y):y^{(g_r)}
 - g_r\cdot f_r(0,y):y^{(g_r)} &= x\cdot h_r.
\end{aligned}
\end{equation}

這個方程組完全建立在給定的多項式對 \(\varphi,\psi\) 之上；
也就是說，多項式 \(f_1,g_1\) 分別被定義為同 \(\varphi,\psi\) 恆等，
但必要時可交換兩者的次序；次序須取為 \((f_1)\geqq(g_1)\)。
由於這一規定，並且只有在這一規定下，\(h_1\) 才是整有理的。

多項式 \(f_2,g_2\) 分別被定義為同 \(h_1,g_1\) 恆等，並附加條件
\((f_2)\geqq(g_2)\)。一般地，多項式對 \(f_r,g_r\) 分別被定義為同
\(h_{r-1},g_{r-1}\) 恆等，並附加條件 \((f_r)\geqq(g_r)\)。

方程組 \eqref{eq:kap-1} 固然可以按所說明的方式無限繼續下去；
但是主定理只要求形成該系統的前 \(t\) 個方程，其中 \(t\) 由以下事實確定：

\smallskip
\textbf{輔助定理 I.} 存在一個自然數 \(t\) [並且 \(t\leqq\mu\)，
這裡 \(\mu\) 表示點 \(x=0,\ y=0\) 在 \(\varphi=0,\ \psi=0\)
中出現的交點重數]，使得由 \eqref{eq:kap-1} 定義的前 \(t\) 個數，
即 \((g_1),(g_2),\ldots,(g_t)\)，全都大於零，而同時
\((g_{t+1})\) 以及其後所有數都等於零。

事實上，考慮結式的乘積定理，點 \(x=0,\ y=0\) 作為
\(f_r=0,\ g_r=0\) 的交點，其重數等於
\[
\mu-(g_1)-(g_2)-\ldots-(g_r),
\]
由此推出至多經過 \(\mu\) 步即會中止。同時得到事實：
\begin{equation}\tag{2}\label{eq:kap-2}
\mu=(g_1)+(g_2)+\ldots+(g_t),
\end{equation}
這個公式之所以值得注意，是因為只要有關零點已知，它就在完全一般的情形下
給出一種有理且實踐上易於執行的確定交點重數的方法。

既知數 \((g_i)\) 之後，數 \((K_i)\) 由下列與
\eqref{eq:kap-1} 相對應的方程組確定：
\begin{equation}\tag{3}\label{eq:kap-3}
\begin{aligned}
K_1\cdot g_1(0,y):y^{(g_1)}
 - g_1\cdot K_1(0,y):y^{(g_1)} &= x\cdot K_1,\\
K_2\cdot g_2(0,y):y^{(g_2)}
 - g_2\cdot K_2(0,y):y^{(g_2)} &= x\cdot K_2,\\
&\vdots \qquad\qquad\vdots\qquad\qquad\vdots\\
K_r\cdot g_r(0,y):y^{(g_r)}
 - g_r\cdot K_r(0,y):y^{(g_r)} &= x\cdot K_{r+1}.
\end{aligned}
\end{equation}

\(K_1\) 被定義為同 \(K\) 恆等。顯然，\(K_2\) 當且僅當
\((K_1)\geqq(g_1)\) 時才是整有理的。若後一條件成立，則進一步有：
\(K_3\) 當且僅當 \((K_2)\geqq(g_2)\) 時才是整有理的。
一般地，若 \(K_r\) 已經是整有理的，則 \(K_{r+1}\) 當且僅當
\((K_r)\geqq(g_r)\) 時才是整有理的。

主定理（定理 II）的證明現在依賴於下面的另一個輔助定理：

\smallskip
\textbf{輔助定理 II.} 若令
\((f_i,g_i,\mathfrak p^\varrho)=\mathfrak q^{(i)}\)，
則對於整數列 \(1,2,\ldots\) 中的每個 \(i\)，
\textit{in inf.} 有下述命題：

同餘式
\[
K_i\equiv0(\mathfrak q^{(i)})
\]
成立的必要且充分條件，是同時滿足下列兩個條件：
\[
(K_i)\geqq(g_i)
\qquad \text{以及} \qquad
K_{i+1}\equiv0(\mathfrak q^{(i+1)}).
\]

輔助定理的證明\textsuperscript{8)}由組合方程組
\eqref{eq:kap-1} 與 \eqref{eq:kap-3} 得到。注意到按輔助定理 I，
\(g_{t+1}\) 在零點處不再消失，也就是說，模
\(\mathfrak q^{(t+1)}\) 由全體多項式組成（即為單位理想），
因而 \(K_{t+1}\equiv0(\mathfrak q^{(t+1)})\) 成為平凡命題；
於是把輔助定理 II 應用 \(t\) 次，便得到主定理的證明。

\smallskip
不作證明（證明見我上面所引的論文），我再給出主定理的兩個補充：

\begin{zusatz}[ I]
主定理中的數 \(t\)，等同於具有
\(x^{v}\equiv0(\mathfrak q)\) 這一性質的最小\textsuperscript{9)}正數 \(\nu\)。
\end{zusatz}

\begin{zusatz}[ II]
方程組 \eqref{eq:kap-1} 與 \eqref{eq:kap-3}，以及由交換 \(x\) 和 \(y\) 從它們得到的方程組，在適當組合時，給出一種有理方法，可精確地數值確定“屬於” \(\mathfrak q\) 的指數
\(\varrho\)，即定理 Ia 中的那個最小指數 \(\varrho\)。
\end{zusatz}

E.\ Bertini\textsuperscript{10)} 曾給出 \(\varrho\) 的一般公式。但是該公式並不總是給出 \(\varrho\) 的最小值，這一點在所引論文中通過例子
\(\varphi=x^3+y^4,\ \psi=x^2+y^3\) 得到說明。最後還要指出，Bertini 公式本身也可以通過我們的主定理以一種新的方式加以論證。

\vspace{0.5em}
\noindent\rule{\textwidth}{0.4pt}
\vspace{0.3em}

{\footnotesize
\textsuperscript{7)} 若對 \(A(x,y)\) 也允許分式有理函數 \(C(x,y):D(x,y)\)，則
\((A)=(C)-(D)\)；在後一種情形中，\((A)\) 因而也可取負值。

\textsuperscript{8)} 也可參看“補充”的備註 I 中的證明。

\textsuperscript{9)} 因而這裡得到的是精確指數，不同於 Grete Hermann 小姐在 ``Die Frage der endlich vielen Schritte in der Theorie der Polynomideale''（使用 K.\ Hentzelt 遺留定理）一文中給出的估計；該文見 \foreign{Math.\ Annalen} 95 (1926)，在那裡，雖處理的問題一般得多，卻只給出上界。

\textsuperscript{10)} E.\ Bertini, \foreign{Math.\ Annalen} 34 (1889).
}

\newpage
\section*{定理 III 與補充}

主定理的一個有趣特殊化來自以下限制性假設：相關
\(\mathfrak q=(\varphi,\psi,\mathfrak p^{\varrho})\) 的零點，至少在兩條曲線
\(\varphi=0,\ \psi=0\) 中的一條上是單點。若例如它在 \(\psi\) 上是單點，則可以證明，主定理中的數 \(t\) 正好等於 \(\mu\)，並且主定理中的 \(t\) 個不等式可由一個唯一的重數條件代替。這導致下列一般定理：

\begin{satz}[定理 III]\label{satz:kap-III}
若 \(\varphi=0,\ \psi=0\) 的所有交點在 \(\psi\) 上都是單點，則在假定 \(K\) 與 \(\psi\) 互素時，同餘式 \(K\equiv0(\varphi,\psi)\) 成立的必要且充分條件是：\(K\) 在所有這些點處消失，而且消失的方式使得在每一點上，曲線對 \(K=0,\ \psi=0\) 的交點重數至少同曲線對 \(\varphi=0,\ \psi=0\) 的交點重數一樣大。
\end{satz}

定理 III 還因一個特別原因而值得注意：它所要求的重數條件，其滿足與否可以通過有限次微分來判定，而不依賴於基於 \eqref{eq:kap-1} 的公式 \eqref{eq:kap-2}。如果我引用下列我在 1923 年提出的“定理”\textsuperscript{11)}，這一點立刻清楚：

\smallskip
\textit{若 \(A\) 和 \(B\) 是 \(x,y\) 中任意兩個互素多項式，\(P\) 是一個公共點，並且是 \(B\) 的單點，則這個點作為交點的重數，當且僅當同一點在曲線對 \(A'=0,\ B=0\) 中的重數恰為 \(\mu-1\) 時，等於數 \(\mu\)；這裡}
\[
A' = \frac{\partial A}{\partial x} \cdot \frac{\partial B}{\partial y} -
\frac{\partial A}{\partial y} \cdot \frac{\partial B}{\partial x}.
\]

最後還要提到，主定理經過小的修改後，也可以轉移到三個變量 \(x,y,z\) 中的齊次多項式上。在上面的特殊情形中，這種轉移做得格外優雅。也就是說，可以證明：

\begin{zusatz}[（對定理 III 的補充）]
定理 III 完全不變，甚至在文字表述上也不變，只要把 \(\varphi,\psi,K\) 理解為三個變量中的齊次多項式\textsuperscript{12)}。
\end{zusatz}

\vspace{0.5em}
\noindent\rule{\textwidth}{0.4pt}
\vspace{0.3em}

{\footnotesize
\textsuperscript{11)} H.\ Kapferer, ``Über die Multiplizität der Schnittpunkte von zwei algebraischen Kurven''. \foreign{Jahresber.\ der D.\ M.-V.} 32 (1923).

\textsuperscript{12)} 關於三個變量中齊次同餘
\(K(x,y,z)\equiv0(\varphi(x,y,z),\psi(x,y,z))\) 成立條件的問題，雖在高度特殊化的假設下，也是 A.\ Ostrowski ``Über die Maxwellsche Erzeugung der Kugelfunktionen'', \foreign{Jahresb.\ der D.M.-V.} 33 (1925) 一文的核心問題。該文假設：\(\varphi\) 分解為全體一次因子；\(\varphi\) 和 \(K\) 同階；\(\psi=x^2+y^2+z^2\)。由於這裡 \(\psi\) 完全無奇點，Ostrowski 筆記中的主定理可以看作我們定理 III 的特殊情形。
}

\newpage
\section*{補充，與 E. Noether 合作\textsuperscript{13)}}

方程組 \eqref{eq:kap-1} 同時給出一種有理方法，用來確定模 \(\mathfrak q\) 的一組完全剩餘系；所謂完全剩餘系，是指相對於係數範圍 \(\mathsf{P}\)，模 \(\mathfrak q\) 線性無關的有限多個剩餘類的一組完整代表元。這樣，主定理證明中的各個步驟變得十分透明；同時還得到一個副結果：結式重數 \(\mu\) 同理想 \(\mathfrak q\) 的長度相一致；

\vspace{0.5em}
\noindent\rule{\textwidth}{0.4pt}
\vspace{0.3em}

{\footnotesize
\textsuperscript{13)} 理想論解釋出自 E.\ Noether，而 (5) 與 (6) 的基本證明出自 H.\ Kapferer。

\textsuperscript{14)} 不妨把域 \(P\) 假定為代數閉的；在這個範圍內，除微分條件外，前述全部定理實際都成立。\(\mathfrak q\) 的長度也定義為由理想組成、從 \(\mathfrak p\) 走向 \(\mathfrak q\) 的一條合成列的長度。參看 E.\ Noether, \foreign{Jahresber.\ d.\ d.\ Math.-Ver.} 34 (1925), S.~101（斜體頁碼）；更詳細見 H.\ Grell, \foreign{Math.\ Annalen} 97 (1927), S.~529。

文獻中已有的、關於這兩個重數（對 \(n\) 個變量）相一致的證明都很複雜；最簡單的是 Hentzelt 尚未發表的遺稿證明。
}

\vspace{1em}
也就是說，同剩餘類環的秩相一致，或同相對於 \(P\) 線性無關的剩餘類數目相一致\textsuperscript{14)}。

用 \eqref{eq:kap-1} 與輔助定理 II 的記號，完全剩餘系由下列定理刻畫：

\begin{satz}[定理 IV]\label{satz:kap-IV}
模 \(\mathfrak q\) 的一組完全剩餘系由下列多項式給出：
\[
h_t,\ h_ty,\ldots,h_ty^{(g_t)-1};\ \ldots;\
h_i,\ h_iy,\ldots,h_iy^{(g_i)-1};\ \ldots;\
h_1,\ h_1y,\ldots,h_1y^{(g_1)-1}.
\]
其中，對每個 \(i\)，\(h_i,h_iy,\ldots,h_iy^{(g_i)-1}\) 給出模
\(\mathfrak q^{(i)}\) 的 \(\mathfrak q^{(i+1)}\) 的一組完全剩餘系。
\end{satz}

由定義立即有 \(\mathfrak q^{(i+1)}=(\mathfrak q^{(i)},h_i)\)，因而
\(\mathfrak q^{(i)}\equiv0(\mathfrak q^{(i+1)})\)。因此，只要證明
\begin{equation}\tag{4}\label{eq:kap-4}
h_i x\equiv0(\mathfrak q^{(i)});\qquad
h_i y^{(g_i)}\equiv0(\mathfrak q^{(i)})
\end{equation}
以及
\begin{equation}\tag{5}\label{eq:kap-5}
\text{由 }h_iF(y)\equiv0(\mathfrak q^{(i)})
\text{ 可推出 }F(y)\equiv0(y^{(g_i)}),
\end{equation}
定理 IV 即得證。

事實上，關係式 \eqref{eq:kap-4} 表明，
\(h_i,h_iy,\ldots,h_iy^{(g_i)-1}\) 的線性組合窮盡模
\(\mathfrak q^{(i)}\) 的 \(\mathfrak q^{(i+1)}\) 的剩餘類；而
\eqref{eq:kap-5} 表示線性無關性。由於
\(\mathfrak q^{(t+1)}\) 成為單位理想——這一點現在由
\(\mathfrak q\) 的有限長度推出——由此便得到定理 IV 的證明，
同時還得到 \(\mathfrak q\) 的長度等於
\((g_1)+\ldots+(g_t)\)，因而按 \eqref{eq:kap-2} 同結式重數一致。

關係式 \eqref{eq:kap-4} 的第一式可從 \eqref{eq:kap-1} 讀出；
第二式如下得到：

有 \(h_ig_i\equiv0(\mathfrak q^{(i)})\)；因此由
\eqref{eq:kap-4} 的第一式，也有
\(h_ig_i(0,y)\equiv0(\mathfrak q^{(i)})\)，或者
\[
h_iy^{(g_i)}\cdot(g_i(0,y):y^{(g_i)})
\equiv0(\mathfrak q^{(i)}).
\]
由於 \(g_i(0,y):y^{(g_i)}\not\equiv0(y)\)，由此推出
\(h_iy^{(g_i)}\equiv0(\mathfrak q^{(i)})\)。

在證明 \eqref{eq:kap-5} 之前，先作一個可由 \eqref{eq:kap-1}
立即讀出的說明：\(f_i\) 與 \(g_i\) 的公共因子只能是
\(y\) 的一個不被 \(y\) 整除的多項式 \(d(y)\)：
\[
f_i=d_i(y)\bar f_i,\qquad g_i=d_i(y)\bar g_i,
\]
並且 \(\bar f_i\) 與 \(\bar g_i\) 互素。由於
\(d(y)\not\equiv0(\mathfrak{p})\)，\(\mathfrak q^{(i)}\) 也等於
\((\bar f_i,\bar g_i,\mathfrak p^\varrho)\)；並且，作為
\(\mathfrak q\) 的因子，對每個 \(\sigma\geqq\varrho\)，它都等於
\((\bar f_i,\bar g_i,\mathfrak p^\sigma)\)。因此
\(\mathfrak q^{(i)}\) 是 \((\bar f_i,\bar g_i)\) 中屬於原點的準素分量。

設 \(\gamma_1,\delta_1;\ldots;\gamma_r,\delta_r\) 是
\((\bar f_i,\bar g_i)\) 中不同於 \(x=0,\ y=0\) 的零點，並設
\(\tau_1,\ldots,\tau_r\) 是 \((\bar f_i,\bar g_i)\) 中相應準素分量的指數；
再在乘積
\[
P=((x-\gamma_1)+u(y-\delta_1))^{\tau_1}\ldots
((x-\gamma_r)+u(y-\delta_r))^{\tau_r}
\]
中把不定元 \(u\) 專門化為係數範圍 \(\mathsf{P}\) 的某個元素，使得仍有
\(P\not\equiv0(\mathfrak p)\)。令 \(G(x,y)\) 等於經此專門化的
\(P\) 同 \(d(y)\) 的乘積；於是
\(G(x,y)\not\equiv0(\mathfrak p)\)，並且
\(\mathfrak q^{(i)}\cdot G(x,y)\equiv0(f_i,g_i)\)。

由假設 \(h_iF(y)\equiv0(\mathfrak q^{(i)})\) 推出
\(h_iF(y)G(x,y)\equiv0(f_i,g_i)\)；將它同 \eqref{eq:kap-1}
合起來，以恆等式形式寫為
\[
h_iF(y)G=A f_i+B g_i;\qquad
h_ix=(g_i(0,y):y^{(g_i)})\cdot f_i
 -(f(0,y):y^{(g_i)})g_i.
\]
由此得到
\[
\bar f_i(xA-F(y)H)\equiv0(\bar g_i);\qquad
H=G(x,y)\cdot g_i(0,y):y^{(g_i)}\not\equiv0(\mathfrak p).
\]

由 \(\bar f_i\) 與 \(\bar g_i\) 互素，進一步推出
\[
xA-F(y)H\equiv0(\bar g_i),
\]
因而
\[
F(y)H\equiv0(\bar g_i,x)
\equiv0(\bar g_i,x,\mathfrak p^\sigma)=(y^{(g_i)},x),
\]
從而 \(F(y)\equiv0(y^{(g_i)})\)。這樣便證明了
\eqref{eq:kap-5}，因而也證明了定理 IV。

\smallskip
\textbf{備註 I.} 關係式 \eqref{eq:kap-4}、\eqref{eq:kap-5} 與
輔助定理 II 合在一起，可以寫成理想商\textsuperscript{15)}之間
相互倒置的關係；其中第二條關係的證明同時給出輔助定理 II 的證明：
\begin{equation}\tag{6}\label{eq:kap-6}
\mathfrak q^{(i)}:\mathfrak q^{(i+1)}
 =(\mathfrak q^{(i)},x)=(y^{(g_i)},x);\qquad
\mathfrak q^{(i)}:(\mathfrak q^{(i)},x)
 =\mathfrak q^{(i+1)}.\qquad\text{\textsuperscript{16)}}
\end{equation}

第一條關係是 \eqref{eq:kap-4} 與 \eqref{eq:kap-5} 的合併；
關於第二條需要說明：在 \eqref{eq:kap-4} 中已經證明
\(x\cdot\mathfrak q^{(i+1)}\equiv0(\mathfrak q^{(i)})\)；
反過來可類似於 \eqref{eq:kap-5} 得到。由
\(xM\equiv0(\mathfrak q^{(i)})\)，如上可推出
\[
xMG\equiv0(f_i,g_i);\quad\text{所以}\quad
xMG\cdot g_i(0,y):y^{(g_i)}\equiv0(xh_i,g_i).
\]

由於按定義（即 \((f_i)\geqq(g_i)\)）有
\(g_i\not\equiv0(x)\)，由此得到
\[
M\cdot(G\cdot g_i(0,y):y^{(g_i)})
\equiv0(h_i,g_i)\equiv0(\mathfrak q^{(i+1)}),
\quad\text{也就是}\quad M\equiv0(\mathfrak q^{(i+1)}).
\]

輔助定理 II 中條件的必要性現在可這樣推出：
\(K_i\equiv0(\mathfrak q_i)\) 給出
\(K_i(0,y)\equiv0(y^{(g_i)},x)\)；因此
\((K_i)\geqq(g_i)\)，並由 \eqref{eq:kap-3} 得到
\[
xK_{i+1}\equiv0(\mathfrak q^{(i)});\quad
\text{於是由 \eqref{eq:kap-6} 的第二條關系}\quad
K_{i+1}\equiv0(\mathfrak q^{(i+1)}).
\]

充分性可直接讀出；而對輔助定理 II 作有限次迭代，就得到主定理。
同樣，對 \eqref{eq:kap-6} 的第二條關係迭代，得到補充 I，即
\(\mathfrak q:x^t=\mathfrak q^{(t+1)}\)，所以等於單位理想；並且
\(t\) 是這樣的最低冪，因為
\(\mathfrak q:x^{t-1}=\mathfrak q^{(t)}\)。

\textbf{備註 II.} 對於一個不被 \(\mathfrak q\) 整除的多項式
\(K\)，還可斷言：若
\(K\equiv0(\mathfrak q^{(i+1)})\)，且
\(K\not\equiv0(\mathfrak q^{(i)})\)，則
\[
(K,\mathfrak q^{(i)})=(h_iy^{(K)},\mathfrak q^{(i)}).
\]
因為按定理 IV，
\(K\equiv h_iy^\lambda c(y)\pmod{\mathfrak q^{(i)}}\)，其中
\(c(y)\not\equiv0(y)\)；因此
\[
(K,\mathfrak q^{(i)})=(h_iy^\lambda,\mathfrak q^{(i)}).
\]

這裡 \(y^{(g_i)-\lambda}K\equiv0(\mathfrak q^{(i)})\)，並且按
\eqref{eq:kap-5}，這正是這樣的最低冪。按輔助定理 II，這個最低冪為
\((g_i)-(K)\)；所以 \(\lambda=(K)\)。

\vfill
\begin{flushright}
（1926 年 10 月 24 日收稿。）
\end{flushright}

\vspace{0.5em}
\noindent\rule{\textwidth}{0.4pt}
\vspace{0.3em}

{\footnotesize
\textsuperscript{15)} 所謂商 \(\mathfrak a:\mathfrak b\)，眾所周知，是指全體多項式 \(c(x,y)\)，使得
\(\mathfrak b c\equiv0(\mathfrak a)\)。

\textsuperscript{16)} 在這兩個互為倒置的關係中，按不可約理想的一般理論，其中一個蘊涵另一個；參看 Macaulay, \textit{Modular Systems}, Nr.\,72 和 73. \foreign{Cambridge Tracts} 19 (1916).
}

%
\clearpage
%

%
%
%
%

\newpage
\part*{第四部\hspace{1em}書目}
\addcontentsline{toc}{part}{第四部：書目}

\section*{書目}
\addcontentsline{toc}{section}{書目}

\begin{enumerate}[label={\small\arabic*.},leftmargin=2em,itemsep=0.3em]
\item 關於三元雙二次型的型系統的形成。\\
\foreign{Sitz.\ Ber.\ d.\ Physikal.-mediz.\ Sozietät in Erlangen} \textbf{39} (1907), S.~176--179

\item 關於三元雙二次型的型系統的形成。\\
\foreign{J.\ Reine Angew.\ Math.}\ \textbf{134} (1908), S.~23--90 及兩張表

\item 關於 \(n\) 個變量的型的不變式理論。\\
\foreign{J.\ Ber.\ d.\ DMV} \textbf{19} (1910), S.~101--104

\item 關於 \(n\) 個變量的型的不變式理論。\\
\foreign{J.\ Reine Angew.\ Math.}\ \textbf{139} (1911), S.~118--154

\item 有理函數域。\\
\foreign{J.\ Ber.\ d.\ DMV} \textbf{22} (1913), S.~316--319

\item 域與有理函數系統。\\
\foreign{Math.\ Ann.}\ \textbf{76} (1915), S.~161--196

\item 有限群不變式的有限性定理。\\
\foreign{Math.\ Ann.}\ \textbf{77} (1916), S.~89--92

\item 關於任意多個基本型的一個系統的不變式的整有理表示。\\
\foreign{Math.\ Ann.}\ \textbf{77} (1916), S.~93--102。（亦見第~16 項）

\item 由整超越數組成的最一般範圍。\\
\foreign{Math.\ Ann.}\ \textbf{77} (1916), S.~103--128。（亦見第~16 項）

\item 同構映射的函數方程。\\
\foreign{Math.\ Ann.}\ \textbf{77} (1916), S.~536--545

\item 具有指定群的方程。\\
\foreign{Math.\ Ann.}\ \textbf{78} (1918), S.~221--229。（亦見第~16 項）

\item 任意微分表達式的不變式。\\
\foreign{Nachr.\ v.\ d.\ Ges.\ d.\ Wiss.\ zu Göttingen} \textbf{1918}, S.~37--44

\item 不變變分問題。\\
\foreign{Nachr.\ v.\ d.\ Ges.\ d.\ Wiss.\ zu Göttingen} \textbf{1918}, S.~235--257

\item 單變量代數函數的算術理論及其同其他理論和數域理論的關係。\\
\foreign{J.\ Ber.\ d.\ DMV} \textbf{28} (1919), S.~182--203

\item 二元型的整係數不變式系統的有限性。\\
\foreign{Nachr.\ v.\ d.\ Ges.\ d.\ Wiss.\ zu Göttingen} \textbf{1919}, S.~138--156

\item 關於型論中的級數展開。\\
\foreign{Math.\ Ann.}\ \textbf{81} (1920), S.~25--30

\item 與 W.\ Schmeidler 合作：非交換範圍中的模，特別是由微分表達式和差分表達式構成的模。\\
\foreign{Math.\ Z.}\ \textbf{8} (1920), S.~1--35

\item 關於陣亡的 K.\ Hentzelt 在消元理論方面的一項工作。\\
\foreign{J.\ Ber.\ d.\ DMV} \textbf{30} (1921), S.~101

\item 環範圍中的理想論。\\
\foreign{Math.\ Ann.}\ \textbf{83} (1921), S.~24--66

\item 絕對不可約性的一個代數判別準則。\\
\foreign{Math.\ Ann.}\ \textbf{85} (1922), S.~26--33

\item 形式變分計算與微分不變式。\\
\foreign{Encyklopädie d.\ math.\ Wiss.}\ III,\ 3 (1922), S.~68--71（載於：
R.\ Weitzenböck, \textit{Differentialinvarianten}）

\item 整理 K.\ Hentzelt：關於多項式理想與結式的理論。\\
\foreign{Math.\ Ann.}\ \textbf{88} (1923), S.~53--79

\item 代數不變量與微分不變量。\\
\foreign{J.\ Ber.\ d.\ DMV} \textbf{32} (1923), S.~177--184

\item 消元理論與一般理想論。\\
\foreign{Math.\ Ann.}\ \textbf{90} (1923), S.~229--261

\item 消元理論與理想論。\\
\foreign{J.\ Ber.\ d.\ DMV} \textbf{33} (1924), S.~116--120

\item 代數數域中理想論的抽象構造。\\
\foreign{J.\ Ber.\ d.\ DMV} \textbf{33} (1924), S.~102

\item 理想論中的 Hilbert 函數。\\
\foreign{J.\ Ber.\ d.\ DMV} \textbf{34} (1925), S.~101

\item 群特徵標與理想論。\\
\foreign{J.\ Ber.\ d.\ DMV} \textbf{34} (1925), S.~144

\item 特徵 \(p\) 的有限線性群不變式的有限性定理。\\
\foreign{Nachr.\ v.\ d.\ Ges.\ d.\ Wiss.\ zu Göttingen} \textbf{1926}, S.~28--35

\item 代數數域和函數域中理想論的抽象構造。\\
\foreign{Math.\ Ann.}\ \textbf{96} (1927), S.~26--61

\item 代數數域或函數域的階的判別式定理。\\
\foreign{J.\ Reine Angew.\ Math.}\ \textbf{157} (1927), S.~82--104

\item 與 R.\ Brauer 合作：關於不可約表示的極小分裂域。\\
\foreign{Sitz.\ Ber.\ d.\ Preuß.\ Akad.\ d.\ Wiss.}\ \textbf{1927}, S.~221--228

\item 算術觀點下的超複量與表示論。\\
\foreign{Atti Congresso Bologna} \textbf{2} (1928), S.~71--73

\item 超複量與表示論。\\
\foreign{Math.\ Z.}\ \textbf{30} (1929), S.~641--692

\item 關於由整值函數組成的極大範圍。\\
\foreign{Rec.\ Soc.\ Math.\ Moscou} \textbf{36} (1929), S.~65--72

\item 理想微分與不同。\\
\foreign{J.\ Ber.\ d.\ DMV} \textbf{39} (1930), S.~17

\item 無高階分歧的域中的正規基。\\
\foreign{J.\ Reine Angew.\ Math.}\ \textbf{167} (1932), S.~147--152

\item 與 R.\ Brauer 和 H.\ Hasse 合作：代數理論中一個主定理的證明。\\
\foreign{J.\ Reine Angew.\ Math.}\ \textbf{167} (1932), S.~399--404

\item 超複系統同交換代數及數論的關係。\\
\foreign{Verhandl.\ Intern.\ Math.-Kongreß Zürich} \textbf{1} (1932), S.~189--194

\item 非交換代數。\\
\foreign{Math.\ Z.}\ \textbf{37} (1933), S.~514--541

\item 相對 Galois 數域的主屬定理。\\
\foreign{Math.\ Ann.}\ \textbf{108} (1933), S.~411--419

\item 分裂交叉積及其極大階。\\
\foreign{Actualités scientifiques et industrielles} \textbf{148} (1934) (15 S.)

\item 理想微分與不同。\\
\foreign{J.\ Reine Angew.\ Math.}\ \textbf{188} (1950), S.~1--21
\end{enumerate}

\noindent\begin{minipage}{\textwidth}
\vspace{1em}
\noindent\rule{\textwidth}{0.4pt}
\vspace{0.5em}

\noindent\textbf{Noether 代數函數基本定理的必要且充分的重數條件：}\\
\textbf{H.\ Kapferer}\\
\textbf{（附一則與 E.\ Noether 合作的補充）。}\\
\foreign{Math.\ Ann.}\ \textbf{97} (1927), S.~559--567
\end{minipage}

\section*{Emmy Noether 的短訊和書評列表}
\addcontentsline{toc}{section}{短訊和書評}

\subsection*{發表於 \foreign{Jahresberichte der Deutschen Mathematiker-Vereinigung} 的短訊}

\begin{itemize}[leftmargin=1.5em,itemsep=0.2em]
\item 任意微分表達式的不變式，27, Teil II, (1918), S.~28
\item 不變變分問題，27, Teil II, (1918), S.~47
\item 整係數二元不變式的有限性，28, Teil II, (1919), S.~29
\item 關於整係數多項式和冪級數，28, Teil II, (1919), S.~29--30
\item 模論的分解定理，29, Teil II, (1920), S.~54
\item 初等因子與一般理想論，30, Teil II, (1921), S.~32
\item 關於 Lipschitz 一個公式的提示，30, Teil II, (1921), S.~48
\item 一般理想論中 Hilbert 定理的一個類似物，32, Teil II, (1923), S.~20--21
\item 關於公理化問題的群論備註，32, Teil II, (1923), S.~21
\item 任意域中的 Galois 理論，33, Teil II, (1925), S.~120
\item 理想論中的 Hilbert 函數，34, Teil II, (1926), S.~101
\item 理想的微分商與分歧理論，38, Teil II, (1929), S.~81
\end{itemize}

\subsection*{發表於 \foreign{Jahresberichte der Deutschen Mathematiker-Vereinigung} 的書評}

\begin{itemize}[leftmargin=1.5em,itemsep=0.2em]
\item E.\ Study, \textit{Einleitung in die Theorie der Invarianten linearer Transformationen
auf Grund der Vektorenrechnung.} Erster Teil. (Die Wissenschaft, Bd.~71)
Braunschweig 1923, Vieweg.\\
36, Teil II, (1927), S.~71

\item Leopold Kroneckers Werke, Bd.~4. Herausgegeben von K.\ Hensel. VI u.
509 S. Leipzig und Berlin 1929, B.\ G.\ Teubner.\\
40, Teil II, (1931), S.~11--12

\item Leopold Kroneckers Werke, Bd.~5. Herausgegeben von K.\ Hensel. 527 S.
Leipzig und Berlin 1930, B.\ G.\ Teubner.\\
41, Teil II, (1932), S.~17

\item Leopold Kroneckers Werke, Bd.\ III.\,2. Herausgegeben von K.\ Hensel.
215 S.\ Leipzig und Berlin 1931, B.\ G.\ Teubner.\\
42, Teil II, (1933), S.~38--39

\item D.\ E.\ Rutherford, \textit{Modular Invariants.} Cambridge Tracts No.~27.
Cambridge 1932, VII und 84 S.\\
43, Teil II, (1934), S.~94--95
\end{itemize}

\newpage
\section*{在 Emmy Noether 參與下產生的書籍}
\addcontentsline{toc}{section}{在 Emmy Noether 參與下產生的書籍}

\begin{itemize}[leftmargin=1.5em,itemsep=0.5em]
\item R.\ Dedekind, \textit{Gesammelte mathematische Werke} I--III. R.\ Fricke、
Emmy Noether 與 Øystein Ore 編。Vieweg, Braunschweig 1930--1932\\
(Bd.~I 1930; Bd.~II 1931; Bd.~III 1932)

\item G.\ Cantor, \textit{Briefwechsel mit R.\ Dedekind.} Emmy
Noether 和 J.\ Cavaillès 編，60 頁。\foreign{Actualités scient.\ et ind.}\ 581,
Paris 1937
\end{itemize}

\vfill

\begin{center}
\rule{0.5\textwidth}{0.4pt}\\[1em]
\textit{選定數學論文結束}\\[0.5em]
\textit{Emmy Noether}\\[0.5em]
\textit{(1882--1935)}
\end{center}

%
\end{document}



